\documentclass[11pt,a4paper]{amsart}

\usepackage{amsmath,amssymb,amsthm,mathtools,bm,mathrsfs}
\usepackage[colorlinks=true,linkcolor=black,citecolor=black,urlcolor=black]{hyperref}
\pdfstringdefDisableCommands{%
  \def\Omega{Omega}%
  \def\Phi{Phi}%
  \def\PsiCal{Psi}%
  \def\RH{RH}%
  \def\corr{corr}%
  \def\toy{toy}%
  \def\val{val}%
  \def\aut{aut}%
  \def\jet{jet}%
  \def\far{far}%
  \def\est{est}%
  \def\rem{rem}%
  \def\sym{sym}%
  \def\curv{curv}%
  \def\tail{tail}%
  \def\mix{mix}%
  \def\core{core}%
  \def\aux{aux}%
  \def\out{out}%
  \def\pair{pair}%
  \def\lin{lin}%
  \def\new{new}%
}
\usepackage{enumitem}
\usepackage{microtype}
\usepackage[margin=1in]{geometry}

\theoremstyle{plain}
\newtheorem{theorem}{Theorem}[section]
\newtheorem{proposition}[theorem]{Proposition}
\newtheorem{lemma}[theorem]{Lemma}
\newtheorem{corollary}[theorem]{Corollary}

\theoremstyle{definition}
\newtheorem{definition}[theorem]{Definition}
\newtheorem{remark}[theorem]{Remark}
\newtheorem{claim}[theorem]{Claim}
\newtheorem{hypothesis}[theorem]{Hypothesis}

\title{A Proof of the Riemann Hypothesis}

\author{John Shields}
\thanks{Preprint.}
\date{April 2026}

\subjclass[2020]{11M26 (primary), 11M06, 11M50 (secondary)}
\keywords{Riemann hypothesis, zeta function, critical line, phase kernel,
  jet normalization, odd-moment cancellation}

\newcommand{\RH}{\mathrm{RH}}
\newcommand{\z}{\zeta}
\newcommand{\Om}{\Omega}
\newcommand{\Ph}{\Phi}
\newcommand{\Del}{\Delta}
\newcommand{\Aut}{\mathrm{aut}}
\newcommand{\sym}{\mathrm{sym}}
\newcommand{\loc}{\mathrm{loc}}
\newcommand{\sm}{\mathrm{sm}}
\newcommand{\far}{\mathrm{far}}
\newcommand{\est}{\mathrm{est}}
\newcommand{\rem}{\mathrm{rem}}
\newcommand{\toy}{\mathrm{toy}}
\newcommand{\val}{\mathrm{val}}
\newcommand{\corr}{\mathrm{corr}}
\newcommand{\main}{\mathrm{main}}
\newcommand{\err}{\mathrm{err}}
\newcommand{\jet}{\mathrm{jet}}
\newcommand{\bg}{\mathrm{bg}}
\newcommand{\curv}{\mathrm{curv}}
\newcommand{\tail}{\mathrm{tail}}
\newcommand{\mix}{\mathrm{mix}}
\newcommand{\core}{\mathrm{core}}
\newcommand{\aux}{\mathrm{aux}}
\newcommand{\out}{\mathrm{out}}
\newcommand{\pair}{\mathrm{pair}}
\newcommand{\lin}{\mathrm{lin}}
\newcommand{\new}{\mathrm{new}}
\newcommand{\off}{\mathrm{off}}
\newcommand{\odd}{\mathrm{odd}}
\newcommand{\amp}{\mathrm{amp}}
\newcommand{\actual}{\mathrm{actual}}
\newcommand{\can}{\mathrm{can}}
\newcommand{\source}{\mathrm{source}}
\newcommand{\base}{\mathrm{base}}
\newcommand{\wh}{\mathrm{wh}}
\newcommand{\phys}{\mathrm{phys}}
\newcommand{\quot}{\mathrm{quot}}
\newcommand{\Qabs}{\mathcal Q^{\mathrm{abs}}}
\newcommand{\F}{\mathrm{F}}
\newcommand{\RePart}{\operatorname{Re}}
\newcommand{\dist}{\operatorname{dist}}
\newcommand{\diag}{\operatorname{diag}}
\newcommand{\Tr}{\operatorname{Tr}}
\newcommand{\Ker}{\operatorname{Ker}}
\newcommand{\sgn}{\operatorname{sgn}}

\begin{document}

\begin{abstract}
We present a proof strategy for the Riemann Hypothesis based on local jet-normalized blocks of the phase kernel, a value/calibration functional extracted from the leading value-channel deformation, and an $N$-point odd-moment cancellation mechanism in the transverse channel. The document is written as a living proof draft: every statement is presented in proof form and the notation is frozen globally, so that future revisions can strengthen any remaining technical points without changing the architecture.
\end{abstract}

\maketitle

\tableofcontents

\newpage

\section{Introduction}

The Riemann Hypothesis, originating in Riemann's 1859 memoir \cite{Riemann1859}, asserts that every nontrivial zero of $\zeta(s)$ lies on the critical line $\RePart s = \tfrac12$ \cite{Bombieri2000,Conrey2003,Edwards2001,Titchmarsh1986}. The argument developed here is local and geometric. It compares two objects:
\begin{enumerate}[label=(\roman*),nosep]
    \item a local jet-whitened block extracted from the zeta/scattering phase on short symmetric windows, and
    \item a local off-line toy model corresponding to a split pair away from the critical line.
\end{enumerate}
The original goal was to show that the toy anomaly has a transverse component
that survives every local normalization, while the zeta-side transverse channel
is annihilated at leading value order and becomes exponentially small after
odd-moment cancellation.  The present draft should now be read more cautiously:
the local analysis has produced a detailed obstruction framework and several
conditional rank-two visibility theorems, but the global proof of \(\RH\) is
not closed.

The present proof-state is layered.  Upstream, ordinary Gate~1 is conditional
on three named actual-zeta nonrecurrence/incidence inputs: fixed-width cell
nonrecurrence, cofactor-vector nonrecurrence, and Wronskian angle
nonalignment.  The deep singular branch is conditional on a singular
explicit-formula anti-design exclusion, while the mild transition region
requires a separate transition-control theorem.  Downstream, Gate~2 has been
reduced to source thickening plus an \(L^2\)-input real-rigidity lemma, and
the row-energy layer has been reduced to gauge-locked wall escape plus
carrier-quotient D4 row conversion.  Thus the draft is a conditional local
obstruction architecture, not a completed proof of \(\RH\).

Classically, $\RH$ sits at the intersection of the zero distribution of $\zeta(s)$, the explicit formula relating zeros and primes, and a number of equivalent or partially equivalent positivity, spectral, and functional-analytic reformulations \cite{Bombieri2000,Conrey2003}.

We introduce a normalized local jet-Gramian associated to the corrected local phase geometry. This object encodes the intrinsic local deformation after removal of the dominant value-channel and background effects, and it is the basic carrier of the odd quotient information used in the argument.

The local framework originally had four main pillars.
\begin{enumerate}[label=\arabic*.,nosep]
    \item \emph{Jet normalization:} replace raw point samples by local value/derivative coordinates.
    \item \emph{Value/curvature splitting:} separate the leading background-subtracted value channel from the curvature and tail corrections.
    \item \emph{Canonical calibration:} extract the correct value scale from the actual leading zeta-side value matrix rather than from an ad hoc scalar functional.
    \item \emph{$N$-point transverse cancellation:} use an odd holomorphic expansion and explicit cancellation coefficients to crush the zeta transverse channel exponentially in the local height parameter.
\end{enumerate}

\begin{remark}[Audit status of the four pillars]
\label{rem:wip-four-pillars-audit-status}
The four pillars above remain part of the local language of the draft, but
they should not be read as a complete proof gate.  The later audits show that
the decisive obstruction is no longer transverse odd-mode detection alone.
The proof-state has separated into named upstream actual-zeta inputs
(O1--O3, S1, and transition control) and downstream local conversion theorems
(Gate~2 real rigidity, wall escape, D4 row conversion, and row mobility).
See Remark~\ref{rem:current-architecture-ledger-unified} for the terminal
dependency ledger.
\end{remark}

This approach is not aimed at discovering new direct relations among primes. Rather, it studies a local, normalized geometric object attached to the zeta function and argues that an off-line zero would force a forbidden transverse odd mode in that object. In that sense, the method is primarily a zeta-side rigidity argument, with prime-distribution consequences entering only indirectly through the arithmetic content encoded in the zeta phase and, more broadly, through the classical zero--prime linkage expressed by the explicit formula \cite{Bombieri2000,Conrey2003,Edwards2001,Titchmarsh1986}. This local geometric viewpoint is not standard among the main classical families of approaches to $\RH$ surveyed in the literature \cite{Bombieri2000,Conrey2003}.

For orientation, it may be helpful to situate the present framework relative to earlier geometric viewpoints on $\RH$. One such direction is the Speiser--Levinson--Montgomery line, in which the geometry of zeros of $\zeta'(s)$ is used to encode information about zeros of $\zeta(s)$ \cite{Speiser1935,LevinsonMontgomery1974}. Another is the de~Branges Hilbert-space program, in which one seeks global kernel/positivity structures attached to entire functions related to $\zeta$ \cite{deBranges1986,ConreyLi2000}. The present approach appears to differ in emphasis: rather than working directly with raw derivative geometry or with a global positivity condition, it constructs a locally normalized transverse invariant after quotienting out the dominant value/gauge directions. This comparison is intended only as heuristic orientation for the reader; no claim is made here that the present framework subsumes those earlier viewpoints.

\paragraph{A Mellin-side interpretation of the $N$-point layer.}
A second point of orientation concerns the $N$-point cancellation operator
introduced later in the proof.  The cancellation coefficients may be viewed as
defining a discrete multiplicative distribution whose Mellin transform preserves
the constant mode while annihilating the first $N-1$ odd jets.  In this
sense, the $N$-point step is not merely a finite-difference device: it may be
read as a local Mellin-side projector acting on the corrected transverse germ.
This is formally reminiscent of the Weil/Connes test-function viewpoint, in
which one probes the zero set through Mellin/Fourier transforms of
multiplicative test functions and explicit- or trace-formula pairings
\cite{Conrey2003,Connes1997,ConnesConsani2021}.  The present argument does
\emph{not} invoke any global trace formula, positivity principle, or spectral
realization theorem from that literature; the comparison is included only as
heuristic orientation.  What is used here is a local, explicit, finite-support
projector whose Mellin symbol can be written in closed form and whose odd part
starts only at order $2N-1$.

\section{Local kernel and jet normalization}

\subsection{Phase kernel}

Let $\Ph$ be a real phase on the local spectral variable $t$, and define the phase kernel
\begin{equation}\label{eq:phase-kernel}
K_{\Ph}(x,y)=
\begin{cases}
\dfrac{\sin(\Ph(x)-\Ph(y))}{\pi(x-y)}, & x\neq y,\\[1ex]
\dfrac{\Ph'(x)}{\pi}, & x=y.
\end{cases}
\end{equation}
We write
\[
q(t):=\Ph'(t).
\]

\begin{definition}[Actual-zeta real phase and canonical Gate--2 normalization]
\label{def:canonical-actual-zeta-phase-curvature}
On ordinary smooth actual-zeta packets we use a real completed critical-line
phase/scalar
\[
\Phi_\zeta.
\]
Concretely, \(\Phi_\zeta\) is defined in a real completed gauge, for instance
from
\[
\Xi(t)=\xi\!\left(\frac12+it\right)\in\mathbb R
\]
or from an equivalent Hardy-normalized real scalar.  It is not a one-sided AFE
branch, not \(P_2\) alone, not \(D_2\) alone, and not a proxy zero kernel.  If a
local real chart is used to define \(\Phi_\zeta\), every chart denominator is
assumed polynomially bounded below on the ordinary packet.  If this chart
condition fails, the packet must be treated in another chart or routed to the
singular/transition branch.
The associated phase derivative is
\[
q_\zeta:=\mathcal D\Phi_\zeta,
\]
where \(\mathcal D\) denotes the normalized packet derivative.  In sections
where ordinary prime notation is used, this means \(q_\zeta=\Phi_\zeta'\) in
the local spectral coordinate.

The canonical source curvature is
\[
\boxed{
q''_{\can}:=\mathcal D^2q_\zeta=\mathcal D^3\Phi_\zeta.
}
\]
Equivalently, in ordinary local prime notation,
\[
q''_{\can}=q_\zeta''=\Phi_\zeta'''.
\]
This is the full actual-zeta phase-derivative curvature in the same-point Gram
normalization.  It is not whitened, projectivized, AFE-normalized, or
affine-subtracted.
Finite zero-side proxies, one-sided AFE truncations, core-only curvatures, and
affine-subtracted complement curvatures are diagnostic objects only unless the
corresponding error has explicitly been included in \(q''_{\can}\) or proved
perturbative at the stated downstream scale.

When no ambiguity is possible, centered rank-two formulas may write
\(q,q'',q^{(4)}\) for
\[
q_\zeta,\qquad q''_{\can},\qquad \mathcal D^4q_\zeta.
\]
If a source/core split is being discussed, the affine-subtracted complement
curvature is instead denoted
\[
q''_{\source,R}(m):=(G^{\new})''(m),
\]
and should not be identified with \(q''_{\can}\) unless the entire finite-\(s\)
row package has explicitly been rebuilt with \(q''_{\source,R}\) in place of
the full actual-zeta curvature.

On ordinary smooth zero-separated packets,
\[
q_\zeta,\qquad q''_{\can},\qquad \Gamma_G,\qquad \Gamma_{G,2}
\]
are real-valued.  If a critical-line zero lies in the unregularized packet
interval, the packet is routed to the singular/transition branch, or else the
corresponding regularized phase convention must be used.
\end{definition}

\begin{lemma}[Reality of the completed critical-line scalar]
\label{lem:completed-critical-line-scalar-real}
Let
\[
\xi(s)=\frac12s(s-1)\pi^{-s/2}\Gamma(s/2)\zeta(s).
\]
Then
\[
\Xi(t):=\xi\!\left(\frac12+it\right)
\]
is real for all real \(t\).
\end{lemma}

\begin{proof}
The completed function satisfies
\[
\xi(s)=\xi(1-s),
\qquad
\overline{\xi(s)}=\xi(\overline{s}).
\]
Thus, for \(s=\frac12+it\),
\[
\overline{\xi\!\left(\frac12+it\right)}
=
\xi\!\left(\frac12-it\right)
=
\xi\!\left(1-\left(\frac12+it\right)\right)
=
\xi\!\left(\frac12+it\right).
\]
Therefore \(\Xi(t)\in\mathbb R\).
\end{proof}

\begin{definition}[Canonical Gate--2 anisotropy functionals]
\label{def:canonical-gate-two-anisotropy-functionals}
In the canonical actual-zeta convention of
Definition~\ref{def:canonical-actual-zeta-phase-curvature}, define
\[
\Gamma_G
:=
q''_{\can}-6q_\zeta+2q_\zeta^3,
\]
and
\[
\Gamma_{G,2}
:=
(1-q_\zeta^2)q''_{\can}
-\frac16\mathcal D^4q_\zeta.
\]
When the local coordinate is clear, we suppress the \(\zeta\)-subscript and
write \(q\) for \(q_\zeta\).
\end{definition}

\begin{remark}[Derivative notation in the canonical Gate--2 identities]
\label{rem:gate-two-derivative-notation}
Throughout the canonical Gate--2 identities, primes denote the same normalized
packet derivative \(\mathcal D\) used in
Definition~\ref{def:canonical-actual-zeta-phase-curvature}.  Thus
\[
        q'= \mathcal D q_\zeta,\qquad
        q''=\mathcal D^2q_\zeta=q''_{\can},
        \qquad
        q^{(4)}=\mathcal D^4q_\zeta.
\]
In physical coordinates \(m\), one has \(\mathcal D=\lambda\partial_m\) on
ordinary packets.  All identities involving \(q'',\Gamma_G,\Gamma_{G,2}\), and
\(\Gamma_G''\) are to be read in this normalized derivative convention unless
explicitly stated otherwise.
\end{remark}

\begin{lemma}[Differentiated form of the second Gate--2 anisotropy]
\label{lem:differentiated-form-second-gate-two-anisotropy}
With the canonical definitions above,
\[
\boxed{
\Gamma_{G,2}
=
2q_\zeta(\mathcal Dq_\zeta)^2
-\frac16\mathcal D^2\Gamma_G.
}
\]
\end{lemma}

\begin{proof}
From
\[
\Gamma_G=q''_{\can}-6q_\zeta+2q_\zeta^3
\]
and \(q''_{\can}=\mathcal D^2q_\zeta\), we compute
\[
\mathcal D^2\Gamma_G
=
\mathcal D^4q_\zeta
-6q''_{\can}
+
12q_\zeta(\mathcal Dq_\zeta)^2
+
6q_\zeta^2q''_{\can}.
\]
Hence
\[
\begin{aligned}
2q_\zeta(\mathcal Dq_\zeta)^2-\frac16\mathcal D^2\Gamma_G
&=
2q_\zeta(\mathcal Dq_\zeta)^2
-\frac16\mathcal D^4q_\zeta
+q''_{\can}
-2q_\zeta(\mathcal Dq_\zeta)^2
-q_\zeta^2q''_{\can}  \\
&=
(1-q_\zeta^2)q''_{\can}
-\frac16\mathcal D^4q_\zeta
=
\Gamma_{G,2}.
\end{aligned}
\]
\end{proof}

\begin{lemma}[Expanded Gate--2 derivative identities]
\label{lem:expanded-gate-two-derivative-identities}
In the normalized canonical convention,
\[
        u=q''_{\can}=q'',
\qquad
        g=\Gamma_G=q''-6q+2q^3,
\]
and
\[
        h=\Gamma_{G,2}.
\]
Then
\[
        h=(1-q^2)q''-\frac16q^{(4)}
        =
        (1-q^2)u-\frac16q^{(4)}.
\]
Moreover,
\[
        g'=q'''-6q'+6q^2q',
\]
and
\[
        h'
        =
        -2qq'q''
        +(1-q^2)q'''
        -\frac16q^{(5)}.
\]
Equivalently,
\[
        h'
        =
        -2qq'u
        +(1-q^2)u'
        -\frac16q^{(5)}.
\]
\end{lemma}

\begin{proof}
The identity for \(h\) follows by expanding
\[
        h=2q(q')^2-\frac16g''
\]
and using
\[
        g''=q^{(4)}-6q''+12q(q')^2+6q^2q''.
\]
Thus
\[
\begin{aligned}
h
&=
2q(q')^2
-\frac16q^{(4)}
+q''
-2q(q')^2
-q^2q''        \\
&=
(1-q^2)q''
-\frac16q^{(4)}.
\end{aligned}
\]
The formulas for \(g'\) and \(h'\) follow by differentiating
\[
        g=q''-6q+2q^3
\]
and
\[
        h=(1-q^2)q''-\frac16q^{(4)}.
\]
\end{proof}

\begin{remark}[Derivative convention and gauge-locking for downstream use]
\label{rem:derivative-convention-gauge-locking-downstream}
All downstream Gate--2, wall-escape, and D4 row-conversion identities use the
same normalized derivative convention
\[
\mathcal D=\lambda\partial_m.
\]
Thus, whenever ordinary prime notation is used in a local packet coordinate in
the downstream sections, it is shorthand for \(\mathcal D\)-differentiation in
that normalized local coordinate.  In particular,
\[
q''=q''_{\can}=\mathcal D^2q_\zeta,
\qquad
\Gamma_G=q''_{\can}-6q_\zeta+2q_\zeta^3,
\]
and
\[
\Gamma_{G,2}
=
2q_\zeta(\mathcal Dq_\zeta)^2-\frac16\mathcal D^2\Gamma_G.
\]

The scalar \(q_\zeta\) used here is also the scalar that appears in the
\(q_0\)-wall and in the D4 source factor.  Thus the wall coordinate is not an
arbitrary affine primitive of \(q''_{\can}\).  The affine freedom in solving
\(\mathcal D^2q=q''_{\can}\) is fixed by the canonical real phase convention
above.  This gauge-locking is necessary because neither
\[
\Gamma_G=q''_{\can}-6q+2q^3
\]
nor the source factor
\[
H_{\mathrm{lead}}=-\frac{q(q''_{\can})^3}{288}
\]
is invariant under affine changes \(q\mapsto q+am+b\).
\end{remark}

For two nearby points around a center $T$, the natural local point basis is not stable under coalescence. One instead passes to the point-to-jet matrix
\begin{equation}\label{eq:point-to-jet}
P_h:=\frac12
\begin{pmatrix}
1 & 1\\[1ex]
-1/h & 1/h
\end{pmatrix},
\end{equation}
which turns two point samples into local value and derivative coordinates.

\subsection{Jet-limit local blocks}

Take two centers $T_1\neq T_2$. In the point basis, let $A_h$ denote the same-point $2\times 2$ block around a single center and $C_h$ the mixed block between the two centers. Passing to jet coordinates yields finite limits as $h\to 0$:
\[
\widetilde A_h=P_hA_hP_h^\top \to J(T),
\qquad
\widetilde C_h=P_hC_hP_h^\top \to J_{12}.
\]
A direct Taylor expansion gives
\begin{equation}\label{eq:J-T}
J(T)=\frac1\pi
\begin{pmatrix}
2q(T) & \frac12 q'(T)\\[1ex]
\frac12 q'(T) & \frac1{12}\bigl(q''(T)+2q(T)^3\bigr)
\end{pmatrix}.
\end{equation}
For two centers $T_1,T_2$, writing
\[
\Delta:=\Ph(T_1)-\Ph(T_2),
\qquad s:=T_1-T_2,
\qquad q_j:=q(T_j),
\]
one obtains the explicit cross-jet matrix
\begin{equation}\label{eq:N12}
N_{12}=
\begin{pmatrix}
\dfrac{2\sin \Delta}{s}
&
\dfrac{\sin \Delta-q_2 s\cos \Delta}{s^2}
\\[2ex]
\dfrac{q_1 s\cos \Delta-\sin \Delta}{s^2}
&
\dfrac{(q_1+q_2)s\cos \Delta+(q_1q_2s^2-2)\sin \Delta}{2s^3}
\end{pmatrix}.
\end{equation}
Writing
\[
J(T_j)=\frac1\pi M_j,
\qquad
J_{12}=\frac1\pi N_{12},
\]
the common $\pi^{-1}$ factor cancels under whitening, and the jet-limit whitened cross block is
\begin{equation}\label{eq:omega-infty}
\Omega_\infty(T_1,T_2)=M_1^{-1/2}N_{12}M_2^{-1/2}.
\end{equation}
This is the canonical local matrix to compare with the toy anomaly.

\section{Finite-\texorpdfstring{$s$}{s} local blocks}

\subsection{Symmetric placement}

Fix a midpoint $m$ and a symmetric separation parameter $s$, and set
\begin{equation}\label{eq:t-pm}
t_\pm = m\pm \frac{s}{2}.
\end{equation}
At finite $s$, the same-point jet Gram blocks are naturally two matrices,
\begin{equation}\label{eq:Gpm}
G_{m,\pm}(s)=\frac1\pi
\begin{pmatrix}
2q(t_\pm) & \frac12 q'(t_\pm)\\[1ex]
\frac12 q'(t_\pm) & \frac1{12}\bigl(q''(t_\pm)+2q(t_\pm)^3\bigr)
\end{pmatrix}.
\end{equation}
The mixed block is obtained from \eqref{eq:N12} by substituting $T_1=t_-$, $T_2=t_+$ and $s=t_--t_+=-s$:
\begin{equation}\label{eq:Nm}
\resizebox{.94\linewidth}{!}{$
N_m(s)=\frac1\pi
\left(\begin{smallmatrix}
-\dfrac{2\sin \Delta_m(s)}{s}
&
\dfrac{\sin \Delta_m(s)+q_+(s)\,s\cos \Delta_m(s)}{s^2}
\\[1.5ex]
-\dfrac{\sin \Delta_m(s)+q_-(s)\,s\cos \Delta_m(s)}{s^2}
&
\dfrac{(q_-(s)+q_+(s))s\cos \Delta_m(s)+(2-q_-(s)q_+(s)s^2)\sin \Delta_m(s)}{2s^3}
\end{smallmatrix}\right)
$}
\end{equation}
where
\[
q_\pm(s):=q(t_\pm),
\qquad
\Delta_m(s):=\Ph(t_-)-\Ph(t_+).
\]
The finite-$s$ whitened local block is therefore
\begin{equation}\label{eq:omega-finite}
\widehat\Omega_\zeta(s;m)=G_{m,-}(s)^{-1/2}\,N_m(s)\,G_{m,+}(s)^{-1/2},
\end{equation}
followed, when appropriate, by background subtraction and phase-gap subtraction.

\subsection{Removable singularity structure}

The powers $1/s$, $1/s^2$, $1/s^3$ in \eqref{eq:Nm} are removable at $s=0$. Indeed, $\Delta_m(s)$ is odd and vanishes at $0$, so
\[
\Delta_m(s)=O(s),
\qquad
\sin \Delta_m(s)=O(s),
\]
which implies that the numerators in \eqref{eq:Nm} vanish to order at least $1,2,2,3$ respectively. Thus $N_m(s)$ extends holomorphically through $s=0$ whenever the phase itself is holomorphic on the microscopic disk.

\section{Zeta-side decomposition}

Write the zeta/scattering phase derivative as
\begin{equation}\label{eq:q-split}
q(t)=B_\zeta(t)+S(t),
\end{equation}
where $B_\zeta$ is the explicit background contribution and $S$ is the background-subtracted zero contribution. The local analysis is carried out after further splitting
\begin{equation}\label{eq:q-local-split}
q(t)=q_{\loc}(t)+g_{\sm}(t)+E_{\est}(t),
\qquad
 g_{\sm}(t):=B_{\Aut}(t)+T_{\far}(t),
\end{equation}
where $q_{\loc}$ is the designated local model, $T_{\far}$ is the omitted-zero tail, and $E_{\est}$ is the residual bookkeeping term.

Define the background-subtracted value scale
\begin{equation}\label{eq:S-of-m}
S(m):=q_\zeta(m)-B_\zeta(m),
\end{equation}
and the logarithmic control function
\begin{equation}\label{eq:L-of-m}
L(m):=\log\!\Bigl(3+|m|+\frac1{S(m)}\Bigr).
\end{equation}

\subsection{Curvature estimate}

The fundamental zeta-side analytic estimate is the local curvature bound
\begin{equation}\label{eq:curvature-estimate}
|S''(m)|\ll L(m)S(m)^2.
\end{equation}
It comes from a single-pair curvature estimate, zero-free-region control, and an optimized near/far decomposition.

\section{Smooth remainder after affine-jet subtraction}

Fix a short interval $I=[a,b]$ of length $L=b-a$ and midpoint $t_*=(a+b)/2$. Let
\[
J_1[g_{\sm}](t):=g_{\sm}(t_*)+g_{\sm}'(t_*)(t-t_*),
\qquad
r(t):=g_{\sm}(t)-J_1[g_{\sm}](t).
\]
Assume
\[
\|g_{\sm}''\|_{L^\infty(I)}\le S_2.
\]
Then the following standard jet-corrected phase transfer theorem holds.

\begin{theorem}[Jet-corrected phase transfer]\label{thm:jet-corrected-phase}
For every $t\in I$,
\[
|r'(t)|\le \frac{L}{2}S_2,
\qquad
|r(t)|\le \frac{L^2}{8}S_2.
\]
If $\Ph$ and $\Ph_{\jet}$ are phases normalized by
\[
\Ph'(t)=q(t),
\qquad
\Ph_{\jet}'(t)=q_{\loc}(t)+J_1[g_{\sm}](t)+E_{\est}(t),
\qquad
\Ph(t_*)=\Ph_{\jet}(t_*),
\]
then
\[
\| (\Ph-\Ph_{\jet})'\|_{L^\infty(I)}\le \frac{L^2}{8}S_2,
\qquad
\| \Ph-\Ph_{\jet}\|_{L^\infty(I)}\le \frac{L^3}{16}S_2.
\]
\end{theorem}

\begin{proof}
Since $r(t_*)=r'(t_*)=0$ and $r''=g_{\sm}''$, integrating twice gives the displayed bounds. Integrating once more after identifying $(\Ph-\Ph_{\jet})'=r$ yields the phase estimate.
\end{proof}

\section{Tail curvature bound}

Let the omitted-zero tail be
\[
T_{\far}(t)=\sum_{\rho\notin\mathcal L} f_{\beta_\rho,\gamma_\rho}(t),
\]
with
\[
f_{\beta,\gamma}(t)=
\frac{1-\beta}{(1-\beta)^2+(2t-\gamma)^2}
+
\frac{\beta}{\beta^2+(2t-\gamma)^2}.
\]
For each omitted zero define
\[
\Delta_\rho:=\dist(2I,\gamma_\rho).
\]

\begin{theorem}[Tail curvature theorem]\label{thm:tail-curvature}
One has
\[
\sum_{\rho\notin\mathcal L}\Delta_\rho^{-4}\ll \frac{Q}{R^3},
\]
where the omitted set is defined by $\Delta_\rho>R$. Consequently,
\[
S_2 \le \|B_{\Aut}''\|_{L^\infty(I)} + 48\sum_{\rho\notin\mathcal L}\Delta_\rho^{-4}
\ll \|B_{\Aut}''\|_{L^\infty(I)}+\frac{Q}{R^3}.
\]
In particular, if $R\to\infty$ with $m$, then $S_2=o(Q)$.
\end{theorem}

\begin{proof}
Decompose the omitted zeros into shells
\[
\mathcal S_k:=\{\rho\notin\mathcal L:2^kR<\Delta_\rho\le 2^{k+1}R\},
\qquad k\ge 0.
\]
Then
\[
\sum_{\rho\notin\mathcal L}\Delta_\rho^{-4}
\le
\sum_{k\ge 0}\#\mathcal S_k\,(2^kR)^{-4}.
\]
Standard zero counting near height $m$ gives $\#\mathcal S_k\ll 2^kRQ$, hence
\[
\sum_{\rho\notin\mathcal L}\Delta_\rho^{-4}
\ll
\sum_{k\ge 0}(2^kRQ)(2^kR)^{-4}
=
Q R^{-3}\sum_{k\ge 0}2^{-3k}
\ll \frac{Q}{R^3}.
\]
The bound for $S_2$ follows from the definition.
\end{proof}

\section{Toy anomaly and transverse obstruction}

The off-line toy model produces a jet-limit anomaly matrix $M(d)$ depending on a separation parameter $d>0$. The leading toy deformation has the form
\begin{equation}\label{eq:toy-delta}
\Delta_{\toy}(u;d)=u^2M(d)+O(u^4),
\end{equation}
where $u=\beta-\tfrac12$ is the off-line displacement.

Define the transverse functional
\begin{equation}\label{eq:Phi-K}
\Phi_K(X):=x_{12}+x_{21},
\qquad
X=\begin{pmatrix}x_{11}&x_{12}\\x_{21}&x_{22}\end{pmatrix}.
\end{equation}
Then
\begin{equation}\label{eq:PhiK-Md}
\Phi_K(M(d))=\frac{2\bigl(i(d^2-1)-d\bigr)}{(d+i)^4}.
\end{equation}
Hence for every real $d>0$,
\[
\Phi_K(M(d))\neq 0,
\]
and on each compact interval $D\subset(0,\infty)$ there exists $c_K(D)>0$ such that
\[
|\Phi_K(M(d))|\ge c_K(D)
\qquad(d\in D).
\]

\section{Leading value matrix and canonical calibration}

We now make the value-channel derivative $A_{\val}$ and its pairing with the
toy anomaly matrix $M(d)$ fully explicit.

Work in the symmetric local regime
\[
T_1=m-\frac{\sigma}{2},
\qquad
T_2=m+\frac{\sigma}{2},
\]
and write
\[
B:=B_\zeta(m),
\qquad
x:=B\,\sigma.
\]
When convenient, we write $A_{\val}(x)$ for $A_{\val}(m,\sigma)$ under the
relation $x=B_\zeta(m)\sigma$.

\paragraph{Pure value-channel perturbation.}
Let the local value-channel parameter be perturbed by
\[
q_1=q_2=B+\alpha,
\]
while all derivative and curvature data are frozen at background level. In the
constant-background symmetric approximation, the phase gap is
\[
\Delta=\Delta_0-\alpha\sigma,
\]
with baseline value
\[
\Delta_0=-B\sigma=-x.
\]

Recall that the jet-limit local block is
\[
\Omega_\infty(T_1,T_2)=M_1^{-1/2}N_{12}M_2^{-1/2}.
\]
Define the pure value-channel derivative by
\begin{equation}\label{eq:A-val-def}
A_{\val}(m,\sigma):=
\left.\frac{\partial}{\partial\alpha}\Omega_\infty(T_1,T_2)\right|_{\alpha=0}.
\end{equation}

\begin{proposition}[Explicit formula for the value-channel derivative]
\label{prop:explicit-Aval}
In the symmetric constant-background regime,
\[
A_{\val}(x)=\frac1B
\begin{pmatrix}
\dfrac{x\cos x-\sin x}{x}
&
\dfrac{\sqrt3\bigl(-x^2\sin x-2x\cos x+2\sin x\bigr)}{x^2}
\\[2.5ex]
\dfrac{\sqrt3\bigl(x^2\sin x+2x\cos x-2\sin x\bigr)}{x^2}
&
\dfrac{3\bigl(x^3\cos x-6x\cos x+3(2-x^2)\sin x\bigr)}{x^3}
\end{pmatrix}.
\]
\end{proposition}

\begin{proof}
In the constant-background symmetric regime, the same-point jet Gram matrices
are diagonal:
\[
M_1=M_2=M=
\begin{pmatrix}
2(B+\alpha) & 0\\[0.5ex]
0 & (B+\alpha)^3/6
\end{pmatrix}.
\]
Hence
\[
M^{-1/2}=
\begin{pmatrix}
(2(B+\alpha))^{-1/2} & 0\\[0.5ex]
0 & \sqrt6\,(B+\alpha)^{-3/2}
\end{pmatrix}.
\]

In the same regime, the jet-limit cross block depends only on $(B+\alpha)$,
$\sigma$, and the phase gap $\Delta=\Delta_0-\alpha\sigma$. Writing out the
explicit jet-limit cross block in this specialization gives
\[
\resizebox{.94\linewidth}{!}{$
N_{12}(\alpha)=
\left(\begin{smallmatrix}
-\dfrac{2\sin\Delta}{\sigma}
&
\dfrac{\sqrt3\bigl(\sin\Delta+(B+\alpha)\sigma\cos\Delta\bigr)}{\sigma^2}
\\[1.5ex]
-\dfrac{\sqrt3\bigl(\sin\Delta+(B+\alpha)\sigma\cos\Delta\bigr)}{\sigma^2}
&
\dfrac{3\bigl(2\sin\Delta-2(B+\alpha)\sigma\cos\Delta-(B+\alpha)^2\sigma^2\sin\Delta\bigr)}{\sigma^3}
\end{smallmatrix}\right),
$}
\]
up to the common normalization factor already absorbed into the definition of
$\Omega_\infty$.

Now differentiate
\[
\Omega_\infty(\alpha)=M(\alpha)^{-1/2}N_{12}(\alpha)M(\alpha)^{-1/2}
\]
at $\alpha=0$, using
\[
\frac{d}{d\alpha}(B+\alpha)^{-1/2}\Big|_{\alpha=0}
=-\frac12 B^{-3/2},
\qquad
\frac{d}{d\alpha}(B+\alpha)^{-3/2}\Big|_{\alpha=0}
=-\frac32 B^{-5/2},
\]
and
\[
\frac{d}{d\alpha}\sin(\Delta_0-\alpha\sigma)\Big|_{\alpha=0}
=-\sigma\cos\Delta_0,
\qquad
\frac{d}{d\alpha}\cos(\Delta_0-\alpha\sigma)\Big|_{\alpha=0}
=\sigma\sin\Delta_0.
\]

After substituting $\Delta_0=-x$, $x=B\sigma$, using
\[
\sin(-x)=-\sin x,\qquad \cos(-x)=\cos x,
\]
and simplifying each entry, one obtains
\[
A_{\val}(x)=\frac1B
\begin{pmatrix}
\dfrac{x\cos x-\sin x}{x}
&
\dfrac{\sqrt3\bigl(-x^2\sin x-2x\cos x+2\sin x\bigr)}{x^2}
\\[2.5ex]
\dfrac{\sqrt3\bigl(x^2\sin x+2x\cos x-2\sin x\bigr)}{x^2}
&
\dfrac{3\bigl(x^3\cos x-6x\cos x+3(2-x^2)\sin x\bigr)}{x^3}
\end{pmatrix}.
\]
\end{proof}

\begin{corollary}[Transverse vanishing on the value channel]
\label{cor:PhiK-Aval-zero}
One has
\[
\Phi_K(A_{\val}(m,\sigma))=0.
\]
\end{corollary}

\begin{proof}
By definition,
\[
\Phi_K(X)=x_{12}+x_{21}.
\]
From Proposition~\ref{prop:explicit-Aval},
\[
A_{\val,12}(x)=
\frac{\sqrt3}{B}\,
\frac{-x^2\sin x-2x\cos x+2\sin x}{x^2},
\]
\[
A_{\val,21}(x)=
\frac{\sqrt3}{B}\,
\frac{x^2\sin x+2x\cos x-2\sin x}{x^2}.
\]
These are exact negatives of each other, hence
\[
A_{\val,12}(x)+A_{\val,21}(x)=0,
\]
and therefore
\[
\Phi_K(A_{\val}(m,\sigma))=0.
\]
\end{proof}

\begin{lemma}[Small-$x$ expansion]
\label{lem:Aval-small-x}
As $x\to 0$,
\[
A_{\val}(x)=\frac1B
\begin{pmatrix}
-\dfrac{x^2}{3}+O(x^4) & -\dfrac{\sqrt3\,x}{3}+O(x^3)\\[1.2ex]
\dfrac{\sqrt3\,x}{3}+O(x^3) & -\dfrac{3x^2}{5}+O(x^4)
\end{pmatrix}.
\]
Consequently,
\[
\|A_{\val}(x)\|_\F\asymp \frac{x}{B}
\qquad (x\to 0).
\]
\end{lemma}

\begin{proof}
Expand
\[
\sin x=x-\frac{x^3}{6}+O(x^5),
\qquad
\cos x=1-\frac{x^2}{2}+O(x^4),
\]
and substitute into the entries of $A_{\val}(x)$.

For the $(1,1)$-entry:
\[
\frac{x\cos x-\sin x}{x}
=
\frac{x(1-\frac{x^2}{2}+O(x^4))-(x-\frac{x^3}{6}+O(x^5))}{x}
=
-\frac{x^2}{3}+O(x^4).
\]

For the $(1,2)$-entry, expanding the numerator gives
\[
-x^2\sin x-2x\cos x+2\sin x
=-x^3+O(x^5)-2x\bigl(1-\tfrac{x^2}{2}+O(x^4)\bigr)+2\bigl(x-\tfrac{x^3}{6}+O(x^5)\bigr)
=-\tfrac{x^3}{3}+O(x^5),
\]
so
\[
\frac{\sqrt3(-x^2\sin x-2x\cos x+2\sin x)}{x^2}
=-\frac{\sqrt3\,x}{3}+O(x^3).
\]

The $(2,1)$-entry is the negative of the $(1,2)$-entry, and the
$(2,2)$-entry similarly simplifies to
\[
-\frac{3x^2}{5}+O(x^4).
\]

Since the off-diagonal entries are linear in $x$ and dominate the quadratic
diagonal terms, the Frobenius norm satisfies
\[
\|A_{\val}(x)\|_\F\asymp \frac{x}{B}.
\]
\end{proof}

\begin{proposition}[Pairing with the toy anomaly]
\label{prop:pairing}
Let $M(d)$ be the toy anomaly matrix. Then, for fixed $d>0$,
\[
B\,\langle A_{\val}(x),M(d)\rangle_\F
=
x\,\frac{2\sqrt3\,d\,(id+3)}{3(d-i)^4}
+
x^2\,\frac{d(-5d^2+46id+77)}{15(d-i)^5}
+
O(x^3).
\]
In particular, for every compact interval
\[
D=[d_-,d_+]\subset(0,\infty),
\]
there exists $x_0(D)>0$ such that for
\[
0<x\le x_0(D),\qquad d\in D,
\]
one has
\[
|\langle A_{\val}(x),M(d)\rangle_\F|\asymp \frac{x}{B}.
\]
\end{proposition}

\begin{proof}
Using Lemma~\ref{lem:Aval-small-x}, write
\[
A_{\val}(x)=\frac1B
\begin{pmatrix}
-\dfrac{x^2}{3}+O(x^4) & -\dfrac{\sqrt3\,x}{3}+O(x^3)\\[1.2ex]
\dfrac{\sqrt3\,x}{3}+O(x^3) & -\dfrac{3x^2}{5}+O(x^4)
\end{pmatrix}.
\]
Now take the Frobenius pairing with the explicit toy anomaly matrix $M(d)$:
\[
\langle A_{\val}(x),M(d)\rangle_\F
=
\sum_{j,k} A_{\val,jk}(x)\,\overline{M_{jk}(d)}.
\]

The linear term in $x$ comes entirely from the off-diagonal entries of
$A_{\val}(x)$, since the diagonal entries are quadratic. Using the explicit
entries of $M(d)$, one finds
\[
B\,\langle A_{\val}(x),M(d)\rangle_\F
=
x\,\frac{\sqrt3}{3}\bigl(\overline{M_{21}(d)}-\overline{M_{12}(d)}\bigr)
+O(x^2).
\]
Substituting the explicit formulas for $M_{12}(d)$ and $M_{21}(d)$ and
simplifying yields
\[
x\,\frac{2\sqrt3\,d\,(id+3)}{3(d-i)^4}.
\]

Including the quadratic contributions from the diagonal entries and the next
terms in the off-diagonal expansion gives
\[
B\,\langle A_{\val}(x),M(d)\rangle_\F
=
x\,\frac{2\sqrt3\,d\,(id+3)}{3(d-i)^4}
+
x^2\,\frac{d(-5d^2+46id+77)}{15(d-i)^5}
+
O(x^3).
\]

The coefficient of $x$ is nonzero for every real $d>0$. Since it depends
continuously on $d$, its modulus is bounded below on every compact interval
$D\subset(0,\infty)$. Shrinking $x_0(D)$ if necessary gives
\[
|\langle A_{\val}(x),M(d)\rangle_\F|\asymp \frac{x}{B}
\]
uniformly for $0<x\le x_0(D)$ and $d\in D$.
\end{proof}

\begin{proposition}[Canonical calibration theorem]
\label{prop:canonical-calibration}
For fixed compact $D\subset(0,\infty)$ and $0<x\le x_0(D)$, define
\[
\Psi_{x,d}(X):=
\frac{\langle A_{\val}(x),X\rangle_\F}
{\langle A_{\val}(x),M(d)\rangle_\F}.
\]
Then:
\begin{enumerate}
\item $\Psi_{x,d}(M(d))=1$;
\item $\Psi_{x,d}(A_{\val}(x))\asymp x/B$;
\item if
\[
\Delta_\zeta=S(m)A_{\val}(x)+R_\zeta,
\]
then
\[
\Psi_{x,d}(\Delta_\zeta)
=
c_{x,d}\,\frac{x}{B}S(m)+\Psi_{x,d}(R_\zeta),
\qquad c_{x,d}\asymp 1;
\]
\item if
\[
\Delta_{\toy}(u;d)=u^2M(d)+O(u^4),
\]
then
\[
\Psi_{x,d}(\Delta_{\toy}(u;d))=u^2+O(u^4).
\]
\end{enumerate}
In particular, whenever
\[
\Psi_{x,d}(R_\zeta)=o\!\left(\frac{x}{B}S(m)\right),
\]
one obtains
\begin{equation}\label{eq:u2-calibration}
u^2\asymp \frac{x}{B}S(m).
\end{equation}
\end{proposition}

\begin{proof}
The first statement is immediate from the definition of $\Psi_{x,d}$.
By Lemma~\ref{lem:Aval-small-x} and Proposition~\ref{prop:pairing},
\[
\Psi_{x,d}(A_{\val}(x))
=
\frac{\|A_{\val}(x)\|_\F^2}{\langle A_{\val}(x),M(d)\rangle_\F}
\asymp
\frac{(x/B)^2}{x/B}
=
\frac{x}{B}.
\]
The remaining statements follow by linearity.
\end{proof}

\begin{proposition}[Calibration remainder under a core-stable cutoff choice]
\label{prop:calibration-remainder-cutoff}
Assume the fixed-core decomposition of
Section~\ref{subsec:fixed-core} and the core-stable cutoff compatibility of
Proposition~\ref{prop:cutoff-compatibility}. If the cutoff $R$ is chosen so
that
\[
R^3\gg \frac{1}{xQ\,S(m)},
\]
then
\[
S_2=o\!\bigl(xQ^2S(m)\bigr),
\]
and consequently
\[
\Psi_{x,d}(R_\zeta)=o\!\left(\frac{x}{B_\zeta(m)}S(m)\right).
\]
\end{proposition}

\begin{proof}
By Theorem~\ref{thm:tail-curvature},
\[
S_2\ll |B_{\Aut}''(m)|+\frac{Q}{R^3}.
\]
The background term is negligible, so under the cutoff condition
\[
R^3\gg \frac{1}{xQ\,S(m)},
\]
one has
\[
\frac{Q}{R^3}=o\!\bigl(xQ^2S(m)\bigr).
\]
Hence
\[
S_2=o\!\bigl(xQ^2S(m)\bigr).
\]
By Corollary~\ref{cor:sharpened-calibration-remainder}, this implies
\[
\Psi_{x,d}(R_\zeta)=o\!\left(\frac{x}{B_\zeta(m)}S(m)\right).
\]
The use of such a cutoff does not alter the toy comparison because, by
Proposition~\ref{prop:cutoff-compatibility}, the core $\mathcal C$ --- and
therefore $M(d)$ and $A_{\val}$ --- is independent of $R$.
\end{proof}

\begin{remark}[Status of the calibration remainder estimate]
\label{rem:status-sharpened-calibration-remainder}
Corollary~\ref{cor:sharpened-calibration-remainder} isolates the quantitative
bound actually used later:
\[
|\Psi_{x,d}(R_\zeta)|\ll_D \frac{S_2}{Q^3}.
\]
Thus the calibration remainder step is no longer a deferred analytic-packaging
item.
\end{remark}

\section{Corrected finite-\texorpdfstring{$s$}{s} holomorphic whitening}

Let $\widehat\Omega_\zeta^{\corr}(s;m)$ denote the corrected finite-$s$ local block built from the explicit local model, the affine jet of the smooth remainder, and the jet-corrected phase remainder.

\begin{proposition}[Microscopic denominator comparability and omitted-smooth holomorphy]
\label{prop:denominator-comparability}
Fix an admissible midpoint $m$ and write
\[
t_\pm=m\pm \frac{s}{2},
\qquad
Q\asymp \log |m|.
\]
For each omitted zero $\rho=\beta_\rho+i\gamma_\rho$, set
\[
D_{\rho,\pm}(s):=\bigl((2m-\gamma_\rho)\pm s\bigr)^2+a_\rho^2,
\qquad
a_\rho\in\{\beta_\rho,1-\beta_\rho\}.
\]
Assume the zero-free region gives
\[
a_\rho\ge \frac{\sigma_0}{Q}
\]
for some absolute $\sigma_0>0$ on the admissible midpoint class. If
\[
|s|\le \frac{\rho_0}{Q},
\qquad
\rho_0\le \frac{\sigma_0}{\sqrt 6},
\]
then
\[
\frac12\Bigl((2m-\gamma_\rho)^2+a_\rho^2\Bigr)
\le
|D_{\rho,\pm}(s)|
\le
\frac32\Bigl((2m-\gamma_\rho)^2+a_\rho^2\Bigr).
\]

Consequently, for every derivative order $r\le r_*$ needed in the corrected
finite-$s$ block entries, the omitted-zero sums defining
\[
T_{\far}(t_\pm),\quad \partial_t T_{\far}(t_\pm),\quad \dots,\quad \partial_t^{\,r_*}T_{\far}(t_\pm)
\]
converge uniformly on $|s|<\rho_0/Q$, provided the corresponding real-line shell
sums converge. In particular,
\[
g_{\sm}(t_\pm)=B_{\Aut}(t_\pm)+T_{\far}(t_\pm)
\]
and its required $t$-derivatives extend holomorphically to
\[
|s|<\rho_0/Q.
\]
\end{proposition}

\begin{proof}
Set
\[
x:=2m-\gamma_\rho,\qquad a:=a_\rho.
\]
Then
\[
D_{\rho,\pm}(s)=(x^2+a^2)\pm 2xs+s^2.
\]
Using
\[
|2xs+s^2|\le 2|x||s|+|s|^2\le \frac12 x^2+3|s|^2,
\]
and the assumptions
\[
a\ge \frac{\sigma_0}{Q},
\qquad
|s|\le \frac{\rho_0}{Q},
\qquad
\rho_0\le \frac{\sigma_0}{\sqrt 6},
\]
we obtain
\[
3|s|^2\le \frac12 a^2.
\]
Hence
\[
|2xs+s^2|
\le
\frac12(x^2+a^2),
\]
which gives the stated lower and upper bounds.

Now every differentiated omitted-zero term is a rational combination of powers of
$D_{\rho,\pm}(s)^{-1}$ with coefficients depending only on the derivative order.
By the above comparability, on $|s|<\rho_0/Q$ these are uniformly comparable to
their real-line counterparts. Therefore the same shell decomposition used on the
real line yields uniform convergence of the differentiated omitted-zero sums on
the complex disk, and hence holomorphy of $T_{\far}(t_\pm)$ and the required
$t$-derivatives there. Adding the holomorphic background term $B_{\Aut}(t_\pm)$
gives the corresponding statement for $g_{\sm}(t_\pm)$.
\end{proof}

\begin{proposition}[Entry-by-entry perturbation estimate for the corrected same-point blocks]
\label{prop:block-perturbation}
In the corrected finite-$s$ normalization, write
\[
G_{m,\pm}^{\corr}(s)
=
\frac1\pi
\begin{pmatrix}
2q(t_\pm) & \frac12 q'(t_\pm)\\[1ex]
\frac12 q'(t_\pm) & \frac1{12}\bigl(q''(t_\pm)+2q(t_\pm)^3\bigr)
\end{pmatrix},
\qquad
t_\pm=m\pm \frac{s}{2}.
\]
Let
\[
q(t)=q_0(t)+\delta q(t),
\]
where $q_0$ is the baseline local-model-plus-affine-jet part and $\delta q$ is the corrected remainder.
Define the baseline same-point blocks
\[
G_{m,\pm}^{(0)}(s)
=
\frac1\pi
\begin{pmatrix}
2q_0(t_\pm) & \frac12 q_0'(t_\pm)\\[1ex]
\frac12 q_0'(t_\pm) & \frac1{12}\bigl(q_0''(t_\pm)+2q_0(t_\pm)^3\bigr)
\end{pmatrix},
\]
and set
\[
R_{m,\pm}(s):=G_{m,\pm}^{\corr}(s)-G_{m,\pm}^{(0)}(s).
\]

Assume that on the microscopic window one has
\[
|\delta q(t_\pm)|\ll L_I^2 S_2,
\qquad
|\delta q'(t_\pm)|\ll L_I S_2,
\qquad
|\delta q''(t_\pm)|\ll S_2,
\]
where $L_I$ is the local window length, and that in the fixed-scaled regime
\[
L_I\asymp \frac{x_0}{Q},
\qquad
|q_0(t_\pm)|\asymp Q.
\]
Then
\[
\boxed{
\|R_{m,\pm}(s)\|\ll S_2
}
\]
uniformly for $|s|<\rho_0/Q$.
\end{proposition}

\begin{proof}
We estimate the entries of $R_{m,\pm}(s)$ one by one.

\medskip
\noindent
\textbf{(1,1)-entry.}
Since
\[
R_{11}
=
\frac1\pi\cdot 2\delta q(t_\pm),
\]
we have
\[
|R_{11}|
\ll
|\delta q(t_\pm)|
\ll
L_I^2S_2
\ll
\frac{S_2}{Q^2}
\ll
S_2.
\]

\medskip
\noindent
\textbf{(1,2)- and (2,1)-entries.}
Since
\[
R_{12}=R_{21}
=
\frac1\pi\cdot \frac12 \delta q'(t_\pm),
\]
we get
\[
|R_{12}|=|R_{21}|
\ll
|\delta q'(t_\pm)|
\ll
L_I S_2
\ll
\frac{S_2}{Q}
\ll
S_2.
\]

\medskip
\noindent
\textbf{(2,2)-entry.}
We have
\[
R_{22}
=
\frac1{12\pi}
\left(
\delta q''(t_\pm)
+
2\Bigl((q_0+\delta q)^3-q_0^3\Bigr)(t_\pm)
\right).
\]
Expand the cubic difference:
\[
(q_0+\delta q)^3-q_0^3
=
3q_0^2\delta q+3q_0(\delta q)^2+(\delta q)^3.
\]
Hence
\[
R_{22}
=
\frac1{12\pi}
\left(
\delta q''
+
6q_0^2\delta q
+
6q_0(\delta q)^2
+
2(\delta q)^3
\right)(t_\pm).
\]

Now estimate each term.

First,
\[
|\delta q''(t_\pm)|\ll S_2.
\]

Second, using $|q_0(t_\pm)|\asymp Q$ and $|\delta q(t_\pm)|\ll L_I^2S_2\ll S_2/Q^2$,
\[
|q_0(t_\pm)^2\delta q(t_\pm)|
\ll
Q^2\cdot \frac{S_2}{Q^2}
=
S_2.
\]

Third,
\[
|q_0(t_\pm)(\delta q(t_\pm))^2|
\ll
Q\left(\frac{S_2}{Q^2}\right)^2
=
\frac{S_2^2}{Q^3}.
\]
Since from the tail-curvature theorem one already has $S_2=o(Q)$, for $Q$
sufficiently large this is $O(S_2/Q^2)$ and hence $\ll S_2$.

Finally,
\[
|(\delta q(t_\pm))^3|
\ll
\left(\frac{S_2}{Q^2}\right)^3
=
\frac{S_2^3}{Q^6}
\ll
S_2.
\]

Therefore
\[
|R_{22}|\ll S_2.
\]

\medskip
\noindent
Collecting the four entrywise bounds, we obtain
\[
|R_{11}|,\ |R_{12}|,\ |R_{21}|,\ |R_{22}|
\ll S_2.
\]
Since all norms on $2\times 2$ matrices are equivalent,
\[
\|R_{m,\pm}(s)\|\ll S_2,
\]
uniformly on the microscopic disk. This proves the proposition.
\end{proof}

\begin{proposition}[Entry-by-entry perturbation estimate for the corrected mixed block]
\label{prop:cross-block-perturbation}
Let
\[
q(t)=q_0(t)+\delta q(t),
\qquad
\Phi(t)=\Phi_0(t)+\delta\Phi(t),
\]
where $q_0$ is the baseline local-model-plus-affine-jet phase derivative and
$\delta q$ is the corrected remainder. Write
\[
t_\pm=m\pm \frac{s}{2},
\qquad
\Delta(s):=\Phi(t_-)-\Phi(t_+),
\qquad
\Delta_0(s):=\Phi_0(t_-)-\Phi_0(t_+),
\]
and set
\[
\delta\Delta(s):=\Delta(s)-\Delta_0(s).
\]

Let $N_m^{\corr}(s)$ be the corrected finite-$s$ mixed block built from
$q,\Delta$, and let $N_m^{(0)}(s)$ be the corresponding block built from
$q_0,\Delta_0$. Define
\[
\widetilde R_m(s):=N_m^{\corr}(s)-N_m^{(0)}(s).
\]

Assume that on the microscopic window one has
\[
|\delta q(t_\pm)|\ll L_I^2S_2,
\qquad
|\delta q'(t_\pm)|\ll L_I S_2,
\qquad
|\delta q''(t_\pm)|\ll S_2,
\]
and
\[
|\delta\Delta(s)|\ll L_I^3S_2,
\]
with
\[
L_I\asymp \frac{x_0}{Q},
\qquad
|q_0(t_\pm)|\asymp Q,
\qquad
|s|\asymp L_I.
\]
Then
\[
\boxed{\ \|\widetilde R_m(s)\|\ll S_2\ }
\]
uniformly for $|s|<\rho_0/Q$.
\end{proposition}

\begin{proof}
Write the corrected mixed block entrywise as
\[
N_m^{\corr}(s)=\frac1\pi
\begin{pmatrix}
-\dfrac{2\sin\Delta}{s}
&
\dfrac{\sin\Delta+q(t_+)s\cos\Delta}{s^2}
\\[2ex]
-\dfrac{\sin\Delta+q(t_-)s\cos\Delta}{s^2}
&
\dfrac{(q(t_-)+q(t_+))s\cos\Delta+(2-q(t_-)q(t_+)s^2)\sin\Delta}{2s^3}
\end{pmatrix},
\]
and analogously for $N_m^{(0)}(s)$ with $(q_0,\Delta_0)$.

We estimate the perturbation entry by entry. Throughout, use
\[
\sin(\Delta_0+\delta\Delta)-\sin\Delta_0=O(|\delta\Delta|),
\qquad
\cos(\Delta_0+\delta\Delta)-\cos\Delta_0=O(|\delta\Delta|).
\]

\medskip
\noindent
\textbf{(1,1)-entry.}
\[
\widetilde R_{11}
=
-\frac{2}{\pi s}\bigl(\sin\Delta-\sin\Delta_0\bigr),
\]
hence
\[
|\widetilde R_{11}|
\ll
\frac{|\delta\Delta|}{|s|}
\ll
\frac{L_I^3S_2}{L_I}
\ll
L_I^2S_2
\ll
S_2.
\]

\medskip
\noindent
\textbf{(1,2)-entry.}
The numerator perturbation is
\[
(\sin\Delta-\sin\Delta_0)
+
\bigl(q(t_+)-q_0(t_+)\bigr)s\cos\Delta
+
q_0(t_+)s(\cos\Delta-\cos\Delta_0).
\]
Therefore
\[
|\widetilde R_{12}|
\ll
\frac{|\delta\Delta|+|\delta q(t_+)|\,|s|+|q_0(t_+)|\,|s|\,|\delta\Delta|}{|s|^2}.
\]
Using
\[
|\delta\Delta|\ll L_I^3S_2,\quad
|\delta q(t_+)|\ll L_I^2S_2,\quad
|q_0(t_+)|\asymp Q,\quad
|s|\asymp L_I\asymp Q^{-1},
\]
we get
\[
|\widetilde R_{12}|\ll \frac{L_I^3S_2}{L_I^2}
+\frac{L_I^3S_2}{L_I^2}
+\frac{Q\,L_I^4S_2}{L_I^2}
\ll
\frac{S_2}{Q}
\ll S_2.
\]
The same argument gives
\[
|\widetilde R_{21}|\ll S_2.
\]

\medskip
\noindent
\textbf{(2,2)-entry.}
Write the numerator as
\[
F(q_-,q_+,\Delta;s):=
(q_-+q_+)s\cos\Delta+(2-q_-q_+s^2)\sin\Delta.
\]
Then
\[
\widetilde R_{22}
=
\frac{F(q(t_-),q(t_+),\Delta;s)-F(q_0(t_-),q_0(t_+),\Delta_0;s)}{2\pi s^3}.
\]
A first-order expansion of $F$ in $(q_-,q_+,\Delta)$ gives
\[
|F-F_0|
\ll
|s|\,|\delta q(t_-)|+|s|\,|\delta q(t_+)|
+|s|^2Q\,|\delta q(t_-)|+|s|^2Q\,|\delta q(t_+)|
+\bigl(|s|Q+1\bigr)|\delta\Delta|.
\]
Using
\[
|\delta q(t_\pm)|\ll L_I^2S_2,\qquad
|\delta\Delta|\ll L_I^3S_2,\qquad
|s|\asymp L_I\asymp Q^{-1},
\]
we obtain
\[
|F-F_0|\ll L_I^3S_2.
\]
Hence
\[
|\widetilde R_{22}|
\ll
\frac{L_I^3S_2}{|s|^3}
\ll
S_2.
\]

Collecting the four entrywise bounds and using equivalence of norms on
$2\times2$ matrices gives
\[
\|\widetilde R_m(s)\|\ll S_2,
\]
uniformly on the microscopic disk.
\end{proof}

\begin{lemma}[Baseline same-point variation]
\label{lem:baseline-variation}
Let $G_{m,\pm}^{(0)}(s)$ be the corrected same-point baseline blocks built
from the explicit local model together with the affine jet $J_1[g_{\sm}]$.
Then, on the microscopic disk $|s|<\rho_0/Q$,
\[
\sup_{|s|<\rho_0/Q}
\|G_{m,\pm}^{(0)}(s)-G_{m,\pm}^{(0)}(0)\|
\ll \rho_0 Q.
\]
\end{lemma}

\begin{proof}
Differentiate the explicit same-point block formula entrywise:
\[
G_{m,\pm}^{(0)}(s)=\frac1\pi
\begin{pmatrix}
2q_0(t_\pm) & \frac12 q_0'(t_\pm)\\[1ex]
\frac12 q_0'(t_\pm) & \frac1{12}\bigl(q_0''(t_\pm)+2q_0(t_\pm)^3\bigr)
\end{pmatrix},
\qquad t_\pm=m\pm \frac{s}{2}.
\]
By the chain rule,
\[
\partial_s G_{m,\pm}^{(0)}(s)
=
\pm \frac1{2\pi}
\begin{pmatrix}
2q_0'(t_\pm) & \frac12 q_0''(t_\pm)\\[1ex]
\frac12 q_0''(t_\pm) &
\frac1{12}\bigl(q_0'''(t_\pm)+6q_0(t_\pm)^2q_0'(t_\pm)\bigr)
\end{pmatrix}.
\]
For the explicit local model plus affine jet, one has the baseline local scales
\[
|q_0(t_\pm)|\ll Q,\qquad
|q_0'(t_\pm)|\ll Q^2,\qquad
|q_0''(t_\pm)|\ll Q^3,\qquad
|q_0'''(t_\pm)|\ll Q^4,
\]
uniformly on the admissible midpoint class. The cubic combination
$
q_0'''(t_\pm)+6q_0(t_\pm)^2q_0'(t_\pm)
$
therefore has size $O(Q^2)$ in the current normalization, and hence
\[
\sup_{|s|<\rho_0/Q}\|\partial_s G_{m,\pm}^{(0)}(s)\|\ll Q^2.
\]
Integrating along the segment from $0$ to $s$ gives
\[
\|G_{m,\pm}^{(0)}(s)-G_{m,\pm}^{(0)}(0)\|
\le
|s|\sup_{|u|<\rho_0/Q}\|\partial_u G_{m,\pm}^{(0)}(u)\|
\ll
\frac{\rho_0}{Q}\cdot Q^2
=
\rho_0 Q.
\]
\end{proof}

\begin{proposition}[Corrected finite-$s$ holomorphic whitening]
\label{prop:corrected-whitening}
The corrected same-point and mixed blocks are holomorphic on
\[
|s|<\rho_0/Q.
\]
Moreover, after shrinking $\rho_0$ if necessary, the same-point blocks
\[
G_{m,\pm}^{\corr}(s)
\]
remain uniformly positive definite on this disk, and the inverse square roots
\[
G_{m,\pm}^{\corr}(s)^{-1/2}
\]
are holomorphic there. Consequently
\[
\widehat\Omega_\zeta^{\corr}(s;m)
=
G_{m,-}^{\corr}(s)^{-1/2}
N_m^{\corr}(s)
G_{m,+}^{\corr}(s)^{-1/2}
\]
is holomorphic on $|s|<\rho_0/Q$.
\end{proposition}

\begin{proof}
Holomorphy of the omitted smooth part and its required $t$-derivatives follows
from Proposition~\ref{prop:denominator-comparability}. The only apparent poles
at $s=0$ in the corrected mixed block come from powers of $1/s$, $1/s^2$,
$1/s^3$, and these are removable because $\Delta_m(s)$ is odd in $s$,
$\sin \Delta_m(s)=O(s)$, and the off-diagonal and $(2,2)$ numerators vanish
to orders $2$ and $3$, respectively. Thus both $G_{m,\pm}^{\corr}(s)$ and
$N_m^{\corr}(s)$ are holomorphic on the disk.

By Proposition~\ref{prop:block-perturbation} and the tail-curvature theorem,
\[
\sup_{|s|<\rho_0/Q}\|R_{m,\pm}(s)\|=o(Q).
\]
By Lemma~\ref{lem:baseline-variation},
\[
\sup_{|s|<\rho_0/Q}
\|G_{m,\pm}^{(0)}(s)-G_{m,\pm}^{(0)}(0)\|
\ll \rho_0 Q.
\]
At $s=0$, the baseline jet-limit block has spectral gap
\[
\lambda_{\min}(G_{m,\pm}^{(0)}(0))\gg Q
\]
on the admissible midpoint class. Therefore
\[
\|G_{m,\pm}^{\corr}(s)-G_{m,\pm}^{(0)}(0)\|
\le
\|G_{m,\pm}^{(0)}(s)-G_{m,\pm}^{(0)}(0)\|
+
\|R_{m,\pm}(s)\|
\ll
\rho_0 Q + o(Q).
\]
For $\rho_0$ sufficiently small and $m$ sufficiently large, this is smaller
than half the baseline spectral gap, hence
\[
\lambda_{\min}(G_{m,\pm}^{\corr}(s))\gg Q
\]
uniformly on the disk. Holomorphic functional calculus then gives the
holomorphic inverse square roots, and the whitened block is holomorphic as a
product of holomorphic factors.
\end{proof}

\subsection{Fixed anomaly core and cutoff-dependent auxiliary bookkeeping}
\label{subsec:fixed-core}

To separate the anomaly-carrying local model from the cutoff-dependent tail
bookkeeping, we fix once and for all a finite local core
\[
\mathcal C
\]
containing the candidate off-line anomaly (for example, the designated off-line
pair together with any symmetry-forced companions needed for the local model).

For each cutoff $R\ge R_0$, we then decompose the zero contribution into:
\begin{enumerate}
\item the fixed core contribution $q_{\core}(t)$, coming from the zeros in
$\mathcal C$;
\item the auxiliary near contribution $q_{\aux,R}(t)$, coming from the zeros
outside $\mathcal C$ but inside the near shell determined by $R$;
\item the far contribution $T_{\far,R}(t)$, coming from the zeros outside the
near shell;
\item the background contribution $B_{\bg}(t)$;
\item the estimation term $E_{\est,R}(t)$.
\end{enumerate}

Thus the full phase derivative is decomposed as
\[
q(t)=q_{\core}(t)+q_{\aux,R}(t)+g_{\sm,R}(t)+E_{\est,R}(t),
\]
where
\[
g_{\sm,R}(t):=B_{\bg}(t)+T_{\far,R}(t).
\]

The local toy comparison is attached only to the fixed core $\mathcal C$. The
auxiliary near part $q_{\aux,R}$, the far part $T_{\far,R}$, and the
estimation term $E_{\est,R}$ are absorbed into the corrected zeta-side scalar
defined below.

\begin{proposition}[Corrected local deformation decomposition with fixed core]
\label{prop:corrected-local-decomposition}
Fix the anomaly-carrying core $\mathcal C$ as above, and for each cutoff
$R\ge R_0$ write
\[
q(t)=q_{\core}(t)+q_{\aux,R}(t)+g_{\sm,R}(t)+E_{\est,R}(t),
\qquad
g_{\sm,R}(t)=B_{\bg}(t)+T_{\far,R}(t).
\]
Let $J_1[g_{\sm,R}]$ denote the affine jet of $g_{\sm,R}$ at the midpoint,
and let $\Phi_{\corr,R}$ be the corrected phase obtained from this
decomposition.

Using $\Phi_{\corr,R}$, define the corrected same-point blocks
\[
G_{m,\pm}^{\corr,R}(s),
\]
the corrected mixed block
\[
N_m^{\corr,R}(s),
\]
and the corrected whitened block
\[
\widehat\Omega_\zeta^{\corr,R}(s;m)
=
G_{m,-}^{\corr,R}(s)^{-1/2}
N_m^{\corr,R}(s)
G_{m,+}^{\corr,R}(s)^{-1/2}.
\]

Let
\[
\widehat\Omega_{\bg}^{\corr,R}(s;m)
\]
denote the corresponding corrected whitened block at background value-channel
parameter $S(m)=0$, and define the corrected local deformation
\[
\Delta_\zeta(m,\sigma)
:=
\widehat\Omega_\zeta^{\corr,R}(s;m)-\widehat\Omega_{\bg}^{\corr,R}(s;m),
\qquad s=\sigma.
\]

Then the corrected local deformation admits the decomposition
\[
\boxed{
\Delta_\zeta(m,\sigma)
=
S(m)\,A_{\val}(m,\sigma)+R_\zeta(m,\sigma),
}
\]
where $A_{\val}(m,\sigma)$ is the pure value-channel derivative of the
jet-limit block associated with the fixed core $\mathcal C$, and where
\[
\boxed{
R_\zeta(m,\sigma)
:=
\Delta_\zeta(m,\sigma)-S(m)\,A_{\val}(m,\sigma).
}
\]
Thus, by definition, $R_\zeta(m,\sigma)$ is exactly the non-leading part of
the corrected local deformation after the explicit value-channel linear term
is removed. In particular, $R_\zeta(m,\sigma)$ is internal to the same full
corrected cutoff-dependent object
\[
H_{m,R}(s):=\Phi_K\!\bigl(\widehat\Omega_\zeta^{\corr,R}(s;m)\bigr),
\]
so once the corrected odd holomorphic package is established for $H_{m,R}(s)$,
no external tail/error term survives outside the resulting $N$-point scalar.
\end{proposition}

\begin{proof}
Consider the analytic whitening map
\[
\mathcal W(G_-,N,G_+):=G_-^{-1/2}NG_+^{-1/2}.
\]
By the corrected finite-$s$ holomorphic-whitening theorem, this map is
holomorphic on the open set where $G_\pm$ remain uniformly positive definite.

For each cutoff $R$, the corrected blocks are built from the decomposition
\[
q=q_{\core}+q_{\aux,R}+g_{\sm,R}+E_{\est,R}.
\]
Thus every auxiliary near contribution from $q_{\aux,R}$, every omitted-zero
tail contribution from $T_{\far,R}$, every affine-jet correction from
$J_1[g_{\sm,R}]$, and every explicit estimation term $E_{\est,R}$ is already
absorbed into the corrected phase $\Phi_{\corr,R}$, hence into the corrected
same-point blocks, the corrected mixed block, the corrected whitened block,
and therefore into the corrected transverse scalar
\[
H_{m,R}(s)=\Phi_K\!\bigl(\widehat\Omega_\zeta^{\corr,R}(s;m)\bigr).
\]

Now expand the corrected whitened block around the background value-channel
parameter $S(m)=0$. By analyticity of $\mathcal W$,
\[
\widehat\Omega_\zeta^{\corr,R}(s;m)
=
\widehat\Omega_{\bg}^{\corr,R}(s;m)
+
S(m)\,D\mathcal W_{(\bg)}[\dot G_-,\dot N,\dot G_+]
+
R_\zeta(m,\sigma),
\]
where the first-order coefficient
\[
D\mathcal W_{(\bg)}[\dot G_-,\dot N,\dot G_+]
\]
is exactly the pure value-channel derivative $A_{\val}(m,\sigma)$ associated
with the fixed core $\mathcal C$. We therefore define
\[
R_\zeta(m,\sigma)
:=
\Delta_\zeta(m,\sigma)-S(m)\,A_{\val}(m,\sigma),
\]
so that $R_\zeta(m,\sigma)$ is, by definition, the full non-leading remainder
inside the same corrected cutoff-dependent object.

Subtracting the background block gives
\[
\Delta_\zeta(m,\sigma)
=
S(m)\,A_{\val}(m,\sigma)+R_\zeta(m,\sigma),
\]
as claimed.
\end{proof}

\begin{proposition}[Core-stable cutoff compatibility]
\label{prop:cutoff-compatibility}
Fix the anomaly-carrying core $\mathcal C$, and for each cutoff $R\ge R_0$
define the corrected scalar
\[
H_{m,R}(s):=\Phi_K\!\bigl(\widehat\Omega_\zeta^{\corr,R}(s;m)\bigr)
\]
from the decomposition
\[
q=q_{\core}+q_{\aux,R}+g_{\sm,R}+E_{\est,R}.
\]

Then enlarging $R$ changes only the auxiliary near part $q_{\aux,R}$ and the
far part $T_{\far,R}$, while leaving the anomaly-carrying core
$q_{\core}$ unchanged. Consequently:
\begin{enumerate}
\item the toy anomaly matrix $M(d)$ and the calibration matrix
$A_{\val}(m,\sigma)$ depend only on the fixed core $\mathcal C$, not on
the cutoff $R$;
\item all $R$-dependent auxiliary, tail, jet-correction, and estimation
contributions are internal to the full corrected scalar $H_{m,R}(s)$;
\item if the corrected odd holomorphic package is proved for $H_{m,R}(s)$,
then the $N$-point odd-moment cancellation theorem applies directly to the
full corrected scalar for that cutoff, with no external remainder term.
\end{enumerate}

In particular, the cutoff $R$ may be chosen for quantitative convenience in
the curvature/tail estimates without changing the fixed toy comparison
attached to the core $\mathcal C$.
\end{proposition}

\begin{proof}
Let $R_1<R_2$. Then the only change in the decomposition
\[
q=q_{\core}+q_{\aux,R}+g_{\sm,R}+E_{\est,R}
\]
is that some non-core zeros move from the far part into the auxiliary near
part:
\[
q_{\aux,R_2}-q_{\aux,R_1}
=
-\bigl(T_{\far,R_2}-T_{\far,R_1}\bigr).
\]
The core contribution $q_{\core}$ is unchanged by construction.

Thus the toy anomaly matrix $M(d)$, which is attached only to the fixed core,
is unchanged, and the value-channel derivative $A_{\val}(m,\sigma)$ associated
with that core is unchanged as well. All cutoff dependence is confined to the
auxiliary near contribution, the far tail, the affine-jet correction built
from $g_{\sm,R}$, and the explicit estimation term $E_{\est,R}$. These enter
only through the corrected phase $\Phi_{\corr,R}$, hence only through the full
corrected scalar $H_{m,R}(s)$.

Therefore the corrected odd holomorphic package, once established for
$H_{m,R}(s)$, applies to the complete cutoff-dependent object. The resulting
$N$-point scalar already contains all auxiliary/tail/error effects for that
cutoff, and no external remainder survives. Since the fixed core is
unchanged, the toy comparison is stable under enlarging $R$. This proves the
claim.
\end{proof}

\begin{remark}[Cutoff choice]
\label{rem:cutoff-choice}
Proposition~\ref{prop:cutoff-compatibility} is a structural statement. It
does not assert any particular optimal choice of $R$; it only says that once
the core $\mathcal C$ is fixed, the cutoff may be adjusted in the perturbative
bookkeeping without altering the underlying toy comparison.
\end{remark}

\begin{proposition}[Whitened mixed-block transfer with preserved $Q^{-3}$ scale]
\label{prop:whitened-mixed-transfer}
Let
\[
\mathcal W(G_-,N,G_+):=G_-^{-1/2}NG_+^{-1/2}.
\]
Write the corrected finite-$s$ blocks as
\[
G_\pm^{\corr}(s)=G_\pm^{(0)}(s)+R_\pm(s),
\qquad
N^{\corr}(s)=N^{(0)}(s)+\widetilde R(s),
\]
and set
\[
\widehat\Omega_\zeta^{\corr}(s;m):=\mathcal W(G_-^{\corr},N^{\corr},G_+^{\corr}),
\qquad
\widehat\Omega_\zeta^{(0)}(s;m):=\mathcal W(G_-^{(0)},N^{(0)},G_+^{(0)}).
\]

Then
\[
\widehat\Omega_\zeta^{\corr}(s;m)-\widehat\Omega_\zeta^{(0)}(s;m)
=
\mathcal T_-(s;m)+\mathcal T_{\mix}(s;m)+\mathcal T_+(s;m)+\mathcal E_2(s;m),
\]
where
\[
\mathcal T_-(s;m):=X_-(s)N^{(0)}(s)G_+^{(0)}(s)^{-1/2},
\]
\[
\mathcal T_{\mix}(s;m):=G_-^{(0)}(s)^{-1/2}\widetilde R(s)G_+^{(0)}(s)^{-1/2},
\]
\[
\mathcal T_+(s;m):=G_-^{(0)}(s)^{-1/2}N^{(0)}(s)X_+(s),
\]
with
\[
X_\pm(s):=D(G^{-1/2})_{G_\pm^{(0)}(s)}[R_\pm(s)],
\]
and $\mathcal E_2$ quadratic in $(R_-,\widetilde R,R_+)$.

Moreover,
\[
\Phi_K(\mathcal T_-(s;m))\ll \frac{S_2}{Q^3},
\qquad
\Phi_K(\mathcal T_{\mix}(s;m))\ll \frac{S_2}{Q^3},
\qquad
\Phi_K(\mathcal T_+(s;m))\ll \frac{S_2}{Q^3},
\]
and
\[
\Phi_K(\mathcal E_2(s;m))\ll \frac{S_2^2}{Q^3}.
\]
Hence
\[
\boxed{
\Phi_K\!\bigl(\widehat\Omega_\zeta^{\corr}(s;m)-\widehat\Omega_\zeta^{(0)}(s;m)\bigr)
\ll
\frac{S_2}{Q^3}.
}
\]
\end{proposition}

\begin{proof}
We first collect the scales used below from results established earlier. The
same-point perturbation bound $\|R_\pm(s)\|\ll S_2$ is
Proposition~\ref{prop:block-perturbation}; the mixed-block perturbation bound
$\|\widetilde R(s)\|\ll S_2$ is
Proposition~\ref{prop:cross-block-perturbation}; and the smallness
$S_2=o(Q)$ is Theorem~\ref{thm:tail-curvature}. The baseline entry scales
\[
G_\pm^{(0)}(s)=
\begin{pmatrix} O(Q) & O(1)\\ O(1) & O(Q^3) \end{pmatrix},
\qquad
G_\pm^{(0)}(s)^{-1/2}=
\begin{pmatrix}
O(Q^{-1/2}) & O(Q^{-3/2})\\
O(Q^{-3/2}) & O(Q^{-3/2})
\end{pmatrix},
\]
\[
N^{(0)}(s)=
\begin{pmatrix} O(Q) & O(Q^2)\\ O(Q^2) & O(Q^3) \end{pmatrix}
\]
are read off directly from the explicit same-point formula \eqref{eq:Gpm} and
the mixed-block formula \eqref{eq:Nm}, given the baseline local scales of
$q_0, q_0', q_0''$ stated in the proof of
Lemma~\ref{lem:baseline-variation}. Uniform positive-definiteness of
$G_\pm^{(0)}(s)$ on the microscopic disk is
Proposition~\ref{prop:corrected-whitening}.

The whitening map admits the first-order expansion
\[
\widehat\Omega_\zeta^{\corr}-\widehat\Omega_\zeta^{(0)}
=
D\mathcal W_{(0)}[R_-,\widetilde R,R_+]+\mathcal E_2,
\]
where
\[
D\mathcal W_{(0)}[R_-,\widetilde R,R_+]
=
X_-N^{(0)}G_+^{(0)-1/2}
+
G_-^{(0)-1/2}\widetilde R\,G_+^{(0)-1/2}
+
G_-^{(0)-1/2}N^{(0)}X_+.
\]

We estimate the three first-order pieces separately.

\medskip
\noindent
\textbf{Step 1: the mixed-block term.}
By Proposition~\ref{prop:cross-block-perturbation} (entrywise form),
\[
\widetilde R_{11}\ll \frac{S_2}{Q^2},
\qquad
\widetilde R_{12},\widetilde R_{21}\ll \frac{S_2}{Q},
\qquad
\widetilde R_{22}\ll S_2.
\]
Combining with the baseline inverse-square-root scales above, the off-diagonal
entries of
\[
\mathcal T_{\mix}:=G_-^{(0)-1/2}\widetilde R\,G_+^{(0)-1/2}
\]
are all $O(S_2/Q^3)$. Hence
\[
\Phi_K(\mathcal T_{\mix})\ll \frac{S_2}{Q^3}.
\]

\medskip
\noindent
\textbf{Step 2: the same-point Fréchet derivative.}
We prove the corresponding estimate for the left same-point term
\[
\mathcal T_-:=X_-N^{(0)}G_+^{(0)-1/2},
\qquad
X_-:=D(G^{-1/2})_{G_-^{(0)}}[R_-].
\]
The proof for $\mathcal T_+$ is identical.

Write one baseline same-point block in the form
\[
G_0=
\begin{pmatrix}
a & b\\
b & c
\end{pmatrix},
\qquad
a\asymp Q,\quad b=O(1),\quad c\asymp Q^3,
\]
and one perturbation as
\[
R=
\begin{pmatrix}
r_{11} & r_{12}\\
r_{12} & r_{22}
\end{pmatrix},
\qquad
r_{11}\ll \frac{S_2}{Q^2},\quad
r_{12}\ll \frac{S_2}{Q},\quad
r_{22}\ll S_2.
\]

Diagonalize
\[
G_0=U\Lambda U^\top,
\qquad
\Lambda=\operatorname{diag}(\lambda_1,\lambda_2),
\]
with $U$ orthogonal. Since
\[
a\asymp Q,\qquad c\asymp Q^3,\qquad b=O(1),
\]
we have
\[
\lambda_1\asymp Q,\qquad \lambda_2\asymp Q^3.
\]
Moreover, if $U$ is a rotation by angle $\theta$, then
\[
\tan(2\theta)=\frac{2b}{c-a}=O(Q^{-3}),
\]
hence
\[
\theta=O(Q^{-3}),
\qquad
U=I+O(Q^{-3}).
\]

Set
\[
H:=U^\top R\,U.
\]
A direct computation gives
\[
H_{11}=r_{11}+2\theta r_{12}+O(\theta^2 r_{22}),
\]
\[
H_{12}=r_{12}+\theta(r_{22}-r_{11})+O(\theta^2 r_{12}),
\]
\[
H_{22}=r_{22}-2\theta r_{12}+O(\theta^2 r_{11}).
\]
Using
\[
\theta=O(Q^{-3}),
\qquad
r_{11}\ll \frac{S_2}{Q^2},
\qquad
r_{12}\ll \frac{S_2}{Q},
\qquad
r_{22}\ll S_2,
\]
we obtain
\[
H_{11}\ll \frac{S_2}{Q^2},
\qquad
H_{12}\ll \frac{S_2}{Q},
\qquad
H_{22}\ll S_2.
\]

Now let $f(x)=x^{-1/2}$. For a diagonal matrix $\Lambda$, the Fréchet
derivative is given entrywise by the Löwner formula
\[
Df_\Lambda(H)_{ij}=f^{[1]}(\lambda_i,\lambda_j)H_{ij},
\]
where
\[
f^{[1]}(u,v)=
\begin{cases}
\dfrac{u^{-1/2}-v^{-1/2}}{u-v}, & u\neq v,\\[1.2ex]
-\dfrac{1}{2u^{3/2}}, & u=v.
\end{cases}
\]
Since
\[
\lambda_1\asymp Q,\qquad \lambda_2\asymp Q^3,
\]
we have
\[
f^{[1]}(\lambda_1,\lambda_1)\asymp Q^{-3/2},
\qquad
f^{[1]}(\lambda_2,\lambda_2)\asymp Q^{-9/2},
\qquad
f^{[1]}(\lambda_1,\lambda_2)\asymp Q^{-7/2}.
\]
Therefore, if
\[
Y:=Df_\Lambda(H),
\]
then
\[
Y_{11}\ll Q^{-3/2}\cdot \frac{S_2}{Q^2}
=\frac{S_2}{Q^{7/2}},
\]
\[
Y_{12}\ll Q^{-7/2}\cdot \frac{S_2}{Q}
=\frac{S_2}{Q^{9/2}},
\]
\[
Y_{22}\ll Q^{-9/2}\cdot S_2
=\frac{S_2}{Q^{9/2}}.
\]

Since
\[
X=U\,Y\,U^\top
\]
and $U=I+O(Q^{-3})$, these scales persist in the original basis:
\[
X_{11}\ll \frac{S_2}{Q^{7/2}},
\qquad
X_{12},X_{21}\ll \frac{S_2}{Q^{9/2}},
\qquad
X_{22}\ll \frac{S_2}{Q^{9/2}}.
\]

\medskip
\noindent
\textbf{Step 3: the left same-point term $\mathcal T_-$.}
The baseline mixed-block scales are
\[
N^{(0)}=
\begin{pmatrix}
O(Q) & O(Q^2)\\
O(Q^2) & O(Q^3)
\end{pmatrix},
\qquad
G_+^{(0)-1/2}=
\begin{pmatrix}
O(Q^{-1/2}) & O(Q^{-3/2})\\
O(Q^{-3/2}) & O(Q^{-3/2})
\end{pmatrix}.
\]
We estimate only the off-diagonal entries, since
\[
\Phi_K(Y)=Y_{12}+Y_{21}.
\]

For $(\mathcal T_-)_{12}$, every summand is $O(S_2/Q^3)$. For example,
\[
X_{11}N_{12}(G_+^{-1/2})_{22}
\ll
\frac{S_2}{Q^{7/2}}\cdot Q^2\cdot Q^{-3/2}
=
\frac{S_2}{Q^3},
\]
and
\[
X_{12}N_{22}(G_+^{-1/2})_{22}
\ll
\frac{S_2}{Q^{9/2}}\cdot Q^3\cdot Q^{-3/2}
=
\frac{S_2}{Q^3}.
\]
All other terms are smaller. Thus
\[
(\mathcal T_-)_{12}\ll \frac{S_2}{Q^3}.
\]
The same argument gives
\[
(\mathcal T_-)_{21}\ll \frac{S_2}{Q^3},
\]
and hence
\[
\Phi_K(\mathcal T_-)\ll \frac{S_2}{Q^3}.
\]

By symmetry, the same estimate holds for the right same-point term:
\[
\Phi_K(\mathcal T_+)\ll \frac{S_2}{Q^3}.
\]

\medskip
\noindent
\textbf{Step 4: the quadratic remainder.}
Every term in $\mathcal E_2$ contains at least two perturbation factors from
$(R_-,\widetilde R,R_+)$, and whitening contributes only polynomial
$Q$-weights. Therefore
\[
\Phi_K(\mathcal E_2)\ll \frac{S_2^2}{Q^3}.
\]
By Theorem~\ref{thm:tail-curvature}, $S_2=o(Q)$, so this is lower order than
$S_2/Q^3$.

Combining the estimates for $\mathcal T_-$, $\mathcal T_{\mix}$,
$\mathcal T_+$, and $\mathcal E_2$ gives
\[
\Phi_K\!\bigl(\widehat\Omega_\zeta^{\corr}(s;m)-\widehat\Omega_\zeta^{(0)}(s;m)\bigr)
\ll
\frac{S_2}{Q^3},
\]
as claimed.
\end{proof}

\begin{corollary}[Sharpened calibration remainder estimate]
\label{cor:sharpened-calibration-remainder}
Fix a compact interval
\[
D\subset(0,\infty),
\]
and work in the fixed-scaled regime
\[
0<x\le x_0(D),\qquad d\in D,\qquad x=B_\zeta(m)\sigma,\qquad |\sigma|\asymp Q^{-1}.
\]
Write the corrected local deformation in the form
\[
\Delta_\zeta(m,\sigma)=S(m)A_{\val}(x)+R_\zeta(m,\sigma).
\]
Then
\[
\bigl|\langle A_{\val}(x),R_\zeta(m,\sigma)\rangle_\F\bigr|
\ll_D
\frac{x}{B_\zeta(m)}\frac{S_2}{Q^3},
\]
and hence
\[
\boxed{
|\Psi_{x,d}(R_\zeta(m,\sigma))|
\ll_D
\frac{S_2}{Q^3}.
}
\]
In particular, if
\[
S_2=o\!\bigl(xQ^2S(m)\bigr),
\]
then
\[
\Psi_{x,d}(R_\zeta(m,\sigma))
=
o\!\left(\frac{x}{B_\zeta(m)}S(m)\right).
\]
\end{corollary}

\begin{proof}
By Proposition~\ref{prop:corrected-local-decomposition}, after subtraction of
the explicit leading value-channel term $S(m)A_{\val}(x)$, the remaining
contribution $R_\zeta(m,\sigma)$ is internal to the same corrected
transferred package controlled by the whitening perturbation analysis.

We therefore use the entry scales already established in the proof of
Proposition~\ref{prop:whitened-mixed-transfer}. First, from
Proposition~\ref{prop:cross-block-perturbation} and the baseline
inverse-square-root scales,
\[
G_\pm^{(0)}(s)^{-1/2}
=
\begin{pmatrix}
O(Q^{-1/2}) & O(Q^{-3/2})\\
O(Q^{-3/2}) & O(Q^{-3/2})
\end{pmatrix},
\qquad
\widetilde R(s)=
\begin{pmatrix}
O(S_2/Q^2) & O(S_2/Q)\\
O(S_2/Q) & O(S_2)
\end{pmatrix},
\]
one gets
\[
\mathcal T_{\mix}(s;m)
=
G_-^{(0)}(s)^{-1/2}\widetilde R(s)G_+^{(0)}(s)^{-1/2}
=
\begin{pmatrix}
O(S_2/Q^3) & O(S_2/Q^3)\\
O(S_2/Q^3) & O(S_2/Q^3)
\end{pmatrix}.
\]

Next, from the same proof,
\[
X_\pm(s)=
\begin{pmatrix}
O(S_2/Q^{7/2}) & O(S_2/Q^{9/2})\\
O(S_2/Q^{9/2}) & O(S_2/Q^{9/2})
\end{pmatrix},
\qquad
N^{(0)}(s)=
\begin{pmatrix}
O(Q) & O(Q^2)\\
O(Q^2) & O(Q^3)
\end{pmatrix},
\]
and multiplying once more by $G_\pm^{(0)}(s)^{-1/2}$ gives
\[
\mathcal T_-(s;m)=
\begin{pmatrix}
O(S_2/Q^3) & O(S_2/Q^3)\\
O(S_2/Q^3) & O(S_2/Q^3)
\end{pmatrix},
\qquad
\mathcal T_+(s;m)=
\begin{pmatrix}
O(S_2/Q^3) & O(S_2/Q^3)\\
O(S_2/Q^3) & O(S_2/Q^3)
\end{pmatrix}.
\]

Finally, every entry of the quadratic remainder $\mathcal E_2(s;m)$ is
\[
O(S_2^2/Q^3),
\]
which is lower order because $S_2=o(Q)$ by
Theorem~\ref{thm:tail-curvature}.

Thus the part of the corrected deformation seen after removal of the leading
value-channel term satisfies
\[
R_\zeta(m,\sigma)
=
\begin{pmatrix}
O(S_2/Q^3) & O(S_2/Q^3)\\
O(S_2/Q^3) & O(S_2/Q^3)
\end{pmatrix}
\]
up to lower-order $O(S_2^2/Q^3)$ terms.

Now use the explicit small-$x$ structure of the value-channel derivative:
by Lemma~\ref{lem:Aval-small-x},
\[
A_{\val}(x)=\frac1{B_\zeta(m)}
\begin{pmatrix}
O(x^2) & O(x)\\
O(x) & O(x^2)
\end{pmatrix}.
\]
Therefore
\[
\bigl|\langle A_{\val}(x),R_\zeta(m,\sigma)\rangle_\F\bigr|
\ll
\frac{x}{B_\zeta(m)}\frac{S_2}{Q^3}.
\]

By Proposition~\ref{prop:pairing},
\[
\bigl|\langle A_{\val}(x),M(d)\rangle_\F\bigr|
\asymp_D
\frac{x}{B_\zeta(m)},
\]
so dividing gives
\[
|\Psi_{x,d}(R_\zeta(m,\sigma))|
\ll_D
\frac{S_2}{Q^3}.
\]

Finally, since $B_\zeta(m)\asymp Q$,
\[
\frac{S_2/Q^3}{(x/B_\zeta(m))S(m)}
\asymp
\frac{S_2}{xQ^2S(m)}\to 0
\]
under the hypothesis $S_2=o(xQ^2S(m))$. This proves the final claim.
\end{proof}

\section{Odd holomorphic transverse channel}

Define the corrected transverse scalar
\begin{equation}\label{eq:Hm-def}
H_m(s):=\Phi_K\!\bigl(\widehat\Omega_\zeta^{\corr}(s;m)\bigr).
\end{equation}
By symmetry under $s\mapsto -s$, this is odd. By
Corollary~\ref{cor:PhiK-Aval-zero}, the leading value channel is
annihilated by $\Phi_K$, so only the curvature, tail, and corrected
finite-$s$ remainder channels contribute.

\begin{proposition}[Boundary estimate]
\label{prop:boundary-estimate}
On the microscopic boundary
\[
|s|=\rho_0/Q,
\]
one has
\[
|H_m(s)|\ll_{\rho_0}\frac{L(m)S(m)^2}{Q^3}.
\]
\end{proposition}

\begin{proof}
By Proposition~\ref{prop:corrected-local-decomposition}, the corrected local
deformation may be written in the form
\[
\Delta_\zeta(m,\sigma)
=
S(m)\,A_{\val}(m,\sigma)+R_\zeta(m,\sigma),
\]
where $R_\zeta(m,\sigma)$ is internal to the full corrected cutoff-dependent
scalar
\[
H_{m,R}(s):=\Phi_K\!\bigl(\widehat\Omega_\zeta^{\corr,R}(s;m)\bigr).
\]
By Proposition~\ref{prop:cutoff-compatibility}, all cutoff-dependent tail,
jet-correction, and estimation contributions are already absorbed into this
full corrected scalar, so no external remainder term is carried separately at
the boundary stage.

Applying $\Phi_K$ to the corrected local decomposition and using
Corollary~\ref{cor:PhiK-Aval-zero},
\[
\Phi_K(A_{\val}(m,\sigma))=0,
\]
gives
\[
H_{m,R}(s)
=
\Phi_K\!\bigl(\widehat\Omega_\zeta^{\corr,R}(s;m)\bigr)
=
\Phi_K\!\bigl(R_\zeta(m,\sigma)\bigr).
\]

Now compare the full corrected whitened block with its baseline counterpart.
By Proposition~\ref{prop:whitened-mixed-transfer},
\[
\Phi_K\!\bigl(\widehat\Omega_\zeta^{\corr,R}(s;m)-\widehat\Omega_\zeta^{(0),R}(s;m)\bigr)
\ll
\frac{S_2}{Q^3}
\qquad (|s|<\rho_0/Q).
\]
On the microscopic boundary $|s|=\rho_0/Q$, the tail-curvature theorem and
the corrected transfer bounds imply
\[
S_2\ll L(m)S(m)^2
\]
in the corrected regime. Therefore
\[
|H_{m,R}(s)|
\ll
\frac{L(m)S(m)^2}{Q^3}
\qquad (|s|=\rho_0/Q).
\]

Since the cutoff bookkeeping has already been absorbed into the full corrected
scalar $H_{m,R}(s)$, this is the boundary estimate for the complete corrected
object, not merely for a leading part plus an external tail/error remainder.
This proves the claim.
\end{proof}

\begin{proposition}[Odd microscopic expansion]
\label{prop:odd-expansion}
The function $H_m(s)$ is odd and holomorphic on $|s|<\rho_0/Q$, and therefore
admits an expansion
\[
H_m(z/Q^2)
=
\sum_{r\ge 0} c_{2r+1}(m)\frac{z^{2r+1}}{Q^{2r+4}},
\qquad |z|<\rho_0Q,
\]
with
\[
|c_{2r+1}(m)|\ll_{\rho_0}L(m)S(m)^2\,\rho_0^{-(2r+1)}.
\]
\end{proposition}

\begin{proof}
Oddness follows from symmetric placement and the jet-basis swap symmetry:
\[
\widehat\Omega_\zeta^{\corr}(-s;m)=J_0\,\widehat\Omega_\zeta^{\corr}(s;m)\,J_0,
\]
with
\[
\Phi_K(J_0XJ_0)=-\Phi_K(X).
\]
Holomorphy follows from Proposition~\ref{prop:corrected-whitening}, and the
coefficient bounds are immediate from Cauchy's estimate applied to
Proposition~\ref{prop:boundary-estimate}.
\end{proof}

\subsection{Coefficient-level corrected-whitening transfer and exterior odd-jet localization}
\label{subsec:differentiated-transfer}

\begin{lemma}[Baseline-scaled matrix gain for corrected whitening]
\label{lem:baseline-scaled-matrix-gain}
Let
\[
\mathcal W(G_-,N,G_+):=G_-^{-1/2}NG_+^{-1/2}.
\]
Assume the baseline same-point and mixed blocks satisfy
\[
G_\pm^{(0)}(s)=
\begin{pmatrix}
O(Q) & O(1)\\
O(1) & O(Q^3)
\end{pmatrix},
\qquad
N^{(0)}(s)=
\begin{pmatrix}
O(Q) & O(Q^2)\\
O(Q^2) & O(Q^3)
\end{pmatrix},
\]
uniformly on $|s|<\rho_0/Q$, and that $G_\pm^{(0)}(s)$ remain uniformly
positive definite there.

Set
\[
S_Q:=\begin{pmatrix}Q^{-1/2}&0\\0&Q^{-3/2}\end{pmatrix},
\qquad
\widehat G_\pm:=S_Q^{-1}G_\pm S_Q^{-1},
\qquad
\widehat N:=S_Q^{-1}NS_Q^{-1}.
\]
Then
\[
\widehat G_\pm^{(0)}(s)=
\begin{pmatrix}
O(1)&O(Q^{-1})\\
O(Q^{-1})&O(1)
\end{pmatrix},
\qquad
\widehat N^{(0)}(s)=
\begin{pmatrix}
O(1)&O(1)\\
O(1)&O(1)
\end{pmatrix},
\]
and
\[
\mathcal W(G_-,N,G_+)
=
S_Q\,
\mathcal W(\widehat G_-,\widehat N,\widehat G_+)\,
S_Q.
\]

In particular, every off-diagonal entry of $\mathcal W(G_-,N,G_+)$ carries an
explicit matrix-level factor $Q^{-2}$ relative to the corresponding rescaled
whitened block. Any additional $Q$-power appearing in the corrected transverse
scalar must therefore come from the scalar normalization used in the definition
of that scalar, not from the matrix identity itself.
\end{lemma}

\begin{proof}
The displayed bounds for $\widehat G_\pm^{(0)}$ and $\widehat N^{(0)}$ are
immediate from the definitions and the baseline scales. Since
\[
G_\pm=S_Q\,\widehat G_\pm\,S_Q,
\]
functional calculus gives
\[
G_\pm^{-1/2}=S_Q^{-1}\,\widehat G_\pm^{-1/2}\,S_Q^{-1}.
\]
Hence
\[
\mathcal W(G_-,N,G_+)
=
G_-^{-1/2}NG_+^{-1/2}
=
S_Q\,
\widehat G_-^{-1/2}\widehat N\,\widehat G_+^{-1/2}\,
S_Q
=
S_Q\,\mathcal W(\widehat G_-,\widehat N,\widehat G_+)\,S_Q.
\]
The off-diagonal entries therefore acquire one factor $Q^{-1/2}$ from the
upper-left entry of $S_Q$ and one factor $Q^{-3/2}$ from the lower-right
entry of $S_Q$, hence a total matrix-level factor $Q^{-2}$.
\end{proof}

\begin{proposition}[Uniform scaled corrected-whitening transfer]
\label{prop:differentiated-whitening-transfer}
Fix $0<c_*<\rho_0/8$. For every $Q$ and every radius $R$ satisfying
\[
0<R\le c_*Q,
\]
set
\[
w:=\frac{z}{R},
\qquad
s=\frac{z}{Q^2}=\frac{Rw}{Q^2}.
\]
Let
\[
\mathfrak D_Q(M):=
\begin{pmatrix}
Q^2 m_{11} & Q m_{12}\\
Q m_{21} & m_{22}
\end{pmatrix},
\qquad
\mathfrak D_Q^{-1}(X)=
\begin{pmatrix}
Q^{-2}x_{11} & Q^{-1}x_{12}\\
Q^{-1}x_{21} & x_{22}
\end{pmatrix}.
\]
For $R'>0$, let
\[
A_{R'}(M_2):=
\left\{
F(z)=\sum_{n\ge 0}F_n z^n:
\|F\|_{A_{R'}}:=\sum_{n\ge 0}\|F_n\|\,R'^n<\infty
\right\},
\]
and write $A_1(M_2)$ for the case $R'=1$. For scalar holomorphic germs,
$A_{R'}(\mathbb C)$ is defined similarly.

Write the baseline corrected same-point and mixed blocks on the $w$-disk as
\[
G^{(0)}_{\pm;Q,R}(w):=G_\pm^{(0)}(Rw/Q^2),
\qquad
N^{(0)}_{Q,R}(w):=N^{(0)}(Rw/Q^2).
\]
For a weighted perturbation triple
\[
X=(X_-,Y,X_+)\in A_1(M_2)^3,
\]
define
\[
\delta G_\pm:=\mathfrak D_Q^{-1}X_\pm,
\qquad
\delta N:=\mathfrak D_Q^{-1}Y,
\]
and, whenever the corrected scalar perturbation
\[
\Phi_K\!\left(
\mathcal W\bigl(
G^{(0)}_{-;Q,R}+\delta G_-,
N^{(0)}_{Q,R}+\delta N,
G^{(0)}_{+;Q,R}+\delta G_+
\bigr)
-
\mathcal W\bigl(
G^{(0)}_{-;Q,R},
N^{(0)}_{Q,R},
G^{(0)}_{+;Q,R}
\bigr)
\right)
\]
vanishes at $w=0$, define
\[
\mathcal T_{Q,R}(X)(w)
:=
\frac{Q}{Rw}\,
\Phi_K\!\left(
\mathcal W\bigl(
G^{(0)}_{-;Q,R}+\delta G_-,
N^{(0)}_{Q,R}+\delta N,
G^{(0)}_{+;Q,R}+\delta G_+
\bigr)
-
\mathcal W\bigl(
G^{(0)}_{-;Q,R},
N^{(0)}_{Q,R},
G^{(0)}_{+;Q,R}
\bigr)
\right).
\]

Then there exist absolute constants
\[
\eta>0,\qquad C>0,\qquad B>0
\]
such that for every $Q$, every $0<R\le c_*Q$, and every
\[
X\in A_1(M_2)^3
\qquad\text{with}\qquad
\|X_-\|_{A_1}+\|Y\|_{A_1}+\|X_+\|_{A_1}\le \eta,
\]
the quotient $\mathcal T_{Q,R}(X)$ is holomorphic on $|w|<1$ and admits an
absolutely convergent homogeneous expansion
\[
\mathcal T_{Q,R}(X)=\sum_{p\ge 1}\mathcal T_p(X),
\]
with
\[
\|\mathcal T_p(X)\|_{A_1}
\le
C B^p
\bigl(\|X_-\|_{A_1}+\|Y\|_{A_1}+\|X_+\|_{A_1}\bigr)^p.
\]
The constants $C,B,\eta$ are independent of $Q$ and $R$.

Moreover, if
\[
X_\pm(w)=\sum_{m\ge 0}X_{\pm,m}w^m,
\qquad
Y(w)=\sum_{m\ge 0}Y_m w^m,
\]
and
\[
\mathcal T_p(X)(w)=\sum_{n\ge 0} T_{p,n}w^n,
\]
then with
\[
\mathbf X_m:=\max\!\bigl(\|X_{-,m}\|,\|Y_m\|,\|X_{+,m}\|\bigr),
\]
one has the coefficient bound
\[
\boxed{
|T_{p,n}|
\le
C B^p
\sum_{M\le n}
\sum_{\substack{m_1+\cdots+m_p=M\\ m_i\ge 0}}
\prod_{i=1}^p \mathbf X_{m_i}.
}
\]
\end{proposition}

\begin{proof}
On $|w|\le 2$,
\[
\left|\frac{Rw}{Q^2}\right|
\le \frac{2R}{Q^2}
\le \frac{2c_*}{Q}
<\frac{\rho_0}{Q},
\]
so the baseline blocks
\[
G^{(0)}_{\pm;Q,R}(w),\qquad N^{(0)}_{Q,R}(w)
\]
are holomorphic and uniformly bounded on $|w|\le 2$, with bounds
independent of $Q$ and $R$. By
Lemma~\ref{lem:baseline-scaled-matrix-gain}, the same-point baseline blocks
remain uniformly positive definite there.

Therefore the corrected-whitening map
\[
(\delta G_-,\delta N,\delta G_+)
\longmapsto
\mathcal W\bigl(
G^{(0)}_{-;Q,R}+\delta G_-,
N^{(0)}_{Q,R}+\delta N,
G^{(0)}_{+;Q,R}+\delta G_+
\bigr)
\]
is analytic on a fixed Banach neighborhood in $A_1(M_2)^3$, uniformly in
$Q,R$. Since $\Phi_K$ is linear, the corrected scalar perturbation depends
analytically on the input triple. For those structured perturbations for which
the corrected scalar perturbation vanishes at $w=0$, the quotient by $w$ in
the definition of $\mathcal T_{Q,R}$ is holomorphic. This gives the
homogeneous expansion with uniform constants $C,B,\eta$.

The coefficient bound follows by testing the homogeneous maps on monomials in
$A_1(M_2)$ and using the Wiener-algebra estimate
\[
|[w^n]F|\le \|F\|_{A_1}
\]
for scalar coefficients.
\end{proof}

\begin{proposition}[Pair-radius kernel-class decomposition for one-pair perturbations]
\label{prop:one-pair-symbolic-coefficient-decomposition}
Fix $0<c_*<\rho_0/8$ and $0<c\ll 1$. For one exterior zero pair $\rho$,
set
\[
\Delta_\rho:=|2m-\gamma_\rho|,
\qquad
R_\rho:=c\,\min(\Delta_\rho,c_*Q).
\]

Let $f_\rho(t)$ denote the one-pair phase-derivative contribution, and remove
its midpoint affine jet:
\[
r_\rho(t):=f_\rho(t)-f_\rho(m)-f_\rho'(m)(t-m).
\]
With
\[
t_\pm:=m\pm \frac{z}{2Q^2},
\]
define
\[
d_\pm(z):=r_\rho(t_\pm),
\qquad
e_\pm(z):=r_\rho'(t_\pm),
\qquad
g_\pm(z):=r_\rho''(t_\pm),
\]
and
\[
\eta(z):=\delta\Delta_\rho(z/Q^2):=\int_{t_+}^{t_-} r_\rho(u)\,du.
\]
Then there exist holomorphic kernels $U_\pm,E_\pm,G_\pm,V$ such that
\[
d_\pm(z)=\frac{z^2}{Q^4}U_\pm(z),
\qquad
e_\pm(z)=\frac{z}{Q^2}E_\pm(z),
\qquad
g_\pm(z)=G_\pm(z),
\qquad
\eta(z)=\frac{z^3}{Q^6}V(z).
\]

For every fixed $H_*>1$, if $\Delta_\rho\ge H_*$, then
\[
U_\pm,\ E_\pm,\ G_\pm,\ V
\]
are holomorphic on $|z|<2R_\rho$ and satisfy
\[
\|U_\pm\|_{A_{R_\rho}}+
\|E_\pm\|_{A_{R_\rho}}+
\|G_\pm\|_{A_{R_\rho}}+
\|V\|_{A_{R_\rho}}
\ll
\Delta_\rho^{-4}.
\]

Let
\[
X_\rho:=(X_{\rho,-},Y_\rho,X_{\rho,+})
\]
be the weighted one-pair perturbation triple
\[
X_{\rho,\pm}:=\mathfrak D_Q(\delta G_{\rho,\pm}),
\qquad
Y_\rho:=\mathfrak D_Q(\delta N_\rho),
\]
where $\delta G_{\rho,\pm}$ and $\delta N_\rho$ are the actual one-pair
perturbations of the corrected same-point and mixed blocks in the draft.

Then, on $|z|<2R_\rho$, every entry of every block of $X_\rho(z)$ is
linear in
\[
U_-,\ U_+,\ E_-,\ E_+,\ G_-,\ G_+,\ V
\]
modulo terms containing at least two such pair kernels. In particular,
\[
\|X_\rho(R_\rho\cdot)\|_{A_1}\ll \Delta_\rho^{-4}.
\]
\end{proposition}

\begin{proof}
The one-pair poles lie at distance $\asymp \Delta_\rho$ in the $s$-plane,
hence at distance $\asymp Q^2\Delta_\rho$ in the $z$-plane. Since
\[
R_\rho\ll \min(\Delta_\rho,Q),
\]
the kernels $U_\pm,E_\pm,G_\pm,V$ are holomorphic on $|z|<2R_\rho$, and
Cauchy gives the displayed $A_{R_\rho}$-bounds.

Now check the actual corrected formulas.

For the corrected same-point block,
\[
G_{m,\pm}^{\corr}(s)=\frac1\pi
\begin{pmatrix}
2q(t_\pm) & \frac12 q'(t_\pm)\\[1ex]
\frac12 q'(t_\pm) & \frac1{12}\bigl(q''(t_\pm)+2q(t_\pm)^3\bigr)
\end{pmatrix}.
\]
Subtracting the baseline and using
\[
q=q_0+d_\pm,\qquad q'=q_0'+e_\pm,\qquad q''=q_0''+g_\pm,
\]
the weighted perturbation entries are:
\[
(1,1)\sim (z/Q)^2U_\pm,
\qquad
(1,2),(2,1)\sim (z/Q)E_\pm,
\]
and
\[
(2,2)\sim
G_\pm
+
\Bigl(\frac{q_0}{Q}\Bigr)^2 (z/Q)^2U_\pm
+
\Bigl(\frac{q_0}{Q}\Bigr)(z/Q)^4U_\pm^2
+
(z/Q)^6U_\pm^3.
\]
On $|z|\le 2R_\rho$, the normalized baseline factors
\[
\frac{q_0 z}{Q^2},\qquad \frac{q_0^2 z^2}{Q^4}
\]
are $O(1)$, so the same-point block is linear in the pair kernels modulo
terms with at least two pair kernels.

For the corrected mixed block,
\[
N_m^{\corr}(s)=\frac1\pi
\begin{pmatrix}
-\dfrac{2\sin\Delta}{s}
&
\dfrac{\sin\Delta+q(t_+)s\cos\Delta}{s^2}
\\[2ex]
-\dfrac{\sin\Delta+q(t_-)s\cos\Delta}{s^2}
&
\dfrac{(q(t_-)+q(t_+))s\cos\Delta+(2-q(t_-)q(t_+)s^2)\sin\Delta}{2s^3}
\end{pmatrix}.
\]
Write
\[
q_\pm=q_{0,\pm}+d_\pm,
\qquad
\Delta=\Delta_0+\eta.
\]

For the $(1,1)$-entry, since
\[
\eta(z)=\frac{z^3}{Q^6}V(z)=s^3V(z),
\]
\[
Q^2\frac{\sin(\Delta_0+\eta)-\sin\Delta_0}{s}
\]
is linear in $V$ modulo higher pair order.

For the off-diagonal entries, the numerator difference is
\[
(\sin(\Delta_0+\eta)-\sin\Delta_0)
+d_\pm s\cos(\Delta_0+\eta)
+q_{0,\pm}s(\cos(\Delta_0+\eta)-\cos\Delta_0),
\]
and after multiplying by $Q/s^2$ the linear terms are bounded-baseline times
$U_\pm$ and $V$.

For the $(2,2)$-entry, define
\[
F(q_-,q_+,\Delta;s):=(q_-+q_+)s\cos\Delta+(2-q_-q_+s^2)\sin\Delta.
\]
Then the perturbation is
\[
\frac{
F(q_{0,-}+d_-,q_{0,+}+d_+,\Delta_0+\eta;s)
-
F(q_{0,-},q_{0,+},\Delta_0;s)
}{s^3}.
\]
Taylor expansion in $(d_-,d_+,\eta)$ around the baseline gives
\[
\partial_{q_-}F_0\,d_-+\partial_{q_+}F_0\,d_++\partial_\Delta F_0\,\eta
\]
plus a remainder quadratic in $(d_-,d_+,\eta)$. After dividing by $s^3$,
the linear part is
\[
\widetilde B_-(z)\,U_-(z)+\widetilde B_+(z)\,U_+(z)+\widetilde B_\Delta(z)\,V(z),
\]
with bounded baseline coefficients $\widetilde B_-,\widetilde B_+,\widetilde
B_\Delta$, and the remainder carries at least two pair kernels.

So every weighted mixed-block entry is also linear in the pair kernels modulo
higher pair order.

Finally, on $|z|\le 2R_\rho$, every factor $z/Q$ is $O(1)$, every
normalized baseline factor is $O(1)$, and each pair kernel has $A_{R_\rho}$
norm $O(\Delta_\rho^{-4})$. Therefore
\[
\|X_\rho(R_\rho\cdot)\|_{A_1}\ll \Delta_\rho^{-4}.
\]
\end{proof}

\begin{definition}[Tagged two-atom pre-whitening source block]
\label{def:tagged-two-atom-prewhitening-source-block}
Introduce formal source tags \(\tau_1,\tau_2\).  For a normalized two-atom
core \((a_1,h_1;a_2,h_2)\), define the tagged corrected source perturbations by
additive substitution into the corrected phase variables:
\[
d_\pm^{(2)}:=\tau_1a_1d_{h_1,\pm}+\tau_2a_2d_{h_2,\pm},
\qquad
e_\pm^{(2)}:=\tau_1a_1e_{h_1,\pm}+\tau_2a_2e_{h_2,\pm},
\]
\[
g_\pm^{(2)}:=\tau_1a_1g_{h_1,\pm}+\tau_2a_2g_{h_2,\pm},
\qquad
\eta^{(2)}:=\tau_1a_1\eta_{h_1}+\tau_2a_2\eta_{h_2}.
\]
Equivalently,
\[
q_{2,\pm}=q_{0,\pm}+d_\pm^{(2)},\qquad
q'_{2,\pm}=q'_{0,\pm}+e_\pm^{(2)},\qquad
q''_{2,\pm}=q''_{0,\pm}+g_\pm^{(2)},\qquad
\Delta_2=\Delta_0+\eta^{(2)}.
\]
Define
\[
\mathcal B_2(a_1,h_1;a_2,h_2)
:=(G_{2,-}^{\corr},N_2^{\corr},G_{2,+}^{\corr})
\]
by the same corrected same-point and mixed block formulas used above for the
one-pair corrected block, with \(q_\pm,\Delta\) replaced by
\(q_{2,\pm},\Delta_2\).  Set
\[
\mathcal B_0:=(G_-^{(0)},N^{(0)},G_+^{(0)}).
\]
\end{definition}

\begin{definition}[Pair-kernel linear part]
\label{def:pair-kernel-linear-part}
Let \(\mathcal K\) be the filtration on entries of
\(\mathfrak D_Q(\mathcal B_2-\mathcal B_0)\) in which each tagged source kernel
\[
\tau_iU_{i,\pm},\qquad \tau_iE_{i,\pm},\qquad
\tau_iG_{i,\pm},\qquad \tau_iV_i
\]
has \(\mathcal K\)-degree \(1\), while baseline factors, the amplitudes \(a_i\),
\(Q\)-scales, and bounded analytic coefficients have degree \(0\).  Let
\(\operatorname{Lin}_{\mathcal K}\) denote projection to
\(\mathcal K\)-degree \(1\).

Define the one-atom source-linear triple by
\[
L_h:=\operatorname{Lin}_{\mathcal K}
\mathfrak D_Q(\mathcal B_1(h)-\mathcal B_0),
\]
where \(\mathcal B_1(h)\) is the one-atom corrected block obtained from the
same corrected block formulas.
\end{definition}

\begin{lemma}[Tagged pair-kernel linearity of the source block]
\label{lem:tagged-pair-kernel-linearity-source-block}
With source tags retained,
\[
\operatorname{Lin}_{\mathcal K}
\mathfrak D_Q\bigl(\mathcal B_2(a_1,h_1;a_2,h_2)-\mathcal B_0\bigr)
=
\tau_1a_1L_{h_1}+\tau_2a_2L_{h_2}.
\]
After this tagged \(\mathcal K\)-linear projection is taken, applying the
augmentation \(\tau_1,\tau_2\mapsto1\) gives
\[
\operatorname{Lin}_{\mathcal K}
\mathfrak D_Q(\mathcal B_2-\mathcal B_0)
=a_1L_{h_1}+a_2L_{h_2}.
\]
\end{lemma}

\begin{proof}
For the same-point blocks \(G_{2,\pm}^{\corr}\), the \((1,1)\) and
off-diagonal entries are linear in \(d_\pm^{(2)}\) and \(e_\pm^{(2)}\).  The
\((2,2)\) entry has the form
\[
g_\pm^{(2)}
+\text{baseline}\cdot d_\pm^{(2)}
+\text{terms of }\mathcal K\text{-degree at least }2,
\]
by the same expansion as in the one-pair corrected block calculation.

For the mixed block \(N_2^{\corr}\), Taylor expansion in
\((d_-^{(2)},d_+^{(2)},\eta^{(2)})\) gives a bounded-baseline linear
combination of those variables and a remainder of \(\mathcal K\)-degree at
least \(2\).  Since the tagged variables are additive sums of the two source
contributions, the \(\mathcal K\)-linear part is the sum of the two one-atom
\(\mathcal K\)-linear parts, with tags still attached.
\end{proof}

\begin{remark}[Scope of the tagged source block]
\label{rem:scope-tagged-source-block}
The definitions and lemma above are pre-whitening statements.  They prove
tagged \(\mathcal K\)-linear additivity of the corrected source block before
the nonlinear corrected-whitening transfer
\[
\mathcal W(G_-,N,G_+)=G_-^{-1/2}NG_+^{-1/2}.
\]
They do not by themselves prove transfer-domain vanishing,
source-bidegree edge divisibility, quotient-output package coincidence,
diagonal merger, the coefficient \(B_7^{\mathfrak f}\), the quotient
\(Q_{7,m}^q\), or the projector \(\Pi_{1,1}\).

In particular, the tagged pre-whitening block does not automatically induce
the homogeneous scalar-grade source tower used in
Corollary~\ref{cor:gate-b-raw-edge-divisibility-order-seven}.  That
identification is a separate source-model theorem for the actual
\(\mathfrak D_Q\)-scaled corrected-whitening package.
\end{remark}

\begin{remark}[Corrected package map for the two-point and mixed four-point frontiers]
\label{rem:corrected-package-map-two-point-four-point}
The current local frontier is best viewed as a single corrected package
problem, not as two independent mixed-geometry problems.  The two-point mixed
branch and the mixed four-point branch both require the same source-provenance
package:
\[
\mathcal B_2
\longrightarrow
C_{\mathcal W}
\longrightarrow
B_7^{\mathfrak f}
\longrightarrow
Q_{7,m}^q
\longrightarrow
\mathcal H_7^q .
\]
Here \(\mathcal B_2\) is the tagged two-atom pre-whitening source block of
Definition~\ref{def:tagged-two-atom-prewhitening-source-block},
\(C_{\mathcal W}\) is the corrected-whitening matrix cross-effect before
\(\Phi_K\) of
Definition~\ref{def:raw-corrected-whitening-cross-effect},
\(B_7^{\mathfrak f}\) is the pre-determinant fixed-sector order-\(7\)
coefficient of
Definition~\ref{def:package-level-order-7-quotient-extractor}, and
\[
Q_{7,m}^q(C):=[B_7^{\mathfrak f}(C)]
\in
\mathfrak f/\mathbf C A_5^{\mathfrak f}(m)
\]
is the fixed midpoint quotient class.  The normalized corrected package
representative \(C_{\mathcal W}^{\corr}\) and its moving fixed-sector
representative \(B_{\mathrm{pkg}}^{\mathfrak f}\) are defined in
Section~\ref{sec:corrected-fixed-target-quotient-package}.

For the two-point branch, the remaining source input is finite: prove swap
symmetry, one-amplitude collapse, and diagonal merger for the actual corrected
two-atom package.  Once these hold, the quadratic collision factorization
\[
\mathcal I_2=a_1a_2(h_1-h_2)^2\mathcal J_2
\]
is formal.  The role of the order-\(7\) quotient package is to make the
corresponding septic quotient defect an actual fixed-target object rather than
an uncontrolled two-atom mixed representative.

For the mixed four-point branch, the current reduction says that a genuinely
mixed clean branch cannot persist independently in the region
\[
\{T_v\ne0\}\cap\{\mathcal R\ne0\}.
\]
Thus every clean mixed branch must exit through one of the scalar walls
\[
T_v=0
\qquad\text{or}\qquad
\mathcal R=0.
\]
The package theorem is needed to prevent the corrected source cross-effect
from creating a new non-scalar fixed-sector direction in the quotient
\[
\mathfrak f/\mathbf C A_5^{\mathfrak f}(m).
\]

Consequently, the local target is not a scalarized identity after \(\Phi_K\).
It is the pre-\(\Phi_K\), pre-determinant quotient-gauge statement
\[
B_7^{\mathfrak f}(C_{\mathcal W}^{\corr})
\in
\mathbf C A_5^{\mathfrak f}(m).
\]
If one keeps a full corrected package before source-bidegree projection, the
equivalent full-package formulation is obtained by applying
\((1-\Pi_{1,1})_{\mathrm{mix}}\) before fixed-sector extraction and
quotienting.
This implication does not by itself prove that the scalar exits
\[
T_v=0
\qquad\text{or}\qquad
\mathcal R=0
\]
are impossible.  Its role is narrower: it removes the possibility that
non-\((1,1)\) corrected source terms create an independent fixed-sector
quotient direction.  After that removal, the existing mixed four-point
capture theorem reduces every clean mixed branch to the scalar exits, and
the two-point residual is pushed into the midpoint quotient edge package.

If this package theorem is proved, the residual two-point mixed obstruction is
absorbed into the midpoint quotient edge package, and the mixed four-point
source-coupled walls have no independent clean local carrier beyond the scalar
exits above.
\end{remark}

\begin{remark}[Conditional role of the corrected package in the mixed branches]
\label{rem:conditional-role-corrected-package-mixed-branches}
The corrected package theorem does not by itself prove \(\RH\).  Its role is
to remove the residual local mixed carrier.  If the package theorem closes,
then the two-point mixed obstruction reduces to the fixed midpoint quotient
edge package, and the mixed four-point branch has no independent clean local
carrier outside the scalar exits
\[
T_v=0
\qquad\text{or}\qquad
\mathcal R=0.
\]
The remaining global proof still requires the finite-core lower-model
transfer or a separate treatment of the near-zero-anomaly branch
\(\mathcal N_k(m,Q)\), hidden-state/\(T\)-horizontal odd-block extraction,
and the downstream endgame reductions, including the source-level
interaction branch controlled by \(J_1=[s]\Phi_K(\mathcal I_{\ge2})\).
\end{remark}

\begin{remark}[Sensitivity of the mixed four-point branch to Gate B]
\label{rem:mixed-four-point-sensitivity-to-gate-b}
The mixed four-point branch is sensitive only to the fixed-sector
quotient class of a Gate~B residual.  Let
\[
E=eI+fS\in\mathfrak f_h.
\]
If
\[
E\in \mathbf C A_5^{\mathfrak f}(h),
\]
then \(E\) is invisible in the quotient package.  On a good patch, this
is equivalently
\[
ev_5-u_5f=0,
\qquad
A_5^{\mathfrak f}(h)=u_5I+v_5S.
\]
If instead
\[
ev_5-u_5f\ne0,
\]
then \(E\) shifts the order-\(7\) quotient scalar:
\[
\Delta_7\longmapsto \Delta_7+(ev_5-u_5f).
\]
Consequently \(\lambda=\Delta_7/v_5\) and the downstream scalar package
\(\mathcal Z,\mathcal R,T_v\) generally change.  Thus a non-gauge Gate~B
residual would create precisely the independent fixed-sector quotient
carrier that the corrected package theorem is designed to remove.

The mixed four-point scalar-exit theorem should therefore be read after
Gate~B has removed non-\((1,1)\) quotient directions.  It is not a
substitute for Gate~B.
\end{remark}

\begin{theorem}[Mixed four-point dependence descends to the quotient determinant]
\label{thm:mixed-four-point-dependence-descends-to-quotient-determinant}
Work on a good fixed-sector patch on which
\[
c(h)\ne0,\qquad A_5^{\mathfrak f}(h)=u_5(h)I+v_5(h)S\ne0.
\]
Let
\[
E=eI+fS\in\mathfrak f_h
\]
be an order-\(7\) fixed-sector residual added to
\[
A_7^{\mathfrak f}=u_7I+v_7S,
\]
with \(c,u_5,v_5\) fixed.  Then the mixed four-point scalar-exit package sees
\(E\) only through the quotient determinant pairing
\[
\omega_h(E,A_5^{\mathfrak f}(h))=ev_5-u_5f.
\]

Equivalently,
\[
\Delta_7=u_7v_5-u_5v_7
\]
is transformed by
\[
\Delta_7
\longmapsto
\Delta_7+\omega_h(E,A_5^{\mathfrak f}(h)),
\]
and every lifted \(v_5\)-patch scalar
\[
Q_v,\qquad T_v,\qquad \mathcal Z,\qquad \mathcal R
\]
changes only through this induced change of \(\Delta_7\).

In particular, if
\[
E\in\mathbf C A_5^{\mathfrak f}(h),
\]
then \(E\) is invisible to the mixed four-point geometry.  If
\[
[E]\ne0
\quad\text{in}\quad
\mathfrak f_h/\mathbf C A_5^{\mathfrak f}(h),
\]
then the resulting mixed four-point effect is exactly the already-defined
quotient-septic obstruction, expressed through \(\Delta_7\).  No independent
mixed four-point carrier is produced.
\end{theorem}

\begin{proof}
The lifted mixed four-point branch is built, on a \(v_5\)-patch, from
\[
x=\frac{v_5}{c},
\qquad
Y=\frac{u_5}{c},
\qquad
S=\sigma_v=\frac{\Delta_7}{c\,v_5},
\]
and hence from
\[
Q_v=(Y,x,S).
\]
The scalar exits \(T_v,\mathcal Z,\mathcal R\) are functions of these data and
their \(x\)-derivatives.  Thus the only possible order-\(7\) input from
\(A_7^{\mathfrak f}\) is through \(\Delta_7\).

Replacing \(A_7^{\mathfrak f}\) by \(A_7^{\mathfrak f}+E\), with
\(E=eI+fS\), changes
\[
\Delta_7=\omega_h(A_7^{\mathfrak f},A_5^{\mathfrak f})
\]
by
\[
\omega_h(E,A_5^{\mathfrak f})
=
ev_5-u_5f.
\]
The same argument on a \(u_5\)-patch uses the parallel quotient coordinate
\[
-\frac{\Delta_7}{c\,u_5}.
\]
Hence no raw representative data \((u_7,v_7)\) or
\(\Phi_K(A_7^{\mathfrak f})\) enter as independent mixed four-point inputs.
\end{proof}

\begin{corollary}[Connection residuals are quotient determinant failures]
\label{cor:connection-residuals-are-quotient-determinant-failures}
For \(M\in\{M_A,M_C\}\), set
\[
q_{\mathrm{tot},7}[M]
:=
R_7^{(M)}
+
\sum_{i=1}^{7}\omega_m(Z_{7-i}[M],A_i).
\]
If
\[
q_{\mathrm{tot},7}[M_A]=q_{\mathrm{tot},7}[M_C]=0,
\]
then the residual monomials \(M_A,M_C\) create no mixed four-point carrier
beyond the scalar exits already present in the lifted branch.

If one of these two scalars is nonzero, the obstruction is exactly a failure
of quotient determinant closure:
\[
\Delta_7\longmapsto \Delta_7+q_{\mathrm{tot},7}[M].
\]
It is not a separate four-point mechanism.
\end{corollary}

\begin{remark}[Three data layers in the mixed four-point audit]
\label{rem:three-data-layers-mixed-four-point-audit}
The mixed four-point carrier audit separates three logically different layers:
\[
\begin{array}{c|c|c}
\text{Layer} & \text{Data} & \text{Role} \\
\hline
\text{pre-}\Phi_K\text{ quotient package}
&
[B_7^{\mathfrak f}]
\in
\mathfrak f_m/\mathbf C A_5^{\mathfrak f}(m)
&
\text{Gate B / UV-026 obstruction object}
\\[0.8ex]
\Phi_K\text{-visible hidden state}
&
A_7^{\mathfrak f}\mapsto
A_7^{\mathfrak f}+\chi A_5^{\mathfrak f},
\quad
\Phi_K(\chi A_5^{\mathfrak f})=2\chi v_5
&
\text{visible to transform extraction, quotient-invisible}
\\[0.8ex]
\text{mixed four-point scalar exits}
&
Q_v,\ T_v,\ \mathcal Z,\ \mathcal R
\text{ built from }
c,u_5,v_5,\Delta_7
&
\text{downstream four-point branch geometry}
\end{array}
\]
Thus quotient invisibility is enough to remove an order-\(7\) corrected source
residual from the mixed four-point branch, but it is not by itself enough to
solve the package-to-transform hidden-state bridge.
\end{remark}

\begin{theorem}[Mixed four-point factorization through quotient determinant]
\label{thm:mixed-four-point-factors-through-delta-seven}
Let
\[
A_5^{\mathfrak f}=u_5I+v_5S,
\qquad
A_7^{\mathfrak f}=u_7I+v_7S,
\qquad
c\ne0,
\]
and define
\[
\Delta_7
=
\omega_h(A_7^{\mathfrak f},A_5^{\mathfrak f})
=
u_7v_5-u_5v_7,
\qquad
T=\frac{v_7}{c}.
\]
On every good mixed four-point patch, the mixed four-point scalar exits depend
only on
\[
c,\quad u_5,\quad v_5,\quad \Delta_7.
\]
They do not depend on \(T\) independently.
\end{theorem}

\begin{proof}
On a \(v_5\ne0\) patch, the mixed four-point branch is written through the
normalized variables
\[
x=\frac{v_5}{c},
\qquad
Y=\frac{u_5}{c},
\qquad
\sigma_v=\frac{\Delta_7}{cv_5}.
\]
All scalar exits are functions of these quantities and their local
derivatives.  Thus they are determined by
\[
c,\quad u_5,\quad v_5,\quad \Delta_7,
\]
and do not require \(u_7\) or \(v_7\) separately.

On a \(u_5\ne0\) patch, one uses the parallel quotient coordinate
\[
\sigma_u=-\frac{\Delta_7}{cu_5},
\]
and the same factorization holds.

For an order-seven fixed-sector residual
\[
E=eI+fS,
\]
adding \(E\) to \(A_7^{\mathfrak f}\) changes the determinant by
\[
\Delta_7\longmapsto
\Delta_7+\omega_h(E,A_5^{\mathfrak f})
=
\Delta_7+ev_5-u_5f.
\]
The hidden coordinate changes by
\[
T\longmapsto T+\frac{f}{c}.
\]
Thus the mixed four-point branch sees only the determinant projection
\(ev_5-u_5f\), while the \(T\)-coordinate records a separate
\(\Phi_K\)-visible projection.
\end{proof}

\begin{remark}[No independent \(T\)-dependency in mixed four-point]
\label{rem:no-t-dependency-mixed-four-point}
A pure quotient-line shift
\[
E=\chi A_5^{\mathfrak f}
\]
has
\[
\omega_h(E,A_5^{\mathfrak f})=0,
\]
and is therefore invisible to mixed four-point geometry.  It changes
\[
T
\]
by
\[
T\longmapsto T+\chi\frac{v_5}{c},
\]
and is generally \(\Phi_K\)-visible on \(v_5\ne0\).  Hence mixed four-point
closure is a quotient-determinant problem, whereas two-point transform closure
also requires \(T\)-horizontal compatibility.
\end{remark}

\begin{corollary}[Single-pair exterior odd-coefficient bound]
\label{cor:single-pair-odd-jet}
Fix $0<c_*<\rho_0/8$. Let $h_\rho(s)$ be the corrected transverse scalar
contributed by one exterior zero pair at absolute distance
\[
\Delta_\rho:=|2m-\gamma_\rho|,
\]
and write
\[
h_\rho(z/Q^2)
=
\sum_{r\ge 0} a_{2r+1}(\rho)\,\frac{z^{2r+1}}{Q^{2r+4}}.
\]
Then for every fixed $H_*>1$, if $\Delta_\rho\ge H_*$, one has
\[
\boxed{
|a_{2r+1}(\rho)|
\ll
\frac{1}{Q}
\begin{cases}
\Delta_\rho^{-(2r+4)}, & \Delta_\rho\le c_*Q,\\[1ex]
(c_*Q)^{-2r}\Delta_\rho^{-4}, & \Delta_\rho>c_*Q.
\end{cases}
}
\]
\end{corollary}

\begin{proof}
Let
\[
R_\rho:=c\,\min(\Delta_\rho,c_*Q).
\]
By Proposition~\ref{prop:one-pair-symbolic-coefficient-decomposition},
\[
\|X_\rho(R_\rho\cdot)\|_{A_1}\ll \Delta_\rho^{-4}.
\]
For $H_*$ large enough, this lies inside the uniform small ball from
Proposition~\ref{prop:differentiated-whitening-transfer}. Since the one-pair
corrected scalar $h_\rho(s)$ is odd in $s$, the corresponding corrected
scalar perturbation vanishes at $w=0$, so the quotient map
$\mathcal T_{Q,R_\rho}$ is holomorphic on the one-pair input. Therefore the
normalized even germ
\[
g_\rho(z):=\frac{Q^3}{z}\,h_\rho(z/Q^2)
\]
satisfies
\[
\|g_\rho(R_\rho\cdot)\|_{A_1}\ll \Delta_\rho^{-4}.
\]
Hence
\[
[z^{2r}]\,g_\rho(z)\ll R_\rho^{-2r}\Delta_\rho^{-4}.
\]
Since
\[
g_\rho(z)=\sum_{r\ge 0}\frac{a_{2r+1}(\rho)}{Q^{2r+1}}\,z^{2r},
\]
this gives
\[
|a_{2r+1}(\rho)|\ll Q^{2r+1}R_\rho^{-2r}\Delta_\rho^{-4}.
\]
Using
\[
R_\rho\asymp \min(\Delta_\rho,c_*Q)
\]
yields the displayed hybrid bound.
\end{proof}

\begin{lemma}[Local zero-counting in dyadic shells]
\label{lem:dyadic-shell-count}
Let
\[
\mathcal S_k:=\{\rho:\ 2^kH<\Delta_\rho\le 2^{k+1}H\},
\qquad
\Delta_\rho:=|2m-\gamma_\rho|,
\qquad
k\ge 0.
\]
Then
\[
\#\mathcal S_k\ll Q\,(2^kH+1),
\]
where $Q\asymp \log|m|$.
\end{lemma}

\begin{proof}
This is the local zero-counting estimate already used elsewhere in the draft,
applied to the interval of ordinates
\[
[\,2m-2^{k+1}H,\;2m+2^{k+1}H\,].
\]
Its length is $O(2^kH+1)$, and the local zero density at height $m$ is
$O(Q)$.
\end{proof}

\begin{corollary}[Exterior higher odd-coefficient bound]
\label{cor:outside-higher-odd-jets}
Let
\[
H_{\out}(z/Q^2)
=
\sum_{r\ge 0} a_{2r+1}^{\out}(m,H)\,\frac{z^{2r+1}}{Q^{2r+4}}
\]
be the corrected transverse scalar contributed by all zeros outside absolute
radius $H$ from the local center. Then for every $r\ge 0$ and every
$H\ge H_*$,
\[
\boxed{
|a_{2r+1}^{\out}(m,H)|
\ll
H^{-(2r+3)}+Q^{-(2r+2)}.
}
\]
\end{corollary}

\begin{proof}
Decompose the exterior zeros into dyadic shells
\[
\mathcal S_k:=\{\rho:\ 2^kH<\Delta_\rho\le 2^{k+1}H\},
\qquad k\ge 0.
\]
By Lemma~\ref{lem:dyadic-shell-count},
\[
\#\mathcal S_k\ll Q(2^kH+1).
\]

For shells with $2^{k+1}H\le c_*Q$, Corollary~\ref{cor:single-pair-odd-jet}
gives
\[
\sum_{\rho\in\mathcal S_k}|a_{2r+1}(\rho)|
\ll
Q(2^kH+1)\cdot \frac{1}{Q(2^kH)^{2r+4}}
\ll
(2^kH)^{-(2r+3)}.
\]
Summing these shells yields
\[
\ll H^{-(2r+3)}.
\]

For shells with $2^kH>c_*Q$, Corollary~\ref{cor:single-pair-odd-jet} gives
\[
\sum_{\rho\in\mathcal S_k}|a_{2r+1}(\rho)|
\ll
Q(2^kH+1)\cdot \frac{1}{Q(c_*Q)^{2r}(2^kH)^4}
\ll
(c_*Q)^{-2r}(2^kH)^{-3}.
\]
Summing the far shells gives
\[
\ll
Q^{-2r-2}.
\]
Adding the near and far ranges proves the claim.
\end{proof}

\begin{corollary}[Core localization of the extracted odd jet]
\label{cor:extracted-odd-jet-core-localization}
Let $N=\lfloor \alpha Q\rfloor$ with $\alpha>0$, and let
$a_{2N-1}^{\le H}(m)$ denote the coefficient contributed by the zeros with
$\Delta_\rho\le H$. Then for every fixed $H\ge H_*$,
\[
\boxed{
a_{2N-1}(m)
=
a_{2N-1}^{\le H}(m)
+
O\!\bigl(H^{-(2N+1)}+Q^{-2N}\bigr).
}
\]
In particular, the first surviving corrected odd coefficient is exponentially
localized to any fixed finite core.
\end{corollary}

\begin{proof}
This is exactly Corollary~\ref{cor:outside-higher-odd-jets} with $r=N-1$.
\end{proof}

\section{\texorpdfstring{$N$}{N}-point odd-moment cancellation}

For $N\le \rho_0Q$, set
\[
s_j=\frac{j}{Q^2},
\qquad j=1,\dots,N,
\]
and define
\[
a_j^{(N)}
:=
(-1)^{j+1}\,
\frac{j\binom{2N}{N+j}}{\binom{2N-2}{N-1}},
\qquad 1\le j\le N.
\]
Define the corrected $N$-point transverse scalar
\begin{equation}\label{eq:Xi-zeta}
\Xi_\zeta^{(N)}(m):=
\sum_{j=1}^N a_j^{(N)}H_m(s_j).
\end{equation}

\begin{lemma}[Explicit odd-moment-canceling coefficients and $\ell^1$ control]
\label{lem:n-point-coefficients}
The coefficients $a_j^{(N)}$ satisfy
\[
\sum_{j=1}^N a_j^{(N)}=1,
\qquad
\sum_{j=1}^N a_j^{(N)}j^{2r+1}=0
\quad (0\le r\le N-2).
\]
Moreover,
\[
\sum_{j=1}^N |a_j^{(N)}|=2N-1.
\]
Consequently, for every integer $k\ge 0$,
\[
\left|\sum_{j=1}^N a_j^{(N)}j^k\right|
\le
(2N-1)N^k.
\]
\end{lemma}

\begin{proof}
Set
\[
D_N:=\binom{2N-2}{N-1}.
\]

We first prove the normalization. Writing
\[
j\binom{2N}{N+j}
=
N\!\left(\binom{2N-1}{N+j-1}-\binom{2N-1}{N+j}\right),
\]
one gets
\[
\sum_{j=1}^N (-1)^{j+1}j\binom{2N}{N+j}
=
N\sum_{j=1}^N (-1)^{j+1}
\left(\binom{2N-1}{N+j-1}-\binom{2N-1}{N+j}\right).
\]
Using the partial alternating-binomial identity
\[
\sum_{k=0}^{m}(-1)^k\binom{n}{k}
=
(-1)^m\binom{n-1}{m},
\]
one finds
\[
\sum_{k=N}^{2N-1}(-1)^k\binom{2N-1}{k}
=
(-1)^N\binom{2N-2}{N-1}.
\]
Substituting this into the preceding display yields
\[
\sum_{j=1}^N (-1)^{j+1}j\binom{2N}{N+j}
=
\binom{2N-2}{N-1}
=
D_N,
\]
hence
\[
\sum_{j=1}^N a_j^{(N)}=1.
\]

Next, for $0\le r\le N-2$, apply the standard finite-difference identity
\[
\sum_{k=0}^{2N}(-1)^k\binom{2N}{k}P(k)=0
\]
to the polynomial
\[
P(k)=(k-N)^{2r+2},
\qquad
\deg P=2r+2<2N.
\]
After the change of variables $j=k-N$, this gives
\[
\sum_{j=-N}^{N}(-1)^j\binom{2N}{N+j}j^{2r+2}=0.
\]
Since $2r+2$ is even and the $j=0$ term vanishes,
\[
\sum_{j=1}^{N}(-1)^j\binom{2N}{N+j}j^{2r+2}=0.
\]
Dividing by $D_N$ yields
\[
\sum_{j=1}^N a_j^{(N)}j^{2r+1}=0
\qquad (0\le r\le N-2).
\]

Finally, since the signs alternate,
\[
\sum_{j=1}^N |a_j^{(N)}|
=
\frac{1}{D_N}\sum_{j=1}^N j\binom{2N}{N+j}.
\]
Using again
\[
j\binom{2N}{N+j}
=
N\!\left(\binom{2N-1}{N+j-1}-\binom{2N-1}{N+j}\right),
\]
this telescopes to
\[
\sum_{j=1}^N j\binom{2N}{N+j}
=
N\binom{2N-1}{N}.
\]
Therefore
\[
\sum_{j=1}^N |a_j^{(N)}|
=
N\frac{\binom{2N-1}{N}}{\binom{2N-2}{N-1}}
=
2N-1.
\]
The final moment bound is immediate from
\[
\left|\sum_{j=1}^N a_j^{(N)}j^k\right|
\le
\sum_{j=1}^N |a_j^{(N)}|\,j^k
\le
(2N-1)N^k.
\qedhere
\]
\end{proof}

\begin{lemma}[Closed-form Mellin symbol of the discrete $N$-point projector]
\label{lem:n-point-mellin-symbol}
Define
\[
A_N(t):=\sum_{j=1}^N a_j^{(N)}\,t^j,
\qquad
\widehat\mu_{N,Q}(z):=A_N(e^{z/Q^2}).
\]
Then
\[
A_N(t)
=
\frac{2(2N-1)}{N+1}\,
t\,{}_2F_1(1-N,\,2;\,N+2;\,t),
\]
so in particular
\[
\widehat\mu_{N,Q}(z)
=
\frac{2(2N-1)}{N+1}\,
e^{z/Q^2}\,
{}_2F_1(1-N,\,2;\,N+2;\,e^{z/Q^2}).
\]
Moreover,
\[
A_N(t)-A_N(t^{-1})
=
\frac{N(-1)^N}{\binom{2N-2}{N-1}}\,
t^{-N}(1-t)^{2N-1}(1+t),
\]
and hence
\[
\widehat\mu_{N,Q}(z)-\widehat\mu_{N,Q}(-z)
=
\frac{N(-1)^N}{\binom{2N-2}{N-1}}\,
e^{-Nz/Q^2}\,
\bigl(1-e^{z/Q^2}\bigr)^{2N-1}
\bigl(1+e^{z/Q^2}\bigr).
\]
In particular,
\[
\widehat\mu_{N,Q}(0)=1,
\qquad
\partial_z^{\,2r+1}\widehat\mu_{N,Q}(0)=0
\quad (0\le r\le N-2),
\]
and the odd part of $\widehat\mu_{N,Q}$ starts only at order $2N-1$.
\end{lemma}

\begin{proof}
Let
\[
B_N(t):=
\frac{2(2N-1)}{N+1}\,
t\,{}_2F_1(1-N,\,2;\,N+2;\,t)
=
\sum_{j=1}^N b_j^{(N)}\,t^j.
\]
Its coefficients satisfy
\[
\frac{b_{j+1}^{(N)}}{b_j^{(N)}}
=
\frac{(1-N+j-1)(2+j-1)}{(N+2+j-1)j}
=
\frac{(j-N)(j+1)}{j(N+j+1)}.
\]
On the other hand, the coefficients $a_j^{(N)}$ satisfy
\[
\frac{a_{j+1}^{(N)}}{a_j^{(N)}}
=
\frac{(-1)^{j+2}(j+1)\binom{2N}{N+j+1}}
     {(-1)^{j+1}j\binom{2N}{N+j}}
=
\frac{(j-N)(j+1)}{j(N+j+1)}.
\]
Thus $a_j^{(N)}$ and $b_j^{(N)}$ agree up to an overall constant.  Matching
the $t^1$-coefficient gives
\[
b_1^{(N)}
=
\frac{2(2N-1)}{N+1}
=
\frac{\binom{2N}{N+1}}{\binom{2N-2}{N-1}}
=
a_1^{(N)},
\]
hence $A_N(t)=B_N(t)$.

For the second identity, define
\[
F_N(t):=\sum_{j=-N}^{N}(-1)^j\binom{2N}{N+j}t^j.
\]
By the binomial theorem,
\[
F_N(t)=(-1)^N t^{-N}(1-t)^{2N}.
\]
Differentiating gives
\[
tF_N'(t)
=
\sum_{j=-N}^{N}j(-1)^j\binom{2N}{N+j}t^j
=
-N(-1)^N t^{-N}(1-t)^{2N-1}(1+t).
\]
Now
\[
A_N(t)-A_N(t^{-1})
=
-\frac{1}{\binom{2N-2}{N-1}}
\sum_{j=1}^N j(-1)^j\binom{2N}{N+j}(t^j-t^{-j}),
\]
while the symmetry $j\mapsto -j$ in the preceding display shows
\[
tF_N'(t)
=
\sum_{j=1}^N j(-1)^j\binom{2N}{N+j}(t^j-t^{-j}).
\]
Substituting the closed form for $tF_N'(t)$ yields
\[
A_N(t)-A_N(t^{-1})
=
\frac{N(-1)^N}{\binom{2N-2}{N-1}}\,
t^{-N}(1-t)^{2N-1}(1+t).
\]
Replacing $t$ by $e^{z/Q^2}$ gives the displayed formula for
$\widehat\mu_{N,Q}(z)-\widehat\mu_{N,Q}(-z)$.  The normalization
$\widehat\mu_{N,Q}(0)=1$ and the low odd-jet vanishing are equivalent to the
moment identities in Lemma~\ref{lem:n-point-coefficients}.
\end{proof}

\begin{lemma}[Exact odd-moment generating function]
\label{lem:n-point-odd-generating-function}
Define
\[
\mu_k^{(N)}:=\sum_{j=1}^N a_j^{(N)}j^k,
\]
and, for $|u|<1/N$,
\[
\mathcal M_N(u)
:=
\sum_{r\ge 0}\mu_{2r+1}^{(N)}u^{2r+1}
=
\sum_{j=1}^N \frac{a_j^{(N)}\,j\,u}{1-j^2u^2}.
\]
Then
\[
\mathcal M_N(u)
=
(-1)^{N+1}(2N-1)N!(N-1)!
\frac{u^{2N-1}}{\prod_{j=1}^N(1-j^2u^2)}.
\]
In particular, for every real $x>N$,
\[
\sum_{r\ge N-1} |\mu_{2r+1}^{(N)}|\,x^{-(2r+1)}
=
(2N-1)N!(N-1)!\,x^2\,
\frac{\Gamma(x-N)}{\Gamma(x+N+1)}.
\]
\end{lemma}

\begin{proof}
The identity
\[
\mathcal M_N(u)
=
\sum_{j=1}^N \frac{a_j^{(N)}\,j\,u}{1-j^2u^2}
\]
is immediate from summing the geometric series
\[
\frac{j\,u}{1-j^2u^2}
=
\sum_{r\ge 0} j^{2r+1}u^{2r+1}
\qquad (|u|<1/j).
\]

Set
\[
D_N(u):=
\prod_{j=1}^N(1-j^2u^2)\,\mathcal M_N(u).
\]
Since $\mathcal M_N$ is a sum of simple fractions, $D_N(u)$ is an odd
polynomial of degree at most $2N-1$. By
Lemma~\ref{lem:n-point-coefficients},
\[
\mu_{2r+1}^{(N)}=0
\qquad (0\le r\le N-2),
\]
so $\mathcal M_N(u)=O(u^{2N-1})$ as $u\to 0$. Hence
\[
D_N(u)=C_N\,u^{2N-1}
\]
for some constant $C_N$.

To determine $C_N$, compare residues at $u=1/N$. The residue of
$\mathcal M_N(u)$ there comes only from the $j=N$ term and equals
\[
\operatorname*{Res}_{u=1/N}\mathcal M_N(u)
=
-\frac{a_N^{(N)}}{2N}.
\]
On the other hand,
\[
\operatorname*{Res}_{u=1/N}
\frac{C_Nu^{2N-1}}{\prod_{j=1}^N(1-j^2u^2)}
=
-\frac{C_N}{2N^2\prod_{k=1}^{N-1}(1-k^2/N^2)}.
\]
Equating these residues gives
\[
C_N
=
a_N^{(N)}\,N\prod_{k=1}^{N-1}(N^2-k^2).
\]
Now
\[
a_N^{(N)}
=
(-1)^{N+1}\frac{N}{\binom{2N-2}{N-1}},
\]
and
\[
\prod_{k=1}^{N-1}(N^2-k^2)
=
\prod_{k=1}^{N-1}(N-k)\prod_{k=1}^{N-1}(N+k)
=
(N-1)!\,\frac{(2N-1)!}{N!}
=
\frac{(2N-1)!}{N}.
\]
Therefore
\[
C_N
=
(-1)^{N+1}\frac{N(2N-1)!}{\binom{2N-2}{N-1}}
=
(-1)^{N+1}(2N-1)N!(N-1)!.
\]
This proves the first formula.

For the tail identity, note that
\[
(-1)^{N+1}\mathcal M_N(u)
=
(2N-1)N!(N-1)!
\frac{u^{2N-1}}{\prod_{j=1}^N(1-j^2u^2)}
\]
has nonnegative Taylor coefficients for $0<u<1/N$. Hence
\[
(-1)^{N+1}\mu_{2r+1}^{(N)}\ge 0
\qquad (r\ge N-1),
\]
and therefore, for $x>N$,
\[
\sum_{r\ge N-1} |\mu_{2r+1}^{(N)}|\,x^{-(2r+1)}
=
\left|\mathcal M_N\!\left(\frac1x\right)\right|
=
(2N-1)N!(N-1)!\,
\frac{x^2}{\prod_{j=1}^N(x^2-j^2)}.
\]
Finally,
\[
\prod_{j=1}^N(x^2-j^2)
=
\prod_{j=1}^N(x-j)\prod_{j=1}^N(x+j)
=
\frac{\Gamma(x)}{\Gamma(x-N)}\cdot
\frac{\Gamma(x+N+1)}{\Gamma(x+1)}
=
\frac{1}{x}\frac{\Gamma(x+N+1)}{\Gamma(x-N)}.
\]
Substituting this identity yields the claimed formula.
\end{proof}

\begin{remark}[Structure of the surviving odd tail]
\label{rem:surviving-odd-tail-structure}
Expanding
\[
\frac{1}{\prod_{j=1}^N(1-j^2u^2)}
=
\sum_{k\ge 0} h_k(1^2,\dots,N^2)\,u^{2k},
\]
where $h_k$ denotes the complete homogeneous symmetric polynomial, the exact
formula of Lemma~\ref{lem:n-point-odd-generating-function} yields
\[
\mu_{2N-1+2k}^{(N)}
=
(-1)^{N+1}(2N-1)N!(N-1)!\,
h_k(1^2,\dots,N^2),
\qquad k\ge 0.
\]
Thus the surviving odd moments are not merely bounded: they are explicit and
have a fixed sign,
\[
(-1)^{N+1}\mu_{2r+1}^{(N)}\ge 0
\qquad (r\ge N-1).
\]
This makes the post-cancellation odd tail completely structured rather than
mysterious.
\end{remark}

\begin{corollary}[Exact action of the $N$-point operator on odd germs]
\label{cor:n-point-exact-odd-action}
Let
\[
F(s)=\sum_{r\ge 0} b_{2r+1}s^{2r+1}
\]
be an odd analytic germ at $s=0$. Then
\[
(\mathcal P_{N,Q}F)
:=
\sum_{j=1}^N a_j^{(N)}\,F\!\left(\frac{j}{Q^2}\right)
=
\sum_{r\ge 0} b_{2r+1}\frac{\mu_{2r+1}^{(N)}}{Q^{4r+2}}.
\]
Hence, by Remark~\ref{rem:surviving-odd-tail-structure},
\[
(\mathcal P_{N,Q}F)
=
(-1)^{N+1}(2N-1)N!(N-1)!
\sum_{k\ge 0}
b_{2N-1+2k}\,
h_k(1^2,\dots,N^2)\,
Q^{-(4N+4k-2)}.
\]
Thus the $N$-point operator does not merely annihilate the first $N-1$
odd jets: on odd germs it acts as an explicit shifted complete-homogeneous
transform of the surviving odd Taylor coefficients.
\end{corollary}

\begin{proposition}[Exact finite-difference factorization on odd germs]
\label{prop:n-point-odd-finite-difference-factorization}
Let
\[
T_Q:=e^{Q^{-2}D},
\qquad
D:=\frac{d}{ds},
\]
so that $T_QF(s)=F(s+Q^{-2})$. Then for every odd analytic germ
$F$ at $s=0$,
\[
(\mathcal P_{N,Q}F)
=
\frac{N(-1)^N}{2\binom{2N-2}{N-1}}
\left[
T_Q^{-N}(1-T_Q)^{2N-1}(1+T_Q)F
\right]_{s=0}.
\]
Equivalently,
\[
(\mathcal P_{N,Q}F)
=
\frac{(-1)^{N+1}N}{2\binom{2N-2}{N-1}}
\left[
T_Q^{-N}(T_Q-1)^{2N-1}(1+T_Q)F
\right]_{s=0}.
\]
In particular, on odd germs the $N$-point operator factors exactly through
the $(2N-1)$-st forward difference.
\end{proposition}

\begin{proof}
Let
\[
P_{N,Q}(D):=\sum_{j=1}^N a_j^{(N)}e^{(j/Q^2)D}
=\sum_{j=1}^N a_j^{(N)}T_Q^j.
\]
Then
\[
(\mathcal P_{N,Q}F)=\bigl(P_{N,Q}(D)F\bigr)\big|_{s=0}.
\]
By Lemma~\ref{lem:n-point-mellin-symbol}, with $t=T_Q$,
\[
P_{N,Q}(D)-P_{N,Q}(-D)
=
\frac{N(-1)^N}{\binom{2N-2}{N-1}}\,
T_Q^{-N}(1-T_Q)^{2N-1}(1+T_Q).
\]
If $F$ is odd, then
\[
\bigl(P_{N,Q}(-D)F\bigr)\big|_{s=0}
=
\sum_{j=1}^N a_j^{(N)}F(-j/Q^2)
=
-\sum_{j=1}^N a_j^{(N)}F(j/Q^2)
=
-(\mathcal P_{N,Q}F).
\]
Hence
\[
2(\mathcal P_{N,Q}F)
=
\bigl((P_{N,Q}(D)-P_{N,Q}(-D))F\bigr)\big|_{s=0},
\]
which gives the claimed formula. The second form follows from
$(1-T_Q)^{2N-1}=-(T_Q-1)^{2N-1}$.
\end{proof}

\begin{remark}[Why the odd tail begins at order $2N-1$]
\label{rem:n-point-odd-tail-begins-by-factorization}
Proposition~\ref{prop:n-point-odd-finite-difference-factorization} shows that
the disappearance of the first $N-1$ odd jets is not merely a moment
identity. It is built into the operator itself: on odd germs, the
$N$-point projector contains the factor $(T_Q-1)^{2N-1}$. The surviving
odd tail therefore begins at order $2N-1$ for operator-theoretic reasons,
not only by coefficient cancellation.
\end{remark}

\begin{lemma}[Integral form of a forward difference]
\label{lem:forward-difference-integral-form}
For $m\ge 1$, $h>0$, and an analytic germ $f$, one has
\[
(e^{hD}-1)^m f(x)
=
h^m\int_0^m f^{(m)}(x+hu)\,\omega_m(u)\,du,
\]
where
\[
\omega_m(u)
:=
\frac{1}{(m-1)!}
\sum_{r=0}^{m}(-1)^r\binom{m}{r}(u-r)_+^{m-1},
\qquad
(u-r)_+:=\max\{u-r,0\}.
\]
The function $\omega_m$ is nonnegative, supported on $[0,m]$, and
satisfies the symmetry
\[
\omega_m(u)=\omega_m(m-u).
\]
\end{lemma}

\begin{proof}
The identity
\[
(e^{hD}-1)f(x)=f(x+h)-f(x)=\int_0^h f'(x+t)\,dt
\]
iterates to
\[
(e^{hD}-1)^m f(x)
=
\int_{[0,h]^m}
f^{(m)}(x+t_1+\cdots+t_m)\,dt_1\cdots dt_m.
\]
Rescaling $t_i=hu_i$ gives
\[
(e^{hD}-1)^m f(x)
=
h^m
\int_{[0,1]^m}
f^{(m)}(x+h(u_1+\cdots+u_m))\,du_1\cdots du_m.
\]
Pushing forward Lebesgue measure on $[0,1]^m$ under
\[
(u_1,\dots,u_m)\mapsto u_1+\cdots+u_m
\]
yields the density $\omega_m$, namely the classical Irwin--Hall density,
which has the stated explicit formula, support, and symmetry. This proves the
claim.
\end{proof}

\begin{proposition}[Positive-kernel representation of the odd projector]
\label{prop:n-point-odd-positive-kernel}
Set
\[
m:=2N-1,
\qquad
h:=Q^{-2},
\]
and define, for $0\le u\le N$,
\[
\kappa_N(u):=\omega_{2N-1}(N-u)+\omega_{2N-1}(N-1-u),
\]
where $\omega_{2N-1}$ is the density from
Lemma~\ref{lem:forward-difference-integral-form}, extended by $0$ outside
$[0,2N-1]$. Then $\kappa_N(u)\ge 0$ on $[0,N]$, and for every odd
analytic germ $F$,
\[
(\mathcal P_{N,Q}F)
=
\frac{N(-1)^{N+1}}{\binom{2N-2}{N-1}}\,
Q^{-(4N-2)}
\int_0^N
\kappa_N(u)\,
F^{(2N-1)}\!\left(\frac{u}{Q^2}\right)\,du.
\]
Moreover, $\kappa_N(u)>0$ for $0<u<N$.
\end{proposition}

\begin{proof}
Apply Proposition~\ref{prop:n-point-odd-finite-difference-factorization} and
Lemma~\ref{lem:forward-difference-integral-form} with $m=2N-1$ and
$h=Q^{-2}$. This gives
\[
(\mathcal P_{N,Q}F)
=
\frac{N(-1)^{N+1}}{2\binom{2N-2}{N-1}}
h^{2N-1}
\int_0^{2N-1}
\omega_{2N-1}(v)\,
\Bigl[
F^{(2N-1)}\bigl((v-N)h\bigr)
+
F^{(2N-1)}\bigl((v-N+1)h\bigr)
\Bigr]\,dv.
\]
Since $F$ is odd, $F^{(2N-1)}$ is even. Using this evenness together with
the symmetry
\[
\omega_{2N-1}(v)=\omega_{2N-1}(2N-1-v),
\]
the two terms may be folded onto $[0,N]$, yielding
\[
(\mathcal P_{N,Q}F)
=
\frac{N(-1)^{N+1}}{\binom{2N-2}{N-1}}\,
h^{2N-1}
\int_0^N
\bigl[
\omega_{2N-1}(N-u)+\omega_{2N-1}(N-1-u)
\bigr]
F^{(2N-1)}(uh)\,du.
\]
This is exactly the stated formula. The strict positivity on $0<u<N$
follows because $N-u\in(0,2N-1)$ there, and
$\omega_{2N-1}$ is strictly positive on the interior of its support.
\end{proof}

\begin{corollary}[Sign transfer from the first surviving derivative]
\label{cor:n-point-sign-transfer}
Let $F$ be an odd analytic germ. If for some $\sigma\in\{\pm1\}$,
\[
\sigma\,F^{(2N-1)}(s)\ge 0
\qquad
\text{for all } 0\le s\le \frac{N}{Q^2},
\]
then
\[
\sigma\,(-1)^{N+1}(\mathcal P_{N,Q}F)\ge 0.
\]
If moreover $F^{(2N-1)}\not\equiv 0$ on $[0,N/Q^2]$, then the inequality
is strict.
\end{corollary}

\begin{proof}
This is immediate from Proposition~\ref{prop:n-point-odd-positive-kernel},
since the kernel $\kappa_N$ is nonnegative and strictly positive on
$(0,N)$.
\end{proof}

\begin{remark}[Why this is the right downstream target]
\label{rem:n-point-positive-kernel-target}
Proposition~\ref{prop:n-point-odd-positive-kernel} shows that the $N$-point
layer is not merely an odd-moment cancellation device: on odd germs it is a
positive average of the first surviving derivative
$F^{(2N-1)}$ over the microscopic interval $[0,N/Q^2]$. Accordingly, any
source-level argument that gives a fixed sign, one-sided lower bound, or
variation control for
\[
H_m^{(2N-1)}(s)
\qquad
\left(0\le s\le \frac{N}{Q^2}\right)
\]
would transfer directly to the sign or size of $\Xi_\zeta^{(N)}(m)$.
This isolates the next natural analytic target in the $N$-point route.
\end{remark}

\begin{remark}[Finite-order recurrence of the normalized surviving tail]
\label{rem:n-point-tail-recurrence}
Define
\[
\nu_k
:=
\frac{(-1)^{N+1}\mu_{2N-1+2k}^{(N)}}{(2N-1)N!(N-1)!}
=
h_k(1^2,\dots,N^2),
\qquad k\ge 0.
\]
Then
\[
\sum_{k\ge 0}\nu_k t^k
=
\frac{1}{\prod_{j=1}^N(1-j^2t)}
=
\frac{1}{1-e_1 t+e_2 t^2-\cdots+(-1)^N e_N t^N},
\]
where $e_r=e_r(1^2,\dots,N^2)$ is the $r$-th elementary symmetric
polynomial in $1^2,\dots,N^2$. Consequently,
\[
\nu_k
=
e_1\nu_{k-1}-e_2\nu_{k-2}+\cdots+(-1)^{N+1}e_N\nu_{k-N}
\qquad (k\ge 1),
\]
with the conventions $\nu_0=1$ and $\nu_{k}=0$ for $k<0$. Thus the
surviving odd tail is governed by an exact order-$N$ linear recurrence.
\end{remark}

\begin{proposition}[Partial-fraction form of the surviving odd tail]
\label{prop:n-point-tail-partial-fractions}
Let
\[
G_N(t):=\frac{1}{\prod_{m=1}^N(1-m^2 t)}.
\]
Then $G_N$ has the exact partial-fraction decomposition
\[
G_N(t)
=
\sum_{j=1}^N \frac{B_{N,j}}{1-j^2 t},
\]
where
\[
B_{N,j}
=
\frac{2(-1)^{N-j}j^{2N}}{(N-j)!(N+j)!}.
\]
Equivalently,
\[
\frac{1}{\prod_{m=1}^N(1-m^2u^2)}
=
\sum_{j=1}^N
\frac{2(-1)^{N-j}j^{2N}}{(N-j)!(N+j)!}\,
\frac{1}{1-j^2u^2}.
\]

Consequently, for every $k\ge 0$,
\[
h_k(1^2,\dots,N^2)
=
\sum_{j=1}^N
\frac{2(-1)^{N-j}j^{2N+2k}}{(N-j)!(N+j)!},
\]
and hence
\[
\mu_{2N-1+2k}^{(N)}
=
(-1)^{N+1}(2N-1)N!(N-1)!
\sum_{j=1}^N
\frac{2(-1)^{N-j}j^{2N+2k}}{(N-j)!(N+j)!}.
\]
\end{proposition}

\begin{proof}
Since the poles of $G_N(t)$ are simple and located at $t=j^{-2}$
($1\le j\le N$), one has
\[
G_N(t)=\sum_{j=1}^N \frac{B_{N,j}}{1-j^2 t}
\]
for uniquely determined coefficients $B_{N,j}$. Multiplying by
$(1-j^2 t)$ and evaluating at $t=j^{-2}$ gives
\[
B_{N,j}
=
\left.
\frac{1}{\prod_{m\ne j}(1-m^2 t)}
\right|_{t=j^{-2}}
=
\frac{1}{\prod_{m\ne j}(1-m^2/j^2)}.
\]
Now
\[
\prod_{m\ne j}\left(1-\frac{m^2}{j^2}\right)
=
j^{-2N+2}\prod_{m\ne j}(j^2-m^2),
\]
and
\[
\prod_{m\ne j}(j^2-m^2)
=
\prod_{m\ne j}(j-m)\prod_{m\ne j}(j+m).
\]
Here
\[
\prod_{m\ne j}(j-m)
=
(j-1)! \,(-1)^{N-j}(N-j)!,
\]
while
\[
\prod_{m\ne j}(j+m)
=
\frac{(N+j)!}{2j\,j!}.
\]
Therefore
\[
\prod_{m\ne j}(j^2-m^2)
=
(-1)^{N-j}\frac{(N-j)!(N+j)!}{2j^2},
\]
and hence
\[
B_{N,j}
=
\frac{2(-1)^{N-j}j^{2N}}{(N-j)!(N+j)!}.
\]
Expanding each geometric series
\[
\frac{1}{1-j^2u^2}=\sum_{k\ge 0} j^{2k}u^{2k}
\]
and comparing coefficients of $u^{2k}$ yields the formula for
$h_k(1^2,\dots,N^2)$. The identity for
$\mu_{2N-1+2k}^{(N)}$ then follows from
Remark~\ref{rem:surviving-odd-tail-structure}.
\end{proof}

\begin{remark}[Dominance of the outermost pole for fixed $N$]
\label{rem:n-point-tail-dominant-outer-pole}
For fixed $N$, the explicit decomposition in
Proposition~\ref{prop:n-point-tail-partial-fractions} shows that the surviving
tail is asymptotically dominated by the outermost node $j=N$. Indeed,
\[
h_k(1^2,\dots,N^2)
=
\frac{2N^{2N+2k}}{(2N)!}
+
O_N\!\bigl((N-1)^{2k}\bigr)
\qquad (k\to\infty),
\]
and therefore
\[
\mu_{2N-1+2k}^{(N)}
=
(-1)^{N+1}\frac{N^{2N+2k}}{\binom{2N-2}{N-1}}
+
O_N\!\bigl((N-1)^{2k}\bigr)
\qquad (k\to\infty).
\]
So for fixed $N$, the far odd tail is eventually controlled by the single
outermost pole $u=\pm 1/N$.
\end{remark}

\begin{proposition}[Kernel form and explicit corrected $N$-point transverse estimate]
\label{prop:corrected-n-point}
Let
\[
H_m(s):=\Phi_K\!\bigl(\widehat\Omega_\zeta^{\corr}(s;m)\bigr)
\]
be the full corrected transverse scalar.
For $N=\lfloor \alpha Q\rfloor$ with $0<\alpha<\rho_0$, define
\[
s_j=\frac{j}{Q^2},
\qquad
\Xi_\zeta^{(N)}(m)
:=
\sum_{j=1}^N a_j^{(N)}\,H_m(s_j).
\]
Then
\[
\Xi_\zeta^{(N)}(m)
=
\frac{1}{2\pi i}
\oint_{|\xi|=\rho_0/Q}
H_m(\xi)\,K_{N,Q}(\xi)\,d\xi,
\]
where
\[
K_{N,Q}(\xi)
:=
2\sum_{j=1}^N a_j^{(N)}\frac{s_j}{\xi^2-s_j^2}
=
2(-1)^{N+1}(2N-1)N!(N-1)! \, Q^4\xi\,
\frac{\Gamma(Q^2\xi-N)}{\Gamma(Q^2\xi+N+1)}.
\]
Moreover, writing $x:=\rho_0Q$, one has the explicit bound
\[
|\Xi_\zeta^{(N)}(m)|
\ll_{\rho_0}
\rho_0^2\,
\frac{(2N-1)N!(N-1)!}{Q}\,
\frac{\Gamma(x-N)}{\Gamma(x+N+1)}\,
L(m)S(m)^2.
\]
\end{proposition}

\begin{proof}
By Proposition~\ref{prop:corrected-whitening}, the full corrected scalar
\[
H_m(s)=\Phi_K\!\bigl(\widehat\Omega_\zeta^{\corr}(s;m)\bigr)
\]
is holomorphic on $|s|<\rho_0/Q$. By
Proposition~\ref{prop:cutoff-compatibility}, all cutoff-dependent tail,
jet-correction, and estimation contributions are already internal to the full
corrected scalar $H_m(s)$; hence no external remainder is carried into
$\Xi_\zeta^{(N)}(m)$.

Since $H_m$ is odd, the odd Cauchy formula gives
\[
H_m(s)
=
\frac{1}{2\pi i}\oint_{|\xi|=\rho_0/Q}
H_m(\xi)\,\frac{2s}{\xi^2-s^2}\,d\xi.
\]
Summing this identity against the coefficients $a_j^{(N)}$ at
$s=s_j=j/Q^2$ yields
\[
\Xi_\zeta^{(N)}(m)
=
\frac{1}{2\pi i}\oint_{|\xi|=\rho_0/Q}
H_m(\xi)\,K_{N,Q}(\xi)\,d\xi
\]
with
\[
K_{N,Q}(\xi)
=
2\sum_{j=1}^N a_j^{(N)}\frac{s_j}{\xi^2-s_j^2}.
\]

Now
\[
\frac{s_j}{\xi^2-s_j^2}
=
\frac{1}{\xi}\,
\frac{j/(Q^2\xi)}{1-j^2/(Q^4\xi^2)},
\]
so, using Lemma~\ref{lem:n-point-odd-generating-function},
\[
K_{N,Q}(\xi)
=
\frac{2}{\xi}\,
\mathcal M_N\!\left(\frac{1}{Q^2\xi}\right)
=
2(-1)^{N+1}(2N-1)N!(N-1)!\,
\frac{Q^2}{\prod_{j=1}^N(Q^4\xi^2-j^2)}.
\]
Finally,
\[
\prod_{j=1}^N(Q^4\xi^2-j^2)
=
\prod_{j=1}^N(Q^2\xi-j)\prod_{j=1}^N(Q^2\xi+j)
=
\frac{1}{Q^2\xi}\,
\frac{\Gamma(Q^2\xi+N+1)}{\Gamma(Q^2\xi-N)},
\]
which gives
\[
K_{N,Q}(\xi)
=
2(-1)^{N+1}(2N-1)N!(N-1)! \, Q^4\xi\,
\frac{\Gamma(Q^2\xi-N)}{\Gamma(Q^2\xi+N+1)}.
\]

For the size bound, Proposition~\ref{prop:boundary-estimate} gives
\[
|H_m(s)|\ll_{\rho_0}\frac{L(m)S(m)^2}{Q^3}
\qquad (|s|=\rho_0/Q).
\]
By Proposition~\ref{prop:odd-expansion}, one may write
\[
H_m(z/Q^2)
=
\sum_{r\ge 0} c_{2r+1}(m)\frac{z^{2r+1}}{Q^{2r+4}},
\qquad
|c_{2r+1}(m)|\ll L(m)S(m)^2\,\rho_0^{-(2r+1)}.
\]
Therefore
\[
\Xi_\zeta^{(N)}(m)
=
\sum_{r\ge 0} c_{2r+1}(m)\frac{\mu_{2r+1}^{(N)}}{Q^{2r+4}}.
\]
By Lemma~\ref{lem:n-point-coefficients} and
Remark~\ref{rem:surviving-odd-tail-structure}, this becomes the exact
surviving expansion
\[
\Xi_\zeta^{(N)}(m)
=
(-1)^{N+1}(2N-1)N!(N-1)!
\sum_{k\ge 0}
c_{2N-1+2k}(m)\,
h_k(1^2,\dots,N^2)\,
Q^{-(2N+2k+2)}.
\]
In particular, the $N$-point operation isolates the first surviving
corrected odd coefficient $c_{2N-1}(m)$ and propagates the higher odd
coefficients through an explicit universal positive tail.

For the bound, the coefficient estimate yields
\[
|\Xi_\zeta^{(N)}(m)|
\ll_{\rho_0}
\frac{L(m)S(m)^2}{Q^3}
\sum_{r\ge N-1}
|\mu_{2r+1}^{(N)}|\,(\rho_0Q)^{-(2r+1)}.
\]
Applying Lemma~\ref{lem:n-point-odd-generating-function} with $x=\rho_0Q$
gives
\[
\sum_{r\ge N-1}
|\mu_{2r+1}^{(N)}|\,(\rho_0Q)^{-(2r+1)}
=
(2N-1)N!(N-1)!(\rho_0Q)^2
\frac{\Gamma(\rho_0Q-N)}{\Gamma(\rho_0Q+N+1)}.
\]
Substituting this proves the stated bound.
\end{proof}

\begin{remark}[Exact surviving-coefficient expansion on the zeta side]
\label{rem:xi-zeta-exact-surviving-expansion}
With the notation of the proof of
Proposition~\ref{prop:corrected-n-point},
\[
\Xi_\zeta^{(N)}(m)
=
\sum_{r\ge 0} c_{2r+1}(m)\frac{\mu_{2r+1}^{(N)}}{Q^{2r+4}}
=
(-1)^{N+1}(2N-1)N!(N-1)!
\sum_{k\ge 0}
c_{2N-1+2k}(m)\,
h_k(1^2,\dots,N^2)\,
Q^{-(2N+2k+2)}.
\]
Thus the $N$-point step does not merely remove the first $N-1$ odd jets:
it converts the corrected odd Taylor expansion into an explicit universal tail
whose leading term is the first surviving corrected odd coefficient
$c_{2N-1}(m)$.
\end{remark}

\begin{corollary}[Finite-core localization of the transformed zeta-side scalar]
\label{cor:xi-zeta-finite-core-localization}
With the notation of Proposition~\ref{prop:corrected-n-point} and
Corollary~\ref{cor:outside-higher-odd-jets}, let
\[
H_{\out}(z/Q^2)
=
\sum_{r\ge 0} a_{2r+1}^{\out}(m,H)\,\frac{z^{2r+1}}{Q^{2r+4}}
\]
be the corrected transverse scalar contributed by all zero pairs with
$\Delta_\rho>H$, and define
\[
\Xi_{\zeta,\out}^{(N)}(m;H)
:=
\sum_{j=1}^N a_j^{(N)}\,H_{\out}(s_j),
\qquad
\Xi_{\zeta,\le H}^{(N)}(m)
:=
\Xi_\zeta^{(N)}(m)-\Xi_{\zeta,\out}^{(N)}(m;H).
\]
If $HQ>N$, then
\[
\boxed{
|\Xi_{\zeta,\out}^{(N)}(m;H)|
\ll
\frac{(2N-1)N!(N-1)!}{Q}\,
\frac{\Gamma(HQ-N)}{\Gamma(HQ+N+1)}
+
(2N-1)N!(N-1)!\,
\frac{\Gamma(Q^2-N)}{\Gamma(Q^2+N+1)}.
}
\]
Consequently,
\[
\boxed{
\Xi_\zeta^{(N)}(m)
=
\Xi_{\zeta,\le H}^{(N)}(m)
+
O\!\left(
\frac{(2N-1)N!(N-1)!}{Q}\,
\frac{\Gamma(HQ-N)}{\Gamma(HQ+N+1)}
\right)
+
O\!\left(
(2N-1)N!(N-1)!\,
\frac{\Gamma(Q^2-N)}{\Gamma(Q^2+N+1)}
\right).
}
\]
\end{corollary}

\begin{proof}
By the exact $N$-point action on odd germs used in the proof of
Proposition~\ref{prop:corrected-n-point},
\[
\Xi_{\zeta,\out}^{(N)}(m;H)
=
\sum_{r\ge 0} a_{2r+1}^{\out}(m,H)\frac{\mu_{2r+1}^{(N)}}{Q^{2r+4}}.
\]
Since $\mu_{2r+1}^{(N)}=0$ for $0\le r\le N-2$,
\[
\Xi_{\zeta,\out}^{(N)}(m;H)
=
\sum_{r\ge N-1} a_{2r+1}^{\out}(m,H)\frac{\mu_{2r+1}^{(N)}}{Q^{2r+4}}.
\]
Taking absolute values and applying
Corollary~\ref{cor:outside-higher-odd-jets},
\[
|\Xi_{\zeta,\out}^{(N)}(m;H)|
\ll
\sum_{r\ge N-1}
|\mu_{2r+1}^{(N)}|
\left(
\frac{H^{-(2r+3)}}{Q^{2r+4}}
+
\frac{Q^{-(2r+2)}}{Q^{2r+4}}
\right).
\]
Thus
\[
|\Xi_{\zeta,\out}^{(N)}(m;H)|
\ll
\frac{1}{H^2Q^3}
\sum_{r\ge N-1}
|\mu_{2r+1}^{(N)}|(HQ)^{-(2r+1)}
+
\frac{1}{Q^4}
\sum_{r\ge N-1}
|\mu_{2r+1}^{(N)}|(Q^2)^{-(2r+1)}.
\]
Now apply Lemma~\ref{lem:n-point-odd-generating-function} first with
$x=HQ$ and then with $x=Q^2$. Since $HQ>N$, both evaluations are valid
for all sufficiently large $Q$, and one obtains
\[
\sum_{r\ge N-1}
|\mu_{2r+1}^{(N)}|(HQ)^{-(2r+1)}
=
(2N-1)N!(N-1)!(HQ)^2
\frac{\Gamma(HQ-N)}{\Gamma(HQ+N+1)},
\]
and
\[
\sum_{r\ge N-1}
|\mu_{2r+1}^{(N)}|(Q^2)^{-(2r+1)}
=
(2N-1)N!(N-1)!Q^4
\frac{\Gamma(Q^2-N)}{\Gamma(Q^2+N+1)}.
\]
Substituting these identities proves the displayed bound for
$\Xi_{\zeta,\out}^{(N)}(m;H)$, and the decomposition
\[
\Xi_\zeta^{(N)}(m)=\Xi_{\zeta,\le H}^{(N)}(m)+\Xi_{\zeta,\out}^{(N)}(m;H)
\]
gives the second display.
\end{proof}

\begin{remark}[Asymptotic finite-core localization at linear depth]
\label{rem:xi-zeta-finite-core-localization-asymptotic}
Let
\[
N=\lfloor \alpha Q\rfloor,
\qquad
0<\alpha<\rho_0,
\]
and fix $H>\alpha$. Then Corollary~\ref{cor:xi-zeta-finite-core-localization}
and Stirling's formula give
\[
|\Xi_{\zeta,\out}^{(N)}(m;H)|
\ll
\exp\!\bigl(-H f(\alpha/H)\,Q+O(\log Q)\bigr)
+
\exp\!\bigl(-2\alpha Q\log Q+O(Q)\bigr),
\]
where
\[
f(\beta)
=
(1+\beta)\log(1+\beta)-2\beta\log\beta-(1-\beta)\log(1-\beta).
\]
In particular,
\[
\Xi_\zeta^{(N)}(m)
=
\Xi_{\zeta,\le H}^{(N)}(m)
+
O\!\bigl(e^{-Hf(\alpha/H)Q+O(\log Q)}\bigr)
+
O\!\bigl(e^{-2\alpha Q\log Q+O(Q)}\bigr).
\]
Thus the full transformed zeta-side scalar is determined, up to exponentially
small error, by a fixed finite core.
\end{remark}

\begin{definition}[Finite-core \(N\)-point scalar]
\label{def:finite-core-npoint-scalar}
Fix a finite same-scale core \(\mathcal C\) of zero pairs around the midpoint
\(m\), and let
\[
H_{\core,m}(s)
:=
\Phi_K\!\bigl(\widehat\Omega_{\core}^{\corr}(s;m)\bigr)
\]
denote the corrected transverse scalar contributed by this core, in the same
finite-\(s\) normalization as \(H_m(s)\).  For \(N\ge1\), define
\[
\Xi_{\core}^{(N)}(m)
:=
\sum_{j=1}^N a_j^{(N)}H_{\core,m}(j/Q^2).
\]
If
\[
H_{\core,m}(z/Q^2)
=
\sum_{r\ge 0}
c_{\core,2r+1}(m)\frac{z^{2r+1}}{Q^{2r+4}},
\]
then the \(N\)-point moment cancellations give
\[
\Xi_{\core}^{(N)}(m)
=
(-1)^{N+1}(2N-1)N!(N-1)!
\sum_{k\ge0}
c_{\core,2N-1+2k}(m)
h_k(1^2,\dots,N^2)
Q^{-(2N+2k+2)}.
\]
Thus the first finite-core coefficient seen by the \(N\)-point projector at
depth \(N\) is
\[
c_{\core,2N-1}(m).
\]
\end{definition}

\begin{remark}[Finite-core localization is not a lower law]
\label{rem:finite-core-localization-not-lower-law}
Corollary~\ref{cor:xi-zeta-finite-core-localization} is an upper/localization
statement.  It does not supply a lower bound for
\[
\Xi_{\core}^{(N)}(m).
\]
In particular, it does not rule out cancellation among the finitely many core
constituents:
\[
c_{\core,2N-1}(m)=0
\]
or cancellation of several consecutive \(N\)-visible odd coefficients.
A separate finite-core lower-transfer theorem is required before the
\(N\)-point contradiction can be applied to genuinely non-pair-like cores.
\end{remark}

\begin{remark}[Exact Laurent tail of the contour kernel]
\label{rem:n-point-kernel-laurent-tail}
The kernel in Proposition~\ref{prop:corrected-n-point} has an exact large-
$\xi$ Laurent expansion. Indeed, using
Lemma~\ref{lem:n-point-odd-generating-function} and
Remark~\ref{rem:surviving-odd-tail-structure},
\[
K_{N,Q}(\xi)
=
\frac{2}{\xi}\,
\mathcal M_N\!\left(\frac{1}{Q^2\xi}\right)
=
2(-1)^{N+1}(2N-1)N!(N-1)!\,
Q^{-4N+2}\xi^{-2N}
\sum_{k\ge 0}
h_k(1^2,\dots,N^2)\,
Q^{-4k}\xi^{-2k}.
\]
So the contour kernel is not merely small: its entire post-cancellation odd
tail is explicit, universal, and sign-rigid term-by-term.
\end{remark}

\begin{corollary}[Optimized zeta-side cancellation depth]
\label{cor:optimized-zeta-cancellation-depth}
Let
\[
N=\lfloor \alpha Q\rfloor,
\qquad
0<\alpha<\rho_0,
\qquad
\beta:=\frac{\alpha}{\rho_0}\in(0,1).
\]
Then
\[
|\Xi_\zeta^{(N)}(m)|
\ll_{\rho_0,\alpha}
\frac{L(m)S(m)^2}{Q}\,
e^{-\delta_*(\alpha,\rho_0)Q},
\]
where
\[
\delta_*(\alpha,\rho_0)
=
\rho_0\Bigl(
(1+\beta)\log(1+\beta)
-2\beta\log\beta
-(1-\beta)\log(1-\beta)
\Bigr).
\]
Moreover, $\delta_*(\alpha,\rho_0)$ is maximized uniquely at
\[
\beta=\frac{1}{\sqrt2},
\qquad\text{i.e.}\qquad
\alpha_{\mathrm{opt}}=\frac{\rho_0}{\sqrt2},
\]
so the optimal zeta-side cancellation depth is
\[
N_{\mathrm{opt}}
=
\left\lfloor \frac{\rho_0}{\sqrt2}\,Q\right\rfloor,
\]
and the corresponding exponential rate is
\[
\delta_{\mathrm{opt}}
=
2\rho_0\log(1+\sqrt2).
\]
\end{corollary}

\begin{proof}
Set $x:=\rho_0Q$, so that $N=\beta x+O(1)$ with $\beta\in(0,1)$ fixed.
By Proposition~\ref{prop:corrected-n-point}, it suffices to estimate
\[
(2N-1)N!(N-1)!\,\frac{\Gamma(x-N)}{\Gamma(x+N+1)}.
\]
A Stirling expansion gives
\[
\log\!\left(
(2N-1)N!(N-1)!\,\frac{\Gamma(x-N)}{\Gamma(x+N+1)}
\right)
=
-x\,f(\beta)+O(\log x),
\]
where
\[
f(\beta)
=
(1+\beta)\log(1+\beta)
-2\beta\log\beta
-(1-\beta)\log(1-\beta).
\]
Hence
\[
(2N-1)N!(N-1)!\,\frac{\Gamma(x-N)}{\Gamma(x+N+1)}
\ll_{\rho_0,\alpha}
e^{-x f(\beta)},
\]
and therefore
\[
|\Xi_\zeta^{(N)}(m)|
\ll_{\rho_0,\alpha}
\frac{L(m)S(m)^2}{Q}\,
e^{-\rho_0 f(\beta)Q}.
\]
This proves the first claim.

For the optimization, differentiate:
\[
f'(\beta)
=
\log\!\left(\frac{1-\beta^2}{\beta^2}\right).
\]
Thus $f'(\beta)=0$ exactly when $\beta=1/\sqrt2$. Since
\[
f''(\beta)
=
-\frac{2\beta}{1-\beta^2}-\frac{2}{\beta}<0
\qquad (0<\beta<1),
\]
this critical point is the unique maximizer. Finally,
\[
f\!\left(\frac{1}{\sqrt2}\right)
=
2\log(1+\sqrt2),
\]
which gives the stated value of $\delta_{\mathrm{opt}}$.
\end{proof}

\begin{remark}
Corollary~\ref{cor:optimized-zeta-cancellation-depth} optimizes the
\emph{zeta-side} cancellation depth. After the exact odd-tail evaluation is
imported into Proposition~\ref{prop:toy-n-point-direct}, the toy side imposes
no asymptotic upper restriction on the linear choice $N=\lfloor \alpha
Q\rfloor$: for every fixed $\alpha>0$, the toy microscopic nodes
$s_j=j/Q^2$ remain inside the toy analyticity disk for all sufficiently
large $Q$, and the residual odd toy tail is in fact super-exponentially
small on the $Q$-scale. Accordingly one may take the exact zeta-side
optimum
\[
\alpha=\frac{\rho_0}{\sqrt2}.
\]
\end{remark}

\begin{remark}[Sharper prefactor in the optimized zeta-side bound]
\label{rem:optimized-zeta-prefactor}
With the notation of Corollary~\ref{cor:optimized-zeta-cancellation-depth},
set
\[
x:=\rho_0Q,
\qquad
N=\beta x+O(1),
\qquad
\beta=\frac{\alpha}{\rho_0}\in(0,1).
\]
A more precise Stirling expansion of the gamma-ratio factor in
Proposition~\ref{prop:corrected-n-point} gives
\[
(2N-1)N!(N-1)!\,\frac{\Gamma(x-N)}{\Gamma(x+N+1)}
=
\frac{4\pi\beta}{\sqrt{1-\beta^2}}\,
e^{-x f(\beta)}\,
(1+o(1)),
\]
where
\[
f(\beta)
=
(1+\beta)\log(1+\beta)
-2\beta\log\beta
-(1-\beta)\log(1-\beta).
\]
Consequently,
\[
|\Xi_\zeta^{(N)}(m)|
\ll
2C_{\partial}(\rho_0)\rho_0\,
\frac{4\pi\beta}{\sqrt{1-\beta^2}}\,
\frac{L(m)S(m)^2}{Q^2}\,
e^{-\delta_*(\alpha,\rho_0)Q}\,
(1+o(1)).
\]
At the zeta-side optimum $\beta=1/\sqrt2$, the prefactor simplifies to
$4\pi$, so
\[
|\Xi_\zeta^{(N_{\mathrm{opt}})}(m)|
\ll
8\pi\,C_{\partial}(\rho_0)\rho_0\,
\frac{L(m)S(m)^2}{Q^2}\,
e^{-2\rho_0\log(1+\sqrt2)\,Q}\,
(1+o(1)).
\]
This sharper prefactor is not needed for the qualitative contradiction, but it
records that, once the boundary constant $C_{\partial}(\rho_0)$ is made
explicit, the zeta-side $N$-point estimate becomes effective with an
explicit leading constant.
\end{remark}

\begin{proposition}[Real profile of the exact $N$-point kernel]
\label{prop:exact-kernel-real-profile}
Let
\[
C_N:=(2N-1)N!(N-1)!,
\qquad
\kappa_N(x):=
x\,\frac{\Gamma(x-N)}{\Gamma(x+N+1)}
\qquad (x>N).
\]
Then
\[
\kappa_N(x)=\frac{1}{\prod_{j=1}^N(x^2-j^2)},
\]
and for real $x>N$,
\[
K_{N,Q}\!\left(\frac{x}{Q^2}\right)
=
2(-1)^{N+1}C_NQ^2\,\kappa_N(x).
\]
Moreover:

\begin{enumerate}
\item $\kappa_N(x)>0$ for all $x>N$;
\item $\kappa_N'(x)<0$ for all $x>N$;
\item for every $x>N$,
\[
\sup_{|\xi|=x/Q^2}|K_{N,Q}(\xi)|
=
2C_NQ^2\,\kappa_N(x)
=
\left|K_{N,Q}\!\left(\frac{x}{Q^2}\right)\right|.
\]
\end{enumerate}
\end{proposition}

\begin{proof}
From Proposition~\ref{prop:corrected-n-point},
\[
K_{N,Q}(\xi)
=
2(-1)^{N+1}C_NQ^4\xi\,
\frac{\Gamma(Q^2\xi-N)}{\Gamma(Q^2\xi+N+1)}.
\]
Substituting $\xi=x/Q^2$ gives
\[
K_{N,Q}\!\left(\frac{x}{Q^2}\right)
=
2(-1)^{N+1}C_NQ^2\,
x\,\frac{\Gamma(x-N)}{\Gamma(x+N+1)}
=
2(-1)^{N+1}C_NQ^2\,\kappa_N(x).
\]

Now
\[
\Gamma(x+N+1)
=
(x+N)(x+N-1)\cdots x\,\Gamma(x),
\]
and
\[
\Gamma(x)
=
(x-1)(x-2)\cdots (x-N)\,\Gamma(x-N).
\]
Hence
\[
x\,\frac{\Gamma(x-N)}{\Gamma(x+N+1)}
=
\frac{1}{(x-N)\cdots(x-1)(x+1)\cdots(x+N)}
=
\frac{1}{\prod_{j=1}^N(x^2-j^2)},
\]
which proves the product formula. Positivity is immediate for $x>N$.

For monotonicity, differentiate the logarithm:
\[
\frac{d}{dx}\log\kappa_N(x)
=
-\sum_{j=1}^N \frac{2x}{x^2-j^2}.
\]
Since $x>N$, each denominator is positive, so
\[
\frac{d}{dx}\log\kappa_N(x)<0,
\]
hence $\kappa_N'(x)<0$.

Finally, using the product form of the kernel,
\[
K_{N,Q}(\xi)
=
2(-1)^{N+1}C_NQ^2\,
\frac{1}{\prod_{j=1}^N(Q^4\xi^2-j^2)},
\]
if $|\xi|=x/Q^2$ then
\[
|Q^4\xi^2-j^2|
\ge \bigl||Q^2\xi|^2-j^2\bigr|
=
x^2-j^2
\qquad (1\le j\le N),
\]
so
\[
|K_{N,Q}(\xi)|
\le
2C_NQ^2\,
\frac{1}{\prod_{j=1}^N(x^2-j^2)}
=
2C_NQ^2\,\kappa_N(x).
\]
Equality is attained at $\xi=x/Q^2$, proving the final claim.
\end{proof}

\begin{remark}[No contour loss in the kernel factor]
\label{rem:no-contour-loss-kernel-factor}
Proposition~\ref{prop:exact-kernel-real-profile} shows that the kernel factor
in Proposition~\ref{prop:corrected-n-point} is controlled exactly by its
positive real value on the boundary circle.  Thus the boundary estimate used in
the zeta-side $N$-point bound carries no hidden contour-loss from the kernel
itself: the worst point is the positive real boundary point
\[
\xi=\frac{\rho_0}{Q}.
\]
Moreover, since $\kappa_N$ is strictly decreasing on $(N,\infty)$, the
kernel factor is largest at the smallest admissible boundary radius.
\end{remark}

\begin{proposition}[Positive smoothing preserves the real kernel profile]
\label{prop:smoothed-real-kernel-profile}
Let $\psi\in C_c^\infty(\mathbf R)$ be even, nonnegative, and normalized by
\[
\int_{\mathbf R}\psi(u)\,du=1.
\]
Set
\[
U_\psi:=\sup\{|u|:u\in \operatorname{supp}\psi\}.
\]
For $x>N+\eta Q^2U_\psi$, define the smoothed real profile
\[
\kappa_{N,\eta}^{\psi}(x)
:=
\int_{\mathbf R}\psi(u)\,\kappa_N(x+\eta Q^2u)\,du,
\]
where
\[
\kappa_N(x):=x\,\frac{\Gamma(x-N)}{\Gamma(x+N+1)}
=\frac{1}{\prod_{j=1}^N(x^2-j^2)}.
\]
Then
\[
\kappa_{N,\eta}^{\psi}(x)>0
\qquad\text{and}\qquad
\frac{d}{dx}\kappa_{N,\eta}^{\psi}(x)<0
\]
for all $x>N+\eta Q^2U_\psi$.
\end{proposition}

\begin{proof}
For $x>N+\eta Q^2U_\psi$ and $u\in \operatorname{supp}\psi$, one has
\[
x+\eta Q^2u>N.
\]
By Proposition~\ref{prop:exact-kernel-real-profile},
\[
\kappa_N(x+\eta Q^2u)>0
\qquad\text{and}\qquad
\kappa_N'(x+\eta Q^2u)<0.
\]
Since $\psi\ge 0$ and $\int \psi=1$, averaging preserves positivity:
\[
\kappa_{N,\eta}^{\psi}(x)>0.
\]
Differentiating under the integral sign gives
\[
\frac{d}{dx}\kappa_{N,\eta}^{\psi}(x)
=
\int_{\mathbf R}\psi(u)\,\kappa_N'(x+\eta Q^2u)\,du<0.
\]
\end{proof}

\begin{remark}[Smoothed real profile]
\label{rem:smoothed-real-kernel-profile}
Proposition~\ref{prop:smoothed-real-kernel-profile} shows that nonnegative
smoothing preserves the qualitative real-axis behavior of the exact kernel:
positivity and strict monotonic decay survive averaging. Thus the Mellin-side
smoothing introduced later does not destroy the basic real-profile structure of
the $N$-point kernel.
\end{remark}

\begin{remark}[Discrete Mellin projector and test-function viewpoint]
\label{rem:discrete-mellin-projector}
Define the additive operator
\[
(\mathcal P_{N,Q}F):=\sum_{j=1}^N a_j^{(N)}\,F\!\left(\frac{j}{Q^2}\right),
\]
and the associated multiplicative distribution on $\mathbb R_+^\times$,
\[
\mu_{N,Q}:=\sum_{j=1}^N a_j^{(N)}\,\delta_{e^{j/Q^2}}.
\]
Its Mellin symbol is precisely
\[
\widehat\mu_{N,Q}(z)=\sum_{j=1}^N a_j^{(N)}e^{zj/Q^2}.
\]

\medskip

Equivalently, if $D:=\frac{d}{ds}$, then
\[
\mathcal P_{N,Q}
=
P_{N,Q}(D)\big|_{s=0},
\qquad
P_{N,Q}(D):=\sum_{j=1}^N a_j^{(N)}e^{(j/Q^2)D}.
\]
Thus $\mathcal P_{N,Q}$ is a polynomial in the translation semigroup,
evaluated at the origin. If
\[
F(s)=\sum_{n\ge 0} c_n s^n,
\]
then
\[
(\mathcal P_{N,Q}F)
=
\sum_{n\ge 0} c_n\,\frac{\mu_n^{(N)}}{Q^{2n}},
\qquad
\mu_n^{(N)}:=\sum_{j=1}^N a_j^{(N)}j^n.
\]
Lemma~\ref{lem:n-point-mellin-symbol} shows that
$\widehat\mu_{N,Q}(0)=1$ and that the odd part of
$\widehat\mu_{N,Q}$ starts only at order $2N-1$; meanwhile,
Corollary~\ref{cor:n-point-exact-odd-action} identifies the exact action of
$\mathcal P_{N,Q}$ on odd germs as a shifted complete-homogeneous transform
of the surviving odd Taylor coefficients.

In this sense, the $N$-point operator may be viewed as a local discrete
Mellin-side projector. This is formally reminiscent of the Weil/Connes
test-function viewpoint, where one probes the zero set through Mellin/Fourier
transforms of multiplicative test functions and explicit- or trace-formula
pairings \cite{Conrey2003,Connes1997,ConnesConsani2021}. The present paper
does not use any global trace formula; the comparison is only structural.

A natural smoothed version is obtained by fixing an even
$\phi\in C_c^\infty(\mathbb R)$ with $\widehat\phi(0)=1$ and setting
\[
f_{N,Q,\eta}(x)
:=
\sum_{j=1}^N a_j^{(N)}
\,\eta^{-1}\phi\!\left(\frac{\log x-j/Q^2}{\eta}\right).
\]
Its Mellin transform is
\[
\widehat f_{N,Q,\eta}(z)
=
\widehat\phi(\eta z)\,\widehat\mu_{N,Q}(z).
\]
Hence the same low odd-jet annihilation persists after smoothing, while the
finite-support projector is replaced by an honest smooth multiplicative test
function. This observation is not used elsewhere in the proof, but it places
the $N$-point cancellation layer in a more canonical operator-theoretic
setting.
\end{remark}

\begin{lemma}[Smoothed Mellin projector]
\label{lem:smoothed-mellin-projector}
Let $\phi\in C_c^\infty(\mathbf R)$ be even and normalized by
\[
\int_{\mathbf R}\phi(u)\,du=1.
\]
Define its bilateral Laplace transform by
\[
\widehat\phi(w):=\int_{\mathbf R}\phi(u)e^{wu}\,du.
\]
For $\eta>0$, define the smoothed additive projector
\[
(\mathcal P_{N,Q,\eta}^{\phi}F)
:=
\sum_{j=1}^N a_j^{(N)}
\int_{\mathbf R}\phi(u)\,
F\!\left(\frac{j}{Q^2}+\eta u\right)\,du,
\]
whenever $F$ is holomorphic on a neighborhood of the integration region.

Then:

\begin{enumerate}
\item for every exponential mode $F(s)=e^{zs}$,
\[
\mathcal P_{N,Q,\eta}^{\phi}(e^{zs})
=
\widehat\phi(\eta z)\,\widehat\mu_{N,Q}(z),
\qquad
\widehat\mu_{N,Q}(z)=\sum_{j=1}^N a_j^{(N)}e^{zj/Q^2};
\]

\item the smoothed symbol preserves the exact low odd-jet annihilation:
\[
\widehat\phi(0)\widehat\mu_{N,Q}(0)=1,
\qquad
\partial_z^{\,2r+1}\!\bigl(\widehat\phi(\eta z)\widehat\mu_{N,Q}(z)\bigr)\Big|_{z=0}=0
\quad (0\le r\le N-2);
\]

\item if $F$ is holomorphic on a neighborhood $\Omega$ containing all points
\[
\frac{j}{Q^2}+\eta u
\qquad
(1\le j\le N,\ u\in \operatorname{supp}\phi),
\]
then
\[
\bigl|\mathcal P_{N,Q,\eta}^{\phi}F-\mathcal P_{N,Q}F\bigr|
\le
\frac{\eta^2}{2}\,(2N-1)\,m_2(\phi)\,
\sup_{\Omega}|F''|,
\]
where
\[
m_2(\phi):=\int_{\mathbf R}|u|^2\,|\phi(u)|\,du.
\]
\end{enumerate}
\end{lemma}

\begin{proof}
For $F(s)=e^{zs}$,
\[
\mathcal P_{N,Q,\eta}^{\phi}(e^{zs})
=
\sum_{j=1}^N a_j^{(N)}e^{zj/Q^2}
\int_{\mathbf R}\phi(u)e^{\eta zu}\,du
=
\widehat\phi(\eta z)\,\widehat\mu_{N,Q}(z),
\]
which proves the first statement.

Since $\phi$ is even, $\widehat\phi$ is an even entire function. In
particular,
\[
\widehat\phi(0)=\int_{\mathbf R}\phi(u)\,du=1,
\qquad
\widehat\phi^{(2r+1)}(0)=0
\quad (r\ge 0).
\]
Lemma~\ref{lem:n-point-mellin-symbol} gives
\[
\widehat\mu_{N,Q}(0)=1,
\qquad
\partial_z^{\,2r+1}\widehat\mu_{N,Q}(0)=0
\quad (0\le r\le N-2).
\]
Applying Leibniz' rule to the product
$\widehat\phi(\eta z)\widehat\mu_{N,Q}(z)$ shows that all odd derivatives up
to order $2N-3$ vanish at $z=0$, proving the second statement.

For the quantitative comparison, expand $F$ at $j/Q^2$:
\[
F\!\left(\frac{j}{Q^2}+\eta u\right)
=
F\!\left(\frac{j}{Q^2}\right)
+
\eta u\,F'\!\left(\frac{j}{Q^2}\right)
+
\frac{\eta^2u^2}{2}\,F''(\xi_{j,u})
\]
for some $\xi_{j,u}\in\Omega$. Since $\phi$ is even,
\[
\int_{\mathbf R}\phi(u)\,u\,du=0.
\]
Also $\int \phi(u)\,du=1$. Therefore
\[
\int_{\mathbf R}\phi(u)\,
F\!\left(\frac{j}{Q^2}+\eta u\right)\,du
-
F\!\left(\frac{j}{Q^2}\right)
=
\frac{\eta^2}{2}
\int_{\mathbf R}\phi(u)\,u^2\,F''(\xi_{j,u})\,du.
\]
Summing against $a_j^{(N)}$ and using
Lemma~\ref{lem:n-point-coefficients},
\[
\sum_{j=1}^N |a_j^{(N)}|=2N-1,
\]
gives
\[
\bigl|\mathcal P_{N,Q,\eta}^{\phi}F-\mathcal P_{N,Q}F\bigr|
\le
\frac{\eta^2}{2}\,(2N-1)\,m_2(\phi)\,
\sup_{\Omega}|F''|.
\qedhere
\]
\end{proof}

\begin{remark}[Smoothing scale]
\label{rem:smoothing-scale-mellin-projector}
Since $N\asymp Q$, Lemma~\ref{lem:smoothed-mellin-projector} gives
\[
\mathcal P_{N,Q,\eta}^{\phi}F
=
\mathcal P_{N,Q}F
+
O\!\bigl(Q\eta^2\sup|F''|\bigr).
\]
Thus any choice $\eta=o(Q^{-1/2})$ makes the smoothing error $o(1)$ while
preserving the exact low odd-jet annihilation. This does not enter the present
proof, but it shows that the discrete $N$-point projector admits a genuinely
smooth Mellin-side avatar with the same exact cancellation pattern.
\end{remark}

\begin{proposition}[Smoothed comparison for the corrected $N$-point zeta scalar]
\label{prop:smoothed-zeta-n-point-comparison}
Let $\phi\in C_c^\infty(\mathbf R)$ be even, normalized by
\[
\int_{\mathbf R}\phi(u)\,du=1,
\]
and set
\[
U_\phi:=\sup\{|u|:u\in \operatorname{supp}\phi\},
\qquad
m_2(\phi):=\int_{\mathbf R}|u|^2\,|\phi(u)|\,du.
\]
Fix $0<\alpha<\rho_0$, let
\[
N=\lfloor \alpha Q\rfloor,
\]
and assume
\[
0<\eta\le \frac{\rho_0-\alpha}{2U_\phi\,Q}.
\]
Define the smoothed corrected scalar
\[
H_{m,\eta}(s):=\int_{\mathbf R}\phi(u)\,H_m(s+\eta u)\,du,
\]
and the corresponding smoothed $N$-point scalar
\[
\Xi_{\zeta,\eta}^{(N)}(m)
:=
\sum_{j=1}^N a_j^{(N)}\,H_{m,\eta}\!\left(\frac{j}{Q^2}\right).
\]
Then:

\begin{enumerate}
\item $H_{m,\eta}$ is odd and holomorphic on
\[
|s|<\frac{\rho_\eta}{Q},
\qquad
\rho_\eta:=\rho_0-U_\phi \eta Q>\alpha.
\]

\item One has the zeta-side exponential estimate
\[
|\Xi_{\zeta,\eta}^{(N)}(m)|
\ll_{\phi,\rho_0,\alpha}
\frac{L(m)S(m)^2}{Q}\,
e^{-\delta_*(\alpha,\rho_\eta)Q},
\]
where
\[
\delta_*(\alpha,\rho_\eta)
=
\rho_\eta\Bigl(
(1+\beta_\eta)\log(1+\beta_\eta)
-2\beta_\eta\log\beta_\eta
-(1-\beta_\eta)\log(1-\beta_\eta)
\Bigr),
\qquad
\beta_\eta:=\frac{\alpha}{\rho_\eta}.
\]

\item The smoothed and discrete $N$-point scalars satisfy
\[
|\Xi_{\zeta,\eta}^{(N)}(m)-\Xi_{\zeta}^{(N)}(m)|
\le
\frac{4\,C_{\partial}(\rho_0)\,m_2(\phi)}{(\rho_0-\alpha)^2}\,
\frac{2N-1}{Q}\,
\eta^2\,L(m)S(m)^2.
\]
In particular,
\[
|\Xi_{\zeta,\eta}^{(N)}(m)-\Xi_{\zeta}^{(N)}(m)|
\ll_{\phi,\rho_0,\alpha}
\eta^2\,L(m)S(m)^2.
\]
\end{enumerate}
\end{proposition}

\begin{proof}
Since $H_m$ is odd and $\phi$ is even,
\[
H_{m,\eta}(-s)
=
\int_{\mathbf R}\phi(u)\,H_m(-s+\eta u)\,du
=
\int_{\mathbf R}\phi(u)\,H_m(-(s+\eta u))\,du
=
-H_{m,\eta}(s),
\]
so $H_{m,\eta}$ is odd. If $|s|<\rho_\eta/Q$ and $u\in \operatorname{supp}\phi$,
then
\[
|s+\eta u|
\le
|s|+\eta U_\phi
<
\frac{\rho_0-U_\phi\eta Q}{Q}+\eta U_\phi
=
\frac{\rho_0}{Q}.
\]
Since $H_m$ is holomorphic on $|s|<\rho_0/Q$, it follows that
$H_{m,\eta}$ is holomorphic on $|s|<\rho_\eta/Q$. This proves (1).

For $|s|=\rho_\eta/Q$ and $u\in \operatorname{supp}\phi$, the same estimate gives
\[
|s+\eta u|\le \frac{\rho_0}{Q}.
\]
Hence Proposition~\ref{prop:boundary-estimate} implies
\[
|H_m(s+\eta u)|
\le
C_{\partial}(\rho_0)\,\frac{L(m)S(m)^2}{Q^3}.
\]
Therefore
\[
|H_{m,\eta}(s)|
\le
\|\phi\|_{L^1}\,C_{\partial}(\rho_0)\,\frac{L(m)S(m)^2}{Q^3}
\qquad
\left(|s|=\frac{\rho_\eta}{Q}\right).
\]
Since $H_{m,\eta}$ is odd and holomorphic on $|s|<\rho_\eta/Q$, the proof of
Corollary~\ref{cor:optimized-zeta-cancellation-depth} applies verbatim with
$\rho_0$ replaced by $\rho_\eta$. This gives
\[
|\Xi_{\zeta,\eta}^{(N)}(m)|
\ll_{\phi,\rho_0,\alpha}
\frac{L(m)S(m)^2}{Q}\,
e^{-\delta_*(\alpha,\rho_\eta)Q},
\]
proving (2).

For (3), note that
\[
\Xi_{\zeta,\eta}^{(N)}(m)-\Xi_{\zeta}^{(N)}(m)
=
\bigl(\mathcal P_{N,Q,\eta}^{\phi}-\mathcal P_{N,Q}\bigr)H_m.
\]
Apply Lemma~\ref{lem:smoothed-mellin-projector} with $F=H_m$. It remains to
bound $H_m''$ on the relevant region. Since
\[
\frac{j}{Q^2}\le \frac{N}{Q^2}\le \frac{\alpha}{Q},
\qquad
\eta U_\phi\le \frac{\rho_0-\alpha}{2Q},
\]
all points
\[
\frac{j}{Q^2}+\eta u
\qquad
(1\le j\le N,\ u\in\operatorname{supp}\phi)
\]
lie in the disk
\[
|s|\le \frac{\alpha}{Q}+\eta U_\phi
\le
\frac{\rho_0+\alpha}{2Q}.
\]
Thus their distance to the boundary circle $|s|=\rho_0/Q$ is at least
\[
\frac{\rho_0-\alpha}{2Q}.
\]
By Cauchy's derivative estimate and Proposition~\ref{prop:boundary-estimate},
\[
\sup_{\Omega}|H_m''|
\le
\frac{2\,C_{\partial}(\rho_0)\,L(m)S(m)^2/Q^3}
{\bigl((\rho_0-\alpha)/(2Q)\bigr)^2}
=
\frac{8\,C_{\partial}(\rho_0)}{(\rho_0-\alpha)^2}\,
\frac{L(m)S(m)^2}{Q}.
\]
Lemma~\ref{lem:smoothed-mellin-projector} therefore yields
\[
|\Xi_{\zeta,\eta}^{(N)}(m)-\Xi_{\zeta}^{(N)}(m)|
\le
\frac{\eta^2}{2}\,(2N-1)\,m_2(\phi)\,
\frac{8\,C_{\partial}(\rho_0)}{(\rho_0-\alpha)^2}\,
\frac{L(m)S(m)^2}{Q},
\]
which is exactly the stated bound.
\end{proof}

\begin{remark}[Microscopic smoothing regime]
\label{rem:microscopic-smoothing-zeta-n-point}
If
\[
\eta\ll Q^{-2},
\]
then Proposition~\ref{prop:smoothed-zeta-n-point-comparison} gives
\[
\Xi_{\zeta,\eta}^{(N)}(m)
=
\Xi_{\zeta}^{(N)}(m)
+
O\!\left(\frac{L(m)S(m)^2}{Q^4}\right).
\]
At the same time,
\[
\rho_\eta=\rho_0+O(Q^{-1}),
\]
so
\[
\delta_*(\alpha,\rho_\eta)
=
\delta_*(\alpha,\rho_0)+O(Q^{-1}).
\]
Thus microscopic smoothing preserves the same zeta-side exponential scale while
replacing the discrete projector by a genuinely smooth Mellin-side test
function.
\end{remark}

\begin{lemma}[Higher-order smoothed Mellin projector]
\label{lem:higher-order-smoothed-mellin-projector}
Fix an integer $M\ge 1$. Let $\phi_M\in C_c^\infty(\mathbf R)$ be even and
satisfy
\[
\int_{\mathbf R}\phi_M(u)\,du=1,
\qquad
\int_{\mathbf R}u^{2\ell}\phi_M(u)\,du=0
\quad (1\le \ell\le M-1).
\]
Define its bilateral Laplace transform by
\[
\widehat\phi_M(w):=\int_{\mathbf R}\phi_M(u)e^{wu}\,du.
\]
For $\eta>0$, define
\[
(\mathcal P_{N,Q,\eta}^{(M)}F)
:=
\sum_{j=1}^N a_j^{(N)}
\int_{\mathbf R}\phi_M(u)\,
F\!\left(\frac{j}{Q^2}+\eta u\right)\,du,
\]
whenever $F$ is holomorphic on a neighborhood of the integration region.

Then:

\begin{enumerate}
\item for every exponential mode $F(s)=e^{zs}$,
\[
\mathcal P_{N,Q,\eta}^{(M)}(e^{zs})
=
\widehat\phi_M(\eta z)\,\widehat\mu_{N,Q}(z),
\qquad
\widehat\mu_{N,Q}(z)=\sum_{j=1}^N a_j^{(N)}e^{zj/Q^2};
\]

\item one has
\[
\widehat\phi_M(0)=1,
\qquad
\widehat\phi_M^{(k)}(0)=0
\quad (1\le k\le 2M-1);
\]

\item the smoothed symbol preserves the exact low odd-jet annihilation:
\[
\partial_z^{\,2r+1}\!\bigl(\widehat\phi_M(\eta z)\widehat\mu_{N,Q}(z)\bigr)\Big|_{z=0}=0
\qquad (0\le r\le N-2);
\]

\item if $F$ is holomorphic on a neighborhood $\Omega$ containing all points
\[
\frac{j}{Q^2}+\eta u
\qquad
(1\le j\le N,\ u\in \operatorname{supp}\phi_M),
\]
then
\[
\bigl|\mathcal P_{N,Q,\eta}^{(M)}F-\mathcal P_{N,Q}F\bigr|
\le
\frac{\eta^{2M}}{(2M)!}\,(2N-1)\,m_{2M}(\phi_M)\,
\sup_{\Omega}|F^{(2M)}|,
\]
where
\[
m_{2M}(\phi_M):=\int_{\mathbf R}|u|^{2M}|\phi_M(u)|\,du.
\]
\end{enumerate}
\end{lemma}

\begin{proof}
For $F(s)=e^{zs}$,
\[
\mathcal P_{N,Q,\eta}^{(M)}(e^{zs})
=
\sum_{j=1}^N a_j^{(N)}e^{zj/Q^2}
\int_{\mathbf R}\phi_M(u)e^{\eta zu}\,du
=
\widehat\phi_M(\eta z)\,\widehat\mu_{N,Q}(z),
\]
which proves the first statement.

Since $\phi_M$ is even, all odd moments vanish. By the imposed even-moment
conditions, one also has
\[
\int_{\mathbf R}u^k\phi_M(u)\,du=0
\qquad (1\le k\le 2M-1).
\]
Differentiating $\widehat\phi_M$ at $0$ gives
\[
\widehat\phi_M^{(k)}(0)=\int_{\mathbf R}u^k\phi_M(u)\,du,
\]
so
\[
\widehat\phi_M(0)=1,
\qquad
\widehat\phi_M^{(k)}(0)=0
\quad (1\le k\le 2M-1).
\]
This proves the second statement.

By Lemma~\ref{lem:n-point-mellin-symbol},
\[
\widehat\mu_{N,Q}(0)=1,
\qquad
\partial_z^{\,2r+1}\widehat\mu_{N,Q}(0)=0
\quad (0\le r\le N-2).
\]
Applying Leibniz' rule to
$\widehat\phi_M(\eta z)\widehat\mu_{N,Q}(z)$ shows that all odd derivatives
up to order $2N-3$ vanish at $z=0$, proving the third statement.

For the quantitative comparison, expand $F$ to order $2M-1$ at $j/Q^2$:
\[
F\!\left(\frac{j}{Q^2}+\eta u\right)
=
\sum_{k=0}^{2M-1}
\frac{(\eta u)^k}{k!}\,
F^{(k)}\!\left(\frac{j}{Q^2}\right)
+
\frac{(\eta u)^{2M}}{(2M)!}\,F^{(2M)}(\xi_{j,u})
\]
for some $\xi_{j,u}\in\Omega$. Integrating against $\phi_M$, all terms of
order $1$ through $2M-1$ vanish by the moment conditions, while the
constant term survives because $\int\phi_M=1$. Hence
\[
\int_{\mathbf R}\phi_M(u)\,
F\!\left(\frac{j}{Q^2}+\eta u\right)\,du
-
F\!\left(\frac{j}{Q^2}\right)
=
\frac{\eta^{2M}}{(2M)!}
\int_{\mathbf R}\phi_M(u)\,u^{2M}F^{(2M)}(\xi_{j,u})\,du.
\]
Summing against $a_j^{(N)}$ and using
Lemma~\ref{lem:n-point-coefficients},
\[
\sum_{j=1}^N |a_j^{(N)}|=2N-1,
\]
gives the claimed bound.
\qedhere
\end{proof}

\begin{proposition}[Higher-order smoothed comparison for the corrected $N$-point zeta scalar]
\label{prop:higher-order-smoothed-zeta-n-point-comparison}
Fix an integer $M\ge 1$, and let $\phi_M\in C_c^\infty(\mathbf R)$ be even,
normalized, and moment-corrected as in
Lemma~\ref{lem:higher-order-smoothed-mellin-projector}. Set
\[
U_M:=\sup\{|u|:u\in \operatorname{supp}\phi_M\},
\qquad
m_{2M}(\phi_M):=\int_{\mathbf R}|u|^{2M}|\phi_M(u)|\,du.
\]
Fix $0<\alpha<\rho_0$, let
\[
N=\lfloor \alpha Q\rfloor,
\]
and assume
\[
0<\eta\le \frac{\rho_0-\alpha}{2U_M\,Q}.
\]
Define
\[
H_{m,\eta}^{(M)}(s):=\int_{\mathbf R}\phi_M(u)\,H_m(s+\eta u)\,du,
\]
and
\[
\Xi_{\zeta,\eta}^{(N,M)}(m)
:=
\sum_{j=1}^N a_j^{(N)}\,H_{m,\eta}^{(M)}\!\left(\frac{j}{Q^2}\right).
\]
Then:

\begin{enumerate}
\item $H_{m,\eta}^{(M)}$ is odd and holomorphic on
\[
|s|<\frac{\rho_{M,\eta}}{Q},
\qquad
\rho_{M,\eta}:=\rho_0-U_M\eta Q>\alpha;
\]

\item one has
\[
|\Xi_{\zeta,\eta}^{(N,M)}(m)|
\ll_{\phi_M,\rho_0,\alpha}
\frac{L(m)S(m)^2}{Q}\,
e^{-\delta_*(\alpha,\rho_{M,\eta})Q};
\]

\item the smoothed and discrete $N$-point scalars satisfy
\[
|\Xi_{\zeta,\eta}^{(N,M)}(m)-\Xi_{\zeta}^{(N)}(m)|
\le
C_{\partial}(\rho_0)\,(2N-1)\,m_{2M}(\phi_M)\,
\left(\frac{2}{\rho_0-\alpha}\right)^{2M}
\eta^{2M}Q^{2M-3}\,
L(m)S(m)^2,
\]
and hence
\[
|\Xi_{\zeta,\eta}^{(N,M)}(m)-\Xi_{\zeta}^{(N)}(m)|
\ll_{\phi_M,\rho_0,\alpha}
(\eta Q)^{2M}\,\frac{L(m)S(m)^2}{Q^2}.
\]
\end{enumerate}
\end{proposition}

\begin{proof}
The proof of oddness and holomorphy is identical to that of
Proposition~\ref{prop:smoothed-zeta-n-point-comparison}, with $\phi$ replaced
by $\phi_M$. Likewise, on $|s|=\rho_{M,\eta}/Q$,
\[
|H_{m,\eta}^{(M)}(s)|
\le
\|\phi_M\|_{L^1}\,C_{\partial}(\rho_0)\,\frac{L(m)S(m)^2}{Q^3}.
\]
Therefore the proof of
Corollary~\ref{cor:optimized-zeta-cancellation-depth} applies verbatim with
$\rho_0$ replaced by $\rho_{M,\eta}$, giving the stated exponential bound.

For the comparison estimate,
\[
\Xi_{\zeta,\eta}^{(N,M)}(m)-\Xi_{\zeta}^{(N)}(m)
=
\bigl(\mathcal P_{N,Q,\eta}^{(M)}-\mathcal P_{N,Q}\bigr)H_m.
\]
Apply Lemma~\ref{lem:higher-order-smoothed-mellin-projector} with $F=H_m$.
Exactly as in the proof of
Proposition~\ref{prop:smoothed-zeta-n-point-comparison}, all relevant points
lie at distance at least
\[
\frac{\rho_0-\alpha}{2Q}
\]
from the boundary circle $|s|=\rho_0/Q$. Hence Cauchy's derivative estimate
and Proposition~\ref{prop:boundary-estimate} give
\[
\sup_{\Omega}|H_m^{(2M)}|
\le
(2M)!\,C_{\partial}(\rho_0)\,
\left(\frac{2Q}{\rho_0-\alpha}\right)^{2M}
\frac{L(m)S(m)^2}{Q^3}.
\]
Substituting this into Lemma~\ref{lem:higher-order-smoothed-mellin-projector}
yields
\[
|\Xi_{\zeta,\eta}^{(N,M)}(m)-\Xi_{\zeta}^{(N)}(m)|
\le
C_{\partial}(\rho_0)\,(2N-1)\,m_{2M}(\phi_M)\,
\left(\frac{2}{\rho_0-\alpha}\right)^{2M}
\eta^{2M}Q^{2M-3}\,
L(m)S(m)^2,
\]
as claimed. Since $N\asymp Q$, the final big-Oh form follows.
\end{proof}

\begin{remark}[Saturation at the natural $Q^{-1}$ smoothing scale]
\label{rem:higher-order-smoothing-saturation}
Proposition~\ref{prop:higher-order-smoothed-zeta-n-point-comparison} shows that
higher-order moment correction improves the comparison error to order
$2M$, but under the largest admissible smoothing scale
\[
\eta\asymp Q^{-1},
\]
one still has
\[
|\Xi_{\zeta,\eta}^{(N,M)}(m)-\Xi_{\zeta}^{(N)}(m)|
\ll
\frac{L(m)S(m)^2}{Q^2}.
\]
Thus higher-order smoothing does not improve the $Q$-power at the natural
admissible scale; its gain is in canonicity and in the ability to make the
smoothed projector arbitrarily close to the discrete one once one works at a
smaller scale. For example, if
\[
\eta=Q^{-1-\sigma}
\qquad (\sigma>0),
\]
then
\[
|\Xi_{\zeta,\eta}^{(N,M)}(m)-\Xi_{\zeta}^{(N)}(m)|
\ll
Q^{-2-2M\sigma}\,L(m)S(m)^2.
\]
\end{remark}

We isolate the toy microscopic expansion in
Appendix~\ref{app:toy-microscopic-expansion}; the only fact used in the main
argument is the following $N$-point preservation statement.

\begin{proposition}[Toy $N$-point transverse preservation]
\label{prop:toy-n-point-direct}
Fix a compact interval
\[
D=[d_-,d_+]\subset(0,\infty),
\]
let $\rho_{\toy}$ and $u_0(D)$ be as in
Lemma~\ref{lem:toy-microscopic-expansion}, and fix
\[
\alpha>0.
\]
For
\[
N=\lfloor \alpha Q\rfloor,
\qquad
s_j=\frac{j}{Q^2},
\]
define
\[
\Xi_{\toy}^{(N)}(u,d)
:=
\sum_{j=1}^N a_j^{(N)}\,
\Phi_K\!\bigl(\Delta_{\toy}(u,d;s_j)\bigr).
\]
Then, for all sufficiently large $Q$ (depending on $D$ and $\alpha$),
for $d\in D$ and $|u|<u_0(D)$,
\[
\Xi_{\toy}^{(N)}(u,d)
=
u^2\Phi_K(M(d))+O(u^4),
\]
uniformly in $d\in D$. In particular,
\[
|\Xi_{\toy}^{(N)}(u,d)|\asymp u^2
\]
uniformly on compact $d$-ranges.
\end{proposition}

\begin{proof}
By Lemma~\ref{lem:toy-microscopic-expansion},
\[
\Delta_{\toy}(u,d;s)=u^2M(d)+u^4E(u,d;s),
\]
where $E(u,d;s)$ is analytic in $s$ on $|s|<\rho_{\toy}$ and uniformly
bounded there for $d\in D$ and $|u|<u_0(D)$. Applying the linear
functional $\Phi_K$, define
\[
R_{u,d}(s):=\Phi_K(E(u,d;s)).
\]
Then $R_{u,d}(s)$ is analytic on $|s|<\rho_{\toy}$, and by Cauchy's
estimate its Taylor coefficients satisfy
\[
R_{u,d}(s)=\sum_{n\ge 0} b_n(u,d)s^n,
\qquad
|b_n(u,d)|\ll \rho_{\toy}^{-n},
\]
uniformly for $d\in D$ and $|u|<u_0(D)$.

Therefore
\[
\Phi_K\!\bigl(\Delta_{\toy}(u,d;s_j)\bigr)
=
u^2\Phi_K(M(d))+u^4R_{u,d}(s_j),
\]
so
\[
\Xi_{\toy}^{(N)}(u,d)
=
u^2\Phi_K(M(d))
+
u^4\sum_{j=1}^N a_j^{(N)}R_{u,d}(s_j).
\]
Expand $R_{u,d}$ at the microscopic nodes $s_j=j/Q^2$:
\[
\sum_{j=1}^N a_j^{(N)}R_{u,d}(s_j)
=
\sum_{n\ge 0} b_n(u,d)\frac{\mu_n^{(N)}}{Q^{2n}},
\qquad
\mu_n^{(N)}:=\sum_{j=1}^N a_j^{(N)}j^n.
\]

The constant term contributes
\[
b_0(u,d)\mu_0^{(N)}=b_0(u,d),
\]
since $\mu_0^{(N)}=\sum_j a_j^{(N)}=1$ by
Lemma~\ref{lem:n-point-coefficients}. For the remaining terms, split into
even and odd parts.

For the even terms, Lemma~\ref{lem:n-point-coefficients} gives
\[
|\mu_{2\ell}^{(N)}|
\le
(2N-1)N^{2\ell},
\]
so
\[
\sum_{\ell\ge 1}
|b_{2\ell}(u,d)|\frac{|\mu_{2\ell}^{(N)}|}{Q^{4\ell}}
\ll
(2N-1)\sum_{\ell\ge 1}
\rho_{\toy}^{-2\ell}\frac{N^{2\ell}}{Q^{4\ell}}
\ll
\frac{1}{Q},
\]
because $N\asymp Q$.

For the odd terms, Lemma~\ref{lem:n-point-coefficients} gives
\[
\mu_{2r+1}^{(N)}=0
\qquad (0\le r\le N-2),
\]
so only the tail $r\ge N-1$ remains. Using the coefficient bound
$|b_{2r+1}(u,d)|\ll \rho_{\toy}^{-(2r+1)}$, we obtain
\[
\sum_{r\ge N-1}
|b_{2r+1}(u,d)|\frac{|\mu_{2r+1}^{(N)}|}{Q^{4r+2}}
\ll
\sum_{r\ge N-1}
|\mu_{2r+1}^{(N)}|\,(\rho_{\toy}Q^2)^{-(2r+1)}.
\]
Now apply Lemma~\ref{lem:n-point-odd-generating-function} with
\[
x=\rho_{\toy}Q^2.
\]
For all sufficiently large $Q$, one has $x>N$, and therefore
\[
\sum_{r\ge N-1}
|\mu_{2r+1}^{(N)}|\,(\rho_{\toy}Q^2)^{-(2r+1)}
=
(2N-1)N!(N-1)!(\rho_{\toy}Q^2)^2
\frac{\Gamma(\rho_{\toy}Q^2-N)}{\Gamma(\rho_{\toy}Q^2+N+1)}.
\]
Since $N=\lfloor \alpha Q\rfloor$ and $x=\rho_{\toy}Q^2$, Stirling's
formula gives
\[
(2N-1)N!(N-1)!(\rho_{\toy}Q^2)^2
\frac{\Gamma(\rho_{\toy}Q^2-N)}{\Gamma(\rho_{\toy}Q^2+N+1)}
=
\exp\!\bigl(-2\alpha Q\log Q+O(Q)\bigr).
\]
In particular, the residual odd toy tail is super-exponentially small on the
$Q$-scale:
\[
\sum_{r\ge N-1}
|b_{2r+1}(u,d)|\frac{|\mu_{2r+1}^{(N)}|}{Q^{4r+2}}
\ll
\exp\!\bigl(-2\alpha Q\log Q+O(Q)\bigr).
\]

Combining the constant, even, and odd contributions,
\[
\sum_{j=1}^N a_j^{(N)}R_{u,d}(s_j)
=
b_0(u,d)+O(Q^{-1})+O(e^{-cQ}),
\]
and in particular
\[
\sum_{j=1}^N a_j^{(N)}R_{u,d}(s_j)=O(1),
\]
uniformly for $d\in D$ and $|u|<u_0(D)$. Therefore
\[
\Xi_{\toy}^{(N)}(u,d)
=
u^2\Phi_K(M(d))+O(u^4),
\]
uniformly in $d\in D$.

Because $\Phi_K(M(d))\neq 0$ for $d>0$, and is bounded away from $0$ on
every compact interval $D\subset(0,\infty)$, the final estimate
\[
|\Xi_{\toy}^{(N)}(u,d)|\asymp u^2
\]
follows.
\end{proof}

\begin{remark}[The toy odd tail is negligible beyond all $e^{-cQ}$-scales]
\label{rem:toy-odd-tail-superexp}
The proof of Proposition~\ref{prop:toy-n-point-direct} shows more than the
$O(u^4)$ preservation statement used later: after $N$-point cancellation,
the residual odd toy tail is of size
\[
\exp\!\bigl(-2\alpha Q\log Q+O(Q)\bigr).
\]
Thus on the toy side the $N$-point odd leakage is not merely exponentially
small in $Q$; it is negligible beyond every $e^{-cQ}$-scale.
\end{remark}

On the toy side, Proposition~\ref{prop:toy-n-point-direct} gives
\begin{equation}\label{eq:Xi-toy}
\Xi_{\toy}^{(N)}(u,d)=u^2\Phi_K(M(d))+O(u^4),
\end{equation}
so on compact $d$-ranges,
\[
|\Xi_{\toy}^{(N)}(u,d)|\asymp u^2.
\]
Using \eqref{eq:u2-calibration},
\begin{equation}\label{eq:Xi-toy-calibrated}
|\Xi_{\toy}^{(N)}(u,d)|\asymp \frac{x_0}{B_\zeta(m)}S(m).
\end{equation}

\begin{remark}[Work in progress: calibrated size versus microscopic toy regime]
\label{rem:wip-calibration-small-u}
As presently written, the final calibration freezes the comparison scale at a
fixed value $x=x_0(D)$, which yields
\[
u^2 \asymp \frac{x_0}{B_\zeta(m)}\,S(m).
\]
With only the crude global bounds used later in the contradiction,
\[
S(m)\ll Q^2,
\qquad
B_\zeta(m)\asymp Q,
\]
this gives at best
\[
u^2\ll x_0 Q,
\]
and therefore does not by itself force the calibrated toy parameter $u$ to
lie in the microscopic toy regime $|u|<u_0(D)$ required by
Proposition~\ref{prop:toy-n-point-direct}.

A plausible repair is to avoid freezing the scale at $x_0(D)$, and instead
to choose the comparison scale adaptively:
\[
x=x(m)\in (0,x_0(D)].
\]
Since the pairing and calibration statements are naturally formulated for
$0<x\le x_0(D)$, one expects to choose $x(m)$ so that
\[
u^2 \asymp \frac{x(m)}{B_\zeta(m)}\,S(m)
\]
is forced to remain below the toy threshold $u_0(D)^2$.  For example, if
the implicit upper bound
\[
u^2 \le C\,\frac{x}{B_\zeta(m)}\,S(m)
\]
is available uniformly, then one may try to take
\[
x(m):=
\min\!\left\{
x_0(D),\;
\frac{u_0(D)^2}{2C}\frac{B_\zeta(m)}{S(m)}
\right\},
\]
which would guarantee
\[
u^2\le \frac{u_0(D)^2}{2}.
\]

To make this repair rigorous, however, the draft must still verify that the
entire calibration chain remains valid uniformly in such a variable
$x$-regime.  In particular, one still needs to check all of the following:

\begin{enumerate}
\item The pairing and denominator estimates remain uniform for all
$0<x\le x_0(D)$, including the regime $x\to 0$, so that the calibration
functional is still well-defined and the main term continues to scale linearly
in $x$.

\item The fixed-scaled statements currently written in terms of
\[
x=B_\zeta(m)\sigma,
\qquad
|\sigma|\asymp Q^{-1},
\]
can be upgraded to a variable/shrinking regime of the form
\[
x=B_\zeta(m)\sigma,
\qquad
|\sigma|\le c\,Q^{-1},
\]
or an equivalent $0<x\le x_0(D)$ formulation, without loss of the transfer,
whitening, and calibration estimates used later.

\item The remainder bound after applying the calibrated functional remains
uniform in the variable-$x$ regime; in particular, one needs the normalized
error estimate to stay independent of $x$, or at least sufficiently mild as
$x\to 0$, so that shrinking $x$ to enforce $|u|<u_0(D)$ does not
destroy the zeta-side control.

\item The cutoff choice required in the calibration remainder estimate can
still be made uniformly when $x=x(m)$ is allowed to shrink.  Concretely, one
must check that the condition
\[
R^3 \gg \frac{1}{xQ\,S(m)}
\]
remains compatible with the admissible cutoff regime under the adaptive choice
of $x(m)$.

\item The final contradiction theorem must be reformulated so that the toy-side
main term is compared at the variable scale $x(m)$, rather than at the fixed
scale $x_0(D)$, and one must verify that the resulting lower bound still
dominates the exponentially small zeta-side bound.
\end{enumerate}

Accordingly, the small-$u$ issue does not presently appear to be a
fundamental obstruction, but the current argument should be read as
conditional on a variable-$x$ calibration rewrite that has not yet been
carried out in the draft.
\end{remark}

The following remark records the current status of the toy model in the
overall proof strategy.

\begin{remark}[Work in progress: pair-like versus finite-core endgame]
\label{rem:wip-pairlike-finitecore}
At the current stage of the draft, the toy anomaly model should be understood
as a \emph{pair-core first-odd-jet normal form}, not yet as a universal normal
form for an arbitrary finite same-scale off-line core.

More precisely, in the pure pair model one has first-odd-jet nondegeneracy:
\[
a_1^{\toy}(u,d)=\kappa(d)u^2+O(u^4),
\qquad
\kappa(d)\neq 0,
\]
so the first odd jet already detects the off-line anomaly.

However, the fixed-core analysis developed later in the draft shows that for
a general finite same-scale core, cancellation may occur at the first odd
order even when a genuine corrected anomaly survives at a higher odd order.
Thus the toy model is presently justified as the correct first-odd-jet
simplification only under a \emph{pair-like core} hypothesis (or a
perturbative enlargement thereof), not yet for an arbitrary finite core.

Accordingly, the argument now has a genuine fork:
\begin{enumerate}
\item \textbf{Pair-like branch.}
If the actual same-scale core is perturbatively pair-like in the corrected
microscopic sense, then the present toy-model first-jet obstruction should
still apply, and the contradiction can be driven by the first odd coefficient.

\item \textbf{Finite-core branch.}
If the actual same-scale core is not perturbatively pair-like, then the
correct endgame must be reformulated in terms of a genuine finite-core
anomaly model, and the obstruction should be attached not specifically to the
first odd jet, but to the \emph{first nonzero odd jet} of the corrected
transverse scalar.
\end{enumerate}

Thus the toy model is best viewed, at present, as the first-odd-jet
simplification of a pair-like core, rather than as a universal model for all
finite cores. This is one of the main places where the present manuscript
should still be read as a work in progress.
\end{remark}

\section{Local projective rigidity of the one-pair and multi-pair families}
\label{sec:projective-rigidity}

Throughout this section, $\Qabs$ denotes the abstract odd-germ quotient by
the value/gauge family (not the height scale $Q\asymp\log|m|$ used
elsewhere in the paper). For the first residual step on the cubic-generic
region, we will also introduce an explicit finite-order model
$\mathcal R_h^{(5)}$ through order $5$. The latter supplements, but does
not globally replace, the ambient abstract quotient $\Qabs$.

\subsection{Cubic leading line and the first residual class for the one-pair family}

Let
\[
F_h(z)\in \Qabs[[z^2]]
\]
denote the corrected quotient odd germ of the normalized one-pair family, where
$\Qabs$ is the abstract odd-germ quotient by the value/gauge family.

\begin{lemma}[First-order value-family generator]
\label{lem:first-order-value-family-generator}
Let
\[
\mathcal W(G_-,N,G_+):=G_-^{-1/2}NG_+^{-1/2}.
\]
Let $G_0$ and $N_0$ denote the coincident background same-point and mixed
blocks, and write
\[
X_0:=G_0^{-1/2}.
\]
Let $L$ be the Fr\'echet derivative of the inverse-square-root map
$G\mapsto G^{-1/2}$ at $G_0$, so that
\[
(G_0+\varepsilon H)^{-1/2}
=
X_0+\varepsilon L(H)+O(\varepsilon^2).
\]

If $\dot G_{-,1}$, $\dot N_1$, and $\dot G_{+,1}$ are the first-order
odd pure-value variations of the left same-point block, the mixed block, and
the right same-point block, respectively, then the induced first-order odd
coefficient in the corrected local block is
\[
B_1
=
L(\dot G_{-,1})\,N_0\,X_0
+
X_0\,\dot N_1\,X_0
+
X_0\,N_0\,L(\dot G_{+,1}).
\]
Equivalently, the first-order family $V_{\val}^{(1)}$ is the image of the
linear map
\[
(\dot G_{-,1},\dot N_1,\dot G_{+,1})
\longmapsto
L(\dot G_{-,1})\,N_0\,X_0
+
X_0\,\dot N_1\,X_0
+
X_0\,N_0\,L(\dot G_{+,1}).
\]
\end{lemma}

\begin{proof}
The corrected local block is obtained by the analytic whitening map
\[
\mathcal W(G_-,N,G_+)=G_-^{-1/2}NG_+^{-1/2}.
\]
Expand the three inputs to first odd order under a pure value deformation:
\[
G_-(s)=G_0+\varepsilon \dot G_{-,1}+O(\varepsilon^2),\qquad
N(s)=N_0+\varepsilon \dot N_1+O(\varepsilon^2),\qquad
G_+(s)=G_0+\varepsilon \dot G_{+,1}+O(\varepsilon^2).
\]
Then
\[
G_\pm(s)^{-1/2}
=
X_0+\varepsilon L(\dot G_{\pm,1})+O(\varepsilon^2),
\]
so multiplying
\[
G_-(s)^{-1/2}\,N(s)\,G_+(s)^{-1/2}
\]
to first order gives
\[
X_0N_0X_0
+
\varepsilon\Bigl(
L(\dot G_{-,1})\,N_0\,X_0
+
X_0\,\dot N_1\,X_0
+
X_0\,N_0\,L(\dot G_{+,1})
\Bigr)
+
O(\varepsilon^2).
\]
The coefficient of $\varepsilon$ is exactly the first-order generator
$B_1$.
\end{proof}

\begin{lemma}[Whitened first-order reduction]
\label{lem:first-order-generator-whitened}
At the coincident background one has
\[
N_0=G_0.
\]
If
\[
\widetilde H:=X_0HX_0
\]
denotes whitening by $X_0=G_0^{-1/2}$, then the first-order generator of
Lemma~\ref{lem:first-order-value-family-generator} becomes
\[
\widetilde B_1
=
\widetilde{\dot N}_1
-\frac12\,\widetilde{\dot G}_{-,1}
-\frac12\,\widetilde{\dot G}_{+,1}.
\]
\end{lemma}

\begin{proof}
Since $X_0G_0X_0=I$, whitening conjugates the background to the identity.
The Fr\'echet derivative of $G\mapsto G^{-1/2}$ at $I$ is
\[
L_I(H)=-\frac12 H.
\]
Therefore, after whitening the formula of
Lemma~\ref{lem:first-order-value-family-generator} becomes
\[
\widetilde B_1
=
L_I(\widetilde{\dot G}_{-,1})\,I\,I
+
I\,\widetilde{\dot N}_1\,I
+
I\,I\,L_I(\widetilde{\dot G}_{+,1})
=
\widetilde{\dot N}_1
-\frac12\,\widetilde{\dot G}_{-,1}
-\frac12\,\widetilde{\dot G}_{+,1}.
\qedhere
\]
\end{proof}

\begin{lemma}[First-order same-point span]
\label{lem:first-order-same-point-span}
The first-order pure-value same-point variations span the full symmetric
subspace:
\[
\operatorname{span}\{\widetilde{\dot G}_{-,1}\}
=
\operatorname{span}\{\widetilde{\dot G}_{+,1}\}
=
\operatorname{Sym}_2(\mathbf C).
\]
\end{lemma}

\begin{proof}
The corrected same-point block has the form
\[
G_{m,\pm}^{\corr}(s)
=
\frac1\pi
\begin{pmatrix}
2q(t_\pm) & \frac12 q'(t_\pm)\\[.5ex]
\frac12 q'(t_\pm) & \frac1{12}\bigl(q''(t_\pm)+2q(t_\pm)^3\bigr)
\end{pmatrix}.
\]
Therefore an infinitesimal pure-value variation at one side gives
\[
\delta G_\pm
=
\frac1\pi
\begin{pmatrix}
2\,\delta q_\pm & \frac12\,\delta q'_\pm\\[.5ex]
\frac12\,\delta q'_\pm &
\frac1{12}\bigl(\delta q''_\pm+6q_{\pm,0}^2\delta q_\pm\bigr)
\end{pmatrix}.
\]
As $\delta q_\pm$, $\delta q'_\pm$, and $\delta q''_\pm$ vary
independently, the off-diagonal entry, the $(1,1)$-entry, and the pure
$(2,2)$-direction may be varied independently. Hence the span is the full
space $\operatorname{Sym}_2(\mathbf C)$. Whitening by $X_0$ preserves the
symmetric subspace and is invertible, so the same conclusion holds for
$\widetilde{\dot G}_{\pm,1}$.
\end{proof}

\begin{lemma}[First-order mixed skew direction]
\label{lem:first-order-mixed-skew}
Let
\[
\mathfrak{so}_2(\mathbf C)
:=
\left\{
\begin{pmatrix}
0&\xi\\
-\xi&0
\end{pmatrix}
:\ \xi\in\mathbf C
\right\}.
\]
Then the first-order pure-value mixed variations contain the full skew line:
\[
\mathfrak{so}_2(\mathbf C)\subset \operatorname{span}\{\widetilde{\dot N}_1\}.
\]
Equivalently, the first-order mixed family contains a nonzero skew direction.
\end{lemma}

\begin{proof}
Freeze the same-point data and vary only the phase gap:
\[
\Delta_m(s)\mapsto \Delta_m(s)+\varepsilon\,\delta\Delta(s),
\qquad
\delta\Delta(s)=\theta s^3+O(s^5),
\quad \theta\neq 0.
\]
Linearizing the mixed block formula in $\varepsilon$, the two off-diagonal
entries vary with opposite sign, while the diagonal entries do not contribute
at first odd order. More precisely,
\[
\delta N_{12}(s)=\frac{\theta}{\pi}s+O(s^3),
\qquad
\delta N_{21}(s)=-\frac{\theta}{\pi}s+O(s^3),
\]
up to the overall sign convention for $\theta$. Hence the first odd mixed
coefficient of this phase-gap variation is a nonzero element of
\[
\mathfrak{so}_2(\mathbf C).
\]
Whitening by $X_0$ is invertible, so the whitened first-order mixed family
also contains a nonzero skew element. Since $\mathfrak{so}_2(\mathbf C)$ is
one-dimensional, this is equivalent to
\[
\mathfrak{so}_2(\mathbf C)\subset \operatorname{span}\{\widetilde{\dot N}_1\}.
\]
\end{proof}

\begin{proposition}[Generic triviality of the first quotient order]
\label{prop:first-order-quotient-trivial}
The first-order pure-value family fills the full matrix space:
\[
V_{\val}^{(1)}=M_2(\mathbf C).
\]
Consequently
\[
\mathcal J_1/V_{\val}^{(1)}=0,
\]
so the generic leading quotient order is at least $3$.
\end{proposition}

\begin{proof}
By Lemma~\ref{lem:first-order-generator-whitened}, the first-order generator is
\[
\widetilde B_1
=
\widetilde{\dot N}_1
-\frac12\,\widetilde{\dot G}_{-,1}
-\frac12\,\widetilde{\dot G}_{+,1}.
\]
By Lemma~\ref{lem:first-order-same-point-span}, the same-point directions
$\widetilde{\dot G}_{-,1}$ and $\widetilde{\dot G}_{+,1}$ span the full
symmetric subspace $\operatorname{Sym}_2(\mathbf C)$. By
Lemma~\ref{lem:first-order-mixed-skew}, the mixed directions contain the skew
line
\[
\mathfrak{so}_2(\mathbf C).
\]
Therefore the image of the generator contains
\[
\operatorname{Sym}_2(\mathbf C)\oplus \mathfrak{so}_2(\mathbf C)=M_2(\mathbf C).
\]
Hence
\[
V_{\val}^{(1)}=M_2(\mathbf C),
\]
and so
\[
\mathcal J_1/V_{\val}^{(1)}=0.
\]
Thus the generic leading quotient order is at least $3$.
\end{proof}

\begin{remark}[Raw versus full value family]
\label{rem:raw-vs-full-value-family}
The first-order surjectivity statement concerns the full value family as it
appears in the three-input derivative
\[
D\mathcal W_{(\bg)}[\dot G_-,\dot N,\dot G_+].
\]
This is larger than the raw swap-odd one-pair slice obtained by tying the left
and right same-point variations together. In particular, at first order the
combined family $V_{\val}^{(1)}$ fills $M_2(\mathbf C)$, while the odd
coefficients of the raw corrected local block are constrained by additional
swap/parity structure. Thus the present notation
\[
V_{\val}^{(2j+1)}
\]
bundles together at least two logically distinct sources of ambiguity: raw
pure-value odd jets and additional normalization/basis-gauge directions. The
explicit finite-order model introduced later makes this split rigorous through
order $5$; outside that local model, the ambient notation
$V_{\val}^{(2j+1)}$ remains schematic.
\end{remark}

\begin{lemma}[Rigidity of the cubic scalar $B$]
\label{lem:B-rigidity}
Let $I$ be a connected parameter interval, and define
\[
B(h):=\frac{q(h)^2q'(h)}{4}-\frac{q'''(h)}{48}.
\]
Then the following are equivalent:
\begin{enumerate}
\item $B\equiv 0$ on $I$;
\item $q'''=12q^2q'$ on $I$;
\item there exists a constant $C_0\in \mathbf C$ such that
\[
q''=4q^3+C_0
\qquad\text{on }I.
\]
\end{enumerate}
In particular, the hypothesis $B\not\equiv 0$ is equivalent to excluding
the special cubic ODE $q''=4q^3+C_0$.
\end{lemma}

\begin{proof}
By definition,
\[
48\,B = 12q^2q' - q'''.
\]
Hence $B\equiv 0$ is equivalent to
\[
q'''=(4q^3)'.
\]
Integrating on the connected interval $I$ gives
\[
q''=4q^3+C_0
\]
for some constant $C_0\in \mathbf C$. The converse is immediate by
differentiation.
\end{proof}

\begin{proposition}[Generic cubic leading quotient order]
\label{prop:cubic-leading-order}
Let $I$ be a connected one-pair parameter interval, and let
$J_3^{\pair}(h)$ denote the cubic jet of the normalized one-pair family.

Assume:
\begin{enumerate}
\item the cubic value-family satisfies
\[
V_{\val}^{(3)}
\subset
\left\{
\begin{pmatrix}
0&*\\ *&0
\end{pmatrix}
\right\};
\]
\item the background same-point block is nondegenerate, so that the diagonal
entries of $X_0=G_0^{-1/2}$ satisfy $x_1\neq x_2$;
\item the scalar $v(h)$ is not identically zero on $I$;
\item $B(h)\not\equiv 0$ on $I$, equivalently, by
Lemma~\ref{lem:B-rigidity}, $q$ does not satisfy
\[
q''=4q^3+C_0
\]
identically on $I$.
\end{enumerate}

Define the cubic-generic region
\[
\mathcal U_{\mathrm{cub}}
:=
\{\,h\in I:\ v(h)B(h)\neq 0\,\}.
\]
Then $\mathcal U_{\mathrm{cub}}$ is open and dense in $I$, and for every
$h\in \mathcal U_{\mathrm{cub}}$,
\[
[J_3^{\pair}(h)]\neq 0
\quad\text{in}\quad
\mathcal J_3/V_{\val}^{(3)}.
\]

More precisely, the image of $J_3^{\pair}(h)$ in the cubic quotient has
nonzero diagonal-traceless component
\[
v(h)\,B(h)\,(x_1-x_2)
\begin{pmatrix}
1&0\\0&-1
\end{pmatrix}.
\]
Consequently,
\[
F_h(z)=v_0(h)z^3+v_1(h)z^5+\cdots,
\qquad
v_0(h):=[J_3^{\pair}(h)]\neq 0
\]
for $h\in \mathcal U_{\mathrm{cub}}$. Since the first quotient order is
generically trivial, the generic leading quotient order is therefore exactly
$3$ on $\mathcal U_{\mathrm{cub}}$.
\end{proposition}

\begin{proof}
By assumption,
\[
\operatorname{diag}J_3^{\pair}(h)
=
v(h)\,B(h)\,(x_1-x_2)
\begin{pmatrix}
1&0\\0&-1
\end{pmatrix}.
\]
Because $V_{\val}^{(3)}$ is confined to the off-diagonal sector, it cannot
cancel any nonzero diagonal component. Hence for every $h$ with
$v(h)B(h)(x_1-x_2)\neq 0$, the cubic quotient class
$[J_3^{\pair}(h)]$ is nonzero.

Since $x_1\neq x_2$ is fixed and both $v$ and $B$ are not identically
zero, the set $\mathcal U_{\mathrm{cub}}$ is nonempty; if the dependence on
$h$ is real-analytic, its complement is a proper analytic subset, hence
$\mathcal U_{\mathrm{cub}}$ is open and dense. Finally,
Proposition~\ref{prop:first-order-quotient-trivial} shows that the leading
quotient order is at least $3$, while the present argument shows that it is
equal to $3$ on $\mathcal U_{\mathrm{cub}}$.
\end{proof}

\begin{remark}[What is genuinely generic here]
\label{rem:cubic-leading-genericity}
The cubic-leading statement depends on three logically distinct inputs:
\begin{enumerate}
\item $x_1\neq x_2$, a background same-point nondegeneracy condition;
\item $v\not\equiv 0$, a one-pair nonfreezing condition;
\item exclusion of the special ODE $q''=4q^3+C_0$, equivalently
$B\not\equiv 0$.
\end{enumerate}
This separates the structural cubic mechanism from the residual quintic
bookkeeping needed later.
\end{remark}

\begin{definition}[Cubic leading line and first residual class]
\label{def:cubic-leading-line}
For $h\in \mathcal U_{\mathrm{cub}}$, the cubic quotient coefficient
\[
v_0(h)=[J_3^{\pair}(h)]
\]
is nonzero by Proposition~\ref{prop:cubic-leading-order}. Define the cubic
leading line by
\[
\ell_h:=\mathbf C\,v_0(h)\subset \Qabs.
\]
Let $v_1(h):=[J_5^{\pair}(h)]$ denote the quintic quotient coefficient.
Its image in the residual quotient
\[
\Qabs/\ell_h
\]
is denoted by
\[
\overline v_1(h)\in \Qabs/\ell_h.
\]
\end{definition}

\subsection{Normalization gauge at the first residual step}

The phrase ``value/gauge family'' bundles together two logically distinct
sources of ambiguity:

\begin{enumerate}
\item the \emph{pure value family}, recorded coefficientwise by the spaces
$
V_{\val}^{(2j+1)}\subset \mathcal J_{2j+1},
$
and

\item the \emph{normalization gauge}, consisting of even scalar
renormalizations and even jet-basis changes.
\end{enumerate}

For the first residual step it is useful to separate these two.

\begin{definition}[Normalization gauge]
\label{def:normalization-gauge}
Let $F(z)$ be a matrix-valued odd germ. The normalization gauge acts by
\[
F(z)\longmapsto a(z)\,P(z)\,F(z)\,P(z)^{-1},
\]
where
\[
a(z)=1+c_2 z^2+c_4 z^4+\cdots
\]
is an even scalar germ with $a(0)=1$, and
\[
P(z)=I+B_2 z^2+B_4 z^4+\cdots
\]
is an even $2\times 2$ matrix germ with $P(0)=I$.
\end{definition}

\begin{proposition}[First-residual transformation law]
\label{prop:first-residual-gauge-law}
Let
\[
F(z)=A_3 z^3+A_5 z^5+O(z^7)
\]
be an odd matrix germ, and let
\[
F^{\new}(z):=a(z)\,P(z)\,F(z)\,P(z)^{-1}
\]
with $a,P$ as in Definition~\ref{def:normalization-gauge}. Then
\[
A_3^{\new}=A_3,
\qquad
A_5^{\new}=A_5+c_2A_3+[B_2,A_3].
\]
\end{proposition}

\begin{proof}
Expand
\[
P(z)^{-1}=I-B_2 z^2+O(z^4).
\]
Then
\[
P(z)F(z)P(z)^{-1}
=
(I+B_2 z^2)\,(A_3 z^3+A_5 z^5)\,(I-B_2 z^2)+O(z^7),
\]
so
\[
P(z)F(z)P(z)^{-1}
=
A_3 z^3+\bigl(A_5+[B_2,A_3]\bigr)z^5+O(z^7).
\]
Multiplying by
\[
a(z)=1+c_2 z^2+O(z^4)
\]
gives
\[
F^{\new}(z)
=
A_3 z^3+\bigl(A_5+c_2A_3+[B_2,A_3]\bigr)z^5+O(z^7),
\]
as claimed.
\end{proof}

\begin{corollary}[Residual effect of normalization gauge]
\label{cor:first-residual-gauge-law}
Let $\ell:=\mathbf C A_3$. Under the normalization gauge,
the quintic coefficient changes by
\[
A_5 \longmapsto A_5 + c_2 A_3 + [B_2,A_3].
\]
Hence its class in the residual quotient modulo the leading line satisfies
\[
[A_5] \longmapsto [A_5+[B_2,A_3]]
\qquad\text{in}\qquad
\mathcal J_5/\ell.
\]
In particular, scalar renormalization is invisible in the first residual
quotient.
\end{corollary}

\begin{proof}
This is immediate from Proposition~\ref{prop:first-residual-gauge-law}, since
$c_2A_3\in \ell$.
\end{proof}

\begin{corollary}[Trace invariance under normalization gauge]
\label{cor:trace-invariant-under-normalization-gauge}
If $A_3$ is diagonal-traceless, then every commutator $[B_2,A_3]$ is
traceless. Consequently
\[
\operatorname{tr}(A_5^{\new})=\operatorname{tr}(A_5)
\]
for every normalization-gauge transform $F^{\new}$ of $F$.
\end{corollary}

\begin{proof}
The scalar renormalization term $c_2A_3$ is traceless because $A_3$ is.
The commutator term is traceless because
\[
\operatorname{tr}[B_2,A_3]=0.
\]
Hence Proposition~\ref{prop:first-residual-gauge-law} gives
\[
\operatorname{tr}(A_5^{\new})=\operatorname{tr}(A_5).
\]
\end{proof}

\begin{proposition}[Second-residual transformation law]
\label{prop:second-residual-gauge-law}
Let
\[
F(z)=A_3 z^3+A_5 z^5+A_7 z^7+O(z^9)
\]
be an odd matrix germ, and let
\[
F^{\new}(z):=a(z)\,P(z)\,F(z)\,P(z)^{-1}
\]
with $a,P$ as in Definition~\ref{def:normalization-gauge}. Then
\[
A_3^{\new}=A_3,
\qquad
A_5^{\new}=A_5+c_2A_3+[B_2,A_3],
\]
and the septic coefficient transforms by
\[
A_7^{\new}
=
A_7+c_2A_5+c_4A_3+[B_2,A_5]+[B_4,A_3]
+c_2[B_2,A_3]+A_3B_2^2-B_2A_3B_2.
\]
\end{proposition}

\begin{proof}
By Proposition~\ref{prop:first-residual-gauge-law}, the formulas for
$A_3^{\new}$ and $A_5^{\new}$ are already known. It remains to compute the
$z^7$-coefficient.

Expand
\[
P(z)^{-1}=I-B_2 z^2+(B_2^2-B_4)z^4+O(z^6).
\]
Then
\[
P(z)F(z)P(z)^{-1}
=
(I+B_2 z^2+B_4 z^4)(A_3 z^3+A_5 z^5+A_7 z^7)
(I-B_2 z^2+(B_2^2-B_4)z^4)+O(z^9).
\]
Collecting the $z^7$-terms gives
\[
A_7+[B_2,A_5]+[B_4,A_3]+A_3B_2^2-B_2A_3B_2.
\]
Multiplying by
\[
a(z)=1+c_2 z^2+c_4 z^4+O(z^6)
\]
adds the scalar contributions
\[
c_2A_5+c_4A_3+c_2[B_2,A_3].
\]
Hence
\[
A_7^{\new}
=
A_7+c_2A_5+c_4A_3+[B_2,A_5]+[B_4,A_3]
+c_2[B_2,A_3]+A_3B_2^2-B_2A_3B_2,
\]
as claimed.
\end{proof}

\begin{lemma}[Quintic trace witness]
\label{lem:quintic-trace-witness}
Let $J_{5,K_1}(h)$ denote the $K_1$-transport contribution to the quintic
jet. At the trace level, the relevant contribution is
\[
\operatorname{tr}J_{5,K_1}(h)
=
\kappa(h)\,\operatorname{tr}\!\Bigl(
X_0JX_4(h)-X_4(h)JX_0
+
X_3(h)JX_1(h)-X_1(h)JX_3(h)
\Bigr),
\]
since
\[
\operatorname{tr}\bigl(X_2(h)JX_2(h)\bigr)=0
\]
for symmetric $X_2(h)$. Therefore
\[
\operatorname{tr}J_{5,K_1}(h)
=
-2\kappa(h)\Bigl((x_1-x_2)r(h)+\beta(h)(u(h)-w(h))+v(h)(\gamma(h)-\alpha(h))\Bigr),
\]
where $K_1=\kappa(h)J$, $v(h)$ is the off-diagonal part of $X_1(h)$,
and $\gamma(h)-\alpha(h)$ is the diagonal-traceless part of $X_3(h)$.

Moreover, the linear $G_3$-contribution to $\gamma(h)-\alpha(h)$ is
\[
(\gamma-\alpha)^{\lin}
=
\frac{1}{2\pi}\left[
x_1^3\frac{q'''}{24}
-
x_2^3\left(
\frac{q'''''}{576}
+\frac{q^2q'''}{96}
+\frac{qq'q''}{16}
+\frac{(q')^3}{48}
\right)
\right].
\]
Since
\[
v(h)=-\frac{1}{4\pi(s_1+s_2)}\,\eta_2(h),
\]
the product $v(h)(\gamma(h)-\alpha(h))^{\lin}$ contains the monomial
\[
\kappa(h)\,\eta_2(h)\,q'''''\!(h)
\]
with nonzero coefficient.
\end{lemma}

\begin{lemma}[Parity of the whitening expansion at a diagonal base]
\label{lem:whitening-parity}
Let
\[
G_0=\diag(\lambda_1,\lambda_2),
\qquad \lambda_1,\lambda_2>0,
\]
and let
\[
X(H):=(G_0+H)^{-1/2}
\]
for $H$ sufficiently small. Write
\[
E:=\left\{
\begin{pmatrix}
*&0\\0&*
\end{pmatrix}
\right\},
\qquad
O:=\left\{
\begin{pmatrix}
0&*\\ *&0
\end{pmatrix}
\right\}.
\]
Then every multilinear Taylor coefficient of $X(H)$ has definite matrix
parity: a coefficient depending on perturbation factors
$H_1,\dots,H_n$ lies in $E$ if an even number of the $H_j$ lie in
$O$, and lies in $O$ if an odd number of the $H_j$ lie in $O$.

In particular, diagonal input gives diagonal output at linear order, and a
bilinear coefficient with one diagonal and one off-diagonal input is
off-diagonal.
\end{lemma}

\begin{proof}
Use the holomorphic functional calculus
\[
(G_0+H)^{-1/2}
=
\frac{1}{2\pi i}\oint_{\Gamma} z^{-1/2}(zI-G_0-H)^{-1}\,dz,
\]
where $\Gamma$ encloses the spectrum of $G_0+H$. Expanding the resolvent,
\[
(zI-G_0-H)^{-1}
=
R(z)+R(z)HR(z)+R(z)HR(z)HR(z)+\cdots,
\]
with
\[
R(z):=(zI-G_0)^{-1},
\]
every factor $R(z)$ is diagonal because $G_0$ is diagonal. Hence the
matrix parity of each word
\[
R(z)H_{i_1}R(z)H_{i_2}\cdots R(z)H_{i_n}R(z)
\]
is exactly the parity of the product $H_{i_1}\cdots H_{i_n}$. Since
\[
E\cdot E\subset E,\qquad
E\cdot O\subset O,\qquad
O\cdot E\subset O,\qquad
O\cdot O\subset E,
\]
the parity rule follows. Integrating over $z$ preserves the subspaces
$E$ and $O$.
\end{proof}

\begin{lemma}[No-cancellation of the quintic trace witness]
\label{lem:no-cancellation-quintic-trace}
The monomial
\[
\kappa(h)\,\eta_2(h)\,q'''''\!(h)
\]
appearing in Lemma~\ref{lem:quintic-trace-witness} is not produced by any
other quintic contribution to $\operatorname{tr}J_5^{\pair}(h)$. More
precisely, on the graded slice corresponding to
\[
(K_1,\eta_2,q^{(5)}),
\]
the only contributing channels in
$\operatorname{tr}J_5^{\pair}(h)$ are
\[
-2\kappa(h)\,(x_1-x_2)r(h)
\qquad\text{and}\qquad
-2\kappa(h)\,v(h)\bigl(\gamma(h)-\alpha(h)\bigr),
\]
and their total coefficient of
$\kappa(h)\eta_2(h)q^{(5)}(h)$ is nonzero. Consequently
\[
\operatorname{tr}J_5^{\pair}(h)\not\equiv 0
\]
on the generic one-pair region.
\end{lemma}

\begin{proof}
Fix $h$ and isolate the graded slice
\[
\kappa(h)\,\eta_2(h)\,q^{(5)}(h).
\]

We first exclude all non-$K_1$ contributions. Same-point quintic terms
carry no $\kappa$-factor, so they cannot contribute. A $K_5$-transport
term leaves no same-point degree available, and a $K_3$-transport term
leaves only same-point degree $2$, whereas the factor $q^{(5)}(h)$
first appears in same-point degree $3$. Hence the target monomial can
occur only in the $K_1$-transport sector.

Now consider the same-point expansion. By the explicit formula
\[
G_{m,\pm}^{\corr}(s)
=
\frac1\pi
\begin{pmatrix}
2q(t_\pm) & \frac12 q'(t_\pm)\\[.5ex]
\frac12 q'(t_\pm) & \frac1{12}(q''(t_\pm)+2q(t_\pm)^3)
\end{pmatrix},
\]
the factor $q^{(5)}(h)$ first appears only in the cubic same-point
coefficient $G_3$, and there only in the diagonal sector, namely through
the cubic Taylor term of the $(2,2)$-entry coming from $q''(t_\pm)$.
Therefore, by Lemma~\ref{lem:whitening-parity}, the linear
$q^{(5)}$-contribution to $X_3$ is diagonal. In particular it contributes
to $\gamma-\alpha$, but not to $\beta$.

Likewise, the $\eta_2(h)$-slice of $X_1(h)$ is off-diagonal, because
\[
v(h)=-\frac{1}{4\pi(s_1+s_2)}\,\eta_2(h).
\]
Applying Lemma~\ref{lem:whitening-parity} again, the mixed
$(\eta_2,q^{(5)})$-slice of $X_4$ is off-diagonal. Hence it contributes
to $r$, but not to the diagonal difference $u-w$.

It follows that on the graded slice
$(K_1,\eta_2,q^{(5)})$, the term
\[
-2\kappa(h)\,\beta(h)\bigl(u(h)-w(h)\bigr)
\]
does not contribute at all. Thus the only surviving $K_1$-channels are
\[
-2\kappa(h)\,(x_1-x_2)r(h)
\qquad\text{and}\qquad
-2\kappa(h)\,v(h)\bigl(\gamma(h)-\alpha(h)\bigr).
\]

To compute their combined coefficient, it suffices by multilinearity to
restrict to the two perturbation directions carrying the relevant tags:
an off-diagonal first-order perturbation $A$ for the $\eta_2$-slice and
a diagonal third-order perturbation $D$ for the $q^{(5)}$-slice.
Accordingly, write
\[
G_0=\diag(a^2,b^2),
\qquad
A=
\begin{pmatrix}
0&u\\u&0
\end{pmatrix},
\qquad
D=
\begin{pmatrix}
0&0\\0&d
\end{pmatrix},
\]
and expand
\[
X(\varepsilon,\delta):=(G_0+\varepsilon A+\delta D)^{-1/2}
=
X_0+\varepsilon U+\delta V+\varepsilon\delta W+O(\varepsilon^2,\delta^2).
\]
Solving the identity
\[
X(\varepsilon,\delta)\,(G_0+\varepsilon A+\delta D)\,X(\varepsilon,\delta)=I
\]
order by order gives
\[
U_{12}=-\frac{u}{ab(a+b)},
\qquad
V_{22}=-\frac{d}{2b^3},
\qquad
W_{12}=\frac{du(a+2b)}{2ab^3(a+b)^2}.
\]
Here $U_{12}$ is the $\eta_2$-slice of $v$, $V_{22}$ is the
$q^{(5)}$-slice of $\gamma-\alpha$, and $W_{12}$ is the mixed
$(\eta_2,q^{(5)})$-slice of $r$. Therefore the total coefficient of the
target slice inside
\[
(x_1-x_2)r+v(\gamma-\alpha)
\]
is
\[
\left(\frac1a-\frac1b\right)W_{12}+U_{12}V_{22}
=
\frac{du}{a^2b^2(a+b)^2},
\]
which is nonzero.

Thus the monomial
$\kappa(h)\eta_2(h)q^{(5)}(h)$ occurs in
$\operatorname{tr}J_5^{\pair}(h)$ with nonzero total coefficient and is not
canceled by any other quintic contribution. Hence
\[
\operatorname{tr}J_5^{\pair}(h)\not\equiv 0.
\]
\end{proof}

\subsection{Explicit finite-order model of the first residual quotient}

For the first residual step it is enough to work with an explicit finite-order
model through order $5$ on the cubic-generic region. In this range the explicit
model is not merely a proxy: Proposition~\ref{prop:finite-order-comparison-qabs}
canonically identifies it with the order-$\le 5$ truncation of the ambient
abstract quotient $\Qabs$ after passage to the cubic line.

Accordingly, the first residual quintic detector is now rigorous at the
finite-order ambient level. What is not addressed here is only a stronger
global identification for later uses of $\Qabs$ beyond order $5$ and beyond the
first residual layer.

\begin{definition}[Finite-order truncation of the ambient abstract quotient]
\label{def:qabs-order-five}
Fix $h\in \mathcal U_{\mathrm{cub}}$, and let $A_3(h)$ represent the cubic
leading line $\ell_h\subset \Qabs$.

Define the order-$\le 5$ truncation of the ambient abstract quotient at $h$
to be the quotient
\[
{\Qabs}_h^{(\le 5)}
:=
\frac{\mathbf C\,A_3(h)z^3\oplus \mathcal J_5 z^5}
{\bigl(V_{\mathrm{raw}}^{(5)}(h)+\mathcal N_h^{(5)}\bigr)z^5}.
\]
Equivalently, ${\Qabs}_h^{(\le 5)}$ is the cubic--quintic odd-germ quotient in
which the quintic ambiguity is generated exactly by the raw pure-value quintic
slice and the realized normalization-gauge shadows through order $5$.

Its first residual quotient is
\[
{\Qabs}_h^{(\le 5)}/\ell_h
\cong
\mathcal J_5/
\bigl(\mathbf C\,A_3(h)+V_{\mathrm{raw}}^{(5)}(h)+\mathcal N_h^{(5)}\bigr).
\]
\end{definition}

\begin{definition}[Swap involution on cubic--quintic germs]
\label{def:sigma-cubic-quintic}
Let
\[
J_0=
\begin{pmatrix}
0&1\\
1&0
\end{pmatrix}.
\]
For a truncated odd germ
\[
F(z)=A_3 z^3+A_5 z^5+O(z^7)
\]
define
\[
\Sigma(F)(z):=J_0\,F(-z)\,J_0
=
-\,J_0A_3J_0\,z^3
-\,J_0A_5J_0\,z^5
+O(z^7).
\]
Thus the induced involution on quintic coefficients is
\[
\Sigma_5(A):=-J_0AJ_0.
\]
\end{definition}

\begin{definition}[Cubic-normalized swap-symmetric raw corrected germ through order $5$]
\label{def:raw-corrected-germ-order-five}
Fix $h\in \mathcal U_{\mathrm{cub}}$, and choose a diagonal-traceless
matrix
\[
A_3(h)
\]
representing the cubic leading line $\ell_h$.

A cubic-normalized swap-symmetric raw corrected germ through order $5$ at
$h$ is a truncated odd germ of the form
\[
R_h(z)=A_3(h)z^3+A_5 z^5+O(z^7)
\]
obtained from the corrected one-pair local block by a cubic normalization that
preserves the swap symmetry inherited from symmetric placement.

Its raw pure-value quintic slice is the coefficient space
\[
V_{\mathrm{raw}}^{(5)}(h)
\]
consisting of all quintic coefficients $A_5$ arising in this way under pure
value deformations at the fixed parameter $h$.
\end{definition}

\begin{lemma}[Raw quintic slice is anti-fixed]
\label{lem:raw-quintic-slice-antifixed}
For every $h\in\mathcal U_{\mathrm{cub}}$,
\[
V_{\mathrm{raw}}^{(5)}(h)\subset
\{A\in M_2(\mathbf C):J_0AJ_0=-A\}.
\]
In particular,
\[
\operatorname{tr}\bigl(V_{\mathrm{raw}}^{(5)}(h)\bigr)=0.
\]
\end{lemma}

\begin{proof}
By Definition~\ref{def:raw-corrected-germ-order-five}, every cubic-normalized
raw corrected representative satisfies
\[
R_h(-z)=J_0\,R_h(z)\,J_0.
\]
Writing
\[
R_h(z)=A_3(h)z^3+A_5 z^5+O(z^7),
\]
coefficient comparison at order $z^5$ gives
\[
J_0A_5J_0=-A_5.
\]
Hence every $A_5\in V_{\mathrm{raw}}^{(5)}(h)$ is anti-fixed. If
\[
A=
\begin{pmatrix}
a&b\\
c&d
\end{pmatrix}
\]
satisfies $J_0AJ_0=-A$, then
\[
\begin{pmatrix}
d&c\\
b&a
\end{pmatrix}
=
\begin{pmatrix}
-a&-b\\
-c&-d
\end{pmatrix},
\]
so $d=-a$, and therefore $\operatorname{tr}(A)=0$.
\end{proof}

\begin{definition}[Quintic ambiguity space and explicit first residual quotient]
\label{def:finite-order-quotient-model}
For $h\in \mathcal U_{\mathrm{cub}}$, let $V_{\mathrm{raw}}^{(5)}(h)$ be
the raw pure-value quintic slice from
Definition~\ref{def:raw-corrected-germ-order-five}.

Define the normalization-gauge quintic shadow space $\mathcal N_h^{(5)}$ to
be the set of all differences
\[
\widetilde A_5-A_5
\]
arising from two cubic-normalized swap-symmetric raw corrected germs
\[
R_h(z)=A_3(h)z^3+A_5 z^5+O(z^7),
\qquad
\widetilde R_h(z)=A_3(h)z^3+\widetilde A_5 z^5+O(z^7),
\]
of the same raw corrected local object through order $5$, when those two
representatives are related by the allowed normalization gauge.

By Corollary~\ref{cor:first-residual-gauge-law}, every such shadow has the
form
\[
\widetilde A_5-A_5=c_2A_3(h)+[B_2,A_3(h)],
\]
but $\mathcal N_h^{(5)}$ is defined to consist only of the actually realized
gauge shadows.

Set
\[
V_{\mathrm{amb}}^{(5)}(h)
:=
V_{\mathrm{raw}}^{(5)}(h)+\mathcal N_h^{(5)}.
\]

Define the explicit first residual quotient by
\[
\mathcal R_h^{(5)}
:=
\mathcal J_5/
\bigl(\mathbf C\,A_3(h)+V_{\mathrm{amb}}^{(5)}(h)\bigr).
\]
\end{definition}

\begin{proposition}[Finite-order comparison theorem]
\label{prop:finite-order-comparison-qabs}
For each $h\in \mathcal U_{\mathrm{cub}}$, the order-$\le 5$ truncation of
the ambient abstract quotient and the explicit first residual model agree
canonically:
\[
{\Qabs}_h^{(\le 5)}/\ell_h
\;\cong\;
\mathcal R_h^{(5)}.
\]

Equivalently, after passing to the cubic-generic region and truncating through
order $5$, the first residual ambiguity in the ambient abstract quotient is
exactly the ambiguity built into
Definition~\ref{def:finite-order-quotient-model}.
\end{proposition}

\begin{proof}
By Definition~\ref{def:qabs-order-five},
\[
{\Qabs}_h^{(\le 5)}
=
\frac{\mathbf C\,A_3(h)z^3\oplus \mathcal J_5 z^5}
{\bigl(V_{\mathrm{raw}}^{(5)}(h)+\mathcal N_h^{(5)}\bigr)z^5}.
\]
Passing to the first residual quotient by the cubic line
$\ell_h=\mathbf C A_3(h)$ gives
\[
{\Qabs}_h^{(\le 5)}/\ell_h
\cong
\mathcal J_5/
\bigl(\mathbf C\,A_3(h)+V_{\mathrm{raw}}^{(5)}(h)+\mathcal N_h^{(5)}\bigr).
\]
By Definition~\ref{def:finite-order-quotient-model},
\[
V_{\mathrm{amb}}^{(5)}(h)
=
V_{\mathrm{raw}}^{(5)}(h)+\mathcal N_h^{(5)},
\]
so the right-hand side is exactly
\[
\mathcal J_5/
\bigl(\mathbf C\,A_3(h)+V_{\mathrm{amb}}^{(5)}(h)\bigr)
=
\mathcal R_h^{(5)}.
\]
This is the claimed canonical identification.
\end{proof}

\begin{lemma}[Gauge shadows are anti-fixed]
\label{lem:gauge-shadows-antifixed}
For every $h\in\mathcal U_{\mathrm{cub}}$,
\[
\mathcal N_h^{(5)}\subset
\{A\in M_2(\mathbf C):J_0AJ_0=-A\}.
\]
In particular,
\[
\operatorname{tr}\bigl(\mathcal N_h^{(5)}\bigr)=0.
\]
\end{lemma}

\begin{proof}
Let
\[
R_h(z)=A_3(h)z^3+A_5 z^5+O(z^7),
\qquad
\widetilde R_h(z)=A_3(h)z^3+\widetilde A_5 z^5+O(z^7)
\]
be two cubic-normalized swap-symmetric representatives of the same raw
corrected local object through order $5$, related by the allowed
normalization gauge. By Definition~\ref{def:raw-corrected-germ-order-five},
both satisfy
\[
R_h(-z)=J_0R_h(z)J_0,
\qquad
\widetilde R_h(-z)=J_0\widetilde R_h(z)J_0.
\]
Comparing coefficients at order $z^5$ gives
\[
J_0A_5J_0=-A_5,
\qquad
J_0\widetilde A_5J_0=-\widetilde A_5.
\]
Subtracting,
\[
J_0(\widetilde A_5-A_5)J_0=-(\widetilde A_5-A_5).
\]
Since $\widetilde A_5-A_5\in \mathcal N_h^{(5)}$ by
Definition~\ref{def:finite-order-quotient-model}, every gauge shadow is
anti-fixed. Tracelessness follows as in
Lemma~\ref{lem:raw-quintic-slice-antifixed}.
\end{proof}

\begin{proposition}[Trace-free quintic ambiguity in the explicit finite-order model]
\label{prop:sigma-stability-quintic-ambiguity}
For every $h\in \mathcal U_{\mathrm{cub}}$,
\[
V_{\mathrm{amb}}^{(5)}(h)\subset
\{A\in M_2(\mathbf C):J_0AJ_0=-A\}.
\]
Equivalently, the involution $\Sigma_5$ preserves the quintic ambiguity
space and acts trivially on it.

In particular,
\[
\operatorname{tr}\bigl(V_{\mathrm{amb}}^{(5)}(h)\bigr)=0.
\]
\end{proposition}

\begin{proof}
Combine Lemma~\ref{lem:raw-quintic-slice-antifixed} and
Lemma~\ref{lem:gauge-shadows-antifixed}; by
Definition~\ref{def:finite-order-quotient-model},
\[
V_{\mathrm{amb}}^{(5)}(h)
=
V_{\mathrm{raw}}^{(5)}(h)+\mathcal N_h^{(5)}.
\]
\end{proof}

\begin{corollary}[Trace descends on the explicit first residual quotient]
\label{cor:trace-descends-explicit-first-residual}
For every $h\in \mathcal U_{\mathrm{cub}}$, trace annihilates
\[
\mathbf C\,A_3(h)+V_{\mathrm{amb}}^{(5)}(h).
\]
Hence trace descends to a well-defined linear functional
\[
\tau_h:\mathcal R_h^{(5)}\to \mathbf C.
\]
\end{corollary}

\begin{proof}
By Proposition~\ref{prop:sigma-stability-quintic-ambiguity}, trace vanishes on
$V_{\mathrm{amb}}^{(5)}(h)$. It also vanishes on $\mathbf C\,A_3(h)$
because $A_3(h)$ is diagonal-traceless.
\end{proof}

\begin{proposition}[Generic first residual nonvanishing in the explicit finite-order model]
\label{prop:first-residual-class-generic}
There exists an open dense subset
\[
\mathcal U_{\mathrm{res}}\subset \mathcal U_{\mathrm{cub}}
\]
such that
\[
\overline v_1(h)\neq 0
\qquad\text{in }\mathcal R_h^{(5)}
\]
for every $h\in \mathcal U_{\mathrm{res}}$.

Hence the one-pair quotient odd germ is not projectively constant on
$\mathcal U_{\mathrm{res}}$ at the explicit first residual level.
\end{proposition}

\begin{proof}
By Corollary~\ref{cor:trace-descends-explicit-first-residual}, trace defines a
linear functional
\[
\tau_h:\mathcal R_h^{(5)}\to\mathbf C.
\]
By Lemma~\ref{lem:no-cancellation-quintic-trace},
\[
\operatorname{tr}J_5^{\pair}(h)\not\equiv 0
\]
on the cubic-generic region. Hence there is an open dense subset
\[
\mathcal U_{\mathrm{res}}\subset \mathcal U_{\mathrm{cub}}
\]
on which
\[
\operatorname{tr}J_5^{\pair}(h)\neq 0.
\]
For such $h$,
\[
\tau_h(\overline v_1(h))
=
\operatorname{tr}J_5^{\pair}(h)\neq 0,
\]
so $\overline v_1(h)\neq 0$ in $\mathcal R_h^{(5)}$.
\end{proof}

\begin{remark}[Status of the explicit first residual model]
\label{rem:status-explicit-first-residual-model}
The first residual step is now closed at the finite-order ambient level.
Indeed, Proposition~\ref{prop:finite-order-comparison-qabs} identifies the
explicit quotient $\mathcal R_h^{(5)}$ canonically with the order-$\le 5$
truncation of the ambient abstract quotient $\Qabs$ modulo the cubic line on
the cubic-generic region.

Thus the remaining issue is no longer a finite-order comparison theorem for
the first residual layer. What remains, if one wants a stronger global ambient identification, is a
separate theorem showing that later uses of $\Qabs$ beyond order $5$ do not
introduce any new ambiguity relevant to the downstream multi-pair and global
arguments.
\end{remark}

\begin{remark}[Beyond the first residual trace]
\label{rem:beyond-first-residual-trace}
The quintic trace is the first scalar residual detector, but it is not the end
of the one-pair local structure. The later fixed-sector refinement
$
A_5^{\mathfrak f},A_7^{\mathfrak f},\Delta_7
$
indicates that the residual germ carries further geometric information beyond
the first trace channel.
\end{remark}

\subsection{Fixed-sector refinement at quintic and septic order}

The trace witness at quintic order suggests a more structured refinement of the
one-pair residual germ than the trace alone records. For bookkeeping purposes,
it is convenient to decompose the quintic and septic coefficients into fixed and
anti-fixed parts with respect to the involution
\[
\iota(A):=J_0AJ_0,
\qquad
J_0=
\begin{pmatrix}
0&1\\
1&0
\end{pmatrix}.
\]

Write
\[
\pi_{\mathfrak f}(A):=\frac12(A+J_0AJ_0),
\qquad
\pi_{\mathfrak a}(A):=\frac12(A-J_0AJ_0).
\]
Then
\[
\mathfrak f=\operatorname{span}\{I,S\},
\qquad
\mathfrak a=\operatorname{span}\{D,J\},
\]
where
\[
I=
\begin{pmatrix}
1&0\\0&1
\end{pmatrix},
\quad
S=
\begin{pmatrix}
0&1\\1&0
\end{pmatrix},
\quad
D=
\begin{pmatrix}
1&0\\0&-1
\end{pmatrix},
\quad
J=
\begin{pmatrix}
0&-1\\1&0
\end{pmatrix}.
\]

\begin{definition}[Fixed-sector quintic and septic coefficients]
\label{def:fixed-sector-quintic-septic}
Let
\[
F_h(z)=A_3(h)z^3+A_5(h)z^5+A_7(h)z^7+\cdots
\]
denote the normalized one-pair quotient odd germ on the cubic-generic region.
Its fixed-sector quintic and septic coefficients are defined by
\[
A_5^{\mathfrak f}(h):=\pi_{\mathfrak f}(A_5(h)),
\qquad
A_7^{\mathfrak f}(h):=\pi_{\mathfrak f}(A_7(h)).
\]
Write
\[
A_5^{\mathfrak f}(h)=u_5(h)\,I+v_5(h)\,S,
\qquad
A_7^{\mathfrak f}(h)=u_7(h)\,I+v_7(h)\,S.
\]
\end{definition}

\begin{definition}[Second residual determinant]
\label{def:second-residual-determinant}
On the cubic-generic region, define
\[
\Delta_7(h):=u_7(h)v_5(h)-u_5(h)v_7(h).
\]
\end{definition}

\begin{remark}[Interpretation]
\label{rem:fixed-sector-residual-package}
The quintic fixed-sector vector $A_5^{\mathfrak f}(h)$ refines the first
residual data beyond the trace detector alone, and the determinant
$\Delta_7(h)$ records the next residual quotient at septic order. In the
working notes, these form the first two nontrivial layers of a one-pair
residual package:
\[
c(h),\qquad A_5^{\mathfrak f}(h),\qquad \Delta_7(h).
\]
The current draft now records explicit coefficient formulas for the quintic
commutator contribution and the septic $K_1$-transport contribution in the
existing notation $(u,v,w,\alpha,\beta,\gamma,r,\dots)$. The canonical
order-$7$ fixed-sector content is the quotient class
\[
[A_7^{\mathfrak f}] \in \mathfrak f/\mathbf C A_5^{\mathfrak f},
\]
equivalently the second residual determinant
\[
\Delta_7=u_7v_5-u_5v_7.
\]
What remains non-unique is only the raw representative pair $(u_7,v_7)$,
which is defined only up to the one-dimensional septic gauge ambiguity along
the quintic fixed-sector line $\mathbf C A_5^{\mathfrak f}$. Thus the later
geometry should be formulated in terms of the canonical quotient-level order-$7$
data $[A_7^{\mathfrak f}]$ and $\Delta_7$, rather than a uniquely chosen
raw pair $(u_7,v_7)$. At the canonical level, these are the natural variables
for the geometric reformulation of the multi-pair problem; raw representatives
$(u_7,v_7)$ may be chosen locally by gauge-fixing when convenient, but are
not themselves canonical.
\end{remark}

\begin{proposition}[Projected septic gauge law]
\label{prop:projected-septic-gauge-law}
Assume the normalization gauge is symmetry-preserving in the sense that
\[
B_2,B_4\in \mathfrak f,
\qquad
A_3\in \mathfrak a,
\]
with
\[
\mathfrak f=\operatorname{span}\{I,S\},
\qquad
\mathfrak a=\operatorname{span}\{D,J\}.
\]
Then the fixed-sector septic coefficient satisfies
\[
(A_7^{\mathfrak f})^{\new}=A_7^{\mathfrak f}+c_2A_5^{\mathfrak f}.
\]
Equivalently, the septic fixed-sector ambiguity is one-dimensional and lies
along the quintic fixed-sector line.
\end{proposition}

\begin{proof}
By Proposition~\ref{prop:second-residual-gauge-law},
\[
A_7^{\new}
=
A_7+c_2A_5+c_4A_3+[B_2,A_5]+[B_4,A_3]
+c_2[B_2,A_3]+A_3B_2^2-B_2A_3B_2.
\]
Apply the fixed-sector projection $\pi_{\mathfrak f}$ term by term.

Since $A_3\in\mathfrak a$, one has
\[
\pi_{\mathfrak f}(c_4A_3)=0.
\]
Because $B_4\in\mathfrak f$ and $A_3\in\mathfrak a$,
\[
[B_4,A_3]\in[\mathfrak f,\mathfrak a]\subset\mathfrak a,
\]
so
\[
\pi_{\mathfrak f}[B_4,A_3]=0.
\]
Likewise,
\[
[B_2,A_3]\in[\mathfrak f,\mathfrak a]\subset\mathfrak a,
\]
hence
\[
\pi_{\mathfrak f}\bigl(c_2[B_2,A_3]\bigr)=0.
\]

Now decompose
\[
A_5=A_5^{\mathfrak f}+A_5^{\mathfrak a}.
\]
Since $\mathfrak f=\operatorname{span}\{I,S\}$ is abelian,
\[
[B_2,A_5^{\mathfrak f}]=0.
\]
And since $B_2\in\mathfrak f$ and $A_5^{\mathfrak a}\in\mathfrak a$,
\[
[B_2,A_5^{\mathfrak a}]\in[\mathfrak f,\mathfrak a]\subset\mathfrak a,
\]
so
\[
\pi_{\mathfrak f}[B_2,A_5]=0.
\]

Finally, $B_2\in\mathfrak f$ implies $B_2^2\in\mathfrak f$. Therefore
\[
A_3B_2^2\in\mathfrak a\mathfrak f\subset\mathfrak a,
\qquad
B_2A_3B_2\in\mathfrak f\mathfrak a\mathfrak f\subset\mathfrak a,
\]
and hence
\[
\pi_{\mathfrak f}(A_3B_2^2-B_2A_3B_2)=0.
\]

Thus the only term surviving fixed projection is
\[
c_2A_5^{\mathfrak f},
\]
and we conclude
\[
(A_7^{\mathfrak f})^{\new}=A_7^{\mathfrak f}+c_2A_5^{\mathfrak f}.
\qedhere
\]
\end{proof}

\begin{corollary}[Gauge invariance of the second residual determinant]
\label{cor:delta7-gauge-invariant}
Under a symmetry-preserving normalization gauge, the scalar
\[
\Delta_7=u_7v_5-u_5v_7
\]
is invariant.
\end{corollary}

\begin{proof}
By Proposition~\ref{prop:projected-septic-gauge-law},
\[
(A_7^{\mathfrak f})^{\new}=A_7^{\mathfrak f}+c_2A_5^{\mathfrak f}.
\]
Writing
\[
A_5^{\mathfrak f}=u_5I+v_5S,
\qquad
A_7^{\mathfrak f}=u_7I+v_7S,
\]
the new septic pair is
\[
(u_7^{\new},v_7^{\new})=(u_7,v_7)+c_2(u_5,v_5).
\]
Therefore
\[
\Delta_7^{\new}
=
(u_7+c_2u_5)v_5-u_5(v_7+c_2v_5)
=
u_7v_5-u_5v_7
=
\Delta_7.
\]
\end{proof}

\begin{definition}[Septic normalization-shadow sectors]
\label{def:septic-shadow-sectors}
In the notation of Proposition~\ref{prop:second-residual-gauge-law}, define
the order-$2$ and order-$4$ septic normalization-shadow sectors by
\[
A_{7,K_5}:=c_2A_5+[B_2,A_5],
\]
and
\[
A_{7,K_3}:=
c_4A_3+[B_4,A_3]+c_2[B_2,A_3]+A_3B_2^2-B_2A_3B_2.
\]
Then the exact septic transformation law may be written
\[
A_7^{\new}=A_7+A_{7,K_5}+A_{7,K_3}.
\]
When the normalization gauge is symmetry-preserving, these are the septic
$K_5$- and $K_3$-transport shadow sectors.
\end{definition}

\begin{proposition}[Projected septic $K_5$-shadow law]
\label{prop:projected-k5-shadow-law}
Assume the normalization gauge is symmetry-preserving in the sense of
Proposition~\ref{prop:projected-septic-gauge-law}. Then
\[
A_{7,K_5}^{\mathfrak f}=c_2A_5^{\mathfrak f}.
\]
In particular,
\[
A_{7,K_5}^{\mathfrak f}\in \mathbf C\,A_5^{\mathfrak f}.
\]
\end{proposition}

\begin{proof}
By Definition~\ref{def:septic-shadow-sectors},
\[
A_{7,K_5}=c_2A_5+[B_2,A_5].
\]
Apply the fixed-sector projection $\pi_{\mathfrak f}$. Decompose
\[
A_5=A_5^{\mathfrak f}+A_5^{\mathfrak a}.
\]
Since $\mathfrak f=\operatorname{span}\{I,S\}$ is abelian,
\[
[B_2,A_5^{\mathfrak f}]=0.
\]
And since $B_2\in\mathfrak f$ and $A_5^{\mathfrak a}\in\mathfrak a$,
\[
[B_2,A_5^{\mathfrak a}]\in[\mathfrak f,\mathfrak a]\subset\mathfrak a,
\]
so its fixed-sector projection vanishes. Therefore
\[
A_{7,K_5}^{\mathfrak f}
=
\pi_{\mathfrak f}(c_2A_5+[B_2,A_5])
=
c_2A_5^{\mathfrak f}.
\qedhere
\]
\end{proof}

\begin{proposition}[Projected septic $K_3$-shadow vanishing]
\label{prop:projected-k3-shadow-vanishing}
Assume the normalization gauge is symmetry-preserving in the sense of
Proposition~\ref{prop:projected-septic-gauge-law}. Then
\[
A_{7,K_3}^{\mathfrak f}=0.
\]
\end{proposition}

\begin{proof}
By Definition~\ref{def:septic-shadow-sectors},
\[
A_{7,K_3}
=
c_4A_3+[B_4,A_3]+c_2[B_2,A_3]+A_3B_2^2-B_2A_3B_2.
\]
Since $A_3\in\mathfrak a$, one has
\[
c_4A_3\in\mathfrak a.
\]
Because $B_4\in\mathfrak f$ and $A_3\in\mathfrak a$,
\[
[B_4,A_3]\in[\mathfrak f,\mathfrak a]\subset\mathfrak a.
\]
Likewise,
\[
[B_2,A_3]\in[\mathfrak f,\mathfrak a]\subset\mathfrak a,
\]
hence
\[
c_2[B_2,A_3]\in\mathfrak a.
\]
Finally, $B_2\in\mathfrak f$ implies $B_2^2\in\mathfrak f$, so
\[
A_3B_2^2\in\mathfrak a\mathfrak f\subset\mathfrak a,
\qquad
B_2A_3B_2\in\mathfrak f\mathfrak a\mathfrak f\subset\mathfrak a.
\]
Thus every term of $A_{7,K_3}$ is anti-fixed, and therefore
\[
A_{7,K_3}^{\mathfrak f}=\pi_{\mathfrak f}(A_{7,K_3})=0.
\qedhere
\]
\end{proof}

\begin{remark}[Interpretation of the septic shadow sectors]
\label{rem:septic-shadow-interpretation}
Propositions~\ref{prop:projected-k5-shadow-law} and
\ref{prop:projected-k3-shadow-vanishing} show that the lower-order
normalization shadows at septic order contribute no new canonical fixed-sector
information beyond the already-known quintic line
$\mathbf C A_5^{\mathfrak f}$.
\end{remark}

\begin{proposition}[Explicit fixed-sector formulas at quintic order]
\label{prop:explicit-quintic-fixed-sector}
Write
\[
X_1=
\begin{pmatrix}
u&v\\
v&w
\end{pmatrix},
\qquad
X_3=
\begin{pmatrix}
\alpha&\beta\\
\beta&\gamma
\end{pmatrix}.
\]
Then the fixed-sector part of the basic quintic commutator channel
\[
X_3JX_1-X_1JX_3
\]
is
\[
\pi_{\mathfrak f}(X_3JX_1-X_1JX_3)
=
\bigl(v(\gamma-\alpha)+\beta(u-w)\bigr)I
+
(\gamma u-\alpha w)S.
\]
Consequently, for the quintic fixed-sector coefficient
\[
A_5^{\mathfrak f}(h)=u_5(h)\,I+v_5(h)\,S,
\]
the commutator-type contribution is explicitly
\[
u_5^{(\mathrm{comm})}
=
v(\gamma-\alpha)+\beta(u-w),
\qquad
v_5^{(\mathrm{comm})}
=
\gamma u-\alpha w.
\]
\end{proposition}

\begin{proof}
A direct computation gives
\[
X_3JX_1-X_1JX_3
=
\begin{pmatrix}
-2\alpha v+2\beta u & -\alpha w+\gamma u\\[.5ex]
-\alpha w+\gamma u & -2\beta w+2\gamma v
\end{pmatrix}.
\]
The fixed-sector projection of a symmetric matrix
\[
\begin{pmatrix}
A&B\\
B&C
\end{pmatrix}
\]
is
\[
\pi_{\mathfrak f}
\begin{pmatrix}
A&B\\
B&C
\end{pmatrix}
=
\frac{A+C}{2}\,I+B\,S.
\]
Applying this to the displayed matrix gives
\[
\pi_{\mathfrak f}(X_3JX_1-X_1JX_3)
=
\frac{(-2\alpha v+2\beta u)+(-2\beta w+2\gamma v)}{2}\,I
+
(\gamma u-\alpha w)S,
\]
which simplifies to
\[
\bigl(v(\gamma-\alpha)+\beta(u-w)\bigr)I
+
(\gamma u-\alpha w)S.
\]
\end{proof}

\begin{proposition}[Explicit septic $K_1$-transport fixed-sector formulas]
\label{prop:explicit-septic-k1-fixed-sector}
Write
\[
X_0=\diag(x_1,x_2),
\qquad
X_1=
\begin{pmatrix}
u&v\\
v&w
\end{pmatrix},
\qquad
X_2=
\begin{pmatrix}
a&b\\
b&c
\end{pmatrix},
\]
\[
X_4=
\begin{pmatrix}
p&r\\
r&s
\end{pmatrix},
\qquad
X_5=
\begin{pmatrix}
\alpha_5&\beta_5\\
\beta_5&\gamma_5
\end{pmatrix},
\qquad
X_6=
\begin{pmatrix}
\mu_6&\nu_6\\
\nu_6&\lambda_6
\end{pmatrix}.
\]
Then the fixed-sector parts of the three septic $K_1$-transport channels are
\[
\pi_{\mathfrak f}(X_5JX_1-X_1JX_5)
=
\bigl(\beta_5(u-w)+v(\gamma_5-\alpha_5)\bigr)I
+
(\gamma_5u-\alpha_5w)S,
\]
\[
\pi_{\mathfrak f}(X_2JX_4-X_4JX_2)
=
\bigl(b(p-s)+r(c-a)\bigr)I
+
(cp-as)S,
\]
and
\[
\pi_{\mathfrak f}(X_0JX_6-X_6JX_0)
=
(x_2-x_1)\nu_6\,I
+
(x_2\mu_6-x_1\lambda_6)S.
\]

Consequently, if the septic $K_1$-transport fixed-sector coefficient is
written
\[
A_{7,K_1}^{\mathfrak f}(h)=u_{7,K_1}(h)\,I+v_{7,K_1}(h)\,S,
\]
then
\[
u_{7,K_1}
=
\kappa\Bigl(
(x_2-x_1)\nu_6
+
\beta_5(u-w)+v(\gamma_5-\alpha_5)
+
b(p-s)+r(c-a)
\Bigr),
\]
\[
v_{7,K_1}
=
\kappa\Bigl(
x_2\mu_6-x_1\lambda_6
+
\gamma_5u-\alpha_5w
+
cp-as
\Bigr).
\]
\end{proposition}

\begin{proof}
For two symmetric matrices
\[
X=
\begin{pmatrix}
a&b\\
b&c
\end{pmatrix},
\qquad
Y=
\begin{pmatrix}
p&q\\
q&r
\end{pmatrix},
\]
a direct calculation gives
\[
XJY-YJX
=
\begin{pmatrix}
-2aq+2bp & -ar+cp\\
-ar+cp & -2br+2cq
\end{pmatrix}.
\]
Taking fixed-sector projection yields
\[
\pi_{\mathfrak f}(XJY-YJX)
=
\bigl(b(p-r)+q(c-a)\bigr)I
+
(cp-ar)S.
\]
Applying this formula to $(X,Y)=(X_5,X_1)$ and $(X,Y)=(X_2,X_4)$ gives the
first two displayed identities.

For the $X_0/X_6$ channel, a direct multiplication gives
\[
X_0JX_6-X_6JX_0
=
\begin{pmatrix}
-2x_1\nu_6 & -x_1\lambda_6+x_2\mu_6\\
-x_1\lambda_6+x_2\mu_6 & 2x_2\nu_6
\end{pmatrix},
\]
whose fixed-sector part is
\[
(x_2-x_1)\nu_6\,I
+
(x_2\mu_6-x_1\lambda_6)S.
\]
Summing the three channels gives the stated formulas for
$u_{7,K_1}$ and $v_{7,K_1}$.
\end{proof}

\begin{definition}[Septic $K_1$-residual determinant]
\label{def:delta7-k1}
If
\[
A_{7,K_1}^{\mathfrak f}(h)=u_{7,K_1}(h)\,I+v_{7,K_1}(h)\,S,
\]
define the corresponding septic $K_1$-residual determinant by
\[
\Delta_{7,K_1}(h):=u_{7,K_1}(h)v_5(h)-u_5(h)v_{7,K_1}(h).
\]
\end{definition}

\begin{corollary}[Direct $q^{(7)}$-driven septic slice]
\label{cor:direct-q7-septic-slice}
On the direct linear $q^{(7)}$-driven slice, the septic $K_1$-transport
fixed-sector contribution is $I$-directional. More precisely,
\[
v_{7,K_1}^{(q^{(7)})}=0,
\]
and
\[
u_{7,K_1}^{(q^{(7)})}
=
\kappa\Bigl(
(x_2-x_1)\nu_6^{(q^{(7)})}
+
v\,(\gamma_5-\alpha_5)^{(q^{(7)})}
\Bigr),
\]
so in particular
\[
\Delta_7^{(q^{(7)})}=u_7^{(q^{(7)})}v_5.
\]
\end{corollary}

\begin{proof}
On the direct $q^{(7)}$-driven slice, the linear $X_5$-input is diagonal,
hence
\[
\beta_5^{(q^{(7)})}=0.
\]
The mixed $X_6$-slice is off-diagonal on this grading, so
\[
\mu_6^{(q^{(7)})}=\lambda_6^{(q^{(7)})}=0.
\]
Moreover the $X_2/X_4$-channel does not lie on the pure
$(K_1,\eta_2,q^{(7)})$-slice. Therefore the $S$-component vanishes and the
displayed formula for $u_{7,K_1}^{(q^{(7)})}$ follows from
Proposition~\ref{prop:explicit-septic-k1-fixed-sector}. Since
\[
\Delta_7=u_7v_5-u_5v_7,
\]
the direct $q^{(7)}$-slice satisfies
\[
\Delta_7^{(q^{(7)})}=u_7^{(q^{(7)})}v_5.
\qedhere
\]
\end{proof}

\begin{proposition}[Noncancellation of the direct septic $K_1$ witness]
\label{prop:direct-q7-k1-witness-nonzero}
On the graded slice
\[
(K_1,\eta_2,q^{(7)}),
\]
the only contributing channels to
\[
u_{7,K_1}^{(q^{(7)})}
\]
are
\[
\kappa(x_2-x_1)\nu_6^{(q^{(7)})}
\qquad\text{and}\qquad
\kappa\,v(\gamma_5-\alpha_5)^{(q^{(7)})},
\]
and their total coefficient of the monomial
\[
\kappa(h)\,\eta_2(h)\,q^{(7)}(h)
\]
is nonzero. Consequently,
\[
u_{7,K_1}^{(q^{(7)})}(h)\not\equiv 0
\]
on the generic one-pair region.
\end{proposition}

\begin{proof}
Fix $h$ and isolate the graded slice
\[
\kappa(h)\,\eta_2(h)\,q^{(7)}(h).
\]

By Corollary~\ref{cor:direct-q7-septic-slice}, the direct $q^{(7)}$-driven
slice is $I$-directional and is given by
\[
u_{7,K_1}^{(q^{(7)})}
=
\kappa\Bigl(
(x_2-x_1)\nu_6^{(q^{(7)})}
+
v(\gamma_5-\alpha_5)^{(q^{(7)})}
\Bigr).
\]
Thus it remains only to show that the two displayed channels do not cancel.

Now consider the same-point expansion. By the explicit formula
\[
G_{m,\pm}^{\corr}(s)
=
\frac1\pi
\begin{pmatrix}
2q(t_\pm) & \frac12 q'(t_\pm)\\[.5ex]
\frac12 q'(t_\pm) & \frac1{12}(q''(t_\pm)+2q(t_\pm)^3)
\end{pmatrix},
\]
the factor $q^{(7)}(h)$ first appears only in the quintic same-point
coefficient $G_5$, and there only in the diagonal sector, namely through
the quintic Taylor term of the $(2,2)$-entry coming from $q''(t_\pm)$.
Therefore, by Lemma~\ref{lem:whitening-parity}, the linear
$q^{(7)}$-contribution to $X_5$ is diagonal. In particular it contributes
to $\gamma_5-\alpha_5$, but not to $\beta_5$.

Likewise, the $\eta_2(h)$-slice of $X_1(h)$ is off-diagonal, because
\[
v(h)=-\frac{1}{4\pi(s_1+s_2)}\,\eta_2(h).
\]
Applying Lemma~\ref{lem:whitening-parity} again, the mixed
$(\eta_2,q^{(7)})$-slice of $X_6$ is off-diagonal. Hence it contributes
to $\nu_6$, but not to $\mu_6$ or $\lambda_6$.

Thus the only contributing channels on the graded slice
\[
(K_1,\eta_2,q^{(7)})
\]
are exactly
\[
\kappa(x_2-x_1)\nu_6^{(q^{(7)})}
\qquad\text{and}\qquad
\kappa\,v(\gamma_5-\alpha_5)^{(q^{(7)})}.
\]

To compute their combined coefficient, it suffices by multilinearity to
restrict to the two perturbation directions carrying the relevant tags:
an off-diagonal first-order perturbation $A$ for the $\eta_2$-slice and
a diagonal fifth-order perturbation $D$ for the $q^{(7)}$-slice.
Accordingly, write
\[
G_0=\diag(a^2,b^2),
\qquad
A=
\begin{pmatrix}
0&u\\u&0
\end{pmatrix},
\qquad
D=
\begin{pmatrix}
0&0\\0&d
\end{pmatrix},
\]
and expand
\[
X(\varepsilon,\delta):=(G_0+\varepsilon A+\delta D)^{-1/2}
=
X_0+\varepsilon U+\delta V+\varepsilon\delta W+O(\varepsilon^2,\delta^2).
\]
Solving the identity
\[
X(\varepsilon,\delta)\,(G_0+\varepsilon A+\delta D)\,X(\varepsilon,\delta)=I
\]
order by order gives
\[
U_{12}=-\frac{u}{ab(a+b)},
\qquad
V_{22}=-\frac{d}{2b^3},
\qquad
W_{12}=\frac{du(a+2b)}{2ab^3(a+b)^2}.
\]
Here $U_{12}$ is the $\eta_2$-slice of $v$, $V_{22}$ is the
$q^{(7)}$-slice of $\gamma_5-\alpha_5$, and $W_{12}$ is the mixed
$(\eta_2,q^{(7)})$-slice of $\nu_6$. Therefore the total coefficient of the
target slice inside
\[
(x_2-x_1)\nu_6+v(\gamma_5-\alpha_5)
\]
is
\[
\left(\frac1a-\frac1b\right)W_{12}+U_{12}V_{22}
=
\frac{du}{a^2b^2(a+b)^2},
\]
which is nonzero.

Thus the monomial
\[
\kappa(h)\,\eta_2(h)\,q^{(7)}(h)
\]
occurs in $u_{7,K_1}^{(q^{(7)})}(h)$ with nonzero total coefficient, so
\[
u_{7,K_1}^{(q^{(7)})}(h)\not\equiv 0.
\]
\end{proof}

\begin{remark}[Septic $K_1$ trace target]
\label{rem:septic-k1-trace-target}
At the graded level, the septic $K_1$-transport channel is the first odd
transport sector capable of coupling the new same-point $q^{(7)}$-slice into
the fixed-sector septic residual. In particular, the expected highest new
contribution to the septic determinant is carried by the $K_1$-channel, in
parallel with the earlier quintic trace witness at order $5$.
\end{remark}

\begin{lemma}[Parity elimination of the pure same-point septic sector]
\label{lem:parity-elimination-pure-septic}
Let $t_\pm=h\pm s/2$, and let the odd two-point germ be defined as the odd
part under the involution
\[
s\longmapsto -s
\qquad
(\text{equivalently }t_+\leftrightarrow t_-).
\]
Let $A_{7,\mathrm{pure}}$ denote the pure same-point septic sector, i.e. the
septic contribution built only from the same-point blocks and their whitening,
with no transport factor coming from the ordered pair $(t_+,t_-)$.

Then $A_{7,\mathrm{pure}}$ is even in $s$, hence vanishes in the odd
two-point germ. In particular,
\[
A_{7,\mathrm{pure}}=0
\quad\text{in the odd germ,}
\qquad
A_{7,\mathrm{pure}}^{\mathfrak f}=0.
\]
Consequently, pure same-point septic terms do not contribute to
$\Delta_7$, and in particular do not contribute to the weight-one part
$[\,\Delta_7\,]^{[1]}$.
\end{lemma}

\begin{proof}
A pure same-point septic term depends only on midpoint jet data at $h$, i.e.
on the same-point expansion at the midpoint, and carries no transport factor
coming from the ordered pair $(t_+,t_-)$. Therefore it is invariant under
exchanging the two points:
\[
t_+\leftrightarrow t_-.
\]
Since this exchange is exactly the involution $s\mapsto -s$, every pure
same-point septic term is an even function of $s$.

By definition, the odd two-point germ is obtained by projecting to the odd part
under $s\mapsto -s$. Hence every even contribution vanishes in that
projection. Therefore
\[
A_{7,\mathrm{pure}}=0
\quad\text{in the odd germ.}
\]
Applying fixed-sector projection gives
\[
A_{7,\mathrm{pure}}^{\mathfrak f}=0.
\]
Since $\Delta_7$ is extracted from the odd pair germ after fixed-sector
projection and quotient normalization, pure same-point septic terms contribute
nothing to $\Delta_7$, and therefore nothing to $[\,\Delta_7\,]^{[1]}$.
\end{proof}

\begin{corollary}[Weight-one septic reduction to the direct $K_1$ slice]
\label{cor:weight-one-septic-reduction}
Let $[\,\cdot\,]^{[1]}$ denote the part linear in the first genuinely new
same-point $q^{(7)}$-slice. Then
\[
[\Delta_7]^{[1]}=[\Delta_{7,K_1}]^{[1]}.
\]
Equivalently, the highest genuinely new $q^{(7)}$-coefficient of the septic
residual determinant comes only from the direct septic $K_1$-transport slice.
\end{corollary}

\begin{proof}
The first genuinely new $q^{(7)}$-slice enters the same-point expansion at
order $5$, so
\[
X_j^{[1]}=0
\qquad (j<5).
\]
Hence every septic transport sector built only from lower same-point orders is
weight zero; in particular, the $K_3$- and $K_5$-transport sectors do not
contribute to $[\Delta_7]^{[1]}$.

By Lemma~\ref{lem:parity-elimination-pure-septic}, pure same-point septic terms
also do not contribute to $[\Delta_7]^{[1]}$, because they vanish already in
the odd two-point germ.

Thus the only septic sector that can see the weight-one $q^{(7)}$-slice is
the direct $K_1$-transport sector, and therefore
\[
[\Delta_7]^{[1]}=[\Delta_{7,K_1}]^{[1]}.
\qedhere
\]
\end{proof}

\begin{proposition}[Generic nonvanishing of the direct septic slice under nontrivial quintic fixed-sector coefficient]
\label{prop:direct-q7-total-delta7-nonzero}
Let $J$ be a connected open interval contained in the quintic-generic region.
Assume
\[
v_5\not\equiv 0
\qquad\text{on }J.
\]
Then the direct $q^{(7)}$-driven slice of the full total septic residual
determinant is generically nonzero on $J$. Equivalently,
\[
\Delta_7^{(q^{(7)})}\not\equiv 0
\qquad\text{on }J.
\]
\end{proposition}

\begin{proof}
By Corollary~\ref{cor:weight-one-septic-reduction},
\[
[\Delta_7]^{[1]}=[\Delta_{7,K_1}]^{[1]}.
\]
Thus the highest genuinely new direct $q^{(7)}$-coefficient of the total
septic residual determinant is the same as that of the direct $K_1$-transport
slice.

By Proposition~\ref{prop:direct-q7-k1-witness-nonzero}, the direct septic
$K_1$ witness is generically nonzero on the quintic-generic region. In
particular,
\[
u_{7,K_1}^{(q^{(7)})}\not\equiv 0
\qquad\text{on }J.
\]
Moreover, by Corollary~\ref{cor:direct-q7-septic-slice}, on the direct
$q^{(7)}$-driven slice one has
\[
v_{7,K_1}^{(q^{(7)})}=0,
\qquad
\Delta_{7,K_1}^{(q^{(7)})}
=
u_{7,K_1}^{(q^{(7)})}\,v_5.
\]
Since $v_5\not\equiv 0$ on $J$, it follows that
\[
\Delta_{7,K_1}^{(q^{(7)})}\not\equiv 0
\qquad\text{on }J.
\]
Therefore
\[
\Delta_7^{(q^{(7)})}\not\equiv 0
\qquad\text{on }J,
\]
as claimed.
\end{proof}

\begin{theorem}[Quotient-septic closure]
\label{thm:quotient-septic-closure}
Let the septic odd-germ coefficient be decomposed as
\[
A_7=A_{7,K_1}+A_{7,K_5}+A_{7,K_3}+A_{7,\mathrm{pure}},
\]
where $A_{7,K_1}$ is the direct $K_1$-transport sector,
$A_{7,K_5}$ and $A_{7,K_3}$ are the septic normalization-shadow sectors
from Definition~\ref{def:septic-shadow-sectors}, and
$A_{7,\mathrm{pure}}$ is the pure same-point septic sector from
Lemma~\ref{lem:parity-elimination-pure-septic}.

Then
\[
A_7^{\mathfrak f}
=
A_{7,K_1}^{\mathfrak f}+A_{7,K_5}^{\mathfrak f},
\]
with
\[
A_{7,K_5}^{\mathfrak f}\in \mathbf C\,A_5^{\mathfrak f}.
\]
Consequently,
\[
[A_7^{\mathfrak f}]
=
[A_{7,K_1}^{\mathfrak f}]
\quad\text{in }\mathfrak f/\mathbf C A_5^{\mathfrak f},
\]
and
\[
\Delta_7=\Delta_{7,K_1}.
\]
\end{theorem}

\begin{proof}
By Proposition~\ref{prop:projected-k3-shadow-vanishing},
\[
A_{7,K_3}^{\mathfrak f}=0.
\]
By Lemma~\ref{lem:parity-elimination-pure-septic},
\[
A_{7,\mathrm{pure}}^{\mathfrak f}=0.
\]
Therefore
\[
A_7^{\mathfrak f}
=
A_{7,K_1}^{\mathfrak f}+A_{7,K_5}^{\mathfrak f}.
\]
By Proposition~\ref{prop:projected-k5-shadow-law},
\[
A_{7,K_5}^{\mathfrak f}\in \mathbf C\,A_5^{\mathfrak f}.
\]
Hence $A_7^{\mathfrak f}$ and $A_{7,K_1}^{\mathfrak f}$ define the same
class in the quotient
\[
\mathfrak f/\mathbf C A_5^{\mathfrak f}.
\]
Writing
\[
A_5^{\mathfrak f}=u_5I+v_5S,
\qquad
A_{7,K_1}^{\mathfrak f}=u_{7,K_1}I+v_{7,K_1}S,
\]
this quotient equality means precisely that
\[
A_7^{\mathfrak f}
=
A_{7,K_1}^{\mathfrak f}+\chi\,A_5^{\mathfrak f}
\]
for some scalar $\chi$. Therefore
\[
\Delta_7
=
(u_{7,K_1}+\chi u_5)v_5-u_5(v_{7,K_1}+\chi v_5)
=
u_{7,K_1}v_5-u_5v_{7,K_1}
=
\Delta_{7,K_1},
\]
where $\Delta_{7,K_1}$ is as in Definition~\ref{def:delta7-k1}.
\end{proof}

\begin{remark}[Status of the order-$7$ fixed-sector package]
\label{rem:status-full-septic-fixed-sector}
The order-$7$ fixed-sector analysis is now closed at the canonical quotient
level. By Theorem~\ref{thm:quotient-septic-closure}, the full total septic
fixed-sector class satisfies
\[
[A_7^{\mathfrak f}]
=
[A_{7,K_1}^{\mathfrak f}]
\quad\text{in }\mathfrak f/\mathbf C A_5^{\mathfrak f},
\]
and the canonical second residual determinant is
\[
\Delta_7=\Delta_{7,K_1}.
\]

What remains non-unique is only the raw representative of $A_7^{\mathfrak f}$
inside the affine line
\[
A_{7,K_1}^{\mathfrak f}+\mathbf C A_5^{\mathfrak f},
\]
as recorded by Proposition~\ref{prop:projected-septic-gauge-law}. Equivalently,
the full lower/mixed septic completion of the raw total pair $(u_7,v_7)$ is
not canonically determined. This remaining gauge-dependent choice does not
affect the quotient class $[A_7^{\mathfrak f}]$ or the residual determinant
$\Delta_7$ used in the later geometric arguments.
\end{remark}

\begin{remark}[Local gauge-fixing for raw septic representatives]
\label{rem:local-gauge-fixing-u7-v7}
Although the raw representative pair $(u_7,v_7)$ is not canonical, local
choices may be made by fixing a gauge section along the quintic fixed-sector
line. Indeed, Proposition~\ref{prop:projected-septic-gauge-law} gives
\[
(u_7^{\new},v_7^{\new})=(u_7,v_7)+c_2(u_5,v_5).
\]
Hence on any region where $v_5\neq 0$, one may choose
\[
c_2=-\frac{v_7}{v_5},
\]
so that
\[
v_7^{\new}=0,
\qquad
u_7^{\new}=\frac{\Delta_7}{v_5}.
\]
Similarly, on any region where $u_5\neq 0$, one may choose
\[
c_2=-\frac{u_7}{u_5},
\]
so that
\[
u_7^{\new}=0,
\qquad
v_7^{\new}=-\frac{\Delta_7}{u_5}.
\]
Thus raw septic representatives can be chosen locally by gauge-fixing, but
such choices are auxiliary and are not canonical.
\end{remark}

\subsection{Normalized planar curve and curvature bridge}

On the region where $c(h)\neq 0$, define the normalized scalars
\[
\rho_1(h):=\frac{v_5(h)}{c(h)},
\qquad
\rho_2(h):=\frac{\Delta_7(h)}{c(h)},
\]
and the associated planar invariant curve
\[
\Gamma(h):=\bigl(\rho_1(h),\rho_2(h)\bigr).
\]

\begin{proposition}[Non-affinity of the second normalized residual scalar]
\label{prop:rho2-not-affine-in-rho1}
Let $J$ be a connected open interval contained in the quintic-generic region
and in $\{h:c(h)\neq 0\}$. Assume
\[
v_5\not\equiv 0
\qquad\text{on }J.
\]
Then there do not exist constants $A,B\in\mathbf C$ such that
\[
\rho_2(h)\equiv A\,\rho_1(h)+B
\]
on $J$.

Equivalently,
\[
\Delta_7(h)\not\equiv A\,v_5(h)+B\,c(h)
\]
on $J$.
\end{proposition}

\begin{proof}
Suppose
\[
\Delta_7(h)\equiv A\,v_5(h)+B\,c(h)
\]
on $J$. Then the direct $q^{(7)}$-driven slice of the right-hand side
vanishes, because $v_5$ and $c$ are governed by the cubic/quintic layers
and carry no direct septic $q^{(7)}$-slice.

On the other hand, by
Proposition~\ref{prop:direct-q7-total-delta7-nonzero}, the direct
$q^{(7)}$-driven slice of $\Delta_7$ is generically nonzero on $J$.
This is a contradiction. Dividing by $c(h)$ yields
\[
\rho_2(h)\not\equiv A\,\rho_1(h)+B.
\qedhere
\]
\end{proof}

\begin{definition}[Normalized derivative numerators]
\label{def:MN-normalized-curve}
Define
\[
M(h):=c(h)v_5'(h)-c'(h)v_5(h),
\qquad
N(h):=c(h)\Delta_7'(h)-c'(h)\Delta_7(h).
\]
Then
\[
\rho_1'(h)=\frac{M(h)}{c(h)^2},
\qquad
\rho_2'(h)=\frac{N(h)}{c(h)^2}.
\]
\end{definition}

\begin{proposition}[Curvature bridge for the normalized planar curve]
\label{prop:curvature-bridge-normalized-curve}
Let
\[
K_\Gamma(h):=\rho_1'(h)\rho_2''(h)-\rho_2'(h)\rho_1''(h)
\]
denote the affine curvature numerator of the planar curve $\Gamma(h)$.
Then
\[
K_\Gamma(h)=\frac{M(h)N'(h)-N(h)M'(h)}{c(h)^4}.
\]
In particular, if $K_\Gamma\equiv 0$ on a connected open interval $J$ on
which
\[
c\neq 0,
\qquad
M\neq 0,
\qquad
v_5\not\equiv 0,
\]
then $\rho_2$ is affine in $\rho_1$ on $J$, contradicting
Proposition~\ref{prop:rho2-not-affine-in-rho1}. Hence on every such interval
one has
\[
K_\Gamma\not\equiv 0.
\]
\end{proposition}

\begin{proof}
By Definition~\ref{def:MN-normalized-curve},
\[
\rho_1'=\frac{M}{c^2},
\qquad
\rho_2'=\frac{N}{c^2}.
\]
Differentiating gives
\[
\rho_1''=\frac{M'c-2Mc'}{c^3},
\qquad
\rho_2''=\frac{N'c-2Nc'}{c^3}.
\]
Therefore
\[
K_\Gamma
=
\frac{M}{c^2}\frac{N'c-2Nc'}{c^3}
-
\frac{N}{c^2}\frac{M'c-2Mc'}{c^3}
=
\frac{MN'-NM'}{c^4}.
\]

If $K_\Gamma\equiv 0$ on a connected open set with $M\neq 0$, then
\[
0=MN'-NM'=M^2\left(\frac{N}{M}\right)',
\]
so $N/M$ is constant there:
\[
\frac{N}{M}=A.
\]
Since
\[
\frac{N}{M}=\frac{\rho_2'}{\rho_1'},
\]
it follows that
\[
\rho_2'=A\,\rho_1',
\]
hence
\[
\rho_2=A\,\rho_1+B
\]
on that connected interval. If in addition
\[
c\neq 0,
\qquad
v_5\not\equiv 0,
\]
there, this contradicts Proposition~\ref{prop:rho2-not-affine-in-rho1}. The
final claim follows.
\end{proof}

\begin{remark}[Geometric consequence]
\label{rem:geometric-consequence-normalized-curve}
Proposition~\ref{prop:curvature-bridge-normalized-curve} is the local-to-global
bridge from the order-$7$ fixed-sector analysis back to the secant/tangent
program. On the open set where
\[
c(h)\neq 0,
\qquad
M(h)\neq 0,
\qquad
M(h)N'(h)-N(h)M'(h)\neq 0,
\]
the normalized planar curve $\Gamma(h)$ decomposes into $C^2$ branches in the
affine chart $(\rho_1,\rho_2)$, each of which is strictly convex or strictly
concave. Thus the remaining global obstruction is not same-branch secancy, but
the possibility of a common four-secant line meeting two separated ECT
branches twice each.
\end{remark}

\begin{definition}[Projective Wronskian of the strengthened one-pair triple]
\label{def:projective-wronskian-c-v5-delta7}
Define
\[
W(h):=
\det\!
\begin{pmatrix}
c(h) & v_5(h) & \Delta_7(h)\\
c'(h) & v_5'(h) & \Delta_7'(h)\\
c''(h) & v_5''(h) & \Delta_7''(h)
\end{pmatrix}.
\]
\end{definition}

\begin{proposition}[Projective-Wronskian bridge]
\label{prop:projective-wronskian-bridge}
With $M(h)$ and $N(h)$ as in
Definition~\ref{def:MN-normalized-curve}, one has
\[
M(h)N'(h)-N(h)M'(h)=c(h)\,W(h).
\]
Consequently,
\[
K_\Gamma(h)=\frac{W(h)}{c(h)^3}
\]
on the region where $c(h)\neq 0$.
\end{proposition}

\begin{proof}
Using
\[
M=c\,v_5'-c'v_5,
\qquad
N=c\,\Delta_7'-c'\Delta_7,
\]
a direct expansion gives
\[
M N'-N M'
=
(c\,v_5'-c'v_5)(c\,\Delta_7''-c''\Delta_7)
-
(c\,\Delta_7'-c'\Delta_7)(c\,v_5''-c''v_5).
\]
Expanding and regrouping yields
\[
M N'-N M'
=
c\,
\det\!
\begin{pmatrix}
c & v_5 & \Delta_7\\
c' & v_5' & \Delta_7'\\
c'' & v_5'' & \Delta_7''
\end{pmatrix}
=
c\,W.
\]
The displayed formula for $K_\Gamma$ now follows from
Proposition~\ref{prop:curvature-bridge-normalized-curve}.
\end{proof}

\begin{lemma}[Exact identities for $W$ and $W'$]
\label{lem:exact-W-and-Wprime-identities}
Let
\[
F(h):=(c(h),v_5(h),\Delta_7(h)),
\qquad
M(h):=c(h)v_5'(h)-c'(h)v_5(h).
\]
Then
\[
W(h)=\det(F(h),F'(h),F''(h))
\]
admits the exact expansion
\[
W
=
M\,\Delta_7''
+
(v_5c''-cv_5'')\,\Delta_7'
+
(c'v_5''-v_5'c'')\,\Delta_7,
\]
and
\[
W'(h)=\det(F(h),F'(h),F'''(h))
\]
admits the exact expansion
\[
W'
=
M\,\Delta_7'''
+
(v_5c'''-cv_5''')\,\Delta_7'
+
(c'v_5'''-v_5'c''')\,\Delta_7.
\]
\end{lemma}

\begin{proof}
Expand the determinant $\det(F,F',F'')$ along the third column. This gives
\[
W
=
\det\!\begin{pmatrix}
c & v_5\\ c' & v_5'
\end{pmatrix}\Delta_7''
-
\det\!\begin{pmatrix}
c & v_5\\ c'' & v_5''
\end{pmatrix}\Delta_7'
+
\det\!\begin{pmatrix}
c' & v_5'\\ c'' & v_5''
\end{pmatrix}\Delta_7,
\]
which is exactly the first identity.

For $W'$, differentiate $W=\det(F,F',F'')$. The terms
$\det(F',F',F'')$ and $\det(F,F'',F'')$ vanish by repeated columns, so
\[
W'=\det(F,F',F''').
\]
Expanding along the third column gives the second formula.
\end{proof}

\begin{proposition}[Exclusion of $W$-flat intervals on the $c\neq 0,\ M\neq 0$ region]
\label{prop:no-W-flat-interval}
Let $J$ be a connected open interval on which
\[
c(h)\neq 0,
\qquad
M(h)\neq 0.
\]
Then
\[
W\not\equiv 0
\qquad\text{on }J.
\]
\end{proposition}

\begin{proof}
Suppose for contradiction that
\[
W\equiv 0
\qquad\text{on }J.
\]
By Proposition~\ref{prop:projective-wronskian-bridge},
\[
K_\Gamma(h)=\frac{W(h)}{c(h)^3}\equiv 0
\qquad\text{on }J.
\]
Since
\[
M=c\,v_5'-c'v_5
\]
and $c\neq 0$ on $J$, the hypothesis $M\neq 0$ on $J$ rules out
$v_5\equiv 0$ on $J$. Hence $v_5\not\equiv 0$ on $J$.

By Proposition~\ref{prop:curvature-bridge-normalized-curve}, it follows that
\[
\rho_2(h)=A\,\rho_1(h)+B
\]
for some constants $A,B\in\mathbf C$, i.e.
\[
\Delta_7(h)=A\,v_5(h)+B\,c(h)
\qquad\text{on }J.
\]
This contradicts Proposition~\ref{prop:rho2-not-affine-in-rho1}. Therefore
\[
W\not\equiv 0
\qquad\text{on }J.
\qedhere
\]
\end{proof}

\begin{remark}[Interpretation of the $\lambda$-reduction]
\label{rem:lambda-reduction-no-W-flatness}
On the direct highest-new slice one has
\[
\Delta_7^{[1]}=\lambda^{[1]}v_5
\]
by Corollary~\ref{cor:direct-q7-septic-slice} together with
Corollary~\ref{cor:weight-one-septic-reduction}. Thus if $W\equiv 0$ on a
connected open interval with $c\neq 0$ and $M\neq 0$, then
\[
v_5\not\equiv 0
\]
on that interval, because $M=c\,v_5'-c'v_5$. The affine dependence
\[
\Delta_7=A\,v_5+B\,c
\]
forced by Proposition~\ref{prop:no-W-flat-interval} implies
\[
\Delta_7^{[1]}=A\,v_5,
\]
hence
\[
\lambda^{[1]}\equiv A
\]
on that interval. In this sense, the present argument excludes $W$-flatness on
an $M\neq 0$ region without requiring a separate differential equation for
$\lambda^{[1]}$; what remains open is not the canonical quotient-level order-$7$
invariant, but only the lower/mixed septic completion of a raw representative
pair $(u_7,v_7)$, which is irrelevant to the quotient class
$[A_7^{\mathfrak f}]$ and to $\Delta_7$.
\end{remark}

\begin{definition}[Projective graph and dual tangent-line variables]
\label{def:mu-lambda-beta-m}
On the region where $v_5\neq 0$, define
\[
\mu:=\frac{c}{v_5},
\qquad
\lambda:=\frac{\Delta_7}{v_5}.
\]
On the subregion where $M\neq 0$, define
\[
L:=v_5\Delta_7'-v_5'\Delta_7,
\qquad
\beta:=\frac{d\lambda}{d\mu},
\qquad
m:=\lambda-\mu\beta.
\]
\end{definition}

\begin{proposition}[Closed minor system and dual tangent-line curve]
\label{prop:closed-minor-system-dual-curve}
Let
\[
F:=(c,v_5,\Delta_7),
\qquad
U:=(L,-N,M).
\]
Then
\[
F\times F'=U,
\qquad
U\times U'=W\,F.
\]
Equivalently,
\[
MN'-NM'=cW,
\qquad
ML'-LM'=v_5W,
\qquad
NL'-LN'=\Delta_7W.
\]

On the region where $v_5M\neq 0$,
\[
\mu'=-\frac{M}{v_5^2},
\qquad
\beta=-\frac{L}{M},
\qquad
m=\frac{N}{M},
\]
and
\[
\beta'=-\frac{v_5W}{M^2},
\qquad
m'=\frac{cW}{M^2}.
\]
Hence the dual tangent-line curve
\[
T(h):=(\beta(h),m(h))
\]
satisfies
\[
T'(h)=\frac{W(h)}{M(h)^2}\,(-v_5(h),c(h)).
\]
Moreover,
\[
\det(T'(h),T''(h))=\frac{W(h)^2}{M(h)^3}.
\]
\end{proposition}

\begin{proof}
The identity
\[
F\times F'=(L,-N,M)
\]
is the direct expansion of the cross product. Differentiating gives
\[
U'=F\times F''.
\]
Applying the vector identity
\[
(A\times B)\times(A\times C)=\det(A,B,C)\,A
\]
with $A=F$, $B=F'$, $C=F''$ yields
\[
U\times U'=(F\times F')\times(F\times F'')=\det(F,F',F'')\,F=W\,F,
\]
which is exactly the three displayed scalar identities.

Next,
\[
\mu'=\left(\frac{c}{v_5}\right)'
=\frac{c'v_5-cv_5'}{v_5^2}
=-\frac{M}{v_5^2},
\]
and
\[
\lambda'=\left(\frac{\Delta_7}{v_5}\right)'
=\frac{v_5\Delta_7'-v_5'\Delta_7}{v_5^2}
=\frac{L}{v_5^2},
\]
so
\[
\beta=\frac{\lambda'}{\mu'}=-\frac{L}{M}.
\]
Also
\[
\mu L+M\lambda
=
\frac{c}{v_5}(v_5\Delta_7'-v_5'\Delta_7)
+
(cv_5'-c'v_5)\frac{\Delta_7}{v_5}
=
c\Delta_7'-c'\Delta_7
=
N,
\]
hence
\[
m=\lambda-\mu\beta=\lambda+\mu\frac{L}{M}=\frac{N}{M}.
\]
The formulas for $\beta'$ and $m'$ follow immediately from
\[
ML'-LM'=v_5W,
\qquad
MN'-NM'=cW.
\]

Finally, write
\[
T'=\alpha V,
\qquad
\alpha:=\frac{W}{M^2},
\qquad
V:=(-v_5,c).
\]
Then
\[
\det(T',T'')=\alpha^2\det(V,V')
=\frac{W^2}{M^4}\det\!\begin{pmatrix}
-v_5 & -v_5'\\
c & c'
\end{pmatrix}
=
\frac{W^2}{M^3},
\]
because
\[
\det(V,V')=cv_5'-c'v_5=M.
\]
\end{proof}

\begin{corollary}[Graph curvature identities in $(\mu,\lambda)$-coordinates]
\label{cor:mu-lambda-curvature-identities}
On the region where $v_5M\neq 0$,
\[
\lambda_{\mu\mu}=\frac{v_5^3}{M^3}W,
\]
and
\[
\lambda_{\mu\mu\mu}
=
-\frac{v_5^5}{M^4}
\left(
W'+3\left(\frac{v_5'}{v_5}-\frac{M'}{M}\right)W
\right).
\]
In particular, at a point where $W=0$, $M\neq 0$, and $v_5\neq 0$,
\[
\lambda_{\mu\mu\mu}
=
-\frac{v_5^5}{M^4}W'.
\]
\end{corollary}

\begin{proof}
Since
\[
\frac{d}{d\mu}=\frac{1}{\mu'}\frac{d}{dh}
=
-\frac{v_5^2}{M}\frac{d}{dh},
\]
the identity
\[
\lambda_{\mu\mu}=\frac{v_5^3}{M^3}W
\]
follows from
\[
\lambda_\mu=\beta,
\qquad
\beta'=-\frac{v_5W}{M^2}.
\]
Differentiating once more gives
\[
\lambda_{\mu\mu\mu}
=
-\frac{v_5^2}{M}\frac{d}{dh}\!\left(\frac{v_5^3}{M^3}W\right),
\]
which expands to the displayed formula. The last statement is the specialisation
$W=0$.
\end{proof}

\begin{proposition}[Projectivized tangent-angle flow and contact identity]
\label{prop:projectivized-angle-contact-identity}
On any simply connected subinterval of the region
\[
v_5\neq 0,
\qquad
M\neq 0,
\]
choose a local lift $\Theta$ of the projectivized tangent direction by
\[
(\cos\Theta,\sin\Theta)
=
\frac{1}{\sqrt{M^2+L^2}}\,(M,-L).
\]
Then
\[
\Theta'=-\frac{v_5W}{M^2+L^2}.
\]

Moreover, with
\[
\mu=\frac{c}{v_5},
\qquad
\beta=-\frac{L}{M},
\qquad
m=\frac{N}{M},
\]
one has the exact contact identity
\[
dm+\mu\,d\beta=0.
\]
Equivalently, on every interval where $\beta'$ does not vanish,
\[
\frac{dm}{d\beta}=-\mu.
\]

In particular, at a generic turning boundary with
\[
M(h_b)=0,
\qquad
N(h_b)\neq 0,
\qquad
v_5(h_b)\neq 0,
\]
the projectivized slope direction is regular: one has
\[
L(h_b)\neq 0,
\]
so the compactified tangent direction extends smoothly across the boundary.
Thus the only critical points of the compactified tangent-direction map are the
zeros of $W$.
\end{proposition}

\begin{proof}
Differentiate the angle of the vector $(M,-L)$:
\[
\Theta'
=
\frac{M(-L)'-(-L)M'}{M^2+L^2}
=
-\frac{ML'-LM'}{M^2+L^2}.
\]
By Proposition~\ref{prop:closed-minor-system-dual-curve},
\[
ML'-LM'=v_5W,
\]
hence
\[
\Theta'=-\frac{v_5W}{M^2+L^2}.
\]

Next, by Definition~\ref{def:mu-lambda-beta-m},
\[
m=\lambda-\mu\beta.
\]
Since $\beta=d\lambda/d\mu$, one has
\[
d\lambda=\beta\,d\mu,
\]
and therefore
\[
dm=d\lambda-\beta\,d\mu-\mu\,d\beta=-\mu\,d\beta.
\]
This is exactly
\[
dm+\mu\,d\beta=0.
\]

Finally, at a generic turning boundary with $v_5(h_b)\neq 0$, one has
\[
\mu(h_b)=\frac{c(h_b)}{v_5(h_b)}\neq 0.
\]
Using
\[
N=\mu L+M\lambda
\]
from Proposition~\ref{prop:closed-minor-system-dual-curve}, and evaluating at
$h_b$ where $M(h_b)=0$, gives
\[
N(h_b)=\mu(h_b)L(h_b).
\]
Since $N(h_b)\neq 0$ and $\mu(h_b)\neq 0$, it follows that
\[
L(h_b)\neq 0.
\]
Hence
\[
M(h_b)^2+L(h_b)^2\neq 0,
\]
so the projectivized tangent direction is regular there.
\end{proof}

\begin{remark}[Determinant interpretation on a fixed-sign $W$-interval]
\label{rem:determinant-interpretation-fixed-sign-W}
On any connected interval where $c$, $M$, and $W$ are nonzero and $W$
keeps one sign, the graph $\lambda=\lambda(\mu)$ is strictly convex or
strictly concave. Consequently, for any ordered triple of distinct points on
that interval, the determinant
\[
D(s,t,u)=\det(F(s),F(t),F(u))
\]
has the corresponding fixed Vandermonde sign. Thus the unresolved global issue
is not intervalwise sign regularity, but propagation across turning seams and
possible nonadjacent recurrence of tangent-slope / tangent-line data.
\end{remark}

\begin{proposition}[Local determinant sign transport across a generic turning boundary]
\label{prop:local-determinant-sign-transport-turning}
Let $h_b$ be a generic turning boundary:
\[
M(h_b)=0,\qquad M'(h_b)\neq 0,\qquad N(h_b)\neq 0.
\]
Then there exists a neighborhood $U$ of $h_b$ and a sign
$\sigma_{\mathrm{turn}}\in\{\pm 1\}$ such that for any three distinct points
$h_1,h_2,h_3\in U$,
\[
\operatorname{sgn} D(h_1,h_2,h_3)
=
\sigma_{\mathrm{turn}}\,
\operatorname{sgn}\bigl((h_2-h_1)(h_3-h_1)(h_3-h_2)\bigr),
\]
where
\[
D(h_1,h_2,h_3):=\det(F(h_1),F(h_2),F(h_3)),
\qquad
F=(c,v_5,\Delta_7).
\]
Equivalently, the triple $(c,v_5,\Delta_7)$ is locally discretely sign-regular
of order $3$ across a generic turning boundary.
\end{proposition}

\begin{proof}
Work in the normalized affine chart
\[
x(h):=\frac{v_5(h)}{c(h)},
\qquad
y(h):=\frac{\Delta_7(h)}{c(h)}.
\]
At a generic turning boundary one has
\[
x'(h_b)=0,\qquad x''(h_b)\neq 0,\qquad y'(h_b)\neq 0.
\]
Hence $y$ is a valid local coordinate near $h_b$, and the normalized curve
may be written as
\[
x=\chi(y)
\]
with
\[
\chi'(y_b)=0,\qquad \chi''(y_b)\neq 0.
\]
After shrinking $U$, the sign of $\chi''$ is constant on $U$.

Now for $h_i\in U$, write
\[
x_i:=x(h_i),\qquad y_i:=y(h_i).
\]
Since
\[
F(h_i)=c(h_i)\,(1,x_i,y_i),
\]
we have
\[
D(h_1,h_2,h_3)
=
c(h_1)c(h_2)c(h_3)\,
\det\!\begin{pmatrix}
1&x_1&y_1\\
1&x_2&y_2\\
1&x_3&y_3
\end{pmatrix}.
\]
Using $x_i=\chi(y_i)$, the $3\times 3$ determinant factors as
\[
\det\!\begin{pmatrix}
1&x_1&y_1\\
1&x_2&y_2\\
1&x_3&y_3
\end{pmatrix}
=
(y_2-y_1)(y_3-y_1)(y_3-y_2)\,\chi[y_1,y_2,y_3],
\]
where $\chi[y_1,y_2,y_3]$ is the second divided difference of $\chi$.
By the scalar divided-difference mean-value theorem,
\[
\chi[y_1,y_2,y_3]=\frac12\,\chi''(\eta)
\]
for some $\eta$ between the $y_i$. Thus $\chi[y_1,y_2,y_3]$ has fixed sign
on $U$.

Also $c$ has fixed sign on a sufficiently small neighborhood of $h_b$, and
since $y'(h_b)\neq 0$, the map $h\mapsto y(h)$ is monotone on $U$. Hence
\[
\operatorname{sgn}\bigl((y_2-y_1)(y_3-y_1)(y_3-y_2)\bigr)
=
\operatorname{sgn}(y')^3\,
\operatorname{sgn}\bigl((h_2-h_1)(h_3-h_1)(h_3-h_2)\bigr),
\]
with $\operatorname{sgn}(y')$ constant on $U$.

Collecting the fixed signs of $c$, $y'$, and $\chi''$ gives the desired
constant $\sigma_{\mathrm{turn}}$.
\end{proof}

\begin{proposition}[Dual asymptote at a generic turning boundary]
\label{prop:dual-asymptote-turning-boundary}
Let $h_b$ be a generic turning boundary:
\[
M(h_b)=0,\qquad M'(h_b)\neq 0,\qquad N(h_b)\neq 0.
\]
Assume $v_5(h_b)\neq 0$, and write
\[
\mu_b:=\frac{c(h_b)}{v_5(h_b)},
\qquad
\lambda_b:=\frac{\Delta_7(h_b)}{v_5(h_b)}.
\]
Let
\[
T(h)=\bigl(\beta(h),m(h)\bigr)
=
\left(-\frac{L(h)}{M(h)},\frac{N(h)}{M(h)}\right)
\]
be the dual tangent-line curve on the neighboring $M\neq 0$ branches.
Then both one-sided dual arcs have the common asymptote
\[
m=-\mu_b\,\beta+\lambda_b,
\]
and, more precisely,
\[
m+\mu_b\beta-\lambda_b
=
\frac{1}{4a\,\beta}+o\!\left(\frac1\beta\right)
\qquad\text{as } |\beta|\to\infty,
\]
where $a\neq 0$ is the quadratic coefficient in the local turning-point normal
form
\[
\mu-\mu_b=a(\lambda-\lambda_b)^2+O((\lambda-\lambda_b)^3).
\]
In particular, the two adjacent dual arcs approach the same asymptote from
opposite sides. This asymptotic separation is the dual-curve explanation of
the local adjacent turning-boundary exclusion theorem proved later.
\end{proposition}

\begin{proof}
At a generic turning boundary one has $\mu'(h_b)=0$ and $\lambda'(h_b)\neq 0$,
so $\lambda$ is a valid local coordinate and
\[
\mu-\mu_b=a(\lambda-\lambda_b)^2+O((\lambda-\lambda_b)^3),
\qquad a\neq 0.
\]
Differentiating gives
\[
\frac{d\mu}{d\lambda}=2a(\lambda-\lambda_b)+O((\lambda-\lambda_b)^2),
\]
hence
\[
\beta=\frac{d\lambda}{d\mu}
=
\frac{1}{2a(\lambda-\lambda_b)}+O(1).
\]
Thus $|\beta|\to\infty$ as $h\to h_b$ from either side.

Now
\[
m=\lambda-\mu\beta,
\]
so
\[
m+\mu_b\beta-\lambda_b
=
(\lambda-\lambda_b)-(\mu-\mu_b)\beta.
\]
Using
\[
\mu-\mu_b=a(\lambda-\lambda_b)^2+O((\lambda-\lambda_b)^3)
\]
and
\[
\beta=\frac{1}{2a(\lambda-\lambda_b)}+O(1),
\]
we obtain
\[
(\mu-\mu_b)\beta
=
\frac12(\lambda-\lambda_b)+O((\lambda-\lambda_b)^2).
\]
Therefore
\[
m+\mu_b\beta-\lambda_b
=
\frac12(\lambda-\lambda_b)+O((\lambda-\lambda_b)^2).
\]
Substituting
\[
\lambda-\lambda_b=\frac{1}{2a\beta}+o\!\left(\frac1\beta\right)
\]
gives
\[
m+\mu_b\beta-\lambda_b
=
\frac{1}{4a\,\beta}+o\!\left(\frac1\beta\right).
\]
Since $\beta\to+\infty$ on one side and $\beta\to-\infty$ on the other, the
correction term changes sign, so the two one-sided dual arcs approach the same
asymptote from opposite sides.
\end{proof}

\begin{proposition}[ECT intervals for the strengthened one-pair triple]
\label{prop:ect-intervals-c-v5-delta7}
Let $I$ be a connected interval on which
\[
c(h)\neq 0,
\qquad
M(h)\neq 0,
\qquad
W(h)\neq 0.
\]
Then the triple
\[
(c,\ v_5,\ \Delta_7)
\]
is an extended complete Chebyshev system on $I$. Equivalently, every
nontrivial linear combination
\[
A\,c(h)+B\,v_5(h)+C\,\Delta_7(h)
\]
has at most two zeros on $I$, counting multiplicity.

In particular, in the affine chart
\[
x=\frac{v_5}{c},\qquad y=\frac{\Delta_7}{c},
\]
every affine line meets $\Gamma(I)$ in at most two points, counting
tangency with multiplicity.
\end{proposition}

\begin{proof}
Define
\[
x(h):=\frac{v_5(h)}{c(h)},
\qquad
y(h):=\frac{\Delta_7(h)}{c(h)}.
\]
Since the one-pair family is real-analytic on the generic region, the
functions $x$ and $y$ are $C^2$ on $I$.

By Definition~\ref{def:MN-normalized-curve},
\[
x'(h)=\frac{M(h)}{c(h)^2},
\qquad
y'(h)=\frac{N(h)}{c(h)^2}.
\]
Because $M(h)\neq 0$ and $c(h)\neq 0$ on $I$, one has
\[
x'(h)\neq 0
\qquad (h\in I),
\]
so $x$ is strictly monotone on $I$ and may be used as a parameter on
$\Gamma(I)$.

Moreover,
\[
\frac{dy}{dx}=\frac{y'(h)}{x'(h)}=\frac{N(h)}{M(h)}.
\]
Differentiating with respect to $x$ gives
\[
\frac{d^2y}{dx^2}
=
\frac{\frac{d}{dh}\bigl(N/M\bigr)}{x'(h)}
=
\frac{(M N'-N M')/M^2}{M/c^2}.
\]
Using Proposition~\ref{prop:projective-wronskian-bridge},
\[
M N'-N M'=c\,W,
\]
hence
\[
\frac{d^2y}{dx^2}
=
\frac{c(h)^3\,W(h)}{M(h)^3}.
\]
Since $c(h)\neq 0$, $M(h)\neq 0$, and $W(h)\neq 0$ on $I$, it follows that
\[
\frac{d^2y}{dx^2}\neq 0
\qquad (h\in I).
\]
Therefore $\Gamma(I)$ is the graph of a strictly convex or strictly concave
$C^2$ function $y=y(x)$.

Now let
\[
L(h):=A\,c(h)+B\,v_5(h)+C\,\Delta_7(h)
\]
be a nontrivial linear combination. Since $c(h)\neq 0$ on $I$,
\[
L(h)=0
\quad\Longleftrightarrow\quad
A+B\,x(h)+C\,y(h)=0.
\]
Thus zeros of $L$ correspond exactly to intersections of $\Gamma(I)$ with the
affine line
\[
A+B\,x+C\,y=0.
\]

If $C=0$, then the equation is
\[
A+B\,x=0,
\]
which has at most one solution because $x$ is strictly monotone on $I$.
If $C\neq 0$, the equation may be written as
\[
y=-\frac{A}{C}-\frac{B}{C}x,
\]
and a strictly convex or strictly concave graph meets an affine line in at
most two points, counting tangency with multiplicity.

Hence every nontrivial linear combination
\[
A\,c+B\,v_5+C\,\Delta_7
\]
has at most two zeros on $I$, counting multiplicity. This is exactly the ECT
property for the triple $(c,v_5,\Delta_7)$. The affine-line statement is the
same conclusion in the chart $(x,y)$.
\end{proof}

\begin{definition}[ECT interval]
\label{def:ect-interval}
An ECT interval is a connected interval
\[
I\subset\{h:\ c(h)\neq 0,\ M(h)\neq 0,\ W(h)\neq 0\}.
\]
Equivalently, by Proposition~\ref{prop:ect-intervals-c-v5-delta7}, it is a
connected interval on which the triple
\[
(c,\ v_5,\ \Delta_7)
\]
is an extended complete Chebyshev system.
\end{definition}

\begin{corollary}[At most two points from one ECT interval]
\label{cor:two-points-per-ect-interval}
Let $I$ be an ECT interval. Then every affine line in the chart
\[
x=\frac{v_5}{c},
\qquad
y=\frac{\Delta_7}{c}
\]
meets $\Gamma(I)$ in at most two points, counting tangency with multiplicity.

In particular, if a finite configuration is cut out by a single affine
relation in the normalized chart, then its contribution from any fixed ECT
interval has cardinality at most two.
\end{corollary}

\begin{proof}
This is exactly the affine-line formulation of
Proposition~\ref{prop:ect-intervals-c-v5-delta7}.
\end{proof}

\begin{remark}[Branch geometry on an ECT interval]
\label{rem:branch-geometry-ect-interval}
On an ECT interval $I$, the quantities $c$, $M$, and $W$ have fixed nonzero
sign. Hence
\[
x=\frac{v_5}{c}
\]
is a monotone parameter on $I$, and in the affine chart
\[
\Gamma(I)=\left\{\left(\frac{v_5(h)}{c(h)},\frac{\Delta_7(h)}{c(h)}\right):h\in I\right\}
\]
is a single $C^2$ branch which is strictly convex or strictly concave
according to the sign of
\[
\frac{c(h)^3W(h)}{M(h)^3}.
\]
The later tangent-separator analysis is therefore an adjacent-branch problem,
not a same-branch secancy problem.
\end{remark}

\begin{definition}[Tangent-separator kernel]
\label{def:tangent-separator-kernel}
Let
\[
F(h):=(c(h),v_5(h),\Delta_7(h)).
\]
For $s,t$ in the generic region, define
\[
K(s,t):=\det\!\bigl(F(s),F'(s),F(t)\bigr).
\]
For fixed $s$, the zeros of $t\mapsto K(s,t)$ are exactly the points
$\Gamma(t)$ lying on the affine tangent line to $\Gamma$ at $\Gamma(s)$.
Equivalently, $K(s,t)=0$ is the tangent-separator incidence condition for
the tangent at $s$ and the point at $t$.
\end{definition}

\begin{lemma}[Affine form of the tangent-separator kernel]
\label{lem:affine-form-of-K}
On the affine chart $c\neq 0$, write
\[
x(h):=\frac{v_5(h)}{c(h)},
\qquad
y(h):=\frac{\Delta_7(h)}{c(h)}.
\]
Then for every $s,t$ with $c(s)c(t)\neq 0$,
\[
K(s,t)
=
c(s)^2\,c(t)\,x'(s)
\Bigl(
\bigl(y(t)-m(s)x(t)\bigr)
-
\bigl(y(s)-m(s)x(s)\bigr)
\Bigr),
\]
where
\[
m(s)=\frac{dy}{dx}(s)=\frac{N(s)}{M(s)}.
\]
Equivalently,
\[
K(s,t)
=
c(s)^2\,c(t)\,x'(s)\,
\bigl(\phi_{m(s)}(t)-\tau_s\bigr),
\]
with
\[
\phi_{m(s)}(t):=y(t)-m(s)x(t),
\qquad
\tau_s:=y(s)-m(s)x(s).
\]
\end{lemma}

\begin{proof}
Write
\[
F(h)=\bigl(c(h),\,c(h)x(h),\,c(h)y(h)\bigr).
\]
Then
\[
F'(s)=\bigl(c'(s),\,c'(s)x(s)+c(s)x'(s),\,c'(s)y(s)+c(s)y'(s)\bigr).
\]
Subtracting $x(s)$ times the first column from the second and $y(s)$ times
the first column from the third gives
\[
K(s,t)
=
\det\!
\begin{pmatrix}
c(s) & 0 & 0\\
c'(s) & c(s)x'(s) & c(s)y'(s)\\
c(t) & c(t)(x(t)-x(s)) & c(t)(y(t)-y(s))
\end{pmatrix}.
\]
Expanding along the first row yields
\[
K(s,t)
=
c(s)^2c(t)\Bigl(x'(s)(y(t)-y(s))-y'(s)(x(t)-x(s))\Bigr).
\]
Since
\[
m(s)=\frac{y'(s)}{x'(s)},
\]
this becomes
\[
K(s,t)
=
c(s)^2c(t)x'(s)
\Bigl((y(t)-m(s)x(t))-(y(s)-m(s)x(s))\Bigr),
\]
as claimed.
\end{proof}

\begin{lemma}[Basic properties of the tangent-separator kernel]
\label{lem:tangent-separator-basic-properties}
For every $s$ in the generic region,
\[
K(s,s)=0,
\qquad
\partial_tK(s,s)=0,
\qquad
\partial_t^2K(s,s)=W(s).
\]
In particular, on any interval where $W$ has fixed sign, the tangent line at
$\Gamma(s)$ is locally supporting.
\end{lemma}

\begin{proof}
The first two identities follow because, at $t=s$, the third column of the
determinant equals $F(s)$, and after differentiating once in $t$ it equals
$F'(s)$. Differentiating a second time gives
\[
\partial_t^2K(s,s)=
\det\!\bigl(F(s),F'(s),F''(s)\bigr)=W(s).
\qedhere
\]
\end{proof}

\begin{lemma}[Compact first-zero reduction on adjacent intervals]
\label{lem:compact-first-zero-reduction}
Let $I_1<I_2$ be adjacent ECT intervals in the affine
$x=v_5/c$ ordering. Assume there are compact subintervals
\[
J_1\Subset I_1,\qquad J_2\Subset I_2
\]
such that every zero of the tangent-separator kernel
\[
K(s,t)=\det(F(s),F'(s),F(t)),
\qquad
F(h)=(c(h),v_5(h),\Delta_7(h)),
\]
with $s\in I_2$, $t\in I_1$, lies in $J_2\times J_1$.

If $K$ has a zero on $J_2\times J_1$, then there exists a zero
\[
(s_*,t_*)\in J_2\times J_1
\]
such that $s_*$ is minimal in the $x=v_5/c$ ordering among all
$s$-coordinates of zeros of $K$ on $J_2\times J_1$.
\end{lemma}

\begin{proof}
Because $K$ is continuous, its zero set in the compact set $J_2\times J_1$
is closed, hence compact. The projection of that zero set to the $s$-axis is
therefore compact as well. Since $x=v_5/c$ is monotone on $I_2$, the
induced $x$-ordering is equivalent to the interval ordering on $J_2$. Thus
the projected zero set has a minimal element $s_*$, and there exists
$t_*\in J_1$ such that $(s_*,t_*)$ is a zero of $K$.
\end{proof}

\begin{lemma}[First adjacent zero implies tangency]
\label{lem:first-adjacent-zero-implies-tangency}
In the setting of Lemma~\ref{lem:compact-first-zero-reduction}, let
\[
(s_*,t_*)\in J_2\times J_1
\]
be a zero with minimal $s$-coordinate in the $x=v_5/c$ ordering. Then
\[
\partial_t K(s_*,t_*)=0.
\]
Consequently,
\[
K(s_*,t_*)=0,
\qquad
\partial_t K(s_*,t_*)=0
\]
force the tangent line at $\Gamma(s_*)$ to be tangent also at
$\Gamma(t_*)$.
\end{lemma}

\begin{proof}
Assume for contradiction that
\[
\partial_tK(s_*,t_*)\neq 0.
\]
Since
\[
(s_*,t_*)\in J_2\times J_1\Subset I_2\times I_1,
\]
the implicit function theorem applies in the open set $I_2\times I_1$ and
gives a local $C^1$ zero curve
\[
t=t(s)
\]
through $(s_*,t_*)$, defined for $s$ in a neighborhood of $s_*$, such that
\[
K(s,t(s))=0.
\]
For $s$ sufficiently close to $s_*$ with $s<s_*$, one still has
\[
(s,t(s))\in I_2\times I_1
\]
and
\[
K(s,t(s))=0.
\]
By the zero-capture hypothesis in
Lemma~\ref{lem:compact-first-zero-reduction}, every zero of $K$ with
$s\in I_2$ and $t\in I_1$ lies in $J_2\times J_1$. Hence these nearby zeros
also lie in $J_2\times J_1$, contradicting the minimality of $s_*$ among all
zeros of $K$ on $J_2\times J_1$.

Therefore
\[
\partial_tK(s_*,t_*)=0.
\]

Finally, the identities
\[
K(s,t)=\det(F(s),F'(s),F(t)),
\qquad
\partial_tK(s,t)=\det(F(s),F'(s),F'(t))
\]
show that $K(s_*,t_*)=0$ places $\Gamma(t_*)$ on the tangent line at
$\Gamma(s_*)$, while $\partial_tK(s_*,t_*)=0$ places the tangent direction
at $t_*$ on that same line. Hence the line is tangent to $\Gamma$ at both
$\Gamma(s_*)$ and $\Gamma(t_*)$; that is, it is a common tangent line.
\end{proof}

\begin{proposition}[Adjacent slope-range separation]
\label{prop:adjacent-slope-range-separation}
Let $I_1<I_2$ be adjacent ECT intervals, meeting at a boundary point $h_b$.

\begin{enumerate}
\item If
\[
M(h_b)=0,\qquad M'(h_b)\neq 0,\qquad N(h_b)\neq 0,
\]
then the interior tangent-slope ranges of $I_1$ and $I_2$ are disjoint.

\item If
\[
W(h_b)=0,\qquad M(h_b)\neq 0,
\]
then $x=v_5/c$ remains a valid local parameter, but no adjacent tangent
separation statement is asserted here. In that case the local affine jet is
governed by
\[
\left.\frac{d^2y}{dx^2}\right|_{h_b}=0,
\qquad
\left.\frac{d^3y}{dx^3}\right|_{h_b}
=
\frac{c(h_b)^5}{M(h_b)^4}W'(h_b),
\]
so the generic behavior is an ordinary affine inflection unless further
special structure forces $W'(h_b)=0$.
\end{enumerate}
\end{proposition}

\begin{proof}
In case (1), $M(h_b)=0$ is a simple zero and $N(h_b)\neq 0$. Therefore
\[
m(h)=\frac{N(h)}{M(h)}
\]
has opposite one-sided blow-up:
\[
m(h)\to \pm\infty
\]
as $h\to h_b^\pm$. Hence the interior tangent-slope ranges on the two sides
are disjoint.

Case (2) is the local affine-jet computation recorded in
Remark~\ref{rem:w0-inflection-obstruction}; it identifies the cubic-inflection
coefficient but does not produce a tangent-separation theorem.
\end{proof}

\begin{lemma}[Endpoint tangent exclusion at a generic turning boundary]
\label{lem:endpoint-tangent-exclusion}
Let $I_1<I_2$ be adjacent ECT intervals meeting at a boundary point $h_b$,
and assume
\[
M(h_b)=0,\qquad M'(h_b)\neq 0,\qquad N(h_b)\neq 0.
\]
Then the boundary tangent line at $\Gamma(h_b)$ is vertical and meets
neither adjacent interval in its interior.
\end{lemma}

\begin{proof}
The boundary tangent line is
\[
x=\rho_1(h_b),
\qquad
\rho_1(h)=\frac{v_5(h)}{c(h)}.
\]
Since $M(h_b)=0$ is a simple zero,
\[
\rho_1'(h_b)=0,\qquad \rho_1''(h_b)\neq 0,
\]
so $\rho_1$ has a strict local extremum at $h_b$. Therefore
\[
\rho_1(h)-\rho_1(h_b)
\]
has fixed sign on each adjacent interval and vanishes only at the boundary
point. Hence the vertical tangent line meets neither adjacent interval in its
interior.
\end{proof}

\begin{theorem}[Conditional adjacent tangent-separator sign regularity across turning boundaries]
\label{thm:adjacent-tangent-separator-sign-regularity}
Let $I_1<I_2$ be adjacent ECT intervals, and assume their common boundary
point $h_b$ is a generic turning boundary:
\[
M(h_b)=0,\qquad M'(h_b)\neq 0,\qquad N(h_b)\neq 0.
\]
Assume in addition the zero-capture hypothesis of
Lemma~\ref{lem:compact-first-zero-reduction} on the adjacent pair
$I_2\times I_1$. Then
\[
K(s,t)\neq 0
\qquad
(s\in I_2,\ t\in I_1),
\]
where
\[
K(s,t)=\det(F(s),F'(s),F(t)),
\qquad
F(h)=(c(h),v_5(h),\Delta_7(h)).
\]
Equivalently, tangent lines from the right interval do not meet the adjacent
left interval across a generic turning boundary.
\end{theorem}

\begin{proof}
By Lemma~\ref{lem:endpoint-tangent-exclusion}, for a generic turning boundary
there is a neighborhood of the boundary pair on which $K(s,t)\neq 0$ for
$s\in I_2$, $t\in I_1$.

Assume for contradiction that some zero $(s_0,t_0)\in I_2\times I_1$
exists. By the zero-capture hypothesis, there are compact subintervals
\[
J_1\Subset I_1,\qquad J_2\Subset I_2
\]
such that every zero of $K$ on $I_2\times I_1$ lies in $J_2\times J_1$.
Replacing them by smaller compact subintervals still containing the zero set,
using the boundary-near zero-free region furnished by
Lemma~\ref{lem:endpoint-tangent-exclusion}, we may also assume that
$J_2\times J_1$ is disjoint from that boundary-near region.
Then the zero set of $K$ in the compact set $J_2\times J_1$ is nonempty
and closed, so by Lemma~\ref{lem:compact-first-zero-reduction} there exists a
zero $(s_*,t_*)\in J_2\times J_1$ whose $s$-coordinate is minimal in the
$x=v_5/c$ ordering among all zeros in $J_2\times J_1$.

By Lemma~\ref{lem:first-adjacent-zero-implies-tangency},
\[
K(s_*,t_*)=0,
\qquad
\partial_tK(s_*,t_*)=0,
\]
so the tangent line at $\Gamma(s_*)$ is also tangent at $\Gamma(t_*)$.

Across a generic turning boundary,
Proposition~\ref{prop:adjacent-slope-range-separation} shows that the
interior tangent-slope ranges of the two adjacent intervals are disjoint.
Hence no common interior tangent slope can occur. This contradicts the
existence of $(s_*,t_*)$. Therefore no such zero exists, and
\[
K(s,t)\neq 0
\qquad
(s\in I_2,\ t\in I_1).
\qedhere
\]
\end{proof}

\begin{remark}[Status of the turning-boundary adjacent theorem]
\label{rem:status-zero-capture-turning-boundary}
Theorem~\ref{thm:adjacent-tangent-separator-sign-regularity} remains a
conditional turning-boundary theorem: the local compact first-zero and
first-zero-implies-tangency mechanism is in place, but it still depends on a
separate zero-capture statement for the directional tangent-separator kernel
$K_{21}(s,t)$ on the full adjacent pair $I_2\times I_1$.

However, by Theorem~\ref{thm:adjacent-pair-pair-exclusion-turning}, full
zero-capture is no longer needed to exclude actual adjacent pair-pair lines.
What remains open here is therefore only the stronger support-ordering /
directional-kernel route, not the local adjacent $2+2$ exclusion itself.
\end{remark}

\begin{lemma}[Boundary-near exclusion for both directional kernels at a generic turning boundary]
\label{lem:boundary-near-directional-kernel-exclusion}
Let $I_1<I_2$ be adjacent ECT intervals meeting at a generic turning boundary
\[
M(h_b)=0,\qquad M'(h_b)\neq 0,\qquad N(h_b)\neq 0.
\]
Define
\[
K_{21}(s,t):=\det(F(s),F'(s),F(t)),
\qquad
K_{12}(t,s):=\det(F(t),F'(t),F(s)),
\]
with
\[
F(h)=(c(h),v_5(h),\Delta_7(h)).
\]
Then there is a neighborhood of the boundary pair on which both directional
kernels are nonzero:
\[
K_{21}(s,t)\neq 0,
\qquad
K_{12}(t,s)\neq 0
\]
for $s\in I_2$, $t\in I_1$ sufficiently close to $h_b$.
\end{lemma}

\begin{proof}
Write
\[
x(h):=\frac{v_5(h)}{c(h)},
\qquad
y(h):=\frac{\Delta_7(h)}{c(h)}.
\]
Since
\[
y'(h)=\frac{N(h)}{c(h)^2},
\]
the hypothesis $N(h_b)\neq 0$ implies that $y$ is a valid local
coordinate near $h_b$. Thus the normalized planar curve may be written
locally as a graph
\[
x=\chi(y).
\]
At $h_b$,
\[
\frac{dx}{dy}=\frac{M(h)}{N(h)},
\]
so
\[
\left.\frac{dx}{dy}\right|_{h_b}=0.
\]
Differentiating once more,
\[
\frac{d^2x}{dy^2}
=
\frac{c(h)^2}{N(h)}\cdot\frac{M'(h)N(h)-M(h)N'(h)}{N(h)^2},
\]
and since $M(h_b)=0$,
\[
\left.\frac{d^2x}{dy^2}\right|_{h_b}
=
\frac{c(h_b)^2\,M'(h_b)}{N(h_b)^2}\neq 0.
\]
Hence, after shrinking to a small neighborhood of $h_b$, the graph
$x=\chi(y)$ is strictly convex or strictly concave.

For such a graph, the tangent line at any point meets the graph only at its
point of tangency. Therefore no tangent line from one side can meet the
opposite side inside a sufficiently small neighborhood of the boundary.
Equivalently, both directional tangent-separator kernels are nonzero near the
boundary pair.
\end{proof}

\begin{lemma}[Mirrored first adjacent zero implies tangency]
\label{lem:mirrored-first-adjacent-zero-implies-tangency}
Let $J_1\Subset I_1$, $J_2\Subset I_2$ be compact interior subintervals
across a generic turning boundary, and define
\[
K_{12}(t,s):=\det(F(t),F'(t),F(s)),
\qquad
t\in I_1,\ s\in I_2.
\]
Assume every zero of $K_{12}$ on $I_1\times I_2$ lies in
$J_1\times J_2$. Let $(t_*,s_*)\in J_1\times J_2$ be a zero whose
$t$-coordinate is minimal in the $x=v_5/c$ ordering among all zeros of
$K_{12}$ on $I_1\times I_2$. Then
\[
\partial_s K_{12}(t_*,s_*)=0.
\]
\end{lemma}

\begin{proof}
Because $M\neq 0$ on the open interval $I_1$, the function
\[
x(h)=\frac{v_5(h)}{c(h)}
\]
is a valid local coordinate on $I_1$. Since every zero of $K_{12}$ lies in
the compact interior rectangle $J_1\times J_2$, the chosen point
$(t_*,s_*)$ is interior in the ambient domain $I_1\times I_2$.

Suppose
\[
\partial_s K_{12}(t_*,s_*)\neq 0.
\]
Then by the implicit function theorem there is a $C^1$ local zero curve
\[
s=s(t)
\]
through $(t_*,s_*)$, defined for $t$ near $t_*$, such that
\[
K_{12}(t,s(t))=0.
\]
Because $t_*$ is interior in $I_1$, there are nearby points with smaller
$x=v_5/c$ ordering still in $I_1$. For such $t$, the point
$(t,s(t))$ is also a zero of $K_{12}$, contradicting the minimality of
$t_*$. Therefore
\[
\partial_s K_{12}(t_*,s_*)=0.
\]
\end{proof}

\begin{theorem}[Conditional symmetric directional-kernel exclusion across turning boundaries]
\label{thm:conditional-symmetric-directional-kernel-exclusion}
Let $I_1<I_2$ be adjacent ECT intervals meeting at a generic turning boundary
\[
M(h_b)=0,\qquad M'(h_b)\neq 0,\qquad N(h_b)\neq 0.
\]
Assume:
\begin{enumerate}
\item the zero-capture hypothesis required in
Theorem~\ref{thm:adjacent-tangent-separator-sign-regularity} for the kernel
$K_{21}$ on $I_2\times I_1$;
\item the symmetric zero-capture hypothesis for $K_{12}$, namely that every
zero of
\[
K_{12}(t,s)=\det(F(t),F'(t),F(s))
\]
on $I_1\times I_2$ lies in some compact interior rectangle
$J_1\times J_2\Subset I_1\times I_2$.
\end{enumerate}
Then both directional kernels are zero-free:
\[
K_{21}(s,t)\neq 0
\qquad (s\in I_2,\ t\in I_1),
\]
and
\[
K_{12}(t,s)\neq 0
\qquad (t\in I_1,\ s\in I_2).
\]
\end{theorem}

\begin{proof}
The first assertion is exactly
Theorem~\ref{thm:adjacent-tangent-separator-sign-regularity}.

For the second assertion, boundary-near exclusion for $K_{12}$ is provided by
Lemma~\ref{lem:boundary-near-directional-kernel-exclusion}. Assume for
contradiction that $K_{12}$ has a zero. By the symmetric zero-capture
hypothesis, every zero lies in a compact interior rectangle $J_1\times J_2$.
Choose a zero $(t_*,s_*)\in J_1\times J_2$ whose $t$-coordinate is minimal
in the $x=v_5/c$ ordering among all zeros of $K_{12}$ on $I_1\times I_2$.

By Lemma~\ref{lem:mirrored-first-adjacent-zero-implies-tangency},
\[
\partial_s K_{12}(t_*,s_*)=0.
\]
Together with
\[
K_{12}(t_*,s_*)=0,
\]
this implies that the tangent line at $\Gamma(t_*)$ is also tangent at
$\Gamma(s_*)$. Hence the two adjacent intervals share a common interior
tangent slope, contradicting
Proposition~\ref{prop:adjacent-slope-range-separation}. Therefore
$K_{12}$ has no zero.
\end{proof}

\begin{remark}[Status of the directional-kernel reduction]
\label{rem:status-directional-kernel-reduction}
Lemma~\ref{lem:pair-pair-gives-directional-kernel-zero} reduces an adjacent
pair-pair line across a generic turning boundary to a compact interior zero of
one of the two directional tangent-separator kernels.
Theorem~\ref{thm:conditional-symmetric-directional-kernel-exclusion}
shows that the remaining turning-boundary issue is no longer orientation
selection between $K_{21}$ and $K_{12}$, but the symmetric zero-capture
problem for these two kernels.
\end{remark}

\begin{remark}[The $W=0$ obstruction and the cubic-inflection coefficient]
\label{rem:w0-inflection-obstruction}
Let
\[
x(h):=\frac{v_5(h)}{c(h)},
\qquad
y(h):=\frac{\Delta_7(h)}{c(h)}.
\]
At a boundary point $h_b$ with
\[
W(h_b)=0,\qquad M(h_b)\neq 0,
\]
the coordinate $x$ remains a valid local parameter. Since
\[
\frac{d^2y}{dx^2}=\frac{c(h)^3W(h)}{M(h)^3},
\]
one has
\[
\left.\frac{d^2y}{dx^2}\right|_{h_b}=0.
\]
Differentiating once more with respect to $x$ gives
\[
\left.\frac{d^3y}{dx^3}\right|_{h_b}
=
\frac{c(h_b)^5}{M(h_b)^4}W'(h_b).
\]
Hence, after translating the point to the origin and subtracting the tangent
line, the normalized curve has local expansion
\[
y
=
\frac{c(h_b)^5W'(h_b)}{6M(h_b)^4}\,x^3+O(x^4).
\]

Therefore a generic $W=0,\ M\neq 0$ boundary is an ordinary affine inflection
unless additional special structure forces
\[
W'(h_b)=0.
\]
In particular, the tangent-separator route is not presently justified across
such boundaries, and the real remaining geometric question is whether the
actual triple $(c,v_5,\Delta_7)$ imposes any non-generic vanishing of
$W'(h_b)$.
\end{remark}

\begin{lemma}[Fixed-slope intercept sets for a strictly convex ECT branch]
\label{lem:fixed-slope-intercept-sets}
Let $I$ be an ECT interval in the affine chart
\[
x=\frac{v_5}{c},
\qquad
y=\frac{\Delta_7}{c},
\]
and suppose $\Gamma(I)$ is a strictly convex $C^2$ graph
\[
\Gamma(I)=\{(x,\psi_I(x)):x\in J_I\}
\]
over an open interval $J_I\subset\mathbf R$.

For each slope $m\in\mathbf R$, define
\[
\mathcal B_I(m):=\{\psi_I(x)-mx:x\in J_I\}.
\]

Then:

\begin{enumerate}
\item $\mathcal B_I(m)$ is an interval.

\item A line $y=mx+b$ meets $\Gamma(I)$ if and only if
\[
b\in \mathcal B_I(m).
\]

\item If $m$ lies in the interior tangent-slope range of $I$, there is a
unique $x_I(m)\in J_I$ such that
\[
\psi_I'(x_I(m))=m,
\]
and the tangent line of slope $m$ has intercept
\[
\tau_I(m):=\psi_I(x_I(m))-m\,x_I(m).
\]
Moreover,
\[
\tau_I(m)=\inf \mathcal B_I(m).
\]

\item For such $m$, a line $y=mx+b$ is $2$-secant to $\Gamma(I)$ if and
only if
\[
b\in \operatorname{int}\mathcal B_I(m).
\]
Equivalently, the $2$-secant intercepts of slope $m$ are exactly the
interior points of the intercept interval.
\end{enumerate}
\end{lemma}

\begin{proof}
For fixed $m$, put
\[
\phi_m(x):=\psi_I(x)-mx.
\]
Since $\phi_m$ is continuous on the interval $J_I$, its image
\[
\phi_m(J_I)=\mathcal B_I(m)
\]
is an interval. This proves (1), and (2) is immediate from
\[
(x,\psi_I(x))\in \{y=mx+b\}
\iff
\psi_I(x)-mx=b.
\]

If $m$ lies in the interior tangent-slope range, strict convexity implies
that $\psi_I'$ is strictly monotone, hence there is a unique $x_I(m)$
with $\psi_I'(x_I(m))=m$. At that point,
\[
\phi_m'(x_I(m))=\psi_I'(x_I(m))-m=0,
\]
and since
\[
\phi_m''(x)=\psi_I''(x)>0,
\]
the point $x_I(m)$ is the unique minimizer of $\phi_m$. Thus
\[
\tau_I(m)=\phi_m(x_I(m))=\inf \mathcal B_I(m),
\]
which proves (3).

Finally, because $\phi_m$ is strictly convex with a unique minimum,
a level $b$ is attained at two points exactly when it lies in the interior
of the interval $\phi_m(J_I)=\mathcal B_I(m)$. This is precisely the
$2$-secant condition, proving (4).
\end{proof}

\begin{proposition}[Support-curve calculus on an ECT branch]
\label{prop:support-curve-calculus-ect}
Let $I$ be as in Lemma~\ref{lem:fixed-slope-intercept-sets}, and assume $m$
lies in the interior tangent-slope range of $I$. Then
\[
x_I'(m)=\frac{1}{\psi_I''(x_I(m))}>0,
\]
\[
\tau_I'(m)=-x_I(m),
\]
and
\[
\tau_I''(m)=-\frac{1}{\psi_I''(x_I(m))}<0.
\]
In particular, the tangent-intercept function $m\mapsto \tau_I(m)$ is strictly
concave on the interior tangent-slope range.
\end{proposition}

\begin{proof}
The defining relation
\[
\psi_I'(x_I(m))=m
\]
implies, by differentiation in $m$,
\[
\psi_I''(x_I(m))\,x_I'(m)=1,
\]
hence
\[
x_I'(m)=\frac{1}{\psi_I''(x_I(m))}>0.
\]

Now
\[
\tau_I(m)=\psi_I(x_I(m))-m\,x_I(m).
\]
Differentiate:
\[
\tau_I'(m)
=
\psi_I'(x_I(m))\,x_I'(m)-x_I(m)-m\,x_I'(m)
=
m\,x_I'(m)-x_I(m)-m\,x_I'(m)
=
-x_I(m).
\]
Differentiating once more gives
\[
\tau_I''(m)=-x_I'(m)=-\frac{1}{\psi_I''(x_I(m))}<0.
\]
\end{proof}

\begin{corollary}[Single crossing and betweenness of support curves]
\label{cor:single-crossing-betweenness-support-curves}
Let $I_i<I_j$ be two ordered strict convex ECT branches in the affine chart.
On every common interior tangent-slope interval,
\[
(\tau_{I_j}-\tau_{I_i})'(m)=-(x_{I_j}(m)-x_{I_i}(m))<0.
\]
Hence $\tau_{I_i}$ and $\tau_{I_j}$ cross at most once; equivalently, the two
ordered branches admit at most one common interior tangent slope.

More generally, let
\[
I_i<I_j<I_k
\]
be three ordered strict convex ECT branches, active on a common interior slope
interval. If $\tau_{I_i}$ and $\tau_{I_j}$ cross at $m_{ij}$ and
$\tau_{I_j}$ and $\tau_{I_k}$ cross at $m_{jk}$ with
\[
m_{ij}<m_{jk},
\]
then there is a unique
\[
m_{ik}\in(m_{ij},m_{jk})
\]
such that
\[
\tau_{I_i}(m_{ik})=\tau_{I_k}(m_{ik}).
\]
So the support curves
\[
b=\tau_I(m)
\]
form a pseudoline arrangement on every common slope window.
\end{corollary}

\begin{proof}
Since the branches are ordered in the $x=v_5/c$ coordinate, one has
\[
x_{I_i}(m)<x_{I_j}(m)
\]
whenever both tangencies exist. Therefore
\[
(\tau_{I_j}-\tau_{I_i})'(m)=-(x_{I_j}(m)-x_{I_i}(m))<0,
\]
so $\tau_{I_j}-\tau_{I_i}$ is strictly decreasing and can vanish at most once.

For the betweenness statement, set
\[
\delta_{ab}(m):=\tau_{I_b}(m)-\tau_{I_a}(m).
\]
Each $\delta_{ab}$ is strictly decreasing. If
\[
\delta_{ij}(m_{ij})=0,\qquad \delta_{jk}(m_{jk})=0,\qquad m_{ij}<m_{jk},
\]
then
\[
\delta_{ik}(m_{ij})
=
\delta_{ij}(m_{ij})+\delta_{jk}(m_{ij})
=
\delta_{jk}(m_{ij})>0,
\]
because $\delta_{jk}$ vanishes later at $m_{jk}$, and similarly
\[
\delta_{ik}(m_{jk})
=
\delta_{ij}(m_{jk})+\delta_{jk}(m_{jk})
=
\delta_{ij}(m_{jk})<0.
\]
By strict monotonicity, $\delta_{ik}$ has a unique zero in $(m_{ij},m_{jk})$.
\end{proof}

\begin{proposition}[Ordered-pair overlap windows are one-shot]
\label{prop:ordered-pair-overlap-windows-one-shot}
Let $I_1<I_2$ be two ordered strict convex ECT branches in the affine chart,
with
\[
\Gamma(I_j)=\{(x,\psi_j(x)):x\in J_j\},
\qquad
J_1<J_2.
\]
For $j=1,2$, let
\[
\mathcal B_{I_j}(m)=\{\psi_j(x)-mx:x\in J_j\},
\qquad
u_j(m):=\sup \mathcal B_{I_j}(m),
\]
and let $\tau_j(m)$ be the tangent intercept on the interior tangent-slope
range.

Then the directional overlap window
\[
\Sigma_{2\to 1}
:=
\{\,m:\tau_1(m)\le \tau_2(m)\le u_1(m)\,\}
\]
is an interval, possibly empty.

Likewise the pair-pair overlap window
\[
\Sigma^{\mathrm{pp}}_{12}
:=
\{\,m:\operatorname{int}\mathcal B_{I_1}(m)\cap
\operatorname{int}\mathcal B_{I_2}(m)\neq\varnothing\,\}
\]
is an interval, possibly empty.

More precisely:

\begin{enumerate}
\item on every slope subinterval where one fixed endpoint support line
\[
m\longmapsto y_\xi-m\xi
\]
realizes $u_1(m)$, one has
\[
\frac{d}{dm}\bigl(u_1(m)-\tau_2(m)\bigr)=x_{I_2}(m)-\xi>0;
\]

\item on every slope subinterval where one fixed endpoint support line
\[
m\longmapsto y_\eta-m\eta
\]
realizes $u_2(m)$, one has
\[
\frac{d}{dm}\bigl(u_2(m)-\tau_1(m)\bigr)=x_{I_1}(m)-\eta<0.
\]
\end{enumerate}

In particular, for a fixed ordered pair, directional or pair-pair overlap
cannot disappear and later reappear.
\end{proposition}

\begin{proof}
By Corollary~\ref{cor:single-crossing-betweenness-support-curves},
\[
\tau_2(m)-\tau_1(m)
\]
is strictly decreasing on the common interior slope range. Therefore the set
\[
\{m:\tau_1(m)\le \tau_2(m)\}
\]
is an interval.

Now fix a slope subinterval on which one endpoint support line
\[
m\mapsto y_\xi-m\xi
\]
realizes $u_1(m)$. Then
\[
u_1(m)=y_\xi-m\xi,
\]
so
\[
\frac{d}{dm}\bigl(u_1(m)-\tau_2(m)\bigr)
=
-\xi-(-x_{I_2}(m))
=
x_{I_2}(m)-\xi.
\]
Because every point of $J_2$ lies strictly to the right of every endpoint of
$J_1$, one has
\[
x_{I_2}(m)-\xi>0.
\]
Hence
\[
u_1(m)-\tau_2(m)
\]
is strictly increasing on that piece. Therefore the set
\[
\{m:\tau_2(m)\le u_1(m)\}
\]
is an interval after patching across the finitely many endpoint-switch slopes.
Thus
\[
\Sigma_{2\to 1}
=
\{m:\tau_1(m)\le \tau_2(m)\}\cap \{m:\tau_2(m)\le u_1(m)\}
\]
is an interval.

Similarly, on a slope subinterval where one endpoint support line
\[
m\mapsto y_\eta-m\eta
\]
realizes $u_2(m)$,
\[
\frac{d}{dm}\bigl(u_2(m)-\tau_1(m)\bigr)
=
-\eta-(-x_{I_1}(m))
=
x_{I_1}(m)-\eta<0,
\]
because every point of $J_1$ lies strictly to the left of every endpoint of
$J_2$.

Finally, by Lemma~\ref{lem:fixed-slope-intercept-sets},
\[
\operatorname{int}\mathcal B_{I_1}(m)\cap
\operatorname{int}\mathcal B_{I_2}(m)\neq\varnothing
\]
is equivalent to
\[
\tau_2(m)<u_1(m)
\qquad\text{and}\qquad
\tau_1(m)<u_2(m).
\]
The first inequality holds on an interval because $u_1-\tau_2$ is piecewise
strictly increasing; the second holds on an interval because $u_2-\tau_1$ is
piecewise strictly decreasing. Their intersection is therefore an interval.
\end{proof}

\begin{lemma}[Support-overlap truncation across a generic turning boundary]
\label{lem:support-overlap-truncation-turning}
Let $I_1<I_2$ be adjacent ECT intervals meeting at a generic turning boundary
\[
c(h_b)\neq 0,
\qquad
M(h_b)=0,
\qquad
M'(h_b)\neq 0,
\qquad
N(h_b)\neq 0.
\]
For a fixed finite slope $m$, define
\[
\phi_m(h):=\frac{\Delta_7(h)-m v_5(h)}{c(h)},
\qquad
\mathcal B_I(m):=\phi_m(I).
\]
Assume
\[
b\in \mathcal B_{I_1}(m)\cap \mathcal B_{I_2}(m).
\]
Then there exists a neighborhood $U$ of $h_b$ and compact subintervals
\[
J_1\Subset I_1\setminus U,
\qquad
J_2\Subset I_2\setminus U
\]
such that
\[
b\in \mathcal B_{J_1}(m)\cap \mathcal B_{J_2}(m).
\]
\end{lemma}

\begin{proof}
Differentiate:
\[
\phi_m'(h)=\frac{N(h)-mM(h)}{c(h)^2}.
\]
At $h=h_b$, since $M(h_b)=0$,
\[
\phi_m'(h_b)=\frac{N(h_b)}{c(h_b)^2}\neq 0.
\]
Hence for each fixed finite $m$, after shrinking to a small neighborhood $U$
of $h_b$, the function $\phi_m$ is strictly monotone on each one-sided piece
\[
I_1\cap U,
\qquad
I_2\cap U.
\]
Therefore, for each $i=1,2$, the equation
\[
\phi_m(h)=b
\]
has at most one solution in $I_i\cap U$.

But the assumption
\[
b\in \mathcal B_{I_i}(m)
\]
means exactly that there is at least one solution of $\phi_m(h)=b$ somewhere
on $I_i$. Therefore, after removing the boundary neighborhood $U$, at least
one such solution
\[
h_i^\sharp\in I_i\setminus U
\]
remains on each side. Choose small compact subintervals
\[
J_i\Subset I_i\setminus U
\]
with $h_i^\sharp\in J_i$. Then by continuity,
\[
b=\phi_m(h_i^\sharp)\in \phi_m(J_i)=\mathcal B_{J_i}(m),
\]
for $i=1,2$. This proves the claim.
\end{proof}

\begin{lemma}[Compact adjacent pair-pair implies a directional kernel zero]
\label{lem:pair-pair-gives-directional-kernel-zero}
Let $I_1<I_2$ be adjacent ECT intervals meeting at a generic turning boundary,
and let
\[
L:\ y=mx+b
\]
be a finite-slope line that is $2$-secant to both $\Gamma(I_1)$ and
$\Gamma(I_2)$.
Then, after removing a sufficiently small boundary neighborhood, there exist
compact subintervals
\[
J_1\Subset I_1,
\qquad
J_2\Subset I_2
\]
and points $s\in J_2$, $t\in J_1$ such that one of the two directional
kernels vanishes:
\[
K_{21}(s,t):=\det(F(s),F'(s),F(t))=0,
\]
or
\[
K_{12}(t,s):=\det(F(t),F'(t),F(s))=0,
\]
where
\[
F(h)=(c(h),v_5(h),\Delta_7(h)).
\]
\end{lemma}

\begin{proof}
Because $L$ is $2$-secant to each branch, Lemma~\ref{lem:fixed-slope-intercept-sets}
shows that $m$ lies in the interior tangent-slope range of each interval.
Hence each branch has a unique interior tangent line of slope $m$, with tangent
intercept
\[
\tau_{I_i}(m)=\inf \mathcal B_{I_i}(m).
\]

At a generic turning boundary,
\[
\frac{dy/dh}{dx/dh}=\frac{N(h)}{M(h)}
\]
has one-sided blow-up, so no finite-slope tangency can occur arbitrarily close
to $h_b$. Therefore, after shrinking away a boundary neighborhood and applying
Lemma~\ref{lem:support-overlap-truncation-turning}, we may choose compact
subintervals
\[
J_1\Subset I_1,
\qquad
J_2\Subset I_2
\]
such that:
\begin{enumerate}
\item $b\in \mathcal B_{J_1}(m)\cap \mathcal B_{J_2}(m)$;
\item the unique slope-$m$ tangency point of each branch lies in the
corresponding $J_i$.
\end{enumerate}
Hence each $\mathcal B_{J_i}(m)$ is a compact interval whose lower endpoint
is the tangent intercept
\[
\tau_i:=\tau_{J_i}(m)=\inf \mathcal B_{J_i}(m).
\]

Now set
\[
b_*:=\inf\bigl(\mathcal B_{J_1}(m)\cap \mathcal B_{J_2}(m)\bigr).
\]
Since the intersection of two compact intervals is again a compact interval,
\[
b_*=\max(\tau_1,\tau_2).
\]
So $b_*$ is equal to one of the two tangent intercepts, and it still lies in
the support set of the other branch.

If $b_*=\tau_2$, then the slope-$m$ tangent line to $\Gamma(J_2)$ meets
$\Gamma(J_1)$. Thus there are points $s\in J_2$, $t\in J_1$ such that
\[
\det(F(s),F'(s),F(t))=0,
\]
i.e.
\[
K_{21}(s,t)=0.
\]

If $b_*=\tau_1$, then the slope-$m$ tangent line to $\Gamma(J_1)$ meets
$\Gamma(J_2)$. Thus there are points $t\in J_1$, $s\in J_2$ such that
\[
\det(F(t),F'(t),F(s))=0,
\]
i.e.
\[
K_{12}(t,s)=0.
\]
This proves the lemma.
\end{proof}

\begin{lemma}[Strip exclusion for $K_{21}$ against a compact opposite branch]
\label{lem:strip-exclusion-k21}
Let $I_1<I_2$ be adjacent ECT intervals meeting at a generic turning boundary
\[
M(h_b)=0,\qquad M'(h_b)\neq 0,\qquad N(h_b)\neq 0.
\]
Let $J_1\Subset I_1$ be a compact subinterval. Then there exists a one-sided
neighborhood $U_2\subset I_2$ of the shared boundary point $h_b$ such that
\[
K_{21}(s,t):=\det(F(s),F'(s),F(t))\neq 0
\qquad
(s\in U_2,\ t\in J_1),
\]
where $F(h)=(c(h),v_5(h),\Delta_7(h))$.
\end{lemma}

\begin{proof}
Work in the affine chart
\[
x(h):=\frac{v_5(h)}{c(h)},
\qquad
y(h):=\frac{\Delta_7(h)}{c(h)}.
\]
At a generic turning boundary one has
\[
x'(h_b)=0,\qquad x''(h_b)\neq 0,\qquad y'(h_b)\neq 0.
\]
Hence $y$ is a valid local coordinate and the normalized curve may be written
near $h_b$ as a graph
\[
x=\chi(y)
\]
with
\[
\chi'(y_b)=0,\qquad \chi''(y_b)\neq 0.
\]
Thus $x$ has a strict local extremum at the boundary, while the tangent slope
\[
m(h)=\frac{dy/dh}{dx/dh}=\frac{N(h)}{M(h)}
\]
blows up one-sidedly as $h\to h_b$ from either adjacent interval.

Because $J_1\Subset I_1$ is compact away from $h_b$, there exists
$\delta>0$ such that
\[
|x(t)-x(h_b)|\ge 2\delta
\qquad
(t\in J_1).
\]
After shrinking to a small one-sided neighborhood $U_2\subset I_2$ of $h_b$,
we may also assume
\[
|x(s)-x(h_b)|\le \delta
\qquad
(s\in U_2).
\]
Hence
\[
|x(t)-x(s)|\ge \delta
\qquad
(s\in U_2,\ t\in J_1).
\]

Fix a preliminary compact neighborhood $U_2^0\Subset I_2$ of $h_b$ and set
\[
B:=\sup\{|y(t)-y(s)|:\ t\in J_1,\ s\in U_2^0\}<\infty.
\]
Since $m(s)\to \pm\infty$ as $s\to h_b$ in $I_2$, we may shrink further to
$U_2\subset U_2^0$ so that
\[
|m(s)|>\frac{2B}{\delta}
\qquad
(s\in U_2).
\]

If $K_{21}(s,t)=0$, then the tangent line at $\Gamma(s)$ passes through
$\Gamma(t)$, equivalently
\[
y(t)-y(s)=m(s)\bigl(x(t)-x(s)\bigr).
\]
Taking absolute values gives
\[
B\ge |y(t)-y(s)|=|m(s)|\,|x(t)-x(s)|>\frac{2B}{\delta}\cdot \delta=2B,
\]
a contradiction. Therefore $K_{21}(s,t)\neq 0$ on $U_2\times J_1$.
\end{proof}

\begin{lemma}[Strip exclusion for $K_{12}$ against a compact opposite branch]
\label{lem:strip-exclusion-k12}
Let $I_1<I_2$ be adjacent ECT intervals meeting at a generic turning boundary
\[
M(h_b)=0,\qquad M'(h_b)\neq 0,\qquad N(h_b)\neq 0.
\]
Let $J_2\Subset I_2$ be a compact subinterval. Then there exists a one-sided
neighborhood $U_1\subset I_1$ of the shared boundary point $h_b$ such that
\[
K_{12}(t,s):=\det(F(t),F'(t),F(s))\neq 0
\qquad
(t\in U_1,\ s\in J_2),
\]
where $F(h)=(c(h),v_5(h),\Delta_7(h))$.
\end{lemma}

\begin{proof}
This is the left-right mirror of Lemma~\ref{lem:strip-exclusion-k21}.
\end{proof}

\begin{lemma}[Compact first-contact for $K_{21}$]
\label{lem:compact-first-contact-k21}
Let $J_1\Subset I_1$, $J_2\Subset I_2$ be compact subintervals across a
generic turning boundary. Assume there is a relatively open collar
$V_2\subset J_2$ at the boundary-facing end of $J_2$ such that
\[
K_{21}(s,t)\neq 0
\qquad
(s\in V_2,\ t\in J_1).
\]
Assume also that the zero set of $K_{21}$ in $J_2\times J_1$ is nonempty.
Then there exists a zero $(s_*,t_*)\in J_2\times J_1$ such that:
\begin{enumerate}
\item $s_*$ is boundary-nearest in the $x=v_5/c$ ordering among all
$s\in J_2$ for which $K_{21}(s,t)=0$ for some $t\in J_1$;
\item
\[
\partial_t K_{21}(s_*,t_*)=0.
\]
\end{enumerate}
\end{lemma}

\begin{proof}
Define
\[
A_{21}:=\{\,s\in J_2:\ \exists\, t\in J_1 \text{ with } K_{21}(s,t)=0\,\}.
\]
Since the zero set of $K_{21}$ in $J_2\times J_1$ is compact and nonempty,
the set $A_{21}$ is compact and nonempty. By hypothesis,
\[
A_{21}\cap V_2=\varnothing.
\]
Hence $A_{21}$ has a boundary-nearest element $s_*$ in the $x=v_5/c$
ordering on $J_2$. Choose $t_*\in J_1$ with
\[
K_{21}(s_*,t_*)=0.
\]

Suppose
\[
\partial_t K_{21}(s_*,t_*)\neq 0.
\]
Then by the implicit function theorem there is a $C^1$ local zero curve
\[
t=t(s)
\]
through $(s_*,t_*)$, defined for $s$ near $s_*$, such that
\[
K_{21}(s,t(s))=0.
\]
Because $s_*\notin V_2$ and $V_2$ is an open collar at the boundary-facing
end of $J_2$, there exist nearby points $s\in J_2$ that are boundary-nearer
than $s_*$ in the $x=v_5/c$ ordering. For such $s$, the point
$(s,t(s))$ is also a zero of $K_{21}$, contradicting the minimality of
$s_*$. Therefore
\[
\partial_t K_{21}(s_*,t_*)=0.
\]
\end{proof}

\begin{lemma}[Compact first-contact for $K_{12}$]
\label{lem:compact-first-contact-k12}
Let $J_1\Subset I_1$, $J_2\Subset I_2$ be compact subintervals across a
generic turning boundary. Assume there is a relatively open collar
$V_1\subset J_1$ at the boundary-facing end of $J_1$ such that
\[
K_{12}(t,s)\neq 0
\qquad
(t\in V_1,\ s\in J_2).
\]
Assume also that the zero set of $K_{12}$ in $J_1\times J_2$ is nonempty.
Then there exists a zero $(t_*,s_*)\in J_1\times J_2$ such that:
\begin{enumerate}
\item $t_*$ is boundary-nearest in the $x=v_5/c$ ordering among all
$t\in J_1$ for which $K_{12}(t,s)=0$ for some $s\in J_2$;
\item
\[
\partial_s K_{12}(t_*,s_*)=0.
\]
\end{enumerate}
\end{lemma}

\begin{proof}
Define
\[
A_{12}:=\{\,t\in J_1:\ \exists\, s\in J_2 \text{ with } K_{12}(t,s)=0\,\}.
\]
Since the zero set of $K_{12}$ in $J_1\times J_2$ is compact and nonempty,
the set $A_{12}$ is compact and nonempty. By hypothesis,
\[
A_{12}\cap V_1=\varnothing.
\]
Hence $A_{12}$ has a boundary-nearest element $t_*$ in the $x=v_5/c$
ordering on $J_1$. Choose $s_*\in J_2$ with
\[
K_{12}(t_*,s_*)=0.
\]

Suppose
\[
\partial_s K_{12}(t_*,s_*)\neq 0.
\]
Then by the implicit function theorem there is a $C^1$ local zero curve
\[
s=s(t)
\]
through $(t_*,s_*)$, defined for $t$ near $t_*$, such that
\[
K_{12}(t,s(t))=0.
\]
Because $t_*\notin V_1$ and $V_1$ is an open collar at the boundary-facing
end of $J_1$, there exist nearby points $t\in J_1$ that are boundary-nearer
than $t_*$ in the $x=v_5/c$ ordering. For such $t$, the point
$(t,s(t))$ is also a zero of $K_{12}$, contradicting the minimality of
$t_*$. Therefore
\[
\partial_s K_{12}(t_*,s_*)=0.
\]
\end{proof}

\begin{theorem}[Adjacent pair-pair exclusion across a generic turning boundary]
\label{thm:adjacent-pair-pair-exclusion-turning}
Let $I_1<I_2$ be adjacent ECT intervals meeting at a generic turning boundary
\[
M(h_b)=0,\qquad M'(h_b)\neq 0,\qquad N(h_b)\neq 0.
\]
Then no finite-slope affine line in the chart
\[
x=\frac{v_5}{c},
\qquad
y=\frac{\Delta_7}{c}
\]
is $2$-secant to both $\Gamma(I_1)$ and $\Gamma(I_2)$.
\end{theorem}

\begin{proof}
Assume for contradiction that a finite-slope line
\[
L:\ y=mx+b
\]
is $2$-secant to both adjacent branches. By
Lemma~\ref{lem:pair-pair-gives-directional-kernel-zero}, after removing a
small boundary neighborhood there exist compact subintervals
\[
J_1\Subset I_1,\qquad J_2\Subset I_2
\]
and a zero of one of the two directional kernels.

First suppose there exist $s_0\in J_2$, $t_0\in J_1$ such that
\[
K_{21}(s_0,t_0)=0.
\]
By Lemma~\ref{lem:strip-exclusion-k21}, there exists a one-sided neighborhood
$U_2\subset I_2$ of $h_b$ such that
\[
K_{21}(s,t)\neq 0
\qquad
(s\in U_2,\ t\in J_1).
\]
Choose a compact interval $\widetilde J_2\Subset I_2$ containing $s_0$ and
whose boundary-facing collar $V_2$ lies in $U_2$. Then the zero set of
$K_{21}$ in $\widetilde J_2\times J_1$ is nonempty, and
Lemma~\ref{lem:compact-first-contact-k21} yields a zero
\[
(s_*,t_*)\in \widetilde J_2\times J_1
\]
such that
\[
\partial_t K_{21}(s_*,t_*)=0.
\]
Together with $K_{21}(s_*,t_*)=0$, this means that the tangent line at
$\Gamma(s_*)$ is also tangent at $\Gamma(t_*)$. Since
$s_*\in I_2$ and $t_*\in I_1$ are interior points, this gives a common
interior tangent slope across the adjacent branches, contradicting
Proposition~\ref{prop:adjacent-slope-range-separation}.

Now suppose instead there exist $t_0\in J_1$, $s_0\in J_2$ such that
\[
K_{12}(t_0,s_0)=0.
\]
By Lemma~\ref{lem:strip-exclusion-k12}, there exists a one-sided neighborhood
$U_1\subset I_1$ of $h_b$ such that
\[
K_{12}(t,s)\neq 0
\qquad
(t\in U_1,\ s\in J_2).
\]
Choose a compact interval $\widetilde J_1\Subset I_1$ containing $t_0$ and
whose boundary-facing collar $V_1$ lies in $U_1$. Then the zero set of
$K_{12}$ in $\widetilde J_1\times J_2$ is nonempty, and
Lemma~\ref{lem:compact-first-contact-k12} yields a zero
\[
(t_*,s_*)\in \widetilde J_1\times J_2
\]
such that
\[
\partial_s K_{12}(t_*,s_*)=0.
\]
Together with $K_{12}(t_*,s_*)=0$, this means that the tangent line at
$\Gamma(t_*)$ is also tangent at $\Gamma(s_*)$. Again this gives a common
interior tangent slope across the adjacent branches, contradicting
Proposition~\ref{prop:adjacent-slope-range-separation}.

Both cases are impossible. Therefore no finite-slope line is $2$-secant to
both adjacent branches across a generic turning boundary.
\end{proof}

\begin{remark}[What the unconditional adjacent turning theorem does and does not yet replace]
\label{rem:status-adjacent-pair-pair-turning}
Theorem~\ref{thm:adjacent-pair-pair-exclusion-turning} gives an unconditional
exclusion of finite-slope pair-pair lines across a single generic turning
boundary. It is stronger than the older local directional-kernel reduction in
that it bypasses full symmetric zero-capture.

However, it does not by itself replace
Theorem~\ref{thm:partial-global-no-secant}. To obtain a full global no-four-secant
theorem from it, one would still need a separate propagation statement reducing
nonadjacent pair-pair lines to the adjacent turning-boundary case.
\end{remark}

\begin{remark}[Status of the directional-kernel reduction]
\label{rem:audit-directional-kernel-orientation}
Lemma~\ref{lem:pair-pair-gives-directional-kernel-zero} reduces an adjacent
pair-pair line across a generic turning boundary to a compact interior zero of
one of the two directional tangent-separator kernels.

There are now two closures of that reduction in the draft:
\begin{enumerate}
\item Theorem~\ref{thm:adjacent-pair-pair-exclusion-turning} gives an
unconditional closure for actual adjacent pair-pair lines, via strip exclusion
and compact first-contact.

\item Theorem~\ref{thm:conditional-symmetric-directional-kernel-exclusion}
gives a stronger zero-freeness statement for both directional kernels, but only
conditionally on symmetric zero-capture.
\end{enumerate}

Thus orientation selection between $K_{21}$ and $K_{12}$ is no longer the
main issue. What remains open only for the stronger tangent-separator /
support-ordering route is the symmetric zero-capture problem on the full
adjacent rectangle.
\end{remark}

\begin{lemma}[Adjacent support-ordering]
\label{lem:adjacent-support-ordering}
Let $I_1<I_2$ be adjacent ECT intervals under the hypotheses of
Theorem~\ref{thm:adjacent-tangent-separator-sign-regularity}. Then, for every
slope $m\in\mathbf R$ such that both fixed-slope intercept sets are nonempty,
\[
\mathcal B_{I_1}(m)\cap \mathcal B_{I_2}(m)=\varnothing.
\]
More precisely, one has one of the two strict alternatives
\[
\sup \mathcal B_{I_1}(m)<\inf \mathcal B_{I_2}(m)
\qquad\text{or}\qquad
\sup \mathcal B_{I_2}(m)<\inf \mathcal B_{I_1}(m),
\]
and the same alternative holds for all such $m$ on the adjacent pair.
\end{lemma}

\begin{proof}
Fix $m$ such that both $\mathcal B_{I_1}(m)$ and $\mathcal B_{I_2}(m)$
are nonempty.

First assume $m$ lies in the interior tangent-slope range of $I_2$. Let
$s\in I_2$ be the unique point with tangent slope $m$. Then by
Lemma~\ref{lem:fixed-slope-intercept-sets},
\[
\tau_{I_2}(m)=\inf \mathcal B_{I_2}(m).
\]
By Theorem~\ref{thm:adjacent-tangent-separator-sign-regularity},
\[
K(s,t)\neq 0
\qquad (t\in I_1).
\]
By Lemma~\ref{lem:affine-form-of-K},
\[
K(s,t)=c(s)^2c(t)x'(s)\bigl(\phi_m(t)-\tau_{I_2}(m)\bigr),
\qquad
\phi_m(t):=y(t)-mx(t).
\]
Since $c(s)^2c(t)x'(s)$ has fixed sign on $I_1$ and $K(s,t)$ is continuous
and nowhere zero there, the quantity
\[
\phi_m(t)-\tau_{I_2}(m)
\]
has constant sign on $I_1$. To determine which sign occurs, move $s$ toward
the common boundary point from inside $I_2$. By
Lemma~\ref{lem:endpoint-tangent-exclusion}, the limiting boundary tangent line
does not meet $I_1$, so $I_1$ lies entirely on one side of that tangent
line. By continuity, the same side persists for nearby interior tangents.
Indeed, if the side changed as the tangency point $s$ moved along $I_2$, then
by the connectedness of $I_2$ and continuity of the tangent-line data there
would exist an intermediate point $s_*\in I_2$ for which the tangent line at
$\Gamma(s_*)$ meets $\Gamma(I_1)$. Equivalently, there would exist
$t_*\in I_1$ such that
\[
K(s_*,t_*)=0,
\]
contradicting Theorem~\ref{thm:adjacent-tangent-separator-sign-regularity}.
Hence the same side holds for all $s\in I_2$. Therefore either
\[
\phi_m(t)<\tau_{I_2}(m)\quad\text{for all }t\in I_1,
\]
or
\[
\phi_m(t)>\tau_{I_2}(m)\quad\text{for all }t\in I_1.
\]
In the first case,
\[
\sup\mathcal B_{I_1}(m)\le \tau_{I_2}(m)=\inf\mathcal B_{I_2}(m),
\]
and since equality would force an intersection with the tangent line, the
inequality is strict. The second case is symmetric.

Now suppose $m$ is not in the interior tangent-slope range of $I_2$, but
both intercept sets are nonempty. Then $m$ lies on one side of the tangent
slope range of $I_2$, so $\phi_m$ is strictly monotone on $I_2$, and
$\mathcal B_{I_2}(m)$ is an interval with endpoint attained in the one-sided
limit at the adjacent boundary. Passing to that boundary-support line and using
Lemma~\ref{lem:endpoint-tangent-exclusion} again shows that $I_1$ lies
strictly on one side of that boundary-support line. Hence the same strict
ordering of intercept sets holds in this case as well.

Thus $\mathcal B_{I_1}(m)$ and $\mathcal B_{I_2}(m)$ are disjoint for every
slope $m$ for which both are nonempty, and the side ordering is constant on
the adjacent pair.
\end{proof}

\begin{theorem}[Partial global no-secant theorem on turning-boundary generic blocks]
\label{thm:partial-global-no-secant}
Assume every adjacent pair of ECT intervals is separated by a generic turning
boundary
\[
M(h_b)=0,\qquad M'(h_b)\neq 0,\qquad N(h_b)\neq 0,
\]
and assume moreover that the zero-capture hypothesis required in
Theorem~\ref{thm:adjacent-tangent-separator-sign-regularity} holds for every
adjacent turning-boundary pair.
Then no affine line in the chart
\[
x=\frac{v_5}{c},
\qquad
y=\frac{\Delta_7}{c}
\]
is $2$-secant to two distinct ECT intervals. In particular, there is no
global four-secant line for the normalized planar curve $\Gamma(h)$ on such
a turning-boundary generic block.
\end{theorem}

\begin{proof}
Under the stated hypotheses, every adjacent boundary is of turning type and
Theorem~\ref{thm:adjacent-tangent-separator-sign-regularity} applies to each
adjacent pair.
For such adjacent pairs, Lemma~\ref{lem:adjacent-support-ordering} applies.
Hence adjacent fixed-slope intercept sets are strictly ordered. Iterating this
ordering over the $x=v_5/c$ ordering of ECT intervals shows that no affine
line can be $2$-secant to two distinct ECT intervals. Therefore no global
four-secant line exists on the block.
\end{proof}

\begin{corollary}[Partial exclusion of pair-pair branch configurations]
\label{cor:pair-pair-exclusion}
Under the hypotheses of Theorem~\ref{thm:partial-global-no-secant}, a bad
finite core cannot contain two distinct ECT intervals each contributing a pair
of points on a turning-boundary generic block.
\end{corollary}

\begin{proof}
Each ECT interval contributes at most two points to the zero set of an affine
line in the normalized chart, and by
Theorem~\ref{thm:partial-global-no-secant} no affine line can be $2$-secant
to two different ECT intervals on such a block. Therefore a pair-pair
configuration across two distinct ECT intervals is impossible there.
\end{proof}

\begin{remark}[Strip picture and clean-block compression]
\label{rem:strip-picture-clean-block-compression}
For each ECT interval $I$ in the normalized affine chart
\[
x=\frac{v_5}{c},
\qquad
y=\frac{\Delta_7}{c},
\]
it is convenient to package the fixed-slope $2$-secant data by the open
secant strip
\[
\mathcal S_I
:=
\{(m,b): b\in \operatorname{int}\mathcal B_I(m)\},
\]
where $\mathcal B_I(m)$ is the intercept interval from
Lemma~\ref{lem:fixed-slope-intercept-sets}. Thus
\[
(m,b)\in \mathcal S_I
\]
if and only if the affine line $y=mx+b$ is $2$-secant to $\Gamma(I)$.

Under the hypotheses of
Theorem~\ref{thm:partial-global-no-secant}, these strips are pairwise disjoint
across distinct ECT intervals on the block:
\[
\mathcal S_{I_i}\cap \mathcal S_{I_j}=\varnothing
\qquad(i\neq j).
\]
Equivalently, a pair-pair state can use at most one ECT interval from such a
turning-boundary generic block. In the later finite-core language, every such
block behaves as a single transport chamber rather than as an independent
source of $2+2$ combinatorics.
\end{remark}

\begin{remark}[Canonical quotient-level closure versus the honest order-$7$ lane]
\label{rem:status-direct-septic-witness}
The direct $q^{(7)}$-slice of total $\Delta_7$ is now closed at the
weight-one level by
Proposition~\ref{prop:direct-q7-total-delta7-nonzero}. Together with
Theorem~\ref{thm:quotient-septic-closure}, this closes the canonical
quotient-level order-$7$ package.

However, this does \emph{not} mean that order $7$ has disappeared from the
later endgame. The honest transverse order-$7$ datum is not raw septic
equality of representatives, but the quotient-septic class
\[
[A_7^{\mathfrak f}]\in \mathfrak f/\mathbf C A_5^{\mathfrak f}.
\]
Once two corrected germs agree through cubic and quintic order and differ only
by odd reparametrization plus the already-quotiented lower scalar/value gauge,
the remaining septic mismatch is automatically a multiple of
$A_5^{\mathfrak f}$, so the quotient-septic defect vanishes. What is not yet
proved is that cubic/quintic merger alone forces that odd-reparametrization
comparability.

Accordingly, the honest order-$7$ theorem target is not ``prove raw septic
equality.'' It is a provenance-sensitive package statement, for example:
\begin{itemize}
\item same reduced image germ at coincidence; or
\item collision-functoriality of the corrected block triple on the
collision/cancellation locus.
\end{itemize}
If such a statement holds, the quotient-septic collapse becomes formal
afterward. So order $7$ is canonically closed, but not yet globally exhausted.
\end{remark}

\begin{remark}[Canonical quotient-septic closure is still locally free]
\label{rem:canonical-quotient-septic-closure-still-locally-free}
The canonical quotient-level order-$7$ package is now closed, but that does not
yet make the quotient-septic geometry rigid.

Indeed, once the lower one-pair data $(c,u_5,v_5)$ are fixed, the canonical
septic residual
\[
\Delta_7=u_7v_5-u_5v_7
\]
is still locally free. On a region where $v_5\neq 0$ one may gauge-fix
\[
v_7^{\new}=0,
\qquad
u_7^{\new}=\frac{\Delta_7}{v_5},
\]
while on a region where $u_5\neq 0$ one may gauge-fix
\[
u_7^{\new}=0,
\qquad
v_7^{\new}=-\frac{\Delta_7}{u_5}.
\]
So after the lower one-pair data are fixed, the canonical quotient-septic curve
\[
\Gamma(h)=\left(\frac{v_5(h)}{c(h)},\frac{\Delta_7(h)}{c(h)}\right)
\]
is still locally determined by an arbitrary analytic choice of $\Delta_7$.

Accordingly, no theorem using only arbitrary canonical quotient-septic geometry
of $(c,v_5,\Delta_7)$ can by itself close the later repeated-source /
discriminant branch. Any successful higher-order hybrid must therefore be
provenance-sensitive, not merely quotient-septic.
\end{remark}

\begin{remark}[Status and ceiling of the highest-new direct provenance slice]
\label{rem:wip-highest-new-scalar-direct-slice}
On the region $v_5\neq 0$, set
\[
\mu:=\frac{c}{v_5},
\qquad
\lambda:=\frac{\Delta_7}{v_5},
\qquad
\phi:=\lambda^{[1]}.
\]
By Corollary~\ref{cor:direct-q7-septic-slice},
Corollary~\ref{cor:weight-one-septic-reduction}, and
Theorem~\ref{thm:quotient-septic-closure}, the highest-new quotient-level
septic scalar is already reduced to the direct $q^{(7)}$-driven $K_1$-slice:
\[
\Delta_7^{[1]}=v_5\,\phi,
\qquad
\phi
=
\kappa\Bigl(
(x_2-x_1)\nu_6^{(q^{(7)})}
+
v(\gamma_5-\alpha_5)^{(q^{(7)})}
\Bigr).
\]

On the subregion where $M\neq 0$, the direct highest-new dual variables are
\[
\beta^{[1]}=\phi_\mu,
\qquad
m^{[1]}=\phi-\mu\phi_\mu,
\]
and by Corollary~\ref{cor:mu-lambda-curvature-identities},
\[
\phi_{\mu\mu}=\frac{v_5^3}{M^3}W^{[1]}.
\]
Thus the direct highest-new quotient-septic geometry is exactly the ordinary
tangent-line geometry of the scalar graph
\[
y=\phi(\mu).
\]

In particular, any bitangent line or four-secant line for the graph
$y=\phi(\mu)$, and any nonadjacent slope re-entry for
\[
\beta^{[1]}=\phi_\mu,
\]
forces at least two zeros of
\[
\phi_{\mu\mu},
\]
equivalently of
\[
W^{[1]}.
\]
So the global direct-slice propagation problem may be viewed as a no-recurrence
problem for the scalar graph $\phi$.

However, the explicit direct provenance formula only shows that
\[
\phi=A\,q^{(7)}
\qquad
(A:=\kappa\,\eta_2),
\]
with $A$ generically nonvanishing on a generic one-pair block. By
Proposition~\ref{prop:full-block-universality-direct-provenance-slice} below,
this still leaves $\phi$ fully universal on each such block. So the direct
graph picture is now a precise dictionary, not by itself a live rigidity
theorem.
\end{remark}

\begin{remark}[Status of the $W=0,\ M\neq 0$ side]
\label{rem:audit-w0-affine-inflection}
On the region $c\neq 0,\ M\neq 0$, Proposition~\ref{prop:no-W-flat-interval}
excludes open-interval flatness $W\equiv 0$. At an isolated boundary
\[
W(h_b)=0,\qquad M(h_b)\neq 0,
\]
the normalized planar curve in the chart
\[
x=\frac{v_5}{c},
\qquad
y=\frac{\Delta_7}{c}
\]
is an ordinary odd-order affine inflection whenever $W$ changes sign, so
adjacent tangent-separator sign regularity fails there in general.

However, Theorem~\ref{thm:finite-order-w-zero-local-sterility} shows that every
finite-order isolated boundary of type $W=0,\ M\neq 0$ is locally sterile for
pair-pair birth. Thus the remaining issue on this side is not local creation of
a bad $2+2$ line, but global continuation / recurrence of determinant-zero or
tangent-slope data across multiple intervals.

The exact identities for $W$ and $W'$ continue to show that this residual
obstruction is still encoded through the canonical order-$7$ package
$\Delta_7$, not by a genuinely new primitive order-$9$ invariant.
\end{remark}

\begin{theorem}[Finite-order $W=0$ boundaries are locally sterile for pair-pair birth]
\label{thm:finite-order-w-zero-local-sterility}
Let $h_b$ satisfy
\[
c(h_b)\neq 0,
\qquad
M(h_b)\neq 0,
\qquad
W(h_b)=0,
\]
and assume that $W$ has finite vanishing order $r\ge 1$ at $h_b$:
\[
W(h)=(h-h_b)^r\,\widetilde W(h),
\qquad
\widetilde W(h_b)\neq 0.
\]
Define
\[
F(h):=(c(h),v_5(h),\Delta_7(h)),
\qquad
D(s,t,u):=\det(F(s),F(t),F(u)).
\]
Then, after shrinking to a sufficiently small neighborhood $U$ of $h_b$,
the following holds.

If $s<t<h_b$ lie in $U$, then:
\begin{enumerate}
\item if $r$ is even, the function $u\mapsto D(s,t,u)$ has no zero on the
opposite side $h_b<u\in U$;
\item if $r$ is odd, the function $u\mapsto D(s,t,u)$ has at most one zero on
the opposite side $h_b<u\in U$.
\end{enumerate}
The same conclusion holds with left and right exchanged. In particular, no
finite-order boundary of type $W=0,\ M\neq 0$ can support a local pair-pair
birth.
\end{theorem}

\begin{proof}
Work in the affine chart
\[
x(h):=\frac{v_5(h)}{c(h)},
\qquad
y(h):=\frac{\Delta_7(h)}{c(h)}.
\]
Since
\[
x'(h)=\frac{M(h)}{c(h)^2},
\]
the hypothesis $M(h_b)\neq 0$ implies that $x$ is a valid local coordinate
near $h_b$. After translating so that $x(h_b)=0$ and subtracting the tangent
line, we may write the normalized curve as
\[
y=f(x),
\qquad
f(0)=f'(0)=0.
\]

By Proposition~\ref{prop:projective-wronskian-bridge},
\[
\frac{d^2y}{dx^2}=\frac{c(h)^3W(h)}{M(h)^3}.
\]
Because $x'(h_b)\neq 0$, the finite vanishing order of $W$ in $h$ is the same
as its finite vanishing order in the local coordinate $x$. Therefore
\[
f''(x)=\lambda x^r+O(x^{r+1})
\qquad(\lambda\neq 0).
\]

Fix two same-side points
\[
x_1<x_2<0
\]
sufficiently close to $0$, and define
\[
\widehat D(x_1,x_2,u):=
\det\!\begin{pmatrix}
1&x_1&f(x_1)\\
1&x_2&f(x_2)\\
1&u&f(u)
\end{pmatrix}.
\]
Then
\[
\widehat D(x_1,x_2,u)
=
(x_2-x_1)(u-x_1)(u-x_2)\,H_{x_1,x_2}(u),
\]
where
\[
H_{x_1,x_2}(u):=f[x_1,x_2,u]
\]
is the second divided difference of $f$.

If $r$ is even, then $f''$ has fixed sign near $0$. By the Hermite--Genocchi
formula for second divided differences, $H_{x_1,x_2}(u)$ has that same fixed
sign for $u>0$ sufficiently small. Hence $\widehat D(x_1,x_2,u)\neq 0$ on the
opposite side.

If $r$ is odd, then
\[
f'''(x)=\lambda r\,x^{r-1}+O(x^r)
\]
has fixed nonzero sign near $0$, because $r-1$ is even. Differentiating with
respect to the repeated node gives
\[
\partial_u H_{x_1,x_2}(u)=f[x_1,x_2,u,u].
\]
By the Hermite--Genocchi formula for repeated-node divided differences,
$f[x_1,x_2,u,u]$ has the same fixed nonzero sign for $u>0$ sufficiently small.
Thus $u\mapsto H_{x_1,x_2}(u)$ is strictly monotone on the opposite side, and
therefore has at most one zero there.

Finally,
\[
D(s,t,u)=c(s)c(t)c(u)\,\widehat D(x(s),x(t),x(u)),
\]
and $c$ is nonzero on a small neighborhood of $h_b$, so the zero set of $D$
is exactly the zero set of $\widehat D$. The theorem follows.
\end{proof}

\begin{remark}[Consequence for the $W=0,\ M\neq 0$ side]
\label{rem:consequence-finite-order-w-zero-local-sterility}
Theorem~\ref{thm:finite-order-w-zero-local-sterility} shows that finite-order
$W=0,\ M\neq 0$ boundaries are not local birthplaces of pair-pair
configurations. Combined with
Lemma~\ref{lem:endpoint-gap-monotonicity-ordered-target} and
Corollary~\ref{cor:turning-boundary-inertness-endpoint-contacts}, this yields a
sharper reduction of the one-pair propagation problem.

Namely, after removing local births at generic turning boundaries and at
finite-order isolated $W=0,\ M\neq 0$ boundaries, any surviving nonadjacent
pair-pair obstruction must come from a \emph{remote} interior zero of one of
the endpoint-gap functions
\[
G_{\xi,J}(h)=y_\xi-y(h)-m(h)\bigl(\xi-x(h)\bigr),
\]
where $(\xi,y_\xi)$ is an endpoint support point and $J$ is an ordered target
branch.

Since each $G_{\xi,J}$ is strictly monotone on $J$, each endpoint/branch pair
can contribute at most one such tangency. Thus the remaining obstruction on the
$W=0$ side is no longer local birth at an affine inflection seam, but a finite
remote-incidence problem.
\end{remark}

\begin{lemma}[Endpoint-gap monotonicity on an ordered target branch]
\label{lem:endpoint-gap-monotonicity-ordered-target}
Let $J$ be an ECT interval in the affine chart
\[
x=\frac{v_5}{c},
\qquad
y=\frac{\Delta_7}{c},
\]
and let $(\xi,y_\xi)$ be a fixed support point with
\[
\xi< x(h)\qquad (h\in J).
\]
Define
\[
G_{\xi,J}(h):=y_\xi-y(h)-m(h)\bigl(\xi-x(h)\bigr),
\qquad
m(h)=\frac{N(h)}{M(h)}.
\]
Then
\[
G_{\xi,J}'(h)
=
\frac{c(h)W(h)}{M(h)^2}\,\bigl(x(h)-\xi\bigr).
\]
In particular, $G_{\xi,J}$ is strictly monotone on $J$.

Equivalently, a fixed support point strictly to the left of $J$ lies on at most
one tangent line to $\Gamma(J)$. The symmetric statement holds for a fixed
support point strictly to the right of $J$.
\end{lemma}

\begin{proof}
Using
\[
y'(h)=m(h)x'(h),
\]
differentiate:
\[
G_{\xi,J}'(h)
=
-y'(h)-m'(h)(\xi-x(h))+m(h)x'(h)
=
m'(h)\bigl(x(h)-\xi\bigr).
\]
By Proposition~\ref{prop:closed-minor-system-dual-curve},
\[
m'(h)=\frac{c(h)W(h)}{M(h)^2},
\]
so
\[
G_{\xi,J}'(h)
=
\frac{c(h)W(h)}{M(h)^2}\,\bigl(x(h)-\xi\bigr).
\]
On an ECT interval, $c$, $W$, and $M$ have fixed nonzero sign, and by
assumption $x(h)-\xi>0$ on $J$, so $G_{\xi,J}'$ has fixed nonzero sign.
Therefore $G_{\xi,J}$ is strictly monotone.
\end{proof}

\begin{corollary}[Turning-boundary inertness for endpoint contacts]
\label{cor:turning-boundary-inertness-endpoint-contacts}
Let $J$ be an ECT interval ending at a generic turning boundary $h_b$ with
\[
M(h_b)=0,
\qquad
M'(h_b)\neq 0,
\qquad
N(h_b)\neq 0,
\qquad
c(h_b)\neq 0.
\]
Assume $v_5(h_b)\neq 0$, and let
\[
x_b:=\frac{v_5(h_b)}{c(h_b)}.
\]
If $(\xi,y_\xi)$ is a fixed support point with $\xi<x_b$, then
\[
G_{\xi,J}(h)\to \pm\infty
\qquad\text{as } h\to h_b \text{ from within } J.
\]
In particular, endpoint-support contact cannot be created or destroyed at a
generic turning-boundary end of an ordered branch.
\end{corollary}

\begin{proof}
By Proposition~\ref{prop:closed-minor-system-dual-curve},
\[
m(h)=\frac{N(h)}{M(h)}\to \pm\infty
\qquad (h\to h_b),
\]
while
\[
x(h)\to x_b,
\qquad
y(h)\to y_b:=\frac{\Delta_7(h_b)}{c(h_b)}.
\]
Hence
\[
G_{\xi,J}(h)
=
\bigl(y_\xi-y(h)\bigr)+m(h)\bigl(x(h)-\xi\bigr).
\]
Since $x(h)-\xi\to x_b-\xi>0$, the second term dominates and has fixed sign
for $h$ sufficiently close to $h_b$. Therefore
\[
G_{\xi,J}(h)\to \pm\infty.
\]
\end{proof}

\begin{remark}[Endpoint purification for bounded strip-overlap components]
\label{rem:endpoint-purification-bounded-strip-overlap}
Let
\[
\mathcal P
:=
\bigcup_{i<j-1}(\mathcal S_{I_i}\cap \mathcal S_{I_j})
\]
denote the nonadjacent pair-pair strip-overlap set on a compact patch, where
the union runs over the ECT intervals in the normalized affine chart.

By Theorem~\ref{thm:adjacent-pair-pair-exclusion-turning} and
Theorem~\ref{thm:finite-order-w-zero-local-sterility}, a bounded connected
component
\[
U\subset \mathcal P
\]
cannot be created or destroyed at a generic turning boundary, and cannot have a
first entrance or last exit at a finite-order boundary of type
\[
W=0,\qquad M\neq 0.
\]

Moreover, by
Lemma~\ref{lem:endpoint-gap-monotonicity-ordered-target}, each fixed endpoint
support point can meet a given ordered target branch at most once. Thus any
surviving bounded strip-overlap component is controlled by a finite collection
of isolated remote endpoint/branch tangencies, rather than by diffuse local
birth along a seam.

So the remaining nonadjacent pair-pair problem is best viewed simply as a
finite remote-incidence problem.
\end{remark}

\begin{remark}[Audit: residual $W=0$ pathology]
\label{rem:audit-residual-w0-pathology}
By Proposition~\ref{prop:no-W-flat-interval}, the $W=0,\ M\neq 0$ side is no
longer obstructed by open-interval flatness $W\equiv 0$ on a component with
$c\neq 0$ and $M\neq 0$. By
Theorem~\ref{thm:finite-order-w-zero-local-sterility}, every finite-order
isolated $W=0$ boundary is locally sterile for pair-pair birth. Thus any
remaining $W=0$ issue would have to come either from non-analytic /
infinite-order flatness or from a genuinely global continuation phenomenon,
not from a generic local inflection birth mechanism.
\end{remark}

\subsection{From the one-pair residual germ to two-pair and minimal-core rigidity}

\begin{remark}[Strengthened organizing principle]
\label{rem:strengthened-organizing-principle}
The original version of this subsection was organized around the cubic leading
line and the first residual trace witness. The later working analysis suggests
a stronger one-pair package:
\[
c(h),\qquad A_5^{\mathfrak f}(h),\qquad \Delta_7(h),
\]
or equivalently the normalized affine curve
\[
h\longmapsto \left(\frac{v_5(h)}{c(h)},\frac{\Delta_7(h)}{c(h)}\right).
\]
The intended two-pair and minimal-core rigidity statements should ultimately be
reformulated in terms of this strengthened one-pair invariant geometry.
\end{remark}

\begin{remark}[Subsection guide]
\label{rem:wip-multi-pair-subsection}
This subsection is best read as a rigidity program built on the strengthened
one-pair local package
\[
c(h),\qquad A_5^{\mathfrak f}(h),\qquad \Delta_7(h),
\]
or equivalently the normalized planar curve
\[
h\longmapsto \left(\frac{v_5(h)}{c(h)},\frac{\Delta_7(h)}{c(h)}\right).
\]
The projective coincidence conditions remain the organizing principle for the
two-pair and minimal-core analysis. The finite-order discussion below now
isolates the exact source-level criterion that would force the two-point
quadratic interaction factorization through order $7$, but the current attack
also exhibits two backup lanes: a weaker quotient-diagonal exact exclusion route
and, on good overlap patches, a mixed $\lambda$-shear / projective-shadow
route. Beyond that, the remaining steps are the verification of one of these
two-point lanes, the genuine multi-pair interaction bookkeeping, and the
extraction into the finite-core endgame.
\end{remark}

\begin{remark}[Interpretation of residual notation in the multi-pair architecture]
\label{rem:residual-notation-multipair}
In the remainder of this section, the symbols
\[
F_h(z),\qquad \ell_h,\qquad \overline v_1(h),\qquad
\Qabs/\ell_h
\]
are used schematically to describe the intended cubic-leading / first-residual
architecture of multi-pair and minimal-core configurations.

The fully explicit rigorous first-residual object established earlier is
\[
\mathcal R_h^{(5)}
=
\mathcal J_5/
\bigl(\mathbf C\,A_3(h)+V_{\mathrm{amb}}^{(5)}(h)\bigr).
\]
Accordingly, whenever a later statement invokes the first residual class in a
theorem-strength way, it should be interpreted through $\mathcal R_h^{(5)}$.
The present two-pair subsection now isolates the exact finite-order source
criterion through order $7$ together with two weaker backup routes. What remains
architectural is the verification of one of these two-point routes for the
actual corrected two-atom package, together with its extension to genuine
multi-pair and minimal-bad-core interaction terms.
\end{remark}

Let $F_h(z)$ denote the corrected quotient odd germ of the normalized one-pair
family. By Proposition~\ref{prop:cubic-leading-order}, on the cubic-generic
region we have
\[
F_h(z)=v_0(h)z^3+v_1(h)z^5+\cdots,
\qquad
v_0(h)\neq 0,
\]
with cubic leading line
\[
\ell_h:=\mathbf C\,v_0(h)\subset \Qabs.
\]
Under the additional hypotheses of
Proposition~\ref{prop:first-residual-class-generic}, there is an open dense
subset $\mathcal U_{\mathrm{res}}\subset \mathcal U_{\mathrm{cub}}$ on which
the explicit first residual class
\[
\overline v_1(h)\in\mathcal R_h^{(5)}
\]
is nonzero.

In the strengthened fixed-sector refinement, the one-pair local package is
recorded by
\[
c(h),\qquad A_5^{\mathfrak f}(h)=u_5(h)I+v_5(h)S,\qquad
\Delta_7(h),
\]
and hence by the normalized planar curve
\[
\Gamma(h):=\left(\frac{v_5(h)}{c(h)},\frac{\Delta_7(h)}{c(h)}\right)
=
\bigl(\rho_1(h),\rho_2(h)\bigr).
\]
By Proposition~\ref{prop:curvature-bridge-normalized-curve}, the affine
curvature numerator of $\Gamma$ is
\[
K_\Gamma(h)
=
\rho_1'(h)\rho_2''(h)-\rho_2'(h)\rho_1''(h)
=
\frac{M(h)N'(h)-N(h)M'(h)}{c(h)^4},
\]
where
\[
M(h)=c(h)v_5'(h)-c'(h)v_5(h),
\qquad
N(h)=c(h)\Delta_7'(h)-c'(h)\Delta_7(h).
\]
Thus the strengthened one-pair geometry is governed by the curve $\Gamma(h)$
and its secant/tangent behavior on the region where
\[
c(h)\neq 0,
\qquad
M(h)\neq 0,
\qquad
M(h)N'(h)-N(h)M'(h)\neq 0.
\]

\begin{definition}[Amplitude-invariant strengthened datum]
\label{def:amplitude-invariant-strengthened-datum}
On the region where
\[
c(h)\neq 0,
\]
define
\[
\widehat\Psi(h):=
\left(
\frac{u_5(h)}{c(h)},
\frac{v_5(h)}{c(h)},
\frac{\Delta_7(h)}{c(h)^2}
\right).
\]
\end{definition}

\begin{lemma}[Scaling law for the strengthened package]
\label{lem:scaling-law-strengthened-package}
If the normalized one-pair odd germ is multiplied by a scalar $\lambda$, then
\[
c\longmapsto \lambda c,
\qquad
A_5^{\mathfrak f}\longmapsto \lambda A_5^{\mathfrak f},
\qquad
A_7^{\mathfrak f}\longmapsto \lambda A_7^{\mathfrak f},
\]
and hence
\[
\Delta_7\longmapsto \lambda^2 \Delta_7.
\]
Consequently the datum
\[
\widehat\Psi(h)=
\left(
\frac{u_5(h)}{c(h)},
\frac{v_5(h)}{c(h)},
\frac{\Delta_7(h)}{c(h)^2}
\right)
\]
is invariant under scalar rescaling of the full one-pair germ.
\end{lemma}

\begin{proof}
The cubic, quintic, and septic coefficients of the odd germ all scale linearly
under multiplication of the germ by $\lambda$. Since
\[
A_5^{\mathfrak f}=u_5I+v_5S,
\qquad
A_7^{\mathfrak f}=u_7I+v_7S,
\]
it follows that
\[
u_5\mapsto \lambda u_5,\qquad v_5\mapsto \lambda v_5,
\qquad
u_7\mapsto \lambda u_7,\qquad v_7\mapsto \lambda v_7.
\]
Therefore
\[
\Delta_7=u_7v_5-u_5v_7
\]
scales as
\[
\Delta_7\mapsto \lambda^2\Delta_7,
\]
which gives the claim.
\end{proof}

\begin{remark}[Two normalizations: geometric versus projective]
\label{rem:two-normalizations-geometric-vs-projective}
The affine curve
\[
\Gamma(h)=\left(\frac{v_5(h)}{c(h)},\frac{\Delta_7(h)}{c(h)}\right)
\]
is the correct object for the one-pair differential geometry, because its
derivatives are governed by the projective Wronskians $M$, $N$, and $W$.

For the multi-pair coincidence problem, however, the natural scalar-invariant
datum is $\widehat\Psi(h)$, not $\Gamma(h)$. The reason is exactly the
scaling law in Lemma~\ref{lem:scaling-law-strengthened-package}: under scalar
rescaling of the full one-pair germ, $\Delta_7/c$ is not invariant, whereas
$\Delta_7/c^2$ is.
\end{remark}

In this formulation, the geometric secant/tangent obstruction is well controlled
on ECT intervals and across generic turning-point boundaries. The remaining
global issue is no longer a purely local $W=0,\ M\neq 0$ obstruction, but the
global propagation / recurrence problem for the determinant-zero and tangent
configurations that survive the local sterility results proved later in the
paper.

\begin{definition}[Two-pair decomposition]
\label{def:two-pair-decomposition}
For a normalized two-pair core with amplitudes $a_1,a_2$ and offsets
$h_1,h_2$, write its quotient odd germ schematically as
\[
F_{(a_1,h_1;a_2,h_2)}(z)
=
a_1F_{h_1}(z)+a_2F_{h_2}(z)+\mathcal I_{12}(z),
\]
where $\mathcal I_{12}(z)$ denotes the pair-interaction germ.

In the explicit finite-order corrected source-level model through order $7$
constructed below, the cubic scalar, quintic fixed-sector, and septic
fixed-sector projections of $\mathcal I_{12}$ are quadratic in the amplitudes
and divisible by $(h_1-h_2)^2$.
\end{definition}

\begin{proposition}[Exact strengthened two-pair coincidence through order $7$]
\label{prop:exact-strengthened-two-pair-coincidence}
Let $(a_1,h_1;a_2,h_2)$ be a normalized two-pair core with
\[
F_{(a_1,h_1;a_2,h_2)}(z)\equiv 0.
\]
Assume moreover that, after passage to the explicit corrected finite-order model
through order $7$, the interaction contribution vanishes at the three
relevant levels, namely:
\[
a_1c(h_1)+a_2c(h_2)=0,
\]
\[
a_1A_5^{\mathfrak f}(h_1)+a_2A_5^{\mathfrak f}(h_2)=0,
\]
and, after the second identity identifies the common quotient line
\[
\mathbf C A_5^{\mathfrak f}(h_1)=\mathbf C A_5^{\mathfrak f}(h_2),
\]
the second residual classes satisfy
\[
a_1\,[A_7^{\mathfrak f}](h_1)+a_2\,[A_7^{\mathfrak f}](h_2)=0
\qquad\text{in}\qquad
\mathfrak f/\mathbf C A_5^{\mathfrak f}(h_1).
\]
Then
\[
\widehat\Psi(h_1)=\widehat\Psi(h_2).
\]
Equivalently,
\[
\frac{u_5(h_1)}{c(h_1)}=\frac{u_5(h_2)}{c(h_2)},
\qquad
\frac{v_5(h_1)}{c(h_1)}=\frac{v_5(h_2)}{c(h_2)},
\qquad
\frac{\Delta_7(h_1)}{c(h_1)^2}=\frac{\Delta_7(h_2)}{c(h_2)^2}.
\]
\end{proposition}

\begin{proof}
Set
\[
\lambda:=-\frac{a_2}{a_1}.
\]
The cubic identity gives
\[
c(h_1)=\lambda\,c(h_2).
\]
The quintic fixed-sector identity gives
\[
A_5^{\mathfrak f}(h_1)=\lambda\,A_5^{\mathfrak f}(h_2).
\]
Writing
\[
A_5^{\mathfrak f}(h_i)=u_5(h_i)I+v_5(h_i)S,
\]
we obtain
\[
u_5(h_1)=\lambda\,u_5(h_2),
\qquad
v_5(h_1)=\lambda\,v_5(h_2).
\]
Hence
\[
\frac{u_5(h_1)}{c(h_1)}=\frac{u_5(h_2)}{c(h_2)},
\qquad
\frac{v_5(h_1)}{c(h_1)}=\frac{v_5(h_2)}{c(h_2)}.
\]

Now the second identity implies
\[
\mathbf C A_5^{\mathfrak f}(h_1)=\mathbf C A_5^{\mathfrak f}(h_2),
\]
so the residual quotient spaces are canonically identified. The order-$7$
residual identity therefore says
\[
[A_7^{\mathfrak f}](h_1)=\lambda\,[A_7^{\mathfrak f}](h_2)
\]
in the common quotient. Equivalently, there exists $\mu\in\mathbf C$ such
that
\[
A_7^{\mathfrak f}(h_1)=\lambda\,A_7^{\mathfrak f}(h_2)+\mu\,A_5^{\mathfrak f}(h_2).
\]
Writing
\[
A_7^{\mathfrak f}(h_i)=u_7(h_i)I+v_7(h_i)S,
\]
and using
\[
\Delta_7(h)=u_7(h)v_5(h)-u_5(h)v_7(h),
\]
we obtain
\[
\Delta_7(h_1)
=
\det\!\bigl(A_7^{\mathfrak f}(h_1),A_5^{\mathfrak f}(h_1)\bigr)
=
\det\!\bigl(\lambda A_7^{\mathfrak f}(h_2)+\mu A_5^{\mathfrak f}(h_2),
\lambda A_5^{\mathfrak f}(h_2)\bigr)
=
\lambda^2\Delta_7(h_2).
\]
Since also $c(h_1)=\lambda c(h_2)$, it follows that
\[
\frac{\Delta_7(h_1)}{c(h_1)^2}
=
\frac{\Delta_7(h_2)}{c(h_2)^2}.
\]
Thus $\widehat\Psi(h_1)=\widehat\Psi(h_2)$.
\end{proof}

\begin{theorem}[Quantitative defect-thickened strengthened coincidence]
\label{thm:quantitative-defect-thickened-strengthened-coincidence}
Let $(a_1,h_1;a_2,h_2)$ be a normalized bad two-pair core, and write
\[
c_i:=c(h_i),
\qquad
A_{5,i}^{\mathfrak f}:=A_5^{\mathfrak f}(h_i)=u_{5,i}I+v_{5,i}S,
\qquad
\Delta_i:=\Delta_7(h_i).
\]
Assume
\[
a_1\neq 0,
\qquad
c_1c_2\neq 0.
\]
Suppose the cubic, quintic, and septic quotient defects are written as
\[
a_1c_1+a_2c_2+E_{12}^{(3)}=0,
\]
\[
a_1A_{5,1}^{\mathfrak f}+a_2A_{5,2}^{\mathfrak f}+E_{12}^{(5)}=0,
\]
and
\[
a_1[A_{7,1}^{\mathfrak f}]+a_2[A_{7,2}^{\mathfrak f}]
+\overline E_{12}^{(7;1)}=0
\qquad\text{in}\qquad
\mathfrak f/\mathbf C A_{5,1}^{\mathfrak f}.
\]
Set
\[
\lambda:=-\frac{a_2}{a_1},
\qquad
e:=\frac{E_{12}^{(3)}}{a_1},
\qquad
E:=\frac{E_{12}^{(5)}}{a_1}=e_u I+e_v S.
\]
Choose any representative $R\in\mathfrak f$ of the quotient class
\[
-\frac{1}{a_1}\,\overline E_{12}^{(7;1)},
\]
so that
\[
A_{7,1}^{\mathfrak f}
=
\lambda A_{7,2}^{\mathfrak f}
+
R
+
\eta A_{5,1}^{\mathfrak f}
\]
for some $\eta\in\mathbf C$. Then the strengthened invariant
\[
\widehat\Psi(h)=
\left(
\frac{u_5(h)}{c(h)},
\frac{v_5(h)}{c(h)},
\frac{\Delta_7(h)}{c(h)^2}
\right)
\]
satisfies the exact identities
\[
\frac{u_{5,1}}{c_1}-\frac{u_{5,2}}{c_2}
=
\frac{u_{5,2}e-c_2e_u}{c_1c_2},
\]
\[
\frac{v_{5,1}}{c_1}-\frac{v_{5,2}}{c_2}
=
\frac{v_{5,2}e-c_2e_v}{c_1c_2},
\]
and
\[
\frac{\Delta_1}{c_1^2}-\frac{\Delta_2}{c_2^2}
=
\frac{
-\lambda c_2^2\det(A_{7,2}^{\mathfrak f},E)
+
c_2^2\det(R,A_{5,1}^{\mathfrak f})
+
\Delta_2(2\lambda c_2 e-e^2)
}{c_1^2c_2^2}.
\]
\end{theorem}

\begin{proof}
The cubic defect identity gives
\[
c_1=\lambda c_2-e.
\]
The quintic defect identity gives
\[
A_{5,1}^{\mathfrak f}=\lambda A_{5,2}^{\mathfrak f}-E,
\]
hence
\[
u_{5,1}=\lambda u_{5,2}-e_u,
\qquad
v_{5,1}=\lambda v_{5,2}-e_v.
\]
Therefore
\[
\frac{u_{5,1}}{c_1}-\frac{u_{5,2}}{c_2}
=
\frac{u_{5,1}c_2-u_{5,2}c_1}{c_1c_2}
=
\frac{(\lambda u_{5,2}-e_u)c_2-u_{5,2}(\lambda c_2-e)}{c_1c_2}
=
\frac{u_{5,2}e-c_2e_u}{c_1c_2},
\]
and similarly
\[
\frac{v_{5,1}}{c_1}-\frac{v_{5,2}}{c_2}
=
\frac{v_{5,2}e-c_2e_v}{c_1c_2}.
\]

For the third coordinate, by the choice of $R$ we have
\[
A_{7,1}^{\mathfrak f}
=
\lambda A_{7,2}^{\mathfrak f}
+
R
+
\eta A_{5,1}^{\mathfrak f}
\]
for some $\eta\in\mathbf C$. Using
\[
\Delta_i=\det(A_{7,i}^{\mathfrak f},A_{5,i}^{\mathfrak f}),
\]
we obtain
\[
\Delta_1
=
\det(A_{7,1}^{\mathfrak f},A_{5,1}^{\mathfrak f})
=
\lambda\det(A_{7,2}^{\mathfrak f},A_{5,1}^{\mathfrak f})
+
\det(R,A_{5,1}^{\mathfrak f}),
\]
since the $\eta A_{5,1}^{\mathfrak f}$ term drops out. Substituting
\[
A_{5,1}^{\mathfrak f}=\lambda A_{5,2}^{\mathfrak f}-E
\]
gives
\[
\det(A_{7,2}^{\mathfrak f},A_{5,1}^{\mathfrak f})
=
\lambda\Delta_2-\det(A_{7,2}^{\mathfrak f},E),
\]
hence
\[
\Delta_1-\lambda^2\Delta_2
=
-\lambda\det(A_{7,2}^{\mathfrak f},E)
+
\det(R,A_{5,1}^{\mathfrak f}).
\]

Also,
\[
c_1=\lambda c_2-e
\]
implies
\[
\lambda^2c_2^2-c_1^2=2\lambda c_2 e-e^2.
\]
Therefore
\[
\frac{\Delta_1}{c_1^2}-\frac{\Delta_2}{c_2^2}
=
\frac{c_2^2\Delta_1-c_1^2\Delta_2}{c_1^2c_2^2}
=
\frac{
c_2^2(\Delta_1-\lambda^2\Delta_2)+\Delta_2(\lambda^2c_2^2-c_1^2)
}{c_1^2c_2^2},
\]
which is exactly the claimed formula.

Finally, the expression
\[
\det(R,A_{5,1}^{\mathfrak f})
\]
is independent of the chosen representative of the quotient defect, since
replacing $R$ by $R+\xi A_{5,1}^{\mathfrak f}$ does not change the
determinant.
\end{proof}

\begin{corollary}[Bad two-pair cores are defect-scale close on good $M\neq 0$ patches]
\label{cor:bad-two-pair-cores-defect-scale-close}
Let $J\Subset\{c\neq 0,\ M\neq 0\}$ be a compact interval. Assume
\[
|a_1|\ge a_*>0
\]
for the normalized bad two-pair core under consideration. Then there exists a
constant $C_J>0$ such that whenever $h_1,h_2\in J$,
\[
|h_1-h_2|
\le
C_J\Bigl(|E_{12}^{(3)}|+\|E_{12}^{(5)}\|\Bigr).
\]
In particular, on good $M\neq 0$ patches, bad two-pair cores force the two
parameters to be quantitatively close once the cubic/quintic interaction
defects are small.
\end{corollary}

\begin{proof}
The second coordinate of $\widehat\Psi$ is
\[
x(h):=\frac{v_5(h)}{c(h)},
\qquad
x'(h)=\frac{M(h)}{c(h)^2}.
\]
Since $J\Subset\{c\neq 0,\ M\neq 0\}$, there exists $\eta_J>0$ such that
\[
|x(h_1)-x(h_2)|\ge \eta_J\,|h_1-h_2|
\qquad
(h_1,h_2\in J).
\]
On the other hand, Theorem~\ref{thm:quantitative-defect-thickened-strengthened-coincidence}
gives the exact identity
\[
x(h_1)-x(h_2)
=
\frac{v_{5,2}e-c_2e_v}{c_1c_2},
\]
with
\[
e=\frac{E_{12}^{(3)}}{a_1},
\qquad
e_v=\text{$S$-coefficient of } \frac{E_{12}^{(5)}}{a_1}.
\]
Since $J$ is compact and $a_1$ is bounded away from $0$, the right-hand
side is bounded by
\[
C'_J\Bigl(|E_{12}^{(3)}|+\|E_{12}^{(5)}\|\Bigr)
\]
for some constant $C'_J$. Combining the two estimates yields the claim.
\end{proof}

\begin{proposition}[No sublinear two-pair defects on good patches]
\label{prop:no-sublinear-two-pair-defects-good-patch}
Let
\[
J\Subset\{c\neq 0,\ M\neq 0\}
\]
be a compact interval. Then there exists a constant $C_J>0$ such that every
normalized bad two-pair core with $h_1,h_2\in J$ and $h_1\neq h_2$
satisfies
\[
|E_{12}^{(3)}|+\|E_{12}^{(5)}\|+\bigl|\overline E_{12}^{(7;1)}\bigr|
\ge
C_J^{-1}\,|h_1-h_2|.
\]
In particular, there is no sequence of nontrivial normalized bad two-pair
cores in $J$ for which
\[
|E_{12}^{(3)}|+\|E_{12}^{(5)}\|+\bigl|\overline E_{12}^{(7;1)}\bigr|
=
o\!\left(|h_1-h_2|\right).
\]
\end{proposition}

\begin{proof}
By Corollary~\ref{cor:bad-two-pair-cores-defect-scale-close}, there exists a
constant $C'_J>0$ such that
\[
|h_1-h_2|
\le
C'_J\Bigl(|E_{12}^{(3)}|+\|E_{12}^{(5)}\|\Bigr)
\]
for every normalized bad two-pair core with $h_1,h_2\in J$. Since
\[
|E_{12}^{(3)}|+\|E_{12}^{(5)}\|
\le
|E_{12}^{(3)}|+\|E_{12}^{(5)}\|+\bigl|\overline E_{12}^{(7;1)}\bigr|,
\]
the same estimate yields
\[
|h_1-h_2|
\le
C'_J\Bigl(
|E_{12}^{(3)}|+\|E_{12}^{(5)}\|+\bigl|\overline E_{12}^{(7;1)}\bigr|
\Bigr).
\]
Renaming $C'_J$ as $C_J$ proves the first claim. The final statement is
immediate: if a sequence satisfied
\[
|E_{12}^{(3)}|+\|E_{12}^{(5)}\|+\bigl|\overline E_{12}^{(7;1)}\bigr|
=
o\!\left(|h_1-h_2|\right),
\]
then the displayed lower bound would force
\[
|h_1-h_2|=o\!\left(|h_1-h_2|\right),
\]
which is impossible for a nontrivial two-point core with $h_1\neq h_2$.
\end{proof}

\begin{lemma}[Minimal finite-order source criterion for quadratic two-point factorization]
\label{lem:minimal-source-level-two-point-criterion}
Let $V$ be a finite-dimensional complex vector space, and let
\[
\mathfrak P_2(a_1,h_1;a_2,h_2)\in V
\]
be analytic on a neighborhood of the small two-atom data. Assume that there is
a one-atom package
\[
F_h\in V
\]
such that the following three identities hold:
\[
\mathfrak P_2(a_1,h_1;a_2,h_2)
=
\mathfrak P_2(a_2,h_2;a_1,h_1),
\]
\[
\mathfrak P_2(0,h_1;a_2,h_2)=a_2F_{h_2},
\qquad
\mathfrak P_2(a_1,h_1;0,h_2)=a_1F_{h_1},
\]
and
\[
\mathfrak P_2(a_1,h;a_2,h)=(a_1+a_2)F_h.
\]
Define the interaction remainder
\[
\mathcal I_2(a_1,h_1;a_2,h_2)
:=
\mathfrak P_2(a_1,h_1;a_2,h_2)-a_1F_{h_1}-a_2F_{h_2}.
\]
Then there exists an analytic map
\[
\mathcal J_2:\ (\text{small two-atom data})\to V
\]
such that
\[
\boxed{
\mathcal I_2(a_1,h_1;a_2,h_2)
=
a_1a_2(h_1-h_2)^2\,\mathcal J_2(a_1,h_1;a_2,h_2).
}
\]
\end{lemma}

\begin{proof}
By the one-amplitude collapse identities,
\[
\mathcal I_2(0,h_1;a_2,h_2)=0,
\qquad
\mathcal I_2(a_1,h_1;0,h_2)=0.
\]
Hence analyticity implies divisibility by $a_1a_2$:
\[
\mathcal I_2(a_1,h_1;a_2,h_2)=a_1a_2\,G(h_1,h_2)
\]
for an analytic
\[
G:\ (h_1,h_2)\longmapsto V.
\]
Swap symmetry gives
\[
G(h_1,h_2)=G(h_2,h_1).
\]
On the diagonal,
\[
\mathcal I_2(a_1,h;a_2,h)
=
\mathfrak P_2(a_1,h;a_2,h)-a_1F_h-a_2F_h
=
(a_1+a_2)F_h-a_1F_h-a_2F_h
=
0,
\]
so
\[
G(h,h)=0.
\]
Now set
\[
m=\frac{h_1+h_2}{2},
\qquad
s=h_2-h_1,
\qquad
\widetilde G(m,s):=
G\!\left(m-\frac{s}{2},m+\frac{s}{2}\right).
\]
Since $G$ is symmetric, $\widetilde G(m,s)$ is even in $s$, and since
$G(h,h)=0$, one has
\[
\widetilde G(m,0)=0.
\]
Therefore $\widetilde G(m,s)$ is divisible by $s^2$:
\[
\widetilde G(m,s)=s^2H(m,s)
\]
for analytic $H$. Returning to $(h_1,h_2)$ gives
\[
\mathcal I_2(a_1,h_1;a_2,h_2)
=
a_1a_2(h_1-h_2)^2\,\mathcal J_2(a_1,h_1;a_2,h_2),
\]
as claimed.
\end{proof}

\begin{corollary}[Conditional quadratic defect estimate and two-point exclusion]
\label{cor:conditional-quadratic-defect-estimate-two-point-exclusion}
Assume that the actual corrected two-pair finite-order package through order $7$
admits a source-level realization
\[
\mathfrak P_2(a_1,h_1;a_2,h_2)
\in
\mathbf C\oplus \mathfrak f\oplus
\bigl(\mathfrak f/\mathbf C A_5^{\mathfrak f}\bigr)
\]
satisfying the three hypotheses of
Lemma~\ref{lem:minimal-source-level-two-point-criterion}, with
\[
F_h=
\bigl(c(h),A_5^{\mathfrak f}(h),[A_7^{\mathfrak f}](h)\bigr).
\]
Then on compacta,
\[
\mathcal E^{(3)}=a_1a_2\,O(\delta^2),
\qquad
\mathcal E^{(5)}=a_1a_2\,O(\delta^2),
\qquad
\overline{\mathcal E}^{(7)}=a_1a_2\,O(\delta^2),
\qquad
\delta:=|h_1-h_2|.
\]
If the normalized defects are obtained by dividing by $a_1$, then on normalized
families with $|a_2|$ bounded one has
\[
E_{12}^{(3)}=O(\delta^2),
\qquad
E_{12}^{(5)}=O(\delta^2),
\qquad
\overline E_{12}^{(7;1)}=O(\delta^2).
\]
Consequently, the two-point bad-core branch is excluded on compact good patches
\[
J\Subset\{c\neq 0,\ M\neq 0\}.
\]
\end{corollary}

\begin{proof}
Apply Lemma~\ref{lem:minimal-source-level-two-point-criterion} to the cubic,
quintic fixed-sector, and septic quotient outputs, then argue exactly as in the
present proof using
Theorem~\ref{thm:quantitative-defect-thickened-strengthened-coincidence}.
\end{proof}

\begin{proposition}[Conditional exact two-point exclusion from quotient diagonal collapse]
\label{prop:conditional-exact-two-point-exclusion-qabs}
Let
\[
J\Subset\{c\neq 0\}
\]
be a compact interval. Assume there exists an actual corrected two-atom quotient
map
\[
\mathcal F_2(a_1,h_1;a_2,h_2)\in \Qabs[[z^2]]
\]
defined on normalized two-pair data, with the following properties:

\begin{enumerate}
\item \textbf{Continuity.}
The map $\mathcal F_2$ is continuous in the normalized two-pair parameters.

\item \textbf{Diagonal collapse.}
For every $h\in J$ and every normalized amplitude pair $(a_1,a_2)$,
\[
\mathcal F_2(a_1,h;a_2,h)=F_h,
\]
where $F_h$ is the normalized one-pair quotient odd germ.

\item \textbf{Precompactness of normalized amplitudes.}
Every sequence of normalized two-pair data with
\[
h_{1,n},h_{2,n}\in J,
\qquad
|h_{1,n}-h_{2,n}|\to 0
\]
admits a subsequence along which the full normalized parameter tuple converges.
\end{enumerate}

Then there is no sequence of nontrivial normalized bad two-pair cores
\[
(a_{1,n},h_{1,n};a_{2,n},h_{2,n})
\]
with
\[
h_{1,n},h_{2,n}\in J,
\qquad
0<|h_{1,n}-h_{2,n}|\to 0.
\]
Equivalently, under the three assumptions above, the exact close two-point
bad-core branch is excluded on compact subsets of $\{c\neq 0\}$.
\end{proposition}

\begin{proof}
Assume for contradiction that there exists a sequence of nontrivial normalized
bad two-pair cores
\[
(a_{1,n},h_{1,n};a_{2,n},h_{2,n})
\]
with
\[
h_{1,n},h_{2,n}\in J,
\qquad
0<|h_{1,n}-h_{2,n}|\to 0.
\]
Since each core is bad, its quotient odd germ vanishes identically, so
\[
\mathcal F_2(a_{1,n},h_{1,n};a_{2,n},h_{2,n})=0
\qquad
\text{for all }n.
\]

By the precompactness assumption, after passing to a subsequence we may assume
that the full normalized parameter tuple converges:
\[
a_{1,n}\to a_1^*,
\qquad
a_{2,n}\to a_2^*,
\qquad
h_{1,n}\to m,
\qquad
h_{2,n}\to m
\]
for some $m\in J$. By continuity of $\mathcal F_2$,
\[
0
=
\lim_{n\to\infty}
\mathcal F_2(a_{1,n},h_{1,n};a_{2,n},h_{2,n})
=
\mathcal F_2(a_1^*,m;a_2^*,m).
\]
Applying diagonal collapse gives
\[
0=\mathcal F_2(a_1^*,m;a_2^*,m)=F_m.
\]

But $m\in J\Subset\{c\neq 0\}$, so the normalized one-pair quotient germ
$F_m$ is nonzero: its cubic leading coefficient is nonzero on $\{c\neq 0\}$.
This contradiction proves the claim.
\end{proof}

\begin{remark}[A weaker exact two-point route]
\label{rem:weaker-exact-two-point-route}
Proposition~\ref{prop:conditional-exact-two-point-exclusion-qabs} is strictly
weaker than the quadratic defect-factorization route of
Lemma~\ref{lem:minimal-source-level-two-point-criterion}. It does not produce
\[
\mathcal I_2=a_1a_2(h_1-h_2)^2\mathcal J_2
\]
and does not yield the quantitative defect estimates used in
Theorem~\ref{thm:quantitative-defect-thickened-strengthened-coincidence}.
Instead, it isolates a smaller alternative input: a genuine corrected two-atom
quotient map, continuity in normalized parameters, and exact diagonal collapse
to the one-pair quotient class.

Thus there are now two distinct two-point routes in the draft:

\begin{enumerate}
\item the stronger finite-order source-level route, which would force quadratic
defect bounds and hence quantitative defect-scale closeness;

\item the weaker quotient-diagonal route, which would already exclude exact
close two-point bad cores on compact subsets of $\{c\neq 0\}$.
\end{enumerate}

At present, the draft does not prove either route outright. The first remains a
finite-order source-level merger problem; the second remains a continuity and
diagonal-collapse problem for the actual corrected two-atom quotient germ.
\end{remark}

\begin{remark}[Exact remaining source-level input on the two-point branch]
\label{rem:wip-source-level-reduction-two-point-branch}
Lemma~\ref{lem:minimal-source-level-two-point-criterion} isolates the exact
finite-order source identities needed for the two-point branch. It is enough to
verify, for the actual corrected two-pair cubic/quintic/septic package, only:

\begin{enumerate}
\item swap symmetry;
\item one-amplitude collapse;
\item diagonal merger
\[
\mathfrak P_2(a_1,h;a_2,h)=(a_1+a_2)F_h.
\]
\end{enumerate}

Once these identities are known, the quadratic factorization
\[
\mathcal I_2=a_1a_2(h_1-h_2)^2\mathcal J_2
\]
is formal. So the remaining issue on the two-point branch is not another
asymptotic estimate, but verification of this finite-order source-level merger
statement for the corrected two-atom package.
\end{remark}

\begin{remark}[Audit: even-amplitude obstruction for the naive source-summed model]
\label{rem:audit-even-amplitude-obstruction-naive-source-summed-two-point}
The most obvious attempt to prove the source-level merger theorem is to insert a
weighted one-atom datum
\[
q_{a,h}=a\,q_h,
\qquad
\Phi_{a,h}=a\,\Phi_h
\]
directly into the corrected same-point and mixed block formulas, whiten, and
then extract the cubic/quintic/septic outputs. This route cannot by itself
prove the required one-atom linearity
\[
\mathfrak P_2(a\delta_h)=a\,F_h.
\]

Indeed the explicit corrected block formulas have the form
\[
G^{\corr}_{a,h,\pm}(s)=a\,\mathcal G_{h,\pm}(a^2;s),
\qquad
N^{\corr}_{a,h}(s)=a\,\mathcal N_h(a^2;s),
\]
so after whitening the overall factor of $a$ cancels:
\[
\widehat\Omega^{\corr}_{a,h}(s)
=
\mathcal G_{h,-}(a^2;s)^{-1/2}\,
\mathcal N_h(a^2;s)\,
\mathcal G_{h,+}(a^2;s)^{-1/2}.
\]
Hence every fixed finite-order coefficient extracted from the whitened block is
an even analytic function of the source weight $a$. In particular, the naive
source-summed corrected block package cannot satisfy analytic one-atom linearity
or coincident-atom additivity unless it vanishes identically. So the remaining
two-point theorem must be proved using a different lift/transport object, not
the raw source-summed whitened coefficients themselves.
\end{remark}

\begin{remark}[Audit: quotient and single-determinant routes are shear-blind]
\label{rem:audit-shear-blind-two-point-determinantal-routes}
On the mixed two-point branch, any route based only on cubic/quintic source
minors,
the septic quotient class,
or the pointwise patch gauges $v_7=0$ and $u_7=0$ is exactly blind to the
repeated-node shear parameter.

For an exact bad two-pair core, the cubic/quintic defect vector has the form
\[
d=\Lambda s_2-s_1,
\qquad
\Lambda:=-\frac{a_2}{a_1},
\qquad
s_i=(c_i,u_{5,i},v_{5,i}),
\]
so
\[
s_2\wedge d=s_1\wedge s_2
\]
is independent of $\Lambda$. Likewise, any pointwise representative extracted
only from the septic quotient class and $A_{5,1}^{\mathfrak f}$ carries no more
information than the same determinant scalar already seen by the quotient. Thus
the remaining mixed two-point burden is genuinely source-coupled rather than
quotient-determinantal.
\end{remark}

\begin{proposition}[Projective compactness of the moving $\lambda$-shear defects]
\label{prop:projective-compactness-moving-lambda-shear-two-point}
Assume a close exact mixed two-point bad branch on a compact good overlap patch
\[
U\Subset\{c\neq 0,\ v_5\neq 0,\ M\neq 0\},
\]
and write
\[
\Lambda_n:=-\frac{a_{2,n}}{a_{1,n}},
\qquad
\delta_n f:=f(h_{1,n})-\Lambda_n f(h_{2,n}),
\qquad
\lambda=\frac{\Delta_7}{v_5}.
\]
Then the mixed branch equations are equivalent to the existence of branch-generated
moving shears $(\eta_n,\xi_n,\gamma_n)$ such that
\[
\delta_n(c)=\eta_n\,\delta_n(\lambda),\qquad
\delta_n(v_5)=\xi_n\,\delta_n(\lambda),\qquad
\delta_n(u_5)=\gamma_n\,\delta_n(\lambda).
\]

Set
\[
d_n:=\bigl(\delta_n(c),\delta_n(v_5),\delta_n(u_5),\delta_n(\lambda)\bigr)\in\mathbf R^4.
\]
Then, after passing to a subsequence, exactly one of the following holds.

\begin{enumerate}
\item
$d_n=0$ for all $n$ in the subsequence. Equivalently, along that subsequence one already has exact repeated-source transport simultaneously for
\[
c,\qquad v_5,\qquad u_5,\qquad \lambda.
\]

\item
$d_n\neq 0$ for all $n$ in the subsequence, and the normalized defect directions
\[
\widehat d_n:=\frac{d_n}{\|d_n\|}\in S^3
\]
converge. Equivalently, the projective classes
\[
[d_n]
=
[\delta_n(c):\delta_n(v_5):\delta_n(u_5):\delta_n(\lambda)]
\in \mathbf P^3
\]
admit a convergent subsequence.
\end{enumerate}

Moreover, on any subsequence for which $\delta_n(\lambda)\neq 0$ and the projective
limit lies in the affine chart with nonzero $\lambda$-coordinate, the affine
$\lambda$-shears converge:
\[
\eta_n=\frac{\delta_n(c)}{\delta_n(\lambda)}\to\eta_*,
\qquad
\xi_n=\frac{\delta_n(v_5)}{\delta_n(\lambda)}\to\xi_*,
\qquad
\gamma_n=\frac{\delta_n(u_5)}{\delta_n(\lambda)}\to\gamma_*.
\]

For the limiting affine shear one has only the asymptotic fixed-shear identities
\[
\delta_n(c-\eta_*\lambda)=(\eta_n-\eta_*)\,\delta_n(\lambda),
\]
\[
\delta_n(v_5-\xi_*\lambda)=(\xi_n-\xi_*)\,\delta_n(\lambda),
\]
\[
\delta_n(u_5-\gamma_*\lambda)=(\gamma_n-\gamma_*)\,\delta_n(\lambda).
\]
Hence a close exact mixed branch automatically reduces only to a
\emph{projectively frozen secant-shadow problem}. In general it does \emph{not}
reduce to exact repeated-source transport for a single fixed triple
\[
(\eta_*,\xi_*,\gamma_*).
\]
\end{proposition}

\begin{proof}
The identities
\[
\delta_n(c)=\eta_n\,\delta_n(\lambda),\qquad
\delta_n(v_5)=\xi_n\,\delta_n(\lambda),\qquad
\delta_n(u_5)=\gamma_n\,\delta_n(\lambda)
\]
are just a rewriting of the branch equations. If $d_n=0$ infinitely often, pass
to that subsequence. Otherwise $d_n\neq 0$ for all large $n$, and compactness of
the unit sphere gives a convergent subsequence of $\widehat d_n$, equivalently
of the associated projective classes in $\mathbf P^3$.

On the affine chart $\delta_n(\lambda)\neq 0$, the projective coordinates are
exactly
\[
\eta_n=\frac{\delta_n(c)}{\delta_n(\lambda)},\qquad
\xi_n=\frac{\delta_n(v_5)}{\delta_n(\lambda)},\qquad
\gamma_n=\frac{\delta_n(u_5)}{\delta_n(\lambda)},
\]
so convergence in that chart is precisely convergence of the affine
$\lambda$-shears. The last three identities are immediate from the definitions.
\end{proof}

\begin{remark}[Affine freezing is stronger than projective freezing on the mixed two-point branch]
\label{rem:affine-freezing-stronger-than-projective-two-point}
Proposition~\ref{prop:projective-compactness-moving-lambda-shear-two-point}
shows that the automatic compactness on the mixed two-point branch is
projective, not affine. Even on a good overlap patch, the affine
$\lambda$-shears can diverge if $\delta_n(\lambda)$ vanishes to higher order
than $\delta_n(c),\delta_n(v_5),\delta_n(u_5)$. And even when
\[
(\eta_n,\xi_n,\gamma_n)\to(\eta_*,\xi_*,\gamma_*),
\]
one gets only asymptotic transport for
\[
c-\eta_*\lambda,\qquad
v_5-\xi_*\lambda,\qquad
u_5-\gamma_*\lambda,
\]
not exact transport at each stage of the branch. Therefore the downstream
derivative-geometric reduction should be formulated in terms of the projective
defect direction
\[
[\delta_n(c):\delta_n(v_5):\delta_n(u_5):\delta_n(\lambda)],
\]
or equivalently in any affine chart on which the projective limit has nonzero
coordinate, rather than as an exact frozen-$\lambda$ repeated-source branch.
\end{remark}

\begin{remark}[Collision-cancellation chart for close two-point cores]
\label{rem:wip-collision-cancellation-chart-two-point}
On a compact good patch
\[
J\Subset\{c\neq 0,\ M\neq 0\},
\]
a close exact two-point bad core is naturally studied not in the raw diagonal
chart $(h_1,h_2)\to(h,h)$, but in the cancellation chart
\[
m:=\frac{h_1+h_2}{2},
\qquad
\delta:=h_2-h_1,
\qquad
\lambda:=-\frac{a_2}{a_1},
\qquad
\omega:=\frac{\lambda-1}{\lambda+1}.
\]
Indeed, the cubic equation
\[
a_1c(h_1)+a_2c(h_2)=0
\]
forces $\lambda\to 1$ whenever $\delta\to 0$ on $J$.

The variable $\omega$ is the natural swap-compatible amplitude coordinate:
under interchange of the two atoms one has
\[
\lambda\longmapsto \lambda^{-1},
\qquad
(m,\omega,\delta)\longmapsto (m,-\omega,-\delta).
\]
Writing
\[
\mathcal D_h:=(c(h),u_5(h),v_5(h)),
\]
and using
\[
\lambda=\frac{1+\omega}{1-\omega}=1+2\omega+O(\omega^2),
\]
one gets the first-order expansion
\[
\mathcal D_{h_1}-\lambda \mathcal D_{h_2}
=
-2\omega\,\mathcal D_m-\delta\,\mathcal D_m'
+
O\!\bigl((|\omega|+|\delta|)^2\bigr).
\]
After polar blow-up
\[
(2\omega,\delta)=\rho(\cos\theta,\sin\theta),
\]
the leading cancellation datum is therefore the projective source direction
\[
\Theta_{m,\theta}
:=
\cos\theta\,\mathcal D_m+\sin\theta\,\mathcal D_m'.
\]

So any future cubic/quintic limit theorem for close cancelling two-point cores
should be phrased in this blown-up collision-cancellation chart rather than as a
Taylor expansion through the singular zero-source point.
\end{remark}

\begin{remark}[Work in progress: parity/projective factorization and the local good-patch frontier]
\label{rem:wip-parity-projective-factorization-collision-blow-up}
The collision-cancellation chart suggests a substantial weakening of the
hypotheses needed for local exclusion on the residual exact fixed-shear lane.

Suppose that an actual corrected defect $E(m,\omega,\delta)$ is analytic in
this swap-compatible chart, is swap-even,
\[
E(m,-\omega,-\delta)=E(m,\omega,\delta),
\]
and vanishes on the collision locus
\[
E(m,0,0)=0.
\]
Then its Taylor expansion has no linear terms, and with
\[
\kappa:=\frac{2\omega}{\delta}
\]
one can write
\[
E(m,\omega,\delta)=\delta^2\,\mathcal H(m,\kappa,\delta^2)
\]
for an analytic germ $\mathcal H$.

So the entire first blown-up interaction layer is parity-sterile: the first
surviving layer is already quadratic and one-projective-parameter. Along an
involutive exact fixed-shear branch this becomes quotient-visible, and on the
quartic--sextic rung the quadratic coefficient is expected to be controlled by
the single quotient transport scalar $Q$.

A sharper downstream target is the following conditional edge law for the
actual corrected cubic/quintic defects:
\[
E_{12}^{(3)}=\delta^2\,\mathcal H_3(m,\kappa,\delta^2),
\qquad
E_{12}^{(5)}=\delta^2\,\mathcal H_5(m,\kappa,\delta^2),
\]
with edge values
\[
\mathcal H_3(m,\kappa,0)=c'(m)-\frac{\kappa}{2}c(m),
\qquad
\mathcal H_5(m,\kappa,0)=(A_5^{\mathfrak f})'(m)-\frac{\kappa}{2}A_5^{\mathfrak f}(m).
\]
On the cubic-killing slope
\[
\kappa_*(m)=2\frac{c'(m)}{c(m)},
\]
this would give the reduced quintic obstruction
\[
\mathfrak R(m,0)=\frac{M(m)}{c(m)},
\qquad
M:=c(A_5^{\mathfrak f})'-c'A_5^{\mathfrak f}.
\]
Hence on every compact good patch
\[
J\Subset\{c\neq 0,\ M\neq 0\},
\]
sufficiently close cancelling bad cores would already be excluded.

Beyond that good-patch region, the same edge-law package would force any
quartic--sextic survivor into the exceptional locus $M=0$, where the residual
problem should reduce, after Weierstrass preparation, to finitely many prepared
branches and then to a one-variable compatibility / resultant / norm test.
This is a reduction of the remaining burden, not a closed theorem at present.
\end{remark}

\paragraph{Gate A and Gate B.}
\label{par:gate-a-gate-b}
\label{subsec:gate-a-gate-b}
The corrected-package audit uses two separate gates.  We call them Gate~A
and Gate~B to prevent the leading homogeneous cubic audit from being
confused with the full raw edge-divisibility theorem.

\begin{remark}[Roadmap of the Gate A/Gate B correction]
\label{rem:roadmap-gate-a-gate-b-correction}
The corrected-package audit separates into two layers.

\emph{Gate A} is the leading homogeneous cubic audit.  It is a finite,
order-\(7\), leading-source calculation.  Under the homogeneous scalar-grade
convention, Gate~A closes at the aggregate fixed-sector level:
\[
F_1^{\mathrm{lead\ cubic}}=0.
\]

\emph{Gate B} is the raw edge-divisibility theorem.  It concerns the full
non-\((1,1)\) projected raw corrected-whitening cross-effect
\[
(1-\Pi_{1,1})_{\mathrm{mix}}C_{\mathcal W}^{\mathrm{raw}}
\]
on the collision chart
\[
h_1=h-\frac{\delta}{2},\qquad
h_2=h+\frac{\delta}{2},\qquad
a_2=-a_1(1+\kappa\delta).
\]
Gate~B is the assertion
\[
[z^7][\delta^1]\,
\pi_{\mathfrak f_h}(1-\Pi_{1,1})_{\mathrm{mix}}C_{\mathcal W}^{\mathrm{raw}}=0.
\]

The proof of Gate~B below uses four inputs:
\begin{enumerate}
\item Gate~A aggregate closure;
\item odd total-grade subset-complement cancellation for amplitude terms;
\item moving Taylor-frame telescoping for the source-motion terms;
\item \(2\)-jet source-frame normalization, which removes the affine lower
boundary.
\end{enumerate}
Only after Gate~B is proved may one form the normalized edge package
\[
-a_1^{-1}\delta^{-2}
(1-\Pi_{1,1})_{\mathrm{mix}}C_{\mathcal W}^{\mathrm{raw}}.
\]
\end{remark}

\begin{definition}[Gate A: leading homogeneous cubic audit]
\label{def:gate-a-leading-homogeneous-cubic-audit}
Work under the homogeneous scalar-grade convention
\[
\operatorname{Gr}_a r
=
\frac{r^{(a+2)}(h)}{(a+2)!}(t-h)^{a+2}.
\]
Gate~A is the audit of the leading homogeneous cubic non-\((1,1)\) source
sector at ordinary order \(7\).  In the active convention, this is the
leading \((1,1,5)\)-type cubic sector.  The Gate~A families are
\[
L_1YR_1,\qquad L_2YR_0,\qquad L_0YR_2,
\]
\[
L_2N_0R_1,\qquad L_1N_0R_2,
\]
and
\[
L_3N_0R_0,\qquad L_0N_0R_3.
\]
Gate~A asks whether these leading homogeneous non-\((1,1)\) cubic terms
produce a fixed-sector order-\(7\) obstruction.
\end{definition}

\begin{definition}[Gate B: raw edge-divisibility]
\label{def:gate-b-raw-edge-divisibility}
Gate~B is the raw edge-divisibility theorem for the full projected
corrected-whitening cross-effect.  Let
\[
C_{\mathcal W}^{\mathrm{raw}}
=
\mathcal W(\mathcal B_2)
-\mathcal W(\mathcal B_1(a_1,h_1))
-\mathcal W(\mathcal B_1(a_2,h_2))
+\mathcal W(\mathcal B_0),
\]
where
\[
\mathcal W(G_-,N,G_+)=G_-^{-1/2}NG_+^{-1/2}.
\]
Apply the source-bidegree projection before fixed-sector extraction,
determinant scalarization, or quotienting:
\[
(1-\Pi_{1,1})_{\mathrm{mix}}C_{\mathcal W}^{\mathrm{raw}}.
\]
On the collision chart
\[
h_1=h-\frac{\delta}{2},\qquad
h_2=h+\frac{\delta}{2},\qquad
a_2=-a_1(1+\kappa\delta),
\]
Gate~B is the assertion
\[
(1-\Pi_{1,1})_{\mathrm{mix}}C_{\mathcal W}^{\mathrm{raw}}
=
O(a_1\delta^2)
\]
through fixed-sector ordinary order \(7\).  Equivalently, the
first-collision coefficient
\[
F_1:=
[z^7][\delta^1]\,
\pi_{\mathfrak f_h}
(1-\Pi_{1,1})_{\mathrm{mix}}C_{\mathcal W}^{\mathrm{raw}}
\]
vanishes.
\end{definition}

\begin{remark}[Why Gate B must be raw]
\label{rem:why-gate-b-must-be-raw}
Gate~B must be proved before forming the normalized package
\[
-a_1^{-1}\delta^{-2}
(1-\Pi_{1,1})_{\mathrm{mix}}C_{\mathcal W}^{\mathrm{raw}}.
\]
If \(F_1\ne0\), this normalized expression has a \(\delta^{-1}\)-pole and
the post-edge package is not holomorphic.  Therefore one must not use the
post-edge quotient statement
\[
Q_{7,m}^q(C)=Q_{7,m}^q(\Pi_{1,1}C)
\]
as the first proof of edge divisibility.  Quotient compatibility is a
post-edge statement unless it is reformulated at the raw edge level.
\end{remark}

\begin{lemma}[Order exhaustion for Gate A]
\label{lem:gate-a-order-exhaustion}
Under the homogeneous scalar-grade convention, every leading homogeneous
cubic non-\((1,1)\) contribution to \(B_7^{\mathfrak f}\) is one of the
seven families
\[
L_1YR_1,\quad L_2YR_0,\quad L_0YR_2,\quad
L_2N_0R_1,\quad L_1N_0R_2,\quad
L_3N_0R_0,\quad L_0N_0R_3.
\]
The three \(Y\)-slot families
\[
L_1YR_1,\qquad L_2YR_0,\qquad L_0YR_2
\]
are absent at \(B_7^{\mathfrak f}\).  The non-\(Y\) families reduce to the
leading placements
\[
L_2^{\{1,1\}}N_0[0]R_1^{[5]},
\qquad
L_2^{\{1,5\}}N_0[0]R_1^{[1]},
\]
and their right-side mirrors, together with the all-left/all-right
placements
\[
L_3^{\{1,1,5\}}N_0[0]R_0,
\qquad
L_0N_0[0]R_3^{\{1,1,5\}}.
\]
\end{lemma}

\begin{proof}[Proof sketch]
The \(Y\)-slot absence is the homogeneous scalar-grade order-count
statement: under this convention the relevant \(Y\)-middle entries do not
reach ordinary order \(7\).  The remaining non-\(Y\) placements are the
only ones whose Fr\'echet degree and homogeneous source grades can sum to
\(7\) with the baseline middle \(N_0[0]\).  This is scoped to the
homogeneous scalar-grade convention and does not apply to an ordinary
matrix-output \(M^{[5]}\) convention.
\end{proof}

\begin{lemma}[Aggregate closure of the leading \(L_2N_0R_1/L_1N_0R_2\) gate]
\label{lem:aggregate-closure-leading-l2n0r1-l1n0r2-gate}
Work under the homogeneous scalar-grade convention.  Let
\[
N:=N_0[0],
\qquad
W:=W_{-,0}=G_0^{-1/2},
\]
and let \(\Sigma_-\) denote the Sylvester inverse
\[
G_0^{1/2}Z+ZG_0^{1/2}=C.
\]
The leading same-point source inputs satisfy
\[
X_{-,1}^{[1]}=-\frac{1}{48\pi Q^2}(I-D),
\]
and
\[
X_{-,5}^{[5]}
=
-\frac{1}{92160\pi Q^{10}}(I-D)
=
\lambda X_{-,1}^{[1]},
\qquad
\lambda:=\frac{1}{1920Q^8}.
\]
Set
\[
A:=S_{-,1}^{[1]}=\Sigma_-(X_{-,1}^{[1]}),
\qquad
M:=WAW,
\]
and
\[
T_2^{1,1}:=\Sigma_-(-2AA),
\qquad
H:=2WAWAW-WT_2^{1,1}W.
\]
Then
\[
S_{-,5}^{[5]}=\lambda A,
\qquad
T_6^{1,5}=\lambda T_2^{1,1},
\qquad
H_{-,6}^{1,5}=\lambda H.
\]
Consequently the leading fixed-sector tag vectors satisfy
\[
C_{112}^{L_2N_0R_1}
=
C_{122}^{L_2N_0R_1}
=
\lambda\,\pi_{\mathfrak f}(HNM),
\]
and
\[
C_{112}^{L_1N_0R_2}
=
C_{122}^{L_1N_0R_2}
=
-\lambda\,\pi_{\mathfrak f}(MNH).
\]
Since \(H,N,M\) are symmetric and
\[
(HNM)^T=MNH,
\]
while \(\pi_{\mathfrak f}\) is transpose-invariant, one has
\[
\pi_{\mathfrak f}(HNM)=\pi_{\mathfrak f}(MNH).
\]
Therefore
\[
C_{112}^{L_2N_0R_1}+C_{112}^{L_1N_0R_2}=0,
\]
and
\[
C_{122}^{L_2N_0R_1}+C_{122}^{L_1N_0R_2}=0.
\]
Thus the leading homogeneous non-\(Y\) \(L_2N_0R_1/L_1N_0R_2\) gate
closes at the aggregate fixed-sector level before quotienting.
\end{lemma}

\begin{lemma}[Aggregate closure of the leading all-left/all-right cubic gate]
\label{lem:aggregate-closure-leading-all-left-all-right-cubic-gate}
Assume the midpoint parity convention extends to the third-Fr\'echet
coefficient table:
\[
G^+_{n;a,b,c}=(-1)^nG^-_{n;a,b,c}.
\]
Under the homogeneous scalar-grade convention, the first
all-left/all-right source multiset is \((1,1,5)\), and it contributes at
ordinary order \(7\).  Let
\[
G:=G^-_{7;1,1,5},
\qquad
W:=W_{-,0}=W_{+,0},
\qquad
N:=N_0[0].
\]
Then
\[
G^+_{7;1,1,5}=-G.
\]
The leading all-left and all-right fixed-sector contributions are
\[
C^{L_3N_0R_0}_{1,1,5}
=
\pi_{\mathfrak f}(GNW),
\]
and
\[
C^{L_0N_0R_3}_{1,1,5}
=
-\pi_{\mathfrak f}(WNG).
\]
Since \(G,N,W\) are symmetric,
\[
(GNW)^T=WNG.
\]
Because \(\pi_{\mathfrak f}\) is transpose-invariant,
\[
\pi_{\mathfrak f}(GNW)=\pi_{\mathfrak f}(WNG).
\]
Hence
\[
C^{L_3N_0R_0}_{1,1,5}
+
C^{L_0N_0R_3}_{1,1,5}
=0
\quad\text{in }\mathfrak f.
\]
Thus the leading all-left/all-right homogeneous third-Fr\'echet gate
closes at the aggregate fixed-sector level before quotienting.
\end{lemma}

\begin{corollary}[Gate A closure]
\label{cor:gate-a-closure}
Under the homogeneous scalar-grade convention and the stated
midpoint-parity conventions,
\[
F_1^{\mathrm{lead\ cubic}}=0.
\]
Equivalently, the leading homogeneous cubic non-\((1,1)\) sector
contributes no fixed-sector \(B_7^{\mathfrak f}\) obstruction and no
first-collision Gate~B obstruction.
\end{corollary}

\begin{remark}[Aggregate closure versus family-by-family gauge]
\label{rem:aggregate-closure-versus-family-by-family-gauge}
The preceding Gate~A closures are aggregate fixed-sector cancellations.
They do not assert that the individual families
\[
L_2N_0R_1,\quad L_1N_0R_2,\quad L_3N_0R_0,\quad L_0N_0R_3
\]
are separately proportional to \(A_5^{\mathfrak f}(h)\).  The package
theorem requires the aggregate corrected-whitening cross-effect, not a
family-by-family gauge statement.  No determinant scalarization, no
\(\Phi_K\)-scalar invisibility, and no downstream diagonal merger is used
in these closures.
\end{remark}

\begin{remark}[Scope of Gate A]
\label{rem:scope-of-gate-a}
Gate~A is scoped to the active homogeneous scalar-grade convention.  It
does not apply to an ordinary matrix-output \(M^{[5]}\) convention, nor to
a bookkeeping system that treats lower-shadow grades \(0,2,4\) as
independent ordinary matrix-output inputs.  Such conventions require a
separate no-double-counting theorem.
\end{remark}

\begin{remark}[Corrected package theorem after the Gate A/Gate B split]
\label{rem:wip-first-cubic-source-transfer-target}
The current corrected-package frontier is no longer the raw seven-family
Fr\'echet shortcut.  The proof architecture is layered as follows.

First, Gate~A closes the leading homogeneous cubic non-\((1,1)\) audit at
the aggregate fixed-sector level.  This uses the homogeneous scalar-grade
convention, the absence of the \(Y\)-slot families at \(B_7\), the aggregate
\(L_2N_0R_1/L_1N_0R_2\) cancellation, and the aggregate
\(L_3N_0R_0/L_0N_0R_3\) cancellation, the last under the stated
third-Fr\'echet midpoint parity convention.

Second, Gate~B proves raw edge-divisibility for the projected raw
corrected-whitening cross-effect:
\[
[z^7][\delta^1]\,
\pi_{\mathfrak f_h}
(1-\Pi_{1,1})_{\mathrm{mix}}C_{\mathcal W}^{\mathrm{raw}}=0.
\]
This is the raw theorem needed before the amplitude-normalized package
\[
-a_1^{-1}\delta^{-2}
(1-\Pi_{1,1})_{\mathrm{mix}}C_{\mathcal W}^{\mathrm{raw}}
\]
can be formed.

Only after Gate~B is proved does the corrected fixed-target package become
available for the quotient-gauge theorem
\[
B_7^{\mathfrak f}(C_{\mathcal W}^{\corr})
\in
\mathbf C A_5^{\mathfrak f}(m),
\]
equivalently
\[
Q_{7,m}^q(C_{\mathcal W}^{\corr})=0.
\]
Thus any statement of the form
\[
Q_{7,m}^q(C)=Q_{7,m}^q(\Pi_{1,1}C)
\]
is post-edge bookkeeping unless it is explicitly reformulated as a raw
edge-divisibility statement.

The older raw seven-family finite-coefficient tests remain useful only as
diagnostic coefficient tools.  They are not the theorem object and should
not be used in place of the corrected package representative.
\end{remark}

\begin{definition}[Homogeneous scalar grade]
\label{def:homogeneous-scalar-grade}
For a residual scalar source \(r\), define
\[
\operatorname{Gr}_a r
:=
\frac{r^{(a+2)}(m)}{(a+2)!}(t-m)^{a+2}.
\]
Thus grade \(1\) corresponds to the cubic scalar jet \(r^{(3)}(m)\), and
grade \(5\) corresponds to the septic scalar jet \(r^{(7)}(m)\).
\end{definition}

\begin{lemma}[Transpose parity of the source-linear mixed input]
\label{lem:transpose-parity-source-linear-mixed-input}
For the full source-linear scaled mixed input \(M(z)\) before \(\Phi_K\), one
has
\[
M(-z)=M(z)^T.
\]
Consequently the diagonal entries of \(M(z)\) are even in \(z\), while the
odd ordinary powers occur in the off-diagonal antisymmetric channel.  Through
ordinary order \(7\), the coefficient \([z^5]M\) is off-diagonal
antisymmetric and is supported by the even source derivatives
\[
r^{(2)}(m),\qquad r^{(4)}(m),\qquad r^{(6)}(m),
\]
not by the homogeneous grade-five source \(r^{(7)}(m)\).
\end{lemma}

\begin{remark}[Verification sketch for the transpose-parity mixed input]
\label{rem:verification-sketch-transpose-parity-mixed-input}
Changing \(z\) to \(-z\) interchanges the two endpoints in the scaled mixed
block.  The mixed block is therefore transformed by endpoint swap, which is
matrix transpose in the left/right basis.  Hence the source-linear mixed
input satisfies
\[
M(-z)=M(z)^T.
\]
The diagonal and symmetric off-diagonal channels are even, while the
antisymmetric off-diagonal channel is odd.  The audited order-\(7\)
expansion places \([z^5]M\) in that antisymmetric channel, supported by
\[
r^{(2)}(m),\qquad r^{(4)}(m),\qquad r^{(6)}(m),
\]
not by the homogeneous grade-five source \(r^{(7)}(m)\).
\end{remark}

\begin{lemma}[\(Y\)-slot families are absent at \(B_7\) under homogeneous grade]
\label{lem:y-slot-families-absent-homogeneous-grade}
Under Definition~\ref{def:homogeneous-scalar-grade}, the \(Y\)-slot cubic
families
\[
L_1YR_1,\qquad
L_2YR_0,\qquad
L_0YR_2
\]
do not contribute to \(B_7^{\mathfrak f}\).

More precisely, the homogeneous grade-five mixed input associated to
\(\operatorname{Gr}_5 r\) has no ordinary \(z^5\) middle coefficient in the
mixed block; its first possible mixed contribution occurs too late to combine
with the corresponding first- or second-Fr\'echet same-point factors at total
ordinary order \(7\).  Consequently the three \(Y\)-slot families above start
at ordinary order at least \(8\).
\end{lemma}

\begin{remark}[Scope of Lemma~\ref{lem:y-slot-families-absent-homogeneous-grade}]
\label{rem:scope-y-slot-absence-homogeneous-grade}
Lemma~\ref{lem:y-slot-families-absent-homogeneous-grade} is an order-count
absence theorem, not a determinant-proportionality theorem.  It is scoped to
the homogeneous scalar grade convention.  If a future theorem defines
\(M^{[5]}\) instead as an ordinary matrix-output order-\(5\) projection, or
as a parity-corrected non-homogeneous source projection, then the \(Y\)-slot
families must be reopened and audited separately.  Such a convention also
requires a no-double-counting theorem against the natural grade \(0,2,4\)
same-point shadows.
\end{remark}

\begin{remark}[Status of the homogeneous-grade lemmas after Gate A/B]
\label{rem:status-homogeneous-grade-lemmas}
Lemmas~\ref{lem:transpose-parity-source-linear-mixed-input}
and~\ref{lem:y-slot-families-absent-homogeneous-grade} are now inputs to
the Gate~A/Gate~B closure mechanism.

Under the homogeneous scalar-grade convention, the \(Y\)-slot families are
absent at \(B_7\), the leading non-\(Y\) Gate~A aggregate closes in
\(\mathfrak f_h\), and Gate~B closes through ordinary order \(7\) after
all-subset midpoint parity, moving source-frame telescoping, and affine
boundary removal are imposed.

This does not mean that lower-shadow grades \(0,2,4\) are discarded.  They
may appear in the moving Taylor tower.  Interior lower-shadow terms are
handled by the telescoping identity of
Theorem~\ref{thm:moving-homogeneous-source-frame-telescoping}; only the
affine grade \(-1\) boundary is removed by the \(2\)-jet source-frame
normalization of
Lemma~\ref{lem:affine-boundary-removal-moving-source-frame}.
\end{remark}

\begin{remark}[Grade-zero sector after the Gate B source-frame audit]
\label{rem:grade-zero-sector-after-gate-b-source-frame-audit}
The homogeneous scalar-grade convention
\[
\operatorname{Gr}_a r
=
\frac{r^{(a+2)}(m)}{(a+2)!}(t-m)^{a+2}
\]
does not discard the grade-zero component
\[
\operatorname{Gr}_0r=\frac{r^{(2)}(m)}{2!}(t-m)^2.
\]
Affine removal eliminates only the constant and affine modes.  Thus
\(r^{(2)}(m)\) is a genuine source-frame direction, not baseline data.

The Gate~B result treats this direction differently from the old
open-status notes.  Grade \(0\) is allowed in the moving homogeneous Taylor
tower.  Its amplitude contributions are covered by the odd total-grade
subset-complement cancellation, including zero-padded labelled multisets.
Its source-motion lower boundary is
\[
(\partial_h\operatorname{Gr}_{0}r)_{\mathrm{lower}}
=
-r^{(2)}(h)(t-h),
\]
which is affine and is removed by the \(2\)-jet source-frame projection
\(\Pi_{\ge2,h}\).

Consequently grade \(0\) is neither an independent leftover Gate~B residual
nor something silently absorbed into \(q_0\).  It is a live source-frame
direction whose first-collision contribution is killed by the Gate~B
source-frame mechanism.
\end{remark}

\begin{remark}[Residual-source tower before highest-new reduction]
\label{rem:residual-source-tower-before-highest-new-reduction}
The homogeneous scalar-grade convention organizes the local residual source
tower, but it does not by itself identify all residual-source coefficients
with the highest-new tower \(q^{(7)}\).  Finite quotient-gauge identities
must first be tested in the residual-source variables supplied by the
corrected source block.  Only a residual that survives this package-level
test should be passed to the later same-tower / highest-new reduction.

Thus the statement that the eventual local bridge should run through the
single highest-new source \(q^{(7)}\) is a downstream package-to-jet bridge
claim.  It is not a substitute for the pre-\(\Phi_K\), pre-determinant
fixed-sector quotient-gauge calculation.
\end{remark}

\begin{remark}[Quarantined matrix-output \(M^{[5]}\) convention]
\label{rem:quarantined-matrix-output-M5-convention}
A nonzero ordinary matrix-output order-\(5\) convention for \(M^{[5]}\) is
not the active homogeneous scalar-grade convention.  If such a convention is
later introduced, then the \(Y\)-slot families
\[
L_1YR_1,\qquad L_2YR_0,\qquad L_0YR_2
\]
must be reopened.  That route requires a no-double-counting theorem
separating the ordinary matrix-output order-\(5\) contribution from the
natural homogeneous scalar grades \(0,2,4\), whose source derivatives also
have same-point shadows.  The earlier raw-seven-family numerical failure
recorded in Remark~\ref{rem:computational-status-raw-package} rejects the raw
matrix-output shortcut; it does not reject the homogeneous scalar-grade gate
of Lemma~\ref{lem:y-slot-families-absent-homogeneous-grade} or the
corrected package representative of
Section~\ref{sec:corrected-fixed-target-quotient-package}.
\end{remark}

\begin{remark}[Historical role of the first non-\(Y\) cubic gate]
\label{rem:first-live-non-y-cubic-gate}
Before the aggregate Gate~A closure was isolated, the first non-\(Y\)
homogeneous cubic gate was the pair
\[
L_2N_0R_1
\qquad\text{and}\qquad
L_1N_0R_2.
\]
The finite coefficient inputs for this diagnostic gate were
\[
[z^2]\mathcal L_-^{\{1,1\}},\qquad
[z^6]\mathcal L_-^{\{1,5\}},
\]
\[
[z^2]\mathcal R_+^{\{1,1\}},\qquad
[z^6]\mathcal R_+^{\{1,5\}},
\]
together with
\[
N_0[0],
\qquad
[z^1]\Lambda_\pm^{[1]},
\qquad
[z^5]\Lambda_\pm^{[5]}.
\]
These data are still useful for auditing the finite coefficient tables.

However, in the current proof they are no longer an open Gate~A target.
Lemma~\ref{lem:aggregate-closure-leading-l2n0r1-l1n0r2-gate} proves that
the \(L_2N_0R_1/L_1N_0R_2\) leading homogeneous gate closes at the
aggregate fixed-sector level.  Therefore this remark records historical
coefficient-tool provenance, not a remaining theorem obligation.
\end{remark}

\begin{remark}[Superseded coefficient test for the first non-\(Y\) gate]
\label{target:leading-coefficient-test-first-non-y-cubic-gate}
The earlier coefficient-level audit asked for the four fixed-sector tag
vectors
\[
C_{112}^{L_2N_0R_1},\qquad
C_{122}^{L_2N_0R_1},\qquad
C_{112}^{L_1N_0R_2},\qquad
C_{122}^{L_1N_0R_2}.
\]
That was the correct finite test before the aggregate closure mechanism
was identified.

In the current draft this test is superseded by
Lemma~\ref{lem:aggregate-closure-leading-l2n0r1-l1n0r2-gate}, which proves
the cancellation in \(\mathfrak f\) directly:
\[
C_{112}^{L_2N_0R_1}+C_{112}^{L_1N_0R_2}=0,
\qquad
C_{122}^{L_2N_0R_1}+C_{122}^{L_1N_0R_2}=0.
\]
Thus the old determinant checks against \(A_5^{\mathfrak f}\) should not
be read as remaining obligations for Gate~A.  They remain only as a
useful coordinate-level audit if the homogeneous scalar-grade convention
is later changed or the aggregate proof is challenged.
\end{remark}

\begin{remark}[All-left/all-right cubic families: aggregate status]
\label{rem:all-left-all-right-families-aggregate-status}
The all-left/all-right families
\[
L_3N_0R_0,
\qquad
L_0N_0R_3
\]
are not closed family-by-family.  However, under the homogeneous
scalar-grade convention and assuming third-Fr\'echet midpoint parity, their
leading \((1,1,5)\) aggregate satisfies
\[
L_3N_0R_0+L_0N_0R_3=0
\quad\text{in }\mathfrak f.
\]
This is stronger than quotient-gauge for the package-level Gate~A target.
It does not assert that either family is separately proportional to
\(A_5^{\mathfrak f}\).

If a coefficient table uses ordered triples rather than symmetric
multilinear coefficients, the aggregate must include all ordered
permutations with the same convention on both left and right sides before
this cancellation is invoked.
\end{remark}

\begin{remark}[Current UV chain and status]
\label{rem:current-uv-chain-and-status}
The current package chain is as follows.

\begin{enumerate}
\item \textbf{UV-022.} Define an invariant source-weight-linear input for the
corrected block and obtain exact one-amplitude collapse at the package level.

\item \textbf{UV-023.} Establish the raw first-collision edge law for the
non-\((1,1)\) projected corrected-whitening cross-effect through ordinary
order \(7\):
\[
[z^7][\delta^1]\,
\pi_{\mathfrak f_h}
(1-\Pi_{1,1})_{\mathrm{mix}}C_{\mathcal W}^{\mathrm{raw}}=0.
\]
In the present draft this is supplied, conditionally, by
Corollary~\ref{cor:gate-b-raw-edge-divisibility-order-seven}.  The
mechanism is not a \(\mathcal T_2\)-only truncation theorem, not
source-weight linearization, and not post-edge quotient compatibility.
It uses Gate~A aggregate closure, odd total-grade subset-complement
cancellation, grade-preserving source-motion cancellation, moving
source-frame telescoping, affine boundary removal, and full-tower
differentiation before ordinary-order truncation.

\item \textbf{UV-024.} Define the pre-determinant fixed-sector coefficient
\[
B_7^{\mathfrak f}(C)
\]
for the actual corrected-whitening matrix cross-effect before \(\Phi_K\), and
define
\[
Q_{7,m}^q(C):=[B_7^{\mathfrak f}(C)]
\in
\mathfrak f/\mathbf C A_5^{\mathfrak f}(m).
\]

\item \textbf{UV-025.} Supply the tagged pre-whitening source block of
Definition~\ref{def:tagged-two-atom-prewhitening-source-block}.  This is now
a guarded source-level input, not a quotient theorem.

\item \textbf{Gate A / UV-026 leading cubic audit.}
Under the homogeneous scalar-grade convention, the leading homogeneous
cubic non-\((1,1)\) sector is closed at the aggregate fixed-sector level.
The \(Y\)-slot families are absent at \(B_7\), and the non-\(Y\)
aggregate pairs
\[
L_2N_0R_1+L_1N_0R_2,
\qquad
L_3N_0R_0+L_0N_0R_3
\]
vanish in \(\mathfrak f\), the latter assuming third-Fr\'echet midpoint
parity.

\item \textbf{Gate B raw edge-divisibility.}
The projected raw corrected-whitening cross-effect satisfies
\[
[z^7][\delta^1]\,
\pi_{\mathfrak f_h}
(1-\Pi_{1,1})_{\mathrm{mix}}C_{\mathcal W}^{\mathrm{raw}}=0
\]
through ordinary order \(7\), provided the all-subset midpoint parity,
leading \(Y\)-middle parity, moving source-frame telescoping, and affine
boundary-removal lemmas are imposed.  This supplies the raw edge
divisibility needed before amplitude normalization and before the lower
package table \(Z_j[M]\) is formed.

\item \textbf{Connection correction subtarget.}
After Gate~B, compute the transported moving line coefficients
\[
A_i=[(h-m)^i]\tau_{h\to m}A_5^{\mathfrak f}(h),
\qquad
1\le i\le7,
\]
and the pre-trivialization lower edge coefficients
\[
Z_j[M_A],\quad Z_j[M_C],
\qquad
0\le j\le6,
\]
for
\[
M_A=\dot r_3r_2\dot r_4,
\qquad
M_C=\dot r_3\dot r_2r_4.
\]
The required cancellation identities are
\[
R_7^{(A)}+\sum_{i=1}^7\omega_m(Z_{7-i}[M_A],A_i)=0,
\]
and
\[
R_7^{(C)}+\sum_{i=1}^7\omega_m(Z_{7-i}[M_C],A_i)=0.
\]
These are package-level pre-trivialization connection identities, not raw
seven-family identities and not post-trivialization \(B_j[M]\) identities.
\end{enumerate}

Shortcut proofs using \(\Phi_K\)-scalar hiding, determinant scalarization,
swap symmetry, one-pair gauge, or the UV-025 source block alone do not prove
UV-026.  The theorem must be proved at the pre-\(\Phi_K\), pre-determinant
fixed-sector quotient level of
Section~\ref{sec:corrected-fixed-target-quotient-package}.

In this formulation, UV-026 is a theorem about the normalized corrected
package representative.  The raw seven-family Fr\'echet sum is only a
diagnostic object; its failure under a matrix-output \(M^{[5]}\) convention
is recorded in Remark~\ref{rem:computational-status-raw-package} and does
not replace the corrected package theorem.
\end{remark}

\begin{remark}[Superseded Gate B status notes]
\label{rem:superseded-gate-b-status-notes}
Earlier WIP notes treated \(p\ge4\) transfer filtration and lower-shadow
absorption as open obstructions.  That was the correct status before the
homogeneous transfer audit.  The current proof replaces those open
obligations as follows:
\begin{enumerate}
\item positive-grade amplitude terms of total grade \(7\) vanish by odd
total-grade subset-complement cancellation;
\item the grade-preserving part of \(X_h\) vanishes by the same
cancellation;
\item the grade-lowering part of \(X_h\) telescopes across adjacent
moving Taylor grades;
\item the sole lower boundary produced by \(\operatorname{Gr}_0\) is
affine and is removed by \(2\)-jet source-frame normalization.
\end{enumerate}
Thus the remaining role of the old notes is as guardrails: do not prove
Gate~B by post-edge quotient compatibility, do not demand
family-by-family vanishing, and do not truncate the source tower before
differentiating.
\end{remark}

\begin{remark}[The old \(d\ge4\) residual partition list is superseded]
\label{rem:old-dge4-residual-list-superseded}
Earlier audits listed positive-grade residual partitions such as
\[
(1,1,1,4),\quad (1,1,2,3),\quad (1,2,2,2),
\]
\[
(1,1,1,1,3),\quad (1,1,1,2,2),\quad
(1,1,1,1,1,2),\quad (1,1,1,1,1,1,1)
\]
as live \(d\ge4\) transfer-filtration candidates.  Those lists are now
superseded by the odd total-grade subset-complement cancellation and the
moving source-frame telescoping theorem.  They should be used only as
historical diagnostics for checking that the aggregate proof has included
all labelled placements, not as remaining theorem obligations.
\end{remark}

\begin{remark}[Lower-shadow sectors and zero-padding after Gate B]
\label{rem:lower-shadow-zero-padding-after-gate-b}
The old lower-shadow obstruction was not a finite missing table.  If
grade \(0\) is treated as an independent source factor, then the
source-grade partition inventory is non-finite:
\[
(0,1,1,5),\qquad
(0,0,1,1,5),\qquad
(\underbrace{0,\ldots,0}_{N},1,1,5),\qquad \ldots
\]
and similarly for
\[
(\underbrace{0,\ldots,0}_{N},7),\qquad
(\underbrace{0,\ldots,0}_{N},1,2,4),\qquad
(\underbrace{0,\ldots,0}_{N},2,2,3).
\]
The Gate~B proof does not close this tower by finite enumeration.  It
closes it uniformly.

More precisely, lower-shadow grades \(0,2,4\) are not reclassified as
baseline data and are not silently discarded.  They remain part of the
moving homogeneous Taylor source tower.  Their amplitude contributions of
total homogeneous grade \(7\) are killed by odd total-grade
subset-complement cancellation for arbitrary labelled transfer degree.
Their source-motion contributions are handled by the moving source-frame
telescoping identity, and the only lower boundary generated by
differentiating \(\operatorname{Gr}_0\) is the affine grade \(-1\) mode,
which is removed by the \(2\)-jet source-frame normalization.

Thus the correct post-Gate~B status is: lower shadows are live
source-frame directions, but not live Gate~B first-collision residuals.
They should not be listed as an infinite family of remaining residual
partitions after
Corollary~\ref{cor:gate-b-raw-edge-divisibility-order-seven}.
\end{remark}

\begin{remark}[Counterterm route if Gate B dependencies fail]
\label{rem:counterterm-route-if-gate-b-dependencies-fail}
The present draft proves Gate~B under the stated midpoint-parity,
source-frame telescoping, and \(2\)-jet normalization hypotheses.  If any
of those hypotheses fails in the final source model, then one cannot
simply subtract the residual after determinant scalarization.

A legitimate counterterm would have to be a canonical pre-\(\Phi_K\),
pre-determinant, source-bidegree-compatible edge counterterm
\(C_{\mathrm{diag}/\mathrm{edge}}\) determined by the raw edge jets
\[
[z^7][\delta^0]\pi_{\mathfrak f_h}
(1-\Pi_{1,1})_{\mathrm{mix}}C_{\mathcal W}^{\mathrm{raw}},
\qquad
[z^7][\delta^1]\pi_{\mathfrak f_h}
(1-\Pi_{1,1})_{\mathrm{mix}}C_{\mathcal W}^{\mathrm{raw}}.
\]
It must be natural under atom swap, must preserve the
\((1-\Pi_{1,1})_{\mathrm{mix}}\) source-bidegree sector, and must be
defined before \(\Phi_K\), determinant pairing, and quotienting.
Otherwise it is merely subtracting the obstruction by hand.
\end{remark}

\begin{remark}[Relation to the package-to-jet bridge]
\label{rem:relation-to-package-to-jet-bridge}
The corrected package program above is not an alternative to the zeta-side
odd-jet extractor.  The extractor side is already available; the missing
bridge is the package-to-jet theorem identifying the first surviving
corrected odd coefficient from finite-core package data.  In the
fixed-shear corner, the finite quotient-visible transport has already been
compressed to the remaining package-level data.  Thus any surviving
obstruction after the corrected package quotient theorem must be treated as
a package-to-transform or hidden-state problem, not as another planar or
scalarized local detector.

In particular, Gate~A and Gate~B are package-construction results.  They do
not by themselves prove the final finite-core contradiction or the
package-to-\(\Xi_\zeta^{(N)}\) hidden-state theorem.  Any residual
\(\Phi_K\)-visible affine-lift or higher odd-block state remains a separate
package-to-transform issue.
\end{remark}

\begin{remark}[Open theorem target: finite-core lower transfer]
\label{rem:open-finite-core-lower-transfer}
The finite-core branch would be closed by a theorem of the following form.

\emph{Finite-core lower-transfer theorem.}
Let \(\mathcal C\) be a normalized genuinely finite same-scale core whose
corrected package anomaly is nonzero after all quotient and hidden-state
identifications relevant to \(\Phi_K\).  Then there exist an admissible
\(N=N(m)\), a sign \(\sigma_{\mathcal C}\in\{\pm1\}\), and a constant
\(c_{\mathcal C}>0\), stable on the compact core-parameter patch, such that
\[
|\Xi_{\core}^{(N)}(m)|
\ge
c_{\mathcal C}\frac{x}{B_\zeta(m)}S(m).
\]

The present draft does not prove this theorem.  The corrected order-\(7\)
package controls the local quotient construction and removes the known
two-point/package quotient residuals, but it does not by itself determine the
\(N\)-visible coefficient
\[
c_{\core,2N-1}(m)
\]
for a genuinely non-pair-like finite core, nor does it rule out cancellation
among finite-core constituents.
\end{remark}

\begin{definition}[Nonsingular actual-source \(k\)-core stratum]
\label{def:nonsingular-actual-source-k-core-stratum}
A compact parameter patch \(\Theta_k\) is called a nonsingular actual-source
\(k\)-core stratum if the following conditions hold.

\begin{enumerate}
\item \(\Theta_k\) is contained in a finite-dimensional real-analytic
parameter manifold with fixed core cardinality \(k\), fixed source
combinatorics, and fixed collision/blow-up chart.

\item The corrected finite-core transverse scalar
\[
H_{\core,\theta}(z)
=
\Phi_K\!\left(
\widehat\Omega_{\core,\theta}^{\corr}(z/Q^2;m)
\right)
\]
is real-analytic jointly in \((\theta,z)\) on a neighborhood of
\(\Theta_k\times\{0\}\).

\item The whitening, inverse-square-root, and midpoint-trivialization maps used
to define \(\widehat\Omega_{\core,\theta}^{\corr}\) are nonsingular and
analytic on the patch.

\item The chosen coordinate chart does not divide by a vanishing quantity.  In
particular, if the chart uses \(c^{-1}\), \(u_5^{-1}\), \(v_5^{-1}\), or a
determinant coordinate relative to \(A_5^{\mathfrak f}\), the corresponding
quantity is nonzero on the patch.
\end{enumerate}
\end{definition}

\begin{theorem}[Finite odd-germ determinacy on compact actual-source strata]
\label{thm:finite-odd-germ-determinacy-actual-source}
Let \(\Theta_k\) be a compact nonsingular actual-source \(k\)-core stratum.
Write
\[
H_{\core,\theta}(z)
=
\sum_{r\ge0}c_{\core,2r+1}(\theta)z^{2r+1}.
\]
Then there exists an integer \(R_{\Theta_k}\ge0\) such that
\[
c_{\core,1}(\theta)=c_{\core,3}(\theta)=\cdots
=c_{\core,2R_{\Theta_k}+1}(\theta)=0
\]
implies
\[
H_{\core,\theta}\equiv0
\]
as an odd germ.

Equivalently, the zero-anomaly locus
\[
\mathcal Z_k
=
\{\theta\in\Theta_k:H_{\core,\theta}\equiv0\}
\]
is cut out on \(\Theta_k\) by finitely many odd coefficient equations.
\end{theorem}

\begin{proof}
By joint analyticity of \(H_{\core,\theta}(z)\), each coefficient
\[
c_{\core,2r+1}(\theta)
\]
is real-analytic on a neighborhood of \(\Theta_k\).

Fix \(\theta_0\in\Theta_k\).  Let
\[
\mathcal I_{\theta_0}
=
\bigl(c_{\core,1},c_{\core,3},c_{\core,5},\ldots\bigr)
\subset
\mathcal O_{\Theta_k,\theta_0}
\]
be the ideal generated by the odd coefficient germs in the local analytic ring
at \(\theta_0\).  The local analytic ring is Noetherian.  Hence
\(\mathcal I_{\theta_0}\) is finitely generated.  Therefore there are indices
\(r_1,\ldots,r_s\) such that
\[
\mathcal I_{\theta_0}
=
\bigl(
c_{\core,2r_1+1},\ldots,c_{\core,2r_s+1}
\bigr)
\]
as germs at \(\theta_0\).

Set
\[
R(\theta_0)=\max(r_1,\ldots,r_s).
\]
After shrinking to a neighborhood \(U_{\theta_0}\), every coefficient
\(c_{\core,2r+1}\) is an analytic linear combination of
\[
c_{\core,1},c_{\core,3},\ldots,c_{\core,2R(\theta_0)+1}.
\]
Thus, on \(U_{\theta_0}\), vanishing of all odd coefficients up to order
\(2R(\theta_0)+1\) forces every odd coefficient to vanish.  Hence
\(H_{\core,\theta}\equiv0\) as an odd germ.

The compactness of \(\Theta_k\) gives a finite subcover
\[
U_{\theta_1},\ldots,U_{\theta_M}.
\]
Taking
\[
R_{\Theta_k}
=
\max_{1\le j\le M}R(\theta_j)
\]
proves the claim uniformly on \(\Theta_k\).
\end{proof}

\begin{remark}[Strata excluded from one finite-determinacy patch]
\label{rem:finite-determinacy-excluded-strata}
Theorem~\ref{thm:finite-odd-germ-determinacy-actual-source} is stratum-local.
The following loci must be separated before applying it.

\begin{enumerate}
\item \emph{Amplitude-drop strata:}
\[
a_j=0,
\]
where a \(k\)-core degenerates to a smaller core.

\item \emph{Collision strata:}
\[
h_i=h_j,
\]
unless a collision/blow-up coordinate chart has been chosen in which the
corrected source construction remains analytic.

\item \emph{Whitening singularities:} loci where the relevant Gram blocks lose
invertibility or the inverse-square-root branches cease to be analytic.

\item \emph{\(c=0\) strata:} these must be excluded from charts using
\[
T=\frac{v_7}{c}
\]
or any other \(c^{-1}\)-normalization.

\item \emph{\(A_5^{\mathfrak f}=0\) strata:} these are rank-jump loci for the
quotient line
\[
\mathbf C A_5^{\mathfrak f}.
\]

\item \emph{\(v_5=0\) and \(u_5=0\) patch boundaries:} these must be separated
when using the corresponding \(v_5\)-patch or \(u_5\)-patch coordinates.

\item \emph{Hidden-state rank jumps:} loci where the augmented package
coordinate
\[
T=\frac{v_7}{c}
\]
or the \(\Phi_K\)-visible lift ceases to be represented by the chosen analytic
chart.
\end{enumerate}

After separating these loci, finite odd-germ determinacy applies on each
compact nonsingular analytic stratum.
\end{remark}

\begin{corollary}[Coefficient floor away from the zero-anomaly locus]
\label{cor:coefficient-floor-away-from-zero-anomaly}
Let \(\Theta_k\) be a compact nonsingular actual-source \(k\)-core stratum, and
let \(R_{\Theta_k}\) be the finite-determinacy depth from
Theorem~\ref{thm:finite-odd-germ-determinacy-actual-source}.  Define
\[
\mathcal Z_k
=
\{\theta\in\Theta_k:H_{\core,\theta}\equiv0\}.
\]
For \(\eta>0\), set
\[
\Theta_{k,\eta}
=
\{\theta\in\Theta_k:
\operatorname{dist}(\theta,\mathcal Z_k)\ge\eta\}.
\]
Then there exists \(\gamma_\eta>0\) such that every
\(\theta\in\Theta_{k,\eta}\) satisfies
\[
\max_{0\le r\le R_{\Theta_k}}
|c_{\core,2r+1}(\theta)|
\ge
\gamma_\eta.
\]
\end{corollary}

\begin{proof}
The function
\[
D(\theta)
=
\max_{0\le r\le R_{\Theta_k}}
|c_{\core,2r+1}(\theta)|
\]
is continuous on \(\Theta_k\).  By
Theorem~\ref{thm:finite-odd-germ-determinacy-actual-source},
\[
D(\theta)=0
\quad\Longleftrightarrow\quad
\theta\in\mathcal Z_k.
\]
Thus \(D>0\) on the compact set \(\Theta_{k,\eta}\), and hence \(D\) has a
positive minimum there.
\end{proof}

\begin{remark}[Finite determinacy is not zeta-scale transfer]
\label{rem:finite-determinacy-not-zeta-scale-transfer}
Theorem~\ref{thm:finite-odd-germ-determinacy-actual-source} is qualitative.
It proves only
\[
H_{\core,\theta}\not\equiv0
\Longrightarrow
\exists N\le R_{\Theta_k}+1:
c_{\core,2N-1}(\theta)\ne0.
\]
It does not prove
\[
|\Xi_{\core}^{(N)}(m)|
\gtrsim
\frac{x}{B_\zeta(m)}S(m).
\]
Nor does it provide a sign law, a tail-domination estimate for the filtered
\(N\)-point coefficient, or an explicit bound for \(R_{\Theta_k}\) depending
only on \(k\).

The remaining quantitative problem is to compare the actual zeta core's
distance from
\[
\mathcal Z_k
\]
with the calibration scale
\[
\frac{x}{B_\zeta(m)}S(m).
\]
That is the calibrated separation problem, not a consequence of
Noetherian finite determinacy.
\end{remark}

\begin{remark}[Later refinement of calibrated separation]
\label{rem:calibrated-separation-later-refinement}
The distance formulation below is an earlier geometric form of the finite-core
quantitative gap.  The current proof-state uses the sharper sampled-vector
formulation in
Subsection~\ref{subsec:finite-core-sampled-vector-nonconcentration}.  In that
formulation the finite-determinacy vector is decomposed as
\[
C_{\rm fd}=P_{\rm fd}+J_{\rm fd},
\]
and the live hypothesis is the sampled vector lower bound
\[
\|P_{\rm fd}(\theta(m))+J_{\rm fd}(\theta(m))\|_\infty
\ge
\Lambda(m,Q).
\]
Thus the branch
\[
\mathcal N_k(m,Q)
\]
defined below should be read as a geometric precursor to the sharper
interaction-aligned branch
\[
\mathcal N_k^{\rm align}(m,Q).
\]
\end{remark}

\begin{hypothesis}[Calibrated finite-core separation]
\label{hyp:calibrated-finite-core-separation}
Let \(\Theta_k\) be a compact nonsingular actual-source \(k\)-core stratum,
with zero-anomaly locus
\[
\mathcal Z_k
=
\{\theta:H_{\core,\theta}\equiv0\}.
\]
For actual zeta finite cores \(\theta(m)\), either
\[
\theta(m)\in\mathcal Z_k,
\]
or \(\theta(m)\) lies in a controlled incidence stratum, or else
\[
\operatorname{dist}(\theta(m),\mathcal Z_k)
\ge
D_{\Theta_k}(m,Q),
\]
where \(D_{\Theta_k}(m,Q)\) is large enough to force
\[
\left|
c_{\core,2N-1}
+
\sum_{j\ge1}
c_{\core,2N-1+2j}
h_j(1^2,\ldots,N^2)Q^{-2j}
\right|
\gtrsim
\frac{Q^{2N+2}}{(2N-1)N!(N-1)!}
\frac{x}{B_\zeta(m)}S(m)
\]
for some
\[
N\le R_{\Theta_k}+1.
\]
\end{hypothesis}

\begin{definition}[Near-zero-anomaly finite-core branch]
\label{def:near-zero-anomaly-finite-core-branch}
On a compact nonsingular actual-source \(k\)-core stratum \(\Theta_k\), define
\[
\mathcal N_k(m,Q)
:=
\left\{
\theta(m)\in\Theta_k:
0<
\operatorname{dist}(\theta(m),\mathcal Z_k)
<
D_{\Theta_k}(m,Q),
\quad
\theta(m)\notin\mathcal E_{\rm geom}
\right\}.
\]
This is the branch where the corrected finite-core anomaly is nonzero but too
close to the zero-anomaly locus for the finite-determinacy--{\L}ojasiewicz
argument to imply
\[
|\Xi_{\core}^{(N)}(m)|
\gtrsim
\frac{x}{B_\zeta(m)}S(m).
\]
\end{definition}

\begin{remark}[Status of the near-zero-anomaly branch]
\label{rem:status-near-zero-anomaly-finite-core-branch}
The present draft does not prove
\[
\mathcal N_k(m,Q)=\varnothing.
\]
In particular, proximity to \(\mathcal Z_k\) is not currently shown to imply
membership in a controlled incidence stratum such as
\[
M=0,\quad W=0,\quad K_v=0,\quad T_v=0,
\]
nor collision, patch-boundary, endpoint-gap incidence, amplitude drop,
whitening singularity, \(c=0\), \(A_5^{\mathfrak f}=0\), or hidden-state rank
jump.  Thus finite-core closure requires either calibrated anti-tangency from
\(\mathcal Z_k\) or a separate analysis of \(\mathcal N_k(m,Q)\).
\end{remark}

\begin{definition}[First interaction branch]
\label{def:first-interaction-branch}
For a normalized same-scale finite core, write
\[
H_{\core}(s)
=
\Phi_K(F_{\core})(s)
=
\sum_{r\ge0}c_{\core,2r+1}s^{2r+1},
\]
and decompose
\[
F_{\core}
=
\sum_{\nu}a_\nu F_{h_\nu}+\mathcal I_{\ge2}.
\]
Define
\[
J_1:=[s]\Phi_K(\mathcal I_{\ge2}).
\]
The first interaction branch is
\[
\mathfrak Z_{\rm int}
=
\left\{
\sum_{\nu}a_\nu K(d_\nu)+J_1=0
\right\}
\setminus\mathcal E_{\rm geom},
\]
where \(\mathcal E_{\rm geom}\) denotes the union of the controlled geometric,
collision, patch-boundary, and endpoint-gap incidence strata.
\end{definition}

\begin{remark}[Interaction rigidity remains a source-level obligation]
\label{rem:interaction-rigidity-remains-source-level-obligation}
The current ECT and lifted-curve geometry controls defect-free affine
incidence.  It does not by itself control the genuinely interacting scalar
\[
J_1=[s]\Phi_K(\mathcal I_{\ge2}).
\]
The known structural properties of \(\mathcal I_{\ge2}\), such as
one-amplitude collapse and collision divisibility, force \(J_1\) to vanish on
some boundary strata, but they do not force
\[
J_1\equiv0,
\]
nor do they impose a sign condition or place \(J_1\) in the radical of the
controlled incidence ideal.

Thus the finite-core proof needs an additional interaction-rigidity theorem:
\(J_1\) must be shown to be zero, sign-controlled, or governed by the
controlled incidence equations.  Without such a theorem,
\[
\mathfrak Z_{\rm int}
\]
remains a separate finite-core branch.
\end{remark}

\begin{remark}[Source-level first-interaction rigidity remains open]
\label{rem:source-level-first-interaction-rigidity-open}
At two-source order, the first interaction coefficient has the form
\[
\mathfrak j_{ij}^{(2)}
=
[s]\Phi_K
\left(
D\mathcal W_0[\mathcal B_{ij}^{(2)}]
+
D^2\mathcal W_0[\mathcal B_i^{(1)},\mathcal B_j^{(1)}]
\right).
\]
The direct Hessian contribution and the bilinear whitening contribution split
into the independent channels
\[
dd,\qquad d\eta+\eta d,\qquad \eta\eta,
\]
together with bilinear whitening products.  No current parity,
endpoint-symmetry, ECT, or controlled-incidence theorem forces their sum to
vanish, gives it a calibrated sign/cone condition, or places it in the radical
of the controlled incidence ideal.

Thus
\[
J_1=[s]\Phi_K(\mathcal I_{\ge2})
\]
remains an independent finite-core obstruction unless a separate
source-level first-interaction rigidity theorem is proved.
\end{remark}

\begin{remark}[Scope after Gate B]
\label{rem:scope-after-gate-b}
Gate~B supplies the raw \(\delta^2\)-edge divisibility needed to construct
the order-\(7\) edge-normalized corrected package.  It does not by itself
prove the final quotient-gauge theorem, diagonal merger,
\(\kappa\)-independence, the connection cancellations for \(M_A,M_C\), or
the package-to-\(\Xi_\zeta^{(N)}\) bridge.

The \(M_A,M_C\) connection branch is now closed under the source-grade
reflection convention by
Theorem~\ref{thm:ma-mc-pretriv-connection-closure-final}.  Thus the remaining
post-Gate-B local obstruction is not the \(Z_j[M_A],Z_j[M_C]\) table.  The
remaining local actual-source/Bottleneck-C obstruction is the promoted
quadratic mixed-block coefficient
\[
E_{A,Y}^{\mathrm{BC}}
\]
of Definition~\ref{def:promoted-quadratic-mixed-block-bc-coefficient}.  Its
quotient test is
\[
e_{A,Y}v_5-u_5f_{A,Y}=0,
\]
and its transform-safe test is
\[
f_{A,Y}=0.
\]
\end{remark}

\begin{remark}[Reduced package equality does not imply hidden-state equality]
\label{rem:reduced-package-equality-not-hidden-state-equality}
The reduced datum \(\widehat\Psi\) forgets the raw septic representative in
the affine line
\[
A_7^{\mathfrak f}\longmapsto A_7^{\mathfrak f}+\chi A_5^{\mathfrak f}.
\]
This forgotten direction is not automatically invisible to the transverse
functional.  Writing
\[
A_5^{\mathfrak f}=u_5 I+v_5 S,
\]
one has
\[
\Phi_K(A_5^{\mathfrak f})=2v_5.
\]
Thus reduced-package equality alone cannot imply agreement modulo
\(\ker\Phi_K\) through the first surviving odd order, except on the special
locus \(v_5=0\) or after an additional theorem showing cancellation in the
actual corrected block.  The missing hidden-state theorem is therefore a
transform-level state-locality statement modulo \(\ker\Phi_K\), not a
formal consequence of quotient-septic closure.
\end{remark}

\begin{definition}[\(\Phi_K\)-visible septic state]
\label{def:phik-visible-septic-state}
On the fixed target fiber
\[
\mathfrak f_m=\mathbf C I\oplus \mathbf C S,
\]
write
\[
A_7^{\mathfrak f}=u_7I+v_7S.
\]
The \(\Phi_K\)-visible septic state is the class
\[
[A_7^{\mathfrak f}]_{\Phi}:=
A_7^{\mathfrak f}\bmod \ker(\Phi_K|_{\mathfrak f_m}),
\]
equivalently the scalar
\[
\phi_7:=\Phi_K(A_7^{\mathfrak f})=2v_7.
\]
If the \(c\)-normalized package is being used, the normalized visible septic
state is
\[
\theta_7:=\frac{\phi_7}{c}=\frac{2v_7}{c}.
\]
\end{definition}

\begin{proposition}[Formal non-descent of the \(\Phi_K\)-visible septic state]
\label{prop:formal-non-descent-phik-visible-septic-state}
On the generic locus \(v_5\neq0\), the quotient class
\[
[A_7^{\mathfrak f}]
\in
\mathfrak f_m/\mathbf C A_5^{\mathfrak f}(m)
\]
does not determine the \(\Phi_K\)-visible septic state
\[
[A_7^{\mathfrak f}]_{\Phi}
\in
\mathfrak f_m/\ker(\Phi_K|_{\mathfrak f_m}).
\]
In particular, quotient package equality does not imply equality of the first
\(\Phi_K\)-visible corrected odd coefficient.
\end{proposition}

\begin{proof}
Write
\[
A_5^{\mathfrak f}=u_5I+v_5S,
\qquad
A_7^{\mathfrak f}=u_7I+v_7S.
\]
The quotient identifies
\[
A_7^{\mathfrak f}
\sim
A_7^{\mathfrak f}+\chi A_5^{\mathfrak f}.
\]
However
\[
\Phi_K(A_7^{\mathfrak f}+\chi A_5^{\mathfrak f})
-
\Phi_K(A_7^{\mathfrak f})
=
\chi\Phi_K(A_5^{\mathfrak f})
=
2\chi v_5.
\]
Thus, if \(v_5\neq0\), the forgotten quotient direction is visible to
\(\Phi_K\).

Equivalently, take the formal model
\[
A_5^{\mathfrak f}=S,
\qquad
A_7^{\mathfrak f,\chi}=\chi S.
\]
Then
\[
[A_7^{\mathfrak f,\chi}]=0
\quad\text{in}\quad
\mathfrak f_m/\mathbf C S
\]
for every \(\chi\), but
\[
\Phi_K(A_7^{\mathfrak f,\chi})=2\chi.
\]
Hence the reduced quotient package does not determine the
\(\Phi_K\)-visible septic state.
\end{proof}

\begin{remark}[Package-to-transform bridge after augmentation]
\label{rem:package-to-transform-bridge-after-augmentation}
The quotient-only package-to-transform bridge is impossible on the generic
patch \(v_5\ne0\), since
\[
A_7^{\mathfrak f}\sim A_7^{\mathfrak f}+\chi A_5^{\mathfrak f}
\]
preserves the quotient class but changes
\[
\Phi_K(A_7^{\mathfrak f})
\]
by
\[
2\chi v_5.
\]
The correct theorem object is therefore the transform-safe augmented package
\[
\widehat\Psi_{\rm red}^{+}
=
(\widehat\Psi_{\rm red},T),
\qquad
T=\frac{v_7}{c}.
\]
With this augmentation, the first \(\Phi_K\)-visible septic coefficient is
read off as
\[
\Phi_K(A_7^{\mathfrak f})=2cT.
\]
The remaining compatibility problem is not quotient-only state-locality; it is
the hidden-horizontal diagonal-merger condition
\[
\partial_\kappa T(m,\kappa,0)=0.
\]
\end{remark}

\begin{definition}[Transform-safe augmented reduced package]
\label{def:transform-safe-augmented-reduced-package}
Work on a good patch with
\[
c\ne0,
\qquad
A_5^{\mathfrak f}=u_5I+v_5S.
\]
Write
\[
A_7^{\mathfrak f}=u_7I+v_7S,
\qquad
\Delta_7:=u_7v_5-u_5v_7.
\]
The transform-safe augmented reduced package is
\[
\widehat\Psi_{\rm red}^{+}
:=
(\widehat\Psi_{\rm red},T),
\qquad
T:=\frac{v_7}{c}.
\]
Equivalently, on the \(v_5\)-patch one may write
\[
\widehat\Psi_{\rm red}^{+}
=
(c,u_5,v_5,\Delta_7,T).
\]
The normalized \(\Phi_K\)-visible coordinate is
\[
\theta_{\Phi,7}
:=
\frac{\Phi_K(A_7^{\mathfrak f})}{c}
=
2T.
\]
\end{definition}

\begin{theorem}[Augmented package-to-transform readout]
\label{thm:augmented-package-to-transform-readout}
Let
\[
\widehat\Psi_{\rm red}^{+}
=
(c,u_5,v_5,\Delta_7,T)
\]
be the transform-safe augmented package on a good patch with \(c\ne0\). Then
the first \(\Phi_K\)-visible septic coefficient of the corrected finite-\(s\)
two-point transform is determined by the augmented package:
\[
\Phi_K(A_7^{\mathfrak f})
=
2v_7
=
2cT.
\]
Equivalently,
\[
\frac{1}{c}
[s^7]\Phi_K(\widehat\Omega^{\corr}(s;m))
=
2T.
\]
Thus the two-point transform-visible septic state factors through
\(\widehat\Psi_{\rm red}^{+}\), not through \(\widehat\Psi_{\rm red}\) alone.
\end{theorem}

\begin{remark}[Why the old quotient-only bridge must be abandoned]
\label{rem:old-quotient-only-bridge-abandoned}
The reduced quotient package remembers only
\[
[A_7^{\mathfrak f}]
\in
\mathfrak f_m/\mathbf C A_5^{\mathfrak f}(m).
\]
It therefore identifies
\[
A_7^{\mathfrak f}
\sim
A_7^{\mathfrak f}+\chi A_5^{\mathfrak f}.
\]
But
\[
\Phi_K(A_5^{\mathfrak f})=2v_5.
\]
Hence, on the generic patch \(v_5\ne0\), the quotient-forgotten line direction
is \(\Phi_K\)-visible:
\[
\Phi_K(A_7^{\mathfrak f}+\chi A_5^{\mathfrak f})
-
\Phi_K(A_7^{\mathfrak f})
=
2\chi v_5.
\]
Therefore no functional of \(\widehat\Psi_{\rm red}\) alone can determine
\(\Phi_K(A_7^{\mathfrak f})\) on the generic patch.  The missing coordinate is
exactly
\[
T=\frac{v_7}{c}.
\]
\end{remark}

\begin{theorem}[Sharp no-descent criterion for the \(\Phi_K\)-visible septic state]
\label{thm:sharp-no-descent-transform-visible-septic}
Let
\[
\mathfrak f_m=\mathbf C I\oplus\mathbf C S,
\qquad
A_5^{\mathfrak f}=u_5I+v_5S.
\]
Let
\[
q:\mathfrak f_m\to \mathfrak f_m/\mathbf C A_5^{\mathfrak f}
\]
be the quotient map.  There exists a functional
\[
\Theta_{\Phi,7}
\]
on the quotient such that
\[
\Phi_K(A_7^{\mathfrak f})
=
\Theta_{\Phi,7}(q(A_7^{\mathfrak f}))
\]
for all \(A_7^{\mathfrak f}\in\mathfrak f_m\) if and only if
\[
\Phi_K(A_5^{\mathfrak f})=0.
\]
Equivalently, since
\[
\Phi_K(A_5^{\mathfrak f})=2v_5,
\]
such descent is possible if and only if
\[
v_5=0.
\]
\end{theorem}

\begin{proof}
The quotient identifies
\[
A_7^{\mathfrak f}
\sim
A_7^{\mathfrak f}+\chi A_5^{\mathfrak f}.
\]
Therefore \(\Phi_K\) descends to the quotient if and only if it is constant on
these affine fibers, i.e.
\[
\Phi_K(A_7^{\mathfrak f}+\chi A_5^{\mathfrak f})
=
\Phi_K(A_7^{\mathfrak f})
\]
for all \(A_7^{\mathfrak f}\) and all \(\chi\).  By linearity this is
equivalent to
\[
\chi\Phi_K(A_5^{\mathfrak f})=0
\]
for all \(\chi\), hence to
\[
\Phi_K(A_5^{\mathfrak f})=0.
\]
Writing
\[
A_5^{\mathfrak f}=u_5I+v_5S
\]
gives
\[
\Phi_K(A_5^{\mathfrak f})=2v_5.
\]
Thus descent holds exactly on the locus \(v_5=0\).
\end{proof}

\begin{theorem}[Augmented Bottleneck C condition]
\label{thm:augmented-bottleneck-c-condition}
On the good patch \(c\ne0\), diagonal merger for the transform-safe package
\[
\widehat\Psi_{\rm red}^{+}
=
(c,u_5,v_5,\Delta_7,T)
\]
requires both
\[
\partial_\kappa Q_{7,m}^q(m,\kappa,0)=0
\]
and
\[
\partial_\kappa T(m,\kappa,0)=0.
\]
The first condition is quotient-horizontal merger.  The second is
hidden-horizontal merger.  On the generic patch \(v_5\ne0\), the first does
not imply the second.
\end{theorem}

\begin{remark}[Scalar Bottleneck C condition]
\label{rem:scalar-bottleneck-c-condition}
Let
\[
P_7(a_1,a_2)
=
[z^7][\delta^1]\tau_{h\to m}\pi_{\mathfrak f_h}
(1-\Pi_{1,1})_{\rm mix}C_{\mathcal W}^{\rm raw}.
\]
The quotient Bottleneck C condition is
\[
\partial_{a_2}P_7(a_1,-a_1)
\in
\mathbf C A_5^{\mathfrak f}(m).
\]
The transform-safe condition is instead
\[
\Phi_K\!\left(\partial_{a_2}P_7(a_1,-a_1)\right)=0.
\]
Equivalently, if
\[
P_{7,\Phi}(a_1,a_2)
:=
\frac{1}{c}\Phi_K(P_7(a_1,a_2)),
\]
then the scalar condition is
\[
P_{7,\Phi}(a_1,a_2)
\in
(a_1+a_2)^2\mathbf C[a_1,a_2].
\]
This scalar condition is independent of quotient-valued double-balanced
divisibility on the generic patch \(v_5\ne0\).
\end{remark}

\begin{lemma}[Hidden-horizontal exceptional divisor condition]
\label{lem:hidden-horizontal-exceptional-divisor-condition}
Work on a good patch with
\[
c(m)\ne0,\qquad
A_5^{\mathfrak f}(m)=u_5I+v_5S,\qquad
v_5\ne0.
\]
Let
\[
F_7(a_1,a_2,\delta)
=
[z^7]\tau_{h\to m}\pi_{\mathfrak f_h}
(1-\Pi_{1,1})_{\rm mix}C_{\mathcal W}^{\rm raw}
=
\delta P_7(a_1,a_2)+O(\delta^2).
\]
Then the augmented hidden coordinate \(T=v_7/c\) is horizontal on the
exceptional divisor if and only if
\[
\Phi_K\!\left(
\partial_{a_2}P_7(a_1,-a_1)
\right)=0.
\]
Equivalently,
\[
\partial_\kappa T(m,\kappa,0)=0
\quad\Longleftrightarrow\quad
\frac1c
\Phi_K\!\left(
\partial_{a_2}P_7(a_1,-a_1)
\right)=0.
\]
\end{lemma}

\begin{proof}
Along the collision chart
\[
a_2=-a_1(1+\kappa\delta),
\]
one has
\[
P_7(a_1,a_2)
=
P_7(a_1,-a_1)
-
a_1\kappa\delta\,\partial_{a_2}P_7(a_1,-a_1)
+
O(\delta^2).
\]
After the standard edge normalization, the exceptional \(\kappa\)-slope is
represented by
\[
\partial_{a_2}P_7(a_1,-a_1).
\]
Applying \(\Phi_K/c\) gives the slope of the normalized transform-visible
coordinate.  Since \(T=\Phi_K(A_7^{\mathfrak f})/(2c)\), the displayed
criterion follows.
\end{proof}

\begin{proposition}[Formal transform-visible line-gauge leakage]
\label{prop:formal-transform-visible-line-gauge-leakage}
On a good patch with
\[
c\ne0,
\qquad
A_5^{\mathfrak f}=u_5I+v_5S,
\qquad
v_5\ne0,
\]
quotient-horizontal Bottleneck~C does not imply \(T\)-horizontal merger.
\end{proposition}

\begin{proof}
Consider an order-\(7\) fixed-sector edge carrier of the form
\[
[\delta^1]A_7^{\mathfrak f,\chi}
=
\lambda(m)\,a_1a_2(a_1+a_2)\,A_5^{\mathfrak f}(m).
\]
Since it is pointwise parallel to \(A_5^{\mathfrak f}(m)\), its quotient class
vanishes:
\[
[A_7^{\mathfrak f,\chi}]=0
\qquad
\text{in}
\qquad
\mathfrak f_m/\mathbf C A_5^{\mathfrak f}(m).
\]
Thus it is invisible to quotient Bottleneck~C.

However,
\[
a_1a_2(a_1+a_2)\big|_{a_2=-a_1}=0,
\]
but
\[
\partial_{a_2}\bigl(a_1a_2(a_1+a_2)\bigr)\big|_{a_2=-a_1}
=
-a_1^2.
\]
Moreover,
\[
\Phi_K(A_5^{\mathfrak f})=2v_5.
\]
Therefore the induced hidden-state slope is
\[
[\delta^1]\partial_{a_2}T_\chi(a_1,-a_1,\delta)
=
-a_1^2\lambda(m)\frac{v_5(m)}{c(m)}.
\]
On the generic patch \(c\ne0,\ v_5\ne0\), this is nonzero unless
\[
\lambda(m)=0.
\]
Thus a quotient-invisible line-gauge edge carrier can be transform-visible.
\end{proof}

\begin{corollary}[First live \(T\)-slope carrier]
\label{cor:first-live-t-slope-carrier}
Let
\[
E_{A,Y}^{\mathrm{BC}}
=
e_{A,Y}I+f_{A,Y}S
\]
be the frozen Bottleneck-C coefficient of the promoted nonlinear mixed-block
carrier.  The amplitude factor of this carrier is
\[
a_1a_2(a_1+a_2),
\]
and
\[
\partial_{a_2}\bigl(a_1a_2(a_1+a_2)\bigr)\big|_{a_2=-a_1}
=
-a_1^2.
\]
Consequently its contribution to the hidden-horizontal slope is
\[
\partial_\kappa T\big|_{E_{A,Y}^{\mathrm{BC}}}
=
-a_1^2\frac{\Phi_K(E_{A,Y}^{\mathrm{BC}})}{2c}
=
-a_1^2\frac{f_{A,Y}}{c}.
\]
Thus the augmented transform-safe package requires
\[
f_{A,Y}=0.
\]
In particular, if
\[
E_{A,Y}^{\mathrm{BC}}=\lambda A_5^{\mathfrak f},
\]
then
\[
\partial_\kappa T\big|_{E_{A,Y}^{\mathrm{BC}}}
=
-a_1^2\lambda\frac{v_5}{c}.
\]
On the generic \(v_5\ne0\) patch, a nonzero line-gauge
\(E_{A,Y}^{\mathrm{BC}}\) closes the quotient condition but still leaks
through \(T\).
\end{corollary}

\begin{remark}[Local diagnostic for the hidden septic state]
\label{rem:local-diagnostic-hidden-septic-state}
On the quintic-good patch
\[
\{c\ne0,\ v_5\ne0\},
\]
one useful local diagnostic enlargement of the reduced package is
\[
\widetilde\Psi_{\min}=(x,Y,S,T)
=
\left(
\frac{v_5}{c},
\frac{u_5}{c},
\frac{\Delta_7}{cv_5},
\frac{v_7}{c}
\right).
\]
This tuple is useful only as normalized fixed-sector septic data.  It is
not a canonical global theorem object and it does not give finite
determination of the first surviving odd jet.  The global target remains
the transform-level first \(\Phi_K\)-visible corrected odd coefficient,
modulo \(\ker\Phi_K\).
\end{remark}

\begin{remark}[Bottleneck C is exceptional-divisor fiber selection]
\label{rem:bottleneck-c-exceptional-divisor-fiber-selection}
For Bottleneck C, the obstruction is not the choice of a local septic
representative such as \(v_7=0\), \(u_7=0\), or \(T=v_7/c\).  Those are
diagnostic data for the hidden-state problem.  The C-side object has
already quotiented raw septic representatives away.  Its remaining
obstruction is the exceptional-divisor trace
\[
B(m,\kappa):=\widetilde\Psi^{\corr}_{\mathrm{red}}(m,\kappa,0).
\]
Analyticity, swap-evenness, fixed codomain, scalar normalization, and the
central value \(B(m,0)=\widehat\Psi(m)\) do not force
\[
\partial_\kappa B=0.
\]
A valid C proof must supply actual diagonal merger, same-reduced-image, or
collision-functoriality for the corrected two-atom package.
\end{remark}

\begin{remark}[Formal C properties do not rule out slope leakage]
\label{rem:formal-c-properties-do-not-rule-out-slope-leakage}
Formal properties such as analyticity, swap-evenness, fixed codomain,
scalar normalization, and the central value at \(\kappa=0\) do not by
themselves force the C-side exceptional-divisor trace to be independent
of \(\kappa\).  A formal model may have
\[
B_3=Z+\varepsilon\kappa,
\]
or, at the quotient-defect level,
\[
R\longmapsto R+R_\kappa,
\qquad
\det(R_\kappa,A_5^{\mathfrak f})=c^2\varepsilon\kappa.
\]
Thus representative-gauge invariance only kills additions parallel to
\(A_5^{\mathfrak f}\).  It does not kill determinant-changing
quotient-defect data.  A C proof must control the actual corrected
exceptional-divisor quotient class, not merely the formal shape of the
reduced package.
\end{remark}

\begin{remark}[Atom swap does not impose \(\kappa\mapsto-\kappa\) on the exceptional divisor]
\label{rem:atom-swap-does-not-kill-kappa-slope}
In the collision chart
\[
h_1=m-\frac{\delta}{2},
\qquad
h_2=m+\frac{\delta}{2},
\qquad
a_2=-a_1(1+\kappa\delta),
\]
interchanging the two labelled atoms gives
\[
\delta\longmapsto-\delta,
\qquad
a_1\longmapsto -a_1(1+\kappa\delta),
\qquad
\kappa\longmapsto \frac{\kappa}{1+\kappa\delta}.
\]
Hence on the exceptional divisor \(\delta=0\), the slope coordinate is fixed:
\[
\kappa\longmapsto\kappa.
\]
Therefore atom-swap symmetry does not force the exceptional trace
\[
B(m,\kappa)
\]
to be even in \(\kappa\), and in particular it does not force
\[
\partial_\kappa B(m,\kappa)=0.
\]
Swap symmetry removes the wrong-parity terms in the transverse collision
variable; it does not remove quotient-visible dependence on the exceptional
slope.
\end{remark}

\begin{remark}[Missing theorem: quotient-valued diagonal merger]
\label{rem:missing-quotient-valued-diagonal-merger}
Bottleneck~C requires a quotient-valued diagonal merger theorem.  On every
good patch on which
\[
c(m)\ne0,
\qquad
A_5^{\mathfrak f}(m)\ne0,
\]
the corrected two-atom collision package must factor through the blow-down of
the exceptional divisor:
\[
(m,\kappa,0)\longmapsto m.
\]
Equivalently,
\[
B(m,\kappa)
=
B(m,0)
=
\widehat\Psi(m)
\qquad
\text{for all sufficiently small }\kappa.
\]
At the order-\(7\) quotient level this means
\[
B_7^{\mathfrak f}(m,\kappa,0)
-
B_7^{\mathfrak f}(m,0,0)
\in
\mathbf C A_5^{\mathfrak f}(m).
\]

This is stronger than choosing a local septic gauge such as \(v_7=0\) or
\(u_7=0\).  Such a gauge selects a representative inside a fixed quotient
class; it does not prove that the quotient class itself is independent of the
exceptional slope \(\kappa\).
\end{remark}

\begin{definition}[Quotient-valued two-atom collision functor]
\label{def:quotient-valued-two-atom-collision-functor}
On a good patch with
\[
A_5^{\mathfrak f}(m)\ne0,
\]
define the order-\(7\) quotient-valued two-atom collision functor by the
ordered composite
\[
\mathcal B_2
\longmapsto
C_{\mathcal W}^{\rm raw}
\longmapsto
(1-\Pi_{1,1})_{\rm mix}C_{\mathcal W}^{\rm raw}
\longmapsto
-a_1^{-1}\delta^{-2}
(1-\Pi_{1,1})_{\rm mix}C_{\mathcal W}^{\rm raw}
\]
\[
\longmapsto
\pi_{\mathfrak f_h}[z^{\le7}]
\longmapsto
\tau_{h\to m}
\longmapsto
[z^7]
\longmapsto
\mathfrak f_m/\mathbf C A_5^{\mathfrak f}(m).
\]
In formulas,
\[
\mathcal C_7^q(m,\kappa,\delta)
:=
\left[
[z^7]\,
\tau_{h\to m}
\pi_{\mathfrak f_h}
\left(
-a_1^{-1}\delta^{-2}
(1-\Pi_{1,1})_{\rm mix}
C_{\mathcal W}^{\rm raw}
\right)
\right].
\]
The exceptional-divisor trace is
\[
B(m,\kappa)=\mathcal C_7^q(m,\kappa,0).
\]
Thus quotient-valued diagonal merger is precisely the statement that
\[
\mathcal C_7^q(m,\kappa,0)
=
\mathcal C_7^q(m,0,0),
\]
or equivalently that \(\mathcal C_7^q\) factors on the exceptional divisor
through the blow-down map
\[
(m,\kappa,0)\longmapsto m.
\]
\end{definition}

\begin{proposition}[Formal quotient-visible exceptional-slope leakage]
\label{prop:formal-quotient-visible-exceptional-slope-leakage}
The formal properties
\[
\text{analyticity},\qquad
\text{atom-swap compatibility},\qquad
\text{fixed codomain},\qquad
\text{scalar normalization},\qquad
B(m,0)=\widehat\Psi(m)
\]
do not imply quotient-valued diagonal merger.
\end{proposition}

\begin{proof}
Work on a good patch and choose
\[
E_m\in\mathfrak f_m
\]
with
\[
\omega_m(E_m,A_5^{\mathfrak f}(m))=1.
\]
Thus \(E_m\) is nonzero in
\[
\mathfrak f_m/\mathbf C A_5^{\mathfrak f}(m).
\]

In the collision chart
\[
h_1=m-\frac{\delta}{2},
\qquad
h_2=m+\frac{\delta}{2},
\qquad
a_2=-a_1(1+\kappa\delta),
\]
atom swap acts by
\[
\delta\longmapsto-\delta,
\qquad
\kappa\longmapsto\frac{\kappa}{1+\kappa\delta}.
\]
Define
\[
H(\kappa,\delta)
:=
\frac12
\left(
\kappa+\frac{\kappa}{1+\kappa\delta}
\right).
\]
Then \(H\) is analytic near \((\kappa,\delta)=(0,0)\), satisfies
\[
H\!\left(\frac{\kappa}{1+\kappa\delta},-\delta\right)
=
H(\kappa,\delta),
\]
and restricts to
\[
H(\kappa,0)=\kappa
\]
on the exceptional divisor.

Let \(Z_0(m)\in\mathfrak f_m\) be any representative of
\(\widehat\Psi(m)\), and define the formal normalized septic edge
representative
\[
\mathcal R_\varepsilon(m,\kappa,\delta;z)
=
z^7
\left(
Z_0(m)+\varepsilon H(\kappa,\delta)E_m
\right).
\]
Equivalently, before edge normalization one may write the raw model as
\[
\mathcal N_\varepsilon
=
-a_1\delta^2z^7
\left(
Z_0(m)+\varepsilon H(\kappa,\delta)E_m
\right).
\]
This raw model is \(O(a_1\delta^2)\), so it is compatible with raw
edge-divisibility.  It is analytic and atom-swap compatible by construction.

At \(\kappa=0\),
\[
B_\varepsilon(m,0)=[Z_0(m)]=\widehat\Psi(m).
\]
But on the exceptional divisor,
\[
B_\varepsilon(m,\kappa)
=
[Z_0(m)+\varepsilon\kappa E_m],
\]
and hence
\[
\partial_\kappa B_\varepsilon(m,\kappa)
=
\varepsilon [E_m]\ne0
\]
in
\[
\mathfrak f_m/\mathbf C A_5^{\mathfrak f}(m).
\]
Therefore the listed formal properties do not force quotient-valued diagonal
merger.
\end{proof}

\begin{remark}[Balancing does not impose hidden \(\kappa\)-horizontality]
\label{rem:balancing-does-not-impose-kappa-horizontality}
The balancing condition
\[
a_2=-a_1(1+\kappa\delta)
\]
gives
\[
\partial_\kappa a_2=-a_1\delta.
\]
This single factor of \(\delta\) does not force the normalized exceptional
slope to vanish after the \(\delta^{-2}\) edge normalization.  A raw term of
the form
\[
-a_1\delta^2\kappa E_m
\]
is still of the allowed Gate-B edge order and survives normalization as
\[
\kappa E_m.
\]
Equivalently, with \(\kappa=2\omega/\delta\), the leakage term
\[
2\varepsilon\omega\delta E_m
=
\varepsilon\kappa\delta^2E_m
\]
is invisible before edge normalization but visible on the exceptional divisor.
Thus balancing kills the zeroth source moment; it does not by itself kill the
projective exceptional slope.
\end{remark}

\begin{definition}[Promoted quadratic mixed-block Bottleneck-C coefficient]
\label{def:promoted-quadratic-mixed-block-bc-coefficient}
Let
\[
Y_{12}^{(2)}:=D^2N_0[\xi_1,\xi_2]
\]
denote the \(\mathcal K^2\), source-bidegree \((1,1)\) nonlinear Taylor
remainder of the actual corrected mixed block
\(N^{\corr}(d_-,d_+,\eta)\).

The promoted channel has amplitude factor
\[
a_1a_2(a_1+a_2).
\]
On the collision chart
\[
a_2=-a_1(1+\kappa\delta),
\]
this amplitude already contributes the first \(\delta\)-edge:
\[
a_1a_2(a_1+a_2)
=
a_1^3\kappa\delta+O(\delta^2).
\]
Therefore the Bottleneck-C coefficient for this promoted channel is the
frozen internal coefficient
\[
E_{A,Y}^{\mathrm{BC}}
:=
\tau_{h\to m}\pi_{\mathfrak f_h}
\operatorname{Sym}
[z^7]\,
D(G_-^{-1/2})[A_-]\,
Y_{12}^{(2)}\,
R_0.
\]
Equivalently,
\[
E_{A,Y}^{\mathrm{BC}}
=
\tau_{h\to m}\pi_{\mathfrak f_h}
\operatorname{Sym}
\sum_{\chi\in\{dd,d\eta+\eta d,\eta\eta\}}
\sum_{g=4}^{7}
\sum_{r+b=7-g}
H^A_{r,0}Y^\chi_{g,0}W_{+,b}.
\]
Write
\[
E_{A,Y}^{\mathrm{BC}}
=
e_{A,Y}I+f_{A,Y}S.
\]
The quotient Bottleneck-C obstruction is
\[
Q_{A,Y}:=e_{A,Y}v_5-u_5f_{A,Y}.
\]
The transform-safe obstruction is
\[
\Phi_K(E_{A,Y}^{\mathrm{BC}})=2f_{A,Y}.
\]

For comparison, the literal source-position edge
\[
E_{A,Y}^{\mathrm{pos}}
:=
\tau_{h\to m}\pi_{\mathfrak f_h}
\operatorname{Sym}
[z^7\delta^1]\,
D(G_-^{-1/2})[A_-]\,
Y_{12}^{(2)}\,
R_0
\]
belongs to the next source-motion layer and should not be confused with the
first Bottleneck-C amplitude-derivative coefficient.
\end{definition}

\begin{definition}[Active Bottleneck-C source-linear index set]
\label{def:active-bc-source-linear-index-set}
Let
\[
\mathcal A_{\mathrm{BC}}
\subseteq\{2,3,4,5,6\}
\]
denote the set of same-point source-linear Fr\'echet primitive indices that
can enter the promoted mixed-block Bottleneck-C coefficient
\[
E_{A,Y}^{\mathrm{BC}}.
\]
In the finite Fr\'echet convention of
Section~\ref{sec:finite-frechet-coefficient-calculus}, the lead top-row
obstruction uses
\[
a=2,
\qquad
\Lambda^-_{2,0}
=
-\frac{1}{24\pi q^3}E_{22},
\]
so \(2\in\mathcal A_{\mathrm{BC}}\).  Older shorthand using
\(a\in\{1,5\}\) should be read only as a deprecated table convention and not
as the active finite Fr\'echet index set.
\end{definition}

\begin{remark}[Executable table for \(E_{A,Y}^{\mathrm{BC}}\)]
\label{rem:executable-table-eay-bc}
The frozen Bottleneck-C coefficient is computed from the rows
\[
\mathcal E^a_{\chi;g,r,b}
:=
\tau_{h\to m}\pi_{\mathfrak f_h}
\operatorname{Sym}
\left(
\Lambda^-_{a,r}
Y^\chi_{g,0}
W_{+,b}
\right),
\]
where
\[
a\in\mathcal A_{\mathrm{BC}},
\qquad
\chi\in\{dd,d\eta+\eta d,\eta\eta\},
\]
\[
4\le g\le7,\qquad r+b=7-g.
\]
Thus
\[
E_{A,Y}^{\mathrm{BC}}
=
\sum_{\chi\in\{dd,d\eta+\eta d,\eta\eta\}}
\sum_{a\in\mathcal A_{\mathrm{BC}}}
\sum_{\substack{4\le g\le7\\ r+b=7-g}}
\mathcal E^a_{\chi;g,r,b}.
\]
The skew row
\[
Y_{d\eta-\eta d}^{(2)}
\]
is excluded because it is killed by fixed-sector projection.
\end{remark}

\begin{remark}[Status of \(E_{A,Y}^{\mathrm{BC}}\)]
\label{rem:status-eay-bc}
The recurrences for
\[
W_{+,0},W_{+,1},W_{+,2},W_{+,3}
\]
and
\[
\Lambda^-_{a,0},\Lambda^-_{a,1},\Lambda^-_{a,2},\Lambda^-_{a,3}
\]
reduce \(E_{A,Y}^{\mathrm{BC}}\) to a finite symbolic coefficient table.
However, the present draft does not yet supply the scalar midpoint source
jets needed to evaluate
\[
E_{A,Y}^{\mathrm{BC}}=e_{A,Y}I+f_{A,Y}S.
\]
The missing scalar jets are
\[
q_0^{(k)}(m),\qquad 0\le k\le9,
\]
and
\[
(r_i^{[a]})^{(k)}(m),
\qquad
i=1,2,\quad a\in\mathcal A_{\mathrm{BC}},\quad 2\le k\le9.
\]
Moreover, the active top row
\[
(a,r,g,b)=(2,0,7,0),\qquad \chi=\eta\eta_{\odd},
\]
recorded in
Subsection~\ref{subsec:concrete-odd-etaeta-top-row-obstruction}, gives a
concrete \(\Phi_K\)-visible survivor unless
\[
\nu_{21,7}=0
\]
or unless it is cancelled by another physical top-row contribution.  Thus the
current local obstruction is sharper than an unevaluated table: the package
must prove cancellation of the odd-\(\eta\eta\) top row, or add a physical
correction cancelling it.
\end{remark}

\begin{remark}[First live promoted nonlinear mixed-block obstruction]
\label{rem:first-live-promoted-y12-obstruction}
The same-point cubic \(q^3\) defect is fixed-sector invisible in same-point-only
Fr\'echet responses by endpoint-swap/transpose symmetry.  The first remaining
actual-source/Bottleneck-C carrier is the promoted nonlinear mixed-block
quadratic remainder
\[
Y_{12}^{(2)}=D^2N_0[\xi_1,\xi_2],
\]
coupled to one same-point linear whitening slot.

The amplitude factor of this channel is
\[
a_1a_2(a_1+a_2).
\]
Since
\[
a_1a_2(a_1+a_2)
=
a_1^3\kappa\delta+O(\delta^2)
\quad
\text{on }a_2=-a_1(1+\kappa\delta),
\]
the first Bottleneck-C coefficient is
\[
E_{A,Y}^{\mathrm{BC}},
\]
not the internal source-position coefficient \(E_{A,Y}^{\mathrm{pos}}\).

Thus quotient-valued Bottleneck C for this channel is exactly
\[
Q_{A,Y}=e_{A,Y}v_5-u_5f_{A,Y}=0.
\]
The augmented transform-safe package requires the stronger scalar condition
\[
f_{A,Y}=0,
\qquad
\text{equivalently}\qquad
\Phi_K(E_{A,Y}^{\mathrm{BC}})=0.
\]
On the generic \(v_5\ne0\) patch, quotient closure together with
\(T\)-horizontal closure forces
\[
E_{A,Y}^{\mathrm{BC}}=0.
\]
\end{remark}

\begin{lemma}[Promoted mixed-block coefficient table]
\label{lem:promoted-y12-coefficient-table}
In the notation of
Definition~\ref{def:promoted-quadratic-mixed-block-bc-coefficient}, the frozen
Bottleneck-C coefficient of the promoted quadratic mixed-block channel is
\[
E_{A,Y}^{\mathrm{BC}}
=
\tau_{h\to m}\pi_{\mathfrak f_h}
\operatorname{Sym}
\sum_{\chi\in\{dd,d\eta+\eta d,\eta\eta\}}
\sum_{g=4}^{7}
\sum_{r+b=7-g}
H^A_{r,0}Y^\chi_{g,0}W_{+,b}.
\]
Here
\[
H^A_{r,0}
=
[z^r]D(G_-^{-1/2})[A_-],
\qquad
W_{+,b}=[z^b]G_+^{-1/2},
\]
and
\[
Y^\chi_{g,0}=[z^g]Y_\chi^{(2)}
\]
for
\[
\chi\in\{dd,d\eta+\eta d,\eta\eta\}.
\]
The skew combination
\[
Y_{d\eta-\eta d}^{(2)}
\]
is not included, since it is killed by fixed-sector projection.

If
\[
E_{A,Y}^{\mathrm{BC}}=e_{A,Y}I+f_{A,Y}S,
\]
then the quotient obstruction is
\[
e_{A,Y}v_5-u_5f_{A,Y},
\]
and the transform-safe obstruction is
\[
f_{A,Y}.
\]
\end{lemma}

\begin{remark}[Current status of the promoted mixed-block coefficient table]
\label{rem:status-promoted-y12-coefficient-table}
The finite Fr\'echet/Sylvester machinery determines the shape of
\(E_{A,Y}^{\mathrm{BC}}\), but the present draft does not yet tabulate the
scalar midpoint jet data needed to decide
\[
E_{A,Y}^{\mathrm{BC}}=0,
\qquad
E_{A,Y}^{\mathrm{BC}}\in\mathbf C A_5^{\mathfrak f}(m),
\qquad
\text{or}\qquad
(E_{A,Y}^{\mathrm{BC}})^q\ne0.
\]
The exact missing finite table is
\[
\left\{
W_{+,0},W_{+,1},W_{+,2},W_{+,3},
\quad
\Lambda^-_{a,0},\Lambda^-_{a,1},\Lambda^-_{a,2},\Lambda^-_{a,3},
\right.
\]
\[
\left.
Y^{dd}_{g,0},\quad
Y^{d\eta+\eta d}_{g,0},\quad
Y^{\eta\eta}_{g,0}
\quad (4\le g\le7)
\right\}.
\]
Until that table is evaluated, \(E_{A,Y}^{\mathrm{BC}}\) remains the first
unresolved local actual-source/Bottleneck-C obstruction.
\end{remark}

\begin{remark}[Frozen Bottleneck-C coefficient versus source-position edge]
\label{rem:frozen-bc-coefficient-versus-source-position-edge}
The coefficient \(E_{A,Y}^{\mathrm{BC}}\) is the first Bottleneck-C
amplitude-ratio carrier.  It uses the frozen internal coefficient because the
amplitude factor
\[
a_1a_2(a_1+a_2)
\]
already supplies the first \(\delta\)-edge on the collision chart
\[
a_2=-a_1(1+\kappa\delta).
\]

The literal source-position edge
\[
E_{A,Y}^{\mathrm{pos}}
:=
\tau_{h\to m}\pi_{\mathfrak f_h}
\operatorname{Sym}
[z^7\delta^1]\,
D(G_-^{-1/2})[A_-]\,
Y_{12}^{(2)}\,
R_0
\]
belongs to the next source-motion layer and should not be used as the first
Bottleneck-C amplitude-derivative coefficient.
\end{remark}

\begin{remark}[Shortcut routes excluded for the corrected package theorem]
\label{rem:shortcut-routes-excluded-current-package-theorem}
The following routes do not prove the corrected package quotient theorem:
\begin{enumerate}
\item \(\Phi_K\)-scalar invisibility or scalar hiding after \(\Phi_K\).
\item Determinant scalarization used as the definition of the quotient
class.
\item Swap symmetry or one-pair \(K_5/K_3\) gauge laws without the actual
two-atom package object.
\item The tagged pre-whitening source block alone.
\item The raw seven-family Fr\'echet sum; tested source-constrained
nonconstant baselines leave nonzero determinant residuals.
\item Cubic/quintic edge laws plus the abstract source criterion alone.
\item Existing one-pair order-\(7\) or centered \(D_2\) formulas promoted
directly to a two-atom order-\(7\) Hessian.
\item Signed source-weight lift alone; it restores oddness but leaves
higher \(a^3,a^5,\ldots\) amplitude terms.
\end{enumerate}
Every positive route must work before \(\Phi_K\), before determinant
scalarization, and inside the fixed-target quotient package.
\end{remark}

\begin{remark}[UV-010 boundary: edge law before descent]
\label{rem:uv010-boundary-edge-law-before-descent}
The midpoint quotient trivialization is available on the good patch
\[
A_5^{\mathfrak f}(m)\ne0,
\]
but this does not construct the order-\(7\) two-atom edge package.  The
UV-010 target is the normalized fixed-midpoint quotient edge law
\[
\overline E_{12}^{(7;1)}
=
\delta^2\mathcal H_7^q(m,\kappa,\delta^2),
\]
with
\[
[R]_{\mathrm{edge}}=-\mathcal H_7^q(m,\kappa,0).
\]
Descent, \(\kappa\)-independence, and diagonal merger are downstream of
this edge-law construction.  The Gate~B edge-divisibility theorem
constructs the holomorphic environment needed for the corrected package
representative; it should not be conflated with the further UV-010 task of
identifying the exceptional-divisor value \([R]_{\mathrm{edge}}\) with the
C-side reduced-package trace.  Determinant scalarization is good-patch
bookkeeping after the quotient-valued edge object exists; it is not a
substitute for the edge package.
\end{remark}

\begin{remark}[What remains in the two-pair rewrite]
\label{rem:what-remains-two-pair-rewrite}
Theorem~\ref{thm:quantitative-defect-thickened-strengthened-coincidence}
still provides the correct algebraic bridge from the two-pair core to the
strengthened invariant $\widehat\Psi$, and
Lemma~\ref{lem:minimal-source-level-two-point-criterion}
still isolates the finite-order source criterion that would close the close
two-point branch on compact good patches.

What remains, however, is no longer raw bookkeeping through orders $3,5,7$ and
no longer a search for another local mixed branch taxonomy. The overlap-side
geometry is now sharply split:
\[
(K_v=0,\ S\neq 0)
\]
is already mixed low-order,
\[
(S=0,\ K_v\neq 0)
\]
is already absorbed by nonmixed $4$-point geometry,
and the only genuinely residual mixed corner is
\[
(S=0,\ K_v=0).
\]
There the branch is quotient-diagonal / involutive, and the remaining burden is
best read through the three-front compression of
Remark~\ref{rem:current-three-front-compression-two-point-endgame}.

Accordingly, the live downstream two-point programs are now the following.

First, a good-patch collision-cancellation theorem for the \emph{actual}
corrected cubic/quintic defects: enough regularity and swap-evenness in the
blown-up cancellation chart to justify the quadratic/projective edge-law
package.

Second, a quotient-diagonal closure theorem on the exact fixed-shear corner:
for example state-locality, transport closure, or a diagonal-collapse result
for the actual corrected defect package after quotient-visible descent.

Third, a provenance-sensitive package theorem --- for example same reduced
image germ at coincidence or collision-functoriality on the
collision/cancellation locus --- strong enough to force the later
quotient-septic collapse formally.

These are conditional downstream programs, not yet closure theorems. What the
present draft already shows is that the residual burden is concentrated there,
rather than in a missing local branch classification.
\end{remark}

\begin{remark}[Current three-front compression of the remaining burden]
\label{rem:current-three-front-compression-two-point-endgame}
The two-point endgame is now substantially narrower than a generic
``one more local sign argument'' picture.

On the nonmixed/global affine side, the remaining burden is an
$\RH$-specific exclusion problem for the isolated remote endpoint/branch
incidences encoded by the monotone gap functions $G_{\xi,J}$.
On the mixed side, the only genuinely residual overlap contribution is the
exact fixed-shear corner
\[
\{S=0\}\cap\{K_v=0\},
\]
already reduced to a quotient-diagonal / involutive transport problem.
On the honest higher-order side, the relevant target is no longer raw
order-$7$ equality of representatives, but a provenance-sensitive package
statement such as same reduced image germ at coincidence or
collision-functoriality on the collision/cancellation locus.

Accordingly, what remains is not a diffuse interaction heuristic.
It is a three-front problem: remote endpoint-gap rigidity, exact fixed-shear
quotient transport, and honest package-level coincidence/functoriality.
\end{remark}

\begin{remark}[Minimal-core reformulation]
\label{rem:minimal-core-reformulation}
The proposition above shows that the correct projective datum for the downstream
multi-pair and minimal-core extraction is the amplitude-invariant strengthened
map
\[
\widehat\Psi(h)=
\left(
\frac{u_5(h)}{c(h)},
\frac{v_5(h)}{c(h)},
\frac{\Delta_7(h)}{c(h)^2}
\right),
\]
not merely the pair $(\ell_h,\overline v_1(h))$.

By contrast, the affine curve
\[
\Gamma(h)=\left(\frac{v_5(h)}{c(h)},\frac{\Delta_7(h)}{c(h)}\right)
\]
remains the correct object for the one-pair secant/tangent geometry. The
finite-core endgame should therefore be viewed as a two-stage problem:
first extract defect-corrected coincidence information for $\widehat\Psi$,
then convert that information into affine-line incidence constraints for
$\Gamma$.
\end{remark}

\begin{remark}[Finite-core compression to a remote endpoint/branch incidence problem]
\label{rem:finite-core-compression-defect-gate-graph}
The two-stage reformulation above suggests a useful downstream combinatorial
compression.

Fix a compact relevant patch $J$. By
Theorem~\ref{thm:partial-global-no-secant}, every maximal
turning-boundary-generic block contributes no pair-pair geometry, so such a
block behaves as a single transport chamber. By
Theorem~\ref{thm:finite-order-w-zero-local-sterility}, the isolated
finite-order $W=0,\ M\neq 0$ boundaries are locally sterile for pair-pair
birth. Moreover, on the repeated-source side,
Proposition~\ref{prop:first-rooting-requires-second-odd-zero-repeated-source},
Corollary~\ref{cor:no-crossing-adjacent-two-odd-zero-under-anchor-sign-preservation},
Proposition~\ref{prop:odd-anchor-passage-no-reset}, and
Corollary~\ref{cor:first-post-anchor-chamber-dangerous} show that odd-order
$\mathcal R$-boundaries are relay-type rather than arbitrary first-root events.

After these absorptions, any surviving minimal bad core contribution from
one-pair incidence must come through the finite remote endpoint/branch
tangencies of
Lemma~\ref{lem:endpoint-gap-monotonicity-ordered-target} and
Remark~\ref{rem:remote-endpoint-gap-reduction}, equivalently through isolated
interior zeros of the strictly monotone endpoint-gap functions
\[
G_{\xi,J}(h)=y_\xi-y(h)-m(h)\bigl(\xi-x(h)\bigr).
\]
On the mixed side, the only residual contribution is the exact fixed-shear
corner $\{S=0\}\cap\{K_v=0\}$, already treated separately as a
quotient-diagonal transport problem.

Thus the remaining global nonadjacent recurrence problem is cleanly a finite
remote-incidence problem for the $G_{\xi,J}$, not a diffuse local-birth problem
on a graph of generic gates.
\end{remark}

\begin{remark}[Ceiling of the remaining affine/global lane]
\label{rem:ceiling-remaining-affine-global-lane}
The derivative-geometric reductions on the nonmixed/global side now have a
clean positive content and a clean ceiling.

The positive content is that fixed-slope tangency/incidence has been reduced to
a finite remote endpoint/branch incidence problem: support curves are strictly
concave, have single-crossing and betweenness, and the endpoint-gap functions
\[
G_{\xi,J}(h)=y_\xi-y(h)-m(h)\bigl(\xi-x(h)\bigr)
\]
are strictly monotone on ordered target branches. Thus each endpoint/branch
pair contributes at most one remote tangency.

The ceiling is that this affine/pseudoline structure is universal for analytic
plane curves. So no further purely planar, interval-graph, pseudoline, or local
convexity argument can by itself finish this lane. Any final contradiction here
must use $\RH$-specific source structure constraining the possible zeros of the
functions $G_{\xi,J}$.
\end{remark}

\begin{definition}[Minimal bad core]
\label{def:minimal-bad-core}
A normalized finite core is called bad if its quotient odd germ vanishes
identically. A minimal bad core is a bad normalized finite core with minimal
number of constituent pairs.
\end{definition}

\begin{theorem}[Defect-free quintic circuit reduction]
\label{thm:defect-free-quintic-circuit-reduction}
Let
\[
F_{\core}(z)=\sum_{j=1}^k a_jF_{h_j}(z)+\mathcal I_{\ge 2}(z)
\]
be a normalized finite core, and assume that the multi-pair interaction germ
$\mathcal I_{\ge 2}(z)$ contributes nothing to the cubic and quintic
coefficients. Write
\[
A_5^{\mathfrak f}(h)=u_5(h)I+v_5(h)S,
\qquad
V(h):=\bigl(c(h),u_5(h),v_5(h)\bigr)\in\mathbf C^3,
\]
and on $\{c\neq 0\}$ define
\[
\Xi(h):=\left(\frac{u_5(h)}{c(h)},\frac{v_5(h)}{c(h)}\right)\in\mathbf C^2.
\]
If
\[
F_{\core}(z)\equiv 0,
\]
then
\[
\sum_{j=1}^k a_j\,V(h_j)=0.
\]
Equivalently, with
\[
b_j:=a_j\,c(h_j),
\]
one has
\[
\sum_{j=1}^k b_j=0,
\qquad
\sum_{j=1}^k b_j\,\Xi(h_j)=0.
\]
Thus the points $\Xi(h_j)$ are affinely dependent.
\end{theorem}

\begin{proof}
Taking the cubic coefficient in the identity
\[
F_{\core}(z)\equiv 0
\]
and using the hypothesis that $\mathcal I_{\ge2}$ contributes nothing at
cubic order gives
\[
\sum_{j=1}^k a_j\,c(h_j)=0.
\]
Taking the quintic fixed-sector coefficient and using the corresponding
defect-free hypothesis at order $5$ gives
\[
\sum_{j=1}^k a_j\,A_5^{\mathfrak f}(h_j)=0.
\]
Since
\[
A_5^{\mathfrak f}(h_j)=u_5(h_j)I+v_5(h_j)S
\]
and $I,S$ are linearly independent, this is equivalent to
\[
\sum_{j=1}^k a_j\,u_5(h_j)=0,
\qquad
\sum_{j=1}^k a_j\,v_5(h_j)=0.
\]
Together these three identities are exactly
\[
\sum_{j=1}^k a_j\,V(h_j)=0.
\]
If $c(h_j)\neq 0$ for all $j$, dividing the last two identities by
$c(h_j)$ after setting $b_j=a_jc(h_j)$ yields
\[
\sum_{j=1}^k b_j=0,
\qquad
\sum_{j=1}^k b_j\,\Xi(h_j)=0,
\]
which is the affine dependence statement.
\end{proof}

\begin{corollary}[Defect-free minimal quintic circuits have size at most four]
\label{cor:defect-free-minimal-quintic-circuits}
In the defect-free order-$\le 5$ model of
Theorem~\ref{thm:defect-free-quintic-circuit-reduction}, every minimal bad
core has size at most $4$. More precisely:
\begin{enumerate}
\item if $k=2$, then $\Xi(h_1)=\Xi(h_2)$;
\item if $k=3$, then the three points $\Xi(h_1),\Xi(h_2),\Xi(h_3)$ are
collinear;
\item if $k=4$, then the four points $\Xi(h_1),\dots,\Xi(h_4)$ form an
affine circuit in $\mathbf C^2$.
\end{enumerate}
\end{corollary}

\begin{proof}
By Theorem~\ref{thm:defect-free-quintic-circuit-reduction}, a minimal
defect-free bad core produces a minimal affine dependence among finitely many
points of the affine plane $\mathbf C^2$. Such a minimal affine dependent
set has size at most $4$. The cases $k=2,3,4$ are the standard affine
configurations in dimension $2$: coincidence, collinear triple, or
four-point affine circuit.
\end{proof}

\begin{definition}[Lifted $v_5$-patch coordinates]
\label{def:lifted-v-patch-coordinates}
On a connected interval $U\subset\{c\neq 0,\ v_5\neq 0\}$, define
\[
x(h):=\frac{v_5(h)}{c(h)},
\qquad
y(h):=\frac{u_5(h)}{c(h)},
\qquad
\sigma_v(h):=\frac{\Delta_7(h)}{c(h)\,v_5(h)}.
\]
The associated lifted planar and space curves are
\[
\Lambda_v(h):=\bigl(x(h),\sigma_v(h)\bigr)\in\mathbf C^2,
\qquad
Q_v(h):=\bigl(y(h),x(h),\sigma_v(h)\bigr)\in\mathbf C^3.
\]
\end{definition}

\begin{proposition}[Defect-free lifted affine dependence on a $v_5$-patch]
\label{prop:defect-free-lifted-affine-dependence-vpatch}
Let
\[
F_{\core}(z)=\sum_{j=1}^k a_jF_{h_j}(z)+\mathcal I_{\ge 2}(z)
\]
be a normalized finite core, and assume:
\begin{enumerate}
\item all parameters $h_j$ lie in one connected interval
$U\subset\{c\neq 0,\ v_5\neq 0\}$;
\item the multi-pair interaction germ $\mathcal I_{\ge2}(z)$ contributes
nothing to the cubic, quintic, or septic coefficients.
\end{enumerate}
Then, after choosing on $U$ the local septic gauge $v_7=0$, the vanishing
identity
\[
F_{\core}(z)\equiv 0
\]
implies
\[
\sum_{j=1}^k b_j=0,
\qquad
\sum_{j=1}^k b_j\,y(h_j)=0,
\qquad
\sum_{j=1}^k b_j\,x(h_j)=0,
\qquad
\sum_{j=1}^k b_j\,\sigma_v(h_j)=0,
\]
where $b_j:=a_jc(h_j)$. Equivalently, the points $Q_v(h_j)\in\mathbf C^3$
are affinely dependent.
\end{proposition}

\begin{proof}
The cubic and quintic identities are exactly those in
Theorem~\ref{thm:defect-free-quintic-circuit-reduction}, so after setting
$b_j=a_jc(h_j)$ we obtain
\[
\sum_{j=1}^k b_j=0,
\qquad
\sum_{j=1}^k b_j\,y(h_j)=0,
\qquad
\sum_{j=1}^k b_j\,x(h_j)=0.
\]
On the interval $U$ we may choose the local septic gauge $v_7=0$. Then
\[
u_7(h)=\frac{\Delta_7(h)}{v_5(h)}=c(h)\,\sigma_v(h).
\]
Since the septic interaction is assumed absent, taking the septic fixed-sector
coefficient of $F_{\core}\equiv 0$ yields
\[
\sum_{j=1}^k a_j\,u_7(h_j)=0.
\]
Substituting $u_7(h_j)=c(h_j)\sigma_v(h_j)$ gives
\[
\sum_{j=1}^k b_j\,\sigma_v(h_j)=0.
\]
Hence the points $Q_v(h_j)$ are affinely dependent.
\end{proof}

\begin{corollary}[Geometric meaning of defect-free lifted affine dependence]
\label{cor:geometric-meaning-defect-free-lifted-affine-dependence-vpatch}
Assume the hypotheses of
Proposition~\ref{prop:defect-free-lifted-affine-dependence-vpatch}. Then:
\begin{enumerate}
\item if $k=2$, the two points $Q_v(h_1)$ and $Q_v(h_2)$ coincide;
\item if $k=3$, the three points $Q_v(h_1),Q_v(h_2),Q_v(h_3)$ are
collinear in $\mathbf C^3$;
\item if $k=4$, the four points $Q_v(h_1),\dots,Q_v(h_4)$ are coplanar in
$\mathbf C^3$.
\end{enumerate}
\end{corollary}

\begin{proof}
By Proposition~\ref{prop:defect-free-lifted-affine-dependence-vpatch}, the
points $Q_v(h_j)$ are affinely dependent. For $k=2$, affine dependence of
two points is exactly coincidence. For $k=3$, affine dependence of three
points is exactly collinearity. For $k=4$, affine dependence of four points
is exactly coplanarity.
\end{proof}

\begin{corollary}[Explicit four-point determinant form on a $v_5$-patch]
\label{cor:explicit-four-point-determinant-form-vpatch}
Assume the hypotheses of
Proposition~\ref{prop:defect-free-lifted-affine-dependence-vpatch} with
$k=4$. Then the coplanarity of the four lifted points is equivalent to
\[
\det\!\begin{pmatrix}
1 & y(h_1) & x(h_1) & \sigma_v(h_1)\\
1 & y(h_2) & x(h_2) & \sigma_v(h_2)\\
1 & y(h_3) & x(h_3) & \sigma_v(h_3)\\
1 & y(h_4) & x(h_4) & \sigma_v(h_4)
\end{pmatrix}
=0.
\]
Equivalently,
\[
\det\!\begin{pmatrix}
1 & \dfrac{u_5(h_1)}{c(h_1)} & \dfrac{v_5(h_1)}{c(h_1)} &
\dfrac{\Delta_7(h_1)}{c(h_1)\,v_5(h_1)}\\[1.2ex]
1 & \dfrac{u_5(h_2)}{c(h_2)} & \dfrac{v_5(h_2)}{c(h_2)} &
\dfrac{\Delta_7(h_2)}{c(h_2)\,v_5(h_2)}\\[1.2ex]
1 & \dfrac{u_5(h_3)}{c(h_3)} & \dfrac{v_5(h_3)}{c(h_3)} &
\dfrac{\Delta_7(h_3)}{c(h_3)\,v_5(h_3)}\\[1.2ex]
1 & \dfrac{u_5(h_4)}{c(h_4)} & \dfrac{v_5(h_4)}{c(h_4)} &
\dfrac{\Delta_7(h_4)}{c(h_4)\,v_5(h_4)}
\end{pmatrix}
=0.
\]
Multiplying row $j$ by $v_5(h_j)$ gives the canonical determinant identity
\[
\det\!\begin{pmatrix}
v_5(h_1) & u_5(h_1) & \dfrac{v_5(h_1)^2}{c(h_1)} &
\dfrac{\Delta_7(h_1)}{c(h_1)}\\[1.2ex]
v_5(h_2) & u_5(h_2) & \dfrac{v_5(h_2)^2}{c(h_2)} &
\dfrac{\Delta_7(h_2)}{c(h_2)}\\[1.2ex]
v_5(h_3) & u_5(h_3) & \dfrac{v_5(h_3)^2}{c(h_3)} &
\dfrac{\Delta_7(h_3)}{c(h_3)}\\[1.2ex]
v_5(h_4) & u_5(h_4) & \dfrac{v_5(h_4)^2}{c(h_4)} &
\dfrac{\Delta_7(h_4)}{c(h_4)}
\end{pmatrix}
=0.
\]
Equivalently, multiplying instead row $j$ by $c(h_j)v_5(h_j)$ gives the
fully polynomial version
\[
\det\!\begin{pmatrix}
c(h_1)v_5(h_1) & u_5(h_1)v_5(h_1) & v_5(h_1)^2 & \Delta_7(h_1)\\
c(h_2)v_5(h_2) & u_5(h_2)v_5(h_2) & v_5(h_2)^2 & \Delta_7(h_2)\\
c(h_3)v_5(h_3) & u_5(h_3)v_5(h_3) & v_5(h_3)^2 & \Delta_7(h_3)\\
c(h_4)v_5(h_4) & u_5(h_4)v_5(h_4) & v_5(h_4)^2 & \Delta_7(h_4)
\end{pmatrix}
=0.
\]
\end{corollary}

\begin{proof}
For $k=4$, Corollary~\ref{cor:geometric-meaning-defect-free-lifted-affine-dependence-vpatch}
shows that the four lifted points
\[
Q_v(h_j)=\bigl(y(h_j),x(h_j),\sigma_v(h_j)\bigr)
\]
are coplanar. This is equivalent to the vanishing of the affine $4\times4$
determinant with first column $1$. The displayed reformulations follow by
substituting the definitions of $x,y,\sigma_v$ and multiplying rows by the
indicated nonzero factors.
\end{proof}

\begin{remark}[Finite-core meaning of the lifted $v_5$-patch geometry]
\label{rem:finite-core-meaning-lifted-vpatch-geometry}
The defect-free finite-core reduction now places the lifted $v_5$-patch
geometry into a sharp global context.

Indeed, by
Theorem~\ref{thm:defect-free-quintic-circuit-reduction},
Corollary~\ref{cor:defect-free-minimal-quintic-circuits}, and
Proposition~\ref{prop:defect-free-lifted-affine-dependence-vpatch},
defect-free bad cores on a good $v_5$-patch reduce to affine dependence of the
lifted points
\[
Q_v(h)=\bigl(y(h),x(h),\sigma_v(h)\bigr),
\]
with effective ceiling
\[
k\le 4.
\]
Thus the three low core sizes have the following exact lifted meanings:
\[
k=2 \Longleftrightarrow Q_v(h_1)=Q_v(h_2),
\]
\[
k=3 \Longleftrightarrow Q_v(h_1),Q_v(h_2),Q_v(h_3)
\ \text{are collinear},
\]
\[
k=4 \Longleftrightarrow Q_v(h_1),\dots,Q_v(h_4)
\ \text{are coplanar}.
\]

So the lifted $3$-point and $4$-point analyses are not auxiliary side
computations: they are exactly the $k=3$ and $k=4$ geometric layers in the
finite-core program. In particular, once the lifted $3$-point branch is
globally rigid on good $v_5$-patches, the remaining finite-core burden shifts
upward to the two-point source-level merger problem and to the mixed/nonmixed
lifted $4$-point mechanisms.
\end{remark}

\begin{remark}[Parallel $u_5$-patch version]
\label{rem:parallel-u5-patch-version}
There is a completely parallel local story on any connected interval in
$\{c\neq 0,\ u_5\neq 0\}$, obtained by choosing the local septic gauge
$u_7=0$ instead of $v_7=0$. Then one defines
\[
\tau_u(h):=-\frac{\Delta_7(h)}{c(h)\,u_5(h)},
\qquad
Q_u(h):=\left(\frac{u_5(h)}{c(h)},\frac{v_5(h)}{c(h)},\tau_u(h)\right),
\]
and the same defect-free affine-dependence conclusions hold with $Q_v$
replaced by $Q_u$. Since the present one-pair geometry already uses
\[
x(h)=\frac{v_5(h)}{c(h)}
\]
as the distinguished ECT coordinate, the $v_5\neq 0$ patch is the natural
first setting and will be used below.
\end{remark}

\begin{proposition}[Centered quintic-patch dichotomy]
\label{prop:centered-quintic-patch-dichotomy}
Assume the centered diagonal family is a genuine analytic branch in the
underlying $h$-geometry, and that the centered scalar package from
Remark~\ref{rem:centered-symmetric-placement-taylor-package} holds on that
branch. Then, after shrinking, there exist one-sided punctured subarcs on which
at least one of the following holds:
\begin{enumerate}
\item
\[
c\neq 0,\qquad M\neq 0,\qquad v_5\neq 0,
\]
so the branch enters a good $v_5$-patch;
\item
\[
c\neq 0,\qquad u_5\neq 0,
\]
so the parallel $u_5$-patch of
Remark~\ref{rem:parallel-u5-patch-version} applies.
\end{enumerate}
In particular, the centered diagonal branch enters some quintic patch on
one-sided punctured arcs.
\end{proposition}

\begin{proof}
By Corollary~\ref{cor:centered-quintic-fixed-sector-nontriviality},
\[
A_5^{\mathfrak f}(s)\not\equiv 0,
\]
so at least one of $u_5$ or $v_5$ is a nonzero analytic germ on the centered
branch.

If $v_5\not\equiv 0$, then its zeros are isolated. Since
\[
c(s)=C_0+C_2s^2+\cdots,\qquad C_0\neq 0,
\]
we also have $c\neq 0$ near the basepoint. On any one-sided punctured arc
where $v_5\neq 0$, define
\[
x=\frac{v_5}{c}.
\]
This is a nonconstant analytic germ. If $M\equiv 0$ on such an arc, then
\[
x'=\left(\frac{v_5}{c}\right)'=\frac{M}{c^2}\equiv 0,
\]
so $x$ would be constant, contradiction. Hence $M$ is not identically zero on
that arc, and by analyticity its zeros are isolated. After shrinking again, one
obtains
\[
c\neq 0,\qquad M\neq 0,\qquad v_5\neq 0,
\]
i.e. a good $v_5$-patch.

If instead $v_5\equiv 0$, then $u_5\not\equiv 0$. Its zeros are isolated, so
after shrinking to a one-sided punctured arc one has
\[
c\neq 0,\qquad u_5\neq 0,
\]
and the $u_5$-patch of Remark~\ref{rem:parallel-u5-patch-version} applies.
\end{proof}

\begin{corollary}[Conditional punctured $v_5$-patch promotion]
\label{cor:conditional-punctured-v5-patch-promotion}
Under the hypotheses of
Proposition~\ref{prop:centered-quintic-patch-dichotomy}, if
\[
v_5\not\equiv 0
\]
on the centered diagonal branch, then after shrinking there exist one-sided
punctured subarcs contained in
\[
\{c\neq 0,\ M\neq 0,\ v_5\neq 0\}.
\]
Hence on such arcs
\[
x=\frac{v_5}{c}
\]
is automatically a valid local coordinate.
\end{corollary}

\begin{proof}
This is exactly the first branch of the proof of
Proposition~\ref{prop:centered-quintic-patch-dichotomy}.
\end{proof}

\begin{proposition}[Linear independence of $1,u_5/c,v_5/c$]
\label{prop:linear-independence-one-u5-over-c-v5-over-c}
Let $J$ be a connected open interval on which
\[
c\neq 0,
\qquad
M\neq 0.
\]
Then the three functions
\[
1,
\qquad
\frac{u_5}{c},
\qquad
\frac{v_5}{c}
\]
are linearly independent on $J$.
\end{proposition}

\begin{proof}
Assume there exist constants $A,B,C\in\mathbf C$, not all zero, such that
\[
A+B\frac{u_5}{c}+C\frac{v_5}{c}\equiv 0
\qquad\text{on }J.
\]
Multiplying by $c$ gives
\[
A\,c+B\,u_5+C\,v_5\equiv 0.
\]

If $B=0$, then
\[
A\,c+C\,v_5\equiv 0,
\]
so $v_5/c$ is constant on $J$. But
\[
\left(\frac{v_5}{c}\right)'=\frac{M}{c^2},
\]
and $M\neq 0$ on $J$. This is impossible.

Hence $B\neq 0$, and after renormalization we may write
\[
u_5=\alpha\,v_5+\beta\,c
\]
for constants $\alpha,\beta$.

Now isolate the graded slice corresponding to
\[
(K_1,\eta_2,q^{(5)}).
\]
By Lemmas~\ref{lem:quintic-trace-witness}
and~\ref{lem:no-cancellation-quintic-trace}, the tagged monomial
\[
\kappa(h)\,\eta_2(h)\,q^{(5)}(h)
\]
occurs with nonzero coefficient in $\operatorname{tr}J_5^{\pair}(h)$.
Since
\[
A_5^{\mathfrak f}(h)=u_5(h)I+v_5(h)S
\]
and $S$ is traceless, this means exactly that the same tagged monomial
occurs in the $I$-coefficient $u_5(h)$.

It cannot occur in $c(h)$, since $c$ is cubic and carries no quintic
$q^{(5)}$-slice. It also cannot occur in $v_5(h)$. Indeed, on the tagged
slice, the only surviving quintic $K_1$-transport channels are those already
isolated in Lemma~\ref{lem:no-cancellation-quintic-trace}. For the channel
\[
X_0JX_4-X_4JX_0,
\]
Lemma~\ref{lem:whitening-parity} shows that the relevant
$(\eta_2,q^{(5)})$-slice is diagonal, hence contributes only to the
$I$-component. For the commutator channel
\[
X_3JX_1-X_1JX_3,
\]
Proposition~\ref{prop:explicit-quintic-fixed-sector} gives
\[
v_5^{(\mathrm{comm})}=\gamma u-\alpha w.
\]
On the tagged slice the $\eta_2$-input in $X_1$ is purely off-diagonal, so
$u=w=0$, while the $q^{(5)}$-input in $X_3$ is diagonal. Hence
\[
v_5^{(\mathrm{comm})}=0
\]
on that tagged slice.

Thus the monomial $\kappa\eta_2 q^{(5)}$ occurs in $u_5$ but in neither
$v_5$ nor $c$, contradicting
\[
u_5=\alpha\,v_5+\beta\,c.
\]
This proves the claim.
\end{proof}

\begin{corollary}[Automatic $u_5$-patch on a good $M\neq 0$ interval]
\label{cor:automatic-u5-patch-on-good-M-interval}
Let $J$ be a connected open interval on which
\[
c\neq 0,
\qquad
M\neq 0.
\]
Then $u_5\not\equiv 0$ on $J$. Consequently, every connected component of
\[
J\setminus u_5^{-1}(0)
\]
is a valid local $u_5$-patch in the sense of
Remark~\ref{rem:parallel-u5-patch-version}.
\end{corollary}

\begin{proof}
If $u_5\equiv 0$ on $J$, then the three functions
\[
1,\qquad \frac{u_5}{c},\qquad \frac{v_5}{c}
\]
would be linearly dependent, contradicting
Proposition~\ref{prop:linear-independence-one-u5-over-c-v5-over-c}. Thus
$u_5\not\equiv 0$ on $J$, so its zeros are isolated and each connected
component of $J\setminus u_5^{-1}(0)$ lies in $\{c\neq 0,\ u_5\neq 0\}$.
The parallel $u_5$-patch geometry then follows from
Remark~\ref{rem:parallel-u5-patch-version}.
\end{proof}

\begin{proposition}[Non-affinity of the lifted septic scalar]
\label{prop:sigmav-not-affine-in-x}
Let $J$ be a connected open interval contained in the quintic-generic region
and in
\[
\{h:c(h)\neq 0,\ v_5(h)\neq 0\}.
\]
Then there do not exist constants $A,B\in\mathbf C$ such that
\[
\sigma_v(h)\equiv A\,x(h)+B
\qquad\text{on }J.
\]
Equivalently,
\[
\Delta_7(h)\not\equiv A\,v_5(h)^2+B\,c(h)v_5(h)
\qquad\text{on }J.
\]
\end{proposition}

\begin{proof}
If
\[
\sigma_v(h)\equiv A\,x(h)+B,
\]
then multiplying by $c(h)v_5(h)$ gives
\[
\Delta_7(h)\equiv A\,v_5(h)^2+B\,c(h)v_5(h).
\]
The right-hand side is built only from the cubic and quintic layers and
therefore has no direct septic $q^{(7)}$-slice. But by
Proposition~\ref{prop:direct-q7-total-delta7-nonzero}, the direct
$q^{(7)}$-slice of $\Delta_7$ is generically nonzero on $J$. This is a
contradiction.
\end{proof}

\begin{definition}[Lifted derivative numerator and curvature numerator]
\label{def:lifted-vpatch-derivative-and-curvature}
On $\{c\neq 0,\ v_5\neq 0\}$, define
\[
P_v(h):=c(h)v_5(h)\Delta_7'(h)-\Delta_7(h)\bigl(c(h)v_5(h)\bigr)'.
\]
Then
\[
\sigma_v'(h)=\frac{P_v(h)}{c(h)^2\,v_5(h)^2}.
\]
Define also the affine curvature numerator of the lifted planar curve
\[
\Lambda_v(h)=\bigl(x(h),\sigma_v(h)\bigr)
\]
by
\[
K_v(h):=x'(h)\sigma_v''(h)-\sigma_v'(h)x''(h).
\]
\end{definition}

\begin{proposition}[Curvature bridge for the lifted $v_5$-patch curve]
\label{prop:curvature-bridge-lifted-vpatch}
On $\{c\neq 0,\ M\neq 0,\ v_5\neq 0\}$,
\[
\frac{d\sigma_v}{dx}
=
\frac{P_v}{M\,v_5^2},
\]
and
\[
K_v(h)
=
\frac{M(h)^2}{c(h)^4}
\frac{d}{dh}\!\left(\frac{P_v(h)}{M(h)\,v_5(h)^2}\right).
\]
In particular, if $K_v\equiv 0$ on a connected open interval $J\subset
\{c\neq 0,\ M\neq 0,\ v_5\neq 0\}$, then $\sigma_v$ is affine in $x$ on
$J$, contradicting Proposition~\ref{prop:sigmav-not-affine-in-x}. Hence
\[
K_v\not\equiv 0
\]
on every such interval.
\end{proposition}

\begin{proof}
Since
\[
x=\frac{v_5}{c},
\qquad
x'=\frac{M}{c^2},
\qquad
\sigma_v'=\frac{P_v}{c^2v_5^2},
\]
we obtain
\[
\frac{d\sigma_v}{dx}
=
\frac{\sigma_v'}{x'}
=
\frac{P_v}{M\,v_5^2}.
\]
Differentiating with respect to $h$ and multiplying by
\[
x'=\frac{M}{c^2}
\]
gives
\[
K_v
=
x'^2\,\frac{d}{dx}\!\left(\frac{d\sigma_v}{dx}\right)
=
\frac{M^2}{c^4}\,
\frac{d}{dh}\!\left(\frac{P_v}{M\,v_5^2}\right).
\]
If $K_v\equiv 0$ on a connected interval with $M\neq 0$, then
\[
\frac{P_v}{M\,v_5^2}
\]
is constant there, so $d\sigma_v/dx$ is constant and $\sigma_v$ is affine
in $x$. This contradicts
Proposition~\ref{prop:sigmav-not-affine-in-x}.
\end{proof}

\begin{corollary}[Local defect-free three-point exclusion on lifted ECT subintervals]
\label{cor:local-three-point-exclusion-lifted-vpatch}
Let $J\subset\{c\neq 0,\ M\neq 0,\ v_5\neq 0\}$ be a connected interval on
which $K_v$ has fixed nonzero sign. Then the triple
\[
1,\qquad x,\qquad \sigma_v
\]
is an ECT system on $J$. In particular, no affine line meets the lifted
planar curve
\[
\Lambda_v(h)=\bigl(x(h),\sigma_v(h)\bigr)
\]
in three points of $J$. Hence the defect-free three-point branch is excluded
on such a subinterval.
\end{corollary}

\begin{proof}
The Wronskian of $1,x,\sigma_v$ is exactly $K_v$. If $K_v$ has fixed
nonzero sign on $J$, the ECT criterion applies.
\end{proof}

\begin{corollary}[Lifted order-$7$ Wronskian is not identically zero]
\label{cor:lifted-vpatch-wronskian-not-identically-zero}
Let $J$ be a connected open interval contained in the quintic-generic region
and in
\[
\{h:c(h)\neq 0,\ M(h)\neq 0,\ v_5(h)\neq 0\}.
\]
Then
\[
T_v(h):=
W_h\!\left(
1,\ \frac{u_5(h)}{c(h)},\ \frac{v_5(h)}{c(h)},\
\frac{\Delta_7(h)}{c(h)\,v_5(h)}
\right)
\]
is not identically zero on $J$.
\end{corollary}

\begin{proof}
Write
\[
y:=\frac{u_5}{c},
\qquad
x:=\frac{v_5}{c},
\qquad
\sigma_v:=\frac{\Delta_7}{c\,v_5}.
\]
If $T_v\equiv 0$ on $J$, then by analyticity the four functions
\[
1,\ y,\ x,\ \sigma_v
\]
are locally linearly dependent on $J$, hence linearly dependent on all of
the connected interval $J$. By
Proposition~\ref{prop:linear-independence-one-u5-over-c-v5-over-c}, the first
three are linearly independent, so there exist constants $A,B,C$ such that
\[
\sigma_v=A+By+Cx.
\]
Multiplying by $c\,v_5$ gives
\[
\Delta_7=A\,c\,v_5+B\,u_5v_5+C\,v_5^2.
\]
The right-hand side is built entirely from the cubic/quintic layers and
therefore carries no direct $q^{(7)}$-slice. This contradicts
Proposition~\ref{prop:direct-q7-total-delta7-nonzero}, which gives
\[
\Delta_7^{(q^{(7)})}\not\equiv 0
\]
on $J$.
\end{proof}

\begin{proposition}[Exact $x$-coordinate identities for the lifted $4$-point invariants]
\label{prop:exact-x-coordinate-identities-lifted-four-point}
On a good $v_5$-patch, write
\[
x=\frac{v_5}{c},
\qquad
Y=\frac{u_5}{c},
\qquad
S=\sigma_v=\frac{\Delta_7}{c\,v_5}.
\]
Then:
\begin{enumerate}
\item the lifted planar curvature numerator satisfies
\[
K_v=(x')^3\,S''(x);
\]
\item the lifted $4$-point Wronskian satisfies
\[
T_v=(x')^6\,W_x(1,Y,x,S)
=(x')^6\bigl(S''Y'''-S'''Y''\bigr);
\]
\item equivalently,
\[
T_v=(x')^6\,S''^{\,2}\,\frac{d}{dx}\!\left(\frac{Y''}{S''}\right)
=K_v^{\,2}\,\frac{d}{dx}\!\left(\frac{Y''}{S''}\right).
\]
\end{enumerate}
In particular, $K_v$ and $S''$ have the same zero set and the same
vanishing order on every good $v_5$-patch.
\end{proposition}

\begin{proof}
By definition,
\[
K_v=x'\sigma_v''-\sigma_v' x''.
\]
Since $S=\sigma_v$ as a function of the local coordinate $x$, we have
\[
\sigma_v'=S' x',
\qquad
\sigma_v''=S''(x')^2+S' x'',
\]
hence
\[
K_v=x'\bigl(S''(x')^2+S'x''\bigr)-S' x' x''=(x')^3S''.
\]

Next,
\[
T_v=W_h(1,Y,x,S).
\]
Under reparametrization $h\mapsto x(h)$, a $4\times4$ Wronskian gains the
factor $(x')^{0+1+2+3}=(x')^6$, so
\[
T_v=(x')^6\,W_x(1,Y,x,S).
\]
Now
\[
W_x(1,Y,x,S)
=
\det\!\begin{pmatrix}
1&Y&x&S\\
0&Y'&1&S'\\
0&Y''&0&S''\\
0&Y'''&0&S'''
\end{pmatrix}
=
S''Y'''-S'''Y'',
\]
which proves the second identity.

Finally,
\[
S''Y'''-S'''Y''
=
S''^{\,2}\,\frac{d}{dx}\!\left(\frac{Y''}{S''}\right),
\]
and using $K_v=(x')^3S''$ gives
\[
(x')^6S''^{\,2}=K_v^{\,2}.
\]
This yields the third identity.
\end{proof}

\begin{proposition}[Canonical scalar invariants $\mathcal Z$ and $\mathcal R$ on a nonmixed lifted patch]
\label{prop:canonical-Z-R-invariants-lifted-four-point}
On a good $v_5$-patch, write
\[
x=\frac{v_5}{c},
\qquad
Y=\frac{u_5}{c},
\qquad
S=\sigma_v=\frac{\Delta_7}{c\,v_5}.
\]
On the nonmixed subregion where $K_v\neq 0$, define
\[
\mathcal Z:=\frac{Y''}{S''},
\qquad
\mathcal R:=\frac{d\mathcal Z}{dx}.
\]
Then:
\begin{enumerate}
\item
\[
\mathcal R=\frac{T_v}{K_v^2};
\]
\item under the affine gauge
\[
Y\longmapsto Y+\alpha+\beta x+\gamma S,
\]
one has
\[
\mathcal Z\longmapsto \mathcal Z+\gamma,
\qquad
\mathcal R\longmapsto \mathcal R.
\]
Hence $\mathcal R$ is the canonical scalar invariant of the lifted
$4$-point branch on every nonmixed patch.
\end{enumerate}
\end{proposition}

\begin{proof}
By Proposition~\ref{prop:exact-x-coordinate-identities-lifted-four-point},
\[
T_v=(x')^6\,S''^{\,2}\,\frac{d}{dx}\!\left(\frac{Y''}{S''}\right)
=K_v^{\,2}\,\frac{d}{dx}\!\left(\frac{Y''}{S''}\right),
\]
so
\[
\mathcal R=\frac{d}{dx}\!\left(\frac{Y''}{S''}\right)=\frac{T_v}{K_v^2}.
\]

For the gauge transformation, one has
\[
(Y+\alpha+\beta x+\gamma S)''=Y''+\gamma S'',
\]
hence
\[
\mathcal Z\longmapsto \frac{Y''+\gamma S''}{S''}=\mathcal Z+\gamma.
\]
Differentiating with respect to $x$ gives $\mathcal R\mapsto \mathcal R$.
\end{proof}

\begin{remark}[Below-lifted Wronskian-ratio reformulation]
\label{rem:below-lifted-wronskian-ratio-reformulation}
Let
\[
x=\frac{v_5}{c},
\qquad
Y=\frac{u_5}{c},
\qquad
S=\frac{\Delta_7}{c\,v_5},
\qquad
\lambda:=\frac{\Delta_7}{v_5},
\qquad
M:=cv_5'-c'v_5.
\]
For any scalar germ $f$, define
\[
J_f(h):=
\det\!\begin{pmatrix}
c(h)&v_5(h)&f(h)\\
c'(h)&v_5'(h)&f'(h)\\
c''(h)&v_5''(h)&f''(h)
\end{pmatrix}.
\]
Then
\[
\frac{d^2}{dx^2}\!\left(\frac{f}{c}\right)
=
\frac{c^3}{M^3}\,J_f.
\]
Applying this to $f=u_5$ and $f=\lambda$ gives
\[
Y''=\frac{c^3}{M^3}J_{u_5},
\qquad
S''=\frac{c^3}{M^3}J_{\lambda},
\]
hence
\[
\boxed{
\mathcal Z:=\frac{Y''}{S''}=\frac{J_{u_5}}{J_{\lambda}}.
}
\]
Differentiating in the $x$-coordinate yields
\[
\boxed{
\mathcal R:=\mathcal Z'=\frac{T_v}{K_v^2}
=
\frac{c^2}{M}\frac{d}{dh}\!\left(\frac{J_{u_5}}{J_{\lambda}}\right).
}
\]
Thus the lifted $4$-point branch may be viewed equivalently as a
parity/divisibility problem for the Wronskian ratio $J_{u_5}/J_\lambda$.
On the nonmixed side, however, the later source-level no-go results show that
the ratio/Wronskian package by itself is not yet restrictive enough: the
remaining burden is a genuine source-coupling theorem, not a generic
Wronskian-ratio argument.
\end{remark}

\begin{proposition}[Source-level Wronskian for the nonmixed $4$-point branch]
\label{prop:source-level-Wronskian-nonmixed-four-point}
Define
\[
\mathcal W_{\mathrm{src}}
:=
J_{\lambda}\,J_{u_5}'-J_{u_5}\,J_{\lambda}'.
\]
Then
\[
\boxed{
\mathcal R
=
\frac{c^2}{M}\,\frac{\mathcal W_{\mathrm{src}}}{J_\lambda^2}.
}
\]
In particular, on every nonmixed arc on which $J_\lambda\neq 0$, the analytic
scalars $\mathcal R$ and $\mathcal W_{\mathrm{src}}$ differ only by a
nonvanishing analytic factor. Hence:

\begin{enumerate}
\item $\mathcal R=0\iff \mathcal W_{\mathrm{src}}=0$;
\item $\operatorname{ord}\mathcal R=\operatorname{ord}\mathcal W_{\mathrm{src}}$;
\item the local parity classification of isolated zeros of $\mathcal R$
transfers verbatim to $\mathcal W_{\mathrm{src}}$.
\end{enumerate}

Moreover, under the affine gauge change
\[
u_5\longmapsto u_5+\alpha c+\beta v_5+\gamma \lambda
\]
one has
\[
J_{u_5}\longmapsto J_{u_5}+\gamma J_\lambda,
\]
and therefore
\[
\mathcal W_{\mathrm{src}}\longmapsto \mathcal W_{\mathrm{src}}.
\]
So $\mathcal W_{\mathrm{src}}$ is the canonical source-level gauge-invariant
nonmixed scalar.
\end{proposition}

\begin{proof}
Differentiate the ratio
\[
\mathcal Z=\frac{J_{u_5}}{J_\lambda}
\]
from Remark~\ref{rem:below-lifted-wronskian-ratio-reformulation}:
\[
\frac{d}{dh}\!\left(\frac{J_{u_5}}{J_\lambda}\right)
=
\frac{J_\lambda J_{u_5}'-J_{u_5}J_\lambda'}{J_\lambda^2}
=
\frac{\mathcal W_{\mathrm{src}}}{J_\lambda^2}.
\]
Substituting this into
\[
\mathcal R
=
\frac{c^2}{M}\frac{d}{dh}\!\left(\frac{J_{u_5}}{J_\lambda}\right)
\]
gives the displayed identity.

On a nonmixed arc, $J_\lambda\neq 0$ and $c^2/M\neq 0$, so $\mathcal R$
and $\mathcal W_{\mathrm{src}}$ differ by a nonvanishing analytic factor.
This proves the zero-set and parity statements.

For gauge invariance, the terms $\alpha c$ and $\beta v_5$ do not change
$J_{u_5}$, while $\gamma\lambda$ adds $\gamma J_\lambda$. Therefore
\[
(J_{u_5}+\gamma J_\lambda)'=J_{u_5}'+\gamma J_\lambda',
\]
and
\[
J_\lambda(J_{u_5}'+\gamma J_\lambda')-(J_{u_5}+\gamma J_\lambda)J_\lambda'
=
J_\lambda J_{u_5}'-J_{u_5}J_\lambda'.
\qedhere
\]
\end{proof}

\begin{remark}[Determinant-level parity on a centered good patch]
\label{rem:determinant-level-parity-centered-good-patch}
Assume the centered scalar package from
Remark~\ref{rem:centered-symmetric-placement-taylor-package} and suppose that
$x=v_5/c$ is being used as the induced odd collision-side coordinate.
Then
\[
u_5,\ \lambda:=\frac{\Delta_7}{v_5}\ \text{are odd in }s,
\qquad
c\ \text{is even in }s,
\qquad
v_5\ \text{is odd in }s.
\]
Hence the rows of the defining determinants for $J_{u_5}$ and $J_\lambda$
have parities
\[
(e,o,o),\qquad (o,e,e),\qquad (e,o,o),
\]
so each determinant monomial is odd. Therefore
\[
J_{u_5}(s)\ \text{and}\ J_\lambda(s)\ \text{are odd in }s.
\]
Consequently,
\[
\mathcal Z=\frac{J_{u_5}}{J_\lambda}
\]
is even and
\[
\mathcal R=\mathcal Z'
\]
is odd, in agreement with the parity package obtained from the direct
$Y''/S''$-analysis.
\end{remark}

\begin{corollary}[Centered parity of the source-level Wronskian]
\label{cor:centered-parity-source-level-Wronskian}
Under the hypotheses of
Remark~\ref{rem:determinant-level-parity-centered-good-patch}, the source-level
Wronskian
\[
\mathcal W_{\mathrm{src}}(s)
=
J_\lambda(s)\,J_{u_5}'(s)-J_{u_5}(s)\,J_\lambda'(s)
\]
is odd in the centered variable $s$.

Consequently, any isolated centered zero of $\mathcal W_{\mathrm{src}}$ has
odd order if and only if it is a potentially branch-seeding zero for the
nonmixed $4$-point geometry; even-order centered zeros are locally sterile.
\end{corollary}

\begin{proof}
By Remark~\ref{rem:determinant-level-parity-centered-good-patch},
\[
J_{u_5}(s)\ \text{and}\ J_\lambda(s)\ \text{are odd in }s.
\]
Hence
\[
J_{u_5}'(s)\ \text{and}\ J_\lambda'(s)\ \text{are even in }s.
\]
Therefore
\[
\mathcal W_{\mathrm{src}}(s)
=
J_\lambda(s)\,J_{u_5}'(s)-J_{u_5}(s)\,J_\lambda'(s)
\]
is odd in $s$. The final statement follows from
Proposition~\ref{prop:source-level-Wronskian-nonmixed-four-point}, which
transfers the parity classification from $\mathcal R$ to
$\mathcal W_{\mathrm{src}}$.
\end{proof}

\begin{corollary}[Mixed low-order exclusion on a centered good $v_5$-patch]
\label{cor:mixed-low-order-exclusion-centered-good-v5-patch}
Under the hypotheses of
Remark~\ref{rem:determinant-level-parity-centered-good-patch}, one has
\[
\mathcal Z(x)\ \text{even},
\qquad
\mathcal R(x)\ \text{odd},
\qquad
K_v(x)\ \text{odd}.
\]
Consequently,
\[
(r,\operatorname{ord}_0\mathcal Z)\notin\{(1,1),(1,3),(2,2)\}.
\]
Equivalently, none of the explicit mixed low-order models
\[
(m,n)=(3,4),\ (3,6),\ (4,6)
\]
can occur on such a centered good patch.
\end{corollary}

\begin{proof}
Since $\mathcal Z$ is even, its finite vanishing order must be even, so the
cases
\[
(r,\operatorname{ord}_0\mathcal Z)=(1,1),\ (1,3)
\]
are impossible. Since $K_v$ is odd, its vanishing order $r=\operatorname{ord}_0K_v$
must be odd, so
\[
(r,\operatorname{ord}_0\mathcal Z)=(2,2)
\]
is impossible as well.
\end{proof}

\begin{theorem}[Local exhaustion of the mixed $4$-point branch near the centered base]
\label{thm:local-exhaustion-mixed-four-point-centered-base}
Assume the centered scalar package from
Remark~\ref{rem:centered-symmetric-placement-taylor-package}, the right-pincer
local reduction from
Corollary~\ref{cor:right-pincer-local-reduction-centered-base}, and the finite-jet
scalar descent from
Corollary~\ref{cor:finite-jet-scalar-descent-centered-branch}. Then every
sufficiently small nearby genuinely new mixed $4$-point obstruction near the
centered diagonal base is excluded.

Equivalently, near the centered base the mixed $4$-point branch is locally
exhausted: the clean branch $V_1\neq 0$ is excluded by the centered
good-$v_5$-patch parity package, while every late-$v_5$ branch is either
locally sterile on the $u_5$-patch or already absorbed into the defect-free
planar quintic-circuit regime.
\end{theorem}

\begin{proof}
By Corollary~\ref{cor:right-pincer-local-reduction-centered-base}, every
sufficiently small nearby genuinely new mixed $4$-point obstruction must lie
on the clean branch
\[
V_1\neq 0,
\qquad\text{equivalently}\qquad
\operatorname{ord}_s v_5=1.
\]
On that branch, Corollary~\ref{cor:finite-jet-scalar-descent-centered-branch}
gives
\[
c(s)=C_0+C_2s^2+C_4s^4+\cdots.
\]
Since $v_5(s)$ is odd with $V_1\neq 0$, the induced patch coordinate
\[
x(s)=\frac{v_5(s)}{c(s)}
\]
is odd and has odd inverse. Since $u_5(s)$ is odd and $\Delta_7(s)$ is
even, it follows that
\[
Y(x)=\frac{u_5}{c}(x)
\qquad\text{and}\qquad
S(x)=\frac{\Delta_7}{cv_5}(x)
\]
are odd in $x$. Hence
\[
\mathcal Z(x)=\frac{Y''}{S''}\ \text{is even},
\qquad
K_v=(x')^3S''
\]
has odd vanishing order, so by
Corollary~\ref{cor:mixed-low-order-exclusion-centered-good-v5-patch}
none of the explicit mixed low-order obstruction patterns
\[
(r,\operatorname{ord}_0\mathcal Z)\in\{(1,1),(1,3),(2,2)\}
\]
can occur on the clean branch.

Thus no sufficiently small nearby genuinely new mixed $4$-point obstruction
remains. The late-$v_5$ branches were already removed by
Corollary~\ref{cor:right-pincer-local-reduction-centered-base}, and the clean
branch is excluded as above.
\end{proof}

\begin{remark}[What the mixed local exhaustion theorem now accomplishes]
\label{rem:what-mixed-local-exhaustion-now-accomplishes}
Theorem~\ref{thm:local-exhaustion-mixed-four-point-centered-base} closes the
mixed $4$-point route locally near the centered diagonal base in a
self-contained way: the clean branch is eliminated by the canonical scalar
descent together with the centered good-$v_5$-patch parity package, while
every late-$v_5$ branch is already removed by the $u_5$-patch right-pincer
package as either locally sterile or already-quintic.
\end{remark}

\begin{remark}[What remains after local mixed exhaustion]
\label{rem:what-remains-after-local-mixed-exhaustion}
Theorem~\ref{thm:local-exhaustion-mixed-four-point-centered-base} exhausts the
mixed $4$-point branch in the centered-base local regime. Any remaining mixed
$4$-point escape must therefore be genuinely global/nonlocal: it cannot come
from a new local mixed septic-jet pattern near the centered diagonal base.

Equivalently, the later mixed/global problem is no longer a local normal-form
problem. It is a finite-core / global propagation problem for the normalized
line-incidence geometry, with the residual mixed two-point derivative geometry
concentrated on the overlap exceptional set and, after the off-$K_v=0$
handoffs, sharply focused on the exact fixed-shear $K_v=0$ slice.
\end{remark}

\begin{remark}[How mixed local exhaustion constrains later hybrid routes]
\label{rem:mixed-local-exhaustion-constrains-later-hybrids}
Theorem~\ref{thm:local-exhaustion-mixed-four-point-centered-base} is useful
beyond the mixed local lane itself. It shows that any later hybrid argument
trying to constrain the nonmixed repeated-source / source-coupling branch
cannot appeal to a new local mixed $4$-point septic jet pattern near the
centered diagonal base, because that possibility is already exhausted. Thus any
remaining mixed input capable of affecting the nonmixed branch must be genuinely
global/nonlocal, or must enter through external canonical/multi-pair structure
rather than through a new local mixed obstruction.
\end{remark}

\begin{remark}[A second below-lifted ratio does not add a new obstruction]
\label{rem:second-below-lifted-ratio-no-new-obstruction}
Let
\[
\Theta:=\frac{W_{u_5}}{W},
\]
where $W$ denotes the corresponding projective-chart Wronskian in the
$(c,v_5,\cdot)$-frame. In the gauge $c\equiv 1$, $v_5=x$, one has the
exact relation
\[
\Theta=\mathcal Z\,\frac{S''}{xS''+2S'},
\qquad
(\mathcal Z-x\Theta)S''=2\Theta S'.
\]
So $\Theta$ does not provide an independent local obstruction beyond
$(\mathcal Z,\mathcal R)$: whenever $2\Theta/(\mathcal Z-x\Theta)$ is
analytic, one can solve locally for $S$ and then recover $Y$ from
$Y''=\mathcal Z S''$. In particular, attacks based only on a generic
differential relation between $\mathcal Z$ and $\Theta$ are too weak to
close the $4$-point branch.
\end{remark}

\begin{corollary}[Local defect-free four-point exclusion on lifted ECT subintervals]
\label{cor:local-four-point-exclusion-lifted-vpatch}
Let $J\subset\{c\neq 0,\ M\neq 0,\ v_5\neq 0\}$ be a connected interval on
which $T_v$ has fixed nonzero sign. Then the quadruple
\[
1,\qquad \frac{u_5}{c},\qquad \frac{v_5}{c},\qquad
\frac{\Delta_7}{c\,v_5}
\]
is an ECT system on $J$. Equivalently, no affine plane meets the lifted
space curve
\[
Q_v(h)=\left(\frac{u_5(h)}{c(h)},\frac{v_5(h)}{c(h)},
\frac{\Delta_7(h)}{c(h)\,v_5(h)}\right)
\]
in four points of $J$. Hence the defect-free four-point branch is excluded
on such a subinterval.
\end{corollary}

\begin{proof}
The Wronskian of the displayed quadruple is exactly $T_v$. If $T_v$ has
fixed nonzero sign on $J$, the ECT criterion applies.
\end{proof}

\begin{lemma}[Osculating determinant expansion for shrinking affine circuits]
\label{lem:osculating-determinant-expansion}
Let $Q:I\to \mathbf C^d$ be real-analytic on an interval $I$, and define
\[
\Delta_{d+1}(t_0,\dots,t_d)
:=
\det\!\begin{pmatrix}
1 & Q(t_0)\\
\vdots & \vdots\\
1 & Q(t_d)
\end{pmatrix}.
\]
Fix $h_*\in I$, and let
\[
t_j=h_*+\varepsilon \lambda_j
\qquad (j=0,\dots,d),
\]
where the $\lambda_j$ are pairwise distinct and $\varepsilon\to 0$. Then
\[
\Delta_{d+1}(t_0,\dots,t_d)
=
\frac{\varepsilon^{d(d+1)/2}}{\prod_{k=1}^d k!}\,
V(\lambda_0,\dots,\lambda_d)\,
\det\!\bigl(Q'(h_*),Q''(h_*),\dots,Q^{(d)}(h_*)\bigr)
+
O\!\left(\varepsilon^{d(d+1)/2+1}\right),
\]
where
\[
V(\lambda_0,\dots,\lambda_d):=\prod_{0\le i<j\le d}(\lambda_j-\lambda_i)
\]
is the Vandermonde determinant.

In particular, if there is a sequence of exact affine $(d+1)$-circuits
\[
\Delta_{d+1}(t_{0,n},\dots,t_{d,n})=0
\]
with pairwise distinct nodes $t_{j,n}\to h_*$, then
\[
\det\!\bigl(Q'(h_*),Q''(h_*),\dots,Q^{(d)}(h_*)\bigr)=0.
\]
\end{lemma}

\begin{proof}
Write the Taylor expansion
\[
Q(h_*+\varepsilon\lambda)
=
Q(h_*)+\sum_{m=1}^d \frac{\varepsilon^m\lambda^m}{m!}\,Q^{(m)}(h_*)
+O(\varepsilon^{d+1}),
\]
uniformly for $\lambda$ in bounded sets. Subtract the first row
$(1,Q(t_0))$ from each of the remaining rows. The first column then becomes
$(1,0,\dots,0)^{\top}$, so $\Delta_{d+1}$ reduces to the determinant of the
$d\times d$ matrix whose $j$-th row ($j=1,\dots,d$) is
\[
Q(t_j)-Q(t_0)
=
\sum_{m=1}^d \frac{\varepsilon^m(\lambda_j^m-\lambda_0^m)}{m!}\,Q^{(m)}(h_*)
+O(\varepsilon^{d+1}).
\]
Factor the column vectors $Q'(h_*),\dots,Q^{(d)}(h_*)$. This gives
\[
\Delta_{d+1}(t_0,\dots,t_d)
=
\det\!\bigl(Q'(h_*),\dots,Q^{(d)}(h_*)\bigr)\,
\det\!\Bigl(\frac{\varepsilon^m(\lambda_j^m-\lambda_0^m)}{m!}\Bigr)_{1\le j,m\le d}
+O\!\left(\varepsilon^{d(d+1)/2+1}\right).
\]
Pulling $\varepsilon^m/m!$ from the $m$-th column yields the factor
\[
\frac{\varepsilon^{1+\cdots+d}}{\prod_{k=1}^d k!}
=
\frac{\varepsilon^{d(d+1)/2}}{\prod_{k=1}^d k!}.
\]
The remaining determinant is the standard Vandermonde determinant
$V(\lambda_0,\dots,\lambda_d)$. This proves the expansion.

If $\Delta_{d+1}(t_{0,n},\dots,t_{d,n})=0$ for a sequence of shrinking
distinct-node circuits, divide by the nonzero Vandermonde factor and let
$\varepsilon\to 0$. The leading term must vanish, giving
\[
\det\!\bigl(Q'(h_*),Q''(h_*),\dots,Q^{(d)}(h_*)\bigr)=0.
\qedhere
\]
\end{proof}

\begin{remark}[Planar curvature and lifted torsion as the first derivative obstructions]
\label{rem:derivative-geometry-curvature-torsion-bridge}
Lemma~\ref{lem:osculating-determinant-expansion} shows that shrinking exact
affine circuits are controlled, at first nontrivial order, by derivative-side
osculating determinants. For a planar curve $q$, shrinking exact collinear
triples force
\[
\det(q',q'')=0,
\]
and for a space curve $Q$, shrinking exact coplanar quadruples force
\[
\det(Q',Q'',Q''')=0.
\]
Thus the defect-free lifted $3$-point and $4$-point branches are not merely
incidence problems: they are rooted at curvature- and torsion-type derivative
invariants of the corresponding lifted curves.
\end{remark}

\begin{proposition}[Shrinking exact lifted $3$-point circuits concentrate into $K_v=0$]
\label{prop:shrinking-exact-lifted-three-point-circuits-force-Kv-zero}
Let $P\Subset\{c\neq 0,\ v_5\neq 0,\ M\neq 0\}$ be compact, and use
\[
x=\frac{v_5}{c},
\qquad
S=\sigma_v=\frac{\Delta_7}{c\,v_5}.
\]
Suppose there is a sequence of distinct triples
\[
x_{1,n},x_{2,n},x_{3,n}\in x(P),
\qquad
x_{1,n},x_{2,n},x_{3,n}\to x_*,
\]
such that
\[
\det\!\begin{pmatrix}
1&x_{1,n}&S(x_{1,n})\\
1&x_{2,n}&S(x_{2,n})\\
1&x_{3,n}&S(x_{3,n})
\end{pmatrix}
=0
\qquad
\text{for every }n.
\]
Then
\[
K_v(x_*)=0.
\]
\end{proposition}

\begin{proof}
Apply Lemma~\ref{lem:osculating-determinant-expansion} with $d=2$ to the
planar lifted curve
\[
q(x):=(x,S(x)).
\]
The exact collinearity condition in the statement is precisely the vanishing of
\[
\det\!\begin{pmatrix}
1&x_1&S(x_1)\\
1&x_2&S(x_2)\\
1&x_3&S(x_3)
\end{pmatrix}.
\]
Hence shrinking exact lifted $3$-point circuits force
\[
\det(q'(x_*),q''(x_*))=0.
\]
Since
\[
q'(x)=(1,S'(x)),
\qquad
q''(x)=(0,S''(x)),
\]
one gets
\[
\det(q'(x_*),q''(x_*))=S''(x_*).
\]
On the $x$-patch, Proposition~\ref{prop:exact-x-coordinate-identities-lifted-four-point}
gives
\[
K_v=(x')^3S'',
\]
and $x'\neq 0$ on a good $v_5$-patch. Therefore $S''(x_*)=0$ implies
\[
K_v(x_*)=0.
\qedhere
\]
\end{proof}

\begin{corollary}[Uniform small-diameter exclusion away from $K_v=0$]
\label{cor:uniform-small-diameter-exclusion-away-from-Kv-zero}
Let $P_0\Subset\{c\neq 0,\ v_5\neq 0,\ M\neq 0\}$ be compact and assume
\[
P_0\cap\{K_v=0\}=\varnothing.
\]
Then there exists $\rho_3>0$ such that no exact lifted $3$-point circuit
contained in $P_0$ has $x$-diameter $<\rho_3$.
\end{corollary}

\begin{proof}
Otherwise there would exist a sequence of exact lifted $3$-point circuits in
$P_0$ with diameters tending to $0$. By compactness, after passing to a
subsequence the three nodes would converge to some $x_*\in x(P_0)$. This
contradicts
Proposition~\ref{prop:shrinking-exact-lifted-three-point-circuits-force-Kv-zero}.
\end{proof}

\begin{proposition}[Shrinking exact lifted $4$-point circuits concentrate into $T_v=0$]
\label{prop:shrinking-exact-lifted-four-point-circuits-force-Tv-zero}
Let $P\Subset\{c\neq 0,\ v_5\neq 0,\ M\neq 0\}$ be compact, and use
\[
x=\frac{v_5}{c},
\qquad
Y=\frac{u_5}{c},
\qquad
S=\sigma_v=\frac{\Delta_7}{c\,v_5}.
\]
Suppose there is a sequence of distinct quadruples
\[
x_{1,n},x_{2,n},x_{3,n},x_{4,n}\in x(P),
\qquad
x_{1,n},x_{2,n},x_{3,n},x_{4,n}\to x_*,
\]
such that
\[
\det\!\begin{pmatrix}
1&Y(x_{1,n})&x_{1,n}&S(x_{1,n})\\
1&Y(x_{2,n})&x_{2,n}&S(x_{2,n})\\
1&Y(x_{3,n})&x_{3,n}&S(x_{3,n})\\
1&Y(x_{4,n})&x_{4,n}&S(x_{4,n})
\end{pmatrix}
=0
\qquad
\text{for every }n.
\]
Then
\[
T_v(x_*)=0.
\]
\end{proposition}

\begin{proof}
Apply Lemma~\ref{lem:osculating-determinant-expansion} with $d=3$ to the
lifted space curve
\[
Q(x):=(Y(x),x,S(x)).
\]
The exact coplanarity condition in the statement is precisely the vanishing of
\[
\det\!\begin{pmatrix}
1&Y(x_1)&x_1&S(x_1)\\
1&Y(x_2)&x_2&S(x_2)\\
1&Y(x_3)&x_3&S(x_3)\\
1&Y(x_4)&x_4&S(x_4)
\end{pmatrix}.
\]
Hence shrinking exact lifted $4$-point circuits force
\[
\det(Q'(x_*),Q''(x_*),Q'''(x_*))=0.
\]
But
\[
\det(Q'(x),Q''(x),Q'''(x))
=
W_x(1,Y,x,S)(x),
\]
and Proposition~\ref{prop:exact-x-coordinate-identities-lifted-four-point}
gives
\[
T_v=(x')^6\,W_x(1,Y,x,S).
\]
Since $x'\neq 0$ on a good $v_5$-patch, it follows that
\[
T_v(x_*)=0.
\qedhere
\]
\end{proof}

\begin{corollary}[Uniform small-diameter exclusion away from $T_v=0$]
\label{cor:uniform-small-diameter-exclusion-away-from-Tv-zero}
Let $P_0\Subset\{c\neq 0,\ v_5\neq 0,\ M\neq 0\}$ be compact and assume
\[
P_0\cap\{T_v=0\}=\varnothing.
\]
Then there exists $\rho_4>0$ such that no exact lifted $4$-point circuit
contained in $P_0$ has $x$-diameter $<\rho_4$.
\end{corollary}

\begin{proof}
Otherwise there would exist a sequence of exact lifted $4$-point circuits in
$P_0$ with diameters tending to $0$. By compactness, after passing to a
subsequence the four nodes would converge to some $x_*\in x(P_0)$. This
contradicts
Proposition~\ref{prop:shrinking-exact-lifted-four-point-circuits-force-Tv-zero}.
\end{proof}

\begin{definition}[Repeated-node tangent-contact determinants on a $v_5$-patch]
\label{def:repeated-node-tangent-contact-determinants-vpatch}
On a $v_5\neq 0$ patch with $x=v_5/c$, $Y=u_5/c$, and
$S=\Delta_7/(c\,v_5)$, define
\[
H_3(a,u):=
\det\!\begin{pmatrix}
1&a&S(a)\\
0&1&S'(a)\\
1&u&S(u)
\end{pmatrix}
=
S(u)-S(a)-S'(a)(u-a),
\]
and
\[
H_4(a,u,v):=
\det\!\begin{pmatrix}
1&Y(a)&a&S(a)\\
0&Y'(a)&1&S'(a)\\
1&Y(u)&u&S(u)\\
1&Y(v)&v&S(v)
\end{pmatrix}.
\]
Thus $H_3(a,u)=0$ means that $(u,S(u))$ lies on the tangent line to
$\Lambda_v$ at $a$, while $H_4(a,u,v)=0$ means that $Q_v(u)$ and
$Q_v(v)$ lie in the tangent plane to $Q_v$ at $a$.
\end{definition}

\begin{proposition}[Repeated-node limits of exact lifted determinants]
\label{prop:repeated-node-limits-of-exact-lifted-determinants}
For the exact lifted determinants
\[
\Delta_3(a,b,u):=
\det\!\begin{pmatrix}
1&a&S(a)\\
1&b&S(b)\\
1&u&S(u)
\end{pmatrix},
\]
and
\[
\Delta_4(a,b,u,v):=
\det\!\begin{pmatrix}
1&Y(a)&a&S(a)\\
1&Y(b)&b&S(b)\\
1&Y(u)&u&S(u)\\
1&Y(v)&v&S(v)
\end{pmatrix},
\]
one has
\[
\frac{\Delta_3(a,b,u)}{b-a}\longrightarrow H_3(a,u)
\qquad (b\to a),
\]
and
\[
\frac{\Delta_4(a,b,u,v)}{b-a}\longrightarrow H_4(a,u,v)
\qquad (b\to a).
\]
\end{proposition}

\begin{proof}
Subtract the first row from the second row in each determinant and divide by
$b-a$. The resulting second row is the first difference quotient of the
corresponding row entries, and the limit $b\to a$ gives the derivative row.
\end{proof}

\begin{theorem}[Local classification of the lifted $3$-point branch near a finite-order $K_v$-zero]
\label{thm:local-classification-lifted-three-point-branch-near-Kv-zero}
Let $x=0$ be an isolated finite-order zero of $K_v$ on a good $v_5$-patch.
After affine normalization of $\Lambda_v=(x,S)$, write
\[
S(x)=x^m+O(x^{m+1}),
\qquad
m=r+2\ge 3,
\]
where $r=\operatorname{ord}_0 K_v$. Then the repeated-node tangent-contact equation
\[
H_3(a,u)=0
\]
with $a<0<u$ has the following local classification:
\begin{enumerate}
\item if $r$ is even (equivalently $m$ is even), there is no
opposite-side tangent-contact branch near $(a,u)=(0,0)$;
\item if $r$ is odd (equivalently $m$ is odd), there exists a unique
real-analytic opposite-side branch
\[
u=\Upsilon(a),
\qquad
a<0,\quad |a|\ll 1,
\]
with asymptotic expansion
\[
\Upsilon(a)=\lambda_m(-a)+O(a^2),
\]
where $\lambda_m>1$ is the unique positive root of
\[
P_m(\lambda):=\lambda^m-m\lambda-(m-1).
\]
\end{enumerate}
\end{theorem}

\begin{proof}
Set $a=-\alpha$ with $\alpha>0$, and write $u=\lambda\alpha$ with
$\lambda>0$. For the model germ $S(x)=x^m$, the tangent-contact equation is
\[
H_3(-\alpha,\lambda\alpha)
=
\alpha^m
\begin{cases}
\lambda^m+m\lambda+(m-1),& m\ \text{even},\\[0.5ex]
\lambda^m-m\lambda-(m-1),& m\ \text{odd}.
\end{cases}
\]
If $m$ is even, the displayed factor is strictly positive for every
$\lambda>0$, so no opposite-side branch exists.

If $m$ is odd, define
\[
P_m(\lambda):=\lambda^m-m\lambda-(m-1).
\]
Then
\[
P_m(1)=-2(m-1)<0,
\qquad
P_m(\lambda)\to+\infty\ \text{as }\lambda\to+\infty,
\]
and
\[
P_m'(\lambda)=m(\lambda^{m-1}-1)
\]
is negative on $(0,1)$ and positive on $(1,\infty)$. Therefore $P_m$
has a unique positive root $\lambda_m>1$.

For the perturbed germ $S(x)=x^m+O(x^{m+1})$, divide $H_3(-\alpha,\lambda\alpha)$
by $\alpha^m$. The resulting equation is
\[
P_m(\lambda)+O(\alpha)=0.
\]
If $m$ is even, positivity of the model factor persists for $\alpha$ small,
so there is still no branch. If $m$ is odd, the root $\lambda_m$ is simple,
so the implicit function theorem yields a unique analytic branch
\[
\lambda=\lambda(\alpha)=\lambda_m+O(\alpha),
\]
equivalently
\[
u=\lambda_m(-a)+O(a^2).
\]
This proves the classification.
\end{proof}

\begin{corollary}[One-sided exclusion for repeated-node tangent contact]
\label{cor:one-sided-exclusion-repeated-node-tangent-contact}
Under the hypotheses of
Theorem~\ref{thm:local-classification-lifted-three-point-branch-near-Kv-zero},
there exists $\varepsilon>0$ such that
\[
H_3(a,u)\neq 0
\]
for every same-side pair with
\[
-\varepsilon<a<u<0
\qquad\text{or}\qquad
0<a<u<\varepsilon.
\]
Equivalently, every sufficiently small off-diagonal solution of
\[
H_3(a,u)=0
\]
is necessarily opposite-side.
\end{corollary}

\begin{proof}
Using
\[
H_3(a,u)=S(u)-S(a)-S'(a)(u-a),
\]
one has the exact integral identity
\[
H_3(a,u)=\int_a^u (u-t)\,S''(t)\,dt.
\]
Since $x=0$ is an isolated finite-order zero of $K_v=(x')^3S''$ on a good
$v_5$-patch, the analytic function $S''$ has constant nonzero sign on each
punctured one-sided interval
\[
(-\varepsilon,0),
\qquad
(0,\varepsilon)
\]
after shrinking. If $a<u$ lie on the same side of $0$, then
\[
u-t>0
\]
for all $t\in(a,u)$, and $S''(t)$ has fixed nonzero sign on the whole interval.
So the integrand has fixed nonzero sign, and therefore
\[
H_3(a,u)\neq 0.
\qedhere
\]
\end{proof}

\begin{corollary}[No infinite-order flatness for $K_v$ on a good $v_5$-patch]
\label{cor:no-infinite-order-flatness-Kv-good-vpatch}
On a good $v_5$-patch, every zero of $K_v$ has finite order.
Equivalently, the lifted curvature $S''(x)$ has no infinite-order zero on such a
patch.
\end{corollary}

\begin{proof}
By Proposition~\ref{prop:exact-x-coordinate-identities-lifted-four-point},
\[
K_v=(x')^3S''(x),
\]
and on a good $v_5$-patch one has $x'\neq 0$. Thus $K_v$ and $S''$ have the
same zero set and the same vanishing order.

If $K_v$ had an infinite-order zero at some interior point $x_*$, then so would
$S''$. Since $S''$ is analytic in the local coordinate $x$, this would force
\[
S''\equiv 0
\]
on a neighborhood of $x_*$. Hence $S$ would be affine in $x$ there, contrary to
Proposition~\ref{prop:sigmav-not-affine-in-x}.
\end{proof}

\begin{lemma}[Slope--intercept parametrization for the lifted $3$-point branch]
\label{lem:slope-intercept-parametrization-lifted-three-point}
Assume the affine normalization from
Theorem~\ref{thm:local-classification-lifted-three-point-branch-near-Kv-zero},
so that
\[
S(0)=S'(0)=0,
\qquad
S''<0\ \text{on the source side }I_-=(\alpha,0).
\]
Let
\[
I_+=(0,\beta)
\]
be the maximal open interval to the right of $0$ on which $S'$ is strictly
increasing. Equivalently, $S''$ does not change sign on $I_+$, and any zeros of
$S''$ there are isolated even-order zeros.

For
\[
m\in J_-:=S'(I_-),
\]
define the inverse source point
\[
a=a(m):=(S'|_{I_-})^{-1}(m),
\]
and the source tangent intercept
\[
\tau_-(m):=S(a(m))-m\,a(m).
\]
For
\[
m\in J_+:=S'(I_+),
\]
define the inverse target tangent point
\[
u_+(m):=(S'|_{I_+})^{-1}(m),
\]
and the target tangent intercept
\[
\tau_+(m):=S(u_+(m))-m\,u_+(m).
\]
Then
\[
\tau_-'(m)=-a(m),
\qquad
\tau_+'(m)=-u_+(m).
\]
Hence on
\[
J:=J_-\cap J_+
\]
the intercept gap
\[
D(m):=\tau_-(m)-\tau_+(m)
\]
satisfies
\[
D'(m)=u_+(m)-a(m)>0.
\]
\end{lemma}

\begin{proof}
Differentiate
\[
\tau_-(m)=S(a(m))-m\,a(m)
\]
to obtain
\[
\tau_-'(m)=S'(a(m))a'(m)-a(m)-m\,a'(m)=-a(m),
\]
because $S'(a(m))=m$. The same computation gives
\[
\tau_+'(m)=-u_+(m).
\]
Subtracting yields
\[
D'(m)=u_+(m)-a(m).
\]
Since $a(m)<0<u_+(m)$, the derivative is strictly positive.
\end{proof}

\begin{theorem}[Global continuation / nonrecurrence of the odd-order lifted $3$-point branch]
\label{thm:global-continuation-nonrecurrence-odd-three-point-branch}
Assume the hypotheses and affine normalization of
Theorem~\ref{thm:local-classification-lifted-three-point-branch-near-Kv-zero},
and suppose that $r=\operatorname{ord}_0K_v$ is odd, so the unique local
opposite-side branch exists.

Let
\[
I_-=(\alpha,0)
\]
be the source interval on which $S''<0$, and let
\[
I_+=(0,\beta)
\]
be the maximal open interval to the right of $0$ on which $S'$ is strictly
increasing. Set
\[
J:=S'(I_-)\cap S'(I_+).
\]
Then there exists a unique maximal open interval
\[
J_{\mathrm{br}}=(0,m_*)\subset J
\]
and a unique real-analytic function
\[
u:J_{\mathrm{br}}\to I_+
\]
such that
\[
H_3(a(m),u(m))=0,
\qquad
a(m):=(S'|_{I_-})^{-1}(m),
\qquad
u(m)>0,
\]
and for small $m>0$ this branch agrees with the local branch from
Theorem~\ref{thm:local-classification-lifted-three-point-branch-near-Kv-zero}.

Moreover:
\begin{enumerate}
\item the intercept gap
\[
D(m):=\tau_-(m)-\tau_+(m)
\]
is strictly positive on all of $J$, and
\[
D'(m)=u_+(m)-a(m)>0;
\]

\item along the branch one has
\[
\partial_u H_3(a(m),u(m))
=
S'(u(m))-m
\neq 0,
\]
so there is no interior common-tangent event;

\item the branch is uniquely continued on $J_{\mathrm{br}}$, with
\[
u'(m)=\frac{u(m)-a(m)}{S'(u(m))-m},
\]
and the source coordinate is strictly monotone,
\[
a'(m)=\frac{1}{S''(a(m))}<0;
\]

\item if $m_*<\sup J$, then
\[
u(m)\longrightarrow \beta
\qquad (m\uparrow m_*).
\]
Hence the branch cannot end in the interior of $I_+$; it can stop only when it
reaches the boundary of the source/target slope-overlap block.

In particular, the odd-order lifted $3$-point branch has no interior
recurrence, no self-return, no interior accumulation, and no infinite-order
flatness. Its only remaining endpoint mechanism is boundary exit of the
relevant good-patch / monotone-slope block.
\end{enumerate}
\end{theorem}

\begin{proof}
For
\[
m\in J,
\]
set
\[
\Phi(m,u):=S(u)-m\,u-\tau_-(m).
\]
Then
\[
\Phi(m,u)=H_3(a(m),u).
\]

\smallskip

\noindent
\textbf{Step 1: positivity of the intercept gap.}
By the odd-order local classification theorem, for all sufficiently small
$m>0$ there exists a point $u(m)\in I_+$ such that
\[
\Phi(m,u(m))=0.
\]
Fix such a small $m$. Since $S'$ is strictly increasing on $I_+$, the function
\[
u\longmapsto \Phi(m,u)
\]
has strictly increasing derivative
\[
\partial_u\Phi(m,u)=S'(u)-m,
\]
hence a unique critical point, namely $u_+(m)$, and that critical point is its
unique minimum on $I_+$. Therefore
\[
\Phi(m,u_+(m))<0.
\]
But
\[
\Phi(m,u_+(m))
=
S(u_+(m))-m\,u_+(m)-\tau_-(m)
=
\tau_+(m)-\tau_-(m)
=
-D(m),
\]
so $D(m)>0$ for all sufficiently small $m>0$.

By Lemma~\ref{lem:slope-intercept-parametrization-lifted-three-point},
\[
D'(m)=u_+(m)-a(m)>0
\]
on all of $J$. Hence $D(m)>0$ for every $m\in J$.

\smallskip

\noindent
\textbf{Step 2: no interior common tangency.}
Suppose
\[
\Phi(m,u)=0
\]
for some $m\in J$ and $u\in I_+$. If also
\[
\partial_u\Phi(m,u)=S'(u)-m=0,
\]
then by uniqueness of the inverse point on $I_+$ one has $u=u_+(m)$, and thus
\[
0=\Phi(m,u_+(m))=-D(m),
\]
contrary to $D(m)>0$. Therefore every branch point satisfies
\[
\partial_u\Phi(m,u)\neq 0.
\]
Equivalently, there is no interior common-tangent point on the odd-order
branch.

\smallskip

\noindent
\textbf{Step 3: uniqueness and analytic continuation.}
For fixed $m\in J$, the derivative
\[
u\longmapsto \partial_u\Phi(m,u)=S'(u)-m
\]
is strictly increasing on $I_+$, and
\[
\partial_u\Phi(m,0)=S'(0)-m=-m<0.
\]
Hence $\Phi(m,\cdot)$ first decreases and then increases, so it has at most one
zero on $I_+$. Thus the local branch from
Theorem~\ref{thm:local-classification-lifted-three-point-branch-near-Kv-zero}
is unique on every connected continuation interval.

Let
\[
E:=\{m\in J:\exists\,u\in I_+\ \text{with}\ \Phi(m,u)=0\}.
\]
By Step 2, every zero is simple in the $u$-variable, so the implicit function
theorem shows that $E$ is open and that the zero depends real-analytically on
$m$. Let
\[
J_{\mathrm{br}}=(0,m_*)
\]
be the connected component of $E$ containing all sufficiently small positive
$m$. This gives the unique maximal analytic continuation
\[
u=u(m),
\qquad
m\in J_{\mathrm{br}}.
\]

Differentiating
\[
\Phi(m,u(m))=0
\]
gives
\[
\bigl(S'(u(m))-m\bigr)u'(m)-u(m)+a(m)=0,
\]
hence
\[
u'(m)=\frac{u(m)-a(m)}{S'(u(m))-m}.
\]
Also
\[
a'(m)=\frac{1}{S''(a(m))}<0
\]
because $S''<0$ on the source interval $I_-$. So the source coordinate is
strictly monotone and cannot return to a previous value.

\smallskip

\noindent
\textbf{Step 4: endpoint control.}
Assume that
\[
m_*<\sup J
\]
but that $u(m)$ does not approach $\beta$ as $m\uparrow m_*$. Then there exists
a compact subinterval
\[
K\Subset I_+
\]
and a sequence $m_n\uparrow m_*$ with
\[
u(m_n)\in K.
\]
After passing to a subsequence,
\[
u(m_n)\to u_*\in K.
\]
By continuity,
\[
\Phi(m_*,u_*)=0.
\]
Since $m_*<\sup J$, one has $m_*\in J$, so by Step 2,
\[
\partial_u\Phi(m_*,u_*)\neq 0.
\]
The implicit function theorem therefore extends the branch analytically past
$m_*$, contradicting maximality of $J_{\mathrm{br}}$.

Hence
\[
u(m)\to\beta
\qquad (m\uparrow m_*).
\]
This proves that the branch cannot terminate by an interior fold, interior
common tangency, or interior accumulation.

Finally, by
Corollary~\ref{cor:no-infinite-order-flatness-Kv-good-vpatch}, every zero of
$K_v$ has finite order. Since even-order zeros do not change the sign of $S''$,
they do not break the strict increase of $S'$ on $I_+$. Therefore the endpoint
$\beta$ can only be the ambient boundary of the good patch or the first
odd-order zero of $K_v$ to the right.
\end{proof}

\begin{remark}[Structural package for the global lifted $3$-point branch]
\label{rem:structural-package-global-lifted-three-point-branch}
The correct global parameter on the odd-order lifted $3$-point branch is the
source tangent slope
\[
m=S'(a).
\]
In that parameter, the branch is governed by the source/target tangent
intercepts
\[
\tau_-(m)=S(a(m))-m\,a(m),
\qquad
\tau_+(m)=S(u_+(m))-m\,u_+(m),
\]
and their gap
\[
D(m)=\tau_-(m)-\tau_+(m),
\qquad
D'(m)=u_+(m)-a(m)>0.
\]
Thus interior common tangency is exactly the equation
\[
D(m)=0,
\]
but this never occurs on the odd-order continuation branch because $D>0$ there.
So the global $3$-point problem is not a hidden recurrence problem inside the
patch: once born at an odd-order $K_v$-zero, the branch is a one-parameter
analytic graph with monotone source coordinate, and even-order $K_v$-zeros are
invisible to it except as isolated points where $S''$ vanishes without changing
the monotone-slope block.
\end{remark}

\begin{remark}[Target-block form of the intercept-gap monotonicity]
\label{rem:target-block-intercept-gap-monotonicity}
The slope--intercept monotonicity from
Lemma~\ref{lem:slope-intercept-parametrization-lifted-three-point}
is not special to the first right-hand target block.

Let $I$ be any open interval to the right of the source interval $I_-$
on which $S'$ is monotone and hence invertible. For
\[
m\in S'(I_-)\cap S'(I),
\]
define
\[
u_I(m):=(S'|_I)^{-1}(m),
\qquad
\tau_I(m):=S(u_I(m))-m\,u_I(m),
\]
and
\[
D_I(m):=\tau_-(m)-\tau_I(m).
\]
Then exactly the same envelope calculation gives
\[
\tau_I'(m)=-u_I(m),
\qquad
D_I'(m)=u_I(m)-a(m)>0.
\]

Thus for any monotone-slope target block, interior common tangency against the
fixed source interval $I_-$ is again the scalar equation
\[
D_I(m)=0.
\]
So any future continuation across successive odd-order $K_v$-zeros must be
organized blockwise through these monotone intercept-gap functions rather than
through a hidden recurrence mechanism inside a block.
\end{remark}

\begin{remark}[Derivative-geometric meaning of the lifted $3$-point branch]
\label{rem:derivative-geometry-meaning-global-three-point-branch}
The lifted $3$-point branch is now best viewed as a derivative-geometric
propagation theorem for the planar lifted curve
\[
\Lambda_v=(x,S).
\]

By Lemma~\ref{lem:osculating-determinant-expansion} and
Proposition~\ref{prop:shrinking-exact-lifted-three-point-circuits-force-Kv-zero},
shrinking exact lifted $3$-point circuits are detected by the osculating
determinant and therefore concentrate into the lifted curvature-zero set
\[
K_v=(x')^3S''(x)=0.
\]
Theorem~\ref{thm:local-classification-lifted-three-point-branch-near-Kv-zero}
then identifies the local repeated-node branch structure at those curvature
zeros, and
Theorem~\ref{thm:global-continuation-nonrecurrence-odd-three-point-branch}
shows that the odd-order branch propagates globally on each monotone-slope
target block by the slope--intercept envelope identities
\[
\tau_-'(m)=-a(m),
\qquad
\tau_+'(m)=-u_+(m),
\qquad
D'(m)=u_+(m)-a(m)>0.
\]

So the lifted $3$-point theory is no longer an ad hoc repeated-node analysis.
It is the global propagation of an osculating-contact defect for the curve
$\Lambda_v$, organized by the derivative geometry of its tangent lines and
curvature.
\end{remark}

\begin{remark}[Exact nonmixed reduction in terms of $\mathcal Z$ and $\mathcal R$]
\label{rem:nonmixed-four-point-zeros-can-carry-preexisting-branches}
At a nonmixed lifted point
\[
K_v(0)\neq 0,
\qquad
T_v(0)=0,
\]
Propositions~\ref{prop:exact-x-coordinate-identities-lifted-four-point}
and~\ref{prop:canonical-Z-R-invariants-lifted-four-point} give
\[
\mathcal Z=\frac{Y''}{S''},
\qquad
\mathcal R=\mathcal Z'=\frac{T_v}{K_v^2}.
\]
Since $K_v(0)\neq 0$, one has $S''(0)\neq 0$, so after analytic
reparametrization and affine target normalization preserving coplanarity one may
write the lifted germ as
\[
Q_v(\xi)=\bigl(\widetilde Y(\xi),\,\xi,\,\xi^2\bigr),
\qquad
\widetilde Y(0)=\widetilde Y'(0)=\widetilde Y''(0)=0.
\]
In these normalized coordinates,
\[
\mathcal Z(\xi)=\frac12\,\widetilde Y''(\xi),
\qquad
\mathcal R(\xi)=\frac12\,\widetilde Y'''(\xi).
\]

Thus the nonmixed lifted $4$-point problem is exactly the jet problem of
excluding odd-order zeros of $\mathcal R$. Equivalently, it is the problem of
excluding even leading degrees of $\widetilde Y$ in the normalized chart.

In particular, the low-order dangerous nonmixed cases are exactly:
\[
\operatorname{ord}_0\mathcal R=1
\iff
\widetilde Y'''(0)=0,\ \widetilde Y^{(4)}(0)\neq 0,
\]
\[
\operatorname{ord}_0\mathcal R=3
\iff
\widetilde Y^{(3)}(0)=\widetilde Y^{(4)}(0)=\widetilde Y^{(5)}(0)=0,\
\widetilde Y^{(6)}(0)\neq 0,
\]
\[
\operatorname{ord}_0\mathcal R=5
\iff
\widetilde Y^{(3)}(0)=\cdots=\widetilde Y^{(7)}(0)=0,\
\widetilde Y^{(8)}(0)\neq 0.
\]

Equivalently, the dangerous nonmixed low-order list is
\[
\operatorname{ord}_0\mathcal Z\in\{2,4,6\}
\qquad\Longleftrightarrow\qquad
\operatorname{ord}_0\mathcal R\in\{1,3,5\}.
\]

So the nonmixed $4$-point endgame has been reduced to a finite jet problem
for the canonical invariant $\mathcal R$.
\end{remark}

\begin{lemma}[Repeated-node tangent-plane determinant in the normalized nonmixed chart]
\label{lem:normalized-nonmixed-H4-divided-difference}
In the normalized nonmixed chart of
Remark~\ref{rem:nonmixed-four-point-zeros-can-carry-preexisting-branches},
define
\[
\widetilde H_4(a,u,v):=
\det\!\begin{pmatrix}
1&\widetilde Y(a)&a&a^2\\
0&\widetilde Y'(a)&1&2a\\
1&\widetilde Y(u)&u&u^2\\
1&\widetilde Y(v)&v&v^2
\end{pmatrix}.
\]
Then
\[
\widetilde H_4(a,u,v)
=
(u-a)^2(v-a)^2\Bigl(\widetilde Y[a,a,v]-\widetilde Y[a,a,u]\Bigr),
\]
where
\[
\widetilde Y[a,a,t]
:=
\frac{\widetilde Y(t)-\widetilde Y(a)-\widetilde Y'(a)(t-a)}{(t-a)^2}
\]
is the repeated-node second divided difference.
Equivalently, for $u\neq v$,
\[
\widetilde H_4(a,u,v)
=
(u-a)^2(v-a)^2(v-u)\,\widetilde Y[a,a,u,v].
\]
\end{lemma}

\begin{proof}
Subtract the first row from the third and fourth rows. Since
\[
u^2-a^2=(u-a)(u+a),
\qquad
v^2-a^2=(v-a)(v+a),
\]
this gives
\[
\widetilde H_4(a,u,v)
=
(u-a)(v-a)\,
\det\!\begin{pmatrix}
1&\widetilde Y(a)&a&a^2\\
0&\widetilde Y'(a)&1&2a\\
0&\dfrac{\widetilde Y(u)-\widetilde Y(a)}{u-a}&1&u+a\\
0&\dfrac{\widetilde Y(v)-\widetilde Y(a)}{v-a}&1&v+a
\end{pmatrix}.
\]
Expanding along the first column yields
\[
\widetilde H_4(a,u,v)
=
(u-a)(v-a)\,
\det\!\begin{pmatrix}
\widetilde Y'(a)&1&2a\\
\dfrac{\widetilde Y(u)-\widetilde Y(a)}{u-a}&1&u+a\\
\dfrac{\widetilde Y(v)-\widetilde Y(a)}{v-a}&1&v+a
\end{pmatrix}.
\]
Now subtract the first row from the second and third rows:
\[
\widetilde H_4(a,u,v)
=
(u-a)(v-a)\,
\det\!\begin{pmatrix}
\widetilde Y'(a)&1&2a\\
(u-a)\widetilde Y[a,a,u]&0&u-a\\
(v-a)\widetilde Y[a,a,v]&0&v-a
\end{pmatrix}.
\]
Expanding along the middle column gives
\[
\widetilde H_4(a,u,v)
=
(u-a)^2(v-a)^2\Bigl(\widetilde Y[a,a,v]-\widetilde Y[a,a,u]\Bigr),
\]
which is the first formula. Since
\[
\widetilde Y[a,a,u,v]
=
\frac{\widetilde Y[a,a,v]-\widetilde Y[a,a,u]}{v-u},
\]
the second formula follows for $u\neq v$.
\end{proof}

\begin{proposition}[Exact nonmixed tangent-plane contacts force nearby zeros of $\mathcal R$]
\label{prop:exact-nonmixed-tangent-plane-contacts-force-R-zero}
In the normalized nonmixed chart of
Remark~\ref{rem:nonmixed-four-point-zeros-can-carry-preexisting-branches},
assume $a,u,v\in\mathbf R$ are distinct and satisfy
\[
\widetilde H_4(a,u,v)=0.
\]
Then there exists
\[
w\in \operatorname{conv}\{a,u,v\}
\]
such that
\[
\mathcal R(w)=0.
\]
Equivalently,
\[
\widetilde Y'''(w)=0.
\]
\end{proposition}

\begin{proof}
By Lemma~\ref{lem:normalized-nonmixed-H4-divided-difference}, since
$a,u,v$ are distinct,
\[
0=\widetilde H_4(a,u,v)
=
(u-a)^2(v-a)^2(v-u)\,\widetilde Y[a,a,u,v].
\]
Hence
\[
\widetilde Y[a,a,u,v]=0.
\]
By the Hermite--Genocchi formula for repeated-node divided differences, there
exists
\[
w\in \operatorname{conv}\{a,u,v\}
\]
such that
\[
\widetilde Y[a,a,u,v]=\frac{\widetilde Y'''(w)}{3!}.
\]
Therefore
\[
\widetilde Y'''(w)=0.
\]
In the normalized nonmixed chart one has
\[
\mathcal R(\xi)=\frac12\,\widetilde Y'''(\xi),
\]
so $\mathcal R(w)=0$.
\end{proof}

\begin{corollary}[Sign-definite nonmixed intervals are exactly-contact sterile]
\label{cor:sign-definite-R-excludes-exact-nonmixed-contacts}
Let $J$ be an open interval in the normalized nonmixed chart such that
\[
\mathcal R \text{ has fixed nonzero sign on }J.
\]
Then there do not exist distinct points $a,u,v\in J$ such that
\[
\widetilde H_4(a,u,v)=0.
\]
Equivalently, there are no exact nonmixed tangent-plane contacts with all three
nodes in $J$.
\end{corollary}

\begin{proof}
If such $a,u,v$ existed, Proposition~\ref{prop:exact-nonmixed-tangent-plane-contacts-force-R-zero}
would produce some
\[
w\in \operatorname{conv}\{a,u,v\}\subset J
\]
with $\mathcal R(w)=0$, contradicting the sign assumption.
\end{proof}

\begin{remark}[Derivative-side interpretation of exact nonmixed contacts]
\label{rem:exact-nonmixed-contacts-are-R-zero-forcing}
Proposition~\ref{prop:exact-nonmixed-tangent-plane-contacts-force-R-zero}
strengthens the shrinking-circuit statement
Proposition~\ref{prop:shrinking-exact-lifted-four-point-circuits-force-Tv-zero}.
It is not merely a limit theorem: every exact nonmixed tangent-plane contact
already forces a zero of the canonical derivative invariant inside the convex
hull of the three nodes,
\[
\mathcal R=\frac12\,\widetilde Y'''.
\]
Thus the nonmixed $4$-point branch is controlled not only asymptotically by
$T_v=0$, but pointwise by the zero set of a single analytic scalar.
\end{remark}

\begin{proposition}[Interval-persistent nonmixed $4$-point degeneracy forces affine planarity]
\label{prop:interval-persistent-nonmixed-four-point-forces-affine-planarity}
Let $J$ be a connected open interval in a good $v_5$-patch such that
\[
K_v\neq 0
\qquad\text{on }J.
\]
Assume
\[
T_v\equiv 0
\qquad\text{on }J.
\]
Then on $J$ one has
\[
\mathcal R\equiv 0,
\qquad
\mathcal Z\equiv \gamma
\]
for some constant $\gamma\in\mathbf C$, and hence
\[
Y=\gamma S+ax+b
\]
for some constants $a,b\in\mathbf C$. Equivalently, the lifted curve
\[
x\longmapsto (Y(x),x,S(x))
\]
lies in an affine plane on $J$.

In particular, any interval-persistent nonmixed exact lifted $4$-point branch
collapses to an affine-planar residual rather than remaining genuinely
three-dimensional.
\end{proposition}

\begin{proof}
On the nonmixed interval $J$, Proposition~\ref{prop:canonical-Z-R-invariants-lifted-four-point}
gives
\[
\mathcal R=\frac{T_v}{K_v^2}.
\]
Since $T_v\equiv 0$ and $K_v\neq 0$ on $J$, it follows that
\[
\mathcal R\equiv 0
\qquad\text{on }J.
\]
By definition,
\[
\mathcal R=\frac{d\mathcal Z}{dx},
\]
so $\mathcal Z$ is constant on $J$:
\[
\mathcal Z\equiv \gamma.
\]
Using
\[
\mathcal Z=\frac{Y''}{S''},
\]
we obtain
\[
Y''=\gamma S''.
\]
Integrating twice with respect to $x$ yields
\[
Y=\gamma S+ax+b
\]
for constants $a,b\in\mathbf C$. Therefore the lifted curve satisfies the
affine linear equation
\[
Y-\gamma S-ax-b=0
\]
throughout $J$, so it lies in an affine plane.
\end{proof}

\begin{corollary}[Recurrent exact nonmixed contacts force affine-planar collapse]
\label{cor:recurrent-exact-nonmixed-contacts-force-affine-planarity}
Let $J$ be a connected open interval in a good $v_5$-patch such that
\[
K_v\neq 0
\qquad\text{on }J.
\]
Assume that every nonempty open subinterval $I\subset J$ contains an exact
nonmixed tangent-plane contact triple. Equivalently, after local nonmixed
normalization on $I$, there exist distinct points $a,u,v\in I$ such that
\[
\widetilde H_4(a,u,v)=0.
\]
Then
\[
\mathcal R\equiv 0
\qquad\text{on }J,
\]
hence
\[
T_v\equiv 0
\qquad\text{on }J,
\]
and therefore the lifted curve is affine-planar on $J$:
\[
Y=\gamma S+ax+b
\]
for some constants $\gamma,a,b\in\mathbf C$.
\end{corollary}

\begin{proof}
Let $I\subset J$ be any nonempty open subinterval. By assumption, $I$
contains an exact nonmixed tangent-plane contact triple. After nonmixed
normalization on $I$, Proposition~\ref{prop:exact-nonmixed-tangent-plane-contacts-force-R-zero}
produces a point $w_I\in I$ such that
\[
\mathcal R(w_I)=0.
\]
Thus every nonempty open subinterval of $J$ contains a zero of $\mathcal R$.
Since $\mathcal R$ is analytic, its zero set cannot have an accumulation point
in $J$ unless
\[
\mathcal R\equiv 0
\qquad\text{on }J.
\]
Because $K_v\neq 0$ on $J$ and
\[
\mathcal R=\frac{T_v}{K_v^2},
\]
it follows that
\[
T_v\equiv 0
\qquad\text{on }J.
\]
Now apply
Proposition~\ref{prop:interval-persistent-nonmixed-four-point-forces-affine-planarity}
to conclude that
\[
Y=\gamma S+ax+b
\]
on $J$.
\end{proof}

\begin{lemma}[Derivative formula for the repeated-node nonmixed profile]
\label{lem:Da-derivative-formula}
In the normalized nonmixed chart, define
\[
D_a(t):=\widetilde Y[a,a,t]
=
\frac{\widetilde Y(t)-\widetilde Y(a)-\widetilde Y'(a)(t-a)}{(t-a)^2},
\qquad
D_a(a):=\frac12\,\widetilde Y''(a).
\]
Then $D_a(t)$ extends real-analytically across $t=a$, and
\[
D_a(t)=\int_0^1 (1-s)\,\widetilde Y''(a+s(t-a))\,ds.
\]
Moreover,
\[
D_a'(t)
=
\int_0^1 s(1-s)\,\widetilde Y'''(a+s(t-a))\,ds
=
2\int_0^1 s(1-s)\,\mathcal R(a+s(t-a))\,ds.
\]
\end{lemma}

\begin{proof}
Taylor's integral formula gives
\[
\widetilde Y(t)
=
\widetilde Y(a)+\widetilde Y'(a)(t-a)
+(t-a)^2\int_0^1 (1-s)\,\widetilde Y''(a+s(t-a))\,ds.
\]
Dividing by $(t-a)^2$ proves the first formula for $t\neq a$, and the
analytic extension across $t=a$ is immediate. Differentiating under the
integral sign gives
\[
D_a'(t)
=
\int_0^1 s(1-s)\,\widetilde Y'''(a+s(t-a))\,ds.
\]
Since $\mathcal R=\frac12\,\widetilde Y'''$ in the normalized nonmixed chart,
the last formula follows.
\end{proof}

\begin{proposition}[Parity classification of isolated nonmixed zeros of $\mathcal R$]
\label{prop:parity-classification-isolated-nonmixed-R-zeros}
Let $\xi_0\in\mathbf R$ be an isolated zero of the nonmixed invariant
\[
\mathcal R(\xi)=\frac12\,\widetilde Y'''(\xi),
\]
and write
\[
\mathcal R(\xi)=(\xi-\xi_0)^m\,\rho(\xi),
\qquad
m\ge 1,
\qquad
\rho(\xi_0)\neq 0.
\]

\smallskip

\noindent
\emph{(i) Even-order sterility.}
If $m$ is even, then there exists a neighborhood $U$ of $\xi_0$ such that
for every $a\in U$, the function
\[
D_a(t):=\widetilde Y[a,a,t]
\]
is strictly monotone on $U$. In particular, there do not exist distinct
$a,u,v\in U$ such that
\[
\widetilde H_4(a,u,v)=0.
\]

\smallskip

\noindent
\emph{(ii) Odd-order birth.}
If $m$ is odd, then there exists a neighborhood $U$ of $\xi_0$ and a
real-analytic function $\tau:U\to U$ such that
\[
\tau(\xi_0)=\xi_0,
\qquad
\tau'(\xi_0)=-1,
\qquad
D_a'(\tau(a))=0
\quad (a\in U).
\]
Moreover,
\[
D_a''(\tau(a))\neq 0
\quad (a\in U\setminus\{\xi_0\}),
\]
so for each $a\in U\setminus\{\xi_0\}$, the function $D_a$ has a strict
local extremum at $t=\tau(a)$. Consequently, for each such $a$, there exist
distinct points $u_-(a),u_+(a)\in U$ with
\[
u_-(a)<\tau(a)<u_+(a),
\qquad
D_a(u_-(a))=D_a(u_+(a)),
\]
and hence
\[
\widetilde H_4(a,u_-(a),u_+(a))=0.
\]
\end{proposition}

\begin{proof}
By translation in the normalized chart, we may assume $\xi_0=0$. Thus
\[
\mathcal R(\xi)=\xi^m\rho(\xi),
\qquad
\rho(0)\neq 0.
\]

For the even-order case, shrink to a neighborhood $U$ on which $\rho$ has
fixed nonzero sign. Since $m$ is even, $\xi^m\ge 0$, so $\mathcal R$ has
fixed sign on $U$, vanishing only at $0$. By
Lemma~\ref{lem:Da-derivative-formula},
\[
D_a'(t)=2\int_0^1 s(1-s)\,\mathcal R(a+s(t-a))\,ds.
\]
For $a\neq t$, the affine segment $a+s(t-a)$ is not identically zero, so the
integrand has fixed sign and is not identically zero. Hence $D_a'(t)$ has
fixed nonzero sign for all $a\neq t$ in $U$, and therefore each $D_a$ is
strictly monotone on $U$. Thus $D_a$ is injective on $U$, and by
Lemma~\ref{lem:normalized-nonmixed-H4-divided-difference} there are no exact
contacts $\widetilde H_4(a,u,v)=0$ with distinct $a,u,v\in U$.

Now assume $m$ is odd. Write $a=h$ and $t=\lambda h$. Then
Lemma~\ref{lem:Da-derivative-formula} gives
\[
D_h'(\lambda h)
=
2h^m\,G(h,\lambda),
\]
where
\[
G(h,\lambda)
:=
\int_0^1 s(1-s)\bigl(1+s(\lambda-1)\bigr)^m
\rho\!\bigl(h(1+s(\lambda-1))\bigr)\,ds.
\]
This is real-analytic in $(h,\lambda)$, and
\[
G(0,\lambda)=\rho(0)\,J_m(\lambda),
\qquad
J_m(\lambda):=\int_0^1 s(1-s)\bigl(1+s(\lambda-1)\bigr)^m\,ds.
\]
Since $m$ is odd,
\[
J_m(-1)=\int_0^1 s(1-s)(1-2s)^m\,ds=0
\]
by the symmetry $s\mapsto 1-s$, while
\[
J_m'(-1)
=
m\int_0^1 s^2(1-s)(1-2s)^{m-1}\,ds>0
\]
because $m-1$ is even. Thus
\[
\partial_\lambda G(0,-1)=\rho(0)J_m'(-1)\neq 0.
\]
By the implicit function theorem, there exists a unique real-analytic function
$\lambda(h)$ near $0$ such that
\[
\lambda(0)=-1,
\qquad
G(h,\lambda(h))=0.
\]
Define
\[
\tau(h):=\lambda(h)h.
\]
Then
\[
\tau(0)=0,
\qquad
\tau'(0)=\lambda(0)=-1,
\]
and
\[
D_h'(\tau(h))=0.
\]

It remains to show that this critical point is a strict local extremum. By
Lemma~\ref{lem:Da-derivative-formula},
\[
D_a''(t)=2\int_0^1 s^2(1-s)\,\mathcal R'(a+s(t-a))\,ds.
\]
Substituting $a=h$, $t=\lambda(h)h$, and using
\[
\mathcal R'(\xi)=m\xi^{m-1}\rho(\xi)+\xi^m\rho'(\xi),
\]
one obtains
\[
D_h''(\tau(h))
=
2h^{m-1}\rho(0)\,J_m'(-1)+O(h^m).
\]
Since $\rho(0)J_m'(-1)\neq 0$, it follows that
\[
D_h''(\tau(h))\neq 0
\]
for $h\neq 0$ sufficiently small. Therefore $t=\tau(h)$ is a strict local
extremum of $D_h$. Hence there exist distinct points $u_-(h)$, $u_+(h)$
arbitrarily close to $\tau(h)$ such that
\[
D_h(u_-(h))=D_h(u_+(h)).
\]
By Lemma~\ref{lem:normalized-nonmixed-H4-divided-difference},
\[
\widetilde H_4(h,u_-(h),u_+(h))
=
(h-u_-(h))^2(h-u_+(h))^2
\bigl(D_h(u_+(h))-D_h(u_-(h))\bigr)=0.
\]
This proves the odd-order birth statement.
\end{proof}

\begin{remark}[Reflection law for odd-order birth]
\label{rem:reflection-law-odd-order-R-zero}
In the odd-order case of
Proposition~\ref{prop:parity-classification-isolated-nonmixed-R-zeros},
the critical-point branch satisfies
\[
\tau(a)=2\xi_0-a+O((a-\xi_0)^2).
\]
Thus the local exact nonmixed contact family born at any odd-order zero of
$\mathcal R$ has the same first-order reflection geometry across that zero as
in the simple-zero case.
\end{remark}

\begin{proposition}[Odd-zero critical branch is uniformly unrooted]
\label{prop:odd-zero-critical-branch-uniformly-unrooted}
Let $\xi_0$ be an isolated odd-order zero of
\[
\mathcal R(\xi)=\frac12\,\widetilde Y'''(\xi),
\]
and, after translation, write
\[
\xi_0=0,
\qquad
\mathcal R(\xi)=\xi^m\rho(\xi),
\qquad
m\ge 1 \text{ odd},
\qquad
\rho(0)\neq 0.
\]
Let
\[
D_a(t):=\widetilde Y[a,a,t],
\qquad
D_a(a)=\frac12\,\widetilde Y''(a)=\mathcal Z(a).
\]

Then there exists $\varepsilon>0$ and a unique real-analytic function
\[
\mu:(-\varepsilon,\varepsilon)\to\mathbf R,
\qquad
\mu(0)=0,
\]
such that for every sufficiently small $h>0$ the unique nontrivial critical pair
\[
a(h):=h(\mu(h)-1),
\qquad
x(h):=h(\mu(h)+1)
\]
satisfies
\[
D_{a(h)}'(x(h))=0,
\]
and every critical pair
\[
a<x,\qquad |a|+|x|<\varepsilon,\qquad D_a'(x)=0
\]
is of this form.

Moreover,
\[
D_{a(h)}(x(h))-D_{a(h)}(a(h))
=
-\frac{\rho(0)}{m+2}\,h^{m+1}+O(h^{m+2}).
\]
Equivalently,
\[
D_{a(h)}(x(h))-D_{a(h)}(a(h))
=
-\frac{\rho(0)}{2^{m+1}(m+2)}(x(h)-a(h))^{m+1}
+o\!\left((x(h)-a(h))^{m+1}\right).
\]

In particular, after shrinking $\varepsilon$ if needed, every nontrivial nearby
critical pair satisfies
\[
D_a(x)\neq D_a(a),
\]
indeed
\[
\operatorname{sgn}\bigl(D_a(x)-D_a(a)\bigr)=-\operatorname{sgn}\rho(0).
\]
So no sheared repeated-node tangent-contact event can occur in a sufficiently
small neighborhood of an isolated odd zero of $\mathcal R$.
\end{proposition}

\begin{proof}
Write every nearby pair as
\[
a=c-h,\qquad x=c+h,\qquad h=\frac{x-a}{2}>0,\qquad \mu=\frac{c}{h}.
\]

By Lemma~\ref{lem:Da-derivative-formula},
\[
D_a'(x)
=
2\int_0^1 s(1-s)\,\mathcal R(a+s(x-a))\,ds.
\]
With the change of variables $u=2s-1$, this becomes
\[
D_{c-h}'(c+h)
=
\frac14\int_{-1}^1 (1-u^2)\,\mathcal R(c+hu)\,du
=
h^m F(h,\mu),
\]
where
\[
F(h,\mu)
:=
\frac14\int_{-1}^1 (1-u^2)(\mu+u)^m
\rho\!\bigl(h(\mu+u)\bigr)\,du.
\]
This is real-analytic in $(h,\mu)$ near $(0,0)$. At $h=0$,
\[
F(0,\mu)=\frac{\rho(0)}{4}K_m(\mu),
\qquad
K_m(\mu):=\int_{-1}^1 (1-u^2)(\mu+u)^m\,du.
\]
Since $m$ is odd,
\[
K_m(0)=0,
\]
and
\[
K_m'(\mu)
=
m\int_{-1}^1 (1-u^2)(\mu+u)^{m-1}\,du>0
\]
for all real $\mu$, because $m-1$ is even. Hence $\mu=0$ is the unique real
root of $K_m$, and it is simple. Therefore
\[
\partial_\mu F(0,0)\neq 0,
\]
so the implicit function theorem gives a unique real-analytic branch
\[
\mu=\mu(h),\qquad \mu(0)=0,
\]
parameterizing all nearby critical pairs.

Now write
\[
\mathcal Z(\xi)=\mathcal Z(0)+\xi^{m+1}\eta(\xi),
\qquad
\eta(0)=\frac{\rho(0)}{m+1},
\]
since $\mathcal Z'=\mathcal R=\xi^m\rho(\xi)$.
Using
\[
D_a(x)-D_a(a)
=
\frac12\int_{-1}^1 (1-u)\bigl(\mathcal Z(c+hu)-\mathcal Z(c-h)\bigr)\,du,
\]
we obtain
\[
D_{c-h}(c+h)-D_{c-h}(c-h)=h^{m+1}G(h,\mu),
\]
where
\[
G(h,\mu):=
\frac12\int_{-1}^1 (1-u)\Bigl(
(\mu+u)^{m+1}\eta\!\bigl(h(\mu+u)\bigr)
-(\mu-1)^{m+1}\eta\!\bigl(h(\mu-1)\bigr)
\Bigr)\,du.
\]
Again $G$ is real-analytic. Evaluating at $(0,0)$ gives
\[
G(0,0)
=
\frac{\eta(0)}{2}\int_{-1}^1 (1-u)(u^{m+1}-1)\,du
=
-\frac{\rho(0)}{m+2}\neq 0.
\]
Therefore along the critical branch,
\[
D_{a(h)}(x(h))-D_{a(h)}(a(h))
=
h^{m+1}G(h,\mu(h))
=
-\frac{\rho(0)}{m+2}h^{m+1}+O(h^{m+2}),
\]
which has fixed nonzero sign for $h$ sufficiently small. This proves the claim.
\end{proof}

\begin{corollary}[Sheared repeated-node tangent-contact is impossible near an isolated odd $\mathcal R$-zero]
\label{cor:sheared-event-impossible-near-isolated-odd-R-zero}
Let $b$ be an isolated odd-order zero of $\mathcal R$ on a genuine opened
nonmixed arc. Then there exists a neighborhood $U$ of $b$ such that no pair
\[
a<x,\qquad a,x\in U,
\]
can satisfy
\[
\widetilde Y_{\mathcal Z(a)}''(a)=0,
\qquad
\widetilde H_3(a,x)=0,
\qquad
\partial_x\widetilde H_3(a,x)=0.
\]
Equivalently, no rooted sheared repeated-node tangent-contact event occurs in
a sufficiently small neighborhood of an isolated odd zero of $\mathcal R$.
\end{corollary}

\begin{proof}
In the normalized nonmixed chart centered at $b$, the rooted sheared event is
exactly the conjunction
\[
D_a'(x)=0,
\qquad
D_a(x)=D_a(a)=\mathcal Z(a).
\]
Proposition~\ref{prop:odd-zero-critical-branch-uniformly-unrooted} shows that
for every sufficiently small nontrivial critical pair one has
\[
D_a(x)-D_a(a)\neq 0.
\]
So the rooted sheared event is impossible near $b$.
\end{proof}

\begin{proposition}[Quadratic-envelope factorization on a nonmixed interval]
\label{prop:quadratic-envelope-factorization-nonmixed}
Let $J$ be a connected open interval in a good $v_5$-patch such that
\[
K_v\neq 0
\qquad\text{on }J.
\]
Write
\[
Q_v(h)=(Y(h),x(h),S(h)),
\qquad
\mathcal Z=\frac{Y''}{S''},
\qquad
\mathcal R=\mathcal Z'=\frac{T_v}{K_v^2}.
\]
For a fixed source point $a\in J$, define
\[
A_a(x):=Y(x)-Y(a)-Y'(a)(x-a),
\]
\[
B_a(x):=S(x)-S(a)-S'(a)(x-a),
\]
and
\[
\Gamma_a(x):=\frac{A_a(x)}{B_a(x)}
\qquad (x\neq a).
\]

Then:

\smallskip

\noindent
\emph{(i) Tangent-contact denominator.}
For $x\neq a$,
\[
B_a(x)=\int_a^x (x-t)\,S''(t)\,dt,
\]
so $B_a(x)\neq 0$ on $J\setminus\{a\}$.

\smallskip

\noindent
\emph{(ii) Exact repeated-source factorization.}
For distinct $a,u,v\in J$,
\[
H_4(a,u,v)
=
-\,B_a(u)\,B_a(v)\,\bigl(\Gamma_a(u)-\Gamma_a(v)\bigr).
\]

\smallskip

\noindent
\emph{(iii) Integral representation.}
For $x\neq a$,
\[
\Gamma_a(x)
=
\frac{\displaystyle\int_a^x (x-t)\,\mathcal Z(t)\,S''(t)\,dt}
{\displaystyle\int_a^x (x-t)\,S''(t)\,dt}.
\]

\smallskip

\noindent
\emph{(iv) Positive-kernel derivative formula.}
For $x\neq a$,
\[
\Gamma_a'(x)
=
\frac{1}{B_a(x)^2}\int_a^x \Omega_{a,x}(r)\,\mathcal R(r)\,dr,
\]
where
\[
\Omega_{a,x}(r)
:=
\int_a^r\!\!\int_r^x (t-s)\,S''(s)\,S''(t)\,dt\,ds.
\]
Moreover, after orienting $x$ so that $S''$ has fixed sign on $J$,
\[
\Omega_{a,x}(r)>0
\qquad (a<r<x).
\]
\end{proposition}

\begin{proof}
Since $K_v=(x')^3S''$ and $x'\neq 0$ on a good $v_5$-patch, the function
$S''$ has fixed nonzero sign on $J$. Therefore
\[
B_a(x)=\int_a^x (x-t)\,S''(t)\,dt
\]
has fixed nonzero sign for $x\neq a$, proving (i).

For (ii), use the repeated-node determinant
\[
H_4(a,u,v)=
\det\!\begin{pmatrix}
1&Y(a)&a&S(a)\\
0&Y'(a)&1&S'(a)\\
1&Y(u)&u&S(u)\\
1&Y(v)&v&S(v)
\end{pmatrix}.
\]
Subtract from the $u$-row and $v$-row the corresponding affine tangent data
at $a$. This gives
\[
H_4(a,u,v)
=
\det\!\begin{pmatrix}
1&Y(a)&a&S(a)\\
0&Y'(a)&1&S'(a)\\
0&A_a(u)&0&B_a(u)\\
0&A_a(v)&0&B_a(v)
\end{pmatrix}
=
A_a(u)B_a(v)-A_a(v)B_a(u),
\]
hence
\[
H_4(a,u,v)
=
-\,B_a(u)B_a(v)\bigl(\Gamma_a(u)-\Gamma_a(v)\bigr).
\]

For (iii), Taylor's integral formula gives
\[
A_a(x)=\int_a^x (x-t)\,Y''(t)\,dt
=
\int_a^x (x-t)\,\mathcal Z(t)\,S''(t)\,dt,
\]
while
\[
B_a(x)=\int_a^x (x-t)\,S''(t)\,dt.
\]
Dividing yields the formula for $\Gamma_a$.

For (iv), set
\[
N_a(x):=\int_a^x (x-t)\,\mathcal Z(t)\,S''(t)\,dt,
\qquad
D_a(x):=\int_a^x (x-t)\,S''(t)\,dt=B_a(x).
\]
Then $\Gamma_a=N_a/D_a$, so
\[
\Gamma_a'(x)=\frac{N_a'(x)D_a(x)-N_a(x)D_a'(x)}{D_a(x)^2}.
\]
A symmetrization gives
\[
N_a'D_a-N_aD_a'
=
\frac12\int_a^x\!\!\int_a^x
(\mathcal Z(t)-\mathcal Z(s))(t-s)\,S''(s)\,S''(t)\,ds\,dt.
\]
Since $\mathcal Z(t)-\mathcal Z(s)=\int_s^t \mathcal R(r)\,dr$, Fubini yields
\[
N_a'D_a-N_aD_a'
=
\int_a^x \Omega_{a,x}(r)\,\mathcal R(r)\,dr
\]
with
\[
\Omega_{a,x}(r)
=
\int_a^r\!\!\int_r^x (t-s)\,S''(s)\,S''(t)\,dt\,ds.
\]
Because $t-s>0$ on the region of integration and $S''$ has fixed sign on
$J$, one has $\Omega_{a,x}(r)>0$ after orienting the $x$-coordinate.
This proves (iv).
\end{proof}

\begin{corollary}[Repeated-source monotonicity on one-sign $\mathcal R$-blocks]
\label{cor:repeated-source-monotonicity-one-sign-R}
Let $I\subset J$ be a subinterval on which
\[
\mathcal R \text{ has fixed nonzero sign.}
\]
Then for every $a\in I$, the function $x\mapsto \Gamma_a(x)$ is strictly
monotone on each component of $I\setminus\{a\}$. In particular, there do not
exist points
\[
a<u<v \quad \text{in } I
\]
such that
\[
H_4(a,u,v)=0.
\]
\end{corollary}

\begin{proof}
This is immediate from Proposition~\ref{prop:quadratic-envelope-factorization-nonmixed}:
the derivative $\Gamma_a'(x)$ is a positive-kernel average of $\mathcal R$,
so it has the same sign as $\mathcal R$ on each one-sign block. The factorization
of $H_4$ then excludes repeated-source exact contacts there.
\end{proof}

\begin{remark}[Below-lifted source-determinant reformulation of $\Gamma_a$]
\label{rem:source-determinant-reformulation-Gammaa}
Let
\[
\lambda:=\frac{\Delta_7}{v_5},
\qquad
x=\frac{v_5}{c},
\qquad
Y=\frac{u_5}{c},
\qquad
S=\frac{\lambda}{c}.
\]
For any scalar germ $f$, define
\[
D_f(a,h):=
\det\!\begin{pmatrix}
c(a)&v_5(a)&f(a)\\
c'(a)&v_5'(a)&f'(a)\\
c(h)&v_5(h)&f(h)
\end{pmatrix}.
\]
A direct row reduction gives
\[
D_f(a,h)
=
c(a)^2c(h)\Bigl(
x'(a)\bigl(F(h)-F(a)\bigr)-F'(a)\bigl(x(h)-x(a)\bigr)
\Bigr),
\qquad
F=\frac{f}{c}.
\]
In particular,
\[
\Psi_a(h):=\frac{D_{u_5}(a,h)}{D_\lambda(a,h)}
=
\Gamma_{x(a)}(x(h)).
\]
Thus the repeated-source nonmixed $4$-point geometry is already encoded by a
finite ratio of source determinants.

Likewise, if one sets
\[
\mathfrak D(a,h):=
\det\!\begin{pmatrix}
c(a)&v_5(a)&\lambda(a)&u_5(a)\\
c'(a)&v_5'(a)&\lambda'(a)&u_5'(a)\\
c(h)&v_5(h)&\lambda(h)&u_5(h)\\
c'(h)&v_5'(h)&\lambda'(h)&u_5'(h)
\end{pmatrix},
\]
then a Pl\"ucker identity yields
\[
\frac{d}{dh}\Psi_a(h)
=
\frac{W_h(c,v_5)(a)\,\mathfrak D(a,h)}{D_\lambda(a,h)^2}.
\]
So the sign of the quadratic-envelope derivative is governed by one explicit
$4\times 4$ source determinant.

Moreover, combining the identity
\[
\Psi_a(h)=\Gamma_{x(a)}(x(h))
\]
with Proposition~\ref{prop:quadratic-envelope-factorization-nonmixed} gives the
exact source-level kernel formula
\[
\mathfrak D(a,h)
=
c(a)^2c(h)^2x'(a)x'(h)
\int_{x(a)}^{x(h)} \Omega_{x(a),x(h)}(r)\,\mathcal R(r)\,dr.
\]
After orienting the $x$-coordinate so that $x'>0$ on the interval, the sign
of $\mathfrak D(a,h)$ is therefore exactly the sign of the corresponding
positive-kernel average of $\mathcal R$. In particular, on every one-sign
$\mathcal R$-block the quadruple
\[
(c,v_5,\lambda,u_5)
\]
is osculatory-ECT in the sense that $\mathfrak D(a,h)$ has fixed nonzero sign
for $a<h$.
\end{remark}

\begin{corollary}[Source-level Wronskian consequences for the nonmixed branch]
\label{cor:source-level-Wronskian-consequences-nonmixed}
On a good nonmixed $v_5$-patch, the following hold.

\begin{enumerate}
\item If there is an exact nonmixed tangent-plane contact, then there exists a
source point $h_*$ in the source convex hull of the contact nodes such that
\[
\mathcal W_{\mathrm{src}}(h_*)=0.
\]

\item If $\mathcal W_{\mathrm{src}}$ has fixed nonzero sign on an open
nonmixed interval, then there are no exact nonmixed tangent-plane contacts on
that interval.

\item If every nonempty open subinterval of a nonmixed interval contains an
exact nonmixed tangent-plane contact, then
\[
\mathcal W_{\mathrm{src}}\equiv 0
\]
on that interval, hence
\[
\mathcal R\equiv 0,
\]
and therefore the lifted curve collapses to affine-planar geometry.
\end{enumerate}
\end{corollary}

\begin{proof}
For (1), Proposition~\ref{prop:exact-nonmixed-tangent-plane-contacts-force-R-zero}
produces a point in the convex hull of the three normalized nodes where
$\mathcal R=0$. Pulling back through the monotone source-to-$x$ coordinate
on the nonmixed patch and using
Proposition~\ref{prop:source-level-Wronskian-nonmixed-four-point} gives
$\mathcal W_{\mathrm{src}}=0$ at the corresponding source point.

For (2), by
Proposition~\ref{prop:source-level-Wronskian-nonmixed-four-point},
$\mathcal W_{\mathrm{src}}$ and $\mathcal R$ have the same sign on a
nonmixed arc. So Corollary~\ref{cor:repeated-source-monotonicity-one-sign-R}
applies.

For (3), every nonempty open subinterval contains a zero of
$\mathcal W_{\mathrm{src}}$ by (1), so analyticity forces
$\mathcal W_{\mathrm{src}}\equiv 0$. Proposition~\ref{prop:source-level-Wronskian-nonmixed-four-point}
then gives $\mathcal R\equiv 0$, and
Corollary~\ref{cor:recurrent-exact-nonmixed-contacts-force-affine-planarity}
gives affine-planar collapse.
\end{proof}

\begin{remark}[Chebyshev formulation of the nonmixed burden]
\label{rem:Chebyshev-formulation-nonmixed}
Corollary~\ref{cor:source-level-Wronskian-consequences-nonmixed} shows that the
nonmixed $4$-point branch is controlled source-level by the pair
\[
(J_\lambda,\ J_{u_5}).
\]
If this pair were a Chebyshev pair on an opened nonmixed arc, i.e.
\[
\mathcal W_{\mathrm{src}}\neq 0
\]
throughout that arc, then the entire nonmixed $4$-point branch would be
sterile there. Thus a one-sign theorem for $\mathcal W_{\mathrm{src}}$ would
still be sufficient. However, the later source-level no-go theorem shows that
generic Wronskian algebra alone is not enough to force such a statement. The
actual remaining burden is therefore a source-coupling theorem that excludes
the analytically possible odd zeros of $\mathcal W_{\mathrm{src}}$ on genuine
opened nonmixed arcs.
\end{remark}

\begin{proposition}[Lift of the source-level Wronskian to a $4\times4$ Wronskian]
\label{prop:source-level-Wronskian-lifts-to-4x4}
With
\[
J_\lambda=W(c,v_5,\lambda),
\qquad
J_{u_5}=W(c,v_5,u_5),
\qquad
M=W(c,v_5)=cv_5'-c'v_5,
\]
one has the exact identity
\[
\boxed{
\mathcal W_{\mathrm{src}}
=
J_\lambda J_{u_5}'-J_{u_5}J_\lambda'
=
W(J_\lambda,J_{u_5})
=
M\,W(c,v_5,\lambda,u_5).
}
\]
Thus, on a nonmixed arc where $M\neq 0$, the odd-zero problem for
$\mathcal W_{\mathrm{src}}$ is equivalent to the odd-zero problem for the
source-level $4\times4$ Wronskian
\[
\mathcal U_{\mathrm{src}}:=W(c,v_5,\lambda,u_5).
\]
\end{proposition}

\begin{proof}
By definition,
\[
\mathcal W_{\mathrm{src}}=W(J_\lambda,J_{u_5}).
\]
Now apply the Sylvester--Wronski identity
\[
W\!\bigl(W(f_1,f_2,f_3),\,W(f_1,f_2,f_4)\bigr)
=
W(f_1,f_2)\,W(f_1,f_2,f_3,f_4)
\]
with
\[
(f_1,f_2,f_3,f_4)=(c,v_5,\lambda,u_5).
\]
This gives
\[
W(J_\lambda,J_{u_5})
=
W(c,v_5)\,W(c,v_5,\lambda,u_5)
=
M\,W(c,v_5,\lambda,u_5),
\]
which is exactly the claimed identity.
\end{proof}

\begin{proposition}[The current source-level Wronskian package is not yet a contradiction engine for nonmixed odd zeros]
\label{prop:source-level-Wronskian-package-not-enough}
Let $I$ be an interval on which
\[
c\neq 0,
\qquad
M=W(c,v_5)\neq 0,
\qquad
J_\lambda=W(c,v_5,\lambda)\neq 0.
\]
Then:

\begin{enumerate}
\item for every analytic forcing $F:I\to\mathbf C$, there exists an analytic
function $u_5$ on $I$ such that
\[
W(c,v_5,\lambda,u_5)=F;
\]

\item consequently, for every analytic forcing $G:I\to\mathbf C$, there exists
an analytic function $u_5$ on $I$ such that
\[
\mathcal W_{\mathrm{src}}=G;
\]

\item equivalently, for every analytic forcing $R:I\to\mathbf C$, there exists
an analytic function $u_5$ on $I$ such that
\[
\mathcal R=R.
\]
\end{enumerate}

In particular, at the current source-level ratio/Wronskian/$4\times4$-determinant
level, the zero set, parity, and odd/even vanishing pattern of
$\mathcal W_{\mathrm{src}}$ and $\mathcal R$ are formally unconstrained.
\end{proposition}

\begin{proof}
Expand the $4\times4$ Wronskian along the $u_5$-column. Since the highest
derivative of $u_5$ appearing in $W(c,v_5,\lambda,u_5)$ is $u_5'''$, with
coefficient exactly
\[
W(c,v_5,\lambda)=J_\lambda,
\]
there exist analytic coefficients $a_0,a_1,a_2$ on $I$ such that
\[
W(c,v_5,\lambda,u_5)
=
J_\lambda\,u_5'''+a_2u_5''+a_1u_5'+a_0u_5.
\]
Because $J_\lambda\neq 0$ on $I$, this is a third-order linear analytic ODE
for $u_5$ with nonvanishing leading coefficient. Standard analytic ODE
existence therefore gives an analytic solution $u_5$ on $I$ for every
analytic forcing $F$. This proves (1).

For (2), given analytic $G$, set
\[
F:=\frac{G}{M}.
\]
Since $M\neq 0$ on $I$, $F$ is analytic. By (1), choose $u_5$ such that
\[
W(c,v_5,\lambda,u_5)=F.
\]
Then Proposition~\ref{prop:source-level-Wronskian-lifts-to-4x4} gives
\[
\mathcal W_{\mathrm{src}}
=
M\,W(c,v_5,\lambda,u_5)
=
G.
\]

For (3), given analytic $R$, set
\[
G:=\frac{M J_\lambda^2}{c^2}\,R.
\]
Since $c\neq 0$, $M\neq 0$, and $J_\lambda\neq 0$ on $I$, this is
analytic. By (2), choose $u_5$ such that $\mathcal W_{\mathrm{src}}=G$.
Then Proposition~\ref{prop:source-level-Wronskian-nonmixed-four-point} gives
\[
\mathcal R
=
\frac{c^2}{M}\frac{\mathcal W_{\mathrm{src}}}{J_\lambda^2}
=
\frac{c^2}{M}\frac{G}{J_\lambda^2}
=
R.
\qedhere
\]
\end{proof}

\begin{remark}[Ceiling of the current source-level reduction]
\label{rem:ceiling-of-current-source-level-reduction}
Propositions~\ref{prop:source-level-Wronskian-lifts-to-4x4} and
\ref{prop:source-level-Wronskian-package-not-enough} show that the current
source-level reduction has been pushed to its natural ceiling.  The nonmixed
$4$-point branch is indeed equivalent to an odd-zero problem for the
gauge-invariant source-level Wronskian
\[
\mathcal W_{\mathrm{src}},
\]
and even to the $4\times4$ Wronskian $W(c,v_5,\lambda,u_5)$.  But at that
level the odd-zero pattern is still formally free.

Accordingly, the missing ingredient can no longer be a generic Wronskian or
determinant argument.  It must be a genuine \emph{source-coupling theorem}:
one that uses additional identities coming from the actual RH source structure
to show that the analytically possible $u_5$-freedom exhibited above is not
realizable on a genuine opened nonmixed arc.
\end{remark}

\begin{remark}[First-loss rigidification of the quadratic-envelope scalar]
\label{rem:first-loss-rigidification-nonmixed}
Set
\[
M(a,x):=B_a(x)^2\,\Gamma_a'(x)=A_a'(x)B_a(x)-A_a(x)B_a'(x).
\]
Then
\[
\partial_x M(a,x)=S''(x)\,B_a(x)\,\bigl(\mathcal Z(x)-\Gamma_a(x)\bigr).
\]
Moreover, at any point where $M(a,x)=0$,
\[
\partial_a M(a,x)=S''(a)\,C_a(x)\,\bigl(\mathcal Z(a)-\Gamma_a(x)\bigr),
\]
where
\[
C_a(x):=S(a)-S(x)+S'(x)(x-a)=\int_a^x (t-a)\,S''(t)\,dt.
\]

Thus any formally attained first loss of monotonicity for $\Gamma_a$ forces
\[
\Gamma_a(x)=\mathcal Z(a).
\]
After the affine shear
\[
\widetilde Y_\gamma:=Y-\gamma S,
\qquad \gamma=\mathcal Z(a),
\]
this yields
\[
\widetilde Y_\gamma''(a)=0,
\qquad
\widetilde H_3(a,x)=0,
\qquad
\partial_x\widetilde H_3(a,x)=0.
\]
So a first loss of quadratic-envelope monotonicity can only occur through a
sheared repeated-node tangent-contact configuration rooted at a curvature-zero
point. In particular, the remaining nonmixed obstruction is diagonal in nature:
it is governed by odd zeros of $\mathcal R$, not by an independent separated
source-target mechanism.
\end{remark}

\begin{proposition}[The current $3$-point package is not by itself a contradiction engine for nonmixed first loss]
\label{prop:three-point-package-not-enough-for-nonmixed-first-loss}
Fix an analytic function $S$ on an interval $I$ with
\[
S''\neq 0
\qquad\text{on }I.
\]
Let $a<x$ be two points of $I$. Then there exist an analytic function $Y$ on
$I$ and a constant $\gamma\in\mathbf R$ such that, with
\[
\widetilde Y_\gamma:=Y-\gamma S,
\]
one has
\[
\widetilde Y_\gamma''(a)=0,
\qquad
\widetilde H_3(a,x)=0,
\qquad
\partial_x\widetilde H_3(a,x)=0,
\]
where
\[
\widetilde H_3(a,u):=
\widetilde Y_\gamma(u)-\widetilde Y_\gamma(a)-\widetilde Y_\gamma'(a)(u-a).
\]
In particular, the rigidified repeated-node tangent-contact event from
Remark~\ref{rem:first-loss-rigidification-nonmixed} is formally compatible with
an arbitrary fixed nonmixed planar curve
\[
\Lambda_v=(x,S),
\]
so the established $3$-point package for $\Lambda_v$ alone cannot exclude that
event.
\end{proposition}

\begin{proof}
Define
\[
T_{a,x}(t):=(t-a)^3(t-x)^2.
\]
Then
\[
T_{a,x}(a)=0,
\qquad
T_{a,x}'(a)=0,
\qquad
T_{a,x}''(a)=0,
\qquad
T_{a,x}(x)=0,
\qquad
T_{a,x}'(x)=0.
\]
Hence
\[
T_{a,x}(x)-T_{a,x}(a)-T_{a,x}'(a)(x-a)=0,
\qquad
T_{a,x}'(x)-T_{a,x}'(a)=0.
\]
Now fix any $\gamma\in\mathbf R$ and set
\[
Y:=\gamma S+T_{a,x}.
\]
Then
\[
\widetilde Y_\gamma=Y-\gamma S=T_{a,x},
\]
so
\[
\widetilde Y_\gamma''(a)=0,
\qquad
\widetilde H_3(a,x)=0,
\qquad
\partial_x\widetilde H_3(a,x)=0.
\]
Since $S$ itself has not changed and $S''\neq 0$ on $I$, the original
$3$-point geometry of $\Lambda_v=(x,S)$ on $I$ is exactly the same as before.
Thus the sheared repeated-node tangent-contact event is not ruled out by the
planar derivative geometry of $\Lambda_v$ alone.
\end{proof}

\begin{proposition}[A sheared repeated-node event forces two interior $\mathcal R$-zeros]
\label{prop:sheared-repeated-node-forces-two-R-zeros}
Let $J$ be a connected opened nonmixed arc, and let
\[
a<x
\qquad\text{lie in }J.
\]
Set
\[
\gamma:=\mathcal Z(a),
\qquad
\widetilde Y_\gamma:=Y-\gamma S,
\]
and let
\[
\widetilde H_3(a,u):=
\widetilde Y_\gamma(u)-\widetilde Y_\gamma(a)-\widetilde Y_\gamma'(a)(u-a).
\]
Assume
\[
\widetilde Y_\gamma''(a)=0,
\qquad
\widetilde H_3(a,x)=0,
\qquad
\partial_x\widetilde H_3(a,x)=0.
\]
Then one of the following holds:
\begin{enumerate}
\item $\mathcal R\equiv 0$ on $[a,x]$;
\item there exist points
\[
a<r_1<r_2<x
\]
such that
\[
\mathcal Z(r_1)=\mathcal Z(r_2)=\mathcal Z(a).
\]
Consequently there exist
\[
\xi_1\in(a,r_1),
\qquad
\xi_2\in(r_1,r_2)
\]
with
\[
\mathcal R(\xi_1)=\mathcal R(\xi_2)=0.
\]
\end{enumerate}
In particular, every nontrivial sheared repeated-node event contains at least two
interior zeros of $\mathcal R$.
\end{proposition}

\begin{proof}
Since
\[
\widetilde Y_\gamma''=Y''-\gamma S''=(\mathcal Z-\gamma)S'',
\]
the choice $\gamma=\mathcal Z(a)$ gives
\[
\widetilde Y_\gamma''=(\mathcal Z-\mathcal Z(a))S''.
\]
Also,
\[
\partial_x\widetilde H_3(a,x)
=
\widetilde Y_\gamma'(x)-\widetilde Y_\gamma'(a)
=
\int_a^x \widetilde Y_\gamma''(t)\,dt,
\]
and
\[
\widetilde H_3(a,x)
=
\int_a^x (x-t)\,\widetilde Y_\gamma''(t)\,dt.
\]
Therefore the hypotheses imply
\[
\int_a^x (\mathcal Z(t)-\mathcal Z(a))\,S''(t)\,dt=0,
\]
and
\[
\int_a^x (x-t)(\mathcal Z(t)-\mathcal Z(a))\,S''(t)\,dt=0.
\]

Choose $\varepsilon\in\{\pm1\}$ so that
\[
\varepsilon S''>0
\qquad\text{on }J,
\]
and define the positive measure
\[
d\mu(t):=\varepsilon S''(t)\,dt,
\qquad
\phi(t):=\mathcal Z(t)-\mathcal Z(a).
\]
Then
\[
\int_a^x \phi\,d\mu=0,
\qquad
\int_a^x (x-t)\phi\,d\mu=0.
\]
Hence also
\[
\int_a^x t\,\phi(t)\,d\mu(t)=0.
\]

If $\phi\equiv 0$ on $[a,x]$, then
\[
\mathcal R=\phi'\equiv 0
\qquad\text{on }[a,x],
\]
which is alternative (1).

Assume now that $\phi\not\equiv 0$. Then $\phi$ cannot have fixed sign on
$[a,x]$, because $\mu$ is positive and
\[
\int_a^x \phi\,d\mu=0.
\]
It also cannot have exactly one sign change. Indeed, if $c\in(a,x)$ were the
unique sign-change point, then
\[
(t-c)\phi(t)\ge 0
\qquad\text{on }[a,x],
\]
with strict inequality on a set of positive $\mu$-measure. Therefore
\[
\int_a^x (t-c)\phi(t)\,d\mu(t)>0.
\]
But
\[
\int_a^x (t-c)\phi(t)\,d\mu(t)
=
\int_a^x t\,\phi(t)\,d\mu(t)
-c\int_a^x \phi(t)\,d\mu(t)
=0,
\]
a contradiction.

So $\phi$ has at least two sign changes on $(a,x)$. Since $\phi(a)=0$ and
$\phi$ is continuous, there exist
\[
a<r_1<r_2<x
\]
such that
\[
\phi(r_1)=\phi(r_2)=0,
\]
i.e.
\[
\mathcal Z(r_1)=\mathcal Z(r_2)=\mathcal Z(a).
\]
Applying Rolle's theorem to $\phi=\mathcal Z-\mathcal Z(a)$ on $[a,r_1]$ and on
$[r_1,r_2]$ gives
\[
\mathcal R(\xi_1)=0
\quad\text{for some }\xi_1\in(a,r_1),
\qquad
\mathcal R(\xi_2)=0
\quad\text{for some }\xi_2\in(r_1,r_2).
\]
This proves (2).
\end{proof}

\begin{corollary}[Isolated $\mathcal R$-zeros are shear-sterile]
\label{cor:isolated-R-zeros-are-shear-sterile}
Let $b\in J$ be an isolated zero of $\mathcal R$. Then there exists a
neighborhood $U\subset J$ of $b$ such that no pair
\[
a<x,
\qquad
a,x\in U,
\]
can satisfy
\[
\gamma=\mathcal Z(a),
\qquad
\widetilde Y_\gamma''(a)=0,
\qquad
\widetilde H_3(a,x)=0,
\qquad
\partial_x\widetilde H_3(a,x)=0.
\]
\end{corollary}

\begin{proof}
Choose $U$ so small that $b$ is the only zero of $\mathcal R$ in $U$. If such a
pair $a<x$ existed in $U$, then
Proposition~\ref{prop:sheared-repeated-node-forces-two-R-zeros} would produce
two interior zeros of $\mathcal R$ in $(a,x)\subset U$, a contradiction.
\end{proof}

\begin{corollary}[Diagonal sheared recurrence forces affine-planar collapse]
\label{cor:diagonal-sheared-recurrence-forces-affine-planar-collapse}
Let $J$ be a connected opened nonmixed arc. Suppose there exists a sequence
\[
a_n<x_n,
\qquad
a_n,x_n\in J,
\qquad
a_n,x_n\to b,
\qquad
x_n-a_n\to 0,
\]
such that, with
\[
\gamma_n:=\mathcal Z(a_n),
\qquad
\widetilde Y_{\gamma_n}:=Y-\gamma_n S,
\]
one has
\[
\widetilde Y_{\gamma_n}''(a_n)=0,
\qquad
\widetilde H_{3,n}(a_n,x_n)=0,
\qquad
\partial_x\widetilde H_{3,n}(a_n,x_n)=0,
\]
where
\[
\widetilde H_{3,n}(a,u):=
\widetilde Y_{\gamma_n}(u)-\widetilde Y_{\gamma_n}(a)-\widetilde Y_{\gamma_n}'(a)(u-a).
\]
Then
\[
\mathcal R\equiv 0
\]
on a neighborhood of $b$. Hence on some nontrivial interval containing $b$ one
has
\[
Y=\gamma S+\alpha x+\beta,
\]
so the lifted curve is locally affine-planar there.
\end{corollary}

\begin{proof}
If for some $n$ the first alternative in
Proposition~\ref{prop:sheared-repeated-node-forces-two-R-zeros} holds, then
$\mathcal R\equiv 0$ on $[a_n,x_n]$ and we are done.

Otherwise, for every $n$ the interval $(a_n,x_n)$ contains at least two zeros of
$\mathcal R$. Since
\[
a_n,x_n\to b,
\qquad
x_n-a_n\to 0,
\]
these zeros accumulate at $b$. Because $\mathcal R$ is analytic, this implies
\[
\mathcal R\equiv 0
\]
on a neighborhood of $b$. Therefore $\mathcal Z$ is locally constant, so
\[
Y''=\gamma S''
\]
there, and integrating twice gives
\[
Y=\gamma S+\alpha x+\beta
\]
on a nontrivial interval around $b$.
\end{proof}

\begin{proposition}[Uniform fixed-level nonrecurrence for the sheared $3$-point branch]
\label{prop:uniform-fixed-level-nonrecurrence-sheared-three-point}
Let $J$ be a connected opened nonmixed arc, fix a real scalar $\gamma$, and set
\[
\widetilde Y_\gamma:=Y-\gamma S=\frac{u_5-\gamma\lambda}{c},
\qquad
f_\gamma:=u_5-\gamma\lambda,
\qquad
J_\gamma:=J_{f_\gamma}=J_{u_5}-\gamma J_\lambda.
\]
Then
\[
\widetilde Y_\gamma''=\frac{c^3}{M^3}J_\gamma.
\]

Assume that $b\in J$ is an isolated finite-order odd zero of $J_\gamma$.
After affine normalization of the planar curve
\[
(x,\widetilde Y_\gamma)
\]
at $b$, let
\[
I_-=(\alpha,0)
\]
be the source interval on which
\[
\widetilde Y_\gamma''<0,
\]
and let
\[
I_+=(0,\beta)
\]
be the maximal open interval to the right of $0$ on which
\[
\widetilde Y_\gamma'
\]
is strictly increasing. Set
\[
\mathcal J_\gamma:=\widetilde Y_\gamma'(I_-)\cap \widetilde Y_\gamma'(I_+).
\]
For
\[
m\in\mathcal J_\gamma,
\]
write
\[
a_\gamma(m):=(\widetilde Y_\gamma'|_{I_-})^{-1}(m),
\]
and define
\[
\widetilde H_{3,\gamma}(a,u):=
\widetilde Y_\gamma(u)-\widetilde Y_\gamma(a)-\widetilde Y_\gamma'(a)(u-a).
\]

Then there exists a unique maximal open interval
\[
\mathcal J_{\gamma,\mathrm{br}}=(0,m_{\gamma,*})\subset \mathcal J_\gamma
\]
and a unique real-analytic function
\[
u_\gamma:\mathcal J_{\gamma,\mathrm{br}}\to I_+
\]
such that
\[
\widetilde H_{3,\gamma}(a_\gamma(m),u_\gamma(m))=0
\qquad
(m\in\mathcal J_{\gamma,\mathrm{br}}),
\]
and this branch agrees near $m=0$ with the unique local opposite-side branch
born at the odd zero $b$.

Moreover, along the whole branch one has
\[
\partial_u\widetilde H_{3,\gamma}(a_\gamma(m),u_\gamma(m))
=
\widetilde Y_\gamma'(u_\gamma(m))-m
\neq 0,
\]
so there is no interior common-tangent event, no interior recurrence, and no
interior accumulation. If
\[
m_{\gamma,*}<\sup\mathcal J_\gamma,
\]
then
\[
u_\gamma(m)\longrightarrow \beta
\qquad (m\uparrow m_{\gamma,*}),
\]
so the branch can stop only by boundary exit of the relevant slope-overlap
block.

Equivalently, using
\[
\widetilde H_{3,\gamma}(a,u)=A_a(u)-\gamma B_a(u)
=
B_a(u)\bigl(\Gamma_a(u)-\gamma\bigr),
\]
one has along the branch
\[
\Gamma_{a_\gamma(m)}(u_\gamma(m))=\gamma,
\qquad
\Gamma_{a_\gamma(m)}'(u_\gamma(m))\neq 0.
\]
In particular, for every fixed $\gamma$, the distinguished level $\gamma$ of
the repeated-source scalar $\Gamma_a$ has no interior critical point on the
sheared odd-zero branch.
\end{proposition}

\begin{proof}
By Remark~\ref{rem:below-lifted-wronskian-ratio-reformulation}, for any scalar
germ $f$ one has
\[
\frac{d^2}{dx^2}\!\left(\frac{f}{c}\right)=\frac{c^3}{M^3}J_f.
\]
Applying this with
\[
f=f_\gamma=u_5-\gamma\lambda
\]
gives
\[
\widetilde Y_\gamma''=\frac{c^3}{M^3}J_\gamma.
\]
Hence $b$ is an isolated finite-order odd zero of $J_\gamma$ if and only if it
is an isolated finite-order odd zero of the planar curvature
\[
\widetilde Y_\gamma''.
\]

Now apply
Theorem~\ref{thm:local-classification-lifted-three-point-branch-near-Kv-zero}
and
Theorem~\ref{thm:global-continuation-nonrecurrence-odd-three-point-branch}
with the affine-normalized analytic function
\[
\widetilde Y_\gamma
\]
playing the role of the planar graphing function $S$. The proofs of those
results use only the planar repeated-node tangent-contact equation for the
chosen analytic function, its curvature parity at the base point, and the
maximal monotone-slope block on the opposite side. Therefore they apply
verbatim to
\[
\widetilde H_{3,\gamma}(a,u)
=
\widetilde Y_\gamma(u)-\widetilde Y_\gamma(a)-\widetilde Y_\gamma'(a)(u-a),
\]
and yield the unique local opposite-side branch, its unique continuation on the
maximal slope-overlap block, the nonvanishing of
\[
\partial_u\widetilde H_{3,\gamma},
\]
and the boundary-exit alternative. This proves the first part.

For the reformulation in terms of $\Gamma_a$, note from the definitions of
$A_a$, $B_a$, and $\Gamma_a$ that
\[
\widetilde H_{3,\gamma}(a,u)
=
\bigl(Y(u)-Y(a)-Y'(a)(u-a)\bigr)
-\gamma\bigl(S(u)-S(a)-S'(a)(u-a)\bigr)
=
A_a(u)-\gamma B_a(u).
\]
Since
\[
A_a(u)=\Gamma_a(u)\,B_a(u),
\]
this becomes
\[
\widetilde H_{3,\gamma}(a,u)
=
B_a(u)\bigl(\Gamma_a(u)-\gamma\bigr).
\]
By Proposition~\ref{prop:quadratic-envelope-factorization-nonmixed},
\[
B_a(u)\neq 0
\qquad (u\neq a),
\]
so along the branch
\[
\widetilde H_{3,\gamma}(a_\gamma(m),u_\gamma(m))=0
\iff
\Gamma_{a_\gamma(m)}(u_\gamma(m))=\gamma.
\]

Differentiating
\[
\widetilde H_{3,\gamma}(a,u)=B_a(u)\bigl(\Gamma_a(u)-\gamma\bigr)
\]
with respect to $u$ gives
\[
\partial_u\widetilde H_{3,\gamma}(a,u)
=
B_a'(u)\bigl(\Gamma_a(u)-\gamma\bigr)+B_a(u)\Gamma_a'(u).
\]
On the branch, where
\[
\Gamma_a(u)=\gamma,
\]
this simplifies to
\[
\partial_u\widetilde H_{3,\gamma}(a,u)=B_a(u)\Gamma_a'(u).
\]
Since the first part already gives
\[
\partial_u\widetilde H_{3,\gamma}(a_\gamma(m),u_\gamma(m))\neq 0
\]
and
\[
B_{a_\gamma(m)}(u_\gamma(m))\neq 0,
\]
it follows that
\[
\Gamma_{a_\gamma(m)}'(u_\gamma(m))\neq 0
\]
along the whole branch. This proves the final statement.
\end{proof}

\begin{remark}[No-go for the fixed-level chamber-exit route]
\label{rem:no-go-fixed-level-chamber-exit-route}
One might hope to strengthen
Proposition~\ref{prop:uniform-fixed-level-nonrecurrence-sheared-three-point}
as follows: if a fixed level $\gamma$ is chosen near a strict local extremum of
$\mathcal Z$, so that
\[
J_\gamma=(\mathcal Z-\gamma)J_\lambda
\]
has two nearby simple zeros
\[
r_-<r_+,
\]
then the corresponding sheared opposite-side $3$-point branch should be unable
to reach the chamber boundary $u=r_+$ without first encountering an interior
common-tangent point
\[
\partial_u \widetilde H_{3,\gamma}=0.
\]
This is false already at the level of the local planar/sheared equations.

Indeed, consider the model functions
\[
S_{\mathrm{mod}}(x):=x^2,
\qquad
\mathcal Z_{\mathrm{mod}}(x):=1-x^2,
\qquad
\mathcal R_{\mathrm{mod}}(x):=\mathcal Z_{\mathrm{mod}}'(x)=-2x,
\]
and fix the level
\[
\gamma:=0.
\]
Set
\[
F_{\mathrm{mod}}(x):=x^2-\frac{x^4}{6}.
\]
Then
\[
F_{\mathrm{mod}}''(x)=2(1-x^2)
=
(\mathcal Z_{\mathrm{mod}}(x)-\gamma)\,S_{\mathrm{mod}}''(x),
\]
so the two simple zeros of the sheared curvature are
\[
r_-=-1,
\qquad
r_+=1.
\]

Let
\[
H_{\mathrm{mod}}(a,u):=
F_{\mathrm{mod}}(u)-F_{\mathrm{mod}}(a)-F_{\mathrm{mod}}'(a)(u-a).
\]
A direct computation gives
\[
H_{\mathrm{mod}}(a,u)
=
-\frac{(u-a)^2}{6}\,
\bigl(3a^2+2au+u^2-6\bigr).
\]
Hence the nontrivial opposite-side branch issued from the odd zero $r_-=-1$ is
\[
u(a)=-a-\sqrt{6-2a^2},
\qquad
-\frac53<a<-1.
\]
Along this branch one has
\[
-1<u(a)<1,
\qquad
u(a)\uparrow 1
\quad\text{as}\quad
a\downarrow -\frac53,
\]
so the branch runs inside the chamber and stops by boundary exit at the second
curvature zero.

Moreover,
\[
\partial_u H_{\mathrm{mod}}(a,u)
=
F_{\mathrm{mod}}'(u)-F_{\mathrm{mod}}'(a)
=
-\frac{u-a}{3}\,\bigl(a^2+au+u^2-3\bigr).
\]
Substituting $u=u(a)$ yields
\[
a^2+a\,u(a)+u(a)^2-3
=
-a^2+a\sqrt{6-2a^2}+3<0
\qquad
\left(-\frac53<a<-1\right),
\]
and therefore
\[
\partial_u H_{\mathrm{mod}}(a,u(a))>0
\qquad
\left(-\frac53<a<-1\right).
\]
So there is \emph{no} interior common-tangent point anywhere on the branch
before boundary exit. In fact, if one continues the same algebraic branch past
the chamber boundary, the first common tangent occurs only later at
\[
(a,u)=(-\sqrt3,\sqrt3),
\]
strictly after the chamber has already been crossed.

Thus fixed-level sheared nonrecurrence cannot be sharpened to a chamber-exit
contradiction by planar arguments alone. The remaining obstruction is therefore
genuinely source-coupled: any successful odd-zero exclusion must use the
source-pinned moving shear
\[
f_a=u_5-\mathcal Z(a)\lambda
\]
or some equivalent RH-specific coupling, rather than only a frozen level
$\gamma$ and the planar branch geometry of
\[
\widetilde Y_\gamma=Y-\gamma S.
\]
\end{remark}

\begin{remark}[Source-level form of the shear transfer]
\label{rem:source-level-form-of-shear-transfer}
For fixed $a\in J$, set
\[
f_a:=u_5-\mathcal Z(a)\lambda.
\]
Then
\[
\widetilde Y_{\mathcal Z(a)}=\frac{f_a}{c},
\qquad
J_{f_a}=J_{u_5}-\mathcal Z(a)J_\lambda
=(\mathcal Z-\mathcal Z(a))J_\lambda,
\]
and
\[
D_{f_a}(a,h)=D_{u_5}(a,h)-\mathcal Z(a)\,D_\lambda(a,h).
\]
Thus the sheared repeated-node conditions are exactly a repeated-node
tangent-contact event for the source triple $(c,v_5,f_a)$:
\[
J_{f_a}(a)=0,
\qquad
D_{f_a}(a,h)=0,
\qquad
\partial_h D_{f_a}(a,h)=0.
\]
Proposition~\ref{prop:sheared-repeated-node-forces-two-R-zeros} shows that such
a source-level event forces either $J_{f_a}\equiv 0$ on the interval or at least
two interior zeros of $J_{f_a}$, equivalently at least two interior zeros of
$\mathcal R$ and hence of $\mathcal W_{\mathrm{src}}$.
\end{remark}

\begin{proposition}[Source-determinant form of the odd-zero unrooted defect]
\label{prop:source-determinant-form-odd-zero-unrooted-defect}
Let $J$ be a good nonmixed $v_5$-patch, and let $b\in J$ be an isolated
odd-order zero of $\mathcal R$.  Use the local $x$-coordinate on $J$ centered
at $b$, so that
\[
x(b)=0,
\qquad
\mathcal R(x)=x^m\rho(x),
\qquad
m\ge 1 \text{ odd},
\qquad
\rho(0)\neq 0.
\]
For each source point $a\in J$, set
\[
\Psi_a(h):=\frac{D_{u_5}(a,h)}{D_\lambda(a,h)},
\qquad
f_a:=u_5-\mathcal Z(x(a))\,\lambda.
\]
Then, after shrinking $J$ around $b$ if necessary, the following hold.

\begin{enumerate}
\item The critical locus
\[
\partial_h\Psi_a(h)=0
\qquad\text{equivalently}\qquad
\mathfrak D(a,h)=0
\]
is the graph of a unique real-analytic branch
\[
h=\sigma(a),
\qquad
\sigma(b)=b,
\qquad
\sigma(a)\neq a
\quad (a\neq b).
\]

\item One has the exact identity
\[
D_{f_a}(a,h)
=
D_{u_5}(a,h)-\mathcal Z(x(a))\,D_\lambda(a,h)
=
D_\lambda(a,h)\,\bigl(\Psi_a(h)-\mathcal Z(x(a))\bigr).
\]

\item Along the critical branch,
\[
\Psi_a(\sigma(a))-\mathcal Z(x(a))
=
-\frac{\rho(0)}{2^{m+1}(m+2)}
\bigl(x(\sigma(a))-x(a)\bigr)^{m+1}
+o\!\left(\bigl|x(\sigma(a))-x(a)\bigr|^{m+1}\right).
\]

\item Consequently,
\[
D_{f_a}(a,\sigma(a))
=
-\frac{c(b)^4J_\lambda(b)\rho(0)}{2^{m+2}(m+2)M(b)^2}
\bigl(x(\sigma(a))-x(a)\bigr)^{m+3}
+o\!\left(\bigl|x(\sigma(a))-x(a)\bigr|^{m+3}\right).
\]
\end{enumerate}

In particular, for all sufficiently small $a\neq b$ one has
\[
D_{f_a}(a,\sigma(a))\neq 0,
\]
and indeed
\[
\operatorname{sgn}D_{f_a}(a,\sigma(a))
=
-\operatorname{sgn}\!\bigl(J_\lambda(b)\rho(0)\bigr).
\]
So near an isolated odd zero of $\mathcal R$, the born repeated-source
critical branch
\[
\mathfrak D(a,h)=0
\]
is uniformly unrooted in source-determinant language:
\[
D_{u_5}(a,h)-\mathcal Z(x(a))\,D_\lambda(a,h)\neq 0.
\]
\end{proposition}

\begin{proof}
By Remark~\ref{rem:source-determinant-reformulation-Gammaa},
\[
\Psi_a(h)=\Gamma_{x(a)}(x(h)).
\]
Therefore the critical locus
\[
\partial_h\Psi_a(h)=0
\]
is exactly the pullback through the source-to-$x$ coordinate of the local
critical branch
\[
\partial_x\Gamma_{x(a)}(x)=0.
\]
So item (1) is the source-coordinate form of
Proposition~\ref{prop:odd-zero-critical-branch-uniformly-unrooted}.

For item (2), linearity of the determinant in the third column gives
\[
D_{f_a}(a,h)
=
D_{u_5}(a,h)-\mathcal Z(x(a))\,D_\lambda(a,h).
\]
Dividing by $D_\lambda(a,h)\neq 0$ on the nonmixed patch yields
\[
\frac{D_{f_a}(a,h)}{D_\lambda(a,h)}
=
\Psi_a(h)-\mathcal Z(x(a)).
\]

For item (3), apply
Proposition~\ref{prop:odd-zero-critical-branch-uniformly-unrooted} to the
$x$-coordinate pair
\[
\bigl(x(a),x(\sigma(a))\bigr).
\]
Since
\[
\Psi_a(\sigma(a))=\Gamma_{x(a)}(x(\sigma(a))),
\]
that proposition gives exactly
\[
\Psi_a(\sigma(a))-\mathcal Z(x(a))
=
-\frac{\rho(0)}{2^{m+1}(m+2)}
\bigl(x(\sigma(a))-x(a)\bigr)^{m+1}
+o\!\left(\bigl|x(\sigma(a))-x(a)\bigr|^{m+1}\right).
\]

For item (4), apply the row-reduction identity from
Remark~\ref{rem:source-determinant-reformulation-Gammaa} to
\[
f=f_a,
\qquad
F=\frac{f_a}{c}=Y-\mathcal Z(x(a))S.
\]
This gives
\[
D_{f_a}(a,h)
=
c(a)^2c(h)x'(a)\,
\widetilde H_{3,\mathcal Z(x(a))}\bigl(x(a),x(h)\bigr),
\]
and since
\[
\widetilde H_{3,\mathcal Z(x(a))}(x(a),x(h))
=
B_{x(a)}(x(h))\bigl(\Psi_a(h)-\mathcal Z(x(a))\bigr),
\]
we obtain
\[
D_{f_a}(a,h)
=
c(a)^2c(h)x'(a)\,B_{x(a)}(x(h))
\bigl(\Psi_a(h)-\mathcal Z(x(a))\bigr).
\]

Along the branch $h=\sigma(a)$, write
\[
\delta(a):=x(\sigma(a))-x(a).
\]
Then
\[
B_{x(a)}(x(\sigma(a)))
=
\frac12\,S''(b)\,\delta(a)^2
+
o\!\left(|\delta(a)|^2\right),
\]
while
\[
x'(b)=\frac{M(b)}{c(b)^2},
\qquad
S''(b)=\frac{c(b)^3}{M(b)^3}J_\lambda(b)
\]
by Remark~\ref{rem:below-lifted-wronskian-ratio-reformulation}. Therefore
\[
c(a)^2c(\sigma(a))x'(a)B_{x(a)}(x(\sigma(a)))
=
\frac{c(b)^4J_\lambda(b)}{2M(b)^2}\,\delta(a)^2
+
o\!\left(|\delta(a)|^2\right).
\]
Multiplying by item (3) gives
\[
D_{f_a}(a,\sigma(a))
=
-\frac{c(b)^4J_\lambda(b)\rho(0)}{2^{m+2}(m+2)M(b)^2}\,
\delta(a)^{m+3}
+
o\!\left(|\delta(a)|^{m+3}\right),
\]
which is the claimed asymptotic. Since $m+3$ is even, the leading term has
fixed nonzero sign for all sufficiently small $a\neq b$.
\end{proof}

\begin{remark}[Ceiling of the local shear-transfer route]
\label{rem:ceiling-local-shear-transfer-route}
Proposition~\ref{prop:source-determinant-form-odd-zero-unrooted-defect} shows
that near an isolated odd zero of $\mathcal R$, the genuinely born local
repeated-source branch is the critical branch
\[
\mathfrak D(a,h)=0
\qquad\text{equivalently}\qquad
\partial_h\Psi_a(h)=0,
\]
but its rooted numerator
\[
D_{f_a}(a,h)
=
D_{u_5}(a,h)-\mathcal Z(x(a))\,D_\lambda(a,h)
\]
has a fixed nonzero leading term there.

So the rooted sheared repeated-node tangent-contact condition is not the local
source avatar of diagonal odd-zero birth.  In particular, no argument using
only the local odd-zero model and the repeated-source critical-branch
equations can prove an implication of the form
\[
\mathfrak D(a,h)=0 \Longrightarrow D_{f_a}(a,h)=0
\]
near odd zeros.

This is a different obstruction from
Remark~\ref{rem:ceiling-of-current-source-level-reduction}: the latter says the
generic source-level Wronskian/$4\times4$ package is formally too free, while
the present remark says that even along the \emph{actual local odd-zero branch},
the rooted numerator fails to vanish with a universal signed defect.  Any
successful shear-transfer theorem must therefore use genuinely additional RH
source coupling beyond the local odd-zero critical-branch analysis.
\end{remark}

\begin{proposition}[Repeated-source first loss is exactly diagonal odd-zero birth]
\label{prop:first-loss-is-diagonal-odd-zero-birth}
Let $J$ be a connected opened nonmixed arc, so
\[
K_v\neq 0
\qquad\text{on }J.
\]
For $a<x$ in $J$, define
\[
A_a(x):=Y(x)-Y(a)-Y'(a)(x-a),
\qquad
B_a(x):=S(x)-S(a)-S'(a)(x-a),
\]
\[
\Gamma_a(x):=\frac{A_a(x)}{B_a(x)},
\qquad
M(a,x):=B_a(x)^2\,\Gamma_a'(x).
\]
Equivalently,
\[
M(a,x)=A_a'(x)B_a(x)-A_a(x)B_a'(x).
\]
Let
\[
\Sigma_M:=\{(a,x)\in J^2:\ a<x,\ M(a,x)=0\}.
\]

Then the following hold.

\begin{enumerate}
\item \emph{No attained separated first loss.}
There is no pair $(a_*,x_*)\in \Sigma_M$ such that
\[
x_*=\inf\{x:\exists\,a<x\text{ with }(a,x)\in \Sigma_M\}.
\]

\item \emph{Any first-loss sequence is diagonal.}
If $(a_n,x_n)\in \Sigma_M$ and
\[
x_n\downarrow x_*:=\inf\{x:\exists\,a<x\text{ with }(a,x)\in \Sigma_M\},
\]
then
\[
x_n-a_n\to 0.
\]

\item \emph{Diagonal accumulation forces a zero of $\mathcal R$.}
If $(a_n,x_n)\in\Sigma_M$ and
\[
a_n,x_n\to b,
\qquad
x_n-a_n\to 0,
\]
then
\[
\mathcal R(b)=0.
\]

\item \emph{Even-order zeros are diagonally sterile.}
If $b$ is an isolated even-order zero of $\mathcal R$, then there exists a
neighborhood $U$ of $b$ such that
\[
\Sigma_M\cap (U\times U)=\varnothing.
\]

\item \emph{Odd-order zeros seed diagonal first loss.}
If $b$ is an isolated odd-order zero of $\mathcal R$, then every
neighborhood $U$ of $b$ contains points
\[
(a,x)\in \Sigma_M
\qquad\text{with}\qquad
a<x,
\qquad
a,x\in U.
\]
\end{enumerate}

Accordingly, repeated-source first loss on a nonmixed arc is not an
independent obstruction mechanism: it occurs if and only if one diagonally
approaches an odd-order zero of $\mathcal R$.
\end{proposition}

\begin{proof}
For (1), suppose $(a_*,x_*)\in\Sigma_M$ were an attained separated first-loss
pair. Then for fixed $a_*$ one has
\[
M(a_*,x)>0
\qquad (a_*<x<x_*),
\]
by the definition of $x_*$. Since $B_{a_*}(x)\neq 0$ on a nonmixed arc,
this implies
\[
\Gamma_{a_*}'(x)>0
\qquad (a_*<x<x_*),
\]
so $\Gamma_{a_*}$ is strictly increasing on $(a_*,x_*)$. Therefore
\[
\Gamma_{a_*}(x_*)>\lim_{x\downarrow a_*}\Gamma_{a_*}(x)=\mathcal Z(a_*).
\]
But Remark~\ref{rem:first-loss-rigidification-nonmixed} shows that every
formally attained first loss satisfies
\[
\Gamma_{a_*}(x_*)=\mathcal Z(a_*),
\]
a contradiction. Thus no attained separated first-loss pair exists.

For (2), suppose $(a_n,x_n)\in\Sigma_M$ and $x_n\downarrow x_*$, but that
$x_n-a_n\not\to 0$. Passing to a subsequence, we may assume
\[
a_n\to a_*<x_*.
\]
By continuity of $M$,
\[
M(a_*,x_*)=0,
\]
so $(a_*,x_*)\in\Sigma_M$, contradicting (1). Hence $x_n-a_n\to 0$.

For (3), since $B_a(x)\neq 0$ on a nonmixed arc, the condition
$(a_n,x_n)\in\Sigma_M$ is equivalent to
\[
\Gamma_{a_n}'(x_n)=0.
\]
In the quadratic-envelope notation from
Proposition~\ref{prop:quadratic-envelope-factorization-nonmixed},
\[
\Gamma_a'(x)
=
\frac{1}{B_a(x)^2}\int_a^x \Omega_{a,x}(r)\,\mathcal R(r)\,dr.
\]
Equivalently, after local nonmixed normalization at $b$, this becomes the
repeated-node profile derivative formula
\[
D_a'(x)=2\int_0^1 s(1-s)\,\mathcal R(a+s(x-a))\,ds.
\]
Letting $a_n,x_n\to b$ with $x_n-a_n\to 0$, continuity yields
\[
0=\lim_{n\to\infty}D_{a_n}'(x_n)
=
2\!\left(\int_0^1 s(1-s)\,ds\right)\mathcal R(b)
=
\frac13\,\mathcal R(b).
\]
Hence $\mathcal R(b)=0$.

For (4), if $b$ is an isolated even-order zero of $\mathcal R$, then
Proposition~\ref{prop:parity-classification-isolated-nonmixed-R-zeros} gives a
neighborhood $U$ of $b$ on which every repeated-node profile $D_a$ is
strictly monotone. Therefore $D_a'(x)\neq 0$ for $a\neq x$ in $U$, hence
$\Gamma_a'(x)\neq 0$ and
\[
\Sigma_M\cap (U\times U)=\varnothing.
\]

For (5), if $b$ is an isolated odd-order zero of $\mathcal R$, then
Proposition~\ref{prop:parity-classification-isolated-nonmixed-R-zeros} gives,
after local normalization, a local repeated-node extremum branch for the
profiles $D_a$. In particular, every neighborhood $U$ of $b$ contains
points $a<x$ with
\[
D_a'(x)=0.
\]
Since $B_a(x)\neq 0$ on the nonmixed arc, this is equivalent to
\[
\Gamma_a'(x)=0,
\]
hence $(a,x)\in\Sigma_M\cap (U\times U)$.

The final statement is the combination of (2), (4), and (5).
\end{proof}

\begin{remark}[Effect on the rigidified first-loss lane]
\label{rem:effect-of-shear-transfer-on-first-loss}
Proposition~\ref{prop:first-loss-is-diagonal-odd-zero-birth} shows that the
repeated-source / first-loss mechanism is not an independent obstruction:
no attained separated first loss can occur, every first-loss sequence is forced
to the diagonal, and diagonal first loss occurs exactly at isolated odd-order
zeros of $\mathcal R$.

At the same time,
Corollary~\ref{cor:sheared-event-impossible-near-isolated-odd-R-zero} shows
that near an isolated odd zero of $\mathcal R$, the rooted sheared
repeated-node tangent-contact event itself is locally absent; and
Corollary~\ref{cor:diagonal-sheared-recurrence-forces-affine-planar-collapse}
shows that any genuine diagonal realization of the sheared event already forces
the affine-planar collapse case
\[
\mathcal R\equiv 0.
\]
So the sheared bottleneck closes in the collapse direction, not by a purely
planar contradiction for $(x,S)$.

Proposition~\ref{prop:uniform-fixed-level-nonrecurrence-sheared-three-point}
shows moreover that for each fixed shear level $\gamma$, the corresponding
sheared opposite-side $3$-point branch is already globally rigid on its maximal
slope-overlap block: it has no interior common-tangent event, no interior
recurrence, and no interior accumulation. Thus the remaining issue is not
sheared $3$-point recurrence at fixed $\gamma$, but the source-coupling /
level-selection step that would force the distinguished level
\[
\gamma=\mathcal Z(b)
\]
arising from a hypothetical odd zero of $\mathcal W_{\mathrm{src}}$ to be hit
by the relevant critical-value configuration of $\Gamma_a$.

Remark~\ref{rem:no-go-fixed-level-chamber-exit-route} shows that this
fixed-level route cannot be pushed further by planar chamber-exit arguments
alone: even with two nearby roots of $\mathcal Z=\gamma$, the sheared
opposite-side branch may run to the second curvature zero by boundary exit
without any interior common-tangent event. Thus the remaining step must be
genuinely source-coupled, involving the source-pinned moving shear
\[
f_a=u_5-\mathcal Z(a)\lambda
\]
or an equivalent RH-specific coupling, rather than only the frozen-level
geometry of
\[
\widetilde Y_\gamma=Y-\gamma S.
\]
\end{remark}

\begin{remark}[The repeated-source lane is now closed]
\label{rem:repeated-source-lane-closed}
Proposition~\ref{prop:first-loss-is-diagonal-odd-zero-birth} shows that the
repeated-source / first-loss mechanism does not define an additional nonmixed
obstruction beyond the odd-zero problem for $\mathcal R$. No attained
separated first loss can occur, every first-loss sequence is forced to the
diagonal, and diagonal first loss occurs exactly at isolated odd-order zeros of
$\mathcal R$. Thus the genuinely remaining nonmixed burden is purely the
global/source-level control of odd zeros of $\mathcal R$, equivalently of the
source-level Wronskian $\mathcal W_{\mathrm{src}}$.
\end{remark}

\begin{proposition}[Single-sign-change source segment cannot realize rigidified first loss]
\label{prop:single-sign-change-cannot-realize-first-loss}
Let $J$ be an opened nonmixed arc, and let $a<x$ in $J$.
After local nonmixed normalization on $[a,x]$, write
\[
R_{a,x}(s):=\mathcal R(a+s(x-a)),
\qquad 0\le s\le 1.
\]
Then
\[
\Gamma_a'(x)
=
2\int_0^1 s(1-s)\,R_{a,x}(s)\,ds,
\]
and
\[
\Gamma_a(x)-\mathcal Z(a)
=
(x-a)\int_0^1 (1-s)^2\,R_{a,x}(s)\,ds.
\]

Assume that $R_{a,x}$ has exactly one sign change on $(0,1)$.
If
\[
\Gamma_a'(x)=0,
\]
then
\[
\operatorname{sgn}\bigl(\Gamma_a(x)-\mathcal Z(a)\bigr)
\]
is the sign of $R_{a,x}$ on the initial side of the sign change.
In particular,
\[
\Gamma_a(x)\neq \mathcal Z(a).
\]

Equivalently, no pair $(a,x)$ for which $\mathcal R$ changes sign exactly once
on $(a,x)$ can satisfy simultaneously
\[
\Gamma_a'(x)=0,
\qquad
\Gamma_a(x)=\mathcal Z(a).
\]
Hence any rigidified first-loss segment must cross at least two sign changes of
$\mathcal R$.
\end{proposition}

\begin{proof}
In the normalized nonmixed chart, $S(t)=t^2$, so
$\Gamma_a = D_a$ with
\[
D_a(t)=\int_0^1 (1-s)\,\widetilde Y''(a+s(t-a))\,ds,
\qquad
D_a(a)=\mathcal Z(a).
\]
Differentiating in $t$ gives
\[
D_a'(t)=\int_0^1 s(1-s)\,\widetilde Y'''(a+s(t-a))\,ds
=2\int_0^1 s(1-s)\,\mathcal R(a+s(t-a))\,ds,
\]
which is the first displayed formula at $t=x$. Integrating
$\widetilde Y''(a+s(t-a))$ by parts in $s$ yields
\[
D_a(t)-D_a(a)
=
(t-a)\int_0^1 (1-s)^2\,\mathcal R(a+s(t-a))\,ds,
\]
which is the second displayed formula at $t=x$.

Let
\[
k_1(s):=s(1-s),\qquad k_0(s):=(1-s)^2,
\]
so that
\[
\frac{k_0(s)}{k_1(s)}=\frac{1-s}{s}
\]
is strictly decreasing on $(0,1)$.

Write $R_{a,x}=\sigma\rho$ with $\rho\ge 0$ and $\sigma$ changing sign exactly
once at some $s_0\in(0,1)$. By hypothesis,
\[
0=\int_0^1 k_1(s)\,\sigma(s)\rho(s)\,ds.
\]
Split
\[
I:=\int_0^1 k_0(s)\,\sigma(s)\rho(s)\,ds
=
\int_0^{s_0}\frac{k_0}{k_1}(s)\,k_1(s)\sigma(s)\rho(s)\,ds
+
\int_{s_0}^1\frac{k_0}{k_1}(s)\,k_1(s)\sigma(s)\rho(s)\,ds.
\]
Let $\epsilon:=\operatorname{sgn}(\sigma|_{(0,s_0)})$. Then
$k_1\sigma\rho\ge 0$ with sign $\epsilon$ on $(0,s_0)$ and sign $-\epsilon$ on
$(s_0,1)$. Because $k_0/k_1$ is strictly decreasing, it is bounded below by
its value at $s_0$ on $(0,s_0)$ and bounded above by its value at $s_0$ on
$(s_0,1)$. Using
\[
\int_0^{s_0}k_1\sigma\rho=-\int_{s_0}^1 k_1\sigma\rho
\]
and strict monotonicity on the support of $\rho$ gives
\[
\epsilon\,I>0.
\]
Hence
\[
\Gamma_a(x)-\mathcal Z(a)=(x-a)\,I\neq 0,
\]
with sign $\epsilon$ (after orienting $x>a$). This proves the claim.
\end{proof}

\begin{corollary}[Source-level two-sign-change obstruction]
\label{cor:source-level-two-sign-change-obstruction}
Let $a<h$ lie on an opened nonmixed arc in source coordinates.
Assume
\[
\Psi_a'(h)=0.
\]
If the source-level Wronskian
\[
\mathcal W_{\mathrm{src}}
=
J_\lambda J_{u_5}'-J_{u_5}J_\lambda'
\]
changes sign exactly once on the source interval $[a,h]$, then
\[
\Psi_a(h)\neq \mathcal Z(a).
\]
Equivalently, the simultaneous source-level first-loss conditions
\[
\Psi_a'(h)=0,
\qquad
\Psi_a(h)=\mathcal Z(a)
\]
force at least two sign changes of $\mathcal W_{\mathrm{src}}$ on $[a,h]$.
\end{corollary}

\begin{proof}
By Remark~\ref{rem:source-determinant-reformulation-Gammaa},
$\Psi_a(h)=\Gamma_{x(a)}(x(h))$, and on the nonmixed arc
$\mathcal W_{\mathrm{src}}$ and $\mathcal R$ differ only by a nonvanishing
analytic factor
(Proposition~\ref{prop:source-level-Wronskian-nonmixed-four-point}). So a
single sign change of $\mathcal W_{\mathrm{src}}$ on $[a,h]$ pulls back to a
single sign change of $\mathcal R$ on $[x(a),x(h)]$. Apply
Proposition~\ref{prop:single-sign-change-cannot-realize-first-loss}.
\end{proof}

\begin{proposition}[Transport equations for a repeated-source critical branch]
\label{prop:transport-equations-repeated-source-critical-branch}
Let $J$ be an opened nonmixed arc, and work in a local normalized nonmixed
chart on which
\[
S(t)=t^2,
\qquad
\Gamma_a=D_a.
\]
Let $U\subset J$ be an interval, and let
\[
\tau:U\to J
\]
be a real-analytic function such that
\[
\tau(a)>a,
\qquad
D_a'(\tau(a))=0
\qquad (a\in U).
\]
Set
\[
V(a):=D_a(\tau(a)).
\]
Then:
\begin{enumerate}
\item
\[
V'(a)=\frac{2}{\tau(a)-a}\,\bigl(V(a)-\mathcal Z(a)\bigr).
\]

\item
\[
D_a''(\tau(a))
=
\frac{2}{(\tau(a)-a)^2}\,\bigl(\mathcal Z(\tau(a))-V(a)\bigr).
\]

\item
At any point $a\in U$ where $D_a''(\tau(a))\neq 0$,
\[
\tau'(a)
=
\frac{V(a)-\mathcal Z(a)}
{\mathcal Z(\tau(a))-V(a)}.
\]
\end{enumerate}
\end{proposition}

\begin{proof}
Write
\[
D_a(t)
=
\frac{\widetilde Y(t)-\widetilde Y(a)-\widetilde Y'(a)(t-a)}{(t-a)^2}.
\]
Differentiating in $a$ gives
\[
\partial_a D_a(t)
=
\frac{2}{t-a}\,\bigl(D_a(t)-\mathcal Z(a)\bigr).
\]
Evaluating at $t=\tau(a)$ and using $D_a'(\tau(a))=0$ yields
\[
V'(a)
=
\frac{d}{da}D_a(\tau(a))
=
\partial_a D_a(\tau(a))
=
\frac{2}{\tau(a)-a}\,\bigl(V(a)-\mathcal Z(a)\bigr),
\]
which proves (1).

Now use the rigidification scalar
\[
M(a,t):=B_a(t)^2\,\Gamma_a'(t).
\]
In the normalized chart,
\[
B_a(t)=(t-a)^2,
\qquad
\Gamma_a=D_a,
\qquad
M(a,t)=(t-a)^4\,D_a'(t).
\]
By Remark~\ref{rem:first-loss-rigidification-nonmixed},
\[
\partial_t M(a,t)
=
2(t-a)^2\,\bigl(\mathcal Z(t)-D_a(t)\bigr).
\]
At $t=\tau(a)$ we have $D_a'(\tau(a))=0$, so
\[
\partial_t M(a,\tau(a))
=
(\tau(a)-a)^4\,D_a''(\tau(a)).
\]
Comparing the two formulas gives
\[
(\tau(a)-a)^4\,D_a''(\tau(a))
=
2(\tau(a)-a)^2\,\bigl(\mathcal Z(\tau(a))-V(a)\bigr),
\]
hence (2).

Finally, still by Remark~\ref{rem:first-loss-rigidification-nonmixed},
at points where $M(a,t)=0$ one has
\[
\partial_a M(a,t)
=
2(t-a)^2\,\bigl(\mathcal Z(a)-D_a(t)\bigr),
\]
in the normalized chart. Since
\[
M(a,\tau(a))=0,
\]
differentiating in $a$ gives
\[
0=\frac{d}{da}M(a,\tau(a))
=
\partial_a M(a,\tau(a))
+
\partial_t M(a,\tau(a))\,\tau'(a).
\]
Substituting the formulas above yields
\[
0
=
2(\tau(a)-a)^2\,\bigl(\mathcal Z(a)-V(a)\bigr)
+
2(\tau(a)-a)^2\,\bigl(\mathcal Z(\tau(a))-V(a)\bigr)\tau'(a).
\]
If $D_a''(\tau(a))\neq 0$, then by (2),
\[
\mathcal Z(\tau(a))-V(a)\neq 0,
\]
so solving for $\tau'(a)$ gives (3).
\end{proof}

\begin{corollary}[One-sign-change chamber is no-return for the born critical branch]
\label{cor:one-sign-change-chamber-no-return}
In the setting of
Proposition~\ref{prop:transport-equations-repeated-source-critical-branch},
assume moreover that for every $a\in U$ the segment
\[
[a,\tau(a)]
\]
contains exactly one sign change of $\mathcal R$.

Then for every $a\in U$:
\begin{enumerate}
\item
\[
\operatorname{sgn}\bigl(V(a)-\mathcal Z(a)\bigr)
\]
is the sign of the initial lobe of $\mathcal R$ on $[a,\tau(a)]$;

\item
\[
\operatorname{sgn}\bigl(\mathcal Z(\tau(a))-V(a)\bigr)
\]
is the sign of the terminal lobe of $\mathcal R$ on $[a,\tau(a)]$;

\item
\[
\bigl(V(a)-\mathcal Z(a)\bigr)\bigl(\mathcal Z(\tau(a))-V(a)\bigr)<0;
\]

\item
\[
D_a''(\tau(a))\neq 0,
\qquad
\tau'(a)<0;
\]

\item
\[
V(a)\notin [\,\mathcal Z(a),\mathcal Z(\tau(a))\,].
\]
In particular,
\[
V(a)\neq \mathcal Z(a),
\]
so the branch cannot realize the rigidified first-loss condition anywhere on
$U$.
\end{enumerate}
\end{corollary}

\begin{proof}
The first sign statement is exactly
Proposition~\ref{prop:single-sign-change-cannot-realize-first-loss} applied to
the critical pair $(a,\tau(a))$.

For the second sign statement, write
\[
R_{a,\tau(a)}(s):=\mathcal R(a+s(\tau(a)-a)).
\]
Since
\[
\mathcal Z(\tau(a))-\mathcal Z(a)
=
(\tau(a)-a)\int_0^1 R_{a,\tau(a)}(s)\,ds
\]
and
\[
V(a)-\mathcal Z(a)
=
(\tau(a)-a)\int_0^1 (1-s)^2\,R_{a,\tau(a)}(s)\,ds,
\]
one gets
\[
\mathcal Z(\tau(a))-V(a)
=
(\tau(a)-a)\int_0^1 (2s-s^2)\,R_{a,\tau(a)}(s)\,ds.
\]
At the critical pair,
\[
0=D_a'(\tau(a))
=
2\int_0^1 s(1-s)\,R_{a,\tau(a)}(s)\,ds,
\]
so
\[
\int_0^1 (2s-s^2)\,R_{a,\tau(a)}(s)\,ds
=
\int_0^1 s^2\,R_{a,\tau(a)}(s)\,ds.
\]
Hence
\[
\mathcal Z(\tau(a))-V(a)
=
(\tau(a)-a)\int_0^1 s^2\,R_{a,\tau(a)}(s)\,ds.
\]
Now set
\[
k_1(s):=s(1-s),
\qquad
k_2(s):=s^2.
\]
Then
\[
\frac{k_2(s)}{k_1(s)}=\frac{s}{1-s}
\]
is strictly increasing on $(0,1)$. Repeating the weighted signed-mass argument
from the proof of
Proposition~\ref{prop:single-sign-change-cannot-realize-first-loss} shows that
\[
\operatorname{sgn}\bigl(\mathcal Z(\tau(a))-V(a)\bigr)
\]
is the sign of the terminal lobe of $\mathcal R$.

The two signs are therefore opposite, which proves (3). Then
Proposition~\ref{prop:transport-equations-repeated-source-critical-branch}(2)
gives
\[
D_a''(\tau(a))\neq 0,
\]
and
Proposition~\ref{prop:transport-equations-repeated-source-critical-branch}(3)
gives
\[
\tau'(a)<0.
\]
Finally, opposite signs of
\[
V(a)-\mathcal Z(a)
\quad\text{and}\quad
\mathcal Z(\tau(a))-V(a)
\]
mean exactly that $V(a)$ lies strictly outside the endpoint-curvature interval
\[
[\,\mathcal Z(a),\mathcal Z(\tau(a))\,].
\]
\end{proof}

\begin{proposition}[First rooting requires a second odd zero along the born repeated-source branch]
\label{prop:first-rooting-requires-second-odd-zero-repeated-source}
Let $b$ be an isolated odd-order zero of $\mathcal R$ on an opened nonmixed
arc, and let
\[
t=\tau(a)
\]
denote an analytic continuation of the local born repeated-source critical
branch, so that
\[
\tau(a)>a,
\qquad
D_a'(\tau(a))=0.
\]
Assume that $U$ is a connected continuation interval on which, for every
$a\in U$, the segment
\[
[a,\tau(a)]
\]
contains exactly one sign change of $\mathcal R$.

Then for every $a\in U$ one has
\[
D_a(\tau(a))\neq \mathcal Z(a).
\]
Equivalently, on the corresponding source branch
\[
h=\sigma(a),
\qquad
\mathfrak D(a,\sigma(a))=0,
\]
one has
\[
\Psi_a(\sigma(a))\neq \mathcal Z(x(a)),
\]
or in determinant form,
\[
D_{u_5}(a,\sigma(a))-\mathcal Z(x(a))\,D_\lambda(a,\sigma(a))\neq 0.
\]

Consequently, any first parameter at which the analytic continuation of the
born repeated-source branch realizes a genuine repeated-source obstruction
inside the opened nonmixed arc must occur only after the corresponding segment
ceases to be one-sign-change. Internally, away from affine-planar collapse and
without leaving the arc, this can happen only when the segment acquires a
second odd-order zero of $\mathcal R$, equivalently of
\[
\mathcal W_{\mathrm{src}}.
\]
\end{proposition}

\begin{proof}
The first claim is exactly
Corollary~\ref{cor:one-sign-change-chamber-no-return}(5): on any continuation
interval on which each segment
\[
[a,\tau(a)]
\]
contains exactly one sign change of $\mathcal R$, the born critical branch
cannot realize the rigidified first-loss condition
\[
D_a(\tau(a))=\mathcal Z(a).
\]

Via Remark~\ref{rem:source-determinant-reformulation-Gammaa},
\[
\Psi_a(h)=\Gamma_{x(a)}(x(h)),
\]
so the same statement in source coordinates is
\[
\Psi_a(\sigma(a))\neq \mathcal Z(x(a)).
\]
Using
Proposition~\ref{prop:source-determinant-form-odd-zero-unrooted-defect}(2),
this is equivalent to
\[
D_{u_5}(a,\sigma(a))-\mathcal Z(x(a))\,D_\lambda(a,\sigma(a))\neq 0.
\]

For the final statement, as long as the continued branch remains inside the
opened nonmixed arc and the corresponding segment remains one-sign-change,
the first part shows that rooting cannot occur. Thus any first rooted point
must be preceded by failure of the one-sign-change condition.

Now even-order zeros of $\mathcal R$ do not change sign, so internally the
sign-change count on the segment can change only when an endpoint of the
segment meets another odd-order zero of $\mathcal R$ (unless the branch exits
the arc or enters the affine-planar collapse case). Finally, on a nonmixed arc
the functions $\mathcal R$ and
\[
\mathcal W_{\mathrm{src}}
=
J_\lambda J_{u_5}'-J_{u_5}J_\lambda'
\]
differ only by a nonvanishing analytic factor
(Proposition~\ref{prop:source-level-Wronskian-nonmixed-four-point}), so odd
zeros of the two objects coincide.
\end{proof}

\begin{proposition}[Discriminant entry into the adjacent-two-odd-zero chamber is not a first rooted point]
\label{prop:discriminant-entry-not-first-root}
Let $J$ be an opened nonmixed arc, and work in a local normalized nonmixed
chart on which
\[
S(t)=t^2,
\qquad
\Gamma_a=D_a.
\]
Let
\[
\tau:(a_-,a_0]\to J
\]
be a real-analytic continuation of the born repeated-source critical branch,
so that
\[
\tau(a)>a,
\qquad
D_a'(\tau(a))=0
\qquad (a\in(a_-,a_0]).
\]
Assume:

\begin{enumerate}
\item for every $a\in(a_-,a_0)$, the segment
\[
[a,\tau(a)]
\]
contains exactly one sign change of $\mathcal R$;

\item the target endpoint reaches the next adjacent odd zero at $a=a_0$, i.e.
\[
q:=\tau(a_0)
\]
is an odd-order zero of $\mathcal R$, there are no odd zeros of $\mathcal R$
in $(\tau(a),q)$ for $a<a_0$ close to $a_0$, and the segment
\[
[a_0,q]
\]
still contains exactly one sign change of $\mathcal R$.
\end{enumerate}

Set
\[
V(a):=D_a(\tau(a)).
\]
Define the sheared discriminant and the RH-selected slice by
\[
\mathcal D:=
\left\{
(a,\gamma,h):a<h,\ D_{f_\gamma}(a,h)=0,\ \partial_h D_{f_\gamma}(a,h)=0
\right\},
\]
\[
\mathcal S_{\mathrm{RH}}
:=
\left\{
(a,\gamma,h):\gamma=\mathcal Z(a)
\right\}.
\]

Then
\[
P(a):=(a,V(a),\tau(a))\in\mathcal D
\qquad (a\in(a_-,a_0]),
\]
but at the adjacent-odd-zero entrance point one has
\[
V(a_0)-\mathcal Z(a_0)\neq 0.
\]
More precisely,
\[
V(a_0)-\mathcal Z(a_0)
=
(q-a_0)\int_0^1 (1-s)^2\,
\mathcal R(a_0+s(q-a_0))\,ds,
\]
and this quantity has the sign of the initial lobe of $\mathcal R$ on
$[a_0,q]$.

Consequently,
\[
P(a_0)\notin \mathcal S_{\mathrm{RH}},
\]
and there exists a neighborhood $N$ of $P(a_0)$ in
\[
J\times\mathbf R\times J
\]
such that
\[
N\cap\mathcal D\cap\mathcal S_{\mathrm{RH}}=\varnothing.
\]

In particular, a first rooted point cannot be created at the moment when the
next adjacent odd zero first enters the active segment. If rooting occurs at
all, it must occur only after the branch has moved into a genuine two-sign-change
chamber of positive width.
\end{proposition}

\begin{proof}
By Remark~\ref{rem:source-determinant-reformulation-Gammaa},
\[
\Psi_a(h)=\Gamma_{x(a)}(x(h)),
\]
and on a nonmixed patch
\[
D_{f_\gamma}(a,h)=D_\lambda(a,h)\bigl(\Psi_a(h)-\gamma\bigr).
\]
Since
\[
D_\lambda(a,h)\neq 0
\qquad (a<h),
\]
we have
\[
D_{f_\gamma}(a,h)=0,\ \partial_hD_{f_\gamma}(a,h)=0
\iff
\gamma=\Psi_a(h),\ \partial_h\Psi_a(h)=0.
\]
Along the continued born branch,
\[
\partial_h\Psi_a(\tau(a))=0,
\qquad
\Psi_a(\tau(a))=V(a),
\]
so indeed
\[
P(a)=(a,V(a),\tau(a))\in\mathcal D.
\]

At $a=a_0$, the integral formula for $D_a$ gives
\[
V(a_0)-\mathcal Z(a_0)
=
(q-a_0)\int_0^1 (1-s)^2\,
\mathcal R(a_0+s(q-a_0))\,ds.
\]
Also, since $D_{a_0}'(q)=0$,
\[
0
=
D_{a_0}'(q)
=
2\int_0^1 s(1-s)\,
\mathcal R(a_0+s(q-a_0))\,ds.
\]

Set
\[
R_0(s):=\mathcal R(a_0+s(q-a_0)),
\qquad
k_0(s):=(1-s)^2,
\qquad
k_1(s):=s(1-s).
\]
Then
\[
\frac{k_0(s)}{k_1(s)}=\frac{1-s}{s}
\]
is strictly decreasing on $(0,1)$. Because $R_0$ has exactly one sign change on
$(0,1)$, the same weighted signed-mass argument used in
Proposition~\ref{prop:single-sign-change-cannot-realize-first-loss} implies
that
\[
\int_0^1 k_0(s)R_0(s)\,ds
\]
has the sign of the initial lobe of $R_0$, hence is nonzero. Therefore
\[
V(a_0)-\mathcal Z(a_0)\neq 0.
\]

Thus
\[
P(a_0)=(a_0,V(a_0),q)\in\mathcal D
\quad\text{but}\quad
P(a_0)\notin\mathcal S_{\mathrm{RH}}.
\]
Since both $\mathcal D$ along the continued branch and the function
\[
(a,\gamma,h)\longmapsto \gamma-\mathcal Z(a)
\]
are continuous, there is a neighborhood $N$ of $P(a_0)$ on which
\[
\gamma-\mathcal Z(a)
\]
keeps the sign of
\[
V(a_0)-\mathcal Z(a_0).
\]
Hence
\[
N\cap\mathcal D\cap\mathcal S_{\mathrm{RH}}=\varnothing.
\]

The final statement is immediate: a rooted point is exactly a point of
\[
\mathcal D\cap\mathcal S_{\mathrm{RH}},
\]
so no such point can be born at the adjacent-odd-zero entrance.
\end{proof}

\begin{remark}[What the discriminant analysis now leaves open]
\label{rem:what-discriminant-adjacent-odd-zero-analysis-still-leaves-open}
Proposition~\ref{prop:discriminant-entry-not-first-root} shows that a rooted
point is not created \emph{at the entry} into the adjacent-two-odd-zero chamber.
Corollary~\ref{cor:no-crossing-adjacent-two-odd-zero-under-anchor-sign-preservation}
shows more: on any continuation interval on which the anchor-side value
\[
\mathcal R(a)
\]
keeps a fixed sign, the RH-selected slice
\[
\gamma=\mathcal Z(a)
\]
does not cross the continued discriminant curve at all.

Thus the adjacent-two-odd-zero chamber is not itself the residual obstruction.
What remains open is narrower and genuinely global: whether a rooted crossing
can occur only after the anchor passes through an odd-order zero of
$\mathcal R$ (equivalently of $\mathcal W_{\mathrm{src}}$), or after branch
exit, or after affine-planar collapse. In other words, the remaining discriminant-theoretic bottleneck is a
\emph{delayed post-anchor no-crossing problem}: not an entry-point singularity
problem, and not crossing at the anchor odd zero itself, but a later crossing on
a dangerous post-anchor chamber.
\end{remark}

\begin{proposition}[Rooted-defect transport and stationary first root]
\label{prop:rooted-defect-transport-stationary-first-root}
In the setting of
Proposition~\ref{prop:transport-equations-repeated-source-critical-branch},
define
\[
L(a):=\tau(a)-a,
\qquad
G(a):=V(a)-\mathcal Z(a)=D_a(\tau(a))-\mathcal Z(a).
\]
Then
\[
G'(a)=\frac{2}{L(a)}\,G(a)-\mathcal R(a).
\]
Hence, if $a_*\in U$ satisfies
\[
G(a_*)=0,
\]
then
\[
G'(a_*)=-\mathcal R(a_*).
\]

If in addition
\[
\mathcal Z(\tau(a_*))\neq \mathcal Z(a_*),
\]
then
\[
D_{a_*}''(\tau(a_*))\neq 0,
\qquad
\tau'(a_*)=0,
\qquad
\tau''(a_*)=
-\frac{\mathcal R(a_*)}{\mathcal Z(\tau(a_*))-\mathcal Z(a_*)}.
\]
\end{proposition}

\begin{proof}
Since
\[
G(a)=V(a)-\mathcal Z(a),
\]
Proposition~\ref{prop:transport-equations-repeated-source-critical-branch}(1)
gives
\[
G'(a)=V'(a)-\mathcal R(a)
=\frac{2}{\tau(a)-a}\,\bigl(V(a)-\mathcal Z(a)\bigr)-\mathcal R(a)
=\frac{2}{L(a)}\,G(a)-\mathcal R(a).
\]
Evaluating at a root $a_*$ of $G$ gives
\[
G'(a_*)=-\mathcal R(a_*).
\]

Now assume also that
\[
\mathcal Z(\tau(a_*))\neq \mathcal Z(a_*).
\]
Since $G(a_*)=0$, this means
\[
\mathcal Z(\tau(a_*))\neq V(a_*).
\]
So Proposition~\ref{prop:transport-equations-repeated-source-critical-branch}(2)
yields
\[
D_{a_*}''(\tau(a_*))
=
\frac{2}{(\tau(a_*)-a_*)^2}\,
\bigl(\mathcal Z(\tau(a_*))-V(a_*)\bigr)
\neq 0.
\]
Then Proposition~\ref{prop:transport-equations-repeated-source-critical-branch}(3)
applies near $a_*$ and gives
\[
\tau'(a)
=
\frac{G(a)}{\mathcal Z(\tau(a))-V(a)}.
\]
Evaluating at $a_*$ gives
\[
\tau'(a_*)=0.
\]

Differentiate this identity once more. Since $G(a_*)=0$ and
\[
\mathcal Z(\tau(a_*))-V(a_*)=\mathcal Z(\tau(a_*))-\mathcal Z(a_*)\neq 0,
\]
one gets
\[
\tau''(a_*)
=
\frac{G'(a_*)}{\mathcal Z(\tau(a_*))-V(a_*)}
=
-\frac{\mathcal R(a_*)}{\mathcal Z(\tau(a_*))-\mathcal Z(a_*)},
\]
as claimed.
\end{proof}

\begin{proposition}[Nondegenerate rooted point is a fold of the moving-shear equation]
\label{prop:nondegenerate-rooted-point-is-moving-shear-fold}
Work in a local normalized nonmixed chart on which
\[
S(t)=t^2.
\]
Define
\[
\mathscr F(a,t):=A_a(t)-\mathcal Z(a)\,B_a(t)
=(t-a)^2\bigl(D_a(t)-\mathcal Z(a)\bigr).
\]
Let
\[
a_*<t_*
\]
satisfy
\[
\mathscr F(a_*,t_*)=0,
\qquad
\partial_t\mathscr F(a_*,t_*)=0,
\]
and assume
\[
\mathcal R(a_*)\neq 0,
\qquad
\mathcal Z(t_*)\neq \mathcal Z(a_*).
\]
Then
\[
\partial_a\mathscr F(a_*,t_*)
=
-(t_*-a_*)^2\,\mathcal R(a_*),
\]
and
\[
\partial_{tt}\mathscr F(a_*,t_*)
=
2\bigl(\mathcal Z(t_*)-\mathcal Z(a_*)\bigr).
\]
In particular,
\[
\partial_a\mathscr F(a_*,t_*)\neq 0,
\qquad
\partial_{tt}\mathscr F(a_*,t_*)\neq 0.
\]
Hence near $(a_*,t_*)$ the zero set
\[
\mathscr F(a,t)=0
\]
is the graph of a unique real-analytic function
\[
a=\alpha(t)
\]
with
\[
\alpha(t_*)=a_*,
\qquad
\alpha'(t_*)=0,
\qquad
\alpha''(t_*)
=
\frac{2\bigl(\mathcal Z(t_*)-\mathcal Z(a_*)\bigr)}
{(t_*-a_*)^2\,\mathcal R(a_*)}.
\]
Equivalently,
\[
a-a_*
=
\frac{\mathcal Z(t_*)-\mathcal Z(a_*)}
{(t_*-a_*)^2\,\mathcal R(a_*)}(t-t_*)^2
+o\!\left((t-t_*)^2\right).
\]
\end{proposition}

\begin{proof}
Since
\[
\mathscr F(a,t)=(t-a)^2\bigl(D_a(t)-\mathcal Z(a)\bigr),
\]
the conditions
\[
\mathscr F(a_*,t_*)=0,
\qquad
\partial_t\mathscr F(a_*,t_*)=0
\]
imply
\[
D_{a_*}(t_*)=\mathcal Z(a_*),
\qquad
D_{a_*}'(t_*)=0.
\]
So $(a_*,t_*)$ is a rooted critical point.

Differentiate $\mathscr F$ with respect to $a$:
\[
\partial_a\mathscr F(a,t)
=
-2(t-a)\bigl(D_a(t)-\mathcal Z(a)\bigr)
+(t-a)^2\bigl(\partial_aD_a(t)-\mathcal R(a)\bigr).
\]
At the rooted point,
\[
D_{a_*}(t_*)-\mathcal Z(a_*)=0,
\]
and by
Proposition~\ref{prop:transport-equations-repeated-source-critical-branch},
\[
\partial_aD_a(t)=\frac{2}{t-a}\bigl(D_a(t)-\mathcal Z(a)\bigr),
\]
so
\[
\partial_a\mathscr F(a_*,t_*)
=
-(t_*-a_*)^2\,\mathcal R(a_*).
\]

Next,
\[
\partial_t\mathscr F(a,t)
=
2(t-a)\bigl(D_a(t)-\mathcal Z(a)\bigr)
+(t-a)^2D_a'(t),
\]
hence
\[
\partial_{tt}\mathscr F(a,t)
=
2\bigl(D_a(t)-\mathcal Z(a)\bigr)
+4(t-a)D_a'(t)
+(t-a)^2D_a''(t).
\]
At the rooted critical point this becomes
\[
\partial_{tt}\mathscr F(a_*,t_*)
=
(t_*-a_*)^2D_{a_*}''(t_*).
\]
Applying
Proposition~\ref{prop:transport-equations-repeated-source-critical-branch}(2)
and using
\[
D_{a_*}(t_*)=\mathcal Z(a_*),
\]
one gets
\[
\partial_{tt}\mathscr F(a_*,t_*)
=
2\bigl(\mathcal Z(t_*)-\mathcal Z(a_*)\bigr).
\]

By the two nonvanishing assumptions,
\[
\partial_a\mathscr F(a_*,t_*)\neq 0.
\]
So the implicit function theorem gives a unique real-analytic graph
\[
a=\alpha(t)
\]
for the zero set of $\mathscr F$ near $(a_*,t_*)$. Since
\[
\partial_t\mathscr F(a_*,t_*)=0,
\]
one has
\[
\alpha'(t_*)=0.
\]
Differentiating
\[
\mathscr F(\alpha(t),t)=0
\]
twice at $t=t_*$ gives
\[
\partial_a\mathscr F(a_*,t_*)\,\alpha''(t_*)
+
\partial_{tt}\mathscr F(a_*,t_*)=0,
\]
hence
\[
\alpha''(t_*)
=
-\frac{\partial_{tt}\mathscr F(a_*,t_*)}
{\partial_a\mathscr F(a_*,t_*)}
=
\frac{2\bigl(\mathcal Z(t_*)-\mathcal Z(a_*)\bigr)}
{(t_*-a_*)^2\,\mathcal R(a_*)}.
\]
The Taylor expansion of $\alpha$ then yields the stated quadratic form.
\end{proof}

\begin{remark}[First rooted point: local fold data versus the actual remaining burden]
\label{rem:first-rooted-point-post-birth-fold-problem}
Let $b$ be an isolated odd-order zero of $\mathcal R$, and let
\[
t=\tau(a)
\]
be an analytic continuation of the born repeated-source critical branch. By
Proposition~\ref{prop:first-rooting-requires-second-odd-zero-repeated-source},
any first rooted point along this continuation can occur only after the segment
has acquired a second odd-order zero of $\mathcal R$.

At such a rooted point, if
\[
\mathcal R(a_*)\neq 0,
\qquad
\mathcal Z(\tau(a_*))\neq \mathcal Z(a_*),
\]
then
Proposition~\ref{prop:rooted-defect-transport-stationary-first-root} shows that
the rooted defect
\[
G(a)=D_a(\tau(a))-\mathcal Z(a)
\qquad
\text{equivalently}\qquad
G(a)=\Psi_a(\sigma(a))-\mathcal Z(a)
\]
crosses transversely and that the target map is stationary there:
\[
\tau'(a_*)=0.
\]
Moreover,
Proposition~\ref{prop:nondegenerate-rooted-point-is-moving-shear-fold} shows
that the moving-shear equation
\[
A_a(x)-\mathcal Z(a)\,B_a(x)=0
\qquad\text{equivalently}\qquad
D_{f_a}(a,h)=0,
\quad
f_a=u_5-\mathcal Z(a)\lambda,
\]
has a local fold at such a nondegenerate rooted point.

Thus the transport/fold package determines the \emph{local shape} of any
prospective rooted crossing. However, the later route-2 results show that this
is not the final obstruction theorem. Proposition~\ref{prop:discriminant-entry-not-first-root}
rules out birth of the first rooted point at entry into the adjacent-two-odd-zero
chamber; Proposition~\ref{prop:anchor-sign-preservation-excludes-first-rooting}
and Corollary~\ref{cor:no-crossing-adjacent-two-odd-zero-under-anchor-sign-preservation}
exclude rooted crossing on continuation intervals where the anchor-side value
keeps fixed sign; Proposition~\ref{prop:odd-anchor-passage-no-reset} rules out
rooted crossing at an odd anchor zero itself or immediately after passage; and
Proposition~\ref{prop:universality-threshold-current-route-two-package}
together with
Corollary~\ref{cor:current-route-two-package-universal-under-normalization}
show that even delayed rooted crossing on dangerous post-anchor chambers cannot
be excluded from the current planar/route-2 package alone.

Accordingly, this remark should be read only as local fold data. The remaining
route-2 burden is no longer a ``first-fold exclusion'' theorem, but the search
for the first genuinely non-universal source-coupled input beyond the normalized
planar model $(\mathcal Z,\mathcal R,D_a)$.
\end{remark}

\begin{proposition}[At most one odd zero implies no genuine repeated-source obstruction]
\label{prop:at-most-one-odd-zero-no-genuine-repeated-source-obstruction}
Let $J$ be an opened nonmixed arc. Assume that $\mathcal R$ has at most one
odd-order zero on $J$. Then there is no pair $a<x$ in $J$ such that
\[
\Gamma_a'(x)=0,
\qquad
\Gamma_a(x)=\mathcal Z(a).
\]
Equivalently, there is no source-level pair $a<h$ on the corresponding source
arc such that
\[
\Psi_a'(h)=0,
\qquad
\Psi_a(h)=\mathcal Z(a).
\]
\end{proposition}

\begin{proof}
On a nonmixed arc, $\mathcal W_{\mathrm{src}}$ differs from $\mathcal R$ only
by a nonvanishing analytic factor
(Proposition~\ref{prop:source-level-Wronskian-nonmixed-four-point}), so odd
zeros and sign changes of the two objects coincide.

If $\mathcal R$ has at most one odd-order zero on $J$, then every segment
contained in $J$ contains at most one sign change of $\mathcal R$ (equivalently
of $\mathcal W_{\mathrm{src}}$), since even-order zeros do not change sign.
Hence Corollary~\ref{cor:source-level-two-sign-change-obstruction} excludes any
simultaneous source-level solution of
\[
\Psi_a'(h)=0,
\qquad
\Psi_a(h)=\mathcal Z(a).
\]
Via Remark~\ref{rem:source-determinant-reformulation-Gammaa}, this is
equivalent to the planar formulation
\[
\Gamma_a'(x)=0,
\qquad
\Gamma_a(x)=\mathcal Z(a).
\]
\end{proof}

\begin{remark}[Minimal unresolved configuration on this route]
\label{rem:minimal-unresolved-configuration-repeated-source}
Proposition~\ref{prop:at-most-one-odd-zero-no-genuine-repeated-source-obstruction}
rules out every opened nonmixed arc carrying at most one odd zero of
$\mathcal R$, and
Corollary~\ref{cor:no-crossing-adjacent-two-odd-zero-under-anchor-sign-preservation}
rules out rooted crossing on the standard adjacent-two-odd-zero chamber as long
as the anchor-side value $\mathcal R(a)$ keeps a fixed sign.

The new odd-anchor passage theorem
(Proposition~\ref{prop:odd-anchor-passage-no-reset})
shows that no rooted crossing can occur at the anchor odd zero itself or
immediately after it is crossed, while
Corollary~\ref{cor:first-post-anchor-chamber-dangerous}
shows that the subsequent post-anchor chambers alternate dangerous and safe,
beginning with a dangerous first post-anchor chamber.

Remark~\ref{rem:no-go-local-scalar-first-dangerous-chamber} shows that the
first dangerous chamber is already formally compatible with the local scalar
rooted equations and fold data. Accordingly, the first genuinely unresolved
configuration is no longer the adjacent-two-odd-zero chamber itself, nor the
anchor odd zero passage, but a \emph{delayed} crossing on a dangerous
post-anchor chamber. Thus the sharpened source-coupling target is a delayed
no-crossing / profile-realizability exclusion theorem for the moving shear
\[
f_a=u_5-\mathcal Z(a)\lambda
\]
on genuine opened nonmixed arcs.
\end{remark}

\begin{remark}[No-go for a pure level-set / barrier theorem for $\mathcal Z$]
\label{rem:no-go-pure-level-set-barrier-for-Z}
A natural next guess would be that the remaining adjacent-odd-zero case could be
excluded by a theorem depending only on the level geometry of
\[
\mathcal Z,\qquad \mathcal R=\mathcal Z',
\]
together with the first-root transport identities
\[
G(a_*):=V(a_*)-\mathcal Z(a_*)=0,
\qquad
G'(a_*)=-\mathcal R(a_*),
\qquad
\tau'(a_*)=0,
\]
and the fact that a rooted point forces two interior recrossings of the level
\[
\mathcal Z(r)=\mathcal Z(a_*).
\]
At the current level of abstraction this is false: the level-set geometry of
$\mathcal Z$ alone is still too free.

Indeed, already in the normalized planar determinant model one can realize all
of the above with exactly two simple odd zeros of $\mathcal R$ and with a
nondegenerate rooted critical pair. Take
\[
S(t)=t^2,
\qquad
\mathcal R(t)=5t^3+5t^2-7t+1,
\qquad
\mathcal Z(t)=\int_0^t \mathcal R(u)\,du
=
t-\frac72 t^2+\frac53 t^3+\frac54 t^4.
\]
Then $\mathcal R$ has exactly two simple zeros in $(0,1)$, numerically
\[
t\approx 0.165726,
\qquad
t\approx 0.660736,
\]
so on $[0,1]$ the sign pattern is $+,-,+$ with exactly two odd sign changes.
Now choose an analytic function $\widetilde Y$ such that
\[
\frac12\,\widetilde Y''=\mathcal Z.
\]
For the normalized determinant function $D_a(x)$ associated to
$(x,\widetilde Y)$, the rooted critical conditions at $(a,x)=(0,1)$ become
\[
D_0'(1)=2\int_0^1 s(1-s)\,\mathcal R(s)\,ds,
\qquad
D_0(1)-\mathcal Z(0)=\int_0^1 (1-s)^2\,\mathcal R(s)\,ds.
\]
For this $\mathcal R$, both moments vanish:
\[
\int_0^1 s(1-s)\,\mathcal R(s)\,ds=0,
\qquad
\int_0^1 (1-s)^2\,\mathcal R(s)\,ds=0.
\]
Hence
\[
D_0'(1)=0,
\qquad
D_0(1)=\mathcal Z(0),
\]
so $(0,1)$ is an exact rooted repeated-source critical pair.

Moreover,
\[
D_0''(1)=2\int_0^1 s^2\,\mathcal R(s)\,ds=\frac56>0,
\]
so the critical branch is nondegenerate and the implicit-function theorem gives
a real-analytic continuation
\[
x=\tau(a),
\qquad
\tau(0)=1,
\qquad
\tau'(0)=0.
\]
Also, with
\[
V(a):=D_a(\tau(a)),
\qquad
G(a):=V(a)-\mathcal Z(a),
\]
one has at the rooted point
\[
G(0)=0,
\qquad
G'(0)=-\mathcal R(0)=-1.
\]

Finally, the level equation
\[
\mathcal Z(t)=\mathcal Z(0)=0
\]
has two interior solutions,
\[
t\approx 0.368001,
\qquad
t\approx 0.851550,
\]
so the rooted point comes with exactly the two interior recrossings of the
moving level demanded by the source-pinned shear picture.

Thus a theorem using only the scalar level-set topology of $\mathcal Z$
together with the first-root identities above cannot suffice to exclude the
remaining adjacent-odd-zero configuration. Any successful next step must use
additional structure: either the provenance of $(a,\tau(a))$ as the
\emph{continued born branch}, or a genuinely source-coupled constraint such as
the moving shear
\[
f_a=u_5-\mathcal Z(a)\lambda,
\]
rather than only the abstract level geometry of $\mathcal Z$.
\end{remark}

\begin{remark}[No-go for branch provenance alone]
\label{rem:no-go-branch-provenance-alone}
One might try to strengthen
Remark~\ref{rem:no-go-pure-level-set-barrier-for-Z}
by using not just the scalar level geometry of $\mathcal Z$, but also the
\emph{provenance} of the rooted point: namely, that $(a,\tau(a))$ lies on the
continued branch born at an earlier odd zero of $\mathcal R$.
At the planar normalized $D_a$-level, this is still not enough. There are
explicit analytic models in which the born branch from one odd zero continues to
a genuine rooted point before leaving the interval bounded by the adjacent odd
zeros.

A concrete example is obtained by taking
\[
S(t)=t^2,
\qquad
\mathcal R(t):=t(t-1)(t^2-4t-2)
=
t^4-5t^3+2t^2+2t,
\]
and
\[
\mathcal Z(t):=\int_0^t \mathcal R(u)\,du
=
\frac15 t^5-\frac54 t^4+\frac23 t^3+t^2.
\]
Then the odd zeros of $\mathcal R$ are
\[
2-\sqrt6,\qquad 0,\qquad 1,\qquad 2+\sqrt6,
\]
so $0$ is an odd zero and its adjacent odd zero on the left is $2-\sqrt6$.

Now choose an analytic function $\widetilde Y$ such that
\[
\frac12\,\widetilde Y''=\mathcal Z,
\]
and let $D_a(x)$ denote the normalized determinant/divided-difference function
for the planar graph $(x,\widetilde Y)$. Near the odd zero $0$ there is the
usual born critical branch
\[
x=\tau(a),
\qquad
a<0<x,
\qquad
D_a'(\tau(a))=0,
\]
with
\[
\tau(a)=-a+O(a^2),
\]
and the local odd-zero analysis gives
\[
G(a):=D_a(\tau(a))-\mathcal Z(a)<0
\]
for $a<0$ sufficiently close to $0$.

However, along this same connected born branch one finds
\[
G(-0.5801)<0,
\qquad
G(-0.5901)>0,
\]
so by continuity there exists
\[
a_*\in(-0.5901,-0.5801)
\]
such that
\[
G(a_*)=0.
\]
Thus the \emph{continued born branch} reaches a rooted point
\[
D_{a_*}'(\tau(a_*))=0,
\qquad
D_{a_*}(\tau(a_*))=\mathcal Z(a_*).
\]
Numerically,
\[
a_*\approx -0.5805673462,
\qquad
\tau(a_*)\approx 0.3888336023,
\]
and
\[
D_{a_*}''(\tau(a_*))>0,
\]
so this rooted point is nondegenerate; in particular,
\[
\tau'(a_*)=0
\]
as predicted by the transport system.

Moreover, the segment
\[
[a_*,\tau(a_*)]
\]
contains exactly the born odd zero $0$ together with the adjacent odd zero
$2-\sqrt6$, and no others. So even the strengthened statement
\[
\text{``continued born-branch provenance plus adjacent-odd-zero geometry
already excludes rooting''}
\]
is false in the planar normalized model.

Therefore the remaining exclusion cannot depend only on abstract branch
continuation data for
\[
D_a'(\tau(a))=0
\]
together with the scalar geometry of $\mathcal Z$. Any successful next step
must use genuinely source-coupled structure, such as the actual quotient
\[
\Psi_a(h)=\frac{D_{u_5}(a,h)}{D_\lambda(a,h)},
\]
the moving shear
\[
f_a=u_5-\mathcal Z(a)\lambda,
\]
or some equivalent RH-specific relation invisible in the planar $D_a$ model.
\end{remark}

\begin{remark}[No-go for disconjugacy/ECT alone]
\label{rem:no-go-disconjugacy-alone}
A further natural idea is to recast the moving-shear family through the
third-order operator
\[
L[y]:=\frac{W(c,v_5,\lambda,y)}{J_\lambda},
\]
and try to exclude rooted moving-shear events by proving that $L$ is
disconjugate, Chebyshev, or sign-regular on opened nonmixed arcs.
At the present level of abstraction this is still not enough.

Indeed, already in the normalized source model
\[
c(t)=1,\qquad v_5(t)=t,\qquad \lambda(t)=t^2,
\]
one has
\[
J_f=W(c,v_5,f)=f'',
\qquad
J_\lambda=2,
\qquad
L[y]=\frac{W(c,v_5,\lambda,y)}{J_\lambda}=\frac12\,y'''.
\]
So the homogeneous equation $L[y]=0$ has basis
\[
1,\ t,\ t^2,
\]
hence is maximally disconjugate / ECT on every interval. Nevertheless, the
planar normalized counterexample from
Remark~\ref{rem:no-go-branch-provenance-alone}
fits exactly into this framework: for a suitable analytic $u_5$ and the
source-pinned family
\[
f_a:=u_5-\mathcal Z(a)\lambda,
\]
one has
\[
L[f_a]=L[u_5],
\qquad
J_{f_a}(a)=0,
\]
and along the continued born branch there is still a rooted point
\[
D_{f_a}(a,h)=0,
\qquad
\partial_h D_{f_a}(a,h)=0.
\]

Thus disconjugacy, ECT, or sign-regularity of the operator $L$ by itself
cannot exclude the remaining rooted moving-shear configuration. Any successful
next step must use additional source-coupled structure beyond the abstract
third-order ODE viewpoint, such as special positivity or monotonicity for the
distinguished quotient
\[
\Psi_a(h)=\frac{D_{u_5}(a,h)}{D_\lambda(a,h)}
\]
or for the particular slice
\[
f_a=u_5-\mathcal Z(a)\lambda.
\]
\end{remark}

\begin{remark}[Ceiling of the present scalar/discriminant analysis]
\label{rem:ceiling-present-scalar-discriminant-analysis}
The preceding results sharply delimit what can and cannot be obtained from the
current planar/scalar reformulation.

On the positive-kernel side, one-sign-change chambers are sterile for genuine
repeated-source first loss:
the born critical branch is simple, has strictly negative slope, and cannot
satisfy
\[
D_a(\tau(a))=\mathcal Z(a)
\]
before the active segment acquires another odd zero of $\mathcal R$
(Proposition~\ref{prop:first-rooting-requires-second-odd-zero-repeated-source}).
On the discriminant side, the adjacent-odd-zero entrance itself cannot create
the first rooted point
(Proposition~\ref{prop:discriminant-entry-not-first-root}).

The new anchor-sign and odd-anchor-passage arguments remove a substantial part
of the remaining freedom. On any continuation interval on which the anchor-side
value
\[
\mathcal R(a)
\]
keeps fixed sign, the RH-selected slice
\[
\gamma=\mathcal Z(a)
\]
cannot cross the continued discriminant curve. When the anchor itself passes
through an odd-order zero of $\mathcal R$, the integrating-factor defect keeps
its incoming sign across that odd anchor zero, so a rooted crossing cannot occur
there or immediately after it. The ensuing post-anchor chambers then alternate
dangerous and safe, beginning with a dangerous first post-anchor chamber.

However, two no-go facts now show that the present scalar/discriminant package
has reached its ceiling. First,
Remark~\ref{rem:no-go-local-scalar-first-dangerous-chamber} shows that already
on the first dangerous post-anchor chamber, the local normalized rooted
equations and fold/transversality data are formally realizable in the planar
scalar model. Second,
Proposition~\ref{prop:universality-threshold-current-route-two-package} and
Corollary~\ref{cor:current-route-two-package-universal-under-normalization}
show that the entire current route-2 source package
\[
\Psi_a,\qquad
\mathfrak D,\qquad
\mathcal W_{\mathrm{src}},\qquad
D_{f_a}
\]
is universal under normalization: every analytic normalized planar model lifts
to a genuine analytic source model for this package.

Accordingly, after the fixed-sign and odd-anchor exclusions, the remaining
problem is \emph{not} a delayed no-crossing / profile-realizability problem
inside the present scalar/discriminant package. It is to find the first
genuinely non-universal source-coupled input, i.e. structure not already
captured by the normalized planar model $(\mathcal Z,\mathcal R,D_a)$.

So any further progress on the nonmixed repeated-source branch must use
genuinely source-coupled structure beyond the present universal package,
possibly through stronger provenance laws, higher canonical/mixed data, or
another RH-specific input invisible in the planar normalized model.
\end{remark}

\begin{proposition}[Discriminant reduction for the moving-shear family]
\label{prop:discriminant-reduction-moving-shear}
Let $J$ be a good nonmixed $v_5$-patch. For $\gamma\in\mathbf R$, set
\[
f_\gamma:=u_5-\gamma\lambda.
\]
Define the sheared discriminant set
\[
\mathcal D:=
\left\{
(a,\gamma,h)\in J\times\mathbf R\times J:
a<h,\
D_{f_\gamma}(a,h)=0,\
\partial_h D_{f_\gamma}(a,h)=0
\right\}.
\]

Then on the region $a<h$ one has the exact equivalence
\[
(a,\gamma,h)\in\mathcal D
\iff
\gamma=\Psi_a(h),
\qquad
\mathfrak D(a,h)=0.
\]
Equivalently,
\[
\mathcal D
=
\left\{
(a,\Psi_a(h),h):\ a<h,\ \mathfrak D(a,h)=0
\right\}.
\]

In particular, on a nonmixed patch the sheared discriminant is not an
independent two-dimensional catastrophe surface in $(a,\gamma,h)$-space:
it is the graph of the repeated-source critical locus
\[
\mathfrak D(a,h)=0
\]
in the extra $\gamma$-direction.
\end{proposition}

\begin{proof}
By Proposition~\ref{prop:source-determinant-form-odd-zero-unrooted-defect}(2),
for every $(a,h)$ with $a<h$ one has
\[
D_{f_\gamma}(a,h)
=
D_\lambda(a,h)\bigl(\Psi_a(h)-\gamma\bigr).
\]
Differentiating in $h$ gives
\[
\partial_h D_{f_\gamma}(a,h)
=
\partial_h D_\lambda(a,h)\bigl(\Psi_a(h)-\gamma\bigr)
+
D_\lambda(a,h)\,\partial_h\Psi_a(h).
\]

On a good nonmixed patch,
\[
D_\lambda(a,h)\neq 0
\qquad (a<h).
\]
Therefore
\[
D_{f_\gamma}(a,h)=0
\iff
\gamma=\Psi_a(h).
\]
At such a point the first term in the formula for $\partial_h D_{f_\gamma}$
vanishes, so
\[
\partial_h D_{f_\gamma}(a,h)=0
\iff
D_\lambda(a,h)\,\partial_h\Psi_a(h)=0
\iff
\partial_h\Psi_a(h)=0.
\]
Finally, by
Remark~\ref{rem:source-determinant-reformulation-Gammaa},
\[
\partial_h\Psi_a(h)
=
\frac{W_h(c,v_5)(a)\,\mathfrak D(a,h)}{D_\lambda(a,h)^2},
\]
and the prefactor is nonvanishing on the nonmixed patch. Hence
\[
\partial_h\Psi_a(h)=0
\iff
\mathfrak D(a,h)=0.
\]
Combining the two equivalences proves the result.
\end{proof}

\begin{proposition}[Generic rooted discriminant point is a transverse crossing]
\label{prop:generic-rooted-discriminant-point-transverse-crossing}
In the setting of
Proposition~\ref{prop:discriminant-reduction-moving-shear},
define the RH-selected slice
\[
\mathcal S_{\mathrm{RH}}
:=
\left\{
(a,\gamma,h)\in J\times\mathbf R\times J:
\gamma=\mathcal Z(x(a))
\right\}.
\]

Let
\[
(a_*,\gamma_*,h_*)\in \mathcal D\cap \mathcal S_{\mathrm{RH}}
\]
and assume
\[
\partial_h^2 D_{f_{\gamma_*}}(a_*,h_*)\neq 0,
\qquad
\mathcal R(x(a_*))\neq 0.
\]
Then there is a neighborhood of $(a_*,\gamma_*,h_*)$ in which
$\mathcal D$ is a unique real-analytic curve
\[
a\longmapsto \bigl(a,\gamma(a),h(a)\bigr),
\]
with
\[
\gamma(a)=\Psi_a(h(a)).
\]

Moreover one has
\[
\gamma'(a_*)=0,
\qquad
h'(a_*)=0,
\]
and
\[
\frac{d}{da}\Bigl(\gamma(a)-\mathcal Z(x(a))\Bigr)\Big|_{a=a_*}
=
-\mathcal R(x(a_*))\,x'(a_*)\neq 0.
\]
Hence the discriminant curve $\mathcal D$ meets the RH-selected slice
$\mathcal S_{\mathrm{RH}}$ transversely at $(a_*,\gamma_*,h_*)$.

In fact the $\gamma$-projection of the discriminant curve has a nondegenerate
local extremum there:
\[
\gamma''(a_*)
=
-\frac{2\,x'(a_*)^2}{x(h_*)-x(a_*)}\,\mathcal R(x(a_*)).
\]
\end{proposition}

\begin{proof}
By Proposition~\ref{prop:discriminant-reduction-moving-shear},
the discriminant is cut out by
\[
\partial_h\Psi_a(h)=0,
\qquad
\gamma=\Psi_a(h).
\]
At the given point,
\[
\partial_h^2 D_{f_{\gamma_*}}(a_*,h_*)\neq 0.
\]
Since
\[
D_{f_\gamma}(a,h)=D_\lambda(a,h)\bigl(\Psi_a(h)-\gamma\bigr),
\]
differentiating twice in $h$ and evaluating at a discriminant point gives
\[
\partial_h^2 D_{f_{\gamma_*}}(a_*,h_*)
=
D_\lambda(a_*,h_*)\,\partial_h^2\Psi_{a_*}(h_*).
\]
Because $D_\lambda(a_*,h_*)\neq 0$, this implies
\[
\partial_h^2\Psi_{a_*}(h_*)\neq 0.
\]
So by the implicit function theorem there is a unique real-analytic function
\[
h=h(a)
\]
with
\[
h(a_*)=h_*,
\qquad
\partial_h\Psi_a(h(a))=0.
\]
Setting
\[
\gamma(a):=\Psi_a(h(a))
\]
gives the asserted local parametrization of $\mathcal D$.

Now pass to the local normalized nonmixed $x$-coordinate and write
\[
\alpha:=x(a),
\qquad
\tau(\alpha):=x(h(a)),
\qquad
V(\alpha):=D_\alpha(\tau(\alpha)).
\]
By Remark~\ref{rem:source-determinant-reformulation-Gammaa},
\[
\gamma(a)=\Psi_a(h(a))=V(\alpha).
\]
Since $(a_*,\gamma_*,h_*)\in\mathcal S_{\mathrm{RH}}$, one has
\[
V(\alpha_*)=\mathcal Z(\alpha_*),
\qquad
\alpha_*:=x(a_*).
\]

By Proposition~\ref{prop:transport-equations-repeated-source-critical-branch},
at a rooted noncollapse point one has
\[
V'(\alpha_*)=0,
\qquad
\tau'(\alpha_*)=0.
\]
Therefore
\[
\gamma'(a_*)
=
V'(\alpha_*)\,x'(a_*)
=
0,
\]
and
\[
x'(h_*)\,h'(a_*)
=
\tau'(\alpha_*)\,x'(a_*)
=
0,
\]
hence
\[
h'(a_*)=0.
\]

For the transverse-crossing derivative, differentiate
\[
\gamma(a)-\mathcal Z(x(a))
=
V(\alpha)-\mathcal Z(\alpha)
\]
to get
\[
\frac{d}{da}\Bigl(\gamma(a)-\mathcal Z(x(a))\Bigr)
=
x'(a)\bigl(V'(\alpha)-\mathcal R(\alpha)\bigr).
\]
At $a=a_*$, since $V'(\alpha_*)=0$, this becomes
\[
\frac{d}{da}\Bigl(\gamma(a)-\mathcal Z(x(a))\Bigr)\Big|_{a=a_*}
=
-\mathcal R(\alpha_*)\,x'(a_*),
\]
which is nonzero by hypothesis. This proves transversality.

Finally, differentiate
\[
V'(\alpha)=\frac{2}{\tau(\alpha)-\alpha}\bigl(V(\alpha)-\mathcal Z(\alpha)\bigr)
\]
from
Proposition~\ref{prop:transport-equations-repeated-source-critical-branch}(1).
At the rooted point,
\[
V(\alpha_*)-\mathcal Z(\alpha_*)=0,
\qquad
\tau'(\alpha_*)=0,
\]
so
\[
V''(\alpha_*)
=
\frac{2}{\tau(\alpha_*)-\alpha_*}\bigl(V'(\alpha_*)-\mathcal R(\alpha_*)\bigr)
=
-\frac{2}{\tau(\alpha_*)-\alpha_*}\,\mathcal R(\alpha_*).
\]
Hence
\[
\gamma''(a_*)
=
V''(\alpha_*)\,x'(a_*)^2
=
-\frac{2\,x'(a_*)^2}{x(h_*)-x(a_*)}\,\mathcal R(x(a_*)),
\]
as claimed.
\end{proof}

\begin{remark}[What the discriminant/catastrophe route actually reduces to]
\label{rem:what-discriminant-catastrophe-route-actually-reduces-to}
The discriminant reformulation sharpens the geometry, but it also shows the
limit of purely local catastrophe arguments on this route.

Indeed, by
Proposition~\ref{prop:discriminant-reduction-moving-shear},
the sheared discriminant on a nonmixed patch is already one-dimensional:
it is the graph of the repeated-source critical locus
\[
\mathfrak D(a,h)=0
\]
in the $\gamma$-direction. And by
Proposition~\ref{prop:generic-rooted-discriminant-point-transverse-crossing},
a generic rooted point is not a tangential meeting of sheets, nor a local
fold/cusp obstruction for the RH-selected slice
\[
\gamma=\mathcal Z(x(a)).
\]
It is a \emph{transverse crossing} of a smooth discriminant curve by that
moving slice, with oriented crossing sign controlled by
\[
-\mathcal R(x(a_*)).
\]

Accordingly, the discriminant/catastrophe route does not reduce the adjacent
odd-zero problem to a local singularity classification in
$(a,\gamma,h)$-space. The remaining issue is global: to exclude the existence
of such a transverse crossing along the continued born branch before the
segment exits the arc, collapses affine-planarly, or acquires the next odd
zero of $\mathcal W_{\mathrm{src}}$.
\end{remark}

\begin{proposition}[Anchor-sign preservation excludes first rooting]
\label{prop:anchor-sign-preservation-excludes-first-rooting}
Let $J$ be an opened nonmixed arc, and work in a local normalized nonmixed
chart on which
\[
S(t)=t^2,
\qquad
\Gamma_a=D_a.
\]
Let
\[
I=(a_1,a_0]\subset J
\]
and let
\[
\tau:I\to J
\]
be a real-analytic continuation of the repeated-source critical branch such that
\[
\tau(a)>a,
\qquad
D_a'(\tau(a))=0
\qquad (a\in I).
\]
Set
\[
V(a):=D_a(\tau(a)),
\qquad
G(a):=V(a)-\mathcal Z(a).
\]

Assume that there exists $\epsilon\in\{\pm1\}$ such that
\[
\epsilon\,\mathcal R(a)>0
\qquad\text{for every }a\in I,
\]
and
\[
\epsilon\,G(a_0)>0.
\]
Then
\[
\epsilon\,G(a)>0
\qquad\text{for every }a\in I.
\]
In particular,
\[
G(a)\neq 0
\qquad (a\in I),
\]
so the branch cannot realize the rigidified first-loss condition anywhere on
$I$.
\end{proposition}

\begin{proof}
Assume for contradiction that $G$ vanishes somewhere on $I$. Since $G$ is
continuous, the set
\[
Z:=\{a\in I:G(a)=0\}
\]
is closed in $I$ and nonempty. Let
\[
a_*:=\sup Z.
\]
Then $a_*<a_0$, and by the choice of $a_*$ one has
\[
\epsilon\,G(a)>0
\qquad\text{for all }a\in(a_*,a_0].
\]
Hence the one-sided derivative at $a_*$ satisfies
\[
\epsilon\,G'(a_*)\ge 0.
\]

On the other hand, by
Proposition~\ref{prop:rooted-defect-transport-stationary-first-root},
\[
G'(a)=\frac{2}{\tau(a)-a}\,G(a)-\mathcal R(a).
\]
Evaluating at the root $a=a_*$ gives
\[
G'(a_*)=-\mathcal R(a_*).
\]
Since $\epsilon\,\mathcal R(a_*)>0$ by hypothesis, this yields
\[
\epsilon\,G'(a_*)
=
-\epsilon\,\mathcal R(a_*)
<0,
\]
contradicting
\[
\epsilon\,G'(a_*)\ge 0.
\]
Therefore $G$ has no zero on $I$. Since
\[
\epsilon\,G(a_0)>0,
\]
the sign of $G$ on all of $I$ is $\epsilon$.
\end{proof}

\begin{corollary}[No crossing on the adjacent-two-odd-zero chamber under anchor-sign preservation]
\label{cor:no-crossing-adjacent-two-odd-zero-under-anchor-sign-preservation}
Let $b<q$ be adjacent odd-order zeros of $\mathcal R$ on an opened nonmixed
arc $J$. Work in a local normalized nonmixed chart, and let
\[
I=(a_1,a_0]\subset J,
\qquad
\tau:I\to J
\]
be a real-analytic continuation of the born repeated-source critical branch such
that
\[
\tau(a)>a,
\qquad
D_a'(\tau(a))=0
\qquad (a\in I),
\]
with
\[
\tau(a_0)=q.
\]
Assume that for every $a\in I$ with $a<a_0$, the segment
\[
[a,\tau(a)]
\]
contains exactly two sign changes of $\mathcal R$, occurring at the odd zeros
$b$ and $q$, and that
\[
\mathcal R(a)
\]
has fixed sign on $I$.

Then
\[
\operatorname{sgn}\bigl(G(a)\bigr)=\operatorname{sgn}\bigl(\mathcal R(a)\bigr)
\qquad (a\in I),
\]
hence
\[
G(a)\neq 0
\qquad (a\in I).
\]
Equivalently, the RH-selected slice
\[
\gamma=\mathcal Z(a)
\]
does not cross the continued discriminant curve
\[
(a,\gamma,h)=\bigl(a,V(a),\tau(a)\bigr)
\]
anywhere on $I$.
\end{corollary}

\begin{proof}
At the entry point $a=a_0$, one has
\[
\tau(a_0)=q,
\]
and the segment
\[
[a_0,q]
\]
contains exactly one sign change of $\mathcal R$, namely at $b$. Therefore
Proposition~\ref{prop:single-sign-change-cannot-realize-first-loss} applied to
the critical pair $(a_0,q)$ gives
\[
\operatorname{sgn}\bigl(G(a_0)\bigr)
=
\operatorname{sgn}\bigl(\mathcal R(a_0)\bigr).
\]
Now apply
Proposition~\ref{prop:anchor-sign-preservation-excludes-first-rooting}
with
\[
\epsilon:=\operatorname{sgn}\bigl(\mathcal R|_I\bigr).
\]
\end{proof}

\begin{remark}[What this really removes from the previous caveat]
\label{rem:what-anchor-sign-preservation-removes}
The preceding corollary shows that the geometric hypothesis
``the anchor stays on the original birth-side outer lobe'' was stronger than
necessary. The proof uses only the intrinsic condition that the anchor-side
value
\[
\mathcal R(a)
\]
keep a fixed sign along the continuation interval.

Accordingly, the remaining unresolved possibility on this route is narrower, but
it is not yet exhausted. A rooted point can occur only after the anchor itself
passes through an odd-order zero of $\mathcal R$, or after branch exit, or after
affine-planar collapse. Moreover,
Proposition~\ref{prop:odd-anchor-passage-no-reset} shows that no rooted crossing
can occur at that anchor odd zero itself or immediately after passage, and
Corollary~\ref{cor:first-post-anchor-chamber-dangerous} shows that the
subsequent post-anchor chambers alternate dangerous and safe. Thus the true residual obstruction is no longer the raw adjacent-two-odd-zero
chamber itself. But by the universality threshold for the current route-2
package, by the local freedom of arbitrary canonical quotient-septic geometry,
and by the full-block universality of the direct provenance slice
\[
\phi=A\,q^{(7)},
\]
even the delayed post-anchor scenario cannot be excluded inside those packages
alone.
\end{remark}

\begin{proposition}[Odd-anchor passage has no sign reset]
\label{prop:odd-anchor-passage-no-reset}
Let $J$ be an opened nonmixed arc, and work in a local normalized nonmixed
chart on which
\[
S(t)=t^2,
\qquad
\Gamma_a=D_a.
\]
Let
\[
q\in J
\]
be an isolated odd-order zero of $\mathcal R$, and let
\[
\tau:(q-\delta,q+\delta)\to J
\]
be a real-analytic continuation of the repeated-source critical branch such that
\[
\tau(a)>a,
\qquad
D_a'(\tau(a))=0
\qquad (a\in(q-\delta,q+\delta)).
\]
Set
\[
G(a):=D_a(\tau(a))-\mathcal Z(a),
\qquad
L(a):=\tau(a)-a,
\]
and define
\[
\mathcal A'(a):=\frac{2}{L(a)},
\qquad
\mathscr G(a):=e^{-\mathcal A(a)}G(a).
\]
Assume that there exists $\epsilon\in\{\pm1\}$ such that
\[
\epsilon\,\mathcal R(a)>0,
\qquad
\epsilon\,G(a)>0
\qquad\text{for every }a\in(q,q+\delta).
\]
Then
\[
\epsilon\,G(q)>0.
\]
In particular, $G(q)\neq 0$. Hence, after possibly shrinking $\delta$,
\[
\epsilon\,G(a)>0
\qquad\text{for every }a\in(q-\delta,q+\delta).
\]
Since $q$ is an odd-order zero of $\mathcal R$, the sign of $\mathcal R$ flips
across $q$; therefore on the outgoing side
\[
\operatorname{sgn}G(a)=-\operatorname{sgn}\mathcal R(a)
\qquad (a\in(q-\delta,q)).
\]
So a rooted crossing cannot occur at the anchor odd zero itself, nor
immediately after passage across it.
\end{proposition}

\begin{proof}
By
Proposition~\ref{prop:rooted-defect-transport-stationary-first-root},
\[
G'(a)=\frac{2}{L(a)}\,G(a)-\mathcal R(a),
\]
hence
\[
\mathscr G'(a)=-e^{-\mathcal A(a)}\mathcal R(a).
\]
Assume for contradiction that $G(q)=0$. Then $\mathscr G(q)=0$, and for
$a\in(q,q+\delta)$,
\[
\mathscr G(a)
=
\mathscr G(q)+\int_q^a \mathscr G'(s)\,ds
=
-\int_q^a e^{-\mathcal A(s)}\mathcal R(s)\,ds.
\]
Since $e^{-\mathcal A(s)}>0$ and $\epsilon\,\mathcal R(s)>0$ on $(q,q+\delta)$,
this gives
\[
\epsilon\,\mathscr G(a)<0
\qquad (a\in(q,q+\delta)),
\]
hence
\[
\epsilon\,G(a)<0
\qquad (a\in(q,q+\delta)),
\]
contradicting the hypothesis. Therefore $G(q)\neq 0$. Since
\[
\epsilon\,G(a)>0
\qquad (a\in(q,q+\delta)),
\]
continuity yields
\[
\epsilon\,G(q)\ge 0,
\]
and the nonvanishing just proved forces
\[
\epsilon\,G(q)>0.
\]
A second continuity argument then gives
\[
\epsilon\,G(a)>0
\qquad (a\in(q-\delta,q+\delta))
\]
after shrinking $\delta$ if necessary. Finally, because $q$ is an odd-order zero
of $\mathcal R$, its sign flips across $q$, so on $(q-\delta,q)$ one has
\[
\operatorname{sgn}G(a)=-\operatorname{sgn}\mathcal R(a).
\]
In particular $G$ cannot vanish there.
\end{proof}

\begin{proposition}[Weighted-defect recursion on post-anchor chambers]
\label{prop:weighted-defect-recursion-post-anchor-chambers}
Assume the setting of
Proposition~\ref{prop:odd-anchor-passage-no-reset}. Let
\[
q_1>q_2>\cdots>q_N
\]
be successive odd-order zeros of $\mathcal R$ crossed by the anchor along the
continued branch after the adjacent-two-odd-zero chamber, and set
\[
I_j:=(q_{j+1},q_j)
\qquad (1\le j\le N-1).
\]
Assume that the continued branch exists real-analytically on
\[
\bigcup_{j=1}^{N-1} I_j.
\]
For each $j$, let
\[
\epsilon_j:=\operatorname{sgn}(\mathcal R|_{I_j})\in\{\pm1\},
\qquad
m_j(a):=\int_a^{q_j} e^{-\mathcal A(s)}|\mathcal R(s)|\,ds
\qquad (a\in I_j),
\]
and
\[
M_j:=m_j(q_{j+1})>0.
\]
Then on each $I_j$,
\[
\mathscr G(a)=\mathscr G(q_j)+\epsilon_j\,m_j(a).
\]
Consequently:
\begin{enumerate}
\item if
\[
\operatorname{sgn}\mathscr G(q_j)=\epsilon_j,
\]
then $\mathscr G$ has no zero on $I_j$;

\item if
\[
\operatorname{sgn}\mathscr G(q_j)=-\epsilon_j,
\]
then $\mathscr G$ has a zero on $I_j$ if and only if
\[
M_j>|\mathscr G(q_j)|,
\]
and in that case the zero is unique;

\item if $\mathscr G$ has no zero on $I_j$, then
\[
\operatorname{sgn}\mathscr G(q_{j+1})
=
\operatorname{sgn}\mathscr G(q_j).
\]
\end{enumerate}
Since $\epsilon_{j+1}=-\epsilon_j$, the chamber type
\[
\operatorname{sgn}\mathscr G(q_j)\ \text{ versus }\ \epsilon_j
\]
alternates from one post-anchor chamber to the next as long as no rooted
crossing has yet occurred.
\end{proposition}

\begin{proof}
On $I_j$ one has
\[
\mathscr G'(a)=-e^{-\mathcal A(a)}\mathcal R(a)
=-\epsilon_j e^{-\mathcal A(a)}|\mathcal R(a)|.
\]
Hence for $a\in I_j$,
\[
\mathscr G(a)
=
\mathscr G(q_j)+\int_{q_j}^a \mathscr G'(s)\,ds
=
\mathscr G(q_j)+\epsilon_j\int_a^{q_j}e^{-\mathcal A(s)}|\mathcal R(s)|\,ds,
\]
which is the stated formula.

Now $m_j(a)$ is strictly increasing as $a$ decreases from $q_j$ to $q_{j+1}$,
with
\[
m_j(q_j)=0,
\qquad
m_j(q_{j+1})=M_j.
\]
So if
\[
\operatorname{sgn}\mathscr G(q_j)=\epsilon_j,
\]
then
\[
\mathscr G(a)=\epsilon_j\bigl(|\mathscr G(q_j)|+m_j(a)\bigr)
\]
has fixed sign on $I_j$ and no zero. If instead
\[
\operatorname{sgn}\mathscr G(q_j)=-\epsilon_j,
\]
then
\[
\mathscr G(a)=\epsilon_j\bigl(-|\mathscr G(q_j)|+m_j(a)\bigr),
\]
so a zero occurs precisely when $m_j(a)=|\mathscr G(q_j)|$, which happens if
and only if
\[
M_j>|\mathscr G(q_j)|.
\]
Because $m_j$ is strictly monotone, that zero is unique.

Finally, if no zero occurs on $I_j$, then $\mathscr G$ keeps one sign
throughout the interval, so in particular
\[
\operatorname{sgn}\mathscr G(q_{j+1})
=
\operatorname{sgn}\mathscr G(q_j).
\]
Since the $\epsilon_j$ alternate at successive odd zeros, the relative chamber
type alternates as claimed.
\end{proof}

\begin{corollary}[The first post-anchor chamber is dangerous]
\label{cor:first-post-anchor-chamber-dangerous}
Let $q_1$ be the first odd-order zero of $\mathcal R$ crossed by the anchor
after the adjacent-two-odd-zero chamber, and let
\[
I_1:=(q_2,q_1)
\]
be the first post-anchor chamber. Assume the continued born branch extends
real-analytically across $q_1$ and on $I_1$. Then:
\begin{enumerate}
\item no rooted crossing occurs at $q_1$ or immediately after crossing it;

\item on $I_1$ one has
\[
\operatorname{sgn}\mathscr G(q_1)=-\epsilon_1,
\]
so $I_1$ is the first \emph{dangerous} post-anchor chamber;

\item as long as no rooted crossing has yet occurred, the post-anchor chambers
alternate dangerous/safe/dangerous/safe\ldots\ .
\end{enumerate}
In particular, the first chamber in which a rooted crossing is even eligible
after the anchor-sign-preservation interval is the first post-anchor chamber.
Without an additional diagonal-capture theorem, however, this eligibility does
not by itself determine the odd-zero count on the active segment.
\end{corollary}

\begin{proof}
The adjacent-two-odd-zero chamber is root-free by
Corollary~\ref{cor:no-crossing-adjacent-two-odd-zero-under-anchor-sign-preservation}.
At its left endpoint $q_1$, the anchor crosses an odd-order zero of
$\mathcal R$. Proposition~\ref{prop:odd-anchor-passage-no-reset} shows that
$\mathscr G$ keeps its incoming sign across $q_1$, while $\mathcal R$ changes
sign there. Thus on the outgoing chamber $I_1$ the signs are opposite:
\[
\operatorname{sgn}\mathscr G(q_1)=-\epsilon_1.
\]
The no-root statement at and immediately after $q_1$ is also exactly
Proposition~\ref{prop:odd-anchor-passage-no-reset}. The alternation of chamber
types then follows from
Proposition~\ref{prop:weighted-defect-recursion-post-anchor-chambers}(3).
\end{proof}

\begin{remark}[No-go for purely local scalar exclusion on the first dangerous post-anchor chamber]
\label{rem:no-go-local-scalar-first-dangerous-chamber}
At a rooted critical point, after normalizing the active interval to
\[
x=a+sL,
\qquad
0\le s\le 1,
\qquad
L=\tau(a)-a,
\]
the profile
\[
R_a(s):=\mathcal R(a+sL)
\]
satisfies the two local scalar moment identities
\[
\int_0^1 (1-s)^2\,R_a(s)\,ds=0,
\qquad
\int_0^1 s(1-s)\,R_a(s)\,ds=0.
\]

Once nonlocal-target landing is admitted, the first dangerous post-anchor
chamber need not correspond to a fixed odd-zero count or to a rigid four-lobe
pattern on the active segment. Even in the simplest three-lobe sign pattern
\[
+,-,+
\]
(or its global negative), the local scalar rooted equations are already
formally realizable.

Indeed, for any $A>0$, define
\[
R_A(s):=
\begin{cases}
A, & 0<s<\frac13,\\[4pt]
-\frac72 A, & \frac13<s<\frac23,\\[4pt]
\frac{11}{2}A, & \frac23<s<1.
\end{cases}
\]
A direct calculation gives
\[
\int_0^1 (1-s)^2\,R_A(s)\,ds=0,
\qquad
\int_0^1 s(1-s)\,R_A(s)\,ds=0,
\]
and also
\[
\int_0^1 R_A(s)\,ds
=
\int_0^1 s\,R_A(s)\,ds
=
\int_0^1 s^2\,R_A(s)\,ds
=
A.
\]
So the local rooted-point equations and the nondegenerate fold data are already
formally compatible with a three-lobe profile.

Thus no argument using only the local normalized rooted profile, the two
vanishing scalar moments, and the local transversality/fold signs can exclude
rooting on the first dangerous post-anchor chamber. Any successful next step
must use genuinely source-coupled structure invisible in the planar normalized
$D_a$ model.
\end{remark}

\begin{remark}[What route 2 now reduces to]
\label{rem:route-two-reduces-to-global-no-crossing}
The discriminant analysis first showed that a rooted point is not born
\emph{at the entry} into the adjacent-two-odd-zero chamber. The stronger
anchor-sign argument showed more: on any continuation interval on which the
anchor-side value
\[
\mathcal R(a)
\]
keeps a fixed sign, the RH-selected slice
\[
\gamma=\mathcal Z(a)
\]
cannot cross the discriminant curve at all.

The odd-anchor analysis sharpened this further: when the anchor itself passes
through an odd-order zero of $\mathcal R$, the integrating-factor defect
\[
\mathscr G(a)=e^{-\mathcal A(a)}\bigl(D_a(\tau(a))-\mathcal Z(a)\bigr)
\]
keeps its incoming sign across that odd anchor zero, so a rooted crossing cannot
occur there or immediately after passage. The subsequent post-anchor chambers
then alternate dangerous and safe, beginning with a dangerous first post-anchor
chamber
(Corollary~\ref{cor:first-post-anchor-chamber-dangerous}).

However, two no-go facts are now clear. First,
Remark~\ref{rem:no-go-local-scalar-first-dangerous-chamber} shows that already
on that first dangerous post-anchor chamber, the local normalized rooted
equations and fold data are formally realizable in the planar scalar model.
Second, the current route-2 source package itself is still universal under the
normalized lift
\[
c=1,\qquad v_5=t,\qquad \lambda=t^2,\qquad u_5''=2\mathcal Z
\]
(Proposition~\ref{prop:universality-threshold-current-route-two-package}),
so even the source-level objects
\[
\Psi_a,\qquad \mathfrak D,\qquad \mathcal W_{\mathrm{src}},\qquad D_{f_a}
\]
do not yet distinguish the RH branch from an arbitrary analytic normalized
model.

So route 2 no longer reduces merely to a delayed no-crossing /
profile-realizability problem \emph{within the current package}. It reduces to
finding the first genuinely \emph{non-universal} source-coupled input, i.e. an
ingredient not already captured by the normalized planar $D_a$ model.

Within the present draft, the first plausible places for such new input lie
beyond the current universal quartet, for example in continuum-strength
provenance constraints on $q$ over whole one-pair blocks, or in external
mixed/canonical structure not already absorbed by the normalized planar model.
\end{remark}

\begin{proposition}[Universality threshold for the current route-2 package]
\label{prop:universality-threshold-current-route-two-package}
Let $I\subset\mathbf R$ be an interval, let
\[
\mathcal Z:I\to\mathbf R
\]
be real-analytic, and set
\[
\mathcal R:=\mathcal Z'.
\]
Define analytic source functions on $I$ by
\[
c(t):=1,\qquad v_5(t):=t,\qquad \lambda(t):=t^2,
\]
and choose $u_5$ so that
\[
u_5''(t)=2\mathcal Z(t).
\]
Then:
\begin{enumerate}
\item
\[
x=\frac{v_5}{c}=t,
\qquad
S=\frac{\lambda}{c}=t^2,
\qquad
Y=\frac{u_5}{c}=u_5;
\]

\item for every analytic $f$,
\[
J_f=W(c,v_5,f)=W(1,t,f)=f'',
\]
hence
\[
J_\lambda=2,
\qquad
J_{u_5}=2\mathcal Z,
\qquad
\mathcal W_{\mathrm{src}}
=
J_\lambda J_{u_5}'-J_{u_5}J_\lambda'
=
4\mathcal R;
\]

\item the repeated-source ratio is exactly the normalized planar one:
\[
\Psi_a(h)=\Gamma_a(h)=D_a(h);
\]

\item the critical determinant equation is exactly the normalized planar
critical-point equation:
\[
\mathfrak D(a,h)=0
\iff
\partial_h \Psi_a(h)=0
\iff
\partial_h D_a(h)=0;
\]

\item for the source-pinned moving shear
\[
f_a:=u_5-\mathcal Z(a)\lambda,
\]
one has
\[
J_{f_a}=2(\mathcal Z-\mathcal Z(a)),
\]
and
\[
D_{f_a}(a,h)=D_\lambda(a,h)\bigl(\Psi_a(h)-\mathcal Z(a)\bigr).
\]
\end{enumerate}
Thus every analytic normalized planar model $(\mathcal Z,\mathcal R,D_a)$
lifts to a genuine analytic source model for the entire current route-2 package.
\end{proposition}

\begin{proof}
The identities
\[
x=\frac{v_5}{c}=t,\qquad
S=\frac{\lambda}{c}=t^2,\qquad
Y=\frac{u_5}{c}=u_5
\]
are immediate. Also,
\[
J_f=W(c,v_5,f)=W(1,t,f)=f'',
\]
so
\[
J_\lambda=2,\qquad J_{u_5}=2\mathcal Z.
\]
Hence
\[
\mathcal W_{\mathrm{src}}
=
J_\lambda J_{u_5}'-J_{u_5}J_\lambda'
=
2(2\mathcal Z')-2\mathcal Z\cdot 0
=
4\mathcal R.
\]

Because the source coordinate and the normalized planar coordinate coincide,
\[
x(t)=t,
\]
the source repeated-source quotient is exactly the planar normalized quadratic
envelope:
\[
\Psi_a(h)=\Gamma_a(h)=D_a(h).
\]
Therefore the source critical determinant equation is exactly the planar
critical-point equation
\[
\mathfrak D(a,h)=0
\iff
\partial_h D_a(h)=0.
\]

Finally,
\[
f_a=u_5-\mathcal Z(a)\lambda
\]
satisfies
\[
J_{f_a}=u_5''-2\mathcal Z(a)=2(\mathcal Z-\mathcal Z(a)),
\]
and
Proposition~\ref{prop:source-determinant-form-odd-zero-unrooted-defect}(2)
gives
\[
D_{f_a}(a,h)=D_\lambda(a,h)\bigl(\Psi_a(h)-\mathcal Z(a)\bigr).
\]
So all route-2 source objects reduce exactly to the normalized planar ones.
\end{proof}

\begin{corollary}[The current route-2 package is universal under normalization]
\label{cor:current-route-two-package-universal-under-normalization}
Any configuration formally realizable by an analytic normalized planar model
\[
(\mathcal Z,\mathcal R,D_a)
\]
is also formally realizable by a genuine analytic source model for the current
route-2 package
\[
(c,v_5,\lambda,u_5),\qquad
\Psi_a,\qquad
\mathfrak D,\qquad
\mathcal W_{\mathrm{src}},\qquad
D_{f_a}.
\]
Accordingly, no contradiction theorem phrased only in terms of the current
route-2 package can distinguish the RH branch from an arbitrary analytic
normalized planar model.
\end{corollary}

\begin{proof}
This is immediate from
Proposition~\ref{prop:universality-threshold-current-route-two-package}.
\end{proof}

\begin{remark}[Explicit source-level countermodel to the current route-2 package]
\label{rem:explicit-source-level-countermodel-current-route-two}
Take
\[
\mathcal R(t):=t(t-1)(t^2-4t-2)
\]
and
\[
\mathcal Z(t):=\int_0^t \mathcal R(u)\,du
=
\frac15 t^5-\frac54 t^4+\frac23 t^3+t^2,
\]
and realize this normalized planar model by
\[
c(t)=1,\qquad v_5(t)=t,\qquad \lambda(t)=t^2,\qquad u_5''(t)=2\mathcal Z(t).
\]
Let
\[
q:=2-\sqrt6.
\]
In this genuine analytic source model, the unique born repeated-source critical
branch from the odd zero $0$ continues analytically down to $a=q$, lands
there with a nonlocal target
\[
\tau(q)\approx 0.3540339256\neq q,
\]
satisfies the anchor-sign safe relation on $(q,0)$, crosses the odd anchor
zero $q$ with no sign reset, and then undergoes a delayed rooted crossing at
\[
a_*\approx -0.5805673462,
\qquad
\tau(a_*)\approx 0.3888336023.
\]
Thus diagonal capture fails, and the whole delayed-crossing scenario is
realized inside a genuine analytic source model for the current route-2
package.
\end{remark}

\begin{remark}[The next live target must be non-universal]
\label{rem:next-live-target-must-be-non-universal}
By
Corollary~\ref{cor:current-route-two-package-universal-under-normalization},
the present lower route-2 package is exhausted: any further progress must use a
genuinely non-universal source-coupled input, i.e. an ingredient not already
absorbed by the normalized planar model $(\mathcal Z,\mathcal R,D_a)$.

However, the thread analysis shows that this non-universal input cannot be
identified merely with arbitrary canonical quotient-septic geometry, nor even
with the explicit direct $q^{(7)}$-slice by itself. Indeed,
Remark~\ref{rem:canonical-quotient-septic-closure-still-locally-free} shows
that arbitrary canonical $\Delta_7$ is still locally free once the lower
one-pair data are fixed, and
Proposition~\ref{prop:full-block-universality-direct-provenance-slice} together
with Corollary~\ref{cor:finite-linear-provenance-too-weak} show that the direct
highest-new provenance slice
\[
\phi=A\,q^{(7)}
\]
is fully universal on each generic one-pair block, even after any finite family
of linear provenance/matching conditions on $q$.

So any viable continuation must use something stronger: either a
\emph{continuum-strength provenance law} for $q$ on a whole block, or an
external mixed/canonical theorem that constrains $q$ beyond finite jet data.
Within the present draft, the first plausible places for such genuinely
non-universal input lie in the mixed / canonical package above arbitrary
quotient-septic geometry.
\end{remark}

\begin{remark}[What remains on the nonmixed $4$-point side]
\label{rem:what-remains-nonmixed-after-R-zero-bridge}
The preceding results now give a full local derivative-side picture of the
nonmixed exact-contact geometry. Exact nonmixed tangent-plane contacts force
zeros of $\mathcal R$
(Proposition~\ref{prop:exact-nonmixed-tangent-plane-contacts-force-R-zero});
intervals on which $\mathcal R$ has fixed nonzero sign are exactly-contact
sterile
(Corollary~\ref{cor:sign-definite-R-excludes-exact-nonmixed-contacts});
recurrent exact contacts force affine-planar collapse
(Corollary~\ref{cor:recurrent-exact-nonmixed-contacts-force-affine-planarity});
and isolated zeros of $\mathcal R$ are classified by parity
(Proposition~\ref{prop:parity-classification-isolated-nonmixed-R-zeros}).

At source level, the same structure is encoded by the gauge-invariant Wronskian
\[
\mathcal W_{\mathrm{src}}
=
J_\lambda J_{u_5}'-J_{u_5}J_\lambda',
\]
which differs from $\mathcal R$ only by a nonvanishing analytic factor on a
nonmixed arc
(Proposition~\ref{prop:source-level-Wronskian-nonmixed-four-point}). Thus exact
nonmixed contacts force zeros of $\mathcal W_{\mathrm{src}}$, one-sign
$\mathcal W_{\mathrm{src}}$-intervals are sterile, and recurrent contacts
force $\mathcal W_{\mathrm{src}}\equiv 0$
(Corollary~\ref{cor:source-level-Wronskian-consequences-nonmixed}).

Moreover, Proposition~\ref{prop:source-level-Wronskian-lifts-to-4x4} lifts
$\mathcal W_{\mathrm{src}}$ to the $4\times4$ Wronskian
\[
W(c,v_5,\lambda,u_5),
\]
while Proposition~\ref{prop:source-level-Wronskian-package-not-enough} shows
that, at the current source-level ratio/Wronskian/$4\times4$-determinant
level, the odd-zero pattern of $\mathcal W_{\mathrm{src}}$ and $\mathcal R$
is still formally unconstrained. Thus generic determinant algebra alone is no
longer enough to close the nonmixed branch.

The repeated-source / first-loss lane is now also closed. For each fixed
source point $a$, repeated-source exact contacts factor through the
quadratic-envelope scalar
\[
\Gamma_a(x)=
\frac{Y(x)-Y(a)-Y'(a)(x-a)}
{S(x)-S(a)-S'(a)(x-a)},
\]
equivalently through the source-determinant quotient
\[
\Psi_a(h)=\frac{D_{u_5}(a,h)}{D_\lambda(a,h)},
\]
with derivative controlled by the osculatory determinant
$\mathfrak D(a,h)$. The formal first-loss event rigidifies to a sheared
repeated-node tangent-contact configuration
(Remark~\ref{rem:first-loss-rigidification-nonmixed}), but this is no longer a
live independent bottleneck. Proposition~\ref{prop:first-loss-is-diagonal-odd-zero-birth}
shows that no attained separated first loss can occur, every first-loss
sequence is forced to the diagonal, and diagonal first loss occurs exactly at
isolated odd-order zeros of $\mathcal R$. Moreover,
Proposition~\ref{prop:uniform-fixed-level-nonrecurrence-sheared-three-point}
shows that for each fixed shear level $\gamma$, the corresponding sheared
opposite-side $3$-point branch is globally rigid on its maximal slope-overlap
block; and
Proposition~\ref{prop:odd-zero-critical-branch-uniformly-unrooted} together
with Corollary~\ref{cor:sheared-event-impossible-near-isolated-odd-R-zero}
show that near an isolated odd zero of $\mathcal R$, the diagonal critical
branch is uniformly unrooted, so the rooted sheared repeated-node
tangent-contact event itself does not occur there. Thus the sheared lane is
closed both in the fixed-$\gamma$ nonrecurrence direction and in the local
isolated-odd-zero direction.

Accordingly, the genuinely live nonmixed $4$-point burden is no longer local
birth, nor repeated-source first loss, nor generic determinant sign.  The
positive-kernel route now shows more: on any chamber where the repeated-source
segment sees exactly one sign change of $\mathcal R$, the born critical branch
is simple, has strictly negative slope, and cannot realize the rigidified
first-loss condition
\[
\Gamma_a(x)=\mathcal Z(a).
\]
At source level, the same branch is characterized by
\[
\mathfrak D(a,h)=0,
\]
but
Proposition~\ref{prop:source-determinant-form-odd-zero-unrooted-defect}
shows that its rooted numerator
\[
D_{u_5}(a,h)-\mathcal Z(x(a))\,D_\lambda(a,h)
\]
has a fixed signed leading term near an isolated odd zero, so the local branch
is uniformly unrooted there. Thus a local branch of
\[
\mathfrak D(a,h)=0
\]
cannot become a genuine obstruction until the source segment acquires a second
odd zero (equivalently a second sign-changing zero) of
\[
\mathcal W_{\mathrm{src}},
\]
or exits the local one-sign-change chamber. By
Proposition~\ref{prop:at-most-one-odd-zero-no-genuine-repeated-source-obstruction},
the positive-kernel route already excludes every opened nonmixed arc carrying
at most one odd zero of $\mathcal W_{\mathrm{src}}$ (equivalently of
$\mathcal R$).

Moreover, once one follows an analytic continuation of the born branch beyond
the local birth chamber, the rooted defect
\[
G(a)=D_a(\tau(a))-\mathcal Z(a)
\]
satisfies an exact first-order transport equation, and any nondegenerate rooted
point is forced to be a stationary point of the target map and a local fold of
the moving-shear equation
\[
A_a(x)-\mathcal Z(a)\,B_a(x)=0
\]
(Propositions~\ref{prop:rooted-defect-transport-stationary-first-root}
and~\ref{prop:nondegenerate-rooted-point-is-moving-shear-fold}).
The discriminant analysis rules out birth of such a rooted point at entry into
the adjacent-two-odd-zero chamber
(Proposition~\ref{prop:discriminant-entry-not-first-root}),
and the anchor-sign argument rules out any rooted crossing on a continuation
interval on which the anchor-side value
\[
\mathcal R(a)
\]
keeps a fixed sign
(Proposition~\ref{prop:anchor-sign-preservation-excludes-first-rooting} and
Corollary~\ref{cor:no-crossing-adjacent-two-odd-zero-under-anchor-sign-preservation}).

Beyond that, Proposition~\ref{prop:odd-anchor-passage-no-reset} shows that even
after the anchor crosses an odd-order zero of $\mathcal R$, no rooted crossing
can occur at that odd anchor zero itself or immediately after passage, while
Corollary~\ref{cor:first-post-anchor-chamber-dangerous} shows that the ensuing
post-anchor chambers alternate dangerous and safe. However,
Remark~\ref{rem:no-go-local-scalar-first-dangerous-chamber} shows that the local
normalized rooted equations and fold data are already formally compatible with
rooting on the first dangerous post-anchor chamber. More strongly,
Proposition~\ref{prop:universality-threshold-current-route-two-package} and
Corollary~\ref{cor:current-route-two-package-universal-under-normalization}
show that the entire current route-2 source package is universal under
normalization, and
Remark~\ref{rem:explicit-source-level-countermodel-current-route-two}
exhibits a genuine analytic source model realizing safe adjacent-two-odd-zero
behavior, nonlocal-target landing at the next odd anchor zero, odd-anchor
no-reset, and delayed rooted crossing.

So the sharpened remaining nonmixed burden is no longer a delayed no-crossing /
profile-realizability theorem \emph{within the current route-2 package}. The
current lower package is already exhausted, and the thread analysis shows that
arbitrary canonical quotient-septic geometry is still too free, while even the
explicit direct provenance slice
\[
\phi=A\,q^{(7)}
\]
is fully universal on each generic one-pair block. More strongly, any finite
family of linear provenance or matching conditions on $q$ is still too weak to
exclude direct-slice recurrence geometry.

Accordingly, any further progress must use a genuinely non-universal
continuum-strength provenance law for $q$, or else an external mixed/canonical
input that constrains $q$ on a whole block and is not already absorbed by the
normalized planar model $(\mathcal Z,\mathcal R,D_a)$. By
Theorem~\ref{thm:local-exhaustion-mixed-four-point-centered-base}, such an
external mixed input cannot come from a new sufficiently small local mixed
$4$-point obstruction near the centered diagonal base; any useful mixed
contribution must therefore be genuinely global/nonlocal, or arise from deeper
canonical/multi-pair structure.
\end{remark}

\begin{remark}[Compact status summary for the nonmixed $4$-point branch]
\label{rem:compact-status-summary-nonmixed-four-point}
The nonmixed $4$-point branch is now reduced to a source-level odd-zero
problem.  Locally, exact nonmixed contacts force zeros of $\mathcal R$, and
isolated zeros are classified by parity: even-order zeros are sterile, while
odd-order zeros seed local exact contacts.  Repeated-source first loss is not a
separate mechanism: it occurs if and only if one diagonally approaches an
odd-order zero of $\mathcal R$.  The sheared repeated-node lane is also closed
in the collapse direction: any genuine diagonal realization of the rigidified
sheared event already forces $\mathcal R\equiv 0$ locally, hence
affine-planar collapse.

On a nonmixed arc, the same structure is encoded by the gauge-invariant
source-level Wronskian
\[
\mathcal W_{\mathrm{src}}=J_\lambda J_{u_5}'-J_{u_5}J_\lambda',
\]
which differs from $\mathcal R$ only by a nonvanishing analytic factor.
Repeated-source first loss is not a separate mechanism: it occurs if and only
if one diagonally approaches an odd-order zero of $\mathcal R$, and the
corresponding rooted sheared repeated-node event is locally absent near an
isolated odd zero of $\mathcal R$. Near such an isolated odd zero, the born
repeated-source branch
\[
\mathfrak D(a,h)=0
\]
is uniformly unrooted at source level:
\[
D_{u_5}(a,h)-\mathcal Z(x(a))\,D_\lambda(a,h)\neq 0
\]
with a fixed signed leading term
(Proposition~\ref{prop:source-determinant-form-odd-zero-unrooted-defect}).
For each fixed shear level $\gamma$, the associated sheared opposite-side
$3$-point branch is globally rigid on its maximal slope-overlap block.
Moreover, the positive-kernel / source-determinant route now excludes every
opened nonmixed arc carrying at most one odd zero of
$\mathcal W_{\mathrm{src}}$:
on any chamber where the repeated-source segment sees exactly one sign change of
$\mathcal R$, the born critical branch is simple, has strictly negative slope,
and cannot realize the rigidified first-loss condition. Beyond that chamber, the
discriminant analysis rules out birth of a rooted point at entry into the
adjacent-two-odd-zero chamber, and the anchor-sign argument rules out any
crossing on a continuation interval where the anchor-side value
\[
\mathcal R(a)
\]
keeps a fixed sign.

The odd-anchor passage theorem shows that once the anchor itself crosses an
odd-order zero of $\mathcal R$, the integrating-factor defect keeps its incoming
sign across that odd anchor zero, so no rooted crossing can occur there or
immediately after it. The subsequent post-anchor chambers alternate dangerous
and safe. However, the local normalized rooted equations and the local
fold/transversality data are already formally realizable on the first dangerous
post-anchor chamber, and more strongly the entire current route-2 source package
is universal under normalization. Even arbitrary canonical quotient-septic
geometry remains too free, and the explicit direct provenance slice
\[
\phi=A\,q^{(7)}
\]
is fully universal on each generic one-pair block; indeed, even after any
finite family of linear provenance/matching conditions on $q$, direct-slice
recurrence geometry is still realizable. Thus the genuinely remaining nonmixed
burden is no longer a problem inside the current route-2 or arbitrary
quotient-septic package at all. Any further progress must use a genuinely
non-universal continuum-strength provenance law for $q$, or an external
mixed/canonical theorem that constrains $q$ on a whole block and is not already
absorbed by the normalized planar model. By the centered-base mixed local
exhaustion theorem, this external mixed input cannot be a new local mixed
$4$-point septic-jet phenomenon near the diagonal base; it must be genuinely
global/nonlocal or come from deeper canonical/multi-pair structure.
\end{remark}

\begin{remark}[Fixed-slope interval graphs, support pseudolines, and the remaining nonmixed global burden]
\label{rem:fixed-slope-interval-graph-monge-support-global-burden}
There is a useful combinatorial reformulation of the nonadjacent pair-pair
geometry on a nonmixed patch.

For fixed slope $m$, Lemma~\ref{lem:fixed-slope-intercept-sets} shows that each
ECT interval $I$ contributes an intercept interval
\[
\mathcal B_I(m)=\{\psi_I(x)-mx:x\in J_I\},
\]
and that a line $y=mx+b$ is $2$-secant to $\Gamma(I)$ exactly when
\[
b\in \operatorname{int}\mathcal B_I(m).
\]
So for fixed $m$, the same-slope pair-pair incidence graph is an interval
graph on the $b$-axis. In particular, any genuinely minimal surviving
nonadjacent recurrence cannot live at one fixed slope; it must be
slope-changing.

There is also a sharper support-curve reformulation. By
Proposition~\ref{prop:support-curve-calculus-ect},
Corollary~\ref{cor:single-crossing-betweenness-support-curves}, and
Proposition~\ref{prop:ordered-pair-overlap-windows-one-shot}, the support
curves
\[
b=\tau_I(m)
\]
are strictly concave pseudolines on common slope windows: each ordered pair has
at most one common interior tangent slope, triple crossings obey the usual
betweenness law, and each ordered pair has at most one directional or pair-pair
overlap window in slope.

Thus any genuinely minimal surviving nonadjacent recurrence cannot be fixed
slope, cannot be pair-repeating, and cannot disappear and later reappear on the
same ordered pair.

There is also a useful endpoint-gap reformulation. For every endpoint support
point $(\xi,y_\xi)$ and every ordered target branch $J$, the endpoint-gap
function
\[
G_{\xi,J}(h):=y_\xi-y(h)-m(h)\bigl(\xi-x(h)\bigr)
\]
is strictly monotone on $J$ by
Lemma~\ref{lem:endpoint-gap-monotonicity-ordered-target}. So each
endpoint/branch pair can contribute at most one remote tangent incidence.

These observations give the sharp current picture of the remaining nonmixed
global burden: any surviving bad core must be slope-changing, pair-nonrepeating,
and supported by a finite table of remote endpoint/branch incidences.

It is equally important to record the ceiling of this reduction. By
Remark~\ref{rem:intrinsic-affine-plane-universality}, intrinsic one-pair affine
geometry is universal for analytic plane curves. Thus pure planarity, parity,
interval-graph combinatorics, fixed-slope support order, or pseudoline
combinatorics by itself cannot close the remaining branch. Any final
contradiction must therefore use RH-specific source/provenance/canonical
structure, not only abstract strip/support geometry.

Accordingly, the natural next global target is not another local normal-form
analysis, but an RH-specific sign theorem for the remote endpoint-gap functions
$G_{\xi,J}$.
\end{remark}

\begin{remark}[Derivative-geometry meaning of the nonmixed invariant $\mathcal R$]
\label{rem:derivative-geometry-meaning-of-R}
On a nonmixed patch,
\[
\mathcal R=\frac{T_v}{K_v^2}
=
\frac{d}{dx}\!\left(\frac{Y''}{S''}\right).
\]
Thus $\mathcal R$ measures the failure of the lifted space curve
\[
x\longmapsto (Y(x),x,S(x))
\]
to be affine-planar at the derivative level: $\mathcal R=0$ means the ratio
$Y''/S''$ is locally constant, equivalently that the curve has vanishing
affine torsion and lies in an affine plane after integrating twice. In this
sense, the nonmixed $4$-point branch is controlled by a single scalar
derivative invariant rather than by raw coplanarity alone.
\end{remark}

\begin{remark}[Exact mixed reduction in terms of $\mathcal Z$ and $\mathcal R$]
\label{rem:mixed-order-bookkeeping-no-generic-parity-obstruction}
At a mixed lifted point
\[
K_v(0)=0,
\qquad
T_v(0)=0,
\]
set
\[
r:=\operatorname{ord}_0 K_v,
\qquad
s:=\operatorname{ord}_0 T_v.
\]
By Proposition~\ref{prop:exact-x-coordinate-identities-lifted-four-point},
\[
K_v=(x')^3S'',
\]
so on a good $v_5$-patch one has
\[
S''(x)=\kappa x^r+O(x^{r+1}),
\qquad \kappa\neq 0.
\]
After removing affine terms and integrating twice,
\[
S(x)=\frac{\kappa}{(r+1)(r+2)}x^{r+2}+O(x^{r+3}).
\]
Thus the mixed exponent $m$ of the monomial model
\[
Q_v(x)\sim (x^n,\ x,\ x^m)
\]
is exactly
\[
m=r+2.
\]

Now let $n$ be the first degree of $Y(x)$ not absorbed by
$\operatorname{span}\{1,x,S\}$, so that
\[
Y(x)=\alpha x^n+O(x^{n+1}),
\qquad \alpha\neq 0,
\qquad n>r+2.
\]
Using
\[
T_v=(x')^6\bigl(S''Y'''-S'''Y''\bigr),
\]
one finds
\[
s=n+r-3,
\qquad
n=s-r+3.
\]

Equivalently, on the nonmixed complement of $K_v=0$,
\[
\mathcal Z=\frac{Y''}{S''},
\qquad
\mathcal R=\mathcal Z'=\frac{T_v}{K_v^2}
\]
satisfy the order bookkeeping
\[
\operatorname{ord}_0\mathcal Z=n-r-2=s-2r+1,
\qquad
\operatorname{ord}_0\mathcal R=s-2r.
\]

Thus the dangerous low-order mixed models are finite and explicit:
\[
(m,n)=(3,4),\ (3,6),\ (4,6),
\]
equivalently
\[
(r,s)=(1,2),\ (1,4),\ (2,5),
\]
equivalently
\[
(r,\operatorname{ord}_0\mathcal Z)=(1,1),\ (1,3),\ (2,2),
\]
or
\[
(r,\operatorname{ord}_0\mathcal R)=(1,0),\ (1,2),\ (2,1).
\]

So the mixed lifted $4$-point endgame has also been reduced to a finite jet
obstruction list, now phrased in terms of the invariant pair
\[
(\mathcal Z,\mathcal R).
\]
\end{remark}

\begin{remark}[Exact finite obstruction coefficients for the low-order $4$-point branch]
\label{rem:exact-finite-obstruction-coefficients-low-order-four-point}
The explicit low-order $4$-point obstruction list can be reduced to a finite
jet package.

\medskip

\noindent
\textbf{Nonmixed side.}
Suppose $K_v(0)\neq 0$, and write
\[
s_k:=S^{(k)}(0),
\qquad
y_k:=Y^{(k)}(0),
\qquad
s_2\neq 0.
\]
Then the Taylor expansion
\[
\mathcal Z(x)=z_0+z_1x+z_2x^2+z_3x^3+z_4x^4+z_5x^5+z_6x^6+O(x^7)
\]
has canonical even coefficients
\[
z_2,\qquad z_4,\qquad z_6,
\]
and the dangerous nonmixed low-order cases are exactly:
\[
n=4 \iff z_2\neq 0,
\qquad
n=6 \iff z_2=0,\ z_4\neq 0,
\qquad
n=8 \iff z_2=z_4=0,\ z_6\neq 0.
\]
Explicitly,
\[
z_2=
\frac{
s_2^2y_4-2s_2s_3y_3+\bigl(2s_3^2-s_2s_4\bigr)y_2
}{
2\,s_2^3
},
\]
\[
z_4=
\frac{
s_2^4y_6
-s_2^3(4s_3y_5+6s_4y_4+4s_5y_3+s_6y_2)
+2s_2^2(6s_3^2y_4+12s_3s_4y_3+4s_3s_5y_2+3s_4^2y_2)
-12s_2s_3^2(2s_3y_3+3s_4y_2)
+24s_3^4y_2
}{
24\,s_2^5
},
\]
and
\[
z_6=\frac{\mathcal P_8(s_2,\dots,s_8;y_2,\dots,y_8)}{720\,s_2^7},
\]
where $\mathcal P_8$ is the explicit degree-$8$ numerator polynomial obtained
by Taylor division of $\mathcal Z=Y''/S''$.

\medskip

\noindent
\textbf{Mixed side.}
Suppose $K_v(0)=T_v(0)=0$, and let
\[
r:=\operatorname{ord}_0 K_v.
\]
For the two low-order mixed cases $r=1,2$, Taylor division of
$\mathcal Z=Y''/S''$ gives the exact obstruction coefficients
\[
\kappa_{1,1},\qquad \kappa_{1,3},\qquad \kappa_{2,2},
\]
defined by the expansions
\[
\mathcal Z(x)
=
\frac{y_2}{s_3}\frac1x
+
\left(\frac{y_3}{s_3}-\frac{s_4y_2}{2s_3^2}\right)
+
\kappa_{1,1}x+\kappa_{1,2}x^2+\kappa_{1,3}x^3+O(x^4)
\qquad (r=1),
\]
and
\[
\mathcal Z(x)
=
\frac{2y_2}{s_4}\frac1{x^2}
+
\left(\frac{2y_3}{s_4}-\frac{2s_5y_2}{3s_4^2}\right)\frac1x
+
\left(
\frac{y_4}{s_4}
-\frac{2s_5y_3}{3s_4^2}
-\frac{s_6y_2}{6s_4^2}
+\frac{2s_5^2y_2}{9s_4^3}
\right)
+\kappa_{2,1}x+\kappa_{2,2}x^2+O(x^3)
\qquad (r=2).
\]
The dangerous mixed low-order list is then exactly:
\[
(m,n)=(3,4)\iff \kappa_{1,1}\neq 0,
\]
\[
(m,n)=(3,6)\iff \kappa_{1,3}\neq 0,
\]
\[
(m,n)=(4,6)\iff \kappa_{2,2}\neq 0.
\]

Thus the full currently explicit \emph{low-order} $4$-point obstruction list is
equivalent to the finite set of coefficients
\[
\boxed{
z_2,\ z_4,\ z_6,\ \kappa_{1,1},\ \kappa_{1,3},\ \kappa_{2,2}.
}
\]

On the mixed side, the three coefficients
\[
\kappa_{1,1},\qquad \kappa_{1,3},\qquad \kappa_{2,2}
\]
exhaust the presently explicit low-order obstruction list.

On the nonmixed side, however, the parity classification from
Proposition~\ref{prop:parity-classification-isolated-nonmixed-R-zeros}
shows that local exact-contact birth occurs at every isolated odd-order zero of
$\mathcal R$. Accordingly,
\[
z_2,\qquad z_4,\qquad z_6
\]
should now be viewed only as the first three explicit low-order nonmixed
obstruction coefficients, not as a complete local nonmixed list.
\end{remark}

\begin{remark}[Wronskian-order reformulation of the mixed low-order list]
\label{rem:wronskian-order-reformulation-mixed-list}
On a good $v_5$-patch one has
\[
Y''=\frac{c^3}{M^3}J_{u_5},
\qquad
S''=\frac{c^3}{M^3}J_\lambda,
\qquad
\mathcal Z=\frac{J_{u_5}}{J_\lambda},
\qquad
\lambda=\frac{\Delta_7}{v_5}.
\]
Thus, if
\[
r:=\operatorname{ord}_0 J_\lambda=\operatorname{ord}_0 K_v,
\qquad
q:=\operatorname{ord}_0 J_{u_5}=\operatorname{ord}_0 Y'',
\]
then
\[
m=r+2,
\qquad
n=q+2,
\qquad
\operatorname{ord}_0\mathcal Z=q-r.
\]
Hence the dangerous mixed models
\[
(m,n)=(3,4),\ (3,6),\ (4,6)
\]
are exactly the Wronskian order-pairs
\[
(\operatorname{ord}_0J_\lambda,\operatorname{ord}_0J_{u_5})=(1,2),\ (1,4),\ (2,4).
\]

If
\[
J_\lambda(x)=x^rA(x),
\qquad
J_{u_5}(x)=x^rB(x),
\qquad
A(0)\neq 0,
\]
then
\[
\mathcal Z(x)=\frac{B(x)}{A(x)}.
\]
For $r=1$, the coefficients $\kappa_{1,1}$ and $\kappa_{1,3}$ are exactly the
$x^1$- and $x^3$-coefficients of the regularized ratio $B/A$; for $r=2$,
$\kappa_{2,2}$ is exactly the $x^2$-coefficient of $B/A$.

Moreover, since
\[
J_c=J_{v_5}=0,
\]
replacing
\[
u_5\longmapsto u_5+a\,c+b\,v_5+\gamma\,\lambda
\]
changes $J_{u_5}$ only by
\[
J_{u_5}\longmapsto J_{u_5}+\gamma J_\lambda.
\]
Thus the nonconstant Taylor coefficients of $\mathcal Z$ are canonical. In
particular, the currently explicit mixed low-order list is already a finite
Wronskian-ratio obstruction list.
\end{remark}

\begin{remark}[Regularized source scalar for the mixed residual quotient]
\label{rem:regularized-source-scalar-mixed-residual}
In the notation of Remark~\ref{rem:wronskian-order-reformulation-mixed-list},
write locally
\[
J_\lambda(x)=x^rA(x),
\qquad
J_{u_5}(x)=x^rB(x),
\qquad
A(0)\neq 0,
\]
and set
\[
Q(x):=\frac{B(x)}{A(x)}=\frac{J_{u_5}(x)}{J_\lambda(x)}.
\]
Then
\[
Q'(x)
=
\frac{J_\lambda J_{u_5}'-J_{u_5}J_\lambda'}{J_\lambda^2}
=
\frac{\mathcal W_{\mathrm{src}}}{J_\lambda^2},
\]
where
\[
\mathcal W_{\mathrm{src}}
=
J_\lambda J_{u_5}'-J_{u_5}J_\lambda'
\]
is the source-level Wronskian from
Proposition~\ref{prop:source-level-Wronskian-nonmixed-four-point}.

Moreover, under the affine gauge change
\[
u_5\longmapsto u_5+\alpha c+\beta v_5+\gamma\lambda
\]
one has
\[
J_{u_5}\longmapsto J_{u_5}+\gamma J_\lambda,
\qquad
Q\longmapsto Q+\gamma,
\]
so
\[
Q'(x)
=
\frac{\mathcal W_{\mathrm{src}}}{J_\lambda^2}
\]
is gauge-invariant. Thus, modulo the trivial additive constant in the
regularized ratio $Q$, the mixed residual quotient germ is already controlled
by the single scalar
\[
\mathcal D:=Q'=\frac{\mathcal W_{\mathrm{src}}}{J_\lambda^2}.
\]

On the nonmixed complement of $K_v=0$,
Remark~\ref{rem:below-lifted-wronskian-ratio-reformulation} and
Proposition~\ref{prop:source-level-Wronskian-nonmixed-four-point} give
\[
\mathcal D
=
\frac{M}{c^2}\,\mathcal R.
\]
So the exact fixed-shear residual lane is governed by the same regularized
source scalar already present on the nonmixed side, rather than by a second
primitive pointwise highest-new field.
\end{remark}

\begin{proposition}[Projective secant-shadow reduction on an overlap patch]
\label{prop:conditional-mixed-lambda-shear-projective-shadow-reduction}
Let
\[
U\Subset\{c\neq 0,\ u_5\neq 0,\ v_5\neq 0,\ M\neq 0\}
\]
be a good overlap patch. Set
\[
\lambda:=\frac{\Delta_7}{v_5},
\qquad
x:=\frac{v_5}{c},
\qquad
Y:=\frac{u_5}{c},
\qquad
S:=\frac{\lambda}{c}=\frac{\Delta_7}{c\,v_5}.
\]

Assume there exists a close exact mixed two-point bad branch
\[
(h_{1,n},h_{2,n},\Lambda_n),
\qquad
h_{1,n}\neq h_{2,n},
\qquad
h_{1,n},h_{2,n}\in U,
\qquad
h_{1,n},h_{2,n}\to m\in U,
\]
and let $(\eta_n,\xi_n,\gamma_n)$ be the branch-generated moving
$\lambda$-shears, so that
\[
(c-\eta_n\lambda)(h_{1,n})=\Lambda_n(c-\eta_n\lambda)(h_{2,n}),
\]
\[
(v_5-\xi_n\lambda)(h_{1,n})=\Lambda_n(v_5-\xi_n\lambda)(h_{2,n}),
\]
\[
(u_5-\gamma_n\lambda)(h_{1,n})=\Lambda_n(u_5-\gamma_n\lambda)(h_{2,n}).
\]
For each $n$, set
\[
A_n:=c-\eta_n\lambda,
\qquad
B_n:=v_5-\xi_n\lambda,
\qquad
F_n:=u_5-\gamma_n\lambda.
\]

Then, after discarding finitely many terms, the branch determines the exact
projective secant shadow
\[
\mathbf\Sigma(h_{1,n},h_{2,n})
:=
[x_1S_2-x_2S_1:\ Y_1S_2-Y_2S_1:\ S_2-S_1],
\]
where
\[
x_i=x(h_{i,n}),
\qquad
Y_i=Y(h_{i,n}),
\qquad
S_i=S(h_{i,n}).
\]
More precisely, for all large $n$ one has the exact identity
\[
[B_n(h_{1,n}):F_n(h_{1,n}):A_n(h_{1,n})]
=
[B_n(h_{2,n}):F_n(h_{2,n}):A_n(h_{2,n})]
=
\mathbf\Sigma(h_{1,n},h_{2,n}).
\]
In particular, the secant shadow is independent of the moving shears
$(\eta_n,\xi_n,\gamma_n)$ and of $\Lambda_n$.

Equivalently, in raw coordinates,
\[
\mathbf\Sigma(h_1,h_2)
=
[v_{5,1}\lambda_2-v_{5,2}\lambda_1:\ u_{5,1}\lambda_2-u_{5,2}\lambda_1:\ c_1\lambda_2-c_2\lambda_1].
\]

As $h_2\to h_1=a$, the secant shadow converges projectively to the tangent
shadow
\[
\mathbf\Pi(a)
:=
[xS_x-S:\ YS_x-SY_x:\ S_x]\big|_{x=x(a)}.
\]
A direct calculation gives
\[
\mathbf\Pi(x)\wedge \partial_x\mathbf\Pi(x)
=
S(x)\,
\bigl(
S_xY_{xx}-S_{xx}Y_x,\ S_{xx},\ -S_{xx}
\bigr).
\]
Hence $\mathbf\Pi$ is projectively immersive wherever
\[
S\neq 0
\qquad\text{and}\qquad
S_{xx}\neq 0.
\]
Since on the $v_5$-patch
\[
K_v=(x')^3S_{xx},
\]
any such close exact mixed two-point bad branch can accumulate only on
\[
\{S=0\}\cup\{K_v=0\}.
\]
\end{proposition}

\begin{proof}
Since
\[
x'=\left(\frac{v_5}{c}\right)'=\frac{M}{c^2},
\]
the hypothesis $M\neq 0$ implies that $x=v_5/c$ is a local coordinate near
$m$. After discarding finitely many terms we may therefore assume
\[
x(h_{1,n})\neq x(h_{2,n}).
\]

Fix such an $n$, and abbreviate
\[
\eta=\eta_n,\qquad \xi=\xi_n,\qquad \gamma=\gamma_n.
\]
Write
\[
p_i:=(x_i,Y_i,1)\in\mathbf R^3,
\qquad
V:=(\xi,\gamma,\eta)\in\mathbf R^3,
\qquad
i=1,2.
\]
Since
\[
A_n=c(1-\eta S),
\qquad
B_n=c(x-\xi S),
\qquad
F_n=c(Y-\gamma S),
\]
the exact repeated-source equations imply
\[
[B_n(h_{1,n}):F_n(h_{1,n}):A_n(h_{1,n})]
=
[p_1-S_1V]
\]
and
\[
[B_n(h_{2,n}):F_n(h_{2,n}):A_n(h_{2,n})]
=
[p_2-S_2V].
\]
Because these two projective points are equal, there is a common projective
value.

If $S_1=S_2$, then
\[
[p_1-S_1V]=[p_2-S_1V]
\]
forces $p_1=p_2$, hence $x_1=x_2$, contradicting the choice of local
coordinate. Therefore
\[
S_1\neq S_2.
\]
Subtracting
\[
p_1-S_1V=p_2-S_2V
\]
gives
\[
p_1-p_2=(S_1-S_2)V,
\]
hence
\[
V=\frac{p_1-p_2}{S_1-S_2}.
\]
Substituting back,
\[
p_1-S_1V
=
p_1-S_1\frac{p_1-p_2}{S_1-S_2}
=
\frac{S_2p_1-S_1p_2}{S_2-S_1}.
\]
Therefore the common projective value is exactly
\[
[S_2p_1-S_1p_2]
=
[x_1S_2-x_2S_1:\ Y_1S_2-Y_2S_1:\ S_2-S_1],
\]
which is the asserted secant shadow and is independent of
$(\eta_n,\xi_n,\gamma_n)$ and of $\Lambda_n$. The raw-coordinate formula is the
same identity multiplied through by $c_1c_2$.

Now divide the homogeneous secant vector by $x_2-x_1$ and let $h_2\to h_1=a$.
The three components converge to
\[
xS_x-S,\qquad YS_x-SY_x,\qquad S_x,
\]
so the secant shadow converges projectively to
\[
\mathbf\Pi(a)
=
[xS_x-S:\ YS_x-SY_x:\ S_x].
\]

Differentiating $\mathbf\Pi$ in the local $x$-coordinate gives
\[
\mathbf\Pi(x)\wedge \partial_x\mathbf\Pi(x)
=
S(x)\,
\bigl(
S_xY_{xx}-S_{xx}Y_x,\ S_{xx},\ -S_{xx}
\bigr).
\]
Therefore $\mathbf\Pi$ is projectively immersive wherever $S\neq 0$ and
$S_{xx}\neq 0$. On a close bad branch the secant shadows converge to tangent
shadows, so accumulation is possible only where immersion fails, i.e. on
\[
\{S=0\}\cup\{S_{xx}=0\}.
\]
Finally, on the $v_5$-patch one has
\[
K_v=(x')^3S_{xx},
\]
so $\{S_{xx}=0\}=\{K_v=0\}$ on $U$.
\end{proof}

\begin{remark}[Useful overlap scalar]
\label{rem:structural-mixed-overlap-scalar-two-point}
On the same overlap patch, set
\[
\tau_u:=-\frac{\Delta_7}{c\,u_5},
\qquad
\mathcal G:=x'S_x+Y'\tau_u'.
\]
Then $\mathcal G$ admits the exact factorization
\[
\mathcal G
=
-\frac{W(u_5,v_5)\,W(u_5,\lambda)}{c^2u_5^2},
\qquad
\lambda=\frac{\Delta_7}{v_5},
\]
where $W(f,g)=fg'-f'g$. Equivalently,
\[
\mathcal G
=
\frac{u_5v_5'-u_5'v_5}{c^2}
\left(\frac{\Delta_7}{u_5v_5}\right)'.
\]
In particular, $\mathcal G\not\equiv 0$ on any connected overlap patch
contained in $\{c\neq 0,\ M\neq 0\}$: the factor $W(u_5,v_5)$ cannot vanish
identically there without forcing linear dependence of
\[
1,\qquad \frac{u_5}{c},\qquad \frac{v_5}{c},
\]
and $W(u_5,\lambda)$ cannot vanish identically there without forcing
\[
\Delta_7=\kappa\,u_5v_5,
\]
which is incompatible with the direct septic slice carried by $\Delta_7$.
This does not yet exclude the mixed two-point branch, but it records the first
genuinely nontrivial $Y$-$S$ overlap coupling scalar exposed by the current
attack.

At highest-new weight, this scalar remains in the same differential tower as
the direct septic source field:
\[
\mathcal G^{[1]}
=
-\frac{W(u_5,v_5)}{c^2u_5^2}\,W(u_5,A\,q^{(7)}),
\qquad
A=\kappa\eta_2.
\]
So $\mathcal G$ does not produce a second primitive pointwise highest-new
generator. It does, however, provide the first useful external-looking rank-$2$
pointwise coupling on the overlap patch: together,
\[
\lambda^{[1]}=A\,q^{(7)}
\qquad\text{and}\qquad
\mathcal G^{[1]}
\]
recover $(q^{(7)},(q^{(7)})')$ generically.
\end{remark}

\begin{proposition}[Same-tower closure for the current explicit highest-new overlap/canonical package]
\label{prop:same-tower-closure-current-explicit-highest-new-overlap-canonical-package}
Work on a good overlap patch
\[
U\Subset\{c\neq 0,\ u_5\neq 0,\ v_5\neq 0,\ M\neq 0\}.
\]
Let $K_{\mathrm{low}}$ denote the differential field generated by the lower
controlled fields on $U$, in particular by
\[
c,\ u_5,\ v_5,\ M,\ A=\kappa\eta_2
\]
and their analytic derivatives, and write
\[
r:=q^{(7)},
\qquad
x:=\frac{v_5}{c},
\qquad
Y:=\frac{u_5}{c},
\qquad
S:=\frac{\Delta_7}{c\,v_5}=\frac{\lambda}{c}.
\]
Then every currently explicit pointwise highest-new field obtained from the
quotient-septic package $(\Delta_7,\lambda)$, the overlap fields
$(x,Y,S,\tau_u)$, the dual tangent-line fields $(\mu,\lambda,\beta,m,W)$, and
the exact fixed-shear overlap constructions by finitely many algebraic
operations, $x$-derivatives, Wronskians, finite jet extractions, and
projectivizations lies in the differential algebra
\[
K_{\mathrm{low}}\langle r\rangle.
\]
More precisely, the weight-one part of each such field is
$K_{\mathrm{low}}$-linear in finitely many $x$-jets
\[
\partial_x^j S^{[1]},
\qquad
j=0,\dots,N,
\]
for some finite $N$, and
\[
S^{[1]}=\frac{A}{c}\,r.
\]

In particular, the currently explicit pointwise highest-new fields
\[
\Delta_7^{[1]},\ \lambda^{[1]},\ \beta^{[1]},\ m^{[1]},\ W^{[1]},\
S^{[1]},\ \mathcal G^{[1]},\ J_\lambda^{[1]},\ \mathcal W_{\mathrm{src}}^{[1]},
\]
the projected fixed-shear germ at weight one, and the quartic--sextic
coefficient $\kappa_{2,2}$ are all same-tower. Hence no second primitive
pointwise highest-new generator appears in the current explicit
overlap/canonical package. This is a statement about the present explicit
pointwise package only, and does not address nonlocal, relational, or
downstream transported-state obstructions.
\end{proposition}

\begin{proof}
On the overlap patch, the lower fields $x$, $Y$, $\mu$, $c$, $u_5$, $v_5$, and
the coordinate-change factors are already fixed at weight zero. Thus all
highest-new weight enters only through
\[
\lambda^{[1]}=A\,r,
\qquad
S^{[1]}=\frac{\lambda^{[1]}}{c}=\frac{A}{c}\,r.
\]
Linearizing any allowed algebraic/differential expression at weight one kills
quadratic interactions among highest-new pieces, so the weight-one part is a
$K_{\mathrm{low}}$-linear combination of finitely many $x$-jets of $S^{[1]}$.
Since
\[
\partial_x=\frac{c^2}{M}\,\partial_h
\]
on $U$, each $\partial_x^jS^{[1]}$ lies in $K_{\mathrm{low}}\langle r\rangle$.
Projectivization only divides by lower nonvanishing factors, and finite jet
coefficients such as $\kappa_{2,2}$ are extracted from finitely many such jets.
The stated list therefore remains same-tower.
\end{proof}

\begin{corollary}[Best current external-looking rank-two coupling]
\label{cor:best-current-external-rank-two-coupling-overlap}
On the same overlap patch,
\[
\begin{pmatrix}
\lambda^{[1]}\\[2pt]
\mathcal G^{[1]}
\end{pmatrix}
=
\begin{pmatrix}
A & 0\\[4pt]
-\dfrac{W(u_5,v_5)}{c^2u_5^2}\bigl(u_5A'-u_5'A\bigr)
&
-\dfrac{W(u_5,v_5)}{c^2u_5}A
\end{pmatrix}
\begin{pmatrix}
r\\[2pt]
r'
\end{pmatrix}.
\]
In particular,
\[
\det
\begin{pmatrix}
A & 0\\[4pt]
-\dfrac{W(u_5,v_5)}{c^2u_5^2}\bigl(u_5A'-u_5'A\bigr)
&
-\dfrac{W(u_5,v_5)}{c^2u_5}A
\end{pmatrix}
=
-\frac{W(u_5,v_5)}{c^2u_5}A^2.
\]
So wherever
\[
A\neq 0,
\qquad
u_5\neq 0,
\qquad
W(u_5,v_5)\neq 0,
\]
the pair $(\lambda^{[1]},\mathcal G^{[1]})$ is a genuine rank-two pointwise
coupling to $(r,r')$. It does not produce a new generator, but it is the best
current external-looking rank-two coupling on the overlap patch.
\end{corollary}

\begin{proof}
This is the explicit highest-new expansion of
Remark~\ref{rem:structural-mixed-overlap-scalar-two-point}, using
\[
\lambda^{[1]}=A\,r,
\qquad
\mathcal G^{[1]}
=
-\frac{W(u_5,v_5)}{c^2u_5^2}
\Bigl(u_5A\,r'+(u_5A'-u_5'A)\,r\Bigr).
\]
\end{proof}

\begin{remark}[Work in progress: explicit-pointwise bridge to the local good-patch detector]
\label{rem:wip-explicit-pointwise-bridge-good-patch-detector}
Proposition~\ref{prop:same-tower-closure-current-explicit-highest-new-overlap-canonical-package}
and
Corollary~\ref{cor:best-current-external-rank-two-coupling-overlap}
show that the present explicit highest-new package has only one currently
visible pointwise source, namely the same tower generated by
\[
r=q^{(7)},
\]
even though the pair
\[
(\lambda^{[1]},\mathcal G^{[1]})
\]
recovers $(r,r')$ generically on the overlap patch.

This shifts the next bridge question. The issue is no longer whether another
explicit pointwise highest-new field appears, but whether this single explicit
pointwise direction is already seen by the local good-patch compensator. If it
is, then the first same-order local compensator should already be
quotient-septic-visible; if it is not, then the entire current explicit
pointwise package is invisible to that local detector, and any surviving
obstruction must already lie in nonlocal transport,
collision-functoriality failure, or other provenance-sensitive two-atom
structure.

At present this is a theorem target, not a proved equivalence.
\end{remark}

\begin{remark}[Ceiling of the current pointwise highest-new package]
\label{rem:ceiling-current-pointwise-highest-new-package}
The present explicit highest-new audit has now reached a stable negative
conclusion. All currently visible pointwise highest-new fields --- including
\[
\Delta_7^{[1]},\ \lambda^{[1]},\ \beta^{[1]},\ m^{[1]},\ W^{[1]},\ S^{[1]},
\ \mathcal G^{[1]}
\]
and the quartic--sextic transported coefficient package --- remain in the
differential algebra generated by the lower data together with
\[
r=q^{(7)}.
\]
So no second primitive explicit pointwise highest-new generator has emerged.

Combined with
Remark~\ref{rem:status-direct-septic-witness}
and
Remark~\ref{rem:audit-even-amplitude-obstruction-naive-source-summed-two-point},
this means that the honest remaining order-$7$/two-point burden is not
``one more local field'' and not the naive source-summed corrected-block route.
What remains is package-level and provenance-sensitive: for example a
reduced-image-germ theorem or collision-functoriality strong enough to force
quotient-septic collapse.
\end{remark}

\begin{proposition}[Exact fixed-shear overlap branches off $K_v=0$ force a fold at $S=0$]
\label{prop:fixed-shear-overlap-branch-off-Kv-forces-fold-at-S-zero}
Work on the overlap patch of
Proposition~\ref{prop:conditional-mixed-lambda-shear-projective-shadow-reduction},
and write
\[
\Phi(x):=
\left(
\frac{x-\xi_* S}{1-\eta_* S},
\frac{Y-\gamma_* S}{1-\eta_* S}
\right),
\qquad
S=\frac{\Delta_7}{c\,v_5}.
\]
Assume there is a close exact fixed-shear mixed $2$-point branch
\[
x_{1,n}\neq x_{2,n},
\qquad
x_{1,n},x_{2,n}\to a,
\qquad
\Phi(x_{1,n})=\Phi(x_{2,n})
\quad(n\ge 1).
\]
If
\[
K_v(a)\neq 0,
\]
then
\[
S(a)=0,
\qquad
\Phi'(a)=0,
\qquad
S_x(a)\neq 0.
\]
More precisely,
\[
0=\Phi_1'(a)=1-(\xi_*-\eta_* x(a))S_x(a),
\]
\[
0=\Phi_2'(a)=Y_x(a)-(\gamma_*-\eta_* Y(a))S_x(a).
\]
In particular, setting
\[
\beta:=\xi_*-\eta_* x(a)=\frac{1}{S_x(a)},
\]
and centering the local $x$-coordinate at $a$, the first-coordinate fold model
\[
U(x):=(x-x(a))-\beta S(x)
\]
satisfies
\[
U(a)=0,
\qquad
U_x(a)=0,
\qquad
U_{xx}(a)=-\beta S_{xx}(a)\neq 0.
\]
Hence, after analytic source reparametrization, the first coordinate of the
fixed-shear projected germ is
\[
t^2.
\]
So any exact fixed-shear mixed branch away from $K_v=0$ is already forced into
the same fold regime that underlies the nonmixed odd-zero analysis. In
particular, the remaining off-$K_v=0$ burden is best understood as an
absorption problem into the existing nonmixed odd-zero / affine-planar package,
rather than as a search for a new mixed local carrier.
\end{proposition}

\begin{proof}
A fixed-shear branch is, in particular, a close exact mixed two-point bad branch
on the same overlap patch. By
Proposition~\ref{prop:conditional-mixed-lambda-shear-projective-shadow-reduction},
any accumulation point of such a branch lies in
\[
\{S=0\}\cup\{K_v=0\}.
\]
So if $K_v(a)\neq 0$, then necessarily
\[
S(a)=0.
\]

Because
\[
\Phi(x_{1,n})=\Phi(x_{2,n}),
\qquad
x_{1,n}\neq x_{2,n},
\qquad
x_{1,n},x_{2,n}\to a,
\]
the mean value theorem applied to each coordinate of $\Phi$ gives points
\[
\zeta_{1,n},\zeta_{2,n}\in(x_{1,n},x_{2,n})
\]
such that
\[
\Phi_1'(\zeta_{1,n})=0,
\qquad
\Phi_2'(\zeta_{2,n})=0.
\]
Passing to the limit yields
\[
\Phi'(a)=0.
\]

Differentiating the two coordinates of $\Phi$ gives
\[
\Phi_1'
=
\frac{1-(\xi_*-\eta_* x)S_x-\eta_* S}{(1-\eta_* S)^2},
\]
\[
\Phi_2'
=
\frac{Y_x-(\gamma_*-\eta_* Y)S_x-\eta_* S\,Y_x}{(1-\eta_* S)^2}.
\]
Evaluating at $x=a$ and using $S(a)=0$ yields the displayed identities
\[
0=\Phi_1'(a)=1-(\xi_*-\eta_* x(a))S_x(a),
\]
\[
0=\Phi_2'(a)=Y_x(a)-(\gamma_*-\eta_* Y(a))S_x(a).
\]
The first identity implies
\[
S_x(a)\neq 0
\qquad\text{and}\qquad
\beta:=\xi_*-\eta_* x(a)=\frac{1}{S_x(a)}.
\]

Now center the local $x$-coordinate at $a$ and set
\[
U(x):=(x-x(a))-\beta S(x).
\]
Then
\[
U(a)=0,
\qquad
U_x(a)=1-\beta S_x(a)=0,
\qquad
U_{xx}(a)=-\beta S_{xx}(a).
\]
Since on the $v_5$-patch
\[
K_v=(x')^3S_{xx},
\qquad
x'=\frac{M}{c^2}\neq 0,
\]
the hypothesis $K_v(a)\neq 0$ implies
\[
S_{xx}(a)\neq 0.
\]
Hence
\[
U_{xx}(a)\neq 0.
\]
By the analytic Morse lemma in one variable, after analytic source
reparametrization the first coordinate is $t^2$.
\end{proof}

\begin{proposition}[Off-septic-zero $K_v$-accumulation is already mixed low-order]
\label{prop:off-septic-zero-Kv-accumulation-is-mixed-low-order}
Work on the overlap patch of
Proposition~\ref{prop:conditional-mixed-lambda-shear-projective-shadow-reduction}.
Assume there exists a close exact mixed two-point bad branch accumulating at a
point $a$ such that
\[
K_v(a)=0,
\qquad
S(a)\neq 0.
\]
Then
\[
Y_{xx}(a)=0.
\]
Equivalently, on the $v_5$-patch one has
\[
J_\lambda(a)=0,
\qquad
J_{u_5}(a)=0.
\]
So any such accumulation already lies in the existing mixed low-order
obstruction regime.
\end{proposition}

\begin{proof}
By Proposition~\ref{prop:conditional-mixed-lambda-shear-projective-shadow-reduction},
any close exact mixed two-point bad branch determines the tangent shadow
\[
\mathbf\Pi(x)=[xS_x-S:\ YS_x-SY_x:\ S_x]
\]
with projective immersion defect
\[
\mathbf\Pi(x)\wedge \partial_x\mathbf\Pi(x)
=
S(x)\bigl(S_xY_{xx}-S_{xx}Y_x,\ S_{xx},\ -S_{xx}\bigr).
\]
At a point $a$ with $K_v(a)=0$, one has $S_{xx}(a)=0$ on the $v_5$-patch,
so this simplifies to
\[
\mathbf\Pi(a)\wedge \partial_x\mathbf\Pi(a)
=
S(a)\,Y_{xx}(a)\,\bigl(S_x(a),\,0,\,S(a)-x(a)S_x(a)\bigr).
\]
If $S(a)\neq 0$ and $Y_{xx}(a)\neq 0$, then the vector on the right is
nonzero, so $\mathbf\Pi$ is projectively immersive at $a$, contradicting
accumulation of secant shadows there. Hence
\[
Y_{xx}(a)=0.
\]
On the $v_5$-patch,
\[
J_\lambda=\frac{M^3}{c^3}S_{xx},
\qquad
J_{u_5}=\frac{M^3}{c^3}Y_{xx},
\]
so $K_v(a)=0$ implies $J_\lambda(a)=0$, and the displayed vanishing of
$Y_{xx}(a)$ gives $J_{u_5}(a)=0$.
\end{proof}

\begin{corollary}[Simple septic-zero fixed-shear branches off $K_v=0$ transfer to nonmixed four-point geometry]
\label{cor:fixed-shear-simple-septic-zero-off-Kv-transfers-to-nonmixed-four-point}
Work on the overlap patch of
Proposition~\ref{prop:fixed-shear-overlap-branch-off-Kv-forces-fold-at-S-zero}.
Assume there is a close exact fixed-shear mixed $2$-point branch
\[
x_{1,n}\neq x_{2,n},
\qquad
x_{1,n},x_{2,n}\to a,
\qquad
\Phi(x_{1,n})=\Phi(x_{2,n}),
\]
and
\[
K_v(a)\neq 0.
\]
Then, after shrinking to a sufficiently small neighborhood of $a$, this lane
produces a local interval-persistent nonmixed $4$-point family for the lifted
curve
\[
Q(x)=(x,Y(x),S(x)).
\]
Consequently it is absorbed by the existing nonmixed $4$-point theory;
in particular, it is not a surviving genuinely mixed $2$-point
derivative-geometric obstruction.
\end{corollary}

\begin{proof}
By Proposition~\ref{prop:fixed-shear-overlap-branch-off-Kv-forces-fold-at-S-zero},
\[
S(a)=0,
\qquad
S_x(a)\neq 0,
\]
and the first coordinate of the fixed-shear projected germ is in fold normal
form. Thus, after analytic source reparametrization, there is a local parameter
$t$ with
\[
\Phi(x(t))=(t^2,\mathcal V(t))
\]
and repeated nearby values force
\[
\mathcal V(t)=\mathcal V(-t)
\]
for all sufficiently small $t$. Equivalently, in the actual local septic
coordinate $s:=S$, there are even analytic germs $E_x,E_Y$ such that
\[
\frac{x(s)-\xi_*s}{1-\eta_*s}=E_x(s),
\qquad
\frac{Y(s)-\gamma_*s}{1-\eta_*s}=E_Y(s),
\qquad
E_x(-s)=E_x(s),\quad E_Y(-s)=E_Y(s).
\]

Define
\[
Q(s):=(x(s),Y(s),s).
\]
If $\eta_*\neq 0$, set
\[
P:=\left(\frac{\xi_*}{\eta_*},\frac{\gamma_*}{\eta_*},\frac1{\eta_*}\right).
\]
Then the parity formulas imply that $Q(s),Q(-s),P$ are collinear for all
small $s$. If $\eta_*=0$, then
\[
Q(s)-Q(-s)=2s(\xi_*,\gamma_*,1),
\]
so all symmetric secants are parallel to one fixed direction, which is the same
projective statement with a point at infinity. Hence for every small
$0<s_1<s_2$, the four points
\[
Q(-s_2),\ Q(-s_1),\ Q(s_1),\ Q(s_2)
\]
are coplanar.

Finally, by Proposition~\ref{prop:fixed-shear-overlap-branch-off-Kv-forces-fold-at-S-zero},
\[
x_s(a)=\frac{1}{S_x(a)}\neq 0,
\]
so $x$ is monotone in $s$ near $a$. Thus the four coplanar points above
lie on one local $x$-arc in the nonmixed ordered pattern, yielding a local
interval-persistent nonmixed $4$-point family. The final claim follows from
Proposition~\ref{prop:interval-persistent-nonmixed-four-point-forces-affine-planarity}.
\end{proof}

\begin{remark}[Exact shadow-adapted freezing on the simple septic-zero slice]
\label{rem:exact-shadow-adapted-freezing-simple-septic-zero}
Work on a good overlap patch and suppose
\[
S=\frac{\Delta_7}{c\,v_5}
\]
is a local coordinate near a simple septic-zero point $a$, so
\[
S(a)=0,
\qquad
S_x(a)\neq 0.
\]
Write $s:=S$ and let
\[
C(s):=(x(s),Y(s)),
\qquad
x=\frac{v_5}{c},
\qquad
Y=\frac{u_5}{c}.
\]
For any distinct nearby secant pair $s_1\neq s_2$, define
\[
\xi(s_1,s_2):=\frac{x(s_1)-x(s_2)}{s_1-s_2},
\qquad
\gamma(s_1,s_2):=\frac{Y(s_1)-Y(s_2)}{s_1-s_2}.
\]
Then one has the exact identities
\[
x(s_1)-\xi(s_1,s_2)\,s_1
=
x(s_2)-\xi(s_1,s_2)\,s_2,
\]
\[
Y(s_1)-\gamma(s_1,s_2)\,s_1
=
Y(s_2)-\gamma(s_1,s_2)\,s_2.
\]
Equivalently, on the simple septic-zero slice every exact secant-shadow pair
already carries an exact fixed-shear normalization with
\[
\eta=0,
\qquad
\xi=\xi(s_1,s_2),
\qquad
\gamma=\gamma(s_1,s_2).
\]

Thus the apparently projective, non-affinely-frozen septic-zero lane is not a
separate carrier mechanism on the simple $S=0$ slice: in the shadow-adapted
chart it freezes exactly and therefore falls under the same exact fixed-shear
analysis used in
Corollary~\ref{cor:fixed-shear-simple-septic-zero-off-Kv-transfers-to-nonmixed-four-point}.
\end{remark}

\begin{remark}[Status of the mixed $\lambda$-shear route]
\label{rem:wip-mixed-lambda-shear-route-status}
Proposition~\ref{prop:conditional-mixed-lambda-shear-projective-shadow-reduction}
shows that the downstream secant-shadow reduction is already exact and fully
projective: a close exact mixed two-point bad branch determines the projective
secant shadow
\[
\mathbf\Sigma(h_1,h_2)
=
[x_1S_2-x_2S_1:\ Y_1S_2-Y_2S_1:\ S_2-S_1],
\qquad
S=\frac{\Delta_7}{c\,v_5},
\]
independently of any affine freezing of the branch-generated $\lambda$-shears.
On the simple septic-zero slice $S=0,\ S_x\neq 0$, there is moreover a
canonical shadow-adapted exact freezing with $\eta=0$; see
Remark~\ref{rem:exact-shadow-adapted-freezing-simple-septic-zero}.

The overlap geometry is now sharply split into three pieces.

First,
Proposition~\ref{prop:off-septic-zero-Kv-accumulation-is-mixed-low-order}
shows that if a close mixed branch accumulates at a point with
\[
K_v=0,
\qquad
S\neq 0,
\]
then
\[
J_\lambda=J_{u_5}=0,
\]
so this part of the $K_v=0$ locus is already absorbed by the existing mixed
low-order obstruction package.

Second,
Corollary~\ref{cor:fixed-shear-simple-septic-zero-off-Kv-transfers-to-nonmixed-four-point}
shows that in the conditional exact fixed-shear overlap lane, the slice
\[
\{S=0\}\cap\{K_v\neq 0\}
\]
is not a surviving mixed $2$-point carrier lane: it transfers into the
established nonmixed $4$-point geometry.

Thus the only genuinely residual exact fixed-shear mixed corner is
\[
\{S=0\}\cap\{K_v=0\}.
\]
On that corner,
Proposition~\ref{prop:finite-transport-jet-quotient-descent-fixed-shear-corner}
shows that every finite parity-normalized transport jet descends analytically to
the quotient variable $u=t^2$. In the quartic--sextic case,
Corollary~\ref{cor:rank-one-transport-state-quartic-sextic-fixed-shear} shows
that every finite quotient-visible transport state is locally an analytic
function of $Q$ alone, and
Corollary~\ref{cor:no-finite-order-hidden-reset-quartic-sextic-fixed-shear}
rules out any third local reset beyond the involution symmetry. So the residual
branch is no longer merely quotient-diagonal / involutive; its finite
transport state has local rank $1$.

If the actual corrected cubic/quintic defects are state-local in this descended
transport state, then
Proposition~\ref{prop:conditional-state-local-defect-closure-fixed-shear-corner}
already excludes the corner. Under the stronger natural merger law,
Corollary~\ref{cor:stronger-state-local-package-collapse-fixed-shear-corner}
shows that the interaction remainder dies identically along the involutive
branch. Thus any still-live mixed two-point obstruction on
\[
\{S=0\}\cap\{K_v=0\}
\]
must come from genuinely relational or provenance-sensitive two-atom data
beyond the descended transport state.
\end{remark}

\begin{remark}[Conditional parity reduction on the exact fixed-shear $K_v=0$ slice]
\label{rem:conditional-parity-reduction-fixed-shear-Kv-zero}
In the additional conditional exact fixed-shear overlap lane, the residual
$K_v=0$ slice should be read as quotient-diagonal / involutive rather than as a
new mixed local branch-classification problem. Indeed, after centering and
analytic source reparametrization, the projected germ has the form
\[
\Phi(t)=(t^m,G(t)),
\]
and exact repeated values force
\[
m\ \text{even},
\qquad
G(-t)=G(t).
\]
So, after quotienting by analytic functions of the first coordinate, the
genuinely mixed residue is an even/even quotient germ
\[
(t^m,\ t^n+\cdots),
\qquad
m,n\ \text{even},\quad n>m.
\]
In particular, among the current explicit low-order mixed models
\[
(m,n)=(3,4),\ (3,6),\ (4,6),
\]
only the quartic--sextic case
\[
(4,6)
\]
can survive on the exact fixed-shear overlap $K_v=0$ slice. The subsequent
quotient-transport reformulation will show that this is the first rung of an
even-support residual ladder
\[
(m,n)=(r+2,\ r+2e+4),
\qquad
r\in 2\mathbf Z_{\ge 0},
\qquad
e=\operatorname{ord}_0\widehat{\mathcal D}.
\]
\end{remark}

\begin{lemma}[Local exact repeated-value locus on the involutive fixed-shear slice]
\label{lem:local-exact-repeated-value-locus-involutive-fixed-shear}
Assume the setup of
Remark~\ref{rem:conditional-parity-reduction-fixed-shear-Kv-zero}, so that in a
centered analytic source coordinate
\[
\Phi(t)=(t^m,G(t)),
\qquad
m\ \text{even},
\qquad
G(-t)=G(t).
\]
After shrinking to a sufficiently small neighborhood of $t=0$, one has
\[
\Phi(t_1)=\Phi(t_2),\qquad t_1\neq t_2
\quad\Longleftrightarrow\quad
t_2=-t_1.
\]
Equivalently, the local exact repeated-value locus is exactly the involution
graph
\[
\{(t,-t):0<|t|<\varepsilon\}.
\]
In particular, on the exact fixed-shear $K_v=0$ slice there is no local hidden
partner choice beyond the involution $t\leftrightarrow -t$.
\end{lemma}

\begin{proof}
If $\Phi(t_1)=\Phi(t_2)$, then comparing first coordinates gives
\[
t_1^m=t_2^m.
\]
Since $m$ is even and the source coordinate is real, this implies
\[
|t_1|=|t_2|.
\]
If $t_1\neq t_2$, then necessarily $t_2=-t_1$. Conversely, because $m$ is even
and $G$ is even,
\[
\Phi(-t)=\Phi(t)
\]
for all sufficiently small $t$. This proves the claim.
\end{proof}

\begin{proposition}[Quotient transport scalar and residual ladder on the exact fixed-shear $K_v=0$ slice]
\label{prop:quotient-transport-scalar-residual-ladder-fixed-shear-Kv-zero}
Assume the additional conditional exact fixed-shear overlap lane from
Remark~\ref{rem:conditional-parity-reduction-fixed-shear-Kv-zero}. Let
\[
r:=\operatorname{ord}_0J_\lambda=\operatorname{ord}_0K_v,
\qquad
q:=\operatorname{ord}_0J_{u_5},
\]
and write in the centered involutive source coordinate $t$
\[
J_\lambda(t)=t^rA(t),
\qquad
J_{u_5}(t)=t^rB(t),
\qquad
A(0)\neq 0.
\]
Set
\[
Q(t):=\frac{B(t)}{A(t)}=\frac{J_{u_5}(t)}{J_\lambda(t)}.
\]
Then the following hold.

\begin{enumerate}
\item $Q$ is even:
\[
Q(-t)=Q(t).
\]
Hence there exists a unique analytic germ $\widehat Q$ in
\[
u:=t^2
\]
such that
\[
Q(t)=\widehat Q(t^2).
\]

\item If
\[
\widehat{\mathcal D}(u):=\widehat Q'(u),
\]
then
\[
\frac{\mathcal W_{\mathrm{src}}(t)}{J_\lambda(t)^2}
=
Q'(t)
=
2t\,\widehat{\mathcal D}(t^2),
\]
and therefore
\[
\mathcal W_{\mathrm{src}}(t)
=
2\,t^{2r+1}A(t)^2\,\widehat{\mathcal D}(t^2).
\]

\item If
\[
e:=\operatorname{ord}_0\widehat{\mathcal D}\in\mathbf Z_{\ge 0},
\]
then
\[
q-r=2e+2,
\qquad
m=r+2,
\qquad
n=r+2e+4,
\]
and
\[
\operatorname{ord}_0\mathcal W_{\mathrm{src}}=2r+1+2e.
\]

\item In particular,
\[
r\ \text{is even},
\qquad
q-r\in 2\mathbf Z_{\ge 2},
\qquad
\operatorname{ord}_0\mathcal W_{\mathrm{src}}\ \text{is odd}.
\]
So the exact fixed-shear residual lane lies on an even-support ladder over
$K_v=0$, while the associated source-level Wronskian vanishes to odd order.

\item The currently surviving quartic--sextic residual model is exactly the
first rung of this ladder:
\[
(m,n)=(4,6)
\iff
(r,e)=(2,0).
\]
\end{enumerate}
\end{proposition}

\begin{proof}
By Remark~\ref{rem:conditional-parity-reduction-fixed-shear-Kv-zero}, the
projected germ has the form
\[
\Phi(t)=(t^m,G(t)),
\qquad
m\ \text{even},
\qquad
G(-t)=G(t).
\]
Thus the quotient-diagonal mixed residue is even/even, so the regularized ratio
$Q=B/A$ is even, giving the factorization
\[
Q(t)=\widehat Q(t^2).
\]
Differentiating yields
\[
Q'(t)=2t\,\widehat Q'(t^2)=2t\,\widehat{\mathcal D}(t^2).
\]
On the other hand,
Remark~\ref{rem:regularized-source-scalar-mixed-residual} gives
\[
Q'(t)=\frac{\mathcal W_{\mathrm{src}}(t)}{J_\lambda(t)^2}.
\]
Since
\[
J_\lambda(t)^2=t^{2r}A(t)^2,
\]
this gives
\[
\mathcal W_{\mathrm{src}}(t)
=
2\,t^{2r+1}A(t)^2\,\widehat{\mathcal D}(t^2).
\]
If $e=\operatorname{ord}_0\widehat{\mathcal D}$, then
\[
\widehat Q(u)-\widehat Q(0)\sim u^{e+1},
\]
so
\[
Q(t)-Q(0)\sim t^{2e+2}.
\]
Hence
\[
q-r=\operatorname{ord}_0(Q(t)-Q(0))=2e+2.
\]
Now
Remark~\ref{rem:wronskian-order-reformulation-mixed-list} gives
\[
m=r+2,
\qquad
n=q+2,
\]
so
\[
n=r+2e+4.
\]
Also
\[
\operatorname{ord}_0\mathcal W_{\mathrm{src}}=2r+1+2e.
\]
Since $m$ is even, $r=m-2$ is even. The final statements follow immediately.
\end{proof}

\begin{proposition}[Sign-run compression and the three-run criterion on a quotient-visible exact fixed-shear branch]
\label{prop:sign-run-compression-three-lobe-criterion-fixed-shear-quotient-branch}
Assume the quotient-visible exact fixed-shear reduction, and let
\[
[a,b]\subset \mathbf R_u
\]
be a connected compact quotient branch on which the descended transport
potential
\[
\widehat Q:[a,b]\to\mathbf R
\]
is analytic and nonconstant, with
\[
\widehat{\mathcal D}=\widehat Q'.
\]
Let
\[
a=s_0<s_1<\cdots<s_\ell=b
\]
be the sign-change partition of $\widehat{\mathcal D}$: on each run
\[
R_j:=(s_j,s_{j+1}),
\qquad
j=0,\dots,\ell-1,
\]
the derivative $\widehat{\mathcal D}$ has fixed nonzero sign, and the sign
alternates from run to run. Define
\[
d_j:=\widehat Q(s_j),
\qquad
\Omega_j:=\bigl(\min\{d_j,d_{j+1}\},\max\{d_j,d_{j+1}\}\bigr),
\]
and the run amplitudes
\[
B_j:=|d_{j+1}-d_j|
=
\int_{R_j}|\widehat{\mathcal D}(u)|\,du.
\]
Then the following hold.

\begin{enumerate}
\item For every regular value
\[
\gamma\notin \widehat Q\bigl(\{\widehat{\mathcal D}=0\}\bigr),
\]
one has the exact multiplicity formula
\[
\#\widehat Q^{-1}(\gamma)=\#\{\,j:\gamma\in\Omega_j\,\}.
\]

\item For every
\[
0\le j\le \ell-3,
\]
one has the exact distance-two overlap criterion
\[
\Omega_j\cap\Omega_{j+2}\neq\varnothing
\iff
B_j+B_{j+2}>B_{j+1}.
\]
Equivalently,
\[
\Omega_j\cap\Omega_{j+2}=\varnothing
\iff
B_{j+1}\ge B_j+B_{j+2}.
\]

\item Two consecutive distance-two separations are impossible. More precisely,
there is no $j$ such that both
\[
\Omega_j\cap\Omega_{j+2}=\varnothing
\qquad\text{and}\qquad
\Omega_{j+1}\cap\Omega_{j+3}=\varnothing.
\]
In particular, if the branch has no nonadjacent overlap of regular values, then
\[
\ell\le 3.
\]
\end{enumerate}
\end{proposition}

\begin{proof}
On each run $R_j$, the function $\widehat Q$ is strictly monotone, so for a
regular value $\gamma$ it has exactly one preimage on $R_j$ if and only if
$\gamma$ lies in the open image interval $\Omega_j$, and otherwise no preimage
there. Summing over runs proves (1).

For (2), fix $j$ and, after replacing $\widehat Q$ by $-\widehat Q$ if needed,
assume that the signs on $R_j,R_{j+1},R_{j+2}$ are
\[
+,\ -, \ +.
\]
Then
\[
d_j<d_{j+1}>d_{j+2}<d_{j+3},
\]
with
\[
B_j=d_{j+1}-d_j,
\qquad
B_{j+1}=d_{j+1}-d_{j+2},
\qquad
B_{j+2}=d_{j+3}-d_{j+2}.
\]
In this configuration,
\[
\Omega_j=(d_j,d_{j+1}),
\qquad
\Omega_{j+2}=(d_{j+2},d_{j+3}),
\]
and since $d_{j+2}<d_{j+1}$ automatically, these open intervals overlap if and
only if
\[
d_j<d_{j+3}.
\]
But
\[
d_{j+3}-d_j
=
(d_{j+1}-d_j)-(d_{j+1}-d_{j+2})+(d_{j+3}-d_{j+2})
=
B_j-B_{j+1}+B_{j+2},
\]
so
\[
d_j<d_{j+3}
\iff
B_j+B_{j+2}>B_{j+1}.
\]
This proves (2).

For (3), if both consecutive separations held, then by (2) one would have
\[
B_{j+1}\ge B_j+B_{j+2},
\qquad
B_{j+2}\ge B_{j+1}+B_{j+3},
\]
which is impossible because all $B_k>0$. Hence among any two consecutive
distance-two comparisons, at least one overlap is forced. If $\ell\ge 4$, such
a consecutive comparison exists, so no-nonadjacent-overlap forces $\ell\le 3$.
\end{proof}

\begin{corollary}[Three-run classification on a quotient-visible exact fixed-shear branch]
\label{cor:three-run-no-reentry-classification-fixed-shear-quotient-branch}
In the setting of
Proposition~\ref{prop:sign-run-compression-three-lobe-criterion-fixed-shear-quotient-branch},
assume that
\[
\ell=3.
\]
Then the branch has no nonadjacent overlap of regular values if and only if
\[
B_1\ge B_0+B_2.
\]
If
\[
B_1>B_0+B_2,
\]
the middle run separates the two outer windows by a genuine gap. If
\[
B_1=B_0+B_2,
\]
the outer windows touch only at a critical value, so there is still no regular
value with nonadjacent re-entry.
\end{corollary}

\begin{proof}
When $\ell=3$, the only nonadjacent comparison is $(\Omega_0,\Omega_2)$, so the
claim is exactly Proposition~\ref{prop:sign-run-compression-three-lobe-criterion-fixed-shear-quotient-branch}(2).
\end{proof}

\begin{remark}[Touching is sterile in the three-run quotient model]
\label{rem:touching-is-sterile-three-run-quotient-model}
In the equality case
\[
B_1=B_0+B_2,
\]
Corollary~\ref{cor:three-run-no-reentry-classification-fixed-shear-quotient-branch}
says that the outer windows meet only at a boundary critical value. Thus the
same-value relation between the two outer runs is zero-dimensional downstairs in
the quotient variable $u=t^2$: there is no open value interval carrying a
same-value branch between the outer runs.

Accordingly, the equality case should be viewed as critical-value touching, not
as a genuine open transport overlap. The genuinely live open-overlap three-run
regime is
\[
B_0+B_2>B_1,
\]
while
\[
B_1>B_0+B_2
\]
is the separation regime.
\end{remark}

\begin{remark}[Separation is not a transport lane in the three-run quotient model]
\label{rem:separation-not-transport-lane-three-run-quotient-model}
In the strict separation regime
\[
B_1>B_0+B_2,
\]
the outer windows are disjoint, so there is no outer-overlap interval
\[
\Omega_0\cap\Omega_2
\]
and hence no inverse-branch-gap transport problem between the two outer runs.
Accordingly, the inverse-branch-gap / $\Delta$ / $\Xi$ setup belongs only to
the genuine outer-overlap regime
\[
B_0+B_2>B_1,
\]
while the strict separation regime is already a no-reentry regime downstairs in
the quotient variable.
\end{remark}

\begin{corollary}[Quartic--sextic quotient injectivity and local no-reentry]
\label{cor:quartic-sextic-quotient-injectivity-local-no-reentry}
Assume the quartic--sextic residual case
\[
(m,n)=(4,6).
\]
Then
\[
Q(t)=Q(0)+\kappa_{2,2}t^2+O(t^4),
\qquad
\widehat Q(u)=Q(0)+\kappa_{2,2}u+O(u^2),
\]
so
\[
\widehat{\mathcal D}(0)=\kappa_{2,2}\neq 0.
\]
Equivalently,
\[
\kappa_{2,2}
=
\frac12\left(\frac{\mathcal W_{\mathrm{src}}}{J_\lambda^2}\right)'(0)
=
\frac{\mathcal W_{\mathrm{src}}^{(5)}(0)}{60\,(J_\lambda''(0))^2}.
\]

Hence $\widehat Q$ is a local analytic diffeomorphism in the quotient variable
$u=t^2$. In particular, after quotienting the involution
$t\leftrightarrow -t$, there is no local exact fixed-shear re-entry: every
nearby fixed-shear level is attained by at most one quotient point.
\end{corollary}

\begin{proof}
In the quartic--sextic case, Proposition~\ref{prop:quotient-transport-scalar-residual-ladder-fixed-shear-Kv-zero}
gives
\[
(r,e)=(2,0).
\]
So $\widehat{\mathcal D}(0)\neq 0$, which is exactly the displayed coefficient
$\kappa_{2,2}$. The explicit formulas follow from
Remark~\ref{rem:wronskian-order-reformulation-mixed-list} and the identity
\[
Q'(t)=\frac{\mathcal W_{\mathrm{src}}}{J_\lambda^2}.
\]
Since
\[
\widehat Q'(0)=\widehat{\mathcal D}(0)\neq 0,
\]
the analytic inverse function theorem implies that $\widehat Q$ is a local
diffeomorphism in $u=t^2$.
\end{proof}

\begin{remark}[The quartic--sextic rung is inert for the lobe-amplitude obstruction]
\label{rem:quartic-sextic-inert-for-lobe-amplitude-obstruction}
On the first residual rung
\[
(m,n)=(4,6),
\]
Proposition~\ref{prop:quotient-transport-scalar-residual-ladder-fixed-shear-Kv-zero}
gives
\[
(r,e)=(2,0),
\]
and
Corollary~\ref{cor:quartic-sextic-quotient-injectivity-local-no-reentry}
gives
\[
\widehat{\mathcal D}(0)=\kappa_{2,2}\neq 0.
\]
So the quartic--sextic basepoint is not a sign-change point of
$\widehat{\mathcal D}$, contributes no compressed node on the quotient branch,
and does not participate in any three-run run-amplitude obstruction. The first
possible
quotient-visible bad core therefore begins only at a higher-rung critical point
of $\widehat{\mathcal D}$.
\end{remark}

\begin{proposition}[Finite transport-jet quotient descent on the exact fixed-shear corner]
\label{prop:finite-transport-jet-quotient-descent-fixed-shear-corner}
Assume the additional conditional exact fixed-shear overlap lane from
Remark~\ref{rem:conditional-parity-reduction-fixed-shear-Kv-zero}, and work
along a local involutive branch
\[
(h_1,h_2)=(h(t),h(-t)).
\]
For each $k\ge 0$, define the parity-normalized transport jet
\[
\Theta_k(t):=
\begin{cases}
Q^{(k)}(t),& k\ \text{even},\\[4pt]
t^{-1}Q^{(k)}(t),& k\ \text{odd},
\end{cases}
\]
where
\[
Q(t)=\frac{J_{u_5}(t)}{J_\lambda(t)}=\widehat Q(t^2).
\]
Then every $\Theta_k$ is even analytic. Equivalently, for each $k\ge 0$
there exists a unique analytic quotient germ $\widetilde\Theta_k(u)$ such that
\[
\Theta_k(t)=\widetilde\Theta_k(t^2).
\]
In particular, the canonical descended scalar
\[
X(t):=\frac{Q'(t)}{t}
\]
is even analytic and satisfies
\[
X(t)=2\,\widehat Q'(t^2)=2\,\widehat{\mathcal D}(t^2).
\]
\end{proposition}

\begin{proof}
Since $Q(t)=\widehat Q(t^2)$ is even, $Q^{(k)}$ has the usual parity:
$Q^{(k)}$ is even for $k$ even and odd for $k$ odd. Every odd analytic
germ is divisible by $t$, so $\Theta_k$ is analytic for odd $k$. In both
cases $\Theta_k$ is even, hence factors uniquely through $u=t^2$. The final
identity is immediate from
\[
Q'(t)=2t\,\widehat Q'(t^2)=2t\,\widehat{\mathcal D}(t^2).
\]
\end{proof}

\begin{corollary}[Rank-one quotient-visible transport state on the quartic--sextic rung]
\label{cor:rank-one-transport-state-quartic-sextic-fixed-shear}
Assume the quartic--sextic residual case
\[
(m,n)=(4,6).
\]
For every $N\ge 0$, set
\[
\mathcal T_N(t):=(\Theta_0(t),\dots,\Theta_N(t)).
\]
Then there exists an analytic map
\[
\mathbf F_N
\]
defined near $Q(0)$ such that
\[
\mathcal T_N(t)=\mathbf F_N(Q(t)).
\]
Equivalently, every finite quotient-visible transport state is locally an
analytic function of the single scalar $Q$ alone.
\end{corollary}

\begin{proof}
By
Corollary~\ref{cor:quartic-sextic-quotient-injectivity-local-no-reentry},
$\widehat Q$ is a local analytic coordinate in $u=t^2$. Since each
$\Theta_k(t)=\widetilde\Theta_k(t^2)$, each descended jet
$\widetilde\Theta_k(u)$ can be re-expressed analytically as a function of
$\widehat Q(u)$, hence of $Q(t)$. Stacking these functions gives
$\mathbf F_N$.
\end{proof}

\begin{corollary}[No finite-order hidden reset on the quartic--sextic residual corner]
\label{cor:no-finite-order-hidden-reset-quartic-sextic-fixed-shear}
Assume the quartic--sextic residual case. Let
\[
\mathcal T_N(t):=(\Theta_0(t),\dots,\Theta_N(t))
\]
be any finite quotient-visible transport state. Then
\[
\mathcal T_N(t_1)=\mathcal T_N(t_2)
\qquad\Longrightarrow\qquad
t_2=\pm t_1.
\]
In particular, no finite quotient-visible transport state produces a third local
reset beyond the involution symmetry $t\leftrightarrow -t$.
\end{corollary}

\begin{proof}
By
Corollary~\ref{cor:rank-one-transport-state-quartic-sextic-fixed-shear},
equality of $\mathcal T_N$ implies equality of $Q$. But
\[
Q(t)=\widehat Q(t^2),
\]
and
Corollary~\ref{cor:quartic-sextic-quotient-injectivity-local-no-reentry}
shows that $\widehat Q$ is locally injective in $u=t^2$. Hence
$t_1^2=t_2^2$, so $t_2=\pm t_1$.
\end{proof}

\begin{proposition}[Conditional state-local defect closure on the residual exact fixed-shear corner]
\label{prop:conditional-state-local-defect-closure-fixed-shear-corner}
Assume the setup of
Proposition~\ref{prop:residual-involutive-defect-criterion-fixed-shear-corner}.
Suppose there exists $N\ge 0$ and analytic functionals
\[
\mathfrak e_3,\mathfrak e_5
\]
defined near $\mathcal T_N(0)$ such that along the involutive branch
\[
E_{12}^{(3)}(t)=\mathfrak e_3(\mathcal T_N(t)),
\qquad
E_{12}^{(5)}(t)=\mathfrak e_5(\mathcal T_N(t)),
\]
and
\[
\mathfrak e_3(\mathcal T_N(0))=\mathfrak e_5(\mathcal T_N(0))=0.
\]
Then
\[
E_{12}^{(3)}(t)=O(t^2),
\qquad
E_{12}^{(5)}(t)=O(t^2),
\]
hence
\[
E_{12}^{(3)}=O((h_1-h_2)^2),
\qquad
E_{12}^{(5)}=O((h_1-h_2)^2),
\]
and no such nontrivial exact residual fixed-shear involutive bad branch can
persist on a compact good patch
\[
J\Subset\{c\neq 0,\ M\neq 0\}.
\]
\end{proposition}

\begin{proof}
By Proposition~\ref{prop:finite-transport-jet-quotient-descent-fixed-shear-corner},
each $\Theta_k$ is even analytic, so
\[
\mathcal T_N(t)=\mathcal T_N(0)+O(t^2).
\]
Analyticity of $\mathfrak e_3,\mathfrak e_5$ and vanishing at $\mathcal T_N(0)$
therefore give
\[
E_{12}^{(3)}(t)=O(t^2),
\qquad
E_{12}^{(5)}(t)=O(t^2).
\]
Since
\[
h_1-h_2=h(t)-h(-t)=2h'(0)t+O(t^3),
\]
this becomes
\[
E_{12}^{(3)}=O((h_1-h_2)^2),
\qquad
E_{12}^{(5)}=O((h_1-h_2)^2).
\]
Now apply Corollary~\ref{cor:bad-two-pair-cores-defect-scale-close}.
\end{proof}

\begin{corollary}[Stronger conditional package collapse on the involutive branch]
\label{cor:stronger-state-local-package-collapse-fixed-shear-corner}
Assume the setup of
Proposition~\ref{prop:residual-involutive-defect-criterion-fixed-shear-corner}.
Suppose that for some $N\ge 0$ the corrected two-pair package along the
involutive branch is state-local in the sense that
\[
\mathcal P_2(t)
=
\mathfrak P_2\bigl(a_1(t),a_2(t);\mathcal T_N(t),\mathcal T_N(t)\bigr),
\]
and the corresponding one-pair package is
\[
F_{h_1(t)}=F_{h_2(t)}=\mathfrak F(\mathcal T_N(t)).
\]
Assume moreover the natural merger law
\[
\mathfrak P_2(a_1,a_2;Y,Y)=(a_1+a_2)\mathfrak F(Y).
\]
Then the interaction remainder
\[
\mathcal I_2(t):=
\mathcal P_2(t)-a_1(t)F_{h_1(t)}-a_2(t)F_{h_2(t)}
\]
vanishes identically along the involutive branch:
\[
\mathcal I_2(t)\equiv 0.
\]
\end{corollary}

\begin{proof}
Substitute the displayed state-local forms into the definition of
$\mathcal I_2(t)$, then apply the merger law with $Y=\mathcal T_N(t)$.
\end{proof}

\begin{remark}[Any remaining residual obstruction is genuinely relational]
\label{rem:any-remaining-residual-obstruction-is-genuinely-relational}
The preceding results show that on the quartic--sextic residual corner the
finite quotient-visible transport state has local rank $1$ and no finite-order
hidden reset. Therefore, if the exact fixed-shear residual corner is still
genuinely live, the obstruction cannot come from:
\begin{itemize}
\item a second local transport variable on the quartic--sextic rung;
\item a third local reset inside any finite quotient-visible transport state;
\item the explicit quotient-visible three-run run-amplitude obstruction on the
quartic--sextic rung, since
Corollary~\ref{cor:quartic-sextic-quotient-injectivity-local-no-reentry}
gives $\widehat{\mathcal D}(0)\neq 0$ there;
\item a two-pair model built only from the descended one-atom transport state,
assuming the natural merger law.
\end{itemize}
So any still-live obstruction on
\[
\{S=0\}\cap\{K_v=0\}
\]
must come from the first genuinely relational or provenance-sensitive two-atom
input that survives after quotient-visible transport descent, or else from
non-finite-order dependence.
\end{remark}

\begin{remark}[Quotient-cusp interpretation and the no-reset issue]
\label{rem:quotient-cusp-interpretation-no-reset-Kv-zero}
The involutive quotient picture on the residual exact fixed-shear
$K_v=0$ slice is best read as a one-branched quotient singularity rather than
as a new branching mixed carrier.

Indeed, in the notation of
Remark~\ref{rem:conditional-parity-reduction-fixed-shear-Kv-zero}, one has
\[
\Phi(t)=(t^m,G(t)),
\qquad
m\ \text{even},
\qquad
G(-t)=G(t),
\]
so after passing to the quotient parameter
\[
u=t^2
\]
the projected germ descends to a single analytic quotient branch
\[
\bar\Phi(u)=\bigl(u^{m/2},\bar G(u)\bigr).
\]
In particular, in the quartic--sextic case
\[
(m,n)=(4,6),
\]
the quotient branch is cusp-like:
\[
\bar\Phi(u)\sim (u^2,u^3+\cdots).
\]

Thus any local quantity that genuinely factors through the quotient-diagonal
projected germ --- equivalently through the quotient-visible fixed-shear
transport data --- is single-valued in $u$ and has trivial local monodromy at
the involutive point. In this sense, the residual exact fixed-shear
$K_v=0$ locus is best viewed as a quotient-cusp splice rather than as a new
local source of cycle rank.

At the same time, this does \emph{not} yet by itself prove global no-reentry on
the $K_v=0$ lane: one still has to show that the full transported overlap /
critical-value state relevant to the later finite-core problem is indeed
quotient-visible. So the live mixed residual question is now a
\emph{no-hidden-reset} question at $K_v=0$, not another local branch
classification problem.
\end{remark}

\begin{remark}[What the quotient-transport package does and does not prove]
\label{rem:what-quotient-transport-package-does-and-does-not-prove}
The quotient-transport package now does more than give a single pointwise
scalar compression of the residual exact fixed-shear $K_v=0$ slice.

First,
Proposition~\ref{prop:finite-transport-jet-quotient-descent-fixed-shear-corner}
shows that every finite parity-normalized transport jet descends analytically to
the quotient variable
\[
u=t^2.
\]

Second, on the quartic--sextic rung,
Corollary~\ref{cor:rank-one-transport-state-quartic-sextic-fixed-shear}
shows that every finite quotient-visible transport state is locally an analytic
function of the single scalar
\[
Q(t)=\widehat Q(t^2)
\]
alone, and
Corollary~\ref{cor:no-finite-order-hidden-reset-quartic-sextic-fixed-shear}
rules out any third local reset inside that descended state beyond the
involution symmetry $t\leftrightarrow -t$.

What is \emph{not} proved here is that the actual corrected two-pair defect
package is already state-local in this descended transport state. If that
state-locality held, then
Proposition~\ref{prop:conditional-state-local-defect-closure-fixed-shear-corner}
would already exclude the residual branch, and under the stronger natural merger
law
Corollary~\ref{cor:stronger-state-local-package-collapse-fixed-shear-corner}
would make the interaction remainder vanish identically.

Accordingly, the present residual mixed burden should now be read as follows:
the exact fixed-shear corner is no longer blocked by a missing local branch
classification or by hidden finite-order transport dynamics. Any remaining
obstruction must come from genuinely relational or provenance-sensitive two-atom
structure, or from non-finite-order dependence, not from the descended
transport state alone.
\end{remark}

\begin{remark}[No finite local rescue on the exact fixed-shear corner]
\label{rem:no-finite-local-rescue-exact-fixed-shear-corner}
The current quotient-transport package leaves very little residual finite local
freedom on the exact fixed-shear corner.

First, every finite parity-normalized transport jet descends to the quotient
variable $u=t^2$. Second, on the quartic--sextic rung the descended transport
state has local rank one and no third local reset beyond the involution
$t\leftrightarrow -t$. Third, the explicit three-run run-amplitude obstruction
is already inert there, because on that rung
\[
\widehat{\mathcal D}(0)=\widehat Q'(0)\neq 0.
\]

So if the residual exact fixed-shear corner still survives, it cannot be
explained by a second local transport variable, by a hidden finite-order reset,
or by the explicit lobe-amplitude obstruction on the quartic--sextic rung.
Any remaining obstruction must therefore be genuinely nonlocal on the transport
side, or package/provenance-sensitive on the two-atom side.
\end{remark}

\begin{proposition}[Local scalar repeated-value classification]
\label{prop:local-scalar-repeated-value-classification}
Let $f$ be a real-analytic function on an open neighborhood of
$a\in\mathbb R$ with $f'(a)=0$, and let $m\ge 2$ be the order of vanishing of
$f-f(a)$ at $a$, so
\[
f(u)-f(a)=c_m(u-a)^m+O((u-a)^{m+1}),
\qquad c_m\neq 0.
\]
Then there exist a sign $\sigma\in\{+1,-1\}$ and a local real-analytic chart
$x$ near $a$ with $x(a)=0$ and $x'(a)\neq 0$ such that
\[
f(u)=f(a)+\sigma\,x(u)^m.
\]
The local repeated-value locus
\[
\{(u_1,u_2):u_1\neq u_2,\ f(u_1)=f(u_2)\}
\]
near $(a,a)$ is:
\begin{itemize}
\item \emph{(Even $m$.)} If $m$ is even, it is the involution graph
\[
u_2=\iota(u_1):=x^{-1}(-x(u_1)),
\]
together with $u_1=u_2$. The involution $\iota$ fixes $a$ and is analytic with
$\iota'(a)=-1$; the repeated-value branch is a single local fold half-arc
attached to $a$.
\item \emph{(Odd $m$.)} If $m$ is odd, then $f$ is a local analytic
homeomorphism near $a$ and the nondiagonal repeated-value locus is empty near
$(a,a)$; the point $a$ is sterile.
\end{itemize}

\begin{proof}
By the Malgrange/Morse-type preparation for analytic functions of one variable
with an isolated critical point, one may write
$f(u)-f(a)=\sigma\,h(u)^m$ for a local analytic chart $h$ with $h(a)=0$ and
$h'(a)\neq 0$, where $\sigma=\operatorname{sgn}(c_m)$ when $m$ is even and
$\sigma=+1$ absorbing the sign into $h$ when $m$ is odd. Setting $x:=h$ gives
the stated normal form.

For even $m$, $f(u_1)=f(u_2)$ near $(a,a)$ is equivalent to
$x(u_1)^m=x(u_2)^m$, i.e.\ $x(u_2)=\pm x(u_1)$; the plus sign is the diagonal,
the minus sign is $u_2=x^{-1}(-x(u_1))=\iota(u_1)$. That $\iota$ is involutive
with $\iota(a)=a$ and $\iota'(a)=-1$ is immediate from $x'(a)\neq 0$ and
$x(a)=0$.

For odd $m$, $x(u_1)^m=x(u_2)^m$ forces $x(u_1)=x(u_2)$, hence $u_1=u_2$ near
$a$ since $x$ is a local homeomorphism.
\end{proof}
\end{proposition}

\begin{corollary}[Local scalar-bud decomposition on the close exact fixed-shear lane]
\label{cor:local-scalar-bud-decomposition-close-fixed-shear}
Apply
Proposition~\ref{prop:local-scalar-repeated-value-classification} to
$f=\widehat Q$ on the nonconstant quotient branch in the exact fixed-shear
corner $\{S=0\}\cap\{K_v=0\}$. Let $a$ be a zero of
$\widehat{\mathcal D}=\widehat Q'$ of order $m-1\ge 1$. Then:
\begin{itemize}
\item If $m$ is even (i.e.\ $a$ is an odd-order zero of $\widehat{\mathcal D}$),
the local close exact repeated-value carrier at $a$ is exactly one fold
half-arc, namely the involution graph $\{u_2=\iota(u_1)\}$ of
Proposition~\ref{prop:local-scalar-repeated-value-classification}.
\item If $m$ is odd (i.e.\ $a$ is an even-order zero of $\widehat{\mathcal D}$),
the local close exact repeated-value carrier at $a$ is empty; the basepoint is
sterile.
\end{itemize}
In particular, on a compact nonconstant quotient branch in the exact
fixed-shear corner, the close exact repeated-value carrier is a finite
disjoint union of local fold half-arcs, one per odd zero of
$\widehat{\mathcal D}$, with even zeros sterile. No connected nonlocal close
carrier remains on that branch.

\begin{proof}
By
Lemma~\ref{lem:local-exact-repeated-value-locus-involutive-fixed-shear}, away
from zeros of $\widehat Q'$ the close exact repeated-value locus is empty, so
any close carrier accumulates only at zeros of $\widehat{\mathcal D}$. At each
such zero $a$ the local structure is dictated by
Proposition~\ref{prop:local-scalar-repeated-value-classification} applied to
$\widehat Q$, with $m=\operatorname{ord}_a(\widehat Q-\widehat Q(a))
=\operatorname{ord}_a\widehat{\mathcal D}+1$; the two cases transcribe to odd
vs.\ even order of $\widehat{\mathcal D}$. Compactness then bounds the number
of such basepoints, and the carrier is their finite disjoint union of local
fold half-arcs.
\end{proof}
\end{corollary}

\begin{remark}[Constant-branch scalarization and the scalar ceiling]
\label{rem:constant-branch-scalarization-and-scalar-ceiling}
There is also a constant-quotient branch. If
\[
\widehat{\mathcal D}\equiv 0,
\]
then the quotient identities give
\[
Y_{xx}=\gamma_0 S_{xx},
\qquad
Y=\gamma_0 S+\beta_0 x+\alpha_0.
\]
After affine gauge one may reduce to $Y\equiv 0$, and the exact fixed-shear
overlap equations then force $\gamma_*=0$. So the constant branch reduces to
the repeated-value problem for the single planar scalar
\[
V(t):=U_{\xi_*,\eta_*}(x(t),S(t)),
\qquad
U_{\xi,\eta}(x,S):=\frac{x-\xi S}{1-\eta S},
\]
to which
Proposition~\ref{prop:local-scalar-repeated-value-classification} applies
pointwise in $t$. At basepoints with
$V'(t_0)\neq 0$, $V$ is a local homeomorphism and close exact repeated-value
pairs are excluded;
Corollary~\ref{cor:quartic-sextic-quotient-injectivity-local-no-reentry}
rules out the basepoint case $\widehat{\mathcal D}(0)\neq 0$, and at basepoints
where $V$ has a critical point the local carrier is again a single fold
half-arc. Thus the constant exact fixed-shear branch also contributes only
local fold half-arcs.

In this form the close exact fixed-shear lane contributes no connected
nonlocal close carrier, only local scalar fold half-arcs.
Corollary~\ref{cor:local-scalar-bud-decomposition-close-fixed-shear} is the
theorem-level scalar ceiling on that lane: the remaining intrinsic obstruction
on the close exact fixed-shear side is now equality of consecutive lobe
amplitudes across these local fold half-arcs, not a hidden nonlocal carrier.
\end{remark}

\begin{remark}[Connected-carrier sterility versus scalar multiplicity on the close exact lane]
\label{rem:connected-carrier-sterility-vs-scalar-multiplicity-close-lane}
Corollary~\ref{cor:local-scalar-bud-decomposition-close-fixed-shear} and
Remark~\ref{rem:constant-branch-scalarization-and-scalar-ceiling} do more than
say that close exact repeated values are locally scalar. They also identify the
global graph type of the close carrier: after shrinking the close threshold,
each connected component is a single local fold bud attached at one odd scalar
seed. In particular, each component is one-sided and pendant, never a cycle
edge, and there is no connected nonlocal close carrier on the exact fixed-shear
lane.

What this closes is the connected-carrier question, not abstract scalar
multiplicity. Pure scalar analyticity can still realize multiplicity-$4$
repeated-value levels, equivalently equal consecutive lobe amplitudes, so no
truthful purely scalar theorem can bound the close multiplicity by $2$ in full
generality. The point is instead that such multiplicity does not organize into a
connected global close carrier on the exact fixed-shear lane.
\end{remark}

\begin{remark}[Work in progress: inverse-branch-gap endgame in the overlap regime on the strict three-run lane]
\label{rem:wip-inverse-branch-gap-endgame-fixed-shear}
Stay in the setting of
Proposition~\ref{prop:sign-run-compression-three-lobe-criterion-fixed-shear-quotient-branch}
and
Corollary~\ref{cor:three-run-no-reentry-classification-fixed-shear-quotient-branch},
and assume the three-run overlap regime
\[
B_0+B_2>B_1.
\]
Then the outer windows overlap on an open value interval
\[
\mathcal I:=\Omega_0\cap\Omega_2,
\]
so there are inverse branches
\[
\phi_0:\mathcal I\to R_0,
\qquad
\phi_2:\mathcal I\to R_2
\]
of $\widehat Q$. Define the inverse-branch gap
\[
\Delta(q):=\phi_2(q)-\phi_0(q),
\]
and the transport-curvature scalar
\[
\Xi(u):=\frac{\widehat Q''(u)}{(\widehat Q'(u))^3},
\qquad
\Upsilon(q):=\Xi(\phi_0(q))-\Xi(\phi_2(q)).
\]
Then
\[
\Delta'(q)=\frac{1}{\widehat Q'(\phi_2(q))}
-\frac{1}{\widehat Q'(\phi_0(q))},
\qquad
\Delta''(q)=\Upsilon(q).
\]

So the strict three-run overlap problem may be studied through the
one-variable gap function $\Delta$. Any interior critical point
$q_*\in\mathcal I$ yields a same-value / same-slope coincidence: with
\[
u_0:=\phi_0(q_*),\qquad u_2:=\phi_2(q_*),
\]
one has
\[
\widehat Q(u_0)=\widehat Q(u_2),
\qquad
\widehat Q'(u_0)=\widehat Q'(u_2).
\]
Thus the overlap obstruction is already a concrete outer-pair witness, not
merely an abstract transport mismatch.

It is natural to package the higher mismatch by
\[
J_1(u):=\frac{1}{\widehat Q'(u)},
\qquad
J_{n+1}(u):=J_1(u)\,\partial_u J_n(u).
\]
The sharp current expectation is that any genuine overlap witness should force a
finite-order even jet mismatch: for some smallest $m\ge 1$,
\[
J_2(u_0)=J_2(u_2),\ \dots,\ J_{2m-1}(u_0)=J_{2m-1}(u_2),
\qquad
J_{2m}(u_0)\neq J_{2m}(u_2),
\]
with sign corresponding to a local minimum of $\Delta$.

There is also a sharper descent fork, now correctly understood as conditional on
actual overlap. If $\Xi$ descends to value space across the two outer runs,
then $\Upsilon\equiv 0$, so $\Delta$ is affine. Any theorem excluding
nonconstant affine same-value transport between the two outer runs would then
rule out the overlap regime. A stronger rigid alternative, if one can justify
the endpoint propagation, would be exact translated replication
\[
\Omega_0=\Omega_2,
\qquad
B_0=B_2,
\qquad
\phi_2(q)-\phi_0(q)\equiv C>0,
\qquad
\widehat Q(u+C)=\widehat Q(u)\qquad (u\in R_0).
\]

So the genuine three-run \emph{overlap} regime now has the following corrected
endgame fork: either $\Xi$-descent fails on an actual outer-overlap interval,
giving the first explicit nonlocal transport obstruction, or $\Xi$-descent
holds and the problem rigidifies to affine, and potentially translated, outer
transport. In particular, the inverse-branch-gap setup belongs only to the
overlap regime
\[
B_0+B_2>B_1,
\]
whereas
\[
B_1>B_0+B_2
\]
is already the separation regime and carries no such outer-transport lane.

Equivalently, the quantitative overlap scalar is
\[
B_0+B_2-B_1,
\]
and one may write
\[
\Delta'(q)=\Delta'(q_0)+\int_{q_0}^q \Upsilon(s)\,ds,
\]
so a transport-mass bound too small to cancel the initial inverse-slope gap
would exclude genuine overlap.

This is a sharpened endgame formulation, not yet a closure theorem.
\end{remark}

\begin{remark}[No-go: overlap geometry alone does not force a witness]
\label{rem:no-go-overlap-geometry-alone-does-not-force-witness}
The corrected strict-overlap transport lane should not be read as if open outer
overlap by itself automatically forces an interior same-value / same-slope
witness in every endpoint configuration. Outside the stricter endpoint-sign-
forcing subcases, the endpoint behavior of
\[
W(q):=\Delta'(q)
=
\frac{1}{\widehat Q'(\phi_2(q))}
-
\frac{1}{\widehat Q'(\phi_0(q))}
\]
depends on which outer branch is singular or regular at each overlap endpoint.
So overlap geometry alone does not determine witness existence in full
generality.

Equivalently, outside the strict sign-forcing cases, one needs extra input:
endpoint slope comparison, a branch-independent closure law, or genuinely
source/package-level structure. Thus the corrected overlap lane has reached a
real ceiling for pure overlap geometry.
\end{remark}

\begin{remark}[Conditional rigidity: $H$-descent forces affine replication]
\label{rem:conditional-rigidity-H-descent-forces-affine-replication}
Still in the corrected overlap regime, let
\[
H(u):=\frac{\widehat Q''(u)}{(\widehat Q'(u))^2},
\qquad
J_1(u):=\frac{1}{\widehat Q'(u)}.
\]
If $H$ descends to value space across the two outer inverse branches, in the
sense that
\[
H(\phi_0(q))=H(\phi_2(q))
\]
on the overlap interval, then the reciprocal-slope branches
\[
y_i(q):=J_1(\phi_i(q))
\]
satisfy the same linear first-order ODE
\[
y'=-h(q)\,y.
\]
Hence
\[
\frac{y_2(q)}{y_0(q)}=\Lambda
\]
is constant, so
\[
J_1(\phi_2(q))=\Lambda\,J_1(\phi_0(q)),
\qquad
\phi_2'(q)=\Lambda\,\phi_0'(q),
\]
and therefore
\[
\phi_2(q)=\Lambda\,\phi_0(q)+B.
\]

So $H$-descent does not merely give a soft transport coincidence: it forces the
outer inverse branches into affine replication. In particular, the translated /
affine-copy exceptional branch is a genuinely rigid symmetry branch, not a soft
hidden-state loophole.
\end{remark}

\begin{proposition}[Residual involutive defect criterion on the exact fixed-shear corner]
\label{prop:residual-involutive-defect-criterion-fixed-shear-corner}
Assume the additional conditional exact fixed-shear overlap lane near a residual
accumulation point
\[
a\in \{S=0\}\cap\{K_v=0\}.
\]
Let a local exact repeated-value branch be parametrized by
\[
(h_1,h_2)=(h(t),h(-t)),
\qquad
h(0)=a,
\qquad
h'(0)\neq 0,
\]
with $h$ analytic. Let
\[
E_{12}^{(3)}(t),
\qquad
E_{12}^{(5)}(t)
\]
denote the cubic and quintic normalized interaction defects along this branch.

Assume:
\[
E_{12}^{(3)}(t)=E_{12}^{(3)}(-t),
\qquad
E_{12}^{(5)}(t)=E_{12}^{(5)}(-t),
\]
and
\[
E_{12}^{(3)}(0)=0,
\qquad
E_{12}^{(5)}(0)=0.
\]

Then
\[
E_{12}^{(3)}(t)=O(t^2),
\qquad
E_{12}^{(5)}(t)=O(t^2).
\]
Since
\[
h_1-h_2=h(t)-h(-t)=2h'(0)t+O(t^3),
\]
it follows that
\[
E_{12}^{(3)}=O((h_1-h_2)^2),
\qquad
E_{12}^{(5)}=O((h_1-h_2)^2).
\]

Consequently no such nontrivial exact residual fixed-shear involutive bad
branch can persist on a compact good patch
\[
J\Subset\{c\neq 0,\ M\neq 0\}.
\]
\end{proposition}

\begin{proof}
By evenness and vanishing at $t=0$, the Taylor expansions have no constant or
odd-order terms:
\[
E_{12}^{(3)}(t)=\alpha_2 t^2+O(t^4),
\qquad
E_{12}^{(5)}(t)=\beta_2 t^2+O(t^4).
\]
So
\[
E_{12}^{(3)}(t)=O(t^2),
\qquad
E_{12}^{(5)}(t)=O(t^2).
\]
Since
\[
h(t)-h(-t)=2h'(0)t+O(t^3),
\]
we have
\[
t=O(h_1-h_2),
\]
and therefore
\[
E_{12}^{(3)}=O((h_1-h_2)^2),
\qquad
E_{12}^{(5)}=O((h_1-h_2)^2).
\]

Now apply Corollary~\ref{cor:bad-two-pair-cores-defect-scale-close} on a
compact good patch
\[
J\Subset\{c\neq 0,\ M\neq 0\}.
\]
It gives
\[
|h_1-h_2|
\le
C_J\Bigl(|E_{12}^{(3)}|+\|E_{12}^{(5)}\|\Bigr)
=
O((h_1-h_2)^2).
\]
For sufficiently small nonzero $|h_1-h_2|$ this is impossible. Hence no such
nontrivial exact residual fixed-shear involutive bad branch can persist.
\end{proof}

\begin{corollary}[Package-level reduction on the residual exact fixed-shear corner]
\label{cor:package-level-reduction-residual-fixed-shear-corner}
In the setting of
Proposition~\ref{prop:residual-involutive-defect-criterion-fixed-shear-corner},
it is enough to verify for the actual corrected cubic and quintic two-pair
defects that:

\begin{enumerate}
\item they are invariant under swapping the two atoms along the involutive
parametrization $t\leftrightarrow -t$; and

\item they vanish on coincident atoms.
\end{enumerate}

Under these two package-level inputs, the exact residual fixed-shear corner
\[
\{S=0\}\cap\{K_v=0\}
\]
is excluded on compact good overlap patches.
\end{corollary}

\begin{proof}
The two hypotheses are exactly the evenness and diagonal-vanishing inputs of
Proposition~\ref{prop:residual-involutive-defect-criterion-fixed-shear-corner}.
\end{proof}

\begin{remark}[The residual exact fixed-shear corner is now a package-level problem]
\label{rem:residual-fixed-shear-corner-now-package-level}
The point of
Proposition~\ref{prop:residual-involutive-defect-criterion-fixed-shear-corner}
is that it does not require construction of the full corrected two-atom quotient
map. Once the residual exact fixed-shear branch is recognized as
quotient-diagonal / involutive, the remaining burden is no longer a new local
mixed branch classification. It is the package-level collapse of the actual
finite-order two-pair defects on coincident atoms, or equivalently a
no-hidden-reset statement for whatever transported overlap state survives on the
quotient side.
\end{remark}

\noindent\emph{Frame convention for the scalar block.}
On the cubic-generic one-pair locus we fix the cubic-normalized frame so that
the distinguished cubic line is the diagonal-traceless line $\mathbf C D$,
with
\[
D=\begin{pmatrix}1&0\\0&-1\end{pmatrix},
\qquad
J_0DJ_0=-D.
\]
Thus the scalar coefficient extracted below is not an auxiliary choice: it is
the canonical coefficient of the cubic line in this fixed cubic-generic frame.

\begin{definition}[Canonical scalar extraction from the raw one-pair stack]
\label{def:canonical-scalar-extraction-raw-one-pair}
Let
\[
R(z)=A_3 z^3 + A_5 z^5 + A_7 z^7 + O(z^9)
\]
be a cubic-generic raw one-pair corrected germ in the cubic-normalized frame.
Write
\[
I=\begin{pmatrix}1&0\\0&1\end{pmatrix},\qquad
D=\begin{pmatrix}1&0\\0&-1\end{pmatrix},\qquad
S=\begin{pmatrix}0&1\\1&0\end{pmatrix},\qquad
K=\begin{pmatrix}0&1\\-1&0\end{pmatrix}.
\]
We use the Frobenius pairing
\[
\langle X,Y\rangle_F:=\operatorname{tr}(X^\ast Y).
\]
The scalar channel of the cubic coefficient is defined by projection to the
diagonal-traceless line $\mathbf C D$:
\[
c(R):=\frac{\langle A_3,D\rangle_F}{\langle D,D\rangle_F}
      =\frac12\,\operatorname{tr}(D A_3).
\]
Equivalently, the scalar normalization equations are
\[
\pi_I(A_3)=0,\qquad
\pi_S(A_3)=0,\qquad
\pi_K(A_3)=0,\qquad
A_3=c(R)\,D,
\]
where $\pi_I,\pi_S,\pi_K$ denote the Frobenius projections to the lines
$\mathbf C I,\mathbf C S,\mathbf C K$.
\end{definition}

\begin{remark}[Compatibility with the existing gauge laws]
\label{rem:scalar-extraction-compatible-with-gauge-laws}
By the already-proved gauge law
\[
A_3^{\new}=A_3,
\]
the cubic coefficient is unchanged by the admissible one-pair normalization
gauge. Hence the scalar extraction in
Definition~\ref{def:canonical-scalar-extraction-raw-one-pair} is gauge
independent on the raw one-pair stack.
\end{remark}

\begin{lemma}[Scalar normalization is $\Sigma$-invariant]
\label{lem:scalar-normalization-sigma-invariant}
Let
\[
\Sigma(F)(z):=J_0\,F(-z)\,J_0
\]
be the centered reflection involution, and let
\[
\Sigma_3(A):=-J_0 A J_0
\]
be its induced action on cubic coefficients. Assume
\[
J_0 D J_0=-D.
\]
Then the scalar extraction
\[
c(R)=\frac12\,\operatorname{tr}(D A_3)
\]
from Definition~\ref{def:canonical-scalar-extraction-raw-one-pair} is
$\Sigma$-invariant:
\[
c(\Sigma R)=c(R).
\]
\end{lemma}

\begin{proof}
If $R(z)=A_3 z^3+O(z^5)$, then
\[
\Sigma(R)(z)=J_0R(-z)J_0
          =(-J_0A_3J_0)z^3+O(z^5)
          =\Sigma_3(A_3)z^3+O(z^5).
\]
Therefore
\[
c(\Sigma R)
=
\frac12\,\operatorname{tr}\!\bigl(D(-J_0A_3J_0)\bigr).
\]
Using cyclicity of trace,
\[
\operatorname{tr}\!\bigl(D(-J_0A_3J_0)\bigr)
=
\operatorname{tr}\!\bigl((-J_0DJ_0)A_3\bigr).
\]
Since $J_0DJ_0=-D$, we get
\[
-J_0DJ_0=D,
\]
hence
\[
c(\Sigma R)=\frac12\,\operatorname{tr}(D A_3)=c(R).
\]
\end{proof}

\begin{lemma}[Uniqueness of the scalar output]
\label{lem:scalar-output-unique}
Let $R$ be a cubic-generic raw one-pair corrected germ. Then the scalar
$c(R)$ from
Definition~\ref{def:canonical-scalar-extraction-raw-one-pair} is uniquely
determined. In particular, if
\[
A_3=cD=\widetilde c\,D,
\]
then $c=\widetilde c$.
\end{lemma}

\begin{proof}
If
\[
(c-\widetilde c)D=0,
\]
then $c-\widetilde c=0$ because $D\neq 0$. Equivalently,
\[
c=\frac12\,\operatorname{tr}(D A_3)
\]
is the unique Frobenius coefficient of $A_3$ along the line $\mathbf C D$.
\end{proof}

\begin{corollary}[Finite-jet scalar descent on the centered branch]
\label{cor:finite-jet-scalar-descent-centered-branch}
Let $R_s$ be the centered one-pair branch, and assume
\[
R_{-s}=\Sigma(R_s).
\]
Then the scalar coefficient
\[
c(s):=c(R_s)
\]
is even:
\[
c(-s)=c(s).
\]
In particular,
\[
C_1=C_3=C_5=0.
\]
\end{corollary}

\begin{proof}
By Lemma~\ref{lem:scalar-normalization-sigma-invariant},
\[
c(R_{-s})=c(\Sigma R_s)=c(R_s).
\]
Thus $c(-s)=c(s)$. The vanishing of the odd Taylor coefficients follows
immediately.
\end{proof}

\begin{remark}[Why the scalar block is the only remaining mixed-local bottleneck]
\label{rem:scalar-block-only-remaining-mixed-local-bottleneck}
Through quintic order, the quotient-level ambiguity is already confined to the
$J_0$-anti-fixed traceless sector, while at septic order the canonical datum
is the gauge-invariant residual $\Delta_7$. Thus, on the clean branch
$V_1\neq 0$, the only remaining mixed-local theorem is the scalar descent
\[
C_1=C_3=C_5=0.
\]
Once this holds, the finite-jet bridge to the good $v_5$-patch coordinate
$x=v_5/c$ is formal and the explicit mixed low-order obstruction list
\[
(r,\operatorname{ord}_0\mathcal Z)\in\{(1,1),(1,3),(2,2)\}
\]
is excluded immediately.
\end{remark}

\begin{remark}[Centered symmetric-placement Taylor package and diagonal-base local exclusion]
\label{rem:centered-symmetric-placement-taylor-package}
Let
\[
t_\pm=h\pm \frac{s}{2},
\]
and regard $s\mapsto -s$ as swapping the two microscopic points. For any
symmetric $2\times 2$ coefficient matrix, the diagonal half-sums are even in
$s$, the diagonal half-differences are odd in $s$, and the off-diagonal
entries are even in $s$. Applying this to the explicit formulas already in the
draft gives
\[
u_5(s),\ v_5(s)\ \text{odd},
\qquad
u_{7,K_1}(s),\ v_{7,K_1}(s)\ \text{odd},
\]
Writing
\[
u_5(s)=U_1s+U_3s^3+U_5s^5+\cdots,
\qquad
v_5(s)=V_1s+V_3s^3+V_5s^5+\cdots,
\]
the first coefficients are
\[
U_1=2\bigl(v_0\,\delta\alpha_1+\beta_0\,\delta u_1\bigr),
\]
\[
U_3=2\bigl(v_0\,\delta\alpha_3+v_2\,\delta\alpha_1+\beta_0\,\delta u_3+\beta_2\,\delta u_1\bigr),
\]
\[
U_5=2\bigl(v_0\,\delta\alpha_5+v_2\,\delta\alpha_3+v_4\,\delta\alpha_1+\beta_0\,\delta u_5+\beta_2\,\delta u_3+\beta_4\,\delta u_1\bigr),
\]
and
\[
V_1=2\bigl(\bar\alpha_0\,\delta u_1+\delta\alpha_1\,\bar u_0\bigr),
\]
\[
V_3=2\bigl(\bar\alpha_0\,\delta u_3+\bar\alpha_2\,\delta u_1+\delta\alpha_1\,\bar u_2+\delta\alpha_3\,\bar u_0\bigr),
\]
\[
V_5=2\bigl(\bar\alpha_0\,\delta u_5+\bar\alpha_2\,\delta u_3+\bar\alpha_4\,\delta u_1+\delta\alpha_1\,\bar u_4+\delta\alpha_3\,\bar u_2+\delta\alpha_5\,\bar u_0\bigr).
\]

Likewise, writing
\[
u_{7,K_1}(s)=\widehat U_1s+\widehat U_3s^3+\cdots,
\qquad
v_{7,K_1}(s)=\widehat V_1s+\widehat V_3s^3+\cdots,
\]
one gets
\[
\widehat U_1
=
2\kappa\bigl(
\delta x_1\,\nu_{6,0}
+\beta_{5,0}\,\delta u_1
+v_0\,\delta\alpha_{5,1}
+b_0\,\delta p_1
+r_0\,\delta a_1
\bigr),
\]
\[
\widehat V_1
=
2\kappa\bigl(
\delta x_1\,\bar\mu_{6,0}
-\bar x_0\,\delta\mu_{6,1}
+\bar\alpha_{5,0}\,\delta u_1
+\delta\alpha_{5,1}\,\bar u_0
+\bar a_0\,\delta p_1
+\delta a_1\,\bar p_0
\bigr),
\]
and
\[
\widehat U_3
=
2\kappa\bigl(
\delta x_1\nu_{6,2}+\delta x_3\nu_{6,0}
+\beta_{5,0}\delta u_3+\beta_{5,2}\delta u_1
+v_0\delta\alpha_{5,3}+v_2\delta\alpha_{5,1}
+b_0\delta p_3+b_2\delta p_1
+r_0\delta a_3+r_2\delta a_1
\bigr),
\]
\[
\widehat V_3
=
2\kappa\bigl(
\delta x_1\bar\mu_{6,2}+\delta x_3\bar\mu_{6,0}
-\bar x_0\delta\mu_{6,3}-\bar x_2\delta\mu_{6,1}
+\bar\alpha_{5,0}\delta u_3+\bar\alpha_{5,2}\delta u_1
+\delta\alpha_{5,1}\bar u_2+\delta\alpha_{5,3}\bar u_0
+\bar a_0\delta p_3+\bar a_2\delta p_1
+\delta a_1\bar p_2+\delta a_3\bar p_0
\bigr).
\]

and therefore
\[
\Delta_7(s)=u_{7,K_1}(s)v_5(s)-u_5(s)v_{7,K_1}(s)
\]
is even in $s$, with Taylor expansion
\[
\Delta_7(s)=D_2s^2+D_4s^4+D_6s^6+\cdots,
\]
where
\[
D_2=\widehat U_1V_1-U_1\widehat V_1,
\]
\[
D_4=\widehat U_1V_3+\widehat U_3V_1-U_1\widehat V_3-U_3\widehat V_1,
\]
\[
D_6=\widehat U_1V_5+\widehat U_3V_3+\widehat U_5V_1-U_1\widehat V_5-U_3\widehat V_3-U_5\widehat V_1.
\]

\medskip

Assume now that
\[
c(s)=C_0+C_2s^2+C_4s^4+\cdots,
\qquad C_0\neq 0.
\]
Then
\[
x(s)=\frac{v_5(s)}{c(s)}
\]
is an odd series
\[
x(s)=X_1s+X_3s^3+X_5s^5+\cdots
\]
with
\[
X_1=\frac{V_1}{C_0},
\]
\[
X_3=\frac{V_3C_0-V_1C_2}{C_0^2},
\]
\[
X_5=\frac{V_5C_0^2-V_3C_0C_2+V_1(C_2^2-C_0C_4)}{C_0^3}.
\]
In particular, if $V_1\neq 0$, then $x$ is an odd local coordinate with odd
inverse
\[
s(x)=\sigma_1x+\sigma_3x^3+\sigma_5x^5+\cdots,
\]
where
\[
\sigma_1=\frac{C_0}{V_1},
\qquad
\sigma_3=-\frac{C_0(V_3C_0-V_1C_2)}{V_1^4},
\]
and so on.

Similarly,
\[
\lambda(s):=\frac{\Delta_7(s)}{v_5(s)}
\]
is an odd series
\[
\lambda(s)=L_1s+L_3s^3+L_5s^5+\cdots
\]
with
\[
L_1=\frac{D_2}{V_1},
\]
\[
L_3=\frac{D_4V_1-D_2V_3}{V_1^2},
\]
\[
L_5=\frac{D_6V_1^2-D_4V_1V_3+D_2(V_3^2-V_1V_5)}{V_1^3}.
\]

So, whenever $x$ is used as the induced odd collision-side coordinate, parity
in $s$ transfers to parity in $x$, and one obtains
\[
Y(x)=\frac{u_5}{c}(x)\ \text{odd},
\qquad
\lambda(x)=\frac{\Delta_7}{v_5}(x)\ \text{odd},
\qquad
S(x)=\frac{\lambda}{c}(x)\ \text{odd}.
\]
Consequently,
\[
Y''(x)\ \text{odd},\qquad
S''(x)\ \text{odd},
\]
hence
\[
\mathcal Z(x)=\frac{Y''}{S''}\ \text{even},
\qquad
\mathcal R(x)=\mathcal Z'(x)\ \text{odd},
\qquad
K_v=(x')^3S''\ \text{odd}.
\]
This is the parity package underlying the provisional mixed-branch exclusion
discussed below.

If $x=v_5/c$ is used as the induced odd collision-side coordinate, then the
centered diagonal germ has the odd/odd form
\[
Q(x)\sim (x^n,\ x,\ x^m),
\qquad
m,n\ \text{odd},\quad n>m\ge 3.
\]
In that regime one obtains the exact factorizations
\[
S''(x)=x^{m-2}\Sigma_m[A](x^2),
\qquad
\mathcal Z(x)=x^{n-m}H(x^2),
\qquad
\mathcal R(x)=x^{n-m-1}J(x^2),
\]
with $H(0)\neq 0$ and $J(0)\neq 0$. Hence there exists $\varepsilon>0$
such that for
\[
0<|x|<\varepsilon
\]
one has
\[
K_v(x)\neq 0,\qquad T_v(x)\neq 0.
\]

So, in the diagonal-based local regime where this odd collision-side coordinate
description is valid, the centered diagonal base has a punctured neighborhood
free of $K_v=0$ and $T_v=0$ points. In particular, the currently explicit
low-order dangerous patterns do not arise at the symmetric diagonal base itself.

\medskip

\noindent
\textit{Work in progress.}
The odd/even transfer from the centered microscopic variable $s$ to a general
good-patch coordinate $x=v_5/c$ should still be treated as a live theorem-level
bridge outside the diagonal-based local regime. What is no longer open at the
same level is total quintic collapse on the centered branch: by
Lemma~\ref{lem:centered-tagged-quintic-witness-survives-specialization} and
Corollary~\ref{cor:centered-quintic-fixed-sector-nontriviality}, one already
has
\[
u_5(s)\not\equiv 0,
\qquad
A_5^{\mathfrak f}(s)\not\equiv 0.
\]
So the remaining mixed-side task is not to prove that some quintic channel
survives, but to promote this centered quintic nontriviality into an actual
punctured $v_5$- or $u_5$-patch.
\end{remark}

\begin{lemma}[Total centered quintic collapse forces source-level degeneration]
\label{lem:total-centered-quintic-collapse}
On the centered symmetric-placement branch one has
\[
u_5=2\bigl(v\,\delta\alpha+\beta\,\delta u\bigr),
\qquad
v_5=2\bigl(\bar\alpha\,\delta u+\bar u\,\delta\alpha\bigr).
\]
If
\[
u_5\equiv 0,
\qquad
v_5\equiv 0,
\]
then
\[
\begin{pmatrix}
\beta(s) & v(s)\\
\bar\alpha(s) & \bar u(s)
\end{pmatrix}
\binom{\delta u(s)}{\delta\alpha(s)}
=0
\]
as an identity of analytic germs. Consequently, either
\[
\Delta_5(s):=\beta(s)\bar u(s)-v(s)\bar\alpha(s)\equiv 0,
\]
or
\[
\delta u(s)\equiv 0,
\qquad
\delta\alpha(s)\equiv 0.
\]
\end{lemma}

\begin{proof}
This is immediate from the displayed formulas for $u_5$ and $v_5$.
If the $2\times 2$ coefficient matrix has nonzero determinant, then the only
solution is $\delta u=\delta\alpha\equiv 0$. If its determinant vanishes
identically, then $\Delta_5\equiv 0$.
\end{proof}

\begin{corollary}[Total centered quintic collapse forces total canonical transport collapse]
\label{cor:total-centered-quintic-collapse-forces-total-transport-collapse}
If
\[
u_5\equiv 0,
\qquad
v_5\equiv 0
\]
on the centered branch, then
\[
A_5^{\mathfrak f}\equiv 0,
\qquad
\Delta_7\equiv 0.
\]
In particular, total centered quintic collapse forces collapse of the canonical
quintic/septic transport package on that branch.
\end{corollary}

\begin{proof}
The identity
\[
A_5^{\mathfrak f}=u_5 I+v_5 S
\]
gives $A_5^{\mathfrak f}\equiv 0$ immediately. Likewise
\[
\Delta_7=u_{7,K_1}v_5-u_5v_{7,K_1}
\]
vanishes identically when $u_5\equiv v_5\equiv 0$.
\end{proof}

\begin{lemma}[Centered tagged quintic witness survives specialization]
\label{lem:centered-tagged-quintic-witness-survives-specialization}
On the tagged slice
\[
(K_1,\eta_2,q^{(5)}),
\]
the first centered odd coefficient
\[
U_1
\]
of
\[
u_5(s)=U_1s+U_3s^3+U_5s^5+\cdots
\]
has a nonzero $\kappa\,\eta_2\,q^{(5)}$-coefficient. In particular
\[
U_1\neq 0
\]
on that tagged slice, and hence
\[
u_5(s)\not\equiv 0
\]
as a centered germ.
\end{lemma}

\begin{proof}
The centered coefficient formula is
\[
U_1=2\bigl(v_0\,\delta\alpha_1+\beta_0\,\delta u_1\bigr).
\]
By the tagged-slice analysis already used in the proof of
Lemma~\ref{lem:no-cancellation-quintic-trace}, the $\eta_2$-slice of $X_1$
is off-diagonal, so
\[
\delta u_1=0,
\]
and the $q^{(5)}$-slice of $X_3$ is diagonal, so
\[
\beta_0=0.
\]
Hence
\[
U_1=2v_0\,\delta\alpha_1
\]
on the tagged slice.

Now use the explicit two-direction computation from the proof of
Lemma~\ref{lem:no-cancellation-quintic-trace}. There one writes
\[
G_0=\operatorname{diag}(a^2,b^2),
\qquad
A=
\begin{pmatrix}
0&u\\u&0
\end{pmatrix},
\qquad
D=
\begin{pmatrix}
0&0\\0&d
\end{pmatrix},
\]
and expands
\[
X(\varepsilon,\delta)
=
X_0+\varepsilon U+\delta V+\varepsilon\delta W+O(\varepsilon^2,\delta^2).
\]
The proof gives
\[
U_{12}=-\frac{u}{ab(a+b)}\neq 0,
\qquad
V_{22}=-\frac{d}{2b^3}\neq 0.
\]
Here $U_{12}$ gives the tagged $\eta_2$-contribution to the off-diagonal
channel $v_0$, while $V_{22}$ gives the tagged $q^{(5)}$-contribution to the
diagonal-difference channel $\gamma-\alpha=2\delta\alpha_1$. Thus the tagged
slice contributes nontrivially to the product $v_0\,\delta\alpha_1$, and hence
to $U_1=2v_0\,\delta\alpha_1$.
\end{proof}

\begin{corollary}[Centered quintic fixed-sector nontriviality]
\label{cor:centered-quintic-fixed-sector-nontriviality}
On the centered diagonal branch,
\[
A_5^{\mathfrak f}(s)=u_5(s)I+v_5(s)S
\]
is not identically zero. Equivalently,
\[
(u_5(s),v_5(s))\not\equiv (0,0).
\]
\end{corollary}

\begin{proof}
By Lemma~\ref{lem:centered-tagged-quintic-witness-survives-specialization},
one has $u_5(s)\not\equiv 0$. Since $I$ and $S$ are linearly independent,
the fixed-sector package $A_5^{\mathfrak f}(s)=u_5(s)I+v_5(s)S$ cannot vanish
identically.
\end{proof}

\begin{lemma}[Valuation form of the septic transport constraint]
\label{lem:valuation-form-septic-transport-constraint}
Assume on the centered branch that
\[
\lambda=\frac{\Delta_7}{v_5}
\]
is analytic odd. Write
\[
\operatorname{ord}_s u_5=2m+1,
\qquad
\operatorname{ord}_s v_5=2r+1,
\qquad
r\ge m.
\]
Then
\[
\operatorname{ord}_s v_{7,K_1}\ge \operatorname{ord}_s v_5-\operatorname{ord}_s u_5+1=2(r-m)+1.
\]
Equivalently, if
\[
U_1=\cdots=U_{2m-1}=0,\qquad U_{2m+1}\neq 0,
\]
and
\[
V_1=\cdots=V_{2r-1}=0,\qquad V_{2r+1}\neq 0,
\]
then
\[
\widehat V_1=\widehat V_3=\cdots=\widehat V_{2(r-m)-1}=0,
\]
and the first potentially surviving even coefficient of $\Delta_7$ is
\[
D_{2r+2}=\widehat U_1V_{2r+1}-U_{2m+1}\widehat V_{2(r-m)+1}.
\]
\end{lemma}

\begin{proof}
Since $\lambda$ is analytic odd and
\[
\Delta_7=\lambda\,v_5,
\]
one has
\[
\operatorname{ord}_s\Delta_7\ge \operatorname{ord}_s v_5+1=2r+2.
\]
Suppose
\[
\operatorname{ord}_s v_{7,K_1}=2d+1<2(r-m)+1.
\]
Then
\[
\operatorname{ord}_s(u_5v_{7,K_1})=(2m+1)+(2d+1)\le 2r,
\]
whereas
\[
\operatorname{ord}_s(u_{7,K_1}v_5)\ge 1+(2r+1)=2r+2
\]
because $u_{7,K_1}$ is odd. Therefore the leading term of
\[
\Delta_7=u_{7,K_1}v_5-u_5v_{7,K_1}
\]
would occur at order at most $2r$, contradicting $\operatorname{ord}_s\Delta_7\ge 2r+2$.
So
\[
\operatorname{ord}_s v_{7,K_1}\ge 2(r-m)+1.
\]

The coefficient statement is the corresponding triangular coefficient form:
vanishing of
\[
D_2,D_4,\dots,D_{2r}
\]
forces
\[
\widehat V_1,\widehat V_3,\dots,\widehat V_{2(r-m)-1}
\]
to vanish successively, and the first potentially surviving coefficient is
\[
D_{2r+2}=\widehat U_1V_{2r+1}-U_{2m+1}\widehat V_{2(r-m)+1}.
\qedhere
\]
\end{proof}

\begin{lemma}[Late-$v_5$ $u_5$-patch affine-normalization transplant]
\label{lem:late-v5-u5-patch-transplant}
Assume the centered scalar package from
Remark~\ref{rem:centered-symmetric-placement-taylor-package}, and assume
\[
U_1\neq 0.
\]
Then
\[
y:=\frac{u_5}{c}
\]
is an odd local coordinate near the centered base. Suppose in addition that
\[
\operatorname{ord}_s v_5=2r+1\ge 3.
\]
On the corresponding $u_5$-patch write
\[
x(y):=\frac{v_5}{c},
\qquad
\tau_u(y):=-\frac{\Delta_7}{c\,u_5},
\qquad
\widetilde\tau_u(y):=\tau_u(y)-\tau_u'(0)\,y.
\]
Then:
\begin{enumerate}
\item $x(y)$ is odd and satisfies
\[
x(y)=O(y^{2r+1});
\]
in particular $\operatorname{ord}_0 x\ge 3$;
\item $\widetilde\tau_u(y)$ is odd and either vanishes identically or has odd
order at least $3$;
\item if $\widetilde\tau_u\not\equiv 0$, then after an ambient affine
normalization consisting of a coordinate permutation and subtraction of a
suitable odd multiple of one higher-order channel from the other, the germ
$Q_u$ is carried to the same odd/odd local normal form
\[
(y^n,\ y,\ y^m),
\qquad
n>m\ge 3,\quad m,n\ \text{odd},
\]
already analyzed in
Remark~\ref{rem:centered-symmetric-placement-taylor-package}.
\end{enumerate}
Consequently, in the subcase $\widetilde\tau_u\not\equiv 0$, the late-$v_5$
corner is locally sterile on punctured $u_5$-patch arcs: no genuinely new
mixed low-order $4$-point branch is born there.
\end{lemma}

\begin{proof}
By Lemma~\ref{lem:centered-tagged-quintic-witness-survives-specialization},
\[
u_5(s)=U_1s+O(s^3),
\qquad
c(s)=C_0+O(s^2),
\qquad
C_0\neq 0,
\]
so
\[
y(s)=\frac{u_5(s)}{c(s)}
\]
is an odd local coordinate with odd inverse $s=s(y)$.

Since $v_5$ is odd and $c$ is even in $s$,
\[
x(y)=\frac{v_5}{c}
\]
is odd in $y$. Because $\operatorname{ord}_s u_5=1$ and
$\operatorname{ord}_s v_5=2r+1$, one has
\[
x(y)=O(y^{2r+1}),
\]
so $\operatorname{ord}_0x\ge 3$.

Likewise $\Delta_7$ is even while $c$ is even and $u_5$ is odd, hence
\[
\tau_u(y)=-\frac{\Delta_7}{c\,u_5}
\]
is odd in $y$. Therefore
\[
\widetilde\tau_u(y):=\tau_u(y)-\tau_u'(0)\,y
\]
is odd and either identically zero or of odd order at least $3$.

Assume now that $\widetilde\tau_u\not\equiv 0$. Then the normalized germ is
\[
Q_u(y)=\bigl(y,\ x(y),\ \widetilde\tau_u(y)\bigr),
\]
with the last two channels odd and of order at least $3$. Apply the ambient
coordinate permutation
\[
(Y,X,Z)\longmapsto (X,Y,Z).
\]
This preserves affine dependence/coplanarity geometry. If the two higher-order
odd channels have the same leading odd order, subtract a suitable odd multiple
of one from the other to separate those leading orders. After this affine
normalization one reaches a germ of the form
\[
(y^nB(y^2),\ y,\ y^mA(y^2)),
\qquad
n>m\ge 3,\quad m,n\ \text{odd},
\]
with $A(0)B(0)\neq 0$, i.e. the same odd/odd local normal form analyzed in
Remark~\ref{rem:centered-symmetric-placement-taylor-package}. The same local
calculation then gives a punctured neighborhood free of nearby mixed and
nonmixed $4$-point obstruction points.
\end{proof}

\begin{lemma}[Planar-lift absorption on the $u_5$-patch]
\label{lem:planar-lift-absorption-u5-patch}
Assume the hypotheses of
Lemma~\ref{lem:late-v5-u5-patch-transplant}, and suppose
\[
\widetilde\tau_u\equiv 0.
\]
Equivalently, there is a constant $a\in\mathbf C$ such that
\[
\Delta_7=-a\,u_5^2
\]
locally on the $u_5$-patch. Then after an ambient affine shear one has
\[
Q_u\sim (y,x,0),
\qquad
y=\frac{u_5}{c},
\quad
x=\frac{v_5}{c}.
\]
In particular, the lifted curve is affine-planar, and the remaining local
content is only the planar affine dependence of
\[
\Xi(h):=\left(\frac{u_5(h)}{c(h)},\frac{v_5(h)}{c(h)}\right).
\]

Moreover, in the defect-free order-$\le 5$ finite-core setting of
Theorem~\ref{thm:defect-free-quintic-circuit-reduction}, this residual corner is
absorbed into the already-known quintic-circuit regime: it does not define a
genuinely new mixed/septic $4$-point branch.
\end{lemma}

\begin{proof}
The identity $\widetilde\tau_u\equiv 0$ means exactly that
\[
\tau_u(y)=a\,y
\]
for some constant $a$. Since
\[
\tau_u=-\frac{\Delta_7}{c\,u_5},
\qquad
y=\frac{u_5}{c},
\]
this is equivalent to
\[
-\frac{\Delta_7}{c\,u_5}=a\,\frac{u_5}{c},
\]
hence
\[
\Delta_7=-a\,u_5^2.
\]
Thus
\[
Q_u=(y,x,ay).
\]
Applying the ambient affine shear
\[
(Y,X,Z)\longmapsto (Y,X,Z-aY)
\]
gives
\[
Q_u\sim (y,x,0).
\]
So the lifted curve is affine-planar.

Now assume additionally the defect-free order-$\le 5$ finite-core hypotheses
from Theorem~\ref{thm:defect-free-quintic-circuit-reduction}. By
Remark~\ref{rem:parallel-u5-patch-version}, the same defect-free
affine-dependence conclusions hold with $Q_v$ replaced by $Q_u$. Since the
third coordinate of $Q_u$ vanishes after the shear, affine dependence of the
lifted points is equivalent to affine dependence of the planar points
\[
\Xi(h_j)=\left(\frac{u_5(h_j)}{c(h_j)},\frac{v_5(h_j)}{c(h_j)}\right)\in\mathbf C^2.
\]
But this is exactly the planar quintic-circuit geometry organized by
Theorem~\ref{thm:defect-free-quintic-circuit-reduction} and
Corollary~\ref{cor:defect-free-minimal-quintic-circuits}. Hence this residual
corner is already-quintic rather than a genuinely new mixed septic branch.
\end{proof}

\begin{corollary}[Right-pincer trichotomy near the centered base]
\label{cor:right-pincer-trichotomy-centered-base}
Near the centered diagonal base, the mixed local picture is reduced to the
following three cases:
\[
V_1\neq 0
\Longrightarrow
\text{clean }v_5\text{-patch branch;}
\]
\[
V_1=0,\ \widetilde\tau_u\not\equiv 0
\Longrightarrow
\text{late-}v_5\text{ branch locally sterile on }Q_u;
\]
\[
V_1=0,\ \widetilde\tau_u\equiv 0
\Longleftrightarrow
\Delta_7=-a\,u_5^2
\Longrightarrow
\text{planar-lift corner, hence already-quintic.}
\]
So every late-$v_5$ branch near the centered base is either locally sterile
or already absorbed into the quintic-circuit regime. The only genuinely new
mixed local branch left to study there is the clean branch
\[
\operatorname{ord}_s v_5=1.
\]
\end{corollary}

\begin{proof}
Combine Lemma~\ref{lem:late-v5-u5-patch-transplant} and
Lemma~\ref{lem:planar-lift-absorption-u5-patch}, together with the fact that
$v_5$ is odd in the centered variable.
\end{proof}

\begin{corollary}[Right-pincer local reduction near the centered base]
\label{cor:right-pincer-local-reduction-centered-base}
Assume the centered scalar package from
Remark~\ref{rem:centered-symmetric-placement-taylor-package}, and assume also
that the centered tagged quintic witness gives
\[
U_1\neq 0.
\]
Then every
sufficiently small nearby genuinely new mixed $4$-point obstruction point near
the centered base must lie on the clean branch
\[
\operatorname{ord}_s v_5=1,
\]
equivalently
\[
V_1\neq 0.
\]
In particular, every late-$v_5$ branch near the centered base is either
locally sterile on the $u_5$-patch or already absorbed into the defect-free
quintic-circuit regime.
\end{corollary}

\begin{proof}
Because $v_5$ is odd in the centered variable, either
\[
\operatorname{ord}_s v_5=1
\qquad\text{or}\qquad
\operatorname{ord}_s v_5\ge 3 \text{ odd}.
\]
If $\operatorname{ord}_s v_5\ge 3$, then
Lemma~\ref{lem:late-v5-u5-patch-transplant} and
Lemma~\ref{lem:planar-lift-absorption-u5-patch}
show that the branch is either locally sterile on the $u_5$-patch or
planar-lift and hence already-quintic. So a sufficiently small nearby
genuinely new mixed $4$-point obstruction cannot lie on such a late-$v_5$
branch. Therefore any such genuinely new mixed local branch must lie on the
only remaining possibility,
\[
\operatorname{ord}_s v_5=1.
\qedhere
\]
\end{proof}

\begin{proposition}[Realization lemma for lifted germs]
\label{prop:realization-lemma-lifted-germs}
Let $\widetilde Y(\xi)$ and $\widetilde S(\xi)$ be arbitrary analytic germs
at $\xi=0$. Define
\[
c(\xi)\equiv 1,
\qquad
v_5(\xi):=1+\xi,
\qquad
u_5(\xi):=\widetilde Y(\xi),
\qquad
\Delta_7(\xi):=(1+\xi)\widetilde S(\xi).
\]
Then on a sufficiently small neighborhood of $0$,
\[
c\neq 0,
\qquad
v_5\neq 0,
\qquad
M=c\,v_5'-c'v_5=1\neq 0,
\]
so this is a good $v_5$-patch. The normalized lifted coordinates are
\[
x=\frac{v_5}{c}=1+\xi,
\qquad
Y=\frac{u_5}{c}=\widetilde Y(\xi),
\qquad
S=\frac{\Delta_7}{c\,v_5}=\widetilde S(\xi).
\]
After the affine shift $x\mapsto x-1$, the lifted germ becomes
\[
Q_v(\xi)=\bigl(\widetilde Y(\xi),\,\xi,\,\widetilde S(\xi)\bigr).
\]

Thus every analytic pair of germs $(\widetilde Y,\widetilde S)$ occurs, after
affine normalization, as the lifted pair $(Y,S)$ on some good $v_5$-patch.
\end{proposition}

\begin{proof}
The displayed formulas are immediate from the definitions. The hypotheses
$c\neq 0$, $v_5\neq 0$, and $M=1\neq 0$ hold on a sufficiently small
neighborhood of $0$, and the formulas for $x,Y,S$ are direct.
\end{proof}

\begin{corollary}[No new obstruction for $(\mathcal Z,\mathcal R)$ is available at the current lifted level]
\label{cor:no-new-jet-obstruction-current-lifted-level}
No identity involving only
\[
c,\qquad v_5,\qquad u_5,\qquad \Delta_7
\]
and their derivatives can force either of the two dream conclusions:
\[
\mathcal R \text{ is even in the normalized nonmixed coordinate},
\]
or the mixed-branch exclusion
\[
(r,\operatorname{ord}_0\mathcal Z)\notin\{(1,1),(1,3),(2,2)\}
\quad\text{at every mixed point.}
\]

In particular, the nonmixed models
\[
Q_v(\xi)=(\xi^n,\xi,\xi^2),
\qquad n\ge 4 \ \text{even},
\]
are all realizable at the present lifted level; the cases
\[
n=4,6,8
\]
are simply the first currently explicit low-order examples.
The dangerous mixed models
\[
Q_v(\xi)=(\xi^n,\xi,\xi^m),
\qquad (m,n)=(3,4),\ (3,6),\ (4,6),
\]
are all realizable at the present lifted level.

Consequently, any RH-specific obstruction to the low-order dangerous
$4$-point patterns must use information not encoded in the current lifted
variables
\[
(c,v_5,u_5,\Delta_7)
\]
alone. It must come from deeper source-level corrected quotient-germ
structure, matrix-level identities behind $(u_5,\Delta_7)$, or an additional
symmetry not yet visible in the present lifted formulation.
\end{corollary}

\begin{proof}
Apply Proposition~\ref{prop:realization-lemma-lifted-germs} with
\[
(\widetilde Y,\widetilde S)=(\xi^n,\xi^2)
\]
for the nonmixed cases, and with
\[
(\widetilde Y,\widetilde S)=(\xi^n,\xi^m)
\]
for the mixed cases. These realize exactly the dangerous low-order patterns for
$(\mathcal Z,\mathcal R)$ listed above.
\end{proof}

\begin{remark}[Updated nonmixed status versus mixed local exhaustion]
\label{rem:wip-dream-parity-identity-four-point}
The local nonmixed lifted route is now closed at the derivative-geometric level.
Exact nonmixed contacts force zeros of $\mathcal R$; one-sign
$\mathcal R$-intervals are sterile; recurrent exact contacts force
affine-planar collapse; and isolated zeros of $\mathcal R$ are classified by
parity. In addition, the repeated-source / first-loss mechanism is now fully
reduced: no attained separated first loss can occur, every first-loss sequence
is forced to the diagonal, and diagonal first loss occurs exactly at odd zeros
of $\mathcal R$. Thus the nonmixed burden is no longer a local parity theorem
or a separate repeated-source obstruction. It is the global/source-level
control of which odd zeros of
\[
\mathcal R=\frac{T_v}{K_v^2},
\qquad\text{equivalently}\qquad
\mathcal W_{\mathrm{src}}=J_\lambda J_{u_5}'-J_{u_5}J_\lambda',
\]
can actually arise along opened nonmixed arcs.

On the mixed side, the decisive pair-level object remains
\[
A_5^{\mathfrak f}(s)=u_5(s)I+v_5(s)S.
\]
By Corollary~\ref{cor:centered-quintic-fixed-sector-nontriviality}, this
package is not identically zero, so after shrinking to one-sided punctured arcs
the centered branch enters either a good $v_5$-patch or the parallel
$u_5$-patch. Near the centered base, however, the theorem package now goes
further: every late-$v_5$ branch is already either locally sterile on the
automatic $u_5$-patch or absorbed into the planar defect-free quintic-circuit
regime, so the only genuinely new mixed local branch left there is the clean
branch
\[
V_1\neq 0.
\]

Accordingly, the mixed-side local problem is no longer a vague codimension-two
obstruction, nor merely a patch-opening issue. Near the centered base it has
been reduced to the single clean branch $V_1\neq 0$, and on that branch the
remaining theorem-level task is the finite scalar descent
\[
c(-s)=c(s)\pmod{s^6},
\qquad\text{equivalently}\qquad
C_1=C_3=C_5=0,
\]
which makes the centered-to-patch finite-jet parity bridge formal.
\end{remark}

\begin{remark}[What the four-point branch now reduces to]
\label{rem:next-algebraic-target-four-point}
The lifted $4$-point branch is now organized by the canonical invariant pair
\[
\mathcal Z=\frac{Y''}{S''}=\frac{J_{u_5}}{J_\lambda},
\qquad
\mathcal R=\mathcal Z'=\frac{T_v}{K_v^2},
\qquad
\lambda=\frac{\Delta_7}{v_5}.
\]

On the nonmixed side, after normalization to
\[
Q_v(\xi)=\bigl(\widetilde Y(\xi),\,\xi,\,\xi^2\bigr),
\]
the local derivative-side picture is now explicit. Exact nonmixed contacts force
zeros of
\[
\mathcal R=\frac12\,\widetilde Y''',
\]
isolated zeros are classified by parity, one-sign $\mathcal R$-intervals are
exactly-contact sterile, and recurrent exact contacts force affine-planar
collapse.

The explicit low-order list
\[
\operatorname{ord}_0\mathcal R\in\{1,3,5\},
\]
equivalently the coefficients
\[
z_2,\qquad z_4,\qquad z_6,
\]
should therefore now be regarded only as the first explicit low-order nonmixed
odd-order cases, not as the full local nonmixed picture.

For fixed source point $a$, repeated-source exact contacts factor through the
quadratic-envelope scalar
\[
\Gamma_a(x)=
\frac{Y(x)-Y(a)-Y'(a)(x-a)}
{S(x)-S(a)-S'(a)(x-a)},
\]
equivalently through the below-lifted quotient
\[
\Psi_a(h)=\frac{D_{u_5}(a,h)}{D_\lambda(a,h)}.
\]
Its derivative is governed by the osculatory $4\times 4$ source determinant
$\mathfrak D(a,h)$. By
Proposition~\ref{prop:first-loss-is-diagonal-odd-zero-birth}, this
repeated-source / first-loss mechanism is now completely reduced: no attained
separated first loss can occur, and any first loss is exactly the diagonal
birth at an isolated odd-order zero of $\mathcal R$.

Thus the remaining nonmixed burden is no longer local parity, nor a separate
first-loss geometry, but purely the global/source-level control of which odd
zeros of $\mathcal W_{\mathrm{src}}$, equivalently of $\mathcal R$, occur
on genuine opened arcs.

On the mixed side, the explicit low-order dangerous list is
\[
(m,n)=(3,4),\ (3,6),\ (4,6),
\]
equivalently
\[
(r,s)=(1,2),\ (1,4),\ (2,5),
\]
equivalently
\[
(r,\operatorname{ord}_0\mathcal Z)=(1,1),\ (1,3),\ (2,2),
\]
equivalently
\[
(r,\operatorname{ord}_0\mathcal R)=(1,0),\ (1,2),\ (2,1).
\]

\noindent
Equivalently, in the below-lifted Wronskian formulation from
Remark~\ref{rem:wronskian-order-reformulation-mixed-list}, the dangerous mixed
models are exactly the order-pairs
\[
(\operatorname{ord}_0J_\lambda,\operatorname{ord}_0J_{u_5})
=(1,2),\ (1,4),\ (2,4).
\]
So the mixed low-order branch is already a finite canonical
Wronskian-ratio obstruction problem, captured in local mixed $x$-coordinates by
\[
\kappa_{1,1},\qquad \kappa_{1,3},\qquad \kappa_{2,2}.
\]

By the realization lemma, no identity involving only
\[
(c,v_5,u_5,\Delta_7)
\]
and their derivatives can eliminate these mixed patterns in general.

At the pair level, the centered fixed-sector quintic package
\[
A_5^{\mathfrak f}(s)=u_5(s)I+v_5(s)S
\]
is now the natural object controlling whether any quintic patch opens at all.
By Corollary~\ref{cor:centered-quintic-fixed-sector-nontriviality},
\[
A_5^{\mathfrak f}(s)\not\equiv 0,
\]
so the centered branch cannot suffer total quintic collapse.

Near the centered base, however, the mixed-side local picture is now sharper
than a mere patch dichotomy. By
Corollary~\ref{cor:right-pincer-local-reduction-centered-base}, every
sufficiently small nearby genuinely new mixed $4$-point obstruction must lie
on the clean branch
\[
\operatorname{ord}_s v_5=1,
\]
while every late-$v_5$ branch is either locally sterile on the automatic
$u_5$-patch or already absorbed into the defect-free planar
quintic-circuit regime. Thus the genuinely new mixed local geometry near the
centered base has been reduced to the single clean branch $V_1\neq 0$.

Accordingly, the remaining mixed-side theorem is no longer another lifted jet
computation, nor a pointwise coefficient nonvanishing statement such as
$V_1\neq 0$ in isolation. It is the finite scalar descent on the clean branch,
namely
\[
c(-s)=c(s)\pmod{s^6},
\qquad\text{equivalently}\qquad
C_1=C_3=C_5=0,
\]
which makes the centered good-$v_5$-patch parity bridge formal. After that
centered local step, the clean-patch mixed-global problem is no longer open
mixed geometry: it is absorbed by the mixed-to-scalar capture theorem
(Theorem~\ref{thm:global-mixed-to-scalar-capture-clean-patch}), with the
source-level viewpoint serving only as the correct language for the lower-core
and source-coupled handoff walls.
\end{remark}

\begin{remark}[What the new nonmixed/source-level machinery changes for the mixed side]
\label{rem:nonmixed-source-machinery-mixed-side}
The new repeated-source/source-determinant viewpoint and the no-go principle
excluding purely planar $(x,S)$-closure do not create a new mixed-local closure
theorem. Rather, they clarify the remaining mixed burden.

First, no argument using only planar $(x,S)$-geometry can close the $4$-point
branch, so any remaining mixed route must use either the centered scalar
descent or genuine coupled $Y$--$S$ source-level information.

Second, if a below-lifted backup route is needed, the natural mixed object is
not an $S$-only invariant, but a coupled source-minor ratio built from
$u_5$ and
\[
\lambda=\frac{\Delta_7}{v_5},
\]
i.e. a below-lifted avatar of
\[
\mathcal Z=\frac{J_{u_5}}{J_\lambda}.
\]

Thus the new nonmixed/source-level machinery functions not only as a strategic
constraint and as a backup language for the mixed side, but also as the
correct external handoff target for the $A=0$ planar collapse; see
Remark~\ref{rem:A-zero-is-nonmixed-handoff}. It does not replace the local
mixed exhaustion theorem near the centered base, but it does absorb the
nonmixed affine-planar lane that remains after the local mixed reduction.
\end{remark}

\begin{remark}[The $A=0$ collapse is already the nonmixed $\mathcal R$-lane]
\label{rem:A-zero-is-nonmixed-handoff}
Work on a good $v_5$-patch and write an exact lifted $4$-point relation in the
form
\[
A\,Y+B\,x+C\,S+D=0,
\qquad
Y=\frac{u_5}{c},\quad x=\frac{v_5}{c},\quad S=\frac{\Delta_7}{c\,v_5}.
\]
If $A=0$, then the relation reduces to
\[
B\,x+C\,S+D=0,
\]
so the four nodes lie on an affine line in the planar $(x,S)$-curve. Thus the
positive-dimensional $A=0$ regime is not a separate mixed branch at all: it is
precisely the nonmixed affine-planar exact-contact regime. In particular, the
relevant theorem-strength input is already on the nonmixed side: exact
nonmixed contacts force zeros of $\mathcal R$
(Proposition~\ref{prop:exact-nonmixed-tangent-plane-contacts-force-R-zero}),
one-sign $\mathcal R$-intervals are sterile
(Corollary~\ref{cor:sign-definite-R-excludes-exact-nonmixed-contacts}), and
interval-persistent / recurrent exact contacts force affine-planar collapse
(Corollary~\ref{cor:recurrent-exact-nonmixed-contacts-force-affine-planarity}).
Accordingly, $A=0$ should be viewed as the external handoff from the mixed side
into the already-handled nonmixed $\mathcal R$-package, not as a fourth
independent mixed-global lane.
\end{remark}

\begin{remark}[Status of the lifted order-$7$ geometry]
\label{rem:wip-lifted-order7-zero-capture}
The lifted geometry is now substantially more stratified than a single
zero-capture problem.

First, shrinking exact lifted three-point and four-point circuits can accumulate
only at
\[
K_v=0
\qquad\text{and}\qquad
T_v=0,
\]
respectively, with uniform small-diameter exclusion away from those zero sets;
equivalently, they are rooted at the curvature- and torsion-type osculating
determinants of the corresponding lifted planar and space curves
(Lemma~\ref{lem:osculating-determinant-expansion},
Propositions~\ref{prop:shrinking-exact-lifted-three-point-circuits-force-Kv-zero}
and~\ref{prop:shrinking-exact-lifted-four-point-circuits-force-Tv-zero}).

Second, repeated-node limits are now explicit:
three-point collapse limits become tangent-line contacts for the
planar lifted curve $\Lambda_v=(x,S)$, and four-point collapse limits become
tangent-plane contacts for the lifted space curve $Q_v=(Y,x,S)$.

Third, the three-point side is now blockwise globally controlled on good
$v_5$-patches.
The curvature zero set $K_v=0$ has no infinite-order flat points; even-order
$K_v$-zeros are sterile; odd-order $K_v$-zeros carry exactly one local
opposite-side tangent-contact branch; and that branch continues uniquely on the
maximal right monotone-slope block, with no interior common-tangent event, no
source-side return, and no interior accumulation. Thus the three-point side is
no longer a live internal recurrence problem on such
patches; only boundary exit of the relevant source/target slope-overlap block
can remain, and in the four-point mixed-to-scalar picture this is used only as
an external lower-core handoff target rather than as a new mixed carrier lane.

On the mixed four-point side, the right-pincer package now gives an analogous
local simplification near the centered base: every late-$v_5$ branch is either
locally sterile on the $u_5$-patch or already absorbed into the planar
quintic-circuit regime, so the only genuinely new mixed local branch left there
is the clean branch $\operatorname{ord}_s v_5=1$.

Fourth, the four-point side splits sharply into nonmixed and mixed regimes.
On the nonmixed side, finite-order odd $T_v$-zeros already support genuine
local exact branch sheets in the model geometry, so a blanket local sterility
theorem is false. On the mixed side, the dangerous low-order obstruction list is
finite and explicit:
\[
(m,n)=(3,4),\ (3,6),\ (4,6),
\]
equivalently
\[
(r,s)=(1,2),\ (1,4),\ (2,5),
\]
equivalently
\[
(r,\operatorname{ord}_0\mathcal Z)=(1,1),\ (1,3),\ (2,2),
\]
equivalently
\[
(r,\operatorname{ord}_0\mathcal R)=(1,0),\ (1,2),\ (2,1).
\]

Fifth, below the lifted variables one has the exact Wronskian-ratio
reformulation
\[
\mathcal Z=\frac{J_{u_5}}{J_\lambda},
\qquad
\mathcal R=\frac{c^2}{M}\frac{d}{dh}\!\left(\frac{J_{u_5}}{J_\lambda}\right),
\qquad
\lambda=\frac{\Delta_7}{v_5},
\]
while the projective curvature ratio $\Theta$ yields no new independent local
obstruction beyond $(\mathcal Z,\mathcal R)$.

Sixth, the centered symmetric-placement Taylor package gives a strong local
collision-side parity picture:
\[
u_5,\ v_5,\ u_{7,K_1},\ v_{7,K_1}\ \text{odd in }s,
\qquad
\Delta_7\ \text{even in }s,
\]
with explicit septic coefficients $D_2,D_4,D_6$. In the diagonal-centered
odd/odd normal form
\[
Q(x)\sim (x^n,\ x,\ x^m),
\qquad
m,n\ \text{odd},\quad n>m\ge 3,
\]
one gets a punctured neighborhood free of $K_v=0$ and $T_v=0$ points.
Thus the explicit low-order dangerous $4$-point patterns do not arise at the
symmetric diagonal base itself.

So the mixed four-point side is no longer waiting for a new local geometry
theorem. On the mixed side, the sufficiently small problem is now organized by
the mixed-to-scalar capture principle: after the source-coupled and patch-loss
reductions, no sufficiently small genuinely mixed branch on a nonplanar good
patch can persist entirely inside
\[
\{T_v\neq 0\}\cap\{\mathcal R\neq 0\},
\]
except after entry into already-absorbed planar, lower-core, repeated-node,
repeated-source, or quintic-collapse regimes. This is formalized below by the
local clean continuation proposition, the positive-scale endpoint corollary, and
Theorem~\ref{thm:global-mixed-to-scalar-capture-clean-patch}. On the still-live
nonmixed side, any genuinely new input would have to be RH-specific source- or
matrix-level structure imposing stronger constraints on
\[
\mathcal R=\frac{T_v}{K_v^2}.
\]
\end{remark}

\begin{remark}[Live attack routes]
\label{rem:live-attack-routes}
At the present stage, the live routes are now quite sharply separated.

\begin{enumerate}
\item
\textbf{Two-point branch: the residual burden is now split more sharply.}
The stronger algebraic route is still
Lemma~\ref{lem:minimal-source-level-two-point-criterion}: prove that the actual
corrected two-pair cubic/quintic/septic package is analytic on small two-atom
data and satisfies swap symmetry, one-amplitude collapse, and diagonal merger
\[
\mathfrak P_2(a_1,h;a_2,h)=(a_1+a_2)F_h.
\]
The weaker exact route is
Proposition~\ref{prop:conditional-exact-two-point-exclusion-qabs}, which would
already remove exact close two-point bad cores on compact subsets of
$\{c\neq 0\}$ from continuity and diagonal collapse of the actual corrected
two-atom quotient germ.

What is now ruled out is the naive source-summed corrected-block route: its
extracted one-atom finite-order coefficients are even in the source weight and
therefore cannot supply the required linear merger law. On the mixed side,
determinant-only and quotient-only arguments are likewise shear-blind. The
genuinely new derivative-geometric lane is the exact projective
$\lambda$-shear / secant-shadow route on a good overlap patch.

By
Proposition~\ref{prop:conditional-mixed-lambda-shear-projective-shadow-reduction},
a close exact mixed two-point bad branch can accumulate only on
\[
\{S=0\}\cup\{K_v=0\},
\qquad
S=\frac{\Delta_7}{c\,v_5}.
\]
By
Proposition~\ref{prop:off-septic-zero-Kv-accumulation-is-mixed-low-order},
the part of $K_v=0$ with $S\neq 0$ is already absorbed by the existing
mixed low-order obstruction package. By
Corollary~\ref{cor:fixed-shear-simple-septic-zero-off-Kv-transfers-to-nonmixed-four-point},
the conditional exact fixed-shear overlap lane on
\[
\{S=0\}\cap\{K_v\neq 0\}
\]
transfers into the established nonmixed $4$-point geometry.

So the only genuinely residual exact fixed-shear mixed corner is
\[
\{S=0\}\cap\{K_v=0\}.
\]
There the branch is quotient-diagonal / involutive, the quartic--sextic first
rung is already quotient-injective, and the remaining exact fixed-shear burden
is now best viewed as split into two sublanes.

The \emph{close} exact fixed-shear lane is now scalar-local at theorem level:
by
Corollary~\ref{cor:local-scalar-bud-decomposition-close-fixed-shear}, on a
compact nonconstant quotient branch the close exact repeated-value carrier is
a finite disjoint union of local fold half-arcs, one per odd zero of
\[
\widehat{\mathcal D}=\widehat Q',
\]
with even zeros sterile; by
Remark~\ref{rem:constant-branch-scalarization-and-scalar-ceiling}, the
constant-quotient branch reduces, after affine gauge, to repeated values of a
single planar scalar and also contributes only local fold half-arcs, with
Corollary~\ref{cor:quartic-sextic-quotient-injectivity-local-no-reentry}
ruling out the $\widehat{\mathcal D}(0)\neq 0$ basepoint. So this lane no
longer admits a connected nonlocal close carrier; the remaining intrinsic
obstruction there is equality of consecutive lobe amplitudes across these
local fold half-arcs.

The \emph{three-run overlap} fixed-shear lane should still be read
conditionally as an inverse-branch-gap transport problem on an actual
outer-overlap interval, namely in the regime
\[
B_0+B_2>B_1.
\]
On such an interval one has outer inverse branches and a gap function
\[
\Delta(q)=\phi_2(q)-\phi_0(q).
\]
Any interior critical point of $\Delta$ then yields a same-value / same-slope
witness, and the current endgame fork is: either the transport-curvature scalar
fails to descend across the outer runs, giving the first explicit nonlocal
transport obstruction, or descent holds and the outer shells rigidify to exact
translated replication. What is not yet settled at theorem level is an
unconditional activation theorem from open overlap alone.

Thus the remaining two-point work is now concentrated in three very specific
targets: finite-order source-level merger, quotient-diagonal
continuity/diagonal collapse, and a transport/package theorem on the three-run
overlap fixed-shear lane; the scalar-local ceiling on the close exact
fixed-shear lane is already captured by
Corollary~\ref{cor:local-scalar-bud-decomposition-close-fixed-shear}.

\item
\textbf{Three-point branch (derivative-geometric layer now stabilized).}
The lifted three-point side is now the most mature local branch. Shrinking
exact lifted $3$-point circuits are tied to the osculating determinant and
hence to
\[
K_v=0,
\]
the local branch structure at finite-order $K_v$-zeros is completely
classified by parity, and the odd-order branch is blockwise globally rigid on
good $v_5$-patches: infinite-order flatness is excluded, the branch continues
uniquely on the maximal right monotone-slope block, and there is no interior
common-tangent recurrence or source-side return. Inside the finite-core
program this stabilizes the lifted $k=3$ collinearity layer; any residual
issue there is only boundary exit of the relevant source/target slope-overlap
block, not an internal three-point propagation mechanism.

\item
\textbf{Four-point branch (mixed side now captured by scalar exits).}
On the nonmixed side, the exact-contact and repeated-source / first-loss
packages are now organized by the scalar
\[
\mathcal R=\frac{T_v}{K_v^2},
\]
with one-sign intervals sterile and interval-persistent exact contacts forcing
affine-planar collapse. On the mixed side, the centered-base local route is
exhausted up to the clean branch, and the source-coupled and patch-loss
frontiers are now reduced to lower-core / shifted-$3$-point /
already-quintic handoffs. Accordingly, the mixed-global four-point side is no longer a search for new
local mixed geometry: by
Theorem~\ref{thm:global-mixed-to-scalar-capture-clean-patch}, a genuinely mixed
clean branch on a nonplanar good patch cannot persist entirely inside
\[
\{T_v\neq 0\}\cap\{\mathcal R\neq 0\},
\]
except after entry into already-absorbed regimes. So the four-point side is now
best viewed as captured by the centered mixed local theorem, the three-point
$K_v$-package, and the nonmixed $\mathcal R$-package, rather than as an
independent mixed-global carrier problem.

\item
\textbf{Late-$v_5$ degeneration is now only a handoff mechanism.}
If the $v_5$-patch fails and $v_5$ starts late, the valuation form of the
septic transport constraint forces matching late vanishing of the septic
transport channel $v_{7,K_1}$. Thus failure of the $v_5$-route is not a free
local cancellation but a coupled quintic--septic degeneration scheme, and in
the current architecture it should be read only as a handoff into lower-core or
already-quintic geometry, not as an independent mixed branch lane.

\item
\textbf{Four-point nonmixed branch via source-level control of odd $\mathcal W_{\mathrm{src}}$-zeros.}
The local nonmixed derivative geometry is now explicit: exact contacts force
zeros of $\mathcal R$, one-sign $\mathcal R$-intervals are sterile,
recurrent exact contacts force affine-planar collapse, and isolated zeros are
classified by parity. At source level, the same structure is encoded by the
canonical Wronskian
\[
\mathcal W_{\mathrm{src}}=J_\lambda J_{u_5}'-J_{u_5}J_\lambda',
\]
which differs from $\mathcal R$ only by a nonvanishing analytic factor on a
nonmixed arc.

For fixed source point $a$, repeated-source exact contacts are controlled by
the quadratic-envelope scalar
\[
\Gamma_a(x)=
\frac{Y(x)-Y(a)-Y'(a)(x-a)}
{S(x)-S(a)-S'(a)(x-a)},
\]
equivalently by the source quotient
\[
\Psi_a(h)=\frac{D_{u_5}(a,h)}{D_\lambda(a,h)},
\]
with derivative governed by the osculatory $4\times 4$ determinant
$\mathfrak D(a,h)$. By
Proposition~\ref{prop:first-loss-is-diagonal-odd-zero-birth},
this repeated-source / first-loss mechanism is not an independent obstruction:
it occurs if and only if one diagonally approaches an odd-order zero of
$\mathcal R$, equivalently of $\mathcal W_{\mathrm{src}}$.

The remaining theorem-level task on this side is therefore the
global/source-level exclusion or absorption of odd zeros of
$\mathcal W_{\mathrm{src}}$ on opened nonmixed arcs. In the unified mixed /
nonmixed picture, this is exactly the scalar $\mathcal R$-lane that absorbs the
interval-persistent $A=0$ collapse from the mixed side.
\end{enumerate}
\end{remark}

\begin{remark}[Stratified lifted local picture]
\label{rem:wip-stratified-lifted-local-picture}
The lifted local obstruction now splits into one closed regime and two genuinely
live regimes.

First, exact shrinking lifted $3$-point circuits concentrate into
\[
K_v=0.
\]
On good $v_5$-patches, this side is now blockwise globally controlled: there is no
infinite-order flatness, even-order $K_v$-zeros are sterile, odd-order
$K_v$-zeros seed exactly one local opposite-side branch, and that branch
continues uniquely on the maximal right monotone-slope block with no interior
common tangent, no source-side return, and no interior accumulation. So the
three-point side contributes only possible boundary exit, not an internal
recurrence mechanism.

Second, exact shrinking lifted $4$-point circuits concentrate into
\[
T_v=0.
\]
On the nonmixed side $K_v\neq 0$, the local derivative-side picture is now
explicit. Exact contacts force zeros of
\[
\mathcal R=\frac{T_v}{K_v^2},
\]
isolated zeros are classified by parity, one-sign $\mathcal R$-intervals are
exactly-contact sterile, and recurrent exact contacts force affine-planar
collapse. So the nonmixed $4$-point side is no longer a vague local
branching problem; what remains there is global/source-level control of which
odd zeros of $\mathcal R$ actually occur on genuine opened arcs.

Third, on the mixed side $K_v=T_v=0$, the low-order obstruction list is now
finite and explicit:
\[
(m,n)=(3,4),\ (3,6),\ (4,6),
\]
equivalently
\[
(r,s)=(1,2),\ (1,4),\ (2,5),
\]
equivalently
\[
(r,\operatorname{ord}_0\mathcal Z)=(1,1),\ (1,3),\ (2,2).
\]
At the pair level, however, the centered quintic fixed-sector package
\[
A_5^{\mathfrak f}(s)=u_5(s)I+v_5(s)S
\]
is now known not to vanish identically on the centered diagonal branch.
Near the centered base, the right-pincer package further reduces the mixed
local picture: every late-$v_5$ branch is either locally sterile on the
automatic $u_5$-patch or already absorbed into the planar quintic-circuit
regime, so the only genuinely new mixed local branch left there is the clean
branch
\[
\operatorname{ord}_s v_5=1.
\]

Thus the remaining lifted issues are now stratified much more sharply. On the
nonmixed side, the live question is the global/source-level control of odd zeros
of $\mathcal R$ on genuine opened arcs. On the mixed side, the centered local
issue is the finite scalar descent
\[
c(-s)=c(s)\pmod{s^6},
\]
while the clean-patch global side has already been compressed to the scalar
capture statement of
Theorem~\ref{thm:global-mixed-to-scalar-capture-clean-patch}. So the mixed
four-point branch should no longer be viewed as an independent open geometric
carrier problem away from the scalar exits.
\end{remark}

\begin{remark}[Current geometric bottleneck]
\label{rem:current-geometric-bottleneck}
The remaining geometric bottleneck is now split into two qualitatively
different layers.

On the one-pair geometric layer, the obstruction is no longer bulk $W=0$
geometry, and it is no longer local adjacent pair-pair birth across a generic
turning boundary. On the $W=0,\ M\neq 0$ side,
Proposition~\ref{prop:no-W-flat-interval} excludes open-interval flatness and
Theorem~\ref{thm:finite-order-w-zero-local-sterility} shows that every
finite-order isolated $W=0$ boundary is locally sterile for pair-pair birth.
On the turning-boundary side,
Theorem~\ref{thm:adjacent-pair-pair-exclusion-turning} removes actual adjacent
pair-pair lines unconditionally.

At the lifted order-$7$ level, the canonical derivative-geometric determinants
\[
K_v
\qquad\text{and}\qquad
T_v
\]
govern shrinking exact lifted $3$- and $4$-point circuits on
$v_5\neq 0$ patches. On the $3$-point side, this layer is stabilized: shrinking
exact circuits concentrate into $\{K_v=0\}$, finite-order isolated zeros are
classified by parity, and the odd-order branch is blockwise globally rigid on
good $v_5$-patches.

On the $4$-point side, exact shrinking lifted $4$-point circuits concentrate
into
\[
\{T_v=0\},
\]
again with uniform small-diameter exclusion away from that zero set.

On the nonmixed side $K_v\neq 0$, generic odd-order $T_v$-zeros can still
carry preexisting close branches, so the nonmixed branch remains genuinely
live and now asks for global/source-level control of which odd zeros of
\[
\mathcal R=\frac{T_v}{K_v^2}
\]
actually arise on opened arcs.

By contrast, on the mixed side $K_v=T_v=0$, the local geometry near the
centered base is no longer the main bottleneck. There the right-pincer package
shows that every late-$v_5$ branch is either locally sterile on the automatic
$u_5$-patch or already absorbed into the defect-free planar quintic-circuit
regime, so the only genuinely new mixed local branch left is the clean branch
\[
V_1\neq 0.
\]
Accordingly, the mixed bottleneck is no longer a diffuse local codimension-two
analysis. At the local level it is the finite scalar descent
\[
c(-s)=c(s)\pmod{s^6}
\]
on that clean branch; at the global level, after the source-coupled and
patch-loss absorptions, the clean-patch mixed problem has already compressed to
the scalar alternatives
\[
T_v=0
\qquad\text{or}\qquad
\mathcal R=0,
\]
as formalized in
Theorem~\ref{thm:global-mixed-to-scalar-capture-clean-patch}.

So the missing lifted ingredient is no longer another generic local geometric
exclusion theorem. On the mixed side, the centered-base local package, the
source-coupled exclusion, the patch-loss reductions, and the nonmixed
$\mathcal R$-handoff now combine to give the clean theorem-level conclusion
formalized below in
Theorem~\ref{thm:global-mixed-to-scalar-capture-clean-patch}: every maximal
genuinely mixed clean branch on a nonplanar good patch is forced either into
the centered-base scalar exit
\[
T_v=0
\]
or into the nonmixed scalar exit
\[
\mathcal R=0,
\]
up to already-absorbed planar, lower-core, repeated-node, repeated-source, and
quintic-collapse regimes. Thus the still-live frontiers are no longer mixed
four-point geometry, but rather the nonmixed/source-level control of odd zeros
of $\mathcal R$ and the separate two-point program.
\end{remark}

\begin{remark}[Current state of the argument]
\label{rem:current-state-of-argument}
The local one-pair analysis now consists of the canonical triple
\[
(c(h),\,v_5(h),\,\Delta_7(h)),
\]
together with the derived quantities
\[
M(h)=c(h)v_5'(h)-c'(h)v_5(h),
\qquad
N(h)=c(h)\Delta_7'(h)-c'(h)\Delta_7(h),
\]
\[
L(h)=v_5(h)\Delta_7'(h)-v_5'(h)\Delta_7(h),
\qquad
W(h)=\det(F(h),F'(h),F''(h)).
\]

On the region where $v_5M\neq 0$, the projective graph variables
\[
\mu=\frac{c}{v_5},
\qquad
\lambda=\frac{\Delta_7}{v_5}
\]
identify the normalized curve with a graph $\lambda=\lambda(\mu)$, and
\[
\lambda_{\mu\mu}=\frac{v_5^3}{M^3}W,
\qquad
\beta:=\lambda_\mu=-\frac{L}{M},
\qquad
T(h)=\left(-\frac{L(h)}{M(h)},\frac{N(h)}{M(h)}\right).
\]
Thus the global line-incidence problem is equivalently a problem about the
recurrence of tangent slopes $\beta$, or about vertical-line re-entry on the
dual tangent-line curve $T$.

Locally, both generic boundary types are now under control. At finite-order
$W=0,\ M\neq 0$ boundaries there is no local pair-pair birth, while at
generic turning boundaries actual adjacent pair-pair lines are excluded
unconditionally. The stronger symmetric directional-kernel theorem remains
available only conditionally on zero-capture and now serves as a sharper but
nonessential alternative route.

Accordingly, the equivalent slope/tangent/dual direct-slice formulations remain
the correct geometric dictionary, but they are no longer themselves the live
theorem target. The explicit direct provenance formula gives
\[
\phi=A\,q^{(7)}
\qquad
(A=\kappa\,\eta_2),
\]
and on any generic one-pair block this direct slice is fully universal
(Proposition~\ref{prop:full-block-universality-direct-provenance-slice}).
Moreover, even after imposing any finite family of linear provenance or
matching conditions on $q$, direct-slice recurrence geometry remains formally
realizable (Corollary~\ref{cor:finite-linear-provenance-too-weak}).

Thus the remaining gap is not an abstract global no-reentry theorem for the
scalar graph $\phi$ by itself. It is the search for a genuinely non-universal
provenance law for $q$, or an external mixed/canonical input, that cuts down
the allowed direct-slice graphs on an entire block.

On the lifted $4$-point side, the mixed low-order obstruction list is now both
finite and canonical:
\[
(r,\operatorname{ord}_0\mathcal Z)\in\{(1,1),(1,3),(2,2)\},
\]
equivalently
\[
(r,s)=(1,2),(1,4),(2,5),
\]
or
\[
(m,n)=(3,4),(3,6),(4,6).
\]
The centered Taylor package already provides the explicit coefficients
\[
U_j,\ V_j,\ \widehat U_j,\ \widehat V_j,\ D_2,\ D_4,\ D_6,
\]
and the tagged quintic witness now shows that
\[
A_5^{\mathfrak f}(s)\not\equiv 0
\]
on the centered diagonal branch.

On the mixed side, however, the local theorem picture is now sharper than a
mere punctured-arc patch dichotomy, and the global side is correspondingly
compressed. Near the centered base, the late-$v_5$
sector is already removed by the $u_5$-patch right-pincer package as either
locally sterile or planar-lift and hence already absorbed into the
defect-free quintic-circuit regime. The only genuinely new mixed local branch
there is the clean branch $V_1\neq 0$, and that branch is excluded by the
canonical scalar descent together with the centered good-$v_5$-patch parity
package. Moreover, after the source-coupled and patch-loss reductions, the
sufficiently small mixed-global problem is no longer a search for a new mixed
carrier lane. At the theorem level, the clean small-branch mixed problem has
already compressed to the scalar alternatives
\[
T_v=0
\quad\text{or}\quad
\mathcal R=0.
\]
Here the $K_v=0$ side no longer survives as an independent mixed carrier on the
clean small-branch four-point side, but on the additional conditional exact
fixed-shear overlap slice it still leaves a narrower residual two-point burden.
That residual slice is not a new mixed local carrier: modulo constant shear,
its pointwise quotient germ is already governed by
\[
\mathcal D=\frac{\mathcal W_{\mathrm{src}}}{J_\lambda^2},
\]
or equivalently on the involution quotient by
\[
\widehat{\mathcal D}(u)=\widehat Q'(u).
\]
Moreover, by
Corollary~\ref{cor:one-sided-exclusion-repeated-node-tangent-contact},
the ordinary local repeated-node tangent-contact mechanism on a good
$v_5$-patch is necessarily opposite-side, so this quotient-diagonal residual
lane cannot locally merge with the standard lifted $3$-point branch geometry.
Thus the genuinely mixed small-branch problem is best read as mixed-to-scalar
capture together with a separate quotient-diagonal residual transport problem,
not as another free mixed carrier geometry.

In parallel, the finite-core endgame now splits into three layers.

First, on the two-pair side, the strengthened invariant
\[
\widehat\Psi(h)=
\left(
\frac{u_5(h)}{c(h)},
\frac{v_5(h)}{c(h)},
\frac{\Delta_7(h)}{c(h)^2}
\right)
\]
still gives the correct algebraic bridge, and on good $M\neq 0$ patches the
quantitative two-pair coincidence theorem forces bad two-point cores into a
defect-scale close regime. The minimal finite-order source criterion
Lemma~\ref{lem:minimal-source-level-two-point-criterion} still isolates the
exact swap / one-amplitude / diagonal-merger package that would imply the
desired quadratic interaction factorization. However, the naive source-summed
corrected block model does not prove this criterion: along a weighted one-atom
path its extracted finite-order coefficients are even analytic in the source
weight, so the required linear merger law cannot come from that model alone.

At present the two-point side is better organized by the three-front
compression of
Remark~\ref{rem:current-three-front-compression-two-point-endgame}
than by the older language of ``three conditional lanes.''

The finite-order source-level merger criterion of
Lemma~\ref{lem:minimal-source-level-two-point-criterion}
remains one possible local route on good patches, but the genuinely residual
mixed burden is now the quotient-diagonal exact fixed-shear corner
\[
(S=0,\ K_v=0).
\]
There the projective $\lambda$-shear compactness and secant-shadow analysis
identify the correct residual geometry, while the collision-cancellation chart,
the no-finite-local-rescue picture, and the quotient-transport package show
that the remaining issue is not a missing local branch taxonomy.

So the live downstream two-point targets are now the following: a good-patch
edge-law theorem for the actual corrected cubic/quintic defects in the
collision-cancellation chart; a quotient-visible transport/state-local closure
on the exact fixed-shear corner; or a provenance-sensitive package theorem
strong enough to force coincidence/functoriality at the
collision/cancellation locus. In particular, the current residual burden should
be read as local blow-up regularity plus quotient transport plus package-level
coincidence, not as a fourth independent mixed lane.

Second, on the lifted three-point side, shrinking exact circuits now
concentrate into $K_v=0$, and this concentration is explicitly
derivative-geometric: by the osculating-determinant expansion, repeated-node
limits become tangent-line contacts for the planar lifted curve
\[
\Lambda_v=(x,S).
\]
On good $v_5$-patches this side is now blockwise globally controlled:
$K_v$ has no infinite-order zeros; even-order zeros are sterile; odd-order
zeros carry exactly one local branch; and that branch continues uniquely on
the maximal right monotone-slope block with no interior common-tangent event
or source-side return. So the three-point side is reduced to boundary exit of
the relevant slope-overlap block, not an internal recurrence mechanism.
Inside the finite-core reduction, this is exactly the stabilization of the
lifted $k=3$ collinearity layer.

Third, on the lifted four-point side, shrinking exact circuits now
concentrate into $T_v=0$, repeated-node limits become tangent-plane
contacts, and the local geometry splits into a nonmixed and a mixed regime.

On the nonmixed side, the local derivative-side theory is now explicit:
exact contacts force zeros of $\mathcal R=T_v/K_v^2$, isolated zeros are
classified by parity, one-sign $\mathcal R$-intervals are exactly-contact
sterile, and recurrent exact contacts force affine-planar collapse. At source
level, the same structure is encoded by the gauge-invariant Wronskian
\[
\mathcal W_{\mathrm{src}}=J_\lambda J_{u_5}'-J_{u_5}J_\lambda',
\]
which differs from $\mathcal R$ only by a nonvanishing analytic factor on a
nonmixed arc. For fixed source point $a$, repeated-source exact contacts
factor through the quadratic-envelope scalar
\[
\Gamma_a(x)=
\frac{Y(x)-Y(a)-Y'(a)(x-a)}
{S(x)-S(a)-S'(a)(x-a)},
\]
equivalently through the source quotient
\[
\Psi_a(h)=\frac{D_{u_5}(a,h)}{D_\lambda(a,h)},
\]
with derivative governed by the osculatory source determinant
$\mathfrak D(a,h)$. So the remaining nonmixed burden is global/source-level
control of which odd zeros of $\mathcal W_{\mathrm{src}}$, equivalently of
$\mathcal R$, occur on a genuine opened nonmixed arc.

On the mixed side, the dangerous low-order obstruction list is the finite
canonical set
\[
\kappa_{1,1},\qquad \kappa_{1,3},\qquad \kappa_{2,2},
\]
equivalently the order patterns
\[
(r,\operatorname{ord}_0\mathcal Z)=(1,1),\ (1,3),\ (2,2).
\]
The remaining theorem-level task there is no longer local patch classification,
and on the clean-patch side it is no longer an open global capture problem
either: the sufficiently small mixed-global closure is now formalized by
Theorem~\ref{thm:global-mixed-to-scalar-capture-clean-patch}. Thus the mixed
side should now be read as split into two parts: the centered clean-branch
scalar descent
\[
c(-s)=c(s)\pmod{s^6},
\]
which is the remaining local centered-base task, and the already-organized
mixed-to-scalar handoff showing that outside absorbed regimes every maximal
genuinely mixed clean branch is forced into
\[
T_v=0
\qquad\text{or}\qquad
\mathcal R=0.
\]
So any genuinely new input beyond the present lifted-variable package would now
have to be RH-specific source-/matrix-level structure on the scalar side, not a
new mixed local geometry theorem.

Moreover, Theorem~\ref{thm:local-exhaustion-mixed-four-point-centered-base}
shows that no new sufficiently small local mixed $4$-point obstruction remains
near the centered diagonal base, so any external mixed input capable of
constraining the direct-slice/provenance lane must be genuinely global/nonlocal
rather than a new local mixed septic-jet pattern.
\end{remark}

\begin{remark}[Post-draft findings: established, provisional, and live]
\label{rem:post-draft-findings-established-provisional-live}
This remark summarizes ledger findings that are now partly formalized in
the surrounding propositions, corollaries, and theorems, and partly
retained as live WIP for the remaining package/frontier obligations.

First, on the two-point branch, the naive source-summed corrected block model
still cannot prove the exact finite-order merger criterion: along a weighted
one-atom path the extracted finite-order coefficients are even analytic in the
source weight, so the needed linear merger law does not come from that model
alone.

At the same time, the residual mixed picture is now much cleaner. The overlap
bad set is fully split:
\[
(K_v=0,\ S\neq 0)
\]
is already absorbed by the mixed low-order package,
\[
(S=0,\ K_v\neq 0)
\]
is already transferred into the established nonmixed $4$-point geometry,
and only the exact fixed-shear corner
\[
\{S=0\}\cap\{K_v=0\}
\]
remains genuinely live.

On that residual corner the branch is quotient-diagonal / involutive,
\[
\Phi(t)=\Phi(-t),
\qquad
Q(t)=\widehat Q(t^2),
\qquad
\widehat{\mathcal D}(u)=\widehat Q'(u).
\]
Moreover,
Proposition~\ref{prop:finite-transport-jet-quotient-descent-fixed-shear-corner}
shows that every finite parity-normalized transport jet descends to the
quotient variable $u=t^2$, while
Corollary~\ref{cor:rank-one-transport-state-quartic-sextic-fixed-shear}
shows that on the quartic--sextic rung the finite quotient-visible transport
state has local rank one, and
Corollary~\ref{cor:no-finite-order-hidden-reset-quartic-sextic-fixed-shear}
rules out any third local reset inside that descended state.

So the residual exact fixed-shear burden is no longer a search for another
local mixed carrier or hidden finite-order transport variable. It has become a
problem of collision-cancellation regularity, quotient transport, and
package/provenance-sensitive coincidence. In particular, any still-live
obstruction must come from the first genuinely relational or
provenance-sensitive two-atom input that survives after quotient-visible
transport descent, or from non-finite-order dependence.

Moreover, by
Corollary~\ref{cor:one-sided-exclusion-repeated-node-tangent-contact},
every sufficiently small off-diagonal repeated-node tangent-contact branch on a
good $v_5$-patch is necessarily opposite-side. Thus the quotient-diagonal
exact fixed-shear residual lane cannot locally merge with the ordinary lifted
$3$-point branch mechanism.

There is now also a sharper internal split of the exact fixed-shear residual
lane itself, and one side of that split has now been promoted to theorem
level. On the close exact side, the scalar-local compression is no longer a
work-in-progress expectation: by
Corollary~\ref{cor:local-scalar-bud-decomposition-close-fixed-shear}, on a
compact nonconstant quotient branch the close exact repeated-value carrier is
a finite disjoint union of local fold half-arcs, one per odd zero of
\[
\widehat{\mathcal D}=\widehat Q',
\]
with even zeros sterile; by
Remark~\ref{rem:constant-branch-scalarization-and-scalar-ceiling}, the
constant-quotient branch also reduces, after affine gauge, to repeated values
of a single planar scalar and contributes only local fold half-arcs. So the
close exact carrier no longer admits a connected nonlocal close carrier, and
the remaining intrinsic obstruction on that sublane is equality of
consecutive lobe amplitudes across these local fold half-arcs.

On the three-run overlap side, by contrast, the endgame is still conditional.
The residual overlap regime is
\[
B_0+B_2>B_1,
\]
not
\[
B_1>B_0+B_2.
\]
In that actual outer-overlap regime, the quotient-visible problem may be
reduced to the inverse-branch-gap transport problem on the two outer runs.
Writing
\[
\Delta(q):=\phi_2(q)-\phi_0(q),
\]
any interior critical point of $\Delta$ gives a same-value / same-slope
witness, and the live endgame fork becomes: either the transport-curvature
scalar fails to descend across the outer runs, giving the first explicit
nonlocal transport obstruction, or descent holds and the outer shells rigidify
to exact translated replication. An unconditional activation theorem from open
overlap alone has not yet been established. So the residual exact fixed-shear
lane is no longer a vague hidden branch problem; it has split into a close
scalar-local lane, now controlled at theorem level by
Corollary~\ref{cor:local-scalar-bud-decomposition-close-fixed-shear}, and a
separate three-run overlap-transport lane that remains the live conditional
target.

Second, on the lifted four-point branch, the explicit low-order obstruction
list has been reduced to the six coefficients
\[
z_2,\ z_4,\ z_6,\ \kappa_{1,1},\ \kappa_{1,3},\ \kappa_{2,2},
\]
or equivalently to the order patterns
\[
\operatorname{ord}_0\mathcal R\in\{1,3,5\},
\qquad
(r,\operatorname{ord}_0\mathcal Z)\in\{(1,1),(1,3),(2,2)\}.
\]

Third, the centered symmetric-placement Taylor package now gives explicit septic
coefficients
\[
D_2,\qquad D_4,\qquad D_6,
\]
explicit bridge coefficients
\[
X_1,X_3,X_5,\qquad L_1,L_3,L_5,
\]
and a punctured-neighborhood exclusion at the diagonal base.

Fourth, the centered fixed-sector quintic package
\[
A_5^{\mathfrak f}(s)=u_5(s)I+v_5(s)S
\]
is now firmly established as the correct pair-level object on the mixed side.
The centered tagged witness survives specialization strongly enough to show
\[
U_1\neq 0,
\qquad
u_5(s)\not\equiv 0,
\qquad
A_5^{\mathfrak f}(s)\not\equiv 0,
\]
so the centered quintic pair cannot totally collapse.

Fifth, the right-pincer package now sharpens the mixed local picture near the
centered base: every late-$v_5$ branch is either locally sterile on the
automatic $u_5$-patch or already absorbed into the planar defect-free
quintic-circuit regime. So the only genuinely new mixed local branch left there
is the clean branch
\[
V_1\neq 0.
\]

Sixth, on that clean branch the remaining theorem has become very small:
the finite scalar descent
\[
c(-s)=c(s)\pmod{s^6},
\qquad\text{equivalently}\qquad
C_1=C_3=C_5=0.
\]
Once this holds, the centered good-$v_5$-patch bridge becomes formal:
\[
Y=\frac{u_5}{c},
\qquad
S=\frac{\Delta_7}{cv_5}
\]
are odd through the required order, and the entire explicit mixed low-order
bad list
\[
(r,\operatorname{ord}_0\mathcal Z)\in\{(1,1),(1,3),(2,2)\}
\]
is excluded there.

Seventh, the relation between the mixed and nonmixed $4$-point sides is now
structurally sharper. On a good $v_5$-patch, an exact mixed relation
\[
A\,Y+B\,x+C\,S+D=0
\]
with
\[
Y=\frac{u_5}{c},
\qquad
x=\frac{v_5}{c},
\qquad
S=\frac{\Delta_7}{c\,v_5}
\]
and $A=0$ is already exact affine-planarity for the planar curve $(x,S)$; by
Remark~\ref{rem:A-zero-is-nonmixed-handoff}, this is precisely the nonmixed
$\mathcal R$-lane. More generally, the sheared four-zero reduction shows that
every exact genuinely mixed $4$-point circuit on a good $v_5$-patch already
forces either
\[
K_v=0
\qquad\text{or}\qquad
\mathcal R=0
\]
in its convex hull. So the mixed-global problem should now be viewed as a
mixed-to-scalar capture problem rather than as a separate branch geometry.

Eighth, the source-coupled and patch-loss frontiers have both been demoted to
handoff mechanisms. In the sufficiently small regime, the source-coupled
lower-core walls
\[
D_\lambda=0,
\qquad
\mathfrak D=0
\]
are not pass-through points. On the patch side, the quintic projective lifts
\[
\widetilde Q_v=[c v_5:u_5v_5:v_5^2:\Delta_7],
\qquad
\widetilde Q_u=[c u_5:u_5^2:u_5v_5:-\Delta_7]
\]
show that single-chart failures are removable after recharting, while the
genuine common boundary
\[
u_5=v_5=0
\]
drops either to pure quintic-direction collapse or to lower-core $3$-point /
repeated-node quotient geometry. Thus
\[
cMv_5=0
\]
should be read as lower-core handoff, not as a fourth independent mixed-global
lane.

Ninth, after the centered-base local exhaustion theorem and the preceding
handoffs are combined, the sufficiently small nonplanar mixed branch has no
independent carrier away from
\[
T_v=0
\quad\text{and}\quad
\mathcal R=0.
\]
Equivalently, there is no sufficiently small genuinely mixed $4$-point branch
contained entirely in
\[
\{T_v\neq 0\}\cap\{\mathcal R\neq 0\}
\]
on a nonplanar good patch, except after entry into already-absorbed
affine-planar or lower-core regimes.

So the present frontier is no longer a generic local geometry problem. On the
residual exact fixed-shear overlap slice
\[
K_v=0,
\]
the pointwise quotient structure has already compressed to the single
transport scalar
\[
\widehat{\mathcal D}(u)=\widehat Q'(u),
\qquad
u=t^2.
\]
The quartic--sextic rung is the first explicit case, and there the quotient
branch is locally one-sheeted. Thus the residual $K_v=0$ lane is best viewed
as a quotient-diagonal / involutive transport problem rather than as a new
mixed local carrier.

At the same time, this pointwise compression does \emph{not} yet prove that the
full transported overlap / critical-value state is quotient-visible. What is
proved inside the quotient-visible exact fixed-shear reduction is narrower and
sharper: on a connected nonconstant quotient branch,
Proposition~\ref{prop:sign-run-compression-three-lobe-criterion-fixed-shear-quotient-branch}
shows that nonadjacent re-entry of regular values can arise only through a
three-run lobe-amplitude pattern, and
Corollary~\ref{cor:three-run-no-reentry-classification-fixed-shear-quotient-branch}
identifies the only surviving inequality
\[
B_1\ge B_0+B_2.
\]
Moreover, on the first live quartic--sextic rung,
Corollary~\ref{cor:quartic-sextic-quotient-injectivity-local-no-reentry}
shows that the basepoint is locally one-sheeted in the quotient variable and
does not itself contribute such a lobe obstruction. Accordingly, the live
two-point burden on the $K_v=0$ slice is now a no-hidden-reset /
state-visibility problem for the actual corrected overlap state, together with
the alternative source-level merger route and the already-separated higher-core
issues.
\end{remark}

\begin{remark}[Two-point burden after mixed four-point closure]
\label{rem:two-point-burden-after-mixed-four-point-closure}
With the clean mixed four-point side now compressed to the scalar exits
\[
T_v=0
\qquad\text{or}\qquad
\mathcal R=0,
\]
the mixed two-point side should no longer be read as an independent frontier of
mixed local geometry.

The overlap bad set has already split as follows:
\[
(K_v=0,\ S\neq 0)
\]
is absorbed by the mixed low-order obstruction package,
\[
(S=0,\ K_v\neq 0)
\]
is transferred into the established nonmixed $4$-point geometry,
and only the exact fixed-shear residual corner
\[
\{S=0\}\cap\{K_v=0\}
\]
remains genuinely live.

On that residual corner, the quotient-visible transport package is now much
more rigid than a single pointwise scalar. By
Proposition~\ref{prop:finite-transport-jet-quotient-descent-fixed-shear-corner},
every finite parity-normalized transport jet descends to the quotient variable
$u=t^2$. On the quartic--sextic rung,
Corollary~\ref{cor:rank-one-transport-state-quartic-sextic-fixed-shear} shows
that every finite quotient-visible transport state is locally a function of
$Q$ alone, and
Corollary~\ref{cor:no-finite-order-hidden-reset-quartic-sextic-fixed-shear}
shows that no third local reset survives there.

So the residual exact fixed-shear corner is no longer best viewed as an
unresolved mixed local curvature-branch problem. If the actual corrected
defect/package is state-local in the descended transport state, then the corner
closes by
Proposition~\ref{prop:conditional-state-local-defect-closure-fixed-shear-corner}
or even by
Corollary~\ref{cor:stronger-state-local-package-collapse-fixed-shear-corner}.
Accordingly, any still-live input on the two-point side would have to come from
the first genuinely relational or provenance-sensitive two-atom datum that
survives after quotient-visible transport descent, or else from non-finite-order
dependence.
\end{remark}

\begin{remark}[Residual mixed two-point lane: structural gain and present limitation]
\label{rem:residual-mixed-two-point-lane-structural-gain-limitation}
The present two-point picture now has three genuine structural gains on the
exact fixed-shear residual corner.

First, after the off-$K_v=0$ absorption results and the transfer of the
simple-septic-zero slice into the nonmixed $4$-point theory, the
derivative-geometric mixed residue is sharply concentrated on the exceptional
corner
\[
\{S=0\}\cap\{K_v=0\}.
\]

Second, by
Proposition~\ref{prop:finite-transport-jet-quotient-descent-fixed-shear-corner},
every finite parity-normalized transport jet is quotient-visible in the
variable
\[
u=t^2.
\]
So the residual corner is no longer a free two-sided $t$-state even at the
finite transport-jet level.

Third, on the quartic--sextic rung,
Corollary~\ref{cor:rank-one-transport-state-quartic-sextic-fixed-shear} shows
that this descended finite transport state has local rank $1$, and
Corollary~\ref{cor:no-finite-order-hidden-reset-quartic-sextic-fixed-shear}
shows that no finite-order hidden reset survives beyond the involution pair
$t\leftrightarrow -t$.

What is \emph{not} yet proved is that the actual corrected two-pair defect or
package is already state-local in this descended transport state. If that
state-locality held, then the residual branch would already be excluded by
Proposition~\ref{prop:conditional-state-local-defect-closure-fixed-shear-corner},
or would collapse identically by
Corollary~\ref{cor:stronger-state-local-package-collapse-fixed-shear-corner}.

Thus the live issue on the mixed residual side is now no longer local
classification and no longer hidden finite-order transport dynamics. The
remaining question is whether the first surviving obstruction is already
quotient-visible, or instead comes from genuinely relational /
provenance-sensitive two-atom structure or from non-finite-order dependence.
\end{remark}

\begin{proposition}[$u_5$-chart exceptional-point handoff]
\label{prop:u5-chart-exceptional-point-handoff}
Work near a collapsing configuration with $u_5\neq 0$, and use the projective
lift
\[
\widetilde Q_u=[c u_5:u_5^2:u_5v_5:-\Delta_7].
\]
Let
\[
q_\nu=[0:1:0:\nu]
\]
be an exceptional limiting point on the exceptional line. Then every limiting
hyperplane through $q_\nu$ induces the exact lower-core relation
\[
D\,c+B\,v_5+C(\lambda-\nu u_5)=0,
\qquad
\lambda=-\frac{\Delta_7}{u_5}.
\]
Accordingly:
\begin{itemize}
\item one-node collapse hands off to exact shifted lower-core $3$-point
geometry;
\item two-node collapse hands off to a repeated-node lower-core wall;
\item if $u_5=0$ as well, one is on the common quintic-collapse locus already
treated separately.
\end{itemize}
So the $u_5$-chart side of patch loss contributes no independent mixed-global
carrier lane.
\end{proposition}

\begin{remark}[Shifted lower-core invariance on the $u_5$-side]
\label{rem:shifted-lower-core-invariance-u5}
The shift
\[
\lambda\longmapsto \lambda-\nu u_5
\]
does not alter the derivative-geometric lower-core package: it changes only the
zeroth-order affine offset, not the relevant repeated-node/tangent-contact
derivative structure. Thus the handoff in
Proposition~\ref{prop:u5-chart-exceptional-point-handoff} is genuinely into the
existing lower-core package, not into new geometry.
\end{remark}

\begin{proposition}[Local clean continuation off the explicit walls]
\label{prop:local-clean-continuation-off-explicit-walls}
Work on a nonplanar good $v_5$-patch and write
\[
Y=\frac{u_5}{c},\qquad x=\frac{v_5}{c},\qquad S=\frac{\Delta_7}{c\,v_5}.
\]
Let
\[
\Phi_\gamma:=Y-\gamma S,
\qquad
F_{\alpha,\beta,\gamma}(x):=\Phi_\gamma(x)-\alpha-\beta x.
\]
Suppose
\[
h_1<h_2<h_3<h_4
\]
is an exact genuinely mixed clean $4$-point configuration, so
\[
F_{\alpha,\beta,\gamma}(x(h_i))=0
\qquad (i=1,2,3,4),
\]
the four roots $x(h_i)$ are distinct and simple, and the carrying hyperplane is
genuinely mixed ($A\neq 0$). Assume moreover that the configuration does not
lie on any of the walls
\[
A=0,\qquad D_\lambda=0,\qquad \mathfrak D=0,\qquad cMv_5=0,\qquad
T_v=0,\qquad \mathcal R=0.
\]
Then this configuration is an interior point of a smooth local solution
manifold of exact genuinely mixed clean $4$-point configurations. In
particular, it cannot be a terminal point of a maximal clean mixed branch.

\begin{proof}
For
\[
\mathcal F_i(h_1,h_2,h_3,h_4,\alpha,\beta,\gamma)
:=
Y(h_i)-\gamma S(h_i)-\alpha-\beta x(h_i),
\]
the Jacobian with respect to $(h_1,h_2,h_3,h_4)$ is diagonal:
\[
\frac{\partial \mathcal F_i}{\partial h_j}
=
\delta_{ij}\bigl(Y'(h_i)-\gamma S'(h_i)-\beta x'(h_i)\bigr).
\]
Since
\[
Y'(h_i)-\gamma S'(h_i)-\beta x'(h_i)
=
x'(h_i)\,F'_{\alpha,\beta,\gamma}(x(h_i)),
\]
the diagonal entries are nonzero on a good $v_5$-patch and at simple roots.
So the Jacobian is invertible, and the claim follows from the implicit
function theorem.
\end{proof}
\end{proposition}

\begin{corollary}[Positive-scale endpoints lie on an absorbed wall or a scalar exit]
\label{cor:positive-scale-endpoint-absorbed-or-scalar}
Let $\Gamma$ be a maximal connected genuinely mixed clean $4$-point branch on a
nonplanar good $v_5$-patch. Suppose $p\in\overline{\Gamma}$ is a terminal point
at which the relevant node scale stays bounded below by a positive constant.
Then $p$ lies on one of
\[
A=0,\qquad D_\lambda=0,\qquad \mathfrak D=0,\qquad cMv_5=0,\qquad
T_v=0,\qquad \mathcal R=0.
\]
In particular, after the previously established handoffs are invoked, a
positive-scale terminal point cannot furnish a new genuinely mixed continuation
mechanism: it must either lie in the centered-base capture set $\{T_v=0\}$,
lie in the nonmixed affine-planar scalar set $\{\mathcal R=0\}$, or already
belong to one of the absorbed planar/lower-core/quintic-collapse regimes.
\end{corollary}

\begin{proof}
If $p$ lay off the listed wall set, then the limiting exact genuinely mixed clean
configuration at $p$ would satisfy the hypotheses of
Proposition~\ref{prop:local-clean-continuation-off-explicit-walls}. Hence the
branch would extend smoothly through $p$, contradicting maximality. Therefore a
positive-scale terminal point must lie on one of
\[
A=0,\qquad D_\lambda=0,\qquad \mathfrak D=0,\qquad cMv_5=0,\qquad
T_v=0,\qquad \mathcal R=0.
\]
The final sentence is then exactly the previously established absorption
package: $A=0$ is the nonmixed affine-planar handoff
(Remark~\ref{rem:A-zero-is-nonmixed-handoff}), the source-coupled walls
$D_\lambda=0$ and $\mathfrak D=0$ are locally non-carrying in the sufficiently
small regime (Remark~\ref{rem:post-draft-findings-established-provisional-live}),
and patch loss $cMv_5=0$ is absorbed into lower-core / shifted-$3$-point /
quintic-collapse geometry
(Proposition~\ref{prop:u5-chart-exceptional-point-handoff} and the surrounding
post-draft findings package).
\end{proof}

\begin{proposition}[Small-scale non-scalar exits are absorbed]
\label{prop:small-scale-non-scalar-exits-absorbed}
Let $\Gamma$ be a sufficiently small genuinely mixed clean $4$-point branch on a
nonplanar good $v_5$-patch. Assume
\[
\Gamma\cap\{T_v=0\}=\varnothing,
\qquad
\Gamma\cap\{\mathcal R=0\}=\varnothing.
\]
Then $\Gamma$ cannot continue through any of the non-scalar exit walls
\[
A=0,\qquad D_\lambda=0,\qquad \mathfrak D=0,\qquad cMv_5=0
\]
as an independent genuinely mixed clean branch.

More precisely:
\begin{enumerate}
\item if $A=0$ persists on an interval, then $\Gamma$ has already entered the
nonmixed affine-planar $\mathcal R$-package
(Remark~\ref{rem:A-zero-is-nonmixed-handoff},
Proposition~\ref{prop:interval-persistent-nonmixed-four-point-forces-affine-planarity},
Corollary~\ref{cor:recurrent-exact-nonmixed-contacts-force-affine-planarity});

\item if $D_\lambda=0$ or $\mathfrak D=0$ occurs in the sufficiently small
regime, then this lies on the source-coupled lower-core marked set, which is
locally non-carrying and not a pass-through mechanism
(Remark~\ref{rem:post-draft-findings-established-provisional-live});

\item if $cMv_5=0$ occurs, then the branch is absorbed into lower-core /
shifted-$3$-point / quotient repeated-node / quintic-collapse geometry
(Remark~\ref{rem:post-draft-findings-established-provisional-live} and
Proposition~\ref{prop:u5-chart-exceptional-point-handoff}).
\end{enumerate}

Accordingly, no sufficiently small genuinely mixed clean branch can remain an
independent mixed carrier while meeting only non-scalar exit walls.
\end{proposition}

\begin{proof}
The $A=0$ case is exactly the nonmixed affine-planar handoff by
Remark~\ref{rem:A-zero-is-nonmixed-handoff}; interval-persistent or recurrent
exact contacts are then absorbed by
Proposition~\ref{prop:interval-persistent-nonmixed-four-point-forces-affine-planarity}
and Corollary~\ref{cor:recurrent-exact-nonmixed-contacts-force-affine-planarity}.

The source-coupled walls $D_\lambda=0$ and $\mathfrak D=0$ are, in the
sufficiently small regime, only source-coupled lower-core marked points and are
locally non-carrying; see
Remark~\ref{rem:post-draft-findings-established-provisional-live}.

Finally, patch loss $cMv_5=0$ has already been reduced to lower-core /
shifted-$3$-point / quotient repeated-node / quintic-collapse handoffs; see
Remark~\ref{rem:post-draft-findings-established-provisional-live} and
Proposition~\ref{prop:u5-chart-exceptional-point-handoff}.

So none of the listed non-scalar walls can provide an independent genuinely
mixed clean continuation mechanism under the stated hypotheses.
\end{proof}

\begin{theorem}[Global mixed-to-scalar capture on the clean patch]
\label{thm:global-mixed-to-scalar-capture-clean-patch}
Let $\Gamma$ be a maximal connected genuinely mixed clean $4$-point branch on a
nonplanar good $v_5$-patch. Then $\Gamma$ cannot persist as an independent
genuinely mixed branch entirely inside
\[
\{T_v\neq 0\}\cap\{\mathcal R\neq 0\}.
\]
Equivalently, every such maximal branch is forced either into
\[
T_v=0
\]
or into
\[
\mathcal R=0,
\]
up to the already-absorbed planar, lower-core, repeated-node, repeated-source,
and quintic-collapse regimes.
\end{theorem}

\begin{proof}
Assume $\Gamma$ is a maximal connected genuinely mixed clean $4$-point branch on
a nonplanar good $v_5$-patch and that
\[
\Gamma\cap\{T_v=0\}=\varnothing,
\qquad
\Gamma\cap\{\mathcal R=0\}=\varnothing.
\]
We show that $\Gamma$ cannot persist as an independent genuinely mixed branch.

If a terminal point of $\Gamma$ has positive scale, then
Corollary~\ref{cor:positive-scale-endpoint-absorbed-or-scalar} places it on one
of
\[
A=0,\qquad D_\lambda=0,\qquad \mathfrak D=0,\qquad cMv_5=0,\qquad
T_v=0,\qquad \mathcal R=0.
\]
The last two are excluded by hypothesis. The remaining three are already
absorbed: $A=0$ is the nonmixed affine-planar $\mathcal R$-lane
(Remark~\ref{rem:A-zero-is-nonmixed-handoff}), the source-coupled walls
$D_\lambda=0$ and $\mathfrak D=0$ are locally dead as pass-through mechanisms
in the sufficiently small regime, and patch loss $cMv_5=0$ is absorbed into the
lower-core / shifted-$3$-point / quintic-collapse package recorded in
Remark~\ref{rem:post-draft-findings-established-provisional-live} and
Proposition~\ref{prop:u5-chart-exceptional-point-handoff}. So no positive-scale
terminal point yields an independent mixed continuation lane.

If instead the relevant scale along $\Gamma$ shrinks, then
Proposition~\ref{prop:small-scale-non-scalar-exits-absorbed} applies. Indeed,
under the standing assumptions
\[
\Gamma\cap\{T_v=0\}=\varnothing,
\qquad
\Gamma\cap\{\mathcal R=0\}=\varnothing,
\]
no sufficiently small genuinely mixed clean branch can continue through any of
the non-scalar exit walls
\[
A=0,\qquad D_\lambda=0,\qquad \mathfrak D=0,\qquad cMv_5=0
\]
as an independent mixed carrier.

Thus in all cases $\Gamma$ is forced either into $T_v=0$, into $\mathcal R=0$,
or into one of the already-absorbed regimes. This contradicts the assumption
that $\Gamma$ persists as an independent genuinely mixed clean branch inside
\[
\{T_v\neq 0\}\cap\{\mathcal R\neq 0\}.
\]
\end{proof}

\begin{remark}
The theorem should be read as the mixed-global four-point closure statement on a
nonplanar good patch: after the source-coupled and patch-loss absorptions, no
independent genuinely mixed clean carrier remains away from the scalar exits
$T_v=0$ and $\mathcal R=0$.
\end{remark}

\begin{remark}[Unified mixed/nonmixed picture after the clean-patch closure]
\label{rem:unified-mixed-nonmixed-picture-after-clean-closure}
Theorem~\ref{thm:global-mixed-to-scalar-capture-clean-patch} should be read as
the point where the mixed and nonmixed $4$-point geometries fully meet.

On the mixed side, the centered-base local theorem
(Theorem~\ref{thm:local-exhaustion-mixed-four-point-centered-base}) closes the
only genuinely new local mixed branch. On the nonmixed side, the affine-planar
lane is already governed by the scalar invariant
\[
\mathcal R=\frac{T_v}{K_v^2},
\]
and interval-persistent or recurrent exact contacts force affine-planar
collapse
(Proposition~\ref{prop:interval-persistent-nonmixed-four-point-forces-affine-planarity}
and Corollary~\ref{cor:recurrent-exact-nonmixed-contacts-force-affine-planarity}).

Accordingly, the sufficiently small clean mixed-global problem is no longer a
search for a new mixed carrier lane. It is a mixed-to-scalar capture theorem:
every independent genuinely mixed clean branch is forced either into the
centered-base scalar exit
\[
T_v=0
\]
or into the already-handled nonmixed scalar exit
\[
\mathcal R=0,
\]
with all source-coupled and patch-loss phenomena absorbed into lower-core,
repeated-node, or already-quintic handoffs.
\end{remark}

\begin{remark}[What remains genuinely open after the clean mixed closure]
\label{rem:what-remains-open-after-clean-mixed-closure}
After
Theorem~\ref{thm:global-mixed-to-scalar-capture-clean-patch},
the clean mixed four-point case is no longer an independent global frontier.
What remains genuinely open is the global nonadjacent propagation problem on the
nonmixed/source-level side, together with the separate two-point program.

More concretely, on a generic component the remaining one-pair theorem may be
read as the absence of same-parity / distance-two $2$-secant recurrence in the
affine chart, equivalently as no same-parity level re-entry for
\[
\beta(h)=-\frac{L(h)}{M(h)},
\]
equivalently as no same-parity vertical-line re-entry for the dual curve
\[
T(h)=\left(-\frac{L(h)}{M(h)},\frac{N(h)}{M(h)}\right).
\]

The current local theory already resolves the finite-order
$W=0,\ M\neq 0$ boundary behavior and the generic turning-boundary blocks.
After the same-parity compression, the genuinely residual affine lane is no
longer arbitrary nonadjacent recurrence but isolated same-parity sign descent,
equivalently discrete peaks of the intercept sequence on triple-active windows.
So the sharpened affine frontier should reduce to middle-limited triples with a
single endpoint sign test, rather than an unrestricted remote recurrence
picture.

On the mixed side, the residual overlap contribution is concentrated on
$K_v=0$ and should be read there as a quotient-diagonal fixed-shear problem:
a close scalar-local lane and a separate three-run overlap-transport lane, not
a new local mixed carrier or hidden finite-order branch type.

Thus the remaining burden is global/nonlocal, but the one-pair derivative
geometry now compresses it substantially: interior common-tangent structure is
pseudoline-ordered, local birth at generic turning and finite-order
$W=0,\ M\neq 0$ boundaries is already sterile, and the affine residual has been
reduced to isolated same-parity endpoint/branch configurations, equivalently to
zeros of the monotone endpoint-gap functions $G_{\xi,J}$ in the middle-limited
triple frontier.

On the mixed residual $K_v=0$ lane, what remains is not another local
branch-classification problem, but the separate scalar-local / transport /
package questions discussed in the fixed-shear section.
\end{remark}

\section{Endgame architecture and remaining propagation problem outside the clean mixed four-point case}

The remaining endgame is now a global propagation problem for the normalized
line-incidence geometry rather than a local boundary-creation problem.

\begin{remark}[Equivalent formulations of the remaining propagation problem outside the clean mixed four-point case]
\label{rem:equivalent-global-propagation-problem}
After Theorem~\ref{thm:global-mixed-to-scalar-capture-clean-patch}, the
remaining missing one-pair theorem on a generic component may be formulated in
any of the following equivalent ways:
\begin{enumerate}
\item no finite-slope affine line in the chart
\[
x=\frac{v_5}{c},
\qquad
y=\frac{\Delta_7}{c}
\]
is $2$-secant to two nonadjacent ECT intervals;

\item the slope function
\[
\beta(h):=-\frac{L(h)}{M(h)},
\qquad
L=v_5\Delta_7'-v_5'\Delta_7,
\]
does not revisit a level on nonadjacent branches;

\item the dual tangent-line curve
\[
T(h)=\left(-\frac{L(h)}{M(h)},\frac{N(h)}{M(h)}\right)
\]
has no nonadjacent same-parity re-entry against a vertical line in the
$(b,m)$-plane;

\item on the highest-new slice, the scalar graph
\[
\phi=\lambda^{[1]}
\]
has no nonadjacent re-entry of its slope $\phi_\mu$.
\end{enumerate}
After the same-parity compression of
Remark~\ref{rem:wip-same-parity-compression-affine-lane}, the equivalent
global formulations in this section should be read with that sharpening built
in: the genuinely residual affine lane is distance-two / same-parity only, and
the remaining global propagation problem is already reduced to isolated
middle-limited triples rather than arbitrary nonadjacent recurrence.

The results proved above resolve the local behavior at finite-order
$W=0,\ M\neq 0$ boundaries and at generic turning boundaries. More precisely,
the one-pair support geometry now shows that fixed-slope overlap windows are
one-shot on each ordered pair, interior common-tangent crossings are
pseudoline-ordered, and endpoint-support incidence is controlled by the
strictly monotone endpoint-gap functions
\[
G_{\xi,J}(h)=y_\xi-y(h)-m(h)\bigl(\xi-x(h)\bigr).
\]
Thus the remaining theorem target is no longer an arbitrary global recurrence
statement. It is the exclusion of remote zeros of these endpoint-gap functions,
equivalently an RH-specific sign theorem for endpoint-support incidence on
ordered nonadjacent branches.
\end{remark}

\begin{remark}[What the present global reduction does and does not buy]
\label{rem:what-global-reduction-does-and-does-not-buy}
The results above now force the remaining propagation problem into a much
smaller combinatorial and analytic class.

On the positive side, fixed-slope incidence is interval-graph geometry,
support curves are strictly concave pseudolines with single-crossing and
betweenness, each ordered pair has at most one directional or pair-pair overlap
window in slope, and each endpoint support point meets a given ordered target
branch at most once. Thus any surviving bad core must be slope-changing and
pair-nonrepeating, and any remaining one-pair obstruction must arise through
some element of a finite table of remote endpoint/branch incidences.

However, it is important to record a limitation. By
Remark~\ref{rem:intrinsic-affine-plane-universality}, intrinsic one-pair affine
geometry is still universal for analytic plane curves. So pure planarity,
parity, interval-graph combinatorics, pseudoline support geometry, or any
similar abstract strip/support formalism does \emph{not} by itself close the
remaining branch. Any final contradiction must therefore use RH-specific
source/provenance/canonical structure, not only generic secant-strip geometry.
\end{remark}

\begin{remark}[Warning: what the current local theory does \emph{not} imply]
\label{rem:warning-no-easy-globalization}
Several tempting globalization statements suggested by the local theory are too
strong in generic models and should not be assumed without further RH-specific
input.

\begin{enumerate}
\item Adjacent exclusion is not enough: exclusion of adjacent pair-pair lines
does not by itself force exclusion for nonadjacent pairs.

\item Local convexity of the dual tangent-line curve, even together with the
shared asymptote structure of
Proposition~\ref{prop:dual-asymptote-turning-boundary}, is not by itself enough
to forbid nonadjacent same-parity re-entry in the dual $(\beta,m)$-plane.

\item The direct highest-new septic coefficient is explicit enough to prove
noncancellation, but not generically one-signed. Thus no trivial sign theorem
for the scalar graph $\phi=\lambda^{[1]}$ should be expected without further
RH-specific structure.

\item By Remark~\ref{rem:intrinsic-affine-plane-universality}, the intrinsic
one-pair affine geometry is universal for analytic plane curves. So no theorem
using only planarity, parity, pseudoline support order, interval-graph overlap,
or similar abstract strip/support data can by itself exclude nonadjacent
recurrence.
\end{enumerate}

Accordingly, the remaining gap is now a genuinely global special-structure
theorem: not another local birth lemma, but an RH-specific sign/ordering law
for the remote endpoint-gap incidences.
\end{remark}

\begin{remark}[Intrinsic affine-plane universality and the ceiling of pure planar arguments]
\label{rem:intrinsic-affine-plane-universality}
It is useful to record a sharp ceiling on what can be proved from the intrinsic
one-pair affine geometry alone.

Indeed, given any real-analytic plane curve
\[
h\longmapsto (x(h),y(h))
\]
on an interval, setting
\[
c\equiv 1,
\qquad
v_5=x,
\qquad
\Delta_7=y
\]
gives
\[
M=x',
\qquad
N=y',
\qquad
W=x'y''-y'x''.
\]
So the normalized affine chart, the ECT/curvature data, the tangent-separator
kernel, and the dual/support geometry reduce exactly to the usual tangent
geometry of the analytic plane curve $(x,y)$.

Accordingly, no theorem using only planarity, parity, support-curve
single-crossing, strip-overlap combinatorics, or any other intrinsic one-pair
affine-topological datum can by itself exclude nonadjacent recurrence. Any
final contradiction must therefore use RH-specific source/provenance/canonical
structure, not only abstract plane-curve or secant-strip geometry.
\end{remark}

\begin{remark}[Remote endpoint-gap reduction of the remaining nonadjacent problem]
\label{rem:remote-endpoint-gap-reduction}
Combine Lemma~\ref{lem:endpoint-gap-monotonicity-ordered-target},
Theorem~\ref{thm:adjacent-pair-pair-exclusion-turning}, and
Theorem~\ref{thm:finite-order-w-zero-local-sterility}. Then every local
boundary-creation mechanism on the generic one-pair region is already sterile:
adjacent generic turning blocks contribute no pair-pair geometry, and
finite-order isolated boundaries of type
\[
W=0,
\qquad
M\neq 0
\]
do not create a new opposite-side pair-pair branch.

Therefore any surviving nonadjacent pair-pair obstruction must appear as a
\emph{remote} interior zero of one of the monotone endpoint-gap functions
\[
G_{\xi,J}(h)=y_\xi-y(h)-m(h)\bigl(\xi-x(h)\bigr),
\]
where $(\xi,y_\xi)$ ranges over endpoint support points of ECT branches and
$J$ ranges over ordered nonadjacent target branches.

Since each $G_{\xi,J}$ is strictly monotone on $J$, each endpoint/branch pair
can contribute at most one such tangency. Thus the remaining propagation
problem reduces to a finite remote-incidence table.

Accordingly, the natural next global target is an RH-specific sign theorem for
the endpoint-gap functions $G_{\xi,J}$, not another local boundary-birth
analysis.
\end{remark}

\begin{remark}[Work in progress: same-parity compression of the residual affine lane]
\label{rem:wip-same-parity-compression-affine-lane}
A useful correction is that the adjacent branch should not be counted in the
nonadjacent envelope: the genuinely residual affine lane is a distance-two /
same-parity propagation problem.

For a common slope slice $m$, write
\[
\Delta_j(m):=\tau_{I_{j+1}}(m)-\tau_{I_j}(m),
\qquad
\varepsilon_j(m):=\operatorname{sgn}\Delta_j(m),
\]
and for a triple
\[
I_{j-1}<I_j<I_{j+1}
\]
active on a common slope interval define
\[
A_j:=\Omega_{j-1}\cap\Omega_j\cap\Omega_{j+1}=(L_j,R_j).
\]
Since each $\Delta_j$ is strictly decreasing, if
\[
c_j:=m_{j,j+1}
\]
denotes the unique adjacent common-tangent slope, then on $A_j$
\[
\Delta_j(m)>0 \iff m<c_j,
\qquad
\Delta_j(m)<0 \iff m>c_j.
\]

Thus the dangerous same-parity sign pattern is exactly
\[
(\varepsilon_{j-1}(m),\varepsilon_j(m))=(+,-),
\]
equivalently
\[
\tau_{I_{j-1}}(m)<\tau_{I_j}(m)>\tau_{I_{j+1}}(m),
\]
equivalently
\[
c_j<m<c_{j-1}.
\]

So the residual affine obstruction is not arbitrary nonadjacent recurrence but
an isolated sign descent / discrete peak of the intercept sequence
$j\mapsto \tau_{I_j}(m)$. The actual triple can be dangerous only if the
active window crosses that descent:
\[
A_j\cap(c_j,c_{j-1})\neq\varnothing.
\]

Writing
\[
e_{j-1}(m):=\sup \mathcal B_{I_{j-1}}(m),
\qquad
B_j(m):=e_{j-1}(m)-\tau_{I_{j+1}}(m),
\]
the residual same-parity obstruction interval is therefore
\[
Q_j:=\{\,m\in A_j\cap(c_j,c_{j-1}) : B_j(m)\ge 0\,\}.
\]

In particular, once the active-window endpoint bookkeeping is added, the
sharpened remaining affine risk should reduce to isolated middle-limited triples
of the form
\[
c_j<r_j<c_{j-1},
\qquad
B_j(r_j)>0,
\]
where $r_j$ is the upper endpoint of the triple-active window. This is the
sharpest current affine formulation, but it remains a theorem target rather
than a proved closure statement.
\end{remark}

\begin{remark}[Further sharpening of the same-parity affine frontier]
\label{rem:further-sharpening-same-parity-affine-frontier}
After the distance-two / same-parity correction, the residual affine lane is
best viewed through the adjacent sign sequence
\[
\varepsilon_j(m):=\sgn\!\bigl(\tau_{j+1}(m)-\tau_j(m)\bigr).
\]
A dangerous same-parity affine defect at the triple
\[
I_{j-1}<I_j<I_{j+1}
\]
is exactly an isolated sign descent
\[
(\varepsilon_{j-1}(m),\varepsilon_j(m))=(+,-),
\]
equivalently a discrete local peak
\[
\tau_{j-1}(m)<\tau_j(m)>\tau_{j+1}(m).
\]
In particular, dangerous affine defect sites are isolated: no consecutive
defects can occur.

Moreover, after imposing the actual triple-active window constraint, the
sharpened residual affine frontier is middle-limited:
\[
c_j<r_j<c_{j-1},
\qquad
B_j(r_j)>0.
\]
So the unresolved affine risk is no longer arbitrary remote recurrence. It is an
isolated same-parity sign descent on an actually active triple, limited on the
right by the middle branch, together with one endpoint sign test.
\end{remark}

\begin{proposition}[Full-block universality of the direct provenance slice]
\label{prop:full-block-universality-direct-provenance-slice}
Let $J$ be a connected generic one-pair block on which
\[
A(h):=\kappa(h)\eta_2(h)
\]
is analytic and nowhere zero. On $J$, the direct provenance slice has the form
\[
\phi(h)=A(h)\,q^{(7)}(h).
\]
Then for every analytic function
\[
\psi:J\to\mathbf C
\]
there exists an analytic function
\[
q:J\to\mathbf C
\]
such that
\[
\phi(h)=\psi(h)
\qquad (h\in J).
\]

More explicitly, for any basepoint $h_0\in J$ and any constants
\[
a_0,\dots,a_6\in\mathbf C,
\]
the function
\[
q_\psi(h)
:=
\sum_{j=0}^6 \frac{a_j}{j!}(h-h_0)^j
+
\int_{h_0}^h \frac{(h-s)^6}{6!}\,\frac{\psi(s)}{A(s)}\,ds
\]
is analytic on $J$ and satisfies
\[
q_\psi^{(7)}(h)=\frac{\psi(h)}{A(h)},
\qquad
\phi(h)=A(h)\,q_\psi^{(7)}(h)=\psi(h).
\]
In particular, on any generic one-pair block, the direct highest-new scalar
graph $\phi$ is an arbitrary analytic scalar graph.
\end{proposition}

\begin{proof}
Since $A$ is analytic and nowhere zero on $J$, the function
\[
\frac{\psi}{A}
\]
is analytic on $J$. The displayed formula for $q_\psi$ is the standard
variation-of-constants solution of
\[
q^{(7)}=\frac{\psi}{A},
\]
with arbitrary initial data at $h_0$. Differentiating seven times gives
\[
q_\psi^{(7)}(h)=\frac{\psi(h)}{A(h)},
\]
hence
\[
\phi(h)=A(h)\,q_\psi^{(7)}(h)=\psi(h).
\]
\end{proof}

\begin{corollary}[Finite linear provenance constraints are too weak]
\label{cor:finite-linear-provenance-too-weak}
Let $J$ be as in
Proposition~\ref{prop:full-block-universality-direct-provenance-slice}.
Fix finitely many linear functionals
\[
\Lambda_r(q)=c_r,
\qquad r=1,\dots,N,
\]
on analytic functions $q$ on $J$, where each $\Lambda_r$ is built from finitely
many point-evaluation jets and/or finitely many linear integral functionals
with analytic kernels.

Fix also two distinct points
\[
\mu_1\neq \mu_2\in J
\]
and an affine line
\[
\ell(\mu)=s\mu+b.
\]
Then there exists a polynomial $q$ on $J$ such that:

\begin{enumerate}
\item all the prescribed linear provenance conditions
\[
\Lambda_r(q)=c_r
\qquad (r=1,\dots,N)
\]
hold;

\item the direct-slice graph
\[
\phi(\mu)=A(\mu)\,q^{(7)}(\mu)
\]
is tangent to the line $y=\ell(\mu)$ at both $\mu_1$ and $\mu_2$, i.e.
\[
\phi(\mu_j)=\ell(\mu_j),
\qquad
\phi_\mu(\mu_j)=s
\qquad (j=1,2).
\]
\end{enumerate}

Accordingly, no theorem using only the explicit direct formula
\[
\phi=A\,q^{(7)}
\]
together with any finite family of linear provenance or matching conditions on
$q$ can exclude bitangent, four-secant, or same-slope recurrence geometry on
the direct slice.
\end{corollary}

\begin{proof}
The tangency conditions at $\mu_j$ are linear jet conditions on $q$:
\[
A(\mu_j)\,q^{(7)}(\mu_j)=\ell(\mu_j),
\]
and, since
\[
\phi_\mu=(Aq^{(7)})'=A' q^{(7)}+A q^{(8)},
\]
also
\[
A'(\mu_j)\,q^{(7)}(\mu_j)+A(\mu_j)\,q^{(8)}(\mu_j)=s.
\]
Because $A(\mu_j)\neq 0$, these are legitimate linear conditions on the jets of
$q$ at the two points $\mu_1,\mu_2$.

Together with the prescribed finitely many provenance constraints
\[
\Lambda_r(q)=c_r,
\]
this is a finite linear system on the coefficients of a polynomial $q$. For a
polynomial of sufficiently large degree, the coefficient space has dimension
larger than the number of imposed linear conditions. Standard Hermite
interpolation / finite-dimensional linear algebra therefore gives a polynomial
$q$ satisfying all of them simultaneously.
\end{proof}

\begin{remark}[What the direct-slice no-go results actually leave open]
\label{rem:what-direct-slice-no-go-results-leave-open}
Proposition~\ref{prop:full-block-universality-direct-provenance-slice} shows
that on a generic one-pair block the direct highest-new scalar graph
\[
\phi
\]
is fully universal. Corollary~\ref{cor:finite-linear-provenance-too-weak} shows
more: even after imposing any finite family of linear provenance / matching
conditions on $q$, the direct slice remains flexible enough to realize
bitangent, four-secant, and same-slope recurrence geometry.

So the next live theorem cannot be another sign/convexity statement about the
graph of $\phi$, nor any theorem using only finitely many endpoint jets or
finitely many linear integral conditions on $q$. Any successful continuation of
the direct-slice lane must instead use a genuinely
\emph{continuum-strength provenance law} for $q$ on a whole block, such as a
lower-order ODE, a sign-definite integral identity, a nonlinear transport law,
or some other infinite-codimension constraint.
\end{remark}

\begin{lemma}[Crude polynomial bound for the local value scale]
\label{lem:crude-S-bound}
Let
\[
S(m)=\sum_\rho f_{\beta_\rho,\gamma_\rho}(m),
\]
where
\[
f_{\beta,\gamma}(m)
=
\frac{1-\beta}{(1-\beta)^2+(2m-\gamma)^2}
+
\frac{\beta}{\beta^2+(2m-\gamma)^2}.
\]
Then
\[
S(m)\ll Q^2,
\qquad Q\asymp \log|m|.
\]
Consequently
\[
L(m)\ll Q
\]
and
\[
B_\zeta(m)L(m)S(m)\ll Q^4.
\]
\end{lemma}

\begin{proof}
By the zero-free region, for zeros at relevant height one has
\[
\min(\beta,1-\beta)\gg \frac1Q.
\]
Hence
\[
f_{\beta,\gamma}(m)\ll \frac{Q}{1+(2m-\gamma)^2}.
\]
Partition the zeros into unit shells
\[
\mathcal S_k:=\{\rho: k\le |2m-\gamma_\rho|<k+1\},
\qquad k=0,1,2,\dots
\]
Standard zero counting gives
\[
\#\mathcal S_k\ll Q
\]
for each unit shell near height $m$. Therefore
\[
S(m)
\ll
\sum_{k\ge 0}\#\mathcal S_k\frac{Q}{1+k^2}
\ll
Q\sum_{k\ge 0}\frac{Q}{1+k^2}
\ll
Q^2.
\]
Since
\[
L(m)=\log\!\Bigl(3+|m|+\frac1{S(m)}\Bigr),
\]
we trivially have $L(m)\ll Q$, and the last estimate follows from
$B_\zeta(m)\asymp Q$.
\end{proof}

The following remark records the present status of the endgame before the
final contradiction theorem.

\begin{remark}[Work in progress: current endgame status]
\label{rem:wip-final-endgame-status}
The contradiction argument in the next theorem is presently written in the
\emph{pair-like branch} of the local model analysis. More precisely, the
current endgame is fully aligned with the case in which the actual same-scale
local core is perturbatively pair-like, so that the toy-model first-odd-jet
obstruction continues to govern the corrected transverse scalar.

The fixed-core analysis developed earlier in the draft shows, however, that
a more general finite-core branch is also possible: the first odd
coefficient may cancel while a genuine corrected anomaly survives at a
higher odd order. Thus the fully general finite-core endgame would need to
be formulated in terms of the \emph{first nonzero odd jet} of the corrected
transverse scalar, rather than specifically the first one.

Accordingly, the theorem below should presently be read together with the
preceding work-in-progress remarks on the pair-like versus finite-core fork.
In particular, the pair-based contradiction is the intended endgame under
the pair-like branch, while the genuinely finite-core branch remains to be
completed in full detail.

At the local algebraic level, the canonical quotient-level order-$7$
fixed-sector package is now closed via
Theorem~\ref{thm:quotient-septic-closure}: the remaining order-$7$ ambiguity is
only the choice of a raw representative inside
$A_{7,K_1}^{\mathfrak f}+\mathbf C A_5^{\mathfrak f}$, equivalently a local
gauge-fixing of $(u_7,v_7)$, and does not affect the quotient class
$[A_7^{\mathfrak f}]$ or the canonical second residual determinant
$\Delta_7$.
\end{remark}

\begin{remark}[Work in progress: asymptotic contradiction versus full $\RH$]
\label{rem:wip-high-height-only}
As presently written, the final contradiction is only asymptotic in the height
parameter $Q$: the comparison closes after reaching a regime in which the
zeta-side upper bound is dominated by the toy-side lower bound, schematically
because
\[
Q^4 e^{-\delta Q}\to 0
\qquad (Q\to\infty).
\]
Accordingly, the current argument excludes sufficiently large pair-like
off-line anomalies, but does not yet by itself rule out a hypothetical
low-lying off-line zero.  In particular, the present proof does not furnish a
uniform-in-height contradiction, nor does it presently prove that a single
off-line zero at one height would force corresponding bad local anomalies at
arbitrarily large heights.

A plausible repair is to split the endgame into an effective high-height
theorem plus a separate finite low-height bridge.

\smallskip

\noindent
\emph{Effective high-height step.}
One should first replace the soft asymptotic comparison by an explicit
inequality.  Concretely, the final contradiction compares a toy-side lower
bound of the form
\[
|\Xi_{\toy}^{(N)}(u,d)|
\ge
c_{\toy}\,\mathcal{L}(m)
\]
against a zeta-side upper bound of the form
\[
|\Xi_{\zeta}^{(N)}(m)|
\le
C_{\zeta}(\rho_0,\alpha)\,
\frac{\mathcal{U}(m)}{Q}\,
e^{-\delta_*(\alpha,\rho_0)Q},
\]
where in the present draft one has schematically
\[
\mathcal{L}(m)\asymp \frac{x}{B_\zeta(m)}S(m)
\]
(or its variable-$x$ replacement from
Remark~\ref{rem:wip-calibration-small-u}),
and
\[
\mathcal{U}(m)\asymp L(m)S(m)^2.
\]
To extract an actual threshold, every use of $O(\cdot)$, $\ll$, and
$\asymp$ in the final comparison must be replaced by explicit constants.
After substituting explicit bounds such as
\[
B_\zeta(m)\le C_B Q,\qquad
L(m)\le C_L Q,\qquad
S(m)\le C_S Q^2,
\]
the contradiction should reduce to an explicit scalar inequality of the form
\[
C_*(\rho_0,\alpha)\,Q^3 e^{-\delta_*(\alpha,\rho_0)Q}<1
\qquad (Q\ge Q_*),
\]
for some computable constant $C_*>0$.  Any explicit $Q_*$ satisfying this
inequality would then yield an \emph{effective} theorem excluding pair-like
off-line anomalies for all $Q\ge Q_*$.

\smallskip

\noindent
\emph{Height extraction.}
Once $Q_*$ is known explicitly, one must convert it into an actual height
threshold.  Since the draft works with a midpoint parameter $m$ satisfying
\[
Q\asymp \log |m|,
\]
this requires making the logarithmic comparability effective, for example in
the form
\[
c_Q \log(3+|m|)\le Q\le C_Q \log(3+|m|).
\]
From such inequalities one obtains an explicit midpoint threshold $m_*$, and
hence an explicit zero-height threshold $T_*$, since the local points under
consideration satisfy
\[
t_\pm = m \pm \frac{\sigma}{2}
\]
with $|\sigma|$ microscopic relative to the height scale.  Thus one would
arrive at a statement of the form:
\[
\text{no pair-like off-line zero can occur for } |\Im \rho|>T_*.
\]

\smallskip

\noindent
\emph{Finite low-height bridge.}
The effective high-height theorem alone still does not imply the full Riemann
Hypothesis unless one also excludes all potential off-line zeros below the
threshold $T_*$.  A natural way to do this is to combine the structural
argument above height $T_*$ with a rigorous finite verification of $\RH$
up to height $T_*$.  Accordingly, if one can arrange that the explicit
threshold $T_*$ lies below a rigorously verified range, then the final step
would be:
\[
\text{high-height exclusion above }T_*
\quad+\quad
\text{rigorous verification below }T_*
\quad\Longrightarrow\quad
\RH.
\]
This would still be an unconditional proof, albeit a computer-assisted one
rather than a purely internal asymptotic contradiction.

\smallskip

\noindent
\emph{What still needs to be made explicit.}
To implement this repair rigorously, the draft must still supply all of the
following ingredients in effective form:
\begin{enumerate}
\item an explicit toy-side lower constant on the admissible compact
$d$-range and microscopic $u$-range;
\item an explicit zeta-side upper constant in the $N$-point cancellation
estimate, including the exponential decay rate $\delta$;
\item explicit growth bounds for the auxiliary quantities appearing in the
final comparison, including $B_\zeta(m)$, $L(m)$, and $S(m)$;
\item an explicit conversion between the logarithmic parameter $Q$ and the
actual height variable;
\item a final theorem stated as an effective exclusion above a concrete
threshold $T_*$, rather than merely ``for $Q$ sufficiently large'';
\item a separate low-height bridge, either by citing a rigorous finite
verification result covering the range up to $T_*$, or by proving an
alternative uniform-in-height contradiction.
\end{enumerate}

\smallskip

\noindent
There is, in principle, a second possible repair: prove a propagation theorem
showing that any off-line zero at one height would force corresponding bad
local anomalies at arbitrarily large heights, so that the asymptotic
contradiction would automatically apply.  However, no such propagation
principle is currently established in the draft, and at present this appears
to be a substantially stronger theorem than the effective-threshold route
described above.

Accordingly, the current endgame should be read as an asymptotic contradiction
theorem in the pair-like branch.  A plausible route to a complete proof of
$\RH$ is available, but it requires an effective extraction of the height
threshold together with a finite low-height bridge, neither of which has yet
been written out in the draft.
\end{remark}

\begin{theorem}[Asymptotic calibrated \(N\)-point contradiction in the pair-like branch]
\label{thm:final-contradiction}
Assume the pair-like branch hypothesis and assume that the pair-like
package-to-\(\Xi\) bridge identifies the calibrated local two-point package
with the toy normal form entering Proposition~\ref{prop:toy-n-point-direct}.
Then, for \(Q\) sufficiently large, no fixed off-line toy anomaly can occur in
the pair-like branch.

This theorem is not a finite-core closure theorem.  The genuinely non-pair-like
finite-core branch remains open pending the lower-transfer theorem described in
Remark~\ref{rem:open-finite-core-lower-transfer}.
\end{theorem}

\begin{proof}
By Proposition~\ref{prop:calibration-remainder-cutoff} and the canonical
calibration theorem, in the fixed small-$x_0$ regime one has
\[
u^2\asymp \frac{x_0}{B_\zeta(m)}S(m).
\]

On the toy side, Proposition~\ref{prop:toy-n-point-direct} gives
\[
\Xi_{\toy}^{(N)}(u,d)=u^2\Phi_K(M(d))+O(u^4),
\]
hence on compact $d$-ranges,
\[
|\Xi_{\toy}^{(N)}(u,d)|\asymp u^2
\asymp
\frac{x_0}{B_\zeta(m)}S(m).
\]

On the zeta side, Corollary~\ref{cor:optimized-zeta-cancellation-depth} gives
\[
|\Xi_\zeta^{(N)}(m)|
\ll_{\rho_0,\alpha}
\frac{L(m)S(m)^2}{Q}\,
e^{-\delta_*(\alpha,\rho_0)Q}.
\]

By Lemma~\ref{lem:crude-S-bound}, one has $S(m)\ll Q^2$, $L(m)\ll Q$, and
$B_\zeta(m)L(m)S(m)\ll Q^4$. Therefore
\[
\frac{(L(m)S(m)^2/Q)e^{-\delta_*(\alpha,\rho_0)Q}}
{(x_0/B_\zeta(m))S(m)}
=
\frac{B_\zeta(m)}{x_0}\,\frac{L(m)S(m)}{Q}\,
e^{-\delta_*(\alpha,\rho_0)Q}
\ll
Q^3 e^{-\delta_*(\alpha,\rho_0)Q}\to 0.
\]
Hence for $Q$ sufficiently large,
\[
|\Xi_\zeta^{(N)}(m)|<|\Xi_{\toy}^{(N)}(u,d)|.
\]

This contradicts the assumed calibrated realization of a fixed off-line toy
anomaly in the pair-like branch, once \(Q\) is sufficiently large.  Thus,
under the pair-like branch hypothesis and the package-to-jet bridge
hypothesis, the pair-like branch is excluded at sufficiently large height.

This theorem is an asymptotic branch exclusion.  To obtain a complete
\(\RH\) proof from it one still needs: an effective extraction of the
threshold \(Q_*\), a conversion from \(Q_*\) to a height threshold \(T_*\),
a finite low-height verification or alternative low-height argument, and
the non-pair-like finite-core branch closure.
\end{proof}

\begin{remark}[Local package results are not yet global transform closure]
\label{rem:local-package-not-global-transform-closure}
The local Gate~B and connection-package results are package-construction
results.  They do not by themselves prove hidden-state equality,
exceptional-divisor fiber selection, finite-core exclusion, or a full
package-to-transform contradiction.  The remaining global obligations are:
\[
\text{state-locality modulo }\ker\Phi_K,
\qquad
\text{quotient-valued diagonal merger},
\qquad
\text{finite-core lower transfer}.
\]
\end{remark}

\begin{remark}[Current status summary after the consolidated round]
\label{rem:current-status-summary-consolidated-round}
Reading
Lemma~\ref{lem:edge-refined-y-middle-parity-final},
Theorem~\ref{thm:ma-mc-pretriv-connection-closure-final},
Definition~\ref{def:promoted-quadratic-mixed-block-bc-coefficient},
Definition~\ref{def:transform-safe-augmented-reduced-package},
Theorem~\ref{thm:finite-odd-germ-determinacy-actual-source}, and
Definition~\ref{def:first-interaction-branch} together, the present status is:
\[
\boxed{
\text{The }M_A,M_C\text{ connection identities close under actual
source-linear }Y\text{-middle edge parity.}
}
\]
\[
\boxed{
\text{The local obstruction has sharpened to the odd-background }
\eta\eta\text{ top row inside }E_{A,Y}^{\mathrm{BC}},
\text{ led by }\mathcal E^{2}_{\eta\eta;7,0,0};
\text{ the present transform-faithful finite-}s\text{ architecture does not
currently remove this row.}
}
\]
\[
\boxed{
\text{The transform-safe package is }
\widehat\Psi_{\rm red}^{+}=(\widehat\Psi_{\rm red},T),
\quad T=v_7/c.
}
\]
\[
\boxed{
\text{Finite determinacy is qualitative; calibrated separation and
interaction rigidity remain open.}
}
\]
\end{remark}

\begin{remark}[Current local obstruction: odd-background \(\eta\eta\) top row]
\label{rem:current-local-obstruction-etaeta-top-row}
After correcting the Bottleneck-C bookkeeping, the first local actual-source
obstruction is no longer merely an unevaluated broad finite table.  The
current leading row is the frozen odd-background \(\eta\eta\) top row
\[
\mathcal E^{2}_{\eta\eta;7,0,0}.
\]
It is governed by
\[
\Lambda^-_{2,0}
\left[
\left(
\frac{\sin\Delta_0}{s^2}
+
\frac{q_-\cos\Delta_0}{s}
\right)
\eta_i\eta_j
\right]_{7}
W_{+,0},
\]
where
\[
\Lambda^-_{2,0}
=
-\frac{1}{24\pi q^3}E_{22}.
\]
Writing
\[
\nu_{21,7}
:=
[z^7]\left[
\left(
\frac{\sin\Delta_0}{s^2}
+
\frac{q_-\cos\Delta_0}{s}
\right)
\eta_i\eta_j
\right],
\]
the top-row \(S\)-coefficient is
\[
f_{\eta\eta}^{\rm top}
=
-\frac{u}{24\pi q^3}\nu_{21,7}.
\]
Thus the local branch is now binary: either prove a zeta-specific identity
\[
\nu_{21,7}=0,
\]
or supply a genuine physical aggregate \(S\)-coefficient cancellation of the
row.  Endpoint-diagonal source realization, even if proved, removes only
endpoint-separation subchannels and does not by itself close this scalar
\(\Delta\)-Hessian \(\eta\eta\) survivor.  The only remaining local rescue is
a genuine physical aggregate \(S\)-coefficient cancellation, not a formal
package counterterm.
\end{remark}

\begin{remark}[Current outcome of the Bottleneck-C repair audits]
\label{rem:current-outcome-bottleneck-c-repair-audits}
The subsequent audits below sharpen the status of
Remark~\ref{rem:current-local-obstruction-etaeta-top-row}.  Replacing the raw
far-tail source by the full non-core zero source is coherent, but the
rederived top row has the same obstruction form.  Replacing the scalar local
readout, quotienting by the \(A_5^{\mathfrak f}\)-line, changing the
finite-\(s\) carrier, or exact endpoint symmetrization likewise does not remove
the \(S\)-visible \(\eta\eta\) row while preserving transform faithfulness.

Consequently, within the present transform-faithful finite-\(s\) architecture,
Bottleneck~C still requires a genuine vanishing of the physical aggregate
\(S\)-coefficient of the \(\eta\eta\) top row.  No such vanishing theorem is
presently proved.
\end{remark}

\section{Downstream local closure package: Gate 2, wall escape, and carrier rows}
\label{sec:downstream-local-closure-package}

This section records the current local downstream replacement for the older
\(D1/D2/D3\) architecture.  The active Gate--2 route is no longer a
Remez--Cartan or pointwise target-smallness route.  It is
\[
\boxed{
\text{Gate 1}
+
|u'|_\infty\text{ source thickening}
+
L^2\text{-input real rigidity}
\Longrightarrow
\text{Gate 2}.
}
\]
The active wall route is
\[
\boxed{
\text{Gate 2}
+
\text{gauge-locked wall escape}
+
\text{D4 carrier-quotient/filtered row conversion}
+
\text{row mobility}
\Longrightarrow
\text{row energy}.
}
\]
The source-robust wall-row search and the scalar \(I\)-coordinate order-nine
incidence wall remain fallback formulations only if the row-output architecture
is forced to be exact-point or scalar-collapsed.

\section{Unified finite-survival architecture}
\label{sec:unified-finite-survival-architecture}

\begin{theorem}[Unified finite-survival theorem]
\label{thm:unified-finite-survival}
The proof is organized around four finite-survival statements:
\begin{enumerate}
\item \textbf{Source survival:} \(\mathrm{Agg}L^2(q''_{\can})\) on a polynomial-weight good packet set.
\item \textbf{Graph survival:} \(\mathrm{Graph/O2O3}_{23}\) after \(R^C\)-normalization and full determinant survival.
\item \textbf{Quotient survival:} the primitive \(Y12\) carrier survives in the same quotient that removes actual-source nuisances.
\item \textbf{Atomic contradiction:} the singular-discard observable is Type-I local-square and is bounded above on the same quotient norm.
\end{enumerate}
If all four hold, the same observable is simultaneously bounded below and above by incompatible polynomial scales, giving the downstream contradiction.
\end{theorem}

\begin{equation}\label{eq:proof-architecture-diagram}
\begin{array}{c}
\boxed{\mathrm{Agg}L^2(q''_{\can})}
\\[3pt]
\Downarrow \quad \text{(packet-core micro-covering + EF completion)}
\\[3pt]
\boxed{\mathrm{Graph/O2O3}_{23}}
\\[3pt]
\Downarrow \quad \text{(\(R^C\)-normalized Riesz basis + thickened Schur survival)}
\\[3pt]
\boxed{Y12\text{-quotient survival}}
\\[3pt]
\Downarrow \quad \text{(retained-space sandwich + labelled transition)}
\\[3pt]
\boxed{\mathrm{AtomicY12Contradiction}}
\end{array}
\end{equation}

\begin{proposition}[Proof sketch / structure]
\label{prop:unified-proof-sketch}
The proof proceeds in four finite-survival layers:
\begin{enumerate}
\item establish aggregate source observability for \(q''_{\can}=B_2+P_2+D_2+E_2\) on a good packet set;
\item establish finite graph nondegeneracy in the \(R^C\)-normalized O2/O3 block and promote it to full determinant survival;
\item pass to the actual-source quotient and verify that the primitive \(Y12\) carrier survives after removing nuisance directions;
\item compare the same quotient observable against the Type-I singular-discard upper and conclude contradiction when the lower and upper scales are separated.
\end{enumerate}
\end{proposition}

\subsection{Shared finite-survival language}
\label{subsec:shared-finite-survival-language}

\begin{definition}[Finite-survival gaps and residuals]
\label{def:finite-survival-gaps}
The downstream proof uses the following shared finite objects.

\paragraph{Source survival.}
Let
\[
\mathcal S_X:=\operatorname{span}\{B_2,P_{2,\le X},D_{2,\le X}\},
\qquad
F_X:=B_2+P_{2,\le X}+D_{2,\le X}.
\]
Define
\[
M:=G_{\mathrm{ref}}^{-1/2}G_{\mathrm{branch}}G_{\mathrm{ref}}^{-1/2},
\qquad
E:=G_{\mathrm{ref}}^{-1/2}\bigl(G_{\mathrm{branch}}(q''_{\can})-G_{\mathrm{branch}}(F_X)\bigr)G_{\mathrm{ref}}^{-1/2},
\]
and the source gap
\[
\delta_{\mathrm{src}}:=\min_i\!\left(M_{ii}-\sum_{j\neq i}|M_{ij}|\right)-\|E\|_{\operatorname{op}}.
\]
Source survival means \(\delta_{\mathrm{src}}\ge Q^{-A}\).

\paragraph{Graph survival.}
Define
\[
K:=G_{00}^{-1/2}G_{0S}G_{SS}^{-1/2},
\qquad
\delta_{\mathrm{graph}}:=1-\|K\|_{\operatorname{op}}.
\]
Graph survival means \(\delta_{\mathrm{graph}}\ge Q^{-A}\), equivalently a Schur complement lower bound or a polynomial principal-angle gap.

\paragraph{Quotient survival.}
Let
\[
U_p:=\operatorname{col}\bigl(T_{\mathrm{par}}^{(\alpha)},T_7,T_{<7},II_{\mathrm{arith},9},T_{E_2}\bigr),
\qquad
\mathcal K_p^Y:=\operatorname{im}
\begin{bmatrix}
R_{B_2,p}&R_{P,p}&R_{D,p}&A_{\alpha,p}&A_{7,p}&A_{<7,p}&A_{II,p}&A_{E_2,p}
\end{bmatrix},
\]
\[
q_p^Y:=I-P_{\mathcal K_p^Y}^{\mathrm{bd}},
\qquad
N_p:=\Pi_{\Sigma(Y12)}T_{\mathrm{tr}\to\mathrm{sing}}\mathcal N_p^{\mathrm{D4}},
\qquad
Y_p:=\Pi_{\Sigma(Y12)}T_{\mathrm{tr}\to\mathrm{sing}}[\mathsf C_{Y12}]^{\mathrm{tr}}.
\]
The quotient residuals are
\[
r_{N,p}:=q_p^Y N_p,
\qquad
r_{Y,p}:=q_p^Y Y_p.
\]
Quotient survival means \(r_{N,p}\approx 0\) and \(r_{Y,p}\neq 0\) with polynomial margin.

\paragraph{Graph/coercivity survival.}
Let
\[
\mathfrak T_p:=(I-P_{\mathrm{sing},p})T_{\mathrm{tr}\to\mathrm{sing},p}J_{\mathrm{tr},p}.
\]
If \(\mathcal R_X=\mathcal R_{X,\mathrm{low},p}\oplus\mathcal R_{X,\mathrm{tail},p}\) is the admissible residual split and \(D_p\) is the low-mode defect matrix for \(\mathfrak T_p\), then
\[
\delta_{\mathrm{coerc}}:=\sigma_{\min}(D_p).
\]
Graph/coercivity survival means \(\delta_{\mathrm{coerc}}\ge Q^{-C}\) and the tail is absorbed by the EF/singular-discard bound.

\paragraph{Extraction covariance survival.}
On each frozen cap \(C_j\), let
\[
B_j:=\sum_{p\in C_j}w_p\,\Gamma_p\Gamma_p^\ast,
\qquad
w_p=\rho_p\frac{E_{1,p}H_p}{N_p},
\qquad
\Gamma_p=\binom{\Gamma_G(p)}{\Gamma_{G,2}(p)}.
\]
Define
\[
\delta_{\mathrm{ext}}:=\frac{\det(B_j)}{(\operatorname{tr}B_j)^2}.
\]
Extraction survival means \(\delta_{\mathrm{ext}}\ge Q^{-A}\).  If \(\delta_{\mathrm{ext}}\ll Q^{-A}\), the wall is effectively static and the cap is handed to the \(Y12\) quotient module.

\paragraph{Atomic contradiction.}
The forbidden observable is
\[
\Lambda_{Y,p}:=q_p^Y\Pi_{\Sigma(Y12),p}T_{\mathrm{tr}\to\mathrm{sing},p}P_{\mathrm{tr},p}.
\]
The final contradiction compares the same quotient norm upper bound for \(\Lambda_{Y,p}J_{\mathrm{tr},p}(q''_{\can})\) with the lower bound carried by \(r_{Y,p}\).
\end{definition}

\begin{remark}[Shared failure labels]
\label{rem:shared-failure-labels}
The proof uses a fixed failure grammar:
\[
\text{branch holes},
\quad
\text{off-diagonal packet correlation too large},
\quad
\text{residual completion failure},
\quad
\text{thickened overlap too coherent},
\quad
\text{quotient killed primitive carrier},
\quad
\text{graph/coercivity collapse},
\quad
\text{wall effectively static}.
\]
Each agent should reduce its module to one of the gaps
\(\delta_{\mathrm{src}},\delta_{\mathrm{graph}},\delta_{\mathrm{coerc}},\delta_{\mathrm{ext}}\), or to the residuals \(r_{N,p},r_{Y,p}\).
\end{remark}

\subsection{\texorpdfstring{\(L^2\)}{L2}-input real rigidity}

\begin{definition}[Ordinary zero-separated source packet]
\label{def:ordinary-zero-separated-source-packet}
Let \(I^\ast\) be an enlarged ordinary packet interval and let
\[
\mathcal D=\lambda\partial_m,\qquad \lambda Q\asymp1.
\]
For \(\rho\ge Q^{-C_\rho}\), define
\[
\mathcal T_\rho(I^\ast)
=
\{m+iy:m\in I^\ast,\ |y|\le \rho\lambda\}.
\]
We say that \(I^\ast\) is ordinary zero-separated if
\[
\operatorname{dist}
\left(
\frac12+i\mathcal T_\rho(I^\ast),
\mathcal Z_\Xi
\right)
\ge
\lambda Q^{-C_0}.
\]
Here \(\mathcal Z_\Xi\) denotes the zero set of the completed scalar underlying
the canonical curvature operator \(\mathfrak C_2[\Xi]\).  It is unrelated to the
finite \(N\)-point odd scalar \(\Xi_\zeta^{(N)}\) used earlier in the
odd-moment cancellation construction.
If this condition fails, the packet is not treated by the ordinary Gate--2
local argument; it is routed to the singular/transition branch, or else a
regularized phase convention must be used.
\end{definition}

\begin{lemma}[Ordinary zero-separated source-complexity transfer]
\label{lem:ordinary-zero-separated-source-complexity-transfer}
Assume \(I^\ast\) is ordinary zero-separated.  Assume also that the canonical
source curvature is the full completed meromorphic curvature
\[
u=q''_{\can}=\mathfrak C_2[\Xi]\left(\frac12+im\right),
\]
where \(\mathfrak C_2\) is a fixed finite expression in the completed
logarithmic derivative
\[
L(s)=\frac{\Xi'}{\Xi}(s)
\]
and finitely many derivatives \(L^{(r)}(s)\), together with regular
gamma/background terms and no denominators other than those implicit in
\(L=\Xi'/\Xi\).
Assume moreover the standard local zero-count estimate for the completed
scalar:
\[
\#\{\rho\in\mathcal Z_\Xi: |s-\rho|\le R_0\}
\le Q^C(1+R_0)
\]
through the polynomial range of \(R_0\) used below.
Then
\[
\boxed{
\|u\|_{L^\infty(I^\ast)}
+
\|\mathcal D u\|_{L^\infty(I^\ast)}
\le Q^C.
}
\]
Equivalently, on \(I\subset I^\ast\),
\[
\boxed{
\|u\|_{L^\infty(I)}+|I|\,\|u'\|_{L^\infty(I)}\le Q^C.
}
\]
\end{lemma}

\begin{proof}
For fixed \(r\), the completed logarithmic derivative has the local
partial-fraction form
\[
L^{(r)}(s)
=
(-1)^r r!\sum_\rho\frac1{(s-\rho)^{r+1}}
+
\mathcal G_r(s),
\]
where \(\mathcal G_r\) denotes the regular gamma/background contribution.  Set
\[
\delta=\lambda Q^{-C_0}.
\]
On the ordinary tube, \(|s-\rho|\ge\delta\).  Split zeros into dyadic annuli
\[
A_\nu=\{\rho:2^\nu\delta\le |s-\rho|\le2^{\nu+1}\delta\}.
\]
The local zero-count estimate gives
\[
\# A_\nu\le Q^C(1+2^\nu\delta).
\]
Hence
\[
\sum_{\rho\in A_\nu}|s-\rho|^{-(r+1)}
\le
Q^C(1+2^\nu\delta)(2^\nu\delta)^{-(r+1)}.
\]
Summing over \(\nu\), the first annuli dominate:
\[
\sum_\rho |s-\rho|^{-(r+1)}
\le Q^C\delta^{-(r+1)}
\le Q^{C_r}.
\]
Thus every fixed \(L^{(r)}\) entering \(\mathfrak C_2[\Xi]\) is polynomially
bounded on the tube.  Since \(u=\mathfrak C_2[\Xi]\) has fixed finite
complexity,
\[
\|u\|_{L^\infty(I^\ast)}\le Q^C.
\]
Applying \(\mathcal D=\lambda\partial_m\) only raises the completed
log-derivative order by one and multiplies by \(\lambda\le1\), so
\[
\|\mathcal D u\|_{L^\infty(I^\ast)}\le Q^C.
\]
The physical-coordinate estimate follows from \(\mathcal D=\lambda\partial_m\)
and \(|I|\asymp\lambda\).
\end{proof}

\begin{remark}[Canonical and denominator guardrails for Gate--2]
\label{rem:canonical-denominator-guardrails-gate-two}
For Gate~2 the inclusive canonical convention is
\[
q''_{\can}=B_2+P_2+D_2+E_2=\mathfrak C_2[\Xi].
\]
Thus \(E_2\) is part of the exact canonical object and no separate
\(E_2\)-smallness is needed for the minimal \(u,\mathcal D u\) transfer.  If
instead \(E_2\) is used as a perturbative truncation error, the minimal Gate--2
transfer requires
\[
\|E_2\|_\infty+\|\mathcal D E_2\|_\infty\le Q^{-A_E}.
\]
The canonical curvature \(q''_{\can}\) may contain \(\Xi'/\Xi\), fixed
completed-log-derivative derivatives, and regular gamma/background terms.  It
must not contain denominators such as
\[
1/q,\qquad 1/u,\qquad 1/\Gamma_G,\qquad
1/\Gamma_{G,2},\qquad 1/|\Gamma|.
\]
Those denominators, if present, belong to later row-projection or quotient
algebra, not to the canonical source curvature used in Gate~2.
\end{remark}

\begin{lemma}[Real-canonical Gate--2 compatibility]
\label{lem:real-canonical-gate-two-compatibility}
Assume the actual-zeta local variables are in the real completed gauge of
Definition~\ref{def:canonical-actual-zeta-phase-curvature}.
If this gauge is implemented through a local chart, assume throughout the
packet that every chart denominator is bounded below by \(Q^{-C}\).
Thus
\[
q=q_\zeta,\qquad u=q''_{\can}=\mathcal D^2q_\zeta,
\]
and the canonical anisotropies are
\[
\Gamma_G=q''_{\can}-6q_\zeta+2q_\zeta^3,
\]
\[
\Gamma_{G,2}
=(1-q_\zeta^2)q''_{\can}-\frac16\mathcal D^4q_\zeta.
\]
Then, on ordinary smooth zero-separated packets,
\[
q,\quad u,\quad \Gamma_G,\quad \Gamma_{G,2}
\]
are real-valued and, with \(g=\Gamma_G\), \(h=\Gamma_{G,2}\),
\[
u=q'',
\]
\[
g=q''-6q+2q^3,
\]
and
\[
h=2q(q')^2-\frac16g''.
\]
\end{lemma}

\begin{proof}
The real-valuedness follows from the real completed critical-line gauge.  The
first two identities are the definitions.  For the final identity, apply
Lemma~\ref{lem:differentiated-form-second-gate-two-anisotropy} after suppressing
the \(\zeta\)-subscript:
\[
\Gamma_{G,2}
=
2q(\mathcal Dq)^2-\frac16\mathcal D^2\Gamma_G.
\]
In ordinary local prime notation this is
\[
h=2q(q')^2-\frac16g''.
\]
\end{proof}

\begin{corollary}[Gauge-locking for the \(q_0\)-wall]
\label{cor:gauge-locking-q0-wall}
In the real-canonical convention above, the \(q_0\)-wall used in wall escape is
the same scalar wall that appears in the Gate--2 anisotropy and in the D4
source factor.  Namely, for \(q_0=q_\zeta(m)\),
\[
u=\mathcal D^2q_0,
\]
\[
\Gamma_G=u-6q_0+2q_0^3,
\]
and the D4 source factor has leading form
\[
H_{\mathrm{lead}}=-\frac{q_0u^3}{288}.
\]
Thus \(q_0\) is not an arbitrary affine primitive of \(u\).  Any affine drift
\(q_0\mapsto q_0+am+b\) changes both
\[
\Gamma_G=u-6q_0+2q_0^3
\]
and
\[
H_{\mathrm{lead}}=-q_0u^3/288.
\]
Consequently, the wall-escape argument is applied only in the canonical gauge
fixed by \(q_\zeta\).
\end{corollary}

\begin{remark}[Residual version of the compatibility package]
\label{rem:gate-two-compatibility-residuals}
If an approximate convention is used, define
\[
        r_C:=u-q'',
\]
\[
        r_G:=\Gamma_G-(q''-6q+2q^3),
\]
and
\[
        r_{G,2}:=
        \Gamma_{G,2}
        -
        \left(2q(q')^2-\frac16\Gamma_G''\right).
\]
The main canonical convention of this draft has
\[
        r_C=r_G=r_{G,2}=0.
\]

If a residual convention is used instead, then the \(L^2\)-input rigidity
argument applies to \(q''\), while the Gate--2 source is \(u\).  Therefore one
must also control the curvature discrepancy \(r_C=u-q''\).  A sufficient residual package
on the source interval \(J\) is
\[
        \|r_C\|_{L^2(J)}
        +\|r_G\|_{L^2(J)}
        +\|r_{G,2}\|_{L^2(J)}
        \le
        \varepsilon |J|^{1/2}.
\]
Under this package, the residual terms enter the kinetic, potential, and
small-\(q\) estimates exactly like \(g\) and \(h\), and the conclusion of
Lemma~\ref{lem:l2-input-real-rigidity-safe} remains valid after increasing the
implicit polynomial loss.
\end{remark}

\begin{lemma}[Optional primitive recovery from \(L^2\)-small \(g\); not used below]
\label{lem:local-primitive-recovery-from-l2-small-g}
Let \(J_1\Subset J\) be intervals with polynomial separation from the boundary.
Assume
\[
u=q'',\qquad g=u-F(q),\qquad F(q)=6q-2q^3,
\]
and
\[
\|u\|_{L^\infty(J)}\le Q^C,
\qquad
\|g\|_{L^2(J)}\le \varepsilon |J|^{1/2}.
\]
Then on \(J_1\),
\[
\|q\|_{L^\infty(J_1)}\le Q^C.
\]
Moreover, on any further polynomially shrunken interval \(J_2\Subset J_1\),
\[
\|q'\|_{L^\infty(J_2)}\le Q^C.
\]
\end{lemma}

\begin{proof}
Since \(F(q)=u-g\),
\[
\|F(q)\|_{L^2(J)}
\le
\|u\|_{L^2(J)}+\|g\|_{L^2(J)}
\le Q^C|J|^{1/2}.
\]
The cubic polynomial \(F(q)=6q-2q^3\) satisfies
\[
|q|^6\le C(1+|F(q)|^2).
\]
Hence
\[
\|q\|_{L^6(J)}\le Q^C|J|^{1/6}.
\]
Since \(q''=u\in L^\infty(J)\), the one-dimensional interpolation estimate on
the polynomially shrunken interval \(J_1\Subset J\) gives
\[
\|q\|_{L^\infty(J_1)}
\le
C\left(
\|q\|_{L^6(J)}
+
\|q\|_{L^6(J)}^{\frac35}
\|q''\|_{L^\infty(J)}^{\frac25}
+
1
\right),
\]
after rescaling \(J\) to a bounded interval and absorbing the polynomial
shrinkage factor into \(Q^C\).  Hence
\[
\|q\|_{L^\infty(J_1)}\le Q^C.
\]
Finally, the Landau estimate gives
\[
\|q'\|_{L^\infty(J_2)}
\le
C\left(
\ell^{-1}\|q\|_{L^\infty(J_1)}
+
\ell\|q''\|_{L^\infty(J_1)}
\right)
\le Q^C,
\]
where \(\ell=\operatorname{dist}(J_2,\partial J_1)\) and the shrinkage is
polynomial.
\end{proof}

\begin{lemma}[\(L^2\)-input real rigidity; safe primitive-bounded form]
\label{lem:l2-input-real-rigidity-safe}
Let
\[
        J=(-2,2),\qquad J'=(-1,1).
\]
Assume \(q,u,g,h\in\mathbb R\) are sufficiently smooth on \(J\), in the
real-canonical convention of
Lemma~\ref{lem:real-canonical-gate-two-compatibility}, and
\[
        u=q'',
\]
\[
        g=q''-6q+2q^3,
\]
\[
        h=2q(q')^2-\frac16g''.
\]
Assume the safe local bounds
\[
        \|q\|_{L^\infty(J)}
        +\|q'\|_{L^\infty(J)}
        +\|q''\|_{L^\infty(J)}
        \le Q^C,
\]
In the active Gate--2 proof this safe bound is assumed directly on the
normalized source interval.  Lemma~\ref{lem:local-primitive-recovery-from-l2-small-g}
records a possible strengthening, but it is not used to justify this lemma in
the main argument.  Assume also the \(L^2\)-smallness condition
\[
        \|g\|_{L^2(J)}+\|h\|_{L^2(J)}
        \le \varepsilon |J|^{1/2}.
\]
Then, after increasing \(C\) if necessary,
\[
        \int_{J'} |u|^2
        \le Q^C\varepsilon^{2/5}|J'|.
\]
\end{lemma}

\begin{proof}
If \(\varepsilon\ge1\), the conclusion follows immediately from
\(\|u\|_\infty=\|q''\|_\infty\le Q^C\).  Assume \(0<\varepsilon<1\), and set
\[
        F(q):=6q-2q^3.
\]
Since \(g=q''-6q+2q^3\), we have
\[
        u=q''=F(q)+g.
\]

Choose
\[
        \phi=\psi^2\in C_c^2(J),
        \qquad
        0\le \phi\le1,
        \qquad
        \phi\equiv1\quad\text{on }J',
\]
with
\[
        |\phi'|^2\le Q^C\phi,
        \qquad
        |\phi''|\le Q^C.
\]
Let \(0<\delta<1\), to be chosen later.  Let
\(\chi_\delta:\mathbb R\to[0,1]\) be smooth with
\[
        \chi_\delta(x)=0\quad (|x|\le\delta),
        \qquad
        \chi_\delta(x)=1\quad (|x|\ge2\delta),
\]
and
\[
        |\chi_\delta'|\le C\delta^{-1},
        \qquad
        |\chi_\delta''|\le C\delta^{-2}.
\]
Define
\[
        a_\delta(q):=q\chi_\delta(q).
\]
Then
\[
        a_\delta'(q)=\chi_\delta(q)+q\chi_\delta'(q)=O(1),
\]
and
\[
        a_\delta''(q)=2\chi_\delta'(q)+q\chi_\delta''(q)=O(\delta^{-1}),
\]
because on the transition region \(\delta\le |q|\le2\delta\) one has
\(|q|\asymp\delta\).

We first prove the weighted kinetic estimate.  From
\[
        h=2q(q')^2-\frac16g'',
\]
multiply by \(\phi a_\delta(q)\) and integrate over \(J\):
\[
        \int_J 2\phi a_\delta(q)q(q')^2
        =
        \int_J \phi a_\delta(q)h
        +
        \frac16\int_J \phi a_\delta(q)g''.
\]
After integrating by parts twice,
\[
        \int_J \phi a_\delta(q)g''
        =
        \int_J g(\phi a_\delta(q))''.
\]
Thus
\[
        \int_J 2\phi\chi_\delta(q)q^2(q')^2
        =
        \int_J \phi a_\delta(q)h
        +
        \frac16\int_J g(\phi a_\delta(q))''.
\]
The first term satisfies
\[
        \left|\int_J \phi a_\delta(q)h\right|
        \le
        \|\phi a_\delta(q)\|_{L^\infty(J)}
        \|h\|_{L^2(J)}
        |J|^{1/2}
        \le
        Q^C\varepsilon |J|.
\]
Next,
\[
        (\phi a_\delta(q))''
        =
        \phi''a_\delta(q)
        +2\phi'a_\delta'(q)q'
        +\phi a_\delta''(q)(q')^2
        +\phi a_\delta'(q)q''.
\]
Using the bounds on \(a_\delta,a_\delta',a_\delta''\), and the safe bounds on
\(q,q',q''\), we get
\[
        \|(\phi a_\delta(q))''\|_{L^2(J)}
        \le Q^C\delta^{-1}|J|^{1/2}.
\]
Therefore
\[
        \left|\int_J g(\phi a_\delta(q))''\right|
        \le
        Q^C\varepsilon\delta^{-1}|J|.
\]
Consequently,
\[
        \int_J\phi\chi_\delta(q)q^2(q')^2
        \le
        Q^C\varepsilon\delta^{-1}|J|.
\]
On \(|q|\ge2\delta\), \(\chi_\delta(q)=1\) and \(q^2\ge4\delta^2\), so
\[
        \int_{\{|q|\ge2\delta\}}\phi(q')^2
        \le
        Q^C\varepsilon\delta^{-3}|J|.
        \tag{1}
\]

Let \(\theta_\delta:\mathbb R\to[0,1]\) be smooth with
\[
        \theta_\delta(q)=0\quad (|q|\le2\delta),
        \qquad
        \theta_\delta(q)=1\quad (|q|\ge3\delta),
\]
and
\[
        |\theta_\delta'|\le C\delta^{-1}.
\]
Define
\[
        X:=\int_J \phi\,\theta_\delta(q)F(q)^2.
\]
Since \(F(q)=q''-g\),
\[
        X
        =
        \int_J \phi\theta_\delta Fq''
        -
        \int_J \phi\theta_\delta Fg.
\]
The \(g\)-term obeys
\[
        \left|\int_J \phi\theta_\delta Fg\right|
        \le
        X^{1/2}\|g\|_{L^2(J)}
        \le
        X^{1/2}\varepsilon |J|^{1/2},
\]
and hence, after absorption,
\[
        \left|\int_J \phi\theta_\delta Fg\right|
        \le
        \frac14X+Q^C\varepsilon^2|J|.
        \tag{2}
\]

For the main term,
\[
        \int_J \phi\theta_\delta Fq''
        =
        -\int_J(\phi\theta_\delta F)'q',
\]
where
\[
        (\phi\theta_\delta F)'
        =
        \phi'\theta_\delta F
        +
        \phi\theta_\delta'(q)q'F
        +
        \phi\theta_\delta F'(q)q'.
\]

For the boundary term, using \(|\phi'|^2\le Q^C\phi\) and (1),
\[
\begin{aligned}
\left|\int_J\phi'\theta_\delta Fq'\right|
&\le
Q^C
\left(\int_J\phi\theta_\delta F^2\right)^{1/2}
\left(
\int_{\{|q|\ge2\delta\}}\phi(q')^2
\right)^{1/2}  \\
&\le
Q^CX^{1/2}(\varepsilon\delta^{-3}|J|)^{1/2}  \\
&\le
\frac14X+Q^C\varepsilon\delta^{-3}|J|.
\end{aligned}
\tag{3}
\]
For the transition term, on \(\operatorname{supp}\theta_\delta'\) one has
\(2\delta\le |q|\le3\delta\), so
\[
        |F(q)|=|6q-2q^3|\le C\delta,
        \qquad
        |\theta_\delta'|\le C\delta^{-1}.
\]
Thus \(|\theta_\delta'F|\le C\), and by (1),
\[
        \left|\int_J\phi\theta_\delta'(q)F(q)(q')^2\right|
        \le
        Q^C\varepsilon\delta^{-3}|J|.
        \tag{4}
\]
Finally, since \(F'(q)=6-6q^2\) and \(\|q\|_{L^\infty(J)}\le Q^C\),
\[
        \left|\int_J\phi\theta_\delta F'(q)(q')^2\right|
        \le
        Q^C\int_{\{|q|\ge2\delta\}}\phi(q')^2
        \le
        Q^C\varepsilon\delta^{-3}|J|.
        \tag{5}
\]
Combining (2)--(5) and absorbing,
\[
        X
        \le
        Q^C(\varepsilon\delta^{-3}+\varepsilon^2)|J|.
        \tag{6}
\]
Since \(\phi=1\) on \(J'\) and \(\theta_\delta=1\) on \(|q|\ge3\delta\),
\[
        \int_{J'\cap\{|q|\ge3\delta\}}F(q)^2
        \le
        Q^C(\varepsilon\delta^{-3}+\varepsilon^2)|J'|.
        \tag{7}
\]

On \(|q|\le3\delta\), \(|F(q)|\le C\delta\).  Since \(u=F(q)+g\),
\[
        \int_{J'\cap\{|q|\le3\delta\}}|u|^2
        \le
        Q^C(\delta^2+\varepsilon^2)|J'|.
        \tag{8}
\]
Combining (7) and (8),
\[
        \int_{J'}|u|^2
        \le
        Q^C(\delta^2+\varepsilon\delta^{-3}+\varepsilon^2)|J'|.
\]
Choose
\[
        \delta=\varepsilon^{1/5}.
\]
Then
\[
        \delta^2=\varepsilon^{2/5},
        \qquad
        \varepsilon\delta^{-3}=\varepsilon^{2/5}.
\]
Since \(\varepsilon^2\le\varepsilon^{2/5}\) for \(0<\varepsilon<1\),
\[
        \int_{J'}|u|^2
        \le
        Q^C\varepsilon^{2/5}|J'|.
\]
\end{proof}

\begin{remark}[Residual-stable form]
\label{rem:l2-rigidity-residual-stable}
The same estimate holds if the identities hold with \(L^2\)-controlled
residuals:
\[
        r_C:=u-q'',
\]
\[
        g=q''-6q+2q^3+r_g,
\]
and
\[
        h=2q(q')^2-\frac16g''+r_h,
\]
where
\[
        \|r_C\|_{L^2(J)}
        +
        \|r_g\|_{L^2(J)}
        +
        \|r_h\|_{L^2(J)}
        \le
        \varepsilon |J|^{1/2}.
\]
Indeed, \(r_h\) enters the weighted kinetic estimate in the same way as \(h\),
while \(r_g\) enters the potential estimate in the same way as \(g\).  The
additional residual \(r_C\) is needed only at the final comparison between the
ideal curvature \(q''\) and the source \(u\).  After increasing \(C\), the same
conclusion follows:
\[
        \int_{J'}|u|^2
        \le Q^C\varepsilon^{2/5}|J'|.
\]
In the exact real-canonical convention used in the main argument,
\[
r_C=r_g=r_h=0.
\]
\end{remark}

\begin{remark}[Primitive recovery not used in the rigidity proof]
\label{rem:l2-rigidity-primitive-free-not-used}
Lemma~\ref{lem:l2-input-real-rigidity-safe} is used only in the safe
primitive-bounded form
\[
\|q\|_\infty+\|q'\|_\infty+\|q''\|_\infty\le Q^C.
\]
Lemma~\ref{lem:local-primitive-recovery-from-l2-small-g} records a possible
way to recover these bounds on shrunken intervals from stronger source and
\(L^2\)-smallness hypotheses, but this recovery is not part of the active
Gate--2 proof.  The proof of
Theorem~\ref{thm:gate-two-from-gate-one-l2-rigidity} uses the displayed safe
primitive bounds as hypotheses on the normalized source interval.
\end{remark}

\begin{remark}[Application on polynomial subintervals]
\label{rem:scale-normalized-use-l2-rigidity}
Lemma~\ref{lem:l2-input-real-rigidity-safe} is stated on the fixed model
intervals
\[
        J=(-2,2),\qquad J'=(-1,1),
\]
for notational convenience.  In applications to a physical subinterval
\(J_{\rm phys}\subset I\), the Gate--2 identities are still interpreted in the
original packet-normalized derivative convention
\[
        \mathcal D=\lambda\partial_m,
\]
not in a newly rescaled derivative unless the variables \(q,u,g,h\) are also
renormalized and the identities are restated.

Thus, when applying the lemma on \(J_{\rm phys}\), we keep
\[
u=q'',\qquad
g=q''-6q+2q^3,\qquad
h=2q(q')^2-\frac16g''
\]
with primes denoting the original \(\mathcal D\)-derivative.  The cutoff
functions used in the proof then have derivatives of size polynomially bounded
in \(|J_{\rm phys}|^{-1}\).  Since the source-thickening interval has
polynomial relative length,
\[
|J_{\rm phys}|\ge Q^{-C}|I|,
\]
these cutoff losses are absorbed into the \(Q^C\) factor in
Lemma~\ref{lem:l2-input-real-rigidity-safe}.  No change of the canonical
Gate--2 derivative convention is made.
\end{remark}

\begin{lemma}[Source thickening from one large source point]
\label{lem:source-thickening-one-large-point}
Let \(I\) be a packet interval and let \(u\) be a \(C^1\), real- or
complex-valued function on \(I\).  Assume
\[
\int_I |u|^2\,dm\ge Q^{-B}|I|
\]
and
\[
\|u'\|_{L^\infty(I)}\le Q^{C_u}|I|^{-1}.
\]
Then there exists an interval \(J\subset I\), one-sided if necessary, such that
\[
|u|\ge \tau:=\frac12Q^{-B/2}
\qquad\text{on }J,
\]
and
\[
|J|\ge Q^{-B/2-C_u-O(1)}|I|.
\]
\end{lemma}

\begin{proof}
Gate~1 gives \(m_\ast\in I\) with \(|u(m_\ast)|\ge Q^{-B/2}\).  Since
\[
\bigl||u(m)|-|u(m_\ast)|\bigr|\le |u(m)-u(m_\ast)|,
\]
the Lipschitz bound gives \(|u|\ge \frac12Q^{-B/2}\) whenever
\[
|m-m_\ast|\le cQ^{-B/2-C_u}|I|
\]
with \(c>0\) sufficiently small.  Intersecting with \(I\) gives the claimed
one-sided or two-sided interval.
\end{proof}

\begin{lemma}[Source-large interval with \(L^2\)-small anisotropy]
\label{lem:source-large-interval-l2-anisotropy}
Assume Gate~1:
\[
\int_I |u|^2\,dm\ge Q^{-B}|I|,
\]
source Lipschitz control
\[
\|u'\|_{L^\infty(I)}\le Q^{C_u}|I|^{-1},
\]
and Gate--2 failure at exponent \(A\):
\[
\int_I |u|^2(g^2+h^2)\,dm\le Q^{-A}|I|.
\]
Then there exists an interval \(J\subset I\) such that
\[
|u|\ge \tau:=\frac12Q^{-B/2}
\qquad\text{on }J,
\]
\[
|J|\ge Q^{-B/2-C_u-O(1)}|I|,
\]
and
\[
\frac1{|J|}\int_J(g^2+h^2)\,dm
\le
Q^{-A+\frac32B+C_u+O(1)}.
\]
\end{lemma}

\begin{proof}
By Lemma~\ref{lem:source-thickening-one-large-point}, there is an interval
\(J\subset I\) with the stated lower bound for \(|u|\) and for \(|J|\).  Since
\(|u|^2\ge \tau^2\) on \(J\),
\[
\tau^2\int_J(g^2+h^2)\,dm
\le
\int_J |u|^2(g^2+h^2)\,dm
\le
\int_I |u|^2(g^2+h^2)\,dm
\le Q^{-A}|I|.
\]
Using \(\tau^{-2}\asymp Q^B\) and
\[
|J|^{-1}\le Q^{B/2+C_u+O(1)}|I|^{-1}
\]
gives the claim.
\end{proof}

\subsection{Gate 2 from source thickening and \(L^2\)-input rigidity}

\begin{theorem}[Gate 2 from Gate 1 and \(L^2\)-input rigidity]
\label{thm:gate-two-from-gate-one-l2-rigidity}
Assume the normalized real-canonical convention of
Definitions~\ref{def:canonical-actual-zeta-phase-curvature} and
\ref{def:canonical-gate-two-anisotropy-functionals}.
This is an ordinary smooth-packet statement.  Packets containing an
unregularized critical-line zero, or packets where the real completed phase is
not smooth in the chosen convention, are routed to the singular/transition
branch and are not covered by this theorem.
Let
\[
u=q''_{\can}=q'',
\qquad
g=\Gamma_G=q''-6q+2q^3,
\qquad
h=\Gamma_{G,2}=2q(q')^2-\frac16g''.
\]
Assume Gate~1:
\[
\int_I |u|^2\,dm\ge Q^{-B}|I|.
\]
Assume source Lipschitz control
\[
\|u'\|_{L^\infty(I)}\le Q^{C_u}|I|^{-1},
\]
Assume moreover that every source interval produced by
Lemma~\ref{lem:source-large-interval-l2-anisotropy} can be placed in the
normalized setting of Lemma~\ref{lem:l2-input-real-rigidity-safe}, with the
safe primitive bounds
\[
\|q\|_\infty+\|q'\|_\infty+\|q''\|_\infty\le Q^C
\]
available on that normalized interval, and with rigidity loss \(Q^{C_R}\).
Then Gate~2 holds with any exponent
\[
\boxed{
A>\frac{13}{2}B+C_u+5C_R+O(1).
}
\]
Equivalently, for such \(A\),
\[
\boxed{
\int_I |u|^2(g^2+h^2)\,dm\ge Q^{-A}|I|.
}
\]
\end{theorem}

\begin{proof}
Assume, toward contradiction, that
\[
\int_I |u|^2(g^2+h^2)\,dm\le Q^{-A}|I|.
\]
By Lemma~\ref{lem:source-large-interval-l2-anisotropy}, there is an interval
\(J\subset I\) such that
\[
|u|\ge \tau:=\frac12Q^{-B/2}\quad\text{on }J,
\]
and
\[
\frac1{|J|}\int_J(g^2+h^2)\,dm
\le
Q^{-A+\frac32B+C_u+O(1)}.
\]
Put
\[
\varepsilon^2
=
Q^{-A+\frac32B+C_u+O(1)}.
\]
Then
\[
\|g\|_{L^2(J)}+\|h\|_{L^2(J)}
\le
Q^{O(1)}\varepsilon |J|^{1/2}.
\]
Absorbing the harmless factor into \(C_R\), Lemma~\ref{lem:l2-input-real-rigidity-safe}
gives on \(J'\Subset J\)
\[
\int_{J'}|u|^2\,dm
\le
Q^{C_R}\varepsilon^{2/5}|J'|.
\]
But \(J'\subset J\) and \(|u|\ge\tau\) on \(J\), so
\[
\int_{J'}|u|^2\,dm\ge Q^{-B-O(1)}|J'|.
\]
Using
\[
\varepsilon^{2/5}
=
Q^{-\frac15A+\frac3{10}B+\frac15C_u+O(1)},
\]
the two estimates are incompatible if
\[
C_R-\frac15A+\frac3{10}B+\frac15C_u+O(1)<-B.
\]
It suffices that
\[
A>\frac{13}{2}B+C_u+5C_R+O(1).
\]
Thus Gate~2 failure is impossible for such \(A\).
\end{proof}

\begin{remark}[Replacement of the old \(D1/D2\) downstream package]
\label{rem:gate-two-replaces-old-d1-d2}
Theorem~\ref{thm:gate-two-from-gate-one-l2-rigidity} replaces the older
``density/Remez propagation'' and ``pointwise derivative propagation''
downstream inputs.  The active local route is
\[
\boxed{
\text{Gate 1}
+
|u'|_\infty\text{ source thickening}
+
L^2\text{-input real rigidity}
\Rightarrow
\text{Gate 2}.
}
\]
No Remez--Cartan theorem, pointwise target-smallness interval, or pointwise
\(\Gamma_G''\)-propagation is used in this route.
\end{remark}

\subsection{Gauge-locked wall escape and row energy}

\begin{definition}[Gauge-locked wall scalar]
\label{def:gauge-locked-wall-scalar}
The wall scalar \(q_0\) is gauge-locked to the Gate--2 and D4 package if the
same scalar \(q_0\) satisfies, in the same normalized derivative convention,
\[
q_0''=u=q''_{\can},
\]
\[
\Gamma_G=u-6q_0+2q_0^3,
\]
and the D4 source factor has leading form
\[
H\sim-\frac{q_0u^3}{288}.
\]
If these identities hold only up to residuals, the residuals must satisfy
\[
|q_0''-u|\le c\tau,
\]
and
\[
\left|
\Gamma_G-(u-6q_0+2q_0^3)
\right|
\le c\tau
\]
on wall-escape neighborhoods.
\end{definition}

\begin{theorem}[Gauge-locked wall escape in normalized coordinates]
\label{thm:gauge-locked-wall-escape-row-energy}
Fix the normalized derivative convention
\[
\mathcal D=\lambda\partial_m.
\]
Let
\[
q=q_0,\qquad u=q''_{\can},
\qquad
\Gamma=(\Gamma_G,\Gamma_{G,2}),
\qquad
G:=|u|^2|\Gamma|^2.
\]
Assume Gate~2 mass on \(I\):
\[
\int_I G(m)\,dm\ge Q^{-A}|I|.
\]
Assume polynomial upper bounds
\[
G\le Q^{C_G},
\qquad
|\Gamma|\le Q^{C_\Gamma},
\qquad
|\mathcal D u|\le Q^{C_{u'}}.
\]
Define the gauge residuals
\[
r_C:=u-\mathcal D^2q_0,
\]
\[
r_G:=
\Gamma_G-\left(\mathcal D^2q_0-6q_0+2q_0^3\right),
\]
and
\[
r_H^\ast:=H+\frac{q_0u^3}{288}.
\]
Assume on wall-escape neighborhoods that
\[
|r_C|\le c\tau,
\qquad
|r_G-r_C|\le c\tau,
\]
and that \(r_H^\ast\) is controlled by the D4 residual estimate.

Choose thresholds
\[
\tau=Q^{-C_\tau},
\qquad
\varepsilon=Q^{-C_\varepsilon},
\]
with
\[
2C_\tau\ge A+2C_\Gamma+O(1),
\]
and
\[
C_\varepsilon
\ge
\max\{3C_\tau+2C_{u'},\ C_\tau+2A+2C_G\}+O(1).
\]

Then either the non-wall region
\[
\{|q_0|\ge\varepsilon\}
\]
already supplies a non-wall D4 sample, or there exists \(m_1\in I\) such that
\[
|q_0(m_1)|=\varepsilon,
\qquad
|u(m_1)|\ge c\tau,
\qquad
|\Gamma_G(m_1)|\ge c\tau.
\]
Consequently
\[
G(m_1)\ge c\tau^4.
\]
If D4 row conversion holds at non-wall samples and row mobility holds, then
\[
\boxed{
\int_I|\mathbf A_N(m)|^2\,dm\ge Q^{-C}|I|.
}
\]
\end{theorem}

\begin{proof}
Split
\[
I_{\rm nd}=\{|q_0|\ge\varepsilon\},
\qquad
I_{\rm wall}=\{|q_0|<\varepsilon\}.
\]
If \(I_{\rm nd}\) carries a fixed polynomial portion of Gate~2 mass, then
\(G\le Q^{C_G}\) gives a point \(m_\ast\in I_{\rm nd}\) with polynomially live
\(G(m_\ast)\), and D4 applies directly.

Assume instead that wall mass dominates.  Remove the low-source set
\[
|u|<\tau.
\]
Choosing
\[
2C_\tau\ge A+2C_\Gamma+O(1)
\]
ensures this low-source set cannot carry all wall mass.  Hence there is a
wall-live point \(m_0\) with
\[
|q_0(m_0)|<\varepsilon,
\qquad
|u(m_0)|\ge\tau,
\]
away from the packet boundary after the standard buffered-packet reduction.

For complex \(q_0\), set
\[
e=\frac{\overline{u(m_0)}}{|u(m_0)|},
\qquad
Q(m):=\Re(e q_0(m)).
\]
Since
\[
\mathcal D^2Q(m)=\Re\bigl(e(\mathcal D^2q_0)(m)\bigr)
=\Re\bigl(e(u-r_C)(m)\bigr),
\]
the assumptions \(|r_C|\le c\tau\) and
\[
|\mathcal D u|\,r_{\mathrm{esc}}\le c\tau
\]
imply
\[
\mathcal D^2Q\ge c\tau
\]
on a normalized interval of radius \(r_{\mathrm{esc}}\).  Taking
\[
r_{\mathrm{esc}}\asymp\sqrt{\varepsilon/\tau}
\]
gives first exit from the strip \(|q_0|<\varepsilon\) within normalized
distance \(O(r_{\mathrm{esc}})\), provided
\[
C_\varepsilon\ge 3C_\tau+2C_{u'}+O(1).
\]
Boundary and buffer loss are controlled by choosing
\[
C_\varepsilon\ge C_\tau+2A+2C_G+O(1).
\]
Thus there exists \(m_1\) with
\[
|q_0(m_1)|=\varepsilon,
\qquad
|u(m_1)|\ge c\tau.
\]

Finally,
\[
\Gamma_G
=
u-6q_0+2q_0^3+(r_G-r_C).
\]
At \(m_1\), if \(\varepsilon\ll\tau\) and
\[
|r_G-r_C|\le c\tau,
\]
then
\[
|\Gamma_G(m_1)|\ge c\tau.
\]
Thus \(G(m_1)\ge c\tau^4\).  D4 converts this non-wall sample to a row lower
bound, and row mobility thickens the pointwise lower bound to interval row
energy.
\end{proof}

\begin{remark}[Absolute residual corridor in wall escape]
\label{rem:absolute-residual-corridor-wall-escape}
In Theorem~\ref{thm:gauge-locked-wall-escape-row-energy}, D4 residual control
is interpreted valuation-wise.  Same-valuation anisotropy-linear terms are
absorbed into the effective D4 row matrix, and valuation-relative residuals are
estimated relative to the leading row scale.  Only absolute residual floors
require a wall-threshold corridor.

Concretely, if the D4 leading valuation is
\[
q_0u^{b_\ast}L_\ast\Gamma
\]
and the wall-escape sample satisfies
\[
|q_0|=\varepsilon,\qquad |u|\gg\tau,\qquad |\Gamma|\gg\tau,
\]
then an absolute residual bound
\[
|\mathbf R_{\rm abs}|\le Q^{-C_R^{\rm abs}}
\]
must satisfy
\[
C_R^{\rm abs}
\ge
C_\varepsilon+(b_\ast+1)C_\tau+C_L+O(1),
\]
where \(Q^{-C_L}\) is the D4 row-rank lower bound.  If D4 uses a filtered
leading block through \(u\)-valuation \(B\), replace \(b_\ast\) by \(B\).
This condition does not apply to residuals carrying useful \(q_0,u,\Gamma\)
valuation.
\end{remark}

\begin{remark}[Packet buffer for wall escape]
\label{rem:packet-buffer-for-wall-escape}
The first-exit wall-escape step moves by a distance
\[
\rho |I|,
\qquad
\rho\asymp\sqrt{\varepsilon/\tau}.
\]
The packet used for row conversion and local-complexity estimates should
therefore be chosen with a buffer of at least this scale.  If the boundary
layer has relative size \(\eta\), and the Gate--2 density is bounded above by
\[
G=|u|^2|\Gamma|^2\le Q^{C_G},
\]
then discarding the edge changes the Gate--2 mass by at most
\(Q^{C_G}\eta |I|\).  Hence a boundary layer with
\[
\eta\le Q^{-A-C_G-O(1)}
\]
is negligible at Gate--2 mass scale \(Q^{-A}|I|\).  This converts packet-edge
effects into exponent bookkeeping rather than a new structural obstruction.
\end{remark}

\begin{remark}[D3 is replaced by wall escape plus D4]
\label{rem:d3-replaced-by-wall-escape-plus-d4}
The old source-robust \(q_0\)-wall row is not the active route in the
row-energy architecture.  Wall escape is the analytic version of saturating by
\(q_0\): instead of dividing by \(q_0\), we move to a nearby point where
\[
|q_0|=\varepsilon
\]
and use the original rows.  The source-robust wall-row search remains a
fallback only if the final contradiction requires row conversion at a fixed
preassigned point.
\end{remark}

\begin{lemma}[Row mobility: pointwise row output thickens to row energy]
\label{lem:row-mobility-point-to-energy}
Let \(\mathbf A_N:I\to\mathbb C^M\) be a finite row vector.  Suppose
\[
|\mathbf A_N(m_\ast)|\ge Q^{-C_A}
\]
for some \(m_\ast\in I\), and
\[
|\mathbf A_N'(m)|\le Q^{C_D}|I|^{-1}
\qquad (m\in I).
\]
Then
\[
\int_I|\mathbf A_N(m)|^2\,dm
\ge Q^{-C'}|I|
\]
for some polynomial exponent \(C'\).
\end{lemma}

\begin{proof}
For
\[
|m-m_\ast|\le \frac12Q^{-C_A-C_D}|I|,
\]
the derivative bound gives
\[
|\mathbf A_N(m)-\mathbf A_N(m_\ast)|\le \frac12Q^{-C_A}.
\]
Thus
\[
|\mathbf A_N(m)|\ge \frac12Q^{-C_A}
\]
on a subinterval of length at least \(Q^{-C_A-C_D-O(1)}|I|\), after taking a
one-sided subinterval if \(m_\ast\) is near the boundary.  Integrating the
square over that subinterval gives the claim.
\end{proof}

\begin{remark}[Quotient-frame version of row mobility]
\label{rem:quotient-frame-row-mobility}
If the row output is measured in a moving quotient row space
\(\mathcal Q_N(m)=\mathcal H_N/\mathcal C_N(m)\), the same argument applies to
\[
P_{\mathcal C_N(m)^\perp}\mathbf A_N(m)
\]
provided the nuisance frame is polynomially stable:
\[
\lambda_{\min}(C(m)^*C(m))\ge Q^{-C},
\qquad
\|C'(m)\|\le Q^C|I|^{-1}.
\]
If raw or carrier-row energy is accepted downstream, no moving-frame derivative
estimate is needed.
\end{remark}

\begin{remark}[Frozen quotient coordinates avoid moving-frame derivatives]
\label{rem:frozen-quotient-coordinates-row-mobility}
If the quotient frame \(\mathcal Q_N(m)\) varies with \(m\), one can avoid
moving-frame derivative estimates by freezing the quotient coordinates at the
D4 sample point \(m_\ast\).  Suppose
\[
|[\mathbf A_N(m_\ast)]|_{\mathcal Q_N(m_\ast)}\ge Q^{-C_A}
\]
and the quotient Gram matrix at \(m_\ast\) is polynomially nondegenerate.  Let
\(\ell_7^\ast,\ell_9^\ast\) be dual quotient coordinates at \(m_\ast\), extended
as fixed linear functionals on the ambient finite row space.  Then
\[
|\ell_7^\ast(\mathbf A_N(m_\ast))|^2
+
|\ell_9^\ast(\mathbf A_N(m_\ast))|^2
\ge Q^{-C}.
\]
If the ambient row vector satisfies
\[
|\mathbf A_N'(m)|\le Q^{C_D}|I|^{-1},
\]
then the two fixed scalar functions
\[
\ell_7^\ast(\mathbf A_N(m)),
\qquad
\ell_9^\ast(\mathbf A_N(m))
\]
obey the same row-mobility argument, yielding
\[
\int_I
\left(
|\ell_7^\ast(\mathbf A_N(m))|^2+
|\ell_9^\ast(\mathbf A_N(m))|^2
\right)\,dm
\ge Q^{-C'}|I|.
\]
Thus carrier-quotient row energy can be proved using fixed dual coordinates at
the sample point, without differentiating the moving quotient frame.
\end{remark}

\begin{theorem}[Local downstream closure from Gate--2 and D4]
\label{thm:local-downstream-closure-from-gate2-d4}
Assume Gate~2 on \(I\):
\[
\int_I |u|^2|\Gamma|^2\,dm\ge Q^{-A}|I|.
\]
Assume the ordinary real-canonical local complexity bounds and the gauge-locking
hypotheses needed for
Theorem~\ref{thm:gauge-locked-wall-escape-row-energy}.  Assume D4 row
conversion at non-wall samples in a chosen row-output space
\[
\mathcal Q_{\rm row},
\]
where \(\mathcal Q_{\rm row}\) is either the carrier quotient
\(\mathcal Q_{\rm car}\), a raw/ambient row space, or another filtered row
space with polynomially equivalent norm.  Finally assume row mobility as in
Lemma~\ref{lem:row-mobility-point-to-energy}, or the quotient-frame version of
Remark~\ref{rem:quotient-frame-row-mobility} if the quotient varies with \(m\).

Then the sampled packet carries polynomial row energy in the chosen row-output
space:
\[
\boxed{
\int_I
\left\|
\mathcal A_N(m)
\right\|_{\mathcal Q_{\rm row}}^2\,dm
\ge Q^{-C}|I|.
}
\]
In the active carrier-quotient formulation this is
\[
\boxed{
\int_I
\left|
[\mathcal A_N^{\rm car}(m)]
\right|_{\mathcal Q_{\rm car}}^2\,dm
\ge Q^{-C}|I|.
}
\]

If the row-output space is represented by a concrete projected coordinate
system, this theorem gives projected coordinate energy only after a polynomial
frame lower bound for that coordinate system.  If the downstream architecture is
exact-point locked or forced to use only the collapsed scalar \(I\)-coordinate,
this theorem is not sufficient and the \(q_0\)-adic row-pencil, source-robust
wall-row fallback, scalar incidence avoidance, or a higher/ambient row
replacement must be used.
\end{theorem}

\section{Current proof state after the local package audits}
\label{sec:current-proof-state-after-local-audits}

\begin{remark}[Audit convention for the current proof-state section]
\label{rem:audit-convention-current-proof-state}
The statements in this section distinguish three levels:
\begin{enumerate}
\item local algebraic lemmas, such as Gate A/B cancellations under fixed
source/core and grading conventions;
\item conditional local geometric infrastructure, such as the reduced-base
rank-two connection and affine forcing equation;
\item missing global proof gates, such as affine transport, boundary,
holonomy/index, or finite-core non-concentration.
\end{enumerate}
A statement at level (1) or (2) should not be read as a proof of RH.  The
current missing theorem is global.
\end{remark}

\begin{remark}[Relation to the final named-input ledger]
\label{rem:current-ledger-vs-final-named-inputs}
This section preserves the local Bottleneck--C and finite-core audit history.
It should not be read as the final dependency graph of the present manuscript.
The final dependency graph is summarized in
Section~\ref{sec:upstream-gate-one-arithmetic-inputs} and
Remark~\ref{rem:current-architecture-ledger-unified}.  In particular, the
ordinary upstream branch is now packaged into O1--O3, while the local
downstream branch is packaged into Gate~2, gauge-locked wall escape,
carrier-quotient D4, and row mobility.
\end{remark}

\begin{remark}[Updated downstream status after the Gate--2 and carrier-quotient audits]
\label{rem:updated-downstream-status-gate2-carrier-quotient}
The downstream local package has changed relative to earlier versions of this
section.

First, the old \(D1/D2\) propagation route is no longer the active Gate--2
route.  In the ordinary real-canonical branch, once Gate~1 is available,
Gate~2 is supplied by Theorem~\ref{thm:gate-two-from-gate-one-l2-rigidity}:
\[
\text{Gate 1}
+
|u'|_\infty\text{ source thickening}
+
L^2\text{-input real rigidity}
\Longrightarrow
\text{Gate 2}.
\]

Second, the old source-robust \(q_0\)-wall row is no longer the active D3
route under row-energy architecture.  It is replaced by the gauge-locked wall
escape theorem, Theorem~\ref{thm:gauge-locked-wall-escape-row-energy}, plus D4.

Third, the scalar order-nine \(I\)-coordinate wall is not the active D4
obstruction if row energy is measured in the carrier quotient.  The order-seven
and order-nine clean carriers survive independently in
\(\mathcal Q_{\rm car}\) by
Theorem~\ref{thm:d4-carrier-quotient-through-order-nine}.

Thus the current downstream implication is
\[
\boxed{
\text{Gate 1}
+
\text{ordinary real-canonical local complexity}
+
\text{carrier-quotient D4 residual control}
\Longrightarrow
\text{row energy}.
}
\]
This remains a downstream theorem only.  It does not discharge the upstream
ordinary/singular actual-zeta inputs needed to obtain Gate~1 globally.
\end{remark}

The remaining subsections of this section record the broader proof state after
the Gate~B, Bottleneck~C, finite-core, and interaction audits.  They are a
status ledger for older and parallel proof obligations and should be read
separately from the ordinary real-canonical downstream closure package above.
In particular, references below to local Bottleneck--C, finite-core sampled
vector non-concentration, or \(dd\)-interaction rigidity are not meant to undo
the updated downstream organization:
\[
\text{Gate 1}
+
|u'|_\infty\text{ source thickening}
+
L^2\text{-input real rigidity}
\Rightarrow
\text{Gate 2},
\]
followed by
\[
\text{Gate 2}
+
\text{gauge-locked wall escape}
+
\text{carrier-quotient D4}
+
\text{row mobility}
\Rightarrow
\text{row energy}.
\]

The closed branches are:
\[
\boxed{
M_A,M_C\text{ connection closure is complete under source-grade reflection,}
}
\]
\[
\boxed{
\text{mixed four-point dependence factors through }
(c,u_5,v_5,\Delta_7),
}
\]
and
\[
\boxed{
\text{finite odd-germ determinacy holds on compact nonsingular
actual-source strata.}
}
\]

The live branches are:
\[
\boxed{
\text{local Bottleneck~C physical edge obstruction,}
}
\]
\[
\boxed{
\text{finite-core sampled vector non-concentration,}
}
\]
and
\[
\boxed{
dd\text{-interaction feature-balance / spectral rigidity.}
}
\]

The local Bottleneck-C branch has now exhausted the scalar, quotient, split,
and line-bundle repair routes.  The remaining viable language is the full
rank-two fixed-sector bundle
\[
\mathcal F_b=\mathfrak f_h(b),
\]
with \(E_{\eta\eta}\) treated as vector-valued affine forcing.  Gate A and
Gate B remain as local algebraic lemmas, but the proof now requires a global
rank-two affine transport, vector-valued non-concentration, projective-plus-scale
index, or comparable non-scalar theorem.

\subsection{Closed branch: \(M_A,M_C\) connection closure}
\label{subsec:closed-ma-mc-connection-closure}

Let
\[
M_A=\dot r_3r_2\dot r_4,
\qquad
M_C=\dot r_3\dot r_2r_4.
\]
Under the homogeneous scalar-grade reflection convention,
\[
\mathsf P_g\circ\mathsf S=(-1)^g\mathsf S\circ\mathsf P_g,
\]
where \(\mathsf S\) is endpoint swap, the corrected mixed block satisfies
\[
N^{\corr}(h_-,h_+)^T=N^{\corr}(h_+,h_-).
\]
Consequently, for the source-linear middle block \(Y_g\),
\[
Y_g(h_-,h_+)^T=(-1)^gY_g(h_+,h_-).
\]
Writing
\[
Y_g(h-\delta/2,h+\delta/2)
=
\sum_{e\ge0}\delta^eY_g^{\langle e\rangle}(h),
\]
one obtains the edge-refined parity
\[
\left(Y_{g,\beta}^{\langle e\rangle}\right)^T
=
(-1)^{g+e}Y_{g,\beta}^{\langle e\rangle}.
\]

This parity closes the labelled \(J=(3,2,4)\) pre-trivialization table:
\[
Z_j[M_A]=Z_j[M_C]=0,
\qquad
0\le j\le7.
\]
Therefore the connection identities
\[
R_7^{(A)}
+
\sum_{i=1}^{7}\omega_m(Z_{7-i}[M_A],A_i)=0,
\]
and
\[
R_7^{(C)}
+
\sum_{i=1}^{7}\omega_m(Z_{7-i}[M_C],A_i)=0
\]
hold under the stated source-grade reflection convention.

Thus the \(M_A,M_C\) connection branch is no longer a live obstruction.

\subsection{Live local obstruction: the frozen promoted mixed-block coefficient}
\label{subsec:live-local-obstruction-eay-bc}

The first live local actual-source/Bottleneck~C carrier is the promoted
quadratic mixed-block channel
\[
Y_{12}^{(2)}:=D^2N_0[\xi_1,\xi_2].
\]
It is coupled to one same-point linear whitening slot.  The promoted channel
has amplitude factor
\[
a_1a_2(a_1+a_2).
\]
On the collision chart
\[
a_2=-a_1(1+\kappa\delta),
\]
one has
\[
a_1a_2(a_1+a_2)=a_1^3\kappa\delta+O(\delta^2).
\]
Thus the amplitude already supplies the first \(\delta\)-edge.  Therefore the
first Bottleneck~C coefficient is the frozen internal coefficient
\[
E_{A,Y}^{\mathrm{BC}}
:=
\tau_{h\to m}\pi_{\mathfrak f_h}
\operatorname{Sym}
[z^7]\,
D(G_-^{-1/2})[A_-]\,
Y_{12}^{(2)}\,
R_0.
\]
Equivalently,
\[
E_{A,Y}^{\mathrm{BC}}
=
\tau_{h\to m}\pi_{\mathfrak f_h}
\operatorname{Sym}
\sum_{\chi\in\{dd,d\eta+\eta d,\eta\eta\}}
\sum_{g=4}^{7}
\sum_{r+b=7-g}
H^A_{r,0}Y^\chi_{g,0}W_{+,b}.
\]
Write
\[
E_{A,Y}^{\mathrm{BC}}=e_{A,Y}I+f_{A,Y}S.
\]
The quotient Bottleneck~C obstruction is
\[
Q_{A,Y}:=e_{A,Y}v_5-u_5f_{A,Y}.
\]
The transform-safe obstruction is
\[
\Phi_K(E_{A,Y}^{\mathrm{BC}})=2f_{A,Y}.
\]

For comparison, the literal source-position edge
\[
E_{A,Y}^{\mathrm{pos}}
:=
\tau_{h\to m}\pi_{\mathfrak f_h}
\operatorname{Sym}
[z^7\delta^1]\,
D(G_-^{-1/2})[A_-]\,
Y_{12}^{(2)}\,
R_0
\]
belongs to the next source-motion layer.  It is not the first
Bottleneck~C amplitude-derivative coefficient.

\begin{remark}[Executable table for \(E_{A,Y}^{\mathrm{BC}}\) (current section)]
\label{rem:executable-table-eay-bc-current}
The frozen Bottleneck~C coefficient is computed from the rows
\[
\mathcal E^a_{\chi;g,r,b}
:=
\tau_{h\to m}\pi_{\mathfrak f_h}
\operatorname{Sym}
\left(
\Lambda^-_{a,r}Y^\chi_{g,0}W_{+,b}
\right),
\]
where
\[
a\in\mathcal A_{\mathrm{BC}},
\qquad
\chi\in\{dd,d\eta+\eta d,\eta\eta\},
\]
\[
4\le g\le7,
\qquad
r+b=7-g.
\]
Thus
\[
E_{A,Y}^{\mathrm{BC}}
=
\sum_{\chi\in\{dd,d\eta+\eta d,\eta\eta\}}
\sum_{a\in\mathcal A_{\mathrm{BC}}}
\sum_{\substack{4\le g\le7\\ r+b=7-g}}
\mathcal E^a_{\chi;g,r,b}.
\]
The skew combination
\[
Y_{d\eta-\eta d}^{(2)}
\]
is excluded because it is killed by fixed-sector projection.
\end{remark}

\begin{remark}[Status of the frozen promoted mixed-block coefficient (current section)]
\label{rem:status-eay-bc-current}
The finite Fr\'echet--Sylvester recurrences reduce
\(E_{A,Y}^{\mathrm{BC}}\) to the finite symbolic table of
Remark~\ref{rem:executable-table-eay-bc-current}.  However, the current draft
does not supply the scalar midpoint source-jet data needed to evaluate
\[
E_{A,Y}^{\mathrm{BC}}=e_{A,Y}I+f_{A,Y}S.
\]
The required scalar data include, safely,
\[
q_0^{(k)}(m),
\qquad
0\le k\le9,
\]
and
\[
(r_i^{[a]})^{(k)}(m),
\qquad
i=1,2,\quad a\in\mathcal A_{\mathrm{BC}},\quad 2\le k\le9.
\]
Equivalently, the paper must supply the midpoint Taylor coefficients of
\[
d_{i,\pm}(z),
\qquad
\eta_i(z),
\]
through the orders contributing to
\[
Y^{dd}_{g,0},
\qquad
Y^{d\eta+\eta d}_{g,0},
\qquad
Y^{\eta\eta}_{g,0},
\qquad
4\le g\le7.
\]

At present no available symmetry, source-convention rebuild, quotient
interpretation, transform-faithful invariant, or endpoint-symmetrized
finite-\(s\) carrier forces
\[
E_{A,Y}^{\mathrm{BC}}=0,
\qquad
e_{A,Y}v_5-u_5f_{A,Y}=0,
\qquad
f_{A,Y}=0.
\]
More specifically, the odd-\(\eta\eta\) top row
\[
\mathcal E^{2}_{\eta\eta;7,0,0}
\]
has the \(S\)-visible coefficient recorded in
Subsection~\ref{subsec:concrete-odd-etaeta-top-row-obstruction}.  Therefore
the first unresolved local actual-source/Bottleneck-C obstruction is now the
top-row cancellation problem
\[
\nu_{21,7}=0
\]
or an equivalent physical cancellation of
\[
\mathcal E^{2}_{\eta\eta;7,0,0}.
\]
\end{remark}

\begin{lemma}[Scalar tests for the frozen promoted mixed-block coefficient]
\label{lem:scalar-tests-eay-bc}
Let
\[
E_{A,Y}^{\mathrm{BC}}=eI+fS,
\qquad
A_5^{\mathfrak f}=u_5I+v_5S.
\]
Then quotient Bottleneck~C for the promoted mixed-block carrier is equivalent
to
\[
ev_5-u_5f=0.
\]
The transform-safe \(T\)-horizontal condition is equivalent to
\[
f=0.
\]
\end{lemma}

\begin{proof}
The quotient line is
\[
\mathbf C A_5^{\mathfrak f}.
\]
Thus
\[
E_{A,Y}^{\mathrm{BC}}\in \mathbf C A_5^{\mathfrak f}
\]
is equivalent, on a good patch, to the vanishing of the determinant
\[
ev_5-u_5f.
\]
On the other hand,
\[
\Phi_K(I)=0,
\qquad
\Phi_K(S)=2.
\]
Hence
\[
\Phi_K(E_{A,Y}^{\mathrm{BC}})=2f,
\]
so the transform-safe condition is exactly \(f=0\).
\end{proof}

\begin{corollary}[First live \(T\)-slope carrier (current section)]
\label{cor:first-live-t-slope-carrier-current}
The amplitude factor of the promoted carrier is
\[
a_1a_2(a_1+a_2),
\]
and
\[
\partial_{a_2}\bigl(a_1a_2(a_1+a_2)\bigr)\big|_{a_2=-a_1}
=
-a_1^2.
\]
Therefore
\[
\partial_\kappa T\big|_{E_{A,Y}^{\mathrm{BC}}}
=
-a_1^2\frac{\Phi_K(E_{A,Y}^{\mathrm{BC}})}{2c}
=
-a_1^2\frac{f}{c}.
\]
If
\[
E_{A,Y}^{\mathrm{BC}}=\lambda A_5^{\mathfrak f},
\]
then
\[
\partial_\kappa T\big|_{E_{A,Y}^{\mathrm{BC}}}
=
-a_1^2\lambda\frac{v_5}{c}.
\]
Thus, on the generic patch \(v_5\ne0\), a nonzero line-gauge
\(E_{A,Y}^{\mathrm{BC}}\) closes the quotient condition but still leaks
through \(T\).
\end{corollary}

\subsection{Concrete odd-\(\eta\eta\) top-row obstruction}
\label{subsec:concrete-odd-etaeta-top-row-obstruction}

The local package audits isolate a concrete top-row survivor inside the frozen
Bottleneck~C coefficient.  The relevant row is
\[
\mathcal E^{2}_{\eta\eta;7,0,0}
=
\tau_{h\to m}\pi_{\mathfrak f_h}
\operatorname{Sym}
\left(
\Lambda^-_{2,0}
\left[
\left(
\frac{\sin\Delta_0}{s^2}
+
\frac{q_-\cos\Delta_0}{s}
\right)
J\,\eta_i\eta_j
\right]_{7}
W_{+,0}
\right),
\]
where
\[
J=
\begin{pmatrix}
0&-1\\
1&0
\end{pmatrix}.
\]
This is the actual finite Fr\'echet source index \(a=2\), not the older
shorthand \(a\in\{1,5\}\).

The row comes from the \((2,1)\)-off-diagonal entry
\[
N_{21}
=
-\frac{\sin\Delta+q_-s\cos\Delta}{s^2}.
\]
Twice differentiating in the scalar phase-gap variable \(\Delta\) gives
\[
\partial_\Delta^2N_{21}\big|_0
=
\frac{\sin\Delta_0}{s^2}
+
\frac{q_-\cos\Delta_0}{s}.
\]
Thus the odd-background \(\eta\eta\) block is not an endpoint-separation
artifact.  It is a scalar \(\Delta\)-Hessian term.

For the active finite Fr\'echet row one has
\[
\Lambda^-_{2,0}
=
-\frac{1}{24\pi q^3}E_{22}.
\]
If
\[
W_{+,0}
=
\begin{pmatrix}
u&w\\
w&v
\end{pmatrix},
\]
then
\[
\pi_{\mathfrak f}\operatorname{Sym}
\left(
\Lambda^-_{2,0}JW_{+,0}
\right)
=
-\frac{1}{24\pi q^3}(wI+uS).
\]
Consequently, if
\[
\nu_{21,7}
:=
[z^7]\left[
\left(
\frac{\sin\Delta_0}{s^2}
+
\frac{q_-\cos\Delta_0}{s}
\right)\eta_i\eta_j
\right],
\]
then this row contributes
\[
\mathcal E^{2}_{\eta\eta;7,0,0}
=
-\frac{\nu_{21,7}}{24\pi q^3}(wI+uS),
\]
up to the harmless good-patch transport \(\tau_{h\to m}\).  In particular,
\[
f_{\eta\eta}^{\top}
=
-\frac{u}{24\pi q^3}\nu_{21,7}.
\]

Under the favorable phase-gap integral model,
\[
\nu_{21,7}
=
\frac{1}{Q^{14}}
\left[
\frac{q_0}{480}
(\rho_{i,2}\rho_{j,4}+\rho_{i,4}\rho_{j,2})
+
\frac{1}{144}
\left(
\frac{q_0''}{6}
-
\frac{2q_0^3}{3}
\right)
\rho_{i,2}\rho_{j,2}
\right].
\]
For the homogeneous source \(r(u)=u^2\), with \(q_0''=0\), this gives
\[
\nu_{21,7}
=
-\frac{q_0^3}{216Q^{14}},
\]
which is nonzero on a good phase patch.

Therefore the row is generically \(T\)-visible:
\[
\partial_\kappa T
=
-a_1^2\frac{f_{\eta\eta}^{\top}}{c}
\ne0.
\]
It is also quotient-visible on a generic \(u_5\ne0\) patch; for example, if
\(w=0\), then
\[
Q_{\eta\eta}^{\top}
=
-u_5 f_{\eta\eta}^{\top}.
\]

\begin{remark}[Endpoint cleanup does not remove this row]
\label{rem:endpoint-cleanup-does-not-remove-row}
Endpoint-diagonal realization, if proved, removes only endpoint-separation
subchannels such as
\[
d\eta^\Delta+\eta^\Delta d,
\qquad
\eta^\Sigma\eta^\Delta.
\]
It does not remove the scalar \(\eta\eta\) term above, since that term comes
from \(\partial_\Delta^2N_0[\eta_i,\eta_j]\).  Thus the older route
\[
[\delta^0]\eta^\Delta=0
\quad\Longrightarrow\quad
E_{A,Y}^{\mathrm{BC}}=0
\]
is no longer sufficient.
\end{remark}

\begin{remark}[No existing package mechanism cancels the row]
\label{rem:no-existing-package-mechanism-cancels-row}
The present package mechanisms do not cancel
\(\mathcal E^{2}_{\eta\eta;7,0,0}\):
source-bidegree projection does not remove the promoted degree-three amplitude
pattern \(a_1a_2(a_1+a_2)\); amplitude normalization only rescales the row;
midpoint trivialization is invertible on a good patch; connection correction is
post-edge and cannot cancel a raw frozen coefficient; and quotienting is
insufficient because the augmented transform-safe package requires the
\(S\)-component itself to vanish.
\end{remark}

Thus the local branch is in the following binary state.  Either one proves a
new source-jet identity
\[
\nu_{21,7}=0,
\]
or one proves a full top-row \(S\)-cancellation
\[
\sum_{\mathrm{top}\ \eta\eta} f_{\eta\eta}^{\top}=0.
\]
The only remaining local rescue is a genuine physical aggregate
\(S\)-coefficient cancellation, not a formal package counterterm.  Without
such cancellation, the current corrected package does not close
Bottleneck~C or the \(T\)-horizontal condition at this row.

\subsection{No-go for a transform-faithful free \(\Theta_\pm\)-normal repair}
\label{subsec:no-go-theta-normal-repair}

The off-diagonal numerator variables
\[
\Theta_+=\sin\Delta+q_+s\cos\Delta,
\qquad
\Theta_-=\sin\Delta+q_-s\cos\Delta
\]
make the off-diagonal entries formally affine:
\[
N_{12}^{\Theta}=\frac{\Theta_+}{s^2},
\qquad
N_{21}^{\Theta}=-\frac{\Theta_-}{s^2}.
\]
As a formal normal package this cancels the original off-diagonal
\(\eta\eta\) Hessian row.  It does not, however, close the physical
Bottleneck~C obstruction.

The distinction is important.  As a formal normal coordinate, \(\Theta_\pm\)
makes the off-diagonal entries affine.  As a physical coordinate induced from
\((q_-,q_+,\Delta)\), however, \(\Theta_\pm\) is constrained and its
source-curvature must either remain in the raw package or be carried by the
hidden physical readout.  The following no-go concerns the latter,
transform-faithful interpretation.

\begin{definition}[\(\Theta\)-local physical edge residual]
\label{def:theta-local-physical-edge-residual}
Let
\[
R_{\off,7}^{\Theta}
\]
be the septic fixed-sector residual produced by the actual constrained source
curvature of the off-diagonal variables \(\Theta_\pm\).  Let
\[
R_{\diag,7}^{\Theta}
\]
be the shifted diagonal \(\eta\eta\) septic residual
\[
\pi_{\mathfrak f}
\operatorname{Sym}
\left(
\Lambda^-_{\bullet,1}Y^{\eta\eta,\diag}_6W_{+,0}
+
\Lambda^-_{\bullet,0}Y^{\eta\eta,\diag}_6W_{+,1}
\right).
\]
Define
\[
R_{\loc,7}^{\Theta}
:=
R_{\off,7}^{\Theta}+R_{\diag,7}^{\Theta}.
\]
Write
\[
R_{\loc,7}^{\Theta}
=
\alpha_{\loc}^{\Theta}I+\beta_{\loc}^{\Theta}S.
\]
The scalar \(\beta_{\loc}^{\Theta}\) is the \(\Phi_K\)-visible physical edge
residual left by the \(\Theta_\pm\)-normal package.
\end{definition}

\begin{remark}[Constrained \(\Theta_\pm\) variables reintroduce curvature]
\label{rem:constrained-theta-reintroduces-curvature}
The variables \(\Theta_\pm\) are not free affine source variables in the
physical finite-\(s\) block.  They are constrained by
\[
\Theta_\sigma=\sin\Delta+q_\sigma s\cos\Delta.
\]
For two source variations,
\[
D_{\source}^2\Theta_\sigma
=
A_\sigma\eta_{ij}
+sC\,\chi_{\sigma,ij}
-\Theta_{\sigma,0}\eta_i\eta_j
-sS(\chi_{\sigma,i}\eta_j+\chi_{\sigma,j}\eta_i).
\]
In the pure \(\eta\eta\) channel, this gives
\[
D_{\source}^2N_{21}^{\Theta,\actual}
=
\left(
\frac{\sin\Delta_0}{s^2}
+
\frac{q_-\cos\Delta_0}{s}
\right)\eta_i\eta_j,
\]
which is the old off-diagonal source-curvature obstruction.  Thus the
\(\Theta_\pm\)-normal form is not a physical coordinate relabeling that
removes the obstruction; it is a formal package subtraction unless this
curvature is carried back into the physical readout.
\end{remark}

\begin{proposition}[No transform-faithful septic normal-form repair]
\label{prop:no-transform-faithful-septic-normal-form-repair}
Let \(K_7\) be a local septic normal-form correction satisfying:
\begin{enumerate}
\item \(K_7\) is invisible below ordinary order \(7\);
\item \(K_7\) is source-quadratic, hence does not alter source-linear Gate~A,
Gate~B, or the \(M_A,M_C\) connection tables;
\item \(K_7\) cancels the local septic residual
\[
R_{\loc,7}^{\Theta}
=
\alpha_{\loc}^{\Theta}I+\beta_{\loc}^{\Theta}S;
\]
\item the physical \(\Phi_K\)-readout is unchanged, so no correction is added
back into the physical hidden coordinate.
\end{enumerate}
Then necessarily
\[
\beta_{\loc}^{\Theta}=0.
\]
Consequently, on any good patch where \(\beta_{\loc}^{\Theta}\ne0\), no such
transform-faithful local normal-form correction exists.
\end{proposition}

\begin{proof}
To cancel \(R_{\loc,7}^{\Theta}\), the correction \(K_7\) must contain the
fixed-sector component
\[
-\beta_{\loc}^{\Theta}S.
\]
But preserving the physical \(\Phi_K\)-readout without adding a correction to
the hidden coordinate requires
\[
[z^7]\Phi_K(K_7)=0.
\]
Since \(\Phi_K(S)=2\), this forces
\[
\beta_{\loc}^{\Theta}=0.
\]
\end{proof}

\begin{proposition}[Physical edge obstruction survives in a good-patch model]
\label{prop:theta-physical-edge-obstruction-good-patch}
There exists a homogeneous good-patch local coefficient model for which
\[
\left.
\partial_\kappa
\left(
\frac{\beta_{\loc}^{\Theta}}{c}
\right)
\right|_{\kappa=0}
\ne0.
\]
\end{proposition}

\begin{proof}
After extracting the frozen Bottleneck-C matrix coefficient, the remaining
amplitude derivative is the derivative of
\[
a_1a_2(a_1+a_2)
\]
in the transverse amplitude direction at the balanced point \(a_2=-a_1\).
Equivalently, using the local edge coordinate
\[
\kappa_{\amp}:=a_2+a_1,
\]
one has
\[
a_1a_2(a_1+a_2)
=
a_1(-a_1+\kappa_{\amp})\kappa_{\amp}
=
-a_1^2\kappa_{\amp}+O(\kappa_{\amp}^2),
\]
so
\[
\left.
\partial_{\kappa_{\amp}}
\bigl(a_1a_2(a_1+a_2)\bigr)
\right|_{\kappa_{\amp}=0}
=
-a_1^2.
\]
This is the same scalar edge factor obtained from the collision chart
\[
a_2=-a_1(1+\kappa\delta)
\]
after the first \(\delta\)-edge has been factored into the frozen
Bottleneck-C coefficient.

Since \(\beta_{\loc}^{\Theta}(0)=0\), differentiating \(1/c\) contributes
nothing at \(\kappa=0\).  Thus
\[
\left.
\partial_\kappa
\left(
\frac{\beta_{\loc}^{\Theta}}{c}
\right)
\right|_{\kappa=0}
=
-\frac{a_1^2}{c_0}B_{\loc}^{\Theta},
\]
where \(B_{\loc}^{\Theta}\) is the frozen physical edge coefficient.

Take a homogeneous source model
\[
r_i=r_j=u^2,
\qquad
q_0''=0,
\]
on a good patch with
\[
u\,q_0\,q\,Q\,c_0\ne0.
\]
The constrained off-diagonal \(\Theta_\pm\)-source curvature contributes a
nonzero fixed-sector \(S\)-coefficient.  In the top-row normalization this
has leading form
\[
B_{\off}^{\Theta,\hom}
=
\frac{u q_0^3}{5184\pi q^3Q^{14}}.
\]
The shifted diagonal term has the form
\[
B_{\diag}^{\Theta,\hom}
=
-\frac{\lambda}{2\pi Q^{12}}
\left(
\frac{q_0\alpha}{72}
+
\frac{q_0^3\beta}{864}
\right),
\qquad
\lambda=
\frac{\alpha\beta}{4\pi Q(p_1+p_2)}.
\]
Choosing
\[
\alpha=\varepsilon\alpha_0,\qquad
\beta=\varepsilon\beta_0,
\]
with \(\varepsilon\ne0\) sufficiently small gives
\[
B_{\diag}^{\Theta,\hom}=O(\varepsilon^3),
\]
while \(B_{\off}^{\Theta,\hom}\) remains nonzero.  Therefore
\[
B_{\loc}^{\Theta,\hom}
=
B_{\off}^{\Theta,\hom}
+
O(\varepsilon^3)
\ne0
\]
for sufficiently small \(\varepsilon\), and the claimed derivative is nonzero.
\end{proof}

\begin{remark}[Status of the \(\Theta_\pm\)-normal repair]
\label{rem:status-theta-normal-repair}
A formal septic package may subtract \(R_{\loc,7}^{\Theta}\) and thereby make
the formal package horizontal.  But the physical hidden coordinate is then
\[
T^{\phys}
=
T^\star+\frac{\beta_{\loc}^{\Theta}}{c}.
\]
Thus the transform-safe physical Bottleneck~C condition is
\[
\partial_\kappa
\left(
T^\star+\frac{\beta_{\loc}^{\Theta}}{c}
\right)=0.
\]
Proposition~\ref{prop:theta-physical-edge-obstruction-good-patch} shows that
the second term is not currently zero.  Therefore local Bottleneck~C cannot be
closed by \(\Theta_\pm\)-normal coordinates or septic normal-form counterterms
alone.  The remaining local route would require a new physical edge identity
\[
\partial_\kappa(\beta_{\loc}^{\Theta}/c)=0
\]
not present in the current package calculus.
\end{remark}

\begin{remark}[\(\Theta_\pm\) no-go versus the frozen Bottleneck-C coefficient]
\label{rem:theta-nogo-versus-frozen-bc-coefficient}
The no-go result for the \(\Theta_\pm\)-normal repair does not replace the
frozen coefficient
\[
E_{A,Y}^{\mathrm{BC}}.
\]
Rather, it says that normal-form coordinates or septic counterterms cannot
make the physical transform-safe Bottleneck-C obstruction disappear.  The
canonical local coefficient to compute remains
\[
E_{A,Y}^{\mathrm{BC}}
=
e_{A,Y}I+f_{A,Y}S,
\]
with quotient test
\[
e_{A,Y}v_5-u_5f_{A,Y}=0
\]
and transform-safe test
\[
f_{A,Y}=0.
\]
The \(\Theta_\pm\) audit shows that if one removes a visible septic residual
from a formal package, the same residual must be carried by the physical
hidden coordinate.  Thus it is an obstruction-isolation result, not a
replacement for the \(E_{A,Y}^{\mathrm{BC}}\) finite table.

The concrete top-row obstruction
\[
\mathcal E^{2}_{\eta\eta;7,0,0}
\]
of Subsection~\ref{subsec:concrete-odd-etaeta-top-row-obstruction} is the
current row-level witness of this phenomenon: removing it from a formal septic
package without changing the physical readout merely transfers the same
obstruction to the hidden \(T\)-coordinate.
\end{remark}

\subsection{Actual exterior-pair source jets and failure of pair-level locking}
\label{subsec:actual-exterior-pair-source-jets-no-locking}

The most direct possible rescue of the Bottleneck-C row would be a hidden
actual-source relation forcing the midpoint post-affine source jets
\[
\rho_2,\qquad \rho_4
\]
to satisfy the ratio required by the \(\eta\eta\) top-row coefficient.  This
does not occur at the one-pair level.

For an exterior functional-equation zero-pair write
\[
f_{\beta,\gamma}(t)
=
\frac{1-\beta}{(1-\beta)^2+(2t-\gamma)^2}
+
\frac{\beta}{\beta^2+(2t-\gamma)^2}.
\]
Put
\[
a:=1-\beta,\qquad b:=\beta,\qquad x:=2m-\gamma,\qquad u:=t-m.
\]
Then
\[
f_{\beta,\gamma}(m+u)
=
\frac{a}{a^2+(x+2u)^2}
+
\frac{b}{b^2+(x+2u)^2}.
\]
Affine subtraction removes only the constant and linear terms.  Thus the
post-affine coefficients \(\rho_2,\rho_4\) are the \(u^2,u^4\) coefficients of
the above expansion.

For one summand,
\[
[u^2]\frac{c}{c^2+(x+2u)^2}
=
\frac{4c(3x^2-c^2)}{(c^2+x^2)^3},
\]
and
\[
[u^4]\frac{c}{c^2+(x+2u)^2}
=
\frac{16c(c^4-10c^2x^2+5x^4)}{(c^2+x^2)^5}.
\]
Therefore
\[
\boxed{
\rho_2(\beta,x)
=
4a\frac{3x^2-a^2}{(a^2+x^2)^3}
+
4b\frac{3x^2-b^2}{(b^2+x^2)^3}
}
\]
and
\[
\boxed{
\rho_4(\beta,x)
=
16a\frac{a^4-10a^2x^2+5x^4}{(a^2+x^2)^5}
+
16b\frac{b^4-10b^2x^2+5x^4}{(b^2+x^2)^5}.
}
\]

\begin{proposition}[No one-pair jet-locking law]
\label{prop:no-one-pair-jet-locking-law}
For \(\rho_2(\beta,x)\ne0\), the ratio
\[
\frac{\rho_4(\beta,x)}{\rho_2(\beta,x)}
=
4
\frac{
a\dfrac{a^4-10a^2x^2+5x^4}{(a^2+x^2)^5}
+
b\dfrac{b^4-10b^2x^2+5x^4}{(b^2+x^2)^5}
}{
a\dfrac{3x^2-a^2}{(a^2+x^2)^3}
+
b\dfrac{3x^2-b^2}{(b^2+x^2)^3}
}
\]
is a nonconstant rational function of the exterior-pair parameters.  In
particular, the actual one-pair source class does not force
\[
\frac{\rho_4}{\rho_2}
=
\frac{10q_0^2}{9}
-
\frac{5q_0''}{18q_0}.
\]
\end{proposition}

\begin{proof}
The displayed formulas are obtained by Taylor expanding the two summands at
\(u=0\).  The ratio \(\rho_4/\rho_2\) depends on the one-pair parameters
\((\beta,x)\), while the Bottleneck-C locking ratio depends on the local
baseline phase jets \(q_0,q_0''\).  There is no source-internal identity
equating these quantities.
\end{proof}

Two useful checks are:
\[
\beta=\frac12,\quad x=0
\quad\Longrightarrow\quad
\rho_2=-64,\qquad \rho_4=1024,\qquad
\frac{\rho_4}{\rho_2}=-16,
\]
and, as \(|x|\to\infty\),
\[
\rho_2=\frac{12}{x^4}+O(x^{-6}),
\qquad
\rho_4=\frac{80}{x^6}+O(x^{-8}),
\]
so
\[
\boxed{
\frac{\rho_4}{\rho_2}
=
\frac{20}{3x^2}+O(x^{-4}).
}
\]
Thus far exterior source atoms approximate the formal bad germ \(r(u)=u^2\)
to finite jet order.  Any cancellation of the aggregate Bottleneck-C
coefficient must therefore be a genuine global zero-tail identity, not a
pair-level source rigidity law.

\begin{remark}[Full symmetry-packet closure does not cancel the source moments]
\label{rem:full-symmetry-packet-does-not-cancel-source-moments}
The paired atom \(f_{\beta,\gamma}\) already includes the
functional-equation partner at the same ordinate:
\[
\beta+i\gamma,
\qquad
1-\beta+i\gamma.
\]
The full symmetry packet also contains the conjugate ordinate pair
\[
\beta-i\gamma,
\qquad
1-\beta-i\gamma,
\]
whose paired contribution is
\[
f_{\beta,-\gamma}(t)
=
\frac{1-\beta}{(1-\beta)^2+(2t+\gamma)^2}
+
\frac{\beta}{\beta^2+(2t+\gamma)^2}.
\]
If
\[
x_+(m)=2m-\gamma,
\qquad
x_-(m)=2m+\gamma,
\]
then the corresponding post-affine coefficients are
\[
\rho_2(\beta,x_+),\quad \rho_4(\beta,x_+),
\qquad
\rho_2(\beta,x_-),\quad \rho_4(\beta,x_-).
\]
The formulas above are even in \(x\).  Thus at a symmetric midpoint
\(m=0\), the \(\gamma\) and \(-\gamma\) packets add:
\[
\rho_2(\beta,-\gamma)=\rho_2(\beta,\gamma),
\qquad
\rho_4(\beta,-\gamma)=\rho_4(\beta,\gamma).
\]
At the midpoint centered on the \(+\gamma\) packet, \(x_+=0\) while
\(x_-=2\gamma\); again the conjugate packet supplies an additional contribution
rather than a sign-reversed copy.  Therefore functional-equation and conjugation
closure do not structurally force
\[
A_2=0,
\qquad
A_4=0,
\qquad
\nu_{\eta\eta}=0.
\]
Any cancellation after summing the full packet is a proper analytic incidence
condition, not a symmetry identity.
\end{remark}

\subsection{Source-convention rebuild and persistence of the \(\eta\eta\) row}
\label{subsec:source-convention-rebuild-etaeta-row}

The preceding obstruction used the tagged omitted-zero tail as the source
object.  One possible repair is to rebuild the source convention by moving all
non-core zero atoms into the source.  The least artificial split is
\[
q
=
q_{\base}^{\new}
+
G^{\new},
\]
where
\[
G^{\new}
:=
q_{\mathrm{aux},R}+T_{\mathrm{far},R},
\qquad
q_{\base}^{\new}
:=
q_{\mathrm{core}}+B_{\mathrm{bg}}+E_{\mathrm{est},R}.
\]
Thus \(G^{\new}\) remains a zero-source object; it contains the auxiliary
non-core zero contribution and the far zero tail, but not the smooth background
or estimation terms.

The affine-subtracted source is
\[
r_G(t)
:=
G^{\new}(t)-G^{\new}(m)-G^{\new\,\prime}(m)(t-m).
\]
Equivalently, the effective baseline in the finite-\(s\) expansion is
\[
q_0^{\new}(t)
:=
q_{\base}^{\new}(t)
+
G^{\new}(m)+G^{\new\,\prime}(m)(t-m),
\]
so that
\[
q(t)=q_0^{\new}(t)+r_G(t),
\]
with
\[
q_0^{\new}(m)=q(m),
\qquad
(q_0^{\new})'(m)=q'(m),
\qquad
(q_0^{\new})''(m)=(q_{\base}^{\new})''(m).
\]
This affine absorption is necessary: otherwise the source would contain value
and linear jets and would reopen lower source-grade gates.

\begin{remark}[Complement-source curvature]
\label{rem:complement-source-curvature-notation}
For the rebuilt source split, define
\[
q''_{\source,R}(m):=(G^{\new})''(m).
\]
Since
\[
r_G(t)=G^{\new}(t)-G^{\new}(m)-G^{\new\,\prime}(m)(t-m),
\]
one has
\[
r_G''(m)=q''_{\source,R}(m).
\]
This is the curvature of the affine-subtracted non-core zero complement.  It is
not the full canonical actual-zeta curvature
\[
q''_{\can}(m)=q_\zeta''(m),
\]
which also includes the core, background, and all remaining completed-phase
contributions in the fixed actual-zeta normalization.
\end{remark}

\begin{remark}[Admissibility conditions for the rebuilt split]
\label{rem:rebuilt-source-split-admissibility-conditions}
The split
\[
G^{\new}=q_{\mathrm{aux},R}+T_{\mathrm{far},R}
\]
is conditionally compatible with the already-closed lower gates under the
following hypotheses:
\begin{enumerate}
\item \(G^{\new}\) is kept as a tagged zero-source sum through the finite-\(s\)
calculation, not collapsed prematurely into an unlabelled aggregate;
\item \(r_G(m)=r_G'(m)=0\) is imposed after the new split;
\item \(q_{\base}^{\new}\) is used consistently in \(q_\pm,\Delta_0,d_\pm,
e_\pm,g_\pm,\eta\);
\item the \(q_{\mathrm{aux},R}\)-atoms satisfy the same endpoint reflection and
source-grade conventions as the \(T_{\mathrm{far},R}\)-atoms;
\item the \(q_{\mathrm{aux},R}\)-atoms have compatible scalar-grade
normalization, or else the affected coefficient tables must be regraded;
\item no \(B_{\mathrm{bg}}\) or \(E_{\mathrm{est},R}\) term is treated as a
zero-source variable.
\end{enumerate}
Under these hypotheses the rebuild does not introduce value or linear source
jets before the Bottleneck-C \(\eta\eta\) row.  By contrast, the more radical
choices
\[
G=q-q_{\mathrm{core}}
\quad\text{or}\quad
G=B_{\mathrm{bg}}+T_{\mathrm{far},R}+E_{\mathrm{est},R}
\]
mix zero atoms with background or error terms.  Those choices introduce
background--source and error--source products and would require a fresh
Gate~A/Gate~B analysis.
\end{remark}

For a tagged exterior-zero source atom \(\alpha\) in \(G^{\new}\), write
\[
r_\alpha(t)
=
g_\alpha(t)-g_\alpha(m)-g_\alpha'(m)(t-m).
\]
Define
\[
d_{\alpha,\pm}^{\new}=r_\alpha(t_\pm),
\qquad
e_{\alpha,\pm}^{\new}=r_\alpha'(t_\pm),
\qquad
g_{\alpha,\pm}^{\new}=r_\alpha''(t_\pm),
\]
and let \(\eta_\alpha^{\new}\) be the corresponding phase-gap source variable.
For two tagged source atoms \(i,j\), the finite-\(s\) substitutions become
\[
q_\pm
=
q_{0,\pm}^{\new}
+
\tau_i a_i d_{i,\pm}^{\new}
+
\tau_j a_j d_{j,\pm}^{\new},
\]
\[
q'_\pm
=
(q_0^{\new})'_\pm
+
\tau_i a_i e_{i,\pm}^{\new}
+
\tau_j a_j e_{j,\pm}^{\new},
\]
\[
q''_\pm
=
(q_0^{\new})''_\pm
+
\tau_i a_i g_{i,\pm}^{\new}
+
\tau_j a_j g_{j,\pm}^{\new},
\]
and
\[
\Delta
=
\Delta_0^{\new}
+
\tau_i a_i\eta_i^{\new}
+
\tau_j a_j\eta_j^{\new}.
\]

The lower-left off-diagonal channel is still
\[
N_{21}
=
-\frac{\sin\Delta+q_-s\cos\Delta}{s^2}.
\]
Therefore its pure \(\eta_i\eta_j\) second phase-gap variation is
\[
\partial_\Delta^2N_{21}\big|_{\Delta=\Delta_0^{\new}}
=
\frac{\sin\Delta_0^{\new}}{s^2}
+
\frac{q_{0,-}^{\new}\cos\Delta_0^{\new}}{s}.
\]
Thus the rebuilt top-row scalar is
\[
\nu_{21,7}^{\new,ij}
:=
[z^7]\left[
\left(
\frac{\sin\Delta_0^{\new}}{s^2}
+
\frac{q_{0,-}^{\new}\cos\Delta_0^{\new}}{s}
\right)
\eta_i^{\new}\eta_j^{\new}
\right].
\]

Let \(\rho_{\alpha,2}^{\new}\) and \(\rho_{\alpha,4}^{\new}\) be the midpoint
source coefficients of \(r_\alpha\) in the same normalization as before.  Then
the rederived pair coefficient has the same form as the old obstruction:
\[
\boxed{
\nu_{21,7}^{\new,ij}
=
\frac{1}{Q^{14}}
\left[
\frac{q_0^{\new}}{480}
\left(
\rho_{i,2}^{\new}\rho_{j,4}^{\new}
+
\rho_{i,4}^{\new}\rho_{j,2}^{\new}
\right)
+
\frac{1}{144}
\left(
\frac{(q_0^{\new})''}{6}
-
\frac{2(q_0^{\new})^3}{3}
\right)
\rho_{i,2}^{\new}\rho_{j,2}^{\new}
\right].
}
\]
Hence the source-convention rebuild does not algebraically kill the top row.

Aggregating over the full non-core zero source, define
\[
A_2^{\new}
:=
\sum_{\alpha\in\mathcal I_{\mathrm{ext}}^{\new}}
w_\alpha\rho_{\alpha,2}^{\new}
=
\frac{(G^{\new})''(m)}{2},
\]
and
\[
A_4^{\new}
:=
\sum_{\alpha\in\mathcal I_{\mathrm{ext}}^{\new}}
w_\alpha\rho_{\alpha,4}^{\new}
=
\frac{(G^{\new})^{(4)}(m)}{24}.
\]
Then
\[
\boxed{
\nu_{21,7}^{\new,\mathrm{agg}}
=
\frac{1}{Q^{14}}
\left[
\frac{q_0^{\new}}{240}A_2^{\new}A_4^{\new}
+
\frac{1}{144}
\left(
\frac{(q_0^{\new})''}{6}
-
\frac{2(q_0^{\new})^3}{3}
\right)
(A_2^{\new})^2
\right].
}
\]

With
\[
W_{+,0}^{\new}
=
\begin{pmatrix}
u^{\new}&w^{\new}\\
w^{\new}&v^{\new}
\end{pmatrix},
\]
the corresponding fixed-sector row is
\[
\mathcal E_{\eta\eta;7,0,0}^{2,\new}
=
-\frac{\nu_{21,7}^{\new}}{24\pi(q_{\wh}^{\new})^3}
\left(
w^{\new}I+u^{\new}S
\right),
\]
so
\[
f_{\eta\eta}^{\new,\top}
=
-\frac{u^{\new}}{24\pi(q_{\wh}^{\new})^3}
\nu_{21,7}^{\new}.
\]
Consequently, on a good patch with \(u^{\new}q_{\wh}^{\new}\ne0\), physical
Bottleneck~C closure still requires
\[
\nu_{21,7}^{\new,\mathrm{agg}}=0.
\]

\begin{remark}[Outcome of the source-convention rebuild]
\label{rem:outcome-source-convention-rebuild}
The split
\[
G^{\new}=q_{\mathrm{aux},R}+T_{\mathrm{far},R}
\]
is cleaner than the raw far-tail source convention because all non-core zero
atoms are placed on the source side.  Nevertheless, the rederived
\(\eta\eta\) top row has the same obstruction form, with
\[
q_0,\ q_0'',\ A_2,\ A_4
\]
replaced by
\[
q_0^{\new},\ (q_0^{\new})'',\ A_2^{\new},\ A_4^{\new}.
\]
Thus the rebuild merely shifts the needed cancellation theorem to the full
non-core zero complement:
\[
\frac{q_0^{\new}}{240}A_2^{\new}A_4^{\new}
+
\frac{1}{144}
\left(
\frac{(q_0^{\new})''}{6}
-
\frac{2(q_0^{\new})^3}{3}
\right)
(A_2^{\new})^2
=
0.
\]
If \(A_2^{\new}\ne0\), this is equivalently
\[
\frac{A_4^{\new}}{A_2^{\new}}
=
\frac{10(q_0^{\new})^2}{9}
-
\frac{5(q_0^{\new})''}{18q_0^{\new}}.
\]
No such global zero-complement identity is presently proved.
\end{remark}

\begin{remark}[Source-tagging and renormalized-complement ODEs]
\label{rem:source-tagging-renormalized-complement-ode-no-go}
The quantities \(A_2^{\new}\) and \(A_4^{\new}\) in the rebuilt
Bottleneck-C coefficient are the second and fourth midpoint jets of the
tagged zero-source object
\[
G^{\new}=q_{\mathrm{aux},R}+T_{\mathrm{far},R},
\]
after affine subtraction.  They are not the corresponding jets of
\[
q-q_{\core},
\qquad
B_{\bg}+T_{\far,R}+E_{\est,R},
\qquad
q,
\]
or any other renormalized complement unless the entire finite-\(s\) row is
rebuilt.

Thus an ODE for a different complement does not affect the already-derived
row.  In the present source convention, the only relevant derivative-ratio
identity is
\[
\frac{(G^{\new})^{(4)}(m)}{(G^{\new})''(m)}
=
\frac{40(q_0^{\new})^2}{3}
-
\frac{10(q_0^{\new})''}{3q_0^{\new}}.
\]
No ODE, meromorphic identity, moving-cutoff identity, or standard
\(\xi'/\xi\) phase-calculus identity currently supplies this condition.

Moreover, near an exterior zero-source pole \(p\) in \(G^{\new}\),
\[
G^{\new}(t)\sim \frac{C}{t-p},
\qquad
(G^{\new})''(t)\sim \frac{2C}{(t-p)^3},
\qquad
(G^{\new})^{(4)}(t)\sim \frac{24C}{(t-p)^5},
\]
so
\[
\frac{(G^{\new})^{(4)}}{(G^{\new})''}
\sim
\frac{12}{(t-p)^2}.
\]
But \(q_{\base}^{\new}=q_{\core}+B_{\bg}+E_{\est,R}\) does not contain that
exterior pole, so the right-hand side
\[
\frac{40(q_{\base}^{\new})^2}{3}
-
\frac{10(q_{\base}^{\new})''}{3q_{\base}^{\new}}
\]
has no matching pole there.  Hence the required equation is not a meromorphic
identity for the source-compatible complement.
\end{remark}

\subsection{Invariant and quotient no-go for the \(S\)-visible defect}
\label{subsec:invariant-quotient-no-go-s-visible-defect}

The source rebuild still leaves a physical \(S\)-visible defect.  One might
therefore ask whether replacing the local invariant or quotienting the fixed
sector can make this defect irrelevant.  This subsection records the negative
answer.

Let
\[
E=eI+fS
\]
be a Bottleneck-C fixed-sector residual.  The current physical readout has
\[
\Phi_K(I)=0,\qquad \Phi_K(S)=2,
\]
so
\[
\Phi_K(E)=2f.
\]
The hidden-state slope is
\[
\partial_\kappa T=-a_1^2\frac{f}{c}.
\]

\begin{proposition}[Local invariant no-go]
\label{prop:local-invariant-no-go-bc}
Let \(\mathcal J\) be a smooth local invariant or readout defined near the
exceptional divisor.  Suppose \(\mathcal J\) is transform-faithful in the sense
that its first variation determines the physical \(T\)-direction.  Then
\[
D\mathcal J[S]\ne0.
\]
Consequently, whenever the Bottleneck-C defect has \(f\ne0\), the first
exceptional-divisor slope of \(\mathcal J\) contains a nonzero contribution
from the same \(S\)-component.  Thus replacing \(\Phi_K\) by
\(\mathcal J\) does not remove the \(\eta\eta\) obstruction.

Conversely, if
\[
D\mathcal J[S]=0,
\]
then \(\mathcal J\) may hide the \(\eta\eta\) row at first order, but it cannot
reconstruct or control the physical hidden-state slope
\(\partial_\kappa T\).  Such a readout is not a valid Bottleneck-C replacement
unless an independent theorem proves that \(T\) is pure gauge.
\end{proposition}

\begin{proof}
At the exceptional divisor,
\[
\partial_\kappa\mathcal J
=
D\mathcal J[E]
=
D\mathcal J[eI+fS].
\]
If \(D\mathcal J[S]\ne0\), then the \(fS\)-component contributes to
\(\partial_\kappa\mathcal J\).  After any normalization that removes scalar
\(I\)-data, the \(S\)-term remains visible.  If \(D\mathcal J[S]=0\), then
\(\mathcal J\) cannot distinguish \(eI\) from \(eI+\varepsilon S\) to first
order.  But this is precisely the distinction measured by
\[
\partial_\kappa T=-a_1^2 f/c.
\]
Thus \(D\mathcal J[S]=0\) loses the physical \(T\)-direction, while
\(D\mathcal J[S]\ne0\) inherits the obstruction.
\end{proof}

\begin{remark}[Examples of non-faithful and faithful invariants]
\label{rem:examples-local-invariant-no-go-bc}
Trace, determinant to first order at the scalar merger, unordered conjugacy
class, and symplectic area are blind to \(f\) at first order.  They hide the
bad row only by discarding the oriented \(S\)-direction.  Oriented eigenvalue
ratios, oriented projective cross-ratios, and oriented curvature or
Schwarzian-type readouts are sensitive to \(f\), and therefore inherit the
same obstruction.
\end{remark}

Next consider the quotient line generated by
\[
A_5^{\mathfrak f}=u_5I+v_5S.
\]
The quotient determinant coordinate is
\[
\Delta_7
=
u_7v_5-u_5v_7.
\]
A residual \(E=eI+fS\) changes this coordinate by
\[
\omega_h(E,A_5^{\mathfrak f})=ev_5-u_5f.
\]
Thus the quotient condition is
\[
ev_5-u_5f=0.
\]

\begin{proposition}[Quotient invisibility is not transform invisibility]
\label{prop:quotient-not-transform-invisibility-bc}
Assume \(v_5\ne0\).  The quotient condition
\[
ev_5-u_5f=0
\]
does not imply the transform-safe condition \(f=0\).  In particular,
\[
E=\chi A_5^{\mathfrak f}
\]
is invisible to the quotient determinant, but
\[
\Phi_K(E)=2\chi v_5
\]
and
\[
T\mapsto T+\chi\frac{v_5}{c}.
\]
Therefore the quotient
\[
\mathfrak f_h/\mathbf C A_5^{\mathfrak f}
\]
is a valid moduli object for mixed four-point scalar exits, but not for the
full Bottleneck-C physical transform state.
\end{proposition}

\begin{proof}
If \(E=\chi A_5^{\mathfrak f}\), then
\[
E=\chi u_5I+\chi v_5S,
\]
so
\[
\omega_h(E,A_5^{\mathfrak f})
=
(\chi u_5)v_5-u_5(\chi v_5)
=
0.
\]
However
\[
\Phi_K(E)=2\chi v_5.
\]
On a \(v_5\ne0\) patch this is nonzero for \(\chi\ne0\).  The hidden coordinate
changes by
\[
T=\frac{v_7}{c}
\quad\mapsto\quad
T+\chi\frac{v_5}{c}.
\]
Thus the quotient identifies states with different \(\Phi_K\)-visible
transform data.
\end{proof}

\begin{remark}[Layer separation]
\label{rem:quotient-layer-separation-bc}
The quotient determinant layer and the physical transform layer see different
functionals:
\[
\omega_h(E,A_5^{\mathfrak f})=ev_5-u_5f,
\qquad
\Phi_K(E)=2f.
\]
The kernel of the former is not contained in the kernel of the latter.  Hence
quotient-horizontal Bottleneck~C can close the mixed four-point scalar-exit
branch while still leaking in the physical \(T\)-coordinate.
\end{remark}

\subsection{Finite-\(s\) carrier and endpoint-symmetrization no-go}
\label{subsec:finite-s-carrier-endpoint-symmetrization-no-go}

The remaining possible repair is to alter the finite-\(s\) carrier before
projection.  This also fails inside the present transform-faithful framework.

Write
\[
N_{12}
=
\frac{n_+}{s^2},
\qquad
n_+
=
\sin\Delta+q_+s\cos\Delta,
\]
and
\[
N_{21}
=
-\frac{n_-}{s^2},
\qquad
n_-
=
\sin\Delta+q_-s\cos\Delta.
\]
For a pure \(\eta_i\eta_j\) source variation, the relevant variation is through
\(\Delta\).  Holding \(q_\pm\) fixed in this Hessian row,
\[
\partial_\Delta n_\pm
=
\cos\Delta-q_\pm s\sin\Delta,
\qquad
\partial_\Delta^2 n_\pm
=
-\sin\Delta-q_\pm s\cos\Delta
=
-n_\pm.
\]
Therefore
\[
\partial_\Delta^2N_{12}
=
-\frac{\sin\Delta}{s^2}
-\frac{q_+\cos\Delta}{s},
\]
whereas
\[
\partial_\Delta^2N_{21}
=
\frac{\sin\Delta}{s^2}
+
\frac{q_-\cos\Delta}{s}.
\]
These are exact finite-\(s\) identities, not artifacts of Taylor truncation.

For a general linear endpoint carrier
\[
C_{a,b}:=aN_{12}+bN_{21},
\]
one obtains
\[
\partial_\Delta^2 C_{a,b}
=
\frac{(b-a)\sin\Delta}{s^2}
+
\frac{(bq_- - aq_+)\cos\Delta}{s}.
\]
Thus identical cancellation requires
\[
b-a=0,
\qquad
bq_- - aq_+=0.
\]
For a nonzero carrier this forces \(q_-=q_+\), which is not a generic
finite-\(s\) identity.

The common endpoint combinations illustrate the obstruction:
\[
N_{12}+N_{21}
=
\frac{(q_+-q_-)\cos\Delta}{s},
\]
so
\[
\partial_\Delta^2(N_{12}+N_{21})
=
-\frac{(q_+-q_-)\cos\Delta}{s}.
\]
This cancels the \(\sin\Delta/s^2\) part but leaves the endpoint-gradient
term.  Similarly,
\[
q_+N_{21}+q_-N_{12}
=
\frac{(q_- - q_+)\sin\Delta}{s^2},
\]
so
\[
\partial_\Delta^2(q_+N_{21}+q_-N_{12})
=
\frac{(q_+ - q_-)\sin\Delta}{s^2}.
\]
This cancels the \(q_\pm\cos\Delta/s\) part but leaves the phase-gap term.
Endpoint balancing can remove one component of the Hessian, but not both.

At the matrix level, define
\[
B_\pm
:=
\frac{\sin\Delta}{s^2}
+
\frac{q_\pm\cos\Delta}{s}.
\]
The off-diagonal second-Hessian matrix is
\[
H_{\eta\eta}
=
\begin{pmatrix}
0&-B_+\\
B_-&0
\end{pmatrix}.
\]
Writing
\[
J=
\begin{pmatrix}
0&-1\\
1&0
\end{pmatrix},
\qquad
X=
\begin{pmatrix}
0&1\\
1&0
\end{pmatrix},
\]
one has
\[
H_{\eta\eta}=AJ+DX,
\]
where
\[
A=\frac{B_++B_-}{2}
=
\frac{\sin\Delta}{s^2}
+
\frac{(q_++q_-)\cos\Delta}{2s},
\]
and
\[
D=\frac{B_- - B_+}{2}
=
\frac{(q_- - q_+)\cos\Delta}{2s}.
\]
Neither \(A\) nor \(D\) is generically zero.  The skew part \(AJ\) is the
endpoint-average obstruction, while the symmetric part \(DX\) removes the
\(\sin\Delta/s^2\) term but retains endpoint-gradient data.

The active left whitening coefficient is
\[
\Lambda^-_{2,0}
=
-\frac{1}{24\pi q^3}E_{22}.
\]
Since
\[
E_{22}J=E_{21},
\qquad
E_{22}X=E_{21},
\]
both endpoint-skew and endpoint-symmetric Hessian pieces feed the same leading
row after left multiplication.  Thus
\[
\Lambda^-_{2,0}(AJ)W_{+,0}
=
-\frac{A}{24\pi q^3}E_{21}W_{+,0},
\]
and
\[
\Lambda^-_{2,0}(DX)W_{+,0}
=
-\frac{D}{24\pi q^3}E_{21}W_{+,0}.
\]
After fixed-sector symmetrization,
\[
\pi_{\mathfrak f}\operatorname{Sym}(E_{21}W_{+,0})
=
wI+uS,
\]
up to the same normalization convention used in
Subsection~\ref{subsec:concrete-odd-etaeta-top-row-obstruction}.  Hence any
nonzero endpoint-symmetrized coefficient remains \(S\)-visible on a patch with
\(u\ne0\).

\begin{proposition}[Endpoint-symmetrized finite-\(s\) carrier no-go]
\label{prop:endpoint-symmetrized-carrier-no-go}
Let a finite-\(s\) endpoint carrier be formed from \(N_{12},N_{21}\) by smooth
endpoint-symmetric linear combinations and endpoint or Gram normalizations
that do not explicitly depend on the source variation direction.  If the
carrier remains transform-faithful, then its pure \(\eta_i\eta_j\) second
\(\Delta\)-Hessian has a nonzero lower-row component.  After the active
left-whitening slot \(\Lambda^-_{2,0}\), this produces a fixed-sector term
\[
-\frac{\mu_{\eta\eta}}{24\pi q^3}(wI+uS)
\]
for some endpoint coefficient \(\mu_{\eta\eta}\).  In particular, unless
\(\mu_{\eta\eta}\) is forced to vanish by a genuine zeta identity, the
physical \(S\)-coefficient
\[
f_{\eta\eta}
=
-\frac{u}{24\pi q^3}\mu_{\eta\eta}
\]
survives.
\end{proposition}

\begin{remark}[Nonlinear carriers do not repair the obstruction]
\label{rem:nonlinear-carriers-do-not-repair-etaeta-row}
Ratio, logarithmic, determinant-normalized, and phase-coordinate carriers do
not supply a transform-faithful repair.  A ratio or logarithmic carrier
normalizes the Hessian but leaves a nonzero endpoint-curvature term.  A phase
coordinate such as \(\Delta\) is affine in the source and therefore has zero
second \(\Delta\)-Hessian, but reconstructing the physical finite-\(s\) kernel
requires applying
\[
\Delta\mapsto \sin\Delta+q_\pm s\cos\Delta,
\]
whose second derivative is exactly the missing Hessian.  Thus the bad row is
either visible in the carrier, reintroduced by reconstruction, or hidden only
by losing the physical \(S\)-visible transform readout.
\end{remark}

\subsection{Local Bottleneck-C ledger after repair attempts}
\label{subsec:local-bottleneck-c-ledger-after-repairs}

Combining
Subsections~\ref{subsec:source-convention-rebuild-etaeta-row},
\ref{subsec:invariant-quotient-no-go-s-visible-defect}, and
\ref{subsec:finite-s-carrier-endpoint-symmetrization-no-go}, the present local
ledger is as follows.

Changing the source convention from the raw far tail to the full non-core
zero source
\[
G^{\new}=q_{\mathrm{aux},R}+T_{\mathrm{far},R}
\]
is coherent under the stated tagging and grading hypotheses, but it reproduces
the same \(\eta\eta\) top-row obstruction form.  Replacing the scalar local
invariant hides the defect only if it discards the physical \(T\)-direction.
Quotienting by \(\mathbf C A_5^{\mathfrak f}\) is legitimate for mixed
four-point scalar exits but not for the full transform state.  Finally,
endpoint symmetrization and alternative finite-\(s\) carriers either preserve
the second-\(\Delta\) Hessian, reintroduce it during reconstruction, or lose
transform faithfulness.

Therefore, within the present transform-faithful finite-\(s\) off-diagonal
kernel architecture, Bottleneck~C closes only if the physical aggregate
\(S\)-visible coefficient of the \(\eta\eta\) top row vanishes.  In the rebuilt
full non-core zero-source convention this means
\[
\nu_{21,7}^{\new,\mathrm{agg}}=0,
\]
or equivalently
\[
\frac{q_0^{\new}}{240}A_2^{\new}A_4^{\new}
+
\frac{1}{144}
\left(
\frac{(q_0^{\new})''}{6}
-
\frac{2(q_0^{\new})^3}{3}
\right)
(A_2^{\new})^2
=
0.
\]
No currently proved identity supplies this vanishing.  The subsequent
subsections record the present interpretation of this obstruction: it is not a
natural boundary term or residue theorem, and it is best viewed as the
vertical defect of the \(T\)-horizontal Bottleneck-C connection.  Thus the
present local route remains open-negative: either one proves a new pointwise
global zero-source cancellation theorem, or else one pivots to a different
architecture, such as global curvature/holonomy control or a quotient-only
no-\(T\) program.

\subsection{Last-chance pointwise flatness in full phase variables}
\label{subsec:last-chance-pointwise-flatness-full-phase-variables}

The rebuilt source convention removes the unstable-cutoff objection by using
the full non-core zero complement
\[
G^{\new}=q_{\aux,R}+T_{\far,R}.
\]
A final possible escape is that the exact decomposition of the full phase
derivative into rebuilt baseline and non-core complement might itself force
the Bottleneck-C flatness identity.  This subsection records that it does not.

Write
\[
Q(t):=q(t)
\]
for the full phase derivative, and decompose
\[
Q(t)=b(t)+G(t),
\]
where
\[
b(t):=q_{\base}^{\new}(t)
=
q_{\core}(t)+B_{\bg}(t)+E_{\est,R}(t),
\]
and
\[
G(t):=G^{\new}(t)=q_{\aux,R}(t)+T_{\far,R}(t).
\]
At the midpoint \(m\), put
\[
b_0:=b(m),
\qquad
b_2:=b''(m),
\]
and
\[
A_2:=\frac{G''(m)}2,
\qquad
A_4:=\frac{G^{(4)}(m)}{24}.
\]
Then
\[
A_2=\frac{Q''(m)-b''(m)}2,
\qquad
A_4=\frac{Q^{(4)}(m)-b^{(4)}(m)}{24}.
\]

The rebuilt Bottleneck-C flatness condition is
\[
\frac{b_0}{240}A_2A_4
+
\frac{1}{144}
\left(
\frac{b_2}{6}
-
\frac{2b_0^3}{3}
\right)
A_2^2
=0.
\]
Substitution gives
\[
\frac{b_0}{11520}
\bigl(Q''-b''\bigr)
\bigl(Q^{(4)}-b^{(4)}\bigr)
+
\frac{1}{576}
\left(
\frac{b''}{6}
-
\frac{2b_0^3}{3}
\right)
\bigl(Q''-b''\bigr)^2
=0,
\]
with all derivatives evaluated at \(m\).  Equivalently, if
\[
D_2:=Q''-b''=G'',
\qquad
D_4:=Q^{(4)}-b^{(4)}=G^{(4)},
\]
then the condition is
\[
D_2
\left[
b_0D_4
+
20
\left(
\frac{b''}{6}
-
\frac{2b_0^3}{3}
\right)
D_2
\right]
=0.
\]

\begin{proposition}[The rebuilt flatness equation is an extra fourth-order constraint]
\label{prop:rebuilt-flatness-extra-fourth-order-constraint}
The identity
\[
Q=b+G
\]
does not imply
\[
\nu_{21,7}^{\new,\mathrm{agg}}=0.
\]
On a good patch with \(G''(m)\ne0\), pointwise flatness is equivalent to the
additional derivative-ratio constraint
\[
\frac{G^{(4)}(m)}{G''(m)}
=
\frac{40b(m)^2}{3}
-
\frac{10b''(m)}{3b(m)}.
\]
Equivalently,
\[
b^{(4)}(m)
=
Q^{(4)}(m)
-
\left(
\frac{40b(m)^2}{3}
-
\frac{10b''(m)}{3b(m)}
\right)
\bigl(Q''(m)-b''(m)\bigr).
\]
Thus the desired flatness is not a consequence of the tautological full-phase
decomposition.  It is an extra nonlinear fourth-order pointwise condition on
the selected baseline/source split.
\end{proposition}

\begin{proof}
The displayed factorization shows that, if \(D_2=G''(m)\ne0\), the vanishing
condition is precisely
\[
b_0D_4
+
20
\left(
\frac{b''}{6}
-
\frac{2b_0^3}{3}
\right)
D_2
=0.
\]
Dividing by \(b_0D_2\) gives
\[
\frac{G^{(4)}(m)}{G''(m)}
=
\frac{40b_0^2}{3}
-
\frac{10b''(m)}{3b_0}.
\]
The equivalent formula for \(b^{(4)}\) follows from
\[
G^{(4)}=Q^{(4)}-b^{(4)}.
\]
The decomposition \(Q=b+G\) supplies only
\[
G''=Q''-b'',
\qquad
G^{(4)}=Q^{(4)}-b^{(4)}.
\]
It supplies no relation fixing \(G^{(4)}/G''\) in terms of \(b,b''\).
\end{proof}

\begin{remark}[Why defining the baseline by flatness is circular]
\label{rem:baseline-by-flatness-is-circular}
One could formally choose \(b\) by requiring the fourth-order equation in
Proposition~\ref{prop:rebuilt-flatness-extra-fourth-order-constraint}.  Such a
choice would kill the Bottleneck-C row by definition.  It is not a proof
unless an independent zeta-side principle selects exactly this baseline, for
example through the functional equation, a canonical scattering law, a
zero-counting conservation principle, or a functorial fixed-core construction.
No such principle is currently present.
\end{remark}

\begin{remark}[Asymptotic scale is not an identity]
\label{rem:asymptotic-scale-not-identity}
If \(b\) is dominated by the smooth phase/background, one expects roughly
\[
b(t)\sim \frac12\log(t/2\pi).
\]
Then the required ratio has scale
\[
\frac{G^{(4)}}{G''}
\sim
\frac{40}{3}b^2
\asymp
\log^2 t.
\]
A non-core zero complement with pole distances of order \(1/\log t\) can have
the same dimensional scale.  This is not enough: the identity requires the
exact constant and sign
\[
\frac{40b^2}{3}-\frac{10b''}{3b}
\]
pointwise.  Local one-pair and finite-tail tests show that source geometry can
produce different fourth-to-second derivative ratios.  Thus a proof of
flatness would have to be a genuinely new global zero-complement rigidity
theorem.
\end{remark}

\subsection{Global obstruction-class interpretation of the Bottleneck-C defect}
\label{subsec:global-obstruction-class-bottleneck-c-defect}

The preceding local ledger shows that the present transform-faithful
finite-\(s\) Bottleneck-C architecture closes only if the rebuilt aggregate
coefficient
\[
\nu_{21,7}^{\new,\mathrm{agg}}
\]
vanishes pointwise.  The next question is whether this coefficient is a hidden
boundary term, residue identity, or global conservation-law density.  The
current evidence is negative: it is better understood as a bulk vertical
connection-defect density for the \(T\)-horizontal Bottleneck-C connection.

For this subsection write, for readability,
\[
G:=G^{\new}=q_{\mathrm{aux},R}+T_{\mathrm{far},R},
\qquad
A_2:=\frac{G''(m)}{2}=\frac12q''_{\source,R}(m),
\qquad
A_4:=\frac{G^{(4)}(m)}{24},
\]
and use the baseline-local notation
\[
q_b:=q_0^{\new},
\qquad
q_b^{(2)}:=(q_0^{\new})''.
\]
The notation \(q_b^{(2)}\) is deliberately not written as \(q''\), because
\(q''\) is reserved in the rank-two arithmetic sections for the full canonical
curvature \(q''_{\can}\).
The rebuilt obstruction is
\[
\nu_{21,7}^{\new,\mathrm{agg}}
=
\frac{1}{Q^{14}}
\left[
\frac{q_b}{240}A_2A_4
+
\frac{1}{144}
\left(
\frac{q_b^{(2)}}{6}
-
\frac{2q_b^3}{3}
\right)
A_2^2
\right].
\]

\begin{proposition}[No natural exactness for the rebuilt Bottleneck-C defect]
\label{prop:no-natural-exactness-rebuilt-bc-defect}
Assume the source set is fixed as \(m\) varies, so that
\[
A_4=\frac{1}{12}A_2''.
\]
Then
\[
2880Q^{14}\nu_{21,7}^{\new,\mathrm{agg}}
=
q_bA_2A_2''
+
20
\left(
\frac{q_b^{(2)}}{6}
-
\frac{2q_b^3}{3}
\right)
A_2^2.
\]
This density is not an exact first-order local derivative in the natural
midpoint variables \((A_2,A_2',q_b,q_b')\).  More precisely,
\[
2880Q^{14}\nu_{21,7}^{\new,\mathrm{agg}}
=
\left(
q_bA_2A_2'
-
\frac12 q_b'A_2^2
\right)'
-
q_b(A_2')^2
+
\left(
\frac{23}{6}q_b^{(2)}
-
\frac{40}{3}q_b^3
\right)
A_2^2.
\]
The final two terms are genuine bulk terms.  Hence
\[
\nu_{21,7}^{\new,\mathrm{agg}}
\]
is not a boundary term or conservation-law density in the ordinary midpoint
calculus.
\end{proposition}

\begin{proof}
Since \(A_2=G''/2\), fixed source tagging gives
\[
A_2''=\frac{G^{(4)}}{2},
\qquad
A_4=\frac{G^{(4)}}{24}=\frac{1}{12}A_2''.
\]
Substitution into the formula for
\(\nu_{21,7}^{\new,\mathrm{agg}}\) gives the first displayed identity.
For the integration-by-parts decomposition, compute
\[
(q_bA_2A_2')'
=
q_bA_2A_2''
+
q_b(A_2')^2
+
q_b'A_2A_2',
\]
and
\[
(q_b'A_2^2)'=q_b^{(2)}A_2^2+2q_b'A_2A_2'.
\]
Therefore
\[
q_bA_2A_2''
=
\left(
q_bA_2A_2'
-
\frac12q_b'A_2^2
\right)'
-
q_b(A_2')^2
+
\frac12q_b^{(2)}A_2^2.
\]
Adding
\[
20
\left(
\frac{q_b^{(2)}}{6}
-
\frac{2q_b^3}{3}
\right)A_2^2
\]
gives the claimed residual term
\[
-q_b(A_2')^2
+
\left(
\frac{23}{6}q_b^{(2)}
-
\frac{40}{3}q_b^3
\right)A_2^2.
\]
\end{proof}

\begin{remark}[Moving source windows]
\label{rem:moving-source-windows-exactness-correction}
If the source set, cutoff, or core window depends on \(m\), write it as
\(G_m(t)\).  Then
\[
A_2(m)=\frac12\partial_t^2G_m(t)\big|_{t=m},
\]
and differentiation gives
\[
A_2''(m)
=
\frac12
\left(
\partial_m^2\partial_t^2
+
2\partial_m\partial_t^3
+
\partial_t^4
\right)G_m(t)\big|_{t=m}.
\]
Consequently
\[
A_4
=
\frac{1}{12}A_2''
-
\frac{1}{24}
\left(
\partial_m^2\partial_t^2
+
2\partial_m\partial_t^3
\right)G_m(t)\big|_{t=m}.
\]
Thus moving cutoffs or moving core windows introduce extra correction terms.
They do not by themselves create a conservation law; they add another
identity requirement.
\end{remark}

\begin{remark}[The induced second-order equation]
\label{rem:induced-second-order-equation-bc-defect}
If \(A_2\ne0\), pointwise vanishing of
\(\nu_{21,7}^{\new,\mathrm{agg}}\) is equivalent, in the fixed-source-set
calculus, to
\[
q_bA_2''
+
20
\left(
\frac{q_b^{(2)}}{6}
-
\frac{2q_b^3}{3}
\right)A_2
=0,
\]
or
\[
A_2''
+
\left(
\frac{10}{3}\frac{q_b^{(2)}}{q_b}
-
\frac{40}{3}q_b^2
\right)A_2
=0.
\]
This should be read as the unproved pointwise equation demanded by the current
Bottleneck-C gate, not as an already available zeta-side ODE.
\end{remark}

\subsection{Meromorphic residue no-go for the rebuilt obstruction}
\label{subsec:meromorphic-residue-no-go-rebuilt-obstruction}

There is also no current residue-theoretic mechanism forcing
\[
\nu_{21,7}^{\new,\mathrm{agg}}(m)=0.
\]
The rebuilt source is a zero-pair pole field.  Locally write
\[
G(m)
=
\sum_{p\in P_{\mathrm{ext}}}\frac{C_p}{m-p},
\]
where \(P_{\mathrm{ext}}\) is the set of non-core zero-source poles and
\(C_p\) includes the residue and source weight.  Then
\[
A_2(m)=\sum_p\frac{C_p}{(m-p)^3},
\qquad
A_4(m)=\sum_p\frac{C_p}{(m-p)^5}.
\]
If
\[
K(m):=\frac{q''(m)}{6}-\frac{2q(m)^3}{3},
\]
then, up to the global factor \(Q^{-14}\),
\[
\nu(m)
:=
\frac{q(m)}{240}A_2A_4+\frac{K(m)}{144}A_2^2
\]
is the quadratic pole-interaction field
\[
\nu(m)
=
\frac{q(m)}{240}
\sum_{p,r}
\frac{C_pC_r}{(m-p)^3(m-r)^5}
+
\frac{K(m)}{144}
\sum_{p,r}
\frac{C_pC_r}{(m-p)^3(m-r)^3}.
\]

\begin{proposition}[The rebuilt obstruction is not residue-killed]
\label{prop:rebuilt-obstruction-not-residue-killed}
The meromorphic differential \(\nu(m)\,dm\) is not, in general, an exact
meromorphic derivative.  Its local residues at zero-source poles are
generically nonzero.  Hence the global residue theorem can at most give
cancellation of the sum of residues; it does not imply the pointwise condition
\[
\nu(m)=0
\]
required by Bottleneck~C.
\end{proposition}

\begin{proof}
Fix a pole \(p\) and write \(\tau=m-p\).  Decompose
\[
A_2=C_p\tau^{-3}+a(\tau),
\qquad
A_4=C_p\tau^{-5}+b(\tau),
\]
where \(a,b\) are holomorphic at \(\tau=0\).  Then
\[
A_2A_4
=
C_p^2\tau^{-8}
+
C_pb\,\tau^{-3}
+
C_pa\,\tau^{-5}
+
ab,
\]
and
\[
A_2^2
=
C_p^2\tau^{-6}
+
2C_pa\,\tau^{-3}
+
a^2.
\]
Thus \(\nu(m)\,dm\) has high-order poles, and its residue depends on Taylor
coefficients of \(q,K,a,b\) at \(p\).  There is no structural reason for this
residue to vanish.

Already in a two-pole constant-\((q,K)\) model, with poles \(p_1,p_2\),
residues \(C_1,C_2\), and \(d=p_1-p_2\), the analytic tail at \(p_1\) gives
\[
\operatorname{Res}_{p_1}(\nu\,dm)
=
C_1C_2
\left(
\frac{q}{8d^7}
+
\frac{K}{12d^5}
\right),
\]
which is generically nonzero.  Therefore \(\nu(m)\,dm\) is not an exact
meromorphic derivative.
\end{proof}

\begin{remark}[What residue calculus can and cannot give]
\label{rem:residue-calculus-limit-bc-obstruction}
Functional-equation pairing and conjugation may pair or cancel residues in a
global sum.  Such cancellation is compatible with the global residue theorem,
but it is much weaker than the pointwise identity
\[
\nu_{21,7}^{\new,\mathrm{agg}}(m)=0.
\]
The Bottleneck-C gate is pointwise because
\[
f_{\eta\eta}^{\new,\mathrm{agg}}
=
-\frac{u^{\new}}{24\pi(q_{\wh}^{\new})^3}
\nu_{21,7}^{\new,\mathrm{agg}},
\qquad
\partial_\kappa T=-a_1^2f/c.
\]
Thus an integrated or cyclic residue cancellation cannot close the present
\(T\)-horizontal gate.
\end{remark}

\subsection{Meromorphic ODE principal-part no-go}
\label{subsec:meromorphic-ode-principal-part-no-go}

A related failed repair is to seek a meromorphic differential equation for the
source complement or for the full phase derivative.  The local flatness
condition would be implied by an equation of the form
\[
\frac{G^{(4)}}{G''}
=
\frac{40q^2}{3}
-
\frac{10q''}{3q},
\]
or by the same equation with \(G=q_\zeta\), the full actual-zeta phase
derivative.  This is incompatible with the principal parts of the actual
zero-pair pole fields.

Let
\[
H(t)=\frac{C}{t-p}+O(1)
\]
have a simple pole.  Then
\[
H''(t)=\frac{2C}{(t-p)^3}+O(1),
\qquad
H^{(4)}(t)=\frac{24C}{(t-p)^5}+O(1),
\]
so
\[
\frac{H^{(4)}}{H''}
=
\frac{12}{(t-p)^2}+O\!\left(\frac1{t-p}\right).
\]

On the other hand, if the full phase derivative has
\[
q_\zeta(t)=\frac{R}{t-p}+O(1),
\]
then
\[
q_\zeta^2(t)=\frac{R^2}{(t-p)^2}+O\!\left(\frac1{t-p}\right),
\qquad
\frac{q''_{\can}}{q_\zeta}
=
\frac{2}{(t-p)^2}+O\!\left(\frac1{t-p}\right).
\]
Hence
\[
\frac{40q^2}{3}-\frac{10q''}{3q}
=
\left(
\frac{40R^2}{3}
-
\frac{20}{3}
\right)
\frac1{(t-p)^2}
+
O\!\left(\frac1{t-p}\right).
\]
Matching the left-hand double-pole coefficient \(12\) would require
\[
R^2=\frac75.
\]

\begin{proposition}[Zero-pair residues rule out the meromorphic ODE repair]
\label{prop:zero-pair-residues-rule-out-meromorphic-ode}
The flatness-forcing ODE
\[
\frac{H^{(4)}}{H''}
=
\frac{40q^2}{3}
-
\frac{10q''}{3q}
\]
cannot hold as a meromorphic identity for the raw non-core zero source, for the
full phase derivative, or for any complement retaining the zero-pair poles.
\end{proposition}

\begin{proof}
For a zero-pair summand
\[
f_a(t)=\frac{a}{a^2+(2t-\gamma)^2},
\]
the poles are
\[
p_a^\pm=\frac{\gamma\pm ia}{2},
\]
with residues
\[
\operatorname{Res}_{p_a^+} f_a=-\frac{i}{4},
\qquad
\operatorname{Res}_{p_a^-} f_a=\frac{i}{4}.
\]
Thus the phase-pole residues are purely imaginary.  For a simple pole,
\[
R=\pm \frac{i}{4},
\]
so
\[
R^2=-\frac1{16},
\]
not \(7/5\).  If several same-sign pole contributions coincide, then
\[
R=\pm \frac{ni}{4},
\]
and still \(R^2<0\).  Therefore the right-hand double-pole coefficient cannot
match the coefficient \(12\) of \(H^{(4)}/H''\).
\end{proof}

\begin{remark}[Pole-cancelled complements are the wrong object]
\label{rem:pole-cancelled-complements-wrong-object}
A pole-cancelled complement can evade the principal-part contradiction at the
cancelled pole.  But Bottleneck-C uses the source moments
\[
A_2=\frac{G''(m)}2,
\qquad
A_4=\frac{G^{(4)}(m)}{24}.
\]
An ODE for a pole-cancelled complement is therefore an ODE for the wrong
object, unless the entire order-seven row is rederived with different moments.
\end{remark}

\subsection{Connection-defect interpretation and possible architecture pivot}
\label{subsec:connection-defect-interpretation-bc}

The preceding no-go results suggest a different interpretation.  The
coefficient
\[
\nu_{21,7}^{\new,\mathrm{agg}}
\]
is not merely a bad scalar coefficient; it is the vertical defect of the
current \(T\)-horizontal Bottleneck-C connection.

Recall that
\[
A_7^{\mathfrak f}=u_7I+v_7S,
\qquad
T=\frac{v_7}{c}.
\]
For a residual
\[
E=eI+fS,
\]
the physical readout satisfies
\[
\Phi_K(E)=2f,
\]
and the Bottleneck-C exceptional slope is
\[
\partial_\kappa T=-a_1^2\frac{f}{c}.
\]
For the rebuilt non-core zero-source convention, the active \(\eta\eta\) row
gives
\[
f_{\eta\eta}^{\new,\mathrm{agg}}
=
-\frac{u^{\new}}{24\pi(q_{\wh}^{\new})^3}
\nu_{21,7}^{\new,\mathrm{agg}}.
\]
Hence, on a good patch,
\[
\partial_\kappa T\big|_{\eta\eta}
=
a_1^2
\frac{u^{\new}}{24\pi c(q_{\wh}^{\new})^3}
\nu_{21,7}^{\new,\mathrm{agg}}.
\]

\begin{definition}[Bottleneck-C vertical defect density]
\label{def:bottleneck-c-vertical-defect-density}
Define
\[
\mathcal K_{\eta\eta;\kappa}
:=
\partial_\kappa T\big|_{\eta\eta}.
\]
Equivalently,
\[
\mathcal K_{\eta\eta;\kappa}
=
C_{\rm geom}\,
\nu_{21,7}^{\new,\mathrm{agg}},
\qquad
C_{\rm geom}
=
a_1^2
\frac{u^{\new}}{24\pi c(q_{\wh}^{\new})^3}.
\]
On a good patch,
\[
C_{\rm geom}\ne0.
\]
\end{definition}

\begin{remark}[Curvature versus second fundamental form]
\label{rem:curvature-versus-second-fundamental-form-bc}
The active row is symmetric in the source directions \(\eta_i,\eta_j\).  Thus
\(\mathcal K_{\eta\eta;\kappa}\) is not ordinary antisymmetric gauge curvature
by itself.  It is better regarded as the vertical second fundamental form, or
quadratic curvature-density, of the diagonal-merger horizontal distribution in
the \(T\)-coordinate.
\end{remark}

\begin{proposition}[Local flatness equivalence]
\label{prop:local-flatness-equivalence-bc}
Within the current transform-faithful Bottleneck-C architecture, local
\(T\)-horizontal closure of the rebuilt \(\eta\eta\) row is equivalent to
\[
\nu_{21,7}^{\new,\mathrm{agg}}=0.
\]
\end{proposition}

\begin{proof}
By Definition~\ref{def:bottleneck-c-vertical-defect-density},
\[
\partial_\kappa T\big|_{\eta\eta}
=
C_{\rm geom}
\nu_{21,7}^{\new,\mathrm{agg}}.
\]
On a good patch \(C_{\rm geom}\ne0\).  Therefore
\[
\partial_\kappa T\big|_{\eta\eta}=0
\]
if and only if
\[
\nu_{21,7}^{\new,\mathrm{agg}}=0.
\]
\end{proof}

\begin{remark}[Possible curvature/holonomy pivot]
\label{rem:possible-curvature-holonomy-pivot}
The current Bottleneck-C architecture demands pointwise flatness:
\[
\mathcal K_{\eta\eta;\kappa}=0.
\]
A different architecture could attempt to treat
\[
C_{\rm geom}\nu_{21,7}^{\new,\mathrm{agg}}\,dm
\]
as a global curvature or holonomy density, replacing pointwise \(T\)-flatness
by a global statement such as
\[
\int \mathcal K_{\eta\eta;\kappa}\,dm=0,
\qquad
\operatorname{Hol}_T(\mathcal K)=0,
\qquad
\text{or}\qquad
\operatorname{Index}(\mathcal K)=0.
\]
This would be a genuine architecture pivot, not a repair of the present local
Bottleneck-C gate.  To make such a pivot meaningful, one would first need to
prove transformation laws showing that
\[
C_{\rm geom}\nu_{21,7}^{\new,\mathrm{agg}}\,dm
\]
is a globally defined connection or curvature density under changes of patch,
core window, source convention, and normalization.
\end{remark}

\subsection{No-\(T\) quotient-determinant architecture as a possible pivot}
\label{subsec:no-t-quotient-determinant-architecture}

A second possible pivot is to abandon the \(T\)-horizontal Bottleneck-C gate
entirely and formulate the local-to-global argument using only quotient
determinant data.  This is not a repair of Bottleneck~C; it is a different
proof architecture.

The fixed-sector data are
\[
A_5^{\mathfrak f}=u_5I+v_5S,
\qquad
A_7^{\mathfrak f}=u_7I+v_7S.
\]
The quotient determinant coordinate is
\[
\Delta_7=u_7v_5-u_5v_7.
\]
A shift
\[
A_7^{\mathfrak f}
\mapsto
A_7^{\mathfrak f}+\chi A_5^{\mathfrak f}
\]
preserves \(\Delta_7\), but changes
\[
T=\frac{v_7}{c}
\quad\mapsto\quad
T+\chi\frac{v_5}{c}.
\]
It also changes
\[
\Phi_K(A_7^{\mathfrak f})
\]
by \(2\chi v_5\).  Thus \(T\) is not quotient-invariant.

A quotient-only architecture would use as local moduli
\[
\mathcal Q(h)
=
\left(
c(h),u_5(h),v_5(h),\Delta_7(h)
\right),
\]
or equivalently the quotient class
\[
[A_7^{\mathfrak f}]
\in
\mathfrak f_h/\mathbf C A_5^{\mathfrak f}.
\]
Its local gate would be
\[
ev_5-u_5f=0,
\]
not
\[
f=0.
\]

\begin{remark}[What a no-\(T\) architecture would require]
\label{rem:no-t-architecture-requirements}
A no-\(T\) architecture can make the \(T\)-leak irrelevant only if every
downstream RH-relevant detection step is reformulated through quotient data:
\[
(c,u_5,v_5,\Delta_7),
\]
mixed four-point scalar exits, and finite-core determinant geometry.  It
cannot use
\[
T=\frac{v_7}{c},
\qquad
\Phi_K(A_7^{\mathfrak f}),
\]
or hidden-state reconstruction as independent data.
\end{remark}

\begin{proposition}[Conditional no-go for quotient-only architecture]
\label{prop:conditional-no-go-quotient-only-architecture}
If a downstream local-to-global step requires reconstruction of \(T\) or
\(\Phi_K(A_7^{\mathfrak f})\), then no quotient-only architecture can replace
the current Bottleneck-C gate.
\end{proposition}

\begin{proof}
The quotient identifies representatives
\[
A_7^{\mathfrak f}
\quad\text{and}\quad
A_7^{\mathfrak f}+\chi A_5^{\mathfrak f}.
\]
These representatives have the same \(\Delta_7\), but their \(T\)-coordinates
differ by
\[
\chi\frac{v_5}{c}.
\]
On a \(v_5\ne0\) patch this is nonzero for \(\chi\ne0\).  Therefore no
function of the quotient data can recover \(T\) or
\(\Phi_K(A_7^{\mathfrak f})\).  Any proof step depending on those quantities
cannot be performed in the quotient-only architecture.
\end{proof}

\begin{remark}[Relation to the Bottleneck-C obstruction]
\label{rem:no-t-relation-to-bc-obstruction}
The no-\(T\) quotient program would make the obstruction
\[
\nu_{21,7}^{\new,\mathrm{agg}}
\]
irrelevant only by removing \(T\)-horizontal closure from the proof.  It would
not prove
\[
\nu_{21,7}^{\new,\mathrm{agg}}=0.
\]
The burden would instead move to quotient finite-core non-concentration and
quotient interaction rigidity using only \((c,u_5,v_5,\Delta_7)\) and the
mixed four-point scalar exits.
\end{remark}

\subsection{Post-quotient audit: split and line-bundle repairs}
\label{subsec:post-quotient-split-line-bundle-audit}

The no-\(T\) quotient architecture keeps the determinant data
\[
D_{2r+1}
=
\omega_h(A_{2r+1}^{\mathfrak f},A_5^{\mathfrak f})
=
u_{2r+1}v_5-u_5v_{2r+1}.
\]
This removes the representative-level \(T\)-coordinate, but it also discards
the full \(A_5^{\mathfrak f}\)-line.  Thus the pure line anomaly
\[
A_{2r+1}^{\mathfrak f}
=
\lambda_{2r+1}A_5^{\mathfrak f}
\]
has
\[
D_{2r+1}=0,
\]
while still giving
\[
\Phi_K(A_{2r+1}^{\mathfrak f})
=
2\lambda_{2r+1}v_5,
\qquad
T_{2r+1}
=
\lambda_{2r+1}\frac{v_5}{c}.
\]
Consequently the quotient scalar
\[
H_m^{\quot}(s)
=
\omega_h\!\left(
\pi_{\mathfrak f}\widehat\Omega_\zeta^{\corr}(s;m),
A_5^{\mathfrak f}
\right)
\]
is analytically coherent but incomplete for the original transform-state
detection problem.

A natural repair is the split readout
\[
\mathcal H_m^{\mathrm{split}}(s)
=
\left(
H_m^{\quot}(s),H_m^\parallel(s)
\right),
\]
where
\[
H_m^\parallel(s)
=
\frac{
\left\langle
\pi_{\mathfrak f}\widehat\Omega_\zeta^{\corr}(s;m),
A_5^{\mathfrak f}
\right\rangle_{\wh}
}{
\left\langle
A_5^{\mathfrak f},A_5^{\mathfrak f}
\right\rangle_{\wh}
}.
\]
For each coefficient
\[
A_{2r+1}^{\mathfrak f}
=
u_{2r+1}I+v_{2r+1}S,
\]
this gives the pair
\[
\left(
D_{2r+1},\lambda_{2r+1}
\right),
\]
with
\[
\lambda_{2r+1}
=
\frac{
\langle A_{2r+1}^{\mathfrak f},A_5^{\mathfrak f}\rangle_{\wh}
}{
\langle A_5^{\mathfrak f},A_5^{\mathfrak f}\rangle_{\wh}
}.
\]
On a good patch, if the whitening metric is fixed and nondegenerate, the pair
\[
(D_{2r+1},\lambda_{2r+1})
\]
recovers the full fixed-sector vector \(A_{2r+1}^{\mathfrak f}\).  Thus the
split readout repairs the fatal incompleteness of the pure quotient readout.

However, this completeness reintroduces the Bottleneck-C obstruction through
the longitudinal coordinate.  For
\[
E_{\eta\eta}
=
-\frac{\nu_{\eta\eta}}{24\pi q^3}(wI+uS),
\]
one has
\[
D(E_{\eta\eta})
=
-\frac{\nu_{\eta\eta}}{24\pi q^3}(wv_5-u u_5),
\]
and
\[
\lambda(E_{\eta\eta})
=
-\frac{\nu_{\eta\eta}}{24\pi q^3}
\frac{
\langle wI+uS,A_5^{\mathfrak f}\rangle_{\wh}
}{
\langle A_5^{\mathfrak f},A_5^{\mathfrak f}\rangle_{\wh}
}.
\]
Thus if the split readout is used as a pointwise vanishing gate, then on a
good patch it requires
\[
E_{\eta\eta}=0,
\]
and hence again requires
\[
\nu_{\eta\eta}=0
\]
unless \(wI+uS\) degenerates.  The split readout is therefore a complete
coordinate system for the fixed-sector vector, not a new closure mechanism.

\begin{remark}[Line-bundle transport attempt]
\label{rem:line-bundle-transport-attempt}
One can try to avoid pointwise vanishing by treating the longitudinal
coordinate as a section of
\[
L=\mathbf C A_5^{\mathfrak f}.
\]
With a fixed whitening metric, the natural local connection is
\[
\alpha_L
=
\frac{
\langle dA_5^{\mathfrak f},A_5^{\mathfrak f}\rangle_{\wh}
}{
\langle A_5^{\mathfrak f},A_5^{\mathfrak f}\rangle_{\wh}
}.
\]
If
\[
s=\lambda A_5^{\mathfrak f},
\]
then
\[
\nabla^Ls=(d\lambda+\lambda\alpha_L)A_5^{\mathfrak f}.
\]
Locally this is coherent: under
\[
A_5^{\mathfrak f}\mapsto \mu A_5^{\mathfrak f}
\]
one has
\[
\alpha_L\mapsto\alpha_L+d\log\mu,
\]
and the covariant derivative transforms correctly.

The obstruction is global.  The construction requires the \(A_5^{\mathfrak f}\)
line and the whitening metric to transform canonically under changes of patch,
source/core window, and whitening convention.  This is not currently proved.
Moreover, no canonical path, boundary condition, or holonomy theorem is
available for \((L,\nabla^L)\), and the longitudinal forcing is not exact in
the natural variables.  Thus line-bundle transport is a useful diagnostic
language but not presently a proof gate.
\end{remark}

\begin{remark}[Current status of quotient, split, and line-bundle repairs]
\label{rem:current-status-quotient-split-line-repairs}
The quotient scalar is invariant but incomplete.  The split readout is complete
but restores the original fixed-sector obstruction.  The line-bundle
connection makes the longitudinal defect locally covariant, but lacks global
covariance, canonical paths, boundary data, exactness, or holonomy control.
Therefore none of these repairs currently replaces pointwise Bottleneck-C
closure.
\end{remark}

\subsection{Reduced core, packet visibility, and the visibility-positive base}
\label{subsec:reduced-core-packet-visibility-base}

The rank-two architecture is only meaningful after fixing the source/core
convention.  Arbitrary transfers of zero atoms between source and core are not
currently functorial.  Moving a single zero contribution \(q_\rho\) from source
to core changes the local corrected object, not merely its coordinates.  At
coefficient level one should expect
\[
A_{2r+1,j}^{\mathfrak f}
=
G_{ij}^{-1}A_{2r+1,i}^{\mathfrak f}
+
B_{\rho,2r+1},
\]
not a pure \(GL_2\)-transition.  Therefore source/core moves are excluded from
the atlas rather than transported.

We use the canonical reduced-core rule
\[
\mathcal C_k(m)
=
\text{the unique symmetry-closed, same-scale \(k\)-packet nearest }m.
\]
Let
\[
\mathcal E_{\rm geom}^{\core}
=
\mathcal E_{\rm geom}
\cup
\mathcal E_{\rm core\ boundary},
\]
where \(\mathcal E_{\rm core\ boundary}\) includes nearest-packet ties,
same-scale boundary, symmetry-closure ambiguity, endpoint-gap collapse, zero
collision or multiplicity, Gram/whitening singularity, and patch-boundary
failure.  Define
\[
\mathcal B_{\rm red}
=
\mathcal B\setminus\mathcal E_{\rm geom}^{\core}.
\]
On \(\mathcal B_{\rm red}\), the selected core is locally constant as a
labelled symmetry-closed packet.

\begin{hypothesis}[Reduced good-patch reach]
\label{hyp:reduced-good-patch-reach}
Every nondegenerate RH-relevant off-line symmetry packet
\[
\mathcal P(\rho)=\{\rho,\bar\rho,1-\rho,1-\bar\rho\},
\qquad
\rho=\beta+i\gamma,\quad \beta\ne\frac12,
\]
admits a nearby reduced good patch
\[
b=(m,s,\mathcal C_k(m))\in\mathcal B_{\rm red},
\]
or else the failure lies in a controlled component of
\(\mathcal E_{\rm geom}^{\core}\).
\end{hypothesis}

\begin{remark}[Status of the reach hypothesis]
\label{rem:status-reduced-good-patch-reach}
The strict-gap, same-scale, endpoint, symmetry-closure, and Gram/whitening
conditions are open once they hold at one point.  Thus
Hypothesis~\ref{hyp:reduced-good-patch-reach} is a non-swallowing assertion:
the enlarged exceptional set
\[
\mathcal E_{\rm geom}^{\core}
\]
must not contain all RH-relevant off-line packet configurations.  This is a
global reach condition for the rank-two route, not a consequence of Gate~A or
Gate~B.
\end{remark}

The RH-relevant local object is not a single zero but the symmetry packet
\[
\mathcal P(\rho)
=
\{\rho,\bar\rho,1-\rho,1-\bar\rho\},
\qquad
\rho=\beta+i\gamma,\quad \beta\ne\frac12.
\]
Packet-level ordering is essential: functional-equation and conjugation
symmetries force internal equalities inside \(\mathcal P(\rho)\), and these
should not be counted as core-selection failures.  The dangerous tie walls are
ties between distinct packets.

For packet visibility, set
\[
\mathcal B_{\rm red}^{+}
=
\mathcal B_{\rm red}
\setminus
\mathcal E_{\rm invis}^{\core,7}.
\]
For the centered one-whitening-slot route, the incidence walls include
\[
\{y=0\}\cup\{q''=0\}\cup\{g=h\},
\]
and, when quotient isolation is used, also \(\{v_5=0\}\).  For the
non-centered \(q^{(7)}\)-transport route, include
\[
\{q^{(7)}=0\}\cup\{\Xi_{K_1}=0\}.
\]
On \(\mathcal B_{\rm red}^{+}\), the off-line packet fixed-sector visibility
vector
\[
C_{\rm fd}^{\vec f}
=
(A_7^{\mathfrak f},A_9^{\mathfrak f},\ldots)
\]
is nonzero modulo these explicit incidence exclusions.

\subsection{Rank-two fixed-sector bundle as the remaining viable language}
\label{subsec:rank-two-fixed-sector-remaining-language}

The preceding failures suggest that the defect should not be scalarized, nor
should one single out the \(A_5^{\mathfrak f}\)-line as the primary geometric
object.  The best remaining local language is the full rank-two fixed-sector
bundle
\[
\mathcal F\to\mathcal B,
\qquad
\mathcal F_b=\mathfrak f_h(b)=\mathbf C I_b\oplus\mathbf C S_b.
\]
This formulation keeps the whole fixed-sector vector
\[
A_{2r+1}^{\mathfrak f}\in\mathcal F_b
\]
and therefore avoids the incompleteness of the quotient scalar and the
noncanonical line-preservation problem of \(L=\mathbf C A_5^{\mathfrak f}\).

The local candidate connection is the projected ambient derivative
\[
\nabla_X^{\mathfrak f}A
=
\Pi_{\mathfrak f}\bigl(d_X^{\mathcal J}A\bigr),
\]
where \(\mathcal J\) is the ambient corrected finite-\(s\) jet/kernel-state
bundle, \(d^{\mathcal J}\) is the ambient derivative, and
\[
\Pi_{\mathfrak f}:\mathcal J_b\to\mathfrak f_h(b)
\]
is the fixed-sector projection.  In a local frame
\[
E=(I,S),
\qquad
A=Ea,
\]
the connection has the form
\[
\nabla^{\mathfrak f}A
=
E(da+\Gamma^{\mathfrak f}a).
\]
If
\[
E=E^{\rm raw}W,
\]
then formally
\[
\Gamma^{\mathfrak f}
=
W^{-1}\Gamma^{\rm raw}W+W^{-1}dW.
\]
In particular, if the raw frame is flat in the chart, then
\[
\Gamma^{\mathfrak f}=W^{-1}dW.
\]

On an overlap with
\[
E_j=E_iG_{ij},
\qquad
G_{ij}\in GL_2(\mathbf C),
\]
the connection matrices must satisfy
\[
\Gamma_j
=
G_{ij}^{-1}\Gamma_iG_{ij}
+
G_{ij}^{-1}dG_{ij}.
\]
Thus the rank-two connection is canonical if the fixed-sector frames glue by
canonical \(GL_2\)-transition maps and the fixed-sector projection is
functorial under the ambient transition maps.

\begin{proposition}[Conditional rank-two connection theorem]
\label{prop:conditional-rank-two-connection}
Assume that on a reduced good-patch base \(\mathcal B_{\rm red}^{+}\):
\begin{enumerate}
\item the core/source convention is fixed and locally stable;
\item no zero moves into or out of \(q_{\core}\);
\item the fixed-sector frames satisfy transition laws
\[
E_j=E_iG_{ij},
\qquad
G_{ij}\in GL_2(\mathbf C),
\qquad
G_{ij}G_{jk}=G_{ik};
\]
\item the ambient corrected jet bundle \(\mathcal J\) and projection
\(\Pi_{\mathfrak f}\) are functorial under these transitions;
\item the whitening branch and Gram normalization are fixed or functorial.
\end{enumerate}
Then
\[
\nabla_X^{\mathfrak f}A
=
\Pi_{\mathfrak f}(d_X^{\mathcal J}A)
\]
defines a canonical connection on
\[
\mathcal F_{\rm red}^{+}\to\mathcal B_{\rm red}^{+}.
\]
\end{proposition}

The reduced-base hypotheses are essential.  On the unreduced base, changing the
source/core convention can produce affine corrections
\[
A_j^{\mathfrak f}
=
G_{ij}^{-1}A_i^{\mathfrak f}
+
B_{ij}.
\]
Such transformations do not define an ordinary vector-bundle connection.  Thus
the present rank-two connection is currently a fixed-convention, reduced-base
object unless an affine torsor transition theorem is proved.

\begin{remark}[Why rank two improves on the line bundle]
\label{rem:rank-two-improves-on-line-bundle}
The \(A_5^{\mathfrak f}\)-line may rotate under legitimate changes of
trivialization, source/core convention, or whitening.  This blocked the
line-bundle route.  The full plane \(\mathfrak f_h=\mathbf C I\oplus\mathbf C S\)
can still transform by a \(GL_2\)-matrix.  Hence the rank-two bundle is the
right level at which to formulate covariance.
\end{remark}

\begin{definition}[Row-normal form for the \(2Y\) live fiber]
\label{def:rank-two-row-normal-form-2Y}
Let \(\ell_7\) be the covariant order-seven live row and let \(\omega_Y\) be a
primitive \(Y12\)-row with \(\omega_Y(e_Y)=1\).  Set
\[
  \kappa:=\omega_Y(e_7),
  \qquad
  \ell_9^Y:=\omega_Y-\kappa\ell_7,
  \qquad
  \beta_9^{\mathrm{tri}}:=\beta_9-\kappa c_7 .
\]
The \(2Y\)-fiber is in row-normal form if
\[
  (\ell_7,\ell_9^Y)(e_7)=(1,0),
  \qquad
  (\ell_7,\ell_9^Y)(e_Y)=(0,1).
\]
All branch classifications below are made in this triangularized row basis.
\end{definition}

\begin{remark}[Triangularized rank-one fallback ratio]
\label{rem:rank-two-triangular-rank-one-ratio}
If the live obstruction collapses to a graph line
\[
  B_Y=\mathbf C(1,\rho_Y)
\]
in the triangularized basis, then the graph-killing row gives the scalar
rank-one wall
\[
  H\left((\beta_9^{\mathrm{tri}}-\rho_Yc_7)\Gamma_G+\alpha_9\Gamma_{G,2}\right).
\]
Hence the corrected rank-one ratio is
\[
  \rho_\Gamma^{\mathrm{tri}}
  :=
  \frac{\beta_9^{\mathrm{tri}}-\rho_Yc_7}{\alpha_9}.
\]
The raw-\(\beta_9\) expression is historical only and is not used in live
row-normal formulas.
\end{remark}

\begin{theorem}[Two-row exact sequence for the \(2Y\) live fiber]
\label{thm:two-row-exact-sequence-2Y}
Let the septic row be the covariant quotient row
\[
  \ell_7=\frac{\Delta_7^\nabla}{\Delta_7^\nabla(C_7)}.
\]
Let the primitive \(Y12\)-whitening row be a covector
\[
  \omega_Y:W\to\mathbf C
\]
such that
\[
  \omega_Y(a_Y)=1,
  \qquad
  W_{\mathrm{dead}}^{\mathrm{fixed}}\subseteq\ker\omega_Y.
\]
Set
\[
  \kappa:=\omega_Y(e_7),
  \qquad
  \ell_9^Y:=\omega_Y-\kappa\ell_7,
  \qquad
  \beta_9^{\mathrm{tri}}:=\beta_9-\kappa c_7.
\]
Assume the septic exact-sequence conditions
\[
  \mathcal A_7^{\loc}\subseteq\ker\Delta_7^\nabla,
  \qquad
  \Delta_7^\nabla(C_7)\ne0,
  \tag{R7}
\]
and the primitive \(Y12\)-whitening conditions
\[
  P_W\mathcal A_Y^{\loc}\subseteq W_{\mathrm{dead}}^{\mathrm{fixed}},
  \qquad
  a_Y\notin W_{\mathrm{dead}}^{\mathrm{fixed}}
  \quad\text{with polynomial separation}.
  \tag{RY}
\]
Assume also the grade-separation condition
\[
  \ell_7(e_Y)=0
  \tag{G}
\]
and the polynomial row bound
\[
  |\omega_Y(e_7)|\le Q^A .
\]
Then \((\ell_7,\ell_9^Y)\) is a polynomially bounded live row pair satisfying
\[
  (\ell_7,\ell_9^Y)(e_7)=(1,0),
  \qquad
  (\ell_7,\ell_9^Y)(e_Y)=(0,1).
\]
Moreover, if
\[
  \ell_7(A_{2Y})=0,
  \qquad
  \omega_Y(A_{2Y})=0,
\]
then
\[
  \ell_7(A_{2Y})=0,
  \qquad
  \ell_9^Y(A_{2Y})=0.
\]
Thus the \(2Y\)-local image is zero in the row-normalized live fiber:
\[
  B_Y=0.
\]
\end{theorem}

\begin{proof}[Proof sketch]
The septic exact sequence gives a bounded row \(\ell_7\) detecting \(C_7\)
and killing all admissible \(s_7\)-dead columns.  The whitening exact sequence
gives a bounded primitive row \(\omega_Y\) detecting \(a_Y\) and killing the
fixed-\(\alpha\) actual \(Y12\)-slot image.  Triangularization removes the
allowed leakage \(\omega_Y(e_7)=\kappa\) by setting
\[
  \ell_9^Y=\omega_Y-\kappa\ell_7.
\]
The condition \(\ell_7(e_Y)=0\) is the non-negotiable grade-separation condition:
the primitive order-nine \(Y12\)-carrier has no covariant order-seven septic
shadow.  The row-normalized conclusions follow immediately.
\end{proof}

\begin{theorem}[Rank-two closure spine]
\label{thm:rank-two-closure-spine}
Assume
\[
  \mathrm{TwoRowExactSequence}_{2Y},
  \qquad
  \mathrm{RowNormalForm}_{2Y},
  \qquad
  \mathrm{NuisancePreservesRows}_{2Y}.
\]
Then
\[
  B_Y=0.
\]
Consequently the row-normal rank selector of
Definition~\ref{def:rank-two-selector-bridge-Y12} satisfies
\[
  r_Y
  =
  \operatorname{rank}(M_{\mathrm{D4}}^{\mathrm{tri}}).
\]
If moreover \(c_7\alpha_9\ne0\), then \(r_Y=2\), and the existing rank-two D4
branch activates.  The D4 matrix in row-normal coordinates is
\[
  M_{\mathrm{D4}}^{\mathrm{tri}}
  =
  \begin{pmatrix}
    c_7&0\\
    \beta_9^{\mathrm{tri}}&\alpha_9
  \end{pmatrix},
\]
and D4 supplies the plain wall
\[
  F_{2Y}^{\mathrm{tri}}
  =
  H
  \begin{pmatrix}
    c_7\Gamma_G\\
    \beta_9^{\mathrm{tri}}\Gamma_G+\alpha_9\Gamma_{G,2}
  \end{pmatrix}.
\]
On coefficient-good cores this is polynomially equivalent to
\[
  H(\Gamma_G,\Gamma_{G,2}).
\]
No derivative wall components \(D\Gamma_G\) or \(D^2\Gamma_G\) are supplied by
this rank-two D4 statement.
\end{theorem}

% === wip-rank-two-selector-bridge-Y12 ===
% Purpose:
%   Connect the row-normal \(2Y\) rank-two package to the paper's rank selector.
% Depends on:
%   RowNormalForm_{2Y}, B_Y, and the triangularized D4 matrix.
% Status:
%   Clarifying bridge / finite-audit convention.

\begin{definition}[Row-normal rank selector for the \(2Y\) branch]
\label{def:rank-two-selector-bridge-Y12}
Let \(B_Y\subset\mathbf C^2\) be the actual/local live-fiber image in the
row-normal coordinates
\[
  (s_7,s_9^Y):=(\ell_7,\ell_9^Y),
  \qquad
  \ell_9^Y=\omega_Y-\omega_Y(e_7)\ell_7.
\]
Let
\[
  M_{\mathrm{D4}}^{\mathrm{tri}}
  :=
  \begin{pmatrix}
    c_7&0\\
    \beta_9^{\mathrm{tri}}&\alpha_9
  \end{pmatrix},
  \qquad
  \beta_9^{\mathrm{tri}}:=\beta_9-\omega_Y(e_7)c_7.
\]
The \(2Y\)-rank selector is
\[
  r_Y
  :=
  \operatorname{rank}
  \left(
    \Pi_{\mathbf C^2/B_Y}M_{\mathrm{D4}}^{\mathrm{tri}}
  \right).
\]
Thus the relevant rank is not the raw rank of formal D4 carriers.  It is the
rank of the row-normal D4 map after quotienting by actual/local source freedom.
\end{definition}

\begin{theorem}[Rank-two bridge to the existing D4 branch]
\label{thm:rank-two-bridge-Y12}
Assume the live fiber is in row-normal form:
\[
  (\ell_7,\ell_9^Y)(e_7)=(1,0),
  \qquad
  (\ell_7,\ell_9^Y)(e_Y)=(0,1),
\]
with polynomially bounded rows.  If the local row-exactness package gives
\[
  B_Y=0,
\]
then
\[
  r_Y
  =
  \operatorname{rank}(M_{\mathrm{D4}}^{\mathrm{tri}})
  =
  2
  \quad\Longleftrightarrow\quad
  c_7\alpha_9\ne0.
\]
In this case the new \(2Y\)-rank-two branch is exactly the existing rank-two
D4 branch, interpreted as two surviving D4 coordinates after quotienting by
actual/local source freedom.  The D4 output is
\[
  F_{2Y}^{\mathrm{tri}}
  =
  H
  \begin{pmatrix}
    c_7\Gamma_G\\
    \beta_9^{\mathrm{tri}}\Gamma_G+\alpha_9\Gamma_{G,2}
  \end{pmatrix},
\]
and on coefficient-good cores,
\[
  |F_{2Y}^{\mathrm{tri}}|^2
  \asymp_{Q^A}
  |H|^2
  \left(
    |\Gamma_G|^2+|\Gamma_{G,2}|^2
  \right).
\]
Equivalently, the honest rank-two wall is the plain wall
\[
  H(\Gamma_G,\Gamma_{G,2}).
\]
No derivative wall components \(D\Gamma_G\) or \(D^2\Gamma_G\) are supplied by
this rank selector.
\end{theorem}

\begin{remark}[Branch classification by \(B_Y\)]
\label{rem:rank-selector-BY-branch-table}
The row-normal selector gives the following branch table:
\[
\begin{array}{c|c|c}
B_Y & r_Y & \text{paper branch}\\
\hline
0
& 2
& \text{rank-two D4 vector wall}\\
\mathbf C(1,\rho_Y)
& 1
& \text{rank-one moving-wall fallback}\\
\mathbf C(0,1)
& 0\text{ for the }Y12\text{ column}
& \text{vertical }Y12\text{ death}\\
\mathbf C^2
& 0
& \text{live fiber absorbed}
\end{array}
\]
In the graph case \(B_Y=\mathbf C(1,\rho_Y)\), the graph-killing row gives the
rank-one scalar wall
\[
  H\left(
    (\beta_9^{\mathrm{tri}}-\rho_Yc_7)\Gamma_G
    +
    \alpha_9\Gamma_{G,2}
  \right).
\]
Thus the corrected triangularized fallback ratio is
\[
  \rho_\Gamma^{\mathrm{tri}}
  =
  \frac{\beta_9^{\mathrm{tri}}-\rho_Yc_7}{\alpha_9}.
\]
The raw-\(\beta_9\) expression is historical only and is not used in live
row-normal formulas.
\end{remark}

\begin{remark}[Do not conflate with the old rank-two norm route]
\label{rem:rank-two-selector-not-old-rank-two-norm-route}
The \(2Y\)-rank-two bridge is not the old abstract rank-two norm route.  It is
the concrete statement
\[
  \operatorname{rank}
  \left(
    \Pi_{\mathbf C^2/B_Y}M_{\mathrm{D4}}^{\mathrm{tri}}
  \right)=2.
\]
It may supersede or bypass older rank-two norm language, but it should not be
identified with any no-go attached to a different rank-two norm-control route
unless the objects are explicitly the same.
\end{remark}

\begin{remark}[Three normal forms for the rank-two branch]
\label{rem:rank-two-three-normal-forms}
The rank-two branch is organized by three internal normal forms.  If the final
contradiction is a global forbidden identity rather than the scalar plain-wall
route, a fourth forbidden-channel normal form is also required.

First, the local live fiber is always put in row normal form:
\[
  \kappa=\omega_Y(e_7),
  \qquad
  \ell_9^Y=\omega_Y-\kappa\ell_7,
  \qquad
  \beta_9^{\mathrm{tri}}=\beta_9-\kappa c_7.
\]
All branch classifications are made in the triangularized basis
\[
  (\ell_7,\ell_9^Y),
\]
not in the raw row pair \((\ell_7,\omega_Y)\).

Second, the analytic wall is put in the source-curvature variable
\[
  p=q'',
\]
so that, up to fixed nonzero normalization constants,
\[
  \Gamma_G=-6p+36q-12q^3,
  \qquad
  \Gamma_{G,2}=p''-6p+6q^2p.
\]
Thus the finite wall spaces are built from a finite source space
\(\mathcal P_I\) for \(p\), together with the affine reconstruction
\[
  \mathcal V_I=\mathcal A_1+\mathcal I^2\mathcal P_I.
\]

Third, weighted transfer is made through predetermined coefficient/smooth-good
cores
\[
  E_I^{\core}\subset I,
  \qquad
  |E_I^{\core}|\ge R^{-C}|I|,
\]
chosen before wall extraction and independent of
\[
  \Gamma_G,\quad \Gamma_{G,2},\quad W_2,
  \qquad
  W_2:=|\Gamma_G|^2+|\Gamma_{G,2}|^2.
\]
This prevents post-hoc wall-dependent selection.

If the route uses a global forbidden channel, the fourth normal form is
forbidden-channel normal form.  One fixes
\[
  M_{\mathrm{forb}}:=\Lambda E_{\mathrm{car}}L_{\mathrm{clean}}
\]
on the D4 live plane and extracts against
\[
  E_{\mathrm{wall}}^{\mathrm{forb}}
  =
  \frac{|H|^2|M_{\mathrm{forb}}\Gamma|^2}{\mathcal N_{\mathrm{D4}}},
  \qquad
  \Gamma=(\Gamma_G,\Gamma_{G,2})^T.
\]
The current concrete Type-I candidate is the transition--singular discard
operator
\[
  \Lambda_{\mathrm{tr},\mathrm{sing},p}
  =
  (I-P_{\mathrm{sing},p})T_{\mathrm{tr}\to\mathrm{sing},p}P_{\mathrm{tr},p}.
\]
\end{remark}

\begin{remark}[Admissibility before quotient, projection, and selection]
\label{rem:rank-two-admissibility-before-projection}
The rank-two branch uses an admissibility-first convention at every layer.

A direction may enter the exact local denominator only after it has been proved
row-dead and live-pair preserving.  A finite coherent direction is routed to
\(\mathcal C_{2Y}^{\mathrm{arith}}\), and a small direction is routed to \(R_X^{2Y}\);
neither is silently quotiented as exact nuisance.

Similarly, finite-atlas Riesz estimates are applied only after wall-variable
completion has produced finite representatives
\[
  G_X\in\mathcal W_G,
  \qquad
  G_{2,X}\in\mathcal W_{G2}
\]
with
\[
  \|\Gamma_G-G_X\|_{H^2_{\mathrm{sc}}(I)}
  +
  \|\Gamma_{G,2}-G_{2,X}\|_{L^2(I)}
  \le Q^{-B}.
\]
Finally, core sets \(E_I^{\core}\) for weighted transfer must be chosen from
coefficient, smooth, and chart data before observing the wall \(W_2\).
\end{remark}

% === wip-current-architecture-l2-gate-one-forbidden-d4 ===
% Purpose:
%   Record the current post-row-normal cleanup architecture.
% Status:
%   Active architecture / finite-audit ledger.

\begin{remark}[Current architecture after the \(L^2\) Gate--1 repair]
\label{rem:current-architecture-l2-gate-one-forbidden-d4}
The active architecture is now
\[
  \boxed{
  \mathrm{Agg}L^2(q''_{\can})
  +
  \mathrm{Finite/O2O3}_{23}
  +
  \mathrm{ForbiddenSingularDiscardD4}.
  }
\]
Here Gate--1 is local scalar \(L^2\)-energy of the full canonical source
\[
  q''_{\can}=B_2+P_2+D_2+E_2.
\]
Old branch observables such as \(T_{\can}X_{\can}\), transition jets,
raw singular sine moments, and projected singular low modes are proof tools or
diagnostics.  They are not the current Gate--1 definition.

The local rank-two/D4 cleanup remains valid: after \(B_Y=0\), D4 supplies only
the plain wall
\[
  H(\Gamma_G,\Gamma_{G,2}),
\]
not derivative wall components.

The finite/O2O3 module uses an analytic \(R^C\)-conditioned Riesz basis for the
finite augmented matrix.  The downstream forbidden channel, if it exists, uses a
separate carrier/source-coherence projection.  These two bases should not be
identified.

The best current downstream forbidden-channel candidate is the local
transition--singular discard operator
\[
  \Lambda_{\mathrm{tr},\mathrm{sing},p}:=(I-P_{\mathrm{sing},p})T_{\mathrm{tr}\to\mathrm{sing},p}P_{\mathrm{tr},p},
\]
not the full three-chart atlas loop.  The finite audit is whether
\[
  (I-P_{\mathrm{sing}})T_{\mathrm{tr}\to\mathrm{sing}}
  \operatorname{span}\{[\mathsf C_7]^{\mathrm{tr}},[\mathsf C_{Y12}]^{\mathrm{tr}}\}
\]
is nonzero.  If it is zero, the singular-discard channel is D4-invisible.  If
it has rank one or two, it gives a genuine local forbidden channel.
\end{remark}

% === audit-D4-needs-forbidden-channel ===
% Purpose:
%   Prevent the carrier-quotient D4 visibility theorem from being mistaken for
%   a contradiction theorem.
% Status:
%   Active no-go / finite-audit warning.

\begin{remark}[D4 visibility is not yet contradiction]
\label{rem:D4-visibility-not-contradiction}
Carrier-quotient D4 supplies local visibility of the row-normalized plain wall:
\[
  F_{2Y}^{\mathrm{tri}}
  =
  H
  \begin{pmatrix}
    c_7\Gamma_G\\
    \beta_9^{\mathrm{tri}}\Gamma_G+\alpha_9\Gamma_{G,2}
  \end{pmatrix}.
\]
This is not by itself a contradiction.  A contradiction requires a global
forbidden channel
\[
  \Lambda:\mathcal H_{\mathrm{row}}\to\mathcal K
\]
such that the global row object satisfies
\[
  \|\Lambda V_{\mathrm{glob}}\|^2\le Q^{-B}\mu(\mathcal C),
\]
or exactly \(\Lambda V_{\mathrm{glob}}=0\), while local D4 produces
\[
  \|\Lambda v_m\|^2\ge Q^{-A}
\]
on a polynomial-weight extracted set.

Thus D4 is currently a carrier-quotient visibility theorem unless a concrete
\(\Lambda\) is supplied and shown to see the live D4 carrier plane.
\end{remark}

\begin{definition}[Covariant septic row]
\label{def:rank-two-covariant-septic-row}
Let
\[
  A_5^{\mathfrak f}(m)=u_5I+v_5S,
  \qquad
  L_5(m):=\mathbf C A_5^{\mathfrak f}(m)\subset \mathfrak f_m .
\]
Define
\[
  \Delta_7^\nabla(u)
  :=
  \omega_m\!\left(Q_7^{\mathrm{pkg},\nabla}u,A_5^{\mathfrak f}(m)\right),
\]
as a covector on the moving quotient
\[
  \mathfrak f_m/L_5(m).
\]
Thus
\[
  \Delta_7^\nabla(u)=0
  \quad\Longleftrightarrow\quad
  Q_7^{\mathrm{pkg},\nabla}u\in \mathbf C A_5^{\mathfrak f}(m).
\]
The normalized septic row is
\[
  \ell_7(u):=\frac{\Delta_7^\nabla(u)}{\Delta_7^\nabla(C_7)}.
\]
Raw determinant drift in \(u_7v_5-u_5v_7\) is not a live \(s_7\)-motion unless
it survives this covariant quotient.
\end{definition}

\begin{definition}[Typed septic denominator]
\label{def:rank-two-typed-septic-denominator}
Define
\[
\begin{aligned}
  \mathcal D_7^{\mathrm{typed}}
  :=
  &\ \mathbf C A_5^{\mathfrak f}
  +\mathcal D_{7,\mathrm{EF}}
  +\mathcal D_{7,\mathrm{chart}}
  +\mathcal D_{7,\mathrm{norm}}
  +\mathcal D_{7,\gamma/\bg}        \\
  &\quad
  +\mathcal D_{7,\tail/\mathrm{res}}
  +\mathcal D_{7,\mathrm{wrong}}
  +\mathcal N_{7,\mathrm{exact}}.
\end{aligned}
\]
Here \(\mathcal N_{7,\mathrm{exact}}\) consists only of exact nuisance directions that
have already passed the row-deadness and live-pair preservation tests.
\end{definition}

\begin{theorem}[Septic exact sequence audit]
\label{thm:rank-two-septic-exact-sequence}
Let
\[
  \mathcal A_7^{\loc}
  :=
  \mathcal I_{\can}^{(2)}
  +T_{\mathrm{lower}}^{\mathrm{arith}}
  +T_{E_2/R_{\mathrm{lab}},\loc}^{s_7}
  +\mathcal N_{7,\mathrm{exact}}.
\]
Assume
\[
  \Delta_7^\nabla(\mathcal A_7^{\loc})=0,
  \tag{SE1}
\]
and
\[
  |\Delta_7^\nabla(C_7)|\ge Q^{-A}.
  \tag{SE2}
\]
Equivalently,
\[
  \operatorname{rank}[\mathcal D_7^{\mathrm{typed}},\mathcal A_7^{\loc}]
  =
  \operatorname{rank}\mathcal D_7^{\mathrm{typed}},
\]
while
\[
  \operatorname{dist}(C_7,\mathcal D_7^{\mathrm{typed}})\ge Q^{-A}\|C_7\|.
\]
Then
\[
  \ell_7(\mathcal A_7^{\loc})=0,
  \qquad
  \ell_7(C_7)=1.
\]
\end{theorem}

\begin{remark}[Actual lower kernel audit]
\label{rem:rank-two-actual-lower-kernel-audit}
The actual-lower part of \((\mathrm{SE1})\) is the finite kernel condition
\[
  \Delta_7^\nabla Q_7^{\mathrm{pkg},\nabla}D\mathcal S_{\mathrm{lower}}=0.
\]
Equivalently, the actual lower parameter space has no corrected septic
package-coordinate summand:
\[
  \Theta_{\mathrm{pkg},7}=0.
\]
Formal \(C_7\)-carriers and formal \(C_7\to Y12\) connections do not enter
\(B_Y\) unless actualized by a genuine actual-source, exact nuisance, or local
fitting tangent.
\end{remark}

\begin{definition}[Actual lower parameter package map]
\label{def:actual-lower-package-map}
Let \(\Theta_{\mathrm{lower}}^{\mathrm{arith}}\) denote the actual lower-source parameter space.
It contains only actual EF/gamma/chart/tail lower deformations, actual
normalization or representative lower freedoms, and genuine exact nuisance
lower directions.  It excludes formal D4 carrier variables, formal \(C_7\)
labels, formal \(C_7\to Y12\) connections without actual tangent, proof
thickening parameters, and free corrected-package coordinates not present in
the actual source.

Define the lower package map
\[
  \mathfrak d_{7,\mathrm{lower}}
  :=
  \Delta_7^\nabla Q_7^{\mathrm{pkg},\nabla}D\mathcal S_{\mathrm{lower}}
  :
  \Theta_{\mathrm{lower}}^{\mathrm{arith}}\to\mathbf C.
\]
\end{definition}

\begin{lemma}[Actual lower no-package-coordinate audit]
\label{lem:actual-lower-no-package-coordinate}
The actual lower contribution to the \(s_7\)-row is zero if and only if
\[
  \mathfrak d_{7,\mathrm{lower}}=0.
\]
Equivalently, the actual lower parameter space has no corrected septic
package-coordinate summand:
\[
  \Theta_{\mathrm{pkg},7}=0.
\]
A formal \(C_7\)-carrier or formal \(C_7\to Y12\) connection contributes to
\(B_Y\) only after actualization by a genuine actual-source, exact nuisance, or
atom-local tangent with nonzero \(\Delta_7^\nabla\).
\end{lemma}

\begin{remark}[Failure routing for actual lower package leakage]
\label{rem:actual-lower-package-leakage-routing}
If \(u_{\mathrm{lower}}\in T_{\mathrm{lower}}^{\mathrm{arith}}\) satisfies
\[
  \Delta_7^\nabla(u_{\mathrm{lower}})\ne0,
\]
then compute its triangularized live-fiber vector
\[
  (\ell_7(u_{\mathrm{lower}}),\ell_9^Y(u_{\mathrm{lower}})).
\]
If all such nonzero vectors span a graph line \(\mathbf C(1,\rho_Y)\), the
rank-one moving-wall fallback remains available.  If a vertical vector
\(\mathbf C(0,1)\), a second slope, or a full rank-two image appears, the
current \(Y12\) route dies locally or requires a different carrier.
\end{remark}

\begin{definition}[Whitening quotient for the primitive \(Y12\) row]
\label{def:whitening-quotient-Y12}
Let \(W\) be the finite same-point whitening-input space seen by the order-nine
singular extractor.  Define
\[
  W_{\mathrm{dead}}^{\mathrm{fixed}}
  :=
  W_{\mathrm{global}}
  +W_{\mathrm{chart}}
  +W_{\mathrm{norm}}
  +W_{\bg}
  +W_{\mathrm{trans}}
  +W_{\mathrm{wrong}}.
\]
Here \(W_{\mathrm{trans}}\) denotes transported same-point whitening data and
\(W_{\mathrm{wrong}}\) denotes wrong-origin terms, such as \(DW[DN_0]\),
\(D^2W[N_0]\), or whitening applied to residual/completion data rather than
to the primitive \(D^2N_0\)-carrier.  Let
\[
  a_Y\in W
\]
be the live same-point whitening input in the primitive carrier
\[
  DW_{\mathrm{same}}[a_Y]\,D^2N_0[\xi_1,\xi_2]
\]
with mixed branch \(12\), same-point incidence, and one whitening slot.
\end{definition}

\begin{theorem}[Whitening exact sequence for fixed-\(\alpha\) \(Y12\)]
\label{thm:whitening-exact-sequence-Y12}
Assume
\[
  P_WV_{\mathrm{arith}}^{\mathrm{fixed}}\subseteq W_{\mathrm{dead}}^{\mathrm{fixed}}
  \tag{W1}
\]
and
\[
  \operatorname{dist}(a_Y,W_{\mathrm{dead}}^{\mathrm{fixed}})
  \ge Q^{-A}|a_Y|.
  \tag{W2}
\]
Then there is a polynomially bounded covector
\[
  \omega_Y:W\to\mathbf C
\]
such that
\[
  \omega_Y(a_Y)=1,
  \qquad
  \omega_Y(W_{\mathrm{dead}}^{\mathrm{fixed}})=0,
  \qquad
  \omega_Y(P_WV_{\mathrm{arith}}^{\mathrm{fixed}})=0.
\]
Consequently fixed-\(\alpha\) actual first and second canonical variations do
not produce a primitive \(Y12\)-row column.

If \(W1\) fails, the fallback audit is the ordered slot cascade:
\[
  \text{same-point whitening slot}
  \longrightarrow
  \text{source slot }1
  \longrightarrow
  \text{source slot }2
  \longrightarrow
  \text{coupled bilinear image}.
\]
Only failure of the final coupled bilinear image creates an actual \(Y12\)-live
column.
\end{theorem}

\begin{remark}[Whitening naturality also controls labelled transition]
\label{rem:whitening-naturality-controls-labelled-transition}
The same whitening quotient controls the \(Y12\)-part of the labelled/raw
transition defect.  Let
\[
  \mathfrak W_Y
  :=
  DW_{\mathrm{sing},\mathrm{same}}
  -
  S_WDW_{\mathrm{tr},\mathrm{same}}S^{-1}.
\]
The \(Y12\)-safe labelled-transition condition is
\[
  D_{\loc}\,
  P_{(W_{\mathrm{dead}}^{\mathrm{fixed}})^\perp}
  \bigl(\mathfrak W_Y[D^2N_0]\bigr)=0.
\]
Equivalently, same-point whitening may be transported or vary coherently, but
it must not be refit atom-by-atom on the primitive carrier.  Fixed coherent
defects are routed to \(\mathcal C_{2Y}^{\mathrm{arith}}\); residual defects are routed
to \(R_X^{2Y}\); atom-local live defects enter \(B_Y\).
\end{remark}

\begin{definition}[Locality filtration for completion and labelled defects]
\label{def:locality-filtration-2Y}
For any \(2Y\)-fiber-valued expression \(X\), write
\[
  X=X_{\mathrm{fixed}}+X_{\mathrm{coh}}+X_{\mathrm{res}}+X_{\loc}.
\]
Only \(X_{\loc}\) contributes to \(B_Y\).  The fixed part is absorbed into
basepoint/recentering data, the coherent part is routed to
\(\mathcal C_{2Y}^{\mathrm{arith}}\), and the residual part is routed to \(R_X^{2Y}\).
\end{definition}

\begin{theorem}[No atom-fitting for \(E_2\) and labelled transition]
\label{thm:no-atom-fitting-2Y-plus}
Assume the following three guardrails.

\begin{enumerate}
\item No local live completion counterterms:
\[
  \operatorname{Loc}\Pi_{\mathrm{fib},Y}E_{\mathrm{model}}=0.
\]
In particular, no local terms of the forms
\[
  \eta_7(p)e_7(p)+\eta_Y(p)e_Y(p),
  \qquad
  \eta(p)(e_7(p)+\rho_Ye_Y(p))
\]
are admitted.

\item Normal form data are not atom-fitted:
\[
D_{\loc}\Pi_{\mathrm{fib},Y}
\left(
D^2N_{0,\mathrm{sing}}\circ(L,L)-S D^2N_{0,\mathrm{tr}}
\right)=0.
\]

\item Same-point whitening is not atom-fitted:
\[
D_{\loc}\Pi_{\mathrm{fib},Y}
\left[
\left(
DW_{\mathrm{sing},\mathrm{same}}-S_WDW_{\mathrm{tr},\mathrm{same}}S^{-1}
\right)[D^2N_0]
\right]=0.
\]
\end{enumerate}

Then
\[
  D_{\loc}\Pi_{\mathrm{fib},Y}(E_2+R_{\mathrm{lab}})=0.
\]
Thus \(E_2/R_{\mathrm{lab}}\) affects \(B_Y\) only through atom-local live-fiber
fitting.  Fixed, coherent, and residual effects do not enter \(B_Y\).
\end{theorem}

\begin{remark}[Amplitude projection is downstream of \(B_Y=0\)]
\label{rem:amplitude-projection-after-BY}
The proof of \(B_Y=0\) must precede any amplitude projection.  A finite coherent
space \(\mathcal C_{2Y}^{\mathrm{arith}}\) may be introduced only after the local
live-fiber quotient has been fixed.  In particular, a projected wall
\[
  (I-P_{\mathcal C_{2Y}^{\mathrm{arith}}})F_{2Y}
\]
must not be used to define \(D_{2Y}\), the live rows, or the rank-two quotient,
unless the projection is itself proved to have bounded local branch complexity.
\end{remark}

\begin{definition}[Exact/coherent/residual nuisance split]
\label{def:rank-two-nuisance-split}
Split the declared singular nuisance space as
\[
  \mathcal N_{\mathrm{sing}}
  =
  \mathcal N_{\mathrm{exact}}
  +
  \mathcal N_{\mathrm{coh}}
  +
  \mathcal N_{\mathrm{res}}.
\]
Only \(\mathcal N_{\mathrm{exact}}\) enters the exact local denominator.  The coherent
part is routed to \(\mathcal C_{2Y}^{\mathrm{arith}}\), and the residual part is routed
to \(R_X^{2Y}\).  Thus
\[
  D_{2Y}:=D_{2Y}^{0}+\mathcal N_{\mathrm{exact}},
\]
not \(D_{2Y}^{0}+\mathcal N_{\mathrm{sing}}\).
\end{definition}

\begin{definition}[Admissible exact nuisance]
\label{def:admissible-exact-nuisance-2Y}
An exact nuisance generator \(n\in\mathcal N_{\mathrm{exact}}\) is admissible for the
rank-two quotient only if
\[
  \Delta_7^\nabla(n)=0,
  \qquad
  \omega_Y(P_Wn)=0.
\]
Any declared nuisance generator failing either row test is not exact nuisance.
It must instead be routed to one of:
\[
  \mathcal C_{2Y}^{\mathrm{arith}},
  \qquad
  R_X^{2Y},
  \qquad
  \text{quotient repair},
  \qquad
  \text{local rank failure}.
\]
\end{definition}

\begin{remark}[Two-grade nuisance homogeneity]
\label{rem:two-grade-nuisance-homogeneity}
The exact nuisance denominator must not identify the two live grades.  In
particular, one must rule out
\[
  e_Y-\lambda e_7\in D_{2Y}^{0}+\mathcal N_{\mathrm{exact}}
\]
with polynomially small angle.  Equivalently, exact nuisance should be
grade-homogeneous,
\[
  \mathcal N_{\mathrm{exact}}
  =
  \mathcal N_{7,\mathrm{exact}}\oplus\mathcal N_{9,\mathrm{exact}},
\]
or at least block-triangular with no live \(7\leftrightarrow9\) relation.
\end{remark}

\begin{lemma}[Nuisance preservation of the live rows]
\label{lem:rank-two-nuisance-preserves-rows}
Assume
\[
  \Delta_7^\nabla(\mathcal N_{\mathrm{exact}})=0,
  \qquad
  \omega_Y(P_W\mathcal N_{\mathrm{exact}})=0.
\]
Assume also the live-pair Schur condition
\[
  \lambda_{\min}
  \left(
  [e_7,e_Y]^*
  (I-P_{D_{2Y}^{0}+\mathcal N_{\mathrm{exact}}})
  [e_7,e_Y]
  \right)
  \ge Q^{-A}.
\]
Then
\[
  (\ell_7,\ell_9^Y)(\mathcal N_{\mathrm{exact}})=0,
\]
and the live pair remains two-dimensional and polynomially conditioned modulo
\(D_{2Y}^{0}+\mathcal N_{\mathrm{exact}}\).
\end{lemma}

\begin{proof}[Proof sketch]
The first row-deadness condition gives \(\ell_7(\mathcal N_{\mathrm{exact}})=0\).
The second gives \(\omega_Y(P_W\mathcal N_{\mathrm{exact}})=0\).  Since
\[
  \ell_9^Y=\omega_Y-\kappa\ell_7,
\]
we get \(\ell_9^Y(\mathcal N_{\mathrm{exact}})=0\).  The Schur condition rules out
\[
  e_7\in D_{2Y}^{0}+\mathcal N_{\mathrm{exact}},
  \qquad
  e_Y\in D_{2Y}^{0}+\mathcal N_{\mathrm{exact}},
\]
and
\[
  e_Y-\lambda e_7\in D_{2Y}^{0}+\mathcal N_{\mathrm{exact}}.
\]
Thus exact nuisance is both row-dead and live-pair preserving.
\end{proof}

\begin{corollary}[Nuisance-safe contribution to \(B_Y=0\)]
\label{cor:rank-two-nuisance-safe-by-rows}
Assume \(\mathcal N_{\mathrm{sing}}\) has been split as
\[
  \mathcal N_{\mathrm{sing}}
  =
  \mathcal N_{\mathrm{exact}}
  +
  \mathcal N_{\mathrm{coh}}
  +
  \mathcal N_{\mathrm{res}},
\]
and that only \(\mathcal N_{\mathrm{exact}}\) enters the exact denominator.  Suppose
\(\mathcal N_{\mathrm{exact}}\) is row-dead:
\[
  \Delta_7^\nabla(\mathcal N_{\mathrm{exact}})=0,
  \qquad
  \omega_Y(P_W\mathcal N_{\mathrm{exact}})=0,
\]
and preserves the live pair:
\[
  e_7\notin D_{2Y}^{0}+\mathcal N_{\mathrm{exact}},
  \qquad
  e_Y\notin D_{2Y}^{0}+\mathcal N_{\mathrm{exact}},
\]
and
\[
  e_Y-\lambda e_7\notin D_{2Y}^{0}+\mathcal N_{\mathrm{exact}}
  \quad\text{for every }\lambda\in\mathbf C,
\]
with polynomial angle.  Then nuisance contributes no local column to the
row-normal live fiber:
\[
  (\ell_7,\ell_9^Y)(\mathcal N_{\mathrm{exact}})=0.
\]
The coherent part \(\mathcal N_{\mathrm{coh}}\) is routed to
\(\mathcal C_{2Y}^{\mathrm{arith}}\), and the residual part \(\mathcal N_{\mathrm{res}}\) is
routed to \(R_X^{2Y}\).  Neither coherent nor residual nuisance may be placed
in the exact denominator merely by name.
\end{corollary}

\begin{theorem}[Row-normal live-fiber avoidance]
\label{thm:rank-two-row-normal-live-fiber-avoidance}
Let the live rows be in row-normal form:
\[
  \ell_7=\frac{\Delta_7^\nabla}{\Delta_7^\nabla(C_7)},
  \qquad
  \ell_9^Y=\omega_Y-\kappa\ell_7,
  \qquad
  \kappa=\omega_Y(e_7).
\]
Assume
\[
  \ell_7(e_7)=1,\qquad \ell_7(e_Y)=0,\qquad
  \omega_Y(e_Y)=1,\qquad |\omega_Y(e_7)|\le Q^A,
\]
and polynomial row bounds.  Assume further that the five row blockers vanish:
\[
  \Delta_7^\nabla(D_{\actual}T\mathcal M_Y^{\mathrm{arith}})=0,
\]
\[
  \omega_Y(P_WD_{\actual}T\mathcal M_Y^{\mathrm{arith}})=0,
\]
\[
  D_{\loc}(\ell_7,\ell_9^Y)(E_2+R_{\mathrm{lab}})=0,
\]
\[
  \Delta_7^\nabla(\mathcal N_{\mathrm{exact}})=0,
\]
and
\[
  \omega_Y(P_W\mathcal N_{\mathrm{exact}})=0.
\]
Then
\[
  B_Y=0.
\]
Consequently, if the triangularized D4 coefficients are polynomially conditioned,
D4 gives the plain rank-two wall
\[
  F_{2Y}^{\mathrm{tri}}
  =
  H
  \begin{pmatrix}
    c_7\Gamma_G\\
    \beta_9^{\mathrm{tri}}\Gamma_G+\alpha_9\Gamma_{G,2}
  \end{pmatrix}
  \asymp_{Q^A}
  H(\Gamma_G,\Gamma_{G,2}).
\]
No derivative wall components are produced by this theorem.
\end{theorem}

\subsection{Order-seven packet visibility on the reduced base}
\label{subsec:order-seven-packet-visibility-reduced-base}

For this and the following subsections, write the centered diagonal Gram
data as
\[
g=\frac{q_0}{\pi},
\qquad
h=\frac1\pi\left(\frac13q_0^3+\frac16q''\right),
\]
so that the centered same-point Gram block at \(s=0\) is, up to the fixed
ambient normalization,
\[
G_0=\operatorname{diag}(2g,2h).
\]

\begin{definition}[First and second Gram-anisotropy functionals]
\label{def:gram-anisotropy-functionals}
On a centered reduced patch, unless explicitly stated otherwise, \(q_0\),
\(q''\), and \(q^{(4)}\) mean the full actual-zeta centered jets
\[
q_0=q_\zeta(m),
\qquad
q''=q''_{\can}(m)=q_\zeta''(m),
\qquad
q^{(4)}=(q_\zeta)^{(4)}(m).
\]
Define
\[
\Gamma_G^{\can}
:=
q''_{\can}-6q_0+2q_0^3 .
\]
For readability we usually suppress the superscript and write
\(\Gamma_G=\Gamma_G^{\can}\).  Then
\[
g-h
=
-\frac{\Gamma_G}{6\pi}.
\]
Thus the Gram-isotropy wall \(g=h\) is equivalently
\[
\Gamma_G=0.
\]

We also define the second Gram-anisotropy functional
\[
\Gamma_{G,2}^{\can}
:=
(1-q_0^2)q''_{\can}-\frac16(q_\zeta)^{(4)}.
\]
Again, when the canonical convention is fixed, we write
\(\Gamma_{G,2}=\Gamma_{G,2}^{\can}\).  This is the anisotropy of the second
centered Gram variation.  In particular, on \(g=h\), the first Gram anisotropy
has vanished, but the second one need not vanish.
\end{definition}

The first admissible visibility coefficient is not \(A_5^{\mathfrak f}\),
which defines the reference line, but \(A_7^{\mathfrak f}\).  The direct raw
off-diagonal order-seven term is not itself the witness.  At a centered even
packet,
\[
N_7^{\off}=C_7J,
\qquad
J=
\begin{pmatrix}
0&1\\
-1&0
\end{pmatrix}.
\]
Since \(X_{-,0}=X_{+,0}=X_0\),
\[
X_0N_7^{\off}X_0=C_7(\det X_0)J,
\]
so
\[
\pi_{\mathfrak f}(X_0N_7^{\off}X_0)=0.
\]

The live centered channel is instead the one-whitening-slot term
\[
T_7^{(1)}
=
X_{-,1}YX_{+,0}+X_{-,0}YX_{+,1},
\qquad
Y=yJ.
\]
At centered parity,
\[
X_{+,0}=X_{-,0}=X_0,
\qquad
X_{+,1}=-X_{-,1}.
\]
Write
\[
X_0=pI+rD,
\qquad
X_{-,1}=BS.
\]
Then
\[
T_7^{(1)}
=
yB(SJX_0-X_0JS)
=
-2yB(pD+rI),
\]
and hence
\[
\boxed{
\pi_{\mathfrak f}(T_7^{(1)})=-2yBr\,I.
}
\]

For the centered off-line packet,
\[
q''=2A_2,
\qquad
A_2=-4(a^{-3}+b^{-3})\ne0,
\]
and
\[
y_{\rm lead}
=
-\frac{q_0A_2^2}{72}.
\]
Moreover
\[
B=
\frac{q''}
{4\pi\sqrt{gh}(\sqrt g+\sqrt h)},
\qquad
r=\frac12(g^{-1/2}-h^{-1/2}).
\]
Thus the one-whitening-slot channel is nonzero outside the explicit incidence
walls
\[
y=0,\qquad q''=0,\qquad g=h.
\]

\begin{lemma}[Centered order-seven fixed-sector row and first Gram anisotropy]
\label{lem:centered-order-seven-gram-anisotropy}
At centered parity,
\[
X_{+,0}=X_{-,0}=X_0,
\qquad
X_{+,1}=-X_{-,1}.
\]
Write
\[
X_0=pI+rD,
\qquad
X_{-,1}=BS,
\qquad
Y=yJ,
\]
where
\[
p=\frac12(g^{-1/2}+h^{-1/2}),
\qquad
r=\frac12(g^{-1/2}-h^{-1/2}).
\]
Then the centered one-whitening-slot row
\[
T_7^{(1)}
:=
X_{-,1}YX_{+,0}+X_{-,0}YX_{+,1}
\]
satisfies
\[
T_7^{(1)}
=
-2yB(pD+rI).
\]
Consequently
\[
\pi_{\mathfrak f}T_7^{(1)}
=
-2yBr\,I.
\]
Equivalently,
\[
\pi_{\mathfrak f}T_7^{(1)}
=
-y\frac{q''(h-g)}
{4\pi gh(\sqrt g+\sqrt h)^2}I.
\]
Thus, away from the source walls \(y=0\) and \(q''=0\), the centered
order-seven fixed-sector signal is nonzero exactly off the first
Gram-isotropy wall \(g=h\), or equivalently off \(\Gamma_G=0\).
\end{lemma}

\begin{remark}[Why order seven alone is insufficient on \(g=h\)]
\label{rem:order-seven-insufficient-on-gram-isotropy}
On \(g=h\) one has \(r=0\).  The order-seven row does not disappear in
the ambient \(2\times2\) algebra; rather, it becomes
\[
T_7^{(1)}\big|_{g=h}
=
-2yBpD.
\]
This component is anti-fixed and is killed by
\(\pi_{\mathfrak f}\).  Hence the centered order-seven fixed-sector
argument requires either quantitative separation from \(g=h\), an
ambient \(D\)-sector enlargement, or a higher-order fixed-sector
replacement.  The next replacement is the order-nine row below.
\end{remark}

\begin{lemma}[Second whitening anisotropy on the Gram-isotropy wall]
\label{lem:second-whitening-anisotropy}
Assume \(g=h\), so that
\[
G_0=gI.
\]
Write
\[
G_-(s)=G_0+sG_1+s^2G_2+O(s^3),
\qquad
X_-(s)=G_-(s)^{-1/2}
=
X_0+sX_1+s^2X_2+O(s^3),
\]
with the convention \(G_-(s)=G(m-s/2)\), so
\[
G_1=-\frac12G'(m),
\qquad
G_2=\frac18G''(m).
\]
At a scalar base,
\[
X_2
=
-\frac12g^{-3/2}G_2
+
\frac38g^{-5/2}G_1^2.
\]
At centered parity, \(G_1\in\mathbf C S\), hence
\[
G_1^2\in\mathbf C I.
\]
Therefore
\[
(X_2)_D
=
-\frac12g^{-3/2}(G_2)_D.
\]

Using the centered Gram model
\[
G(t)
=
\frac1\pi
\begin{pmatrix}
q(t) & \frac12q'(t)\\[1mm]
\frac12q'(t) &
\frac13q(t)^3+\frac16q''(t)
\end{pmatrix},
\]
and \(q'(m)=q'''(m)=0\), one has
\[
G''(m)
=
\frac1\pi
\begin{pmatrix}
q'' & 0\\[1mm]
0 & q_0^2q''+\frac16q^{(4)}
\end{pmatrix}.
\]
Thus
\[
(G_2)_D
=
\frac{1}{16\pi}
\left(
(1-q_0^2)q''-\frac16q^{(4)}
\right)
=
\frac{\Gamma_{G,2}}{16\pi}.
\]
Consequently
\[
\boxed{
(X_2)_D
=
-\frac{1}{32\pi}g^{-3/2}\Gamma_{G,2}.
}
\]
\end{lemma}

\begin{proposition}[Order-nine fixed-sector survivor on \(g=h\)]
\label{prop:order-nine-gram-isotropy-survivor}
Assume \(g=h\).  At centered parity,
\[
X_{+,1}=-X_{-,1},
\qquad
X_{+,2}=X_{-,2}.
\]
Let
\[
X_1:=X_{-,1}=BS,
\qquad
Y=yJ,
\qquad
X_2:=X_{-,2}=aI+dD,
\]
with \(d=(X_2)_D\).  Consider the order-nine row
\[
T_9^{(1,2)}
:=
X_{-,1}YX_{+,2}
+
X_{-,2}YX_{+,1}.
\]
Then
\[
T_9^{(1,2)}
=
X_1YX_2-X_2YX_1.
\]
Using \(SJ=-D\) and \(JS=D\), we obtain
\[
T_9^{(1,2)}
=
-2yB(aD+dI).
\]
Therefore
\[
\pi_{\mathfrak f}T_9^{(1,2)}
=
-2yBd\,I.
\]
At \(g=h\),
\[
B=\frac{q''}{8\pi g^{3/2}}
\]
and, by Lemma~\ref{lem:second-whitening-anisotropy},
\[
d
=
-\frac{1}{32\pi}g^{-3/2}\Gamma_{G,2}.
\]
Hence
\[
\boxed{
\pi_{\mathfrak f}T_9^{(1,2)}
=
\frac{yq''\Gamma_{G,2}}{128\pi^2g^3}I.
}
\]
In particular, on the source-good locus \(yq''\ne0\), the order-nine
fixed-sector row survives on \(g=h\) whenever
\[
\Gamma_{G,2}\ne0,
\]
modulo the correction incidence defined below.
\end{proposition}

\begin{lemma}[The second anisotropy wall is not forced by \(g=h\)]
\label{lem:second-anisotropy-proper}
The first Gram-isotropy relation
\[
g=h
\]
is equivalent to
\[
q''=6q_0-2q_0^3.
\]
It does not imply
\[
\Gamma_{G,2}=0.
\]
Indeed, on \(g=h\),
\[
\Gamma_{G,2}
=
(1-q_0^2)(6q_0-2q_0^3)-\frac16q^{(4)}.
\]
Thus \(\Gamma_{G,2}=0\) imposes the additional fourth-jet relation
\[
q^{(4)}
=
6(1-q_0^2)(6q_0-2q_0^3).
\]
This is an independent higher-jet incidence condition.

For a centered paired packet
\[
f_{\beta,\gamma}(t)
=
\frac{a}{a^2+x^2}
+
\frac{b}{b^2+x^2},
\qquad
x=2t-\gamma,
\]
evaluated at \(x=0\), one has
\[
q_0=a^{-1}+b^{-1},
\qquad
q''=-8(a^{-3}+b^{-3}),
\qquad
q^{(4)}=384(a^{-5}+b^{-5}).
\]
Hence
\[
\Gamma_{G,2}
=
(1-q_0^2)q''-\frac16q^{(4)}.
\]
For example, for \(a=\frac14\), \(b=\frac34\),
\[
q_0=\frac{16}{3},
\qquad
q''=-\frac{14336}{27},
\qquad
q^{(4)}=\frac{31981568}{81},
\]
and therefore
\[
\Gamma_{G,2}
=
-\frac{12449792}{243}\ne0.
\]
Thus the order-nine wall \(\Gamma_{G,2}=0\) is a proper additional
analytic wall, not a consequence of centered parity or of \(g=h\).
\end{lemma}

\begin{proposition}[Raw order-nine \(I\)-row audit]
\label{prop:raw-order-nine-row-audit}
Under the centered raw two-point convention, the raw odd carriers are skew
off-diagonal:
\[
N_{2j+1}=F_{2j+1}J.
\]
In particular,
\[
N_9=F_9J,
\qquad
N_7=F_7J,
\qquad
N_5=F_5J,
\qquad
N_3=F_3J,
\qquad
N_1=F_1J.
\]
Consequently their \(I\)- and \(D\)-components vanish:
\[
N_{2j+1}^I=N_{2j+1}^D=0.
\]

The raw even carriers are fixed under endpoint swap.  In the centered raw
two-point convention, their symmetric off-diagonal \(S\)-parts do not
contribute to the order-nine \(I\)-coordinate:
\[
N_4^S=N_2^S=N_0^S=0
\]
for the raw lower-carrier row relevant here.

Therefore the raw order-nine \(I\)-coordinate receives no contribution from
the raw carriers other than the order-nine survivor
\[
T_9^{(1,2)}
=
X_{-,1}YX_{+,2}+X_{-,2}YX_{+,1}.
\]
Equivalently,
\[
R_{\le4}^{I,\mathrm{raw}}=0,
\]
and
\[
R_9^I[N_9]=R_7^I[N_7]=R_5^I[N_5]=0.
\]
\end{proposition}

\begin{remark}[Convention caveat for corrected rows]
\label{rem:raw-versus-corrected-row-caveat}
Proposition~\ref{prop:raw-order-nine-row-audit} concerns the raw
pre-whitening two-point carriers \(N_b\).  If a later notation folds
normalization, subtraction, background/core, line, or connection correction
terms into a symbol also called \(N_b\), those terms should be separated from
the raw carrier and placed into the correction remainder below.  The raw
matrix-parity audit does not by itself eliminate such correction rows.
\end{remark}

\begin{definition}[Order-nine correction incidence]
\label{def:order-nine-correction-incidence}
Let \(R_{\rm corr}^{I,(9)}\) denote the residual order-nine \(I\)-coordinate
coming from correction rows not included in the raw two-point carrier table.
After covariantization and the fixed-convention local gates, this residual may
only include line, background/core, or normalization/subtraction correction
data:
\[
R_{\rm corr}^{I,(9)}
=
R_{\rm line}^{I,(9)}
+
R_{\rm bg/core}^{I,(9)}
+
R_{\rm norm/sub}^{I,(9)}.
\]
Connection/frame/projection/whitening-gauge terms are absorbed by
\(\nabla^{\mathfrak f}\); Gate~A/B rows are closed under the fixed-convention
hypotheses; shadow rows are projected away; and pure same-point odd rows
vanish by centered parity.

For the collapsed scalar \(I\)-coordinate fallback, define
\[
F_9^I
:=
R_{\rm corr}^{I,(9)}
+
\frac{yq''\Gamma_{G,2}}{128\pi^2g^3}.
\]
Thus the order-nine scalar \(I\)-coordinate is
\[
(A_9^{\mathfrak f})_I=F_9^I.
\]
Define the order-nine correction incidence wall by
\[
\mathcal E_{\rm corr}^{(9)}
:=
\left\{
F_9^I=0,
\quad
yq''\Gamma_{G,2}\ne0
\right\}.
\]
Equivalently,
\[
\mathcal E_{\rm corr}^{(9)}
=
\left\{
R_{\rm corr}^{I,(9)}
=
-\frac{yq''\Gamma_{G,2}}{128\pi^2g^3},
\quad
yq''\Gamma_{G,2}\ne0
\right\}.
\]
On the complement of this wall, the collapsed order-nine \(I\)-coordinate is
nonzero whenever \(yq''\Gamma_{G,2}\ne0\).
\end{definition}

\begin{theorem}[Layered fixed-sector visibility through order nine]
\label{thm:layered-fixed-sector-visibility-order-nine}
Work on a centered reduced good patch in the fixed-core convention.  Assume
\[
g>0,\qquad h>0,\qquad yq''\ne0.
\]
Let
\[
\Gamma_G=q''-6q_0+2q_0^3,
\qquad
\Gamma_{G,2}=(1-q_0^2)q''-\frac16q^{(4)}.
\]
Assume also that the point is not in the order-seven cross-channel incidence
wall and not in the order-nine correction incidence wall:
\[
\theta\notin
\mathcal E_{\rm cc}^{(7)}
\cup
\mathcal E_{\rm corr}^{(9)}.
\]
If
\[
\Gamma_G\ne0
\quad\text{or}\quad
\Gamma_{G,2}\ne0,
\]
then
\[
(A_7^{\mathfrak f},A_9^{\mathfrak f})\ne0.
\]

More precisely:
\[
\Gamma_G\ne0
\quad\Longrightarrow\quad
A_7^{\mathfrak f}\ne0,
\]
where the nonzero fixed-sector row is the centered order-seven row
\[
\pi_{\mathfrak f}T_7^{(1)}
=
-2yBr\,I.
\]
On the first Gram-isotropy wall
\[
\Gamma_G=0,
\]
if
\[
\Gamma_{G,2}\ne0
\]
and
\[
\theta\notin\mathcal E_{\rm corr}^{(9)},
\]
then
\[
A_9^{\mathfrak f}\ne0,
\]
with order-nine \(I\)-coordinate
\[
(A_9^{\mathfrak f})_I
=
\frac{yq''\Gamma_{G,2}}{128\pi^2g^3}
+
R_{\rm corr}^{I,(9)}.
\]
\end{theorem}

\begin{proof}
If \(\Gamma_G\ne0\), then \(g-h\ne0\).  By
Lemma~\ref{lem:centered-order-seven-gram-anisotropy},
\[
\pi_{\mathfrak f}T_7^{(1)}
=
-y\frac{q''(h-g)}
{4\pi gh(\sqrt g+\sqrt h)^2}I,
\]
which is nonzero under \(yq''\ne0\), \(g,h>0\), and \(g-h\ne0\), outside the
named order-seven cross-channel incidence.

It remains to consider \(\Gamma_G=0\), equivalently \(g=h\).  By
Proposition~\ref{prop:order-nine-gram-isotropy-survivor},
\[
\pi_{\mathfrak f}T_9^{(1,2)}
=
\frac{yq''\Gamma_{G,2}}{128\pi^2g^3}I.
\]
By Proposition~\ref{prop:raw-order-nine-row-audit}, the raw order-nine
\(I\)-row has no additional contribution from \(N_9,N_7,N_5\), or from the
raw lower carriers.  Thus the only possible remaining \(I\)-coordinate
correction is \(R_{\rm corr}^{I,(9)}\).  Outside
\(\mathcal E_{\rm corr}^{(9)}\), this correction does not cancel the displayed
candidate row.  Hence \(A_9^{\mathfrak f}\ne0\).
\end{proof}

\begin{remark}[What the layered theorem does and does not prove]
\label{rem:layered-visibility-scope}
Theorem~\ref{thm:layered-fixed-sector-visibility-order-nine} is a local
visibility theorem.  It improves the centered order-seven story by replacing
the single excluded wall
\[
\Gamma_G=0
\]
with the layered incidence
\[
\Gamma_G=0,\qquad \Gamma_{G,2}=0,
\]
together with the named correction incidence.  In particular, the order-seven
fixed-sector signal detects off first Gram isotropy, while the order-nine
fixed-sector signal detects on the first Gram-isotropy wall unless the second
Gram anisotropy also vanishes.

This is still not an RH-scale quantitative lower bound.  For a global proof
one would need a non-concentration estimate of the form
\[
|yq''|
\max\{|\Gamma_G|,|\Gamma_{G,2}|\}
\ge Q^{-A},
\]
or an equivalent lower bound for
\[
\|(A_7^{\mathfrak f},A_9^{\mathfrak f})\|
\]
outside the layered incidence and correction walls.  Qualitative
nonvanishing alone does not rule out super-polynomial approach to the layered
incidence set.
\end{remark}

\begin{remark}[No-go ledger for the centered order-seven/nine analysis]
\label{rem:no-go-ledger-order-seven-nine}
The following points should be kept separate.

\begin{enumerate}
\item The centered order-seven fixed-sector row is structurally proportional
to the first Gram anisotropy \(g-h\).  Hence it cannot detect the wall
\(g=h\) inside the fixed sector.

\item The ambient order-seven \(D\)-component survives on \(g=h\), but using it
would require an ambient, non-fixed-sector enlargement.  The present theorem
instead stays within the rank-two fixed sector and uses order nine.

\item The order-nine survivor is not a raw-carrier artifact.  It arises from
the channel
\[
Y_JX_{1,S}X_{2,D},
\]
namely
\[
X_{-,1}YX_{+,2}+X_{-,2}YX_{+,1}.
\]

\item The raw carriers \(N_9,N_7,N_5,N_3,N_1\) are skew off-diagonal at
centered parity:
\[
N_{2j+1}=F_{2j+1}J.
\]
They do not contribute to the order-nine \(I\)-coordinate.  The raw even lower
carriers are likewise harmless for the order-nine \(I\)-row under the centered
raw convention.

\item Any remaining order-nine cancellation is therefore a correction-row or
cross-channel analytic incidence issue, not a sector/parity identity of the
raw row table.
\end{enumerate}
\end{remark}

\begin{definition}[Carrier-graded row module through order nine]
\label{def:carrier-graded-row-module-order-nine}
Let \(\mathcal M_{\rm car}\) be the pre-collapse carrier row module through
order nine, freely generated over the local coefficient ring by
\[
\mathsf C_7,\quad
\mathsf C_{Y12},\quad
\mathsf C_{N7,I},\quad
\mathsf C_{N7,D},\quad
\mathsf C_{N5,I},\quad
\mathsf C_{N5,D},
\]
together with line, background/core, and normalization/subtraction carriers
\[
\mathsf C_{\rm line},\qquad
\mathsf C_{\rm bg/core},\qquad
\mathsf C_{\rm norm/sub}.
\]
Here \(\mathsf C_7\) denotes the order-seven live carrier
\(T_7^{(1)}\), and \(\mathsf C_{Y12}\) denotes the clean centered order-nine
carrier
\[
T_9^{(1,2)}
=
X_{-,1}YX_{+,2}+X_{-,2}YX_{+,1}.
\]
This is the same carrier denoted \(T_9^{(1,2)}\) in
Proposition~\ref{prop:raw-order-nine-row-audit} and
Theorem~\ref{thm:layered-fixed-sector-visibility-order-nine}.

Define the nuisance submodule
\[
\mathcal C_{\rm car}
:=
\operatorname{span}
\{
\mathsf C_{N7,I},\mathsf C_{N7,D},
\mathsf C_{N5,I},\mathsf C_{N5,D},
\mathsf C_{\rm line},\mathsf C_{\rm bg/core},\mathsf C_{\rm norm/sub}
\}.
\]
The carrier quotient is
\[
\mathcal Q_{\rm car}:=\mathcal M_{\rm car}/\mathcal C_{\rm car}.
\]
\end{definition}

\begin{lemma}[Survival of the order-nine clean carrier in the carrier quotient]
\label{lem:order-nine-carrier-survives-quotient}
In \(\mathcal Q_{\rm car}\),
\[
[\mathsf C_{Y12}]\ne0,
\]
and
\[
[\mathsf C_7]\wedge[\mathsf C_{Y12}]\ne0
\qquad
\text{in }\wedge^2\mathcal Q_{\rm car}.
\]
\end{lemma}

\begin{proof}
By definition, \(\mathcal M_{\rm car}\) is free on the displayed carrier
basis.  The nuisance submodule \(\mathcal C_{\rm car}\) is generated by the
\(N_7\), \(N_5\), line, background/core, and normalization/subtraction
carriers.  It does not contain \(\mathsf C_{Y12}\) or \(\mathsf C_7\).
Therefore their quotient classes are nonzero and linearly independent, unless
an additional row relation is imposed identifying \(\mathsf C_{Y12}\) or
\(\mathsf C_7\) with the nuisance carriers.  No such relation is part of the
carrier-graded associated row module.
\end{proof}

\begin{remark}[Carrier freeness is the row-output choice]
\label{rem:carrier-freeness-row-output-choice}
Lemma~\ref{lem:order-nine-carrier-survives-quotient} is a statement in the
pre-collapse carrier module.  It uses the row-output choice that carrier labels
are retained until after quotienting nuisance directions.  If one instead
defines the row output only after applying the scalar collapse map
\[
S_I:\mathcal M_{\rm car}\to\mathcal O\cdot I,
\]
then the freeness assertion is no longer available and the scalar
cross-channel incidence problem for \(F_9^I\) returns.  Thus carrier freeness
is not a hidden scalar non-cancellation theorem; it is the explicit
carrier-quotient row-energy architecture.
\end{remark}

\begin{lemma}[Leading source factor]
\label{lem:leading-source-factor-H}
In the centered reduced-patch normalization,
\[
H:=yq'',
\]
and the leading source term satisfies
\[
H_{\mathrm{lead}}=-\frac{q_0(q'')^3}{288}.
\]
Equivalently, with \(u=q''_{\can}=q''\),
\[
H_{\mathrm{lead}}=-\frac{q_0u^3}{288}.
\]
\end{lemma}

\begin{proof}
The centered source calculation gives
\[
y_{\mathrm{lead}}=-\frac{q_0A_2^2}{72},
\qquad
q''=2A_2.
\]
Hence
\[
H_{\mathrm{lead}}=y_{\mathrm{lead}}q''
=
-\frac{q_0A_2^2}{72}\cdot 2A_2
=
-\frac{q_0A_2^3}{36}.
\]
Since \(A_2=q''/2\),
\[
H_{\mathrm{lead}}
=
-\frac{q_0(q'')^3}{288}.
\]
\end{proof}

\begin{theorem}[D4 carrier-quotient row conversion through order nine]
\label{thm:d4-carrier-quotient-through-order-nine}
This theorem is a quotient row-energy theorem.  It does not assert coercivity
of the collapsed scalar \(I\)-coordinate.  If the final downstream norm is
forced to be the scalar \(I\)-coordinate norm, then one must instead control
the scalar incidence wall in
Definition~\ref{def:order-nine-correction-incidence}, or pass to a higher or
ambient row architecture.

Assume the carrier-quotient row-energy formulation is admissible downstream.
Let
\[
W_{\rm car}:=\operatorname{span}\{[\mathsf C_7],[\mathsf C_{Y12}]\}
\subset \mathcal Q_{\rm car}.
\]
Assume the quotient Gram matrix of the ordered pair
\[
[\mathsf C_7],\qquad [\mathsf C_{Y12}]
\]
has determinant bounded below by \(Q^{-C}\).

Assume that, in \(\mathcal Q_{\rm car}\), the order-seven and clean order-nine
live carrier classes have the leading-symbol form, modulo nuisance directions
and D4 residual terms,
\[
[\mathcal A_7]
=
H\Gamma_G[\mathsf C_7]+\mathcal E_7,
\]
\[
[\mathcal A_9]
=
H\Gamma_{G,2}[\mathsf C_{Y12}]+\mathcal E_9,
\]
and that the quotient projection of the residual into \(W_{\rm car}\) is
half-leading:
\[
\left|
P_{W_{\rm car}}
\left(
\mathcal E_7+\mathcal E_9
\right)
\right|_{\mathcal Q_{\rm car}}
\le
\frac12
|H|
\left(\Gamma_G^2+\Gamma_{G,2}^2\right)^{1/2}.
\]
Then
\[
\left|
[\mathcal A_7]+[\mathcal A_9]
\right|_{\mathcal Q_{\rm car}}
\ge
Q^{-C}
|H|
\left(\Gamma_G^2+\Gamma_{G,2}^2\right)^{1/2}.
\]
Consequently, combined with Gate~2, gauge-locked wall escape, and
Lemma~\ref{lem:row-mobility-point-to-energy}, one obtains carrier-quotient row
energy
\[
\int_I
\left|[\mathcal A_N^{\rm car}(m)]\right|_{\mathcal Q_{\rm car}}^2\,dm
\ge Q^{-C}|I|.
\]
\end{theorem}

\begin{proof}
The Gram determinant lower bound gives a polynomial lower frame bound for the
two quotient carrier directions.  Therefore
\[
\left|
H\Gamma_G[\mathsf C_7]+H\Gamma_{G,2}[\mathsf C_{Y12}]
\right|_{\mathcal Q_{\rm car}}
\ge
Q^{-C}
|H|
\left(\Gamma_G^2+\Gamma_{G,2}^2\right)^{1/2}.
\]
The residual bound is half-leading in the only quotient directions that can
cancel the leading response.  The triangle inequality gives the pointwise row
lower bound.  Row mobility upgrades the pointwise lower bound to row energy.

This proves carrier-quotient row conversion only.  It does not assert
coercivity of any collapsed scalar coordinate such as the \(I\)-coordinate
order-nine readout.
\end{proof}

\begin{remark}[Finite D4 audit ledger]
\label{rem:finite-d4-audit-ledger}
The carrier-quotient theorem separates the structural rank statement from the
finite row-algebra audit.  The finite D4 audit consists of the following
checks.

\begin{enumerate}[label=(\roman*),nosep]
\item Verify the principal source factor
\[
        \operatorname{gr}(H)=-\frac1{288}q_0u^3.
\]

\item Verify the order-seven carrier detecting \(\Gamma_G\):
\[
        [\mathcal A_7]\sim H\Gamma_G[\mathsf C_7].
\]

\item Verify the clean order-nine carrier detecting \(\Gamma_{G,2}\):
\[
        [\mathcal A_9]\sim H\Gamma_{G,2}[\mathsf C_{Y12}].
\]

\item Identify the nuisance carrier subspace being quotiented.

\item Prove quotient coercivity:
\[
        [\mathsf C_7]\wedge[\mathsf C_{Y12}]\ne0
\]
with polynomial quantitative lower bound in the chosen quotient norm.

\item Bound only the residual components that can cancel the leading response:
\[
        \left|
        P_{W_{\rm car}}(\mathcal E_7+\mathcal E_9)
        \right|_{\mathcal Q_{\rm car}}
        \le
        \frac12
        |H|(\Gamma_G^2+\Gamma_{G,2}^2)^{1/2}.
\]
\end{enumerate}

In scalar \(I\)-coordinate output, additional cross-channel incidence walls may
appear after independent carrier directions are collapsed into the same scalar
coordinate.  These are scalar-output artifacts unless the downstream theorem
explicitly requires scalar \(I\)-coordinate energy.
\end{remark}

\subsection{Forbidden singular-discard D4 candidate}
\label{subsec:forbidden-singular-discard-d4}

The best current finite-D4 forbidden operator is the local singular-discard
operator, not the full three-chart atlas loop.

\begin{definition}[Singular retained/coherent space and discard operator]
\label{def:singular-retained-space-and-discard-lambda}
Let \(H_{\mathrm{sing},p}\) be the singular Hilbert/module space in which
\[
  T_{\mathrm{tr}\to\mathrm{sing},p}J_{\mathrm{tr},p}(q''_{\can})
\]
is represented.  Let
\[
  R_{\mathrm{sing},p}\subset H_{\mathrm{sing},p}
\]
be the retained singular coherence space, and let \(P_{\mathrm{sing},p}\) be the
orthogonal projection onto \(R_{\mathrm{sing},p}\).

The intended choice is not the Taylor-jet-complete singular space.  A
Taylor-jet-complete retained space transports formal D4 carriers functorially
and gives rank zero.  The intended choice is instead an actual-source
EF-coherent singular space generated by the localized finite canonical model
\[
  F_X=B_2+P_{2,\le X}+D_{2,\le X},
\]
together with mandatory gamma/background/coherent directions.

Define
\[
  \Lambda_{\mathrm{tr},\mathrm{sing},p}
  :=
  (I-P_{\mathrm{sing},p})T_{\mathrm{tr}\to\mathrm{sing},p}P_{\mathrm{tr},p}.
\]
\end{definition}

\begin{definition}[Combined nuisance block and projected transition operator]
\label{def:combined-nuisance-block-y12}
For each packet \(p\), define the combined \(Y12\) nuisance block
\[
U_p := \operatorname{span}\!\bigl(T_{\mathrm{par}}^{(\alpha)},\,T_7,\,T_{<7},\,II_{\mathrm{arith},9},\,T_{E_2}\bigr),
\]
with the understanding that \(T_{E_2}\) is included only when \(E_2\) is admitted
in the current quotient layer.

Define the projected transition operator
\[
\mathfrak T_p := (I-P_{\mathrm{sing},p})\,T_{\mathrm{tr}\to\mathrm{sing},p}\,J_{\mathrm{tr},p}.
\]
\end{definition}

\begin{lemma}[Finite-codimension \(Y12\) graph/coercivity lemma]
\label{lem:finite-codim-y12-graph}
Assume that the \(Y12\) quotient annihilates the combined nuisance block,
\[
U_p \subseteq \ker \Lambda_{Y,p},
\]
and that the admissible residual class \(\mathcal R_X\) splits as a finite-codimension graph space
\[
\mathcal R_X = \mathcal R_{X,\mathrm{low},p}\oplus \mathcal R_{X,\mathrm{tail},p},
\qquad
\dim \mathcal R_{X,\mathrm{low},p}<\infty.
\]
Let \(\{e_{1,p},\dots,e_{r,p}\}\) be a basis of \(\mathcal R_{X,\mathrm{low},p}\), and define the defect matrix
\[
D_p := \begin{bmatrix}
\mathfrak T_p e_{1,p} & \cdots & \mathfrak T_p e_{r,p}
\end{bmatrix}.
\]
If
\[
\sigma_{\min}(D_p)\ge Q^{-C}
\qquad\text{and}\qquad
\sum_p \|\mathfrak T_p(R_X)\|^2 \le Q^{-B}\mu(\mathcal C),
\]
then \(\mathfrak T_p\) is coercive on \(\mathcal R_X\) up to the tail norm, and the
\(Y12\) discarded observable survives on a polynomial-weight good set.
\end{lemma}

\begin{remark}[wip-zeta-bridge-proof-y12-graph]
\label{rem:wip-zeta-bridge-proof-y12-graph}
What remains is the finite-codimension graph/coercivity statement for
\(\mathfrak T_p=(I-P_{\mathrm{sing},p})T_{\mathrm{tr}\to\mathrm{sing},p}J_{\mathrm{tr},p}\):
a finite low-mode defect matrix and a controlled singular-discard tail.
\end{remark}

\begin{remark}[Defect matrix and \(D4\) transversality]
\label{rem:defect-matrix-d4-transversality}
The finite defect matrix \(D_p\) is the low-mode representation of the same
discarded operator encoded by
\[
K_{\mathrm{D4},p}
=
(I-P_{\mathrm{sing},p})T_{\mathrm{tr}\to\mathrm{sing},p}E_{\mathrm{car},p}^{\mathrm{tr}}L_{\mathrm{clean},p}^{\mathrm{tr}}.
\]
Thus the coercivity of \(\mathfrak T_p\) on the low block and the rank condition
for \(K_{\mathrm{D4},p}\) are the same finite test in two coordinate systems.
\end{remark}

\begin{remark}[Graph-space no-go]
\label{rem:graph-space-no-go}
If \(U_p\not\subseteq\ker\Lambda_{Y,p}\), or if \(\mathcal R_X\) fails to admit
a finite-codimension graph decomposition with tail norm
\(\|\cdot\|_{\mathrm{EF}}\), then the \(Y12\) discarded observable is not
coercive and the downstream atomic contradiction cannot be derived in the
current norm.
\end{remark}

\begin{hypothesis}[Singular-discard upper from EF coherence]
\label{hyp:singular-discard-upper-from-EF-coherence}
The actual canonical rows satisfy
\[
  \boxed{
  \sum_p
  \left\|
  (I-P_{\mathrm{sing},p})T_{\mathrm{tr}\to\mathrm{sing},p}J_{\mathrm{tr},p}(q''_{\can})
  \right\|^2
  \le
  Q^{-B}\mu(\mathcal C).
  }
\]
A sufficient proof route is:
\[
  q''_{\can}=F_X+R_X,
  \qquad
  T_{\mathrm{tr}\to\mathrm{sing},p}J_{\mathrm{tr},p}(F_X)\in R_{\mathrm{sing},p},
\]
and
\[
  \sum_p
  \|T_{\mathrm{tr}\to\mathrm{sing},p}J_{\mathrm{tr},p}(R_X)\|^2
  \le
  Q^{-B}\mu(\mathcal C).
\]
Thus the same EF coherence estimates used in
Hypothesis~\ref{hyp:agg-l2-qcan-gate-one} may also supply the forbidden upper,
provided they hold in the singular overlap norms.
\end{hypothesis}

\begin{definition}[\(\mathrm{EFCompletion}_{K9}\)]
\label{def:EFCompletion-K9}
The local \(L^2\) completion used for Gate--1 is not automatically strong
enough for the singular-discard forbidden upper.  Define
\(\mathrm{EFCompletion}_{K9}\) to be the estimate
\[
  \boxed{
  \sum_p
  \left\|
  \Lambda_{\mathrm{tr},\mathrm{sing},p}J_{\mathrm{tr},p}(R_X)
  \right\|^2
  \le Q^{-B}\mu(\mathcal C),
  }
\]
Here
\[
  \Lambda_{\mathrm{tr},\mathrm{sing},p}
  =
  (I-P_{\mathrm{sing},p})T_{\mathrm{tr}\to\mathrm{sing},p}P_{\mathrm{tr},p}
\]
already includes the transition-to-singular transport.  Equivalently, for a
normal residual functional \(K_{9,p}\), the same estimate may be written as
\[
  \sum_p
  \left|
    K_{9,p}T_{\mathrm{tr}\to\mathrm{sing},p}J_{\mathrm{tr},p}(R_X)
  \right|^2
  \le Q^{-B}\mu(\mathcal C).
\]
Here
\[
  q''_{\can}=F_X+R_X,
  \qquad
  F_X=B_2+P_{2,\le X}+D_{2,\le X}.
\]
Equivalently, the residual must be small not only in the Gate--1
\(L^2\)-window norm, but also after transport to the singular chart and
application of the discarded singular projection.
\end{definition}

\begin{definition}[\(\mathrm{EFCompletion}_{Y12}\)]
\label{def:EFCompletion-Y12}
The targeted \(Y12\)-signature completion estimate is
\[
  \boxed{
  \sum_p
  \left|
    \Pi_{\Sigma(Y12)}
    T_{\mathrm{tr}\to\mathrm{sing},p}J_{\mathrm{tr},p}(R_X)
  \right|^2
  \le
  Q^{-B}\mu(\mathcal C).
  }
\]
This is weaker than \(\mathrm{EFCompletion}_{K9}\), which controls the full
order-nine residual after applying \(\mathcal K_9\).  It is sufficient for the
\(Y12\)-signature D4 contradiction if
Hypothesis~\ref{hyp:second-order-Y12-signature-separation} holds.
\end{definition}

\begin{remark}[Three completion ledgers]
\label{rem:three-EF-completion-ledgers}
The proof now uses three different completion strengths:
\[
\boxed{\mathrm{EFCompletion}_{L^2}}
\]
for aggregate Gate--1 source energy,
\[
\boxed{\mathrm{EFCompletion}_{K9}}
\]
for a full order-nine residual upper, and
\[
\boxed{\mathrm{EFCompletion}_{Y12}}
\]
for the targeted primitive \(Y12\)-signature channel.  These must not be
conflated.  In particular,
\[
  \mathrm{EFCompletion}_{L^2}\not\Rightarrow
  \mathrm{EFCompletion}_{K9}
\]
unless the operator
\[
  \mathcal K_9T_{\mathrm{tr}\to\mathrm{sing}}J_{\mathrm{tr}}
\]
is polynomially bounded from the \(L^2\)-completion norm.  The targeted
\(Y12\)-estimate may be substantially easier than full \(K_9\)-completion.
\end{remark}

\begin{remark}[\(\mathrm{EFCompletion}_{L^2}\) does not automatically imply
\(\mathrm{EFCompletion}_{K9}\)]
\label{rem:EFCompletion-L2-not-K9}
The implication
\[
  \mathrm{EFCompletion}_{L^2}
  \Longrightarrow
  \mathrm{EFCompletion}_{K9}
\]
requires a polynomial operator bound for
\[
  \Lambda_{\mathrm{tr},\mathrm{sing}}J_{\mathrm{tr}}
\]
from the \(L^2\)-completion norm to the discarded singular norm.  If such a
bound is unavailable, the singular-discard upper remains open even when
aggregate \(L^2\) Gate--1 is well-defined.  This is a separate completion
ledger and should not be merged with the source-energy Gate--1 completion.
\end{remark}

\begin{definition}[D4 singular-discard transversality matrix]
\label{def:d4-singular-discard-transversality-matrix}
On the D4 live plane define
\[
  K_{\mathrm{D4},p}
  :=
  (I-P_{\mathrm{sing},p})T_{\mathrm{tr}\to\mathrm{sing},p}
  E_{\mathrm{car},p}^{\mathrm{tr}}L_{\mathrm{clean},p}^{\mathrm{tr}}.
\]
Since \(L_{\mathrm{clean},p}^{\mathrm{tr}}\) is invertible on the live plane, the decisive
rank test is equivalent to the rank of
\[
  (I-P_{\mathrm{sing},p})T_{\mathrm{tr}\to\mathrm{sing},p}
  \Big|_{\operatorname{span}\{[\mathsf C_7]^{\mathrm{tr}},[\mathsf C_{Y12}]^{\mathrm{tr}}\}}.
\]
Write
\[
  v_7^\perp
  :=
  (I-P_{\mathrm{sing}})T_{\mathrm{tr}\to\mathrm{sing}}[\mathsf C_7]^{\mathrm{tr}},
  \qquad
  v_Y^\perp
  :=
  (I-P_{\mathrm{sing}})T_{\mathrm{tr}\to\mathrm{sing}}[\mathsf C_{Y12}]^{\mathrm{tr}}.
\]
Then the rank is \(0\) if both vanish, \(1\) if they are nonzero and dependent,
and \(2\) if they are independent.
\end{definition}

\begin{definition}[Singular order-nine connection residual]
\label{def:singular-order-nine-K9-residual}
Let \(s=(s_0,\ldots,s_9)\) denote the singular coefficient vector in the
transition-to-singular overlap chart.  Suppose the actual EF-coherent singular
source satisfies an order-nine connection relation of the form
\[
  s_9
  =
  \mathcal R_{7\to9}s_7
  +
  \mathcal R_{<7\to9}(s_{<7})
  +
  \operatorname{Tail}_9 .
\]
Define the singular order-nine residual
\[
  \mathcal K_9(s)
  :=
  s_9-\mathcal R_{7\to9}s_7-\mathcal R_{<7\to9}(s_{<7}).
\]
The associated forbidden operator is
\[
  \Lambda_{K9,p}
  :=
  \mathcal K_9\circ T_{\mathrm{tr}\to\mathrm{sing},p}P_{\mathrm{tr},p}.
\]
It is a valid singular-discard channel only if actual EF-coherent rows satisfy
\[
  \sum_p
  \left|
    \mathcal K_9\bigl(T_{\mathrm{tr}\to\mathrm{sing},p}J_{\mathrm{tr},p}(q''_{\can})\bigr)
  \right|^2
  \le Q^{-B}\mu(\mathcal C),
\]
and if the D4 live plane has nonzero normal image:
\[
  \operatorname{rank}
  \left(
    \mathcal K_9
    T_{\mathrm{tr}\to\mathrm{sing},p}
    E_{\mathrm{car},p}^{\mathrm{tr}}
    L_{\mathrm{clean},p}^{\mathrm{tr}}
  \right)\ge1
\]
on a polynomial-weight set of packets.
\end{definition}

\begin{definition}[Lower-order tangent observability for \(K_9\)]
\label{def:lower-order-tangent-observability-K9}
Let the actual-source tangent matrix in the singular order-nine chart be
written in block form
\[
  T=
  \begin{pmatrix}
    T_L\\
    T_9
  \end{pmatrix},
\]
where \(T_L\) records the lower-order coordinates and \(T_9\) records the
order-nine coordinate.  The residual
\[
  K_9(s)=s_9-\mathcal R_{7\to9}s_7-\mathcal R_{<7\to9}(s_{<7})
\]
is polynomially conditioned if the lower-order block observes every
actual-source tangent:
\[
  \boxed{
  |T_La|\ge Q^{-A}|T_9a|
  \qquad\text{for every actual-source tangent }a.
  }
\]
Equivalently, the connection operator
\[
  \mathcal R:=T_9T_L^\dagger
\]
exists on the relevant tangent image with
\[
  \|\mathcal R\|\le Q^A.
\]
\end{definition}

\begin{remark}[Why \(K_9\) conditioning is separate from \(Y12\)-nonvanishing]
\label{rem:K9-conditioning-separate-from-Y12-nonvanishing}
The condition
\[
  K_9(T_{\mathrm{tr}\to\mathrm{sing}}[\mathsf C_{Y12}]^{\mathrm{tr}})\ne0
\]
is a signature nonvanishing test.  It does not by itself ensure that the normal
functional \(K_9\) is quantitatively usable.  If an actual-source tangent is
almost purely order nine, then \(K_9\) may have super-polynomial norm after
quotienting lower-order directions.  Thus the \(K_9\)-route requires both:
\[
  \boxed{K_9(T_{\mathrm{tr}\to\mathrm{sing}}[\mathsf C_{Y12}]^{\mathrm{tr}})\ne0}
  \qquad\text{and}\qquad
  \boxed{\mathrm{LowerOrderTangentObservability}_{K9}}.
\]
\end{remark}

\begin{remark}[The \(Y12\) channel is the preferred \(K_9\) test]
\label{rem:Y12-preferred-K9-test}
For the rank-two D4 output
\[
  H
  \begin{pmatrix}
    c_7\Gamma_G\\
    \beta_9^{\mathrm{tri}}\Gamma_G+\alpha_9\Gamma_{G,2}
  \end{pmatrix},
\]
the cleanest singular-discard test is
\[
  \mathcal K_9\!\left(
    T_{\mathrm{tr}\to\mathrm{sing}}[\mathsf C_{Y12}]^{\mathrm{tr}}
  \right)\neq0.
\]
If this holds, then the pure \(\Gamma_{G,2}\) channel survives in the discarded
singular normal direction, because \(\alpha_9\neq0\) on D4 coefficient-good
charts.  If it fails, one should next test the \(\Gamma_G\)-channel through
\[
  \mathcal K_9\!\left(
    T_{\mathrm{tr}\to\mathrm{sing}}[\mathsf C_7]^{\mathrm{tr}}
  \right).
\]
If both tests vanish, the order-nine singular residual \(\mathcal K_9\) is
rank-zero for finite D4 and cannot serve as the downstream forbidden channel.
\end{remark}

\begin{remark}[\(E_2^{Y12}\) honesty audit]
\label{rem:E2-Y12-honesty-audit}
Before claiming that the \(Y12\)-normal direction is forbidden, one must audit
the primitive \(Y12\)-signature of the actual completion/curvature term:
\[
  \Pi_{Y12}E_2.
\]
There are three cases.

\[
\boxed{\text{Case A: }\Pi_{Y12}E_2=0\text{ or }Q^{-B}\text{-small}.}
\]
Then \(Y12\) is a strong forbidden-channel candidate, and the actual-row upper
requires only
\[
  \Pi_{Y12}T_{\mathrm{tr}\to\mathrm{sing}}J_{\mathrm{tr}}(R_X)
\]
to be small.

\[
\boxed{\text{Case B: }\Pi_{Y12}E_2\ne0\text{ but fixed/coherent}.}
\]
Then this piece must be included in the retained actual-source model or in
\(R_{\mathrm{sing}}\).  The audit becomes whether its admissible variations span
the same primitive \(Y12\) direction as the D4 carrier.

\[
\boxed{\text{Case C: }\Pi_{Y12}E_2\text{ supplies a free primitive }Y12
\text{ amplitude}.}
\]
Then the \(Y12\)-normal D4 channel is tangent/coherent and the rank-one
\(Y12\)-forbidden route is dead.

Thus the finite audit is not merely
\[
  K_9(T_{\mathrm{tr}\to\mathrm{sing}}[\mathsf C_{Y12}])\ne0,
\]
but also whether actual \(E_2\)-coherence already contains that same primitive
direction.
\end{remark}

\begin{remark}[Coherent \(E_2^Y\) and the final \(Y12\)-normal wall]
\label{rem:coherent-E2Y-final-Y12-normal-wall}
The \(Y12\)-normal D4 channel is safe only after the \(E_2^Y\) component is
classified.  There are three cases.

First, \(E_2^Y\) may be absent, fixed, or residual-controlled.  Then it does
not modify the local forbidden wall.

Second, \(E_2^Y\) may be coherent and graph-like.  In this case it modifies the
fiber graph slope
\[
  \rho_Y\longmapsto \rho_Y^{\mathrm{final}},
\]
and hence modifies the rank-one wall ratio to
\[
  \rho_\Gamma^{\mathrm{final}}
  =
  \frac{\beta_9^{\mathrm{tri}}-\rho_Y^{\mathrm{final}}c_7}{\alpha_9}.
\]
This is safe only if \(\rho_Y^{\mathrm{final}}\) remains low-order/coherent and
does not depend on the actual \(\Gamma_{G,2}\)-wall value.

Third, \(E_2^Y\) may generate an independent finite-dimensional coherent
amplitude space
\[
  \mathcal C_Y^{E_2}.
\]
Then the actual forbidden wall is the Schur projection
\[
  F_Y^{\mathrm{eff}}
  =
  (I-P_{\mathcal C_Y^{E_2}}^{\mathrm{bd}})
  \left[
    H\bigl(
      (\beta_9^{\mathrm{tri}}-\rho_Yc_7)\Gamma_G
      +
      \alpha_9\Gamma_{G,2}
    \bigr)
  \right].
\]
The branch-charge bandwidth \(d_0\) used in finite/O2/O3 must be computed from
\(F_Y^{\mathrm{eff}}\), not from the unprojected wall.  If \(E_2^Y\) is locally
free at extraction scale, then \(Y12\) is actual-source coherent and the
normal-D4 channel is killed.
\end{remark}

\begin{hypothesis}[AmplitudeSchur\(_{Y12}\)]
\label{hyp:amplitude-schur-Y12}
Let \(\mathcal C_Y^{\mathrm{arith}}\) be the finite coherent \(Y12\)-amplitude
space remaining after exact nuisance and residual routing.  Under the
bounded-degree Stack weight \(d\nu_{\mathrm{bd}}\), the effective \(Y12\)-wall
satisfies
\[
  \operatorname{dist}_{L^2(d\nu_{\mathrm{bd}})}
  \left(
    H\bigl(
      (\beta_9^{\mathrm{tri}}-\rho_Yc_7)\Gamma_G
      +
      \alpha_9\Gamma_{G,2}
    \bigr),
    \mathcal C_Y^{\mathrm{arith}}
  \right)^2
  \ge
  Q^{-A}
  \left\|
    H\bigl(
      (\beta_9^{\mathrm{tri}}-\rho_Yc_7)\Gamma_G
      +
      \alpha_9\Gamma_{G,2}
    \bigr)
  \right\|_{L^2(d\nu_{\mathrm{bd}})}^2.
\]
If this fails, the coherent amplitude model absorbs the \(Y12\)-normal wall.
\end{hypothesis}

\begin{definition}[Order-nine \(Y12\) normal quotient]
\label{def:Y12-normal-quotient-Q9-forb}
Let \(Y_9\) denote the order-nine singular carrier space containing the
transitioned primitive \(Y12\)-carrier.  Define the forbidden order-nine
quotient
\[
  \mathcal Q_9^{\mathrm{forb}}
  :=
  Y_9/
  \left(
    T_{\mathrm{par}}
    +
    T_7
    +
    T_{<7}
  \right),
\]
where \(T_{\mathrm{par}}\) denotes actual parameter tangents, \(T_7\) denotes
actual or lower order-seven connection tangents, and \(T_{<7}\) denotes lower
typed dead directions.  The primitive \(Y12\)-normal D4 channel survives exactly
when
\[
  [\mathsf C_{Y12}]\ne0
  \qquad\text{in}\qquad
  \mathcal Q_9^{\mathrm{forb}}.
\]
\end{definition}

\begin{remark}[Relation between singular-discard and \(K_9/Y12\)]
\label{rem:singular-discard-vs-K9-Y12}
The \(K_9/Y12\) channel is a specialization of the singular-discard candidate.
The finite test
\[
  (I-P_{\mathrm{sing}})T_{\mathrm{tr}\to\mathrm{sing}}[\mathsf C_{Y12}]^{\mathrm{tr}}\ne0
\]
is equivalent to the existence of some normal \(K_9\) detecting the transitioned
\(Y12\)-class.  The quotient \(\mathcal Q_9^{\mathrm{forb}}\) records precisely
which actual-source and lower connection directions must be divided out before
\(Y12\) is considered a genuine forbidden normal carrier.

If \([\mathsf C_{Y12}]=0\) in \(\mathcal Q_9^{\mathrm{forb}}\), then the
rank-one \(Y12\)-normal route is dead, even if D4 remains locally visible in a
larger carrier quotient.  If \([\mathsf C_{Y12}]\ne0\), then finite D4 has a
concrete atomic forbidden channel.
\end{remark}

\begin{hypothesis}[Moving \(Y12\)-wall anti-concentration]
\label{hyp:moving-Y12-wall-anti-concentration}
Let
\[
  \ell_p(\Gamma):=a_p\Gamma_G+b_p\Gamma_{G,2}
\]
be the moving linear form obtained from the \(Y12\)-signature D4 channel, with
\[
  |b_p|\ge Q^{-A}
\]
on the extracted set.  Under the post-Gate/O2/O3 measure \(\nu\), assume
\[
  \boxed{
  \int |\ell_p(\Gamma)|^2\,d\nu
  \ge
  Q^{-A}
  \int
  \left(|\Gamma_G|^2+|\Gamma_{G,2}|^2\right)\,d\nu.
  }
\]
Equivalently, the post-Gate/O2/O3 measure is not concentrated on the moving
sublevel set
\[
  \{p:|a_p\Gamma_G+b_p\Gamma_{G,2}|\le Q^{-A}|\Gamma|\}.
\]
In the normalized case \(a_p=0\), this reduces to anti-concentration of
\(\Gamma_{G,2}\).
\end{hypothesis}

\begin{lemma}[Extraction covariance dichotomy]
\label{lem:extraction-covariance-dichotomy}
On each frozen cap \(C_j\), let
\[
B_j=\sum_{p\in C_j} w_p\,\Gamma_p\Gamma_p^\ast,
\qquad
w_p=\rho_p\frac{E_{1,p}H_p}{N_p},
\qquad
\Gamma_p=\binom{\Gamma_G(p)}{\Gamma_{G,2}(p)}.
\]
Then either
\[
\det(B_j)\ge Q^{-A}(\operatorname{tr}B_j)^2,
\]
equivalently
\[
\sum_{p\in C_j}\lambda_p|r_p-\bar r_j|^2\ge Q^{-A},
\]
where \(r_p=\Gamma_{G,2}(p)/\Gamma_G(p)\) on the chart
\(|\Gamma_G|\ge \tfrac12|\Gamma|\), or else \(B_j\) is nearly rank one and the
wall is effectively static.  In the static-line case, one must hand the cap to
the \(Y12\) quotient test:
\[
\eta_j:=|(I-P_{\mathcal K_j^Y})u_j|^2\ge Q^{-A}.
\]
\end{lemma}

\begin{remark}[wip-zeta-bridge-proof-extraction]
\label{rem:wip-zeta-bridge-proof-extraction}
What remains is the \(2\times2\) covariance test on each frozen cap:
\[
\det(B_j)\ge Q^{-A}(\operatorname{tr}B_j)^2,
\]
or else the cap is static and must be handed to the \(Y12\) quotient module.
\end{remark}

\begin{theorem}[Normal D4 through the \(Y12\)-signature \(K_9\) channel]
\label{thm:normal-D4-K9-Y12}
Assume the following.

\begin{enumerate}
\item \(\mathrm{EFCompletion}_{Y12}\) holds:
\[
  \sum_p
  \left|
    \Pi_{\Sigma(Y12)}
    T_{\mathrm{tr}\to\mathrm{sing},p}J_{\mathrm{tr},p}(R_X)
  \right|^2
  \le
  Q^{-B}\mu(\mathcal C).
\]

\item Lower-order tangent observability holds, so the covariant residual
\(K_9\) is polynomially conditioned.

\item The \(Y12\)-signature separation
Hypothesis~\ref{hyp:second-order-Y12-signature-separation} holds.

\item The singular actual-source manifold has polynomial tubular control as in
Definition~\ref{def:tubular-normality-singular-manifold}.

\item On a polynomial-weight extracted set, the D4 coefficient package is good:
\[
  |H|,\ |\alpha_9|,\ |\alpha_9^{-1}|,\ |c_7|,\ |c_7^{-1}|,
  |\beta_9^{\mathrm{tri}}|\le Q^A,
\]
with \(|H|\ge Q^{-A}\).

\item The moving \(Y12\)-wall anti-concentration estimate
Hypothesis~\ref{hyp:moving-Y12-wall-anti-concentration} holds.
\end{enumerate}

Then the D4 displacement has a forbidden \(Y12\)-signature component.  More
precisely,
\[
  \Pi_{\Sigma(Y12)}
  T_{\mathrm{tr}\to\mathrm{sing}}
  E_{\mathrm{car}}^{\mathrm{tr}}L_{\mathrm{clean}}^{\mathrm{tr}}\Gamma
  =
  H\,(a\,\Gamma_G+b\,\Gamma_{G,2})
\]
with
\[
  \boxed{b\ne0}
\]
on the extracted set, up to \(Q^{-A}\)-controlled normalization.  Hence the
forbidden wall
\[
  \boxed{
  E_{\mathrm{wall}}^{Y12}
  :=
  \frac{|H|^2|a\,\Gamma_G+b\,\Gamma_{G,2}|^2}
       {\mathcal N_{\mathrm{D4}}}
  }
\]
has polynomial mass.  Summing over packets contradicts the
\(\mathrm{EFCompletion}_{Y12}\) upper, after the tubular quadratic error is
absorbed.
\end{theorem}

\begin{proof}[Proof sketch]
The \(Y12\)-signature separation removes actual-source tangents, lower
connection residuals, and actual second-curvature effects from the primitive
\(Y12\)-signature quotient.  Therefore
\[
  \Pi_{\Sigma(Y12)}
  T_{\mathrm{tr}\to\mathrm{sing}}[\mathsf C_{Y12}]
\]
is nonzero.  Applying the triangular D4 matrix gives
\[
  H\left(
  \beta_9^{\mathrm{tri}}\Gamma_G+\alpha_9\Gamma_{G,2}
  \right)
  \Pi_{\Sigma(Y12)}T_{\mathrm{tr}\to\mathrm{sing}}[\mathsf C_{Y12}],
\]
so the \(\Gamma_{G,2}\)-coefficient is nonzero because \(\alpha_9\ne0\).
Tubular control turns the linear normal component into an actual distance from
the singular actual-source manifold.  The \(Y12\)-completion upper bounds the
same normal component for actual rows.  The moving-wall anti-concentration
gives positive mass for the D4 normal component, producing the contradiction.
\end{proof}

\begin{hypothesis}[ForbiddenSingularDiscardD4]
\label{hyp:forbidden-singular-discard-d4}
On a polynomial-weight good set,
\[
  \boxed{
  \operatorname{rank}K_{\mathrm{D4},p}\ge1.
  }
\]
When the rank is nonzero, define the forbidden wall energy
\[
  E_{\mathrm{wall}}^{\mathrm{forb}}(p)
  :=
  \frac{
    |H(p)|^2
    \left\|
      K_{\mathrm{D4},p}\Gamma(p)
    \right\|^2
  }{
    \mathcal N_{\mathrm{D4}}(p)
  },
  \qquad
  \Gamma=(\Gamma_G,\Gamma_{G,2})^T.
\]
Rank \(1\) gives one forbidden linear form in
\((\Gamma_G,\Gamma_{G,2})\).  Rank \(2\) gives full raw
\(\Gamma\)-anisotropy.  Rank \(0\) means this \(\Lambda\) is D4-invisible.
\end{hypothesis}

\begin{remark}[Failure routes for the singular-discard channel]
\label{rem:singular-discard-failure-routes}
The singular-discard D4 route fails in exactly the following ways:
\begin{enumerate}
\item \(R_{\mathrm{sing}}\) is Taylor-jet-complete or otherwise too large, so
\[
  T_{\mathrm{tr}\to\mathrm{sing}}
  \operatorname{span}\{[\mathsf C_7],[\mathsf C_{Y12}]\}
  \subset R_{\mathrm{sing}}.
\]
Then the rank is zero.
\item The EF-coherent upper fails:
\[
  (I-P_{\mathrm{sing}})T_{\mathrm{tr}\to\mathrm{sing}}J_{\mathrm{tr}}(q''_{\can})
\]
is not small for actual rows.
\item Transition stability fails:
\[
  \|T_{\mathrm{tr}\to\mathrm{sing}}\|>Q^C
\]
in the singular overlap norms.
\item The D4 live plane is not coefficient-good, so \(E_{\mathrm{wall}}^{\mathrm{forb}}\) loses
polynomial conditioning.
\end{enumerate}
If none of these failures occurs, finite D4 has a genuine Type-I positive
atomic forbidden channel.
\end{remark}

\begin{definition}[Tubular normality for the singular actual-source manifold]
\label{def:tubular-normality-singular-manifold}
Let \(\mathcal M_{\mathrm{sing}}^{\mathrm{arith}}\) be the finite-dimensional
actual-source singular coherence manifold in the retained singular chart.  We
say it has polynomial tubular control on a packet family if its normal
projection is defined on a \(Q^{-C}\)-tube and the second fundamental form
satisfies
\[
  \boxed{
  \|\mathrm{II}_{\mathcal M_{\mathrm{sing}}^{\mathrm{arith}}}\|\le Q^C.
  }
\]
Equivalently, for every tangent point \(u\in\mathcal M_{\mathrm{sing}}^{\mathrm{arith}}\)
and small displacement \(D\),
\[
  \operatorname{dist}(u+D,\mathcal M_{\mathrm{sing}}^{\mathrm{arith}})
  \ge
  \|\Pi_u^\perp D\|
  -
  Q^C\|D\|^2
  -
  \mathrm{Err}_{\mathrm{chart}}.
\]
\end{definition}

\begin{remark}[Why tubular control is needed]
\label{rem:why-tubular-control-needed}
A nominally normal D4 displacement can be absorbed by a nearby nonlinear
actual-source parameter shift if the singular actual-source manifold has large
curvature at the relevant scale.  Therefore the singular-discard contradiction
requires both a linear transversality audit
\[
  \operatorname{rank}
  \left(
  \Lambda_{\mathrm{tr},\mathrm{sing}}E_{\mathrm{car}}^{\mathrm{tr}}L_{\mathrm{clean}}^{\mathrm{tr}}
  \right)\ge1
\]
and polynomial tubular control for
\(\mathcal M_{\mathrm{sing}}^{\mathrm{arith}}\).  The latter is finite-geometric,
not an arithmetic Gate--1 statement.
\end{remark}

\begin{remark}[Partial explicit-formula residuals are upstream completion tools]
\label{rem:partial-EF-residual-role-separation}
The full canonical explicit-formula identity
\[
  q''_{\can}=B_2+P_2+D_2+E_2
\]
is definitional for the canonical source.  Its full residual
\[
  q''_{\can}-(B_2+P_2+D_2+E_2)
\]
is therefore expected to be D4-invisible if all terms are varied coherently.  It
should not be used as the downstream forbidden channel \(\Lambda\).

Partial residuals such as
\[
  R_X
  :=
  q''_{\can}-(B_2+P_{2,\le X}+D_{2,\le X})
\]
are useful upstream for source completion and for proving aggregate
\(L^2\)-observability of \(q''_{\can}\).  A typical use is:
\[
  q''_{\can}=F_X+R_X,\qquad
  F_X:=B_2+P_{2,\le X}+D_{2,\le X},
\]
with
\[
  \|R_X\|_{L^2(W)}\le \frac12\|F_X\|_{L^2(W)}
\]
on the weighted branch window \(W\).  Then the full canonical source inherits
the finite approximant's \(L^2\)-mass.

A partial residual becomes a downstream forbidden channel only if both of the
following are proved:
\[
  \|\Lambda_{X,\psi}V_{\mathrm{glob}}\|^2\le Q^{-B}\|\psi\|^2
\]
for localized cutoffs \(\psi\), and
\[
  \operatorname{rank}
  \left(
    \Lambda_{X,\psi}E_{\mathrm{car}}L_{\mathrm{clean}}
  \right)\ge1
\]
on the D4 live plane.  In the absence of this upper/rank pair, partial EF
residuals are completion-coherence inputs, not the proof-closing \(\Lambda\).
\end{remark}

\begin{remark}[Final D4 residual status in the carrier quotient]
\label{rem:agent-six-final-D4-residual-status-carrier-quotient}
The D4 residual statement is a carrier-quotient row-energy statement for the
canonical scalar source
\[
u=q''_{\can}.
\]
It is not a scalar \(I\)-coordinate coercivity statement.

In the carrier quotient, the D4 row has the effective form
\[
P_{W_{\rm car}}([\mathcal A_7]+[\mathcal A_9])
=
H\,E_{\rm car}L_{\rm eff}
\begin{pmatrix}
\Gamma_G\\
\Gamma_{G,2}
\end{pmatrix}
+
\varepsilon_{\rm rem}.
\]
The clean row-normal matrix is
\[
L_{\rm clean}^{\mathrm{tri}}
=
\begin{pmatrix}
c_7&0\\
\beta_9^{\mathrm{tri}}&\alpha_9
\end{pmatrix},
\qquad
\beta_9^{\mathrm{tri}}:=\beta_9-\omega_Y(e_7)c_7,
\]
with \(c_7,\alpha_9\) as in
Lemma~\ref{lem:clean-triangular-d4-carrier-matrix}.  The same-scale
\(\beta_9^{\mathrm{tri}}\Gamma_G\) term is part of \(L_{\rm eff}\), not a
residual.  All D4 coefficient-good and rank-one fallback statements in this
branch use the triangularized coefficient \(\beta_9^{\mathrm{tri}}\), not the
raw \(\beta_9\).

The remaining correction-row issue is purely carrier-theoretic.  If, before
scalar collapse, the order-seven and order-nine line, background/core, and
normalization/subtraction correction rows satisfy
\[
R_{\nu}^{(b)}\in\mathcal C_{\rm car},
\qquad
\nu\in\{\mathrm{line},\mathrm{bg/core},\mathrm{norm/sub}\},
\qquad
b\in\{7,9\},
\]
then
\[
P_{W_{\rm car}}R_{\nu}^{(b)}=0
\]
for all six correction rows.  Consequently
\[
L_{\rm eff}=L_{\rm clean}^{\mathrm{tri}},
\qquad
\varepsilon_{\rm rem}=0.
\]
Thus the half-leading residual estimate closes automatically in the active
carrier quotient.
\end{remark}

\begin{lemma}[Nuisance correction rows vanish in the carrier quotient]
\label{lem:agent-six-final-nuisance-correction-rows-vanish}
Assume that the finite correction rows through order nine satisfy
\[
R_{\rm line}^{(b)}\in\mathcal C_{\rm car},
\qquad
R_{\rm bg/core}^{(b)}\in\mathcal C_{\rm car},
\qquad
R_{\rm norm/sub}^{(b)}\in\mathcal C_{\rm car},
\]
for \(b\in\{7,9\}\).  Then, in
\[
\mathcal Q_{\rm car}
=
\mathcal M_{\rm car}/\mathcal C_{\rm car},
\]
one has
\[
[R_{\rm line}^{(b)}]
=
[R_{\rm bg/core}^{(b)}]
=
[R_{\rm norm/sub}^{(b)}]
=
0.
\]
Therefore
\[
P_{W_{\rm car}}R_{\nu}^{(b)}=0,
\qquad
\nu\in\{\mathrm{line},\mathrm{bg/core},\mathrm{norm/sub}\},
\quad
b\in\{7,9\}.
\]
Equivalently, in the residual-table notation,
\[
L_{\nu}^{(b)}=0,
\qquad
R_{\nu,\mathrm{val}}^{(b)}=0,
\qquad
R_{\nu,\mathrm{abs}}^{(b)}=0.
\]
\end{lemma}

\begin{proof}
The quotient map
\[
\mathcal M_{\rm car}\longrightarrow
\mathcal Q_{\rm car}
=
\mathcal M_{\rm car}/\mathcal C_{\rm car}
\]
kills every element of \(\mathcal C_{\rm car}\).  The stated carrier-routing
hypothesis places all six correction rows in \(\mathcal C_{\rm car}\), so their
classes vanish in \(\mathcal Q_{\rm car}\).  Since \(W_{\rm car}\) is the live
carrier plane in the quotient, their projection to \(W_{\rm car}\) is also
zero.  The residual-table identities follow by taking all same-scale matrices
and residual components equal to zero.
\end{proof}

\begin{corollary}[D4 residual audit closes in the formal carrier quotient]
\label{cor:agent-six-final-D4-residual-closes-carrier-quotient}
Assume Lemma~\ref{lem:agent-six-final-nuisance-correction-rows-vanish}, the
Gram-good bounds
\[
Q^{-C_0}\le g_{\rm Gr},h_{\rm Gr}\le Q^{C_0},
\]
and the carrier-frame lower bound
\[
\sigma_{\min}(E_{\rm car})\ge Q^{-C_0}.
\]
Then
\[
L_{\rm eff}=L_{\rm clean}^{\mathrm{tri}}
=
\begin{pmatrix}
c_7&0\\
\beta_9^{\mathrm{tri}}&\alpha_9
\end{pmatrix},
\qquad
\beta_9^{\mathrm{tri}}:=\beta_9-\omega_Y(e_7)c_7,
\]
and
\[
\varepsilon_{\rm rem}=0.
\]
Moreover
\[
\det L_{\rm eff}
=
-\frac{1}
{1536\pi^4
g_{\rm Gr}^{3/2}h_{\rm Gr}^{3}
(\sqrt{g_{\rm Gr}}+\sqrt{h_{\rm Gr}})^3},
\]
so
\[
\sigma_{\min}(E_{\rm car}L_{\rm eff})\ge Q^{-C}
\]
on Gram-good patches.  Finally,
\[
\|\varepsilon_{\rm rem}\|_{\mathcal Q_{\rm car}}
=
0
\le
\frac12 |H|\,
\sigma_{\min}(E_{\rm car}L_{\rm eff})
\left(\Gamma_G^2+\Gamma_{G,2}^2\right)^{1/2}.
\]
Thus \(\mathrm{D4}_{\rm res}\) is Green in the formal carrier quotient for the
canonical scalar source \(u=q''_{\can}\).
\end{corollary}

\begin{proof}
By Lemma~\ref{lem:agent-six-final-nuisance-correction-rows-vanish}, all six
correction-row projections vanish in the carrier quotient.  Hence
\[
L_{\rm eff}
=
L_{\rm clean}^{\mathrm{tri}}
+
\sum_{\nu,b}L_{\nu}^{(b)}
=
L_{\rm clean}^{\mathrm{tri}},
\qquad
\varepsilon_{\rm rem}=0.
\]
The determinant formula is the clean triangular determinant from
Lemma~\ref{lem:clean-triangular-d4-carrier-matrix}.  On the
Gram-good patch this determinant is polynomially bounded below, while the
entries of \(L_{\rm clean}^{\mathrm{tri}}\) are polynomially bounded.
Therefore
\[
\sigma_{\min}(L_{\rm clean}^{\mathrm{tri}})\ge Q^{-C}.
\]
Combining this with
\[
\sigma_{\min}(E_{\rm car})\ge Q^{-C_0}
\]
gives
\[
\sigma_{\min}(E_{\rm car}L_{\rm eff})\ge Q^{-C}.
\]
The residual estimate is immediate from \(\varepsilon_{\rm rem}=0\).
\end{proof}

\begin{remark}[Scalar \(I\)-coordinate D4 is dead]
\label{rem:agent-six-final-scalar-I-D4-quarantined}
The scalar \(I\)-coordinate D4 formulation is not an active proof route.  It
collapses distinct carrier classes and reintroduces cancellations that are
absent in the carrier quotient.  The active D4 statement is the
carrier-quotient row-energy statement for the canonical scalar source
\[
u=q''_{\can}.
\]
Any future scalar \(I\)-coordinate argument would be a new theorem, not a
fallback use of the present D4 package.
\end{remark}

\begin{remark}[Compatibility with high-dimensional Gate--1 redesigns]
\label{rem:agent-six-final-D4-family-valued-compatibility}
The D4 carrier-quotient mechanism is tied to the canonical scalar curvature
source
\[
u=q''_{\can}.
\]
It is not a generic converter for arbitrary high-dimensional or family-valued
Gate--1 observables.

Thus a redesigned upstream Gate--1 may feed the present downstream D4 route
only if it certifies the same canonical source \(q''_{\can}\), for example by
implying the required canonical jet or row lower bound for \(u=q''_{\can}\).
Arbitrary projected observables
\[
G_\nu=L_\nu X
\]
are not themselves D4 sources unless each \(G_\nu\) comes with its own
canonical phase, Gram matrix, whitening tower, carrier quotient, and order-seven
and order-nine D4 identities.  No such generic transfer is asserted here.
\end{remark}

\begin{theorem}[Effective-matrix carrier-quotient D4 conversion]
\label{thm:effective-matrix-carrier-quotient-d4}
Let
\[
E_{\rm car}
:=
\bigl([\mathsf C_7]\ [\mathsf C_{Y12}]\bigr)
\]
be the two-column carrier response map into \(\mathcal Q_{\rm car}\).  Assume
the carrier frame is polynomially conditioned:
\[
\sigma_{\min}(E_{\rm car})\ge Q^{-C}.
\]
Assume the carrier-row class satisfies
\[
P_{W_{\rm car}}([\mathcal A_7]+[\mathcal A_9])
=
H\,E_{\rm car}L_{\rm eff}
\begin{pmatrix}
\Gamma_G\\[1mm]
\Gamma_{G,2}
\end{pmatrix}
+
\varepsilon_{\rm rem},
\]
where
\[
  \sigma_{\min}(L_{\rm eff})\ge Q^{-C}
\]
and
\[
  \|\varepsilon_{\rm rem}\|_{\mathcal Q_{\rm car}}
  \le
  \frac12
  |H|\,
  \sigma_{\min}(E_{\rm car}L_{\rm eff})
  \left(\Gamma_G^2+\Gamma_{G,2}^2\right)^{1/2}.
\]
Then
\[
\left|
P_{W_{\rm car}}([\mathcal A_7]+[\mathcal A_9])
\right|_{\mathcal Q_{\rm car}}
\ge
Q^{-C}
|H|
\left(\Gamma_G^2+\Gamma_{G,2}^2\right)^{1/2}.
\]
Consequently, combined with Gate--2, wall escape, and row mobility, the local
carrier-quotient D4 output gives interval carrier-quotient row energy.

In the formal carrier quotient of
Corollary~\ref{cor:agent-six-final-D4-residual-closes-carrier-quotient},
one has
\[
  L_{\rm eff}=L_{\rm clean}^{\mathrm{tri}},
  \qquad
  \varepsilon_{\rm rem}=0,
\]
so the hypotheses above reduce to the clean triangular determinant and carrier
frame lower bounds.
\end{theorem}

\begin{proof}
The polynomial lower bound for \(E_{\rm car}\) and \(L_{\rm eff}\) gives
\[
  \left\|
  H\,E_{\rm car}L_{\rm eff}
  \begin{pmatrix}
  \Gamma_G\\
  \Gamma_{G,2}
  \end{pmatrix}
  \right\|_{\mathcal Q_{\rm car}}
  \ge
  Q^{-C}
  |H|
  \left(\Gamma_G^2+\Gamma_{G,2}^2\right)^{1/2}.
\]
The residual estimate is half-leading, so the triangle inequality gives the
claimed lower bound.  Row mobility then converts the pointwise or packet-level
carrier lower bound into interval carrier-quotient energy.
\end{proof}

\begin{remark}[Why the unit-coordinate D4 form is no longer used]
\label{rem:no-unit-coordinate-d4-form}
The carrier coordinates should not be normalized as if the D4 map were
\[
  H I_2
  \begin{pmatrix}
  \Gamma_G\\
  \Gamma_{G,2}
  \end{pmatrix}.
\]
The clean finite computation gives the triangular row-normal matrix
\[
  L_{\rm clean}^{\mathrm{tri}}
  =
  \begin{pmatrix}
  c_7&0\\
  \beta_9^{\mathrm{tri}}&\alpha_9
  \end{pmatrix},
  \qquad
  \beta_9^{\mathrm{tri}}:=\beta_9-\omega_Y(e_7)c_7.
\]
The same-scale \(\Gamma_G\)-term in the order-nine row is part of the effective
leading matrix.  It is not a residual.
\end{remark}

\begin{remark}[Scalar \(I\)-coordinate incidence is a collapse artifact]
\label{rem:scalar-i-incidence-collapse-artifact}
The scalar obstruction
\[
(A_9^{\mathfrak f})_I
=
\frac{yq''\Gamma_{G,2}}{128\pi^2g^3}
+
R_9^I
\]
arises after the collapse map
\[
S_I:\mathcal M_{\rm car}\to \mathcal O\cdot I
\]
has sent several independent carrier directions to the same scalar
\(I\)-coordinate.
The active D4 theorem in this draft is the carrier-quotient theorem
\ref{thm:d4-carrier-quotient-through-order-nine}; the scalar \(I\)-coordinate
formulation below is retained only as a fallback audit.  In the carrier quotient, the lower-carrier, line,
background/core, and normalization/subtraction contributions are nuisance
directions.  Thus the scalar incidence wall
\[
F_9^I=0
\]
is relevant only if the downstream row theorem insists on collapsed scalar
\(I\)-coordinate energy.  It is not an obstruction to carrier/quotient row
energy.

If one nevertheless imposes scalar \(I\)-coordinate row energy as the downstream
output, the remaining finite obstruction is the incidence wall
\[
\mathcal E_{\rm cc}^{(9)}
=
\left\{
F_9^I=0,\quad yq''\Gamma_{G,2}\ne0
\right\},
\]
where
\[
F_9^I
:=
R_9^I+
\frac{yq''\Gamma_{G,2}}{128\pi^2g^3}.
\]
Closing that scalar formulation would require either proving
\(F_9^I\not\equiv0\) together with a quantitative avoidance theorem for
\(\mathcal E_{\rm cc}^{(9)}\), proving that \(R_9^I\) is zero,
\(H\)-divisible, or nuisance-quotient, or replacing the scalar readout by a
higher-order or ambient row.  Under the carrier-quotient formulation this
scalar incidence is not a D4 obstruction.
We refer to this fallback scalar obstruction as the order-nine scalar
cross-channel incidence problem.  It is not part of the active
carrier-quotient D4 route.
\end{remark}

\begin{remark}[Admissibility of carrier-quotient row energy]
\label{rem:carrier-quotient-row-energy-admissible}
The downstream row obstruction is formulated at the row-module or quotient-row
level, not only after scalar collapse to the \(I\)-coordinate.  Therefore a
lower bound of the form
\[
\int_I
\left|
[\mathcal A_N^{\rm car}(m)]
\right|_{\mathcal Q_{\rm car}}^2\,dm
\ge Q^{-C}|I|
\]
is an admissible row-energy output.  The scalar coordinate
\[
S_I(\mathcal A_N)
\]
is only one coordinate readout and is not faithful on the carrier quotient.  If
one were to require scalar \(I\)-coordinate energy only, then the scalar
cross-channel wall would remain a genuine obstruction.  Under the carrier
quotient formulation it is quotiented away.
\end{remark}

\begin{definition}[D4 residual valuation ledger]
\label{def:d4-residual-valuation-ledger}
After extracting the carrier-quotient leading rows, write the remaining D4 row
contribution schematically as
\[
\mathbf R
=
\mathbf R_{\rm same}
+
\mathbf R_{\rm val}
+
\mathbf R_{\rm abs}.
\]
Terms with the same anisotropy-linear valuation as the leading row are not
automatically errors.  They must first be absorbed into the effective leading
matrix \(L_{\rm eff}\).  Only the component remaining after this absorption is
eligible to be estimated as a residual.
A residual term \(T\) has valuation
\[
T=q_0^a u^b\Gamma^cE_T,
\qquad
\Gamma=(\Gamma_G,\Gamma_{G,2}),
\]
where \(E_T\) is polynomially bounded or separately controlled.

Terms with the same anisotropy-linear valuation as the chosen leading D4 row
are not errors.  In the expected carrier-quotient case this valuation is
\[
(a,b,c)=(1,3,1),
\]
corresponding to \(q_0u^3\Gamma\).  If the exact row algebra contains a lower
or delayed \(u\)-valuation anisotropy-linear layer, the leading valuation is
chosen according to the filtered row-observability rule in
Remark~\ref{rem:filtered-d4-row-observability}.  They are absorbed into the
effective leading matrix or the carrier-quotient coordinate map.

Terms with higher \(q_0\)-valuation, higher anisotropy valuation, or higher
source valuation belong to \(\mathbf R_{\rm val}\) and are controlled by the
chosen non-wall and wall-escape thresholds.

Terms without sufficient factors of \(q_0,u,\Gamma\) belong to
\(\mathbf R_{\rm abs}\).  These are the only residuals requiring a hard
absolute smallness condition in the wall-escape corridor.

Equivalently, for a schematic term
\[
T=q_0^a u^b \Gamma^\ell D,
\qquad |D|\le Q^C,
\]
compare to the leading scale
\[
q_0u^3\Gamma.
\]
The relative size is bounded by
\[
\frac{|T|}{|q_0u^3\Gamma|}
\le
Q^C
|q_0|^{a-1}
|u|^{b-3}
|\Gamma|^{\ell-1}.
\]
Terms with \((a,b,\ell)=(1,3,1)\) belong in the effective leading matrix \(L\),
not in the residual.  Terms with extra useful \(q_0,u,\Gamma\)-valuation are
relative residuals.  Only terms without enough useful valuation are absolute
floors and must satisfy the wall-escape corridor estimates.
\end{definition}

\begin{remark}[Filtered D4 row observability]
\label{rem:filtered-d4-row-observability}
The notation \(q_0u^3L\Gamma\) records the expected leading row scale.  It
should not be hard-coded if the exact row algebra contains lower
\(u\)-valuation anisotropy-linear terms.  In general one should write
\[
\mathbf A_N^{\rm lin}
=
q_0\sum_{b\ge b_0}u^bL_b\Gamma.
\]
D4 closes once the stacked filtered response
\[
\begin{pmatrix}
L_{b_0}\\
uL_{b_0+1}\\
\vdots\\
u^{B-b_0}L_B
\end{pmatrix}
\]
has second singular value bounded below by \(Q^{-C}\) on the sampled branch.
The expected carrier-quotient case is \(b_0=B=3\) and
\[
L_3=-\frac1{288}I_2.
\]
If a lower \(u\)-valuation anisotropy-linear layer appears, it should be
promoted into the leading filtered matrix rather than treated as an error.
\end{remark}

\begin{lemma}[Clean triangular D4 carrier matrix]
\label{lem:clean-triangular-d4-carrier-matrix}
In the D4 finite carrier calculation, write the centered Gram diagonal scalars as
\[
g_{\rm Gr}:=\frac{q_0}{\pi},
\qquad
h_{\rm Gr}:=
\frac1\pi\left(\frac13q_0^3+\frac16u\right),
\qquad
u=q''_{\can}.
\]
Let \(\kappa=\omega_Y(e_7)\), and define the row-normal coefficient
\[
\beta_9^{\mathrm{tri}}(g,h,u):=\beta_9(g,h,u)-\kappa\,c_7(g,h).
\]
Then
\[
g_{\rm Gr}-h_{\rm Gr}
=
-\frac{\Gamma_G}{6\pi}.
\]
In this lemma set \(g=g_{\rm Gr}\), \(h=h_{\rm Gr}\).  The clean order-seven
carrier coordinate is
\[
\pi_{\mathfrak f}T_7^{(1)}
=
H\,c_7(g,h)\Gamma_G I,
\qquad
H=yu,
\]
where
\[
c_7(g,h)
=
-\frac{1}
{24\pi^2gh(\sqrt g+\sqrt h)^2}.
\]

For the clean order-nine channel, write
\[
X_{-,1}=BS,
\qquad
B=
\frac{u}
{4\pi\sqrt{gh}(\sqrt g+\sqrt h)}.
\]
Let
\[
G_2=\operatorname{diag}(K_1,K_2),
\qquad
K_1=\frac{u}{8\pi},
\qquad
K_2=
\frac{q_0^2u+\frac16q^{(4)}}{8\pi}.
\]
If \(X_{-,2}=aI+dD\), then
\[
d=
\frac14
\left[
\frac{B^2h+2B^2\sqrt{gh}-K_1/g}{\sqrt g}
-
\frac{B^2g+2B^2\sqrt{gh}-K_2/h}{\sqrt h}
\right].
\]
Equivalently,
\[
d
=
-\frac{\Gamma_{G,2}}{32\pi h^{3/2}}
+
\Gamma_G\,\rho_9(g,h,u),
\]
where
\[
\rho_9(g,h,u)
=
\frac{B^2(g+3\sqrt{gh}+h)}
{24\pi(\sqrt h+\sqrt g)\sqrt{gh}}
-
\frac{K_1(g+\sqrt{gh}+h)}
{24\pi(\sqrt h+\sqrt g)g^{3/2}h^{3/2}}.
\]
Since
\[
\pi_{\mathfrak f}T_9^{(1,2)}
=
-2yBd\,I,
\]
in the row-normal triangular coordinates the clean order-nine carrier
coordinate has the form
\[
\pi_{\mathfrak f}T_9^{(1,2)}
=
H\left(
\beta_9^{\mathrm{tri}}(g,h,u)\Gamma_G
+
\alpha_9(g,h)\Gamma_{G,2}
\right)I,
\]
where
\[
\alpha_9(g,h)
=
\frac{1}
{64\pi^2\sqrt{gh}\,h^{3/2}(\sqrt g+\sqrt h)}
\]
and
\[
\beta_9^{\mathrm{tri}}(g,h,u)
=
-\frac{\rho_9(g,h,u)}
{2\pi\sqrt{gh}(\sqrt g+\sqrt h)}
-\kappa\,c_7(g,h).
\]
In particular, on \(g=h\),
\[
\alpha_9(g,g)=\frac1{128\pi^2g^3},
\]
recovering the clean wall formula
\[
\pi_{\mathfrak f}T_9^{(1,2)}
=
\frac{H\Gamma_{G,2}}{128\pi^2g^3}I.
\]

Thus, in the clean carrier coordinates,
\[
\begin{pmatrix}
B_7\\
B_9
\end{pmatrix}
=
H
\begin{pmatrix}
c_7&0\\
\beta_9^{\mathrm{tri}}&\alpha_9
\end{pmatrix}
\begin{pmatrix}
\Gamma_G\\
\Gamma_{G,2}
\end{pmatrix}.
\]
The same-scale \(\beta_9^{\mathrm{tri}}\Gamma_G\) term belongs to the effective
leading matrix, not a residual.  Moreover
\[
\det
\begin{pmatrix}
c_7&0\\
\beta_9^{\mathrm{tri}}&\alpha_9
\end{pmatrix}
=
c_7\alpha_9
=
-\frac{1}
{1536\pi^4
g^{3/2}h^3(\sqrt g+\sqrt h)^3}.
\]
Consequently this clean D4 carrier matrix is polynomially invertible on every
Gram-good patch on which
\[
Q^{-C}\le g_{\rm Gr},h_{\rm Gr}\le Q^C.
\]
\end{lemma}

\begin{remark}[Associated-graded D4 determinant certificate]
\label{rem:associated-graded-d4-determinant-certificate}
\emph{Historical pre-triangular bookkeeping.} This remark is retained only as a
reference to the pre-normalization carrier. All live row-normal formulas use the
triangularized coefficient \(\beta_9^{\mathrm{tri}}\) from the row-normal D4 branch.

In the clean carrier normalization of
Lemma~\ref{lem:clean-triangular-d4-carrier-matrix},
\[
L_{\rm clean}
=
\begin{pmatrix}
c_7&0\\
\beta_9^{\mathrm{tri}}&\alpha_9
\end{pmatrix},
\]
and the corresponding clean carrier image is
\[
F_{2Y}^{\mathrm{tri}}
=
H
\begin{pmatrix}
c_7\Gamma_G\\
\beta_9^{\mathrm{tri}}\Gamma_G+\alpha_9\Gamma_{G,2}
\end{pmatrix}.
\]
The same-scale \(\beta_9^{\mathrm{tri}}\Gamma_G\) contribution belongs to
\(L_{\rm eff}\), not to the residual.  Consequently,
\[
\det L_{\rm clean}
=
c_7\alpha_9
=
-\frac{1}
{1536\pi^4
g_{\rm Gr}^{3/2}h_{\rm Gr}^{3}
(\sqrt{g_{\rm Gr}}+\sqrt{h_{\rm Gr}})^3},
\]
which is polynomially bounded below on Gram-good patches.  Hence the clean D4
carrier matrix is triangularly invertible before correction-row residuals are
added.
\end{remark}

\begin{remark}[Nonzero D4 root traps are harmless under energetic sampling]
\label{rem:nonzero-d4-root-trap}
If a triangularized row coordinate has the form
\[
B_7=X+E_7,
\qquad
B_9=Y R_Y(Y)+X(\cdots)+E_9,
\]
where
\[
X=\Gamma_G,\qquad Y=\Gamma_{G,2},
\]
then a possible finite obstruction is a nonzero root
\[
R_Y(\zeta)=0,\qquad \zeta\ne0.
\]
The corresponding bad branch is
\[
X\approx0,\qquad Y\approx\zeta.
\]
This branch cannot persist on a polynomial-size real interval under the
Gate--2 ODE identities.  Indeed,
\[
X=q''-6q+2q^3,
\qquad
Y=2q(q')^2-\frac16X''.
\]
Testing the second identity against a cutoff and integrating \(X''\) by parts
twice converts the apparent derivative obstruction into an \(L^2\)-small
\(X\)-term, just as in the \(L^2\)-input rigidity proof.  Thus an interval on
which \(X\approx0\) and \(Y\approx\zeta\ne0\) would force \(Y\) small on a
shrunken interval, contradicting \(|\zeta|\ge Q^{-C}\).

Consequently, under existential or energetic row sampling, nonzero row roots
are not fatal: isolated occurrences can be avoided, and interval traps are
ruled out by the same local dynamics behind Gate--2 rigidity.  If the proof
were exact-point locked, such roots would again become relevant.
\end{remark}

\begin{remark}[Raw skew coefficients \(F_5\) and \(F_7\)]
\label{rem:raw-skew-coefficients-F5-F7}
For the centered raw off-diagonal scalar
\[
F(s)
=
\frac{\sin\Delta(s)+q_-(s)s\cos\Delta(s)}{s^2},
\]
with
\[
\Delta(s)
=
q_0s+\frac{q_2}{24}s^3+\frac{q_4}{1920}s^5
+\frac{q_6}{322560}s^7+\cdots,
\]
and
\[
q_-(s)
=
q_0+\frac{q_2}{8}s^2+\frac{q_4}{384}s^4
+\frac{q_6}{46080}s^6+\cdots,
\]
the raw odd carriers are skew:
\[
N_{2j+1}=F_{2j+1}J.
\]
In the normalization used here,
\[
F_5
=
-\frac{q_0^7}{630}
+\frac{q_0^4q_2}{72}
-\frac{q_0^2q_4}{480}
-\frac{q_0q_2^2}{144}
+\frac{q_6}{40320},
\]
and
\[
\begin{aligned}
F_7
={}&
\frac{q_0^9}{36288}
-\frac{q_0^6q_2}{1728}
+\frac{q_0^4q_4}{4608}
+\frac{5q_0^3q_2^2}{3456} \\
&-\frac{q_0^2q_6}{64512}
-\frac{q_0q_2q_4}{4608}
-\frac{5q_2^3}{41472}
+\frac{q_8}{9289728}.
\end{aligned}
\]
These coefficients are useful for auditing the raw row table, but they multiply
\(J\) and hence do not contribute directly to the fixed \(I\)-coordinate.
\end{remark}

There is also a non-centered visibility route.  The local construction uses
\[
t_\pm=m\pm\frac{s}{2},
\]
so \(m\) is a sampling midpoint, not necessarily the zero-packet center.  For
the paired packet
\[
f_{\beta,\gamma}(t)
=
\frac{a}{a^2+x^2}+\frac{b}{b^2+x^2},
\qquad x=2t-\gamma,
\]
one has, near \(x_m=2m-\gamma=0\),
\[
q^{(7)}(m)
=
5160960(a^{-9}+b^{-9})x_m+O(x_m^3).
\]
Thus \(x_m=0\) is a simple zero of \(q^{(7)}(m)\), and every sufficiently small
nonzero perturbation of \(m\) activates the septic \(K_1\)-transport slice,
outside
\[
q^{(7)}=0,
\qquad
\Xi_{K_1}=0.
\]

\begin{proposition}[Packet visibility on \(\mathcal B_{\rm red}^{+}\)]
\label{prop:packet-visibility-on-bred-plus}
On \(\mathcal B_{\rm red}^{+}\), the finite fixed-sector vector
\[
C_{\rm fd}^{\vec f}
=
(A_7^{\mathfrak f},A_9^{\mathfrak f},\ldots)
\]
is nonzero for the reduced off-line packet family.
\end{proposition}

\begin{proof}
For the centered route, the direct skew \(N_7^{\off}\)-term projects to zero,
but the one-whitening-slot \(K_1\)-channel has fixed-sector projection
\[
-2yBr\,I.
\]
This vector is nonzero outside
\[
\{y=0\}\cup\{q''=0\}\cup\{g=h\}.
\]
To rule out cancellation by the \(K_5\)-line sector using quotient isolation,
we also exclude
\[
v_5=0.
\]
Indeed,
\[
\omega_h(-2yBr\,I,A_5^{\mathfrak f})
=
-2yBr\,v_5,
\]
so \(K_5\)-line terms cannot cancel the \(K_1\)-channel on the generic patch
\[
yBrv_5\ne0.
\]
The branch \(v_5=0\) is therefore either included in
\(\mathcal E_{\rm invis}^{\core,7}\) or must be treated by a separate
full-vector \(K_5\)-coefficient computation.

For the non-centered route, \(m\) is a free sampling midpoint.  The expansion
of \(q^{(7)}(m)\) shows that \(q^{(7)}=0\) is a proper analytic wall near the
packet center.  Away from this wall and from \(\{\Xi_{K_1}=0\}\), the septic
\(K_1\)-transport slice gives \(A_7^{\mathfrak f}\ne0\).
\end{proof}

\begin{remark}[Order-seven channel independence ledger]
\label{rem:order-seven-channel-independence-ledger}
The full order-seven corrected-whitening coefficient is
\[
A_{7,\core}^{\mathfrak f}
=
\pi_{\mathfrak f}
\sum_{a+g+b=7}
X_{-,a}N_gX_{+,b}.
\]
The direct skew \(N_7^{\off}\)-term is killed by \(\pi_{\mathfrak f}\).  The
live one-whitening-slot term belongs to the \(K_1\)-transport sector.  \(K_3\)
shadow terms are projected away, pure same-point septic terms vanish in the
odd two-point germ, and \(K_5\)-terms lie in
\[
\mathbf C A_5^{\mathfrak f}.
\]
Thus exact cancellation of the live \(-2yBr\,I\) channel can only come from a
same \(K_1\)-transport contribution or a direct diagonal \(N_7\)-term with the
same jet monomial.  No structural identity forcing such cancellation is
currently present; such cancellation is included in the visibility incidence
locus.
\end{remark}

\subsection{The \(\eta\eta\) row as rank-two affine forcing}
\label{subsec:etaeta-row-rank-two-affine-forcing}

In the rank-two formulation, the active Bottleneck-C row is no longer a scalar
defect by itself.  It is a vector-valued section:
\[
E_{\eta\eta}
=
-\frac{\nu_{\eta\eta}}{24\pi q^3}(wI+uS)
\in
\Gamma(\mathcal F_{\rm red}^{+}).
\]
Under a frame change
\[
E_j=E_iG_{ij},
\]
the coordinate vector of \(wI+uS\) transforms by
\[
[w,u]_j^T
=
G_{ij}^{-1}[w,u]_i^T.
\]
The scalar coefficient
\[
-\frac{\nu_{\eta\eta}}{24\pi q^3}
\]
is fixed-convention scalar data.  Therefore \(E_{\eta\eta}\) transforms as a
section of the rank-two fixed-sector bundle, provided the reduced atlas is
functorial.

This row should not be identified with ordinary curvature of
\(\nabla^{\mathfrak f}\).  Curvature is antisymmetric in two tangent
directions:
\[
F_{\nabla^{\mathfrak f}}(X,Y)
=
[\nabla_X^{\mathfrak f},\nabla_Y^{\mathfrak f}]
-
\nabla_{[X,Y]}^{\mathfrak f}.
\]
The active row is symmetric in the source directions \(\eta_i,\eta_j\).  It is
better classified as a symmetric source-Hessian forcing term, or a
second-fundamental-form component, in the fixed-sector transport equation.

The natural local replacement for scalar Bottleneck-C flatness is therefore an
affine transport equation of the form
\[
\nabla_\kappa^{\mathfrak f}A_7^{\mathfrak f}
=
E_{\eta\eta}
+
E_{\rm lower}.
\]
Here \(E_{\rm lower}\) denotes all lower, line, shadow, pure, connection, and
background/core terms.  The expected decomposition is
\[
E_{\rm lower}
=
E_{\rm GateA}
+
E_{\rm GateB}
+
E_{\rm line}
+
E_{K_3}
+
E_{\rm pure}
+
E_{\rm conn}
+
E_{\rm bg/core}
+
E_{K_1}^{\rm comp}.
\]
Under the current local ledger,
\[
E_{\rm GateA}=0,
\qquad
E_{\rm GateB}=0
\]
under the accepted Gate A/B hypotheses,
\[
E_{\rm conn}=0
\]
after passing to the covariant derivative,
\[
E_{K_3}=0
\]
after fixed-sector projection,
\[
E_{\rm pure}=0
\]
in the odd two-point germ, and
\[
E_{\rm line}\in \mathbf C A_5^{\mathfrak f}.
\]
The remaining table dependency is to verify that no additional uncontrolled
same-grade \(K_1\) source-Hessian companion survives outside the already
identified \(E_{\eta\eta}\) row.

The current local row table is:
\[
\begin{array}{c|c|c|c}
\text{row class} & \text{source} & \text{fixed-sector status}
& \text{present role}\\ \hline
E_{\rm GateA}
& \text{aggregate homogeneous scalar-grade rows}
& 0
& \text{closed by Gate A}\\
E_{\rm GateB}
& \text{ordinary order }\le 7\text{ local rows}
& 0
& \text{closed by Gate B hypotheses}\\
E_{\eta\eta}
& \text{symmetric second-source Hessian / mixed-block row}
& -\dfrac{\nu_{\eta\eta}}{24\pi q^3}(wI+uS)
& \text{retained as affine forcing}\\
E_{\rm line}
& K_5\text{ normalization shadows}
& \in\mathbf C A_5^{\mathfrak f}
& \text{transported line term}\\
E_{K_3}
& K_3\text{ shadow sector}
& 0
& \text{projected away}\\
E_{\rm pure}
& \text{pure same-point septic sector}
& 0
& \text{vanishes in odd two-point germ}\\
E_{\rm conn}
& \text{frame/projection/whitening gauge motion}
& 0
& \text{absorbed by }\nabla^{\mathfrak f}\\
E_{\rm bg/core}
& \text{fixed background/core baseline terms}
& \text{fixed analytic section}
& \text{included in affine forcing}
\end{array}
\]

The visibility-positive incidence exclusions used to define
\(\mathcal B_{\rm red}^{+}\) do not alter Gate A/B; they only restrict the
base to points where the packet-visibility channel is active.

\begin{theorem}[Conditional local rank-two affine forcing equation]
\label{thm:conditional-local-rank-two-affine-forcing}
On \(\mathcal B_{\rm red}^{+}\), assume the rank-two connection theorem and
the row-table classifications above.  Then
\[
\nabla_\kappa^{\mathfrak f}A_7^{\mathfrak f}
=
E_{\eta\eta}+E_{\rm lower},
\]
where
\[
E_{\eta\eta}
=
-\frac{\nu_{\eta\eta}}{24\pi q^3}(wI+uS)
\]
is a tensorial section of \(\mathcal F_{\rm red}^{+}\), and \(E_{\rm lower}\)
contains only Gate A/B-closed, connection-absorbed, line, shadow, pure,
background/core, and already-classified analytic forcing terms.
\end{theorem}

\begin{proof}
By definition,
\[
A_7^{\mathfrak f}
=
[s^7]\Pi_{\mathfrak f}\widehat\Omega_\zeta^{\corr}(s;m).
\]
The projected connection gives
\[
\nabla_\kappa^{\mathfrak f}A_7^{\mathfrak f}
=
\Pi_{\mathfrak f}d_\kappa^{\mathcal J}
\left([s^7]\widehat\Omega_\zeta^{\corr}\right),
\]
so frame, projection, and whitening-gauge variations are absorbed into
\(\nabla^{\mathfrak f}\).  The \(\eta_i\eta_j\) second-source Hessian row of
the promoted mixed block gives \(E_{\eta\eta}\).  All remaining terms are
classified by the displayed table.
\end{proof}

\begin{remark}[The equation is local, not yet global]
\label{rem:local-affine-forcing-not-global}
Theorem~\ref{thm:conditional-local-rank-two-affine-forcing} does not prove a
global contradiction.  It changes the local interpretation of the
Bottleneck-C row: the row is no longer a scalar defect demanding pointwise
vanishing, but an inhomogeneous forcing term in a rank-two transport equation.
A global boundary, holonomy, exactness, index, or transported finite-core
non-concentration theorem is still required.
\end{remark}

\begin{remark}[What the affine equation changes]
\label{rem:what-affine-equation-changes}
The old scalar Bottleneck-C gate demanded local pointwise vanishing of the
\(S\)-visible component.  The rank-two equation instead allows the same row to
be an inhomogeneous forcing term in a vector-valued transport equation.  This
is a genuine conceptual improvement only if a later global theorem controls
the forced transport.  If one still demands
\[
E_{\eta\eta}=0
\]
pointwise, then the rank-two formulation collapses back to the old obstruction
\[
\nu_{\eta\eta}=0
\]
on a good patch.
\end{remark}

\subsection{Global affine transport dependency}
\label{subsec:global-affine-transport-dependency}

Assuming the local affine equation
\[
\nabla^{\mathfrak f}A=E
\]
is available, parallel transport along a path \(\Gamma:a\to b\) gives
\[
A(b)
=
U_{a\to b}A(a)
+
\int_a^b U_{t\to b}E(t)\,dt.
\]
For the order-seven fixed-sector section this becomes
\[
A_7^{\mathfrak f}(b)
=
U_{a\to b}A_7^{\mathfrak f}(a)
+
\int_a^b
U_{t\to b}
\left(
E_{\eta\eta}+E_{\rm lower}
\right)(t)\,dt.
\]
Thus the local Bottleneck-C obstruction becomes an affine displacement
problem.

At present no canonical global proof gate for this displacement is available.
The natural path and boundary candidates have the following status:
\begin{itemize}
\item \(\kappa\)-paths are locally natural but have no known endpoint values for
\(A_7^{\mathfrak f}\).
\item midpoint paths \(m\mapsto b(m)\) are natural locally, but good-patch
boundaries are exceptional walls rather than known boundary conditions.
\item source/core deformation paths are excluded on the reduced base, since
source/core transfer is not functorial as an ordinary vector-bundle
transition.
\item Gram or whitening loops may measure gauge/atlas data rather than zeta
geometry unless a global atlas theorem is proved.
\item zero-counting contours are canonical for \(\xi\) or \(\zeta\), but no
map from those contours into the rank-two fixed-sector base is currently
constructed.
\end{itemize}

Nor is the forcing currently exact.  In the fixed-source midpoint calculus, the
active scalar density satisfies
\[
2880Q^{14}\nu_{\eta\eta}
=
qAA''
+
20
\left(
\frac{q''}{6}
-
\frac{2q^3}{3}
\right)A^2,
\qquad
A=A_2.
\]
Integration by parts yields
\[
qAA''
+
20
\left(
\frac{q''}{6}
-
\frac{2q^3}{3}
\right)A^2
=
\left(qAA'-\frac12 q'A^2\right)'
-
q(A')^2
+
\left(
\frac{23}{6}q''
-
\frac{40}{3}q^3
\right)A^2.
\]
The last two terms are genuine bulk terms.  Therefore the forcing is not
presently a boundary term in the natural variables.

For a loop \(\Gamma\) based at \(b\), the corresponding affine holonomy data are
\[
\operatorname{AffHol}_\Gamma
=
\left(
U_\Gamma,
\Delta_\Gamma^{\rm aff}
\right),
\]
where
\[
\Delta_\Gamma^{\rm aff}
=
\int_\Gamma
U_{t\to b}
\left(E_{\eta\eta}+E_{\rm lower}\right)(t)\,dt.
\]
A holonomy proof gate would need all of:
\begin{enumerate}
\item a canonical loop \(\Gamma\subset\mathcal B_{\rm red}^{+}\);
\item a global theorem constraining \(U_\Gamma\);
\item a global zeta-theoretic law forcing \(\Delta_\Gamma^{\rm aff}=0\), or a
prescribed value;
\item an off-line-packet theorem forcing an incompatible value of
\(\Delta_\Gamma^{\rm aff}\).
\end{enumerate}
None of these is presently available.

\begin{remark}[Affine holonomy is not presently a proof gate]
\label{rem:affine-holonomy-not-proof-gate}
Loops around core-boundary walls, packet-visibility incidence walls, Gram
singularities, or whitening branch loci are diagnostic but not currently
canonical RH objects.  A loop around a core-boundary wall may change labels or
source/core convention; a Gram loop may measure only gauge or atlas monodromy;
and zero-counting contours live in the complex zeta plane, not yet in the
rank-two fixed-sector base \(\mathcal B_{\rm red}^{+}\).  Thus affine holonomy
cannot presently replace the local affine forcing equation.  The most
plausible global successors are real path/boundary transport or a transported
finite-core non-concentration theorem.
\end{remark}

\begin{remark}[Current global dependency]
\label{rem:current-global-affine-transport-dependency}
The rank-two affine equation is currently descriptive rather than
contradictory.  To turn it into a proof gate one needs one of:
\begin{enumerate}
\item canonical boundary values for \(A_7^{\mathfrak f}\);
\item a canonical path or loop with prescribed affine holonomy;
\item a path-independent affine displacement theorem;
\item a projective/scale index theorem derived from the rank-two transport;
\item a transported finite-core non-concentration theorem for the full
fixed-sector vector.
\end{enumerate}
None of these is presently proved.
\end{remark}

\begin{hypothesis}[Global rank-two affine transport gate]
\label{hyp:global-rank-two-affine-transport-gate}
There exists a canonical path, loop, boundary condition, index, or transported
finite-core non-concentration theorem on \(\mathcal B_{\rm red}^{+}\) that
controls
\[
\Delta_\Gamma^{\rm aff}
=
\int_\Gamma
U_{t\to b}
\left(E_{\eta\eta}+E_{\rm lower}\right)(t)\,dt.
\]
Equivalently, the rank-two affine forcing equation must be coupled to a global
law that distinguishes RH-compatible configurations from off-line-packet
configurations.  Without such a law, the rank-two architecture is locally
coherent but globally non-closing.
\end{hypothesis}

\subsection{Derived diagnostics: projective, vector-valued, and energy data}
\label{subsec:derived-diagnostics-projective-vector-energy}

Several objects remain useful as diagnostics, but none currently replaces the
rank-two affine transport theorem.

First, projectivizing the fixed-sector vector gives
\[
\ell_{2r+1}(m)
=
[A_{2r+1}^{\mathfrak f}(m)]
\in
\mathbf P(\mathfrak f_h(m)).
\]
The quotient determinant is exactly projective incidence:
\[
D_{2r+1}
=
\omega_h(A_{2r+1}^{\mathfrak f},A_5^{\mathfrak f}),
\]
and
\[
D_{2r+1}=0
\iff
\ell_{2r+1}=\ell_5.
\]
Thus projective/Maslov/spectral-flow data recover the quotient determinant
geometry.  They remain incomplete by themselves because pure line anomalies
\[
A_{2r+1}^{\mathfrak f}
=
\lambda A_5^{\mathfrak f}
\]
are projectively invisible but still visible to \(T\) and \(\Phi_K\).  A
projective index would therefore need to be paired with scale, energy, or the
full rank-two connection.

Second, the complete local fixed-sector object is the vector-valued odd germ
\[
\mathcal A_{\odd}^{\mathfrak f}(s;m)
=
\sum_{r\ge0}
A_{2r+1}^{\mathfrak f}(m)s^{2r+1}
\in
\mathfrak f_h(m).
\]
It retains all fixed-sector information, preserves componentwise holomorphy
and Cauchy recovery, and avoids the incompleteness of quotient scalarization.
However, it is essentially the full fixed-sector version of the split readout.
It does not by itself prove the required finite-core lower bound
\[
\theta(m)\notin
\mathcal Z_k^{\vec f}\cup\mathcal E_{\rm geom}^{\vec f}
\Longrightarrow
\left\|
P_{\rm fd}^{\vec f}(\theta(m))
+
J_{\rm fd}^{\vec f}(\theta(m))
\right\|_\infty
\ge
\Lambda_{\vec f}(m,Q).
\]
The aligned branch
\[
J_{\rm fd}^{\vec f}\approx -P_{\rm fd}^{\vec f}
\]
therefore remains a genuine downstream dependency.

Third, one may define fixed-sector energies such as
\[
\mathcal E_7
=
\|A_7^{\mathfrak f}\|_{\wh}^2
\]
or
\[
\mathcal E_{\odd}
=
\sum_r \omega_r
\|A_{2r+1}^{\mathfrak f}\|_{\wh}^2.
\]
For
\[
E_{\eta\eta}
=
\gamma(wI+uS),
\]
the first variation of \(\mathcal E_7\) contains
\[
2\operatorname{Re}\left[
\gamma\,
\langle A_7^{\mathfrak f},wI+uS\rangle_{\wh}
\right],
\]
which is sign-indefinite.  The second variation is positive but tautological:
\[
2|\gamma|^2\|wI+uS\|_{\wh}^2\ge0.
\]
No zero-energy, monotonicity, or endpoint theorem is presently available.
Thus energy detects the defect but does not yet provide a contradiction.

\begin{remark}[Role of Gate A and Gate B after the downstream repairs]
\label{rem:role-of-gate-a-b-after-audits}
Gate A and Gate B remain useful local algebraic lemmas:
\[
\text{Gate A: aggregate homogeneous cancellation,}
\]
\[
\text{Gate B: conditional local quotient/fixed-sector algebra through order }7.
\]
They simplify the local expansion and classify lower terms, but they are not
proof-closing global gates.

After the \(L^2\)-rigidity repair, Gate~2 source--anisotropy reintegration is
no longer an independent downstream input in the ordinary real-canonical
branch; it is supplied by
Theorem~\ref{thm:gate-two-from-gate-one-l2-rigidity}, conditional on the
real-canonical definitions and source local complexity.  The remaining burden
is therefore sharper: prove Gate~1 source-curvature observability from the
actual-zeta ordinary/singular inputs, and complete the D4 carrier-quotient row
conversion and residual-control audit.
\end{remark}

\begin{remark}[Coherent graph versus forbidden vertical residual]
\label{rem:coherent-graph-versus-forbidden-vertical}
The current architecture repeatedly has the same form:
\[
  \boxed{\text{coherent graph}}
  \qquad\text{versus}\qquad
  \boxed{\text{forbidden vertical residual}.}
\]

For the \(Y12\)-fiber, the coherent graph has the form
\[
  s_9^Y=\rho_Ys_7+\eta_{\mathrm{coh}},
\]
and the forbidden normal component is
\[
  s_9^Y-\rho_Ys_7.
\]

For finite/O2/O3, the coherent graph is the near-pure determinant span
\[
  D_{-n+1},\ldots,D_{-n+d_0},
\]
and the forbidden residual is the Schur component
\[
  D_{-n}
  \notin
  \operatorname{span}\{D_{-n+1},\ldots,D_{-n+d_0}\}.
\]

For Gate--1 finite AFE anti-lock, the coherent graph is the fold-locked
prime/dual cancellation model, and the forbidden residual is the remaining
canonical source energy
\[
  B_2+P_{2,\le X}+D_{2,\le X}
  \not\approx0
\]
on \(W_{\mathrm{branch}}\).

This language should be used to keep the three live modules aligned:
\[
  \mathrm{Agg}L^2(q''_{\can}),
  \qquad
  \mathrm{Graph/O2O3}_{23},
  \qquad
  \mathrm{Y12NormalD4}.
\]
\end{remark}

The finite-core problem below should therefore be read in a further updated
role.  It is no longer merely a downstream step after a closed Bottleneck-C
gate, nor is it only a generic vector-valued non-concentration problem.  The
source-collapse audits show that anisotropy non-concentration alone is
insufficient, because the known fixed-sector rows carry the common source
factor
\[
H=yq''_{\can}.
\]
The local downstream package supplies source--anisotropy reintegration from
Gate~1 by Theorem~\ref{thm:gate-two-from-gate-one-l2-rigidity}.  Thus the
remaining proof-relevant finite-core theorem must include Gate~1 source-level
nonflatness,
\[
\int_{I_\rho}|q''_{\can}(m)|^2\,dm\ge Q^{-B}|I_\rho|,
\]
and then connect the resulting Gate~2 source-weighted anisotropy to the D4
carrier-quotient row-energy output.

\begin{remark}[Visibility is not non-concentration]
\label{rem:visibility-not-nonconcentration}
The order-seven packet-visibility result proves that the relevant fixed-sector
vector is not identically zero on the reduced visibility-positive base.  It
does not prove the sampled finite-core lower bound.  The aligned branch
\[
J_{\rm fd}^{\vec f}\approx -P_{\rm fd}^{\vec f}
\]
remains possible until a vector-valued non-concentration, cone-control,
interaction-rigidity, or transport-index theorem rules it out.  Thus
visibility closes the local detection problem but not the global sampling
problem.
\end{remark}

\begin{remark}[Transported finite-core data do not remove alignment]
\label{rem:transported-finite-core-data-do-not-remove-alignment}
One may try to strengthen the finite-core vector by adding covariant transport
data, for instance
\[
C_{\rm fd}^{\rm tr}
=
(A_7^{\mathfrak f},
\nabla^{\mathfrak f}A_7^{\mathfrak f},
A_9^{\mathfrak f},
\nabla^{\mathfrak f}A_9^{\mathfrak f},
\ldots).
\]
Using the local affine equation, this is equivalent to adding
\[
E_{\eta\eta}+E_{\rm lower}
\]
as additional coordinates.  This is coherent, but it does not by itself prove
sampled non-concentration.  The old aligned branch simply becomes
\[
J_{\rm fd}^{\rm tr}\approx -P_{\rm fd}^{\rm tr}.
\]
Adding derivative or forcing coordinates enlarges the finite vector but does
not impose a sign, cone, monotonicity, or endpoint constraint that forbids
alignment.  A transported finite-core theorem would therefore require a new
coercivity, boundary, incidence, or cone-control input.
\end{remark}

\subsection{Source-weighted interval non-concentration after the Gate--2 repair}
\label{subsec:layered-finite-core-nonconcentration-gate}

The local fixed-sector visibility result is layered through order nine:
\[
\Gamma_G\ne0
\Longrightarrow
A_7^{\mathfrak f}\ne0,
\]
while
\[
\Gamma_G=0,\quad \Gamma_{G,2}\ne0
\Longrightarrow
A_9^{\mathfrak f}\ne0
\]
outside the named correction and cross-channel incidence walls.  This local
visibility theorem is not the remaining proof gate.

After Theorem~\ref{thm:gate-two-from-gate-one-l2-rigidity}, the old two-stage
downstream problem should be read as follows.  First one needs Gate~1
source-curvature observability,
\[
\int_{I_\rho}|q''_{\can}(m)|^2\,dm\ge Q^{-B}|I_\rho|.
\]
In the ordinary real-canonical branch, Gate~2 source--anisotropy reintegration
is then supplied locally:
\[
\int_{I_\rho}
|q''_{\can}(m)|^2
(\Gamma_G(m)^2+\Gamma_{G,2}(m)^2)\,dm
\ge Q^{-A}|I_\rho|.
\]
The remaining downstream conversion is D4: convert this source-weighted
anisotropy into carrier-quotient row energy.  In scalar language the target
scale is
\[
\boxed{
\sup_{m\in I_\rho}
|H(m)|
\max\{|\Gamma_G(m)|,|\Gamma_{G,2}(m)|\}
\ge Q^{-A},
\qquad
H=yq''_{\can}.
}
\]
In the carrier quotient, this is replaced by the quotient row-energy object in
Remark~\ref{rem:carrier-quotient-source-weighted-residual}.

In the transported version, one may instead work with
\[
C_{\rm fd}^{\rm tr}
=
\left(
A_7^{\mathfrak f},
\nabla^{\mathfrak f}A_7^{\mathfrak f},
A_9^{\mathfrak f},
\nabla^{\mathfrak f}A_9^{\mathfrak f},
\ldots
\right),
\]
or minimally, using the affine equation,
\[
C_{\rm fd}^{\rm tr}
=
\left(
A_7^{\mathfrak f},
E_{\eta\eta}+E_{\rm lower},
A_9^{\mathfrak f},
\ldots
\right).
\]
The remaining global obstruction is that pair and interaction pieces may still
align:
\[
J_{\rm fd}^{\rm tr}\approx -P_{\rm fd}^{\rm tr}.
\]
Thus the next proof-closing theorem must be a quantitative layered
non-concentration, incidence-distance, cone, or arithmetic separation result.

\begin{definition}[Source-weighted layered residual]
\label{def:source-weighted-layered-residual}
On a centered reduced good patch, in the canonical rank-two convention define
\[
H:=yq''_{\can}.
\]
The order-seven/order-nine layered fixed-sector rows have the schematic form
\[
A_7^{\mathfrak f}\sim H\Gamma_G,
\qquad
A_9^{\mathfrak f}\sim H\Gamma_{G,2},
\]
where
\[
\Gamma_G=q''_{\can}-6q_0+2q_0^3
\]
and
\[
\Gamma_{G,2}=(1-q_0^2)q''_{\can}-\frac16(q_\zeta)^{(4)}.
\]
Thus the relevant norm-scale residual is not merely
\[
(\Gamma_G,\Gamma_{G,2}),
\]
but the source-weighted vector
\[
\mathcal R_{\rm sw}
:=
\left(
H\Gamma_G,\,
H\Gamma_{G,2}
\right).
\]
Equivalently, the scalar lower-bound target is
\[
\boxed{
\sup_{m\in I_\rho}
|H(m)|
\max\{|\Gamma_G(m)|,|\Gamma_{G,2}(m)|\}
\ge Q^{-A}.
}
\]
At leading centered source order,
\[
y_{\rm lead}=-\frac{q_0A_2^2}{72},
\qquad
q''_{\can}=2A_2,
\]
and hence
\[
\boxed{
H=yq''_{\can}
\sim
-\frac{q_0(q''_{\can})^3}{288}
=
-\frac{q_0A_2^3}{36}.
}
\]
This cubic source factor is the main vulnerability of the rank-two norm route.
\end{definition}

\begin{remark}[Carrier quotient version of the source-weighted residual]
\label{rem:carrier-quotient-source-weighted-residual}
Definition~\ref{def:source-weighted-layered-residual} records the scalar
source-weighted scale
\[
(H\Gamma_G,\ H\Gamma_{G,2}).
\]
In the carrier-quotient row formulation, the corresponding object is
\[
\mathcal R_{\rm sw}^{\rm car}
:=
\left(
H\Gamma_G[\mathsf C_7],
\,
H\Gamma_{G,2}[\mathsf C_{Y12}]
\right)
\in \mathcal Q_{\rm car}^{\oplus 2}.
\]
This is the proof-relevant row-energy object whenever the downstream
contradiction is stated before scalar \(I\)-coordinate collapse.
\end{remark}

\begin{lemma}[Source-curvature completion for the rank-two plain wall]
\label{lem:rank-two-source-curvature-completion}
Let \(I\Subset I^*\) be a chart-good interval with \(R^{-C}\)-scale buffer, and set
\[
  p:=q''.
\]
Assume there is a finite source space \(\mathcal P_I\) and \(p_X\in\mathcal P_I\)
such that
\[
  \|p-p_X\|_{H^2_{\mathrm{sc}}(I^*)}\le Q^{-B_0}.
\]
Let
\[
  \mathcal V_I:=\mathcal A_1+\mathcal I^2\mathcal P_I,
  \qquad
  \mathcal A_1:=\operatorname{span}\{1,t\},
\]
and choose \(q_X\in\mathcal V_I\) with \(q_X''=p_X\), with affine modes included
or normalized.  Define
\[
  \mathcal W_G
  :=
  \mathcal P_I+\mathcal V_I+\operatorname{Sym}^3\mathcal V_I,
\]
and
\[
  \mathcal W_{G2}
  :=
  D^2\mathcal P_I+\mathcal P_I+
  \operatorname{Sym}^2\mathcal V_I\cdot\mathcal P_I.
\]
If
\[
  \|q_X\|_{C^4_{\mathrm{sc}}(I^*)}\le Q^{C_0},
\]
then
\[
  G_X:=\Gamma_G(q_X)\in\mathcal W_G,
  \qquad
  G_{2,X}:=\Gamma_{G,2}(q_X)\in\mathcal W_{G2},
\]
and, after increasing \(B_0\),
\[
  \|\Gamma_G(q)-G_X\|_{H^2_{\mathrm{sc}}(I)}
  +
  \|\Gamma_{G,2}(q)-G_{2,X}\|_{L^2(I)}
  \le Q^{-B}.
\]
\end{lemma}

\begin{proof}[Proof sketch]
Use
\[
  \Gamma_G(q)=-6q''+36q-12q^3=-6p+36q-12q^3,
\]
and
\[
  \Gamma_{G,2}(q)=q^{(4)}-6(1-q^2)q''=p''-6p+6q^2p.
\]
Since \(q_X''=p_X\), the claimed finite-space inclusions are immediate.

Set \(r=q-q_X\).  Then \(r''=p-p_X\), and affine reconstruction gives
\[
  \|r\|_{H^4_{\mathrm{sc}}(I)}
  \le Q^C\|p-p_X\|_{H^2_{\mathrm{sc}}(I^*)}.
\]
The residuals are
\[
  R_G=-6r''+(36-36q_X^2)r-36q_Xr^2-12r^3,
\]
and
\[
\begin{aligned}
  R_{G2}
  &=
  r^{(4)}-6r''
  +6q_X^2r''
  +12q_Xr q_X''
  +12q_Xr r''
  +6r^2q_X''
  +6r^2r'' .
\end{aligned}
\]
One-dimensional product estimates yield the stated
\(H^2_{\mathrm{sc}}\times L^2\) residual bound.
\end{proof}

\begin{remark}[Wall-variable completion contract]
\label{rem:rank-two-wall-variable-completion-contract}
The completion needed for the plain rank-two wall is a wall-variable completion,
not necessarily a primitive \(H^4(q)\)-completion.  It is enough to produce
finite wall representatives
\[
  G_X\in\mathcal W_G,
  \qquad
  G_{2,X}\in\mathcal W_{G2},
\]
such that
\[
  \|\Gamma_G-G_X\|_{H^2_{\mathrm{sc}}(I)}
  +
  \|\Gamma_{G,2}-G_{2,X}\|_{L^2(I)}
  \le Q^{-B}.
\]
The finite-atlas Riesz/Markov estimates are applied to \(G_X,G_{2,X}\), and
the residual is inserted only after the finite-atlas estimates are made.

A sufficient source-level route is to work with
\[
  p=q'',
\]
to approximate \(p\) in \(H^2_{\mathrm{sc}}(I^*)\) by \(p_X\in\mathcal P_I\),
and to reconstruct
\[
  q_X\in \mathcal V_I:=\mathcal A_1+\mathcal I^2\mathcal P_I
\]
with affine modes included or normalized.  Then, up to fixed harmless
normalization constants,
\[
  \Gamma_G=-6p+36q-12q^3,
  \qquad
  \Gamma_{G,2}=p''-6p+6q^2p,
\]
and the finite wall spaces are
\[
  \mathcal W_G
  =
  \mathcal P_I+\mathcal V_I+\operatorname{Sym}^3\mathcal V_I,
\]
\[
  \mathcal W_{G2}
  =
  D^2\mathcal P_I+\mathcal P_I+
  \operatorname{Sym}^2\mathcal V_I\cdot\mathcal P_I.
\]
This is the source-curvature price of the rank-two plain-wall route.
\end{remark}

\begin{remark}[Two-derivative tax for \(B_2\) and background curvature]
\label{rem:rank-two-b2-two-derivative-tax}
If the canonical gamma/background curvature term \(B_2\) is included in
\[
  p=q'',
\]
then the second rank-two wall component
\[
  \Gamma_{G,2}=p''-6p+6q^2p
\]
sees two scaled derivatives of \(B_2\).  Hence the wall-variable completion
ledger must include \(B_2\), \(D B_2\), and \(D^2B_2\), either exactly or to
\(Q^{-B}\)-accuracy in the scaled norms used on \(I^*\).  A zeroth-order
approximation to \(B_2\) is not sufficient for the rank-two wall completion.
\end{remark}

\begin{theorem}[FiniteAtlasRiesz in \(p=q''\) form]
\label{thm:finite-atlas-riesz-2Y-p}
Assume on each chart-good \(I\Subset I^*\):
\begin{enumerate}
\item \(\dim\mathcal P_I\le R^C\) and \(\mathcal P_I\) has a \(Q^C\)-Riesz basis in \(L^2(I^*)\).
\item
\[
  \|D^2p_X\|_{L^2(I)}\le Q^C\|p_X\|_{L^2(I^*)},
  \qquad p_X\in\mathcal P_I.
\]
\item Affine reconstruction into \(\mathcal V_I=\mathcal A_1+\mathcal I^2\mathcal P_I\)
is polynomially controlled.
\item \(\operatorname{Sym}^3\mathcal V_I\) and
\(\operatorname{Sym}^2\mathcal V_I\cdot\mathcal P_I\) have \(Q^C\)-conditioned
quotient bases and bounded leverage.
\item All relevant chart denominators are \(Q^{-C}\)-bounded below on \(I^*\).
\item For every predetermined \(E\subset I\) with \(|E|\ge R^{-C}|I|\), Remez
holds for
\[
  \mathcal W_G
  =
  \mathcal P_I+\mathcal V_I+\operatorname{Sym}^3\mathcal V_I
\]
and
\[
  \mathcal W_{G2}
  =
  D^2\mathcal P_I+\mathcal P_I+
  \operatorname{Sym}^2\mathcal V_I\cdot\mathcal P_I.
\]
\end{enumerate}
Then
\[
  \|w\|_{H^2_{\mathrm{sc}}(I)}
  \le Q^C\|w\|_{L^2(I^*)},
  \qquad w\in\mathcal W_G.
\]
Moreover, after Lemma~\ref{lem:rank-two-source-curvature-completion},
\[
  |E|\ge R^{-C}|I|
  \quad\Longrightarrow\quad
  \int_E W_2\ge Q^{-C}\int_I W_2-Q^{-B},
\]
where
\[
  W_2:=|\Gamma_G|^2+|\Gamma_{G,2}|^2.
\]
\end{theorem}

\begin{definition}[Product-Riesz audit for the \(p=q''\) wall atlas]
\label{def:product-riesz-audit-2Y-p}
Let \(\{p_i\}\) be a basis of \(\mathcal P_I\) and let \(\{v_i\}\) be the
corresponding reconstructed basis of
\[
  \mathcal V_I=\mathcal A_1+\mathcal I^2\mathcal P_I.
\]
The product-Riesz audit consists of the raw product spanning lists
\[
  v_iv_jv_k
  \quad\text{for}\quad
  \operatorname{Sym}^3\mathcal V_I,
\]
and
\[
  v_iv_jp_k
  \quad\text{for}\quad
  \operatorname{Sym}^2\mathcal V_I\cdot\mathcal P_I.
\]
For each list, form the \(L^2(I^*)\)-Gram matrix, remove singular directions
below \(Q^{-C}\), and retain a quotient basis.  The audit passes if the retained
quotient basis is \(Q^C\)-Riesz and has \(Q^C\)-bounded leverage on \(I\).
\end{definition}

\begin{remark}[FiniteAtlasRiesz failure modes]
\label{rem:finite-atlas-riesz-failure-modes-2Y}
The theorem \(\mathrm{FiniteAtlasRiesz}_{2Y}^{p}\) fails only through explicit
finite-atlas mechanisms:
\[
\boxed{\text{\(D^2\mathcal P_I\) not controlled,}}
\]
\[
\boxed{\text{affine reconstruction unstable,}}
\]
\[
\boxed{\text{product-Riesz or product-leverage spike,}}
\]
\[
\boxed{\text{chart denominator singularity on \(I^*\),}}
\]
\[
\boxed{\text{lack of \(R^{-C}\)-Remez on predetermined cores.}}
\]
These are finite-model failures, not rank-one wall finite-type failures.
\end{remark}

\begin{proposition}[Source-divisibility of the established fixed-sector tower]
\label{prop:source-divisibility-established-fixed-sector-tower}
The established fixed-sector escape tower generated by the channels
\[
X_1YX_{2k}-X_{2k}YX_1,
\qquad
Y=yJ,
\qquad
X_1=BS,
\]
is divisible by the common source factor \(H=yq''_{\can}\).

More precisely, write
\[
X_{2k}=a_k I+d_kD+\cdots .
\]
Then
\[
X_1YX_{2k}-X_{2k}YX_1
=
yB(SJX_{2k}-X_{2k}JS).
\]
Using
\[
SJ=-D,\qquad JS=D,
\]
one obtains
\[
X_1YX_{2k}-X_{2k}YX_1
=
-yB(DX_{2k}+X_{2k}D).
\]
The \(I\)-component of \(DX_{2k}+X_{2k}D\) is \(2d_kI\).  Hence
\[
\boxed{
\pi_{\mathfrak f}
\left(
X_1YX_{2k}-X_{2k}YX_1
\right)
=
-2yB\,d_k\,I.
}
\]
Since
\[
B\propto q''_{\can},
\]
every row in this tower contains
\[
\boxed{
yB\sim yq''_{\can}=H.
}
\]
Thus increasing the anisotropy order replaces
\[
\Gamma_G,\Gamma_{G,2},\ldots
\]
by higher anisotropy factors, but it does not remove the common source factor
\(H\).
\end{proposition}

\begin{remark}[Why the anisotropy tower does not fix source collapse]
\label{rem:anisotropy-tower-does-not-fix-source-collapse}
The order-seven row detects first Gram anisotropy:
\[
A_7^{\mathfrak f}\sim H\Gamma_G.
\]
The order-nine row detects the second anisotropy on the first wall:
\[
A_9^{\mathfrak f}\sim H\Gamma_{G,2}.
\]
A higher row in the same tower would have the schematic form
\[
A_{2k+7}^{\mathfrak f}\sim H\Gamma_{G,k+1}.
\]
This improves the local layered visibility against the walls
\[
\Gamma_G=0,\quad \Gamma_{G,2}=0,\quad \ldots,
\]
but it leaves the source-collapse wall
\[
H=yq''_{\can}=0
\]
untouched.  Therefore the rank-two norm route requires either a
source-robust row not divisible by \(H\), or an external quantitative theorem
preventing source collapse.
\end{remark}

\begin{proposition}[Source-collapse countermodel for the rank-two norm route]
\label{prop:source-collapse-countermodel-rank-two-norm-route}
The source-weighted lower bound
\[
\sup_{m\in I_\rho}
|H(m)|
\max\{|\Gamma_G(m)|,|\Gamma_{G,2}(m)|\}
\ge Q^{-A}
\]
is not a consequence of the current local rank-two structure.

Indeed, consider the local jet target
\[
q_0=2,\qquad q''_{\can}=0,\qquad (q_\zeta)^{(4)}=0.
\]
Then
\[
g=\frac{q_0}{\pi}=\frac2\pi>0,
\]
and
\[
h=
\frac1\pi
\left(
\frac13q_0^3+\frac16q''_{\can}
\right)
=
\frac{8}{3\pi}>0.
\]
Moreover
\[
\Gamma_G
=
q''_{\can}-6q_0+2q_0^3
=
0-12+16
=
4,
\]
while
\[
\Gamma_{G,2}
=
(1-q_0^2)q''_{\can}-\frac16(q_\zeta)^{(4)}
=
0.
\]
Thus the layered anisotropy vector is live:
\[
\max\{|\Gamma_G|,|\Gamma_{G,2}|\}=4.
\]
However,
\[
H=yq''_{\can}
\sim
-\frac{q_0(q''_{\can})^3}{288}=0.
\]
Perturbing the target to
\[
\sup_{m\in I_\rho}|q''_{\can}(m)|\le Q^{-C}
\]
while keeping
\[
q_0(m)\asymp1,\qquad
g(m),h(m)>0,
\qquad
\max\{|\Gamma_G(m)|,|\Gamma_{G,2}(m)|\}\asymp1
\]
gives
\[
\sup_{m\in I_\rho}
|H(m)|
\max\{|\Gamma_G(m)|,|\Gamma_{G,2}(m)|\}
\ll Q^{-3C}.
\]
Choosing \(C>A/3\) defeats the desired \(Q^{-A}\) lower bound.
\end{proposition}

\begin{remark}[Interpretation of the source-collapse countermodel]
\label{rem:interpretation-source-collapse-countermodel}
Proposition~\ref{prop:source-collapse-countermodel-rank-two-norm-route}
does not assert that actual zeta realizes this jet configuration.  It says
that the present local rank-two row algebra, Gram positivity, and layered
anisotropy structure do not prevent it.  A proof must therefore supply a new
input such as
\[
\sup_{m\in I_\rho}|q''(m)|\ge Q^{-B},
\]
or a stronger source-jet / interval nonflatness theorem for actual zeta data.
\end{remark}

\begin{definition}[D4 coefficient-good core]
\label{def:d4-coeff-core-2Y}
For each routed block \(I\), choose before wall extraction a core
\[
  E_I^{\core}\subset I
\]
depending only on
\[
  H,\ c_7,\ \alpha_9,\ \beta_9^{\mathrm{tri}},\ N_{\sm},
  \quad\text{and chart data,}
\]
not on
\[
  \Gamma_G,\ \Gamma_{G,2},\ W_2,\ \text{hard atoms, or live residuals}.
\]
The core is admissible if
\[
  |E_I^{\core}|\ge R^{-C}|I|,
\]
and on \(E_I^{\core}\),
\[
  |H|\ge Q^{-A},
\]
\[
  |c_7|,\ |c_7^{-1}|,\ |\alpha_9|,\ |\alpha_9^{-1}|,\ |\beta_9^{\mathrm{tri}}|
  \le Q^A,
\]
and
\[
  N_{\sm}\ge Q^{-A}.
\]
\end{definition}

\begin{theorem}[D4 coefficient-core routing]
\label{thm:d4-coeff-core-routing-2Y}
For every selected block \(I\), suppose there is a predetermined core
\[
  E_I^{\core}\subset I,
  \qquad
  |E_I^{\core}|\ge R^{-C}|I|,
\]
chosen before wall extraction.  The core may depend on
\[
  H,\quad c_7,\quad \alpha_9,\quad \beta_9^{\mathrm{tri}},\quad N_{\sm},
  \quad\text{and chart data},
\]
but not on
\[
  \Gamma_G,\quad \Gamma_{G,2},\quad
  W_2:=|\Gamma_G|^2+|\Gamma_{G,2}|^2,
\]
hard extracted atoms, or live residuals.  Assume on \(E_I^{\core}\) that
\[
  |H|\ge Q^{-A},
\]
\[
  |c_7|,\ |c_7^{-1}|,\ |\alpha_9|,\ |\alpha_9^{-1}|,
  \ |\beta_9^{\mathrm{tri}}|
  \le Q^A,
\]
and
\[
  N_{\sm}\ge Q^{-A}.
\]
Then on the core
\[
  \left\|
  H
  \begin{pmatrix}
    c_7\Gamma_G\\
    \beta_9^{\mathrm{tri}}\Gamma_G+\alpha_9\Gamma_{G,2}
  \end{pmatrix}
  \right\|^2
  \asymp_{Q^A}
  |H|^2
  \left(
    |\Gamma_G|^2+|\Gamma_{G,2}|^2
  \right).
\]
After multiplying by \(N_{\sm}\), the same core supplies the weighted wall
\[
  N_{\sm}|H|^2
  \left(
    |\Gamma_G|^2+|\Gamma_{G,2}|^2
  \right).
\]
Thus coefficient-good routing is a core-selection statement, not a post-hoc
selection of a wall-large set.
\end{theorem}

\begin{lemma}[Core weighted plain-wall transfer]
\label{lem:core-weighted-plain-wall-transfer-2Y}
Assume admissible cores and vector Remez:
\[
  |E|\ge R^{-C}|I|
  \quad\Longrightarrow\quad
  \int_E W_2\,d\nu\ge Q^{-C}\int_IW_2\,d\nu.
\]
Then
\[
  \int_I N_{\sm}|H|^2W_2\,d\nu
  \ge Q^{-A'}\int_IW_2\,d\nu.
\]
\end{lemma}

\begin{remark}[Non-circularity of coefficient-good cores]
\label{rem:rank-two-core-transfer-noncircularity}
The core
\[
  E_I^{\core}
\]
must be selected before wall extraction.  It may depend on
\[
  H,\quad c_7,\quad \alpha_9,\quad \beta_9^{\mathrm{tri}},\quad N_{\sm},
  \quad\text{and chart data,}
\]
but not on
\[
  \Gamma_G,\quad \Gamma_{G,2},\quad
  W_2=|\Gamma_G|^2+|\Gamma_{G,2}|^2,
\]
nor on hard-extracted atoms or live residuals.  This is the transfer analogue
of the local no-atom-fitting rule.  If \(E_I^{\core}\) is chosen after seeing
\(W_2\), the weighted transfer becomes circular.
\end{remark}

\begin{remark}[Symbolic rank-two O2/O3 degree]
\label{rem:rank-two-symbolic-o2o3-degree}
The scalar O2/O3 input is degree-parametric.  After scalar component selection,
the weight is
\[
  W_\star=N_{\sm}|H|^2|G_\star|^2,
  \qquad
  G_\star\in\{\Gamma_G,\Gamma_{G,2}\}.
\]
The required branch-charge bandwidth is
\[
  d_{\max}^{(2)}
  =
  d_N+2d_H+2\max(d_{\Gamma_G},d_{\Gamma_{G,2}}).
\]
One should not hard-code \(d_{\max}=14\) unless this degree audit has been
performed in the active convention.  If \(H\) is a smooth coefficient or
normalization denominator absorbed into the measure, then \(d_H=0\).  If \(H\)
is branch-bearing, its fixed branch degree must be counted.  If \(H\) is
locally adjustable, the present D4 normalization is not valid without
renormalization or an alternate chart.
\end{remark}

\begin{lemma}[Scalar component selection]
\label{lem:rank-two-scalar-component-selection}
If
\[
  \int N_{\sm}|H|^2\left(|\Gamma_G|^2+|\Gamma_{G,2}|^2\right)d\nu\ge Q^{-A},
\]
then there exists a fixed component
\[
  G_\star\in\{\Gamma_G,\Gamma_{G,2}\}
\]
chosen at the smooth aggregate stage such that
\[
  \int N_{\sm}|H|^2|G_\star|^2\,d\nu\ge Q^{-A'}.
\]
\end{lemma}

\begin{definition}[Rank-two wall degree]
\label{def:rank-two-wall-degree}
For
\[
  G_\star\in\{\Gamma_G,\Gamma_{G,2}\},
\]
define
\[
  d_0(G_\star):=d_N+2d_H+2d_{G_\star}.
\]
Set
\[
  d_{\max}^{(2)}
  :=
  d_N+2d_H+2\max(d_{\Gamma_G},d_{\Gamma_{G,2}}).
\]
Conservatively,
\[
  d_{\Gamma_G},d_{\Gamma_{G,2}}\le 3d_q,
\]
so
\[
  d_{\max}^{(2)}\le d_N+2d_H+6d_q.
\]
If \(H\) is smooth-renormalizable and carries no branch charge, set \(d_H=0\).
\end{definition}

\begin{theorem}[Scalar O2/O3 interface for rank two]
\label{thm:rank-two-scalar-o2o3-interface}
Assume weighted rank-two plain wall mass has been transferred, and assume
\[
  \mathbb C_{d_{\max}^{(2)}}[y]\subset V_{\mathrm{tr}}.
\]
Assume the scalar O2/O3 near-pure separation holds uniformly for every fixed
\(d\le d_{\max}^{(2)}\):
\[
  \operatorname{dist}_{L^2(\mu_{\mathrm{sub}})}
  \left(
  1,
  \operatorname{span}\{R_{\mathrm{fold}},R_{\mathrm{fold}}^2,\ldots,R_{\mathrm{fold}}^{d}\}
  \right)^2
  \ge Q^{-A(d_{\max}^{(2)})}.
\]
Then the selected scalar wall
\[
  W_\star=N_{\sm}|H|^2|G_\star|^2
\]
is a valid bounded-charge O2/O3 input.
\end{theorem}

\begin{remark}[Rank-two no-go and guardrail register]
\label{rem:rank-two-no-go-register}
The following rules are active guardrails for the rank-two branch.

\begin{enumerate}
\item D4 supplies only the plain wall
\[
  H(\Gamma_G,\Gamma_{G,2}).
\]
The derivative family
\[
  \{\Gamma_G,D\Gamma_G,D^2\Gamma_G,\Gamma_{G,2}\}
\]
is analytic auxiliary data only.  It is not a D4/O2/O3 wall unless a separate
\(\mathrm{DerivativeD4Channels}_{2Y}\) or block-difference derivative-channel
theorem is proved.

\item All branch classification is made in row-normal coordinates:
\[
  \ell_9^Y=\omega_Y-\omega_Y(e_7)\ell_7,
  \qquad
  \beta_9^{\mathrm{tri}}=\beta_9-\omega_Y(e_7)c_7.
\]
Any rank-one fallback must use
\[
  \rho_\Gamma^{\mathrm{tri}}
  =
  \frac{\beta_9^{\mathrm{tri}}-\rho_Yc_7}{\alpha_9}.
\]

\item The rank-two selector is the quotient rank
\[
  r_Y=
  \operatorname{rank}
  \left(
    \Pi_{\mathbf C^2/B_Y}M_{\mathrm{D4}}^{\mathrm{tri}}
  \right),
\]
not the raw rank of formal D4 carriers.  Therefore formal \(C_7\to Y12\)
connections, fixed basepoint offsets, and coherent amplitude remnants do not
reduce \(r_Y\) unless they produce actual/local columns in \(B_Y\).

\item Formal \(C_7\) carriers and formal \(C_7\to Y12\) connections do not
contribute to \(B_Y\) unless actualized by a genuine actual-source, exact
nuisance, or local fitting direction.

\item Exact nuisance must be row-dead and live-pair preserving before it is
placed in the denominator.  Coherent nuisance and residual nuisance are routed
to \(\mathcal C_{2Y}^{\mathrm{arith}}\) and \(R_X^{2Y}\), respectively.

\item \(E_2\) and \(R_{\mathrm{lab}}\) affect \(B_Y\) only through atom-local live-fiber
fitting.  Fixed, coherent, and residual terms do not define local columns of
\(B_Y\).

\item Amplitude projection occurs only after \(B_Y=0\) is proved.  A projected
wall is not fed back into the local quotient, the row definitions, the live
pair, or the D4 chart unless its branch complexity has been separately
controlled.

\item The weighted-transfer core \(E_I^{\core}\) must be selected from
coefficient, smooth, and chart data before observing
\[
  W_2=|\Gamma_G|^2+|\Gamma_{G,2}|^2.
\]
Post-hoc wall-dependent core selection is not allowed.

\item The scalar O2/O3 theorem is degree-parametric.  Do not hard-code a
specific value such as \(14\) until the rank-two wall degree audit has fixed
\[
  d_{\max}^{(2)}
  =
  d_N+2d_H+2\max(d_{\Gamma_G},d_{\Gamma_{G,2}}).
\]

\item An \(\alpha\)-twisted determinant average is not proof-critical unless
\(\mathrm{AlphaAdmissibility}^{+}\) is proved.  In particular, one needs
regular zero persistence, badness persistence, twist nondegeneracy, Stack
uniformity, and \(\alpha\)-D4 compatibility:
\[
  [\mathsf C_{Y12}]\ne0
  \quad\text{in}\quad
  \mathcal Q_9^{\mathrm{forb},\alpha}.
\]

\item A pure branch determinant is not a scalar O2/O3 theorem unless it
survives inside the full determinant of the canonical matrix:
\[
  \int |D_{\mathrm{full}}|^2\,d\nu
  \ge
  Q^{-A}\int |D_{\mathrm{dom}}|^2\,d\nu.
\]
The Stack measure is
\[
  d\nu=F_1F_{\mathrm{wall}}^{\mathrm{forb}}N\,d\rho,
\]
so neutral-measure branch orthogonality is not enough.

\item A forbidden \(\Lambda\) must be non-functorial on at least one D4 live
carrier.  If every natural carrier lift is functorial through order \(9\), then
\[
  \Lambda[\mathsf C_7]=\Lambda[\mathsf C_{Y12}]=0,
\]
and finite D4 has no contradiction channel.

\item Carrier-quotient D4 is a visibility theorem, not a contradiction theorem,
until a forbidden channel \(\Lambda\) is supplied.  The finite audit is
\[
  \operatorname{rank}
  \left(
    \Lambda E_{\mathrm{car}}L_{\mathrm{clean}}
  \right)\ge1.
\]
\end{enumerate}
\end{remark}

\begin{remark}[Interval-local source collapse]
\label{rem:interval-local-source-collapse}
The pointwise source-collapse target becomes an interval-local obstruction
because the proof requires a supremum over an admissible packet interval
\[
I_\rho=[m_0-\Delta,m_0+\Delta].
\]
The scale audit gives the natural estimates
\[
\Omega_{\rm eff}(q'')\sim Q,
\qquad
\Delta\sim Q^{-1},
\qquad
\Omega_{\rm eff}(q'')\Delta\sim1.
\]
Thus \(q''\) has only \(O(1)\) effective oscillatory complexity on the
natural microscopic packet interval.

Let
\[
z(m):=q''(m).
\]
To force
\[
\sup_{|m-m_0|\le c/Q}|z(m)|\le Q^{-C},
\]
it is enough to flatten finitely many derivatives:
\[
z^{(r)}(m_0)\approx0,
\qquad
0\le r\le R.
\]
A Taylor--Bernstein estimate gives, schematically,
\[
\sup_{|u|\le c/Q}|z(m_0+u)|
\lesssim
Q^{M_0}
\frac{C_c^{R+1}}{(R+1)!}.
\]
Hence a safe choice is
\[
\boxed{
R\asymp (C+M_0)\frac{\log Q}{\log\log Q},
}
\]
or more conservatively \(R=O(\log Q)\).  This order is far below the
recurrence-danger threshold \(Q/\log Q\).  Therefore the interval-local
source-collapse countermodel is not defeated by interval length or bandwidth
alone.
\end{remark}

\begin{remark}[Fixed-prime and AFE derivative-flat surrogates]
\label{rem:fixed-prime-afe-derivative-flat-surrogates}
A fixed-prime surrogate for the source-collapse mechanism is
\[
L_P(m)
=
\sum_{p\le P}(\log p)p^{-1/2}e^{-im\log p}.
\]
Define
\[
J_0(m)=B_0+\Re L_P(m),
\]
\[
J_2(m)=B_2+\Re L_P''(m),
\]
\[
J_4(m)=B_4+\Re L_P^{(4)}(m).
\]
The model residuals are
\[
R_1=J_2-6J_0+2J_0^3,
\]
\[
R_2=(1-J_0^2)J_2-\frac16J_4,
\]
and
\[
H_{\rm model}=-\frac{J_0J_2^3}{288}.
\]
The interval-local source-collapse target is
\[
J_2^{(r)}(m_0)\approx0,\qquad 0\le r\le R,
\]
with side constraints such as
\[
J_0(m_0)\approx2,\qquad J_4(m_0)\approx0.
\]
At the exact target
\[
(J_0,J_2,J_4)=(2,0,0)
\]
one has
\[
R_1=4,\qquad R_2=0,
\]
and
\[
g=\frac2\pi>0,\qquad h=\frac{8}{3\pi}>0.
\]
Thus the source is collapsed while anisotropy and Gram positivity remain live.

The phase-relaxed finite-prime problem has no rank obstruction: the rows
\[
J_2^{(r)}(m_0)
\]
are governed by powers of the frequencies
\[
\lambda_j=\log p_j
\quad\text{or}\quad
\lambda_j=k\log p,
\]
so the derivative-jet matrix has a Vandermonde-type structure.  The remaining
arithmetic issue is whether the actual one-parameter prime-log flow
\[
m_0\mapsto (m_0\log p_j)_j\pmod{2\pi}
\]
recurs to such a derivative-flat target.  The expected codimension is
approximately
\[
K=R+3,
\]
so a heuristic recurrence cost at tolerance \(\varepsilon=Q^{-C}\) is
\[
T_{\rm hit}\sim Q^{CK}.
\]
For \(R=O(\log Q)\), this is
\[
T_{\rm hit}
=
\exp(O((\log Q)^2))
\ll e^Q.
\]
Thus scale considerations do not presently block the fixed-prime
source-collapse surrogate.

The phase-relaxed derivative-flat target can be sharpened further.  In the
fixed-amplitude prime model
\[
L_P(m)=\sum_{p\le P}a_p e^{-im\log p},
\qquad
a_p=\frac{\log p}{\sqrt p},
\]
one obtains
\[
(J_0,J_2,J_4)
=
\sum_{p\le P}
x_p a_p
\left(
1,\,-(\log p)^2,\,(\log p)^4
\right),
\qquad
x_p=\cos(m\log p).
\]
Thus the phase-relaxed attainable jet values form a fixed-amplitude prime
zonotope.  The source-collapse target
\[
(J_0,J_2,J_4)=(2,0,0)
\]
is compatible with good-patch positivity:
\[
g=\frac2\pi>0,\qquad h=\frac{8}{3\pi}>0,
\]
and has live anisotropy
\[
R_1=4,\qquad R_2=0.
\]
Finite-dimensional feasibility tests for the moment system
\[
\sum_{p\le P}a_px_p=2,
\qquad
\sum_{p\le P}a_p(\log p)^{2k}x_p=0
\]
show interior solutions for low and moderate finite orders.  For example:
\[
\begin{array}{c|c|c|c}
\text{moments imposed} & P & \#\{p\le P\} & \min \max_{p\le P}|x_p|\\
\hline
0,2,4 & 100 & 25 & 0.699\\
0,2,4,6 & 300 & 62 & 0.679\\
0,2,4,6,8 & 1000 & 168 & 0.637\\
0,2,4,6,8,10 & 1000 & 168 & 0.888\\
0,2,4,6,8,10,12 & 5000 & 669 & 0.718.
\end{array}
\]
So the target lies in the interior of the fixed-prime phase-relaxed cosine
zonotope, not on a boundary.  The remaining issue is therefore not amplitude
feasibility.  It is the one-parameter recurrence problem
\[
m_0\log p\equiv \theta_p\pmod{2\pi}
\]
for selected phases \(\theta_p\).

AFE-compatible models add prime powers, smoothing, a dual sum, and gamma
background terms.  These shift or couple the target, but no algebraic
obstruction is presently known that forces
\[
\max_{0\le r\le R}|J_2^{(r)}(m_0)|\ge Q^{-B}.
\]
Therefore the exact missing Gate~1 theorem is actual-zeta source-curvature
observability, and it must later be paired with Gate~2 source--anisotropy
reintegration.  It is not another local row identity.
\end{remark}

\begin{remark}[Prime-log recurrence to derivative-flat tubes]
\label{rem:prime-log-recurrence-derivative-flat-tubes}
Let
\[
\lambda=(\log p_1,\ldots,\log p_N),
\qquad
\Theta(m)=m\lambda\pmod{2\pi}.
\]
Let \(\mathcal T_{\flat}(\varepsilon)\subset\mathbb T^N\) denote the
\(\varepsilon\)-tube around the derivative-flat target
\[
|J_2^{(r)}|\le\varepsilon,\qquad 0\le r\le R,
\]
with side constraints
\[
|J_0-2|\le\varepsilon,
\qquad
|J_4|\le\varepsilon.
\]
If this target is regular, then
\[
\mu(\mathcal T_{\flat}(\varepsilon))\asymp \varepsilon^K,
\qquad
K\sim R+3.
\]
This is a positive-measure analytic tube, not a prescribed point in
\(\mathbb T^N\).

For polynomial tolerance
\[
\varepsilon=Q^{-C},
\]
the random/equidistribution hitting scale is
\[
T_{\rm hit}\sim \varepsilon^{-K}=Q^{CK}.
\]
For the derivative orders relevant to Taylor--Bernstein source flattening,
\[
R\asymp (C+M_0)\frac{\log Q}{\log\log Q}
\]
or even \(R=O(\log Q)\), this gives
\[
T_{\rm hit}
=
\exp\!\left(O\!\left(\frac{(\log Q)^2}{\log\log Q}\right)\right)
\quad\text{or}\quad
T_{\rm hit}=\exp(O((\log Q)^2)).
\]
Both are far below the available height scale \(e^Q\).

Thus recurrence cost does not presently block the interval-local source-collapse
surrogate.  A proof of the desired source-weighted lower bound would need a
nonstandard quantitative nonrecurrence theorem for the prime-log flow against
these derivative-flat tubes, or actual-zeta Gate~1 source-curvature
observability together with Gate~2 source--anisotropy reintegration.
\end{remark}

\begin{remark}[Current Gate--1 attack ledger: finite-jet and Wronskian forms]
\label{rem:current-gate-one-finite-jet-wronskian-ledger}
The current Gate--1 target is
\[
\int_{I_\rho}|q''_{\can}(m)|^2\,dm\ge Q^{-B}|I_\rho|.
\]
In normalized coordinates
\[
m=m_0+\lambda x,\qquad
F_\lambda(x)=\lambda^3q''_{\can}(m_0+\lambda x),
\]
it is enough, subject to the finite-jet-to-energy transfer, to prove a
finite-jet anti-design estimate of the form
\[
\max_{0\le r\le R(Q)}
|\partial_x^rF_\lambda(x_0)|
\ge Q^{-B_{\rm jet}}\lambda^3.
\]
The natural depth is
\[
R(Q)\asymp\frac{\log Q}{\log\log Q},
\]
with \(R=O(\log Q)\) as a safer but more expensive depth.

A promising formulation is a normalized Wronskian/transversality package.  Let
\[
G_0=Z_\alpha,\qquad
G_{j+1}=\Delta^j\partial_m^j q''_{\can},
\quad 0\le j\le R.
\]
The associated normalized Wronskian is
\[
\widetilde{\mathcal W}_R(m,\alpha)
=
\det\bigl(\mathcal D^\ell G_j(m,\alpha)\bigr)_{0\le j,\ell\le R+1},
\qquad
\mathcal D=\Delta\partial_m.
\]
A polynomial lower bound
\[
|\widetilde{\mathcal W}_R(m,\alpha)|\ge Q^{-A}
\]
at admissible centers would imply finite-jet anti-design and hence Gate~1.

At present this route is not Green.  The \(Z_\alpha\)-symbol appears to add
genuine rank to the \(q''_{\can}\)-jet system, but the Wronskian expands into
a fixed-coefficient subset-product exponential sum.  Cauchy--Binet
nonvanishing proves only that the determinant is not identically zero; it does
not give the required pointwise lower bound on the actual prime-log orbit.
Thus the current missing theorem is a fixed actual-zeta determinant or
nonrecurrence lower bound, not another formal rank calculation.
\end{remark}

\begin{remark}[\(Z_\alpha\)-symbol transversality subproblem]
\label{rem:zalpha-symbol-transversality-subproblem}
The first finite-dimensional test for the Wronskian route is whether the
off-line equation contributes a row independent from the \(q''_{\can}\)-jet
rows.  In a simplified AFE block, write
\[
Z_\alpha(m)\sim\sum_n z_n(\alpha)e^{-im\log n},
\]
and
\[
\partial_m^r q''_{\can}(m)
\sim
\sum_n c_n(-i\log n)^r e^{-im\log n}.
\]
On an active block, the relevant symbolic test is whether
\[
\frac{z_n(\alpha)}{c_n}
\]
is approximable by a polynomial in \(\log n\) of degree at most \(R\).  The
expected model is
\[
\frac{z_n(\alpha)}{c_n}
\approx
\frac{n^{-\alpha}}{(\log n)^3}
\frac{V_\alpha(n/N)}{W_2(n/N)},
\]
or, with \(u=\log n\),
\[
f_\alpha(u)\approx e^{-\alpha u}u^{-3}A_\alpha(u).
\]
This is not polynomial-like on a thick active block unless the weights are
degenerate.  Thus \(Z_\alpha\) likely adds genuine symbolic rank.  The missing
step is quantitative: one needs a lower bound
\[
\inf_{P\in\mathcal P_{\le R}}
\max_{n\in\mathcal B}
\left|
\frac{z_n(\alpha)}{c_n}-P(\log n)
\right|
\ge Q^{-C}
\]
for the actual AFE coefficients and active block.  Symbolic transversality
alone still does not imply the required pointwise Wronskian lower bound.
\end{remark}

\begin{remark}[Derivative-order and error budget for the Wronskian route]
\label{rem:wronskian-derivative-order-error-budget}
The normalized Wronskian route is sensitive to derivative order.  Since
\[
G_{j+1}=\Delta^j\partial_m^j q''_{\can},
\]
and the Wronskian differentiates \(G_j\) up to order \(R+1\), the argument
requires derivatives of \(q''_{\can}\) through order approximately \(2R+1\),
equivalently derivatives of \(q_\zeta\) through order \(2R+3\).  The
\(Z_\alpha\)-row requires derivatives through order \(R+1\).

Therefore the canonical error term in
\[
q''_{\can}=B_2+P_2+D_2+E_2
\]
must be controlled through the same derivative order.  The safe choice
\[
R\asymp\frac{\log Q}{\log\log Q}
\]
keeps the normalized factorial and AFE derivative losses polynomial.  Taking
\(R=O(\log Q)\) is dangerous unless a separate lemma removes the resulting
factorial and differentiated-error losses.
\end{remark}

\begin{remark}[Zero-side same-scale source flattening]
\label{rem:zero-side-same-scale-source-flattening}
The source-collapse mechanism is not merely a prime-side artefact.  On the
zero side, one can seek same-scale positive packet configurations for which
the aggregate source curvature is flat to high finite order.

For a centered Cauchy packet
\[
k_c(t)=\frac{c}{c^2+4(t-m)^2},
\]
the centered jets are
\[
k_c(m)=c^{-1},
\]
\[
k_c''(m)=-8c^{-3},
\]
and
\[
k_c^{(4)}(m)=384c^{-5}.
\]
A single centered packet cannot realize
\[
q''=q^{(4)}=0
\]
with \(q_0>0\).  However, once off-center same-scale packets and positive
linear combinations are allowed, the problem becomes a finite moment-matching
problem.

In particular, one can use same-scale packet families to impose finitely many
conditions of the form
\[
q^{(2)}(m_0)=q^{(3)}(m_0)=\cdots=q^{(R+2)}(m_0)=0
\]
while preserving
\[
q_0(m_0)>0.
\]
At the target
\[
q_0=2,\qquad q''=0,\qquad q^{(4)}=0,
\]
Gram positivity and live anisotropy remain:
\[
g=\frac2\pi>0,
\qquad
h=\frac{8}{3\pi}>0,
\qquad
\Gamma_G=4.
\]
Thus same-scale packet geometry alone does not provide Gate~1 source-curvature
observability, and a fortiori does not provide the full two-gate package.  A
proof must use information not already contained in finite same-scale positive
packet geometry.
\end{remark}

\begin{remark}[Local-kernel notation in the same-scale model]
\label{rem:local-kernel-notation-same-scale-model}
In the same-scale packet model below, the unadorned \(q,q'',q^{(4)}\) denote
the local positive source-kernel aggregate being constructed, not the full
actual-zeta canonical curvature \(q''_{\can}\).  The model is used only as a
local/Hadamard-side no-go: it shows that finite same-scale positive packet
geometry alone does not force nonflatness of the canonical actual-zeta
curvature.
\end{remark}

\begin{theorem}[Continuous same-scale source-flattening model]
\label{thm:continuous-same-scale-source-flattening-model}
Fix \(d>0\).  For
\[
c=d\sin\theta,\qquad x=d\cos\theta,\qquad 0<\theta<\pi,
\]
consider the same-distance packet kernel
\[
K_{\theta,d}(u)
=
\frac{c}{c^2+(x+2u)^2}.
\]
Let
\[
d\mu(\theta)=\frac{2}{\pi}\sin\theta\,d\theta.
\]
Then, for \(|2u/d|<1\),
\[
\boxed{
\int_0^\pi K_{\theta,d}(u)\,d\mu(\theta)=\frac1d.
}
\]
Consequently,
\[
q_\mu(u):=\int_0^\pi K_{\theta,d}(u)\,d\mu(\theta)
\]
satisfies
\[
q_\mu(0)=\frac1d>0,
\qquad
q_\mu^{(r)}(0)=0
\quad (r\ge1).
\]
In particular,
\[
q_\mu^{(2)}(0)=q_\mu^{(4)}(0)=q_\mu^{(6)}(0)=\cdots=0.
\]
\end{theorem}

\begin{proof}
Put
\[
t=\frac{2u}{d}.
\]
Then
\[
K_{\theta,d}(u)
=
\frac1d
\frac{\sin\theta}{1+2t\cos\theta+t^2}.
\]
Hence
\[
\int_0^\pi K_{\theta,d}(u)\,d\mu(\theta)
=
\frac1d
\int_0^\pi
\frac{2}{\pi}
\frac{\sin^2\theta}{1+2t\cos\theta+t^2}
\,d\theta.
\]
For \(|t|<1\), the elementary Poisson-kernel identity
\[
\int_0^\pi
\frac{2}{\pi}
\frac{\sin^2\theta}{1+2t\cos\theta+t^2}
\,d\theta
=
1
\]
gives the claim.  Therefore \(q_\mu(u)=1/d\) on a neighborhood of the
origin, and all positive-order derivatives at the origin vanish.
\end{proof}

\begin{corollary}[Finite same-scale source flattening to arbitrary finite jet order]
\label{cor:finite-same-scale-source-flattening}
For every \(M\ge1\) and every \(\epsilon>0\), there exist finitely many angles
\[
0<\theta_1,\ldots,\theta_N<\pi
\]
and positive weights \(w_i>0\) such that, with
\[
c_i=d\sin\theta_i,\qquad x_i=d\cos\theta_i,
\]
one has
\[
x_i^2+c_i^2=d^2
\]
for every \(i\), and
\[
\left|
\sum_i w_iK_{\theta_i,d}^{(r)}(0)
\right|
\le \epsilon,
\qquad
1\le r\le M,
\]
while
\[
\sum_i w_iK_{\theta_i,d}(0)>0.
\]
Equivalently, finite positive same-distance packet clusters can make
\[
q^{(r)}(0),\qquad 1\le r\le M,
\]
arbitrarily small while retaining
\[
q_0=q(0)>0.
\]
In particular, taking \(M\ge R+2\) gives the finite source-flatness target
\[
(q'')^{(r)}(0)\approx0,
\qquad
0\le r\le R.
\]
\end{corollary}

\begin{proof}
The continuous measure \(\mu\) in
Theorem~\ref{thm:continuous-same-scale-source-flattening-model} satisfies
\[
q_\mu(u)=\int K_{\theta,d}(u)\,d\mu(\theta)=\frac1d
\]
near \(u=0\).  Hence all derivatives \(q_\mu^{(r)}(0)\), \(r\ge1\), vanish.

Apply Tchakaloff's theorem, or finite-dimensional Carath\'eodory reduction, to
the positive measure \(\mu\) and the finite list of continuous functions
\[
K_{\theta,d}(0),\quad
K_{\theta,d}^{(1)}(0),\quad
K_{\theta,d}^{(2)}(0),\ldots,K_{\theta,d}^{(M)}(0).
\]
This gives a finite positive quadrature rule reproducing the required finite
moment vector on the closed support.  If an atom occurs at \(\theta=0\) or
\(\theta=\pi\), replace \(\mu\) by its restriction to
\[
[\varepsilon,\pi-\varepsilon]
\]
and renormalize.  Since
\[
d\mu(\theta)=\frac{2}{\pi}\sin\theta\,d\theta
\]
has vanishing endpoint mass, the corresponding finite moment vector converges
to the original one as \(\varepsilon\downarrow0\).  Applying finite quadrature
on \([\varepsilon,\pi-\varepsilon]\) and choosing \(\varepsilon\) sufficiently
small gives interior atoms \(0<\theta_i<\pi\) and the stated finite-order
approximation.  This approximation is the form needed below for interval
source-collapse, where the target tolerance is polynomial in \(Q^{-1}\).
\end{proof}

\begin{remark}[Uniform same-scale caveat]
\label{rem:uniform-same-scale-caveat}
The construction above is same-scale in the strong fixed-distance sense
\(x_i^2+c_i^2=d^2\).  If a later argument requires the additional uniform
angular condition
\[
\epsilon\le \frac{c_i}{d}\le 1-\epsilon
\]
with \(\epsilon>0\) independent of \(M\), then one still needs a strengthened
annular positive-cone lemma.  The present result is already enough to show
that finite source-jet flattening is not excluded by positive same-distance
packet geometry.  Thus any proof-closing lower bound for \(q''\) must use
additional actual-zeta arithmetic, not merely same-scale positivity.
\end{remark}

\begin{remark}[Source-kernel versus actual-zeta packet caveat]
\label{rem:source-kernel-versus-actual-zeta-packet-caveat}
The same-scale flattening construction is a statement about positive local
Cauchy source kernels at a fixed distance scale.  It should not be read as a
claim that actual zeta zeros can realize the same quadrature configuration.
Actual zeta imposes functional-equation symmetry, conjugation, discreteness,
zero-counting constraints, and possibly Euler-product arithmetic not present
in the continuous source-kernel model.

The point of
Theorem~\ref{thm:continuous-same-scale-source-flattening-model} and
Corollary~\ref{cor:finite-same-scale-source-flattening} is therefore negative
but precise: local same-scale positivity does not by itself force
source-curvature nonflatness.  Any proof of
\[
\sup_{m\in I_\rho}|q''_{\can}(m)|\ge Q^{-B}
\]
or of the stronger interval estimate
\[
\int_{I_\rho}|q''_{\can}(m)|^2\,dm\ge Q^{-B}|I_\rho|
\]
must use actual-zeta arithmetic or a stronger packet-realizability theorem,
not merely same-scale positive kernel geometry.
\end{remark}

\begin{corollary}[Interval-local source collapse from finite quadrature]
\label{cor:interval-local-source-collapse-from-quadrature}
All unadorned \(q,q'',q^{(4)}\) in this corollary denote the local
same-scale source-kernel aggregate from
Remark~\ref{rem:local-kernel-notation-same-scale-model}, not the actual-zeta
canonical curvature.  In the actual-zeta notation, this is a model for the
hypothetical regime \(q''_{\can}\approx0\), but it is not a construction of
actual zeta data.

Fix \(A>0\).  Assume the packet scale is microscopic,
\[
d\asymp Q^{-1},
\]
and let
\[
I=[m_0-c/Q,m_0+c/Q]
\]
with \(c>0\) fixed sufficiently small.  For every \(C>A/3\), a finite
same-distance positive packet cluster can be chosen so that
\[
\sup_{m\in I}|q''(m)|\le Q^{-C},
\]
while
\[
q_0(m)\asymp1,\qquad g(m)>0,\qquad h(m)>0,
\]
and
\[
\max\{|\Gamma_G(m)|,|\Gamma_{G,2}(m)|\}\asymp1
\]
on \(I\), after a harmless normalization of the total weight.

Consequently, at leading source order,
\[
H_{\rm loc}=yq''\sim-\frac{q_0(q'')^3}{288}
\]
satisfies
\[
\sup_{m\in I}|H_{\rm loc}(m)|
\max\{|\Gamma_G(m)|,|\Gamma_{G,2}(m)|\}
\ll Q^{-3C},
\]
and therefore can be made smaller than \(Q^{-A}\).
\end{corollary}

\begin{proof}
By Corollary~\ref{cor:finite-same-scale-source-flattening}, choose the packet
cluster so that, for a prescribed polynomial tolerance,
\[
q^{(r)}(m_0)=O(Q^{-C'})
\qquad
1\le r\le M,
\]
with \(M\) large enough relative to the desired interval bound.  In particular,
\[
(q'')^{(r)}(m_0)=O(Q^{-C'})
\qquad
0\le r\le R
\]
for the required finite order \(R\).

The standard Taylor--Bernstein estimate at scale \(d\asymp Q^{-1}\) then gives
\[
\sup_I |q''|\le Q^{-C}
\]
after taking \(C'\) and \(M\) sufficiently large.  For example, the usual
Taylor bookkeeping permits
\[
M\asymp \frac{C\log Q}{\log\log Q}
\]
or the safer choice \(M=O(\log Q)\), with constants depending on the fixed
interval parameter and same-scale annulus.
Normalizing the total mass so that \(q_0\) is close to \(2\), one obtains
\[
g=\frac{q_0}{\pi}>0,
\qquad
h=\frac1\pi\left(\frac13q_0^3+\frac16q''\right)>0,
\]
and
\[
\Gamma_G=q''-6q_0+2q_0^3=4+O(Q^{-C}).
\]
Hence the anisotropy factor remains bounded below while the source factor
collapses cubically in \(q''\).

This is a finite-order local obstruction model.  It shows that source-collapse
is compatible with positive same-scale packet geometry to every prescribed
polynomial accuracy.  It is not a claim that actual zeta zeros realize such a
quadrature configuration.
\end{proof}

\begin{proposition}[No-go for the current rank-two norm route]
\label{prop:no-go-current-rank-two-norm-route}
The current rank-two norm route cannot close RH from the established local
row theory and same-scale positive packet geometry alone.

Here the displayed \(q''\) and \(H=yq''\) first refer to the local
same-scale source-kernel model of
Corollary~\ref{cor:interval-local-source-collapse-from-quadrature}.  The
proposition's conclusion for the actual rank-two route is obtained by
comparison: if actual zeta permits the analogous canonical regime
\[
\sup_{m\in I_\rho}|q''_{\can}(m)|\le Q^{-C},
\]
then the actual source factor
\[
H=yq''_{\can}
\]
collapses in the same way.

More precisely, the following data are compatible with the current local
rank-two structure:
\[
q_0\asymp1,
\qquad
g,h>0,
\qquad
\max\{|\Gamma_G|,|\Gamma_{G,2}|\}\asymp1,
\]
and
\[
\sup_{m\in I_\rho}|q''(m)|\le Q^{-C}.
\]
For such data,
\[
H=yq''
\sim
-\frac{q_0(q'')^3}{288}
\]
satisfies
\[
\sup_{m\in I_\rho}|H(m)|\ll Q^{-3C}.
\]
Consequently,
\[
\sup_{m\in I_\rho}
|H(m)|
\max\{|\Gamma_G(m)|,|\Gamma_{G,2}(m)|\}
\ll Q^{-3C}.
\]
Taking \(C>A/3\) violates the desired \(Q^{-A}\) source-weighted lower bound.

Thus the local rank-two norm route requires the actual-zeta Gate~1 theorem
excluding canonical source-flat intervals,
\[
\int_{I_\rho}|q''_{\can}(m)|^2\,dm\ge Q^{-B}|I_\rho|,
\]
and then the Gate~2 source--anisotropy reintegration theorem connecting that
source mass back to the row vector.
By Corollary~\ref{cor:interval-local-source-collapse-from-quadrature}, the
source-flat interval model is compatible with finite positive same-distance
packet quadrature to arbitrary finite order.  Therefore the missing input
cannot be a finite local packet-geometry lemma; it must use actual zeta
arithmetic, a source-robust row, or a genuinely global transport/projective
constraint.

This proposition is not a counterexample to \(\RH\) or to actual zeta; it is a
no-go for closing the present rank-two norm route from local packet geometry
alone.
\end{proposition}

\begin{remark}[What would be needed to save the norm route]
\label{rem:needed-source-jet-nonflatness-theorem}
A necessary new input is actual-zeta source-curvature observability, for
example
\[
\boxed{
\int_{I_\rho}|q''_{\can}(m)|^2\,dm\ge Q^{-B}|I_\rho|.
}
\]
A finite-jet version would be
\[
\boxed{
\max_{0\le r\le R(Q)}
|\partial_m^r q''_{\can}(m_0)|
\ge Q^{-B}.
}
\]
The currently available local framework does not prove either statement.  The
distinguished packet contributes nontrivially to \(q''_{\can}\), but
rest/core/background/dual data can cancel that contribution.  Same-scale
selection, strict core gap, endpoint separation, Gram positivity, finite
Hadamard symmetry, and the existing fixed-sector row algebra do not prevent
such cancellation.

Even this source-curvature input is not sufficient by itself.  The norm route
also needs source--anisotropy reintegration:
\[
\boxed{
\int_{I_\rho}
|q''_{\can}|^2(\Gamma_G^2+\Gamma_{G,2}^2)\,dm
\ge Q^{-A}|I_\rho|.
}
\]
Thus the required new input is either the two-gate pair above, a single
stronger theorem implying both, a genuinely source-robust packet-visible row,
or a global projective/trace/transport theorem.
\end{remark}

\begin{remark}[No-go ledger after the source-collapse push]
\label{rem:no-go-ledger-after-source-collapse-push}
The current push rules out several tempting repairs of the source-collapse
obstruction.

\begin{enumerate}
\item \emph{Higher anisotropy rows.}
The tower
\[
X_1YX_{2k}-X_{2k}YX_1
\]
is divisible by \(H=yq''_{\can}\).  Higher order can detect higher Gram anisotropies,
but it does not remove source collapse.

\item \emph{Ambient \(D\)-sector at order seven.}
At \(g=h\), the ambient row
\[
T_7^{(1)}=-2yBpD
\]
survives the fixed-sector Gram-isotropy failure, but still contains
\[
yB\sim yq''_{\can}.
\]
It therefore does not solve Regime~B source collapse.

\item \emph{Transported derivatives.}
If
\[
A_j=HG_j,
\]
then
\[
\nabla A_j=(\nabla H)G_j+H\nabla G_j.
\]
Since
\[
H\sim-\frac{q_0(q''_{\can})^3}{288},
\]
one has
\[
\nabla H=O((q''_{\can})^2)
\]
under good-patch derivative control.  Thus finitely many transported
derivatives do not presently defeat an interval where \(q''_{\can}\) and finitely
many of its derivatives are small.

\item \emph{Projectivization.}
The source-normalized projective ratio
\[
[A_7:A_9]=[\Gamma_G:\Gamma_{G,2}]
\]
formally cancels \(H\), but it loses norm-scale information, is singular at
\[
\Gamma_G=\Gamma_{G,2}=0,
\]
and requires correction-row control plus a new global projective index or
transport theorem.  No such theorem is presently known.

\item \emph{Correction rows.}
Line, background/core, and normalization/subtraction rows are not currently
certified as source-robust off-line packet detectors.  They are either
gauge/baseline/normalization data or must be carried as enlarged projective
coordinates.  They do not provide a clean replacement for the source-divisible
rows.

\item \emph{Explicit formula / prime-side coercivity.}
The explicit formula sees the linear ingredients
\[
q_0,\quad q''_{\can},\quad (q_\zeta)^{(4)},
\]
but the residual
\[
(yq''_{\can})^2(\Gamma_G^2+\Gamma_{G,2}^2)
\]
is nonlinear.  Tensorizing the explicit formula produces oscillatory
multi-prime sums with no identified positivity or coercive lower bound.
\end{enumerate}

Accordingly, the current route is non-closing unless one proves the two-gate
actual-zeta package
\[
\int_{I_\rho}|q''_{\can}|^2\,dm\ge Q^{-B}|I_\rho|
\quad\text{and}\quad
\int_{I_\rho}|q''_{\can}|^2(\Gamma_G^2+\Gamma_{G,2}^2)\,dm
\ge Q^{-A}|I_\rho|,
\]
or else proves a source-robust row theorem or a genuinely global
arithmetic/projective/transport index theorem.
\end{remark}

\begin{remark}[Status ledger for candidate routes]
\label{rem:status-ledger-candidate-routes}
The source-collapse push leaves the following status for the main candidate
routes considered so far.

\begin{enumerate}
\item \emph{Order eleven and the finite anisotropy tower.}
A further local row computation may identify the next anisotropy factor
\(\Gamma_{G,3}\) and may be useful for documentation.  However, within the
established tower
\[
X_1YX_{2k}-X_{2k}YX_1,
\qquad
Y=yJ,\qquad
X_1=BS,\qquad B\propto q''_{\can},
\]
the fixed-sector row is always divisible by
\[
H=yq''_{\can}.
\]
Thus an order-eleven row of the same kind would at best have the schematic
form
\[
A_{11}^{\mathfrak f}\sim H\Gamma_{G,3}.
\]
It may help with further anisotropy walls, but it does not address the
dominant source-collapse regime \(H\approx0\).  Order eleven is therefore
diagnostic / documentary, not presently proof-closing.

\item \emph{Norm-based rank-two finite-core route.}
The known rows have the form
\[
A_7^{\mathfrak f}\sim H\Gamma_G,
\qquad
A_9^{\mathfrak f}\sim H\Gamma_{G,2}.
\]
Since
\[
H\sim-\frac{q_0(q''_{\can})^3}{288},
\]
an interval on which \(q''_{\can}\) is small collapses the norm vector even when
\(\Gamma_G\) remains bounded away from zero.  This route is blocked unless one
proves both source-curvature observability and source--anisotropy
reintegration, or finds a source-robust row.

\item \emph{Transported derivative tower.}
Differentiation gives
\[
\nabla A_j=(\nabla H)G_j+H\nabla G_j.
\]
Near \(q''=0\),
\[
\nabla H=O((q'')^2)
\]
under good-patch derivative control.  Thus first derivatives still collapse.
For every fixed finite derivative order, finite local models can plausibly
flatten \(H,\nabla H,\ldots,\nabla^R H\).  A derivative-tower rescue would
require a genuinely new growing-order source-jet theorem.

\item \emph{Projectivization / source-normalized ratios.}
The ratio
\[
[A_7:A_9]=[\Gamma_G:\Gamma_{G,2}]
\]
formally cancels \(H\) in the idealized factorized model.  However, this loses
magnitude, remains singular at
\[
\Gamma_G=\Gamma_{G,2}=0,
\]
and is not cleanly defined for the corrected row vector unless correction rows
vanish or are \(H\)-divisible.  No global projective index, endpoint, winding,
or holonomy theorem is presently available.  The projective branch is therefore
diagnostic only.

\item \emph{Correction rows as source-robust detectors.}
The current audits did not identify a correction row that is simultaneously
non-\(H\)-divisible, packet-visible, and absent in RH-compatible/on-line
models.  Line rows are gauge/reference-line terms, background/core rows are
baseline data, and normalization/subtraction rows are convention-dependent
unless separately packet-isolated.  Correction rows therefore do not currently
provide the missing source-robust detector.

\item \emph{Explicit formula / prime-side coercivity.}
The explicit formula represents the linear ingredients
\[
q_0,\qquad q'',\qquad q^{(4)},
\]
but the source-weighted residual
\[
(yq'')^2(\Gamma_G^2+\Gamma_{G,2}^2)
\]
is nonlinear.  Tensorizing the explicit formula produces oscillatory
multi-prime sums with no identified positivity or coercive lower bound.
Thus the explicit formula is representational here, not currently coercive.

\item \emph{Fixed-prime and AFE-style adversaries.}
Fixed-prime and AFE-compatible surrogate models support the interval-local
source-collapse mechanism.  In such models one attempts to force
\[
J_2^{(r)}(m_0)\approx0,\qquad 0\le r\le R,
\]
with \(J_2\sim q''\), while maintaining
\[
J_0\asymp1,\qquad g,h>0,\qquad \Gamma_G\asymp1.
\]
The AFE main sum, dual sum, smoothing, and gamma/background terms shift or
couple the finite target, but no algebraic blocker is currently known.

\item \emph{Same-scale packet geometry.}
The continuous same-scale source-flattening model and its finite quadrature
corollary show that positive same-distance packet geometry can make finite
source jets vanish while preserving \(q_0>0\), Gram positivity, and live
\(\Gamma_G\).  Therefore source-collapse cannot be excluded by local
same-scale packet geometry alone.

\item \emph{Interval \(L^2\) / source-curvature observability.}
The sharpest remaining positive target is not the full nonlinear residual but
the source-curvature estimate
\[
\int_{I_\rho}|q''_{\can}(m)|^2\,dm\ge Q^{-B}|I_\rho|,
\]
or at least
\[
\sup_{m\in I_\rho}|q''_{\can}(m)|\ge Q^{-B}.
\]
This is the least nonlinear form of the missing arithmetic input.

\item \emph{Source--anisotropy reintegration.}
Even if the interval source-curvature estimate
\[
\int_{I_\rho}|q''_{\can}(m)|^2\,dm\ge Q^{-B}|I_\rho|
\]
is proved, it does not by itself imply the source-weighted rank-two lower
bound.  One still needs to exclude concentration of \(q''_{\can}\)-mass near
\[
\Gamma_G=0,
\qquad
\Gamma_{G,2}=0.
\]
Equivalently, the architecture needs the vector estimate
\[
\int_{I_\rho}
|q''_{\can}|^2
(\Gamma_G^2+\Gamma_{G,2}^2)\,dm
\ge Q^{-A}|I_\rho|.
\]
The identity
\[
\Gamma_{G,2}
=
2q(q')^2-\frac16\Gamma_G''
\]
turns this into a local analytic/ODE rigidity problem coupled to the arithmetic
source-observability problem.  This bridge is currently a missing theorem, not
a consequence of source observability alone.

\item \emph{Global trace / projective index / Connes-style imbalance.}
A global trace, projective index, or spectral accounting theorem is not ruled
out.  However, it would be a new global proof architecture, not a continuation
of the present local rank-two norm route.
\end{enumerate}

Thus the current proof state is: the local rank-two machinery is a precise
obstruction and diagnostic framework, but it is non-closing unless one supplies
both actual-zeta source-curvature observability and source--anisotropy
reintegration, or else supplies a genuinely source-robust row or a new global
arithmetic/projective/trace theorem.
\end{remark}

\begin{remark}[How the downstream repairs modify the source-collapse ledger]
\label{rem:downstream-repairs-modify-source-collapse-ledger}
The source-collapse ledger remains valid as a no-go statement for purely local
rank-two norm arguments and for collapsed scalar \(I\)-coordinate row energy.
It should not be read as a no-go for the repaired downstream package.  The
repaired package first proves the source--anisotropy Gate~2 estimate from
Gate~1 by \(L^2\)-input rigidity, then uses wall escape and carrier-quotient
row energy to convert the resulting interval mass into row energy.  Thus the
source-collapse examples remain diagnostics and negative controls, but they no
longer constitute the active downstream obstruction once the hypotheses of
Theorems~\ref{thm:gate-two-from-gate-one-l2-rigidity},
\ref{thm:gauge-locked-wall-escape-row-energy}, and
\ref{thm:d4-carrier-quotient-through-order-nine} are assumed.
\end{remark}

\begin{remark}[Corrected projective transport obstruction]
\label{rem:corrected-projective-transport-obstruction}
One possible response to source divisibility is to pass to the projective ratio.
In the idealized factorized model
\[
A_7=HG_1,\qquad A_9=HG_2,\qquad H=yq''_{\can},
\]
one has
\[
[A_7:A_9]=[G_1:G_2].
\]
In the present setting \(G_1\) and \(G_2\) are represented by
\[
G_1\sim \Gamma_G,
\qquad
G_2\sim \Gamma_{G,2}.
\]
Thus projectivization formally cancels the common source factor \(H\).

However, the actual corrected rows may have the form
\[
A_7=HG_1+R_7,
\qquad
A_9=HG_2+R_9,
\]
where \(R_7,R_9\) denote line, background/core, normalization, or other
correction-row contributions not known to be \(H\)-divisible.  On the chart
\(A_7\ne0\), put
\[
Z=\frac{A_9}{A_7}.
\]
Then
\[
\nabla Z
=
\frac{A_7\nabla A_9-A_9\nabla A_7}{A_7^2},
\]
and substitution gives
\[
\begin{aligned}
A_7\nabla A_9-A_9\nabla A_7
={}&
H^2(G_1\nabla G_2-G_2\nabla G_1)\\
&+H(G_1\nabla R_9-G_2\nabla R_7)\\
&+H(R_7\nabla G_2-R_9\nabla G_1)\\
&+(\nabla H)(R_7G_2-R_9G_1)\\
&+(R_7\nabla R_9-R_9\nabla R_7).
\end{aligned}
\]
Therefore the clean source-normalized identity
\[
\nabla\left(\frac{A_9}{A_7}\right)
=
\nabla\left(\frac{G_2}{G_1}\right)
\]
is valid only if the correction rows vanish or are themselves
\(H\)-divisible.  If non-\(H\)-divisible corrections remain, then near
\(H=0\)
\[
[A_7:A_9]\to [R_7:R_9],
\]
not
\[
[\Gamma_G:\Gamma_{G,2}].
\]
Thus the projective branch is coherent only after a correction-divisibility or
correction-absorption theorem, or after enlarging the projective coordinate to
\[
[A_7:A_9:R_7:R_9:\cdots].
\]
No global projective endpoint, holonomy, or index theorem is presently known
for that enlarged object.
\end{remark}

\subsection{Finite-core sampled vector non-concentration}
\label{subsec:finite-core-sampled-vector-nonconcentration}

\begin{remark}[Updated role of finite-core non-concentration]
\label{rem:updated-role-finite-core-nonconcentration}
The finite-core non-concentration problem below is now secondary to the
source-collapse obstruction.  The known fixed-sector rows are source-divisible,
so a lower bound for a transported finite-core vector cannot be proved from
anisotropy non-concentration alone.  Any proof-relevant finite-core theorem
must either include the two-gate actual-zeta package
\[
\int_{I_\rho}|q''_{\can}(m)|^2\,dm\ge Q^{-B}|I_\rho|
\]
and
\[
\int_{I_\rho}
|q''_{\can}(m)|^2
\bigl(\Gamma_G(m)^2+\Gamma_{G,2}(m)^2\bigr)\,dm
\ge Q^{-A}|I_\rho|,
\]
or else use a genuinely source-robust row/invariant not divisible by
\[
H=yq''_{\can}.
\]
Without that additional source-level input or source-robust replacement, the
aligned/non-concentration language below remains a diagnostic framework rather
than a proof-closing theorem.
\end{remark}

The notation in this subsection originally arose in the scalar
\(\Phi_K\)-sampled formulation.  In the current rank-two architecture, the
corresponding fixed-sector vector is
\[
C_{\rm fd}^{\vec f}(\theta)
=
(A_7^{\mathfrak f},A_9^{\mathfrak f},\ldots,
A_{2R_{\Theta_k}+1}^{\mathfrak f})(\theta).
\]
The old scalar vector
\[
C_{\rm fd}(\theta)
=
(c_{\core,1},c_{\core,3},\ldots,c_{\core,2R_{\Theta_k}+1})(\theta)
\]
should be regarded as its \(\Phi_K\)-readout, not as the full fixed-sector
object.  The rank-two non-concentration theorem needed by the present
architecture is therefore a lower bound for
\[
P_{\rm fd}^{\vec f}(\theta)+J_{\rm fd}^{\vec f}(\theta),
\]
or for a transported enhancement of this vector, rather than merely for the
scalar \(\Phi_K\)-projection.

Let \(\Theta_k\) be a compact nonsingular actual-source \(k\)-core stratum.
Define the finite-determinacy vector
\[
C_{\rm fd}(\theta)
=
(c_{\core,1},c_{\core,3},\ldots,c_{\core,2R_{\Theta_k}+1})(\theta).
\]
Decompose it as
\[
C_{\rm fd}(\theta)=P_{\rm fd}(\theta)+J_{\rm fd}(\theta),
\]
where \(P_{\rm fd}\) is the pair-like or one-body finite vector and
\[
J_{\rm fd}
=
(J_1,J_3,\ldots,J_{2R_{\Theta_k}+1}),
\qquad
J_{2r+1}=[s^{2r+1}]\Phi_K(\mathcal I_{\ge2}),
\]
is the finite interaction vector.

For \(1\le N\le R_{\Theta_k}+1\), define
\[
F_N(\theta,Q)
=
c_{\core,2N-1}(\theta)
+
\sum_{j\ge1}
c_{\core,2N-1+2j}(\theta)
h_j(1^2,\ldots,N^2)Q^{-2j}.
\]
Set
\[
\Lambda(m,Q)
=
C_*
\max_{1\le N\le R_{\Theta_k}+1}
\left(
\frac{Q^{2N+2}}
{(2N-1)N!(N-1)!}
\frac{x}{B_\zeta(m)}S(m)
\right).
\]

\begin{hypothesis}[Bounded-depth projective detection]
\label{hyp:bounded-depth-projective-detection}
On the compact nonsingular actual-source stratum \(\Theta_k\), there exist
constants
\[
Q_0(\Theta_k),\qquad c_0>0
\]
such that for every \(Q\ge Q_0(\Theta_k)\) and every
\(\theta\notin\mathcal Z_k\), there is an index
\[
1\le N\le R_{\Theta_k}+1
\]
with
\[
|F_N(\theta,Q)|
\ge
c_0
\|C_{\rm fd}(\theta)\|_\infty.
\]
This is the bounded-depth projective detection and \(Q^{-2}\)-tail-control
input.  It is separate from sampled vector non-concentration: it detects the
finite-determinacy vector once that vector has calibrated size, but it does
not prove that calibrated size.
\end{hypothesis}

\begin{hypothesis}[Sampled finite-core vector non-concentration]
\label{hyp:sampled-finite-core-vector-nonconcentration}
For every actual zeta finite-core parameter
\[
\theta(m)\in\Theta_k,
\]
one needs the following alternative:
\[
\theta(m)\in\mathcal Z_k,
\]
or
\[
\theta(m)\in\mathcal E_{\rm geom},
\]
or
\[
\|P_{\rm fd}(\theta(m))+J_{\rm fd}(\theta(m))\|_\infty
\ge
\Lambda(m,Q).
\]
\end{hypothesis}

\begin{definition}[Interaction-aligned finite-core branch]
\label{def:interaction-aligned-finite-core-branch}
Define
\[
\mathcal N_k^{\rm align}(m,Q)
\]
to be the set of actual zeta finite-core parameters
\[
\theta(m)\in\Theta_k\setminus(\mathcal Z_k\cup\mathcal E_{\rm geom})
\]
such that
\[
0<
\|P_{\rm fd}(\theta(m))+J_{\rm fd}(\theta(m))\|_\infty
<
\Lambda(m,Q).
\]
Equivalently, Hypothesis~\ref{hyp:sampled-finite-core-vector-nonconcentration}
asserts
\[
\mathcal N_k^{\rm align}(m,Q)=\varnothing.
\]
\end{definition}

\begin{remark}[Relation among finite-core bad branches]
\label{rem:relation-among-finite-core-bad-branches}
The branch
\[
\mathcal N_k^{\rm align}(m,Q)
\]
is the current sharp form of the finite-core obstruction.  Earlier
formulations used either a distance-to-\(\mathcal Z_k\) tube
\[
\mathcal N_k(m,Q),
\]
or a coefficient-small branch
\[
\mathcal N_k^{\rm coef}(m,Q)
=
\left\{
\theta(m)\notin\mathcal Z_k\cup\mathcal E_{\rm geom}:
\max_N |F_N(\theta(m),Q)|<\Lambda(m,Q)
\right\}.
\]
For \(Q\ge Q_0(\Theta_k)\), bounded-depth projective detection and
\(Q^{-2}\)-tail control identify the coefficient-small branch with radial
smallness of the finite-determinacy vector, up to fixed compact-patch
constants:
\[
0<
\|C_{\rm fd}(\theta(m))\|_\infty
<
\Lambda(m,Q).
\]
Since
\[
C_{\rm fd}=P_{\rm fd}+J_{\rm fd},
\]
this is exactly the interaction-aligned branch
\[
\mathcal N_k^{\rm align}(m,Q).
\]
Accordingly, subsequent finite-core arguments should target
\(\mathcal N_k^{\rm align}\), not repeat finite determinacy or projective
detection.
\end{remark}

\begin{lemma}[Finite-core lower transfer from sampled vector non-concentration]
\label{lem:finite-core-lower-transfer-from-svnc-current}
Assume Hypothesis~\ref{hyp:sampled-finite-core-vector-nonconcentration} and
Hypothesis~\ref{hyp:bounded-depth-projective-detection}.
Then every actual zeta finite-core sample
\[
\theta(m)\notin\mathcal Z_k\cup\mathcal E_{\rm geom}
\]
has some \(1\le N\le R_{\Theta_k}+1\) satisfying
\[
|\Xi_{\core}^{(N)}(m)|
\ge
c
\frac{x}{B_\zeta(m)}S(m).
\]
\end{lemma}

\begin{proof}
Choose \(N\) by projective detection.  Since
\[
C_{\rm fd}=P_{\rm fd}+J_{\rm fd},
\]
Hypothesis~\ref{hyp:sampled-finite-core-vector-nonconcentration} gives
\[
|F_N(\theta(m),Q)|
\ge
c_0\Lambda(m,Q).
\]
By the definition of \(\Lambda(m,Q)\),
\[
\Lambda(m,Q)
\ge
C_*
\frac{Q^{2N+2}}
{(2N-1)N!(N-1)!}
\frac{x}{B_\zeta(m)}S(m).
\]
Substituting this into
\[
\Xi_{\core}^{(N)}(m)
=
(-1)^{N+1}
(2N-1)N!(N-1)!
Q^{-(2N+2)}
F_N(\theta(m),Q)
\]
gives the asserted lower bound.
\end{proof}

\begin{remark}[Why sampled vector non-concentration is the exact finite-core dependency]
\label{rem:why-svnc-is-exact-finite-core-dependency}
Lemma~\ref{lem:finite-core-lower-transfer-from-svnc-current} separates the
finite-core problem into a qualitative part and a sampled-size part.  The
qualitative part is finite odd-germ determinacy: on a compact nonsingular
actual-source stratum,
\[
C_{\rm fd}(\theta)=0
\quad\Longleftrightarrow\quad
H_{\core,\theta}\equiv0.
\]
Together with bounded-depth \(Q^{-2}\)-tail control, this gives projective
detection of every nonzero finite-core germ by some
\[
1\le N\le R_{\Theta_k}+1.
\]
What remains is the sampled finite-vector lower bound
\[
\|C_{\rm fd}(\theta(m))\|_\infty
=
\|P_{\rm fd}(\theta(m))+J_{\rm fd}(\theta(m))\|_\infty
\ge
\Lambda(m,Q)
\]
outside
\[
\mathcal Z_k\cup\mathcal E_{\rm geom}.
\]
Equivalently, the remaining finite-core branch is the interaction-aligned
small-vector branch
\[
\mathcal N_k^{\rm align}(m,Q).
\]
Thus any further finite-core proof must either prove
Hypothesis~\ref{hyp:sampled-finite-core-vector-nonconcentration}, or add an
independent theorem forcing \(P_{\rm fd}+J_{\rm fd}\)-smallness into a
controlled incidence stratum.
\end{remark}

\begin{remark}[Status of the finite-core branch]
\label{rem:status-finite-core-branch-current}
Finite odd-germ determinacy supplies qualitative detection on compact
nonsingular strata.  It does not supply the sampled vector lower bound in
Hypothesis~\ref{hyp:sampled-finite-core-vector-nonconcentration}.  The
live finite-core obstruction is calibrated interaction alignment:
\[
J_{\rm fd}(\theta(m))\approx -P_{\rm fd}(\theta(m)),
\]
making \(C_{\rm fd}=P_{\rm fd}+J_{\rm fd}\) nonzero but smaller than the
\(N\)-point scale \(\Lambda(m,Q)\), without entering any currently listed
controlled incidence stratum.
\end{remark}

\subsection{\(dd\)-interaction feature-balance obstruction}
\label{subsec:dd-interaction-feature-balance-obstruction}

In the covariance-weighted \(dd\)-interaction model, write
\[
K_{ij}^{dd}
=
(u_i-u_j)^2\langle v_i,v_j\rangle,
\qquad
v_i=v(u_i).
\]
For balanced amplitudes
\[
\sum_i a_i=0,
\]
the associated quadratic form is
\[
Q_{dd}(a)
=
2\left\langle
\sum_i a_i u_i^2v_i,\sum_j a_jv_j
\right\rangle
-
2\left|
\sum_i a_i u_iv_i
\right|^2.
\]
The second term is one-sided.  The first term has no fixed sign unless one has
feature balance:
\[
\sum_i a_iv_i=0.
\]
Ordinary amplitude balance does not imply feature balance.  Even rank-one
locking
\[
v_i=\lambda_i v_0
\]
only changes the missing condition to
\[
\sum_i a_i\lambda_i=0.
\]

\begin{lemma}[Feature balance is stronger than amplitude balance]
\label{lem:feature-balance-stronger-than-amplitude-balance}
Suppose that, for a fixed finite core,
\[
\sum_i a_i v_i=0
\]
for every balanced amplitude vector
\[
\sum_i a_i=0.
\]
Then
\[
v_1=v_2=\cdots=v_k.
\]
Thus ordinary amplitude balance implies feature balance only when the
transverse \(dd\)-features are constant on the core.
\end{lemma}

\begin{proof}
Take \(a=e_i-e_j\).  Then \(\sum_\ell a_\ell=0\), so feature balance gives
\(v_i-v_j=0\).  Since \(i,j\) are arbitrary, all \(v_i\) are equal.
\end{proof}

\begin{remark}[Rank-one locking does not close feature balance]
\label{rem:rank-one-locking-does-not-close-feature-balance}
If
\[
v_i=\lambda_i v_0,
\]
then feature balance becomes
\[
\sum_i a_i\lambda_i=0.
\]
This is not implied by
\[
\sum_i a_i=0
\]
unless the weights \(\lambda_i\) are constant on the core.  Thus rank-one
transverse locking merely transfers the missing condition from the vector
features \(v_i\) to the scalar weights \(\lambda_i\).
\end{remark}

\begin{definition}[\(dd\)-interaction branch]
\label{def:dd-interaction-branch-current}
The centered \(dd\)-interaction branch is
\[
\mathfrak Z_{\rm int}^{dd}
=
\left\{
\sum_\nu a_\nu K(d_\nu)
+
\sum_{i,j}a_i a_j
(u_i-u_j)^2\langle v_i,v_j\rangle
=0
\right\}
\setminus\mathcal E_{\rm geom}.
\]
\end{definition}

\begin{definition}[Transported \(dd\)-kernel Taylor coefficients]
\label{def:transported-dd-kernel-taylor}
Assume the favorable transported-feature reduction
\[
K_{dd}(t;m)
=
t^2
\left\langle
v\!\left(m+\frac t2\right),
v\!\left(m-\frac t2\right)
\right\rangle.
\]
Write
\[
v(m+h)
=
v_0+h v_1+\frac{h^2}{2}v_2+O(h^3).
\]
Then
\[
K_{dd}(t;m)
=
\kappa_2t^2+\kappa_4t^4+O(t^6),
\]
where
\[
\kappa_2=\langle v_0,v_0\rangle,
\]
and
\[
\kappa_4
=
\frac14
\left(
\langle v_0,v_2\rangle-\langle v_1,v_1\rangle
\right).
\]
\end{definition}

\begin{remark}[First Taylor obstruction to \(dd\)-kernel conditional definiteness]
\label{rem:first-taylor-obstruction-dd-conditional-definiteness}
A positive spectral representation of the form
\[
\operatorname{Re}(e^{-i\theta_0}K_{dd}(t;m))
=
\alpha(m)t^2+\int_{\mathbb R}(1-\cos(t\xi))\,d\mu_m(\xi),
\qquad
\mu_m\ge0,
\]
or its negative, imposes sign restrictions on the Taylor coefficients.  In
particular, if the phase is chosen so that
\[
\operatorname{Re}(e^{-i\theta_0}\kappa_2)>0,
\]
then the first nonquadratic coefficient must satisfy
\[
\operatorname{Re}(e^{-i\theta_0}\kappa_4)\le0.
\]
By Definition~\ref{def:transported-dd-kernel-taylor}, this requires
\[
\operatorname{Re}
\left(
e^{-i\theta_0}
\left(
\langle v_0,v_2\rangle-\langle v_1,v_1\rangle
\right)
\right)
\le0.
\]
The current one-pair incidence package does not prove this transverse
feature-curvature inequality, nor the higher moment/Hankel positivity
conditions needed for a full positive spectral representation.
\end{remark}

\begin{proposition}[Unit transported feature curves need not be spectrally positive]
\label{prop:unit-feature-curve-not-spectrally-positive}
Let
\[
v(u)=(\cos\phi(u),\sin\phi(u))
\]
be a unit transported transverse feature curve, and set
\[
K_{dd}(t)=
t^2
\left\langle
v\!\left(\frac t2\right),
v\!\left(-\frac t2\right)
\right\rangle.
\]
If
\[
\phi(u)=au+bu^3,
\]
then
\[
K_{dd}(t)
=
t^2
-\frac{a^2}{2}t^4
+
\left(
\frac{a^4}{24}-\frac{ab}{4}
\right)t^6
+
O(t^8).
\]
Consequently, even unit normalization of the transverse feature curve does
not imply the coefficient signs required by a conditionally one-sided spectral
representation.  In the negative-type convention, after
\(\kappa_2>0\) and \(\kappa_4\le0\), one still needs
\[
\kappa_6
=
\frac{a^4}{24}-\frac{ab}{4}
\ge0,
\]
which fails whenever
\[
b>\frac{a^3}{6}.
\]
\end{proposition}

\begin{proof}
Since
\[
\left\langle
v\!\left(\frac t2\right),
v\!\left(-\frac t2\right)
\right\rangle
=
\cos\left(
\phi\!\left(\frac t2\right)-\phi\!\left(-\frac t2\right)
\right),
\]
and
\[
\phi\!\left(\frac t2\right)-\phi\!\left(-\frac t2\right)
=
at+\frac b4t^3,
\]
we have
\[
K_{dd}(t)=t^2\cos\left(at+\frac b4t^3\right).
\]
Expanding the cosine through order \(t^6\) gives the stated formula.
\end{proof}

\begin{remark}[Needed input for the \(dd\)-interaction branch]
\label{rem:feature-balance-or-spectral-positivity-dd-needed}
On each compact nonsingular same-scale finite-core patch away from the
controlled degeneracy set, let \(v(u)\) be the transverse \(dd\)-source feature
after same-side whitening normalization.  To close the \(dd\)-interaction
branch, one needs at least one of the following:
\begin{enumerate}
\item \emph{Feature balance:}
\[
\sum_i a_i v(u_i)=0
\]
for every admissible balanced finite-core amplitude vector
\[
\sum_i a_i=0.
\]
By Lemma~\ref{lem:feature-balance-stronger-than-amplitude-balance}, this is
much stronger than ordinary amplitude balance.

\item \emph{Conditional spectral rigidity:} the reduced kernel
\[
K_{dd}(t;m)
=
t^2
\left\langle
v\!\left(m+\frac t2\right),
v\!\left(m-\frac t2\right)
\right\rangle
\]
admits, after phase rotation, a conditionally one-sided spectral
representation, or its negative:
\[
\operatorname{Re}(e^{-i\theta_0}K_{dd}(t;m))
=
\alpha(m)t^2+
\int_{\mathbb R}(1-\cos(t\xi))\,d\mu_m(\xi),
\qquad
\alpha(m)\ge0,\quad \mu_m\ge0,
\]
or the negative of such a kernel.  At minimum this route requires the Taylor
coefficient sign condition of
Remark~\ref{rem:first-taylor-obstruction-dd-conditional-definiteness}.
Proposition~\ref{prop:unit-feature-curve-not-spectrally-positive} shows that
unit normalization of \(v(u)\) alone is insufficient.

\item \emph{Incidence containment:} all feature-balance failures and centered
kernel signature-changing loci lie in the controlled incidence set
\(\mathcal E_{\rm geom}\).
\end{enumerate}
Absent one of these inputs, the branch
\[
\mathfrak Z_{\rm int}^{dd}
\]
remains independent.  In particular, pair-like positivity of the one-body
kernel is not by itself a centered-kernel spectral theorem for the
covariance-weighted \(dd\)-interaction.
\end{remark}

\subsection{The remaining global gate: quantitative transported non-concentration}
\label{subsec:current-global-proof-gate}

The reduced rank-two construction should be regarded as local infrastructure.
Under the reach, atlas, and row-table hypotheses, it gives a visibility-positive
fixed-sector bundle
\[
\mathcal F_{\rm red}^{+}\to \mathcal B_{\rm red}^{+}
\]
with connection
\[
\nabla_X^{\mathfrak f}A=\Pi_{\mathfrak f}(d_X^{\mathcal J}A),
\]
and a local affine forcing equation
\[
\nabla_\kappa^{\mathfrak f}A_7^{\mathfrak f}
=
E_{\eta\eta}+E_{\rm lower}.
\]
This does not by itself imply RH.

The remaining global problem is quantitative.  For a transported finite-core
vector
\[
C_{\rm fd}^{\rm tr}
=
(A_7^{\mathfrak f},
\nabla^{\mathfrak f}A_7^{\mathfrak f},
A_9^{\mathfrak f},
\nabla^{\mathfrak f}A_9^{\mathfrak f},\ldots),
\]
one must prove a lower bound of the form
\[
\|C_{\rm fd}^{\rm tr}(\theta(m))\|_\infty
\ge Q^{-B}
\]
outside controlled transported incidence.

A qualitative nonvanishing statement is insufficient.  Indeed, in the
pair/interaction decomposition
\[
C_{\rm fd}^{\rm tr}
=
P_{\rm fd}^{\rm tr}
+
J_{\rm fd}^{\rm tr},
\]
the aligned branch
\[
J_{\rm fd}^{\rm tr}\approx -P_{\rm fd}^{\rm tr}
\]
is not excluded by the local rank-two formalism.  Equivalently, one needs a
quantitative incidence-distance theorem
\[
\operatorname{dist}
\left(
\theta(m),
\mathcal Z_k^{\rm tr}\cup\mathcal E_{\rm geom}^{\rm tr}
\right)
\ge Q^{-A},
\]
A more concrete form of the same obstruction is the source-weighted layered
lower bound
\[
\boxed{
\sup_{m\in I_\rho}
|y(m)q''_{\can}(m)|
\max\{|\Gamma_G(m)|,|\Gamma_{G,2}(m)|\}
\ge Q^{-A}.
}
\]
The source-collapse countermodels above show that this estimate cannot be
deduced from local rank-two row algebra, finite Hadamard/quartet symmetry,
same-scale packet geometry, Gram positivity, or finite-prime/AFE surrogate
structure alone.  It now requires the two-gate actual-zeta package:
\[
\int_{I_\rho}|q''_{\can}(m)|^2\,dm\ge Q^{-B}|I_\rho|
\]
together with
\[
\int_{I_\rho}
|q''_{\can}(m)|^2
\bigl(\Gamma_G(m)^2+\Gamma_{G,2}(m)^2\bigr)\,dm
\ge Q^{-A}|I_\rho|.
\]
Alternatively, a zeta-specific rigidity, cone-separation, source-robust row,
or global arithmetic non-concentration theorem implying the displayed
source-weighted lower bound would suffice.

\begin{remark}[Numerical fixed-prime and AFE-style evidence]
\label{rem:numerical-fixed-prime-afe-source-collapse-evidence}
Fixed-prime and AFE-style surrogate experiments support the source-collapse
mechanism.  In these models one defines
\[
L_P(m)=\sum_{p\le P}(\log p)p^{-1/2}e^{-im\log p},
\]
\[
J_0=B_0+\Re L_P,\qquad
J_2=B_2+\Re L_P'',\qquad
J_4=B_4+\Re L_P^{(4)},
\]
and
\[
H_{\rm model}=-\frac{J_0J_2^3}{288}.
\]
The interval-local source-collapse search minimizes
\[
J_2^{(r)}(m_0)\approx0,\qquad 0\le r\le R,
\]
with side constraints
\[
J_0(m_0)\approx2,\qquad J_4(m_0)\approx0,
\]
and then verifies on
\[
I=[m_0-c/Q,m_0+c/Q]
\]
that
\[
\sup_I |J_2|
\]
is small while
\[
\max(|R_1|,|R_2|)
\]
remains bounded below.

Such experiments are not evidence about actual zeta by themselves, because
they use simplified backgrounds, finite prime cutoffs, and surrogate AFE
terms.  They do, however, show that the source-collapse mechanism is not a
mere symbolic artefact: it survives fixed-prime phases, prime powers,
smoothing, dual oscillatory blocks, and interval verification in the surrogate
models.  Therefore the positive route must identify a genuine actual-zeta
arithmetic obstruction, not another local rank-two algebraic identity.

For orientation, the latest normalized fixed-prime source-collapse run found an
interval with
\[
\sup_I |J_2|\approx 5.3\times 10^{-3},
\]
\[
\inf_I\max(|R_1|,|R_2|)\approx4.49,
\]
and
\[
\sup_I |H|\max(|R_1|,|R_2|)
\approx4.7\times10^{-9}.
\]
Adding prime powers did not remove the mechanism; a reported prime-power run
gave interval product about
\[
8.52\times10^{-8}.
\]
A smoothed main sum plus frozen-dual AFE surrogate gave a best interval product
near
\[
8.77\times10^{-9}.
\]
Including moving differentiated \(\chi(t)\), followed by exact
\(\chi(1/2+it)\)-derivatives, still left source-collapse examples; the
reported exact-\(\chi\) interval product was about
\[
1.93\times10^{-7}.
\]
These values are retained only as diagnostics.  They remove several easy
objections to the source-collapse mechanism, but they are not a theorem about
actual \(\zeta\) or its exact AFE normalization.

The numerical evidence is not used as proof.  Its role is to identify the
precise missing theorem: actual-zeta source-curvature observability on
\(I_\rho\), as formulated in
Subsection~\ref{subsec:arithmetic-source-curvature-observability-next-attack}.
\end{remark}

\begin{remark}[Exact \(\chi\)-derivatives in AFE-style source-collapse tests]
\label{rem:exact-chi-derivatives-used-in-afe-surrogates}
For the exact critical-line factor
\[
\chi(t):=\chi\!\left(\frac12+it\right)
=
\pi^{it}
\frac{\Gamma\left(\frac14-\frac{it}{2}\right)}
{\Gamma\left(\frac14+\frac{it}{2}\right)},
\]
write \(g(t)=\log\chi(t)\).  Then
\[
g(t)
=
it\log\pi+
\log\Gamma\left(\frac14-\frac{it}{2}\right)
-
\log\Gamma\left(\frac14+\frac{it}{2}\right),
\]
and, for \(n\ge1\),
\[
g^{(n)}(t)
=
\mathbf 1_{n=1}\,i\log\pi
+
\left(-\frac i2\right)^n
\psi^{(n-1)}\left(\frac14-\frac{it}{2}\right)
-
\left(\frac i2\right)^n
\psi^{(n-1)}\left(\frac14+\frac{it}{2}\right).
\]
If
\[
E_0=1,\qquad
E_n=\sum_{k=0}^{n-1}\binom{n-1}{k}E_k\,g^{(n-k)},
\]
then
\[
\chi^{(n)}=\chi E_n.
\]
Consequently,
\[
\frac{d^k}{dt^k}
\left[
\chi(t)\sum_n b_ne^{it\log n}
\right]
=
\sum_{a=0}^k
\binom{k}{a}
\chi^{(a)}(t)
\sum_n b_n(i\log n)^{k-a}e^{it\log n}.
\]
In the staged AFE-style tests, replacing the leading asymptotic phase by this
exact \(\chi\)-derivative calculation did not remove interval-local source
collapse.  It changed the finite derivative target, but did not create a Gate~1
source-curvature observability mechanism, let alone the full two-gate package.
\end{remark}

\begin{remark}[Gamma/background terms as shifted derivative targets]
\label{rem:gamma-background-shifted-derivative-targets}
Let
\[
J_2(t)=B_2(t)+P_2(t),
\]
where \(B_2\) is the gamma/Riemann--Siegel/far-tail background contribution and
\(P_2\) is the oscillatory prime/AFE contribution.  Interval source collapse
requires
\[
J_2^{(r)}(m_0)\approx0,
\qquad 0\le r\le R.
\]
Equivalently,
\[
P_2^{(r)}(m_0)\approx -B_2^{(r)}(m_0),
\qquad 0\le r\le R.
\]
Thus exact gamma/background terms shift the finite derivative-flat target.
They do not change the type of the problem unless one proves a range,
sign-cone, or arithmetic nonflatness theorem of the form
\[
\max_{0\le r\le R}
|B_2^{(r)}(m_0)+P_2^{(r)}(m_0)|
\ge Q^{-B}.
\]
No such theorem is currently supplied by the local rank-two package.  On the
microscopic interval \( |t-m_0|\le c/Q \), the background is smooth and is
captured by finite Taylor data, so it behaves as a shifted target for the
oscillatory AFE part.
\end{remark}

\begin{remark}[Next numerical test: actual-zeta source curvature]
\label{rem:next-numerical-test-actual-zeta-source-curvature}
The surrogate computations should now be followed by an actual-zeta numerical
scan of the canonical \(q''\) defined in
Definition~\ref{def:canonical-source-curvature-representation-problem}.  The
diagnostics should be interval-local:
\[
\sup_{m\in I}|q''_{\can}(m)|,
\qquad
\int_I |q''_{\can}(m)|^2\,dm,
\]
on intervals
\[
I=[m_0-c/Q,m_0+c/Q].
\]
The purpose is not to prove anything numerically, but to decide whether
source-flat intervals are common in actual zeta data or are an artefact of the
fixed-prime/AFE surrogates.  If source-flat intervals appear in actual data,
Hypothesis~\ref{hyp:offline-packet-source-curvature-observability} is likely
false in its present form.  If they do not appear, the first theorem to attack
is the interval \(L^2\) source-curvature observability estimate for the full
canonical object
\[
q''_{\can}=B_2+P_2+D_2+E_2.
\]
In the ordinary real-canonical branch, the source--anisotropy step is no
longer a separate next theorem; it is supplied by
Theorem~\ref{thm:gate-two-from-gate-one-l2-rigidity}.
\end{remark}

\begin{remark}[Current branch decision]
\label{rem:current-branch-decision-source-collapse}
This is a historical source-collapse fork retained for the interval/local norm
route.  After the \(L^2\)-Gate--1 repair, the active proof-closing upstream
branch is the aggregate observability theorem
\[
  \mathrm{Agg}L^2(q''_{\can}).
\]
The interval statements below should be read as localized proof tools or
diagnostics for that theorem.  In that older interval formulation, the sharp
fork was:

\begin{enumerate}
\item Prove actual-zeta source-curvature observability,
\[
\int_{I_\rho}|q''_{\can}(m)|^2\,dm\ge Q^{-B}|I_\rho|.
\]
In the ordinary real-canonical branch, source--anisotropy reintegration then
follows from Theorem~\ref{thm:gate-two-from-gate-one-l2-rigidity}.  The
remaining downstream burden is to convert the resulting Gate--2 mass into row
energy through D4 carrier-quotient row conversion and residual control;

\item or construct/verify an actual-zeta source-collapse model satisfying
\[
0<
\sup_{m\in I_\rho}
|yq''_{\can}|
\max\{|\Gamma_G|,|\Gamma_{G,2}|\}
<Q^{-A};
\]

\item or find a genuinely source-robust packet-visible row/invariant not
divisible by \(H=yq''_{\can}\);

\item or prove a global projective/holonomy/trace/index theorem that gives a
contradiction without a local norm lower bound.
\end{enumerate}

Absent one of these, the present rank-two norm/projective architecture remains
a strong local obstruction framework but not a proof of RH.
\end{remark}

\begin{remark}[Current proof frontier after the local package audits]
\label{rem:current-proof-frontier-after-local-package-audits}
This is the proof frontier after the local package audits and before the final
source-collapse reduction.  It is retained as a ledger of local dependencies;
the terminal status is given in
Remark~\ref{rem:terminal-status-current-architecture}.

\begin{enumerate}
\item The \(M_A,M_C\) connection branch is closed under the source-grade
reflection convention.

\item The promoted quadratic mixed-block Bottleneck~C coefficient
\[
E_{A,Y}^{\mathrm{BC}}
\]
is now blocked at the concrete odd-background top row
\[
\mathcal E^{2}_{\eta\eta;7,0,0}.
\]
The associated \(S\)-coefficient is
\[
f_{\eta\eta}^{\rm top}
=
-\frac{u}{24\pi q^3}\nu_{21,7}.
\]
Unless one proves
\[
\nu_{21,7}=0,
\]
finds a cancelling active top-row term, or adds a principled second-order
correction, this row is generically \(T\)-visible and can also be quotient
visible on a generic \(u_5\ne0\) patch.

\item The \(\Theta_\pm\)-normal form and septic normal-form counterterms do
not close the physical local obstruction.  They can move the obstruction into
a formal package, but the physical hidden coordinate carries it back through
\[
T^{\phys}
=
T^\star+\frac{\beta_{\loc}^{\Theta}}{c}.
\]
The known good-patch models have
\[
\partial_\kappa(\beta_{\loc}^{\Theta}/c)\ne0.
\]

\item The finite-core lower-transfer branch is now subordinate to the
source-collapse problem.  Earlier formulations reduced it to sampled vector
non-concentration, for example
\[
\mathcal N_k^{\rm align}(m,Q)=\varnothing.
\]
The later audits show that such a theorem must include source-level
nonflatness.  In particular, anisotropy or vector non-concentration alone
cannot control the known rows when
\[
H=yq''_{\can}\approx0.
\]
Thus the remaining source-side input is Gate~1 for the full canonical object
on \(I_\rho\), not merely finite determinacy, compactness, pairwise
zero-spacing, endpoint-gap control, or Remez--{\L}ojasiewicz volume estimates.
In the ordinary real-canonical branch, the former source--anisotropy step is
supplied downstream by Theorem~\ref{thm:gate-two-from-gate-one-l2-rigidity}.

\item The source-weighted layered route is blocked by source collapse.  The
known order-seven/order-nine rows have the form
\[
A_7^{\mathfrak f}\sim H\Gamma_G,
\qquad
A_9^{\mathfrak f}\sim H\Gamma_{G,2},
\qquad
H=yq''_{\can}.
\]
At leading order,
\[
H\sim-\frac{q_0(q''_{\can})^3}{288}.
\]
Thus an admissible interval on which
\[
q''_{\can}(m)\approx0
\]
while
\[
q_0(m)\asymp1,\qquad g(m),h(m)>0,\qquad
\max\{|\Gamma_G(m)|,|\Gamma_{G,2}(m)|\}\asymp1
\]
makes the norm vector
\[
(A_7^{\mathfrak f},A_9^{\mathfrak f})
\]
super-polynomially small.  Fixed-prime and AFE-style derivative-flat
surrogates, as well as same-scale packet moment models, do not presently
supply an obstruction.  The missing positive source-side input is Gate~1:
\[
\int_{I_\rho}|q''_{\can}|^2\,dm\ge Q^{-B}|I_\rho|.
\]
Once this is available in the ordinary real-canonical branch, the source-weighted
anisotropy estimate
\[
\int_{I_\rho}
|q''_{\can}|^2(\Gamma_G^2+\Gamma_{G,2}^2)\,dm
\ge Q^{-A}|I_\rho|
\]
is supplied by Theorem~\ref{thm:gate-two-from-gate-one-l2-rigidity}.

\item The \(dd\)-interaction branch remains independent unless feature
balance, conditional spectral positivity, or incidence containment is proved.
In the transported-kernel reduction, the first nonquadratic Taylor coefficient
\[
\kappa_4
=
\frac14
(\langle v_0,v_2\rangle-\langle v_1,v_1\rangle)
\]
already gives a necessary sign test for conditional definiteness.
\end{enumerate}
\end{remark}

\subsection{Arithmetic source-curvature observability: the next attack}
\label{subsec:arithmetic-source-curvature-observability-next-attack}

The preceding source-collapse countermodels show that the current rank-two
norm route cannot close from local row algebra, same-scale packet geometry, or
finite Hadamard/quartet symmetry alone.  The first missing input is an
arithmetic source-curvature theorem for actual zeta data.  In the notation
fixed below, the minimal useful pointwise form is
\[
\boxed{
\sup_{m\in I_\rho}|q''_{\can}(m)|\ge Q^{-B}.
}
\]
A stronger and more stable form is the interval energy estimate
\[
\boxed{
\int_{I_\rho}|q''_{\can}(m)|^2\,dm
\ge
Q^{-B}|I_\rho|.
}
\]
This estimate would block the pure source-collapse branch, but it is now only
the first gate.  It does not by itself recover the source-weighted rank-two
lower bound, because the \(q''_{\can}\)-mass could still concentrate near the
anisotropy base locus
\[
\Gamma_G\approx0,
\qquad
\Gamma_{G,2}\approx0.
\]
The second missing input is therefore a source--anisotropy reintegration
theorem, formulated below.

This subsection is the positive-theorem counterpart of
Remark~\ref{rem:status-ledger-candidate-routes}: all currently explored local,
projective, correction-row, higher-anisotropy, and finite-Hadamard routes have
either been classified as diagnostic or shown vulnerable to source collapse.

\begin{definition}[Canonical source-curvature representation problem]
\label{def:canonical-source-curvature-representation-problem}
In this subsection and in every source-observability statement below, the
canonical source-curvature object is
\[
\boxed{
q''_{\can}(m)=q_\zeta''(m)=\Phi_\zeta'''(m).
}
\]
This is the full actual-zeta phase-derivative curvature in the same Gram
normalization used by the local rank-two rows.  It is not the affine-subtracted
complement curvature.  The latter is denoted separately by
\[
q''_{\source,R}(m):=(G^{\new})''(m),
\qquad
G^{\new}=q_{\aux,R}+T_{\far,R}.
\]
The affine-subtracted complement curvature is a component in a chosen
source/core split; it is not the canonical arithmetic object whose nonflatness
is required.

The arithmetic representation problem is to write \(q''_{\can}\) on both sides
of the explicit formula.  On the zero side one needs a representation
\[
q''_{\can}(m)
=
B_2(m)+\sum_{\rho}^{\rm reg}K_2(m;\rho),
\]
with all background, tail, source/core, and regularization conventions fixed.
The background curvature is fixed canonically by
\[
\boxed{
B_2(t)=\theta'''(t)
=
-\frac18\Re\psi^{(2)}
\!\left(\frac14+\frac{it}{2}\right).
}
\]
Equivalently, for large \(t\),
\[
B_2(t)
=
-\frac1{2t^2}
-\frac1{8t^4}
-\frac7{64t^6}
+O(t^{-8}).
\]
The background term is included by default in zero-side proxy and canonical
scans.  A \(B_2\)-only mode is diagnostic only; it is never a Gate~1 candidate
because Gate~1 concerns the full object \(q''_{\can}\), not the
gamma/background term alone.
On the prime/AFE side one needs a representation of the form
\[
q''_{\can}(m)
=
B_2(m)
+
P_2(m)
+
D_2(m)
+
E_2(m),
\]
where
\[
P_2(m)
=
\Re\sum_{n\le X}
b_n(\log n)^2 n^{-1/2-im}W(n/T)
\]
has fixed actual-zeta coefficients, \(D_2\) denotes the dual or
functional-equation contribution, and \(E_2\) denotes a controlled
differentiated error term.  The coefficients \(b_n\), smoothing, dual term,
gamma/background contribution, and error scale must be specified in the same
normalization used by the rank-two local rows.  The point of this convention
is to forbid replacing \(q''_{\can}\) by an arbitrary Dirichlet polynomial or
by the affine-subtracted complement \(q''_{\source,R}\).
The displayed \((\log n)^2\)-weight corresponds to the literal
\(q''_{\can}\) curvature convention.  Some numerical determinant stress tests
also use a script-convention source row with \(k_q=3\), hence a
\(\Lambda(n)(\log n)^3n^{-1/2}\)-type coefficient.  That \(k_q=3\) convention
is a numerical stress-test convention recorded in
Remark~\ref{rem:log-derivative-afe-stress-test-convention}; it is not a
replacement for the literal \(q''_{\can}\) representation displayed here.
\end{definition}

\begin{remark}[Log-derivative AFE convention for numerical stress tests]
\label{rem:log-derivative-afe-stress-test-convention}
The packet-frequency determinant and local \(Z_\alpha\)-amplitude stress tests
use the logarithmic-derivative convention
\[
A(s)=-\frac{\zeta'}{\zeta}(s),
\qquad
C_k(s)=(-1)^k\partial_s^k A(s).
\]
The reflected functional-equation sign is
\[
\sigma_k=(-1)^{k+1}.
\]
For the current \(Z_\alpha\)-row convention one has \(k_Z=0\), so the plus and
minus branch coefficients are
\[
z_n^+(\alpha)
=
\Lambda(n)n^{-1/2-\alpha}V_Z^+(n/N_Z^+),
\]
\[
z_n^-(\alpha)
=
-\Lambda(n)n^{-1/2+\alpha}V_Z^-(n/N_Z^-).
\]
For the script-convention source row one uses \(k_q=3\):
\[
c_n^+
=
\Lambda(n)(\log n)^3n^{-1/2}V_q^+(n/N_q^+),
\]
\[
c_n^-
=
+\Lambda(n)(\log n)^3n^{-1/2}V_q^-(n/N_q^-).
\]
For a literal \(q''_{\can}\)-row test, replace \(k_q=3\) by \(k_q=2\), giving
\[
c_n^+
=
\Lambda(n)(\log n)^2n^{-1/2}V_q^+,
\qquad
c_n^-
=
-\Lambda(n)(\log n)^2n^{-1/2}V_q^-.
\]
This convention switch must not be made inside a single proof-level
Wronskian.  A given proof run must fix one convention for \(q''_{\can}\), and
then use the corresponding \(Z_\alpha\)-symbol ratios, dual branch, \(E_2\)
budget, and local row identities throughout.

The Wronskian operator node is the signed differential-frequency node, not the
folded conductor diagnostic node:
\[
\text{plus operator node }=+\log n,
\qquad
\text{minus operator node }=-\log n.
\]
Folded quantities such as
\[
\log\frac{m}{2\pi n}
\]
may be recorded as conductor/saddle diagnostics, but they are not substituted
for the eigenvalue of \(\lambda\partial_m\) unless a different folded
differential operator has been explicitly defined.

These conventions are numerical guardrails.  Proof-level canonical use still
requires the exact Mellin smoothing, the paired \(V^\pm\)-weights, derivatives
of \(m\)-dependent weights, the \(E_2\)-tail budget, and the selected real/phase
projection.
\end{remark}

\begin{hypothesis}[Gate 1: off-line-packet source-curvature observability]
\label{hyp:offline-packet-source-curvature-observability}
Let \(I_\rho\) be an admissible reduced good-patch interval attached to a
hypothetical off-line packet \(\rho\).  Actual zeta data satisfy
\[
\int_{I_\rho}|q''_{\can}(m)|^2\,dm
\ge
Q^{-B}|I_\rho|
\]
for some fixed \(B\), outside the already named reduced-base and correction
incidence loci.
Equivalently, in a weaker pointwise form,
\[
\sup_{m\in I_\rho}|q''_{\can}(m)|\ge Q^{-B}.
\]
This is a necessary source-level nonflatness input for the current rank-two
norm route.  It is not sufficient by itself to recover the desired lower bound
for
\[
|H|\max\{|\Gamma_G|,|\Gamma_{G,2}|\}.
\]
\end{hypothesis}

\begin{remark}[Gate 2 is no longer an independent ordinary-branch hypothesis]
\label{rem:gate-two-no-longer-independent-ordinary-branch}
In earlier versions, the source--anisotropy lower bound
\[
\int_{I_\rho}
|q''_{\can}(m)|^2
\left(
\Gamma_G(m)^2+\Gamma_{G,2}(m)^2
\right)\,dm
\ge
Q^{-A}|I_\rho|
\]
was stated as a separate Gate--2 hypothesis.  In the ordinary smooth
real-canonical branch, this is now supplied by
Theorem~\ref{thm:gate-two-from-gate-one-l2-rigidity}, assuming Gate~1,
source Lipschitz control for \(u=q''_{\can}\), and the \(L^2\)-input real
rigidity lemma.

Thus the displayed source--anisotropy estimate is retained as the Gate--2
conclusion, not as an independent ordinary-branch input.  Singular packets,
transition packets, or packets outside the real-canonical smooth convention
still require their own routing.
\end{remark}

\begin{remark}[ODE form of the reintegration problem]
\label{rem:ode-form-source-anisotropy-reintegration}
Let
\[
u=q''_{\can},
\qquad
\Gamma_G=u-6q+2q^3,
\qquad
\Gamma_{G,2}=(1-q^2)u-\frac16q^{(4)}.
\]
Differentiating \(\Gamma_G\) gives
\[
\Gamma_G''
=
u''-6u+12q(q')^2+6q^2u.
\]
Since \(u''=q^{(4)}\), one obtains the identity
\[
\boxed{
\Gamma_{G,2}
=
2q(q')^2-\frac16\Gamma_G''.
}
\]
Thus the historical reintegration problem has a local analytic/ODE
formulation: one must prevent source mass for \(u=q''_{\can}\) from
concentrating where both
\[
\Gamma_G\approx0
\qquad\text{and}\qquad
2q(q')^2-\frac16\Gamma_G''\approx0.
\]
In the present ordinary real-canonical branch this prevention is supplied by
the \(L^2\)-input rigidity mechanism of
Lemma~\ref{lem:l2-input-real-rigidity-safe} and
Theorem~\ref{thm:gate-two-from-gate-one-l2-rigidity}.  The ODE formulation is
therefore retained as explanatory structure and as a fallback diagnostic, not
as an independent active proof gate.
\end{remark}


\begin{remark}[Status of the Gate--2 analytic package after the \(L^2\)-rigidity repair]
\label{rem:status-gate-two-analytic-good-patch-package}
The previous conditional route through thick-set-to-interval propagation and
pointwise control of \(\Gamma_G''\) is no longer the active Gate--2 route.
In the ordinary real-canonical branch, Gate~2 is supplied by
Theorem~\ref{thm:gate-two-from-gate-one-l2-rigidity}.

The remaining local hypotheses are now sharply reduced:
\begin{enumerate}
\item the real-canonical definition of \(q_\zeta\), \(q''_{\can}\),
\(\Gamma_G\), and \(\Gamma_{G,2}\);

\item the source Lipschitz bound
\[
\|u'\|_{L^\infty(I)}\le Q^{C_u}|I|^{-1},
\qquad
u=q''_{\can};
\]

\item the safe primitive-bounded hypotheses in
Lemma~\ref{lem:l2-input-real-rigidity-safe}, or a later replacement proving
the primitive-free strengthening;

\item residual-stable identities if the displayed formulas for \(g\) and \(h\)
are used with controlled \(L^2\)-errors.
\end{enumerate}

Thus the old \(D1/D2\) status should be read as historical/fallback for this
branch, not as the current downstream proof gate.
\end{remark}

\begin{remark}[Off-line-packet conditional form]
\label{rem:offline-packet-conditional-source-observability}
A more targeted version of
Hypothesis~\ref{hyp:offline-packet-source-curvature-observability} uses the
hypothetical off-line packet itself.  Write
\[
q''_{\can}=q''_P+q''_{\mathrm{rest}},
\]
where \(q''_P\) is the contribution of the distinguished off-line packet and
\(q''_{\mathrm{rest}}\) is the contribution of all remaining actual-zeta data:
other core packets, auxiliary/tail terms, background, dual terms, and
controlled errors in the fixed convention.  This \(q''_{\mathrm{rest}}\)
should not be confused with the affine-subtracted complement curvature
\(q''_{\source,R}=(G^{\new})''\) unless the source/core split has explicitly
identified the two.

The interval observability claim may be sharpened to
\[
\int_{I_\rho}|q''_P(m)+q''_{\mathrm{rest}}(m)|^2\,dm
\ge
Q^{-B}|I_\rho|.
\]
Equivalently, the rest term cannot satisfy
\[
q''_{\mathrm{rest}}(m)=-q''_P(m)+O(Q^{-B})
\]
throughout the packet interval.  This formulation uses the off-line-packet
hypothesis directly and is stronger than a generic short-interval lower bound
for \(q''_{\can}\).  The same-scale source-flattening construction shows that
this cannot be proved from arbitrary positive same-scale packet geometry
alone; it must use actual-zeta arithmetic or a restriction on how actual zero
packets may be arranged around \(\rho\).
\end{remark}

\begin{remark}[Why this theorem is genuinely arithmetic]
\label{rem:source-curvature-observability-is-arithmetic}
Hypothesis~\ref{hyp:offline-packet-source-curvature-observability} cannot
follow from the local packet geometry developed above.  The same-scale
source-flattening theorem and finite quadrature corollary show that positive
same-distance packet configurations can make
\[
q''_{\can}(m_0),\ (q_\zeta)^{(4)}(m_0),\ldots,(q_\zeta)^{(M)}(m_0)
\]
vanish to arbitrary finite order while keeping
\[
q_0>0,\qquad g>0,\qquad h>0,
\]
and while keeping
\[
\Gamma_G=q''_{\can}-6q_0+2q_0^3
\]
bounded away from zero after normalization, for example near \(q_0=2\).
Thus any proof of
Hypothesis~\ref{hyp:offline-packet-source-curvature-observability} must use
actual zeta arithmetic beyond finite same-scale positive packet geometry.
\end{remark}

\begin{remark}[Zero-side quadrature obstruction to local proofs]
\label{rem:zero-side-quadrature-obstruction-to-local-proofs}
The same-scale source-flattening model gives a concrete obstruction to any
purely local zero-side proof of
Hypothesis~\ref{hyp:offline-packet-source-curvature-observability}.  The
continuous measure in
Theorem~\ref{thm:continuous-same-scale-source-flattening-model} lives on the
same-distance curve
\[
x^2+c^2=d^2
\]
and has positive density
\[
d\mu(\theta)=\frac{2}{\pi}\sin\theta\,d\theta.
\]
It makes the averaged packet profile locally constant, and finite quadrature
then approximates arbitrary finite source jets while keeping positive weights
and same-distance packet scale.  Therefore actual zeta must exclude, or fail
to approximate, this quadrature pattern if the rank-two norm route is to close.

A plausible arithmetic formulation is: actual zeta zeros near a hypothetical
off-line packet cannot realize the finite same-scale quadrature needed to make
\[
q''(m_0),q^{(4)}(m_0),\ldots,q^{(2M)}(m_0)
\]
all small while \(q_0>0\) and the Gram block remains positive.  This is the
zero-side form of the source-curvature observability problem.
\end{remark}

\begin{remark}[Prime-side observability formulation]
\label{rem:prime-side-source-curvature-observability-formulation}
After fixing the representation in
Definition~\ref{def:canonical-source-curvature-representation-problem}, the
source-curvature problem becomes a short-interval fixed-coefficient
non-approximation problem.  Writing
\[
q''_{\can}(m)=B_2(m)+P_2(m)+D_2(m)+E_2(m),
\]
source collapse means
\[
P_2(m)\approx -B_2(m)-D_2(m)
\quad\text{throughout }I_\rho,
\]
up to the controlled error \(E_2\).  The desired arithmetic theorem is
therefore an observability estimate of the form
\[
\int_{I_\rho}|B_2(m)+P_2(m)+D_2(m)|^2\,dm
\ge Q^{-B}|I_\rho|,
\]
with \(E_2\) absorbed into the final error budget.

This is not a generic Dirichlet-polynomial problem.  The coefficients in
\(P_2\) are fixed by the actual zeta/AFE normalization.  Standard localized
large-sieve, Montgomery--Vaughan, explicit-formula energy, or uncertainty
principles are useful diagnostics, but no theorem currently proves this
fixed-coefficient lower bound.  If interval estimates remain out of reach, the
equivalent finite-jet formulation below is the sharper arithmetic target.
\end{remark}

\begin{remark}[Finite-jet version of the arithmetic problem]
\label{rem:finite-jet-source-nonflatness-problem}
The interval energy estimate may also be attacked through finite jets.  On an
interval
\[
I_\rho=[m_0-c/Q,m_0+c/Q],
\]
source collapse requires simultaneous smallness of
\[
\partial_m^r q''_{\can}(m_0),
\qquad
0\le r\le R,
\]
where the scale audit gives
\[
R\asymp \frac{\log Q}{\log\log Q}
\]
or more conservatively \(R=O(\log Q)\).  In an AFE model this becomes a finite
system
\[
\partial_m^r\bigl(B_2+D_2+P_2\bigr)(m_0)
\approx0,
\qquad
0\le r\le R,
\]
or equivalently
\[
P_2^{(r)}(m_0)
\approx
-\left(B_2+D_2\right)^{(r)}(m_0).
\]
The positive theorem would be a fixed-coefficient prime-log derivative-jet
nonrecurrence statement:
\[
\max_{0\le r\le R}
\left|
\partial_m^r q''_{\can}(m_0)
\right|
\ge Q^{-B}.
\]
The recurrence heuristics developed above suggest that this is a difficult
arithmetic theorem and may be false in fixed-prime or AFE surrogates.  It is
therefore not currently proved by large-sieve or generic
Dirichlet-polynomial approximation estimates.
\end{remark}

\begin{remark}[Terminal status of the current architecture]
\label{rem:terminal-status-current-architecture}
The present proof architecture is conditional and non-closing as a proof of
\(\RH\).  The active stack is now
\[
  \boxed{
  \mathrm{Agg}L^2(q''_{\can})
  +
  \mathrm{Graph/O2O3}_{23}
  +
  \mathrm{NormalD4}_{K9,Y12}.
  }
\]

The active Gate--1 object is not a pointwise ordinary Wronskian, transition
jet, singular sine moment, or off-line interval diagnostic.  It is scalar
local \(L^2\)-energy of the full canonical source
\[
  q''_{\can}=B_2+P_2+D_2+E_2.
\]
Equivalently, the primary upstream target is
\[
  \boxed{
  \|q''_{\can}\|_{L^2(W_{\mathrm{branch}})}^2
  \ge
  Q^{-A}
  \|q''_{\can}\|_{L^2(W_{\mathrm{ref}})}^2
  }
\]
plus
\[
  \boxed{
  \|q''_{\can}\|_{L^2(W_{\mathrm{ref}})}^2
  \ge
  Q^{-A}\sum_{b,m}\rho_b(m).
  }
\]
The interval source-curvature estimates in
Subsection~\ref{subsec:arithmetic-source-curvature-observability-next-attack}
are retained as localized proof routes and source-collapse diagnostics for this
\(L^2\)-Gate--1 theorem.

In the ordinary real-canonical branch, Gate--2 is no longer an independent
downstream missing input.  It is supplied by
Theorem~\ref{thm:gate-two-from-gate-one-l2-rigidity}, assuming the canonical
definition package, source local complexity, and the safe \(L^2\)-input
rigidity lemma.

The finite/O2/O3 module is separate from Gate--2.  Its current analytic form is
the \(R^C\)-normalized graph problem
\[
  \mathbf C\cdot1+m_{\can}\mathbf C_R[x]
  =
  m_{\can}\left(\mathbf C\cdot(1/m_{\can})+\mathbf C_R[x]\right),
\]
together with full determinant survival inside the scalar canonical matrix.
The pure graph estimate for \(1/m_{\can}\) is only a branch witness; it must
survive in the full determinant under the Stack weight.

The current downstream contradiction candidate is the \(Y12\)-signature
specialization of the singular-discard channel:
\[
  [\mathsf C_{Y12}]\ne0
  \quad\text{in}\quad
  \mathcal Q_{9,Y12}^{\mathrm{forb},\alpha}.
\]
The most economical downstream stack is therefore
\[
  \boxed{
  \mathrm{Agg}L^2(q''_{\can})
  +
  \mathrm{FullDetSurvival/RieszNorm}_{23}
  +
  \mathrm{SecondOrderY12SignatureSeparation}
  +
  \mathrm{EFCompletion}_{Y12}
  +
  \mathrm{MovingWallAntiConcentration}_{Y12}
  +
  \mathrm{TubularNormality}_{\mathrm{sing}}.
  }
\]
If the \(Y12\)-signature vanishes after adding actual parameter tangents,
lower connection effects, and actual second curvature, the present
\(K_9/Y12\) finite-D4 route fails and one must move to a different residual,
higher order, or fixed-\(\alpha\) O2/O3 route.
\end{remark}

\section{How the older framing maps to the new framing}
\label{sec:old-to-new-framing}

\begin{center}
\renewcommand{\arraystretch}{1.15}
\begin{tabular}{p{0.30\linewidth} p{0.30\linewidth} p{0.30\linewidth}}
\hline
Older framing & New framing & Status \\
\hline
Finite-row O1/O2/O3 Gate--1 & \(\mathrm{Agg}L^2(q''_{\can})\) source survival & historical tool only \\
S1/T1 transition / singular diagnostics & proof tools for scalar source observability & not the live Gate--1 definition \\
Scalar \(I\)-coordinate D4 & carrier quotient / row-normal plain wall & dead fallback \\
Derivative-wall D4 & separate no-go / future derivative-channel theorem & not active \\
Optimized \(\Xi_N\) route & rank-zero on the current finite D4 live plane & no-go \\
One-parameter row mobility & not a forbidden identity & no-go \\
Full \(K_9\)-or-bust route & stale & replaced by \(Y12\)-signature quotient \\
\hline
\end{tabular}
\end{center}

The current live chain is
\[
\mathrm{Agg}L^2(q''_{\can})
\;\Rightarrow\;
\mathrm{Graph/O2O3}_{23}
\;\Rightarrow\;
Y12\text{-quotient survival}
\;\Rightarrow\;
\mathrm{AtomicY12Contradiction}.
\]

\section{No-go and audit ledger}
\label{sec:no-go-and-audit-ledger}

All later \(\mathrm{wip}\)-remarks in this section use the shared finite-survival
notation of
Subsection~\ref{subsec:shared-finite-survival-language}
(Definition~\ref{def:finite-survival-gaps}).

\begin{remark}[Unified bad-zero/nonrecurrence obstruction]
\label{rem:unified-bad-zero-nonrecurrence-obstruction}
After the ordinary-route audit, the remaining upstream burden is not local
row algebra.  O1a, O1b, O2--BZ, and O3 are all pointwise actual-orbit
bad-zero or incidence exclusions at zeros of \(Z_\alpha\).  Continuous
mean-square estimates, phase-relaxed torus models, generic determinant
nonvanishing, exterior coefficient mass, and averaged large-sieve estimates
do not imply the needed pointwise lower bounds.  The missing input is a
prime-log nonrecurrence theorem at specified \(Z_\alpha\)-zeros, formulated
for the full canonical source
\[
q''_{\can}=B_2+P_2+D_2+E_2.
\]
\end{remark}

\begin{remark}[Historical route register]
\label{rem:historical-route-register}
The following routes are historical or no-go unless explicitly reactivated:
\begin{itemize}
\item O1/O2/O3, S1/T1, and ordinary finite-row Gate--1 as active definitions;
\item scalar \(I\)-coordinate D4 and derivative-wall D4;
\item optimized \(\Xi_N\) as a finite D4 contradiction channel;
\item one-parameter row mobility as a forbidden upper;
\item functorial carrier lifts when \(R_{\mathrm{sing}}\) is Taylor-complete or carrier-complete;
\item pure-branch determinant-only proofs for full O2/O3.
\end{itemize}
All live statements should be phrased in the current architecture:
\[
\mathrm{Agg}L^2(q''_{\can})+\mathrm{Graph/O2O3}_{23}+\mathrm{NormalD4}_{K9,Y12}.
\]
\end{remark}

\begin{remark}[Unified live/dead dictionary]
\label{rem:unified-live-dead-dictionary}
Whenever the text says \(\mathrm{Gate\text{--}1}\), it means \(\mathrm{Agg}L^2(q''_{\can})\).  Whenever it says \(\Lambda\), it means the transition--singular discarded operator \((I-P_{\mathrm{sing}})T_{\mathrm{tr}\to\mathrm{sing}}P_{\mathrm{tr}}\) unless a local section explicitly defines a different quotient.  D4 is only local visibility unless the forbidden channel has been established in the same quotient norm as the upper bound.
\end{remark}

\begin{lemma}[\(Y12\) quotient descent and primitive survival]
\label{lem:Y12-quotient-descent}
For each packet \(p\), define
\[
\mathcal K_p^Y
:=
\operatorname{im}
\begin{bmatrix}
R_{B_2,p}&R_{P,p}&R_{D,p}&A_{\alpha,p}&A_{7,p}&A_{<7,p}&A_{II,p}&A_{E_2,p}
\end{bmatrix},
\qquad
q_p^Y:=I-P_{\mathcal K_p^Y}^{\mathrm{bd}},
\]
and let
\[
\Lambda_{Y,p}
:=
q_p^Y\Pi_{\Sigma(Y12),p}T_{\mathrm{tr}\to\mathrm{sing},p}P_{\mathrm{tr},p}.
\]
Assume the \(D4\) nuisance image is killed by the same quotient:
\[
\Pi_{\Sigma(Y12),p}T_{\mathrm{tr}\to\mathrm{sing},p}\mathcal N_p^{\mathrm{D4}}
\subseteq
\mathcal K_p^Y.
\]
Then \(\Lambda_{Y,p}\) factors through the \(D4\) quotient \(V_p/N_p^{\mathrm{D4}}\).
Moreover, if
\[
Y_{Y,p}^{\mathrm{tar}}
:=
\Pi_{\Sigma(Y12),p}T_{\mathrm{tr}\to\mathrm{sing},p}[\mathsf C_{Y12}]^{\mathrm{tr}},
\]
then the Schur residual
\[
r_{Y,p}
:=
(I-P_{\mathcal K_p^Y}^{\mathrm{bd}})Y_{Y,p}^{\mathrm{tar}}
\]
vanishes if and only if \(Y_{Y,p}^{\mathrm{tar}}\in \mathcal K_p^Y\).
\end{lemma}

\begin{proof}
If \(v-v'\in N_p^{\mathrm{D4}}\), then by hypothesis
\[
\Pi_{\Sigma(Y12),p}T_{\mathrm{tr}\to\mathrm{sing},p}(v-v')\in \mathcal K_p^Y,
\]
hence
\[
q_p^Y\Pi_{\Sigma(Y12),p}T_{\mathrm{tr}\to\mathrm{sing},p}(v-v')=0.
\]
Therefore \(\Lambda_{Y,p}(v)=\Lambda_{Y,p}(v')\), so \(\Lambda_{Y,p}\) is constant
on cosets of \(N_p^{\mathrm{D4}}\) and descends to a unique induced map on
\(V_p/N_p^{\mathrm{D4}}\).

For the residual statement, let \(A_p\) be any matrix whose columns span
\(\mathcal K_p^Y\).  Then
\[
r_{Y,p}=(I-P_{\mathcal K_p^Y}^{\mathrm{bd}})Y_{Y,p}^{\mathrm{tar}}
\]
and
\[
\|r_{Y,p}\|^2
=
\|Y_{Y,p}^{\mathrm{tar}}\|^2
-
(Y_{Y,p}^{\mathrm{tar}})^*A_p(A_p^*A_p)^\dagger A_p^*Y_{Y,p}^{\mathrm{tar}}.
\]
Hence \(r_{Y,p}=0\) precisely when \(Y_{Y,p}^{\mathrm{tar}}\in \mathcal K_p^Y\).
\end{proof}

\begin{corollary}[Rank form]
\label{cor:Y12-quotient-rank}
With the notation of Lemma~\ref{lem:Y12-quotient-descent},
\[
\operatorname{rank}
\begin{bmatrix}
\mathcal K_p^Y & \Pi_{\Sigma(Y12),p}T_{\mathrm{tr}\to\mathrm{sing},p}\mathcal N_p^{\mathrm{D4}}
\end{bmatrix}
=
\operatorname{rank}(\mathcal K_p^Y)
\]
if and only if the \(Y12\) observable descends through the \(D4\) quotient, and
\[
\operatorname{rank}
\begin{bmatrix}
\mathcal K_p^Y & Y_{Y,p}^{\mathrm{tar}}
\end{bmatrix}
=
\operatorname{rank}(\mathcal K_p^Y)+1
\]
if and only if the primitive \(Y12\) carrier survives the quotient.
\end{corollary}

\noindent
In particular, the current \(D4\to Y12\) route reduces to the two finite checks
\[
\Pi_{\Sigma(Y12)}T_{\mathrm{tr}\to\mathrm{sing}}\mathcal N_p^{\mathrm{D4}}\subseteq \mathcal K_p^Y
\qquad\text{and}\qquad
Y_{Y,p}^{\mathrm{tar}}\notin \mathcal K_p^Y.
\]

\begin{remark}[wip-zeta-bridge-proof-y12-descent]
\label{rem:wip-zeta-bridge-proof-y12-descent}
What remains is the quotient compatibility test:
\[
\Pi_{\Sigma(Y12)}T_{\mathrm{tr}\to\mathrm{sing}}\mathcal N_p^{\mathrm{D4}}\subseteq \mathcal K_p^Y,
\]
together with primitive carrier survival:
\[
Y_{Y,p}^{\mathrm{tar}}\notin \mathcal K_p^Y.
\]
\end{remark}

\begin{remark}[No-go: Taylor-complete singular retained space]
\label{rem:no-go-taylor-complete-singular-retained-space}
If the retained singular space \(R_{\mathrm{sing}}\) contains all Taylor/jet
data through the D4 carrier orders, then
\[
  T_{\mathrm{tr}\to\mathrm{sing}}
  \operatorname{span}\{[\mathsf C_7]^{\mathrm{tr}},[\mathsf C_{Y12}]^{\mathrm{tr}}\}
  \subset R_{\mathrm{sing}}.
\]
Consequently
\[
  (I-P_{\mathrm{sing}})T_{\mathrm{tr}\to\mathrm{sing}}[\mathsf C_7]^{\mathrm{tr}}=0,
  \qquad
  (I-P_{\mathrm{sing}})T_{\mathrm{tr}\to\mathrm{sing}}[\mathsf C_{Y12}]^{\mathrm{tr}}=0,
\]
and the singular-discard forbidden channel has rank zero.  Thus
\(R_{\mathrm{sing}}\) must be an actual EF-coherent retained space, not a
Taylor-complete carrier space, if this \(\Lambda\)-route is to remain viable.
\end{remark}

\begin{remark}[Active no-go boundary]
\label{rem:active-no-go-boundary}
The current \(D4\to Y12\) bridge is only valid if the forbidden observable
descends through the same quotient that kills the \(D4\) nuisance block and
the primitive \(Y12\) carrier survives that quotient.  If either condition
fails, the bridge is structurally incompatible, not merely unfinished.
\end{remark}

\begin{remark}[Active rank-two closure register]
\label{rem:active-rank-two-closure-register}
The active rank-two branch is the following closure pipeline:
\[
  \mathrm{RowNormalForm}_{2Y}
  +
  \mathrm{RowExactness}_{2Y}
  \Longrightarrow
  B_Y=0,
\]
\[
  \mathrm{D4CoeffCoreRouting}_{2Y}
  \Longrightarrow
  F_{2Y}\asymp H(\Gamma_G,\Gamma_{G,2})
  \quad\text{on predetermined cores},
\]
\[
  \mathrm{WallVariableCompletion}_{2Y}
  +
  \mathrm{FiniteAtlasRiesz}_{2Y}^{p}
  +
  \mathrm{RankTwoCoercivity}
  +
  \mathrm{Agg}L^2(q''_{\can})
  \Longrightarrow
  \int W_2\ge Q^{-A},
\]
Here \(\mathrm{Agg}L^2(q''_{\can})\) is the branch-complete aggregate
\(L^2\)-source input for the full canonical source
\[
  q''_{\can}=B_2+P_2+D_2+E_2.
\]
It is not the archived pointwise ordinary \(\mathrm{O1/O2/O3}\) finite-row
route.
\[
  W_2=|\Gamma_G|^2+|\Gamma_{G,2}|^2.
\]
Then
\[
  \mathrm{CoreWeightedPlainWallTransfer}_{2Y}
  \Longrightarrow
  \int N_{\sm}|H|^2W_2\ge Q^{-A'},
\]
and a fixed scalar component
\[
  G_\star\in\{\Gamma_G,\Gamma_{G,2}\}
\]
is selected at the smooth aggregate stage for scalar O2/O3.

The unresolved local blockers are now:
\[
  T_{\mathrm{lower}}^{\mathrm{arith}}\text{ actual corrected-package coordinate exclusion},
\]
and
\[
  \mathcal N_{\mathrm{sing}}\text{ row-dead/live-pair preserving exact nuisance}.
\]
The unresolved analytic blockers are:
\[
  p=q''\text{ source completion in }H^2_{\mathrm{sc}},
\]
\[
  \text{ProductRiesz/Remez for }\mathcal W_G,\mathcal W_{G2},
\]
and
\[
  \text{predetermined }R^{-C}\text{-large coefficient/smooth-good cores}.
\]
\end{remark}

\begin{remark}[Next finite audits after the \(L^2\)-Gate--1 repair]
\label{rem:next-finite-audits-after-l2-gate-one}
After the \(L^2\)-Gate--1 repair, the next finite audits are:
\begin{enumerate}
\item \textbf{Finite AFE anti-lock.}
Prove
\[
  \|B_2+P_{2,\le X}+D_{2,\le X}\|_{L^2(W_{\mathrm{branch}})}^2
  \ge
  Q^{-A}
  \|B_2+P_{2,\le X}+D_{2,\le X}\|_{L^2(W_{\mathrm{ref}})}^2.
\]

\item \textbf{Central tube routing.}
For \(\delta=R^{-C}\), route
\[
  \Omega_{\mathrm{cen}}=\{|x-\tfrac12|<\delta\}
\]
as a rescaled branch and prove both the \(L^2\)-anti-lock and graph/O2/O3
lower bounds there.

\item \textbf{Canonical quadratic multiplier.}
Prove
\[
  m_{\can,+}(x)=\kappa(x-z_0)^2u(x)
\]
with \(R^C\)-controlled divided differences for \(u\), giving
\[
  \operatorname{dist}(1/m_{\can,+},\mathbf C_R[x])\ge Q^{-A}.
\]

\item \textbf{Near-pure Schur block.}
Prove
\[
  D_{-n}
  \notin_{Q^{-A}}
  \operatorname{span}\{D_{-n+1},\ldots,D_{-n+d_0}\}
\]
under the active Stack weight.

\item \textbf{EF residual split.}
Prove both
\[
  \mathrm{EFCompletion}_{L^2}
\]
and, separately,
\[
  \mathrm{EFCompletion}_{K9}.
\]

\item \textbf{Y12 honesty.}
Classify
\[
  \Pi_{Y12}E_2
\]
as absent/small, fixed/coherent, or freely tangent.

\item \textbf{Y12 normal quotient.}
Compute whether
\[
  [\mathsf C_{Y12}]\ne0
  \quad\text{in}\quad
  \mathcal Q_9^{\mathrm{forb}}.
\]

\item \textbf{Transition--singular leakage.}
Compute
\[
  (I-P_{\mathrm{sing}})T_{\mathrm{tr}\to\mathrm{sing}}[\mathsf C_7]^{\mathrm{tr}},
  \qquad
  (I-P_{\mathrm{sing}})T_{\mathrm{tr}\to\mathrm{sing}}[\mathsf C_{Y12}]^{\mathrm{tr}}.
\]

\item \textbf{\(K_9\)-conditioning.}
Prove lower-order tangent observability:
\[
  |T_La|\ge Q^{-A}|T_9a|
  \qquad
  \text{for every actual-source tangent }a.
\]

\item \textbf{Tubular singular coherence.}
Prove a \(Q^C\)-bounded second fundamental form / tubular radius estimate for
\[
  \mathcal M_{\mathrm{sing}}^{\mathrm{arith}}.
\]
\end{enumerate}
\end{remark}

\begin{remark}[Final unresolved finite decisions]
\label{rem:final-unresolved-finite-decisions}
After the current cleanup, the remaining decisive finite decisions are:
\begin{enumerate}
\item Define \(R_{\mathrm{sing}}\) as an actual-source coherence space and test
\[
  (I-P_{\mathrm{sing}})T_{\mathrm{tr}\to\mathrm{sing}}
  \operatorname{span}\{[\mathsf C_7]^{\mathrm{tr}},[\mathsf C_{Y12}]^{\mathrm{tr}}\}.
\]
Rank \(0\) kills the singular-discard \(\Lambda\)-route; rank \(1\) gives a
rank-one forbidden wall; rank \(2\) gives the full rank-two forbidden wall.

\item Decide whether \(T_{\mathrm{tr}\to\mathrm{sing}}\) transports formed labelled
carriers or recomputes \(N_0/\mathcal W\).  Transport preserves the
\(Y12\)-origin filter; recomputation forces a direct Schur separation theorem.

\item Decide whether O2/O3 will be proved by full graph-block Riesz survival,
by determinant branch survival, or by an \(\alpha\)-admissible determinant
average.  Pure-branch graph non-polynomiality alone does not certify the full
canonical augmented matrix.

\item Decide whether \(Y12\)-normality is proved by the full \(K_9\)-residual
or by the targeted \(Y12\)-signature quotient
\[
  \mathcal Q_{9,Y12}^{\mathrm{forb},\alpha}.
\]
The targeted route is weaker and should be attempted first.
\end{enumerate}
\end{remark}

% === audit-no-Lambda-no-D4-contradiction ===
% Purpose:
%   Record the current downstream no-go boundary.
% Status:
%   Active audit ledger.

\begin{remark}[No forbidden channel, no D4 contradiction]
\label{rem:no-Lambda-no-D4-contradiction}
Carrier-quotient D4 is a local visibility statement.  It becomes a contradiction
only after a globally forbidden channel \(\Lambda\) is identified.  A total
Bessel or row-energy upper of the form
\[
  \|V_{\mathrm{glob}}\|^2\le Q^C\mu(\mathcal C)
\]
does not contradict a lower bound of size \(Q^{-C}\mu(\mathcal C)\).

The acceptable forbidden-channel mechanisms are:
\begin{enumerate}
\item a positive atomic upper
\[
  \sum_a|\Lambda_a V_{\mathrm{glob}}|^2\le Q^{-B}\mu(\mathcal C);
\]
\item a localizable identity family
\[
  |\Lambda_\psi V_{\mathrm{glob}}|^2\le Q^{-B}\|\psi\|^2;
\]
\item a coboundary/curl identity such as \(\delta^2=0\);
\item an approximate identity
\[
  \Lambda V_{\mathrm{glob}}=\operatorname{Err},
  \qquad
  \|\operatorname{Err}\|^2\le Q^{-B}\mu(\mathcal C).
\]
\end{enumerate}
Without one of these, finite D4 remains a detector, not a proof-closing
contradiction.  The current best finite-D4 candidate is
\[
  \Lambda_{\mathrm{tr},\mathrm{sing}}=(I-P_{\mathrm{sing}})T_{\mathrm{tr}\to\mathrm{sing}}P_{\mathrm{tr}}.
\]
\end{remark}

% === wip-current-postclosure-architecture-ledger ===
% Purpose:
%   Record the current post-rank-two-cleanup proof architecture.
% Status:
%   Active architecture ledger.

\begin{remark}[Current post-closure architecture]
\label{rem:current-postclosure-architecture}
After the rank-two closure cleanup, the live proof architecture is best recorded as
\[
  \boxed{
  \mathrm{Agg}L^2(q''_{\can})
  +
  \mathrm{Graph/O2O3}_{23}
  +
  \mathrm{ForbiddenSingularDiscardD4}.
  }
\]
The first term is branch-complete local scalar source energy:
\[
  E_{1,b}^{L^2}(m)
  =
  \int_{I_b}
  |\lambda_b^3 q''_{\can}(m+\lambda_bx)|^2\,dx.
\]
The second term is the finite O2/O3 nondegeneracy package, preferably expressed
through a Riesz graph model for
\[
  \mathbf C\cdot1+m_{\can}\mathbf C_R[x]
  \simeq
  m_{\can}\bigl(\mathbf C\cdot(1/m_{\can})+\mathbf C_R[x]\bigr).
\]
The third term is a forbidden D4 channel, currently represented by the
transition--singular discard operator
\[
  \Lambda_{\mathrm{tr},\mathrm{sing}}=(I-P_{\mathrm{sing}})T_{\mathrm{tr}\to\mathrm{sing}}P_{\mathrm{tr}}.
\]

The older row-normal rank-two chain remains a finite-audit submodule proving
that D4 supplies the plain wall \(H(\Gamma_G,\Gamma_{G,2})\).  The present
post-closure architecture adds the missing contradiction requirement: the
plain wall must be visible through an actual forbidden channel.
\end{remark}

The following negative results and audit boundaries are part of the current
proof state.

\begin{enumerate}[label=(\roman*)]
\item \emph{Pointwise base-locus rigidity is false.}
At a single point one may have
\[
\Gamma_G=0,\qquad \Gamma_{G,2}=0,\qquad u\ne0.
\]
Gate~2 must use interval or \(L^2\) structure.

\item \emph{Finite same-scale packet geometry does not prove Gate~1.}
Positive finite quadrature can annihilate finitely many source-curvature jets.
Gate~1 must use actual-zeta arithmetic, not only local Hadamard packet geometry.

\item \emph{Band-limited singular probes do not avoid long-prime control.}
In the deep singular branch, detecting the \(R\) sine moments
\[
\sin(k\theta),\qquad 3\le k\le R+3,
\]
requires analytic or band-limited probes with bandwidth at least polynomial in
\(R\).  The explicit formula then sees primes with
\[
\log n\lesssim R^C/c.
\]
Since \(cQ\ll(\log Q)^{-C_0}\), this is generally longer than the ordinary
packet-frequency range \(\log n\asymp Q\).  Therefore S1 is a separate
long-prime actual-zeta nonrecurrence theorem, not a refinement of the ordinary
Wronskian input.

\item \emph{Proxy source curvature is insufficient.}
A zero-kernel proxy, flexible Dirichlet polynomial, \(B_2\)-only diagnostic, or
affine-subtracted complement curvature is not the proof object.  Gate~1 and
Gate~2 use
\[
q''_{\can}=B_2+P_2+D_2+E_2.
\]

\item \emph{Generic determinant nonvanishing is insufficient.}
A Wronskian or finite determinant can be non-identically zero while becoming
too small on a one-parameter prime-log orbit.  The missing theorem is an
actual-zeta nonrecurrence/incidence lower bound, not formal rank.

\item \emph{Filtered-minor isolation is not locally available at current depth.}
Isolating one subset-product frequency would require an \(m\)-space filter
larger than the packet scale or derivative depth beyond the current budget.

\item \emph{Trace positivity is not coercivity.}
A positive explicit-formula or trace identity gives nonnegativity, but Gate~1
requires a polynomial lower bound on the packet interval.

\item \emph{The scalar \(I\)-coordinate is not a faithful carrier quotient.}
The order-nine scalar incidence wall appears after independent carrier
directions are collapsed into the same scalar coordinate.  Carrier or quotient
row energy should be used unless scalar \(I\)-coordinate output is explicitly
required.

\item \emph{Wall escape depends on row-output architecture.}
Wall escape supplies a non-wall sample somewhere in the packet, not a row
lower bound at a preassigned wall point.  Exact-point architectures require a
separate \(q_0\)-adic wall-row audit.
\end{enumerate}

\begin{remark}[Numerical and surrogate evidence archive]
\label{rem:numerical-surrogate-status}
The numerical and surrogate experiments in this draft are diagnostic only.
They do not prove Gate~1, O1, O2, O3, or any global incidence theorem.
Their value is to rule out obvious relaxed countermechanisms and to guide the
next arithmetic inputs.

The current evidence archive is:
\[
\begin{array}{c|c|c|c}
\text{branch} & \text{scope} & \text{result} & \text{interpretation}\\
\hline
\text{packet-frequency determinant stress}
&
R\le16,\ m_{\min}=50,\ \text{reflected/actual-like dual and }E_2
&
\text{Yellow}=0,\ \text{Red}=0
&
\text{no determinant-incidence collapse detected}\\
\hline
\text{local }Z_\alpha\text{-amplitude v1}
&
1152\text{ rows}
&
\text{Yellow}=0,\ \text{Red}=0
&
\text{no local amplitude/jet flatness detected}\\
\hline
\text{v2 direct jet/cell-sum hardening}
&
\text{adversarial targets from v1/determinant runs}
&
\text{no Yellow/Orange/Red in the recorded runs}
&
\text{no cell-mass collapse found at }Z_\alpha\text{-near zeros}\\
\hline
\text{multiblock }K=8\text{ cell test}
&
794\text{ rows with log-uniform support}
&
\text{Yellow}=0,\ \text{Orange}=0,\ \text{Red}=0
&
\text{supports multiblock fixed-width cell nonrecurrence}\\
\hline
\text{small }Z_\alpha\text{ adversarial rows}
&
|Z_\alpha|\sim10^{-15}\text{--}10^{-12}
&
\text{amplitude, jets, and cell mass stayed non-small}
&
Z_\alpha\approx0\not\Rightarrow\text{ local collapse}
\end{array}
\]

The strongest empirical message is:
\[
\boxed{
Z_\alpha(m_0)\approx0
\not\Rightarrow
\text{amplitude / jet / cell collapse}
}
\]
in the tested packet-frequency surrogates.

The numerical evidence also has clear limits:
\begin{enumerate}
\item It does not prove fixed-width cell nonrecurrence.
\item It does not prove cofactor-vector nonrecurrence.
\item It does not prove Wronskian angle nonalignment.
\item It does not prove singular anti-design exclusion.
\item It does not replace canonical transfer from model rows to the actual
object \(q''_{\can}=B_2+P_2+D_2+E_2\).
\end{enumerate}

The first constraint not broken by the adversarial relaxed models remains full
actual-zeta compatibility:
\[
q''_{\can}=B_2+P_2+D_2+E_2
\]
with fixed actual coefficients, completed dual/gamma/tail structure, and
certified error control.  Thus the numerical evidence supports prioritizing
the actual-zeta incidence/nonrecurrence theorem; it is not a substitute for
that theorem.
\end{remark}

\appendix

\section{Toy microscopic expansion}
\label{app:toy-microscopic-expansion}

\begin{lemma}[Uniform microscopic expansion of the toy local deformation]
\label{lem:toy-microscopic-expansion}
Fix a compact interval
\[
D=[d_-,d_+]\subset(0,\infty),
\]
and let
\[
\Delta_{\toy}(u,d;s)
\]
denote the toy local deformation matrix in the symmetric microscopic regime,
with off-line parameter $u$, separation parameter $d$, and microscopic node
$s$.

Then there exist constants $\rho_{\toy}>0$ and $u_0(D)>0$ such that for all
\[
d\in D,
\qquad
|s|<\rho_{\toy},
\qquad
|u|<u_0(D),
\]
one has the expansion
\[
\Delta_{\toy}(u,d;s)=u^2M(d)+u^4E(u,d;s),
\]
where $M(d)$ is the toy anomaly matrix, $E(u,d;s)$ is analytic in $s$
for $|s|<\rho_{\toy}$, and
\[
\sup_{\substack{d\in D\\ |s|<\rho_{\toy}\\ |u|<u_0(D)}} \|E(u,d;s)\|\ll 1.
\]
In particular, for the microscopic nodes
\[
s_j=\frac{j}{Q^2},
\qquad
j=1,\dots,N,
\]
with $N\le \rho_{\toy}Q^2$, one has uniformly
\[
\Delta_{\toy}(u,d;s_j)=u^2M(d)+u^4E_j(u,d),
\qquad
\sup_j \|E_j(u,d)\|\ll 1.
\]
\end{lemma}

\begin{proof}
The toy local deformation is obtained from an explicit finite-dimensional
off-line perturbation of the toy block. By construction it depends smoothly
(in fact analytically) on the parameters $(u,s,d)$ in a neighborhood of
\[
(u,s,d)=(0,0,d_0),
\qquad d_0\in D.
\]
At $u=0$, the deformation vanishes, and by the even symmetry of the off-line
split there is no linear term in $u$. Hence the first nontrivial term is
quadratic:
\[
\Delta_{\toy}(u,d;s)=u^2M(d,s)+O(u^4),
\]
uniformly on compact $d$-ranges and for $|s|$ in a fixed microscopic disk.

By the jet-limit computation for the toy model, the leading quadratic
coefficient is independent of the microscopic node $s$ and equals the toy
anomaly matrix $M(d)$. Thus
\[
M(d,s)\equiv M(d),
\]
and the expansion becomes
\[
\Delta_{\toy}(u,d;s)=u^2M(d)+u^4E(u,d;s),
\]
with $E$ analytic in $s$ and uniformly bounded on compact parameter
ranges. Restricting to the microscopic nodes $s_j=j/Q^2$ gives the stated
uniform bound for $E_j(u,d)$.
\end{proof}

\begin{remark}[Work in progress: toy microscopic expansion]
\label{rem:wip-toy-microscopic-expansion}
The proof above is schematic: it uses smoothness in $(u,s,d)$, vanishing at
$u=0$, and even symmetry in the off-line split to obtain the quadratic
leading term, and invokes the jet-limit computation to identify the
coefficient as the toy anomaly matrix $M(d)$. A later pass should make the
analyticity and uniform-bound steps quantitative by recording an explicit
Taylor remainder with a constant depending only on $D$.
\end{remark}

\section{Finite Fr\'echet coefficient calculus for the corrected-whitening block}
\label{sec:finite-frechet-coefficient-calculus}

This section records the finite Fr\'echet/Sylvester coefficient recurrences
and finite gates used by the corrected-whitening package frontier of
Remark~\ref{rem:wip-first-cubic-source-transfer-target}.  These are table
contracts, not completed source tables.  Their role is to make the finite
computations auditable by agents and scripts without re-opening convention
choices already fixed in the main text.

\begin{remark}[Fr\'echet coefficient normalization]
\label{rem:frechet-coefficient-normalization}
The finite tables in this section use homogeneous Taylor coefficients in the
ordinary \(z\)-expansion of the inverse-square-root transfer.  Thus
\(\Lambda\), \(H\), and \(G\) denote the actual coefficient contributions to
the expanded matrix series, not raw multilinear derivatives \(D,D^2,D^3\)
unless explicitly stated.

Equivalently, if a separate computation uses derivative-normalized tensors,
then the dictionary is
\[
H_{\mathrm{Taylor}}=\frac{1}{2!}H_{\mathrm{deriv}},
\qquad
G_{\mathrm{Taylor}}=\frac{1}{3!}G_{\mathrm{deriv}}.
\]
All determinant tests in this paper use the Taylor-coefficient convention
unless the contrary is explicitly stated.
\end{remark}

\begin{remark}[Stage 1 source-table status]
\label{rem:stage-one-source-table-status}
The finite Fr\'echet recurrences in this section become a proof tool only
after the required source tables are supplied in the same ordinary-\(z\),
pre-\(\Phi_K\), tagged normalization.  The current Stage 1 infrastructure
reduces the needed source input to a finite midpoint-jet theorem:
\[
\mathsf{Jet}_{\mathrm{Stage1}}
=
\left\{
q_0^{(0)},\ldots,q_0^{(9)};
\ (r_i^{[a]})^{(k)}
:
i=1,2,\quad a\in\mathcal A_{\mathrm{BC}},\quad 2\le k\le9
\right\}_{m}.
\]
The lead local obstruction currently uses \(a=2\).  Hence any Stage~1 table
omitting the \(a=2\) source-linear primitive cannot evaluate the present
Bottleneck-C frontier.

The table generator is infrastructure, not a zeta-source proof.  A promotion
requires a source theorem supplying these jets, or a reduction of them to
named baseline objects, in the same normalization used for
\(B_7^{\mathfrak f}\).  Scale estimates are not a substitute for these
finite coefficient tables.
\end{remark}

\subsection{Source-linear same-point primitives}
\label{subsec:source-linear-same-point-primitives}

For \(\varepsilon\in\{-,+\}\), identified with \(\varepsilon=-1,+1\), write
\[
X_{\varepsilon,n}
:=
[z^n]\mathfrak D_Q(\delta G_\varepsilon^{\lin})
=
\sum_{a=2}^{6}X_{\varepsilon,n}^a r_a.
\]
For \(n\le2\), the nonzero source-constrained coefficients are
\[
X_{\varepsilon,0}^2=\frac{1}{24\pi}(I-D),
\]
\[
X_{\varepsilon,1}^2=\frac{\varepsilon}{4\pi Q}S,
\qquad
X_{\varepsilon,1}^3=\frac{\varepsilon}{48\pi Q^2}(I-D),
\]
\[
X_{\varepsilon,2}^2
=
\left(
\frac{1}{8\pi Q^2}
+
\frac{q^2}{32\pi Q^4}
\right)I
+
\left(
\frac{1}{8\pi Q^2}
-
\frac{q^2}{32\pi Q^4}
\right)D,
\]
\[
X_{\varepsilon,2}^3=\frac{1}{16\pi Q^3}S,
\qquad
X_{\varepsilon,2}^4=\frac{1}{192\pi Q^4}(I-D).
\]
All other \(X_{\varepsilon,n}^a\) with \(n\le2\), \(2\le a\le6\), vanish.

\begin{remark}
\label{rem:source-linear-primitives-residual-normal-form}
These coefficients are not free symbolic inputs.  They are the residual-source
normal form of the source-linear same-point perturbation.
\end{remark}

\subsection{Baseline mixed block coefficients}
\label{subsec:baseline-mixed-block-coefficients}

Let
\[
s=\frac{z}{Q^2},\qquad
q_\pm(z)=q_0\left(m\pm\frac{z}{2Q^2}\right),
\]
\[
\Delta_0(z)=\Phi_0(t_-)-\Phi_0(t_+),
\qquad
C(z)=\cos\Delta_0(z),\quad S(z)=\sin\Delta_0(z).
\]
The baseline mixed block is
\[
N^{(0)}(z)=\frac1\pi
\begin{pmatrix}
-\dfrac{2S}{s}
&
\dfrac{S+q_+sC}{s^2}
\\[1.4ex]
-\dfrac{S+q_-sC}{s^2}
&
\dfrac{(q_-+q_+)sC+(2-q_-q_+s^2)S}{2s^3}
\end{pmatrix}.
\]
For \(j\ge0\), write
\[
N^{(0)}_j=[z^j]N^{(0)}(z).
\]
Then \(N^{(0)}_{0..4}\) are obtained entrywise from the coefficient formulas
\[
(N^{(0)}_j)_{11}
=
-\frac{2Q^2}{\pi}[z^{j+1}]S(z),
\]
\[
(N^{(0)}_j)_{12}
=
\frac{Q^4}{\pi}
\left(
[z^{j+2}]S(z)
+
\frac1{Q^2}[z^{j+1}](q_+(z)C(z))
\right),
\]
\[
(N^{(0)}_j)_{21}
=
-\frac{Q^4}{\pi}
\left(
[z^{j+2}]S(z)
+
\frac1{Q^2}[z^{j+1}](q_-(z)C(z))
\right),
\]
\[
(N^{(0)}_j)_{22}
=
\frac{Q^6}{2\pi}
\left(
\frac1{Q^2}[z^{j+2}]((q_-+q_+)C)
+
2[z^{j+3}]S
-
\frac1{Q^4}[z^{j+1}](q_-q_+S)
\right).
\]

\subsection{First Fr\'echet coefficient table}
\label{subsec:first-frechet-coefficient-table}

Let \(P_{\varepsilon,0}\) be the zeroth square-root coefficient of
\(G_\varepsilon^{(0)}\), and let \(\Sigma_\varepsilon(C)\) denote the unique
solution of
\[
P_{\varepsilon,0}Z+ZP_{\varepsilon,0}=C.
\]
For each source index \(a\), define \(S_{\varepsilon,n}^a\) recursively by
\[
S_{\varepsilon,n}^a
=
\Sigma_\varepsilon
\left(
X_{\varepsilon,n}^a
-
\sum_{j=1}^{n}
\bigl(
P_{\varepsilon,j}S_{\varepsilon,n-j}^a
+
S_{\varepsilon,n-j}^aP_{\varepsilon,j}
\bigr)
\right).
\]
Then
\[
\Lambda_{\varepsilon,n}^{a}
=
-\sum_{i+j+k=n}
W_{\varepsilon,i}S_{\varepsilon,j}^aW_{\varepsilon,k}.
\]
The required ranges used below include
\[
\Lambda^-_{2,0..4},\qquad
\Lambda^-_{3,1..5},\qquad
\Lambda^-_{4,2..6}.
\]

\subsection{Second Fr\'echet coefficient table}
\label{subsec:second-frechet-coefficient-table}

For an ordered source pair \((a,b)\), define \(T_{\varepsilon,n}^{ab}\) by
\[
T_{\varepsilon,n}^{ab}
=
\Sigma_\varepsilon
\left(
-\sum_{i+j=n}
(S_{\varepsilon,i}^aS_{\varepsilon,j}^b+
 S_{\varepsilon,i}^bS_{\varepsilon,j}^a)
-
\sum_{j=1}^{n}
(P_{\varepsilon,j}T_{\varepsilon,n-j}^{ab}
+
T_{\varepsilon,n-j}^{ab}P_{\varepsilon,j})
\right).
\]
The corresponding second Fr\'echet coefficient is
\[
H_{\varepsilon,n}^{ab}
=
\sum_{i+j+k+\ell+m=n}
\left(
W_{\varepsilon,i}S_{\varepsilon,j}^bW_{\varepsilon,k}
S_{\varepsilon,\ell}^aW_{\varepsilon,m}
+
W_{\varepsilon,i}S_{\varepsilon,j}^aW_{\varepsilon,k}
S_{\varepsilon,\ell}^bW_{\varepsilon,m}
\right)
-
\sum_{i+j+k=n}
W_{\varepsilon,i}T_{\varepsilon,j}^{ab}W_{\varepsilon,k}.
\]
The five-family endpoint gate uses
\[
H^-_{32,1..5},\qquad
H^-_{34,3..7},\qquad
H^-_{24,2..6}.
\]

\subsection{Third Fr\'echet coefficient table}
\label{subsec:third-frechet-coefficient-table}

For source triple \((a,b,c)\), define \(U_{\varepsilon,n}^{abc}\) by the third
Sylvester recursion
\[
P_{\varepsilon,0}U_{\varepsilon,n}^{abc}
+
U_{\varepsilon,n}^{abc}P_{\varepsilon,0}
=
-\sum_{i+j=n}
\Bigl(
S_{\varepsilon,i}^{a}T_{\varepsilon,j}^{bc}
+S_{\varepsilon,i}^{b}T_{\varepsilon,j}^{ac}
+S_{\varepsilon,i}^{c}T_{\varepsilon,j}^{ab}
+
T_{\varepsilon,i}^{bc}S_{\varepsilon,j}^{a}
+T_{\varepsilon,i}^{ac}S_{\varepsilon,j}^{b}
+T_{\varepsilon,i}^{ab}S_{\varepsilon,j}^{c}
\Bigr)
-
\sum_{j=1}^{n}
\bigl(
P_{\varepsilon,j}U_{\varepsilon,n-j}^{abc}
+
U_{\varepsilon,n-j}^{abc}P_{\varepsilon,j}
\bigr).
\]
The all-left third Fr\'echet coefficient \(G^-_{abc,n}\) is then obtained from
the coefficient of the cubic source monomial in
\[
(G_-^{(0)}+\text{source})^{-1/2},
\]
equivalently by the coefficient recursion determined by
\(U_{\varepsilon,n}^{abc}\), the lower \(S\)- and \(T\)-coefficients, and the
inverse coefficients \(W_{\varepsilon,n}\).

The seven-family audit for the first surviving residual requires
\[
G^-_{324,3..7}
\quad\text{and}\quad
G^+_{324,3..7}.
\]

\begin{remark}[Status of the third-Fr\'echet table]
\label{rem:status-third-frechet-table}
The recursion above defines the finite third-Fr\'echet coefficient table
that a raw seven-family audit would require.  It does not assert that the
raw \(L_3N_0R_0/L_0N_0R_3\) completion is quotient-gauge.  Under the
current package architecture, this table is a diagnostic tool for raw local
coefficients.  The theorem object is the normalized corrected package
\(C_{\mathcal W}^{\corr}\) of
Definition~\ref{def:normalized-corrected-whitening-package-representative},
not the raw seven-family Fr\'echet sum.
\end{remark}

\begin{remark}[Convention: transfer degree versus source bidegree]
\label{conv:transfer-degree-versus-source-bidegree}
We write \(d\) for homogeneous transfer degree in
\[
\mathcal T(X)=\sum_{d\ge1}\mathcal T_d(X),
\]
and reserve \((p,q)\) for source bidegree in the source-tagged expansion
\[
C(\tau_1,\tau_2)=\sum_{p,q}\tau_1^p\tau_2^qC_{p,q}.
\]
\end{remark}

\begin{remark}[Fixed-sector projection in the moving fiber]
\label{rem:fixed-sector-projection-moving-fiber}
When the midpoint is \(h\), the fixed-sector projection is denoted
\(\pi_{\mathfrak f_h}\).  The transpose-invariance identity
\[
\pi_{\mathfrak f_h}(Z^T)=\pi_{\mathfrak f_h}(Z)
\]
is applied in the moving midpoint fiber before midpoint trivialization to
\(m\).
\end{remark}

\begin{lemma}[Universal first-collision formula by transfer degree]
\label{lem:universal-first-collision-formula-transfer-degree}
Let
\[
\mathcal T(X)=\sum_{d\ge1}\mathcal T_d(X)
\]
be the homogeneous transfer expansion of the corrected-whitening map.  On
the collision chart
\[
h_1=h-\frac{\delta}{2},\qquad
h_2=h+\frac{\delta}{2},\qquad
a_2=-a_1(1+\kappa\delta),
\]
the first-collision contribution of homogeneous transfer degree \(d\) has
the form
\[
\boxed{
F_{1,d}
=
[z^7]\pi_{\mathfrak f_h}
\left[
(-1)^{d+1}d\kappa\,\mathcal T_d(X^d)
+
\frac d2(1-(-1)^d)\mathcal T_d(X^{d-1},X_h)
\right].
}
\]
Thus even \(d\) contributes only the amplitude term
\[
[z^7]\pi_{\mathfrak f_h}\mathcal T_d(X^d),
\]
whereas odd \(d\) contributes both the amplitude term and the
source-motion term
\[
[z^7]\pi_{\mathfrak f_h}\mathcal T_d(X^{d-1},X_h).
\]
\end{lemma}

\begin{proof}[Proof sketch]
The factor \(a_2=-a_1(1+\kappa\delta)\) produces the amplitude derivative
term proportional to \(d\kappa\).  The endpoint displacement
\(h_1=h-\delta/2\), \(h_2=h+\delta/2\) produces the midpoint-motion
insertion.  The parity of \(d\) controls whether the endpoint-motion
contributions cancel or add: for even \(d\) the \(X_h\)-insertion cancels,
while for odd \(d\) it appears with coefficient \(d\).
\end{proof}

\subsection{Fourth-transfer coefficient table for Gate B}
\label{subsec:fourth-transfer-coefficient-table-gate-b}

The corrected-whitening map is
\[
\mathcal W(G_-,N,G_+)=G_-^{-1/2}NG_+^{-1/2}.
\]
Let a labelled source input be
\[
X_j=(A_j,Y_j,B_j),
\]
where
\[
A_j=\delta G_{-,j}^{\lin},\qquad
Y_j=\delta N_j^{\lin},\qquad
B_j=\delta G_{+,j}^{\lin}.
\]
Write
\[
(G_-^{(0)}+A)^{-1/2}
=
L_0+L_1[A]+L_2[A,A]+L_3[A,A,A]+L_4[A,A,A,A]+\cdots,
\]
and
\[
(G_+^{(0)}+B)^{-1/2}
=
R_0+R_1[B]+R_2[B,B]+R_3[B,B,B]+R_4[B,B,B,B]+\cdots.
\]
The middle block is source-linear:
\[
N=N_0+Y.
\]
For a labelled four-source set \(J=\{1,2,3,4\}\), the fourth homogeneous
corrected-whitening transfer is
\[
\boxed{
\mathcal T_4(X_1,X_2,X_3,X_4)
=
\sum_{I\subseteq J}
L^-_I\,N_0\,R^+_{J\setminus I}
+
\sum_{m\in J}
\sum_{I\subseteq J\setminus\{m\}}
L^-_I\,Y_m\,R^+_{J\setminus(I\cup\{m\})}.
}
\]
The first sum consists of
\[
L_4N_0R_0,\quad
L_3N_0R_1,\quad
L_2N_0R_2,\quad
L_1N_0R_3,\quad
L_0N_0R_4,
\]
and the second sum consists of
\[
L_3YR_0,\quad
L_2YR_1,\quad
L_1YR_2,\quad
L_0YR_3.
\]
The ordinary \(z^7\)-coefficient is
\[
\boxed{
[z^7]\mathcal T_4(J)
=
\sum_{I\subseteq J}
\sum_{\alpha+\beta+\gamma=7}
L^-_{I,\alpha}N^{(0)}_{\beta}R^+_{J\setminus I,\gamma}
+
\sum_{m\in J}
\sum_{I\subseteq J\setminus\{m\}}
\sum_{\alpha+\beta+\gamma=7}
L^-_{I,\alpha}Y_{m,\beta}R^+_{J\setminus(I\cup\{m\}),\gamma}.
}
\]
For a source multiset, \(\operatorname{Sym}\mathcal T_4\) denotes the
labelled sum divided by the automorphism factor of the multiset.

\begin{lemma}[All-subset midpoint parity]
\label{lem:all-subset-midpoint-parity}
Let \(K\) be any labelled subset of homogeneous source inputs and define
\[
|K|:=\sum_{g\in K}g.
\]
At leading homogeneous order, the left and right same-point
inverse-square-root coefficients satisfy
\[
\boxed{
R_K^+=(-1)^{|K|}L_K^-.
}
\]
\end{lemma}

\begin{remark}[Scope of all-subset midpoint parity]
\label{rem:scope-all-subset-midpoint-parity}
Lemma~\ref{lem:all-subset-midpoint-parity} is a leading homogeneous
coefficient statement.  It uses the midpoint chart
\[
t_\pm=h\pm\frac{z}{2Q^2}
\]
and the Taylor-coefficient Fr\'echet convention.  If one works with
derivative-normalized Fr\'echet tensors, ordered unsymmetrized source
slots, or non-leading ordinary coefficients, the parity factors and
multiplicities must be translated back to this convention before applying
the subset-complement cancellation.
\end{remark}

\begin{lemma}[Symmetry of same-point subset coefficients]
\label{lem:symmetry-same-point-subset-coefficients}
For every labelled subset \(K\),
\[
(L_K^-)^T=L_K^-,
\qquad
(R_K^+)^T=R_K^+.
\]
Moreover \(N_0[0]^T=N_0[0]\).
\end{lemma}

\begin{lemma}[Leading \(Y\)-middle parity]
\label{lem:leading-y-middle-parity}
The leading \(Y\)-middle input satisfies
\[
Y^{[a]}_{\mathrm{lead}}=0\qquad(a\ \mathrm{odd}),
\]
and
\[
(Y^{[a]}_{\mathrm{lead}})^T=Y^{[a]}_{\mathrm{lead}}
\qquad(a\ \mathrm{even}).
\]
\end{lemma}

\begin{remark}[Scope of the leading \(Y\)-middle parity table]
\label{rem:scope-leading-y-middle-parity-table}
Lemma~\ref{lem:leading-y-middle-parity} concerns the leading homogeneous
ordinary-order coefficient of the source-linear mixed middle input in the
active homogeneous scalar-grade convention.  It is not a statement about
every ordinary coefficient of a non-homogeneous matrix-output \(M^{[5]}\)
table.

In particular, older coefficient audits that found antisymmetric
\(J\)-channel middle terms belong to the quarantined non-homogeneous
matrix-output convention.  They do not contradict the leading homogeneous
\(Y\)-middle parity table used in Gate~B.
\end{remark}

\begin{theorem}[Odd total-grade subset-complement cancellation]
\label{thm:odd-total-grade-subset-complement-cancellation}
Let
\[
J=(g_1,\ldots,g_d)
\]
be a labelled homogeneous source multiset with total grade
\[
|J|=\sum_i g_i=7.
\]
Then the leading ordinary-order amplitude contribution of the
corrected-whitening transfer vanishes in fixed sector:
\[
\boxed{
[z^7]\pi_{\mathfrak f_h}
\mathcal T_d(X^{[g_1]},\ldots,X^{[g_d]})=0.
}
\]
\end{theorem}

\begin{proof}
For the \(N_0\)-middle part, the leading contribution is
\[
\mathcal A_{N_0}(J)
=
[z^7]\pi_{\mathfrak f_h}
\sum_{I\subseteq J}L_I^-N_0[0]R^+_{J\setminus I}.
\]
By Lemma~\ref{lem:all-subset-midpoint-parity},
\[
R^+_{J\setminus I}=(-1)^{|J|-|I|}L^-_{J\setminus I}.
\]
Since \(|J|=7\), this sign is \(-(-1)^{|I|}\).  Pair \(I\) with
\(J\setminus I\).  The paired contribution is
\[
(-1)^{|I|}
\pi_{\mathfrak f_h}
\left(
-L_I^-N_0[0]L^-_{J\setminus I}
+
L^-_{J\setminus I}N_0[0]L_I^-
\right).
\]
The two matrix products are transposes, and \(\pi_{\mathfrak f_h}\) is
transpose-invariant.  Hence the pair cancels.

For the \(Y\)-middle part, fix the source label \(m\in J\) in the middle
and write \(a=g_m\), \(K=J\setminus\{m\}\).  If \(a\) is odd, then
\(Y_m=0\) by Lemma~\ref{lem:leading-y-middle-parity}.  If \(a\) is even,
then \(Y_m\) is symmetric and \(|K|=7-a\) is odd, so the same
subset-complement pairing cancels the fixed-sector projection.
\end{proof}

\begin{lemma}[Grade-preserving source-motion cancellation]
\label{lem:grade-preserving-source-motion-cancellation}
Let \(J=(g_1,\ldots,g_d)\) be a labelled homogeneous source multiset with
\[
|J|=\sum_i g_i=7.
\]
Replace one source slot by the grade-preserving part of its midpoint
derivative,
\[
(\partial_h\operatorname{Gr}_{g_i}r)_{\mathrm{same}},
\]
and keep all other source slots fixed.  Then the corresponding fixed-sector
source-motion contribution vanishes:
\[
[z^7]\pi_{\mathfrak f_h}
\mathcal T_d
\left(
X^{[g_1]},\ldots,
X^{[g_i]}_{\mathrm{same},h},
\ldots,
X^{[g_d]}
\right)
=0.
\]
\end{lemma}

\begin{proof}[Proof sketch]
The grade-preserving derivative has the same ordinary source parity and
the same endpoint transpose parity as the undifferentiated homogeneous
grade in that slot.  Therefore the same subset-complement pairing used in
Theorem~\ref{thm:odd-total-grade-subset-complement-cancellation} applies:
each labelled subset contribution is paired with its complement, the two
matrix products are transposes, and the transpose-invariant projection
\(\pi_{\mathfrak f_h}\) kills the aggregate fixed-sector contribution.
\end{proof}

\begin{corollary}[Amplitude closure of higher transfers]
\label{cor:amplitude-closure-higher-transfers}
The amplitude part of every Gate~B transfer of total homogeneous grade
\(7\) vanishes in \(\mathfrak f_h\).  In particular, the aggregate
fourth-transfer positive-grade contribution
\[
\sum_{J\in\{(1,1,1,4),(1,1,2,3),(1,2,2,2)\}}
[z^7]\pi_{\mathfrak f_h}\mathcal T_4(J)
\]
vanishes.

The same argument also kills the first lower-shadow fourth-transfer
aggregate
\[
[z^7]\pi_{\mathfrak f_h}\mathcal T_4(0,1,1,5)=0.
\]
More generally, the subset-complement argument applies to every labelled
multiset of total homogeneous grade \(7\), including zero-padded
lower-shadow multisets, provided the full Taylor source tower is
differentiated before ordinary-order truncation.

Thus the amplitude parts also vanish for
\[
d=5:\quad (1,1,1,1,3),\ (1,1,1,2,2),
\]
\[
d=6:\quad (1,1,1,1,1,2),
\]
and
\[
d=7:\quad (1,1,1,1,1,1,1).
\]
\end{corollary}

\begin{lemma}[Raw midpoint derivative of homogeneous scalar grades]
\label{lem:raw-midpoint-derivative-homogeneous-scalar-grades}
For
\[
\operatorname{Gr}_g r(h;t)
=
\frac{r^{(g+2)}(h)}{(g+2)!}(t-h)^{g+2},
\]
one has
\[
\partial_h\operatorname{Gr}_g r
=
(\partial_h\operatorname{Gr}_g r)_{\mathrm{same}}
+
(\partial_h\operatorname{Gr}_g r)_{\mathrm{lower}},
\]
where
\[
(\partial_h\operatorname{Gr}_g r)_{\mathrm{same}}
=
\frac{r^{(g+3)}(h)}{(g+2)!}(t-h)^{g+2},
\]
and
\[
(\partial_h\operatorname{Gr}_g r)_{\mathrm{lower}}
=
-\frac{r^{(g+2)}(h)}{(g+1)!}(t-h)^{g+1}.
\]
\end{lemma}

\begin{theorem}[Moving homogeneous source-frame telescoping]
\label{thm:moving-homogeneous-source-frame-telescoping}
For every interior grade \(g\ge1\),
\[
\boxed{
(\partial_h\operatorname{Gr}_{g-1}r)_{\mathrm{same}}
+
(\partial_h\operatorname{Gr}_{g}r)_{\mathrm{lower}}
=0.
}
\]
Consequently, after labelled symmetrization, the grade-lowering part of
the raw source-motion insertion cancels source-slot by source-slot
against the same-basis derivative of the adjacent lower grade.
\end{theorem}

\begin{proof}
By Lemma~\ref{lem:raw-midpoint-derivative-homogeneous-scalar-grades},
\[
(\partial_h\operatorname{Gr}_{g-1}r)_{\mathrm{same}}
=
\frac{r^{(g+2)}(h)}{(g+1)!}(t-h)^{g+1},
\]
whereas
\[
(\partial_h\operatorname{Gr}_{g}r)_{\mathrm{lower}}
=
-\frac{r^{(g+2)}(h)}{(g+1)!}(t-h)^{g+1}.
\]
The signs and factorials agree exactly.  The cancellation is
labelled-slot local, so it survives the symmetrized multilinear transfer
expansion.
\end{proof}

\begin{lemma}[Affine boundary removal in the moving source frame]
\label{lem:affine-boundary-removal-moving-source-frame}
In the \(\mathfrak D_Q\)-scaled corrected source block, admissible source
inputs are \(2\)-jet normalized at the moving midpoint:
\[
f(h)=0,\qquad f'(h)=0.
\]
Equivalently, the source-frame projection is
\[
\Pi_{\ge2,h}f(t)=f(t)-f(h)-f'(h)(t-h).
\]
Thus constant and affine modes are not independent raw source inputs.

This statement is a source-admissibility assertion for the holomorphic
\(\mathfrak D_Q\)-scaled corrected source variables.  It is not a license to
discard lower homogeneous grades.  It uses the residual-source convention in
which the moving midpoint affine jet has already been subtracted; without
that convention an affine component would not be killed by a formal
coefficient argument alone.

This is not a blanket assertion that all lower-shadow grades are discarded.
The grade \(0,2,4\) source directions may appear in the moving Taylor tower.
They are handled by the interior telescoping identity when they arise from
source-motion, not by reclassifying them as baseline data.  The only
boundary mode removed here is the affine grade \(-1\) term.

For the homogeneous scalar-grade basis, the remaining lower boundary
\[
(\partial_h\operatorname{Gr}_{0}r)_{\mathrm{lower}}
=
-r^{(2)}(h)(t-h)
\]
is affine and is killed by \(\Pi_{\ge2,h}\).  Therefore no grade \(-1\)
source term enters
\[
[z^7][\delta^1]\pi_{\mathfrak f_h}
(1-\Pi_{1,1})_{\mathrm{mix}}C_{\mathcal W}^{\mathrm{raw}}.
\]
The Taylor source tower is differentiated before ordinary-order
truncation; in particular, grade \(8\) is included when proving ordinary
order \(7\) cancellation and truncation is applied only after
\(\partial_h\).
\end{lemma}

\begin{proof}[Proof sketch]
An affine perturbation \(\psi(t)=t-h\) satisfies
\[
\psi(t_\pm)=\pm\frac{z}{2Q^2},\qquad
\psi'(t_\pm)=1,\qquad
\psi''(t_\pm)=0.
\]
In the \(\mathfrak D_Q\)-scaled source variables
\[
d_\pm(z)=\frac{z^2}{Q^4}U_\pm(z),
\qquad
e_\pm(z)=\frac{z}{Q^2}E_\pm(z),
\]
an affine source would force \(U_\pm\) and \(E_\pm\) to have poles in
\(z\).  Hence affine modes are outside the admissible holomorphic source
sector.  This is the source-frame reason the homogeneous scalar-grade
tower starts at \(g=0\), i.e. at the second derivative.
\end{proof}

\begin{remark}[No premature ordinary-order truncation]
\label{rem:no-premature-ordinary-order-truncation-gate-b}
The cancellation in
Theorem~\ref{thm:moving-homogeneous-source-frame-telescoping} is a statement
about the full moving Taylor source tower.  For ordinary order \(7\), one
must include the next homogeneous grade before differentiating: the lower
part of \(\partial_h\operatorname{Gr}_{8}r\) cancels the apparent upper
boundary coming from the same-basis derivative of
\(\operatorname{Gr}_{7}r\).  If the tower is truncated at grade \(7\) before
\(\partial_h\) is applied, this cancellation is artificially destroyed.

Thus ``through ordinary order \(7\)'' means: differentiate the full
homogeneous source tower first, perform the adjacent-grade telescoping, and
only then discard ordinary powers above \(7\).
\end{remark}

\begin{proposition}[Decomposition of the Gate B first-collision coefficient]
\label{prop:decomposition-gate-b-first-collision}
Through ordinary order \(7\), the Gate~B first-collision coefficient
decomposes as
\[
F_1
=
F_1^{\mathrm{GateA}}
+
F_1^{\mathrm{amp}}
+
F_1^{X_h,\mathrm{same}}
+
F_1^{X_h,\mathrm{lower}}
+
F_1^{\mathrm{bdry}},
\]
where:
\begin{enumerate}
\item \(F_1^{\mathrm{GateA}}\) is the leading homogeneous cubic
contribution;
\item \(F_1^{\mathrm{amp}}\) is the amplitude contribution of higher
transfer terms \(\mathcal T_d(X^d)\) with total homogeneous grade \(7\);
\item \(F_1^{X_h,\mathrm{same}}\) is the grade-preserving part of the
source-motion insertion;
\item \(F_1^{X_h,\mathrm{lower}}\) is the grade-lowering part of the
source-motion insertion;
\item \(F_1^{\mathrm{bdry}}\) consists of the lower affine boundary and
the upper truncation boundary of the moving Taylor source frame.
\end{enumerate}
Under the hypotheses of Gate~A closure, all-subset midpoint parity,
leading \(Y\)-middle parity, moving source-frame telescoping, and
\(2\)-jet source-frame normalization, every term on the right vanishes.
\end{proposition}

\begin{proof}[Proof sketch]
The decomposition follows from
Lemma~\ref{lem:universal-first-collision-formula-transfer-degree}.  Gate~A
handles the leading cubic term.
Theorem~\ref{thm:odd-total-grade-subset-complement-cancellation} handles
the amplitude part.  The grade-preserving component of \(X_h\) is handled
by Lemma~\ref{lem:grade-preserving-source-motion-cancellation}.  The
grade-lowering part of \(X_h\) is
handled by Theorem~\ref{thm:moving-homogeneous-source-frame-telescoping}.
The lower affine boundary is removed by
Lemma~\ref{lem:affine-boundary-removal-moving-source-frame}, and the
upper boundary is avoided by differentiating the full Taylor source tower
before truncation.
\end{proof}

\begin{corollary}[Conditional Gate B raw edge-divisibility through ordinary order \(7\)]
\label{cor:gate-b-raw-edge-divisibility-order-seven}
Assume that the following source-frame hypotheses have been verified for the
actual \(\mathfrak D_Q\)-scaled corrected source model, not merely for a
formal homogeneous Taylor model:
\begin{enumerate}
\item Gate~A aggregate closure;
\item all-subset midpoint parity for same-point whitening coefficients;
\item symmetry of same-point whitening subset coefficients and of
\(N_0[0]\);
\item transpose-invariance of \(\pi_{\mathfrak f_h}\);
\item the leading \(Y\)-middle parity table;
\item moving homogeneous source-frame telescoping;
\item affine boundary removal by \(2\)-jet source-frame normalization;
and
\item differentiation of the full Taylor source tower before
ordinary-order truncation.
\end{enumerate}
Then
\[
\boxed{
[z^7][\delta^1]\,
\pi_{\mathfrak f_h}
(1-\Pi_{1,1})_{\mathrm{mix}}C_{\mathcal W}^{\mathrm{raw}}=0.
}
\]
Under these same hypotheses, equivalently,
\[
\boxed{
(1-\Pi_{1,1})_{\mathrm{mix}}C_{\mathcal W}^{\mathrm{raw}}
=
O(a_1\delta^2)
}
\]
through ordinary order \(7\) in fixed sector.
\end{corollary}

\begin{remark}[Dependency ledger for Gate B]
\label{rem:dependency-ledger-gate-b}
Corollary~\ref{cor:gate-b-raw-edge-divisibility-order-seven} is a raw
edge-divisibility result, not a post-edge quotient compatibility result.
Its inputs are logically independent and should not be silently replaced by
quotient-level language.

The dependencies are:
\begin{enumerate}
\item Gate~A aggregate closure, which removes the leading homogeneous
cubic non-\((1,1)\) contribution.
\item All-subset midpoint parity for the same-point whitening
coefficients, which pairs labelled subsets with their complements.
\item Symmetry of the relevant same-point subset coefficients and of
\(N_0[0]\).
\item Transpose-invariance of the moving fixed-sector projection
\(\pi_{\mathfrak f_h}\).
\item Leading \(Y\)-middle parity, which removes odd \(Y\)-middle entries
and makes the even entries compatible with subset-complement
cancellation.
\item Moving homogeneous source-frame telescoping, which cancels the
grade-lowering source-motion term against the same-basis derivative of
the adjacent lower grade.
\item \(2\)-jet source-frame normalization, which removes the affine
lower boundary produced by differentiating \(\operatorname{Gr}_0\).
\item Differentiation of the full Taylor source tower before
ordinary-order truncation, which prevents an artificial upper-boundary
residual.
\end{enumerate}

If any one of these inputs is removed, the statement should be demoted
back to a conditional Gate~B target rather than used as a proved package
construction.  Thus the corollary should be cited as a conditional raw
edge-divisibility theorem: it closes Gate~B for the corrected package only
after the listed parity, source-frame, and full-tower-before-truncation
hypotheses have been identified with the actual source model.
\end{remark}

\begin{remark}[Gate B remains conditional on source-model hardening]
\label{rem:gate-b-source-model-hardening}
Corollary~\ref{cor:gate-b-raw-edge-divisibility-order-seven} is a proved
implication from its stated source-frame hypotheses.  It should not yet be
read as an unconditional theorem of
Definition~\ref{def:tagged-two-atom-prewhitening-source-block} alone.

The tagged source block gives pre-whitening \(\mathcal K\)-linear additivity.
The remaining source-model obligation is to prove that the actual
\(\mathfrak D_Q\)-scaled corrected source variables feed the nonlinear
corrected-whitening transfer
\[
\mathcal W(G_-,N,G_+)=G_-^{-1/2}NG_+^{-1/2}
\]
with the same homogeneous scalar grades, midpoint subset parities,
source-motion telescoping, affine-boundary removal, and
full-tower-before-truncation used in Gate~B.

Thus the corollary is presently conditional in exactly this sense: the raw
edge-divisibility conclusion follows once these source-frame hypotheses are
verified for the actual corrected package, rather than only for the formal
homogeneous Taylor tower.
\end{remark}

\begin{remark}[Negative findings from the Gate A/Gate B audit]
\label{rem:negative-findings-gate-a-gate-b-audit}
The Gate A/Gate B audit rules out several tempting shortcuts.

\begin{enumerate}
\item \textbf{Placement-level vanishing is false.}
Individual transfer placements can have nonzero fixed-sector components.
The cancellation is only after labelled subset-complement aggregation and
transpose-invariant fixed-sector projection.

\item \textbf{Family-by-family gauge is too strong.}
The package target is aggregate fixed-sector cancellation or quotient
invisibility.  It does not require each Fr\'echet family to be separately
proportional to \(A_5^{\mathfrak f}\).

\item \textbf{Post-edge quotient compatibility is too late.}
A statement such as
\[
Q_{7,m}^q(C)=Q_{7,m}^q(\Pi_{1,1}C)
\]
cannot be used to prove raw edge-divisibility unless it is reformulated at
the raw edge layer.

\item \textbf{Determinant scalarization is only a good-patch coordinate.}
The quotient object is
\[
\mathfrak f_m/\mathbf C A_5^{\mathfrak f}(m).
\]
The scalar determinant test is a coordinate on a nondegenerate patch, not
the definition of the quotient.

\item \textbf{The matrix-output \(M^{[5]}\) convention is quarantined.}
The active proof uses homogeneous scalar grades.  Replacing the grade-five
source by an ordinary matrix-output order-\(5\) coefficient reopens the
\(Y\)-slot families and requires a separate no-double-counting theorem
against grades \(0,2,4\).

\item \textbf{The \(Z_j\)-table is post-Gate-B, not the Gate-B
obstruction.} After raw edge-divisibility, \(Z_j[M]\) is the next
pre-trivialization table for the connection calculation.  It is not a
missing proof of the \(\delta^2\)-edge law itself.
\end{enumerate}
\end{remark}

\begin{remark}[Mechanism of the Gate B cancellation]
\label{rem:mechanism-gate-b-cancellation}
The Gate~B cancellation is not a brute-force finite coefficient accident.
It is the combination of three structural symmetries: odd total-grade
subset-complement cancellation for amplitude terms, moving Taylor-frame
telescoping for source-motion lower terms, and \(2\)-jet source-frame
normalization for the affine boundary.  These mechanisms act before
\(\Phi_K\), before determinant scalarization, and before quotient
extraction.  The resulting raw edge-divisibility is exactly the input
needed to construct the normalized order-\(7\) corrected package.
\end{remark}

\begin{remark}[What remains attackable in Gate B]
\label{rem:what-remains-attackable-in-gate-b}
After Corollary~\ref{cor:gate-b-raw-edge-divisibility-order-seven}, the
leading cubic inventory, the \(d\ge4\) amplitude terms, and the raw
source-motion lower-shadow terms should not be reopened as independent
obstructions unless one of the stated hypotheses is rejected.  A valid
attack on Gate~B must now target one of the hypotheses:
\[
\text{all-subset midpoint parity},\quad
\text{leading \(Y\)-middle parity},\quad
\text{moving source-frame telescoping},\quad
\text{affine boundary removal},\quad
\text{full-tower-before-truncation}.
\]
In particular, a nonzero unsymmetrized transfer placement is not a Gate~B
counterexample unless it survives the labelled aggregate and the
source-frame telescoping procedure.
\end{remark}

\begin{remark}[Gate A is not Gate B]
\label{rem:gate-a-is-not-gate-b}
Gate~A is the leading homogeneous cubic audit.  Gate~B is the raw
edge-divisibility theorem for the projected corrected-whitening package.
Gate~A closure is a necessary local input, but it is not the same
statement as edge-normalized package construction.
Corollary~\ref{cor:gate-b-raw-edge-divisibility-order-seven} records the
additional transfer and source-frame mechanisms that promote the Gate~A
closure to the order-\(7\) raw edge-divisibility statement.
\end{remark}

\begin{remark}[Unsymmetrized subchannels versus aggregate cancellation]
\label{rem:unsymmetrized-subchannels-versus-aggregate-cancellation}
Individual transfer placements can have nonzero fixed-sector components.
This does not contradict the aggregate cancellations above.  The
cancellations pair each labelled subset with its complement, use midpoint
parity, and then project to the transpose-invariant fixed sector.  No
placement-by-placement vanishing is asserted.
\end{remark}

\begin{remark}[Raw placement-level vanishing is false]
\label{rem:raw-placement-level-vanishing-false}
The aggregate cancellations above must not be read as placement-level
vanishing.  A single fourth-transfer subchannel, for instance a diagonal
same-point placement in the \((1,1,1,4)\) class, can produce a nonzero
fixed-sector \(I\)-component.  The cancellation occurs only after summing
the labelled subset-complement aggregate and applying the
transpose-invariant projection \(\pi_{\mathfrak f}\).  Thus the correct
theorem is aggregate fixed-sector cancellation, not termwise support-zero
or family-by-family gauge.
\end{remark}

\begin{remark}[Layering of the corrected package]
\label{rem:layering-of-corrected-package}
The raw edge-divisibility statement is pre-normalization.  Only after
\[
(1-\Pi_{1,1})_{\mathrm{mix}}C_{\mathcal W}^{\mathrm{raw}}
=
O(a_1\delta^2)
\]
has been established may one form the amplitude-normalized package
\[
-a_1^{-1}\delta^{-2}
(1-\Pi_{1,1})_{\mathrm{mix}}C_{\mathcal W}^{\mathrm{raw}}.
\]
Post-edge quotient compatibility statements involving \(Q_{7,m}^q\)
should not be used as the first proof of edge divisibility.
\end{remark}

\begin{remark}[Post-edge quotient compatibility is not a proof of edge divisibility]
\label{rem:post-edge-quotient-compatibility-not-edge-proof}
The formal projection identity
\[
\Pi_{1,1}\mathcal T_{d\ge3}=0
\]
does not by itself prove Gate~B.  Gate~B uses the retained complement
\[
(1-\Pi_{1,1})_{\mathrm{mix}}C_{\mathcal W}^{\mathrm{raw}},
\]
and must be proved before the normalized quotient package exists.
Therefore a statement such as
\[
Q_{7,m}^q(C)=Q_{7,m}^q(\Pi_{1,1}C)
\]
is a post-edge compatibility statement unless it is reformulated at the
raw edge level.  It cannot be used as the first proof that
\[
(1-\Pi_{1,1})_{\mathrm{mix}}C_{\mathcal W}^{\mathrm{raw}}=O(a_1\delta^2).
\]
\end{remark}

\begin{definition}[Pre-trivialization edge table after Gate B]
\label{def:pre-trivialization-edge-table-after-gate-b}
After Corollary~\ref{cor:gate-b-raw-edge-divisibility-order-seven}, the
edge-normalized package is holomorphic through ordinary order \(7\).  For
a source monomial \(M\), define the pre-trivialization edge coefficient
\[
\boxed{
Z_j[M](z)
:=
-[\epsilon^j\delta^2M]\,
a_1^{-1}\,
\tau_{h\to m}\pi_{\mathfrak f_h}[z^{\le7}]
(1-\Pi_{1,1})_{\mathrm{mix}}C_{\mathcal W}^{\mathrm{raw}},
}
\]
where \(\epsilon=h-m\).  Equivalently,
\[
Z_j[M](z)
=
[\epsilon^jM]\,
\tau_{h\to m}\pi_{\mathfrak f_h}[z^{\le7}]
\widehat C_{\mathcal W}^{\mathrm{amp}},
\]
once the edge-normalized package
\(\widehat C_{\mathcal W}^{\mathrm{amp}}\) has been formed.
\end{definition}

\begin{remark}[Polynomial versus vector notation for {\(Z_j[M]\)}]
\label{rem:polynomial-versus-vector-notation-zj}
Definition~\ref{def:pre-trivialization-edge-table-after-gate-b} first
defines the pre-trivialization edge table as a polynomial
\[
Z_j[M](z)\in\mathfrak f_m[z]_{\le7}.
\]
When \(Z_j[M]\) appears without the argument \(z\) in the moving-line
connection formula, it denotes its ordinary \(z^7\)-coefficient:
\[
Z_j[M]:=[z^7]Z_j[M](z)\in\mathfrak f_m.
\]
Thus
\[
q_{\mathrm{conn},7}[M]
=
\sum_{i=1}^{7}\omega_m(Z_{7-i}[M],A_i)
\]
uses the vector coefficient \([z^7]Z_{7-i}[M](z)\), not the full
polynomial.
\end{remark}

\begin{remark}[Pre-trivialization versus post-trivialization coefficients]
\label{rem:pre-trivialization-versus-post-trivialization-coefficients}
The coefficients \(Z_j[M]\) are pre-trivialization coefficients.  If
\(B_j[M]\) denotes coefficients after applying the representative-level
midpoint trivialization, then \(B_j[M]\) already includes the moving-line
connection contribution.  Therefore the connection sum
\[
q_{\mathrm{conn},7}[M]
=
\sum_{i=1}^{7}\omega_m(Z_{7-i}[M],A_i)
\]
must use \(Z_j[M]\), not post-trivialization \(B_j[M]\), or the
connection is counted twice.
\end{remark}

\begin{remark}[Next object after Gate B]
\label{rem:next-object-after-gate-b}
After Gate~B, the next finite table is not a raw seven-family coefficient
table.  It is the pre-trivialization \(\delta^2\)-edge table
\[
Z_j[M]
=
-[\epsilon^j\delta^2M]\,
a_1^{-1}
\tau_{h\to m}\pi_{\mathfrak f_h}[z^{\le7}]
(1-\Pi_{1,1})_{\mathrm{mix}}C_{\mathcal W}^{\mathrm{raw}}.
\]
This table is defined only after raw edge-divisibility makes the
\(\delta^{-2}\)-normalized package holomorphic.  The finite
Fr\'echet/Sylvester calculus supplies the raw coefficient infrastructure
for computing \(Z_j[M]\); it does not replace the edge-divisibility
theorem.
\end{remark}

\subsection{Midpoint parity}
\label{subsec:midpoint-parity}

\begin{lemma}[Midpoint parity of Fr\'echet coefficients]
\label{lem:midpoint-parity-frechet-coefficients}
In the midpoint chart,
\[
t_\pm=m\pm\frac{z}{2Q^2},
\]
the plus-side and minus-side coefficients satisfy
\[
\Lambda^+_{a,n}=(-1)^n\Lambda^-_{a,n},
\]
and
\[
H^+_{ab,n}=(-1)^nH^-_{ab,n},
\]
for the source indices used in the finite coefficient tables.  The same
parity convention applies to the third Fr\'echet coefficients once the
corresponding \(G^\pm_{abc,n}\) table is defined.
\end{lemma}

\subsection{Fixed-sector projection rule}
\label{subsec:fixed-sector-projection-rule}

\begin{lemma}[Fixed-sector projection of a \(J\)-middle product]
\label{lem:fixed-sector-projection-j-middle-product}
Let
\[
H=aI+bS+cD,\qquad W=pI+qS+rD.
\]
Then
\[
\pi_{\mathfrak f}(HJW)
=
(cq-br)I+(cp-ar)S,
\]
and
\[
\pi_{\mathfrak f}(WJH)
=
-(cq-br)I-(cp-ar)S.
\]
Consequently,
\[
\pi_{\mathfrak f}(WJH)=-\pi_{\mathfrak f}(HJW).
\]
\end{lemma}

\subsection{Computed finite gates}
\label{subsec:computed-finite-gates}

\begin{remark}[Scope of the computed finite gates subsection]
\label{rem:scope-computed-finite-gates-subsection}
The finite gates in this subsection record diagnostic computations from
the ordinary matrix-output / raw-family audit.  They are useful as
coefficient checks and as warnings against termwise or raw-package
shortcuts, but they are not the active Gate~A/Gate~B proof route.

In the active homogeneous scalar-grade proof, the \(Y\)-slot families
are absent at \(B_7\), Gate~A closes by aggregate fixed-sector
cancellation, and Gate~B closes through ordinary order \(7\) by odd
total-grade subset-complement cancellation, moving source-frame
telescoping, and affine boundary removal.  Therefore the older
\(Y^5/Y^6\), \(d^-_{324}\), and five-family \(N^{(0)}\)-weighted
residual formulas below should be read as historical diagnostics for
the raw package, not as remaining theorem obligations.
\end{remark}

\begin{remark}[Source-coefficient parity in the \(Y\)-slot]
\label{rem:y-slot-source-coefficient-parity}
For the grade-five mixed slot \(Y^a\), the source-coefficient table gives
\[
Y^5=0,
\]
while
\[
Y^6=y_6J,\qquad y_6=\frac{1}{53760\pi Q^9}.
\]
Thus the \(Y^5\)-slot residuals vanish individually, and the \(Y^6\)-slot
residuals cancel in left-right pairs by
Lemma~\ref{lem:fixed-sector-projection-j-middle-product}.
\end{remark}

\begin{remark}[First invisible endpoint residual]
\label{rem:first-invisible-endpoint-residual}
For the invisible endpoint pair \((3,4)\) with middle source \(q=2\), the
three-family \(C_2\) gate reduces to one scalar
\[
d^-_{324}
=
y_4\Xi_{324},
\qquad
y_4=-\frac{q^2}{640\pi Q^9},
\]
where
\[
\Xi_{324}
=
\mathcal D(H^-_{2;32}JW_0-H^-_{1;32}JW_1).
\]
The accompanying \(d^+_{324}\) term vanishes by \(Y^3=0\), and the
\(d^+_{234}\) term vanishes by source-degree support
\(H^-_{0..2;34}=0\).
\end{remark}

\begin{remark}[Five-family \(N^{(0)}\)-weighted endpoint gate]
\label{rem:five-family-n0-weighted-endpoint-gate}
In the older raw-family diagnostic audit, the \(N_0\)-middle families did
not collapse to the \(r+t=2\) range.  They involved the full baseline
middle factor and hence the convolution
\[
n+s+t=7.
\]
For the first endpoint residual define
\[
E^-_{324}
=
\mathcal D
\sum_{\substack{1\le n\le5\\2\le t\le6\\n+t\le7}}
(-1)^tH^-_{32,n}N^{(0)}_{7-n-t}\Lambda^-_{4,t},
\]
\[
E^+_{234}
=
\mathcal D
\sum_{\substack{0\le n\le4\\3\le t\le7\\n+t\le7}}
(-1)^t\Lambda^-_{2,n}N^{(0)}_{7-n-t}H^-_{34,t},
\]
\[
E^+_{324}
=
\mathcal D
\sum_{\substack{1\le n\le5\\2\le t\le6\\n+t\le7}}
(-1)^t\Lambda^-_{3,n}N^{(0)}_{7-n-t}H^-_{24,t}.
\]
The five-family residuals are
\[
R^{(A)}_5=2y_4\Xi_{324}+2E^-_{324}+E^+_{234},
\]
\[
R^{(C)}_5=y_4\Xi_{324}+E^-_{324}+2E^+_{324}.
\]
These are finite coefficient identities, not consequences of support or
generic Fr\'echet algebra.

In the active Gate~B proof these raw-family residual formulas are
superseded by Corollary~\ref{cor:gate-b-raw-edge-divisibility-order-seven}.
They should not be listed as live post-Gate-B obstructions unless one
rejects one of the Gate~B hypotheses in
Remark~\ref{rem:what-remains-attackable-in-gate-b}.
\end{remark}

\subsection{Raw-package failure status}
\label{subsec:raw-package-failure-status}

\begin{remark}[Computational status of the raw package shortcut]
\label{rem:computational-status-raw-package}
Source-constrained nonconstant baseline tests show that the raw five-family
and raw seven-family Fr\'echet packages are not generally quotient-gauge
under the provisional fixed-sector line.  The constant baseline is
quotient-inconclusive because \(A_5^{\mathfrak f}=0\).  Nonconstant linear
baselines give nonzero determinant residuals.

The possible Fr\'echet-normalization rescue has also been checked: both the
derivative-normalized weights
\[
(\lambda_2,\lambda_3)=(1,1)
\]
and the homogeneous Taylor weights
\[
(\lambda_2,\lambda_3)=\left(\frac12,\frac16\right)
\]
leave nonzero residuals in the tested nonconstant baselines.

These tests reject the raw UV-026 shortcut.  They do not test UV-026
itself, which is the package quotient-gauge statement for
\[
C_{\mathcal W}^{\corr}.
\]
The corrected package representative is the object defined in
Definition~\ref{def:normalized-corrected-whitening-package-representative};
the finite recurrences and raw gates are a coefficient toolkit, not a
closure of the package theorem.

The word ``inconclusive'' for the constant baseline is meant in the precise
sense of Remark~\ref{rem:good-patch-convention-quotient-coordinate}: when
\(A_5^{\mathfrak f}=0\), the determinant coordinate does not define the
fixed-target quotient.  The nonconstant baseline failures occur on
nondegenerate good patches and therefore reject the raw Fr\'echet shortcut
under the tested convention.  They do not replace the corrected package
theorem, whose object is \(C_{\mathcal W}^{\corr}\).

In particular, these computations belong to the old non-homogeneous
matrix-output \(M^{[5]}\) / fixed-midpoint shortcut lane; they are not
tests of the homogeneous scalar-grade Gate~A/Gate~B mechanism above.
\end{remark}

\section{Corrected fixed-target quotient package}
\label{sec:corrected-fixed-target-quotient-package}

This section defines the corrected package object targeted by the package
frontier of Remark~\ref{rem:wip-first-cubic-source-transfer-target}.  The
sequence of definitions is: raw matrix cross-effect before \(\Phi_K\),
source-bidegree projection, amplitude normalization, representative-level
midpoint trivialization, the normalized package representative, the order-\(7\)
quotient extractor, and the lower moving package coefficients with the
order-\(7\) connection correction.  The closing remark separates definition
material from the theorem obligations that remain.  Throughout this section,
the shorthand \(\mathrm{raw}\) labels the pre-projection cross-effect and
\(\mathrm{amp}\) labels the amplitude-normalized package; these are not
new macros, only superscripts.

The definitions in this section are package definitions.  Raw
edge-divisibility through ordinary order \(7\) is supplied by
Corollary~\ref{cor:gate-b-raw-edge-divisibility-order-seven}, conditional on
the stated source-frame, parity, and boundary hypotheses.  These definitions
do not by themselves prove the later quotient-gauge theorem
\[
B_7^{\mathfrak f}(C_{\mathcal W}^{\corr})
\in
\mathbf C A_5^{\mathfrak f}(m),
\]
nor do they prove the post-Gate-B connection cancellation identities for
\(M_A\) and \(M_C\).  Those obligations are listed in
Remark~\ref{rem:definition-level-material-and-theorem-obligations}.

\begin{remark}[Naming convention for UV-026]
\label{rem:naming-convention-uv026}
In what follows, ``UV-026'' means the corrected package quotient-gauge
target:
\[
B_7^{\mathfrak f}(C_{\mathcal W}^{\corr})
\in
\mathbf C A_5^{\mathfrak f}(m),
\]
where \(C_{\mathcal W}^{\corr}\) is the normalized,
source-bidegree-projected, fixed-target corrected-whitening package
representative.

The older statement that the raw seven-family cubic Fr\'echet inventory is
already \(A_5^{\mathfrak f}(m)\)-gauge is called the \emph{raw UV-026
shortcut}.  It is not the theorem object.

Thus, in the current draft, ``UV-026'' always refers to the corrected
package quotient-gauge theorem.  The live coefficient gates are merely
local tests for that theorem, not separate theorem statements.
\end{remark}

\begin{theorem}[Corrected UV-026 quotient-gauge theorem on the quintic-good patch]
\label{thm:uv026-corrected-package-quotient-gauge}
Work on a good patch on which
\[
A_5^{\mathfrak f}(m)\ne0.
\]
Assume that Gate~B raw edge-divisibility has supplied the holomorphic
edge-normalized object
\[
\widehat C_{\mathcal W}^{\mathrm{amp}}
=
-a_1^{-1}\delta^{-2}
(1-\Pi_{1,1})_{\mathrm{mix}}
C_{\mathcal W}^{\mathrm{raw}}
\]
through ordinary order \(7\) on the collision chart.  Let
\(C_{\mathcal W}^{\corr}\) denote the fixed-target fixed-sector package
representative obtained from this object by
\[
C_{\mathcal W}^{\corr}
=
B_{\mathrm{pkg}}^{\mathfrak f}(h;z)
=
\sum_{n=0}^{7}
z^n\operatorname{Triv}_{h\to m} Z_n(h),
\]
where
\[
\pi_{\mathfrak f_h}[z^{\le7}]
\widehat C_{\mathcal W}^{\mathrm{amp}}(h)
=
\sum_{n=0}^{7}z^n Z_n(h).
\]
Then the corrected UV-026 statement is
\[
B_7^{\mathfrak f}(C_{\mathcal W}^{\corr})
:=
[z^7]B_{\mathrm{pkg}}^{\mathfrak f}(m;z)
\in
\mathbf C A_5^{\mathfrak f}(m),
\]
equivalently
\[
Q_{7,m}^q(C_{\mathcal W}^{\corr})
=
0
\quad\text{in}\quad
\mathfrak f_m/\mathbf C A_5^{\mathfrak f}(m).
\]
\end{theorem}

\begin{remark}[UV-026 theorem object versus proof data]
\label{rem:uv026-object-versus-proof-data}
The theorem object in UV-026 is the quotient class
\[
Q_{7,m}^q(C_{\mathcal W}^{\corr})
\in
\mathfrak f_m/\mathbf C A_5^{\mathfrak f}(m).
\]
The source-bidegree projection, amplitude normalization, fixed-sector
extraction, and midpoint trivialization are part of the definition of
\(C_{\mathcal W}^{\corr}\).  They should not be applied again to
\(C_{\mathcal W}^{\corr}\).

The pre-trivialization edge table \(Z_j[M]\) is not extra output data in the
UV-026 statement.  It is proof data needed to verify the connection
cancellation for
\[
M_A=\dot r_3r_2\dot r_4,
\qquad
M_C=\dot r_3\dot r_2r_4.
\]
The required proof identities are
\[
R_7^{(A)}
+
\sum_{i=1}^{7}\omega_m(Z_{7-i}[M_A],A_i)=0,
\]
and
\[
R_7^{(C)}
+
\sum_{i=1}^{7}\omega_m(Z_{7-i}[M_C],A_i)=0.
\]
Replacing \(Z_j[M]\) by post-trivialization \(B_j[M]\) inside this connection
sum double-counts the moving-line correction.
\end{remark}

\begin{definition}[Raw corrected-whitening matrix cross-effect before \(\Phi_K\)]
\label{def:raw-corrected-whitening-cross-effect}
Let
\[
\mathcal B_0=(G_-^{(0)},N^{(0)},G_+^{(0)})
\]
be the baseline corrected block, and let
\[
\mathcal B_1(a_i,h_i)
\]
denote the one-source tagged corrected block obtained by substituting only the
\(i\)-th tagged source perturbation into the corrected same-point and mixed
block formulae.  Let
\[
\mathcal B_2(a_1,h_1;a_2,h_2)
\]
be the tagged two-source corrected block.

Let
\[
\mathcal W(G_-,N,G_+):=G_-^{-1/2}NG_+^{-1/2}
\]
be the corrected-whitening transfer, defined on the good analytic patch where
the inverse square roots are holomorphic.

Define the matrix cross-effect before applying \(\Phi_K\), determinant
scalarization, or quotient extraction by
\[
\begin{aligned}
C_{\mathcal W}^{\mathrm{raw}}(a_1,h_1;a_2,h_2)
:={}&
\mathcal W(\mathcal B_2(a_1,h_1;a_2,h_2))
-\mathcal W(\mathcal B_1(a_1,h_1))
\\
&-\mathcal W(\mathcal B_1(a_2,h_2))
+\mathcal W(\mathcal B_0).
\end{aligned}
\]
This cross-effect is the correct interaction layer for one-amplitude
collapse, but it does not by itself imply diagonal or collision vanishing
of the order-\(7\) quotient component; that remains the merger,
renormalization, or same-reduced-image theorem.
\end{definition}

\begin{definition}[Source-bidegree projection before quotient extraction]
\label{def:source-bidegree-projection-before-quotient}
Work with formal source tags \(\tau_1,\tau_2\) retained through the matrix
cross-effect.  For a tagged matrix series
\[
F(\tau_1,\tau_2)=\sum_{p,q\ge0}\tau_1^p\tau_2^qF_{p,q},
\]
define
\[
\Pi_{1,1}F:=F_{1,1}.
\]
The non-\((1,1)\) mixed source component is
\[
(1-\Pi_{1,1})_{\mathrm{mix}}F
:=
\sum_{\substack{p,q\ge1\\(p,q)\ne(1,1)}}F_{p,q}.
\]
This projection is applied before fixed-sector extraction, before
\(B_7^{\mathfrak f}\), before determinant pairing, and before quotienting.
\end{definition}

\begin{definition}[Amplitude-normalized non-\((1,1)\) package]
\label{def:amplitude-normalized-non-11-package}
On the collision chart
\[
h_1=h-\frac{\delta}{2},
\qquad
h_2=h+\frac{\delta}{2},
\qquad
a_2=-a_1(1+\kappa\delta),
\]
define the formally amplitude-normalized package by
\[
\widehat C_{\mathcal W}^{\mathrm{amp}}
:=
-a_1^{-1}\delta^{-2}
(1-\Pi_{1,1})_{\mathrm{mix}}
C_{\mathcal W}^{\mathrm{raw}}
\!\left(a_1,h-\frac{\delta}{2};
        -a_1(1+\kappa\delta),h+\frac{\delta}{2}\right).
\]
This is a matrix-valued object before \(\Phi_K\), before determinant
scalarization, and before quotient extraction.  Its use below is conditional
on holomorphic extension to \(\delta=0\).
\end{definition}

\begin{definition}[Representative-level midpoint trivialization]
\label{def:representative-level-midpoint-trivialization}
Let
\[
A(h):=A_5^{\mathfrak f}(h)\in\mathfrak f_h
\]
be the moving quintic fixed-sector line.  Fix a good patch on which
\[
A(m)\ne0.
\]
Fix an analytic identification
\[
\tau_{h\to m}:\mathfrak f_h\longrightarrow\mathfrak f_m.
\]
Write
\[
A_h^m:=\tau_{h\to m}A(h),
\qquad
A_m:=A(m).
\]
Let
\[
\omega_m:\mathfrak f_m\times\mathfrak f_m\to\mathbf C
\]
be the fixed determinant form in the \(I,S\)-basis.

Choose a local transverse section \(E_m\in\mathfrak f_m\) satisfying
\[
\omega_m(E_m,A_m)=1,
\]
and an analytic linear functional
\[
\ell_h:\mathfrak f_m\to\mathbf C
\]
satisfying
\[
\ell_h(A_h^m)=1.
\]
For \(Z(h)\in\mathfrak f_h\), set
\[
\widetilde Z(h):=\tau_{h\to m}Z(h),
\]
and define
\[
\operatorname{Triv}_{h\to m}Z(h)
:=
\ell_h(\widetilde Z(h))A_m
+
\omega_m(\widetilde Z(h),A_h^m)E_m.
\]
Then
\[
\omega_m\!\left(\operatorname{Triv}_{h\to m}Z(h),A_m\right)
=
\omega_m(\widetilde Z(h),A_h^m).
\]
Changing \(E_m\) or \(\ell_h\) changes the representative only by a multiple of
\(A_m\), and hence does not change its class in
\[
\mathfrak f_m/\mathbf C A_m.
\]
\end{definition}

\begin{definition}[Normalized corrected-whitening fixed-sector package representative]
\label{def:normalized-corrected-whitening-package-representative}
Assume the amplitude-normalized package
\[
\widehat C_{\mathcal W}^{\mathrm{amp}}
\]
extends holomorphically to \(\delta=0\) through ordinary order \(7\).

First extract the moving fixed-sector \(z\)-polynomial
\[
Z_{\mathcal W}^{\mathfrak f}(h;z)
:=
\pi_{\mathfrak f_h}[z^{\le7}]
\widehat C_{\mathcal W}^{\mathrm{amp}}(h)
\in
\mathfrak f_h[z]_{\le7}.
\]
Write
\[
Z_{\mathcal W}^{\mathfrak f}(h;z)
=
\sum_{n=0}^{7}z^nZ_n(h),
\qquad
Z_n(h)\in\mathfrak f_h .
\]
The midpoint trivialization is applied coefficientwise:
\[
B_{\mathrm{pkg}}^{\mathfrak f}(h;z)
:=
\sum_{n=0}^{7}
z^n\operatorname{Triv}_{h\to m}Z_n(h)
\in
\mathfrak f_m[z]_{\le7}.
\]
Thus ordinary \(z\)-coefficients are extracted only after fixed-target
trivialization.

Equivalently, expand the fixed-target \(z\)-polynomial in the moving
midpoint parameter:
\[
B_{\mathrm{pkg}}^{\mathfrak f}(h;z)
=
\sum_{j=0}^{7}(h-m)^j\mathcal B_j(z)
+
O((h-m)^8),
\qquad
\mathcal B_j(z)\in\mathfrak f_m[z]_{\le7}.
\]
Although the coefficients still depend on the moving parameter \(h\), the
representatives themselves have already been transported to the fixed
midpoint fiber \(\mathfrak f_m\).

When we write \(C_{\mathcal W}^{\corr}\) below, we mean this normalized
fixed-target fixed-sector package representative
\(B_{\mathrm{pkg}}^{\mathfrak f}(h;z)\), not the full unprojected matrix
cross-effect.
\end{definition}

\begin{definition}[Package-level order-\(7\) quotient extractor]
\label{def:package-level-order-7-quotient-extractor}
For the normalized fixed-target package representative, define
\[
B_7^{\mathfrak f}(C_{\mathcal W}^{\corr})
:=
[z^7]\,B_{\mathrm{pkg}}^{\mathfrak f}(m;z)
\in \mathfrak f_m .
\]
The quotient-valued order-\(7\) extractor is
\[
Q_{7,m}^q(C_{\mathcal W}^{\corr})
:=
\left[
B_7^{\mathfrak f}(C_{\mathcal W}^{\corr})
\right]
\in
\mathfrak f_m/\mathbf C A_5^{\mathfrak f}(m).
\]
\end{definition}

\begin{remark}[Good-patch convention for the quotient coordinate]
\label{rem:good-patch-convention-quotient-coordinate}
The quotient
\[
Q_{7,m}^q(C_{\mathcal W}^{\corr})
\in
\mathfrak f_m/\mathbf C A_5^{\mathfrak f}(m)
\]
is the invariant object.  On a good patch where
\[
A_5^{\mathfrak f}(m)=u_5I+v_5S\ne0,
\]
the determinant coordinate
\[
\mathcal D_m(Z):=(v_5\pi_I-u_5\pi_S)\pi_{\mathfrak f_m}(Z)
\]
identifies this one-dimensional quotient with \(\mathbf C\).  This
coordinate is not available on the rank-jump locus
\[
A_5^{\mathfrak f}(m)=0.
\]
Consequently, baseline tests for which \(A_5^{\mathfrak f}(m)=0\) are
quotient-inconclusive, not evidence for or against the quotient-gauge
theorem.
\end{remark}

\begin{remark}[Rank-jump locus is not covered by the good-patch quotient theorem]
\label{rem:rank-jump-locus-not-covered-uv026}
The corrected UV-026 theorem is a good-patch statement.  The
representative-level midpoint trivialization uses a nonzero quintic line
\[
A_5^{\mathfrak f}(m)\ne0
\]
to identify the quotient
\[
\mathfrak f_m/\mathbf C A_5^{\mathfrak f}(m).
\]
On the rank-jump locus
\[
A_5^{\mathfrak f}(m)=0,
\]
the determinant coordinate is unavailable.  Moreover, the literal inclusion
\[
B_7^{\mathfrak f}(C_{\mathcal W}^{\corr})
\in
\mathbf C A_5^{\mathfrak f}(m)
\]
would reduce to the stronger assertion
\[
B_7^{\mathfrak f}(C_{\mathcal W}^{\corr})=0.
\]
That zero statement is not supplied by the good-patch quotient theorem.
Rank-jump cases therefore require a separate extension, blow-up, or limiting
quotient-line theorem.
\end{remark}

\begin{remark}[Two coefficient indices and two lower tables]
\label{rem:two-coefficient-indices-package}
There are two independent indexing conventions in the corrected package.

First, the subscript \(7\) in
\[
B_7^{\mathfrak f}
\]
denotes the ordinary \(z\)-order coefficient extracted from the fixed-target
package polynomial.

Second, the subscript \(j\) in a lower moving table denotes the coefficient
of \(\epsilon^j=(h-m)^j\) after the ordinary \(z^7\) coefficient has been
selected.

There are also two lower tables.  The pre-trivialization table is
\[
Z_j[M]
=
-[z^7][\epsilon^j\delta^2M]\,
a_1^{-1}\tau_{h\to m}
\pi_{\mathfrak f_h}
(1-\Pi_{1,1})_{\mathrm{mix}}C_{\mathcal W}^{\mathrm{raw}},
\]
and it is the table that enters the moving-line connection formula:
\[
q_{\mathrm{conn},7}[M]
=
\sum_{i=1}^{7}\omega_m(Z_{7-i}[M],A_i).
\]
The post-trivialization table is
\[
B_j[M]=[z^7][\epsilon^jM]\,B_{\mathrm{pkg}}^{\mathfrak f}(h;z),
\]
after applying the representative-level midpoint trivialization.  It records
the chosen final representative, but it must not be inserted into the
pre-trivialization connection formula.
\end{remark}

\begin{definition}[Lower moving package coefficients and connection data]
\label{def:lower-moving-package-coefficients}
Let \(M\) be a residual-source monomial and set \(\epsilon=h-m\).  After
Gate~B raw edge-divisibility, define the pre-trivialization edge
coefficients by
\[
Z_j[M]
:=
-[z^7][\epsilon^j\delta^2M]\,
a_1^{-1}\tau_{h\to m}
\pi_{\mathfrak f_h}
(1-\Pi_{1,1})_{\mathrm{mix}}C_{\mathcal W}^{\mathrm{raw}},
\qquad
0\le j\le 6.
\]
Equivalently,
\[
Z_j[M]
=
[z^7][\epsilon^jM]\,
\tau_{h\to m}\pi_{\mathfrak f_h}
\widehat C_{\mathcal W}^{\mathrm{amp}}
\]
once the edge-normalized package has been formed.

By contrast, let
\[
B_{\mathrm{pkg}}^{\mathfrak f}(h;z;M)
\]
be the coefficient of \(M\) in the fixed-target package representative
after representative-level midpoint trivialization, and write
\[
B_{\mathrm{pkg}}^{\mathfrak f}(h;z;M)
=
\sum_{j\ge0}\epsilon^j\mathcal B_j[M](z),
\qquad
\mathcal B_j[M](z)\in\mathfrak f_m[z]_{\le7}.
\]
The post-trivialization lower representatives are
\[
B_j[M]:=[z^7]\mathcal B_j[M](z)\in\mathfrak f_m.
\]
These coefficients are useful for recording the final representative, but
they must not be inserted into the moving-line connection formula, because
the trivialization has already transported along the moving quotient
line.

For the connection formula, all moving-line coefficients are first
transported to the fixed target \(\mathfrak f_m\).  Thus set
\[
A_h^m:=\tau_{h\to m}A_5^{\mathfrak f}(h),
\qquad
A_h^m=\sum_{i\ge0}\epsilon^iA_i,
\qquad
A_i=u_i^AI+v_i^AS\in\mathfrak f_m.
\]
Equivalently, every occurrence of
\[
A_i=[\epsilon^i]A_5^{\mathfrak f}(h)
\]
inside \(\omega_m(\,\cdot\,,A_i)\) means
\[
A_i=[\epsilon^i]\tau_{h\to m}A_5^{\mathfrak f}(h).
\]

If
\[
Z_j[M]=u_j^Z[M]I+v_j^Z[M]S,
\]
then the order-\(7\) quotient-line connection correction on \(M\) is
\[
q_{\mathrm{conn},7}[M]
=
\sum_{i=1}^{7}
\omega_m(Z_{7-i}[M],A_i).
\]
In coordinates,
\[
q_{\mathrm{conn},7}[M]
=
\sum_{i=1}^{7}
\left(
u_{7-i}^Z[M]v_i^A
-
u_i^Av_{7-i}^Z[M]
\right).
\]
If \(A_i,Z_j[M]\) are derivative-normalized rather than
coefficient-normalized, replace this formula by
\[
q_{\mathrm{conn},7}[M]
=
\sum_{i=1}^{7}\binom{7}{i}
\omega_m(Z_{7-i}[M],A_i).
\]
\end{definition}

\begin{remark}[Representative choices and lower coefficient tables]
\label{rem:representative-choices-lower-coefficient-tables}
There are two different kinds of lower coefficients.

The pre-trivialization coefficients \(Z_j[M]\) are extracted before applying
\(\operatorname{Triv}_{h\to m}\).  These are the coefficients that enter the
moving-line connection formula
\[
q_{\mathrm{conn},7}[M]
=
\sum_{i=1}^{7}\omega_m(Z_{7-i}[M],A_i).
\]

The post-trivialization coefficients \(B_j[M]\) record the chosen
fixed-target representative after representative-level midpoint
trivialization.  They are useful for describing the final normalized
representative, but they already contain the moving-line correction.

Consequently, one may not arbitrarily replace
\[
Z_j[M]\mapsto Z_j[M]+c_jA_0
\]
inside the connection sum unless the direct term and the moving-line
expansion are transformed coherently as part of the same analytic
representative choice.  Likewise, post-trivialization shifts
\[
B_j[M]\mapsto B_j[M]+c_jA_0
\]
are harmless only when \(B_j[M]\) is used solely as a quotient
representative; they are not harmless if the shifted \(B_j[M]\) is then
inserted into a separate connection formula.

Thus coefficient tables may choose an analytic section, but they must keep
the section fixed through the final quotient-coordinate calculation.
Changing representatives in the middle of a finite table is not a valid
simplification.
\end{remark}

\begin{remark}[Pre- and post-trivialization lower coefficients]
\label{rem:pre-post-trivialization-lower-coefficients}
Definition~\ref{def:lower-moving-package-coefficients} uses
pre-trivialization edge coefficients \(Z_j[M]\) in the connection formula
\[
q_{\mathrm{conn},7}[M]
=
\sum_{i=1}^{7}\omega_m(Z_{7-i}[M],A_i).
\]
Post-trivialization coefficients \(B_j[M]\) already include the
moving-line connection contribution.  Inserting \(B_j[M]\) into the
displayed pre-trivialization connection formula would double-count the
moving-line correction.
\end{remark}

\begin{remark}[Next package table required for the connection correction]
\label{rem:next-package-table-connection-correction}
The next computable package table is not a raw Fr\'echet-family table and
not a post-trivialization representative table.  The connection correction
requires the transported moving-line table
\[
A_i=[\epsilon^i]\tau_{h\to m}A_5^{\mathfrak f}(h),
\qquad
1\le i\le7,
\qquad
\epsilon=h-m,
\]
together with the pre-trivialization edge coefficients
\[
Z_j[M_A]
=
-[z^7][\epsilon^j\delta^2M_A]\,
a_1^{-1}\tau_{h\to m}
\pi_{\mathfrak f_h}
(1-\Pi_{1,1})_{\mathrm{mix}}C_{\mathcal W}^{\mathrm{raw}},
\]
and
\[
Z_j[M_C]
=
-[z^7][\epsilon^j\delta^2M_C]\,
a_1^{-1}\tau_{h\to m}
\pi_{\mathfrak f_h}
(1-\Pi_{1,1})_{\mathrm{mix}}C_{\mathcal W}^{\mathrm{raw}},
\qquad
0\le j\le6,
\]
where
\[
M_A=\dot r_3r_2\dot r_4,
\qquad
M_C=\dot r_3\dot r_2r_4.
\]
The connection-cancellation identities to prove are
\[
R_7^{(A)}
+
\sum_{i=1}^{7}
\omega_m(Z_{7-i}[M_A],A_i)=0,
\]
and
\[
R_7^{(C)}
+
\sum_{i=1}^{7}
\omega_m(Z_{7-i}[M_C],A_i)=0.
\]
The post-trivialization coefficients \(B_j[M]\) may be recorded after
applying \(\operatorname{Triv}_{h\to m}\), but inserting them into the
displayed connection formula double-counts the moving quotient-line
correction.
\end{remark}

\begin{remark}[Status of the \(M_A,M_C\) connection identities]
\label{rem:status-ma-mc-connection-identities}
The identities
\[
R_7^{(A)}
+
\sum_{i=1}^{7}\omega_m(Z_{7-i}[M_A],A_i)=0,
\]
and
\[
R_7^{(C)}
+
\sum_{i=1}^{7}\omega_m(Z_{7-i}[M_C],A_i)=0
\]
are not proved in the present draft by parity alone.  The available
antisymmetry and subset-complement mechanisms are Gate~B edge mechanisms;
they do not by themselves evaluate the post-Gate-B moving-line connection
sum.

The viable closure mechanism is cancellation between the direct corrected
order-\(7\) residuals \(R_7^{(A)},R_7^{(C)}\) and the connection
contributions computed from the lower pre-trivialization table
\(Z_j[M_A],Z_j[M_C]\).  Therefore the current status is:
\[
\boxed{
\text{connection closure depends on the missing finite }Z_j\text{-tables
and corrected }R_7\text{-tables.}
}
\]
In particular, \(R_7^{(A)}\) and \(R_7^{(C)}\) must be corrected-package
residuals in the same pre-\(\Phi_K\), pre-determinant convention.  Old raw
seven-family residuals must not be substituted for these quantities.
\end{remark}

\begin{lemma}[Source-grade reflection]
\label{lem:source-grade-reflection}
Let \(\mathsf S\) denote endpoint swap \(h_-\leftrightarrow h_+\), equivalently
\(\delta\mapsto-\delta\).  Under the homogeneous scalar-grade convention
\[
\operatorname{Gr}_g r(t)
=
\frac{r^{(g+2)}(m)}{(g+2)!}(t-m)^{g+2},
\]
the source-linear grade extractor \(\mathsf P_g\) satisfies
\[
\mathsf P_g\circ\mathsf S
=
(-1)^g\mathsf S\circ\mathsf P_g.
\]
\end{lemma}

\begin{proof}
Reflection about the midpoint sends \(t-m\mapsto-(t-m)\).  Hence
\[
(t-m)^{g+2}\mapsto(-1)^{g+2}(t-m)^{g+2}=(-1)^g(t-m)^{g+2}.
\]
This is exactly the asserted commutation law.
\end{proof}

\begin{lemma}[Corrected mixed-block endpoint-swap identity]
\label{lem:corr-mixed-block-endpoint-swap}
The corrected mixed block satisfies
\[
N^{\corr}(h_-,h_+)^T=N^{\corr}(h_+,h_-).
\]
\end{lemma}

\begin{proof}
Writing the corrected mixed block as
\[
N^{\corr}(q_-,q_+,\Delta;s)
=
\frac1\pi
\begin{pmatrix}
-\dfrac{2\sin\Delta}{s}
&
\dfrac{\sin\Delta+q_+s\cos\Delta}{s^2}
\\[2ex]
-\dfrac{\sin\Delta+q_-s\cos\Delta}{s^2}
&
\dfrac{(q_-+q_+)s\cos\Delta+(2-q_-q_+s^2)\sin\Delta}{2s^3}
\end{pmatrix},
\]
endpoint swap sends
\[
(q_-,q_+,\Delta,s)\mapsto(q_+,q_-,-\Delta,-s).
\]
Using \(\sin(-\Delta)=-\sin\Delta\), \(\cos(-\Delta)=\cos\Delta\),
and \((-s)^2=s^2\), \((-s)^3=-s^3\), the swapped matrix is exactly the
transpose of the original matrix.
\end{proof}

\begin{lemma}[Edge-refined \(Y\)-middle parity]
\label{lem:edge-refined-y-middle-parity-final}
Let \(Y_g\) be the source-linear grade-\(g\) middle block extracted from
\(N^{\corr}\).  Then
\[
Y_g(h_-,h_+)^T=(-1)^gY_g(h_+,h_-).
\]
Consequently, if
\[
Y_g(h-\delta/2,h+\delta/2)
=
\sum_{e\ge0}\delta^eY_g^{\langle e\rangle}(h),
\]
then
\[
\left(Y_{g,\beta}^{\langle e\rangle}\right)^T
=
(-1)^{g+e}Y_{g,\beta}^{\langle e\rangle}.
\]
\end{lemma}

\begin{proof}
Apply Lemma~\ref{lem:corr-mixed-block-endpoint-swap} and then extract the
source-linear homogeneous grade using Lemma~\ref{lem:source-grade-reflection}.
The coefficient of \(\delta^e\) gains the additional factor \((-1)^e\)
because endpoint swap sends \(\delta\mapsto-\delta\).
\end{proof}

\begin{theorem}[\(M_A,M_C\) pre-trivialization connection closure]
\label{thm:ma-mc-pretriv-connection-closure-final}
Let
\[
M_A=\dot r_3r_2\dot r_4,
\qquad
M_C=\dot r_3\dot r_2r_4.
\]
Assume all-subset midpoint parity, transpose-invariance of
\(\pi_{\mathfrak f_h}\), symmetry of \(N^{(0)}\), and the source-grade
reflection convention of Lemma~\ref{lem:source-grade-reflection}.  Then the
pre-trivialization table satisfies
\[
Z_j[M_A]=Z_j[M_C]=0,
\qquad 0\le j\le7.
\]
Consequently,
\[
R_7^{(A)}
+
\sum_{i=1}^{7}\omega_m(Z_{7-i}[M_A],A_i)=0,
\]
and
\[
R_7^{(C)}
+
\sum_{i=1}^{7}\omega_m(Z_{7-i}[M_C],A_i)=0.
\]
\end{theorem}

\begin{proof}
For the \(N_0\)-middle block, the labelled source multiset is
\(J=(3,2,4)\), hence \(|J|=9\).  Complementing a placement
\(I\subseteq J\) gives a factor \((-1)^{|J|}=-1\), while the
\(\delta^2\)-edge weights of \(M_A\) and \(M_C\) are complement-invariant
because both monomials have two dotted endpoint sources.  Hence every
\(N_0\)-middle complement pair cancels.

For the \(Y\)-middle block, fix a middle grade \(m\in J\), set
\(K=J\setminus\{m\}\), and let \(e\in\{0,1\}\) be the number of edge
derivatives carried by the middle source.  The complement placement gains
the matrix factor
\[
(-1)^{|K|}(-1)^{m+e}.
\]
The source-edge weight gains the factor \((-1)^e\).  Since
\[
|K|+m=|J|=9,
\]
the total complement factor is
\[
(-1)^e(-1)^{|K|}(-1)^{m+e}
=
(-1)^{|J|}
=
-1.
\]
Thus every \(Y\)-middle complement pair cancels.  Therefore
\[
Z_j[M_A]=Z_j[M_C]=0
\]
for \(0\le j\le7\).  The displayed connection identities follow immediately
from
\[
R_7^{(X)}=\omega_m(Z_7[M_X],A_0),
\qquad X\in\{A,C\}.
\]
\end{proof}

\begin{remark}[Scope of the \(M_A,M_C\) connection closure]
\label{rem:scope-ma-mc-connection-closure-final}
Theorem~\ref{thm:ma-mc-pretriv-connection-closure-final} closes the
\(M_A,M_C\) pre-trivialization connection branch under the stated
source-grade reflection convention.  It does not prove the full
actual-source Gate~B upgrade
\[
\Theta_7^{\actual}=0,
\]
and it does not address the promoted nonlinear mixed-block Bottleneck-C
coefficient
\[
E_{A,Y}^{\mathrm{BC}}
=
\tau_{h\to m}\pi_{\mathfrak f_h}
\operatorname{Sym}
[z^7]\,
D(G_-^{-1/2})[A_-]\,
Y_{12}^{(2)}\,
R_0,
\qquad
Y_{12}^{(2)}=D^2N_0[\xi_1,\xi_2],
\]
which remains a separate actual-source/Bottleneck-C obstruction.
\end{remark}

\begin{remark}[Connection table reduction for the residual monomials]
\label{rem:connection-table-reduction-residual-monomials}
The post-Gate-B connection identities for
\[
M_A=\dot r_3r_2\dot r_4,
\qquad
M_C=\dot r_3\dot r_2r_4
\]
are finite pre-trivialization table identities.  They are not raw
seven-family identities and they must not be tested using the
post-trivialization coefficients \(B_j[M]\).

Write
\[
\epsilon=h-m,
\qquad
A_h^m:=\tau_{h\to m}A_5^{\mathfrak f}(h)
       =\sum_{i\ge0}\epsilon^iA_i,
\]
so that
\[
A_i=[\epsilon^i]\tau_{h\to m}A_5^{\mathfrak f}(h),
\qquad 1\le i\le7.
\]
If
\[
A_i=\alpha_i I+\beta_iS,
\]
and
\[
Z_j[M_X]=u_j^X I+v_j^X S,
\qquad X\in\{A,C\},
\]
then the coefficient-normalized connection contribution is
\[
q_{\mathrm{conn},7}[M_X]
=
\sum_{i=1}^{7}
\omega_m(Z_{7-i}[M_X],A_i)
=
\sum_{i=1}^{7}
\left(
u_{7-i}^X\beta_i-\alpha_i v_{7-i}^X
\right).
\]
Thus the two required scalar identities are
\[
R_7^{(A)}
+
\sum_{i=1}^{7}
\left(
u_{7-i}^A\beta_i-\alpha_i v_{7-i}^A
\right)
=0,
\]
and
\[
R_7^{(C)}
+
\sum_{i=1}^{7}
\left(
u_{7-i}^C\beta_i-\alpha_i v_{7-i}^C
\right)
=0.
\]

If instead the table uses derivative-normalized coefficients
\[
\mathsf A_i
=
\partial_\epsilon^i
\left(\tau_{m+\epsilon\to m}A_5^{\mathfrak f}(m+\epsilon)\right)
\big|_{\epsilon=0},
\qquad
\mathsf Z_j[M]
=
\partial_\epsilon^j Z[M](m+\epsilon)\big|_{\epsilon=0},
\]
then
\[
A_i=\frac{\mathsf A_i}{i!},
\qquad
Z_j[M]=\frac{\mathsf Z_j[M]}{j!}.
\]
Consequently the coefficient-normalized order-\(7\) term is
\[
q_{\mathrm{conn},7}[M_X]
=
\frac{1}{7!}
\sum_{i=1}^{7}
\binom{7}{i}
\omega_m(\mathsf Z_{7-i}[M_X],\mathsf A_i).
\]
Equivalently, after applying \(\partial_\epsilon^7|_{\epsilon=0}\) to the
whole order-\(7\) identity, the derivative-normalized connection term is
\[
\sum_{i=1}^{7}
\binom{7}{i}
\omega_m(\mathsf Z_{7-i}[M_X],\mathsf A_i).
\]
The coefficient-normalized and derivative-normalized conventions must not be
mixed.
\end{remark}

\begin{definition}[Finite \(Z_j\)-table target for \(M_A\) and \(M_C\)]
\label{def:finite-zj-table-target-ma-mc}
For \(X\in\{A,C\}\), set
\[
M_A=\dot r_3r_2\dot r_4,
\qquad
M_C=\dot r_3\dot r_2r_4.
\]
The required pre-trivialization table entries are
\[
Z_j[M_X]
=
-[z^7][\epsilon^j\delta^2M_X]\,
a_1^{-1}\tau_{h\to m}
\pi_{\mathfrak f_h}
(1-\Pi_{1,1})_{\mathrm{mix}}C_{\mathcal W}^{\mathrm{raw}},
\qquad
0\le j\le6.
\]
Equivalently, after Gate~B edge divisibility has supplied the holomorphic
edge-normalized package,
\[
Z_j[M_X]
=
[z^7][\epsilon^jM_X]\,
\tau_{h\to m}\pi_{\mathfrak f_h}
\widehat C_{\mathcal W}^{\mathrm{amp}}.
\]

To compute these entries from the finite Fr\'echet/Sylvester calculus, one
must keep the labelled source slots through the computation.  For the
labelled source triple \(J=(3,2,4)\), the \(N_0\)-middle contribution has
the schematic contract
\[
\sum_{I\subseteq J}
\sum_{\alpha+\beta+\gamma=7}
\pi_{\mathfrak f_h}
\left(
L^-_{I,\alpha}N^{(0)}_{\beta}
R^+_{J\setminus I,\gamma}
\right),
\]
and the \(Y\)-middle contribution has the schematic contract
\[
\sum_{m_0\in J}
\sum_{I\subseteq J\setminus\{m_0\}}
\sum_{\alpha+\beta+\gamma=7}
\pi_{\mathfrak f_h}
\left(
L^-_{I,\alpha}Y_{m_0,\beta}
R^+_{J\setminus(I\cup\{m_0\}),\gamma}
\right).
\]
The operator
\[
-[\epsilon^j\delta^2M_X]\,a_1^{-1}\tau_{h\to m}
(1-\Pi_{1,1})_{\mathrm{mix}}
\]
is then applied to the labelled aggregate before determinant scalarization
and before midpoint trivialization.

The current finite Fr\'echet recurrences supply the formal \(L\)- and
\(R\)-coefficient infrastructure and the baseline \(N^{(0)}\)-coefficient
infrastructure.  Before the edge-refined \(Y\)-middle parity theorem was
isolated, this table was the direct finite computation needed for the
\(M_A,M_C\) branch.

After Lemma~\ref{lem:edge-refined-y-middle-parity-final} and
Theorem~\ref{thm:ma-mc-pretriv-connection-closure-final}, this table should
be read as an audit fallback rather than as an open primary dependency.  Under
the stated actual source-grade reflection convention, the labelled
\(J=(3,2,4)\) \(N_0\)-middle and \(Y\)-middle complement pairs cancel
before midpoint trivialization:
\[
Z_j[M_A]=Z_j[M_C]=0,\qquad 0\le j\le7.
\]
Thus explicit evaluation of the finite \(Z_j\)-table is no longer required
for the \(M_A,M_C\) connection branch unless the actual source-grade
reflection convention or edge-refined \(Y\)-middle parity is later reopened.
\end{definition}

\begin{remark}[Definition-level material and theorem obligations]
\label{rem:definition-level-material-and-theorem-obligations}
Definitions~\ref{def:raw-corrected-whitening-cross-effect}--\ref{def:lower-moving-package-coefficients}
specify the intended corrected package object.  They do not prove its
quotient-gauge property.

The following are theorem obligations.

\begin{enumerate}
\item \textbf{Source-weight-linear package input.}
One must prove that the tagged pre-whitening source block induces the actual
source-weight-linear input used by the corrected-whitening cross-effect.  In
particular, the construction must not reintroduce the rejected signed-lift
shortcut in which multiplying an even whitened package by the source
amplitude leaves higher \(a^3,a^5,\ldots\) terms.  The package input must be
invariant, functorial through the corrected-whitening transfer, and
compatible with the quotient extractor \(Q_{7,m}^q\).
For the order-\(7\) Gate~B theorem, this requirement is realized concretely
by the moving homogeneous source-frame hypotheses used in
Corollary~\ref{cor:gate-b-raw-edge-divisibility-order-seven}; future
attacks should therefore target those hypotheses directly rather than
reintroducing the discarded signed-lift shortcut.

\item \textbf{Homogeneous scalar-grade source theorem.}
One must prove that the residual scalar source convention
\[
\operatorname{Gr}_a r
=
\frac{r^{(a+2)}(m)}{(a+2)!}(t-m)^{a+2}
\]
is the correct source-grade decomposition for the actual corrected package.
In particular, the absence of the \(Y\)-slot families at
\(B_7^{\mathfrak f}\) relies on this homogeneous scalar-grade convention.
If a different convention defines \(M^{[5]}\) as an ordinary matrix-output
order-\(5\) projection, then the \(Y\)-slot families
\[
L_1YR_1,\qquad L_2YR_0,\qquad L_0YR_2
\]
must be reopened and audited separately, as noted in
Remark~\ref{rem:quarantined-matrix-output-M5-convention}.
The homogeneous scalar-grade convention is also the convention under which
moving source-frame telescoping is stated; changing the source convention
requires re-proving the telescoping and affine-boundary lemmas.

\item \textbf{Edge divisibility.}
Corollary~\ref{cor:gate-b-raw-edge-divisibility-order-seven} proves
\[
[z^7][\delta^1]\pi_{\mathfrak f_h}
(1-\Pi_{1,1})_{\mathrm{mix}}C_{\mathcal W}^{\mathrm{raw}}=0
\]
under its stated source-frame, parity, and full-tower-before-truncation
hypotheses.  Equivalently,
\[
(1-\Pi_{1,1})_{\mathrm{mix}}C_{\mathcal W}^{\mathrm{raw}}
=
O(a_1\delta^2)
\]
through ordinary order \(7\) in fixed sector.

The remaining obligation is not to re-prove the leading cubic or
higher-transfer audits, but to ensure that the hypotheses of
Corollary~\ref{cor:gate-b-raw-edge-divisibility-order-seven} are part of
the actual corrected source package, and to extend the edge law if a
later argument requires orders beyond \(7\).

\item \textbf{Source-bidegree compatibility.}
The source-bidegree projector must be applied before fixed-sector
coefficient extraction, determinant pairing, or quotienting.  At the raw
edge level, Corollary~\ref{cor:gate-b-raw-edge-divisibility-order-seven}
proves the required first-collision divisibility through ordinary order
\(7\) under its stated hypotheses:
\[
[z^7][\delta^1]\pi_{\mathfrak f_h}
(1-\Pi_{1,1})_{\mathrm{mix}}C_{\mathcal W}^{\mathrm{raw}}=0.
\]
After this holomorphic edge normalization is available, the corresponding
post-edge quotient formulation may be stated as
\[
Q_{7,m}^q(C)=Q_{7,m}^q(\Pi_{1,1}C)
\]
for the full corrected package, or equivalently as the quotient-gauge
condition for the already-projected representative
\(C_{\mathcal W}^{\corr}\).  The post-edge statement should not be used as
the proof of raw edge divisibility.

\item \textbf{Midpoint trivialization independence.}
The quotient class in
\[
\mathfrak f_m/\mathbf C A_5^{\mathfrak f}(m)
\]
must be independent of the auxiliary representative choices \(E_m,\ell_h\).

\item \textbf{Connection cancellation.}
For the residual monomials
\[
M_A=\dot r_3r_2\dot r_4,
\qquad
M_C=\dot r_3\dot r_2r_4,
\]
one must prove the pre-trivialization connection identities
\[
R_7^{(A)}
+
\sum_{i=1}^{7}\omega_m(Z_{7-i}[M_A],A_i)
=
0,
\]
and
\[
R_7^{(C)}
+
\sum_{i=1}^{7}\omega_m(Z_{7-i}[M_C],A_i)
=
0.
\]
Here \(Z_j[M]\) denotes the pre-trivialization edge coefficient of
Definition~\ref{def:lower-moving-package-coefficients}
(equivalently Definition~\ref{def:pre-trivialization-edge-table-after-gate-b});
it is not the post-trivialization coefficient \(B_j[M]\).

\item \textbf{UV-026 package quotient gauge.}
With the present definitions, \(C_{\mathcal W}^{\corr}\) already denotes the
amplitude-normalized, non-\((1,1)\) mixed package after midpoint
trivialization.  Thus UV-026 is the quotient-gauge statement
\[
B_7^{\mathfrak f}(C_{\mathcal W}^{\corr})
\in
\mathbf C A_5^{\mathfrak f}(m),
\]
equivalently
\[
Q_{7,m}^q(C_{\mathcal W}^{\corr})=0.
\]

If one instead works with an unprojected full corrected package
\(C_{\mathcal W}^{\corr,\mathrm{full}}\), then the equivalent full-package
formulation is
\[
B_7^{\mathfrak f}
\left(
(1-\Pi_{1,1})_{\mathrm{mix}}
C_{\mathcal W}^{\corr,\mathrm{full}}
\right)
\in
\mathbf C A_5^{\mathfrak f}(m),
\]
or
\[
Q_{7,m}^q(C_{\mathcal W}^{\corr,\mathrm{full}})
=
Q_{7,m}^q(\Pi_{1,1}C_{\mathcal W}^{\corr,\mathrm{full}}).
\]
The source-bidegree projection must be applied before fixed-sector
extraction, determinant pairing, or quotienting.
\end{enumerate}
\end{remark}

\section{Finite coefficient table contracts for the corrected-whitening package}
\label{app:finite-coefficient-table-contracts}

\begin{remark}[Active convention for this appendix]
\label{rem:active-convention-finite-coefficient-appendix}
This appendix is written under the homogeneous scalar-grade convention
\[
\operatorname{Gr}_a r
=
\frac{r^{(a+2)}(m)}{(a+2)!}(t-m)^{a+2}.
\]
Under this convention the ordinary matrix-output \(M^{[5]}\) route is
quarantined.  In particular, the \(Y\)-slot families
\[
L_1YR_1,\qquad L_2YR_0,\qquad L_0YR_2
\]
are absent at \(B_7\).  If the project later changes to an ordinary
matrix-output \(M^{[5]}\) convention, then those families must be reopened
and a no-double-counting theorem against grades \(0,2,4\) is required.
\end{remark}

This appendix records the finite coefficient tables needed by the current
corrected-whitening package frontier.  These are table contracts, not
completed source tables.  Their role is to make the finite computations
auditable by agents and scripts without re-opening convention choices
already fixed in the main text.

All tables below are understood:
\begin{enumerate}
\item before \(\Phi_K\);
\item before determinant scalarization;
\item in ordinary \(z\);
\item with source tags retained through the relevant \(112/122\) extraction;
\item in the same fixed-sector normalization used for
\(B_7^{\mathfrak f}\) and \(Q_{7,m}^q\).
\end{enumerate}

The finite Fr\'echet recurrences in
Section~\ref{sec:finite-frechet-coefficient-calculus} become proof tools
only after the source tables in this appendix are supplied.  They are not,
by themselves, quotient-gauge proofs.

\subsection{Stage 1 source-table contract}
\label{app:stage-one-source-table-contract}

The Stage 1 source-table generator requires the following scalar midpoint
jets:
\[
q_0^{(k)}(m),
\qquad
0\le k\le9,
\]
and
\[
(r_i^{[a]})^{(k)}(m),
\qquad
i=1,2,\quad a\in\mathcal A_{\mathrm{BC}},\quad 2\le k\le9.
\]
Thus the input table is
\[
\mathsf{Jet}_{\mathrm{Stage1}}
=
\left\{
q_0^{(0)},\ldots,q_0^{(9)};
(r_1^{[1]})^{(2)},\ldots,(r_1^{[1]})^{(9)};
(r_1^{[5]})^{(2)},\ldots,(r_1^{[5]})^{(9)};
(r_2^{[1]})^{(2)},\ldots,(r_2^{[1]})^{(9)};
(r_2^{[5]})^{(2)},\ldots,(r_2^{[5]})^{(9)}
\right\}_{m}.
\]

The current source supports ungraded one-pair source jets after affine
removal, but it does not yet supply the graded source jets above as a
proved theorem.  In particular, the baseline \(q_0\)-jet theorem
\[
q_0^{(k)}(m),\qquad 0\le k\le9,
\]
is still a separate source obligation.  Scale estimates do not substitute
for this finite table.

\subsection{Homogeneous scalar-grade status table}
\label{app:homogeneous-scalar-grade-status-table}

Under the homogeneous scalar-grade convention
\[
\operatorname{Gr}_a r
=
\frac{r^{(a+2)}(m)}{(a+2)!}(t-m)^{a+2},
\]
the current \(B_7^{\mathfrak f}\)-status of the seven cubic families is:
\[
\begin{array}{c|c|c}
\text{Family} & \text{Status at }B_7^{\mathfrak f}
& \text{Reason / hypothesis}
\\
\hline
L_1YR_1 & \text{absent} &
\text{middle homogeneous grade-five input starts too late}
\\
L_2YR_0 & \text{absent} &
\text{second Fr\'echet terms cannot lower ordinary }z\text{-order}
\\
L_0YR_2 & \text{absent} &
\text{right-hand mirror of }L_2YR_0
\\
L_2N_0R_1 & \text{order-admissible placement, aggregate closed} &
\text{cancels with }L_1N_0R_2\text{ by transpose-pairing}
\\
L_1N_0R_2 & \text{order-admissible placement, aggregate closed} &
\text{cancels with }L_2N_0R_1\text{ by transpose-pairing}
\\
L_3N_0R_0 & \text{order-admissible placement, aggregate closed conditionally} &
\text{cancels with }L_0N_0R_3\text{ under third-Fr\'echet midpoint parity}
\\
L_0N_0R_3 & \text{order-admissible placement, aggregate closed conditionally} &
\text{right-hand mirror of }L_3N_0R_0
\end{array}
\]

The three \(Y\)-slot absences are order-count statements, not determinant
proportionality statements.  The four non-\(Y\) rows are not asserted to
vanish placement by placement.  They close only after the indicated
aggregate pairing and fixed-sector projection.  If a future theorem
defines \(M^{[5]}\) as an ordinary matrix-output order-\(5\) projection,
then the \(Y\)-slot families must be reopened.  That route requires a
no-double-counting theorem against the natural homogeneous scalar grades
\(0,2,4\).

\subsection{Mixed-input parity table}
\label{app:mixed-input-parity-table}

The full source-linear scaled mixed input satisfies
\[
M(-z)=M(z)^T.
\]
Consequently:
\[
\begin{array}{c|c}
\text{Component of }M(z) & \text{ordinary }z\text{-parity}
\\
\hline
\text{diagonal entries} & \text{even}
\\
\text{off-diagonal symmetric channel} & \text{even}
\\
\text{off-diagonal antisymmetric channel} & \text{odd}
\end{array}
\]

Through ordinary order \(7\), the coefficient \([z^5]M\) lies in the
off-diagonal antisymmetric channel and is source-supported by
\[
r^{(2)}(m),\qquad r^{(4)}(m),\qquad r^{(6)}(m),
\]
not by the homogeneous grade-five source \(r^{(7)}(m)\).  Therefore
\([z^5]M\) should not be identified with homogeneous
\(\operatorname{Gr}_5 r\) without a separate parity-corrected
source-projection theorem.

\subsection{Source primitive support table}
\label{app:source-primitive-support-table}

For a homogeneous scalar source grade
\[
\operatorname{Gr}_a r
=
\frac{r^{(a+2)}(m)}{(a+2)!}(t-m)^{a+2},
\]
the same-point input first appears at ordinary pre-whitening order \(a\).
The mixed input must also respect the transpose-parity table above, and
need not have a nonzero coefficient at ordinary order \(a\).
\[
\begin{array}{c|c|c}
\text{Source grade} & \text{same-point input} & \text{mixed input}
\\
\hline
1
&
\delta G_\pm^{[1]}\text{ starts at }z^1
&
\text{use Lemma~\ref{lem:transpose-parity-source-linear-mixed-input}
for the corresponding mixed coefficient}
\\
5
&
\delta G_\pm^{[5]}\text{ starts at }z^5
&
[z^5]\mathfrak D_Q\delta N^{[5]}=0
\text{ under the homogeneous convention}
\end{array}
\]

The grade-zero sector corresponding to \(r^{(2)}(m)\) is not removed by
affine normalization and should not be erased from the source table.  Its
correct status after the Gate~B source-frame audit is more precise: grade
\(0\) remains a live moving Taylor source direction, but its
first-collision contribution is handled uniformly.  Amplitude terms of
total homogeneous grade \(7\), including zero-padded lower-shadow
multisets, vanish by odd total-grade subset-complement cancellation;
source-motion lower terms telescope against adjacent grades; and the sole
lower boundary
\[
(\partial_h\operatorname{Gr}_{0}r)_{\mathrm{lower}}
=
-r^{(2)}(h)(t-h)
\]
is affine and is removed by the \(2\)-jet source-frame projection.

\subsection{Baseline coefficient table contract}
\label{app:baseline-coefficient-table-contract}

The baseline tables needed for the current cubic gates are:
\[
G_\pm^{(0)}(z),\qquad
W_\pm(z):=G_\pm^{(0)}(z)^{-1/2},
\qquad
N_0(z):=N^{(0)}(z).
\]
Write
\[
W_\pm(z)=\sum_{n\ge0}W_{\pm,n}z^n,
\qquad
N_0(z)=\sum_{n\ge0}N_0[n]z^n.
\]

For the historical first non-\(Y\) coefficient audit
\(L_2N_0R_1/L_1N_0R_2\), the leading aggregate calculation uses only the
baseline mixed coefficient
\[
N_0[0].
\]
For the Stage 1 generator and later all-left/all-right audits, one needs
the ordinary-\(z\) baseline tables through the ranges specified in
Appendix~\ref{app:frechet-table-range-contract}.

\subsection{Fr\'echet table range contract}
\label{app:frechet-table-range-contract}

The finite Fr\'echet recurrences of
Section~\ref{sec:finite-frechet-coefficient-calculus} require the following
coefficient ranges for the active UV-026 gates.

\[
\begin{array}{c|c|c}
\text{Object} & \text{Range} & \text{Use}
\\
\hline
\Lambda_\pm^{[1]} & [z^1] & \text{grade-one first Fr\'echet input}
\\
\Lambda_\pm^{[5]} & [z^5] & \text{grade-five first Fr\'echet input}
\\
\mathcal L_-^{\{1,1\}} & [z^2] &
L_2N_0R_1\text{ with }(1,1)\text{ on the left}
\\
\mathcal L_-^{\{1,5\}} & [z^6] &
L_2N_0R_1\text{ with }(1,5)\text{ on the left}
\\
\mathcal R_+^{\{1,1\}} & [z^2] &
L_1N_0R_2\text{ with }(1,1)\text{ on the right}
\\
\mathcal R_+^{\{1,5\}} & [z^6] &
L_1N_0R_2\text{ with }(1,5)\text{ on the right}
\\
G^-_{324} & [z^3],\ldots,[z^7] &
\text{third-Fr\'echet parity audit / convention check}
\\
G^+_{324} & [z^3],\ldots,[z^7] &
\text{right-side mirror in the same convention}
\end{array}
\]

Here \(\mathcal L\) and \(\mathcal R\) denote the left and right second
Fr\'echet inverse-square-root outputs in the same fixed
\(B_7^{\mathfrak f}\) normalization.  The third-Fr\'echet rows are now a
normalization check, not a pending Gate~A determinant table.  Under the
symmetric multilinear third-Fr\'echet convention used in
Lemma~\ref{lem:aggregate-closure-leading-all-left-all-right-cubic-gate},
the all-left/all-right pair closes in aggregate.  If an ordered-triple
coefficient convention is used instead, the ordered-permutation aggregate
must be included consistently on both sides before invoking the
transpose-pairing argument.

\subsection{Diagnostic leading table for \(L_2N_0R_1/L_1N_0R_2\)}
\label{app:l2n0r1-leading-table}

Before the aggregate Gate~A closure was isolated, the first homogeneous
non-\(Y\) coefficient table to audit was
\[
L_2N_0R_1
\qquad\text{and}\qquad
L_1N_0R_2.
\]
The table below is retained as a diagnostic coefficient contract.  It is
not a remaining live Gate~A obstruction and it is not a family-by-family
determinant-gauge requirement.
The leading table is:
\[
\begin{array}{c|c}
\text{Tag vector} & \text{leading fixed-sector coefficient}
\\
\hline
C_{112}^{L_2N_0R_1}
&
\pi_{\mathfrak f}
\left(
[z^2]\mathcal L_-^{\{1,1\}}\,
N_0[0]\,
[z^5]\Lambda_+^{[5]}
\right)
\\[1.2ex]
C_{122}^{L_2N_0R_1}
&
\pi_{\mathfrak f}
\left(
[z^6]\mathcal L_-^{\{1,5\}}\,
N_0[0]\,
[z^1]\Lambda_+^{[1]}
\right)
\\[1.2ex]
C_{112}^{L_1N_0R_2}
&
\pi_{\mathfrak f}
\left(
[z^1]\Lambda_-^{[1]}\,
N_0[0]\,
[z^6]\mathcal R_+^{\{1,5\}}
\right)
\\[1.2ex]
C_{122}^{L_1N_0R_2}
&
\pi_{\mathfrak f}
\left(
[z^5]\Lambda_-^{[5]}\,
N_0[0]\,
[z^2]\mathcal R_+^{\{1,1\}}
\right)
\end{array}
\]

Each entry is a fixed-sector vector in
\[
\mathfrak f=\operatorname{span}\{I,S\}.
\]
These four vectors are retained as a coefficient-table contract, but they
are no longer required to pass an \(A_5^{\mathfrak f}\)-line test
family-by-family.  Under the homogeneous scalar-grade convention,
Lemma~\ref{lem:aggregate-closure-leading-l2n0r1-l1n0r2-gate} proves the
aggregate cancellation
\[
C_{112}^{L_2N_0R_1}+C_{112}^{L_1N_0R_2}=0,
\qquad
C_{122}^{L_2N_0R_1}+C_{122}^{L_1N_0R_2}=0
\]
in \(\mathfrak f\).  This is stronger than quotient-gauge and does not use
the \(A_5^{\mathfrak f}\)-line.

If this table is later used outside the homogeneous scalar-grade
convention, then the old determinant coordinate
\[
uv_5-u_5v=0
\]
is only a good-patch audit coordinate after the fixed-sector vector has
been computed; it is not the definition of the quotient object and not the
current Gate~A closure mechanism.

\begin{remark}[Aggregate closure versus coefficient-table diagnostics]
\label{rem:aggregate-closure-versus-coefficient-table-diagnostics}
The coefficient tables in this appendix are diagnostic contracts.  They
are designed to let agents reproduce the finite entries used by the
aggregate proofs.  They should not be read as replacing the aggregate
statements in the main text.  In particular, a nonzero placement-level
vector is not a counterexample to Gate~A unless the full aggregate
specified in the relevant lemma also survives in \(\mathfrak f\).
\end{remark}

\subsection{All-left/all-right cubic aggregate table}
\label{app:all-left-all-right-cubic-aggregate-table}

The all-left and all-right families
\[
L_3N_0R_0,
\qquad
L_0N_0R_3
\]
are live placements by order count, but they are not remaining
independent Gate~A residuals under the third-Fr\'echet midpoint-parity
convention.  Let
\[
G:=G^-_{7;1,1,5},
\qquad
W:=W_{-,0}=W_{+,0},
\qquad
N:=N_0[0].
\]
The two aggregate fixed-sector contributions are
\[
\pi_{\mathfrak f}(GNW)
\]
and
\[
-\pi_{\mathfrak f}(WNG).
\]
Since \(G,N,W\) are symmetric,
\[
(GNW)^T=WNG,
\]
and since \(\pi_{\mathfrak f}\) is transpose-invariant,
\[
\pi_{\mathfrak f}(GNW)=\pi_{\mathfrak f}(WNG).
\]
Therefore
\[
L_3N_0R_0+L_0N_0R_3=0
\quad\text{in }\mathfrak f
\]
at the aggregate fixed-sector level, conditional on the third-Fr\'echet
midpoint-parity convention.

Thus the required table is not a pending quotient-gauge determinant
table.  The remaining obligation is to verify that the third-Fr\'echet
coefficient normalization used by any coefficient agent is the same
symmetric multilinear normalization used in
Lemma~\ref{lem:aggregate-closure-leading-all-left-all-right-cubic-gate}.
If an ordered-triple convention is used instead, the ordered-permutation
aggregate must be included consistently on both sides before applying the
transpose-pairing argument.

\subsection{Validation checklist for coefficient agents}
\label{app:validation-checklist-coefficient-agents}

A valid coefficient-table submission must specify:
\[
\begin{array}{c|c}
\text{Item} & \text{Required declaration}
\\
\hline
\text{ordinary variable} & z\text{ after }s=z/Q^2
\\
\text{stage} & \text{pre-}\Phi_K
\\
\text{source tags} & \tau_1,\tau_2\text{ retained through }112/122
\\
\text{same-point input} & \text{raw vs. }\mathfrak D_Q\text{ convention stated}
\\
\text{mixed input} & M_i^{[a]}=\operatorname{Gr}_a(\mathfrak D_Q\delta N_i^{\lin})
\\
\text{baseline table} & N_0[0]\text{ and needed }W_{\pm,n}
\\
\text{projection} & \pi_{\mathfrak f}\text{ before determinant test}
\\
\text{determinant coordinate}
&
\text{allowed only after the fixed-sector vector is computed; never as the definition}
\\
\text{ordered vs symmetric source slots}
&
\text{state whether Fr\'echet coefficients are ordered-labelled or
symmetric;
if ordered, include all labelled permutations before applying aggregate
transpose cancellation}
\end{array}
\]

For aggregate Gate~A tables, determinant proportionality to
\(A_5^{\mathfrak f}\) is not the closure mechanism unless explicitly
stated.  The primary object is the fixed-sector vector in \(\mathfrak f\),
and the primary test is the aggregate identity in \(\mathfrak f\).  The
determinant coordinate
\[
uv_5-u_5v
\]
is only a good-patch coordinate after the relevant fixed-sector vector or
aggregate has already been formed.

A submission is not valid if it:
\begin{enumerate}
\item argues after \(\Phi_K\);
\item uses determinant scalarization as the definition of the quotient;
\item replaces the active homogeneous grade by a matrix-output \(M^{[5]}\)
without the no-double-counting theorem;
\item treats recurrence formulas as quotient-gauge proofs without
computing the fixed-sector vectors;
\item omits the baseline \(q_0\)-jet convention or the grade-zero sector.
\end{enumerate}

\begin{remark}[Final status of the Gate A/Gate B package layer]
\label{rem:final-status-gate-a-gate-b-package-layer}
At this stage the corrected-package construction has the following status.

Gate~A is closed at the aggregate fixed-sector level under the active
homogeneous scalar-grade convention and the stated third-Fr\'echet
midpoint parity convention.  It should not be reopened as a
family-by-family determinant-gauge problem.

Gate~B is closed only in the conditional raw edge-divisibility sense of
Corollary~\ref{cor:gate-b-raw-edge-divisibility-order-seven}: its use
depends on the stated midpoint-parity, source-frame telescoping,
\(2\)-jet source-frame normalization, and full-Taylor-tower-before-truncation
hypotheses.  If one of those hypotheses is not verified in the final
source model, Gate~B reverts to an open raw edge-divisibility target, not
to a post-edge quotient-compatibility statement.

After Gate~B, the next local package table is the pre-trivialization
\(Z_j[M]\)-table, and the next local identities are
\[
R_7^{(A)}+\sum_{i=1}^{7}\omega_m(Z_{7-i}[M_A],A_i)=0,
\qquad
R_7^{(C)}+\sum_{i=1}^{7}\omega_m(Z_{7-i}[M_C],A_i)=0.
\]
The post-trivialization coefficients \(B_j[M]\) must not be substituted
into these connection sums.

Finally, none of Gate~A, Gate~B, or the \(Z_j[M]\)-connection package by
itself proves the global package-to-\(\Xi_\zeta^{(N)}\) hidden-state
bridge or the final finite-core contradiction.  Those remain downstream
theorem obligations.
\end{remark}

\section{Upstream Gate--1 arithmetic inputs after the ordinary-route audit}
\label{sec:upstream-gate-one-arithmetic-inputs}

The downstream local package of Section~\ref{sec:downstream-local-closure-package}
reduces row-energy production to Gate~1 and finite D4 row conversion.  The
present section records the current upstream status for Gate~1.  This is a
conditional arithmetic ledger: the inputs below are not proved in this draft.

\subsection{Branch-complete local \(L^2\) Gate--1}
\label{subsec:branch-complete-local-l2-gate-one}

The active Gate--1 input is local scalar \(L^2\)-source energy of
\[
  u=q''_{\can}.
\]
For a branch \(b\in\{\mathrm{ord},\mathrm{tr},\mathrm{sing}\}\), define
\[
  F_{b,m}(x):=\lambda_b^3q''_{\can}(m+\lambda_bx),
\]
and
\[
  E_{1,b}^{L^2}(m)
  :=
  \int_{I_b}|F_{b,m}(x)|^2w_b(x)\,dx,
\]
where \(w_b\ge0\) is a fixed smooth local branch window.  Gate--1 on the packet
\(m\) means
\[
  \boxed{
  E_{1,b}^{L^2}(m)\ge Q^{-A}.
  }
\]
This definition is branch-independent after normalization and is the scalar
source input consumed by the Gate--2 \(L^2\)-real-rigidity theorem.

\begin{hypothesis}[Aggregate scalar \(L^2\) observability for \(q''_{\can}\)]
\label{hyp:agg-l2-qcan-gate-one}
Let
\[
  W_{\mathrm{branch}}(t)
\]
be the routed branch-window weight obtained from the ordinary, transition, and
singular packet windows:
\[
  W_{\mathrm{branch}}(t)
  :=
  \sum_b\sum_{m\in\mathcal C_b}
  \rho_b(m)\lambda_b^5
  w_b\!\left(\frac{t-m}{\lambda_b}\right),
\]
with the evident modification if \(w_b\) has already been normalized after
change of variables.  The aggregate Gate--1 input is
\[
  \boxed{
  \int |q''_{\can}(t)|^2W_{\mathrm{branch}}(t)\,dt
  \ge
  Q^{-A}\mathcal N_{\can}(W_{\mathrm{branch}}).
  }
\]
The reference normalization must also satisfy
\[
  \boxed{
  \mathcal N_{\can}(W_{\mathrm{branch}})
  \ge
  Q^{-A}
  \sum_b\sum_{m\in\mathcal C_b}\rho_b(m).
  }
\]
Both estimates are for the full canonical source
\[
  q''_{\can}=B_2+P_2+D_2+E_2.
\]
\end{hypothesis}

\begin{lemma}[Oscillatory packet kernel decay]
\label{lem:packet-oscillatory-kernel-decay}
Let \(\psi_p(t)=a_p(t)e^{i\omega_p t}\) and \(\psi_q(t)=a_q(t)e^{i\omega_q t}\),
with \(a_p,a_q\in C^N\) supported on packet cores and satisfying
\[
\|\partial_t^j a_p\|_\infty+\|\partial_t^j a_q\|_\infty \le Q^{Cj}
\qquad (0\le j\le N).
\]
Then
\[
\langle \psi_p,\psi_q\rangle_{W_{\mathrm{branch}}}
=\widehat{F_{pq}}(\omega_p-\omega_q),
\qquad
F_{pq}(t)=a_p(t)\overline{a_q(t)}W_{\mathrm{branch}}(t),
\]
and hence
\[
|\langle \psi_p,\psi_q\rangle_{W_{\mathrm{branch}}}|
\le C_N Q^{C_N}(1+|\omega_p-\omega_q|)^{-N}.
\]
In particular, if \(|\omega_p-\omega_q|\ge Q^{-\delta}\) outside resonant blocks,
then the off-diagonal correlation is \(Q^{-A}\)-small for \(N\) large enough.
\end{lemma}

\begin{lemma}[Source frame coercivity for Gate--1]
\label{lem:source-frame-coercivity}
Let
\[
\mathcal S_X=\operatorname{span}\{B_2,P_{2,\le X},D_{2,\le X}\},
\qquad
F_X=B_2+P_{2,\le X}+D_{2,\le X}.
\]
Let \(\{\psi_i\}_{i=1}^N\) be packet-adapted source atoms spanning \(\mathcal S_X\),
normalized in the reference metric:
\[
\|\psi_i\|_{L^2(W_{\mathrm{ref}})}=1.
\]
Define the normalized packet Gram matrix
\[
M_{ij}:=\langle \psi_i,\psi_j\rangle_{W_{\mathrm{branch}}},
\qquad
M=G_{\mathrm{ref}}^{-1/2}G_{\mathrm{branch}}G_{\mathrm{ref}}^{-1/2},
\]
and the residual correction
\[
E:=G_{\mathrm{ref}}^{-1/2}\bigl(G_{\mathrm{branch}}(q''_{\can})-G_{\mathrm{branch}}(F_X)\bigr)G_{\mathrm{ref}}^{-1/2}.
\]

Assume the packet-core micro-covering and oscillatory-kernel hypotheses already
recorded in this subsection, namely:

\begin{enumerate}
\item the packet cores form a polynomially thick sampling set at scale
\(h\asymp \log Q/Q\), with bounded overlap and
\[
W_{\mathrm{branch}}\ge Q^{-C}W_{\mathrm{ref}}
\quad\text{on the packet cores};
\]

\item the packet amplitudes satisfy polynomial derivative bounds so that the
off-diagonal packet kernels obey polynomial oscillatory decay, and the
resonant clusters are finite-dimensional;

\item the residual term satisfies the EF completion bound
\[
\|E\|_{\operatorname{op}}\le Q^{-B};
\]

\item the reference mean-square does not collapse, so the canonical finite
source has nontrivial \(W_{\mathrm{ref}}\)-mass.
\end{enumerate}

Then, after grouping indices into resonant packet clusters \(B\),
\[
\min_B \lambda_{\min}(M_B)-\|M_{\off}\|_{\operatorname{op}}-\|E_B\|_{\operatorname{op}}\ge Q^{-A}
\]
for some absolute \(A>0\), and in particular
\[
\delta_{\mathrm{src}}
:=
\min_i\!\left(M_{ii}-\sum_{j\neq i}|M_{ij}|\right)-\|E\|_{\operatorname{op}}
\ge Q^{-A}.
\]
Equivalently,
\[
\lambda_{\min}(M+E)\ge Q^{-A},
\qquad
\|f\|_{L^2(W_{\mathrm{branch}})}^2\ge Q^{-A}\|f\|_{L^2(W_{\mathrm{ref}})}^2
\quad(\forall f\in\mathcal S_X).
\]
\end{lemma}

\begin{proof}
Normalize the packet atoms so that \(\|\psi_i\|_{L^2(W_{\mathrm{ref}})}=1\).  The
micro-covering hypothesis gives a diagonal lower bound on each packet core:
\[
M_{ii}=\langle \psi_i,\psi_i\rangle_{W_{\mathrm{branch}}}\ge 1-Q^{-A_1},
\]
since \(W_{\mathrm{branch}}\ge Q^{-C}W_{\mathrm{ref}}\) on the core and the EF
completion bound controls the tail outside the core.

For \(i\neq j\), write \(\psi_i(t)=a_i(t)e^{i\omega_i t}\) and
\(\psi_j(t)=a_j(t)e^{i\omega_j t}\).  The oscillatory-kernel hypothesis gives
\[
|M_{ij}|
=
\left|\int a_i(t)\overline{a_j(t)}W_{\mathrm{branch}}(t)e^{i(\omega_i-\omega_j)t}\,dt\right|
\le C_NQ^{C_N}(1+|\omega_i-\omega_j|)^{-N}.
\]
After grouping the packet indices into resonant blocks, only finitely many
near-resonant neighbors remain in each row, and the row-sum bound becomes
\[
\sum_{j\neq i}|M_{ij}|\le Q^{-A_2}.
\]
Choosing \(N\) large enough and using the finite block size yields
\[
\min_i\!\left(M_{ii}-\sum_{j\neq i}|M_{ij}|\right)\ge 1-Q^{-A_3}.
\]

Now incorporate the residual perturbation \(E\).  The EF completion bound gives
\[
\|E\|_{\operatorname{op}}\le Q^{-B}.
\]
Taking \(B\) larger than the constants from the diagonal and off-diagonal
estimates, we obtain
\[
\delta_{\mathrm{src}}
=
\min_i\!\left(M_{ii}-\sum_{j\neq i}|M_{ij}|\right)-\|E\|_{\operatorname{op}}
\ge Q^{-A}.
\]

By Gershgorin, this implies
\[
\lambda_{\min}(M+E)\ge Q^{-A}.
\]
Since \(M+E\) is exactly the normalized source frame operator on \(\mathcal S_X\),
\[
\|f\|_{L^2(W_{\mathrm{branch}})}^2\ge Q^{-A}\|f\|_{L^2(W_{\mathrm{ref}})}^2
\quad\text{for all }f\in\mathcal S_X.
\]
This is the desired source-survival inequality.
\end{proof}

\begin{corollary}[Gate--1 source survival]
\label{cor:gate-one-source-survival}
The aggregate source layer satisfies
\[
\int |q''_{\can}(t)|^2W_{\mathrm{branch}}(t)\,dt
\ge
Q^{-A}\int |q''_{\can}(t)|^2W_{\mathrm{ref}}(t)\,dt,
\]
equivalently \(\mathrm{Agg}L^2(q''_{\can})\) holds with polynomial loss.
\end{corollary}

\begin{remark}[Failure labels]
\label{rem:source-failure-labels}
If the source frame bound fails, the only remaining failure modes are:
branch holes, off-diagonal packet correlation too large, residual completion
failure, or reference mean-square collapse.
\end{remark}

\begin{remark}[Resonant block conditioning]
\label{rem:source-resonant-block-conditioning}
If \(|\omega_p-\omega_q|<Q^{-\delta}\), the pair must be handled inside a
resonant block \(B\).  It is enough to prove
\[
G_B\ge Q^{-A}G_{B,\mathrm{ref}}
\]
for every resonant block.
\end{remark}

\begin{lemma}[Fubini form of branch-complete \(L^2\)-Gate--1]
\label{lem:branch-complete-l2-gate-one-fubini}
Let \(u=q''_{\can}\), and let
\[
  E_{1,b}^{L^2}(m)
  =
  \int
  \left|
    \lambda_b(m)^3 u(m+\lambda_b(m)x)
  \right|^2
  w_b(x)\,dx.
\]
Then
\[
  E_{1,b}^{L^2}(m)
  =
  \int
  |u(t)|^2
  \lambda_b(m)^5
  w_b\!\left(\frac{t-m}{\lambda_b(m)}\right)
  dt.
\]
Consequently
\[
  \sum_b\sum_{m\in\mathcal C_b}
  \rho_b(m)E_{1,b}^{L^2}(m)
  =
  \int |u(t)|^2 W_{\mathrm{branch}}(t)\,dt,
\]
where
\[
  W_{\mathrm{branch}}(t)
  =
  \sum_b\sum_{m\in\mathcal C_b}
  \rho_b(m)\lambda_b(m)^5
  w_b\!\left(\frac{t-m}{\lambda_b(m)}\right).
\]
Thus \(\mathrm{Agg}L^2(q''_{\can})\) is exactly the weighted source
mean-square lower bound
\[
  \int |q''_{\can}(t)|^2W_{\mathrm{branch}}(t)\,dt
  \ge
  Q^{-A}\mathcal N_{\can}(W_{\mathrm{branch}}).
\]
In particular, the proof of Gate--1 splits into:
\[
  \boxed{\text{coverage of the }|q''_{\can}|^2\text{-energy measure}}
  \qquad\text{and}\qquad
  \boxed{\text{canonical source mean-square nontriviality}.}
\]
\end{lemma}

\begin{remark}[Branch observables are proof tools, not Gate--1 definitions]
\label{rem:branch-observables-proof-tools-only-upstream}
The older branch observables
\[
  T_{\can}X_{\can},
  \qquad
  \mathcal J_R^{\mathrm{tr}},
  \qquad
  \mathbf P_R-\mathbf G_R-\mathbf T_R,
  \qquad
  P_R[\lambda^3q''_{\can}]
\]
are retained only as possible proof tools for
Hypothesis~\ref{hyp:agg-l2-qcan-gate-one}.  They are not the Gate--1
definition.  In particular, raw singular sine moments and projected singular
low modes are diagnostics unless they are used to prove local scalar
\(L^2\)-energy of the actual canonical source.

The safe branch-to-downstream rule is:
\[
  \boxed{
  \text{every branch must certify local scalar }L^2\text{ energy of }q''_{\can}.
  }
\]
\end{remark}

\begin{remark}[Finite approximation and \(L^2\)-anti-lock]
\label{rem:finite-approximation-l2-antilock}
The aggregate Gate--1 theorem may be attacked through finite canonical
approximants such as
\[
  F_X:=B_2+P_{2,\le X}+D_{2,\le X},
\]
with a completion estimate
\[
  \|q''_{\can}-F_X\|_{L^2(W_{\mathrm{branch}})}^2
  +
  \|q''_{\can}-F_X\|_{L^2(W_{\mathrm{ref}})}^2
  \le
  Q^{-A_{\rm big}}
  \|F_X\|_{L^2(W_{\mathrm{ref}})}^2.
\]
The hard finite theorem is then a weighted \(L^2\) observability estimate
\[
  \boxed{
  \|F_X\|_{L^2(W_{\mathrm{branch}})}^2
  \ge
  Q^{-A}
  \|F_X\|_{L^2(W_{\mathrm{ref}})}^2.
  }
\]
The principal failure mode is \(L^2\)-level anti-lock:
\[
  B_2+P_{2,\le X}+D_{2,\le X}\approx0
\]
on the observed branch windows.  This is a lower-frame/spectral observability
problem for the actual canonical source, not a pointwise finite-row problem.
\end{remark}

\begin{remark}[Finite AFE anti-lock and the central square-root tube]
\label{rem:finite-afe-antilock-central-tube}
After finite truncation
\[
  F_X=B_2+P_{2,\le X}+D_{2,\le X},
\]
the first nontrivial \(L^2\)-Gate--1 obstruction is finite AFE anti-lock:
\[
  \|B_2+P_{2,\le X}+D_{2,\le X}\|_{L^2(W_{\mathrm{branch}})}^2
  \ge
  Q^{-A}
  \|B_2+P_{2,\le X}+D_{2,\le X}\|_{L^2(W_{\mathrm{ref}})}^2.
\]
For the prime--dual component this is controlled by the fold
\[
  x\longmapsto 1-x.
\]
The only serious fold-degeneracy region is the central square-root tube
\[
  \Omega_{\mathrm{cen}}(\delta)
  :=
  \left\{x:\left|x-\frac12\right|<\delta\right\},
  \qquad
  \delta=R^{-C}.
\]
On the complement
\[
  \Omega_{\mathrm{out}}(\delta)
  :=
  \left\{x:\left|x-\frac12\right|\ge\delta\right\},
\]
the ratios
\[
  \frac{x^2}{(1-x)^2},
  \qquad
  \frac{(1-x)^2}{x^2}
\]
are \(R^{-C}\)-transverse.  Thus finite AFE anti-lock outside
\(\Omega_{\mathrm{cen}}\) reduces to a bounded-degree fold-transversality
estimate.

The central tube is not an exception to be discarded.  It is a rescaled branch.
Writing
\[
  x=\frac12+\delta y,\qquad |y|\le1,
\]
the graph obstruction for
\[
  f(x):=\frac1{m_{\can}(x)}
\]
becomes
\[
  \operatorname{dist}
  \left(
    f\!\left(\frac12+\delta y\right),
    \mathbf C_R[y]
  \right).
\]
For the model multiplier \(m_{\can}(x)=\kappa x^2\),
\[
  \frac1{m_{\can}(\frac12+\delta y)}
  =
  \frac4\kappa(1+2\delta y)^{-2},
\]
and the degree-\(R\) tail is of size
\[
  \asymp_R \delta^{R+1}.
\]
Thus choosing \(\delta=R^{-C}\) retains only polynomial loss:
\[
  \delta^R=R^{-CR}=Q^{-O(1)}.
\]
The same central-tube routing is used both for
\(\mathrm{Agg}L^2(q''_{\can})\) and for the graph/O2/O3 branch-survival audit.
\end{remark}

\begin{remark}[Degree-one fold anti-lock as the model for determinant branch survival]
\label{rem:degree-one-fold-antilock-models-determinant-survival}
The \(L^2\)-Gate--1 anti-lock theorem is the degree-one model for the
finite/O2/O3 determinant branch-survival theorem.  At degree one, the relevant
branches are
\[
  B_2,\qquad P_{2,\le X},\qquad D_{2,\le X}.
\]
At determinant degree \(n=R+2\), the full determinant decomposes into bounded
branch-charge sectors
\[
  D_{\mathrm{full}}=\sum_r D_r.
\]
After dyadic freezing of smooth denominators and the Stack-C2 weight, only
finitely many near-pure sectors can interact with the pure plus sector:
\[
  D_{-n},\quad D_{-n+1},\ldots,D_{-n+d_0},
  \qquad d_0=O(1).
\]
Thus fixed-\(\alpha\) determinant survival may be formulated as a finite Schur
complement
\[
  D_{-n}\notin
  \operatorname{span}\{D_{-n+1},\ldots,D_{-n+d_0}\},
\]
rather than as an uncontrolled \(O(R)\)-branch forest.  This is the determinant
analogue of prime--dual fold anti-lock.
\end{remark}

\begin{remark}[Failure modes for \(\mathrm{Agg}L^2(q''_{\can})\)]
\label{rem:agg-l2-qcan-failure-modes}
The branch-complete \(L^2\) Gate--1 theorem can fail only through:
\begin{enumerate}
\item branch holes: \(W_{\mathrm{branch}}\) misses where \(q''_{\can}\) has energy;
\item finite packet anti-lock:
\[
  B_2+P_{2,\le X}+D_{2,\le X}\approx0
\]
on \(W_{\mathrm{branch}}\);
\item residual failure:
\[
  \|q''_{\can}-F_X\|_{L^2(W_{\mathrm{branch}})}^2
  +
  \|q''_{\can}-F_X\|_{L^2(W_{\mathrm{ref}})}^2
  \not\ll
  \|F_X\|_{L^2(W_{\mathrm{ref}})}^2;
\]
\item reference collapse:
\[
  \mathcal N_{\can}(W_{\mathrm{branch}})\ll
  Q^{-A}\sum_{b,m}\rho_b(m);
\]
\item genuine local source flatness on most branch windows.
\end{enumerate}
These replace the older finite-row pointwise projection obstruction as the
active Gate--1 failure modes.
\end{remark}

\subsection{Ordinary packet-frequency Gate--1 inputs}
\label{subsec:ordinary-packet-frequency-wronskian-route}

This subsection records the historical finite-row ordinary Wronskian route.
It is no longer the primary Gate--1 definition once
Hypothesis~\ref{hyp:agg-l2-qcan-gate-one} is adopted.  The old O1/O2/O3
package remains useful as an alternate proof tool and as a no-go ledger for
pointwise zero-conditioned nonrecurrence.  The local scalar \(L^2\)-source
input for Gate--2 should not be replaced by \(T_{\can}X_{\can}\), transition
jets, or singular sine/projection diagnostics unless the proof explicitly
returns to this finite-row route.

In the ordinary packet-frequency regime
\[
\lambda Q\asymp1,
\qquad
R\asymp\frac{\log Q}{\log\log Q},
\qquad
\mathcal D=\lambda\partial_m,
\]
the ordinary Gate--1 route is retained only as the named arithmetic input package
\[
\boxed{
\mathrm{O1}+\mathrm{O2}+\mathrm{O3}
\Longrightarrow
\text{ordinary Gate--1}.
}
\]
The source object is always the full canonical curvature
\[
q''_{\can}=B_2+P_2+D_2+E_2.
\]

\begin{definition}[O1: fixed-width cell nonrecurrence]
\label{def:O1-fixed-width-cell-nonrecurrence}
For each admissible ordinary-scale zero \(Z_\alpha(m_0)=0\), let \(B_r\) denote
the full canonical fixed-width packet cell vector, including main branch,
folded dual branch, gamma/background, and error contributions in the selected
AFE convention.  O1 is
\[
\boxed{
\sum_r |B_r(m_0,\alpha)|^2\ge Q^{-A}.
}
\]
This is an actual-orbit nonrecurrence input.  It is not implied by local cell
geometry, diagonal mass, or phase-relaxed coefficient magnitudes.
\end{definition}

\begin{definition}[O1b finite-frame folded audit]
\label{def:O1b-finite-frame-folded-audit}
In the historical ordinary finite-row route, write the main and folded-dual
cell amplitudes as
\[
  a_r=B_r^+(m_0,\alpha),
  \qquad
  b_r=B_r^-(m_0,\alpha),
  \qquad
  e=e^{i\Theta_Z(m_0,\alpha)}.
\]
Let
\[
  c=(a_1,\ldots,a_M,eb_1,\ldots,eb_M)\in\mathbf C^{2M}.
\]
Let
\[
  \tau_1,\ldots,\tau_{2M}
  =
  \sigma_1^+,\ldots,\sigma_M^+,\sigma_1^-,\ldots,\sigma_M^-
\]
be the plus and folded-dual branch-symbol nodes.  The \(O1b\) finite-frame
audit is the assertion that
\[
  \sum_{\ell=0}^{R+1}
  \left|
    \sum_{j=1}^{2M}p_\ell(\tau_j)c_j
  \right|^2
  \ge
  Q^{-A}\sum_{j=1}^{2M}|c_j|^2 .
\]
At an admissible zero \(Z_\alpha(m_0)=0\), the value row may vanish, but the
derivative rows still carry the folded branch mass if this finite-frame audit
and the cell-symbol bridge hold.
\end{definition}

\begin{lemma}[Coarse paired cells discharge finite-row \(O1b\)]
\label{lem:O1b-coarse-paired-cells-discharge}
Assume
\[
  2M\le R+2.
\]
Assume also that the branch-symbol nodes are polynomially separated, that the
row polynomials \(p_0,\ldots,p_{R+1}\) are polynomially equivalent on the two
branch-symbol intervals to a stable polynomial basis, and that the cell-symbol
bridge is valid.  Then
\[
  \sum_{\ell=0}^{R+1}
  \left|
    \sum_{r=1}^{M}
    \left(
      p_\ell(\sigma_r^+)B_r^+
      +
      p_\ell(\sigma_r^-)e^{i\Theta_Z}B_r^-
    \right)
  \right|^2
  \ge
  Q^{-A}
  \sum_{r=1}^{M}\left(|B_r^+|^2+|B_r^-|^2\right).
\]
Thus folded value anti-lock is not an independent arithmetic obstruction in
the coarse-cell finite-row setting.
\end{lemma}

\begin{proof}[Proof sketch]
Since \(2M\le R+2\), the row space contains enough polynomial degrees to
interpolate arbitrary data on the \(2M\) branch nodes.  The interpolation
loss is at most
\[
  \exp(O(M\log M)).
\]
For
\[
  M\asymp R\asymp \frac{\log Q}{\log\log Q},
\]
this is \(Q^{O(1)}\).  Hence the evaluation matrix
\[
  V_{\ell j}=p_\ell(\tau_j)
\]
has polynomial lower singular value.  The cell-symbol bridge transfers this
model finite-frame estimate to the actual canonical cell sums.
\end{proof}

\begin{remark}[Cell-symbol bridge for \(O1b\)]
\label{rem:O1b-cell-symbol-bridge}
The deterministic \(O1b\) finite-frame lemma applies to symbol-compatible cell
amplitudes.  For hard cells, the derivative row contains
\[
  \sum_{n\in I_r^\pm}p_\ell(\sigma_n^\pm)A_n^\pm,
\]
not automatically
\[
  p_\ell(\sigma_r^\pm)\sum_{n\in I_r^\pm}A_n^\pm.
\]
At depth \(\ell\sim R\), symbol variation over cells of width \(1/R\) is not
automatically perturbative.  Therefore one must either use cell-dual polynomial
observables or prove a depth-\(R\) hard-cell oscillation lemma.
\end{remark}

\begin{definition}[O2: cofactor-vector nonrecurrence]
\label{def:O2-cofactor-vector-nonrecurrence}
Let \(\mathbf M(m_0,\alpha)\) be the cofactor/Hodge-dual vector of the canonical
\(q''_{\can}\)-jet block at an ordinary zero \(Z_\alpha(m_0)=0\).  O2 is
\[
\boxed{
|\mathbf M(m_0,\alpha)|\ge Q^{-A}.
}
\]
Averaged exterior mass or large-sieve estimates do not imply this pointwise
zero-conditioned statement.
\end{definition}

\begin{definition}[O3: prime-log Wronskian angle nonalignment]
\label{def:O3-wronskian-angle-nonalignment}
Let \(\mathbf z'(m_0,\alpha)\) be the \(Z_\alpha\)-row derivative vector at an
ordinary zero, and let \(\mathbf M(m_0,\alpha)\) be the canonical cofactor
vector.  O3 is
\[
\boxed{
|\langle \mathbf z'(m_0,\alpha),\mathbf M(m_0,\alpha)\rangle|
\ge Q^{-A}|\mathbf z'(m_0,\alpha)|\,|\mathbf M(m_0,\alpha)|.
}
\]
O3 is an actual prime-log incidence theorem.  It is not implied by O1 and O2.
\end{definition}

\begin{theorem}[Ordinary finite-row Gate--1 from O1--O3, historical route]
\label{thm:ordinary-gate-one-from-O1-O2-O3}
Assume O1, O2, and O3 for every admissible ordinary-scale zero.  Then the
ordinary packet-frequency branch supplies Gate--1 for the full canonical source:
\[
\boxed{
\int_I |q''_{\can}(m)|^2\,dm\ge Q^{-A}|I|.
}
\]
\end{theorem}

\begin{remark}[Status of the finite-row ordinary route]
\label{rem:ordinary-finite-row-route-historical-status}
The finite-row ordinary route is not deleted: it remains a valid conditional
sub-architecture.  Its current status is that O1, O2, and O3 are
zero-conditioned actual-orbit nonrecurrence inputs.  In particular:
\[
  \text{O1 same-branch mass}\not\Rightarrow\text{folded nonalignment},
\]
However, the folded subproblem \(O1b\) is finite-frame rather than
fundamentally arithmetic once the paired-cell count is chosen with
\[
  2M\le R+2
\]
and the cell-symbol bridge of Remark~\ref{rem:O1b-cell-symbol-bridge} is
available.  This does not repair the main ordinary Gate--1 dimension problem:
the remaining decisive ordinary obstruction is still the actual-orbit
\(O1a/O2/O3\) package for the full canonical source.
\[
  \text{exterior coefficient mass}\not\Rightarrow\text{O2},
\]
and
\[
  \text{O1+O2}\not\Rightarrow\text{O3}.
\]
The active Gate--1 route is instead the branch-complete \(L^2\) theorem
\[
  \mathrm{Agg}L^2(q''_{\can}).
\]
\end{remark}

\begin{remark}[Ordinary-route no-go summary]
\label{rem:ordinary-route-no-go-summary}
The following implications are not available:
\[
\text{diagonal mass}\not\Rightarrow\mathrm{O1},
\qquad
\mathrm{O1}+\mathrm{O2}\not\Rightarrow\mathrm{O3},
\]
\[
\text{generic determinant nonvanishing}\not\Rightarrow\mathrm{O3},
\qquad
\text{averaged large sieve}\not\Rightarrow\mathrm{O2},
\]
\[
\text{phase-relaxed model}\not\Rightarrow\text{actual-orbit nonrecurrence}.
\]
Low-prime or proxy blocks are diagnostic only.  Proof-level ordinary Gate--1
requires the full canonical source \(q''_{\can}=B_2+P_2+D_2+E_2\).
\end{remark}

\begin{remark}[Macroscopic smoothed zero-Gram as a proof tool]
\label{rem:macroscopic-smoothed-zero-gram-proof-tool}
Although the pointwise ordinary finite-row route is archived, a macroscopic
smoothed zero-Gram calculation remains useful for proving aggregate
\(L^2\)-source estimates.

Let
\[
  w_T(t)=w\!\left(\frac{t-T}{H}\right),
  \qquad
  H\gg T^{c_++\delta},
  \qquad
  n,n'\le T^{c_+},
  \qquad
  c_+<\tfrac12.
\]
For
\[
  K_H(u):=\sum_\gamma w_T(\gamma)e^{-i\gamma u},
\]
the smoothed explicit formula localizes the prime peaks to
\[
  u=\pm\log k.
\]
If
\[
  u=\log(n/n')
\]
and \(n\ne kn'\), then
\[
  \left|\log\frac{n}{kn'}\right|\gg T^{-c_+}.
\]
Thus
\[
  H\left|\log\frac{n}{kn'}\right|\gg T^\delta,
\]
and rapid decay of \(\widehat w\) kills non-exact ratios.  Hence only
\[
  n=kn'\qquad\text{or}\qquad n'=kn
\]
contribute to the leading zero-Gram.

With the natural \(n^{-1/2}\)-normalization, the Landau term uses
\[
  \sum_k\frac{\Lambda(k)}{k}\delta_{\log k}\approx du.
\]
Thus the leading scalar zero-Gram has rank-one form
\[
  G_Z\approx I-\frac1Q e_Z\otimes e_Z.
\]
For smooth local cells of log-width \(\Delta\asymp\log Q\), the corresponding
cell Gram is block-local:
\[
  G_{\mathrm{cell}}
  \approx
  \bigoplus_a
  \left(I_a-\frac1Qe_a\otimes e_a\right),
\]
and remains positive since
\[
  \|e_a\|^2/Q\asymp\Delta/Q\ll1.
\]

This is a possible proof tool for aggregate \(L^2\)-source sampling.  It is
not the Gate--1 definition, and it must use the full canonical source
\[
  q''_{\can}=B_2+P_2+D_2+E_2
\]
with smooth packets.  Hard cells, \(P_2\)-only packets, or scalar value packets
are not sufficient proof objects.
\end{remark}

\subsection{Finite/O2O3\(_{23}\): analytic Riesz graph and exterior-pivot input}
\label{subsec:finite-o2o3-23}

O2/O3 is now separated from Gate--2.  It should be handled in an analytic
\(R^C\)-conditioned Riesz basis, not in the downstream singular-discard carrier
basis.

Let
\[
  A_{\le X,m}
\]
be the finite augmented matrix built from
\[
  F_X=B_2+P_{2,\le X}+D_{2,\le X}
\]
and the augmented column space
\[
  \mathbf C\cdot1+m_{\can}\mathbf C_R[x].
\]
The analytic normalization requirement is
\[
  \boxed{
  \|A_{\le X,m}\|_{\operatorname{op}}\le R^C.
  }
\]

\begin{definition}[Riesz graph model for O2/O3]
\label{def:finite-o2o3-riesz-graph-model}
The augmented column space has the graph form
\[
  \mathbf C\cdot1+m_{\can}\mathbf C_R[x]
  =
  m_{\can}
  \bigl(\mathbf C\cdot(1/m_{\can})+\mathbf C_R[x]\bigr).
\]
Thus the pure branch obstruction is the non-polynomiality of \(1/m_{\can}\):
\[
  \operatorname{dist}\bigl(1/m_{\can,+},\mathbf C_R[x]\bigr)\ge Q^{-A}
\]
in the \(R^C\)-conditioned Riesz norm.  In the favorable model
\[
  m_{\can,+}(x)=\kappa x^2(1+\varepsilon(x)),
\]
this is an analytic divided-difference/Riesz problem for \(1/m_{\can,+}\).
\end{definition}

\begin{remark}[Graph/O2/O3 route and near-pure Schur block]
\label{rem:graph-o2o3-near-pure-schur-block}
In the \(R^C\)-normalized analytic basis, the pure branch graph model is
\[
  \mathbf C\cdot1+m_{\can}\mathbf C_R[x]
  =
  m_{\can}
  \left(
    \mathbf C\cdot(1/m_{\can})+\mathbf C_R[x]
  \right).
\]
Thus pure-branch O2/O3 reduces to non-polynomiality of \(1/m_{\can}\) in the
finite Riesz norm:
\[
  \operatorname{dist}_{\mathcal R_R}
  \left(1/m_{\can},\mathbf C_R[x]\right)\ge Q^{-A}.
\]
The full determinant problem is not solved by the pure branch alone.  Under a
bounded branch-charge Stack weight, one must also prove survival of the pure
branch inside the full determinant.  Equivalently, after branch-charge
decomposition,
\[
  D_{\mathrm{full}}
  =
  \sum_{r=0}^{n}D_{-n+r},
  \qquad n=R+2,
\]
one needs the finite near-pure Schur complement
\[
  D_{-n}
  \notin_{Q^{-A}}
  \operatorname{span}\{D_{-n+1},\ldots,D_{-n+d_0}\},
\]
where \(d_0\) is the bounded branch-charge degree of the smooth Stack weight.
This statement should be proved before capped extraction; hard extraction
destroys the branch-charge/Fourier structure needed for determinant survival.
\end{remark}

\begin{definition}[Bounded-charge near-pure Schur block]
\label{def:bounded-charge-near-pure-schur-block}
Let \(n=R+2\), and decompose the full determinant into branch-charge sectors
\[
  D_{\mathrm{full}}
  =
  \sum_r D_{-n+r}.
\]
Let \(d_0\) be the branch-charge bandwidth of the bounded-degree Stack weight
used before capped extraction:
\[
  d_0
  :=
  d_N+2d_H+2d_{\Gamma,\mathrm{eff}}.
\]
Here \(d_{\Gamma,\mathrm{eff}}\) is computed from the final effective forbidden
wall linear form after coherent amplitude projection.  Then the pure plus
sector \(D_{-n}\) can interact, under the pre-extraction Stack weight, only
with the finite block
\[
  D_{-n+1},\ldots,D_{-n+d_0}.
\]
The fixed-\(\alpha\) full determinant survival condition is therefore the
finite Schur condition
\[
  \boxed{
  D_{-n}
  \notin_{Q^{-A}}
  \operatorname{span}
  \{D_{-n+1},\ldots,D_{-n+d_0}\}.
  }
\]
This is the default finite/O2/O3 branch-survival theorem unless an
AlphaAdmissibility package is installed.
\end{definition}

\begin{remark}[Do not cap before determinant survival]
\label{rem:do-not-cap-before-determinant-survival}
The finite/O2/O3 determinant-survival theorem must be proved at the smooth
aggregate level.  Capped extraction destroys the branch-charge and Fourier
structure used to prove
\[
  D_{-n}
  \notin_{Q^{-A}}
  \operatorname{span}\{D_{-n+1},\ldots,D_{-n+d_0}\}.
\]
Thus the proof order is
\[
  \boxed{
  \text{smooth aggregate}
  \to
  \text{branch survival/O2/O3}
  \to
  \text{capped extraction}
  \to
  \text{atomic D4 contradiction}.
  }
\]
In particular, the Stack weight used for determinant survival should be a
bounded-degree pre-extraction weight.  Capped factors such as
\(F_{\mathrm{wall}}^{\mathrm{forb}}\) may be used only after
determinant-survival and aggregate lower bounds are secured.
\end{remark}

\begin{hypothesis}[Canonical quadratic multiplier theorem]
\label{hyp:canonical-quadratic-multiplier}
On each ordinary branch subpacket used for finite/O2/O3, the canonical
curvature multiplier has the form
\[
  m_{\can,+}(x)=\kappa (x-z_0)^2u(x),
\]
where
\[
  |\kappa|,\ |\kappa|^{-1}\le Q^C,
  \qquad
  |u(x)|,\ |u(x)|^{-1}\le Q^C,
\]
and \(u\) has \(R^C\)-controlled divided differences through order \(R+1\) in
the active \(R^C\)-conditioned Riesz norm.  Moreover, \(z_0\) is polynomially
close to the active branch scale, so that \(1/m_{\can,+}\) has a controlled
second-order pole visible to the finite polynomial space.

Consequently
\[
  \operatorname{dist}_{\mathcal R_R}
  \left(
    1/m_{\can,+},
    \mathbf C_R[x]
  \right)
  \ge e^{-O(R\log R)}
  \ge Q^{-C}.
\]
\end{hypothesis}

\begin{remark}[Why the quadratic multiplier is the O2/O3 keystone]
\label{rem:canonical-quadratic-multiplier-keystone}
The pure-plus graph theorem
\[
  \operatorname{dist}(1/m_{\can,+},\mathbf C_R[x])\ge Q^{-A}
\]
is the analytic source of pure-plus determinant mass.  In the model
\[
  m_{\can,+}(x)\sim \kappa x^2,
\]
this is the elementary non-polynomiality of \(x^{-2}\), with loss
\(e^{-O(R\log R)}=Q^{-O(1)}\).

The same multiplier also produces the branch-substitution ratio
\[
  r(x)=\frac{(1-x)^2}{x^2}
\]
which enters the bounded-charge near-pure Schur block.  Thus the same finite
theorem supplies both:
\[
  \boxed{\text{pure-plus graph non-polynomiality}}
  \qquad\text{and}\qquad
  \boxed{\text{near-pure substitution-ratio separation}.}
\]
\end{remark}

\begin{definition}[Harmonic exterior O2/O3 quantity]
\label{def:finite-o2o3-harmonic-exterior}
For an \(n\times n\) augmented matrix \(A\), \(n=R+2\), define
\[
  \mathcal H_{23}(A)
  :=
  \frac{\det(A^*A)}{e_{n-1}(A^*A)}.
\]
Equivalently,
\[
  \mathcal H_{23}(A)
  =
  \frac{|\det A|^2}
       {\sum_{i,j}|\operatorname{Cof}_{ij}(A)|^2},
\]
with the value \(0\) when \(A\) is singular.  If
\(\sigma_{\min}(A)^2\) is the least eigenvalue of \(A^*A\), then
\[
  \frac1n\sigma_{\min}(A)^2
  \le
  \mathcal H_{23}(A)
  \le
  \sigma_{\min}(A)^2.
\]
Thus \(\mathcal H_{23}\) is the scale-invariant aggregate substitute for raw
determinant energy.
\end{definition}

\begin{hypothesis}[Weighted exterior-pivot anti-trough for finite/O2O3]
\label{hyp:weighted-exterior-pivot-anti-trough-23}
Let
\[
  d\nu:=F_1F_{\mathrm{wall}}^{\mathrm{forb}}N\,d\rho
\]
be the Stack-C2 weight after dyadic normalization and removal of geometric bad
sets.  For
\[
  D(m):=\det A_{\le X,m},
  \qquad
  C_{ij}(m):=\operatorname{Cof}_{ij}(A_{\le X,m}),
\]
define the dominant-cofactor classes
\[
  \mathcal C_{ij}
  :=
  \left\{
    m:\ |C_{ij}(m)|^2
    \ge
    n^{-2}\sum_{p,q}|C_{pq}(m)|^2
  \right\}.
\]
For every dominant cofactor class,
\[
  \boxed{
  \int_{\mathcal C_{ij}}
  \min\!\left(
    Q^A\left|\frac{D(m)}{C_{ij}(m)}\right|^2,
    1
  \right)\,d\nu(m)
  \ge
  Q^{-B}\nu(\mathcal C_{ij}).
  }
\]
Equivalently, one may prove the logarithmic resultant form
\[
  \int_{\mathcal C_{ij}}
  \log_+
  \frac{|C_{ij}(m)|^2}{|D(m)|^2}\,d\nu(m)
  \le
  C(\log Q)\nu(\mathcal C_{ij}).
\]
\end{hypothesis}

\begin{definition}[Full determinant survival under Stack weight]
\label{def:full-determinant-survival-stack-weight}
Let \(A_{\mathrm{full}}(m)\) denote the full scalar canonical O2/O3 augmented matrix
built from
\[
  q''_{\can}=B_2+P_2+D_2+E_2.
\]
Let
\[
  D_{\mathrm{full}}:=\det A_{\mathrm{full}}.
\]
If a branch-count decomposition is available, write
\[
  D_{\mathrm{full}}=\sum_{\beta}D_\beta.
\]
A phase-aligned branch determinant \(D_{\mathrm{dom}}\) is a proof witness only if it
survives inside the full determinant under the active Stack measure
\[
  d\nu=F_1F_{\mathrm{wall}}^{\mathrm{forb}}N\,d\rho:
\]
\[
  \boxed{
  \int |D_{\mathrm{full}}|^2\,d\nu
  \ge
  Q^{-A}
  \int |D_{\mathrm{dom}}|^2\,d\nu.
  }
\]
A lower bound for \(D_{\mathrm{dom}}\) alone is a projected branch certificate and does
not prove scalar O2/O3 for the canonical source.
\end{definition}

\begin{remark}[Cauchy--Binet determinant coefficient and divided differences]
\label{rem:cauchy-binet-divided-difference-o2o3}
For a single branch with frequencies \(x_n=\lambda\log n\), write the augmented
columns as
\[
  h_0(x)=1,\qquad h_{j+1}(x)=m_{\can}(x)x^j,\quad 0\le j\le R.
\]
If
\[
  A_m=(L_m(h_jx^k))_{0\le k,j\le R+1},
\]
then Cauchy--Binet gives
\[
  \det A_m
  =
  \sum_{\substack{S\\ |S|=R+2}}
  C_S e^{-im\log N_S},
  \qquad
  N_S:=\prod_{n\in S}n.
\]
For a pure branch one has
\[
  C_S
  =
  \left(\prod_{n\in S}a_n\right)
  (\det V_S)^2
  \left(\prod_{n\in S}m_{\can}(x_n)\right)
  [x_n:n\in S]\left(\frac1{m_{\can}}\right),
\]
where \(V_S\) is the Vandermonde matrix of the nodes \(x_n\).  Thus the
analytic O3 content of the exterior determinant is the divided difference
\[
  [x_n:n\in S]\left(\frac1{m_{\can}}\right).
\]

In the favorable branch-normalized elliptic model
\[
  m_{\can,+}(x)=\kappa x^2(1+\varepsilon(x)),
\]
one compares \(1/m_{\can,+}\) to \((\kappa x^2)^{-1}\).  Since
\[
  \left(x^{-2}\right)^{(R+1)}
  =
  (-1)^{R+1}(R+2)!x^{-R-3},
\]
the model divided differences have fixed phase and non-small mass on \(x>0\).
Therefore the analytic multiplier obstruction reduces to the finite audit that
\[
  \left(
  \frac1{m_{\can,+}}
  -
  \frac1{\kappa x^2}
  \right)^{(R+1)}
\]
is perturbative relative to
\[
  \left(\frac1{\kappa x^2}\right)^{(R+1)}
\]
on the chosen branch-routed subpacket.
\end{remark}

\begin{remark}[Collision and branch-survival guardrail for determinant methods]
\label{rem:collision-branch-survival-determinant-methods}
A pure-branch determinant lower bound is not, by itself, a scalar finite/O2/O3
theorem for the full canonical matrix.  The determinant expansion has product
frequencies
\[
  \log N_S=\sum_{n\in S}\log n,
\]
with \(|S|=R+2\), and therefore effective length \(e^{O(RQ)}\).  Exact and
near multiplicative collisions can occur among the products \(N_S\).  Moreover,
the full determinant also contains pure dual, mixed, background, and completion
exterior terms.

Thus any determinant route must prove both:
\[
  \boxed{\text{collision/near-collision noncancellation for the relevant
  exterior coefficients}}
\]
and
\[
  \boxed{\text{survival of the branch determinant inside the full scalar
  canonical determinant}.}
\]
Prime-labelled or otherwise collision-safe subpackets are useful only as
witnesses inside a sector-compatible component of the full determinant.  A
projected prime determinant alone certifies a projected map, not the scalar
canonical augmented matrix.
\end{remark}

\begin{lemma}[Riesz-normalized determinant controls \(\mathcal H_{23}\)]
\label{lem:riesz-normalized-determinant-controls-H23}
Let \(A\) be an \(n\times n\) matrix with \(n=R+2\), written in an
\(R^C\)-conditioned domain/codomain basis so that
\[
  \|A\|\le R^C.
\]
Then
\[
  \mathcal H_{23}(A)
  \ge
  Q^{-C}|\det A|^2.
\]
\end{lemma}

\begin{proof}
If \(s_1,\ldots,s_n\) are the singular values of \(A\), then
\[
  e_{n-1}(A^*A)
  =
  \sum_{i=1}^n\prod_{j\ne i}s_j^2
  \le
  n\|A\|^{2(n-1)}.
\]
Thus
\[
  \mathcal H_{23}(A)
  \ge
  \frac{|\det A|^2}{n\|A\|^{2(n-1)}}.
\]
Since \(n=R+2\), \(\|A\|\le R^C\), and
\[
  R^{O(R)}=Q^{O(1)},
\]
the claimed polynomial loss follows.
\end{proof}

\begin{theorem}[Thickened overlap Schur criterion]
\label{thm:thickened-o2o3-schur}
Let
\[
K:=G_{00}^{-1/2}G_{0S}G_{SS}^{-1/2}
\]
after Fej\'er--Riesz thickening of the Stack weight.  Then
\[
G_{00}-G_{0S}G_{SS}^{\dagger}G_{S0}\ge Q^{-A}G_{00}
\]
if and only if
\[
\|K\|_{\operatorname{op}}\le 1-Q^{-A}.
\]
A sufficient finite criterion is
\[
\Bigl(\max_i\sum_j |K_{ij}|\Bigr)\Bigl(\max_j\sum_i |K_{ij}|\Bigr)<(1-Q^{-A})^2.
\]
Equivalently, if \(K\) is viewed as a weighted bipartite overlap graph, then a
coherence bound \(d\mu<1-Q^{-A}\) suffices.
\end{theorem}

\begin{lemma}[Packetwise thickened overlap decay]
\label{lem:packetwise-thickened-overlap-decay}
Let \(\{u_i\}\) be the thickened pure packet family and \(\{v_j\}\) the thickened
contaminant packet family after Fej\'er--Riesz thickening.  Define
\[
K_{ij}:=\big\langle G_{00}^{-1/2}u_i,\;G_{SS}^{-1/2}v_j\big\rangle.
\]
Assume the packet centers satisfy the oscillatory-kernel separation hypothesis
from the source-side packet lemma, with decay profile \(\mu(\Delta)\).  Then
\[
|K_{ij}|\le \mu(\Delta_{ij}),
\]
and in particular if the interacting packet pairs satisfy
\(\Delta_{ij}\ge \Delta_0\), then
\[
|K_{ij}|\le \mu_0:=\mu(\Delta_0).
\]
\end{lemma}

\begin{proof}
This is the packetwise form of the oscillatory-kernel estimate already used in
the source layer.  After Fej\'er--Riesz thickening, the normalized overlap
entry is exactly the inner product between the normalized pure and contaminant
packets:
\[
K_{ij}=\langle \tilde u_i,\tilde v_j\rangle,
\qquad
\tilde u_i:=G_{00}^{-1/2}u_i,
\qquad
\tilde v_j:=G_{SS}^{-1/2}v_j.
\]
Writing the packets in oscillatory form
\(u_i=a_i e^{i\omega_i t}\), \(v_j=b_j e^{i\nu_j t}\), the overlap is an
oscillatory integral.  The oscillatory-kernel hypothesis gives the stated decay
in packet separation \(\Delta_{ij}\).  That is exactly the bound claimed.
\end{proof}

\begin{lemma}[Row/column coherence bound]
\label{lem:row-column-coherence-bound}
Let \(K=(K_{ij})\) be the normalized overlap matrix in
Lemma~\ref{lem:packetwise-thickened-overlap-decay}.  Define
\[
r_i:=\sum_j |K_{ij}|,
\qquad
c_j:=\sum_i |K_{ij}|.
\]
If
\[
\Bigl(\max_i r_i\Bigr)\Bigl(\max_j c_j\Bigr)<(1-Q^{-A})^2,
\]
then
\[
\|K\|_{\operatorname{op}}<1-Q^{-A}.
\]
Equivalently, if each row has at most \(d_0\) nonnegligible entries and each
column at most \(d_1\), with \(|K_{ij}|\le \mu\), then
\[
d_0d_1\mu^2<(1-Q^{-A})^2
\quad\Longrightarrow\quad
\|K\|_{\operatorname{op}}<1-Q^{-A}.
\]
\end{lemma}

\begin{proof}
This is a standard Schur test estimate.  For any vector \(x\),
\[
|(Kx)_i|\le \sum_j |K_{ij}|\,|x_j|.
\]
By Cauchy--Schwarz and the row/column sum bounds,
\[
\|Kx\|_2^2\le \Bigl(\max_i r_i\Bigr)\Bigl(\max_j c_j\Bigr)\|x\|_2^2.
\]
Therefore
\[
\|K\|_{\operatorname{op}}^2\le \Bigl(\max_i r_i\Bigr)\Bigl(\max_j c_j\Bigr),
\]
and the stated implication follows.  The bounded-degree form is the special
case \(r_i\le d_0\mu\), \(c_j\le d_1\mu\).
\end{proof}

\begin{lemma}[Schur complement / principal-angle equivalence]
\label{lem:schur-principal-angle-equivalence}
With
\[
K:=G_{00}^{-1/2}G_{0S}G_{SS}^{-1/2},
\]
the following are equivalent:
\begin{enumerate}
\item \(\|K\|_{\operatorname{op}}<1-Q^{-A}\);
\item the principal angle between the thickened pure and contaminant subspaces
satisfies
\[
\sin\theta_{\min}(U,V)\ge Q^{-A};
\]
\item the Schur complement satisfies
\[
G_{00}-G_{0S}G_{SS}^{\dagger}G_{0S}^\ast\ge Q^{-A}G_{00}.
\]
\end{enumerate}
In particular, any of the bounds in
Lemma~\ref{lem:row-column-coherence-bound} closes the
\(\mathrm{Graph/O2O3}_{23}\) gap.
\end{lemma}

\begin{proof}
By definition, \(K\) is the normalized cross-Gram block between the pure and
contaminant subspaces.  Hence \(\|K\|_{\operatorname{op}}\) is the cosine of
the smallest principal angle:
\[
\|K\|_{\operatorname{op}}=\cos\theta_{\min}(U,V).
\]
Therefore \(\|K\|_{\operatorname{op}}<1-Q^{-A}\) is equivalent to
\[
\sin\theta_{\min}(U,V)\ge Q^{-A},
\]
up to the standard trigonometric identity.

For the Schur complement, write the block Gram matrix
\[
\begin{pmatrix}
G_{00} & G_{0S}\\
G_{0S}^\ast & G_{SS}
\end{pmatrix}.
\]
After conjugating by \(G_{00}^{-1/2}\oplus G_{SS}^{-1/2}\), the positivity of
the block matrix is equivalent to
\[
I-KK^\ast \succ 0.
\]
A quantitative lower bound \(I-KK^\ast\ge Q^{-A}I\) is equivalent to
\[
\|K\|_{\operatorname{op}}<1-Q^{-A},
\]
which is equivalent to the Schur complement lower bound
\[
G_{00}-G_{0S}G_{SS}^{\dagger}G_{0S}^\ast\ge Q^{-A}G_{00}.
\]
\end{proof}

\begin{corollary}[Exact target for \(\mathrm{Graph/O2O3}_{23}\)]
\label{cor:graph-o2o3-target}
To close \(\mathrm{Graph/O2O3}_{23}\), it is enough to prove one of the following:
\[
\Bigl(\max_i\sum_j |K_{ij}|\Bigr)\Bigl(\max_j\sum_i |K_{ij}|\Bigr)<(1-Q^{-A})^2,
\]
or
\[
d_0d_1\mu^2<(1-Q^{-A})^2,
\]
or
\[
G_{00}-G_{0S}G_{SS}^{\dagger}G_{0S}^\ast\ge Q^{-A}G_{00}.
\]
If none of these hold, the precise failure label is:
\[
\text{thickened overlap too coherent}.
\]
\end{corollary}

\begin{remark}[wip-zeta-bridge-proof-o2o3]
\label{rem:wip-zeta-bridge-proof-o2o3}
What remains is the thickened Schur / overlap test:
\[
\|K\|_{\operatorname{op}}<1-Q^{-A}.
\]
Equivalently, the Stack-thickened pure/contaminant overlap must be
spectrally subcritical.
\end{remark}

\begin{remark}[O2/O3 failure modes]
\label{rem:o2o3-thickened-overlap-failure}
The only local failure mode at this stage is thickened overlap too coherent.
The upstream prerequisites are \(R^C\)-normalization and
genuine-pole/divided-difference visibility; if either fails, the overlap matrix
is not in the correct finite regime.
\end{remark}

\begin{remark}[What finite/O2O3 does and does not use]
\label{rem:finite-o2o3-basis-and-no-go-register}
Finite/O2O3 uses the analytic Riesz basis for \(A_{\le X,m}\).  It does not use
the singular retained projection \(P_{\mathrm{sing}}\) except through the weight
\(F_{\mathrm{wall}}^{\mathrm{forb}}\).  Conversely, the singular-discard carrier projection is
not required to make \(A_{\le X,m}\) functorial.

The following routes are no-gos as standalone proofs of finite/O2O3:
\begin{enumerate}
\item fixed-vector aggregate frames, since moving null directions may make every
local \(A_{\le X,m}\) singular;
\item raw determinant energy without \(R^C\)-normalization, since the
determinant-to-\(\sigma_{\min}\) loss is super-polynomial;
\item generic Jensen/resultant/type arguments, since \(\det A_{\le X,m}\) has
effective type \(QR\) while the zero density is only \(Q\);
\item pure branch determinant lower bounds without full determinant survival;
\item the \(\Lambda\)-wall factor alone, since it does not force
\(\mathcal H_{23}(A_{\le X,m})\) to be non-small.
\end{enumerate}
The actual remaining theorem is weighted exterior-pivot anti-trough for the
full finite canonical matrix, not merely for a pure branch determinant.
\end{remark}

\begin{definition}[\(\alpha\)-admissibility for determinant survival]
\label{def:alpha-admissibility}
An \(\alpha\)-twisted determinant survival argument is admissible only if the
following package holds.

\begin{enumerate}
\item \textbf{Zero persistence.}
Every fixed-target bad zero \(Z_0(m_0)=0\) in the regular stratum extends to a
zero curve
\[
  Z_\alpha(m(\alpha))=0
\]
on an interval \(I_\alpha\) of polynomial length.  Quantitatively,
\[
  |\partial_mZ_\alpha(m(\alpha))|\ge Q^{-C},
  \qquad
  \frac{dm}{d\alpha}
  =
  -\frac{\partial_\alpha Z_\alpha}{\partial_m Z_\alpha}.
\]

\item \textbf{Badness persistence.}
The Stack quantities
\[
  E_1,\qquad E_{23},\qquad E_{\mathrm{wall}}^{\mathrm{forb}},\qquad
  D_{\mathrm{full}},\qquad D_{\mathrm{dom}},\qquad N
\]
have polynomial log-Lipschitz variation along the zero curve, away from capped
sublevel boundaries.

\item \textbf{Twist nondegeneracy.}
If
\[
  D_{\mathrm{full}}(m,\alpha)
  =
  \sum_{r=0}^{R+2}e^{ir\Theta(m,\alpha)}D_r(m,\alpha),
\]
then along the zero curve
\[
  \boxed{
  \left|
  \frac{d}{d\alpha}\Theta(m(\alpha),\alpha)
  \right|
  =
  \left|
  \partial_\alpha\Theta
  -
  \partial_m\Theta
  \frac{\partial_\alpha Z_\alpha}{\partial_m Z_\alpha}
  \right|
  \ge Q^{1-C}.
  }
\]

\item \textbf{Stack uniformity.}
Gate--2, O2/O3, D4, and the forbidden functional \(\Lambda\) are defined with
uniform polynomial constants on the \(\alpha\)-tube.
\end{enumerate}

Without these four properties, \(\alpha\)-averaging is heuristic and cannot be
used to prove the fixed \(\alpha=0\) target.
\end{definition}

\begin{remark}[Regular-zero branch and near-multiple branch]
\label{rem:alpha-regular-versus-near-multiple-branch}
The \(\alpha\)-twist route applies only on the regular zero stratum
\[
  |\partial_mZ_\alpha|\ge Q^{-C}.
\]
The complementary branch
\[
  |\partial_mZ_\alpha|<Q^{-C}
\]
must not be ignored.  It should be routed to a separate near-multiple-zero
module.  Such zeros may already seed Gate--1, since
\[
  Z_\alpha(m)\approx0,\qquad
  \partial_mZ_\alpha(m)\approx0
\]
forces local curvature unless the canonical source is flat to higher order.
The required auxiliary theorem is
\[
  \boxed{
  \mathrm{NearMultipleGate1}:
  \quad
  |Z_\alpha|+|\partial_mZ_\alpha|\le Q^{-A}
  \Longrightarrow
  \int_I |\lambda^3q''_{\can}(m+\lambda x)|^2\,dx\ge Q^{-A'}.
  }
\]
If this theorem fails, near-multiple zeros require a separate singular/flatness
branch.
\end{remark}

\begin{definition}[\(\alpha\)-D4 compatibility]
\label{def:alpha-d4-compatibility}
Let \(\mathcal K_9\) be the covariant order-nine singular coherence residual,
with residual codomain \(\mathcal Y_9\).  Let
\[
  \mathcal T_{\mathrm{par}}^{(\alpha)}
  :=
  \operatorname{span}\{D_p\mathcal K_9(\partial_{p_a})\},
\]
where \(p_a\) ranges over the allowed actual-source parameters, including
\(\alpha\) if \(\alpha\) is admitted.  Let \(\mathcal T_7\) and
\(\mathcal T_{<7}\) be the order-seven and lower-order connection residual
subspaces.  Define
\[
  \boxed{
  \mathcal Q_9^{\mathrm{forb},\alpha}
  :=
  \mathcal Y_9/
  \bigl(
  \mathcal T_{\mathrm{par}}^{(\alpha)}
  +
  \mathcal T_7
  +
  \mathcal T_{<7}
  \bigr).
  }
\]
The \(\alpha\)-twist is compatible with the \(Y12\) D4 channel if
\[
  \boxed{
  [\mathsf C_{Y12}]\ne0
  \quad\text{in}\quad
  \mathcal Q_9^{\mathrm{forb},\alpha}.
  }
\]
Equivalently,
\[
  T_{\mathrm{tr}\to\mathrm{sing}}[\mathsf C_{Y12}]
  \notin
  \mathcal T_{\mathrm{par}}^{(\alpha)}+\mathcal T_7+\mathcal T_{<7}.
\]
\end{definition}

\begin{definition}[\(Y12\)-signature quotient]
\label{def:Y12-signature-quotient}
Let \(\Sigma(Y12)\) denote the primitive order-nine mixed-block signature of
the carrier
\[
  Y_{12}^{(2)}=D^2N_0[\xi_1,\xi_2]
\]
coupled to one same-point whitening slot.  Let
\[
  \Pi_{\Sigma(Y12)}:\mathcal Y_9\to \mathcal Y_{\Sigma(Y12)}
\]
be the corresponding signature projection in the residual codomain of the
covariant order-nine coherence residual \(\mathcal K_9\).

Define the \(Y12\)-signature forbidden quotient by
\[
  \boxed{
  \mathcal Q_{9,Y12}^{\mathrm{forb},\alpha}
  :=
  \mathcal Y_{\Sigma(Y12)}
  \Big/
  \Pi_{\Sigma(Y12)}
  \bigl(
    \mathcal T_{\mathrm{par}}^{(\alpha)}
    +\mathcal T_7
    +\mathcal T_{<7}
    +\mathcal{II}_{\mathrm{arith},9}
  \bigr).
  }
\]
Here \(\mathcal{II}_{\mathrm{arith},9}\) denotes the second-order actual-source
curvature/tangent image at order nine.  The \(Y12\)-channel survives iff
\[
  \boxed{
  \Pi_{\Sigma(Y12)}T_{\mathrm{tr}\to\mathrm{sing}}[\mathsf C_{Y12}]
  \ne0
  \quad\text{in}\quad
  \mathcal Q_{9,Y12}^{\mathrm{forb},\alpha}.
  }
\]
\end{definition}

\begin{hypothesis}[Second-order \(Y12\)-signature separation]
\label{hyp:second-order-Y12-signature-separation}
The actual-source tangent, lower-connection, and actual-source second-curvature
images have no primitive \(Y12\)-signature:
\[
  \boxed{
  \Pi_{\Sigma(Y12)}
  \bigl(
    \mathcal T_{\mathrm{par}}^{(\alpha)}
    +\mathcal T_7
    +\mathcal T_{<7}
    +\mathcal{II}_{\mathrm{arith},9}
  \bigr)=0.
  }
\]
Moreover the transported D4 carrier has nonzero \(Y12\)-signature:
\[
  \boxed{
  \Pi_{\Sigma(Y12)}T_{\mathrm{tr}\to\mathrm{sing}}[\mathsf C_{Y12}]\ne0.
  }
\]
If \(\alpha\) is admitted as an actual parameter, it is required to be
source-linear in the sense that
\[
  \Pi_{\Sigma(Y12)}\partial_\alpha u=0,
  \qquad
  \Pi_{\Sigma(Y12)}\partial_\alpha^2 u=0,
\]
unless a separate source-nonlinear \(\alpha\)-audit is supplied.
\end{hypothesis}

\begin{remark}[Why the \(Y12\)-signature test is sharper than full \(K_9\)]
\label{rem:Y12-signature-sharper-than-full-K9}
The full \(K_9\)-upper may be stronger than needed.  The D4 channel currently
used by the rank-two route is the primitive \(Y12\)-signature component of the
order-nine residual.  Therefore it is enough to control the \(Y12\)-signature
projection of actual EF residuals and to prove that the D4 carrier survives in
\(\mathcal Q_{9,Y12}^{\mathrm{forb},\alpha}\).  This avoids proving unnecessary
completion estimates for unrelated order-nine residual sectors.
\end{remark}

\begin{remark}[\(\alpha\) can help O2/O3 while hurting D4]
\label{rem:alpha-o2o3-d4-tradeoff}
The \(\alpha\)-twist is a possible determinant-survival tool because it may
orthogonalize folded-dual branch counts.  However, if \(\alpha\) is admitted as
an actual-source parameter, the actual singular tangent space expands by
\(\partial_\alpha u\).  The D4 contradiction survives this enlargement only if
\[
  [\mathsf C_{Y12}]\ne0
  \quad\text{in}\quad
  \mathcal Q_9^{\mathrm{forb},\alpha}.
\]
Thus the proof-critical admissibility package is
\[
  \boxed{
  \mathrm{AlphaAdmissibility}^{+}
  :=
  \mathrm{AlphaAdmissibility}
  +
  \mathrm{AlphaD4Compatibility}.
  }
\]
There are two logically distinct uses of \(\alpha\).

First, \(\alpha\) may be used as a proof-thickening or averaging parameter for
determinant survival.  In that use, one proves persistence back to the
fixed-target problem at \(\alpha=0\), and \(\partial_\alpha u\) need not be
added to the downstream actual-source tangent space.

Second, \(\alpha\) may be admitted as an actual source deformation.  In that
case \(\partial_\alpha u\) belongs to the tangent space in
\(\mathcal T_{\mathrm{par}}^{(\alpha)}\), and one must verify
\[
  [\mathsf C_{Y12}]\ne0
  \quad\text{in}\quad
  \mathcal Q_9^{\mathrm{forb},\alpha}.
\]
Therefore the proof should not automatically insert \(\partial_\alpha u\) into
the downstream tangent space merely because \(\alpha\) is used as an auxiliary
averaging parameter.

If \(\alpha\)-admission kills \([\mathsf C_{Y12}]\), then \(\alpha\) may still
serve as an auxiliary averaging parameter, but it must not be added to the
downstream actual-source manifold without a replacement forbidden channel.
\end{remark}

\begin{remark}[Two basis layers: analytic Riesz versus carrier projection]
\label{rem:two-basis-layers-o2o3-lambda}
The analytic basis used for finite/O2/O3 and the carrier projection used for
\(\Lambda\) have different purposes and should not be identified.

The finite/O2/O3 basis should be an \(R^C\)-conditioned analytic Riesz basis in
which the augmented matrix satisfies
\[
  \|A_{\le X,m}\|\le R^C.
\]
This is what makes
\[
  \mathcal H_{23}(A_{\le X,m})
  \ge Q^{-C}|\det A_{\le X,m}|^2
\]
a polynomial-loss implication.

The forbidden-channel projection \(P_{\mathrm{sing}}\), in contrast, should
encode actual EF-coherent singular source retention.  It must be non-functorial
on at least one formal D4 carrier:
\[
  (I-P_{\mathrm{sing}})T_{\mathrm{tr}\to\mathrm{sing}}[\mathsf C_7]\ne0
  \quad\text{or}\quad
  (I-P_{\mathrm{sing}})T_{\mathrm{tr}\to\mathrm{sing}}[\mathsf C_{Y12}]\ne0.
\]
Using one fully functorial basis for both jobs risks making O2/O3 well
conditioned while killing the forbidden D4 channel.
\end{remark}

\subsection{Singular and transition proof tools after scalar \(L^2\) Gate--1}
\label{subsec:singular-transition-gate-one-branches}

The regimes
\[
\mu\gtrsim1,
\qquad
(\log Q)^{-C_{\mathrm{tr}}}\lesssim\mu\lesssim1,
\qquad
\mu\ll(\log Q)^{-C_{\mathrm{tr}}},
\qquad
\mu:=\lambda Q,
\]
still govern the construction of packet windows, chart coordinates, proof
tools, and downstream carrier transitions.  They no longer define three
different Gate--1 objects.  All branches feed the same scalar downstream input:
\[
  \int_{I_b}|\lambda_b^3q''_{\can}(m+\lambda_bx)|^2\,dx\ge Q^{-A}.
\]

\begin{definition}[Deprecated pointwise S1 diagnostic]
\label{def:S1-singular-explicit-formula-anti-design}
The old S1 vector
\[
  \mathbf P_R-\mathbf G_R-\mathbf T_R
\]
is a singular diagnostic/proof tool only.  It is not a proof-level Gate--1
definition unless one proves that it gives
\[
  \int_{I_b}|\lambda_b^3q''_{\can}(m+\lambda_bx)|^2\,dx\ge Q^{-A}.
\]
Without that transfer, S1 is only a proxy anti-design statement.
\end{definition}

\begin{definition}[Deprecated pointwise T1 diagnostic]
\label{def:T1-weighted-transition-source-observability}
The weighted transition jet
\[
  \mathcal J_R^{\mathrm{tr}}(m)
  :=
  \left(
  \mu^{-r}\partial_x^r[
  \lambda^3q''_{\can}(m+\lambda x)]_{x=0}
  \right)_{0\le r\le R},
  \qquad
  \mu=\lambda Q,
\]
is a transition proof tool.  If
\[
  \|\mathcal J_R^{\mathrm{tr}}(m)\|\ge Q^{-A},
\]
then a weighted finite-jet-to-energy lemma may yield local scalar Gate--1.
However, the proof-level Gate--1 quantity remains
\[
  E_{1,\mathrm{tr}}^{L^2}(m)
  =
  \int_{I_{\mathrm{tr}}}|\lambda^3q''_{\can}(m+\lambda x)|^2\,dx.
\]
\end{definition}

\begin{lemma}[Weighted transition chart loss]
\label{lem:weighted-transition-chart-loss}
Let
\[
F_\lambda(x)=\sum_{n\in\mathcal B} a_n e^{i\lambda(\log n)x}
+F_{\rm bg}(x)+F_{\rm tail}(x),
\]
with \(c_-Q\le\log n\le c_+Q\).  Set
\[
\tau_n:=\lambda\log n,
\qquad
\mu:=\lambda Q.
\]
Then \(\tau_n/\mu\in[c_-,c_+]\), and
\[
\mu^{-r}\partial_x^r e^{i\tau_nx}\big|_{x=0}
=
i^r(\tau_n/\mu)^r.
\]
Thus the weighted transition jet uses a fixed-window moment vector in the
bounded variable \(\tau_n/\mu\).  Passing from weighted to unweighted
derivatives costs at most \(\mu^{-R}\) at depth \(R\), rather than the
Vandermonde determinant loss \(\mu^{-O(R^2)}\).  Since
\[
\mu\ge(\log Q)^{-C_{\mathrm{tr}}},
\qquad
R\asymp\frac{\log Q}{\log\log Q},
\]
the loss is polynomial in \(Q\).
\end{lemma}

\begin{remark}[Finite transition jets are not complete source frames]
\label{rem:transition-point-jets-not-complete-frame}
The weighted transition jet
\[
  \left(\mu^{-r}\partial_x^rF_\lambda(0)\right)_{0\le r\le R}
\]
records only \(R+1\) bounded-window moments of the local coefficient measure.
For arbitrary transition packets, one may choose nonzero coefficients with
\[
  \sum_n a_n(\tau_n/\mu)^r=0,
  \qquad 0\le r\le R,
\]
while
\[
  \|F_\lambda\|_{L^2}>0.
\]
Thus weighted point jets alone are not a deterministic lower frame for local
\(L^2\)-source energy.  They may be used as one coordinate block of a larger
transition projection family, or as a diagnostic, but not as a standalone
deterministic Gate--1 lower frame.
\end{remark}

\begin{remark}[Regime routing after scalar \(L^2\) normalization]
\label{rem:gate-one-regime-routing}
The ordinary, transition, and singular regimes now determine how to prove
local scalar \(L^2\)-energy and how to construct chart/carrying maps.  The
Gate--1 object itself is always
\[
  E_{1,b}^{L^2}(m)
  =
  \int_{I_b}|\lambda_b^3q''_{\can}(m+\lambda_bx)|^2\,dx.
\]
\end{remark}

\begin{theorem}[Projected source-frame branch certificate]
\label{thm:projected-source-frame-branch-certificate}
Let
\[
  F_m(x):=\lambda_m^3q''_{\can}(m+\lambda_m x).
\]
Assume there are nonnegative smooth branch weights
\[
  \chi_{\mathrm{ord}}(m)+\chi_{\mathrm{tr}}(m)+\chi_{\mathrm{sing}}(m)=1
\]
and branch observables
\[
  E_{\mathrm{ord}}(m),\qquad E_{\mathrm{tr}}(m),\qquad E_{\mathrm{sing}}(m)
\]
with
\[
  E_{\mathrm{ord}}(m)=\|T_{\can}X_{\can}(m)\|^2
\]
in the ordinary packet-frequency chart, while
\[
  E_{\mathrm{tr}}(m)=\sum_j\|P_{\mathrm{tr},j}F_m\|_{L^2}^2,
  \qquad
  E_{\mathrm{sing}}(m)=\sum_j\|P_{\mathrm{sing},j}F_m\|_{L^2}^2.
\]
If the multichart source-frame inequality
\[
  \chi_{\mathrm{ord}}(m)E_{\mathrm{ord}}(m)
  +
  \chi_{\mathrm{tr}}(m)E_{\mathrm{tr}}(m)
  +
  \chi_{\mathrm{sing}}(m)E_{\mathrm{sing}}(m)
  \ge
  Q^{-A}\|F_m\|_{L^2}^2
\]
holds, then the branch observables certify the same scalar canonical source
norm used by \(\mathrm{Agg}L^2(q''_{\can})\).

In particular, transition point jets and singular sine moments are not
standalone Gate--1 definitions.  They are valid proof tools only after being
embedded in projection families satisfying the lower-frame inequality above.
\end{theorem}

\begin{remark}[Point jets and sine diagnostics are not complete frames]
\label{rem:transition-singular-diagnostics-not-complete-frames}
The transition point jet
\[
  \left(\mu^{-r}\partial_x^rF_m(0)\right)_{0\le r\le R}
\]
records only \(R+1\) bounded-window moments.  It is not, by itself, a lower
frame for \(\|F_m\|_{L^2}\).  Similarly, a raw singular sine-moment vector is
not a proof-level Gate--1 input unless it is proved to be an actual coefficient
or projection frame for \(F_m\).

Thus the active transition and singular proof objects are projected source
energies:
\[
  \sum_j\|P_{\mathrm{tr},j}F_m\|_{L^2}^2,
  \qquad
  \sum_j\|P_{\mathrm{sing},j}F_m\|_{L^2}^2.
\]
\end{remark}

\subsection{Ordinary Wronskian no-go archive}
\label{subsec:no-go-archive-ordinary-wronskian}

\begin{remark}[Ordinary route archive]
\label{rem:ordinary-route-archive}
The ordinary Wronskian and finite-row O1/O2/O3 discussions are now archived.
They are proof tools or historical alternate routes only, unless they are
explicitly transferred to the scalar local \(L^2\) source energy of
\(q''_{\can}=B_2+P_2+D_2+E_2\).
The active upstream statement is \(\mathrm{Agg}L^2(q''_{\can})\), not the old
finite-row route.
\end{remark}


\subsection{Canonical audit for the named Gate--1 inputs}
\label{subsec:canonical-audit-named-gate-one-inputs}

\begin{remark}[Canonical audit for named Gate--1 inputs]
\label{rem:canonical-audit-named-gate-one-inputs}
Every named Gate--1 input must be stated for the same full canonical source
\[
q''_{\can}=B_2+P_2+D_2+E_2.
\]
The following substitutions are not allowed:
\[
B_2\text{-only},\qquad
P_2\text{-only},\qquad
\text{zero-kernel-only},\qquad
\text{flexible Dirichlet-polynomial proxies}.
\]
The finalized gamma/background curvature \(B_2\) is part of the canonical
source.  Diagnostics using \(B_2\) alone are not Gate--1 candidates.

For the downstream rank-two analytic route, if \(B_2\) is included in
\[
p=q''_{\can},
\]
then the wall-completion module must include or approximate
\[
B_2,\qquad DB_2,\qquad D^2B_2
\]
in the scale-normalized \(H^2\)-ledger.
\end{remark}

\section{Conditional status after the downstream closure repairs}
\label{sec:final-status-after-downstream-closure-repairs}

\begin{remark}[Current architecture ledger]
\label{rem:current-architecture-ledger-unified}
The current proof architecture is
\[
\mathrm{Agg}L^2(q''_{\can})+\mathrm{Finite/O2O3}_{23}+\mathrm{ForbiddenSingularDiscardD4}.
\]
The unresolved inputs are:
\begin{enumerate}
\item aggregate \(L^2\)-observability for \(q''_{\can}\),
\item finite/O2O3 exterior-pivot anti-trough,
\item construction of \(R_{\mathrm{sing}}\) and the singular-discard upper,
\item D4 transversality to the discarded singular channel.
\end{enumerate}
Any sentence implying that finite D4 alone closes the proof should be replaced
by ``visibility, not contradiction unless \(\Lambda\).''
\end{remark}

\begin{remark}[wip-zeta-bridge-proof-atomic]
\label{rem:wip-zeta-bridge-proof-atomic}
The last step is the same-norm atomic comparison:
the forbidden upper must be Type-I local-square in the same quotient norm
as the surviving lower bound.
\end{remark}

\begin{remark}[No-go registry]
\label{rem:no-go-registry-unified}
The dead routes are:
\begin{itemize}
\item scalar \((I)\)-coordinate D4;
\item derivative D4 wall;
\item optimized \(\Xi_N\) as a finite D4 forbidden channel;
\item one-parameter row mobility as a contradiction mechanism;
\item functorial carrier lifts when \(R_{\mathrm{sing}}\) is Taylor-complete or carrier-complete;
\item low-dimensional pointwise finite-row \(O1/S1/T1\) as the active Gate--1 definition.
\end{itemize}
\end{remark}

\begin{remark}[Active input status table]
\label{rem:active-input-status-table}
The active inputs are now as follows.
\[
\begin{array}{c|c|c}
\text{module} & \text{current status} & \text{main obstruction}\\
\hline
\mathrm{Agg}L^2(q''_{\can})
& \text{open}
& \text{finite canonical }L^2\text{ anti-lock and branch coverage}\\
\mathrm{EFCompletion}_{L^2}
& \text{finite-audit}
& \text{localized }R_X\text{ bounds in }W_{\mathrm{branch}},W_{\mathrm{ref}}\\
\mathrm{EFCompletion}_{K9}
& \text{strong/open}
& \text{full order-nine residual after }\mathcal K_9 T_{\mathrm{tr}\to\mathrm{sing}}J_{\mathrm{tr}}\\
\mathrm{EFCompletion}_{Y12}
& \text{targeted/open}
& \text{residual in primitive }Y12\text{-signature sector}\\
\mathrm{LowerOrderTangentObservability}
& \text{finite-audit}
& \text{conditioning of the covariant }K_9\text{ row}\\
\mathrm{SecondOrderY12SignatureSeparation}
& \text{finite-audit/open}
& \text{actual tangents and second curvature must have no }Y12\text{ signature}\\
\mathrm{CanonicalQuadraticMultiplier}
& \text{finite-audit/open}
& m_{\can,+}=\kappa(x-z_0)^2u(x)
  \text{ with divided-difference control}\\
\mathrm{NearPureSchur}_{23}
& \text{open}
& D_{-n}\notin
  \operatorname{span}\{D_{-n+1},\ldots,D_{-n+d_0}\}\\
\mathrm{FullDetSurvival}
& \text{open}
& \text{branch-count cancellation under Stack weight}\\
\mathrm{RieszNorm}_{23}
& \text{finite-audit}
& \|A_{\le X,m}\|_{\operatorname{op}}\le R^C\\
\mathrm{AlphaAdmissibility}^{+}
& \text{conditional route}
& \text{zero persistence, twist speed, Stack uniformity, D4 compatibility}\\
\mathrm{NormalD4}_{K9,Y12}
& \text{conditional target}
& Y12\text{-signature survival plus wall anti-concentration}\\
\mathrm{TubularNormality}_{\mathrm{sing}}
& \text{finite-audit}
& \text{polynomial second fundamental form of }\mathcal M_{\mathrm{sing}}^{\mathrm{arith}}\\
\mathrm{MovingWallAntiConcentration}_{Y12}
& \text{open}
& \text{sublevel control for }a\Gamma_G+b\Gamma_{G,2}\\
\mathrm{O1/O2/O3}
& \text{historical alternate route}
& \text{zero-conditioned actual-orbit nonrecurrence}
\end{array}
\]

The proof must not claim closure from the historical O1/O2/O3 package unless it
explicitly returns to the ordinary finite-row Wronskian route.  In the current
route, the decisive unresolved finite downstream question is
\[
  T_{\mathrm{tr}\to\mathrm{sing}}
  \operatorname{span}\{[\mathsf C_7]^{\mathrm{tr}},[\mathsf C_{Y12}]^{\mathrm{tr}}\}
  \subset R_{\mathrm{sing}}
  \quad\text{or not.}
\]
In the current \(Y12\)-signature specialization, this rank test is refined to
whether
\[
  [\mathsf C_{Y12}]\ne0
  \quad\text{in}\quad
  \mathcal Q_{9,Y12}^{\mathrm{forb},\alpha}.
\]

In particular, the downstream status should not be summarized as ``D4 closed''
without qualification.  The accurate statement is:
\[
\boxed{
\text{D4 is closed as local carrier-quotient visibility for }u=q''_{\can},
\text{ but proof closure still requires a Type I/II forbidden channel.}
}
\]
\end{remark}

\begin{remark}[Empirical support for O1: local amplitude and cell tests]
\label{rem:empirical-support-o1-local-amplitude-cell-tests}
The current numerical evidence for O1 is purely diagnostic but consistent.
The local \(Z_\alpha\)-amplitude stress tests checked both plus-only and
reflected-dual models over \(1152\) target/window rows and found
\[
\text{Yellow}=0,\qquad \text{Red}=0.
\]
A hardened v2 run selected adversarial targets by smallest local amplitude,
smallest \(J_{\max}\), combined threat score, minima crowding, determinant
small-\(\eta_Z\) rows, and direct jet minimizers.  Its cell-mass diagnostic
used
\[
\mathrm{CellMassRatio}_M
=
\frac{\sum_r|B_r|^2}{\sum_n|z_n|^2},
\qquad
M\in\{R,2R,4R\}.
\]
In the recorded multiblock \(K=8\) cell test, \(794\) rows produced
\[
\text{Yellow}=0,\qquad \text{Orange}=0,\qquad \text{Red}=0.
\]
These tests support the heuristic
\[
Z_\alpha(m_0)\approx0
\not\Rightarrow
\text{polylog-local amplitude, jet, or cell collapse}.
\]
A particularly useful adversarial subfamily consisted of \(52\) reflected-dual
direct-jet-minimizer rows for which the optimizer drove
\[
|Z_\alpha(m_0)|\sim10^{-15}\text{--}10^{-12}.
\]
In these rows, the rescue diagnostics remained non-small: the local amplitude
at polylogarithmic radius, the jet size \(J_{\max}\), and the multiscale cell
mass ratios did not collapse.  This is the empirical behavior predicted by
\[
Z_\alpha(m_0)=0
\Longrightarrow
|\mathbf z'(m_0,\alpha)|\ge Q^{-A}.
\]

The \(K=8\) multiblock cell test is especially relevant because it addresses
the main weakness of earlier single-block diagnostics.  It used log-uniform
support coverage and recorded actual support width and active-cell count per
row.  Rows with undercovered support, too few active cells, or out-of-band
widened support were demoted to diagnostic status rather than counted as
Yellow/Orange/Red evidence.  In the recorded \(K=8\) run, \(794\) rows gave
\[
\text{Yellow}=0,\qquad
\text{Orange}=0,\qquad
\text{Red}=0.
\]

They do not prove O1, and O1 is no longer the active Gate--1 definition in the
current architecture.  The active Gate--1 object is local scalar \(L^2\)-energy
of \(q''_{\can}\), as in
Hypothesis~\ref{hyp:agg-l2-qcan-gate-one}.  The computations above should be
kept only as historical adversarial diagnostics for the old ordinary
finite-row route.  They should not be cited as evidence for
\(\mathrm{Agg}L^2(q''_{\can})\), finite/O2/O3, \(Y12\)-normal D4, or any
uniform polynomial theorem over admissible heights.
\end{remark}

\begin{remark}[Zero-side Gate--1 proxy and watch-rerun evidence]
\label{rem:zero-side-gate-one-proxy-watch-evidence}
The zero-side Gate--1 proxy scans did not produce a stable source-flat
candidate.

In the canonical regularized zero-side pilot using \(N=180\) actual zeta zeros,
\(260\) intervals, and \(121\) sample points per interval, the scan found
\[
\min_I \sup_I |q''_{\can,N}|\approx 1.6854,
\]
and
\[
\min_I |I|^{-1}\int_I |q''_{\can,N}(m)|^2\,dm
\approx 2.4954.
\]
The flattest interval occurred near
\[
m_0\approx22.8958,\qquad Q\approx3.8241,
\]
with
\[
B_I^{\rm eff}
=
-\frac{\log E_I}{\log Q}
\approx -0.6817.
\]
Thus the normalized interval energy stayed order-one or larger in this
canonical zero-side pilot; no source-flat interval was observed.

The targeted watch-rerun evidence is also favorable.  In one targeted rerun of
the watch centers, \(8\) centers generated \(1152\) rows:
\[
\text{healthy}=272,\qquad
\text{watch}=880,\qquad
\text{candidate\_flat}=0,\qquad
\text{serious}=0.
\]
The smallest proxy energy was
\[
E_I\approx9.8\times10^{-5},
\qquad
B_I^{\rm eff}\approx23.36,
\]
but this occurred only under an aggressive regularization setting and remained
classified as watch-level because the regularization drift was large,
approximately
\[
D_{\rm reg}\approx0.99.
\]

A larger watch rerun with \(8640\) rows gave
\[
\text{reject\_tail\_edge}=6480,\quad
\text{healthy}=1728,\quad
\text{watch}=432,
\]
and again
\[
\text{candidate\_proxy\_flat}=0,\qquad
\text{serious\_proxy\_warning}=0.
\]
The best stable watch rows had low interval energy but failed the finite-jet
criterion; for example near
\[
m_0\approx109.099,\qquad Q\approx2.85,
\]
one had
\[
B_I^{\rm eff}\approx7.25,\qquad
E_I\approx5.0\times10^{-4},\qquad
J_{\max}\approx0.022,
\]
whereas
\[
Q^{-5}\approx5.3\times10^{-3}.
\]
Thus these were smooth gap-midpoint or regularization artifacts, not
quadrature-like source-flat candidates.
\end{remark}

\begin{remark}[Broad proxy search top-target status]
\label{rem:broad-proxy-top-target-status}
A broad zero-proxy extraction produced
\[
46408\text{ input rows},\qquad 46396\text{ valid rows},
\]
and the top \(50\) ranked rows had only
\[
B_I^{\rm eff}\in[0.83,0.87].
\]
All top rows were classified as healthy, with no tail-edge failure.  Their
finite-jet diagnostics were far from flat:
\[
J_{\max}\sim10^3 Q^{-5}
\]
for the leading rows.  The ranking also came in physical-center pairs because
the labels `gap\_midpoint' and `large\_gap\_midpoint' can refer to the same
center.  Thus the top \(50\) rows cover roughly \(25\) unique physical targets.

These rows are useful canonical cross-check targets, but they are not Gate--1
candidates.
\end{remark}

\begin{remark}[Same-scale quadrature negative-control benchmark]
\label{rem:same-scale-quadrature-negative-control-numerics}
The same-scale quadrature negative control was verified to produce the expected
source-flat signature.  In a representative run at
\[
Q\approx3.078,\qquad d=1/Q,\qquad
I=[-0.2/Q,0.2/Q],
\]
the positive same-scale quadrature model gave
\[
\sup_I |q''_{\rm quad}|\approx5.7\times10^{-15},
\]
\[
|I|^{-1}\int_I |q''_{\rm quad}(m)|^2\,dm
\approx3.5\times10^{-30},
\]
and
\[
B_I^{\rm eff}\approx60.3.
\]
This confirms that the numerical pipeline can detect the known
same-scale source-flattening adversary when it is present.  The actual-zero
proxy scans above did not resemble this negative control.
\end{remark}

\begin{remark}[Packet-frequency determinant stress-test status]
\label{rem:packet-frequency-determinant-stress-status}
The ordinary packet-frequency determinant stress tests use active blocks
\[
\mathcal B_{\rm pf}
=
\{n:cQ\le\log n\le cQ+Q/R\},
\qquad
0<c<1/2,
\]
so that \(\lambda\log n\asymp1\).  This avoids the low-prime node-compression
failure.  The v3 stress tests ran through
\[
R=4,6,8,10,12,16
\]
with reflected dual, actual-like dual, and \(E_2\)-perturbation stages.  The
completed larger-toy run recorded
\[
\text{Yellow}=0,\qquad
\text{Red}_{\rm pre\text{-}E2}=0,\qquad
\text{Red}=0.
\]
The old \(m=0,\alpha=0\) artifacts disappeared after imposing
\[
m_{\min}=50.
\]
This evidence does not prove the Wronskian incidence theorem, but it found no
surviving packet-frequency determinant collapse under the current diagnostics
\[
\eta_Z,\qquad \eta_W,\qquad \rho_{\rm rank},\qquad
\kappa_{\rm basis}.
\]
Hence no numerical stop rule is currently triggered for the ordinary
packet-frequency Wronskian route.

The detailed determinant stress archive also found no stop-rule event.  The
larger-toy run covered
\[
R=4,6,8,10,12,16
\]
with reflected dual, actual-like dual, and \(E_2\)-perturbation stages, and
reported
\[
\text{Yellow}=0,\qquad
\text{Red}_{\rm pre\text{-}E2}=0,\qquad
\text{Red}=0.
\]
The \(R=16\) row count was lower than the smaller-\(R\) counts because some
blocks or \(\alpha\)-cases were filtered before emission, not because the run
failed.

The best actual-like non-collapsing candidate occurred at roughly
\[
R=6,\qquad Q=58,\qquad c=0.2,\qquad \alpha=0,\qquad m\approx5377.16,
\]
with
\[
\eta_Z\approx2.41\times10^{-5},
\qquad
\eta_W\approx1.35\times10^{-5},
\]
\[
\rho_{\rm rank,2}\approx2.89\times10^{-4},
\qquad
\kappa_{\rm basis}\approx4.02\times10^{4}.
\]
It was classified as `no\_collapse'.  Some rows had very small determinants,
for instance
\[
\eta_W\sim10^{-13},
\]
but these did not co-occur with sufficiently small \(\eta_Z\), or were filtered
as degenerate.  Thus determinant-smallness alone did not produce an incidence
collapse.
\end{remark}

\begin{remark}[Numerical classification rule for Gate--1 warnings]
\label{rem:numerical-classification-rule-gate-one}
A numerical row is not a serious Gate--1 warning unless all of the following
hold:
\[
E_I\text{ is small},
\]
\[
J_{\max}\text{ is small through the declared jet depth},
\]
the tail, quadrature, regularization, and width drifts are stable, and either
the AFE cross-check passes or the row is explicitly labeled
\[
\texttt{serious\_warning\_zero\_side\_only}.
\]
A route-threatening conclusion requires a family, not a single interval.  A
single low-energy row is only watch-level unless it also exhibits finite-jet
flatness and stability.
\end{remark}

\begin{remark}[Gate--1 numerical validation sequence]
\label{rem:gate-one-numerical-validation-sequence}
Future Gate--1 numerical runs should follow this order:
\[
\text{same-scale quadrature negative control}
\to
\text{random packet positive/control test}
\to
\text{singular dominance and cancellation tests}
\to
\text{actual zero-side scan}
\to
\text{candidate refinement}
\to
\text{AFE cross-check}
\to
\text{surrogate comparison and height trend}.
\]
Each run should report at least
\[
E_I,\quad S_I,\quad B_I^{\rm eff},\quad J_{\max},\quad J_{\rm rms},
\]
together with tail, quadrature, regularization, width, and AFE stability
metrics.  The flattest \(0.1\%\) of rows should be rerun with higher precision,
larger zero cutoff, more quadrature nodes, finite jets, and AFE cross-checks.
\end{remark}

\begin{remark}[No Green \(\mathrm{Agg}L^2(q''_{\can})\) proof yet]
\label{rem:no-green-gate-one-proof-yet}
At the current stage, no argument in the draft proves the active
\(\mathrm{Agg}L^2(q''_{\can})\) theorem for the full canonical object
\(q''_{\can}\).  Formal Wronskian or Cauchy--Binet nonvanishing is insufficient
unless it is upgraded to a weighted \(L^2\)-source theorem for
\[
q''_{\can}=B_2+P_2+D_2+E_2.
\]
A proof for a zero-kernel proxy, a flexible Dirichlet polynomial, a
\(P_2\)-only source, a \(B_2\)-only diagnostic, or the affine-subtracted
complement curvature \(q''_{\source,R}\) does not prove the active input.

Equivalently, the present draft has not proved
\[
  \mathrm{Coverage}_{q''},
  \qquad
  \mathrm{MeanSq}_{q''_{\can}},
\]
nor the projected source-frame versions of the ordinary, transition, and
singular proof tools.  The historical \(\mathrm{O1/O2/O3}\) route is also not
proved, but it is no longer the active Gate--2 input.
\end{remark}

\begin{thebibliography}{99}

\bibitem{Riemann1859}
B.~Riemann,
\newblock Ueber die Anzahl der Primzahlen unter einer gegebenen Gr\"osse,
\newblock \emph{Monatsberichte der Berliner Akademie}, 1859, 671--680.

\bibitem{Bombieri2000}
E.~Bombieri,
\newblock \emph{Problems of the Millennium: The Riemann Hypothesis},
\newblock Clay Mathematics Institute, 2000.

\bibitem{Conrey2003}
J.~B.~Conrey,
\newblock The Riemann Hypothesis,
\newblock \emph{Notices of the American Mathematical Society} \textbf{50} (2003), no.~3, 341--353.

\bibitem{Edwards2001}
H.~M.~Edwards,
\newblock \emph{Riemann's Zeta Function},
\newblock Dover Publications, Mineola, NY, 2001.

\bibitem{Titchmarsh1986}
E.~C.~Titchmarsh,
\newblock \emph{The Theory of the Riemann Zeta-Function},
\newblock 2nd ed., revised by D.~R.~Heath-Brown,
\newblock Oxford University Press, Oxford, 1986.

\bibitem{Speiser1935}
A.~Speiser,
\newblock Geometrisches zur Riemannschen Zetafunktion,
\newblock \emph{Mathematische Annalen} \textbf{110} (1935), 514--521.

\bibitem{LevinsonMontgomery1974}
N.~Levinson and H.~L.~Montgomery,
\newblock Zeros of the derivatives of the Riemann zeta-function,
\newblock \emph{Acta Mathematica} \textbf{133} (1974), 49--65.

\bibitem{deBranges1986}
L.~de~Branges,
\newblock The Riemann hypothesis for Hilbert spaces of entire functions,
\newblock \emph{Bulletin of the American Mathematical Society (N.S.)}
\textbf{15} (1986), no.~1, 1--17.

\bibitem{ConreyLi2000}
J.~B.~Conrey and X.-J.~Li,
\newblock A note on some positivity conditions related to zeta and L-functions,
\newblock \emph{International Mathematics Research Notices} \textbf{2000}
(2000), no.~18, 929--940.

\bibitem{Connes1997}
A.~Connes,
\newblock Trace formula in noncommutative geometry and the zeros of the Riemann zeta function,
\newblock \emph{Journ\'ees \'Equations aux d\'eriv\'ees partielles} (1997), Art.~No.~4, 1--28.

\bibitem{ConnesConsani2021}
A.~Connes and C.~Consani,
\newblock Weil positivity and trace formula: the archimedean place,
\newblock arXiv:2006.13771, 2021.

\end{thebibliography}

\end{document}
